% Edisi Inggris: 32 bab + seluruh materi belakang (materi depan, lampiran A--D,
% indeks, kolofon, daftar rujukan). Materi belakang ada di belakang/*-en.tex;
% sumber Indonesianya disalin apa adanya ke belakang/_sumber/ untuk gerbang paritas.
\documentclass[11pt,a4paper,oneside]{book}
\usepackage[T1]{fontenc}
\usepackage[utf8]{inputenc}
\usepackage{lmodern}
\usepackage{microtype}
\usepackage[a4paper,top=2.6cm,bottom=2.8cm,left=3.0cm,right=2.6cm]{geometry}
\usepackage{amsmath}
\usepackage{amssymb}
\usepackage{amsthm}
\usepackage{enumitem}
\usepackage{booktabs}
\usepackage{array}
\usepackage{xcolor}
\usepackage{graphicx}
\usepackage{tikz}
\usepackage{pgfplots}
\pgfplotsset{compat=1.18}
\usetikzlibrary{arrows.meta,positioning,calc,patterns}
\usepackage[font=small,labelfont=bf,justification=centering]{caption}
\usepackage[most]{tcolorbox}
\usepackage{makeidx}
\usepackage[hidelinks,bookmarksnumbered,bookmarksopenlevel=1,plainpages=false]{hyperref}

\hypersetup{
  pdftitle={Why Machines Can Be Wrong: From Linear Algebra to Machine-Learning Security (First Edition)},
  pdfauthor={Afiq Andico Pangimpian},
  pdfsubject={A tiered textbook: mathematics, computing, machine learning, and security},
  pdfkeywords={linear algebra, real analysis, probability, optimisation, neural networks, adversarial ML, Lean 4},
  pdfcreator={pdfLaTeX},
}

\definecolor{ink}{HTML}{1A1A1A}
\definecolor{accent}{HTML}{8C2F39}
\definecolor{layerorange}{HTML}{B45309}
\definecolor{layerteal}{HTML}{0F766E}
\definecolor{layergrey}{HTML}{374151}
\definecolor{layerindigo}{HTML}{4338CA}
\definecolor{layerjade}{HTML}{15803D}
\definecolor{layerbrown}{HTML}{92400E}
\definecolor{layerplum}{HTML}{86198F}
\color{ink}

% Environment names stay Indonesian so Indonesian and English chapters share one
% structural checker; only the printed labels are translated.
\theoremstyle{definition}
\newtheorem{definisi}{Definition}[chapter]
\theoremstyle{plain}
\newtheorem{teorema}{Theorem}[chapter]
\newtheorem{lemma}[teorema]{Lemma}
\newtheorem{akibat}[teorema]{Corollary}
\newtheorem{proposisi}[teorema]{Proposition}
\theoremstyle{remark}
\newtheorem{contoh}{Example}[chapter]
\newtheorem{latihan}{Exercise}[chapter]
\newtheorem{catatan}{Note}[chapter]
\newtheorem*{ringkasan}{Summary}

\renewcommand{\chaptername}{Chapter}
\renewcommand{\partname}{Part}
\renewcommand{\contentsname}{Contents}
\renewcommand{\bibname}{References}
\renewcommand{\indexname}{Index}
\renewcommand{\figurename}{Figure}
\renewcommand{\tablename}{Table}
\renewcommand{\listfigurename}{List of Figures}

\setlist[itemize]{topsep=0.35em,itemsep=0.15em,parsep=0pt}
\setlist[enumerate]{topsep=0.35em,itemsep=0.15em,parsep=0pt}

% ---------------- four reader layers ----------------
\newtcolorbox{untukanak}{colback=orange!6,colframe=layerorange,
  title={\small\bfseries For younger readers (age 10 and up)},
  fonttitle=\bfseries,breakable,left=6pt,right=6pt,top=4pt,bottom=4pt}
\newtcolorbox{untukdosen}{colback=teal!5,colframe=layerteal,
  title={\small\bfseries Note for instructors},
  fonttitle=\bfseries,breakable,left=6pt,right=6pt,top=4pt,bottom=4pt}
\newtcolorbox{untukpeneliti}{colback=black!4,colframe=layergrey,
  title={\small\bfseries For researchers: limits and next steps},
  fonttitle=\bfseries,breakable,left=6pt,right=6pt,top=4pt,bottom=4pt}
% ---- learning-style channels: auditory, kinesthetic, procedural, problem solving ----
\newtcolorbox{bacakeras}{colback=violet!5,colframe=layerindigo,
  title={\small\bfseries Read it aloud (auditory learner)},
  fonttitle=\bfseries,breakable,left=6pt,right=6pt,top=4pt,bottom=4pt}
\newtcolorbox{tangan}{colback=green!5,colframe=layerjade,
  title={\small\bfseries Do it by hand (kinesthetic learner)},
  fonttitle=\bfseries,breakable,left=6pt,right=6pt,top=4pt,bottom=4pt}
\newtcolorbox{resep}{colback=orange!4,colframe=layerbrown,
  title={\small\bfseries Step-by-step recipe (procedural learner)},
  fonttitle=\bfseries,breakable,left=6pt,right=6pt,top=4pt,bottom=4pt}
\newtcolorbox{tantangan}{colback=magenta!4,colframe=layerplum,
  title={\small\bfseries Challenge (problem solver)},
  fonttitle=\bfseries,breakable,left=6pt,right=6pt,top=4pt,bottom=4pt}

% In the English edition the first argument is the English term and the second is
% the Indonesian original, so the printed pair stays the same way round.
\newcommand{\kamus}[3]{%
  \par\smallskip\noindent
  \colorbox{blue!5}{\parbox{\dimexpr\linewidth-2\fboxsep}{%
    \small\textbf{#1} (\emph{#2}) --- #3}}\par\smallskip}
\newcommand{\tingkat}[1]{%
  \par\noindent{\footnotesize\scshape Level: #1}\par\smallskip}
\newcommand{\kertas}[1]{%
  \par\medskip\noindent
  \fbox{\parbox{\dimexpr\linewidth-2\fboxsep-2\fboxrule}{\small #1}}%
  \par\medskip}
\newcommand{\peta}[1]{%
  \par\medskip\noindent{\small\itshape Map to source texts: #1}\par\medskip}

\newenvironment{lab}{\par\medskip\noindent\textbf{Computational lab}\par\nopagebreak}{\par\medskip}
\makeatletter
\renewcommand{\verbatim@font}{\footnotesize\ttfamily}
\makeatother

\makeindex
\setlength{\emergencystretch}{3em}

% ---------------------------------------------------------------------------
% Running head untuk judul bab yang panjang.
%
% Kepala halaman hanya memuat "CHAPTER n. <judul>". Kalau judul penuh terlalu
% lebar, kepala halaman meluber (Overfull \hbox) dan itu bukan hal yang boleh
% dibiarkan. Menyingkat judulnya dilarang: halaman bab dan daftar isi harus
% memuat terjemahan utuh. Jadi yang dipendekkan hanya mark header, dan
% pemendekannya berlaku SEKLALI untuk bab berikutnya:
%
%   \judulpendek{Riemann Integrals}
%   \chapter{Riemann Integrals and the Fundamental Theorem of Calculus}
%
% Tanpa mekanisme sekali-pakai ini, \renewcommand{\chaptermark} di dalam satu
% berkas bab akan bocor ke semua bab sesudahnya - pernah terjadi: header bab
% 14-18 ikut berbunyi "CHAPTER 13. RIEMANN INTEGRALS".
\makeatletter
% \@chapapp (bukan \chaptername): setelah \appendix, \@chapapp menjadi "Appendix",
% sedangkan \chaptername tetap "Chapter", sehingga kepala halaman lampiran dulu
% berbunyi "CHAPTER A. TABLE OF NOTATION" tanpa satu pun peringatan LaTeX.
\newcommand*\judulpendek[1]{\gdef\kepala@pendek{#1}}
\newcommand*\kepala@pendek{}
\renewcommand*\chaptermark[1]{%
  \markright{\MakeUppercase{%
    \ifnum\c@secnumdepth>\m@ne\if@mainmatter\@chapapp\ \thechapter.\ \fi\fi
    \ifx\kepala@pendek\@empty #1\else\kepala@pendek\fi}}%
  \global\let\kepala@pendek\@empty}
\makeatother


\begin{document}

\frontmatter

% ---------------------------- title page ----------------------------
\hypersetup{pageanchor=false}
\begin{titlepage}
\centering
\vspace*{1.6cm}
{\color{accent}\rule{\textwidth}{1.2pt}}\\[0.45cm]
{\large\scshape Layered Textbook}\\[0.15cm]
{\itshape Mathematics $\cdot$ Computation $\cdot$ Machine Learning $\cdot$ Security}\\[0.45cm]
{\color{accent}\rule{\textwidth}{0.6pt}}\\[1.8cm]
{\Huge\bfseries Why Machines\\[0.25cm] Can Be Wrong?}\\[0.7cm]
{\Large\itshape From Linear Algebra\\[0.18cm] to Machine-Learning Security}\\[0.7cm]
{\large\itshape First Edition}\\[0.5cm]
{\large Six parts, thirty-two chapters, four reading tracks}\\[0.35cm]
{\normalsize\itshape for readers aged 10, students, lecturers, and researchers}\\[2.0cm]
{\large Compiled by Afiq Andico Pangimpian}\\[0.3cm]
{\large October 2026}
\vfill
{\footnotesize\itshape A textbook about how we know that an answer is true --- and
how machines fail to know. Full sources are in the References.}
\end{titlepage}
\hypersetup{pageanchor=true}

\tableofcontents

% ---------------------------- front matter ----------------------------
\chapter*{Preface}
\addcontentsline{toc}{chapter}{Preface}
\markboth{Preface}{Preface}

The first edition of this book grew out of one simple problem: the collection of
textbooks I had gathered ran to more than ten thousand pages, but those pages were
never read in sequence. Opening Rudin at chapter three without ever touching the
definition of a supremum produces nothing but frustration. What was missing was not
material, but order.

Before this manuscript there was a thirteen-chapter version. Read again, it showed two
weaknesses. First, thirteen chapters went too fast, so the hardest parts went by in a
glance. Second, that book could be read by only one kind of reader: a student already
used to formal notation. A curious middle-school student had no door in, and a lecturer
found no teaching notes inside it.

This version repairs both. Its scope has been widened to six parts and thirty-two
chapters, with material the thirteen-chapter version did not have: sets and
cardinality, infinite series, systems of linear equations and determinants, multiple
integrals together with the vector field theorems, differential equations and dynamical
systems, regression, and model generalization and regularization. Every chapter is now
built in layers.

The first layer is the \emph{For younger readers} box. It holds analogies to everyday
objects, without a single mathematical symbol, and closes with one question that makes
the reader curious. The second layer is the body of the chapter itself: definitions,
theorems, proofs, worked examples, and exercises, as in an ordinary textbook. The third
layer is the \emph{Note for instructors} box: the teaching order I recommend, the
misconceptions that almost always surface, and one classroom activity. The fourth layer
is the \emph{For researchers} box: what is really proved, what is only popularized, the
primary sources, and the questions that are still open.

The four layers do not replace one another: a younger reader can read a chapter from the
orange box and the summary alone, an undergraduate reads everything except the
researcher box, an instructor uses the green box as a plan for the meeting, and a
researcher starts from the gray box. One book, four paths.

Those layers set the \emph{depth}. What sets the \emph{way in} is the second set of four
boxes: \emph{Read it aloud} for readers who understand by hearing, \emph{Do it by hand}
for those who understand by moving and touching, \emph{Step-by-step recipe} for those who
understand by memorizing a procedure first, and \emph{Challenge} for those who understand
precisely when the problem is not yet clear. The two sets live together: one and the same
chapter, with six doors.

Three principles of writing have not changed. \emph{No claim without a testing tool}:
every theorem comes with a proof, or with an honest admission that the proof is skipped
together with where it is skipped; every number comes from a real execution that the
reader can repeat. \emph{Notation is introduced before it is used}. \emph{Honesty in
evaluation is part of mathematics}: the last chapters do not stop at how to build a
model, but at how to know whether a claim about that model is proved.

One thing is deliberately absent from this book: complete answer keys. Giving a final
number without a trace of the check would violate the first principle. In its place,
Appendix~C gives grading criteria and a verification tool for every kind of exercise, so
that readers can check their own work in the same way the author checked this manuscript.

This book is not a substitute for the source texts it points to, but an ordered way in:
if one chapter makes a reader open Rudin, Axler, Boyd, or a paper by Athalye and then
read the section it names, this book has succeeded. Whatever errors remain belong to the
author, not to those texts.

\enlargethispage{3\baselineskip}
\vspace{0.35cm}
\begin{flushright}
Afiq Andico Pangimpian\\
October 2026
\end{flushright}

\chapter*{How to Read This Book: Four Paths}
\addcontentsline{toc}{chapter}{How to Read This Book: Four Paths}
\markboth{How to Read This Book}{How to Read This Book}

Every chapter has the same shape, so the reader knows what to expect and may skip ahead
without losing the way.

\begin{center}
\small
\begin{tabular}{@{}p{3.1cm}p{4.6cm}p{4.4cm}@{}}
\toprule
\textbf{Path} & \textbf{What is read} & \textbf{What you get} \\
\midrule
Younger reader (10-14) & the \emph{For younger readers} box, the Summary, one worked
  example & the core idea without notation; you know what is being computed \\
Undergraduate (S1) & the whole body of the chapter, the Computational lab, the Exercises;
  skip the researcher box & definitions, proofs, the ability to compute and to check \\
Instructor & the \emph{Note for instructors} box, the source map, the Exercises & a plan
  for the meeting, misconceptions to anticipate, blackboard material \\
Researcher & the \emph{For researchers} box, the source map, Appendices~A and~B & the
  limits of the claims, primary sources, open questions, consistent notation \\
\bottomrule
\end{tabular}
\end{center}

\medskip
The \textsc{level} mark at the start of every section shows how hard the reading is:
\emph{basic} can be read without any earlier chapter other than the ones named;
\emph{intermediate} demands the algebra and the notation of this book;
\emph{advanced} may be skipped on a first reading without breaking the flow of the
chapter.

\medskip
\noindent\textbf{Six ways of learning, six doors in.} People do not learn in the same
way. So every chapter carries more than one route to the same idea, and no route is
compulsory: pick the one that makes you understand fastest, then go back to the body of
the chapter to see why that way works.

\begin{center}
\small
\begin{tabular}{@{}p{3.4cm}p{4.3cm}p{4.6cm}@{}}
\toprule
\textbf{Type} & \textbf{Look for this in every chapter} & \textbf{How to use it} \\
\midrule
Visual & 2-3 numbered figures (in the List of Figures), tables, and the box colors & read
  the figure before its text; the figure carries the same numbers as the computation in
  the text \\[2pt]
Auditory & the \emph{Read it aloud} box (indigo) & read the text aloud or record your
  voice; repeat the passages that have a rhythm \\[2pt]
Kinesthetic & the \emph{Do it by hand} box (green) & do it with real objects: paper,
  cards, the floor, money, or your own body \\[2pt]
Concept-first & definitions, theorems, proofs, the \emph{For researchers} box & do not
  use a formula before you know where it comes from; the proof is where that origin is
  written \\[2pt]
Procedural memorizer & the \emph{Step-by-step recipe} box (brown) and the Exercises &
  memorize the order of the recipe cards, then apply it to different numbers until it is
  automatic \\[2pt]
Problem solver & the \emph{Challenge} box (purple) and the starred exercise items & take
  the challenge first; the hints are tiered, do not open them all at once \\
\bottomrule
\end{tabular}
\end{center}

The four paths below set \emph{how much} is read; the six types above set \emph{which
door} you come in through. The two may be combined: a visual reader who pursues the
researcher path still reads every proof, only starting from the figure.

Three learning paths that have already been tried, from the lightest:

\begin{enumerate}
  \item \textbf{Curiosity path (8-12 weeks).} Read the orange box, the Summary, and one
        worked example from every chapter, in order from Chapter~1 to Chapter~32. No
        exercises. The aim is one thing: to know the map of the terrain.
  \item \textbf{Student path (2 semesters).} First semester: Chapters~1-10 and
        Chapter~17. Second semester: Chapters~11-16, then 18-32. That is twenty-one
        chapters, more than an ordinary semester; the way to trim it without breaking the
        prerequisites is to postpone Chapters~31-32 (cryptography and privacy) to an
        elective course, so that nineteen chapters remain. Do the Computational lab of
        every chapter; at least four exercise items per chapter. This order keeps the
        promise about prerequisites: Chapter~18 demands Chapters~13 and~17, so the three
        must not be far apart; Chapter~21 demands Chapters~10 and~20; Chapter~23 demands
        Chapter~12.
  \item \textbf{Machine-learning security researcher path (9 weeks).} Nineteen chapters,
        not seven: Chapters~3, 4, 9, 10, 12, 13, 18, 19, 20, 22, 23, 25, then 26-30,
        closing with Chapters~31 and~32. Chapters~13, 18, 22, and~23 are included
        because each is a prerequisite of Chapter~20 (hypothesis testing), Chapter~29
        (adversarial), and Chapter~25 (neural networks); Chapters~31 and~32 complete the
        cryptography and privacy side that the title of the book demands. Reading this
        path in six weeks means skipping the prerequisites and losing the reason why the
        claims hold.
\end{enumerate}

Symbols are used consistently throughout the book: $\mathbb{N}$, $\mathbb{Z}$,
$\mathbb{Q}$, $\mathbb{R}$ for the number sets; $\mathbf{x}$ for a vector and $x_i$ for
its components; $\|\mathbf{x}\|$ for the Euclidean norm; $\nabla f$ for the gradient;
$\varepsilon$ and $\delta$ always small positive numbers; $\mathbb{E}[X]$ for the
expected value; $\mathcal{L}(V,W)$ for the space of linear transformations; $\top$ for
the transpose. Appendix~A contains the complete list.

\chapter*{Chapter Dependency Map}
\addcontentsline{toc}{chapter}{Chapter Dependency Map}
\markboth{Chapter Dependency Map}{Chapter Dependency Map}

This book may be read out of order, but not every chapter stands on its own. The table
below names the minimum prerequisites of each chapter; every other prerequisite may be
ignored on a first reading.

\begin{center}
\small
\begin{tabular}{@{}r p{9.0cm} p{2.8cm}@{}}
\toprule
\textbf{Chapter} & \textbf{Title} & \textbf{Prerequisites} \\
\midrule
1 & The Language of Mathematics and the Art of Proof & --- \\
2 & Sets, Relations, Functions, and Cardinality & 1 \\
3 & Axioms of the Real Numbers and Supremum & 1, 2 \\
4 & Sequences and Limits & 3 \\
5 & Infinite Series & 4 \\
6 & Vector Spaces and Subspaces & 2 \\
7 & Basis, Dimension, and Linear Transformations & 6 \\
8 & Linear Systems, Determinants, and Inverses & 6 \\
9 & Eigenvalues and Diagonalization & 7, 8 \\
10 & Orthogonality, SVD, and Least Squares & 7, 9 \\
11 & Continuity and Uniform Continuity & 3, 4 \\
12 & Derivatives: One and Several Variables & 11, 7 \\
13 & Riemann Integrals and the Fundamental Theorem of Calculus & 11 \\
14 & Series of Functions and the Exchange of Limits & 5, 11, 13 \\
15 & Multiple Integrals, Change of Variables, and Vector Field Theorems & 10, 12, 13 \\
16 & Differential Equations and Dynamical Systems & 12, 14, 15 \\
17 & Probability and Conditional Probability & 2, 5 \\
18 & Random Variables and Distributions & 17, 13 \\
19 & The Law of Large Numbers and the Central Limit Theorem & 18, 5 \\
20 & Estimation, Confidence Intervals, and Hypothesis Testing & 19 \\
21 & Regression and Applied Statistics & 10, 20 \\
22 & Floating-Point Arithmetic and Error Analysis & 3, 5 \\
23 & Convex Optimization and Duality & 6, 12 \\
24 & Linear and Integer Programming & 8, 23 \\
25 & Neural Networks and Backpropagation & 10, 12, 23 \\
26 & Generalization, Regularization, and Training & 18, 25 \\
27 & Honest Model Evaluation & 20, 26 \\
28 & Formal Methods with Lean 4 & 1, 2 \\
29 & Adversarial Machine Learning: Attacks and Defenses & 25, 22 \\
30 & The Limits of Attack: Obfuscated Gradients and NIDS Evasion & 27, 29 \\
31 & Cryptography: Modular Integers to the Quantum Threat & 2, 3, 28 \\
32 & Differential Privacy and the Mathematical Foundations of QML & 10, 17, 18 \\
\bottomrule
\end{tabular}
\end{center}

\phantomsection
\pdfbookmark[0]{Contents}{toc}
\tableofcontents
\cleardoublepage
\phantomsection
\pdfbookmark[0]{List of Figures}{lof}
\listoffigures


\mainmatter

% ---------------------------- part I ----------------------------
\part{Foundations: Language, Sets, and Numbers}

This part answers a question that is usually skipped: what the words ``holds for all''
mean, and how a sentence can be proved rather than merely believed. The only prerequisite
is school algebra. After five chapters you can write direct proofs, proofs by contradiction,
and induction, and you know when an argument is still incomplete.
\chapter{The Language of Mathematics and the Art of Proof}

In a classroom, a student writes on the board: ``$n^{2}+n+41$ is prime for every
integer $n \geq 0$''. He then checks them one at a time. For $n=0$ he obtains
$41$; for $n=1$ he obtains $43$; for $n=2$ he obtains $47$. He continues up to
$n=39$, obtains $1601$, and $1601$ is indeed prime. Forty values, forty
failures that never came. At $n=40$ the claim breaks:
$40^{2}+40+41 = 1681 = 41^{2}$. One extra line cancels the whole preceding page.

That event is not an arithmetic accident but a grammatical mistake. ``For every
$n \geq 0$'' speaks about infinitely many numbers, whereas a list touches only
finitely many; no rule of logic moves truth from the finite to the infinite. This
chapter therefore offers a way of speaking such that checking is not needed.

The need is not a luxury. In Chapter~29 we will conclude that a
classifier can be encircled by shifting one input by $\varepsilon$,
and that conclusion holds for all inputs. In Chapter~27 a good accuracy report on
thousands of test data proves nothing about tomorrow's data. This chapter prepares
the apparatus for both.

\begin{untukanak}
On the shore of a lake lived a flock of geese. A child came every morning for
almost a year and counted the geese he saw: white, all of them, for two hundred
and seventy mornings. So he wrote in his notebook: all geese on this lake are white.

His grandfather asked how many geese the child had never seen. The child admitted he
did not know, and his grandfather said that the notebook told about the geese already
seen, not about all the geese on the lake.

This story is not about geese but about the difference between ``I have checked
everything'' and ``it holds for everything''. The first is only as strong as our memory. The second does not depend on who does the checking. One black goose
suffices to bring down a notebook of any thickness, whereas a thousand white geese
are never enough to set it up.

The black goose of this chapter is called forty. At that number the formula that looks tidy
suddenly yields a number that can be split. In this chapter we will learn why one
example brings down a claim, why many examples do not establish it, and how the gap
is closed without inspecting the last goose.
\end{untukanak}

\section{Statements, connectives, and truth tables}
\tingkat{basic}

A sentence can be judged only after we know what it promises. ``Seven is a
prime number'' promises something that can be checked, and that promise is true. ``Close
the window'' promises nothing. ``$x+1=3$'' cannot yet be judged, because its value
depends on $x$. The three look the same to spelling, but only the first
counts as mathematical material.

\begin{definisi}[Statement]
\index{statement}
A \emph{statement} is a sentence that is either true or false, but not
both at once.
\end{definisi}

A single statement becomes useful once it is combined. Each connective has one meaning
that is stipulated, not inferred from habits of speech.

\kamus{logical connective}{penghubung logika}{like a button on a calculator: two
lamps in, one lamp out, and the relation between the lamps is fixed at the factory,
not guessed by its user.}

With the picture of that button in mind, we stipulate the meaning of the five connectives officially.

\begin{definisi}[Logical connectives]
Let $P$ and $Q$ be statements.

We use the following five connectives.
\begin{itemize}
  \item $\neg P$ (\emph{negation}, \emph{not}): true exactly when $P$ is false.
  \item $P \wedge Q$ (\emph{conjunction}, \emph{and}): true exactly when both are true.
  \item $P \vee Q$ (\emph{disjunction}, \emph{or}): true exactly when at least
        one is true; the word ``or'' here is inclusive.
  \item $P \Rightarrow Q$ (\emph{implication}, \emph{if--then}): false exactly when
        $P$ is true and $Q$ is false; true in the other three cases.
  \item $P \Leftrightarrow Q$ (\emph{biconditional}, \emph{iff}): true exactly when
        $P$ and $Q$ have the same value.
\end{itemize}
\end{definisi}

All the content of that definition can be compressed into a single table: the rows are all combinations
of values, the columns are the value of each connective.

{\centering
\begin{tabular}{cc|ccccc}
\hline
$P$ & $Q$ & $\neg P$ & $P \wedge Q$ & $P \vee Q$ & $P \Rightarrow Q$ & $P \Leftrightarrow Q$ \\
\hline
T & T & F & T & T & T & T \\
T & F & F & F & T & F & F \\
F & T & T & F & T & T & F \\
F & F & T & F & F & T & T \\
\hline
\end{tabular}\par}

\begin{figure}[htbp]
\centering
\begin{tikzpicture}[font=\small]
  \fill[layerjade!16] (0,0) rectangle (3,1.8);
  \fill[accent!16] (3,0) rectangle (6,1.8);
  \fill[layerjade!16] (0,-1.8) rectangle (3,0);
  \fill[layerjade!16] (3,-1.8) rectangle (6,0);
  \draw[layergrey!70,thick] (0,-1.8) grid[xstep=3,ystep=1.8] (6,1.8);
  \draw[accent,very thick] (3,0) rectangle (6,1.8);
  \node at (1.5,0.9) {T};
  \node at (4.5,0.9) {F};
  \node at (1.5,-0.9) {T};
  \node at (4.5,-0.9) {T};
  \node[above,font=\footnotesize] at (1.5,1.8) {$Q$ true};
  \node[above,font=\footnotesize] at (4.5,1.8) {$Q$ false};
  \node[left,font=\footnotesize] at (0,0.9) {$P$ true};
  \node[left,font=\footnotesize] at (0,-0.9) {$P$ false};
  \node[below,font=\footnotesize,text=accent] at (4.5,-1.8) {the only false cell};
\end{tikzpicture}
\caption{The truth grid of $P \Rightarrow Q$: rows hold the value of $P$, columns hold the value of $Q$.
Three cells hold T and one cell holds F, namely when $P$ is true while $Q$ is false.
Anyone who wants to bring the implication down need only point at the cell framed in thick lines.}
\label{fig:kisi-implikasi}
\end{figure}

The third and fourth rows are the ones that most often surprise. In the fourth row,
$P$ is false and $Q$ is false, and the implication is judged true. In the third row, $P$ is false
and $Q$ is true, and the implication is true as well. So the only way to bring down
$P \Rightarrow Q$ is to exhibit $P$ true while $Q$ is false; the grid in
Figure~\ref{fig:kisi-implikasi} contains that cell as the only false cell.

That stipulation is not a matter of style. Consider ``if $n$ is divisible by $4$, then $n$
is even'', which must hold for every integer $n$, including $n=3$. There the
antecedent is false; if the implication were judged false, the statement that is plainly true would
collapse because of one irrelevant number. An implication records a relation of values, not
cause and effect.

\kamus{implication}{implikasi}{like the promise ``if it rains, the road is wet'': the promise is
broken only when rain falls and the road stays dry; when there is no rain, the one who made the
promise has done nothing wrong.}

Two compound statements can have the same value in every case; such an equality
has a name of its own.

\begin{definisi}[Tautology and equivalence]
\index{tautology}
A compound statement is called a \emph{tautology} if it is true under every assignment
of values to its constituent statements. Two compound statements $A$ and $B$ are called
\emph{logically equivalent}, written $A \equiv B$, if $A \Leftrightarrow B$ is a tautology.
\end{definisi}

For two constituent statements there are $2^{2}=4$ assignments of values; for three constituents there are
$2^{3}=8$. An equivalence that involves only two or three statements can therefore be
proved by checking every row, and that check is complete. This is the only place in this chapter where
exhaustive checking is accepted as proof, because the set of cases is finite and
has been counted.

\begin{proposisi}
For every statement $P$ and $Q$ we have $P \Rightarrow Q \equiv \neg P \vee Q$.
\end{proposisi}
\par\smallskip\noindent\textit{Proof.} We compare the column $P \Rightarrow Q$ with the column $\neg P \vee Q$
in the four rows of the truth table. In the first row both are true. In the second row $\neg P$ is false and
$Q$ is false, so both are false. In the third and fourth rows $\neg P$ is true,
so both are true. The four rows give the same value, hence the two statements are
equivalent.
\hfill$\blacksquare${\parfillskip=0pt\par}

Two derived equivalences are used throughout this book. The first lets us reverse
the direction of the talk without changing the value.

\begin{akibat}[Contraposition]
\index{contraposition}
For every statement $P$ and $Q$ we have
$P \Rightarrow Q \equiv \neg Q \Rightarrow \neg P$.
\end{akibat}
\par\smallskip\noindent\textit{Proof.} By the preceding proposition, both sides equal $\neg P \vee Q$ and
$\neg(\neg Q) \vee \neg P$. Since $\neg(\neg Q) \equiv Q$ and a disjunction pays no
attention to order, the two forms are identical.
\hfill$\blacksquare${\parfillskip=0pt\par}

What is \emph{not} equivalent is the converse, namely $Q \Rightarrow P$. In the third row of
the table, $Q \Rightarrow P$ is true while $P \Rightarrow Q$ is false. Swapping antecedent
for consequent is the most frequent mistake, and it is often invisible because the sentence
still reads smoothly.

\begin{proposisi}
For every statement $P$ and $Q$ we have
$\neg(P \Rightarrow Q) \equiv P \wedge \neg Q$.
\end{proposisi}
\par\smallskip\noindent\textit{Proof.} By the definition of negation and the first proposition,
$\neg(P \Rightarrow Q) \equiv \neg(\neg P \vee Q)$. De Morgan's law, proved in
Chapter~2, gives $\neg(\neg P \vee Q) \equiv \neg(\neg P) \wedge \neg Q$, and
$\neg(\neg P) \equiv P$. Hence $\neg(P \Rightarrow Q) \equiv P \wedge \neg Q$.
A reader who has not yet accepted De Morgan's law can check the four rows of the table
directly.
\hfill$\blacksquare${\parfillskip=0pt\par}

\begin{tangan}
Prepare eight small cards: four bearing T, four bearing F. Lay out four rows,
each row with one card for $P$ and one for $Q$: T-T, T-F, F-T, F-F. To the right of each
row place a third card that you write yourself: T if you judge ``if $P$ then $Q$'' to be
true, F if false.

Work without looking at the table, then compare. Only one row gets the F card, namely the
row whose $P$ is T and whose $Q$ is F. Slide that card up and down; it remains
the only one. Now swap the P and Q cards in the T-F row so that it reads F-T, then evaluate
again: the card changes from F to T. That one-card move is the converse.
\end{tangan}
\section{Quantifiers and how to negate them}
\tingkat{basic}

Most mathematical sentences speak about many objects, and that multiplicity is
marked by quantifiers. There are two quantifiers, and both read order as part of the meaning.

\kamus{quantifier}{kuantor}{like two kinds of announcement at school: ``all pupils
must bring a book'' and ``some pupil brings a book''; the two sound similar
but demand very different proofs.}

We write down their meaning officially.

\begin{definisi}[Universal and existential quantifiers]
\index{quantifier}
Let $P(x)$ be a statement that depends on $x$ inside the set $A$.
\begin{itemize}
  \item $(\forall x \in A)\,P(x)$ means $P(x)$ is true for every $x \in A$.
  \item $(\exists x \in A)\,P(x)$ means there is at least one $x \in A$
        with $P(x)$ true.
\end{itemize}
The first statement is false if there is one $x \in A$ with $P(x)$ false; one such
object is called a \emph{counterexample}.
\end{definisi}

The existential quantifier demands at least one, not exactly one; adding a demand
that is not written in the definition is the most frequent source of error.

The order of two quantifiers of different kinds may not be swapped. Inside $\mathbb{R}$,
the statement
\[
  (\forall x)(\exists y)\; y > x
\]
is true, whereas
\[
  (\exists y)(\forall x)\; y > x
\]
is false. The proofs of both are short.

\begin{proposisi}
Inside $\mathbb{R}$ we have $(\forall x)(\exists y)\, y > x$, but we do not have
$(\exists y)(\forall x)\, y > x$.
\end{proposisi}
\par\smallskip\noindent\textit{Proof.} For the first part, take any $x \in \mathbb{R}$
and choose $y = x+1$. Then $y > x$, so $y$ is a witness for that $x$.
Since $x$ was arbitrary, every $x$ has a witness; the witness may differ for different
$x$.

For the second part, suppose there is $y_{0} \in \mathbb{R}$ with $y_{0} > x$ for
every $x \in \mathbb{R}$. Take $x = y_{0}+1$, which is also a real number. Then
$y_{0} > y_{0}+1$, that is $0 > 1$, which is impossible. So no such $y_{0}$
exists.
\hfill$\blacksquare${\parfillskip=0pt\par}

\index{negation}Negation works from the outside in: every quantifier changes kind, and the negation travels inward
until it sits on the innermost statement. For compound statements, De Morgan's laws
move the negation past conjunctions and disjunctions.

{\centering
\begin{tabular}{ll}
\hline
Statement & Its negation \\
\hline
$(\forall x \in A)\,P(x)$ & $(\exists x \in A)\,\neg P(x)$ \\
$(\exists x \in A)\,P(x)$ & $(\forall x \in A)\,\neg P(x)$ \\
$P \wedge Q$ & $\neg P \vee \neg Q$ \\
$P \vee Q$ & $\neg P \wedge \neg Q$ \\
$P \Rightarrow Q$ & $P \wedge \neg Q$ \\
\hline
\end{tabular}\par}

The last row is the one most often violated. Denying ``if it rains then the road
is wet'' is not ``if it rains then the road is dry'', but ``it rains and the road is not
wet''. The same mistake appears when people deny a definition: the verb inside it is
negated, the quantifier is left intact.

\begin{proposisi}
For every statement $P(x)$ and $Q(x)$ we have
\[
  \neg(\forall x)\bigl(P(x) \Rightarrow Q(x)\bigr)
  \;\equiv\; (\exists x)\bigl(P(x) \wedge \neg Q(x)\bigr).
\]
\end{proposisi}
\par\smallskip\noindent\textit{Proof.} A universal quantifier is negated into an existential
quantifier over the negation of the body:
\[
  \neg(\forall x)R(x) \;\equiv\; (\exists x)\neg R(x).
\]
The body $R(x)$ here is $P(x) \Rightarrow Q(x)$, whose negation is $P(x) \wedge
\neg Q(x)$ by the proposition of the previous part. So the negation has the form
$(\exists x)(P(x) \wedge \neg Q(x))$.
\hfill$\blacksquare${\parfillskip=0pt\par}

That form sets a clear task when one wants to reject a quantified statement: exhibit
one object that satisfies the antecedent and fails the consequent.

The next two negations are used repeatedly; the first is the material of Chapter~4.

\begin{contoh}[Negation of the definition of a limit]
The definition of a limit reads: the sequence $(a_{n})$ tends to $L$ if
\[
  (\forall \varepsilon > 0)(\exists N \in \mathbb{N})(\forall n \geq N)\;
  |a_{n} - L| < \varepsilon .
\]
Negating from the outside in gives
\[
  (\exists \varepsilon > 0)(\forall N \in \mathbb{N})(\exists n \geq N)\;
  |a_{n} - L| \geq \varepsilon .
\]
The reading of the second sentence: there is one tolerance $\varepsilon$ that no
$N$ can repair, because there is always a term after $N$ that deviates by at least
$\varepsilon$. To prove that $(a_{n})$ does not tend to $L$, we must find
that $\varepsilon$, then for every $N$ find $n \geq N$ that violates it; $n$ may
depend on $N$.
\end{contoh}

The second negation concerns the existence of primes. Euclid proved that for
every natural number $n$ there is a prime larger than $n$:
\[
  (\forall n \in \mathbb{N})(\exists p)\bigl(p \text{ prime} \wedge p > n\bigr).
\]
Its negation is
\[
  (\exists n \in \mathbb{N})(\forall p)\bigl(p \text{ prime} \Rightarrow p \leq n\bigr),
\]
that is: there is one bound $n$ such that all primes do not exceed $n$. Denying
Euclid's claim demands such a bound, not merely failing to find a large prime.

\begin{untukanak}
Picture a class with thirty pupils. The teacher writes on the board: ``all pupils
bring a pencil.'' To prove that sentence true, the teacher must check thirty desks; to
prove it false, one pupil without a pencil is enough.

Now the teacher writes another sentence: ``some pupil brings a pencil.'' This time
the situation is reversed. One pupil with a pencil is already enough to prove
that sentence. But to prove it false, the teacher must check thirty desks
and make sure there is not a single pencil.

The two sentences only swap a word at the front, but the work of checking them swaps
entirely. Mathematicians are careful with ``all'' and ``some'' because both
count how many things must be checked.

There is one more trap. If some pupil brings two pencils, is ``some pupil brings
a pencil'' still true? Of course it is true: that sentence demands at least one,
not exactly one. In this chapter we will learn why a demand that is not written may not
be added on one's own.
\end{untukanak}

\section{Reading a definition literally}
\tingkat{intermediate}

A mathematical definition is not a short explanation of a word; it is a contract that
holds in both directions. If a number is called odd when it has the form $2k+1$ for
some integer $k$, we are bound to two things: every number of the form $2k+1$
must be called odd, and every number called odd must have the form $2k+1$.

To read a definition means to read a single sentence of the form ``$x$ is called $B$ if and only
if $C(x)$'', even though the words ``only if'' are rarely written; the phrase ``is called'' already
contains both directions. The first direction grants permission to use the name $B$, the second grants permission
to use the condition $C(x)$, and reversing either one wrecks the whole argument.

\kamus{necessary condition}{syarat perlu}{like the condition for entering a building: without a membership
card we may not enter, but having a membership card does not yet mean we are already
seated inside.}

Compare two statements about an integer $n$: ``if $6$ divides $n$ then $3$
divides $n$'' is true, whereas ``if $3$ divides $n$ then $6$ divides $n$'' is false,
with counterexample $n=3$. So ``$3$ divides $n$'' is a necessary condition for ``$6$ divides
$n$'', but not a sufficient one. A good definition names both, since only
in that way can the name and the condition be exchanged.

\begin{bacakeras}
Read it aloud as a dialogue for two voices, alternating, with one breath at every full stop.

First voice: ``If it rains, the road is wet.''
Second voice: ``The road is wet. So it rains.''
First voice: ``Not necessarily. The road is wet when it is hosed down, too.''
Second voice: ``Still wet. Still raining.''
First voice: ``What I promised was one direction only. Rain first, then wet. The reverse direction was never promised.''

Repeat that dialogue once more, faster, until the ear catches the difference. The first
sentence puts the rain in front; the second moves the rain to the back. The sound turns
around, the promise does not.

The line to repeat three times out loud: ``if it rains then it is wet, not if
it is wet then it rains.''
\end{bacakeras}

Now consider the definition used most often in this book.

\begin{definisi}[Prime number]
\index{prime number}
An integer $p$ is called \emph{prime} if $p \geq 2$ and the only positive divisors
of $p$ are $1$ and $p$ itself.
\end{definisi}

In everyday words: a prime is larger than one and cannot be split
into a product of two smaller numbers; the two clauses are tested separately.

\begin{proposisi}
The number $1$ is not prime.
\end{proposisi}
\par\smallskip\noindent\textit{Proof.} The definition demands $p \geq 2$. For $p = 1$
that condition fails, so $1$ does not satisfy the definition, whatever we can
say about its divisors. Note that $1$ in fact passes the second clause:
the only positive divisors of $1$ are indeed $1$ and $1$ itself. So what sets it aside
is the first clause, not the second.
\hfill$\blacksquare${\parfillskip=0pt\par}

This kind of event explains why a definition must be read to the last clause.
A reader who skips the first clause will conclude that $1$ is prime, and
the uniqueness of the factorization $12 = 2^{2}\cdot 3$ in Chapter~28 will be lost, because $12$ may
also be written $1 \cdot 2^{2} \cdot 3$.

A definition also determines how to deny membership. To show that a
number is not prime, we must find a positive divisor that is neither $1$ nor
the number itself. Here is one example that often escapes a quick check.

\begin{proposisi}
The number $91$ is not prime. More generally, the statement ``every odd number
that is not divisible by $3$ and not divisible by $5$ is prime'' is false.
\end{proposisi}
\par\smallskip\noindent\textit{Proof.} In fact $91 = 7 \cdot 13$, and both factors
lie between $1$ and $91$, so $91$ has a positive divisor other than $1$ and
itself. The number $91$ is odd, is not divisible by $3$ because $9+1 = 10$ is not a multiple
of $3$, and is not divisible by $5$ because its last digit is neither $0$ nor $5$. So $91$
is a counterexample to that general statement. For comparison, $97$ is indeed prime; among the integers $2$ and $99$ there are
exactly $25$ primes, a count produced by the program in the lab section.
\hfill$\blacksquare${\parfillskip=0pt\par}

\begin{definisi}[Perfect square]
An integer $n$ is called a \emph{perfect square} if there is an integer $k$
with $n = k^{2}$.
\end{definisi}

By this definition, $0$ is a perfect square, because $0 = 0^{2}$, and so is $1$,
because $1 = 1^{2}$. Meanwhile $2$ is not a perfect square, and that statement needs
a proof, not an admission that we failed to find $k$.

\begin{proposisi}
The number $2$ is not a perfect square.
\end{proposisi}
\par\smallskip\noindent\textit{Proof.} Suppose $2 = k^{2}$ for some $k \in
\mathbb{Z}$. Since $k^{2} = 2 > 0$, we have $k \neq 0$. If $|k| = 1$ then
$k^{2} = 1 \neq 2$. If $|k| \geq 2$ then $k^{2} \geq 4 > 2$. Since $|k|$ is a
positive integer, only those two possibilities remain, and both fail. Hence there is no
such $k$.
\hfill$\blacksquare${\parfillskip=0pt\par}

\index{counterexample}A habit worth keeping: every time you meet a
definition, ask three questions. What are its clauses? Which objects satisfy
all of them? Which objects are often thought to satisfy them although they fail one clause?
The third question produces counterexamples, and a counterexample can be neither
argued with nor debated.

\section{Four patterns of proof}
\tingkat{intermediate}

A proof is a transport of truth: from statements we have already accepted to the
statement we want to accept, through steps that can each be checked.
The four patterns below cover most needs, and
Figure~\ref{fig:peta-tiga-bukti} maps the first three routes. The fifth common pattern,
proof by cases, is merely the first pattern applied to cases that exclude one another.

\begin{figure}[htbp]
\centering
\begin{tikzpicture}[font=\small,>=Stealth,
  jalur/.style={draw,rounded corners=2pt,align=center,inner sep=5pt,text width=4.6cm},
  ujung/.style={draw,rounded corners=2pt,align=center,inner sep=4pt,text width=2.2cm}]
  \node[jalur,fill=layerorange!12] (klaim) at (0,0) {claim: $P \Rightarrow Q$};
  \node[jalur,fill=layerteal!12] (lgs) at (6.1,2.4)
    {\textbf{direct}\\ unpack $P$, compute, out comes $Q$};
  \node[jalur,fill=layerindigo!10] (ktp) at (6.1,0)
    {\textbf{contraposition}\\ prove $\neg Q \Rightarrow \neg P$; the value is the same};
  \node[jalur,fill=layerplum!10] (ktr) at (6.1,-2.4)
    {\textbf{contradiction}\\ suppose $P$ true and $Q$ false, derive a clash};
  \node[ujung,fill=layergrey!10] (selesai) at (11.6,0)
    {no assignment\\brings it\\down};
  \draw[->,thick] (klaim.east) -- (lgs.west);
  \draw[->,thick] (klaim.east) -- (ktp.west);
  \draw[->,thick] (klaim.east) -- (ktr.west);
  \draw[->,thick] (lgs.east) -- (selesai.north west);
  \draw[->,thick] (ktp.east) -- (selesai.west);
  \draw[->,thick] (ktr.east) -- (selesai.south west);
\end{tikzpicture}
\caption{Road map of three forms of proof for an implication. All routes start from the same claim
and end at the fact that no assignment of values makes $P$ true while
$Q$ is false. A proof by contraposition is valid because $P \Rightarrow Q \equiv \neg Q \Rightarrow \neg P$.
The fourth form, a constructive proof with a witness, does not enter this map because it proves
a statement of the form $\exists$, not an implication.}
\label{fig:peta-tiga-bukti}
\end{figure}

\paragraph{Pattern one: direct proof.} We start from the hypothesis, unpack it
by definition, then compute until the conclusion appears. The next example is chosen
because its conclusion is stronger than the one people usually prove.

\begin{teorema}
For every odd integer $n$, the number $n^{2}-1$ is divisible by $8$.
\end{teorema}
\par\smallskip\noindent\textit{Proof.} Since $n$ is odd, $n = 2k+1$ for some
$k \in \mathbb{Z}$. Then
\[
  n^{2} - 1 = (2k+1)^{2} - 1 = 4k^{2} + 4k = 4k(k+1).
\]
Among two consecutive integers $k$ and $k+1$, exactly one is even,
so $k(k+1) = 2m$ for some $m \in \mathbb{Z}$. Substitution gives
$n^{2}-1 = 4 \cdot 2m = 8m$, which means $n^{2}-1$ is divisible by $8$.
\hfill$\blacksquare${\parfillskip=0pt\par}

An immediate consequence: the square of an odd number always leaves remainder $1$ when divided by $8$. For
$n = 1,3,5,7,9,11$ we obtain in turn $n^{2}-1 = 0, 8, 24, 48, 80, 120$, and
all six are divisible by $8$. That check is not a proof, but it makes sure that
the proof above contains no typo.

\begin{figure}[htbp]
\centering
\begin{tikzpicture}[font=\small,>=Stealth,
  tahap/.style={draw,rounded corners=2pt,align=left,inner sep=5pt,text width=6.6cm},
  nama/.style={font=\footnotesize\itshape,text=layergrey,align=right,text width=2.4cm}]
  \node[tahap,fill=layerteal!10] (asumsi) at (0,0)
    {$n$ is odd, so $n = 2k+1$ for some $k \in \mathbb{Z}$};
  \node[nama] at (-4.9,0) {assumption (given)};
  \node[tahap,fill=layerorange!10] (langkah) at (0,-2.1)
    {$n^{2}-1 = 4k^{2}+4k = 4k(k+1) = 8m$};
  \node[nama] at (-4.9,-2.1) {step (computed)};
  \node[tahap,fill=layerjade!10] (kesimpulan) at (0,-4.2)
    {$8$ divides $n^{2}-1$};
  \node[nama] at (-4.9,-4.2) {conclusion (required)};
  \draw[->,thick] (asumsi.south) -- (langkah.north);
  \draw[->,thick] (langkah.south) -- (kesimpulan.north);
\end{tikzpicture}
\caption{Anatomy of one direct proof of the statement ``if $n$ is odd then $8$ divides
$n^{2}-1$''. The top box is the only one that may be used without reason; the middle box holds
its reason, namely the definition of an odd number and algebra; the bottom box is what is required.
The arrows may only run downward.}
\label{fig:anatomi-bukti}
\end{figure}

Figure~\ref{fig:anatomi-bukti} names the three parts of that proof: assumption, step,
conclusion.

\paragraph{Pattern two: proof by contraposition.} We prove $\neg Q \Rightarrow \neg P$
instead of $P \Rightarrow Q$. By the corollary in the first section the two are equivalent,
so this is valid. The pattern is chosen when the negation of the conclusion is easier to hold than
its hypothesis.

\begin{proposisi}
For every integer $n$, if $5n+3$ is even then $n$ is odd.
\end{proposisi}
\par\smallskip\noindent\textit{Proof.} The contrapositive reads: if $n$ is even then
$5n+3$ is odd. Let $n = 2k$. Then $5n+3 = 10k+3 = 2(5k+1)+1$, which has the form two
times an integer plus one, hence is odd. Since the contrapositive is proved, the
original implication is proved.
\hfill$\blacksquare${\parfillskip=0pt\par}

A direct proof would demand dividing $5n+3 = 2m$ by $5$; not impossible,
but longer. The pattern of contraposition is chosen because it saves work.

Before the third pattern, we need one lemma about parity. This lemma will be used twice:
once here, once in the proof that the square root of two is irrational.

\begin{lemma}
\index{parity lemma}
For every integer $n$, if $n^{2}$ is even then $n$ is even.
\end{lemma}
\par\smallskip\noindent\textit{Proof.} We prove the contrapositive: if $n$ is odd
then $n^{2}$ is odd. By the first theorem of this section, for odd $n$ we have
$n^{2} = 8m+1$ for some $m$, so $n^{2}$ leaves remainder $1$ when divided by $2$, that is,
is odd.
\hfill$\blacksquare${\parfillskip=0pt\par}

\paragraph{Pattern three: proof by contradiction.} We suppose the conclusion false, then derive two statements that cannot be true
together. This pattern is used when the object being denied has no convenient form
to compute with, for instance when what is denied is the existence of something.

\begin{teorema}
\index{square root of two}
There is no rational number $r$ with $r^{2} = 2$.
\end{teorema}
\par\smallskip\noindent\textit{Proof.} Suppose there is $r \in \mathbb{Q}$ with
$r^{2} = 2$. Every rational number can be written as $p/q$ with $p,q$ integers,
$q \neq 0$, and $p,q$ having no common factor other than $1$; we
remove that common factor by dividing both, so a form without a common
factor can always be reached.

From $p^{2}/q^{2} = 2$ we obtain $p^{2} = 2q^{2}$, so $p^{2}$ is even. By
the lemma above, $p$ is even; write $p = 2m$ with $m \in \mathbb{Z}$. Substitution gives
$4m^{2} = 2q^{2}$, that is $q^{2} = 2m^{2}$, so $q^{2}$ is even and, once again
by the lemma, $q$ is even. Now $p$ and $q$ are both even, so both are divisible by
$2$, contradicting the choice of $p/q$ without a common factor. The supposition is false,
hence there is no rational $r$ with $r^{2} = 2$.
\hfill$\blacksquare${\parfillskip=0pt\par}

The contradiction is not between $r^{2}=2$ and ordinary arithmetic, but between ``$p,q$
both even'' and the opening promise ``$p,q$ without a common factor''. That is why the step
of removing common factors is not decoration; it is part of the hypothesis that is later violated.

\paragraph{Pattern four: constructive proof with a witness.} To prove a statement
of the form $\exists$, we point at one object and check that the object satisfies
the condition. The object pointed at is called a \emph{witness}.

\begin{teorema}
For every natural number $n$ there are $n$ consecutive integers that are all
not prime.
\end{teorema}
\par\smallskip\noindent\textit{Proof.} Take $N = (n+1)!$ and consider the $n$ numbers
\[
  N+2,\; N+3,\; \dots,\; N+(n+1).
\]
For every $k$ with $2 \leq k \leq n+1$, the number $k$ divides $(n+1)!$ because $k$
is one of its factors, so $k$ also divides $N+k$. Moreover
$N+k > k > 1$, so $N+k$ has a positive divisor $k$ that is neither $1$ nor
itself. By the definition, $N+k$ is not prime. We have pointed at the $n$ numbers
required, and there is no need to look for the smallest ones.
\hfill$\blacksquare${\parfillskip=0pt\par}

For $n = 5$ that proof points at $6! = 720$, then $722, 723, 724, 725, 726$. All five
indeed split: $722 = 2 \cdot 19^{2}$, $723 = 3 \cdot 241$, $724 = 2^{2} \cdot 181$,
$725 = 5^{2} \cdot 29$, and $726 = 2 \cdot 3 \cdot 11^{2}$. The number $721 = 7 \cdot 103$ splits as well, but it is not on our list.

A constructive proof does not promise the smallest or the fastest witness; what is judged is only
whether the object pointed at satisfies the condition. Denying a statement of the form $\exists$
demands showing that no object satisfies it, and that is where a contradiction
is needed.

\begin{resep}
Four standard steps to prove ``if $P$ then $Q$''.
\begin{enumerate}
  \item Write $P$ and $Q$ by their definitions, one line each.
  \item Open with ``suppose $P$ is true''. Only $P$ may be supposed; $Q$ may not
        yet be used.
  \item Derive $Q$ by substitution and algebra, one step per line, and name
        the definitions used.
  \item Close with a sentence that states $Q$ exactly as it appears in step 1.
\end{enumerate}
A warning about order: step 2 may not be swapped with step 3. Scratch work from the conclusion
backward is allowed, but the proof submitted must run from $P$ to $Q$; otherwise
what is proved is the converse.
\end{resep}
\section{Induction and strong induction}
\tingkat{intermediate}

Induction is a proof pattern for statements that come in a row, one for each natural
number. We set the first card upright and make sure that every card that falls knocks down
the next card; after both, no card may remain standing.

\kamus{induction}{induksi}{like a row of dominoes: it is enough to make sure the first card falls
and that every card knocks down the card in front of it, and then the whole row falls without needing
to be pushed one by one.}

We take the principle of induction as a basic property of $\mathbb{N} = \{1,2,3,\dots\}$, equivalent
to the statement that every non-empty subset of $\mathbb{N}$ has a
least member; nothing is proved circularly, we merely choose a starting point as in Chapter~3.

\begin{catatan}[Principle of induction]
Let $P(n)$ be a statement about natural numbers. If $P(1)$ is true and, for every
$k \in \mathbb{N}$, the truth of $P(k)$ forces the truth of $P(k+1)$, then $P(n)$ is true
for every $n \in \mathbb{N}$.
\end{catatan}

With that principle, we prove the formula that was already guessed in the opening of this section.

\begin{teorema}
For every natural number $n$ we have
\[
  1 + 3 + 5 + \cdots + (2n-1) = n^{2}.
\]
\end{teorema}
\par\smallskip\noindent\textit{Proof.} Base: for $n = 1$ the left side equals $1$ and
the right side equals $1^{2} = 1$. Step: suppose the formula is true for $n = k$, that is
$1+3+\cdots+(2k-1) = k^{2}$. The next term is $2(k+1)-1 = 2k+1$, so
\[
  1+3+\cdots+(2k-1)+(2k+1) = k^{2} + 2k + 1 = (k+1)^{2}.
\]
Hence the formula is true for $k+1$. By the principle of induction, the formula is true for every
$n \in \mathbb{N}$.
\hfill$\blacksquare${\parfillskip=0pt\par}

The second example has the form of an inequality. From $2^{k} \geq k+1$ and the condition $k \geq 0$
we obtain $2^{k+1} = 2 \cdot 2^{k} \geq 2(k+1) \geq k+2$, so with base $n = 0$
we have $2^{n} \geq n+1$ for every $n \geq 0$; the hypothesis and the side condition are used
at once.

Ordinary induction demands one step back. Sometimes the conclusion for $n$ follows only from
many earlier values at once, and there ordinary induction is not enough.

\begin{catatan}[Principle of strong induction]
If for every $n \in \mathbb{N}$ the truth of $P(1), P(2), \dots, P(n-1)$ forces
the truth of $P(n)$, then $P(n)$ is true for every $n \in \mathbb{N}$.
\end{catatan}

\kamus{strong induction}{induksi kuat}{like building a tower of blocks: to
place the tenth block we may use every block beneath it, not only
the ninth block.}

With that permission we can prove something ordinary induction cannot reach.

\begin{teorema}
Every integer $n \geq 2$ can be written as a product of prime
numbers.
\end{teorema}
\par\smallskip\noindent\textit{Proof.} We use strong induction on $n$. For $n = 2$:
the number $2$ is itself prime, so it is a product with one factor. Now take
$n \geq 3$ and suppose the statement true for all integers $m$ with
$2 \leq m < n$. If $n$ is prime, we are done. If $n$ is not prime, by the definition there is a
positive divisor $a$ with $a \neq 1$ and $a \neq n$; write $n = ab$ with
$b = n/a$. Since $a \neq 1$ and $a \neq n$, we have $2 \leq a \leq n-1$, and likewise
$2 \leq b \leq n-1$. The induction hypothesis applies to $a$ and to $b$, so both are
products of primes, and the product of both is $n$.
\hfill$\blacksquare${\parfillskip=0pt\par}

The strong induction hypothesis is used twice, on $a$ and on $b$. That is where ordinary induction
is awkward: both may be much smaller than $n-1$, so the statement for
$n-1$ alone does not cover them.

That theorem closes the promise of the previous section: a list of primes is never complete.

\begin{akibat}[Infinitely many primes]
\index{infinitely many primes}
For every finite list of primes $p_{1}, p_{2}, \dots, p_{k}$ there is
a prime that is not on that list.
\end{akibat}
\par\smallskip\noindent\textit{Proof.} Let $N = p_{1}p_{2}\cdots p_{k} + 1$. Since
$N \geq 2$, the previous theorem gives that $N$ is a product of primes;
take one prime factor $q$ of $N$. Suppose $q$ is on the list, say
$q = p_{i}$. Then $q$ divides the product $p_{1}\cdots p_{k}$, and since $q$ also
divides $N$, the number $q$ divides the difference $N - p_{1}\cdots p_{k} = 1$. No
prime divides $1$, so the supposition is false. Hence $q$ is not on
the list.
\hfill$\blacksquare${\parfillskip=0pt\par}

That proof does not prove that $N$ is prime, only that $N$ has a prime factor outside the
list: for $2,3,5,7$ we obtain $N = 211$, which happens to be prime, for $2,3,5,7,11,13$
we obtain $N = 30031 = 59 \cdot 509$. The proof is constructive; ``suppose there were only
finitely many primes'' is merely its rhetorical wrapping.

\begin{untukanak}
Picture a row of a thousand boxes, and we want to prove that every box contains
a marble. We do not need to open a thousand boxes. It is enough to open the first box, then prove
one thing: if a box contains a marble, the box beside it contains a marble too.
If both things are true, a thousand boxes are checked without being opened one by one.

This is induction, and it feels like magic. The key lies in the second promise: if the relation
between boxes leaks in even one place, the whole row collapses.

Now picture the first box turning out empty. The relation between boxes remains true, but
it does not help, because there is nothing first to fall. This is the mistake that
happens most often.

In this chapter we will learn why the two steps may not be swapped, and why
checking many boxes is no substitute for either of them.
\end{untukanak}

\section{When a proof fails}
\tingkat{advanced}

Back to the blackboard from the opening of the chapter. The first forty values are prime, and at
$n = 40$ the pattern breaks:
\[
  40^{2} + 40 + 41 = 1600 + 40 + 41 = 1681 = 41^{2}.
\]
That failure is no coincidence, and its shape can be explained without a computer. For
every integer $p \geq 2$, the substitution $n = p-1$ gives
\[
  (p-1)^{2} + (p-1) + p = p^{2} - 2p + 1 + p - 1 + p = p^{2},
\]
whereas the substitution $n = p$ gives $p^{2} + p + p = p(p+2)$. So every number
$p \geq 2$ produces two points of failure, namely $n = p-1$ and $n = p$. With $p = 41$
both produce $1681 = 41^{2}$ and $1763 = 41 \cdot 43$.

The same pattern appears in smaller polynomials. For $n^{2}+n+11$ the first point of failure
is $n = 10$, which gives $121 = 11^{2}$; for $n^{2}+n+17$ the first point of failure
is $n = 16$, which gives $289 = 17^{2}$; for $n^{2}+n+29$ the failure comes at $n = 2$,
because $35 = 5 \cdot 7$ is not prime. The larger the number in the constant position,
the longer the clean initial stretch; that is the whole secret of the reputation of $41$.

That failure yields three lessons that are often confused.

The first lesson is asymmetric. One counterexample brings down a statement
of the form ``for all'', whereas forty successful examples prove nothing
about the forty-first value. A computer check is therefore a tool for finding conjectures,
not a tool for proving.

The second lesson concerns the location of the defect. The first defect that most often escapes notice appears
in an induction proof. Consider the following false claim: for every $n \geq 1$
we have $1+3+\cdots+(2n-1) = n^{2}+1$. Its induction step, when inspected, turns out to be
valid:
\[
  \bigl(n^{2}+1\bigr) + (2n+1) = n^{2} + 2n + 2 = (n+1)^{2} + 1 .
\]
So from the truth of the formula at $n$ the truth at $n+1$ does indeed flow. What fails
is the base: for $n = 1$ the left side equals $1$ while the right side equals $2$.
The induction step moves truth, it does not create it.

The third lesson concerns the limits of checking. The usual divisibility test: $m > 1$ is prime if
it is not divisible by any prime smaller than $\sqrt{m}$. For $m = 1681$ we get $\sqrt{1681} = 41$, and the primes
smaller than $41$ are $2, 3, 5, 7, 11, 13, 17, 19, 23, 29, 31, 37$; not
one of them divides $1681$. By that rule $1681$ passes as prime. The error lies
in the word ``smaller'': the correct rule uses ``smaller than or equal to'',
because $1681 = 41 \cdot 41$ and the divisor is exactly the size of the root. One loose inequality sign
produces one false prime, precisely at $n = 40$.

Three other defects recur in longer writing. The defect of circular reasoning:
``since $n^{2}+n+41$ is prime for all $n$, it is also prime for $n = 40$'' uses
what must be proved as its reason. The defect of dividing by zero: from $a = b$ people
conclude $a+b = b$ after dividing by $a-b$, although $a-b = 0$. The defect of swapping
direction: from $P \Rightarrow Q$ and $Q$ people conclude $P$; in the truth table,
the third row is where that error hides.

All four defects share one root: a step is used outside its condition.
A useful habit is not to memorize patterns but to check the condition of every step:
in a contradiction the supposition must be the negation of the conclusion, in induction there is a base and a
step, in a division the divisor is not zero.

\begin{tantangan}
Claim: for every integer $k \geq 0$, the number $2^{2^{k}}+1$ is prime. For
$k = 0,1,2,3,4$ the values are $3$, $5$, $17$, $257$, $65537$, and all five are indeed prime.
Find the smallest $k$ that brings the claim down, then exhibit a divisor; the claim is
false, and the error can be shown without a computer.

Hint 1: five values are not a proof. Check $k = 5$; the number is $4294967297$, and its smallest
divisor is far smaller than its square root.

Hint 2: that divisor has the form $5 \cdot 2^{7}+1 = 641$, and
$4294967297 = 641 \cdot 6700417$.

If you get stuck, go back to the section ``When a proof fails''.
\end{tantangan}

\section{Worked examples}
\tingkat{intermediate}

The three computations below support this chapter and can be repeated on paper. The first is
the sum of the first odd numbers, term by term from $n = 1$.

{\centering
\begin{tabular}{cccc}
\hline
$n$ & terms & sum & $n^{2}$ \\
\hline
$1$ & $1$ & $1$ & $1$ \\
$2$ & $1+3$ & $4$ & $4$ \\
$3$ & $1+3+5$ & $9$ & $9$ \\
$4$ & $1+3+5+7$ & $16$ & $16$ \\
$5$ & $1+3+5+7+9$ & $25$ & $25$ \\
$6$ & $1+3+5+7+9+11$ & $36$ & $36$ \\
\hline
\end{tabular}\par}

The continuation: $36+13 = 49 = 7^{2}$, then $49+15 = 64$, $64+17 = 81$, $81+19 = 100$.
The difference of two consecutive squares is indeed $(n+1)^{2}-n^{2} = 2n+1$, and that is where induction
takes its stand.

The second is a primality test on $n^{2}+n+41$, only at the interesting values.

{\centering
\begin{tabular}{ccc}
\hline
$n$ & $n^{2}+n+41$ & remark \\
\hline
$0$ & $41$ & prime \\
$1$ & $43$ & prime \\
$39$ & $1601$ & prime \\
$40$ & $1681 = 41^{2}$ & not prime \\
$41$ & $1763 = 41 \cdot 43$ & not prime \\
\hline
\end{tabular}\par}

An exhaustive check for $n = 0$ up to $n = 40$ gives $40$ prime values out of $41$;
the only one that fails is $n = 40$, and $\sqrt{1681} = 41$ exactly, so the decisive
divisor sits right at the boundary of the check.

The third is the contraposition check on the truth table, written out in full.

{\centering
\begin{tabular}{cc|cc|c}
\hline
$P$ & $Q$ & $P \Rightarrow Q$ & $\neg Q \Rightarrow \neg P$ & equal \\
\hline
T & T & T & T & yes \\
T & F & F & F & yes \\
F & T & T & T & yes \\
F & F & T & T & yes \\
\hline
\end{tabular}\par}

The four rows give equality: the contrapositive pattern is not a rhetorical trick but a replacement
of form that does not change the value. The column $Q \Rightarrow P$ would differ in the third row,
and that difference is what makes the converse unusable.

\kertas{Compute $1+3+\cdots+(2n-1)$ for $n = 11$ and $n = 12$ by starting from
$100 = 10^{2}$, that is by adding the next two odd numbers one at a time. Then
compute $n^{2}+n+17$ for $n = 16$ and check by division whether $289$ equals
$17^{2}$. Finally, find the first point of failure of $n^{2}+n+11$ by testing
$n = 0, 1, 2, \dots$ and explain why that point can be predicted from the shape of
its polynomial.}

\section{Computational lab}
\tingkat{advanced}

The following script uses only the Python standard library and can be rerun without
any installation; every output line is genuine output.

\begin{lab}
\begin{verbatim}
from itertools import product

def prime(n):
    if n < 2:
        return False
    d = 2
    while d * d <= n:
        if n % d == 0:
            return False
        d += 1
    return True

pairs = product([False, True], repeat=2)
print("contraposition:", all((not P or Q) == (Q or not P)
                            for P, Q in pairs))

D = range(1, 11)
print("y>x:", sum(y > x for x in D for y in D), "of 100")
print("x without a witness:", [x for x in D if all(y <= x for y in D)])

values = [(n, n*n + n + 41) for n in range(41)]
print("prime:", sum(prime(v) for _, v in values), "of 41")
print("failed:", [(n, v, int(v**0.5)) for n, v in values if not prime(v)])
print("41*41, 41*43:", 41*41, 41*43)

residues = {(n*n - 1) % 8 for n in range(1, 40, 2)}
print("residues mod 8:", residues)

step = all(n*n + 1 + 2*n + 1 == (n+1)**2 + 1 for n in range(1, 500))
print("false step:", step)
\end{verbatim}
\end{lab}

The output obtained:

\begin{verbatim}
contraposition: True
y>x: 45 of 100
x without a witness: [10]
prime: 40 of 41
failed: [(40, 1681, 41)]
41*41, 41*43: 1681 1763
residues mod 8: {0}
false step: True
\end{verbatim}

The first line checks all four rows of the table at once. The second and third lines compute
what is awkward to do on paper: out of a hundred pairs $(x,y)$ in
$\{1,\dots,10\}$, $45$ satisfy $y > x$, and the only $x$ without a witness is $10$;
that is the failure of $(\forall x)(\exists y)$ on a finite set.

The fourth and fifth lines give $40$ primes out of $41$ values, failing only at
$(40,1681,41)$; that last $41$ is $\sqrt{1681}$ truncated. The sixth line
confirms $1681 = 41^{2}$ and $1763 = 41 \cdot 43$. The seventh states that every odd
$n$ below $40$ gives residue $0$ modulo $8$, in agreement with the theorem of the first pattern. The last
line is a warning: the induction step of that false claim is valid for the first five hundred values,
so a numerical check does not find its defect.

\begin{untukdosen}
A teaching order that has worked: start from $n = 40$ and $1681$, not from the
definitions. After the students are annoyed that forty checks were in vain, then
enter the truth table, with emphasis on the third and fourth rows. Continue to
the negation of quantifiers through the definition of the limit of a sequence, then to literal definitions with the example
$91 = 7 \cdot 13$; close with induction, a failing base, and the lab.

Two misconceptions are most frequent. First, an implication is read as cause and effect, so a row
with a false antecedent is judged false and a vacuous case proof is rejected. Second, induction
is thought finished once the step is proved, and strong induction is thought to be ordinary induction
repeated; as a result the base is not checked and the hypothesis is used without naming
its range.

Blackboard activity: split the class into two groups. The first group writes a quantified claim,
the second group writes its negation, then each names one object that
refutes the other side's claim. The game stops when a negation can no longer be refuted.

Students are ready to go on if they can negate the definition of a limit without opening their notes
and can explain why the base in induction may not be skipped.
\end{untukdosen}

\section{Summary}
\tingkat{basic}

\begin{ringkasan}
\begin{itemize}
  \item A statement is true or false; $P \Rightarrow Q$ is false only when $P$ is
        true and $Q$ is false.
  \item $P \Rightarrow Q \equiv \neg Q \Rightarrow \neg P$, but the converse is not
        equivalent; $\neg(P \Rightarrow Q) \equiv P \wedge \neg Q$.
  \item Negation works from the outside in: every quantifier changes kind, so
        denying ``for all'' means pointing at one counterexample.
  \item A definition holds in both directions: the name may be used only if every clause
        is satisfied, and every clause is used to check it.
  \item Four proof patterns: direct, contraposition, contradiction, constructive. Choose
        the shortest, not the most magnificent.
  \item Induction demands a base and a step; strong induction demands all earlier
        values. A step without a base proves nothing.
\end{itemize}
\end{ringkasan}

\section{Exercises}
\tingkat{intermediate}

\begin{latihan}
\begin{enumerate}
  \item Write the negation of the following statement, then turn the result into a
        plain sentence that states what must be shown in order to
        refute it: ``for every $\varepsilon > 0$ there is a natural number $N$
        such that for every $n \geq N$ we have $|a_{n} - L| < \varepsilon$''.

  \item Write the negation of ``every prime larger than $2$ is
        odd'', then name one object that refutes the original statement.
        Do the same for ``there is a natural number $n$ such that
        $n^{2}+n+41$ is not prime''.

  \item Prove that for every integer $n$, the number $n^{2}+n$ is always
        even. Do it in two ways: by cases on parity, and by contraposition.

  \item Prove by induction that for every natural number $n$ we have
        \[
          1^{2} + 2^{2} + \cdots + n^{2} = \frac{n(n+1)(2n+1)}{6}.
        \]
        Check the formula for $n = 1, 2, 3$ before writing the proof.

  \item Prove by strong induction that every natural number $n$ can be written
        in the form $n = 2^{k}m$ with $k \geq 0$ an integer and $m$ an odd
        number. Show where the strong induction hypothesis is used.

  \item Judge the following claims: true, false, or not enough evidence.
        \begin{enumerate}
          \item If $P(1)$ is true and $P(k) \Rightarrow P(k+1)$ for every
                $k \in \mathbb{N}$, then $P(n)$ is true for every $n \in
                \mathbb{N}$.
          \item If $P(n)$ is true for $n = 1, 2, \dots, 10^{6}$, then $P(n)$ is true
                for every $n \in \mathbb{N}$.
          \item If $n^{2}+n+41$ is prime for $n = 0, 1, \dots, 39$, then the claim
                is true for every $n \geq 0$.
          \item If $3$ divides $n$, then $6$ divides $n$.
        \end{enumerate}
        For every item that is false, give one concrete counterexample.

  \item Using Python, count how many primes occur among $n^{2}+n+41$ for
        $n = 0$ up to $n = 100$, then determine all $n$ in that range whose
        value is not prime. Check whether each point of failure has the form
        $n = p-1$ or $n = p$ for some prime $p$.

  \item For $n = 1$ up to $n = 12$, compute $2^{n}$ and $n^{2}$, then determine
        the smallest natural number $n_{0}$ such that $2^{n} > n^{2}$ holds for every
        $n \geq n_{0}$. Prove that claim by induction, with base $n_{0}$.

  \item Prove that the statement ``there is an integer $n > 1$ with
        $n^{2}-1$ prime'' is false. First state the factorization of $n^{2}-1$
        for $n > 1$, then explain why the factors cannot equal $1$.

  \item Show that the statement ``if $n^{2}$ is divisible by $4$ then $n$ is divisible
        by $4$'' is false, then repair the conclusion into a true claim and
        prove it.
\end{enumerate}

\par\smallskip\noindent\textit{Hints for the three hardest items.} Item 5: separate the
case of odd $n$, which is settled at once with $k=0$ and $m=n$, from the case of even $n$,
which uses the induction hypothesis on $n/2$. Item 7: also compute $\sqrt{n^{2}+n+41}$
for every point of failure, and compare it with the smallest divisor. Item 9: $n^{2}-1 = (n-1)(n+1)$; for $n > 1$ both factors lie between $1$
and $n^{2}-1$.
\end{latihan}

\begin{untukpeneliti}
What is genuinely proved in this chapter is three classical facts that survive without a computer:
the irrationality of $\sqrt{2}$ by a parity proof, the infinitude of primes through
Euclid's construction, and the uniqueness of the form of every natural number as a product of primes.
What is merely popularized is the story that $41$ is ``almost'' a prime formula; what happens
is only a long initial stretch, and its length can be predicted from $n^{2}+n+p$.

Primary sources: Abbott, \emph{Understanding Analysis}, Chapter~1 Section~1.1 (irrationality
of $\sqrt{2}$) and Section~1.2 (introduction, including induction); Avigad and Massot,
\emph{Mathematics in Lean}, Chapter~3 (logic) and Chapter~5 Sections~5.1 and 5.2
(irrational roots, induction and recursion); Euclid, \emph{Elements}, Book~IX
Proposition~20 for infinitely many primes.

Open questions: there is no polynomial of positive degree with integer coefficients
that is prime for all $n \geq 0$ (a one-line proof: $f(0)$ divides
$f(t\,f(0))$), but how long a prime stretch a polynomial of high degree can reach
is still being studied, and the distribution of primes is tied to the Riemann hypothesis,
which remains unproved.

Limits of applicability: all results in this chapter are elementary and do not depend on the choice of set
axioms. Formal verification in Lean (Chapter~28) checks the validity of every deductive step,
but it does not check whether the formal statement proved is the one its author
intended.
\end{untukpeneliti}

\peta{Chapter 1 of \emph{Understanding Analysis} (Abbott), in particular Section 1.1 on the
irrationality of $\sqrt{2}$ and Section 1.2 on the introduction; Chapter 3 (logic), Chapter 5
Sections 5.1 and 5.2 (irrational roots as well as induction and recursion) of
\emph{Mathematics in Lean} (Avigad and Massot); and Book IX Proposition 20 of
\emph{Elements} (Euclid) for infinitely many primes.}

\chapter{Sets, Relations, Functions, and Cardinality}

A security tester receives a model specification: a classifier with thirty binary
features, each taking only the value zero or one. To test exhaustively, it must try
every possible input. Each feature doubles the number of possibilities, so the answer
is $2^{30} = 1{,}073{,}741{,}824$. At one million inputs per second, that test takes
$1073.74$ seconds, about $17.9$ minutes. Add ten more features: the number becomes
$2^{40} = 1{,}099{,}511{,}627{,}776$, and a machine a thousand times faster needs
$1099.51$ seconds, about $18.3$ minutes. We multiply the problem by a thousand and
speed the machine up a thousandfold, and we are back at the starting point.

The same difficulty appears more sharply once features may take an arbitrary real
value. The number of inputs stops being "many" and starts being "infinitely many",
and the question changes: can two infinite sets still be compared in size? Instinct
says there are more rational numbers than natural numbers, since between any two
natural numbers there always lie infinitely many rationals. That instinct is wrong,
and this chapter shows why.

The previous chapter gave us a language: quantifiers, implication, proof patterns.
What was missing is the object being talked about. This chapter installs its basic
vocabulary, namely sets and their operations, relations, functions, and one way of
comparing the size of sets that does not depend on the ability to count. That
vocabulary is used throughout the book, from subspaces in Chapter 6 to formal proofs
in Chapter 28.

\begin{untukanak}
Picture two baskets of marbles, one red and one blue. We want to know which basket
holds more, but we are not allowed to count. Take one red marble, pair it with one
blue marble, set both aside, and repeat until one basket is empty: if both empty at
the same moment they hold the same number, if blue is left over then blue holds more.
This method never mentions a number; it only pairs.

Now suppose the baskets never end, so that we can never finish moving marbles. Can
pairing still be used? It can, and that is the surprise of this chapter: with the
right pairing rule, all the counting numbers can be paired off with all the even
numbers, with nothing left over on either side, and pairs of numbers can be paired
with a single number. We will understand why pairing is stronger than counting, and
why some endless sets still cannot be paired with the counting numbers, however hard
we try.
\end{untukanak}

\section{Operations on sets and De Morgan's laws}
\tingkat{basic}

A single set is of no use until we can build new sets from it. Four ways of building
are defined by the same pattern: state a membership condition, and the new set is the
collection of all objects satisfying it. Let $X$ be a universal set and
$A, B \subseteq X$; from Chapter 1 we recall that $A \subseteq B$ means every member
of $A$ is also a member of $B$, and that equality of sets is proved by two inclusions.

\begin{definisi}[Union, intersection, difference, complement]
The union $A \cup B$ is the set of all $x \in X$ that are members of $A$ or members
of $B$. The intersection $A \cap B$ is the set of all $x \in X$ that are members of
$A$ and at the same time members of $B$. The difference $A \setminus B$ is the set of
all $x \in X$ that are members of $A$ but not members of $B$. The complement
$A^{c} = X \setminus A$ is the set of all $x \in X$ that are not members of $A$.
\end{definisi}

The word "or" in a union is inclusive: $x$ may lie in both sets, and that is the
condition under which $A \cap B \subseteq A \cup B$ is always true, as Hoffman and
Kunze note in the Appendix. The empty set $\varnothing$ is needed so that the
intersection of two disjoint sets is still a set.

\begin{proposisi}[De Morgan's laws]\index{De Morgan's laws}
For any $A, B \subseteq X$,
\[
  X \setminus (A \cup B) = (X \setminus A) \cap (X \setminus B),
  \qquad
  X \setminus (A \cap B) = (X \setminus A) \cup (X \setminus B).
\]
\end{proposisi}
\par\smallskip\noindent\textit{Proof.} Take any $x \in X$. By the definition of
complement and union,
\[
  x \in X \setminus (A \cup B)
  \iff x \in X \ \text{and}\ \text{not}\ (x \in A \ \text{or}\ x \in B).
\]
By De Morgan's law for statements from Chapter 1, the negation of "at least one" is
"both fail", so
\[
  x \in X \setminus (A \cup B)
  \iff x \in X \ \text{and}\ x \notin A \ \text{and}\ x \notin B
  \iff x \in (X \setminus A) \cap (X \setminus B).
\]
Since $x$ was arbitrary, the two sets are equal. The second equality is obtained by
replacing every "and" by "or" in the same chain.
\hfill$\blacksquare${\parfillskip=0pt\par}

The law holds for two sets; what we need later is the version for any number of sets.

\begin{definisi}[Indexed family]
Let $I$ be an index set and let $A_i \subseteq X$ be given for each $i \in I$.
\begin{align*}
  \bigcup_{i \in I} A_i &= \{\, x \in X : x \in A_i \ \text{for some}\ i \in I \,\}, \\
  \bigcap_{i \in I} A_i &= \{\, x \in X : x \in A_i \ \text{for every}\ i \in I \,\}.
\end{align*}
\end{definisi}

The general version reads $X \setminus \bigcup_{i} A_i = \bigcap_{i} (X \setminus A_i)$
and $X \setminus \bigcap_{i} A_i = \bigcup_{i} (X \setminus A_i)$. Its proof adds no
new idea: only "for some" turns into "for every" when the negation is pulled inside,
exactly as in Chapter 1.

Numbers make the definition checkable. For $X = \{1,\dots,10\}$,
$A = \{1,\dots,5\}$, and $B = \{4,\dots,8\}$, Figure~\ref{fig:venn} shows the three
operations on the same sets: the union has $3 + 2 + 3 = 8$ members, the intersection
has $2$, and the complement of $B$ has $3 + 2 = 5$.

\begin{figure}[htbp]
\centering
\begin{tikzpicture}[scale=0.82,every node/.style={font=\footnotesize}]
  \begin{scope}
    \begin{scope}
      \clip (-0.5,0) circle (0.85);
      \fill[layerteal!30] (-1.55,-1.25) rectangle (1.55,1.25);
    \end{scope}
    \begin{scope}
      \clip (0.5,0) circle (0.85);
      \fill[layerteal!30] (-1.55,-1.25) rectangle (1.55,1.25);
    \end{scope}
    \draw[ink,thick] (-0.5,0) circle (0.85);
    \draw[ink,thick] (0.5,0) circle (0.85);
    \draw[ink,thick] (-1.55,-1.25) rectangle (1.55,1.25);
    \node at (-1.0,0) {3};
    \node at (0,0) {2};
    \node at (1.0,0) {3};
    \node at (0,-1.0) {2};
    \node[above] at (-0.5,0.86) {$A$};
    \node[above] at (0.5,0.86) {$B$};
    \node[below] at (0,-1.25) {$A \cup B$: 8};
  \end{scope}
  \begin{scope}[xshift=3.9cm]
    \begin{scope}
      \clip (-0.5,0) circle (0.85);
      \clip (0.5,0) circle (0.85);
      \fill[accent!35] (-1.55,-1.25) rectangle (1.55,1.25);
    \end{scope}
    \draw[ink,thick] (-0.5,0) circle (0.85);
    \draw[ink,thick] (0.5,0) circle (0.85);
    \draw[ink,thick] (-1.55,-1.25) rectangle (1.55,1.25);
    \node at (-1.0,0) {3};
    \node at (0,0) {2};
    \node at (1.0,0) {3};
    \node at (0,-1.0) {2};
    \node[above] at (-0.5,0.86) {$A$};
    \node[above] at (0.5,0.86) {$B$};
    \node[below] at (0,-1.25) {$A \cap B$: 2};
  \end{scope}
  \begin{scope}[xshift=7.8cm]
    \fill[layergrey!22] (-1.55,-1.25) rectangle (1.55,1.25);
    \fill[white] (0.5,0) circle (0.85);
    \draw[ink,thick] (-0.5,0) circle (0.85);
    \draw[ink,thick] (0.5,0) circle (0.85);
    \draw[ink,thick] (-1.55,-1.25) rectangle (1.55,1.25);
    \node at (-1.0,0) {3};
    \node[ink!40] at (0,0) {2};
    \node[ink!40] at (1.0,0) {3};
    \node at (0,-1.0) {2};
    \node[above] at (-0.5,0.86) {$A$};
    \node[above] at (0.5,0.86) {$B$};
    \node[below] at (0,-1.25) {$B^{c}$: 5};
  \end{scope}
\end{tikzpicture}
\caption{Three Venn diagrams over the same sets, with the number of members in each
region. Left: $A \cup B$ has $3 + 2 + 3 = 8$ members; middle: $A \cap B$ has $2$
members; right: $B^{c}$ has $3 + 2 = 5$ members, namely the shaded region, while the
members inside $B$ are printed faded.}
\label{fig:venn}
\end{figure}

\begin{tangan}
Materials: twenty used name cards, two pieces of string two meters long, an open floor.

Form two string loops overlapping by about a quarter of their width. Write the names
of five friends who like coffee and five friends who like tea. Put the coffee cards
in the first loop, the tea cards in the second, then move the card of a friend who
likes both into the overlap: that is the intersection. Pick one name who likes
neither and put that card outside both loops.

Count three numbers: cards inside both loops, cards in the overlap, and cards outside
the second loop. What must become visible: the number of cards in the union is not the
number of coffee cards plus the number of tea cards, because the cards in the overlap
are counted twice, and the excess is exactly the intersection.
\end{tangan}

The set we most often build from $A$ is the set of all its subsets, and its size grows
far faster than $A$ itself.

\begin{definisi}[Power set]\index{power set}
The power set $\mathcal{P}(A)$ is the set of all subsets of $A$, that is,
$\mathcal{P}(A) = \{\, B : B \subseteq A \,\}$.
\end{definisi}

\kamus{power set}{himpunan kuasa}{A list of every way of choosing part of the members of a bag, including choosing nothing and choosing everything.}

For example, $\mathcal{P}(\varnothing) = \{\varnothing\}$ has one member, and
$\mathcal{P}(\{1,2\}) = \{\varnothing, \{1\}, \{2\}, \{1,2\}\}$ has four.

\begin{proposisi}
If $A$ is finite with $|A| = n$, then $|\mathcal{P}(A)| = 2^{n}$.
\end{proposisi}
\par\smallskip\noindent\textit{Proof.} For $B \subseteq A$ define the characteristic
function $\chi_B : A \to \{0,1\}$ by $\chi_B(a) = 1$ if $a \in B$ and
$\chi_B(a) = 0$ if $a \notin B$. The map $B \mapsto \chi_B$ is a bijection between
$\mathcal{P}(A)$ and the set of all functions $A \to \{0,1\}$: from $\chi$ we recover
$B = \{a : \chi(a) = 1\}$, and two different subsets give two different functions.
Hence $|\mathcal{P}(A)|$ equals the number of functions $A \to \{0,1\}$.

We count those functions by induction on $n$. For $n = 0$ there is only the empty
function, and $2^{0} = 1$. Suppose $|A| = n+1$, choose $a_0 \in A$, and write
$A' = A \setminus \{a_0\}$. Every $f : A \to \{0,1\}$ is determined uniquely by the
restriction $f|_{A'}$ and the value $f(a_0)$; the induction hypothesis gives $2^{n}$
choices of restriction and $2$ choices of value, $2^{n+1}$ in all.
\hfill$\blacksquare${\parfillskip=0pt\par}

That proof never mentions combinations: the bijection moves the count, and induction
splits it into two free choices.

\begin{bacakeras}
Read in pairs, one voice for member and one voice for subset. Start slowly, then
speed up.

Member: One orange goes into the basket. I am inside.
Subset: One basket goes into the store. I am inside.
Member: I am cargo.
Subset: I am container.
Member: I am pointed at one by one.
Subset: I am proved all at once.
Member: When I go in, I am a single mark.
Subset: When I go in, I am a single proof.
Member: I am small and many.
Subset: I am large and sometimes only one.

Repeat until the last line comes out without thinking. The line repeated three times
together: cargo goes into bag, bag goes into bag, and the two use different marks.
\end{bacakeras}

\section{Cartesian products and relations}
\tingkat{basic}

Order matters. The set $\{1,2\}$ equals $\{2,1\}$, but the coordinates of a point in
the plane may not be swapped without moving the point. We therefore need an object
that remembers order, and its definition must be built from sets alone.

\begin{definisi}[Ordered pair and Cartesian product]\index{Cartesian product}
The ordered pair $(a,b)$ is defined as $\{\{a\}, \{a,b\}\}$. The Cartesian product
$A \times B$ is the set of all ordered pairs $(a,b)$ with $a \in A$ and $b \in B$.
\end{definisi}

The definition looks strange, but only one property is needed, namely
$(a,b) = (c,d)$ if and only if $a = c$ and $b = d$. Checking that property requires
four cases, and we leave it as Exercise 4.

The size of a product is easy to guess, and its proof uses a cut along the first
coordinate.

\begin{proposisi}
If $A$ and $B$ are finite, then $|A \times B| = |A| \cdot |B|$.
\end{proposisi}
\par\smallskip\noindent\textit{Proof.} For each $a \in A$ the set
$\{a\} \times B$ has $|B|$ members, since $b \mapsto (a,b)$ is a bijection between $B$
and $\{a\} \times B$. The family $\{\{a\} \times B\}_{a \in A}$ consists of mutually
disjoint sets, because $(a,b) = (a',b')$ forces $a = a'$. Their union is $A \times B$,
so $|A \times B|$ is the sum of $|A|$ terms each equal to $|B|$.
\hfill$\blacksquare${\parfillskip=0pt\par}

Now the object that gives this chapter its title. A relation is merely a statement
that may be true or false for each pair of objects, and so it can be stored as the set
of pairs for which the statement is true.

\begin{definisi}[Relation]\index{relation}
A relation from $A$ to $B$ is any $R \subseteq A \times B$; when $A = B$ we call it a
relation on $A$. We write $a \mathrel{R} b$ for $(a,b) \in R$. The domain of $R$ is
$\{\, a \in A : (a,b) \in R\ \text{for some}\ b \in B \,\}$ and its range is
$\{\, b \in B : (a,b) \in R\ \text{for some}\ a \in A \,\}$.
\end{definisi}

Since a relation on $A$ is a subset of $A \times A$, the number of relations on a set
with $n$ members is $2^{n^{2}}$. For $n = 3$ the number is $2^{9} = 512$ and for
$n = 4$ it becomes $2^{16} = 65{,}536$; both are checked by the script in the lab
section.

Four properties filter them into useful objects. A relation $R$ on $A$ is called
reflexive if $a \mathrel{R} a$ for every $a \in A$; symmetric if
$a \mathrel{R} b$ implies $b \mathrel{R} a$; transitive if $a \mathrel{R} b$ and
$b \mathrel{R} c$ imply $a \mathrel{R} c$; and antisymmetric if
$a \mathrel{R} b$ together with $b \mathrel{R} a$ is possible only when $a = b$.

\begin{definisi}[Partial order and equivalence relation]\index{equivalence relation}
A relation on $A$ that is reflexive, antisymmetric, and transitive is called a partial
order. A relation on $A$ that is reflexive, symmetric, and transitive is called an
equivalence relation.
\end{definisi}

The relation $\le$ on $\mathbb{R}$ is a partial order, as is the relation "divides" on
the positive natural numbers, and $\subseteq$ on $\mathcal{P}(A)$ gives an example
that will later become the framework of subspaces in Chapter 6. The word "partial" is
used because two members need not be comparable: on $\mathcal{P}(\{1,2,3\})$ the sets
$\{1\}$ and $\{2\}$ are not contained in one another. An order in which every two
members are comparable is called a total order.

\section{Equivalence relations and partitions}
\tingkat{intermediate}

An equivalence relation cuts a set into groups that do not overlap. Before reaching
the theorem, we need a name for the groups.

\begin{definisi}[Equivalence class]\index{equivalence class}
Let $\sim$ be an equivalence relation on $A$ and $x \in A$. The equivalence class of
$x$ is $[x] = \{\, y \in A : y \sim x \,\}$.
\end{definisi}

Since $x \sim x$, no class is empty. The more interesting property: two classes cannot
overlap even slightly.

\begin{lemma}
Let $\sim$ be an equivalence relation on $A$ and $x, y \in A$. Then (i) $x \in [x]$;
(ii) $[x] = [y]$ if and only if $x \sim y$; and (iii) if
$[x] \cap [y] \neq \varnothing$, then $[x] = [y]$.
\end{lemma}
\par\smallskip\noindent\textit{Proof.} Part (i) is the statement $x \sim x$, true
because the relation is reflexive. For (ii), suppose $[x] = [y]$. By (i),
$y \in [y] = [x]$, so $y \sim x$ and by symmetry $x \sim y$. Conversely, suppose
$x \sim y$ and take $z \in [x]$, meaning $z \sim x$. From $z \sim x$ and
$x \sim y$ we get $z \sim y$, so $z \in [y]$; this gives
$[x] \subseteq [y]$, and the argument with the roles exchanged gives
$[y] \subseteq [x]$. For (iii), suppose $z \in [x] \cap [y]$. Then $z \sim x$ and
$z \sim y$; from $z \sim x$ and symmetry we get $x \sim z$, and together with
$z \sim y$ and transitivity we get $x \sim y$, so $[x] = [y]$ by (ii).
\hfill$\blacksquare${\parfillskip=0pt\par}

Consequently the set of all classes divides $A$ into pieces that are mutually disjoint
and cover all of $A$.

\begin{definisi}[Partition]\index{partition}
A partition of a set $A$ is a family $\mathcal{B}$ of nonempty subsets of $A$ such
that every $x \in A$ lies in exactly one member of $\mathcal{B}$.
\end{definisi}

\begin{teorema}
Every equivalence relation on $A$ determines a partition of $A$, every partition of
$A$ determines an equivalence relation on $A$, and the two processes invert each
other.
\end{teorema}
\par\smallskip\noindent\textit{Proof.} Let $\sim$ be an equivalence relation and
$\mathcal{B} = \{\, [x] : x \in A \,\}$. By the Lemma, each $[x]$ is nonempty and two
of its members are either disjoint or identical; every $x$ lies in $[x]$. Hence
$\mathcal{B}$ is a partition.

Conversely, let $\mathcal{B}$ be a partition of $A$ and set $x \sim y$ when $x$ and
$y$ lie in the same member of $\mathcal{B}$. Reflexivity and symmetry are clear. For
transitivity, suppose $x \sim y$ and $y \sim z$, and let $B$ be the member of
$\mathcal{B}$ containing $y$; then $x \in B$ and $z \in B$, so $x \sim z$.

The two processes invert each other: from $\sim$ we obtain the partition into classes,
and that partition returns the original $\sim$ because $x \sim y$ exactly when
$x \in [y]$; from a partition we obtain $\sim$, and the class $[x]$ is the member of
$\mathcal{B}$ containing $x$.
\hfill$\blacksquare${\parfillskip=0pt\par}

The number of classes is called the index. The set of all classes is written
$A/{\sim}$ and called the quotient set, while $q : A \to A/{\sim}$ with $q(x) = [x]$
is called the quotient map. Hoffman and Kunze note that every equivalence relation can
be viewed as coming from a function: two members are equivalent exactly when that
function gives them the same value.

\begin{contoh}
On the integers, set $x \equiv y \pmod{n}$ when $n$ divides $x-y$. This relation is
reflexive because $n \mid 0$, symmetric because $n \mid (x-y)$ implies
$n \mid (y-x)$, and transitive because $n \mid (x-y)$ together with
$n \mid (y-z)$ implies $n \mid (x-z)$. Its classes are the sets of numbers with the
same remainder, and there are exactly $n$ classes, namely $[0], \dots, [n-1]$, since
every integer can be written $x = qn + r$ with $0 \le r \le n-1$ uniquely.
\end{contoh}

An example closer to our own work: for $f : A \to B$, set $x \sim y$ when
$f(x) = f(y)$. The class $[x]$ is the set of inputs with the same output; in machine
learning that is "one prediction file", and the number of classes equals the number of
distinct labels.

\begin{untukanak}
The wardrobe at home is usually not sorted by any official rule. White shirts pile up
on one shelf, printed shirts on another, sports shirts in a third basket. Nobody
writes the rule down, yet every shirt sits in exactly one place.

That is what an equivalence relation does: it decides when two items count as
roommates, then groups all the items. The result is a tidy division, every item lands
in exactly one group, and different groups share no members.

Now the question that makes one curious. What if one shirt belongs to the white group
under one rule, but to the sports group under another rule? Is the grouping still
tidy? The tidiness does not depend on the rule chosen, but on three properties: every
item is a roommate of itself, roommatehood works both ways, and roommatehood can be
passed along through a third item.
\end{untukanak}

\section{Functions: injective, surjective, bijective}
\tingkat{intermediate}

We now fix the meaning of the word "function", because this chapter depends on the
fine distinction between codomain and range.

\begin{definisi}[Function]\index{injective function}
A function from $A$ to $B$ is a rule $f$ that pairs every $x \in A$ with exactly one
$f(x) \in B$. The set $A$ is called the domain, $B$ is called the codomain, and
$f(A) = \{\, f(x) : x \in A \,\}$ is called the range.
\end{definisi}

Codomain and range are different, and the difference matters. The function
$f : \mathbb{R} \to \mathbb{R}$ with $f(x) = x^{2}$ has codomain $\mathbb{R}$ but
range $[0,\infty)$. Hoffman and Kunze use three names interchangeably for the same
object: function, mapping, transformation.

For $S \subseteq B$ we define the preimage
$f^{-1}(S) = \{\, x \in A : f(x) \in S \,\}$. The notation $f^{-1}$ here does not
presuppose that $f$ has an inverse: a preimage is always defined, an inverse need not
exist.

\begin{definisi}[Injective, surjective, bijective]\index{surjective function}\index{bijection}
A function $f : A \to B$ is called injective if $f(x) = f(y)$ is possible only when
$x = y$; surjective if $f(A) = B$; and bijective if both.
\end{definisi}

\kamus{bijection}{bijeksi}{A perfect matching: every item in the first box has exactly one partner in the second box, with nothing left over on either side.}

Figure~\ref{fig:bijeksi} compares the three on two columns of nodes: the left panel
leaves no node over on either side, the middle panel leaves $d$ over on the right, and
the right panel sends two left nodes to one right node.

\begin{figure}[htbp]
\centering
\begin{tikzpicture}[scale=0.95,every node/.style={font=\footnotesize},>=Stealth,
  kotak/.style={draw=ink,thick,circle,inner sep=1.2pt,minimum size=4.4mm}]
  \begin{scope}
    \foreach \i/\y in {1/1.0,2/0,3/-1.0} {
      \node[kotak] (a\i) at (0,\y) {\i};
    }
    \node[kotak] (ba) at (1.9,1.0) {$a$};
    \node[kotak] (bb) at (1.9,0) {$b$};
    \node[kotak] (bc) at (1.9,-1.0) {$c$};
    \draw[->,thick,layerjade] (a1) -- (ba);
    \draw[->,thick,layerjade] (a2) -- (bb);
    \draw[->,thick,layerjade] (a3) -- (bc);
    \node[below] at (0.95,-1.75) {bijective};
  \end{scope}
  \begin{scope}[xshift=3.4cm]
    \foreach \i/\y in {1/1.2,2/0.4,3/-0.4} {
      \node[kotak] (c\i) at (0,\y) {\i};
    }
    \node[kotak] (da) at (1.9,1.2) {$a$};
    \node[kotak] (db) at (1.9,0.4) {$b$};
    \node[kotak] (dc) at (1.9,-0.4) {$c$};
    \node[kotak] (dd) at (1.9,-1.2) {$d$};
    \draw[->,thick,layerindigo] (c1) -- (da);
    \draw[->,thick,layerindigo] (c2) -- (db);
    \draw[->,thick,layerindigo] (c3) -- (dc);
    \node[below] at (0.95,-1.75) {injective only};
  \end{scope}
  \begin{scope}[xshift=6.9cm]
    \foreach \i/\y in {1/1.2,2/0,3/-1.2} {
      \node[kotak] (e\i) at (0,\y) {\i};
    }
    \node[kotak] (fa) at (1.9,0.8) {$a$};
    \node[kotak] (fb) at (1.9,-0.8) {$b$};
    \draw[->,thick,layerbrown] (e1) -- (fa);
    \draw[->,thick,layerbrown] (e2) -- (fa);
    \draw[->,thick,layerbrown] (e3) -- (fb);
    \node[below] at (0.95,-1.75) {surjective only};
  \end{scope}
\end{tikzpicture}
\caption{Three mappings between two columns of nodes. Left: every node is used exactly
once, hence bijective. Middle: the right column has four members and the node $d$
receives no incoming arrow, hence injective but not surjective. Right: two left nodes
go to $a$, hence surjective but not injective.}
\label{fig:bijeksi}
\end{figure}

The word "bijective" is only a short name for the combination of the two previous
properties. Its power appears in the following proposition: a bijection is the exact
condition for a function to be invertible.

\begin{proposisi}
$f : A \to B$ is bijective if and only if there is $g : B \to A$ with
$g \circ f = \mathrm{id}_A$ and $f \circ g = \mathrm{id}_B$. If $g$ exists, it is
unique.
\end{proposisi}
\par\smallskip\noindent\textit{Proof.} Suppose $g$ satisfies both equalities. If
$f(x) = f(y)$ then $x = g(f(x)) = g(f(y)) = y$, so $f$ is injective. For any
$b \in B$ we have $b = f(g(b))$, so $f$ is surjective.

Conversely, suppose $f$ is bijective. For each $b \in B$ there is exactly one
$x \in A$ with $f(x) = b$: existence from surjectivity, uniqueness from injectivity.
Define $g(b) = x$ for that $x$, then $g(f(x)) = x$ and $f(g(b)) = b$.

For uniqueness, suppose $g$ and $h$ satisfy both equalities. Then
\[
  g = g \circ \mathrm{id}_B = g \circ (f \circ h) = (g \circ f) \circ h
  = \mathrm{id}_A \circ h = h,
\]
which uses only the associativity of composition and the two equalities.
\hfill$\blacksquare${\parfillskip=0pt\par}

Composition behaves well with respect to all three properties: if $f$ and $g$ are
injective then $g \circ f$ is injective, and similarly for surjective and bijective.
The converse fails: with $f : \{1\} \to \{1,2\}$, $f(1) = 1$, and
$g : \{1,2\} \to \{1\}$, $g(1) = g(2) = 1$, the composition $g \circ f$ is bijective,
even though $g$ is not injective and $f$ is not surjective.

On finite sets the three properties lock together, and that is why "one-to-one" and
"onto" can be distinguished only on infinite sets.

\begin{lemma}
Let $A$ and $B$ be finite, $f : A \to B$, and $C \subseteq B$. Then
$|f(A)| \le |A|$ with equality if and only if $f$ is injective, and
$|C| \le |B|$ with equality if and only if $C = B$.
\end{lemma}
\par\smallskip\noindent\textit{Proof.} We prove the second part by induction on
$n = |B|$. For $n = 0$ we have $B = \varnothing$, hence $C = B$. Assume the claim for
sets with $n$ members and suppose $|B| = n+1$; choose $b \in B$ and write
$B' = B \setminus \{b\}$.

If $b \notin C$, then $C \subseteq B'$ and the induction hypothesis gives
$|C| \le |B'| = n < n+1$, with equality impossible. If $b \in C$, write
$C' = C \setminus \{b\} \subseteq B'$. Then $|C| = |C'| + 1 \le |B'| + 1 = |B|$;
equality holds exactly when $|C'| = |B'|$, which by the induction hypothesis means
$C' = B'$, that is $C = B$.

The first part is obtained by the same induction on $|A|$: write
$A = A' \cup \{a\}$ with $a \notin A'$ and observe that
$f(A) = f(A') \cup \{f(a)\}$, so $|f(A)| \le |A'| + 1 = |A|$. Equality demands
$f(a) \notin f(A')$ and $|f(A')| = |A'|$; the latter means $f|_{A'}$ is injective by
the induction hypothesis, and the former means $f$ does not repeat a value at $a$.
\hfill$\blacksquare${\parfillskip=0pt\par}

\begin{proposisi}
Let $A$ and $B$ be finite with $|A| = |B|$, and let $f : A \to B$. Then $f$
injective, $f$ surjective, and $f$ bijective are three equivalent statements.
\end{proposisi}
\par\smallskip\noindent\textit{Proof.} If $f$ is injective, the Lemma gives
$|f(A)| = |A| = |B|$; since $f(A) \subseteq B$ and both have the same size, the Lemma
gives $f(A) = B$, so $f$ is surjective. If $f$ is surjective, then $f(A) = B$ so that
$|f(A)| = |B| = |A|$, and the Lemma forces $f$ to be injective. Since bijective means
both properties at once, the equivalence of the three is complete.
\hfill$\blacksquare${\parfillskip=0pt\par}

A warning to carry into the next chapter: that proposition uses the word "finite"
essentially. On $\mathbb{N}$ the function $f(x) = 2x$ is injective but not
surjective, and the function mapping $1$ to $1$ and $n \ge 2$ to $n-1$ is surjective
but not injective.

\section{Finite, countable, and uncountable sets}
\tingkat{advanced}

Now we return to the opening question. To compare two sets without counting we pair
their members, and all we need to require is that the pairing uses an invertible
function.

\begin{definisi}[Equinumerous]
The sets $A$ and $B$ are equinumerous, written $A \sim B$, if there is a bijection
$A \to B$.
\end{definisi}

The relation $\sim$ is reflexive (the identity is a bijection), symmetric (the inverse
of a bijection is a bijection), and transitive (the composition of two bijections is a
bijection). So it behaves as an equivalence relation, but not on "the set of all sets",
because that collection is not a set; this is discussed in the box for researchers.

\begin{definisi}[Finite, countable, uncountable]\index{countable set}
A set $A$ is finite if $A = \varnothing$ or there is $n \in \mathbb{N}$ and a
bijection $A \to \{1,2,\dots,n\}$; in the second case we write $|A| = n$. A set $A$
is countable if $A \sim \mathbb{N}$. An infinite set that is not countable is called
uncountable.
\end{definisi}

\kamus{countable}{terhitung}{All its members can be given serial numbers with none skipped, as when every passenger receives a distinct ticket number.}

The name "countable" is apt: if $A \sim \mathbb{N}$, the bijection is a list. We can
name the first member, the second, and so on, and every member appears exactly once.
What is surprising is how many sets can be listed that way. Pairs of numbers, for
instance.

\begin{proposisi}[Cantor pairing]\index{Cantor pairing}
$\mathbb{N} \times \mathbb{N}$ is countable.
\end{proposisi}
\par\smallskip\noindent\textit{Proof.} Define
$\pi(m,n) = \frac{(m+n)(m+n+1)}{2} + n$ for $m, n \in \mathbb{N} =
\{0,1,2,\dots\}$. Write $T_s = \frac{s(s+1)}{2}$, the number of pairs with
$m + n < s$. For a pair with $m + n = s$ we have $0 \le n \le s$, so
$\pi(m,n) = T_s + n$ takes the values $T_s, \dots, T_s + s$ exactly once each. Since
$T_s + s = T_{s+1} - 1$, the block of values for diagonal $s$ is
$\{T_s, \dots, T_{s+1}-1\}$. The blocks join without gaps and without overlap
starting from $T_0 = 0$. Hence every counting number lies in exactly one block, and
within its block comes from exactly one pair; therefore $\pi$ is bijective.
\hfill$\blacksquare${\parfillskip=0pt\par}

The inverse is explicit: $s = \lfloor (\sqrt{8q+1}-1)/2 \rfloor$, $n = q - T_s$, and
$m = s - n$. That square root is the one place where machine arithmetic can deceive
us, and the example is in the lab section.

Figure~\ref{fig:kisi} draws that grid: the cell in row $m$ and column $n$ is labelled
$\pi(m,n)$, and the color marks the diagonal $m + n = s$, so that reading by color
gives $0$ through $14$ without gaps and without repetition. The right panel applies
the same reading order to a list of four real numbers, with row $d$ flipping the
diagonal digits so that $d$ is not in the list.

\begin{figure}[htbp]
\centering
\begin{tikzpicture}[scale=0.95,every node/.style={font=\footnotesize}]
  \foreach \s/\warna in {0/layerteal,1/layerorange,2/layerindigo,3/layerjade,4/layerplum} {
    \foreach \n in {0,...,\s} {
      \pgfmathtruncatemacro{\m}{\s-\n}
      \pgfmathtruncatemacro{\v}{int(\s*(\s+1)/2+\n)}
      \draw[fill=\warna!30,draw=ink!45] (\n*0.86,-\m*0.72) rectangle ++(0.8,0.66);
      \node at (\n*0.86+0.4,-\m*0.72+0.33) {\v};
    }
  }
  \draw[->,thick] (4.35,0.33) -- (4.95,0.33);
  \node[right] at (4.95,0.33) {$n$};
  \draw[->,thick] (-0.45,-0.05) -- (-0.45,-3.0) node[below] {$m$};
  \begin{scope}[xshift=6.6cm,yshift=0.45cm]
    \foreach \i/\j/\b in {1/1/0,1/2/1,1/3/1,1/4/0, 2/1/1,2/2/0,2/3/0,2/4/1,
                          3/1/0,3/2/0,3/3/1,3/4/0, 4/1/1,4/2/1,4/3/0,4/4/0} {
      \ifnum\i=\j
        \fill[accent!28] (\j*0.52-0.52,-\i*0.42) rectangle ++(0.5,0.38);
      \fi
      \draw[ink!45] (\j*0.52-0.52,-\i*0.42) rectangle ++(0.5,0.38);
    }
    \draw[accent,thick] (0.25,-0.23) -- (1.79,-1.49);
    \foreach \i/\j/\b in {1/1/0,1/2/1,1/3/1,1/4/0, 2/1/1,2/2/0,2/3/0,2/4/1,
                          3/1/0,3/2/0,3/3/1,3/4/0, 4/1/1,4/2/1,4/3/0,4/4/0} {
      \node at (\j*0.52-0.27,-\i*0.42+0.19) {\b};
    }
    \foreach \i in {1,...,4} {
      \node[left] at (-0.62,-\i*0.42+0.19) {$x_{\i}$};
    }
    \foreach \j/\b in {1/1,2/1,3/0,4/1} {
      \fill[layergrey!25] (\j*0.52-0.52,-2.34) rectangle ++(0.5,0.38);
      \draw[ink!45] (\j*0.52-0.52,-2.34) rectangle ++(0.5,0.38);
      \node at (\j*0.52-0.27,-2.34+0.19) {\b};
    }
    \node[left] at (-0.10,-2.15) {$d$};
    \node[below] at (0.52,-2.42) {digit flipped};
  \end{scope}
\end{tikzpicture}
\caption{Left: the grid of pairs $(m,n)$ with $m + n \le 4$, each cell labelled
$\pi(m,n)$, with color marking the diagonal $m + n = s$; the first five diagonals
carry the values $0$ through $14$. Right: the digit table of four real numbers, with
the diagonal cells colored and row $d$ flipping the diagonal digits $0,0,1,0$ into
$1,1,0,1$.}
\label{fig:kisi}
\end{figure}

Since every positive rational is a pair of natural numbers already in lowest terms,
the set of rationals is countable too. All we must guard against is listing a number
twice.

\begin{proposisi}
The set $\mathbb{Q}_{>0}$ of positive rational numbers is countable.
\end{proposisi}
\par\smallskip\noindent\textit{Proof.} Every $q \in \mathbb{Q}_{>0}$ can be written
uniquely as $q = p/r$ with $p, r$ natural numbers and $\gcd(p,r) = 1$. Define the
height of $q$ as $p + r$. For each $k \ge 2$ there are only finitely many rationals of
height $k$, since $p$ and $r$ lie between $1$ and $k-1$. The following list therefore
enumerates $\mathbb{Q}_{>0}$ without repetition: for $k = 2, 3, 4, \dots$ write down
all $p/r$ of height $k$ with $\gcd(p,r) = 1$, increasing in $p$. Every positive
rational appears exactly once, namely at level $k = p + r$ in its unique reduced form,
so that list is a bijection $\mathbb{N} \to \mathbb{Q}_{>0}$.
\hfill$\blacksquare${\parfillskip=0pt\par}

\begin{akibat}
$\mathbb{Q}$ and $\mathbb{Z}$ are countable.
\end{akibat}
\par\smallskip\noindent\textit{Proof.} The list $0, 1, -1, 2, -2, \dots$ is a
bijection $\mathbb{N} \to \mathbb{Z}$. For $\mathbb{Q}$, enumerate each nonzero
rational as $\pm p/r$ in reduced form by $(p+r)$ with alternating signs, then place
$0$ at the front; one extra member does not destroy countability.
\hfill$\blacksquare${\parfillskip=0pt\par}

So $\mathbb{N}$, $\mathbb{Z}$, and $\mathbb{Q}$ are equinumerous, even though
$\mathbb{Q}$ feels far more crowded. Density is not a measure of size; the measure is
whether or not a list exists.

\begin{tantangan}
A hotel on the seafront has rooms numbered $1, 2, 3, \dots$ without end, and that
night every room holds exactly one guest. A traveler arrives without a reservation.
The receptionist says the hotel is full; the manager says there is still a room for
him without sending a single guest home.

Show a rule for moving guests that keeps the manager's promise, then explain why that
rule is a bijection between the old guests plus one new guest and all the rooms. Why
can the same hotel not accommodate as many new guests as there are real numbers?

Hint 1: do not look for an empty room, because there is none. Create an empty room by
moving its occupant.

Hint 2: shift the occupant of room $n$ to room $n+1$ for every $n$, then put the
traveler in room $1$; check that every room holds exactly one person and that no
guest is lost.

If you get stuck, return to the section Finite, countable, and uncountable sets, in
particular the definition of equinumerous.
\end{tantangan}

Cantor showed that $\mathcal{P}(A)$ is always larger than $A$, by a trick that does not
depend on the size of $A$.

\begin{teorema}[Cantor]\index{Cantor diagonal}
For every set $A$ there is no surjective function $A \to \mathcal{P}(A)$.
\end{teorema}
\par\smallskip\noindent\textit{Proof.} Let $f : A \to \mathcal{P}(A)$ be arbitrary
and define $D = \{\, a \in A : a \notin f(a) \,\}$. Then $D \subseteq A$, so
$D \in \mathcal{P}(A)$. Suppose $D = f(a_0)$ for some $a_0 \in A$. If
$a_0 \in D$, then by the definition of $D$ we have $a_0 \notin f(a_0) = D$, a
contradiction. If $a_0 \notin D$, then $a_0 \notin f(a_0)$, which is precisely the
membership condition for $D$, so $a_0 \in D$, a contradiction again. Hence $D$ is not
in the range of $f$.
\hfill$\blacksquare${\parfillskip=0pt\par}

The set $D$ is formed by asking, for each $a$, whether $a$ is a member of the set
paired with it, then flipping the answer; that flip is what makes $D$ differ from
every value of $f$. For $A = \mathbb{N}$ the trick appears as reading the diagonal of
a table with every entry changed, which is why it is called the diagonal argument.

\begin{akibat}
$\mathcal{P}(\mathbb{N})$ is uncountable, and $\mathbb{R}$ is uncountable.
\end{akibat}
\par\smallskip\noindent\textit{Proof.} By Cantor's Theorem with $A = \mathbb{N}$,
there is no surjection $\mathbb{N} \to \mathcal{P}(\mathbb{N})$, so
$\mathcal{P}(\mathbb{N})$ is uncountable.

For $\mathbb{R}$ it suffices to show that $[0,1]$ is uncountable. Every $x \in [0,1]$
has a binary expansion $x = \sum_{k \ge 1} b_k 2^{-k}$ with $b_k \in \{0,1\}$; its
existence is proved in Chapter 3 and Chapter 5. The map from the set of binary
sequences to $[0,1]$ is therefore surjective, and the set of binary sequences is
equinumerous with $\mathcal{P}(\mathbb{N})$ because a binary sequence is a function
$\mathbb{N} \to \{0,1\}$. If $[0,1]$ were countable, its preimage in the set of binary
sequences would be countable as well, contradicting the previous paragraph.
\hfill$\blacksquare${\parfillskip=0pt\par}

One note of honesty: binary expansions are not unique, since
$0.0111\dots = 0.1000\dots$ in base two. The argument above is undisturbed because
it uses only the surjectivity, but $|\mathbb{R}| = |\mathcal{P}(\mathbb{N})|$ demands
the following tool.

\begin{proposisi}[Cantor--Bernstein--Schroeder]
If there is an injective function $A \to B$ and an injective function $B \to A$, then
$A \sim B$.
\end{proposisi}
\par\smallskip\noindent\textit{Proof sketch.} Let $f : A \to B$ and $g : B \to A$ be
injective. Let $A_0$ be the set of points of $A$ that have no predecessor in $B$, then
$A_{k+1} = g(f(A_k))$, and write $A^{+}$ for the union of all $A_k$ with $k$ even.
Define $h(x) = f(x)$ if $x \in A^{+}$ and $h(x) = g^{-1}(x)$ if $x \notin A^{+}$. On
the first part $h$ is injective because $f$ is injective, on the second because
$g^{-1}$ is injective, and the images of the two parts are separated by the way
$A^{+}$ is built, so $h$ is injective. For surjectivity of the image, take
$y \in B$. If $g(y) \in A^{+}$, then $g(y)$ lies in some $A_k$ with $k$ even so that
$g(y) = g(f(x))$ for $x \in A_{k-1}$ and $h(x) = y$; if $g(y) \notin A^{+}$, then
$h(g(y)) = y$. Hence $h$ is surjective.

The bookkeeping of chain depth is the delicate part, and its details are skipped here;
the full formal version is in \emph{Mathematics in Lean}, section 4.3.
\hfill$\blacksquare${\parfillskip=0pt\par}

With that tool we may conclude $|\mathbb{R}| = |\mathcal{P}(\mathbb{N})|$: an
injection from $\mathcal{P}(\mathbb{N})$ to $\mathbb{R}$ is given by binary expansion,
and the injection the other way by the tangent function. The size of $\mathbb{R}$ is
usually written $2^{\aleph_0}$, with $\aleph_0$ the size of $\mathbb{N}$.

\begin{untukanak}
Picture the ticket machine at the entrance to an exhibition. Every visitor receives one
numbered ticket: 1, 2, 3, and so on, no number is used twice, and every possible
number will eventually be given out. A collection of things that can be ticketed like
that is called countable.

It looks as though every endless collection can be ticketed. But not so. Some
collections have so many members that however many tickets we print, there is always a
member left without one. The ticket machine is not broken; that collection really is
bigger than the supply of numbers.

The question that makes one curious: how can we know that a collection cannot be
ticketed, when not all methods have been tried? The answer is Cantor's flipping trick.
Each time someone claims to have ticketed all the members, we build one new member
deliberately different from every ticketed member; that member is certainly left out.
\end{untukanak}

\section{Inequalities, absolute value, and upper bounds}
\tingkat{intermediate}

The last section installs two tools for Chapter 3. The first measures the distance
between two numbers regardless of direction, namely the absolute value.

\begin{definisi}[Absolute value]\index{absolute value}
For $x \in \mathbb{R}$, $|x| = x$ if $x \ge 0$ and $|x| = -x$ if $x < 0$.
\end{definisi}

A route through a third point is never shorter than the direct route.

\begin{proposisi}[Triangle inequality]
For every $x, y \in \mathbb{R}$, $|x + y| \le |x| + |y|$ and
$\bigl| |x| - |y| \bigr| \le |x - y|$.
\end{proposisi}
\par\smallskip\noindent\textit{Proof.} Note that $-|x| \le x \le |x|$ and
$-|y| \le y \le |y|$. Adding the two gives
$-(|x|+|y|) \le x + y \le |x|+|y|$, which by the definition of absolute value means
$|x+y| \le |x|+|y|$. For the second, write $x = (x-y) + y$ and apply the first result:
$|x| \le |x-y| + |y|$, so $|x| - |y| \le |x-y|$. Exchanging the roles of $x$ and $y$
gives $|y| - |x| \le |x-y|$, hence $\bigl| |x| - |y| \bigr| \le |x-y|$.
\hfill$\blacksquare${\parfillskip=0pt\par}

Now bounds. A set of numbers can be bounded above without having a largest member, and
precisely that case is the most important one.

\begin{definisi}[Upper bound and supremum]\index{supremum}
Let $S \subseteq \mathbb{R}$ be nonempty. A number $u$ is called an upper bound of $S$
if $s \le u$ for every $s \in S$; if one exists, $S$ is called bounded above. A number
$u$ is called the supremum of $S$, written $u = \sup S$, if $u$ is an upper bound of
$S$ and $u \le v$ for every upper bound $v$ of $S$.
\end{definisi}

\kamus{supremum}{batas atas terkecil}{The tightest fence: the highest boundary that cannot be lowered any further without leaving a member above it.}

The supremum, when it exists, is unique, since two least upper bounds dominate each
other. What the definition does not guarantee is existence:
$\{x \in \mathbb{Q} : x^{2} < 2\}$ is bounded above in $\mathbb{Q}$ but has no
supremum there.

\begin{catatan}[Completeness axiom]
Every nonempty set $S \subseteq \mathbb{R}$ that is bounded above has a supremum in
$\mathbb{R}$. This axiom is not proved here; its construction together with its
consequences is the content of Chapter 3, including the epsilon characterization that
turns a statement about a least bound into a statement about $\varepsilon$, and the
example of a set that has a supremum without a maximum.
\end{catatan}

On a finite set the supremum and the maximum coincide, and swapping the two often
slips through on small examples.

\section{Worked examples}
\tingkat{intermediate}

The following four computations are done by hand down to the number, then checked
again with a script. All of them are small, so arithmetic errors show up immediately.

\begin{resep}
For two finite sets $A$ and $B$:
\begin{enumerate}
  \item Compute $|A|$, the number of members of $A$.
  \item Compute $|B|$ the same way.
  \item Compute $|A \cap B|$, the members appearing in both lists.
  \item Add and then subtract: $|A \cup B| = |A| + |B| - |A \cap B|$.
\end{enumerate}
A warning about order: step 3 may not be postponed until after step 4. If the addition
is done first, the members of the intersection are counted twice and the result
overshoots by exactly $|A \cap B|$. For $A = \{1,\dots,5\}$ and $B = \{4,\dots,8\}$:
$5 + 5 - 2 = 8$, the same number as in the Venn diagram above.
\end{resep}

First, the power set. For $A = \{1,2,3\}$ the subsets have sizes $0,1,2,3$ with counts
$1, 3, 3, 1$, in total $1 + 3 + 3 + 1 = 8 = 2^{3}$. The row $n = 4$ gives
$1+4+6+4+1 = 16$, $n = 2$ gives $1+2+1 = 4$, and $n = 1$ gives $1+1 = 2$: the same
Pascal triangle as in the Proposition about $2^{n}$.

Second, the Cantor pairing. With $\pi(m,n) = T_{m+n} + n$ and $T_0 = 0$, $T_1 = 1$,
$T_2 = 3$, $T_3 = 6$: diagonal $s = 0$ contains $(0,0) \mapsto 0$; $s = 1$ contains
$(1,0) \mapsto 1$ and $(0,1) \mapsto 2$; $s = 2$ contains $(2,0) \mapsto 3$,
$(1,1) \mapsto 4$, $(0,2) \mapsto 5$; $s = 3$ contains $(3,0) \mapsto 6$,
$(2,1) \mapsto 7$, $(1,2) \mapsto 8$, $(0,3) \mapsto 9$. Those ten values are exactly
$0, \dots, 9$ without repetition and without gaps.

The opposite direction is easier to get wrong. Take $q = 1000$. We look for $s$ with
$T_s \le 1000 < T_{s+1}$. Since $T_{44} = 44 \cdot 45/2 = 990$ and $T_{45} = 1035$, we
get $s = 44$, then $n = 1000 - 990 = 10$ and $m = 44 - 10 = 34$. Check:
$\pi(34,10) = 44 \cdot 45/2 + 10 = 990 + 10 = 1000$, correct.

Third, the rational numbers. By the height rule, level $k = 2$ contains $1/1$;
$k = 3$ contains $1/2, 2/1$; $k = 5$ contains $1/4, 2/3, 3/2, 4/1$. The number of terms
at each level is the value of Euler's totient function, $1, 2, 2, 4, 2, 6$ for
$k = 2, \dots, 7$. The cumulative count up to $k = 7$ is $1+2+2+4+2+6 = 17$, so $6/1$
occupies position 17. Since the cumulative count up to $k = 9$ is $27$ and level
$k = 10$ contains $1/9, 3/7, 7/3, 9/1$, the number $7/3$ occupies position 30.

Fourth, the number of relations. On a set with $3$ members there are $2^{9} = 512$
relations, and only $5$ are equivalence relations: $1$ with one class, $3$ with two
classes, $1$ with three classes. On a set with $4$ members there are $2^{16} = 65{,}536$
relations and $15$ that are equivalence relations. Those four numbers appear again as
the last four lines of the script output; the sequence $1, 2, 5, 15, 52$ is known as
the Bell numbers.

\kertas{Four numbers already checked again. (i) $|\mathcal{P}(A)| = 2^{n}$ for
$n = 0,1,2,3,4$ gives $1, 2, 4, 8, 16$.
(ii) $\pi(34,10) = 990 + 10 = 1000$, and $\pi$ on diagonals $s = 0,1,2,3$ gives
$0,1,2,3,4,5,6,7,8,9$. (iii) The number of positive rationals of height $k$ is the
totient value at $k$; the cumulative count up to $k = 9$ is $27$, so $7/3$ sits at
position 30. (iv) Relations on a set with $3$ members: $512$ in all, $5$ equivalence
relations.}

\section{Computational lab}
\tingkat{advanced}

The following script runs the four computations and adds one check that cannot be done
by hand: it compares the inverse of the Cantor pairing computed with machine floating
point against the one computed with integer square roots. Both should give the same
answer; at one point the answers differ.

\begin{lab}
\begin{verbatim}
from itertools import combinations
from math import gcd, isqrt, sqrt, floor

# (1) number of subsets for small n
for n in range(5):
    A = range(n)
    subs = [S for k in range(n + 1) for S in combinations(A, k)]
    print("n =", n, "|P(A)| =", len(subs), "2^n =", 2 ** n)

# (2) Cantor pairing pi(m, n) = (m+n)(m+n+1)/2 + n
def pi(m, n):
    s = m + n
    return s * (s + 1) // 2 + n

pasangan = [(0, 0), (1, 0), (0, 1), (2, 0), (1, 1),
            (0, 2), (3, 0), (2, 1), (1, 2), (0, 3)]
print("pi:", [pi(m, n) for m, n in pasangan])
keys = {pi(m, n) for m in range(61) for n in range(61)}
print("unique keys over 61x61 pairs:", len(keys) == 61 * 61)

# (3) inverse: float vs isqrt
s, n = 2 * 10**8, 2 * 10**8 - 1
q = s * (s + 1) // 2 + n
print("q =", q, "s correct =", s)
print("s from float:", int(floor((sqrt(8 * q + 1) - 1) / 2)))
print("s from isqrt:", (isqrt(8 * q + 1) - 1) // 2)

# (4) enumerate positive Q by height k = p + q
seq = []
for k in range(2, 10):
    row = [(p, k - p) for p in range(1, k) if gcd(p, k - p) == 1]
    print("k =", k, "count =", len(row), row)
    seq += row
print("terms 1..8:", seq[:8])
print("total up to k = 9:", len(seq))

# (5) relations on an n-element set: 2^(n*n) in all, Bell numbers equivalence
def ekuivalen(R, n):
    refl = all((a, a) in R for a in range(n))
    sim = all((b, a) in R for a, b in R)
    trans = all((a, c) in R for a, b in R for c in range(n)
                if (b, c) in R)
    return refl and sim and trans

for n in range(1, 5):
    semua = 2 ** (n * n)
    hitung = 0
    for bit in range(semua):
        R = {(a, b) for a in range(n) for b in range(n)
             if bit >> (a * n + b) & 1}
        if ekuivalen(R, n):
            hitung += 1
    print("n =", n, "relations =", semua, "equivalences =", hitung)
\end{verbatim}
\end{lab}

The output as printed, verbatim:

\begin{verbatim}
n = 0 |P(A)| = 1 2^n = 1
n = 1 |P(A)| = 2 2^n = 2
n = 2 |P(A)| = 4 2^n = 4
n = 3 |P(A)| = 8 2^n = 8
n = 4 |P(A)| = 16 2^n = 16
pi: [0, 1, 2, 3, 4, 5, 6, 7, 8, 9]
unique keys over 61x61 pairs: True
q = 20000000299999999 s correct = 200000000
s from float: 200000001
s from isqrt: 200000000
k = 2 count = 1 [(1, 1)]
k = 3 count = 2 [(1, 2), (2, 1)]
k = 4 count = 2 [(1, 3), (3, 1)]
k = 5 count = 4 [(1, 4), (2, 3), (3, 2), (4, 1)]
k = 6 count = 2 [(1, 5), (5, 1)]
k = 7 count = 6 [(1, 6), (2, 5), (3, 4), (4, 3), (5, 2), (6, 1)]
k = 8 count = 4 [(1, 7), (3, 5), (5, 3), (7, 1)]
k = 9 count = 6 [(1, 8), (2, 7), (4, 5), (5, 4), (7, 2), (8, 1)]
terms 1..8: [(1, 1), (1, 2), (2, 1), (1, 3), (3, 1), (1, 4), (2, 3), (3, 2)]
total up to k = 9: 27
n = 1 relations = 2 equivalences = 1
n = 2 relations = 16 equivalences = 2
n = 3 relations = 512 equivalences = 5
n = 4 relations = 65536 equivalences = 15
\end{verbatim}

What must be read from that output: the first five lines confirm $2^{n}$ for
$n = 0$ through $4$; the ten values of $\pi$ are consecutive without gaps, and
$61 \times 61 = 3721$ pairs give $3721$ unique keys, which is numerical evidence, not a
mathematical proof, because the mathematical proof is the $T_s$ block argument. The
last four lines sweep all $2^{n^{2}}$ relations for $n = 1$ through $4$ and count those
that are reflexive, symmetric, and transitive: the result is $1, 2, 5, 15$. There is no
line for $n = 5$ because sweeping $2^{25}$ relations takes too long; its value is $52$.

Third, the inverse with faulty floating point. At $q = 20{,}000{,}000{,}299{,}999{,}999$ the
formula $\lfloor (\sqrt{8q+1}-1)/2 \rfloor$ with 64-bit machine floating point gives
$s = 200{,}000{,}001$, whereas the true value is $200{,}000{,}000$. The error is one, and that
is enough to return a completely different pair: with $s = 200{,}000{,}001$ the value
$n = q - T_s$ becomes negative. The number $8q+1$ is of order $10^{17}$, whereas
64-bit floating point stores only about $16$ decimal digits exactly; \texttt{isqrt}
works on integers and gives $200{,}000{,}000$. The lesson is general: inverting an exact
numbering formula demands integer arithmetic, as discussed in Chapter 22.

\begin{untukdosen}
Four meetings. First: set operations and De Morgan, with each proof written as a chain
of equivalences, not as two separate inclusions. Second: Cartesian products,
relations, then the filtering into partial orders and equivalence relations. Third:
classes, partitions, and functions, with the Lemma on classes proved at the board.
Fourth: cardinality up to Cantor's Theorem, plus a sketch of Cantor--Bernstein if time
allows.

Two misconceptions appear most often. First, students think a symmetric and transitive
relation is automatically reflexive; the counterexample
$R = \{(1,1),(2,2),(1,2),(2,1)\}$ on $\{1,2,3\}$ is usually enough, and Exercise 2 asks
for a verdict on that claim. Second, students think $\mathbb{Q}$ is larger than
$\mathbb{N}$ because of its density; ask for the first ten terms of the height
enumeration at the board, then ask for the position of $7/3$.

An activity that works: divide the board into "natural numbers" and "positive
rationals", then have two volunteers pair members for two minutes. The naive pairing
will fail, and then together build the pairing by height. Students are ready to move on
when they can explain why that list misses no rational, and why a similar list fails
for $[0,1]$.
\end{untukdosen}

\section{Summary}
\tingkat{basic}

\begin{ringkasan}
\begin{itemize}
  \item Union, intersection, difference, and complement are defined through
        membership conditions; De Morgan's laws are laws of logic transferred to
        sets, and their version holds for any indexed family.
  \item $|\mathcal{P}(A)| = 2^{|A|}$ for finite $A$, via a bijection to the
        characteristic functions $A \to \{0,1\}$.
  \item A relation is a subset of $A \times B$; on a set with $n$ members there are
        $2^{n^{2}}$ relations, and four properties filter them into equivalence
        relations and partial orders.
  \item Equivalence relations and partitions are two faces of the same object; the
        classes are always nonempty, mutually disjoint, and cover the whole set.
  \item A bijective function is exactly one that has a unique two-sided inverse, and
        on finite sets of equal size the three properties lock together.
  \item $\mathbb{N} \times \mathbb{N}$ and $\mathbb{Q}$ are countable through explicit
        numberings, whereas $\mathcal{P}(\mathbb{N})$ and $\mathbb{R}$ are uncountable
        through Cantor's diagonal argument.
  \item The supremum is the least upper bound; its existence is guaranteed by the
        completeness axiom, while the epsilon characterization and its construction
        continue in Chapter 3.
\end{itemize}
\end{ringkasan}

\section{Exercises}
\tingkat{intermediate}

\begin{latihan}
\begin{enumerate}
  \item Prove $A \cap (B \cup C) = (A \cap B) \cup (A \cap C)$ and
        $A \setminus (B \setminus C) = (A \setminus B) \cup (A \cap C)$ with chains of
        membership equivalences.
  \item Assess this claim: true, false, or not enough evidence. "Every relation on a
        nonempty set that is symmetric and transitive is necessarily reflexive." If it
        is false, give the smallest possible counterexample; if there is not enough
        evidence, state what further evidence would be needed.
  \item Prove that if $f : A \to B$ is injective, then $S \mapsto f(S)$ is an
        injective function from $\mathcal{P}(A)$ to $\mathcal{P}(B)$. Explain where
        injectivity is used and why the conclusion fails without it.
  \item Prove Kuratowski's ordered pair lemma:
        $\{\{a\},\{a,b\}\} = \{\{c\},\{c,d\}\}$ if and only if $a = c$ and
        $b = d$. Handle the case $a = b$ separately.
  \item Complete the first part of the counting Lemma: for $f : A \to B$ with $A$
        finite, $|f(A)| \le |A|$ and equality holds only if $f$ is injective.
  \item Count the equivalence relations on a set with $5$ members that have exactly
        two classes, then check the answer with Python.
  \item Compute $\pi(7,5)$ with the Cantor pairing formula, then find $(m,n)$ with
        $\pi(m,n) = 1000$. Check both answers with Python.
  \item Determine the position of $7/3$ in the enumeration of positive rationals by
        height, then determine the number at positions 30 and 41. Check with Python.
  \item Prove that if $A_1, A_2, A_3, \dots$ are each countable and the numbering of
        every $A_n$ is given explicitly, then $\bigcup_{n \ge 1} A_n$ is countable.
  \item Show that $(0,1)$ is equinumerous with $\mathbb{R}$ through one explicit
        bijection, and explain why the function we choose must have its domain
        restricted.
\end{enumerate}
\end{latihan}

Hints for the three hardest items. Exercise 5: in the induction step write
$A = A' \cup \{a\}$ with $a \notin A'$, then compare $f(A)$ with
$f(A') \cup \{f(a)\}$; equality of sizes demands $f(a) \notin f(A')$ and
$|f(A')| = |A'|$ at the same time. Exercise 6: the number of ways to split $n$ members
into two nonempty parts is $(2^{n}-2)/2$, since choosing the first part or the second
gives the same split. Exercise 9: arrange the numbering of all the $A_n$ as a two-column
table, read it diagonal by diagonal as in the proof that $\mathbb{N} \times \mathbb{N}$
is countable, then discard repeated terms.

\begin{untukpeneliti}
Four things are genuinely proved here. Cantor's Theorem uses no axiom of choice and is
constructive: from any $f$ we build $D$ that lies outside its range. The
uncountability of $\mathbb{R}$ depends on the binary expansion of every point of
$[0,1]$; that fact requires completeness and here is only referred to Chapter 3 and
Chapter 5.

The Cantor--Bernstein--Schroeder theorem in this chapter is only a sketch. A full
proof requires careful bookkeeping of chain depth; the complete formal version is in
Avigad--Massot, \emph{Mathematics in Lean}, section 4.3, while the classical treatment
of sets, functions, and equivalence relations is in Hoffman--Kunze,
\emph{Linear Algebra}, Appendix A.1--A.3. The primary references for the diagonal and
for uncountability are Cantor, "Ueber eine Eigenschaft des Inbegriffes aller
reellen algebraischen Zahlen" (1874) and "Ueber eine elementare Frage der
Mannigfaltigkeitslehre" (1891).

The most famous open question: is there a cardinal between $|\mathbb{N}|$ and
$|\mathbb{R}|$? The continuum hypothesis states that there is none. Godel (1940) showed
that the hypothesis is consistent with ZFC, Cohen (1963) showed that its negation is
consistent too. So the question cannot be answered inside ZFC.

Two technical limits. First, $\sim$ may not be called an equivalence relation on "the
set of all sets", because that collection is not a set; the correct treatment uses
classes as metatheoretic objects. Second, the corollary "a countable union of
countable sets is countable" requires countable choice when the numberings are not
given, and Exercise 9 deliberately requires an explicit numbering.
\end{untukpeneliti}

\peta{Chapter 2 maps to Abbott, \emph{Understanding Analysis}, sections "Some
Preliminaries" (sets, functions, cardinality) and "Cantor's Theorem"; Hoffman--Kunze,
\emph{Linear Algebra}, Appendix A.1 (sets), A.2 (functions and inverses), A.3
(equivalence relations, classes, quotients); Avigad--Massot, \emph{Mathematics in
Lean}, Chapter 4 ("Sets and Functions") for extensionality, preimages, and
Cantor--Bernstein--Schroeder, and Chapter 6.1 for finite sets.}

\chapter{Axioms of the Real Numbers and Supremum}

A program searches for the length of the diagonal of a square with side one. It keeps
two numbers, $a=1$ and $b=2$, with $a^2-2<0<b^2-2$. At every step it computes the
midpoint $m=(a+b)/2$, evaluates $m^2-2$, and then replaces $a$ or $b$ by $m$ so that
the signs at the two ends stay different. After 40 steps the width of the interval is
below $10^{-12}$, after 100 steps below $10^{-30}$. But the program never produces an
answer: no step yields $x^2-2=0$ exactly, and the two endpoints keep squeezing
something that is not a rational number.

That failure is not a mistake in the algorithm. It shows that the rational numbers
have holes, and that the set $\{q\in\mathbb{Q}: q^2<2\}$ has no least upper bound
inside $\mathbb{Q}$.

Chapter 1 taught us to write proofs. Chapter 2 gave us the language of sets and
functions, and used $|x|$ and $\sup$ as computing tools without asking where they come
from. What is still missing is the object that makes those tools legitimate. This
chapter builds it, not out of decimals, but out of three groups of statements accepted
without proof: the field laws, the order laws, and completeness. Everything else,
including the existence of $\sqrt{2}$, is derived from there.

\begin{untukanak}
Imagine a ruler that has only fraction marks: one half, one quarter, one eighth, and
ever closer together. There are no other numbers on that ruler. Now measure the
diagonal of a square whose side is one hand-span. The end of the diagonal falls
between two marks, and it always does, no matter how close the marks are. We can get
as near as we like, but never exactly.

The question is: does that length really exist? The ruler clearly has no answer. Yet
the diagonal, if you carry it around or lay it on the floor, has a definite length.
That length does not disappear merely because our measuring tool does not contain it.

The same story happens with numbers. If we have only fractional numbers, there is a
point on the line that has no name. Mathematicians close that hole with a single
promise: every collection of numbers that is bounded above has a lowest roof. In this
chapter we will learn why that promise is needed, what its consequences are, and how
it gives a name to the diagonal length above, namely root two.
\end{untukanak}

\section{Field axioms}
\tingkat{basic}

We do not define the real numbers as endless decimals: that approach makes addition
complicated, and makes the definition of a least upper bound circular, since
``the supremum of all decimal cuts'' already uses the numbers being built. It is
cleaner to state the properties first and then derive the rest.

The first layer is algebraic. We assume two operations that obey the rules we have
used since school; the following definition only gives those rules a name.

\begin{definisi}[Field]\index{field}
A set $F$ is called a \emph{field} if addition and multiplication are defined on its
members so that: the result of both operations always lies in $F$; both operations are
commutative and associative; there is an element $0$ with $0+x=x$ and an element
$1\neq 0$ with $1\cdot x=x$; every $x$ has a negative $-x$ with $x+(-x)=0$; every
$x\neq 0$ has a reciprocal $1/x$ with $x\cdot(1/x)=1$; and $x(y+z)=xy+xz$ holds.
\end{definisi}

Those axioms mention neither order, nor distance, nor completeness. Other sets
therefore satisfy them too. The rational numbers with ordinary addition and
multiplication form a field. A reasonable question: what is already contained inside
those rules?

Three algebraic statements we usually memorize turn out to be consequences, not
axioms; the one most often used is multiplication by zero, which destroys everything.

\begin{teorema}
In an arbitrary field, for all $x,y$: (a) $0\cdot x=0$;
(b) $(-x)y=-(xy)$; (c) $(-x)(-y)=xy$; (d) if $xy=0$ then $x=0$ or $y=0$.
\end{teorema}

\begin{proof}
(a) From $0+0=0$ and the distributive law, $0\cdot x=(0+0)x=0\cdot x+0\cdot x$.
Add $-(0\cdot x)$ to both sides. The left side becomes $0$, the right side becomes
$0\cdot x$. Hence $0=0\cdot x$.

(b) Compute $(-x)y+xy=(-x+x)y=0\cdot y=0$, using (a). Hence $(-x)y$ is the negative
of $xy$, and that negative is unique.

(c) Apply (b) twice: $(-x)(-y)=-\bigl(x(-y)\bigr)=-\bigl(-(xy)\bigr)=xy$.
The last step uses $--a=a$, which is also a consequence of the uniqueness of the
negative.

(d) Suppose $x\neq 0$ and $xy=0$. Multiply by $1/x$: $(1/x)(xy)=0$, so that
$y=0$. Hence one of the factors is zero.
\end{proof}

Not one of those four statements is assumed; all of them descend from the distributive
law and the existence of negatives. This pattern is used throughout the book: state as
few rules as possible, then harvest the consequences.

\kamus{field}{lapangan}{A set with two arithmetic operations whose rules stay fixed, like
a toolbox whose lid always shuts tight: the result of a computation never leaves the
box.}

Chapter 2 used $\mathbb{Q}$ as a countable set. Here it appears in a different role:
as an ordered field. Those two properties are not yet enough, and that is the content
of the next section.

\section{Order axioms}
\tingkat{basic}

A field alone cannot compare two numbers. We need a relation ``less than'' that
respects addition and multiplication, because without it inequalities cannot be
manipulated. The order axioms say three things: any two numbers can be compared, the
comparison is transitive, and adding or multiplying by a positive number does not
break it.

\begin{definisi}[Ordered field]\index{ordered field}
A field $F$ is called \emph{ordered} if there is a relation $<$ on $F$ such that for
all $a,b,c$: exactly one of $a<b$, $a=b$, $b<a$ holds; if $a<b$ and $b<c$ then
$a<c$; if $a<b$ then $a+c<b+c$; and if $a<b$ and $c>0$ then $ac<bc$.
\end{definisi}

Almost every sign rule we memorize is a consequence of that definition, not an
addition to it. The following proof derives five of them at once; the pattern is the
same: add something to both sides, or multiply by a number whose sign is known.

\begin{proposisi}
In every ordered field: (a) if $x>0$ then $-x<0$, and conversely;
(b) if $x>0$ and $y<z$ then $xy<xz$; (c) if $x<0$ and $y<z$ then $xy>xz$;
(d) if $x\neq 0$ then $x^2>0$, in particular $1>0$; (e) if $0<x<y$ then
$0<1/y<1/x$.
\end{proposisi}

\begin{proof}
(a) If $x>0$, add $-x$ to both sides: $0>-x$. Conversely, if $x<0$ then
$0-x>0$, that is, $-x>0$.

(b) Since $z>y$, we have $z-y>0$. Multiply by $x>0$: $x(z-y)>0$. Hence
$xz=x(z-y)+xy>xy$.

(c) From (a), $-x>0$. Apply (b) to $-x$ and $y<z$: $(-x)y<(-x)z$. By the previous
Theorem, the left side is $-(xy)$ and the right side is $-(xz)$. Hence
$-(xy)<-(xz)$, and by (a) again $xy>xz$.

(d) If $x>0$, the order axiom applied to $y=x>0$ gives $x^2>0$. If $x<0$ then
$-x>0$, so that $(-x)^2>0$; yet $(-x)^2=x^2$. For $x=0$ clearly $x^2=0$. Since
$1=1^2$, we obtain $1>0$.

(e) Since $y>0$ and $1/y\neq 0$, from (d) $1/y>0$; the same holds for $1/x$.
Multiply $x<y$ by the positive number $(1/x)(1/y)$: the result is $1/y<1/x$.
\end{proof}

Two consequences that are used often: $a<b$ is equivalent to $b-a>0$, and $a<b$ is
equivalent to $-b<-a$; the second explains why multiplying by a negative number
reverses the direction, because multiplying by $-1$ reflects the number line about
zero. Another consequence: if $0<a<b$ then $a^2<b^2$, simply by multiplying $a<b$ by
$a>0$ and then by $b>0$ and adding the two. Hence comparing squares is equivalent to
comparing the original positive numbers; every time we write $q^2<2$ for a rational
$q$, we are comparing $q$ with something that need not be rational.

\kamus{ordered field}{lapangan terurut}{A field whose members can be arranged like the
steps of a staircase, and whose arithmetic rules do not disturb that arrangement.}

Note what we do not have yet. We do not know whether there is a positive number smaller
than every $1/n$, and we do not know whether $\{q\in\mathbb{Q}: q^2<2\}$ has a roof.
That is not a question of order but a question of completeness.

\section{Absolute value and the triangle inequality}
\tingkat{intermediate}

In Chapter 2 we wrote $|x-y|$ as a distance, and used it to compare errors, without
proving that the symbol behaves like a distance. Now its properties are derived from
the order. Absolute value discards the sign and leaves the distance to the point zero;
the following definition writes two branches, and it is the second branch that is
usually used.

\begin{definisi}[Absolute value]\index{absolute value}
For $x\in\mathbb{R}$, $|x|=x$ if $x\ge 0$ and $|x|=-x$ if $x<0$. Equivalently,
$|x|=\max\{x,-x\}$, and for real numbers $|x|=\sqrt{x^2}$ with the root discussed in
the next section.
\end{definisi}

From that definition it follows directly: $|x|\ge 0$, $|{-x}|=|x|$, $|x|^2=x^2$, and
$|xy|=|x|\,|y|$. The last is proved by checking the four sign combinations, or by
noting that both sides are nonnegative and their squares agree.

Inequalities containing an absolute value are easiest to handle after being converted
into two ordinary inequalities. The following lemma is the reason that conversion is
legitimate.

\begin{lemma}
For $c\ge 0$, $|x|\le c$ if and only if $-c\le x\le c$.
\end{lemma}

\begin{proof}
If $|x|\le c$: for $x\ge 0$ we have $x=|x|\le c$ and $-c\le 0\le x$; for
$x<0$ we have $-x=|x|\le c$, so that $x\ge -c$, and $x<0\le c$. Hence
$-c\le x\le c$.

Conversely, if $-c\le x\le c$: for $x\ge 0$, $|x|=x\le c$; for $x<0$,
$|x|=-x\le c$ because $x\ge -c$. In both cases $|x|\le c$.
\end{proof}

The triangle inequality is the only tool we need to control the difference of two
numbers. Its name comes from a picture: the length of one side of a triangle does not
exceed the sum of the lengths of the other two sides.

\begin{teorema}[Triangle inequality]\index{triangle inequality}
For all $x,y\in\mathbb{R}$, $|x+y|\le |x|+|y|$. Moreover
$\bigl||x|-|y|\bigr|\le |x-y|$.
\end{teorema}

\begin{proof}
From the definition, $-|x|\le x\le |x|$ and $-|y|\le y\le |y|$. Add the two pairs:
$-(|x|+|y|)\le x+y\le |x|+|y|$. By the previous lemma with $c=|x|+|y|$, we obtain
$|x+y|\le |x|+|y|$.

For the second part, write $x=(x-y)+y$ and apply the first part:
$|x|\le |x-y|+|y|$, so that $|x|-|y|\le |x-y|$. Swapping the roles of $x$ and $y$
gives $|y|-|x|\le |x-y|$. Since $c\le d$ and $-c\le d$ is equivalent to $|c|\le d$,
we have $\bigl||x|-|y|\bigr|\le |x-y|$.
\end{proof}

As a check, take $x=3$ and $y=-5$: $|x+y|=2$ while $|x|+|y|=8$, so the inequality is
slack. For $x=3$ and $y=5$ both sides equal $8$ and the inequality becomes an equation;
that happens exactly when the signs agree or one of them is zero.

Its practical use: bounds on distance become bounds on functions. If
$|x-\sqrt2|<\delta$, then
\[
|x^2-2|=|x-\sqrt2|\,|x+\sqrt2|<\delta\,(2\sqrt2+\delta),
\]
because $|x+\sqrt2|\le|x-\sqrt2|+2\sqrt2$. For $\delta=10^{-3}$ that bound equals
$0.002829427$, and the true value at $x=\sqrt2+10^{-3}$ is also $0.002829427$ since
$x^2-2=2\sqrt2\,\delta+\delta^2$ exactly: the bound is tight. For $\delta=10^{-6}$ the
same bound shrinks to $2.828428\times10^{-6}$.

\kamus{absolute value}{nilai mutlak}{The distance of a point from zero on the number line,
regardless of which side of zero the point lies on.}

The triangle inequality uses only the order axioms, not completeness, so it also holds
in $\mathbb{Q}$ and cannot distinguish the two. The distinguishing tool appears in the
next section.

\section{Completeness: the least upper bound}
\tingkat{intermediate}

Order gives comparison, but it does not give a roof. The set
$A=\{1-1/n: n\in\mathbb{N}\}$ contains numbers that come ever closer to one without
ever equaling one: one is not a member, yet it is the tightest bound. The set
$B=\{q\in\mathbb{Q}: q^2<2\}$ is bounded, yet it has no tight bound in
$\mathbb{Q}$. Both show that a slack bound and a tightest bound need different names.

\begin{untukanak}
Imagine we must put up a tent roof for a group of children. A roof that is too low:
some head sticks out. A roof that is too high: the tent wastes material and wobbles.
What we are looking for is the lowest roof that still lets all the children fit. Call
it the right roof.

Now something strange happens. In that group there is no tallest child; there is always
someone a little taller. Yet the right roof still exists, because their heights are
bounded: nobody is as tall as a tree. The right roof is determined by the whole group,
not by a single child.

In this chapter the right roof has a name: the least upper bound. The property that
guarantees a right roof always exists is called completeness. The fractional numbers
are incomplete: there are groups of bounded height whose right roof is not among the
fractional numbers. A question you can take home: if there is no tallest child, how do
you prove that the right roof really is the lowest one?
\end{untukanak}

An upper bound is a ceiling that contains every member; the least upper bound is the
lowest ceiling among them. The following definition names both, together with their
mirror image below.

\begin{definisi}[Upper bound and supremum]\index{upper bound}\index{supremum}
Let $A\subseteq\mathbb{R}$ and $A\neq\varnothing$. A number $u$ is called an
\emph{upper bound} of $A$ if $a\le u$ for every $a\in A$. If $A$ has at least one upper
bound, $A$ is called bounded above. A number $s$ is called the \emph{supremum} of $A$,
written $s=\sup A$, if $s$ is an upper bound of $A$ and $s\le u$ for every upper bound
$u$. Dually, $t$ is called the \emph{infimum} of $A$, written $t=\inf A$, if $t$ is a
lower bound and $t\ge v$ for every lower bound $v$.
\end{definisi}

The two conditions in the definition of supremum are often confused: the first says
that $s$ is still an upper bound, the second that no upper bound is smaller. For
$A=\{1-1/n\}$, the number $1$ satisfies both while $2$ satisfies only the first; that
set also has no largest member, so ``maximum'' and ``supremum'' are not synonyms.

\begin{bacakeras}
Two voices. The first voice asks, the second answers. Read them aloud alternately with
a friend.

``What is that fence?'' ``The fence is an upper bound: every member of the collection
lies below it.'' ``How many fences are there?'' ``Many. Every higher fence is also a
fence.'' ``Then which one do we use?'' ``The lowest one. Its name is the least upper
bound.'' ``What if that lowest fence happens to be occupied by a member of the
collection?'' ``Then it is also a maximum. If it is not occupied, it is still the least
upper bound. Two names, two jobs.''

Repeat three times, faster each time: a tall fence is not the answer, the lowest fence
is the answer.

The line the reader must repeat, once, in a clear voice: an upper bound is a fence; the
least upper bound is the lowest fence that still holds all the members in.
\end{bacakeras}

Figure~\ref{fig:sup-garis} places both ideas on a single line. The members of
$A=\{1-1/n\}$ crowd toward one direction without ever reaching the end; every number to
the right of one is an upper bound, and the leftmost among them is the supremum.

\begin{figure}[htbp]
\centering
\begin{tikzpicture}[x=1.5cm,y=1cm,font=\small]
  \draw[-{Latex[length=2mm]},ink,thick] (-0.35,0) -- (5.5,0);
  \foreach \t in {0,1,2,3,4,5}{
    \draw[ink] (\t,0.07) -- (\t,-0.07);
    \node[below=2pt,ink] at (\t,0) {$\t$};
  }
  \foreach \n in {1,...,6}{
    \fill[accent] ({1-1/\n},1.0) circle (1.7pt);
  }
  \draw[accent,thick,fill=white] (1,1.0) circle (1.7pt);
  \node[anchor=west,accent,font=\footnotesize] at (1.18,1.0) {members of $A=\{1-1/n\}$};
  \draw[layergrey,dashed] (1,0.95) -- (1,0);
  \node[accent,anchor=west,font=\footnotesize] at (1.12,0.45)
    {$\sup A=1$, and $1$ is not a member of $A$};
  \fill[layerteal!14] (1,-0.95) rectangle (5.35,-0.6);
  \draw[layerteal,thick] (1,-0.95) rectangle (5.35,-0.6);
  \node[layerteal!80!black,font=\footnotesize] at (3.2,-0.775) {all upper bounds of $A$};
  \draw[accent,line width=1.4pt] (1,-1.02) -- (1,-0.52);
\end{tikzpicture}
\caption{The set $A=\{1-1/n:n\in\mathbb{N}\}$ on the number line. Every $u\ge 1$
is an upper bound; the smallest, $u=1$, is $\sup A$. The number $1$ is not a member of
$A$, so $A$ has no maximum.}
\label{fig:sup-garis}
\end{figure}

\kamus{supremum}{supremum (batas atas terkecil)}{The lowest roof that still covers every
member of the set; that roof need not belong to the set.}

A supremum, if it exists, is unique; the second condition can be written with epsilon,
and that form is the one most often used in proofs.

\begin{lemma}
(i) If $A$ has a supremum, then its supremum is unique. (ii) A number $s$ is the
supremum of $A$ if and only if $s$ is an upper bound of $A$ and for every
$\varepsilon>0$ there exists $a\in A$ with $a>s-\varepsilon$.
\end{lemma}

\begin{proof}
(i) Suppose $s$ and $s'$ are both suprema. Since $s$ is an upper bound and $s'$ is a
supremum, $s'\le s$. Since $s'$ is an upper bound and $s$ is a supremum, $s\le s'$.
Hence $s=s'$.

(ii) Suppose $s$ is an upper bound and the epsilon condition holds, but there is an
upper bound $u<s$. Take $\varepsilon=s-u>0$. By the epsilon condition there is
$a\in A$ with $a>s-\varepsilon=u$, contradicting that $u$ is an upper bound.
Conversely, if $s$ is a supremum and there is $\varepsilon>0$ such that
$a\le s-\varepsilon$ for all $a\in A$, then $s-\varepsilon$ is an upper bound smaller
than $s$, contradicting that $s$ is the least upper bound.
\end{proof}

\begin{resep}
To prove that a number $s$ is the supremum of a set $A$:
\begin{enumerate}
\item Show $A\neq\varnothing$ by naming one of its members.
\item Prove that $s$ is an upper bound: take arbitrary $a\in A$, show $a\le s$.
\item Prove that no upper bound is smaller: take $\varepsilon>0$, produce
      $a\in A$ with $a>s-\varepsilon$, for instance $a=1-1/n$ with $n>1/\varepsilon$.
\item Conclude: $s$ is an upper bound and every number below $s$ is not an upper bound,
      so $s=\sup A$.
\end{enumerate}
Order warning: step 2 must be finished before step 3. Choosing $a>s-\varepsilon$ first
proves nothing, since $s$ is not yet known to be an upper bound. For
$A=\{1-1/n\}$, the number $1/2$ passes step 3 but fails step 2.
\end{resep}

The following axiom distinguishes $\mathbb{R}$ from $\mathbb{Q}$. It cannot be proved
from the field and order laws, because it is a statement about the completeness of the
line.

\begin{teorema}[Least-upper-bound property]\index{completeness axiom}
Every nonempty set $A\subseteq\mathbb{R}$ that is bounded above has a supremum in
$\mathbb{R}$.
\end{teorema}

\begin{proof}[Proof]
This statement is accepted as the completeness axiom. Other equivalent formulations:
nested interval intersections, and convergence of Cauchy sequences. We choose the
supremum because it is the most directly usable; the construction of $\mathbb{R}$ from
$\mathbb{Q}$ that proves it exists is in the Appendix to Chapter 1 of Rudin.
\end{proof}

The infimum need not be axiomatized: it is the supremum seen in a mirror.

\begin{proposisi}
If $A$ is nonempty and bounded below, then $\inf A=-\sup(-A)$, with
$-A=\{-a: a\in A\}$.
\end{proposisi}

\begin{proof}
The set $-A$ is nonempty and bounded above, so $s=\sup(-A)$ exists. For every
$a\in A$ we have $-a\le s$, hence $a\ge -s$: so $-s$ is a lower bound of $A$. Now
take $t>-s$. Then $-t<s$, so $-t$ is not an upper bound of $-A$; there is $a\in A$
with $-a>-t$, that is, $a<t$. Hence $t$ is not a lower bound of $A$. Therefore $-s$
is the greatest lower bound.
\end{proof}

First consequence of completeness: the natural numbers have no ceiling. It sounds
trivial, but without completeness it cannot be proved.

\begin{teorema}[Archimedean property]\index{archimedean property}
For every $x>0$ and every $y\in\mathbb{R}$ there exists $n\in\mathbb{N}$ with
$nx>y$.
\end{teorema}

\begin{proof}
Suppose the contrary: there are $x>0$ and $y$ such that $nx\le y$ for all $n$. Then the
set $S=\{nx: n\in\mathbb{N}\}$ is nonempty and bounded above by $y$, so
$\alpha=\sup S$ exists. Since $x>0$, the number $\alpha-x$ is smaller than $\alpha$ and
therefore not an upper bound of $S$. Hence there is $m$ with $mx>\alpha-x$, which gives
$(m+1)x>\alpha$. Yet $(m+1)x\in S$, and $\alpha$ is an upper bound of $S$. Contradiction.
\end{proof}

The reciprocal version is the one most often used: for every $\varepsilon>0$ there is
$n$ with $1/n<\varepsilon$, simply by applying the theorem to $x=\varepsilon$ and
$y=1$. The phrase ``for $n$ sufficiently large'' is shorthand for that consequence, and
now we know it is legitimate.

That property also destroys very small positive numbers: there is no $\delta>0$ that is
smaller than $1/n$ for all $n$, because $1/\delta$ would become an upper bound of
$\mathbb{N}$. Real analysis therefore needs no ``infinitesimals''.

Second consequence: every open interval, however small, contains a rational number.

\begin{teorema}[Density of $\mathbb{Q}$]\index{density}
If $x<y$ are real numbers, then there exists $p\in\mathbb{Q}$ with $x<p<y$.
\end{teorema}

\begin{proof}
Since $y-x>0$, the Archimedean property gives $n\in\mathbb{N}$ with $n(y-x)>1$. Apply
Archimedes again to obtain integers $m_1>nx$ and $m_2>-nx$; so
$-m_2<nx<m_1$. Since $nx$ is trapped between two integers, there is an integer
$m$ with $m-1\le nx<m$. Then $nx<m$, and from $m-1\le nx$ we get
$m\le nx+1<ny$. Since $n>0$, division gives $x<m/n<y$.
\end{proof}

That proof is constructive: it gives a recipe for finding $p$, namely choosing $n$
large enough and then rounding $nx$ upward. It does not give a small denominator, so
efficient approximation still demands other ideas.

\section{Why the rational numbers are not enough}
\tingkat{advanced}

The program in the opening fails because the set it squeezes has no roof in
$\mathbb{Q}$. Let us prove that, not merely display the numbers. Take
$A=\{q\in\mathbb{Q}: q^2<2\}$. This set is nonempty, since $1\in A$, and it is bounded
above inside $\mathbb{Q}$: if $q\ge 3/2$ then $q^2\ge 9/4>2$, so $3/2$ is an upper
bound.

The intervals that program keeps are drawn as a stack of ever shorter bars
(Figure~\ref{fig:selang-bersarang}). Each bar still contains root two and its length
halves, while all its endpoints stay rational; the bars narrow onto a single point
without ever touching it.

\begin{figure}[htbp]
\centering
\begin{tikzpicture}[x=9cm,y=1cm,font=\small]
  \foreach \y/\a/\b/\lab/\pan in {%
    0/1/2/{$[1,\; 2]$}/{$1$},
    0.55/1/1.5/{$[1,\; \tfrac32]$}/{$1/2$},
    1.1/1.25/1.5/{$[\tfrac54,\; \tfrac32]$}/{$1/4$},
    1.65/1.375/1.5/{$[\tfrac{11}{8},\; \tfrac32]$}/{$1/8$},
    2.2/1.375/1.4375/{$[\tfrac{11}{8},\; \tfrac{23}{16}]$}/{$1/16$},
    2.75/1.40625/1.4375/{$[\tfrac{45}{32},\; \tfrac{23}{16}]$}/{$1/32$}}{
    \fill[layerindigo!12] (\a,\y) rectangle (\b,\y+0.3);
    \draw[layerindigo,thick] (\a,\y) rectangle (\b,\y+0.3);
    \node[layerindigo,anchor=east,font=\footnotesize] at (\a-0.018,\y+0.15) {\lab};
    \node[layergrey,anchor=west,font=\footnotesize] at (\b+0.015,\y+0.15) {\pan};
  }
  \draw[accent,dashed,thick] (1.41421,-0.25) -- (1.41421,3.15);
  \node[accent,font=\footnotesize,above] at (1.41421,3.15) {$\sqrt2$};
  \draw[-{Latex[length=1.6mm]},layergrey] (1.0,-0.12) -- (2.0,-0.12);
  \node[layergrey,font=\footnotesize,below] at (1.5,-0.14) {length $=1$};
\end{tikzpicture}
\caption{The intervals kept by the opening program: $[1,2]$, $[1,\tfrac32]$,
$[\tfrac54,\tfrac32]$, $[\tfrac{11}{8},\tfrac32]$, $[\tfrac{11}{8},\tfrac{23}{16}]$,
$[\tfrac{45}{32},\tfrac{23}{16}]$. Their lengths fall from $1$ to $1/32$; continued
without end, the intersection shrinks to a single point, namely $\sqrt2$, which is not
a rational number.}
\label{fig:selang-bersarang}
\end{figure}

First step: there is no largest member. Every rational with $q^2<2$ can still be nudged
a little upward without passing two, and the size of the nudge has a formula.

\begin{lemma}
If $q\in\mathbb{Q}$, $q>0$, and $q^2<2$, then there exists $h\in\mathbb{Q}$, $h>0$,
with $(q+h)^2<2$. In particular $q+h\in A$ and $q+h>q$.
\end{lemma}

\begin{proof}
Take $h=\dfrac{2-q^2}{2q+2}$. Since $q^2<2$ and $2q+2>0$, we have $h>0$ and $h$
rational. Also $h<1$, because $2-q^2<2q+2$ is equivalent to $q^2+2q>0$, which is true
for $q>0$. Then $h^2\le h$, so
\[
(q+h)^2=q^2+2qh+h^2\le q^2+h(2q+1)
=q^2+(2-q^2)\frac{2q+1}{2q+2}<q^2+(2-q^2)=2,
\]
because $\frac{2q+1}{2q+2}<1$.
\end{proof}

That formula can be checked numerically. For $q=1393/985$ we have $2-q^2=1/970225$, so
$h=1/4684660$ and
\[
q+h=\frac{6625109}{4684660}=1.414213411432206\ldots,
\]
whose square is $1.999999573074719\ldots<2$. So $1393/985$ is not the largest member,
even though its squared error is already about one part in a million. This procedure
can be repeated without end: there is no last step.

\begin{tangan}
Materials: 1.5 m of rope, six pieces of tape, a tape measure.

Stick the tape on the floor perpendicular to the wall, at distances $60/n$ cm for
$n=1,\ldots,6$: 60, 30, 20, 15, 12, and 10 cm. You have only stuck six, but the rule
does not stop. Stretch the rope across as a fence, then slide it toward the wall. Each
time you stop, check one thing: are all the tape marks on the same side of the rope as
the wall? As long as they are, the rope is still a fence.

The rope can be brought arbitrarily close to the wall and still be a fence: 5 cm, 1 cm,
1 mm. What stops it is not the tape but the wall. The collection of tape marks has no
smallest member, and not one of them touches the wall. The lowest fence still exists
--- the wall --- and it is not a member of the collection. That is the supremum.
\end{tangan}

It is not only the maximum that is missing. The least upper bound is missing from the
rational numbers as well, and the proof requires one fact about root two.

\begin{proposisi}
The set $A=\{q\in\mathbb{Q}: q^2<2\}$ has no supremum in $\mathbb{Q}$.
\end{proposisi}

\begin{proof}
Suppose $p\in\mathbb{Q}$ is the supremum of $A$. Since $1\in A$, we have $p\ge 1>0$. We
show that $p^2=2$, which is impossible because no rational number has square two.

If $p^2<2$, the previous lemma gives $q=p+h\in A$ with $q>p$, contradicting that $p$
is an upper bound. If $p^2>2$, take $k=(p^2-2)/(2p)>0$, which is rational and smaller
than $p$ since $p^2-2<2p^2$. For every $q\in A$,
\[
(p-k)^2=p^2-2pk+k^2>p^2-2pk=2>q^2,
\]
so $q<p-k$ for $q>0$, and for $q\le 0$ this is clear because $p-k>0$: the number
$p-k$ is still an upper bound of $A$, but smaller than $p$.
That contradicts $p$ being the least upper bound. Hence $p^2=2$.

Now $p=m/n$ with $m,n$ integers having no common factor. From $m^2=2n^2$, $m^2$ is
even, so $m$ is even; write $m=2r$. Then $4r^2=2n^2$, that is, $n^2=2r^2$, so $n$
is even as well. Both are even, contradicting the choice of $m,n$.
\end{proof}

Inside $\mathbb{R}$ that does not happen. The same set, viewed this time as a set of
real numbers, is bounded above by $2$, so by the completeness axiom it has a supremum.
That supremum is the number we name $\sqrt{2}$. The following proof closes the two
remaining possibilities, and its shape imitates the earlier lemma with $q$ replaced by
a real number.

\begin{teorema}[Existence of root two]\index{square root of two}
There is exactly one $s>0$ with $s^2=2$.
\end{teorema}

\begin{proof}
Let $E=\{t\in\mathbb{R}: t>0,\ t^2<2\}$. This set is nonempty because $1\in E$, and it
is bounded above by $2$: if $t>2$ then $t^2>4>2$, so $t\notin E$. By the
least-upper-bound property, $s=\sup E$ exists, and $s\ge 1>0$ because $1\in E$.

Suppose $s^2<2$. Take $h=\frac12\min\{1,(2-s^2)/(2s+1)\}$. Then $0<h<1$ and
$h(2s+1)<2-s^2$: if the first value is the minimum, $h=\frac12$ and
$\frac12(2s+1)\le\frac12(2-s^2)<2-s^2$; if the second is, $h(2s+1)=\frac12(2-s^2)$.
Since $h<1$, we have $h^2<h$, so
\[
(s+h)^2=s^2+2sh+h^2<s^2+h(2s+1)<s^2+(2-s^2)=2 .
\]
Hence $s+h\in E$ although $s+h>s$, contradicting that $s$ is an upper bound of $E$.
Therefore $s^2\ge 2$.

Suppose $s^2>2$. Take $k=(s^2-2)/(2s)>0$; then $k<s$ because $s^2-2<2s^2$. For every
$t\in E$ with $t\ge s-k$ we have
\[
t^2\ge(s-k)^2=s^2-2sk+k^2>s^2-2sk=2,
\]
so $t\notin E$. Hence $s-k$ is an upper bound of $E$ smaller than $s$, contradicting
that $s$ is the least upper bound. Since $s^2<2$ and $s^2>2$ are both impossible,
$s^2=2$.

Uniqueness: if $s,t>0$ and $s^2=t^2=2$, then $(s-t)(s+t)=0$. Since $s+t>0$, $s=t$.
\end{proof}

The number $s$ is irrational, by the parity argument above: the hole in $\mathbb{Q}$ is
real, and $\mathbb{R}$ closes it with exactly one number. What closes it is not order
and not the field laws, but completeness. The opening program fails not from lack of
precision, but because its number system does not contain the object it is looking for.

One thing must be stated honestly: we have not built $\mathbb{R}$ from $\mathbb{Q}$, we
have only assumed that $\mathbb{R}$ satisfies three groups of axioms. The construction
is known. Dedekind cuts $\mathbb{Q}$ in two: a \emph{cut} is a subset
$\alpha\subset\mathbb{Q}$ that is nonempty, not all of $\mathbb{Q}$, closed downward,
and without a largest member. That last property is exactly our lemma above. Cantor
uses equivalence classes of rational Cauchy sequences; both appeared in 1872. The
details are in the Appendix to Chapter 1 of Rudin.

That construction explains the standing of decimals. For $x>0$, take $n_0$ the largest
integer with $n_0\le x$, then $n_1$ the largest digit with $n_0+n_1/10\le x$, then
$n_2$ with $n_0+n_1/10+n_2/100\le x$, and so on. That cut is bounded above by $x$, so
its supremum exists and equals $x$. A decimal is not the raw material of the real
numbers but a way of naming a supremum; completeness guarantees that the name exists,
not that the name can be written with fractions.

\kamus{completeness axiom}{aksioma kelengkapan}{The promise that the number line has no
holes: every collection that is bounded above always has a lowest roof.}

The existence of $\sqrt2$ makes the computational problem sensible: $s$ can be
computed without knowing its digits in advance. Newton's iteration,
$x_{k+1}=\frac12(x_k+2/x_k)$, stays rational when $x_k$ is rational, so the whole
process runs inside $\mathbb{Q}$. From $s^2=2$ we get
\[
x_{k+1}-s=\frac{(x_k-s)^2}{2x_k},\qquad e_{k+1}=e_k^2\cdot\frac{s}{2x_k},
\]
with $e_k=|x_k-s|/s$. The right side of the first form is nonnegative, so every step
stays above $s$; there the factor $\frac{s}{2x_k}$ does not exceed $1/2$, so
$e_{k+1}\le e_k^2/2$ and the number of correct digits roughly doubles at each step.

One warning: we have not proved that $x_k\to s$. That needs the tools of Chapter 4.
What can be said now: since $e_1<0.061$, from $e_{k+1}\le e_k^2/2$ one can derive
$e_k\le 2^{-(2^k-1)}$ for $k\ge 1$, and that is enough to design an honest numerical
test.

\section{Intervals, accumulation points, and open sets}
\tingkat{advanced}

With the supremum in hand, we can name the pieces of the number line, like the finish
rules of a race: the line is the same, only whether the line itself counts differs.
Those names are used again in Chapter 11 and Chapter 13, so we take only what we need
without building a topology in full.

An interval is a piece of the line with no holes. The following definition writes the
nine possible forms; all are determined by two endpoints and by the decision whether
each endpoint counts.

\begin{definisi}[Interval]
Let $a<b$ be real numbers. The set $(a,b)=\{x: a<x<b\}$ is called an open interval,
$[a,b]=\{x: a\le x\le b\}$ a closed interval, and $[a,b)$, $(a,b]$ are mixed forms. The
sets $(a,\infty)=\{x: x>a\}$, $[a,\infty)$, $(-\infty,b)$,
$(-\infty,b]$, and $\mathbb{R}$ itself are called unbounded intervals. All nine forms
are called intervals.
\end{definisi}

The feature that marks them all: if two points lie in the interval, the whole segment
between them lies there too, so intervals are called hole-free. The list above is not
proved complete; the classification only inspects upper bound, lower bound, and whether
that bound is included. The notion of a point being ``near'' is expressed by a
neighborhood, that is, a small open interval around one point.

\begin{definisi}[Neighborhood, interior point, open, closed]\index{open set}\index{closed set}\index{accumulation point}
For $\varepsilon>0$, the neighborhood of a point $x$ is
$V_\varepsilon(x)=(x-\varepsilon,x+\varepsilon)$. A point $x$ is called an interior
point of $S\subseteq\mathbb{R}$ if $V_\varepsilon(x)\subseteq S$ for some
$\varepsilon>0$; $S$ is called open if every member is an interior point. A point $x$ is
called an accumulation point of $S$ if every $V_\varepsilon(x)$ contains a member of
$S$ different from $x$. A set $S$ is called closed if its complement is open.
\end{definisi}

A few examples organize the intuition. The interval $(0,1)$ is open and not closed,
because every neighborhood of $0$ meets it. The interval $[0,1]$ is closed and not open,
while $[0,1)$ is neither. The set $\mathbb{Q}$ is also neither, because every
neighborhood contains rationals and irrationals. The set $\{1/n:n\in\mathbb{N}\}$ is not
closed: $0$ lies in its complement, but every neighborhood of $0$ contains some $1/n$.
The set $\mathbb{Z}$ is closed, because its complement is the union of the open
intervals $(k,k+1)$.

So ``closed'' is not the negation of ``open'': $\mathbb{R}$ and $\varnothing$ are both
open and closed at once, and the complement of an open set is always closed. The density
of $\mathbb{Q}$ makes these examples sharper. Every neighborhood contains rationals and
irrationals, so no point is an interior point of $\mathbb{Q}$ or of its complement: that
set is dense and hole-ridden everywhere at once. Completeness therefore cannot be
replaced by density.

A neighborhood is itself open, and that is what makes the definition of an interior
point reusable.

\begin{proposisi}
(i) Every neighborhood $V_\varepsilon(x)$ is open. (ii) The union of an arbitrary
family of open sets is open. (iii) The intersection of finitely many open sets is open.
\end{proposisi}

\begin{proof}
(i) Take $y\in V_\varepsilon(x)$, so $|y-x|<\varepsilon$. Set
$\delta=\varepsilon-|y-x|>0$. For every $z\in V_\delta(y)$,
$|z-x|\le|z-y|+|y-x|<\delta+|y-x|=\varepsilon$ by the triangle inequality. Hence
$V_\delta(y)\subseteq V_\varepsilon(x)$.

(ii) If $x$ lies in the union, it lies in one of its members, and since that member is
open there is a neighborhood of $x$ contained in it, hence also contained in the union.

(iii) If $x$ lies in the finite intersection $S_1\cap\cdots\cap S_n$, then for each $i$
there is $\varepsilon_i>0$ with $V_{\varepsilon_i}(x)\subseteq S_i$. Take
$\varepsilon=\min\{\varepsilon_1,\ldots,\varepsilon_n\}>0$. Then $V_\varepsilon(x)$ is
contained in every $S_i$, hence contained in the intersection.
\end{proof}

Item (iii) holds only for finite intersections: $\bigcap_n(-1/n,1/n)=\{0\}$ is not open,
because for $x\neq0$ the Archimedean property gives $n$ with $1/n<|x|$.

\section{Worked examples}
\tingkat{intermediate}

Three computations promised earlier are carried out here: Newton's iteration, rational
approximation to error below $10^{-6}$, and a check that $\{q:q^2<2\}$ has no maximum.
All the numbers are computed with exact fractions.

Starting from $x_0=1$, the iteration $x_{k+1}=(x_k+2/x_k)/2$ gives the following
fractions; the denominators grow, and that is the price paid for accuracy.

\begin{center}\small
\begin{tabular}{rlll}
\hline
$k$ & $x_k$ & decimal & relative error \\
\hline
0 & $1$ & $1$ & $2.92893\times10^{-1}$ \\
1 & $3/2$ & $1.5$ & $6.06602\times10^{-2}$ \\
2 & $17/12$ & $1.41666666667$ & $1.73461\times10^{-3}$ \\
3 & $577/408$ & $1.41421568627$ & $1.50183\times10^{-6}$ \\
4 & $665857/470832$ & $1.41421356237$ & $1.12774\times10^{-12}$ \\
5 & $886731088897/627013566048$ & $1.41421356237$ & $6.35896\times10^{-25}$ \\
\hline
\end{tabular}
\end{center}

The relative error is measured against $\sqrt{2}=1.414213562373095\ldots$. The reduction
factor climbs by itself: from step 0 to 1 about $4.8$, then $35$, $1155$, and
$1.3\times10^{6}$ --- that is the signature of squaring. The ratio $e_{k+1}/e_k^2$ takes
the values $0.7071$, $0.4714$, $0.4991$, $0.49999925$, and $0.5$, in agreement with the
formula $e_{k+1}=e_k^2\cdot\frac{s}{2x_k}$ because $x_k$ approaches $s$ from above. At
steps 4 and 5 the decimals agree to eleven digits; the next digits do not appear because
we print twelve significant digits, not because the error stops falling. The fraction at
step 5 has twelve digits with error $6.36\times10^{-25}$, and one additional digit
demands roughly one more step.

The second method uses rational approximations that spend denominators more
economically, namely the continued fraction of $\sqrt2$:
\[
1,\quad \frac32,\quad \frac75,\quad \frac{17}{12},\quad \frac{41}{29},\quad \ldots
\]
with $p_{k+1}=p_k+2q_k$ and $q_{k+1}=p_k+q_k$ starting from $(p_0,q_0)=(1,1)$. The
following table lists their absolute errors.

\begin{center}\small
\begin{tabular}{rlll}
\hline
$p/q$ & decimal & absolute error & $p^2-2q^2$ \\
\hline
$1/1$ & $1$ & $4.14214\times10^{-1}$ & $-1$ \\
$3/2$ & $1.5$ & $8.57864\times10^{-2}$ & $1$ \\
$7/5$ & $1.4$ & $1.42136\times10^{-2}$ & $-1$ \\
$17/12$ & $1.41666666667$ & $2.45310\times10^{-3}$ & $1$ \\
$41/29$ & $1.41379310345$ & $4.20459\times10^{-4}$ & $-1$ \\
$99/70$ & $1.41428571429$ & $7.21519\times10^{-5}$ & $1$ \\
$239/169$ & $1.41420118343$ & $1.23789\times10^{-5}$ & $-1$ \\
$577/408$ & $1.41421568627$ & $2.12390\times10^{-6}$ & $1$ \\
$1393/985$ & $1.41421319797$ & $3.64404\times10^{-7}$ & $-1$ \\
$3363/2378$ & $1.41421362489$ & $6.25218\times10^{-8}$ & $1$ \\
\hline
\end{tabular}
\end{center}

The first approximation whose error is below $10^{-6}$ is $1393/985$ with error
$3.64\times10^{-7}$; the one before it, $577/408$, is still $2.12\times10^{-6}$. The last
column explains why: since $p^2-2q^2=\pm1$,
\[
\left|\frac pq-\sqrt2\right|
=\frac{|p^2-2q^2|}{q^2\left(\frac pq+\sqrt2\right)}
<\frac{1}{2\sqrt2\,q^2},
\]
so the error falls like the inverse square of the denominator. For $q=985$ that bound is
$3.644\times10^{-7}$ and the true error is $3.64404\times10^{-7}$: almost exactly,
because $\frac pq$ is already very close to $\sqrt2$. The denominator 985 is far smaller
than $627013566048$ at Newton step 5, although its error is only about a million times
larger --- that is the difference between a method chasing quadratic error and one
chasing the smallest rational error.

Figure~\ref{fig:tangga-akar} displays that error column as a descending staircase. The
sign of $p^2-2q^2$ alternates, so the fractions jump from below to above root two, and
the ratio of two successive errors approaches $5.8284=(1+\sqrt2)^2$.

\begin{figure}[htbp]
\centering
\begin{tikzpicture}[x=1.02cm,y=1cm,font=\small]
  \foreach \lv in {1,...,7}{
    \draw[layergrey!20] (0.55,\lv) -- (10.75,\lv);
    \node[layergrey,anchor=east,font=\scriptsize] at (0.45,\lv) {$10^{-\lv}$};
  }
  \node[layergrey,anchor=east,font=\scriptsize] at (0.45,0) {$10^{0}$};
  \draw[-{Latex[length=2mm]},ink,thick] (0.55,0) -- (10.95,0);
  \draw[-{Latex[length=2mm]},ink,thick] (0.55,0) -- (0.55,7.85);
  \foreach \k/\lab in {%
    1/{$1/1$},2/{$3/2$},3/{$7/5$},4/{$17/12$},5/{$41/29$},
    6/{$99/70$},7/{$239/169$},8/{$577/408$},9/{$1393/985$},10/{$3363/2378$}}{
    \draw[ink] (\k,0.07) -- (\k,-0.07);
    \node[rotate=32,anchor=north east,font=\scriptsize,layergrey] at (\k,-0.12) {\lab};
  }
  \foreach \k/\v in {1/0.3828,3/1.8473,5/3.3763,7/4.9073,9/6.4384}{
    \fill[accent] (\k,\v) circle (1.8pt);
  }
  \foreach \k/\v in {2/1.0666,4/2.6103,6/4.1418,8/5.6729,10/7.2040}{
    \fill[layerteal] (\k,\v) circle (1.8pt);
  }
  \draw[layergrey,dashed,thick] (0.55,6) -- (10.75,6);
  \node[layergrey,font=\scriptsize,anchor=south east] at (10.7,6.02) {threshold $10^{-6}$};
  \fill[accent] (0.95,7.55) circle (1.8pt);
  \node[right,font=\scriptsize] at (1.05,7.55) {below $\sqrt2$};
  \fill[layerteal] (4.35,7.55) circle (1.8pt);
  \node[right,font=\scriptsize] at (4.45,7.55) {above $\sqrt2$};
  \node[font=\scriptsize,anchor=north] at (5.7,-1.18) {successive fractions};
  \node[font=\scriptsize,rotate=90,anchor=south] at (-0.35,3.9) {$|p/q-\sqrt2|$};
\end{tikzpicture}
\caption{A staircase of rational approximations: the error $|p/q-\sqrt2|$ for ten
successive fractions, logarithmic scale. The $10^{-6}$ threshold is first crossed by
$1393/985$, with error $3.644\times10^{-7}$; after that each step falls by a factor of
$5.8284$.}
\label{fig:tangga-akar}
\end{figure}

The third check concerns the absence of a maximum. One more step from the previous $q$
gives
\[
q'=\frac{149856914148961}{105964844887080}=1.41421349985133266\ldots,
\]
and its square is still below two. Both numbers are rational, both lie in $A$, and each
is larger than the one before: there is no reason for this process to stop.

\begin{tantangan}
Consider $S=\{q\in\mathbb{Q}: q\ge 0,\ q^2\le 2\}$. At first glance this set has a
largest member: the number whose square is exactly two, since the sign $\le$ seems to
include that bound in the set. Prove that $S$ has no maximum, then determine $\sup S$ in
$\mathbb{R}$ and explain why the answer is not a member of $S$.

Hint 1: members of $S$ must be rational. First check whether your candidate for largest
is even allowed to be a member of the set.

Hint 2: for every $q\in S$ with $q>0$ we have $q^2<2$ strictly, because $q^2=2$ is
impossible for a rational. So the shift lemma $h=(2-q^2)/(2q+2)$ applies and produces a
larger member of $S$; the unattainable bound is $\sqrt2$.

If you get stuck, go back to the section ``Why the rational numbers are not enough'' and
to the lemma on shifting.
\end{tantangan}

Finally, the Archimedean property in numbers. For $\varepsilon=10^{-6}$ we need
$n>10^{6}$. At $n=10^{6}$ the reciprocal is exactly $10^{-6}$, which is not less than
$\varepsilon$; at $n=10^{6}+1$ the reciprocal is $9.99999\times10^{-7}$, which does
qualify. One unit in $n$ decides whether the choice is legitimate, and that is why
strict inequalities must be respected.

\kertas{For $x=1393/985$, compute $h=(2-x^2)/(2x+2)$ and check by hand that
$(x+h)^2<2$. Then compute the relative error $|x-\sqrt2|/\sqrt2$ and compare it with
$1/(2\sqrt2\,x^2)$. Explain why the two are almost equal, and why that fact does not
contradict the absence of a maximum in $A$.}

\section{Computational lab}
\tingkat{advanced}

The following code runs with Python 3 and \texttt{mpmath} for high-precision
arithmetic; Newton's iteration uses \texttt{Fraction}, so there is no rounding before
printing.

\begin{lab}
\begin{verbatim}
from fractions import Fraction as F
import mpmath as mp
mp.mp.dps = 40
root2 = mp.sqrt(2)

print("Newton for root two, x0 = 1")
x = F(1)
for k in range(6):
    v = mp.mpf(x.numerator) / x.denominator
    print(k, f"{x.numerator}/{x.denominator}", mp.nstr(v, 12),
          mp.nstr(abs(v - root2) / root2, 6))
    x = (x + F(2) / x) / 2

print()
print("Successive rational approximations, error < 1e-6 ?")
p, q = 1, 1
for k in range(10):
    v = mp.mpf(p) / q
    print(f"{p}/{q}", mp.nstr(v, 12), mp.nstr(abs(v - root2), 6),
          abs(v - root2) < mp.mpf(10) ** -6)
    p, q = p + 2 * q, p + q

print()
print("Test: {rational q : q^2 < 2} has no maximum")
q = F(1393, 985)
for k in range(2):
    h = (2 - q * q) / (2 * q + 2)
    q = q + h
    print(q, mp.nstr(mp.mpf(q.numerator) / q.denominator, 18), q * q < 2)
\end{verbatim}
\end{lab}

The output obtained, verbatim:

\begin{verbatim}
Newton for root two, x0 = 1
0 1/1 1.0 0.292893
1 3/2 1.5 0.0606602
2 17/12 1.41666666667 0.00173461
3 577/408 1.41421568627 1.50183e-6
4 665857/470832 1.41421356237 1.12774e-12
5 886731088897/627013566048 1.41421356237 6.35896e-25

Successive rational approximations, error < 1e-6 ?
1/1 1.0 0.414214 False
3/2 1.5 0.0857864 False
7/5 1.4 0.0142136 False
17/12 1.41666666667 0.0024531 False
41/29 1.41379310345 0.000420459 False
99/70 1.41428571429 7.21519e-5 False
239/169 1.41420118343 1.23789e-5 False
577/408 1.41421568627 2.1239e-6 False
1393/985 1.41421319797 3.64404e-7 True
3363/2378 1.41421362489 6.25218e-8 True

Test: {rational q : q^2 < 2} has no maximum
6625109/4684660 1.41421341143220639 True
149856914148961/105964844887080 1.41421349985133266 True
\end{verbatim}

Three things are readable from that output. The Newton error stops being printed at
twelve digits at step 4, even though the true value still falls to $6.4\times10^{-25}$:
what limits it is the print width, not the mathematics. The $10^{-6}$ threshold is
crossed exactly at $1393/985$, where the column of \texttt{True} begins. Both numbers in
the maximum test are also \texttt{True}: both are rational and both have squares below
two.

That code proves nothing. It checks that the formulas in this chapter give the numbers
stated and that the direction of the inequalities is right; the proofs remain in the
lemmas and theorems. Program output is data, not argument.

\begin{untukdosen}
The teaching order I suggest: the root-finding program in the opening, then the field
axioms with $0\cdot x=0$ proved on the board, then order, absolute value, and
completeness. Do not postpone completeness; the entire content of the chapter depends on
it.

Two misconceptions appear often. First, students think a supremum must be a member of
the set, so they reject $\sup\{1-1/n\}=1$. Second, students think completeness can be
proved from the field and order laws, whereas $\mathbb{Q}$ satisfies both and fails
completeness. Both disappear once students see those two examples side by side.

A blackboard activity that works: each student takes $q=1.4$, computes
$h=(2-q^2)/(2q+2)$ and checks $(q+h)^2<2$; the differing results are precisely what
shows the formula works for every $q$. A readiness signal: students can explain why $A$
has no maximum without computing.
\end{untukdosen}

\section{Summary}
\tingkat{basic}

\begin{ringkasan}
\begin{itemize}
\item The real numbers rest on three groups of axioms: field, order, completeness.
      The first gives algebra, the second comparison, the third closes the holes.
\item From the field laws we derive $0\cdot x=0$, $(-x)y=-(xy)$, and the absence of
      zero divisors.
\item From the order laws we derive the sign rules, including nonnegative squares and
      $1>0$.
\item Absolute value turns distance into inequality: $|x|\le c$ is equivalent to
      $-c\le x\le c$, and $|x+y|\le|x|+|y|$.
\item Completeness guarantees a supremum for every nonempty set that is bounded above,
      and then yields the Archimedean property and the density of $\mathbb{Q}$.
\item $\mathbb{Q}$ satisfies the field and order laws but is not complete:
      $\{q: q^2<2\}$ has no maximum; in $\mathbb{R}$ its supremum is $\sqrt{2}$.
\end{itemize}
\end{ringkasan}

\section{Exercises}
\tingkat{intermediate}

\begin{latihan}
\begin{enumerate}
  \item Prove that in an ordered field, $a<b$ is equivalent to $-b<-a$, and conclude
        that multiplying an inequality by $-1$ reverses its direction.
  \item Prove $a<\frac{a+b}{2}<b$ for $a<b$, then explain why that result does not
        require completeness whereas $\sup\{q:q^2<2\}$ does.
  \item Compute $\sup$ and $\inf$ of the set $S=\{2-1/n : n\in\mathbb{N}\}$, then check
        both conditions of the definition. Also determine whether $S$ has a maximum.
  \item Assess this claim: ``Since every positive number is greater than $1/n$ for some
        $n$, there is no smallest positive real number.'' True, false, or not yet
        proved? Name the part that the proof supports.
  \item With $h=(2-q^2)/(2q+2)$, compute $q+h$ for $q=1$ and $q=7/5$, then check that
        both still lie in $A=\{q\in\mathbb{Q}:q^2<2\}$. Compare the sizes of the
        increases.
  \item Show that if $|x|\le 2$ and $|y|\le 2$ then
        $|x^3-y^3|\le 12\,|x-y|$, and say where the triangle inequality is used.
  \item Prove that if $S$ is open then its complement is closed, and conversely, using
        the definition of a neighborhood.
  \item Compute five steps of Newton's iteration for $\sqrt{3}$ starting from $x_0=2$,
        and report the relative error at each step. Check whether the ratio
        $e_{k+1}/e_k^2$ approaches $1/2$.
  \item Show $\bigcap_{n=1}^{\infty}(-1/n,1/n)=\{0\}$ and explain why that does not
        contradict the closedness of finite intersections.
  \item Prove that between any two real numbers there is always an irrational number.
\end{enumerate}

\noindent\textbf{Hints.} Item 4: the claim contains two different statements, one about
every $\varepsilon$ and one about the existence of a minimum. Item 8: use
$e_{k+1}=e_k^2\cdot\frac{\sqrt3}{2x_k}$ and $x_k\ge\sqrt3$. Item 10: apply the density
of $\mathbb{Q}$ to $x/\sqrt2$ and $y/\sqrt2$, then multiply back.
\end{latihan}

\begin{untukpeneliti}
What is genuinely proved: the completeness axiom does not follow from the field and
order laws, because $\mathbb{Q}$ satisfies both and fails the third. The existence of a
complete ordered field is proved by Rudin via Dedekind cuts, and there it is also stated
without proof that any two complete ordered fields are isomorphic; that claim is not
used here.

Primary sources: Rudin, \emph{Principles of Mathematical Analysis}, Chapter 1
(Definition 1.12, Proposition 1.18, Theorems 1.19--1.21, the Appendix on cuts); Abbott,
\emph{Understanding Analysis}, Chapter 1 (The Axiom of Completeness and its
consequences).

What remains open: the effective content of Newton's iteration. Without an explicit
initial bound, $e_{k+1}\le e_k^2/2$ does not yet give a step count; that requires
Chapter 4 and, for general functions, Kantorovich theory. The chapter's scope: topology
only as far as intervals and neighborhoods, and $\mathbb{R}$ assumed rather than built.
\end{untukpeneliti}

\peta{Chapter 3 in Rudin, \emph{Principles of Mathematical Analysis}: Chapter 1 (Ordered
Sets, Fields, The Real Field, The Least-Upper-Bound Property, the Archimedean property,
density of $\mathbb{Q}$, existence of $n$-th roots, Appendix on Dedekind cuts). In
Abbott, \emph{Understanding Analysis}: Chapter 1 (The Axiom of Completeness and
Consequences of Completeness); open sets and accumulation points from Chapter 3, not
developed here.}

\chapter{Sequences and Limits}

A small program prints one number per line. The rule is only one: the next term is
the square root of two plus the previous term, starting from one. What appears on
the screen: 1.000000, then 1.732051, 1.931852, 1.982890, 1.995718, 1.998929, and
so on. On the twelfth line the screen shows 1.999999738545. That number appears to
stop moving. But appearing and being true are two different things: the program
does not know whether the thousandth line will still shift, and it also does not
know which number this run is heading toward.

The second example is more mischievous. The sequence $(1+1/n)^n$ for $n=10$ has
value 2.5937424601, for $n=100$ has value 2.7048138294, for $n=1000$ has value
2.7169239322, and for $n=10^7$ has value 2.7182816941. The first six digits appear
settled at 2.71828. But the true value at $n=10^7$ is 2.7182816925, so the last
seven digits in that machine output are rounding error, not a sign of closeness.
A table can deceive in two directions at once: it makes a divergent sequence look
calm, and it makes a convergent sequence look off course. Chapter 22 will explain
why machine arithmetic behaves that way; this chapter prepares the tools that let
us not depend on tables.

Chapter 3 already supplied the raw material: the field axioms, the order axioms,
and the supremum property. But the supremum only works once the set of terms has
been collected; it does not tell us whether the sequence stops moving. What we
need is a definition that can be checked on paper, and then theorems that turn
that check into estimated work: the $\varepsilon$-$N$ definition, the algebraic
properties of limits, the squeeze theorem, monotone sequences, subsequences, the
Cauchy criterion, and the limit superior.

\begin{untukanak}
Have you ever walked toward a door under a strange rule: each step covers only
half of the distance that remains? The first step takes you halfway. The second
cuts half of what is left, so you are already three quarters of the way. The third
adds another eighth. Every step moves you forward, yet there is always a remainder
you have not reached. If that rule is kept up, you never actually touch the door.

Even so, there are questions that can be answered with certainty. How many steps
are needed so that the distance to the door is less than a hand's width? The answer
can be computed before you set out, and it does not change when the room becomes
twice as long: you only need a few extra steps. So there are two different things
here. ``Never arrives'' and ``as close as one likes can be approached'' are not two
contradictory statements. In fact both are true at once.

Now the question that makes one curious. What if the rule is changed slightly, say
each step covers a third of the remaining distance, or half plus a little? Can you
still guess where that journey ends, and can its speed still be computed? In this
chapter we will know why such questions have definite answers, and why we never
obtain the answers merely by watching numbers on a screen.
\end{untukanak}

\section{The definition of a limit and how to use it}
\tingkat{basic}

\kamus{sequence}{barisan}{a run of numbers given consecutive numbers, like the house
numbers along one side of a street that never ends}

A limit answers one question: how close the later terms gather, and from when on.
In everyday language, a sequence tends to $L$ when, after a few initial terms are
discarded, all the remaining terms are as close as one likes to $L$.

\begin{definisi}[Convergence]
A sequence $(x_n)$ of real numbers \emph{converges} to $L \in \mathbb{R}$ if for
every $\varepsilon>0$ there is $N \in \mathbb{N}$ such that for every $n>N$ we have
$|x_n-L|<\varepsilon$. We write $x_n \to L$, and $L$ is called the
\emph{limit}\index{limit of a sequence} of the sequence. A sequence that does not
converge to any real number is called \emph{divergent}\index{divergent sequence}.
\end{definisi}

The order of the quantifiers in that definition decides everything. The value
$\varepsilon$ is chosen first, then $N$ follows and may depend on $\varepsilon$.
Reversing that order produces a far stronger and usually false statement: there are
sequences that satisfy ``for every $n$ there is $N$'' vacuously, because $N=n$ can
always be used. A healthy way to read it is as a game between two parties. The
challenger names an $\varepsilon$, however small. The responder must answer with an
$N$ that works for all later terms. The sequence converges when the responder
always has an answer, for every challenge.

Figure~\ref{fig:pita-epsilon} captures that game for $x_n=1+1/n$. Two bands centered
at $L=1$ are drawn with different widths, and each has its own threshold: $N=2$ for
$\varepsilon=0.5$, $N=10$ for $\varepsilon=0.1$. Nothing about the sequence changes;
only the challenger's demand changes.

\begin{figure}[htbp]
\centering
\begin{tikzpicture}[x=0.72cm,y=3.0cm,>=Stealth,font=\footnotesize]
  \fill[layerteal!14] (0,0.5) rectangle (14,1.5);
  \fill[layerorange!20] (0,0.9) rectangle (14,1.1);
  \draw[->,ink] (-0.1,0.35) -- (14.6,0.35) node[right] {$n$};
  \draw[->,ink] (0,0.35) -- (0,2.15) node[above] {$x_n$};
  \draw[dashed,layergrey,thick] (0,1) -- (14.1,1);
  \node[layergrey,anchor=south west] at (0.08,1.02) {$L=1$};
  \foreach \n/\y in {1/2.0000,2/1.5000,3/1.3333,4/1.2500,5/1.2000,6/1.1667,%
      7/1.1429,8/1.1250,9/1.1111,10/1.1000,11/1.0909,12/1.0833,13/1.0769,14/1.0714}
    \fill[accent] (\n,\y) circle (2.2pt);
  \draw[dashed,layerteal,thick] (2,0.35) -- (2,1.52);
  \draw[dashed,layerorange,thick] (10,0.35) -- (10,1.12);
  \node[anchor=north] at (2,0.30) {$N_1=2$};
  \node[anchor=north] at (10,0.30) {$N_2=10$};
  \node[anchor=north west,align=left] at (4.6,2.28)
    {outer band: $\varepsilon=0.5$\\ inner band: $\varepsilon=0.1$};
\end{tikzpicture}
\caption{The sequence $x_n=1+1/n$ entering the $\varepsilon$ band around $L=1$. With
$\varepsilon=0.5$ every term after $n=2$ is already inside the band; with
$\varepsilon=0.1$ the threshold shifts to $n=10$. A narrower band demands a more
distant threshold.}
\label{fig:pita-epsilon}
\end{figure}

\begin{bacakeras}
Two voices take turns. The first asks, the second answers. Read it aloud slowly,
pause at every period, and do not add words in between.

``Longer and longer?'' \ ``Closer and closer.'' \ ``How close?'' \ ``As close as
requested.'' \ ``What if one thousandth is requested?'' \ ``Wait a thousand steps,
then measure.'' \ ``What if one hundredth of a millionth?'' \ ``Wait one hundred
million steps. Patience costs nothing.'' \ ``So it will arrive?'' \ ``No. It never
arrives.'' \ ``Then what is the use of waiting?'' \ ``Because the distance shrinks,
and the one who asks is free to choose how small.'' \ ``And if the requests stop?''
\ ``Then there is nothing to prove.''

The line repeated three times, more slowly each time: close is not arrived; close
is a distance we are allowed to set ourselves.
\end{bacakeras}

For $x_n=1/n$ the answer is immediate: take $N$ a natural number with
$N>1/\varepsilon$, which exists by the Archimedean property in Chapter 3. For $n>N$
we have $|1/n-0|=1/n<1/N<\varepsilon$. If the challenger asks for
$\varepsilon=10^{-3}$, the answer is $N=1000$; at $n=1000$ the error is exactly
$10^{-3}$ and does not yet meet the strict condition, whereas at $n=1001$ the error
is $0.000999$. This one-term difference is not pedantry: the definition uses $n>N$,
not $n \ge N$, so the $N$ that is chosen must genuinely pass the threshold.

\begin{tangan}
Materials: duct tape, a tape measure, chalk, and one sheet of notes. Mark a number
line on the floor two meters long: zero at the left end, two at the right end, and
the number one exactly in the middle. Stick down two horizontal strips as wide as a
foot, centered at the number one: a wide band from one half to one and a half
meters, a narrow band from nine tenths to one point one meters.

Now walk to the position one plus one over n, for n equal to one, two, three, up to
fourteen. Each time, note whether your heel lands in the wide band, in the narrow
band, or outside both, then walk back so the marks are not stepped on.

What should be visible: the first two steps are still outside the narrow band, the
tenth step stops exactly at its edge, and after that no heel comes out again. A
narrower band demands more steps, while the number one is never touched by a heel.
\end{tangan}

The use of a definition is felt most strongly when we want to prove that something
does \emph{not} converge. The negation of the definition reads: there is
$\varepsilon>0$ such that for every $N$ there is $n>N$ with $|x_n-L| \ge
\varepsilon$. Note that $L$ is still a free name there; to refute convergence we
must do it for every $L$.

\begin{contoh}
The sequence $x_n=(-1)^n$ is divergent.
\end{contoh}

\begin{proof}
Suppose $x_n \to L$ for some $L$. Take $\varepsilon=1$; there is $N$ such that
$|x_n-L|<1$ for all $n>N$. Choose an even $n>N$; then $x_n=1$ and $|1-L|<1$. The
next term $n+1$ also exceeds $N$ and is odd, so $x_{n+1}=-1$ and $|-1-L|<1$. The
triangle inequality gives
\[
2 = |1-(-1)| \le |1-L| + |L-(-1)| < 1+1 = 2,
\]
an impossibility. So no $L$ satisfies the definition.
\end{proof}

The proof just given shows a pattern that will appear often: to refute, we choose an
$\varepsilon$ smaller than the distance between two groups of terms. The even and
odd groups are a fixed distance $2$ apart, so no single $N$ can separate them when
$\varepsilon<1$.

The first property worth recording is uniqueness. In everyday words: a sequence
cannot stop in two different places at once.

\begin{teorema}[Uniqueness of the limit]
If $x_n \to L$ and $x_n \to M$, then $L=M$.
\end{teorema}

\begin{proof}
Suppose $L \neq M$ and let $\varepsilon=|L-M|/2>0$. There is $N_1$ with
$|x_n-L|<\varepsilon$ for $n>N_1$, and $N_2$ with $|x_n-M|<\varepsilon$ for
$n>N_2$. For $n>\max(N_1,N_2)$ both inequalities hold at once, so
\[
|L-M| \le |L-x_n| + |x_n-M| < 2\varepsilon = |L-M| .
\]
The inequality $|L-M|<|L-M|$ is impossible, so the supposition is false.
\end{proof}

The phrase ``for $n$ sufficiently large'' means: there is $N$ such that the stated
property holds for all $n>N$. With that vocabulary, the definition of a limit in
brief: $x_n$ is as close as one likes to $L$ when $n$ is sufficiently large.

The second property: a convergent sequence cannot run off to infinity.

\begin{teorema}
Every convergent sequence is bounded.
\end{teorema}

\begin{proof}
Suppose $x_n \to L$. Take $\varepsilon=1$; there is $N$ with $|x_n-L|<1$ for
$n>N$, so $|x_n| \le |x_n-L|+|L| < |L|+1$. Only the first term up to term $N$ is
still uncontrolled, and there are finitely many of those. Hence
\[
M=\max\{|x_1|,|x_2|,\dots,|x_N|,|L|+1\}
\]
is a bound that works for every $n$.
\end{proof}

The converse does not hold: $(-1)^n$ is bounded but divergent. Boundedness is a
necessary condition, not a sufficient one; from a table that looks calm we may only
conclude boundedness.

\section{Algebraic properties of limits}
\tingkat{intermediate}

In everyday language: a limit is polite toward addition, multiplication, and
division, as long as the divisor does not tend to zero.

Before proving those properties we need one small control: a sequence tending to a
nonzero number eventually stays away from zero.

\begin{lemma}
If $y_n \to y$ with $y \neq 0$, then there is $N$ such that $|y_n|>|y|/2$ for all
$n>N$.
\end{lemma}

\begin{proof}
Take $\varepsilon=|y|/2>0$. There is $N$ with $|y_n-y|<|y|/2$ for $n>N$. Then
$|y_n| \ge |y|-|y_n-y| > |y|-|y|/2 = |y|/2$.
\end{proof}

The consequence for computation: limits can be computed term by term, and the only
exception is division by zero.

\begin{teorema}[Algebraic properties of limits]
If $x_n \to x$ and $y_n \to y$, then
$x_n+y_n \to x+y$, \ $x_n y_n \to xy$, and $c x_n \to cx$ for every
$c \in \mathbb{R}$. If $y \neq 0$ and $y_n \neq 0$ for all $n$, then
$x_n/y_n \to x/y$.
\end{teorema}

\begin{proof}
For the sum, take $\varepsilon>0$. Choose $N_1$ with $|x_n-x|<\varepsilon/2$ for
$n>N_1$ and $N_2$ with $|y_n-y|<\varepsilon/2$ for $n>N_2$. For
$n>\max(N_1,N_2)$,
$|(x_n+y_n)-(x+y)| \le |x_n-x|+|y_n-y| < \varepsilon$. The scalar multiple case is
similar and shorter: for $c=0$ the sequence is constant; for $c \neq 0$ use
$\varepsilon/|c|$.

For the product, write the difference as two parts,
\[
x_n y_n - xy = x_n(y_n-y) + (x_n-x)y .
\]
Because $(x_n)$ converges it is bounded: there is $M>0$ with $|x_n| \le M$ for all
$n$. Choose $N$ so that for $n>N$ we have $|y_n-y|<\varepsilon/(2M)$ and
$|x_n-x|<\varepsilon/(2(|y|+1))$. Then
\[
|x_n y_n - xy| \le M|y_n-y| + |x_n-x|\,|y| < \frac{\varepsilon}{2}
+ \frac{\varepsilon}{2} = \varepsilon .
\]
The denominator $|y|+1$ is used so that it is never zero when $y=0$; it stays a
positive constant, so it does not change the conclusion.

For the quotient it suffices to show $1/y_n \to 1/y$, then use the product result.
By the lemma, for $n>N_1$ we have $|y_n|>|y|/2$, so
\[
\left|\frac{1}{y_n}-\frac{1}{y}\right| = \frac{|y-y_n|}{|y_n y|}
\le \frac{2}{|y|^2}\,|y_n-y| .
\]
The right side is smaller than $\varepsilon$ as long as $|y_n-y|<\varepsilon|y|^2/2$,
and that can be guaranteed by choosing $N_2$.
\end{proof}

Note the technique that has just appeared: adding and subtracting the same term to
split a difference into two small parts. Almost every convergence proof involving a
product uses it, and the threshold taken is always the largest one.

The next property concerns order, not operations. In everyday language: a limit
cannot jump over a fence that all terms obey.

\begin{teorema}[Order property]
If $x_n \to x$ and $x_n \ge 0$ for all $n$, then $x \ge 0$. More generally, if
$x_n \le y_n$ for all sufficiently large $n$, $x_n \to x$, and $y_n \to y$, then
$x \le y$.
\end{teorema}

\begin{proof}
Suppose $x<0$ and take $\varepsilon=-x>0$. There is $N$ with $|x_n-x|<\varepsilon$
for $n>N$, so in particular $x_n < x+\varepsilon = 0$. That contradicts
$x_n \ge 0$. For the second statement, apply the first part to $y_n-x_n \ge 0$,
whose limit is $y-x$ by the previous theorem.
\end{proof}

Note that the legitimate conclusion is $x \ge 0$, not $x>0$: the sequence $1/n$ is
entirely positive but its limit is zero. The order property also must not be used
backwards. From $x_n<y_n$ we only obtain $x \le y$; a strict gap can vanish in the
limit.

With the algebraic properties, limits of rational sequences can be computed as
usual. For
\[
x_n=\frac{3n^2+2n}{n^2-5},
\]
divide numerator and denominator by $n^2$:
$x_n=\dfrac{3+2/n}{1-5/n^2}$. Because $1/n \to 0$ and $1/n^2 \to 0$, the numerator
tends to $3$ and the denominator to $1 \neq 0$, so $x_n \to 3$. The numerical values
do creep there: $x_{10}=3.368421$, $x_{100}=3.021511$. Another example,
\[
y_n=\frac{2n^2-n}{3n^2+7} \longrightarrow \frac{2}{3},
\]
with $y_{10}=0.618893$ and $y_{100}=0.663179$; both are still far from $2/3$ at
$n=100$ because the error is about $1/n$.

The algebraic properties speak only about limits that already exist; applying them
to a sequence not yet known to converge is an expensive mistake.

\section{The squeeze theorem and comparison sequences}
\tingkat{intermediate}

Sometimes the limit of a sequence cannot be computed directly, but it can be trapped
from above and from below by two sequences with the same limit. In everyday
language: when two guards walk close together with a child between them, the child
must stop in the same place as they do.

\begin{teorema}[Squeeze theorem]
\index{squeeze theorem}
Suppose $a_n \le x_n \le b_n$ for all sufficiently large $n$. If
$a_n \to L$ and $b_n \to L$, then $x_n \to L$.
\end{teorema}

\begin{proof}
Take $\varepsilon>0$. Choose $N_1$ so that $|a_n-L|<\varepsilon$ for $n>N_1$, that
is $a_n>L-\varepsilon$, and $N_2$ so that $b_n<L+\varepsilon$ for $n>N_2$. Also take
$N_3$ so that the squeezing holds for $n>N_3$. For $n>\max(N_1,N_2,N_3)$,
\[
L-\varepsilon < a_n \le x_n \le b_n < L+\varepsilon,
\]
so $|x_n-L|<\varepsilon$.
\end{proof}

Before using it we need one comparison sequence that will keep appearing: powers that
shrink to zero.

\begin{lemma}
If $0<q<1$, then $q^n \to 0$.
\end{lemma}

\begin{proof}
Write $1/q=1+h$ with $h>0$. Bernoulli's inequality, proved by induction in Chapter
1, gives $(1+h)^n \ge 1+nh$. Since both sides are positive,
\[
0 < q^n = \frac{1}{(1+h)^n} \le \frac{1}{1+nh} .
\]
Given $\varepsilon>0$, the Archimedean property gives $n$ with $1+nh>1/\varepsilon$,
that is $1/(1+nh)<\varepsilon$. By the squeeze theorem, $q^n \to 0$.
\end{proof}

Now an example that cannot be settled by dividing by the highest power. For
$x_n=(2^n+3^n)^{1/n}$, the largest term inside the brackets is $3^n$; factor that
power out:
\[
x_n = 3\left(1+\left(\tfrac{2}{3}\right)^{n}\right)^{1/n}.
\]
Because $(2/3)^n \to 0$ and $1 \le 1+(2/3)^n \le 2$, we obtain
$1 \le x_n/3 \le 2^{1/n}$. It remains to show $2^{1/n} \to 1$: write
$2^{1/n}=1+a_n$ with $a_n \ge 0$, then
$2=(1+a_n)^n \ge 1+na_n$, so $0 \le a_n \le 1/n$. The squeeze theorem closes both
ends, and $x_n \to 3$. Python numbers confirm the direction:
$x_1=5$, $x_2=3.605551$, $x_5=3.075152$, $x_{10}=3.005162$, and $x_{100}=3$.

\begin{untukanak}
Two school guards walk side by side along a corridor, and a pupil walks between them.
The guard on the left steps toward the exit, the guard on the right does too, and the
two get closer until their shoulders almost touch. The pupil squeezed in the middle
does not need to know the way; he only needs to stay between the two guards. Wherever
those guards stop, there the pupil stops as well.

Now the question that makes one curious. What if the left guard stops earlier than
the right guard: can we still guess where the pupil stops? And what if their distance
is never zero, only smaller from day to day? Does the pupil still have a definite
place to stop, or does he keep walking with no end?

It turns out the answer does not depend on the faces of the guards, only on two
things: where the left guard is heading, and where the right guard is heading. If
both head for the same place, the pupil in the middle has no other choice. That way
of thinking, ``squeezed between two parties'', is what mathematicians use for
sequences whose formula is too complicated to compute directly.
\end{untukanak}

\section{Monotone sequences and the monotone convergence theorem}
\tingkat{intermediate}

In everyday language: a sequence that moves in only one direction, and never crosses
a certain fence, must stop somewhere. This is where completeness works most cheaply,
because the limit need not be guessed in advance.

\kamus{monotone sequence}{barisan monoton}{a staircase whose steps only go up, never
down}

The everyday language is simple: one direction only, never turning back.

\begin{definisi}[Monotone sequence]
A sequence $(x_n)$ is called \emph{increasing}\index{monotone sequence} if
$x_n \le x_{n+1}$ for all $n$, and \emph{decreasing} if $x_n \ge x_{n+1}$ for all
$n$. A sequence is called monotone if it is increasing or decreasing.
\end{definisi}

A rough picture: what moves in one direction without crossing a fence must stop, and
it stops exactly at the nearest fence it never crosses.

\begin{teorema}[Monotone convergence]
\index{monotone convergence theorem}
An increasing sequence bounded above converges, and its limit is
$\sup\{x_n : n \in \mathbb{N}\}$.
\end{teorema}

\begin{proof}
Suppose $(x_n)$ is increasing and bounded above. Its set of terms is nonempty and
bounded above, so by the supremum property in Chapter 3 it has a supremum; call it
$s=\sup\{x_n\}$. Take $\varepsilon>0$. Because $s-\varepsilon$ is not an upper bound,
there is $N$ with $x_N>s-\varepsilon$. For $n>N$, monotonicity gives
$x_n \ge x_N>s-\varepsilon$, while $x_n \le s$ by the definition of supremum. So
$|x_n-s|<\varepsilon$.
\end{proof}

The immediate consequence: for such a sequence the limit can be read off from its set
of terms, without computing term by term.

\begin{akibat}
A decreasing sequence bounded below converges, and its limit is the infimum of its
set of terms.
\end{akibat}

\begin{proof}
Apply the theorem to $(-x_n)$, which is increasing and bounded above, then use the
scalar multiple property of limits and the identity
$\sup\{-x_n\}=-\inf\{x_n\}$.
\end{proof}

Note that this theorem does not say the limit equals the bound we use. The bound $3$
in the following example is only a proof device; the limit is in fact $e$.

The first application is the recursive sequence that opens this chapter. Suppose
$x_1=1$ and $x_{n+1}=\sqrt{2+x_n}$.

First, all terms lie between $1$ and $2$. For $n=1$ this is clear. If
$1 \le x_n \le 2$, then $3 \le 2+x_n \le 4$, so
$1 \le \sqrt{3} \le x_{n+1} \le 2$. By induction, $1 \le x_n \le 2$ for all $n$.

Second, the sequence is increasing. Because $x_n \le 2$,
\[
x_{n+1} \ge x_n \iff 2+x_n \ge x_n^2 \iff (2-x_n)(1+x_n) \ge 0,
\]
and the last inequality is true since $x_n \le 2$ and $x_n \ge 1$.

By the monotone convergence theorem, $(x_n)$ converges; call its limit $L$. Because
$x_n \le 2$ for all $n$, the order property gives $L \le 2$. Now square the relation:
$x_{n+1}^2=2+x_n$. The left side tends to $L^2$ by the product property, the right
side to $2+L$, so $L^2=2+L$, or $(L-2)(L+1)=0$. Since $L \ge 1$, the only
possibility is $L=2$.

Note that we never use continuity of the square root; the algebraic properties of
limits and the order property are enough, so the machinery of Chapter 11 is not
needed here.

The second application is the sequence $y_n=(1+1/n)^n$. The binomial expansion gives
\[
y_n = \sum_{k=0}^{n} \binom{n}{k}\frac{1}{n^k}
= \sum_{k=0}^{n} \frac{1}{k!}\prod_{j=1}^{k-1}\left(1-\frac{j}{n}\right),
\]
with the empty product equal to $1$ for $k=0$ and $k=1$. Every factor $1-j/n$
increases as $n$ grows, so every term with index $k$ increases; the sum also gains
one new positive term. Hence $(y_n)$ is increasing. Every factor is at most $1$, so
\[
y_n \le \sum_{k=0}^{n}\frac{1}{k!} \le 1+1+\frac12+\frac14+\cdots = 3,
\]
because $1/k! \le 1/2^{k-1}$ for $k \ge 1$, which is checked by induction. The
sequence is increasing and bounded above by $3$, so it converges. We give its limit
the name $e$.

Figure~\ref{fig:monoton-e} displays that sequence for $n=1$ up to $n=30$. The points
rise without ever turning back, the dashed line at $3$ is the fence used in the proof,
and the dotted line at $2.71828$ is where it actually stops. The gap between fence and
destination stays there throughout the figure; the fence works because it is never
crossed, not because it is reached.

\begin{figure}[htbp]
\centering
\begin{tikzpicture}[x=0.32cm,y=3.0cm,>=Stealth,font=\footnotesize]
  \draw[->,ink] (-0.6,1.85) -- (31.2,1.85) node[right] {$n$};
  \draw[->,ink] (0,1.85) -- (0,3.12) node[above] {$y_n$};
  \draw[dashed,layergrey,thick] (0,3.0) -- (30.6,3.0);
  \node[layergrey,anchor=south] at (20,3.02) {proof bound $B=3$};
  \draw[densely dotted,layerindigo,very thick] (0,2.71828) -- (30.4,2.71828);
  \node[layerindigo,anchor=west] at (31.3,2.71828) {limit $e\approx2.71828$};
  \foreach \n/\y in {1/2.0000,2/2.2500,3/2.3704,4/2.4414,5/2.4883,6/2.5216,%
      7/2.5465,8/2.5658,9/2.5812,10/2.5937,11/2.6042,12/2.6130,13/2.6206,%
      14/2.6272,15/2.6329,16/2.6379,17/2.6424,18/2.6464,19/2.6500,20/2.6533,%
      21/2.6563,22/2.6590,23/2.6615,24/2.6637,25/2.6658,26/2.6678,27/2.6696,%
      28/2.6713,29/2.6728,30/2.6743}
    \fill[accent] (\n,\y) circle (1.7pt);
  \draw[->,layerteal,thick] (9.5,2.42) -- (15.5,2.60);
  \node[layerteal,anchor=north] at (12,2.40) {increasing, never turning back};
\end{tikzpicture}
\caption{The sequence $y_n=(1+1/n)^n$ for $n=1$ up to $n=30$. The points rise and
never turn back; the fence $B=3$ is never touched; the value approached is
$e\approx2.71828$, not the fence.}
\label{fig:monoton-e}
\end{figure}

\section{Subsequences, Bolzano--Weierstrass, and the Cauchy criterion}
\tingkat{advanced}

In everyday language: from a bounded sequence one can always pick a run of children
that stops somewhere, even when the parent run is chaotic.

\kamus{subsequence}{subbarisan}{sorting out some letters from a pile without changing
their order}

A subsequence is formed by choosing only some terms: the indices used grow larger,
but any number of terms may be skipped.

\begin{definisi}[Subsequence]
Let $n_1<n_2<n_3<\cdots$ be increasing natural numbers. The sequence $(x_{n_k})_{k
\in \mathbb{N}}$ is called a \emph{subsequence}\index{subsequence} of $(x_n)$.
\end{definisi}

Because $n_k \ge k$ for every $k$, which is checked by induction, a subsequence
cannot be ``slower'' than its parent sequence in terms of index.

In everyday language: choosing some terms does not change the stopping place, it only
slows down or speeds up its arrival.

\begin{proposisi}
If $x_n \to L$, then every subsequence $(x_{n_k})$ also converges to $L$.
\end{proposisi}

\begin{proof}
Take $\varepsilon>0$ and $N$ with $|x_n-L|<\varepsilon$ for $n>N$. For $k>N$ we have
$n_k \ge k>N$, so $|x_{n_k}-L|<\varepsilon$.
\end{proof}

The consequence is simple but most often used to refute: two different stopping places
are already enough.

\begin{akibat}
If $(x_n)$ has two subsequences with different limits, then $(x_n)$ is divergent.
\end{akibat}

\begin{proof}
If $(x_n)$ converged, all its subsequences would converge to the same limit, and
uniqueness of the limit forbids two different values.
\end{proof}

This corollary explains the example $(-1)^n$ without an argument by impossibility: the
even subsequence tends to $1$, the odd subsequence to $-1$. The way it works also
explains why divergence is more often provable than convergence.

To prove Bolzano--Weierstrass we need the following lemma. Its content: every
sequence, however chaotic, contains a monotone subsequence.

\begin{lemma}
Every real sequence has a monotone subsequence.
\end{lemma}

\begin{proof}
Call an index $n$ a \emph{peak} if $x_n \ge x_m$ for all $m>n$, that is, if no later
term is larger. There are two possibilities.

If there are infinitely many peaks, list them as $n_1<n_2<\cdots$. Because $n_1$ is a
peak, $x_{n_1} \ge x_{n_2}$; because $n_2$ is a peak, $x_{n_2} \ge x_{n_3}$; and so
on. The result is a decreasing subsequence.

If there are only finitely many peaks, take $n_1$ beyond the last peak. Because $n_1$
is not a peak, there is $n_2>n_1$ with $x_{n_2}>x_{n_1}$. Because $n_2$ also lies
beyond the last peak, it is not a peak either, so there is $n_3>n_2$ with
$x_{n_3}>x_{n_2}$. The process continues without end and produces an increasing
subsequence.
\end{proof}

The content of this theorem: from a trapped sequence one can always pick some terms
that genuinely tend to a single point.

\begin{teorema}[Bolzano--Weierstrass]
\index{bolzano-weierstrass}
Every bounded real sequence has a convergent subsequence.
\end{teorema}

\begin{proof}
By the lemma, $(x_n)$ has a monotone subsequence. A subsequence of a bounded sequence
is also bounded. A bounded monotone sequence converges by the monotone convergence
theorem. Hence that subsequence converges.
\end{proof}

An alternative proof uses repeated bisection of intervals and the nested interval
property; both proofs lean on completeness, which we borrowed above through the
supremum.

\kamus{Cauchy sequence}{barisan Cauchy}{passengers crowding ever tighter in a train
corridor, without knowing where the exit door is}

Now consider a different problem. The definition of a limit requires us to name $L$
first, whereas in many computations it is $L$ that we are seeking. All we have is a
list of numbers. We therefore need a measure of closeness between the terms
themselves, without naming a destination.

\begin{definisi}[Cauchy sequence]
A sequence $(x_n)$ is called \emph{Cauchy}\index{cauchy sequence} if for every
$\varepsilon>0$ there is $N$ such that for all $m,n>N$ we have
$|x_n-x_m|<\varepsilon$.
\end{definisi}

Note the difference in the order of ideas. In the definition of a limit all terms are
compared with one fixed number $L$; in the Cauchy definition terms are compared with
other terms, with no outside number. That is why the Cauchy definition makes sense in
any space that has a distance, whereas the limit definition demands a candidate
destination.

\begin{teorema}[Cauchy criterion]
\index{cauchy sequence}
A real sequence converges if and only if it is Cauchy.
\end{teorema}

\begin{proof}
Suppose $x_n \to L$ and take $\varepsilon>0$. There is $N$ with
$|x_n-L|<\varepsilon/2$ for $n>N$. For $m,n>N$,
\[
|x_n-x_m| \le |x_n-L| + |L-x_m| < \varepsilon .
\]
So convergence gives Cauchy.

Conversely, suppose $(x_n)$ is Cauchy. Take $\varepsilon=1$; there is $N_0$ such that
$|x_n-x_m|<1$ for $m,n>N_0$. Choosing $m=N_0+1$ gives
$|x_n|<|x_{N_0+1}|+1$ for $n>N_0$, so $(x_n)$ is bounded. By Bolzano--Weierstrass
there is a subsequence $x_{n_k} \to L$. Take $\varepsilon>0$; choose $N_1$ with
$|x_n-x_m|<\varepsilon/2$ for $m,n>N_1$, and $K$ with $|x_{n_k}-L|<\varepsilon/2$ for
$k>K$. Choose $k>K$ that also satisfies $n_k>N_1$; then for $n>N_1$,
\[
|x_n-L| \le |x_n-x_{n_k}| + |x_{n_k}-L| < \varepsilon .
\]
So $x_n \to L$.
\end{proof}

The second direction uses Bolzano--Weierstrass, which uses the monotone convergence
theorem, which uses the supremum. That chain cannot be cut: inside $\mathbb{Q}$ this
theorem fails. The sequence of decimal truncations
$1;\ 1.4;\ 1.41;\ 1.414;\ 1.4142;\ \dots$ tending to the square root of two is Cauchy,
since the difference of two terms is at most $10^{-k}$ when both pass term $k$, yet no
rational number is its limit. If that sequence converged to $L$, then $a_k^2 \to L^2$;
on the other hand $a_k^2 \to 2$, so $L^2=2$. Chapter 3 already showed that no rational
has square $2$. So ``Cauchy'' and ``convergent'' are two different properties; what
makes them agree is completeness, and that is the content of the supremum axiom.

\begin{tantangan}
The sequence $h_n=1+\frac12+\cdots+\frac1n$ grows very slowly: $h_{10}=2.928968$,
$h_{100}=5.187378$, $h_{1000}=7.485471$, and $h_{10^6}=14.392727$. The difference of
two consecutive terms, namely $1/(n+1)$, tends to zero, so this sequence passes the
shallowest test, the one that most often deceives: its terms get closer together. The
question is whether $(h_n)$ converges. Prove your answer, and explain why
$h_{10^6}<15$ is not a proof that the sequence is bounded.

Hint 1: do not inspect terms one by one. Sum in blocks whose lengths double, that is
one term, two terms, four terms, eight terms, and so on.

Hint 2: the $k$-th block contains $2^{k-1}$ terms, and the smallest term in it is at
least $1/2^k$; so every block contributes at least $1/2$. Compare the result with
$h_{2^m}$ and see what happens as $m$ grows.

If you are stuck, return to the Monotone Convergence Theorem: there an increasing
sequence is proved convergent precisely because a fence exists, and here it is
precisely the fence that is missing.
\end{tantangan}

\section{Limit superior and limit inferior}
\tingkat{advanced}

In everyday language: for a sequence that oscillates, we can still name the highest and
the lowest place it oscillates to, even though the sequence itself has no limit.

\kamus{limit superior}{limit superior}{the highest peak still visited by the sequence
as often as one likes after passing any point}

Throughout this section $(x_n)$ is bounded. For every $n$ form the tail set
$\{x_k : k \ge n\}$ and write
\[
s_n=\sup\{x_k : k \ge n\}, \qquad t_n=\inf\{x_k : k \ge n\}.
\]
Because the tails shrink, $(s_n)$ is decreasing and $(t_n)$ is increasing; both are
bounded, so both converge by the corollary to the monotone convergence theorem.

In everyday language: what we track is not where the sequence is, but how high and how
low it still visits after any point.

\begin{definisi}[Limit superior and inferior]
For a bounded sequence $(x_n)$,
\[
\limsup_{n\to\infty} x_n = \lim_{n\to\infty} s_n
\qquad \text{and} \qquad
\liminf_{n\to\infty} x_n = \lim_{n\to\infty} t_n .
\]
The first number is called the \emph{limit superior}\index{limit superior}, the second
the \emph{limit inferior}\index{limit inferior}.
\end{definisi}

Because $s_n \ge t_n$ for every $n$ and limits preserve order, we always have
$\liminf x_n \le \limsup x_n$. The two can differ greatly, and precisely that
difference is the measure of the chaos of the sequence.

\begin{contoh}
For $x_n=(-1)^n+1/n$, the tail contains the even terms $1+1/k$ and the odd terms
$-1+1/k$. The supremum of the $n$-th tail comes from the smallest even term still
available, so $s_1=1.5$, $s_2=1.5$, $s_3=1.25$, $s_4=1.25$, $s_5=1.166667$, and so on
decreasing toward $1$. The infimum of the tail behaves differently: the odd terms
$-1+1/k$ are always present in every tail and decrease toward $-1$ without ever
reaching it, so $t_n=-1$ for all $n$, even though not a single term equals $-1$. Hence
$\limsup x_n=1$ and $\liminf x_n=-1$, while the sequence is divergent.
\end{contoh}

Figure~\ref{fig:divergen-ayun} plots the $n$-th term against $n$ for the same sequence.
The run of even points clings near $1$, the run of odd points clings near $-1$, and the
distance between the two groups is fixed at $2$. No single horizontal line can be
approached by all the points, however far we walk.

\begin{figure}[htbp]
\centering
\begin{tikzpicture}[x=0.62cm,y=1.9cm,>=Stealth,font=\footnotesize]
  \draw[->,ink] (-0.3,0) -- (13.4,0) node[right] {$n$};
  \draw[->,ink] (0,-1.3) -- (0,2.0) node[above] {$x_n$};
  \draw[densely dotted,layerindigo,very thick] (0,1) -- (13.0,1);
  \draw[densely dotted,layerplum,very thick] (0,-1) -- (13.0,-1);
  \node[layerindigo,anchor=south west] at (0.1,1.03) {$\limsup x_n=1$};
  \node[layerplum,anchor=north west] at (0.1,-1.03) {$\liminf x_n=-1$};
  \foreach \n/\y in {1/0.0000,2/1.5000,3/-0.6667,4/1.2500,5/-0.8000,6/1.1667,%
      7/-0.8571,8/1.1250,9/-0.8889,10/1.1000,11/-0.9091,12/1.0833}
    \fill[accent] (\n,\y) circle (2.2pt);
  \draw[<->,layerbrown,thick] (12.5,-1) -- (12.5,1) node[midway,right] {distance $2$};
\end{tikzpicture}
\caption{Graph of the $n$-th term against $n$ for $x_n=(-1)^n+1/n$. The even points
cling near $1$, the odd points near $-1$, and the distance between the two groups is
fixed at $2$. No horizontal band contains all the terms, so this sequence is
divergent.}
\label{fig:divergen-ayun}
\end{figure}

In everyday language: the highest value still visited by the sequence as often as one
likes really is visited by some of its terms, not merely approached from afar.

\begin{teorema}
For a bounded sequence $(x_n)$, its limit superior is the largest subsequential limit:
$\limsup x_n$ equals the supremum of the set of all limits of convergent subsequences,
and that value itself is attained by some subsequence.
\end{teorema}

\begin{proof}
Write $S=\limsup x_n$. First, every subsequential limit is at most $S$. If
$x_{n_k} \to L$, then for every $k$ we have $x_{n_k} \le s_{n_k}$, and because
$n_k \ge k$ and $(s_n)$ is decreasing, $s_{n_k} \le s_k$. So $x_{n_k} \le s_k$, and
taking limits gives $L \le S$.

Second, we build a subsequence tending to $S$, starting from $n_0=0$. For every $k$,
because $s_n \to S$ and $(s_n)$ is decreasing, there is $M \ge k$ with $s_M<S+1/k$.
The tail $\{x_n : n \ge M\}$ has supremum $s_M$ and contains indices as large as we
like, so there is $n_k>\max(k,n_{k-1})$ with $n_k \ge M$ and $x_{n_k}>s_M-1/k$. For
that choice
\[
S - \frac1k \le s_M - \frac1k < x_{n_k} \le s_M < S + \frac1k ,
\]
so $|x_{n_k}-S|<1/k$. Hence $x_{n_k} \to S$.
\end{proof}

The consequence: equality of limit inferior and limit superior becomes the cheapest way
to check convergence without guessing the limit in advance.

\begin{akibat}
A bounded sequence $(x_n)$ converges to $L$ if and only if
$\liminf x_n = \limsup x_n = L$.
\end{akibat}

\begin{proof}
If $x_n \to L$, every subsequence also tends to $L$, so the set of subsequential limits
contains only $L$ and its supremum is $L$; in the same way its infimum is $L$.
Conversely, if both values equal $L$, then $t_n \le x_n \le s_n$ and both fences tend to
$L$, so the squeeze theorem gives $x_n \to L$.
\end{proof}

For unbounded sequences the convention $\limsup x_n=+\infty$ is used when $(x_n)$ is not
bounded above; the arithmetic laws for $\infty$ do not hold automatically, and we do not
use them here.

\section{Worked examples}
\tingkat{intermediate}

We return to the two opening sequences, this time with numbers that can be rechecked and
with the intermediate steps written down.

\begin{resep}
\begin{enumerate}
  \item Write the difference. Form $|x_n-L|$, then simplify until a single fraction in
        $n$ is left; for the example below, $|x_n-3|=5/(n+2)$.
  \item Solve for $n$. Solve $5/(n+2)<\varepsilon$ for $n>5/\varepsilon-2$, and name
        that threshold $g(\varepsilon)$.
  \item Choose $N$. Take $N=\lfloor g(\varepsilon)\rfloor$, the smallest admissible
        integer, then check the error at $n=N$ and $n=N+1$ before writing anything.
  \item Close with the sentence of the definition: for every $\varepsilon>0$ that $N$
        is chosen for, every $n>N$ then satisfies $|x_n-3|<\varepsilon$.
\end{enumerate}
Order warning: $\varepsilon$ is named first, $N$ answers. Setting up $N$ and then
looking for a matching $\varepsilon$ reverses the game; the result is usually an $N$
that depends on $n$, not on $\varepsilon$, and at that point the proof collapses.
\end{resep}

For $x_n=(3n+1)/(n+2)$,
\[
|x_n-3| = \left|\frac{3n+1-3n-6}{n+2}\right| = \frac{5}{n+2}.
\]
The inequality $5/(n+2)<\varepsilon$ is equivalent to $n>5/\varepsilon-2$. Because the
condition is strict, the largest integer still below that threshold is already
admissible: $N=\lfloor 5/\varepsilon-2 \rfloor$ is the smallest choice that satisfies the
definition. For $\varepsilon=10^{-2}$ the threshold is $498$ and $N=498$. Test the
boundary: at $n=498$ the error is $5/500=0.01$, which is \emph{not} smaller than
$\varepsilon$; at $n=499$ the error is $5/501=0.00998004$, which just meets it. For
$\varepsilon=10^{-4}$ we obtain $N=49998$, with error $0.0001$ at $n=49998$ and
$9.9998 \times 10^{-5}$ at $n=49999$.

The terms move very slowly: $n=10$ gives $2.583333$, $n=1000$ gives $2.995010$, and an
error of $10^{-4}$ is only reached at term $49999$. Four digits before the decimal point
are already correct long before the error is small.

\kertas{For $x_n=(3n+1)/(n+2)$, compute a valid $N$ for $\varepsilon=10^{-3}$ with the
formula $N=\lfloor 5/\varepsilon-2 \rfloor$, then check the error at $n=N$ and $n=N+1$.
Explain why the error at $n=N$ may be exactly equal to $\varepsilon$ while the definition
is still satisfied.}

For $y_n=(1+1/n)^n$, 40-digit arithmetic gives error $0.124539$ at $n=10$, $0.00135790$
at $n=10^3$, and $1.35914 \times 10^{-7}$ at $n=10^7$. The product of $n$ with its error
is nearly constant: $1.245394$, $1.357896$, then creeping to $1.3591409 = e/2$. So the
error falls like $1/n$, not $1/n^2$; one extra digit after the decimal point demands an
$n$ ten times larger. This is where the machine table at the opening of the chapter
breaks: at $n=10^7$ float64 reports $1.34327 \times 10^{-7}$, because a small relative
error in $1+1/n$ is amplified by the power $n$.

For the recursion $x_{n+1}=\sqrt{2+x_n}$ with $x_1=1$, the error relative to $2$ is
$2.679 \times 10^{-1}$ at $n=2$, $1.711 \times 10^{-2}$ at $n=4$, and
$2.615 \times 10^{-7}$ at $n=12$. The ratio of two consecutive errors settles at
$0.250000$.

The number $1/4$ is no coincidence. From the relation,
\[
x_{n+1}-2 = \frac{(2+x_n)-4}{\sqrt{2+x_n}+2}
= \frac{x_n-2}{\sqrt{2+x_n}+2},
\]
so
\[
|x_{n+1}-2| = \frac{|x_n-2|}{\sqrt{2+x_n}+2}
\le \frac{|x_n-2|}{2+\sqrt{2}} = 0.292893\,|x_n-2| ,
\]
because $x_n \ge 1$ gives $\sqrt{2+x_n} \ge \sqrt{3}$. The bound $0.292893$ holds for
all $n$, and the actual ratio tends to $1/4$ because the denominator tends to $4$. The
practical consequence: error $10^{-6}$ is reached at $n=12$, $10^{-9}$ at $n=17$, and
machine precision around $n=27$; far cheaper than bisecting an interval again and again.

\section{Computational lab: four checks}
\tingkat{advanced}

\kamus{truncation error}{galat pemotongan}{the difference between a portrait and the
actual face}

Four claims in this chapter can be checked with one short script: the explicit threshold
$N$, the error rate of $(1+1/n)^n$, machine precision, and the error ratio of the
recursion.

\begin{lab}
\begin{verbatim}
import math, mpmath as mp

def N_for(eps):
    N = int(5/eps - 2)
    while 5/(N+2) >= eps: N += 1
    while N > 0 and 5/(N+2) < eps: N -= 1
    return N

print("== 1. explicit N for (3n+1)/(n+2) -> 3 ==")
for eps in (1e-2, 1e-4):
    N = N_for(eps)
    print("eps=%.0e  N=%d  err(N)=%.6g  err(N+1)=%.6g"
          % (eps, N, abs((3*N+1)/(N+2)-3), abs((3*N+4)/(N+3)-3)))

print("== 2. (1+1/n)^n in float64 ==")
for k in range(1, 8):
    n = 10**k
    print("n=1e%d  %.16f  e-y=%.3e" % (k, (1+1/n)**n, math.e-(1+1/n)**n))

print("== 3. true error (mpmath, 40 digits) ==")
mp.mp.dps = 40
for k in (5, 6, 7):
    n = 10**k
    y = (1+mp.mpf(1)/n)**n
    print("n=1e%d  e-y=%s  (float64: %.3e)"
          % (k, mp.nstr(mp.e-y, 6), math.e-(1+1/n)**n))

print("== 4. recursion x_{n+1}=sqrt(2+x_n), x_1=1 ==")
x, xs = 1.0, []
for _ in range(30):
    xs.append(x); x = math.sqrt(2+x)
for n in (1, 2, 3, 4, 5, 6, 11, 12, 13, 14):
    e = abs(xs[n-1]-2)
    r = "" if n == 1 else "  ratio=%.6f" % (e/abs(xs[n-2]-2))
    print("n=%2d  x=%.12f  error=%.3e%s" % (n, xs[n-1], e, r))
e6 = next(n for n in range(1, 31) if abs(xs[n-1]-2) < 1e-6)
e9 = next(n for n in range(1, 31) if abs(xs[n-1]-2) < 1e-9)
print("error < 1e-6 first at n=%d ; error < 1e-9 at n=%d" % (e6, e9))
\end{verbatim}
\end{lab}

\newpage
The output, as it stands:

\begin{verbatim}
== 1. explicit N for (3n+1)/(n+2) -> 3 ==
eps=1e-02  N=498  err(N)=0.01  err(N+1)=0.00998004
eps=1e-04  N=49998  err(N)=0.0001  err(N+1)=9.9998e-05
== 2. (1+1/n)^n in float64 ==
n=1e1  2.5937424601000023  e-y=1.245e-01
n=1e2  2.7048138294215285  e-y=1.347e-02
n=1e3  2.7169239322355936  e-y=1.358e-03
n=1e4  2.7181459268249255  e-y=1.359e-04
n=1e5  2.7182682371922975  e-y=1.359e-05
n=1e6  2.7182804690957534  e-y=1.359e-06
n=1e7  2.7182816941320818  e-y=1.343e-07
== 3. true error (mpmath, 40 digits) ==
n=1e5  e-y=1.35913e-5  (float64: 1.359e-05)
n=1e6  e-y=1.35914e-6  (float64: 1.359e-06)
n=1e7  e-y=1.35914e-7  (float64: 1.343e-07)
== 4. recursion x_{n+1}=sqrt(2+x_n), x_1=1 ==
n= 1  x=1.000000000000  error=1.000e+00
n= 2  x=1.732050807569  error=2.679e-01  ratio=0.267949
n= 3  x=1.931851652578  error=6.815e-02  ratio=0.254333
n= 4  x=1.982889722748  error=1.711e-02  ratio=0.251074
n= 5  x=1.995717846477  error=4.282e-03  ratio=0.250268
n= 6  x=1.998929174953  error=1.071e-03  ratio=0.250067
n=11  x=1.999998954179  error=1.046e-06  ratio=0.250000
n=12  x=1.999999738545  error=2.615e-07  ratio=0.250000
n=13  x=1.999999934636  error=6.536e-08  ratio=0.250000
n=14  x=1.999999983659  error=1.634e-08  ratio=0.250000
error < 1e-6 first at n=12 ; error < 1e-9 at n=17
\end{verbatim}

Three things need to be read from that output. First, the error at $n=N$ is exactly
equal to $\varepsilon$ in both lines; that is why the definition uses the strict
condition $n>N$. Second, the $e-y$ column shrinks by a factor of ten each time $n$ grows
tenfold up to $n=10^6$, then stops following the pattern; part 3 shows that what stops
is float64 precision, not the convergence of the sequence. The truncation error keeps
shrinking like $e/(2n)$, but the rounding error takes over as soon as the two meet.
Third, the error ratio in part 4 falls to $0.250000$ and stays there, exactly as in the
hand computation before. The program proves nothing: it only checks claims we have
already proved.

\begin{untukdosen}
Teaching order: start with the $\varepsilon$-$N$ game in class, not with the theorems.
Point to two rows of students; one names $\varepsilon$, one answers $N$ for
$x_n=(3n+1)/(n+2)$, and the instructor writes it down until the formula appears by
itself. After that come uniqueness, the algebraic properties, the squeeze theorem,
monotone sequences, then subsequences and Cauchy. The limit superior is placed at the
end, as ``a limit for sequences that have no limit''.

Two typical misconceptions. First, students think $N$ may depend on $n$, then write
``for every $n$ there is $N=n$'' and believe the definition is met; the symptom shows
when they cannot produce a single $N$ for the whole tail. Second, they identify ``terms
getting closer together'' with ``convergent'' and use the Cauchy criterion without
checking the space; the decimal truncations of the square root of two inside
$\mathbb{Q}$ are an effective cure.

Board activity: ask students to write the negation of the definition of a limit, then
prove the divergence of $(-1)^n$ with that negation, not by contradiction. The class is
ready to move on when they can produce an explicit $N$ without being given a formula
and name one bounded sequence that is divergent.
\end{untukdosen}

\section{Summary}
\tingkat{basic}

\begin{ringkasan}
\begin{itemize}
  \item $x_n \to L$ means that for every $\varepsilon>0$ there is $N$ such that
        $|x_n-L|<\varepsilon$ for all $n>N$. The $\varepsilon$ is chosen first; $N$
        answers, and may depend on $\varepsilon$ alone.
  \item The limit is unique, and a convergent sequence is bounded. The converse fails:
        boundedness is only a necessary condition.
  \item Limits respect $+$, $\times$, and division by a nonzero denominator; order is
        never jumped over, only weakly preserved.
  \item The squeeze theorem turns trapping into a limit, and the lemma $q^n \to 0$
        supplies a comparison used repeatedly. A bounded monotone sequence always
        converges, and its limit is the supremum or the infimum.
  \item Every bounded sequence has a convergent subsequence; a real sequence converges
        if and only if it is Cauchy, and that equivalence is completeness.
  \item $\liminf \le \limsup$, and the two are equal to $L$ exactly when the sequence
        converges to $L$.
\end{itemize}
\end{ringkasan}

\section{Exercises}
\tingkat{intermediate}

\begin{latihan}
\begin{enumerate}
  \item Prove from the definition that $x_n=(2n-1)/(3n+4)$ converges to $2/3$, and give
        an explicit $N$ as a function of $\varepsilon$. Compare it with the smallest
        $n$ that satisfies the condition.
  \item Prove: if $x_n \to L$ then $|x_n| \to |L|$. Show by example that the converse is
        false.
  \item Compute $\limsup$ and $\liminf$ of $x_n=1+(-1)^n+1/n$ and of
        $y_n=(-1)^n(1+1/n)$. Check your answers with Python for $n$ up to $1000$.
  \item Evaluate this claim: ``Because $(1+1/n)^n$ is increasing and bounded above by
        $3$, its limit is $3$.'' True, false, or not yet enough evidence? Say which part
        is right and which part is wrong.
  \item Prove that $(2^n+3^n)^{1/n} \to 3$ by writing out both fences completely,
        including a proof that $2^{1/n} \to 1$.
  \item For the recursion $x_{n+1}=\sqrt{2+x_n}$ with $x_1=1$, use Python to find the
        smallest term with error below $10^{-6}$ and below $10^{-12}$. Compare with the
        guess from the ratio $1/4$.
  \item Prove that a Cauchy sequence is bounded, without using Bolzano--Weierstrass.
        Then prove that a bounded monotone sequence is Cauchy.
  \item For $x_n=(3n+1)/(n+2)$ and $\varepsilon=10^{-3}$, determine the smallest
        admissible $N$, then compute the error at $n=N$, $N+1$, and $N+2$ with Python.
        Explain why $N-1$ is not admissible.
  \item Prove: if every subsequence of $(x_n)$ has a subsequence converging to $L$, then
        $x_n \to L$.
\end{enumerate}

Hints for the three hardest items. Item 5: use
$1 \le (1+(2/3)^n)^{1/n} \le 2^{1/n}$ and then the lemma $2^{1/n} \to 1$ via Bernoulli's
inequality. Item 7: for the second part, compare $|x_n-x_m|$ with $x_m-x_n$ when $m>n$
and use monotone convergence. Item 9: suppose $x_n$ does not tend to $L$, build a
subsequence staying a fixed distance from $L$, then use Bolzano--Weierstrass.
\end{latihan}

\begin{untukpeneliti}
What is really proved here is only the following equivalence, and all of it rests on one
axiom: the supremum property. In an ordered field, the following statements are
equivalent: every nonempty set bounded above has a supremum; every bounded increasing
sequence converges; every Cauchy sequence converges; every bounded sequence has a
convergent subsequence. The chain of proofs is in Rudin, \emph{Principles of Mathematical
Analysis}, Chapter 3 (the sections Convergent Sequences, Subsequences, Cauchy Sequences)
and Abbott, \emph{Understanding Analysis}, Chapter 2 (sections 2.2 to 2.6). What is
popular but not proved: ``a Cauchy sequence always converges''. That is false without
completeness, and the decimal truncations of the square root of two inside
$\mathbb{Q}$ are the counterexample.

An open question: how fast a sequence converges, and whether that rate can always be
computed. The answer is no. A classical result of Specker (1949) gives a computable,
increasing, bounded sequence of rationals whose limit is not computable; so
``convergent'' does not automatically give an error bound that can be programmed.

Limits of validity: this whole chapter works in $\mathbb{R}$ with the usual distance. In
a general metric space, Bolzano--Weierstrass becomes a statement about compactness, and
the Cauchy criterion remains a necessary condition but no longer sufficient. The squeeze
theorem is also silent when the two fences do not tend to the same value; there we need a
finer comparison, and that is the material of Chapter 5.
\end{untukpeneliti}

\peta{Chapter 4 in Rudin, \emph{Principles of Mathematical Analysis}: Chapter 3 (Numerical
Sequences and Series, with the sections Convergent Sequences, Subsequences, Cauchy
Sequences, Upper and Lower Limits, and Some Special Sequences). In Abbott,
\emph{Understanding Analysis}: Chapter 2 (The Limit of a Sequence, The Algebraic and
Order Limit Theorems, The Monotone Convergence Theorem and a First Look at Infinite
Series, Subsequences and the Bolzano--Weierstrass Theorem, The Cauchy Criterion).
The axioms and completeness consequences used here are built in Chapter 3 of this book,
not repeated.}

\chapter{Infinite Series}

A team trains a classifier on one million examples. The objective function is the sum of
the per-example losses, and training stops automatically when the increase in the total
falls below $10^{-7}$ per epoch. At the fortieth epoch the process stops by itself. The
model remains bad: each new example contributes only about $10^{-6}$ to the total, so the
total hardly moves, while that total itself is not necessarily heading toward a finite
value.

The suspicion can be tested with the simplest series. The sum
$H_n = 1 + \tfrac{1}{2} + \cdots + \tfrac{1}{n}$ at $n=10^6$ equals
$14.3927267229$. Adding another million terms raises it only to
$15.0858736534$; the addition is $0.6931469306$, almost exactly $\ln 2 = 0.6931471806$.
Each time the number of terms doubles, the harmonic series rises by about $\ln 2$, without
bound: the estimate $\ln n + \gamma$ gives $H_{10^9} \approx 21.3004815$. A stopping rule
based on ``the total has stopped moving'' therefore cannot tell convergence from
divergence; what tells them apart is the rate at which the terms shrink.

Chapter 4 already supplied limits, the Cauchy criterion, monotone convergence, and the
limit superior. Series add no new idea, because a partial sum is merely a sequence. But
partial sums rarely have a closed formula, and computing $s_{10^6}$ does not tell us
whether $s_{10^7}$ is still moving. This chapter turns that question into one that can be
answered by inspecting $a_n$ alone.

\begin{untukanak}
Imagine two jars of coins. In the first jar you save under a strange rule:
one coin on the first day, half a coin on the second, a quarter on the third, and so on,
always half of the day before. You save every day, never taking a day off, for as long as
you like. How full is that jar after a hundred years? The answer is surprising: never more
than two coins. Every addition does shrink, and the sum stops at a number you can name
exactly.

In the second jar the rule is different. The coin on the first day is large, then smaller,
and smaller still. After a month the jar holds only a small heap. After a year it is a
little thicker. You might conclude that this jar too will stop at some number, because the
additions grow thinner. It does not. If you save long enough, the contents will pass any
number you name, even though climbing by a single unit takes years. Two jars, two rules
that look alike, two opposite fates.

The question: without waiting for years, how do we know which jar stops and which does
not? Is there a sign in the coins themselves, rather than in the final result? This
chapter answers it, and the answer is used every day by people who write programs.
\end{untukanak}

\section{Series, partial sums, and series we can compute directly}
\tingkat{basic}

A series\index{series} is a sum that never finishes. Its definition therefore does not use
the word ``sum'' directly, but rather the sequence of running totals\index{partial sums}.
We read the result after some number of terms, then ask where that reading is heading.

\begin{definisi}
Given a sequence $(a_n)$ in $\mathbb{R}$. The $n$-th partial sum is
\[
  s_n = a_1 + a_2 + \cdots + a_n = \sum_{k=1}^{n} a_k .
\]
The symbol $\sum_{n=1}^{\infty} a_n$ is called a series, and that series \emph{converges}
to $s$ when the sequence $(s_n)$ converges to $s$; in that case we write
$\sum_{n=1}^{\infty} a_n = s$ and call $s$ the sum of the series. If $(s_n)$ diverges, the
series is called \emph{divergent}.
\end{definisi}

\kamus{series}{deret}{a sum that never finishes, whose value is read from the running total after some number of terms.}

Note that $s_n - s_{n-1} = a_n$. So a series and a sequence are the same object; only the
question changes: for a sequence the formula for $s_n$ is known, for a series we hold only
$a_n$.

The Cauchy criterion from Chapter 4 can be translated directly, and its translation is the
most basic tool in this chapter.

\begin{teorema}[Cauchy criterion for series]
$\sum a_n$ converges if and only if for every $\varepsilon > 0$ there is an integer $N$
such that
\[
  m \ge n \ge N \quad \Longrightarrow \quad
  \left| \sum_{k=n}^{m} a_k \right| < \varepsilon .
\]
\end{teorema}

\noindent\textbf{Proof.} Since $\sum_{k=n}^{m} a_k = s_m - s_{n-1}$, the condition above is
nothing but the Cauchy condition for the sequence $(s_n)$. By Chapter 4, a real sequence
converges if and only if it is a Cauchy sequence.

This criterion speaks about a cut-off tail of the series, not about a single term. Taking
$m = n$ gives a necessary condition that is very cheap to check.

\begin{akibat}[The $n$-th term test]
If $\sum a_n$ converges, then $a_n \to 0$.
\end{akibat}

\noindent\textbf{Proof.} Take $m=n$ in the Cauchy criterion: for every $\varepsilon>0$
there is $N$ such that $n \ge N$ implies $|a_n| < \varepsilon$. That is exactly the
definition of $a_n \to 0$.

This test only filters. It says that a term that does not tend to zero certainly wrecks
convergence, and it says nothing about terms that do tend to zero. Two quick examples for
the easy side: $\sum \frac{n}{n+1}$ diverges because $\frac{n}{n+1} \to 1$, and
$\sum (-1)^n$ diverges because its terms do not tend to zero. The hard side is the harmonic
series: $\frac{1}{n} \to 0$ but $\sum \frac{1}{n}$ diverges. We prove that in the second
section.

\begin{bacakeras}
Two voices. Read them in turn, then swap roles.

\textbf{Voice A:} I shrink. Every day I am smaller than yesterday.

\textbf{Voice B:} Good. But shrinking does not yet mean stopping.

\textbf{A:} I shrink by half. Half of half. Half again.

\textbf{B:} Then you have a fence. Your sum will stop.

\textbf{A:} I shrink, but only a little. One third. One quarter. One fifth.

\textbf{B:} Careful. Thin does not mean used up.

\textbf{A:} However small my additions, I keep adding.

\textbf{B:} And however long it takes, your sum keeps rising.

Repeat three times, faster each time: shrinking is the business of terms, stopping is the
business of sums. Memorize: what decides is not the term that shrinks, but the remainder
that never runs out.
\end{bacakeras}

The first series that can be computed directly is the geometric series\index{geometric
series}. Each of its terms is obtained by multiplying the previous term by a fixed number.

\begin{teorema}[Geometric series]
For $|r| < 1$,
\[
  \sum_{k=0}^{\infty} r^k = \frac{1}{1-r},
\]
and for $|r| \ge 1$ the series diverges.
\end{teorema}

\noindent\textbf{Proof.} For $r \neq 1$ the identity
\[
  (1-r)\left(1 + r + \cdots + r^n\right) = 1 - r^{n+1},
\]
holds, so $s_n = \frac{1-r^{n+1}}{1-r}$. If $|r|<1$ then $r^{n+1} \to 0$ (Chapter 4,
since $|r|^{n} \to 0$ for $0 \le |r| < 1$), hence $s_n \to \frac{1}{1-r}$. If
$|r| \ge 1$ then $|r|^k \ge 1$ for all $k$, so $s_n$ cannot converge: for $r=1$ clearly
$s_n = n+1$, and for $|r|\ge 1$ with $r\neq 1$ the term $r^k$ does not tend to zero.

With $r = \tfrac{1}{2}$ we get $1 + \tfrac{1}{2} + \tfrac{1}{4} + \cdots = 2$, and with
$r = -\tfrac{1}{2}$ we get $1 - \tfrac{1}{2} + \tfrac{1}{4} - \cdots = \tfrac{2}{3}$. The
second is interesting: its terms alternate in sign and the sum is smaller than twice the
first term, because opposite signs wipe each other out.

The case $r=\tfrac{1}{2}$ has a picture that explains itself: cut the unit square in half,
cut the remainder again, and keep going. Figure~\ref{fig:potong-geometri} shows the pieces
$\tfrac12$, $\tfrac14$, $\tfrac18$ and the remainder $2^{-5} = 0.03125$ after five cuts,
with used area $1 - 2^{-5} = 0.96875$. That remainder shrinks without vanishing; that is
the geometric face of the sentence ``the sum stops''.

\begin{figure}[htbp]
\centering
\begin{tikzpicture}[scale=6.0]
  \draw[fill=layerteal!45,draw=ink] (0.5,0) rectangle (1,1);
  \draw[fill=layerindigo!40,draw=ink] (0,0.5) rectangle (0.5,1);
  \draw[fill=layerorange!50,draw=ink] (0.25,0) rectangle (0.5,0.5);
  \draw[fill=layerjade!45,draw=ink] (0,0.25) rectangle (0.25,0.5);
  \draw[fill=layerplum!40,draw=ink] (0.125,0) rectangle (0.25,0.25);
  \draw[fill=white,draw=ink,densely dashed] (0,0) rectangle (0.125,0.25);
  \draw[ink,thick] (0,0) rectangle (1,1);
  \node[font=\small] at (0.75,0.5) {$\tfrac{1}{2}$};
  \node[font=\small] at (0.25,0.75) {$\tfrac{1}{4}$};
  \node[font=\small] at (0.375,0.25) {$\tfrac{1}{8}$};
  \node[font=\small] at (0.125,0.375) {$\tfrac{1}{16}$};
  \draw[ink,thin,-{Stealth}] (0.06,0.10) -- (-0.30,-0.05);
  \node[font=\scriptsize,anchor=east] at (-0.30,-0.05) {remainder $\tfrac{1}{32}$};
\end{tikzpicture}
\caption{The geometric series of one half as pieces of the unit square. Pieces of size
$\tfrac12$, $\tfrac14$, $\tfrac18$, $\tfrac1{16}$ are laid down in order, and the
remainder still uncovered after four pieces is $\tfrac1{16}$; after five pieces the
remainder is $2^{-5} = 0.03125$ and the used area $1 - 2^{-5} = 0.96875$. The remainder
always shrinks, it never vanishes.}
\label{fig:potong-geometri}
\end{figure}

\begin{tangan}
Materials: one sheet of square paper, a pencil, a ruler.

\begin{enumerate}
\item Fold the paper in half, open it, shade one half. That is the first piece.
\item Fold the unshaded part in half again, open it, shade half of it.
\item Repeat three more times, always only on the part not yet shaded.
\item Check the white part that remains: after five folds only one thirty-second of the
      original paper is left.
\end{enumerate}

What you should see: every new shading is always half of the previous one, and the white
part never vanishes, it only shrinks toward one corner. If the paper is thin enough to fold
ten times, the white part is only $\tfrac{1}{1024} = 0.0009765625$; by the computation it
tends to zero, so the shading finally covers the whole sheet. Now change the rule: shade
$\tfrac12$ of the paper, then $\tfrac13$, then $\tfrac14$. The folds can no longer follow,
and the paper runs out before the series stops.
\end{tangan}

The second family that can be computed directly is the telescoping series\index{telescoping
series}, that is, a series whose terms can be written as the difference of two consecutive
terms of another sequence, so that almost all subtractions cancel each other.

\begin{contoh}
For $n \ge 1$,
\[
  \frac{1}{n(n+1)} = \frac{1}{n} - \frac{1}{n+1}.
\]
Hence
\[
  s_N = \sum_{n=1}^{N} \left(\frac{1}{n} - \frac{1}{n+1}\right)
      = 1 - \frac{1}{N+1},
\]
because all the intermediate terms cancel and only the first and last terms remain. So
$\sum_{n=1}^{\infty} \frac{1}{n(n+1)} = 1$. At $N=100$ that formula gives
$s_{100} = \tfrac{100}{101} = 0.9900990099$, and the remaining tail is
$\tfrac{1}{101} = 0.0099009901$; the two add up to exactly $1$.
\end{contoh}

Notice that in the telescoping example we need no test at all. A closed formula for $s_N$
gives the limit directly. This is what rarely happens. For almost every series, including
the harmonic series and $\sum 1/n^2$, no such formula exists, and we must decide
convergence without ever knowing the sum exactly.

Before turning to the tests, we record that series behave linearly, just as limits of
sequences do.

\begin{proposisi}
If $\sum a_n = A$ and $\sum b_n = B$ converge, then for every $c, d \in \mathbb{R}$ the
series $\sum (c a_n + d b_n)$ converges with sum $cA + dB$. Adding or discarding finitely
many terms does not change convergence, it only changes the sum.
\end{proposisi}

\noindent\textbf{Proof.} The partial sums of the combined series are
$c s_n^{(a)} + d s_n^{(b)}$, and by the algebraic properties of limits in Chapter 4 that
sequence converges to $cA + dB$. For the second statement, suppose terms are discarded
starting at index $M$. The partial sums of the new series are $s_{N} - s_{M-1}$ for some
$N$, that is, the difference of a convergent sequence and a constant, so the convergence is
the same.

As an application, $\sum_{n=1}^{\infty} \left(\frac{1}{2^n} + \frac{1}{n(n+1)}\right) = 1
+ 1 = 2$, because the first part is a geometric series with $r=\tfrac{1}{2}$ starting at
$n=1$.

\section{Series of nonnegative terms: comparison and the $p$-series}
\tingkat{intermediate}

When all terms are nonnegative, the partial sums can only rise or stay put. The question
of convergence then becomes a far simpler one: do the partial sums stop below some fence?

\begin{teorema}
If $a_n \ge 0$ for all $n$, then $\sum a_n$ converges if and only if its sequence of
partial sums is bounded.
\end{teorema}

\noindent\textbf{Proof.} With $a_n \ge 0$ we have $s_{n+1} = s_n + a_{n+1} \ge s_n$, so
$(s_n)$ is nondecreasing. By the monotone convergence theorem in Chapter 4, a monotone
sequence bounded above converges, and a monotone sequence not bounded above diverges to
$+\infty$.

This theorem looks simple, but this whole section is only a variation on it. All we need is
a way to compare an unknown series with a known one, and the stock of known series is still
thin: the geometric series.

\begin{teorema}[Comparison test]
Suppose $0 \le a_n \le c_n$ for $n \ge N_0$.
\begin{enumerate}
\item If $\sum c_n$ converges, then $\sum a_n$ converges.
\item If $a_n \ge d_n \ge 0$ for $n \ge N_0$ and $\sum d_n$ diverges, then $\sum a_n$
      diverges.
\end{enumerate}
\end{teorema}

\noindent\textbf{Proof.} For (1), take $\varepsilon>0$. By the Cauchy criterion there is
$N \ge N_0$ such that $m \ge n \ge N$ implies $\sum_{k=n}^{m} c_k < \varepsilon$. Since
every term is nonnegative and $a_k \le c_k$,
\[
  \left| \sum_{k=n}^{m} a_k \right| = \sum_{k=n}^{m} a_k \le \sum_{k=n}^{m} c_k
  < \varepsilon ,
\]
so $\sum a_n$ satisfies the Cauchy criterion. For (2), if $\sum a_n$ converged then by (1)
and $d_n \le a_n$ the series $\sum d_n$ would converge as well, contradicting the
hypothesis.

The comparison test\index{comparison test} demands that we pick the right comparison. For
terms in the form of rational fractions, the natural comparison is the largest term in the
numerator divided by the largest term in the denominator. For instance
$\frac{1}{n^2 + 3n} \le \frac{1}{n^2}$, so $\sum \frac{1}{n^2+3n}$ converges provided
$\sum \frac{1}{n^2}$ converges. Hence the whole question moves to a single family of
series.

That family is the $p$-series\index{p-series}, and the tool for handling it is grouping the
terms into blocks whose size is a power of two. This grouping idea is used repeatedly,
including in the Cauchy condensation test\index{Cauchy condensation} below.

\begin{lemma}[Cauchy condensation]
Suppose $a_1 \ge a_2 \ge a_3 \ge \cdots \ge 0$. Then $\sum_{n=1}^{\infty} a_n$ converges
if and only if $\sum_{k=0}^{\infty} 2^k a_{2^k}$ converges.
\end{lemma}

\noindent\textbf{Proof.} Since the terms are nonnegative, both series converge exactly when
their partial sums are bounded (the first theorem of this section). We compare blocks. For
$k \ge 0$, the block $a_{2^k+1} + \cdots + a_{2^{k+1}}$ contains $2^k$ terms, each no
larger than $a_{2^k}$ and no smaller than $a_{2^{k+1}}$. Thus
\[
  \sum_{n=1}^{2^m} a_n \;\le\; a_1 + \sum_{k=0}^{m-1} 2^k a_{2^k},
  \qquad
  \sum_{n=1}^{2^m} a_n \;\ge\; \frac{1}{2}\sum_{k=1}^{m} 2^k a_{2^k}.
\]
The left inequality shows that the partial sums of $\sum a_n$ are bounded when
$\sum 2^k a_{2^k}$ is bounded. The right inequality shows the converse, because the partial
sums of $\sum 2^k a_{2^k}$ are bounded when $\sum a_n$ is bounded. Both directions are
done.

\begin{teorema}[The $p$-series]
$\sum_{n=1}^{\infty} \frac{1}{n^p}$ converges if $p > 1$ and diverges if $p \le 1$.
\end{teorema}

\noindent\textbf{Proof.} For $p \le 0$ the term $\frac{1}{n^p} = n^{-p}$ does not tend to
zero, so the series diverges by the $n$-th term test. For $p>0$ the sequence
$a_n = n^{-p}$ decreases to zero, so the condensation lemma applies and
\[
  \sum_{k=0}^{\infty} 2^k a_{2^k}
  = \sum_{k=0}^{\infty} 2^k \cdot 2^{-kp}
  = \sum_{k=0}^{\infty} \left(2^{1-p}\right)^{k}.
\]
That is a geometric series with ratio $r = 2^{1-p}$. Since $p>1$ exactly when $1-p<0$,
that is exactly when $r<1$, the geometric series theorem gives: convergence for $p>1$,
divergence for $p \le 1$.

The case $p=1$ is the harmonic series\index{harmonic series}, and its proof is now a single
line: the geometric ratio is $2^{1-1} = 1$, so the condensed series $\sum 1$ diverges. We
can read its fragility straight off the blocks: every block of size $2^k$ contributes at
least $2^k \cdot \frac{1}{2^{k+1}} = \tfrac{1}{2}$. Hence
\[
  H_{2^m} \;\ge\; 1 + \frac{m}{2},
\]
so $H_n \to \infty$, but raising $H$ by one requires doubling $n$ many times over: the
growth rate is logarithmic. This is why the number $14.39$ at the opening of this chapter
looks like a bound although it is not.

Figure~\ref{fig:batang-suku} puts the two kinds of terms side by side. The harmonic bars
stay plump: the tenth term is still $0.1$, whereas the tenth geometric term is
$2^{-10} = 0.0009765625$. Sum of the first ten terms: $H_{10} = 2.928968254$ against
$1 - 2^{-10} = 0.9990234375$. The harmonic term is larger at every index, so the size of
the term misleads; what decides is whether the shrinking leaves a finite tail.

\begin{figure}[htbp]
\centering
\begin{tikzpicture}
\begin{axis}[
  width=0.93\linewidth, height=6.0cm,
  ybar, bar width=3.4pt,
  xlabel={index $n$}, ylabel={size of term $a_n$},
  xmin=0.3, xmax=10.7, ymin=0, ymax=1.12,
  xtick={1,2,3,4,5,6,7,8,9,10},
  scaled y ticks=false,
  tick label style={font=\scriptsize},
  label style={font=\small},
  ymajorgrids, grid style={layergrey!20},
  legend style={at={(0.98,0.96)},anchor=north east,draw=none,fill=none,
                font=\scriptsize},
]
\addplot[fill=layerteal!85,draw=layerteal] coordinates {
  (1,1) (2,0.5) (3,0.3333333) (4,0.25) (5,0.2)
  (6,0.1666667) (7,0.1428571) (8,0.125) (9,0.1111111) (10,0.1)};
\addlegendentry{harmonic $1/n$}
\addplot[fill=layerorange!85,draw=layerorange] coordinates {
  (1,0.5) (2,0.25) (3,0.125) (4,0.0625) (5,0.03125)
  (6,0.015625) (7,0.0078125) (8,0.0039063) (9,0.0019531) (10,0.0009766)};
\addlegendentry{geometric $2^{-n}$}
\end{axis}
\end{tikzpicture}
\caption{Bars of the harmonic terms $1/n$ and the geometric terms $2^{-n}$ for
$n=1,\ldots,10$. The harmonic bars are taller at every index and the tenth term is still
$0.1$, whereas the tenth geometric term is $2^{-10} = 0.0009765625$. Sum of the first ten
terms: $H_{10} = 2.928968254$ against $1 - 2^{-10} = 0.9990234375$.}
\label{fig:batang-suku}
\end{figure}

\kamus{harmonic series}{deret harmonik}{a pile that keeps growing although each addition is thinner than the last, so the total height climbs very slowly without bound.}

The $p$-series is not the final boundary. With the same condensation, Rudin shows that
$\sum \frac{1}{n(\log n)^p}$ converges for $p>1$ and diverges for $p<1$, and that adding
one more logarithm shifts the boundary. Those series differ very little yet lie on opposite
sides: there is no tidy boundary line, only tests that hold on a particular class.

\begin{tantangan}
Take $k(n)$ the integer with $2^{k-1} \le n < 2^k$, then set $a_n = \frac{1}{k(n)}$. So
$a_1 = 1$, $a_2 = a_3 = \tfrac12$, $a_4 = \cdots = a_7 = \tfrac13$, and so on. Show that
$a_n \to 0$ but $\sum a_n$ diverges, then explain why the ratio test gives limit superior
$1$ here and is therefore useless.

\textbf{Hint 1.} Do not add the terms one by one. Add one block, and first count how many
terms that block contains.

\textbf{Hint 2.} The $k$-th block contains $2^{k-1}$ terms, all of them equal to
$\frac{1}{k}$, so its contribution is $\frac{2^{k-1}}{k}$. The first ten blocks already
add up to $118.653968$.

If you are stuck, go back to the section on series of nonnegative terms, where blocks of
power-of-two size are used in the Cauchy condensation lemma.
\end{tantangan}

\begin{untukanak}
Try to imagine stacking coins on a table. The first coin is thick. The second is half as
thick. The third a third, the fourth a quarter, and so on: the millionth coin is very thin,
as thin as one millionth of the first. If you stack a million coins, the height of the
stack is about fourteen times the first coin, not a million times. It looks as though the
stack will stop growing, because the next coin adds almost nothing.

Now the question that trips up many people: how many coins does it take for the stack to
reach twenty times the first coin? The answer is about two hundred and seventy-two million
coins. For twenty-one times, one needs about nine billion coins. Each level demands many
times more coins, but no level is impossible. The pile never stops, it only slows down.

This is why computer programs are easily fooled. If you stop stacking after a million
coins, the stack looks as though it no longer moves, when in fact it is not finished. In
this chapter we will know why slow does not mean stopped, and how to tell the difference
before spending the time to stack.
\end{untukanak}

\section{The ratio test and the root test}
\tingkat{intermediate}

The comparison test demands that we already have a comparison series. For terms containing
factorials or powers of $n$, such a comparison is hard to guess, but the comparison between
consecutive terms is easy to compute. The idea: if every term is at most $\beta$ times the
previous one with $\beta<1$, the series shrinks at least as fast as a geometric series, so
it converges.

\begin{teorema}[Ratio test]
Suppose $a_n \neq 0$ for all sufficiently large $n$.
\begin{enumerate}
\item If $\limsup_{n\to\infty} \left| \frac{a_{n+1}}{a_n} \right| < 1$, then
      $\sum a_n$ converges, and in fact converges absolutely.
\item If there is $n_0$ such that $\left| \frac{a_{n+1}}{a_n} \right| \ge 1$ for all
      $n \ge n_0$, then $\sum a_n$ diverges.
\end{enumerate}
\end{teorema}

\noindent\textbf{Proof.} For (1), write $L = \limsup |a_{n+1}/a_n|$ and choose $\beta$ with
$L < \beta < 1$. By the properties of the limit superior in Chapter 4, only finitely many
$n$ satisfy $|a_{n+1}/a_n| > \beta$; so there is $N$ such that $n \ge N$ implies
$|a_{n+1}| \le \beta |a_n|$. By induction,
\[
  |a_{N+k}| \le \beta^{k} |a_N| \qquad (k \ge 0).
\]
The series $\sum \beta^k |a_N|$ converges because $0<\beta<1$, so by the comparison test
$\sum |a_n|$ converges, and by the absolute convergence theorem in the fifth section of
this chapter, absolute convergence implies convergence.

For (2), the condition $|a_{n+1}| \ge |a_n|$ for $n \ge n_0$ makes $(|a_n|)$ nondecreasing
and nonnegative, so it cannot tend to zero. Hence $a_n$ does not tend to zero, and the
$n$-th term test rules out convergence.

\kamus{ratio test}{uji rasio}{measures whether each term shrinks by at least a fixed fraction of the previous one, so that the series shrinks like a geometric series.}
\index{ratio test}

The bound $1$ in the ratio test is a blind spot. For $1/n$ the ratio is
$\frac{n}{n+1} \to 1$ and the series diverges; for $1/n^2$ the ratio is
$\left(\frac{n}{n+1}\right)^2 \to 1$ and the series converges. So a ratio tending to $1$
gives no information; shrinking rates that differ greatly can produce the same ratio.

The next two series show the ratio test working in opposite directions. For
$a_n = \frac{n!}{10^n}$ we have $\frac{a_{n+1}}{a_n} = \frac{n+1}{10}$, which tends to
$+\infty$; the series diverges. For $a_n = \frac{10^n}{n!}$ we have
$\frac{a_{n+1}}{a_n} = \frac{10}{n+1}$, which tends to $0$; the series converges. The sum
of the second series is exactly $e^{10} = 22026.4657948$, whereas the sum of its first $20$
terms is only $21950.3788494$.

The root test\index{root test} uses a different quantity: the $n$-th root of the $n$-th
term. Intuitively it measures the ``average geometric ratio'' of all the terms, not the
ratio between two consecutive terms.

\begin{teorema}[Root test]
\begin{enumerate}
\item If $\limsup_{n\to\infty} \sqrt[n]{|a_n|} < 1$, then $\sum a_n$ converges, and in
      fact converges absolutely.
\item If $\limsup_{n\to\infty} \sqrt[n]{|a_n|} > 1$, then $\sum a_n$ diverges.
\end{enumerate}
\end{teorema}

\noindent\textbf{Proof.} For (1), choose $\beta$ with
$\limsup \sqrt[n]{|a_n|} < \beta < 1$. Only finitely many $n$ satisfy
$\sqrt[n]{|a_n|} > \beta$, so there is $N$ such that $|a_n| \le \beta^n$ for $n \ge N$. The
comparison test with the geometric series $\sum \beta^n$ closes this case.

For (2), if the limit superior is larger than $1$, then $\sqrt[n]{|a_n|} > 1$, that is
$|a_n| > 1$, for infinitely many $n$. Those terms do not tend to zero, so the series
diverges by the $n$-th term test.

The root test is stronger than the ratio test. Let $a_n = 2^{-n}$ for even $n$ and
$a_n = 3^{-n}$ for odd $n$. Then $\sqrt[n]{a_n}$ equals $\tfrac{1}{2}$ or $\tfrac{1}{3}$,
so its limit superior is $\tfrac{1}{2}$ and the root test declares the series convergent.
The ratio swings: at $n=100$, $\frac{a_{101}}{a_{100}}$ equals about
$3.05 \times 10^{17}$, with limit superior $+\infty$ and limit inferior $0$. The ratio test
concludes nothing exactly when the terms swing on different scales.

The ratio test stays popular because factorials and powers of $n$ cancel in the ratio,
whereas the $n$-th root of a factorial demands Stirling's approximation. Order of use: the
$n$-th term test, comparison with a $p$-series or a geometric series, then the ratio test;
the root test is used when the ratio swings or gives $1$.

Now the tests pile up, and the order of checking becomes a skill of its own.

\begin{resep}
\begin{enumerate}
\item Check $a_n \to 0$. If not, the series diverges; done.
\item Recognize the shape: factorials call for the ratio test, powers of $n$ call for the
      root test, rational fractions call for comparison with a $p$-series.
\item If the test you chose gives the bound $1$, conclude nothing. Move to comparison or
      condensation.
\item If the signs alternate and $|a_n|$ decreases to zero, use Leibniz and report the
      error bound $b_{N+1}$.
\end{enumerate}

Warning on order: do not swap step 1 with step 2. The ratio test on a series whose terms do
not tend to zero can give the bound $1$ and mislead, whereas the $n$-th term test settles
it with no computation.
\end{resep}

\section{Alternating series and the Leibniz test}
\tingkat{intermediate}

Up to here all series have nonnegative terms, and there convergence means the partial sums
stop rising. Once the signs alternate, the partial sums can rise and fall, and cancellation
between signs can produce convergence that the size of the terms does not guarantee.

\begin{teorema}[Leibniz test]
Suppose $(b_n)$ is a nonnegative nonincreasing sequence, that is
$b_1 \ge b_2 \ge \cdots \ge 0$, and $b_n \to 0$. Then
$\sum_{n=1}^{\infty} (-1)^{n+1} b_n$ converges.
\end{teorema}

\noindent\textbf{Proof.} Write $s_N$ for the $N$-th partial sum. Consider the partial sums
with even index and with odd index separately. Since
\[
  s_{2m+2} - s_{2m} = b_{2m+1} - b_{2m+2} \ge 0,
\]
the sequence $(s_{2m})$ is nondecreasing. Since
\[
  s_{2m+1} - s_{2m-1} = -b_{2m} + b_{2m+1} \le 0,
\]
the sequence $(s_{2m-1})$ is nonincreasing. Besides, $s_{2m} \le s_{2m-1}$ and
$s_{2m-1} - s_{2m} = b_{2m} \to 0$. So $(s_{2m})$ is nondecreasing and bounded above by
$s_1 = b_1$, hence converges to some $s$; $(s_{2m-1})$ is nonincreasing and bounded below
by $s_2$, hence converges to some $t$; and $s - t = \lim (s_{2m} - s_{2m-1}) = 0$. Both
subsequences tend to the same value, and since the even and odd subsequences together cover
all indices, the whole sequence $(s_N)$ converges to that $s$.

The Leibniz test\index{Leibniz test} gives more than convergence: an error bound that can
be used without knowing the true sum, and that is what makes alternating series useful for
computing constants.

\begin{teorema}[Leibniz error bound]
Under the hypotheses of the Leibniz test and with $s = \sum (-1)^{n+1} b_n$,
\[
  \left| s - s_N \right| \le b_{N+1}.
\]
Sharper: $s$ lies between $s_N$ and $s_{N+1}$, so the sign of the error $s - s_N$ is the
sign of the first discarded term, namely $(-1)^{N+2}$.
\end{teorema}

\noindent\textbf{Proof.} From the earlier proof, $(s_{2m})$ rises toward $s$ and
$(s_{2m-1})$ falls toward $s$. So for every $N$, the number $s$ lies between $s_N$ and
$s_{N+1}$. Since $|s_{N+1} - s_N| = b_{N+1}$, the distance from $s$ to $s_N$ does not
exceed $b_{N+1}$, and the direction of the shift is set by the sign of
$s_{N+1} - s_N = (-1)^{N+2} b_{N+1}$.

The classic example is the alternating harmonic series,
\[
  1 - \frac{1}{2} + \frac{1}{3} - \frac{1}{4} + \cdots = \ln 2 .
\]
The equality with $\ln 2$ is proved via the Taylor series for $\ln(1+x)$ in Chapter 14;
here we only use its convergence, which the Leibniz test guarantees because
$b_n = \frac{1}{n}$ decreases to zero. The number $\ln 2$ appears again, now as the sum of
a convergent series; the difference $H_{2n} - H_n \to \ln 2$ is another version of the same
computation.

Figure~\ref{fig:tangga-leibniz} draws the partial sums of that series as two staircases:
the odd ones descend from $s_1 = 1$, the even ones ascend from $s_2 = 0.5$, and both
bracket $\ln 2 = 0.6931471806$. At $n = 10$ and $n = 11$ the bracket is
$0.6456349206 < \ln 2 < 0.7365440115$, of width $b_{11} = \tfrac{1}{11} = 0.0909090909$,
exactly as the error bound theorem promises.

\begin{figure}[htbp]
\centering
\begin{tikzpicture}
\begin{axis}[
  width=0.93\linewidth, height=6.2cm,
  xlabel={number of terms $n$}, ylabel={partial sum $s_n$},
  xmin=0.4, xmax=15.0, ymin=0.45, ymax=1.06,
  xtick={1,2,3,4,5,6,7,8,9,10,11,12,13,14},
  tick label style={font=\scriptsize},
  label style={font=\small},
  ymajorgrids, grid style={layergrey!20},
  legend style={at={(0.98,0.04)},anchor=south east,draw=layergrey!40,fill=white,
                fill opacity=0.9,text opacity=1,font=\scriptsize},
]
\addplot[color=layerindigo,mark=*,mark size=1.6pt,thick] coordinates {
  (1,1.0) (3,0.8333333) (5,0.7833333) (7,0.7595238)
  (9,0.7456349) (11,0.7365440) (13,0.7301338)};
\addlegendentry{$n$ odd, descending}
\addplot[color=layerteal,mark=square*,mark size=1.6pt,thick] coordinates {
  (2,0.5) (4,0.5833333) (6,0.6166667) (8,0.6345238)
  (10,0.6456349) (12,0.6532107) (14,0.6587052)};
\addlegendentry{$n$ even, ascending}
\draw[dashed,accent,thick] (axis cs:0.4,0.6931472) -- (axis cs:15.0,0.6931472);
\node[font=\scriptsize,text=accent,fill=white,inner sep=1pt]
  at (axis cs:6.4,0.668) {$\ln 2 = 0.6931471806$};
\end{axis}
\end{tikzpicture}
\caption{Partial sums of the alternating harmonic series. The odd subsequence descends,
the even subsequence ascends, and both bracket the limit line $\ln 2 = 0.6931471806$.
At $n = 10$ and $n = 11$: $s_{10} = 0.6456349206$ and $s_{11} = 0.7365440115$, of width
$b_{11} = \tfrac{1}{11} = 0.0909090909$.}
\label{fig:tangga-leibniz}
\end{figure}

The Leibniz bound is useful and is often estimated too optimistically. With $1000$ terms,
the partial sum is $0.7851481635$ whereas $\pi/4 = 0.7853981634$; the actual error is
$2.5000 \times 10^{-4}$, the promised bound
$b_{1001} = \tfrac{1}{2001} = 4.9975 \times 10^{-4}$. The bound is correct but twice the
real error, because the actual error is about $\tfrac{1}{4N}$ whereas the bound is
$\tfrac{1}{2N}$. For an error below $10^{-6}$ we need $N \ge 500000$, and at that $N$ the
error is $5.0000 \times 10^{-7}$.

\section{Absolute convergence, grouping, and rearrangement}
\tingkat{advanced}

Series that converge because their signs cancel are far more fragile than series that
converge because their terms shrink fast enough. That fragility has a name, and this
chapter owes one explanation of it.

\begin{definisi}
The series $\sum a_n$ \emph{converges absolutely} if $\sum |a_n|$ converges. If
$\sum a_n$ converges but $\sum |a_n|$ diverges, the series is called
\emph{conditionally convergent}\index{conditional convergence}.
\end{definisi}

\kamus{absolute convergence}{konvergensi mutlak}{a series that still converges after all its signs are forced positive, so it does not depend on cancellation.}
\index{absolute convergence}

\begin{teorema}
If $\sum a_n$ converges absolutely, then $\sum a_n$ converges and
$\left| \sum a_n \right| \le \sum |a_n|$.
\end{teorema}

\noindent\textbf{Proof.} For every $n \le m$ we have
$\left| \sum_{k=n}^{m} a_k \right| \le \sum_{k=n}^{m} |a_k|$. Since $\sum |a_n|$
converges, the rightmost expression can be made smaller than $\varepsilon$ by the Cauchy
criterion, so $\sum a_n$ satisfies the Cauchy criterion as well. The absolute value
inequality for the sum is obtained by taking limits in
$\left| s_n \right| \le \sum_{k=1}^{n} |a_k|$.

The converse does not hold: the alternating harmonic series converges conditionally,
because $\sum \frac{1}{n}$ diverges. For such series, the order of summation begins to
matter. Before going there, we first look at an operation that is safe, namely
grouping\index{grouping}.

\begin{proposisi}[Grouping]
Suppose $\sum a_n$ converges to $s$, and $n_1 < n_2 < n_3 < \cdots$ is a sequence of
indices. Then the series whose terms are
\[
  A_k = a_{n_{k-1}+1} + a_{n_{k-1}+2} + \cdots + a_{n_k}
\]
(with $n_0 = 0$) also converges to $s$. The converse fails: a convergent grouping does not
guarantee that the original series converges.
\end{proposisi}

\noindent\textbf{Proof.} The partial sums of the grouped series are
$S_K = A_1 + \cdots + A_K = s_{n_K}$, that is, a subsequence of $(s_n)$. A subsequence of a
convergent sequence tends to the same limit (Chapter 4), so $S_K \to s$.

For the converse, take $1 - 1 + 1 - 1 + \cdots$. This series diverges because its terms do
not tend to zero, but grouping in pairs gives $(1-1) + (1-1) + \cdots = 0 + 0 + \cdots$,
which converges to $0$.

That counterexample matters in practice: parentheses may be inserted into a series already
known to converge, but they may not be used to prove it.

Now we change the order, not merely the grouping: a rearrangement renumbers the same terms
so that each term appears exactly once, but at a different position.

\begin{definisi}
Let $\sigma$ be a bijection from the set of positive integers to itself. The series
$\sum a_{\sigma(n)}$ is called a \emph{rearrangement} of $\sum a_n$.
\end{definisi}

\kamus{rearrangement}{susunan ulang}{swapping the order in which goods arrive without changing the goods, so the running total can differ even though the warehouse contents are the same.}
\index{rearrangement}

For an absolutely convergent series, rearrangement is harmless.

\begin{teorema}
If $\sum a_n$ converges absolutely, then every rearrangement of it converges, and all of
them tend to the same sum.
\end{teorema}

\noindent\textbf{Proof.} Write $s = \sum a_n$ and let $\sum a_{\sigma(n)}$ be a
rearrangement with partial sums $s'_N$. Take $\varepsilon > 0$. By the Cauchy criterion
there is $N_0$ such that $m \ge n \ge N_0$ implies $\sum_{k=n}^{m} |a_k| < \varepsilon$; in
particular $\sum_{k=N_0}^{\infty} |a_k| \le \varepsilon$.

Choose $P$ so that the numbers $1, 2, \ldots, N_0-1$ all appear among
$\sigma(1), \ldots, \sigma(P)$. For $N > P$, the terms $a_1, \ldots, a_{N_0-1}$ cancel each
other in the difference $s_N - s'_N$ (both contain those terms), so all that remains are
terms with index $k \ge N_0$. Hence
\[
  \left| s_N - s'_N \right| \;\le\; \sum_{k \ge N_0} |a_k| \;\le\; \varepsilon
  \qquad (N > P).
\]
Since $s_N \to s$ and $\left| s'_N - s \right| \le \left| s'_N - s_N \right| +
\left| s_N - s \right|$, we get $s'_N \to s$.

For series that converge only conditionally, that conclusion collapses completely. The next
theorem is stated without proof; the proof is long and technical, and can be read in Rudin,
in the chapter on sequences and series of numbers, the section ``Rearrangements''.

\begin{teorema}[Riemann rearrangement theorem, statement]
Suppose $\sum a_n$ converges but does not converge absolutely, and suppose $\alpha \le
\beta$ are two real numbers (possibly $\pm\infty$). Then there is a rearrangement
$\sum a_{\sigma(n)}$ whose partial sums have limit inferior $\alpha$ and limit superior
$\beta$. In particular, for every $c \in \mathbb{R}$ there is a rearrangement that
converges to $c$, and there is a rearrangement that diverges.
\end{teorema}

The idea of the proof is this. If $\sum a_n$ converges conditionally, the series of its
positive terms and the series of absolute values of its negative terms both diverge. We
rearrange by rule: add positive terms until we pass $\beta$, then negative terms until we
drop below $\alpha$, then repeat. Because both series diverge, that step can always be
carried out, and because $a_n \to 0$, the overshoot past the boundary gets smaller. The
result oscillates with limit inferior $\alpha$ and limit superior $\beta$. The complete
construction is in Rudin, Theorem 3.54.

The best-known case is the alternating harmonic series rearranged in the pattern two
positive terms then one negative term,
\[
  1 + \frac{1}{3} - \frac{1}{2} + \frac{1}{5} + \frac{1}{7} - \frac{1}{4} + \cdots
\]
This series converges, but not to $\ln 2$: its sum is $\tfrac{3}{2}\ln 2 = 1.0397207708$,
and the computation is in the next section. Changing the order changes the answer even
though the terms are exactly the same. Addition of real numbers is not commutative without
conditions, and that is a concrete warning for anyone who sums fractions on a computer.

This also explains why the tests in the second and third sections require nonnegative terms.
When the signs alternate, shrinking is not enough: the shrinking must be fast enough to beat
the cancellation.

\begin{untukanak}
Imagine two piles of cards: green marked ``in'', red marked ``out'', each with a number.
You move the cards one at a time, and the balance on the screen changes: green adds, red
subtracts. If you change the order of the moves, does the final balance change?

If the greens and the reds are both large so that they cancel each other, the order
matters. Moving the greens first makes the balance spike before it falls; moving the reds
first makes the balance go negative for a while. The final balance can differ too.

There is one safe situation: if the total of all the numbers, ignoring color, stays bounded,
the order has no effect. That is absolute convergence. In this chapter we will know why, and
why a computer can give different answers for sums that look the same.
\end{untukanak}

\section{Worked examples}
\tingkat{intermediate}

This section works through three computations that form the backbone of the chapter. All the
numbers are checked with Python in the lab section; here the intermediate steps are written
so that they can be repeated with paper and pencil.

\medskip\noindent\textbf{1. The harmonic series at one million terms}

We want $H_{10^6}$ without adding a million terms. The tool is the asymptotic form
\[
  H_n = \ln n + \gamma + \frac{1}{2n} - \frac{1}{12n^2} + O\left(\frac{1}{n^4}\right),
\]
where $\gamma = 0.5772156649$ is the Euler--Mascheroni constant. For $n = 10^6$:
\[
  \ln 10^6 = 13.8155105580, \qquad \ln 10^6 + \gamma = 14.3927262229 .
\]
The first correction term $\frac{1}{2n} = 5 \times 10^{-7}$ gives $14.3927267229$, and the
second correction $-\frac{1}{12n^2} = -8.33 \times 10^{-14}$ is far below the precision we
need. So
\[
  H_{10^6} = 14.3927267228657\ldots
\]
Direct summation with \texttt{math.fsum} gives $14.392726722865724$, in agreement to
$10^{-15}$.

More important is how many terms are needed to reach a given bound. From
$H_n \approx \ln n + \gamma$ we get $n \approx e^{H - \gamma}$. For $H = 20$:
\[
  n \approx e^{20 - 0.5772156649} = e^{19.4227843351} = 2.724006 \times 10^{8}.
\]
A direct search gives $n = 272400600$ as the first index with $H_n > 20$, and $H$ at that
point is $20.0000000016$. For $H = 21$ one needs about $7.4 \times 10^8$ terms; the estimate
at $n = 10^9$ gives $21.3004815$. Each added unit demands multiplying $n$ by
$e \approx 2.718$.

\medskip\noindent\textbf{2. The sum $\sum 1/n^2$ and $\pi^2/6$}

With $N = 10^6$ terms, direct summation gives
\[
  s_{10^6} = 1.6449330668487263, \qquad \frac{\pi^2}{6} = 1.6449340668482264 .
\]
The difference is $9.9999950 \times 10^{-7}$, very close to $\frac{1}{N} = 10^{-6}$; that
makes sense because the tail $\sum_{n>N} n^{-2}$ is about $\frac{1}{N}$. The next order
correction is $-\frac{1}{2N^2} = -5 \times 10^{-13}$, so
\[
  s_{10^6} + \frac{1}{N} - \frac{1}{2N^2} = 1.6449340668482262,
\]
which matches $\pi^2/6$ to $2 \times 10^{-16}$, the precision limit of 64-bit fractions. One
million terms plus one correction already gives machine accuracy.

\medskip\noindent\textbf{3. Rearranging the alternating harmonic series}

Take the pattern two positive terms and one negative: in the $k$-th block we add
$\frac{1}{4k-3} + \frac{1}{4k-1} - \frac{1}{2k}$. After $K$ blocks, all the positive terms
with denominator $1, 3, \ldots, 4K-1$ have been used, and the negative terms with
denominator $2, 4, \ldots, 2K$ have been used as well. So
\[
  S_K = \sum_{j=1}^{2K} \frac{1}{2j-1} - \sum_{j=1}^{K} \frac{1}{2j}
      = \left( H_{4K} - \frac{1}{2}H_{2K} \right) - \frac{1}{2}H_K .
\]
With $H_m \approx \ln m + \gamma$,
\[
  S_K \approx \left(\ln 4K - \tfrac{1}{2}\ln 2K\right) - \tfrac{1}{2}\ln K
  = \ln \frac{4K}{\sqrt{2K}\sqrt{K}} = \ln\left(2\sqrt{2}\right) = \tfrac{3}{2}\ln 2 .
\]
Every term $\gamma$ cancels, and the result does not contain $K$. So
$S_K \to 1.0397207708$, not $\ln 2 = 0.6931471806$. At $K = 10^6$ blocks the numerical
value is $1.0397205208$, an error of $2.5 \times 10^{-7}$ against $\tfrac{3}{2}\ln 2$,
consistent with order $\frac{1}{4K}$ as for the Leibniz series.

\kertas{Numerical summary. $H_{10^6} = 14.3927267229$ and
$\ln 10^6 + \gamma = 14.3927262229$, the difference is $\frac{1}{2n} = 5 \times 10^{-7}$.
$H_n > 20$ is first reached at $n = 272400600$, close to $e^{20-\gamma} = 2.724006 \times
10^8$. The Leibniz series with $1000$ terms gives an error of $2.5000 \times 10^{-4}$
against $\pi/4$, with bound $\frac{1}{2001} = 4.9975 \times 10^{-4}$. The sum
$\sum 1/n^2$ up to $10^6$ terms falls short of $\pi^2/6$ by $9.9999950 \times 10^{-7}$; one
Euler--Maclaurin correction closes it to machine precision. The rearrangement two positive
one negative tends to $\frac{3}{2}\ln 2 = 1.0397207708$, not $\ln 2$.}

\section{Computational lab}
\tingkat{advanced}

The following code runs as is on Python 3.14.4 with NumPy 2.5.2. The first part tests the
accuracy of the harmonic sum and of its asymptotic estimate.

\begin{lab}
\begin{verbatim}
import math
import numpy as np

n = 10**6
a = 1.0 / np.arange(1, n + 1)
vals = a.tolist()
naive = 0.0
for x in vals:
    naive += x

print("np.sum   ", repr(float(np.sum(a))))
print("for-loop ", repr(naive))
print("fsum     ", repr(math.fsum(vals)))
small = a.astype(np.float32)
print("float32  ", repr(float(np.sum(small, dtype=np.float32))))
g = 0.5772156649015328606
print("ln n + g ", repr(math.log(n) + g))
print("1/(2n)   ", 1 / (2 * n))
\end{verbatim}
\end{lab}

The output obtained:
\begin{verbatim}
np.sum    14.392726722865726
for-loop  14.392726722864989
fsum      14.392726722865724
float32   14.392724990844727
ln n + g  14.392726222865807
1/(2n)    5e-07
\end{verbatim}

The first four results should be the same, and three of them differ in the last digit:
rounding error. Adding a million numbers of size $10^{-6}$ into an accumulator of size $14$
throws away low bits. \texttt{float32} loses $1.7 \times 10^{-6}$, larger than the
correction $\frac{1}{12n^2}$ we want to measure; when the error we are after is smaller
than the error of the summation itself, the result means nothing.

The second part checks three numbers from the previous section.

\begin{lab}
\begin{verbatim}
import math
import numpy as np

pi4 = math.pi / 4
s = 0.0
for k in range(1000):
    s += (-1)**k / (2*k + 1)
print("Leibniz 1000 terms", repr(s))
print("error", repr(pi4 - s))
print("bound 1/2001", repr(1/2001))

N = 10**6
m = np.arange(1, N + 1, dtype=np.float64)
S = float(np.sum(1.0 / m**2))
print("sum 1/n^2", repr(S))
print("pi^2/6", repr(math.pi**2/6))
print("corrected", repr(S + 1/N - 1/(2*N**2)))

S = 0.0
for k in range(1, 10**6 + 1):
    S += 1.0/(4*k-3) + 1.0/(4*k-1) - 1.0/(2*k)
print("rearrangement", repr(S))
print("1.5 ln 2", repr(1.5*math.log(2)))
\end{verbatim}
\end{lab}

Its output:
\begin{verbatim}
Leibniz 1000 terms 0.7851481634599485
error 0.00024999993749974525
bound 1/2001 0.0004997501249375312
sum 1/n^2 1.6449330668487263
pi^2/6 1.6449340668482264
corrected 1.6449340668482262
rearrangement 1.0397205208399414
1.5 ln 2 1.0397207708399179
\end{verbatim}

Three things need to be read from that output. First, the Leibniz error
$2.5000 \times 10^{-4}$ is exactly half the bound $4.9975 \times 10^{-4}$; that is no
accident, because the remainder of a monotone alternating series approaches half of the
first discarded term. Second, the Euler--Maclaurin correction closes the gap to
$2 \times 10^{-16}$, one rounding unit at magnitude $1.64$. Third, the rearrangement gives
$1.0397205$, not $0.6931472$, and its error against $\tfrac{3}{2}\ln 2$ is only
$2.5 \times 10^{-7}$.

\begin{untukdosen}
The teaching order I recommend: start from the epoch problem in the opening, then the
definition of partial sums and the two series that can be computed (geometric,
telescoping). The tests come in as a sequence of checks, not a list to memorize: the
$n$-th term, comparison, ratio, root. The $p$-series follows after condensation, which
explains why the harmonic series diverges and why it does so slowly.

Two misconceptions appear most often. First, students conclude convergence from
$a_n \to 0$; the harmonic series is the cure, and it is best shown with a table
$H_{10} = 2.9290$, $H_{100} = 5.1874$, and $H_{1000} = 7.4855$ so that the slowness is
visible. Second, students take ``convergent series'' to mean ``partial sums stop changing'',
when the change per term can be far below machine precision while the series diverges. This
is the right place to mention rounding error.

A class activity that works: split the class into groups, give each group four unnamed
series, and the only task is to decide convergence or divergence together with the test,
in ten minutes. The chosen series must include one trap with ratio $=1$. Discuss groups that
use different tests and are both right. Students are ready to move on when they can explain
why $\sum 1/n^2$ converges while $\sum 1/n$ does not, using only block grouping.
\end{untukdosen}

\section{Summary}
\tingkat{basic}

\begin{ringkasan}
\begin{itemize}
\item A series is a sequence of partial sums; converging means that sequence converges.
\item For nonnegative terms, convergence is equivalent to bounded partial sums. From this
      come the comparison test, condensation, and the classification of $p$-series:
      convergent exactly when $p>1$.
\item The ratio test and the root test measure the speed of shrinking relative to a
      geometric series. Both are blind at the bound $1$, and the root test is stronger than
      the ratio test.
\item The Leibniz test gives convergence for alternating series that decrease to zero,
      along with an error bound $b_{N+1}$ whose sign is known.
\item Absolute convergence is immune to grouping and rearrangement. Conditional
      convergence is not: a Riemann rearrangement can turn the sum into any number, or make
      it diverge.
\item The harmonic series diverges logarithmically: $H_{10^6} \approx 14.39$ and it takes
      $272$ million terms to reach $20$. A stopping rule based on ``the total no longer
      changes'' cannot be trusted.
\end{itemize}
\end{ringkasan}

\section{Exercises}
\tingkat{intermediate}

\begin{latihan}
\begin{enumerate}
\item Compute $s_{20}$ for the geometric series $\sum_{k=0}^{\infty}
      \left(\frac{3}{5}\right)^k$, then compare it with the true sum. What is the error,
      and what error bound can you promise from the formula $\frac{r^{21}}{1-r}$?

\item Prove that $\sum_{n=1}^{\infty} \frac{1}{n(n+2)} = \frac{3}{4}$ by splitting the
      term into partial fractions and showing that the series telescopes.

\item Prove: if $a_n \ge 0$ and $\sum a_n$ converges, then $\sum a_n^2$ converges. Show by
      example that the converse is false.

\item Evaluate this claim: true, false, or not yet enough evidence. ``Because
      $\frac{1}{n} \to 0$ and every term is smaller than the terms of the geometric series
      with ratio $\frac{1}{2}$ for $n$ sufficiently large, the harmonic series converges.''
      State exactly where the reasoning breaks.

\item Decide convergence or divergence, naming the test used:
      (a) $\sum \frac{n^2}{2^n}$, (b) $\sum \frac{2^n}{n^2}$,
      (c) $\sum \frac{n!}{n^n}$, (d) $\sum \frac{1}{n^{1/3}}$.

\item Compute the sum $\sum_{n=1}^{10^6} \frac{1}{n^2}$ with NumPy, then estimate the tail
      with $\frac{1}{N} - \frac{1}{2N^2}$ and compare the result with $\pi^2/6$. How many
      digits after the decimal point are correct before the correction, and after it?

\item The series $\sum \frac{(-1)^{n+1}}{\sqrt{n}}$ converges. Determine the smallest $N$
      such that the Leibniz bound guarantees an error below $10^{-3}$, then compute the
      actual error at that $N$ with Python. Is the bound loose or sharp?

\item Prove that if $\sum a_n$ converges absolutely and $(b_n)$ is bounded, then
      $\sum a_n b_n$ converges absolutely.

\item Rearrange the alternating harmonic series in the pattern three positive terms and one
      negative. Guess the sum with the method of the worked examples section, then check it
      numerically. Is the result still a multiple of $\ln 2$?

\item Prove that if $\sum a_n$ converges conditionally, then the series of its positive
      terms diverges and the series of absolute values of its negative terms diverges as
      well. Explain why this fact is the fuel of the Riemann rearrangement theorem.
\end{enumerate}

\noindent\emph{Hints for the three hardest items.} Item 7: the Leibniz bound uses the
first discarded term, so solve $\frac{1}{\sqrt{N+1}} < 10^{-3}$ and check whether that
inequality is sharp enough. Item 9: after $K$ blocks, the positive terms are used up to
denominator $6K-1$, the negative ones up to $\frac{1}{2K}$; write the partial sums in terms
of $H_{6K}$, $H_{3K}$, and $H_{K}$, then take the limit. Item 10: write
$p_n = \max(a_n, 0)$ and $q_n = \max(-a_n, 0)$; then $a_n = p_n - q_n$ and
$|a_n| = p_n + q_n$, so convergence of one of those series says something about
$\sum |a_n|$.
\end{latihan}

\begin{untukpeneliti}
What is really proved in this chapter are the statements with complete proofs: the Cauchy
criterion for series, the comparison test, Cauchy condensation, the classification of
$p$-series, the ratio test and the root test in limit superior form, the Leibniz test
together with its error bound, and the closure of absolutely convergent series under
rearrangement. I quote the Riemann rearrangement statement without proof. The primary
source is Rudin, \emph{Principles of Mathematical Analysis}, chapter ``Numerical Sequences
and Series'': Theorems 3.22--3.23, 3.25, 3.27--3.28, 3.33--3.34, 3.42--3.43, 3.45--3.46,
and 3.54--3.55. Abbott, \emph{Understanding Analysis}, chapter ``Sequences and Series'',
discusses the same material with emphasis on rearrangement in the opening section of that
chapter.

What is often popularized beyond what is proved is the idea of a threshold between
convergent and divergent series. Rudin shows that no such threshold exists: two series that
differ very little can lie on opposite sides, and the procedure can be continued without
end.

An open question: how reliable is a convergence decision when $a_n$ is available only as a
noisy estimate, such as the sum of losses in machine learning? For such series there is not
yet a tidy theory.

Limits of validity: the whole chapter treats series of real numbers. In other normed spaces
the definition is the same but the tests differ; series of functions require uniform
convergence (Chapter 14). All numbers come from NumPy 2.5.2, 64-bit arithmetic.
\end{untukpeneliti}

\peta{Chapter 5 in Rudin, \emph{Principles of Mathematical Analysis}, chapter ``Numerical
Sequences and Series'': series and partial sums (3.21--3.24), series of nonnegative terms
and comparison (3.25--3.29), the root and ratio tests (3.33--3.37), absolute convergence and
rearrangement (3.45--3.55). In Abbott, \emph{Understanding Analysis}, chapter ``Sequences
and Series'': rearrangement in the opening section of that chapter. This chapter does not
reuse the prose of Chapter 4; from there it refers only to the definition of a limit, the
Cauchy criterion, and monotone convergence.}


% ---------------------------- part II ----------------------------
\part{Linear Algebra}

This part installs the framework used by almost every later chapter: vector spaces, bases,
determinants, eigenvalues, and the singular value decomposition. The prerequisite is
Chapters~1--2. After five chapters you can turn a problem into a matrix, solve it, and
explain why its answer is unique or not.
\chapter{Vector Spaces and Subspaces}

The following table holds three records of network traffic; each row is one record, each column
one feature (incoming packets, outgoing packets, duration, failed connections). Read from top to
bottom, the columns are
\[
  \mathbf{c}_{1}=(1,2,3),\quad \mathbf{c}_{2}=(2,4,6),\quad
  \mathbf{c}_{3}=(0,1,2),\quad \mathbf{c}_{4}=(1,3,5),
\]
while the rows are $\mathbf{r}_{1}=(1,2,0,1)$, $\mathbf{r}_{2}=(2,4,1,3)$, and
$\mathbf{r}_{3}=(3,6,2,5)$.

How many columns actually carry new information? Counting the nonzero columns gives four, and
that count is wrong: the second column is twice the first, the fourth is the sum of the first
and the third. The table holds only two independent directions of data. The sentence ``can be
recomputed'' acquires meaning only once we can measure how few vectors build it up.

Chapters 1 through 5 prepared the language of sets, the axioms of the real numbers, and then
limits of sequences and series; all of them work with closeness. The problem above contains no
closeness at all: its entries are exact, and the question is purely algebraic. Which sets are
closed under both operations, and how few vectors build them. One habit from Chapter 5 stays in
use: here every combination involves finitely many terms.

There is a second question that demands the same tool. A linear scorer assigns the score
$s(\mathbf{x})=\mathbf{w}\cdot\mathbf{x}$ with $\mathbf{w}=(1,1,1)$. A direction
$\boldsymbol{\delta}$ that does not change the score satisfies
$\mathbf{w}\cdot\boldsymbol{\delta}=0$, which is a plane; every perturbation splits into a part
that changes the score and a part that does not. Is that split unique? Yes, and the property is
called a direct sum.

\begin{untukanak}
Picture a box of poster paint holding four cans: red, yellow, blue, and orange. We may mix the
contents of the cans in any proportion, and the mixture may itself be used to mix again. Orange
can be made from red and yellow. So the orange can adds not a single capability that the box did
not already have; it only takes up space.

The nagging question: how few cans must remain so that every color we can make stays a color we
can make? The answer is not the same as the number of cans left after we throw one away, and the
order of removal must not change that number. That is no accident, and proving it is not as
simple as it looks.

There is a stranger question. Suppose two people mix according to different recipes, and the
results are exactly the same. Are the recipes the same? Sometimes yes, sometimes no. In this
chapter we will learn why that number has a name of its own, why it does not depend on the order
of removal, and why two different mixtures sometimes yield the same color.
\end{untukanak}

\section{The vector space axioms and their examples}
\tingkat{basic}

A vector space is a set whose members may be added and may be enlarged or shrunk by a number,
with the result always staying inside that set, and with the same rules of computation as for
ordinary pairs of numbers.

\begin{definisi}[Vector space]
\index{vector space}
A nonempty set $V$ together with an addition $V\times V\to V$ and a scalar multiplication
$\mathbb{R}\times V\to V$ is called a \emph{vector space over $\mathbb{R}$} if for all
$\mathbf{u},\mathbf{v},\mathbf{w}\in V$ and all $a,b\in\mathbb{R}$ the following hold:
\begin{enumerate}
  \item $\mathbf{u}+\mathbf{v}=\mathbf{v}+\mathbf{u}$;
  \item $(\mathbf{u}+\mathbf{v})+\mathbf{w}=\mathbf{u}+(\mathbf{v}+\mathbf{w})$;
  \item there is $\mathbf{0}\in V$ with $\mathbf{v}+\mathbf{0}=\mathbf{v}$ for every
        $\mathbf{v}$;
  \item there is $-\mathbf{v}\in V$ with $\mathbf{v}+(-\mathbf{v})=\mathbf{0}$;
  \item $a(\mathbf{u}+\mathbf{v})=a\mathbf{u}+a\mathbf{v}$;
  \item $(a+b)\mathbf{v}=a\mathbf{v}+b\mathbf{v}$;
  \item $a(b\mathbf{v})=(ab)\mathbf{v}$;
  \item $1\mathbf{v}=\mathbf{v}$.
\end{enumerate}
\end{definisi}

Both operations carry the closure requirement inside their type: the result always lies in $V$,
so this definition contains no separate closure axiom. The eight axioms are not independent of
one another either, and we do not chase the shortest list; what we need are their small
consequences, all derived without looking at the shape of the members of $V$.

\begin{lemma}
In any vector space, the zero vector is unique, the additive inverse of every vector is unique,
$0\mathbf{v}=\mathbf{0}$, $(-1)\mathbf{v}=-\mathbf{v}$, and $c\mathbf{v}=\mathbf{0}$ implies
$c=0$ or $\mathbf{v}=\mathbf{0}$.
\end{lemma}

\par\medskip\noindent\textbf{Proof.}\ 
For uniqueness of zero, suppose $\mathbf{0}$ and $\mathbf{0}'$ are both neutral. Then
$\mathbf{0}=\mathbf{0}+\mathbf{0}'=\mathbf{0}'$. For uniqueness of inverses, suppose
$\mathbf{v}+\mathbf{a}=\mathbf{0}$ and $\mathbf{v}+\mathbf{b}=\mathbf{0}$. Then
\[
  \mathbf{a}=\mathbf{a}+\mathbf{0}=\mathbf{a}+(\mathbf{v}+\mathbf{b})
  =(\mathbf{a}+\mathbf{v})+\mathbf{b}=\mathbf{0}+\mathbf{b}=\mathbf{b}.
\]
Next $0\mathbf{v}=(0+0)\mathbf{v}=0\mathbf{v}+0\mathbf{v}$. Adding $-(0\mathbf{v})$ to both
sides gives $\mathbf{0}=0\mathbf{v}$. From there
$\mathbf{v}+(-1)\mathbf{v}=1\mathbf{v}+(-1)\mathbf{v}=(1-1)\mathbf{v}=0\mathbf{v}=\mathbf{0}$,
so $(-1)\mathbf{v}$ is an inverse of $\mathbf{v}$, and since inverses are unique,
$(-1)\mathbf{v}=-\mathbf{v}$. Finally, if $c\mathbf{v}=\mathbf{0}$ and $c\neq 0$, then
\[
  \mathbf{v}=1\mathbf{v}=(c^{-1}c)\mathbf{v}=c^{-1}(c\mathbf{v})=c^{-1}\mathbf{0}=\mathbf{0}.
\]
\hfill$\square$\par\medskip

The following four examples suffice for this entire book.

\emph{(i) Coordinate space.} $\mathbb{R}^{n}$ with addition and scalar multiplication coordinate
by coordinate. Here $\mathbf{0}=(0,\dots,0)$, and the eight axioms reduce to the field properties
of the real numbers already discussed in Chapter 3.

\emph{(ii) Matrix space.} $M_{2\times3}(\mathbb{R})$, the set of $2\times3$ matrices with real
entries, with operations entry by entry and the zero matrix as $\mathbf{0}$. Nothing changes if
$2\times3$ is replaced by $m\times n$.

\emph{(iii) Polynomial space.} $P_{3}$ is the set of polynomials with real coefficients of degree
at most $3$, together with the zero polynomial, with operations term by term. A trap worth
noting: polynomials of degree \emph{exactly} $3$ do not form a vector space, since $x^{3}$ and
$-x^{3}+x$ both have degree three but their sum $x$ has degree one. The words ``at most'' are not
decoration.

\emph{(iv) Function space.} The set of all functions $\mathbb{R}\to\mathbb{R}$, written
$\mathbb{R}^{\mathbb{R}}$, with $(f+g)(x)=f(x)+g(x)$ and $(af)(x)=af(x)$. The zero function is
$\mathbf{0}$, and the vector space axioms are inherited from $\mathbb{R}$ at each point $x$.

Example (iv) yields a subtler example. The set $C[0,1]$ of continuous functions on $[0,1]$ is a
vector space, because sums and scalar multiples of continuous functions remain continuous. So is
the set of convergent real sequences, using the limit-of-sums property from Chapter 4. Both are
parts of a larger space, and parts of that kind we now define.

Back to $\mathbb{R}^{2}$, where both operations are visible directly.
Figure~\ref{fig:ch06-dua-operasi} uses $\mathbf{u}=(2,1)$ and $\mathbf{v}=(1,2)$: the sum
$\mathbf{u}+\mathbf{v}=(3,3)$ is reached along two orders, and $2\mathbf{u}$ as well as
$-\mathbf{u}$ stay collinear with $\mathbf{u}$.

\begin{figure}[htbp]
\centering
\begin{tikzpicture}[scale=0.85,>=Stealth,font=\small]
  \draw[->,layergrey] (-2.8,0) -- (5.0,0) node[below] {$x$};
  \draw[->,layergrey] (0,-1.8) -- (0,4.2) node[left] {$y$};
  \draw[dashed,layergrey] (2,1) -- (3,3);
  \draw[dashed,layergrey] (1,2) -- (3,3);
  \draw[->,thick,accent] (0,0) -- (2,1) node[below right] {$\mathbf{u}$};
  \draw[->,thick,layerindigo] (0,0) -- (1,2) node[above left] {$\mathbf{v}$};
  \draw[->,thick,layerteal] (0,0) -- (3,3) node[right] {$\mathbf{u}+\mathbf{v}$};
  \draw[->,layerorange] (0,0) -- (4,2) node[right] {$2\mathbf{u}$};
  \draw[->,layerbrown] (0,0) -- (-2,-1) node[left] {$-\mathbf{u}$};
  \fill[ink] (0,0) circle (1.2pt);
  \node[below left] at (0,0) {$\mathbf{0}$};
\end{tikzpicture}
\caption{The two operations in $\mathbb{R}^{2}$. The thick arrows show $\mathbf{u}=(2,1)$ and
$\mathbf{v}=(1,2)$; the dashed lines complete the parallelogram, so that
$\mathbf{u}+\mathbf{v}=(3,3)$ is reached along two identical orders. The thin arrows show
$2\mathbf{u}=(4,2)$ and $-\mathbf{u}=(-2,-1)$, both staying on one line with
$\mathbf{u}$.}
\label{fig:ch06-dua-operasi}
\end{figure}

\kamus{vector space}{ruang vektor}{a mixing box: its contents may be added and scaled, and the result never leaves the box.}
\begin{bacakeras}
Eight axioms, read as a checklist. Two voices, alternating; the second voice always answers in
three words.

First voice: Swap the places. Second voice: The result is the same.
First voice: Move the brackets. Second voice: The result is the same.
First voice: There is a zero vector. Second voice: It adds nothing.
First voice: There is an opposite vector. Second voice: Back to zero.
First voice: A scalar spreads over two terms. Second voice: Left and right agree.
First voice: Two scalars add up. Second voice: One vector carries it.
First voice: Two scalars multiply. Second voice: One multiplication only.
First voice: The scalar one. Second voice: Silent, unchanged.

The last two lines sound almost alike although they differ: one adds proportions, the other
multiplies proportions. That is where the most frequent arithmetic error is born.

Say it once more, slowly, in a single breath: swap the places, move the brackets, there is a
zero, there is an opposite, it spreads, it adds up, it multiplies, it is silent.
\end{bacakeras}


\section{Subspaces and the subspace test}
\tingkat{basic}

A subspace is a part of a vector space that still behaves as a vector space even when we use only
the members of that part.

\begin{definisi}[Subspace]
\index{subspace}
A set $W\subseteq V$ is called a \emph{subspace} of $V$ if $W$ is a vector space with the same
operations as on $V$.
\end{definisi}

That definition demands eight axioms, and checking eight things every time is too expensive. Six
of them are free: commutativity, associativity, distributivity hold in $W$ because they hold in
$V$. Three things remain, and they are the content of the next theorem.

If $W$ is nonempty and closed under both operations, then $W$ is automatically a subspace.

\begin{teorema}[Subspace test]
\index{subspace test}
Let $V$ be a vector space and $W\subseteq V$ with $W\neq\varnothing$. Then $W$ is a subspace of
$V$ if and only if $W$ is closed under addition and scalar multiplication.
\end{teorema}

\par\medskip\noindent\textbf{Proof.}\ 
If $W$ is a subspace, the result of both operations on members of $W$ must be a member of $W$, so
both closures hold. Conversely, suppose $W$ is nonempty and closed under both operations. Take
$\mathbf{w}\in W$. By scalar closure, $\mathbf{0}=0\mathbf{w}\in W$, so $W$ contains the neutral
element. Also $-\mathbf{w}=(-1)\mathbf{w}\in W$, so every member of $W$ has its inverse inside
$W$. The remaining axioms hold because they hold in $V$ and the operations on $W$ are merely the
restrictions of the operations on $V$.
\hfill$\square$\par\medskip

In practice we check three things: $\mathbf{0}\in W$ and the two closures. All three are
necessary; the next three examples show that failure of one condition cannot be covered up by
another.
\begin{resep}
\begin{enumerate}
  \item Write the membership condition of $W$ as an equation, not as a picture or a name.
  \item Substitute $\mathbf{v}=\mathbf{0}$ into that condition. If it fails, stop: not a subspace.
  \item Take arbitrary $\mathbf{u},\mathbf{v}\in W$, check $\mathbf{u}+\mathbf{v}$ using the same
        condition, not using numerical examples.
  \item Take arbitrary $c\in\mathbb{R}$ and $\mathbf{u}\in W$, check $c\mathbf{u}$; include
        $c<0$ and $0<c<1$.
\end{enumerate}
A warning about order: do not swap step 2 with steps 3 and 4. A set that does not contain
$\mathbf{0}$ need not be tested for closure, but stopping after step 2 is also wrong: containing
$\mathbf{0}$ is not enough, as the following examples show.
\end{resep}


\emph{A line not through the origin.} $W=\{(x,y)\in\mathbb{R}^{2}:x+y=1\}$ does not contain
$\mathbf{0}$, so it is not a subspace; it is also closed under neither addition nor scalar
multiplication, since $(1,0)+(0,1)$ and $2(1,0)$ both leave $W$.

\emph{The union of two axes.} $W=\{(x,0)\}\cup\{(0,y)\}$ contains $\mathbf{0}$ and is closed
under scalar multiplication, but $(1,0)+(0,1)=(1,1)$ leaves $W$.

\emph{The first quadrant.} $W=\{(x,y):x\ge 0,\ y\ge 0\}$ is closed under addition, but
$(-1)(1,1)=(-1,-1)$ leaves $W$.

All three examples fail a different condition; Figure~\ref{fig:ch06-garis} shows the simplest
pair, a line through the origin and a line that misses it.

\begin{figure}[htbp]
\centering
\begin{tikzpicture}[scale=0.72,>=Stealth,font=\small]
  \draw[->,layergrey] (-1.6,0) -- (3.2,0) node[below] {$x$};
  \draw[->,layergrey] (0,-1.6) -- (0,3.2) node[left] {$y$};
  \draw[very thick,layerteal] (-1.4,-1.4) -- (2.8,2.8);
  \fill[ink] (0,0) circle (1.3pt) node[below left] {$\mathbf{0}$};
  \fill[accent] (2,2) circle (1.5pt) node[above right] {$(2,2)$};
  \node at (0.8,-2.4) {\emph{a line through the origin}};
  \begin{scope}[xshift=7cm]
    \draw[->,layergrey] (-1.6,0) -- (3.2,0) node[below] {$x$};
    \draw[->,layergrey] (0,-1.6) -- (0,3.2) node[left] {$y$};
    \draw[very thick,accent] (-0.8,1.8) -- (2.4,-1.4);
    \fill[accent] (1,0) circle (1.5pt) node[below right] {$(1,0)$};
    \fill[accent] (0,1) circle (1.5pt) node[above left] {$(0,1)$};
    \draw[ink,thick] (0,0) circle (3pt);
    \node[below left] at (0,0) {$\mathbf{0}$};
    \node at (0.8,-2.4) {\emph{a line without the origin}};
  \end{scope}
\end{tikzpicture}
\caption{Two sets in $\mathbb{R}^{2}$. Left: $\{(x,y):y=x\}$ contains $\mathbf{0}$ and
$(2,2)$, so it is closed under both operations. Right: $\{(x,y):x+y=1\}$ contains $(1,0)$ and
$(0,1)$, but the origin is marked with an empty circle because it is not included; a single
missing point is enough to disqualify a subspace.}
\label{fig:ch06-garis}
\end{figure}



Conversely, the next two classes of examples always pass.

\emph{Null space.} For an $m\times n$ matrix $A$, the set
$\{\mathbf{x}:A\mathbf{x}=\mathbf{0}\}$ is a subspace, since $A\mathbf{0}=\mathbf{0}$,
$A(\mathbf{x}+\mathbf{y})=A\mathbf{x}+A\mathbf{y}=\mathbf{0}$, and
$A(c\mathbf{x})=cA\mathbf{x}=\mathbf{0}$. That set is called the null space or kernel of $A$,
written $\ker A$, and is discussed further in Chapter 7.

\emph{Polynomials vanishing at one point.} The set
$\{p\in P_{3}:p(1)=0\}$ is a subspace of $P_{3}$, because evaluation at the point $1$ is linear:
$(p+q)(1)=p(1)+q(1)$ and $(cp)(1)=c\,p(1)$.

Both examples preserve the solution set of a homogeneous equation. That is no accident.
\begin{tantangan}
The first-quadrant example fails because of one negative scalar. Find a set
$W\subseteq\mathbb{R}^{2}$ that contains $\mathbf{0}$, is closed under addition, is not closed
under scalar multiplication, and is neither the union of two axes, nor a quadrant, nor a set with
integer coordinates.

Hint 1: do not regulate the direction of the arrows; regulate which numbers may appear in the sum
of the two coordinates.

Hint 2: require $x+y$ to be an integer. Then $(1,0)+(0,1)=(1,1)$ stays inside, whereas
$\tfrac{1}{2}(1,0)$ leaves.

If you get stuck, go back to the section ``Subspaces and the subspace test'' and compare with the
line example $x+y=1$; the only difference is the right-hand side of the equation.
\end{tantangan}


\kamus{subspace}{subruang}{a corridor inside a building: we may walk back and forth in that corridor, but we never leave it.}

\begin{untukanak}
Picture two straight corridors in a building that meet at a single point, like the letter X. Each
corridor, if we walk inside it, is still a corridor: we can go forward, backward, and stop
anywhere without leaving it. The same holds if we walk twice as far along the same corridor; we
are still in that corridor.

Now consider the union of the two corridors. Take one step in the first corridor, then one step in
the second. The place where we stop lies in neither corridor. So the union of two corridors is not
a corridor, even though each corridor is a corridor. This is the thing most often forgotten about
the word ``closed'': what is tested is not the members one by one, but the result of mixing them.

The nagging question: if three corridors meet at one point, is their union a corridor? Sometimes
yes, when all three are the same corridor. In this chapter we will learn why ``closed under
mixing'' is far stronger than ``contains the meeting point'', and why a single example of a step
that leaves is enough to disqualify everything.
\end{untukanak}

\section{Linear combinations and span}
\tingkat{intermediate}

A linear combination is the result of mixing several vectors in certain proportions. The span, in
turn, is the set of all those mixtures, for every conceivable proportion.

\begin{definisi}[Linear combination and span]
\index{linear combination}
\index{span}
A vector $\mathbf{v}$ is a \emph{linear combination} of
$\mathbf{v}_{1},\dots,\mathbf{v}_{k}$ if there are scalars $c_{1},\dots,c_{k}$ with
$\mathbf{v}=c_{1}\mathbf{v}_{1}+\dots+c_{k}\mathbf{v}_{k}$. The \emph{span} of a set
$S\subseteq V$, written $\operatorname{span}(S)$, is the set of all finite linear combinations
of members of $S$. For $S=\varnothing$ we stipulate
$\operatorname{span}(S)=\{\mathbf{0}\}$.
\end{definisi}

A span is always a subspace, and it is the smallest subspace containing $S$. The first statement
says the span deserves to be called a space; the second, that no other space is more economical.

\begin{lemma}
For every $S\subseteq V$, $\operatorname{span}(S)$ is a subspace of $V$, and every subspace
containing $S$ also contains $\operatorname{span}(S)$.
\end{lemma}

\par\medskip\noindent\textbf{Proof.}\ 
For $S=\varnothing$ the statement is true because $\{\mathbf{0}\}$ is a subspace. For $S$
nonempty, take $\mathbf{u}=\sum a_{i}\mathbf{s}_{i}$ and $\mathbf{v}=\sum b_{j}\mathbf{t}_{j}$,
two finite linear combinations of members of $S$. Then
$\mathbf{u}+\mathbf{v}=\sum a_{i}\mathbf{s}_{i}+\sum b_{j}\mathbf{t}_{j}$ and
$c\mathbf{u}=\sum (ca_{i})\mathbf{s}_{i}$ are also finite linear combinations of members of $S$,
so the span is closed under both operations; it also contains $\mathbf{0}=0\mathbf{s}$ for any
$\mathbf{s}\in S$. If $W$ is a subspace and $S\subseteq W$, then every $\mathbf{s}\in S$ lies in
$W$, and closure of $W$ forces every finite linear combination of members of $S$ to lie in $W$
as well.
\hfill$\square$\par\medskip

Therefore $\operatorname{span}(S)$ is the intersection of all subspaces containing $S$, and
adding a vector already inside the span changes nothing:
$\operatorname{span}(\operatorname{span}(S))=\operatorname{span}(S)$.

If $\mathbf{b}$ lies in the span of the columns of $A$, then the system $A\mathbf{x}=\mathbf{b}$
has a solution.

\begin{proposisi}
Let $A$ be an $m\times n$ matrix with columns $\mathbf{a}_{1},\dots,\mathbf{a}_{n}$ and let
$\mathbf{b}\in\mathbb{R}^{m}$. The system $A\mathbf{x}=\mathbf{b}$ is consistent if and only if
$\mathbf{b}\in\operatorname{span}(\mathbf{a}_{1},\dots,\mathbf{a}_{n})$.
\end{proposisi}

\par\medskip\noindent\textbf{Proof.}\ 
Writing $A\mathbf{x}=\mathbf{b}$ coordinate by coordinate is the same as writing
$x_{1}\mathbf{a}_{1}+\dots+x_{n}\mathbf{a}_{n}=\mathbf{b}$. So a solution $\mathbf{x}$ exists
exactly when $\mathbf{b}$ is a linear combination of the columns of $A$ with coefficients
$x_{1},\dots,x_{n}$.
\hfill$\square$\par\medskip

The set $\operatorname{span}(\mathbf{a}_{1},\dots,\mathbf{a}_{n})$ is called the \emph{column
space} of $A$, and it is a subspace of $\mathbb{R}^{m}$ by the lemma above. For the opening
table, that column space is built by four columns, and two of them are redundant.

As a small example, in $\mathbb{R}^{3}$ consider $\mathbf{s}_{1}=(1,0,1)$ and
$\mathbf{s}_{2}=(0,1,1)$. Every linear combination has the form $(a,b,a+b)$, so the span is the
plane $z=x+y$. The vector $(1,2,4)$ lies outside that plane because $4\neq 1+2$, whereas
$(1,2,3)$ lies inside it with $a=1$, $b=2$.

\begin{tangan}
Materials: two pencils of equal length and one flat table. Mark a point on the table with a
fingertip; call it the origin.

Place the first pencil with its tail at the origin. Place the second pencil with its tail at the
tip of the first, at a right angle to it. Now slide the second pencil without rotating it until
its tail is back at the origin. Its tip moves to another place, and that place is the sum of the
two pencils: two orders of addition stop at the same point.

Now the scalars. Reverse one pencil, and its direction becomes negative. Join two identical
pencils from the origin, and the length doubles. What the reader must see: all those tips stay on
one straight line, however far the pencils are joined. One pencil covers only one line on the
table; two nonparallel pencils cover the entire table surface. That is the difference between a
line and a plane, and no third pencil adds a new surface.
\end{tangan}



\section{Linear independence and dependence}
\tingkat{intermediate}

A set of vectors is called linearly independent if not one of its members can be mixed from the
other members. It is linearly dependent if some member is redundant.

\begin{definisi}[Linear independence and dependence]
\index{linear independence}
A set $\{\mathbf{v}_{1},\dots,\mathbf{v}_{k}\}\subseteq V$ is \emph{linearly independent} if the
only choice of scalars satisfying
$c_{1}\mathbf{v}_{1}+\dots+c_{k}\mathbf{v}_{k}=\mathbf{0}$ is
$c_{1}=\dots=c_{k}=0$. If some other choice exists, the set is \emph{linearly dependent}.
\end{definisi}

Two small conventions: the empty set is linearly independent because there are no coefficients to
check, and $\{\mathbf{0}\}$ is linearly dependent because $1\cdot\mathbf{0}=\mathbf{0}$ is a
nontrivial relation.

In a linearly dependent set there is always one vector that lies in the span of the preceding
vectors.

\begin{lemma}[Dependence]
Let $\mathbf{v}_{1}\neq\mathbf{0}$ and $S=\{\mathbf{v}_{1},\dots,\mathbf{v}_{k}\}$.
The set $S$ is linearly dependent if and only if there is an index $j\ge 2$ with
$\mathbf{v}_{j}\in\operatorname{span}(\mathbf{v}_{1},\dots,\mathbf{v}_{j-1})$.
\end{lemma}

\par\medskip\noindent\textbf{Proof.}\ 
If $\mathbf{v}_{j}=\sum_{i<j}c_{i}\mathbf{v}_{i}$, then
$\sum_{i<j}c_{i}\mathbf{v}_{i}-\mathbf{v}_{j}=\mathbf{0}$ is a nontrivial relation, because the
coefficient of $\mathbf{v}_{j}$ is $-1$; so $S$ is dependent. Conversely, if $S$ is dependent then
there is a nontrivial relation $\sum_{i=1}^{k}c_{i}\mathbf{v}_{i}=\mathbf{0}$. Choose the largest
$j$ with $c_{j}\neq 0$. Since $\mathbf{v}_{1}\neq\mathbf{0}$, the index $j$ is not $1$, so
$j\ge 2$ and
\[
  \mathbf{v}_{j}=-\frac{1}{c_{j}}\sum_{i<j}c_{i}\mathbf{v}_{i}
  \in\operatorname{span}(\mathbf{v}_{1},\dots,\mathbf{v}_{j-1}).
\]
\hfill$\square$\par\medskip

Three consequences we use constantly: a set containing $\mathbf{0}$ is always dependent; two
vectors are dependent exactly when one is a multiple of the other; and every subset of a linearly
independent set is also linearly independent.

As an example, consider $\mathbf{v}_{1}=(1,2,-1)$, $\mathbf{v}_{2}=(0,1,3)$, and
$\mathbf{v}_{3}=(2,5,1)$ in $\mathbb{R}^{3}$. Since
$2\mathbf{v}_{1}+\mathbf{v}_{2}=(2,5,1)=\mathbf{v}_{3}$, the relation
$2\mathbf{v}_{1}+\mathbf{v}_{2}-\mathbf{v}_{3}=\mathbf{0}$ is nontrivial and the three are
dependent. Not one of those vectors is zero, and no two are multiples; the dependence appears only
at three vectors.

Figure~\ref{fig:ch06-bidang} places the two side by side: the span of two vectors as a plane, and
three vectors one of which already lies in the span of the other two.

\begin{figure}[htbp]
\centering
\begin{tikzpicture}[scale=0.8,>=Stealth,font=\small]
  \begin{scope}[x={(1cm,0cm)}, y={(0.40cm,0.28cm)}, z={(0cm,1cm)}]
    \fill[layerteal!10] (0,0,0) -- (2.4,0,2.4) -- (2.4,1.8,4.2) -- (0,1.8,2.4) -- cycle;
    \draw[layerteal] (0,0,0) -- (2.4,0,2.4) -- (2.4,1.8,4.2) -- (0,1.8,2.4) -- cycle;
    \draw[->,layergrey] (0,0,0) -- (3.0,0,0) node[below] {$x$};
    \draw[->,layergrey] (0,0,0) -- (0,2.2,0) node[right] {$y$};
    \draw[->,layergrey] (0,0,0) -- (0,0,2.6) node[left] {$z$};
    \draw[->,thick,accent] (0,0,0) -- (1,0,1)
      node[above left,font=\scriptsize,inner sep=1pt] {$\mathbf{s}_{1}$};
    \draw[->,thick,layerindigo] (0,0,0) -- (0,1,1)
      node[right,font=\scriptsize,inner sep=1pt] {$\mathbf{s}_{2}$};
    \fill[layerjade] (1,2,3) circle (1.6pt);
    \node[right,font=\scriptsize] at (1,2,3) {$(1,2,3)$};
    \draw[ink] (1,2,4) circle (2pt);
    \node[right,font=\scriptsize] at (1,2,4) {$(1,2,4)$};
  \end{scope}
  \begin{scope}[xshift=6.4cm]
    \draw[->,layergrey] (-0.4,0) -- (3.2,0) node[below] {$x$};
    \draw[->,layergrey] (0,-0.4) -- (0,2.6) node[left] {$y$};
    \draw[dashed,layergrey] (0,0) -- (2,1);
    \draw[dashed,layergrey] (2,1) -- (2.4,2);
    \draw[->,thick,accent] (0,0) -- (1,0.5) node[below right] {$\mathbf{v}_{1}$};
    \draw[->,thick,layerindigo] (0,0) -- (0.4,1) node[above left] {$\mathbf{v}_{2}$};
    \draw[->,very thick,layerorange] (0,0) -- (2.4,2) node[right] {$\mathbf{v}_{3}$};
    \fill[ink] (0,0) circle (1.2pt);
  \end{scope}
\end{tikzpicture}
\caption{Left: $\operatorname{span}\{(1,0,1),(0,1,1)\}$ is the plane $z=x+y$; the point
$(1,2,3)$ lies inside it, the point $(1,2,4)$ outside it. Right: a two-dimensional sketch of
$\mathbf{v}_{1}=(1,2,-1)$ and $\mathbf{v}_{2}=(0,1,3)$ with
$\mathbf{v}_{3}=2\mathbf{v}_{1}+\mathbf{v}_{2}=(2,5,1)$; the dashed lines show
$2\mathbf{v}_{1}$ plus $\mathbf{v}_{2}$ landing exactly at the tip of $\mathbf{v}_{3}$, so that
the three are dependent.}
\label{fig:ch06-bidang}
\end{figure}



In $\mathbb{R}^{2}$ any three vectors are always dependent, without computing anything.

\begin{proposisi}
Every set of three vectors in $\mathbb{R}^{2}$ is linearly dependent.
\end{proposisi}

\par\medskip\noindent\textbf{Proof.}\ 
Let $S=\{\mathbf{u},\mathbf{v},\mathbf{w}\}\subseteq\mathbb{R}^{2}$. If two of them are multiples,
$S$ is dependent. If not, $\mathbf{u}$ and $\mathbf{v}$ are linearly independent. Take
$\mathbf{w}=(w_{1},w_{2})$ and look for $a,b$ with
$a\mathbf{u}+b\mathbf{v}=\mathbf{w}$. Writing out the components gives two linear equations in
$a$ and $b$, and since $\mathbf{u}$, $\mathbf{v}$ are not multiples, that system has a unique
solution. Hence $\mathbf{w}\in\operatorname{span}(\mathbf{u},\mathbf{v})$.
\hfill$\square$\par\medskip

In $\mathbb{R}^{3}$ a similar statement holds for four vectors, but the method above would become
a long case check; we need a uniform argument.

To test independence computationally, arrange the vectors as the columns of a matrix and then
solve $A\mathbf{c}=\mathbf{0}$; the set is linearly independent exactly when that system has only
the trivial solution. The determinant offers a shortcut in $\mathbb{R}^{3}$, but that tool is only
built in Chapter 8.

\section{Basis and dimension}
\tingkat{advanced}

A basis is a list of vectors that is enough to build the whole space and contains not one
redundant vector. The words ``enough'' and ``no redundancy'' are exactly the two conditions we
already have.

\begin{definisi}[Basis and dimension]
\index{basis}
A set $B\subseteq V$ is called a \emph{basis} of $V$ if $B$ is linearly independent and
$\operatorname{span}(B)=V$. If $V$ has a finite basis, $V$ is called
\emph{finite-dimensional}; if not, $V$ is infinite-dimensional.
\end{definisi}

The definition of dimension cannot yet be written, because it depends on a fact not yet proved:
two bases of the same space always have the same size.

First step: a representation in a basis is always unique.

\begin{lemma}[Uniqueness of coordinates]
If $B=\{\mathbf{b}_{1},\dots,\mathbf{b}_{n}\}$ is a basis of $V$, then every $\mathbf{v}\in V$
can be written uniquely as $\mathbf{v}=a_{1}\mathbf{b}_{1}+\dots+a_{n}\mathbf{b}_{n}$.
\end{lemma}

\par\medskip\noindent\textbf{Proof.}\ 
Existence of the representation comes from $\operatorname{span}(B)=V$. Suppose there are two
representations with coefficients $a_{i}$ and $a_{i}'$. Subtracting them gives
$\sum_{i}(a_{i}-a_{i}')\mathbf{b}_{i}=\mathbf{0}$. Because $B$ is linearly independent, all those
difference coefficients are zero, so $a_{i}=a_{i}'$ for every $i$.
\hfill$\square$\par\medskip

The numbers $a_{1},\dots,a_{n}$ are called the coordinates of $\mathbf{v}$ with respect to $B$;
change of basis is the subject of Chapter 7, and here we only need uniqueness.

Second step: one vector can be exchanged for another without destroying the spanning property.
Replacing one vector in a spanning set by a combination that contains it does not reduce what can
be built.

\begin{lemma}[Exchange]
\index{exchange lemma}
Let $\mathbf{u}=\sum_{\mathbf{s}\in S}a_{\mathbf{s}}\mathbf{s}$ be a finite sum with
$a_{\mathbf{s}_{0}}\neq 0$ for some $\mathbf{s}_{0}\in S$. If $S\cup\{\mathbf{u}\}$
spans $V$, then $(S\setminus\{\mathbf{s}_{0}\})\cup\{\mathbf{u}\}$ also spans $V$.
\end{lemma}

\par\medskip\noindent\textbf{Proof.}\ 
From that relation,
$\mathbf{s}_{0}=a_{\mathbf{s}_{0}}^{-1}\bigl(\mathbf{u}-\sum_{\mathbf{s}\neq\mathbf{s}_{0}}a_{\mathbf{s}}\mathbf{s}\bigr)$,
so $\mathbf{s}_{0}$ is a linear combination of members of the smaller set
$T=(S\setminus\{\mathbf{s}_{0}\})\cup\{\mathbf{u}\}$. Every member of $S\cup\{\mathbf{u}\}$
therefore lies in $\operatorname{span}(T)$, so
$\operatorname{span}(S\cup\{\mathbf{u}\})\subseteq\operatorname{span}(T)$. The reverse inclusion
is clear because $T\subseteq S\cup\{\mathbf{u}\}$. Hence both spans equal $V$.
\hfill$\square$\par\medskip

The third step is the heart of this chapter.

If a space can be built by $n$ vectors, then it is impossible to find $n+1$ linearly independent
vectors inside it.

\begin{teorema}[Exchange]
\index{exchange theorem}
If $V$ is spanned by $n$ vectors and $\mathbf{w}_{1},\dots,\mathbf{w}_{m}$ are linearly
independent in $V$, then $m\le n$.
\end{teorema}

\par\medskip\noindent\textbf{Proof.}\ 
Let $S_{0}=\{\mathbf{v}_{1},\dots,\mathbf{v}_{n}\}$ span $V$. For $j=1$:
$\mathbf{w}_{1}\neq\mathbf{0}$ is a linear combination of members of $S_{0}$ with at least one
nonzero coefficient, so the exchange lemma gives $T_{1}\subseteq S_{0}$ of size $n-1$ with
$\{\mathbf{w}_{1}\}\cup T_{1}$ spanning.

Suppose $\{\mathbf{w}_{1},\dots,\mathbf{w}_{j}\}\cup T_{j}$ spans with $|T_{j}|=n-j$. The vector
$\mathbf{w}_{j+1}$ is a linear combination of that set. If all coefficients on $T_{j}$ are zero,
then $\mathbf{w}_{j+1}$ becomes a linear combination of $\mathbf{w}_{1},\dots,\mathbf{w}_{j}$,
which is impossible because that set is linearly independent. So some member of $T_{j}$ has a
nonzero coefficient, and the exchange lemma gives $T_{j+1}$ of size $n-j-1$ that still spans
together with $\mathbf{w}_{1},\dots,\mathbf{w}_{j+1}$.

After $m$ steps we obtain a spanning set of size $n-m$; the size of a set cannot be negative, so
$m\le n$.
\hfill$\square$\par\medskip

Now the size of a space can be written: the number of vectors in its most economical list, and
that number does not depend on the list chosen.

Two bases of a finite-dimensional space always have the same number of elements.

\begin{akibat}
If $V$ is finite-dimensional with a basis $B_{1}$ of size $n$ and a basis $B_{2}$ of size $m$,
then $n=m$.
\end{akibat}

\par\medskip\noindent\textbf{Proof.}\ 
$B_{1}$ spans and $B_{2}$ is linearly independent, so the exchange theorem gives $m\le n$.
Conversely $B_{2}$ spans and $B_{1}$ is linearly independent, hence $n\le m$.
\hfill$\square$\par\medskip

The number of vectors in that most economical list is important enough to deserve a name of its
own.

\begin{definisi}[Dimension]
\index{dimension}
For finite-dimensional $V$, $\dim V$ is the number of members of any basis of $V$. The null space
$\{\mathbf{0}\}$ is stipulated to have dimension $0$.
\end{definisi}

The size of a subspace never exceeds the size of its space; a list whose length matches the
dimension is immediately a basis; and a linearly independent list can always be completed to full
length.

\begin{akibat}
Let $\dim V=n$. (a) Every linearly independent set of $n$ vectors is a basis.
(b) Every spanning set of $n$ vectors is a basis. (c) Every linearly independent set
can be extended to a basis. (d) If $U$ is a subspace of $V$, then $\dim U\le\dim V$, and
equality holds only when $U=V$.
\end{akibat}

\par\medskip\noindent\textbf{Proof.}\ 
For (a), a linearly independent set of size $n$ that does not span can be enlarged by one vector
outside its span, and the result is still linearly independent; so there are $n+1$ linearly
independent vectors, and the exchange theorem forbids that. For (b), a spanning set of size $n$
that is dependent can lose one vector through the dependence lemma and still span with $n-1$
vectors; again this contradicts exchange. For (c), add a vector outside the span as long as one
exists; the process stops after at most $n$ steps, and when it stops the span is already all of
$V$. For (d), a basis of $U$ is a linearly independent set in $V$, so its size is at most $n$; if
its size is $n$, part (a) makes it a basis of $V$, hence $U=V$.
\hfill$\square$\par\medskip

Thus $\dim\mathbb{R}^{n}=n$, $\dim P_{n}=n+1$ with basis $1,x,\dots,x^{n}$,
$\dim M_{m\times n}(\mathbb{R})=mn$ with basis the matrices having exactly one entry equal to $1$,
and $\dim\{\mathbf{0}\}=0$. For function spaces such as $\mathbb{R}^{\mathbb{R}}$ there is no
finite basis; the proof requires the axiom of choice and lies outside the scope of this book.

\kamus{basis}{basis}{the most economical list of paint cans: enough to make every color, and not one can in it can be thrown away.}

\kamus{dimension}{dimensi}{the number of cans in that most economical list; the number belongs to the space, not to the list we happen to choose.}

\section{Direct sums and complementary subspaces}
\tingkat{advanced}

The sum of two subspaces is the set of all results of adding one member of each subspace.

\begin{definisi}[Sum of subspaces]
\index{sum of subspaces}
For subspaces $U,W\subseteq V$, the \emph{sum} $U+W$ is
$\{\mathbf{u}+\mathbf{w}:\mathbf{u}\in U,\ \mathbf{w}\in W\}$.
\end{definisi}

It differs from the union, which is almost never a subspace, as the two-axes example shows.

The sum is itself a subspace, and no smaller subspace contains both.

\begin{proposisi}
For subspaces $U,W\subseteq V$, the set $U+W$ is a subspace of $V$, and it is contained in every
subspace that contains $U\cup W$.
\end{proposisi}

\par\medskip\noindent\textbf{Proof.}\ 
The set $U+W$ contains $\mathbf{0}=\mathbf{0}+\mathbf{0}$. If
$\mathbf{u}_{1}+\mathbf{w}_{1}$ and $\mathbf{u}_{2}+\mathbf{w}_{2}$ are two of its members, their
sum $(\mathbf{u}_{1}+\mathbf{u}_{2})+(\mathbf{w}_{1}+\mathbf{w}_{2})$ has its first term in $U$
and its second in $W$, and $c(\mathbf{u}_{1}+\mathbf{w}_{1})=c\mathbf{u}_{1}+c\mathbf{w}_{1}$ is
of the same form; so $U+W$ is a subspace. If $Z$ is a subspace containing $U$ and $W$, then $Z$
contains every $\mathbf{u}+\mathbf{w}$ because it is closed under addition, so $U+W\subseteq Z$.
\hfill$\square$\par\medskip

The size of a sum is not the sum of the sizes, because the part that overlaps is counted twice.

$\dim(U+W)$ equals $\dim U$ plus $\dim W$ minus the dimension of the intersection.

\begin{teorema}[Dimension of a sum]
\index{dimension of a sum}
If $U$ and $W$ are subspaces of a finite-dimensional space $V$, then
\[
  \dim(U+W)=\dim U+\dim W-\dim(U\cap W).
\]
\end{teorema}

\par\medskip\noindent\textbf{Proof.}\ 
The intersection $U\cap W$ is a subspace, so it is finite-dimensional. Take a basis
$\{\mathbf{z}_{1},\dots,\mathbf{z}_{m}\}$ of $U\cap W$. Since it is linearly independent in $U$,
that set is extended to a basis of $U$ by adding $\mathbf{u}_{1},\dots,
\mathbf{u}_{j}$, so that $\dim U=m+j$; in the same way, adding
$\mathbf{w}_{1},\dots,\mathbf{w}_{k}$ yields a basis of $W$, so that $\dim W=m+k$.

Claim: $\{\mathbf{z}_{1},\dots,\mathbf{z}_{m},\mathbf{u}_{1},\dots,\mathbf{u}_{j},
\mathbf{w}_{1},\dots,\mathbf{w}_{k}\}$ is a basis of $U+W$. The set contains a basis of $U$ and a
basis of $W$, so its span contains $U+W$; the reverse inclusion is clear. For independence,
suppose
\[
  \sum_{i}a_{i}\mathbf{z}_{i}+\sum_{r}b_{r}\mathbf{u}_{r}
  +\sum_{s}c_{s}\mathbf{w}_{s}=\mathbf{0}.
\]
Then $\sum_{s}c_{s}\mathbf{w}_{s}=-\sum_{i}a_{i}\mathbf{z}_{i}-\sum_{r}b_{r}\mathbf{u}_{r}$
lies in $U$ and at the same time in $W$; so it is a member of $U\cap W$ and can be written
$\sum_{i}d_{i}\mathbf{z}_{i}$. Then $\sum_{s}c_{s}\mathbf{w}_{s}-\sum_{i}d_{i}\mathbf{z}_{i}
=\mathbf{0}$ is a linear relation in the basis of $W$, so all $c_{s}$ and $d_{i}$ are zero; the
original equation shrinks to a linear relation in the basis of $U$, so that $a_{i}$ and $b_{r}$
are also zero.

Hence $\dim(U+W)=m+j+k=(m+j)+(m+k)-m=\dim U+\dim W-\dim(U\cap W)$.
\hfill$\square$\par\medskip

Now the second question from the opening of the chapter: when is the decomposition into two terms
unique?

The sum $U+W$ is called a \emph{direct sum}, written $U\oplus W$, if every member of
$U+W$ can be written as $\mathbf{u}+\mathbf{w}$ in exactly one way.

\begin{definisi}[Direct sum]
\index{direct sum}
If every $\mathbf{x}\in U+W$ has exactly one representation
$\mathbf{x}=\mathbf{u}+\mathbf{w}$ with $\mathbf{u}\in U$ and $\mathbf{w}\in W$, then
$U+W$ is called a direct sum and is written $U\oplus W$.
\end{definisi}

Uniqueness of all representations can be checked by looking only at the representation of zero.

$U+W$ is a direct sum exactly when $U\cap W=\{\mathbf{0}\}$.

\begin{teorema}[Criterion for a direct sum]
Let $U$ and $W$ be subspaces of $V$. Then $U+W=U\oplus W$ if and only if
$U\cap W=\{\mathbf{0}\}$.
\end{teorema}

\par\medskip\noindent\textbf{Proof.}\ 
Suppose $U+W$ is a direct sum and take $\mathbf{x}\in U\cap W$. Then
$\mathbf{0}=\mathbf{x}+(-\mathbf{x})$ with $\mathbf{x}\in U$ and $-\mathbf{x}\in W$, and also
$\mathbf{0}=\mathbf{0}+\mathbf{0}$. Uniqueness gives $\mathbf{x}=\mathbf{0}$.

Conversely, suppose $U\cap W=\{\mathbf{0}\}$. Take $\mathbf{x}\in U+W$ with two representations
$\mathbf{x}=\mathbf{u}_{1}+\mathbf{w}_{1}=\mathbf{u}_{2}+\mathbf{w}_{2}$. Then
$\mathbf{u}_{1}-\mathbf{u}_{2}=\mathbf{w}_{2}-\mathbf{w}_{1}$ lies in $U$ and in $W$, so it equals
$\mathbf{0}$. Consequently $\mathbf{u}_{1}=\mathbf{u}_{2}$ and $\mathbf{w}_{1}=\mathbf{w}_{2}$.
\hfill$\square$\par\medskip

For two subspaces that criterion is enough. For three subspaces or more, checking pairwise
intersections is not enough, and the next example shows why.

In $\mathbb{R}^{3}$ take $V_{1}=\operatorname{span}\{(1,0,0)\}$,
$V_{2}=\operatorname{span}\{(0,1,0)\}$, and $V_{3}=\operatorname{span}\{(1,1,0)\}$. Every
pairwise intersection contains only $\mathbf{0}$, but $V_{1}+V_{2}+V_{3}$ is not a direct sum:
$(1,1,0)$ can be written $(1,0,0)+(0,1,0)+\mathbf{0}$ and also $\mathbf{0}+\mathbf{0}+(1,1,0)$.
Even zero has two representations, namely $\mathbf{0}+\mathbf{0}+\mathbf{0}$ and
$(1,0,0)+(0,1,0)+(-1,-1,0)$.

For many terms, the criterion returns to the definition:
$\mathbf{0}=\mathbf{v}_{1}+\dots+\mathbf{v}_{k}$ with $\mathbf{v}_{i}\in V_{i}$ is possible only
when all the $\mathbf{v}_{i}$ are zero.

Finally, every subspace can be paired with another subspace so that the two decompose the whole
space.

\begin{teorema}[Complement]
\index{complement}
If $V$ is finite-dimensional and $U$ is a subspace of $V$, then there is a subspace
$W\subseteq V$ with $V=U\oplus W$. In that case $\dim V=\dim U+\dim W$.
\end{teorema}

\par\medskip\noindent\textbf{Proof.}\ 
If $U=\{\mathbf{0}\}$, take $W=V$. If not, take a basis
$\{\mathbf{u}_{1},\dots,\mathbf{u}_{m}\}$ of $U$ and extend it to a basis
$\{\mathbf{u}_{1},\dots,\mathbf{u}_{m},\mathbf{w}_{1},\dots,\mathbf{w}_{k}\}$ of $V$;
this extension is possible by consequence (c) in the previous section. Set
$W=\operatorname{span}\{\mathbf{w}_{1},\dots,\mathbf{w}_{k}\}$.

Every $\mathbf{x}\in V$ is a linear combination of the entire basis, and separating the
$\mathbf{u}$ terms from the $\mathbf{w}$ terms exhibits it as a member of $U+W$. For the
intersection, take $\mathbf{x}\in U\cap W$. Then
$\mathbf{x}=\sum a_{i}\mathbf{u}_{i}=\sum b_{s}\mathbf{w}_{s}$, so that
$\sum a_{i}\mathbf{u}_{i}-\sum b_{s}\mathbf{w}_{s}=\mathbf{0}$. This is a linear relation in the
basis of $V$, so all its coefficients are zero, in particular $\mathbf{x}=\mathbf{0}$. By the
criterion above, $V=U\oplus W$, and from the dimension-of-a-sum theorem,
$\dim V=\dim U+\dim W$.

The complement is not unique. In $\mathbb{R}^{2}$ with $U$ the $x$-axis, every line through the
origin other than the $x$-axis is a complement of $U$; what is unique is not the choice but the
dimension. In Chapter 10, orthogonal complements make that decomposition the shortest one as well.

Now the opening question can be answered. For $\mathbf{w}=(1,1,1)$, let
$U=\{\boldsymbol{\delta}:\mathbf{w}\cdot\boldsymbol{\delta}=0\}$ and
$W=\operatorname{span}\{\mathbf{w}\}$. Their intersection is only $\{\mathbf{0}\}$, since
$\mathbf{w}\cdot(t\mathbf{w})=3t$ is zero only when $t=0$; with $\dim U=2$ and
$\dim W=1$ we obtain $U\oplus W=\mathbb{R}^{3}$. Every perturbation splits uniquely into a part
that does not change the score, namely a member of $U$, and a part that changes it, a member of
$W$.

\kamus{direct sum}{jumlah langsung}{two corridors that meet only at the origin, so that every place in the building is reachable by one route that cannot be disputed.}

\section{Worked examples}
\tingkat{intermediate}

The following four computations use the entire content of the chapter; all of them are rechecked
with Python in the lab section.

\subsection*{(a) Rank and a basis of the column space}

Back to the opening table. Row elimination: $R_{2}\leftarrow R_{2}-2R_{1}$ turns the second row
into $(0,0,1,1)$, and $R_{3}\leftarrow R_{3}-3R_{1}$ turns the third row into
$(0,0,2,2)$. Then $R_{3}\leftarrow R_{3}-2R_{2}$ eliminates the third row, so that what remains
is $\mathbf{r}_{1}=(1,2,0,1)$ and $(0,0,1,1)$.

That form is already reduced echelon, with pivot columns $1$ and $3$. Hence
$\operatorname{rank}A=2$ and its column space has dimension $2$ with basis
$\{\mathbf{c}_{1},\mathbf{c}_{3}\}=\{(1,2,3),(0,1,2)\}$.

A direct check against the discarded columns:
\[
  \mathbf{c}_{2}=(2,4,6)=2\mathbf{c}_{1},
  \qquad
  \mathbf{c}_{4}=(1,3,5)=\mathbf{c}_{1}+\mathbf{c}_{3}.
\]
Both are exact, with no rounding. This is the answer ``two'' to the opening question: the pivot
columns give a linearly independent set that spans the column space, and a column without a pivot
is a combination of the pivot columns to its left.

The null space of $A$ is also readable from the reduced form. From $x_{1}+2x_{2}+x_{4}=0$ and
$x_{3}+x_{4}=0$ we obtain $x_{1}=-2x_{2}-x_{4}$ and $x_{3}=-x_{4}$, so that
\[
  \mathbf{x}=x_{2}\,(-2,1,0,0)+x_{4}\,(-1,0,-1,1),
\]
written as row vectors for brevity. Hence $\dim\ker A=2$, and
$\operatorname{rank}A+\dim\ker A=2+2=4$ equals the number of columns; the general statement is
rank-nullity, proved in Chapter 7.

\subsection*{(b) Three dependent vectors}

Consider $\mathbf{v}_{1}=(1,2,-1)$, $\mathbf{v}_{2}=(0,1,3)$, $\mathbf{v}_{3}=(2,5,1)$ as
the columns of $V$. Cofactor expansion along the first row:
\[
  \det V=1(1\cdot1-3\cdot5)+0+2\bigl(2\cdot3-1\cdot(-1)\bigr)=(1-15)+2(6+1)=0.
\]
A zero determinant means the columns are dependent; the relation is guessed and then checked:
$2\mathbf{v}_{1}+\mathbf{v}_{2}=(2,5,1)=\mathbf{v}_{3}$, so
$2\mathbf{v}_{1}+\mathbf{v}_{2}-\mathbf{v}_{3}=\mathbf{0}$ and $\operatorname{rank}V=2$.

\subsection*{(c) Dimension of a polynomial space}

$\dim P_{3}=4$, because $\{1,x,x^{2},x^{3}\}$ is a basis. The set
$B=\{1,\ 1+x,\ 1+x+x^{2},\ 1+x+x^{2}+x^{3}\}$ is also a basis: its coordinate matrix with respect
to the standard basis is upper triangular with diagonal $1$, so $B$ is linearly independent, and
since its size is $4=\dim P_{3}$, consequence (a) makes it a basis.

The coordinates of $p=2-x+3x^{3}$ with respect to $B$ are computed from the highest degree
downward: $a_{4}=3$, then $a_{3}+a_{4}=0$ gives $a_{3}=-3$, $a_{2}+a_{3}+a_{4}=-1$ gives
$a_{2}=-1$, and the constant term gives $a_{1}=3$. Hence
$p=3-(1+x)-3(1+x+x^{2})+3(1+x+x^{2}+x^{3})$, which upon expansion gives $2-x+3x^{3}$.

\subsection*{(d) Decomposition as a direct sum}

Take $U=\{(x,y,z):x+y+z=0\}$ and $W=\operatorname{span}\{(1,1,1)\}$ in $\mathbb{R}^{3}$.
A basis of $U$ is $\{(1,-1,0),(1,0,-1)\}$, so $\dim U=2$; with $\dim W=1$ and
$U\cap W=\{\mathbf{0}\}$ we obtain $\mathbb{R}^{3}=U\oplus W$.

Decompose $\mathbf{v}=(3,0,1)$ as $\mathbf{u}+t(1,1,1)$. The condition that the sum of the
coordinates of $\mathbf{u}$ be zero gives $3+0+1-3t=0$, so $t=4/3$. Then
\[
  \mathbf{v}=\underbrace{\Bigl(\tfrac53,-\tfrac43,-\tfrac13\Bigr)}_{\in U}
  +\underbrace{\Bigl(\tfrac43,\tfrac43,\tfrac43\Bigr)}_{\in W},
\]
and $\tfrac53-\tfrac43-\tfrac13=0$ confirms that the first term really lies in $U$. By the
criterion theorem, there is no other decomposition.

\kertas{Redo (a) without looking: carry out the elimination on $A$, mark the pivot columns, then
write $\mathbf{c}_{2}$ and $\mathbf{c}_{4}$ as combinations of $\mathbf{c}_{1}$ and
$\mathbf{c}_{3}$. Then compute $\dim\ker A$ from the reduced form and check that its sum with the
rank equals the number of columns. Finally, for the matrix $V$ in (b), replace entry $(3,1)$ by
$0$ and decide whether the columns are still dependent.}

\section{Computational lab}
\tingkat{advanced}

The following code runs as is with numpy. The function \texttt{rref} is written by hand so that the
pivot columns are visible; its last line cleans up leftover rounding and negative zeros.

\begin{lab}
\begin{verbatim}
import numpy as np

np.set_printoptions(precision=3, suppress=True)

def rref(M):
    M = M.astype(float).copy()
    r = 0
    for c in range(M.shape[1]):
        i = np.argmax(np.abs(M[r:, c])) + r
        if abs(M[i, c]) < 1e-12:
            continue
        M[[r, i]] = M[[i, r]]
        M[r] = M[r] / M[r, c]
        for j in range(M.shape[0]):
            if j != r:
                M[j] = M[j] - M[j, c] * M[r]
        r = r + 1
        if r == M.shape[0]:
            break
    M[np.abs(M) < 1e-12] = 0.0
    return M

A = np.array([[1, 2, 0, 1],
              [2, 4, 1, 3],
              [3, 6, 2, 5]], float)
print("RREF(A):"); print(rref(A))
print("rank(A) =", np.linalg.matrix_rank(A))
print("k2-2k1 =", A[:, 1] - 2 * A[:, 0])
print("k4-k1-k3 =", A[:, 3] - A[:, 0] - A[:, 2])

V = np.array([[1, 0, 2],
              [2, 1, 5],
              [-1, 3, 1]], float)
print("det(V) =", np.linalg.det(V))
print("2v1+v2-v3 =", 2 * V[:, 0] + V[:, 1] - V[:, 2])

x = np.array([0., 1, 2, 3])
M = np.vander(x, 4, increasing=True)
print("det(M) =", np.linalg.det(M))
print("rank(M) =", np.linalg.matrix_rank(M))

B = np.array([[1, 1, 1, 1],
              [0, 1, 1, 1],
              [0, 0, 1, 1],
              [0, 0, 0, 1]], float)
p = np.array([2, -1, 0, 3], float)
k = np.linalg.solve(B, p)
print("koordinat p =", k)
print("sisa =", B @ k - p)
\end{verbatim}
\end{lab}

The printed output:
\begin{verbatim}
RREF(A):
[[1. 2. 0. 1.]
 [0. 0. 1. 1.]
 [0. 0. 0. 0.]]
rank(A) = 2
k2-2k1 = [0. 0. 0.]
k4-k1-k3 = [0. 0. 0.]
det(V) = 0.0
2v1+v2-v3 = [0. 0. 0.]
det(M) = 12.0
rank(M) = 4
koordinat p = [ 3. -1. -3.  3.]
sisa = [0. 0. 0. 0.]
\end{verbatim}

Four things are readable from that output. The pivot columns $1$ and $3$ as well as
$\operatorname{rank}A=2$ agree with the hand computation. The differences
$\mathbf{c}_{2}-2\mathbf{c}_{1}$ and $\mathbf{c}_{4}-\mathbf{c}_{1}-\mathbf{c}_{3}$ are exactly
zero; on real data those differences are small but nonzero, and that is where the limits of
applying rank numerically lie. $\det V=0$ confirms the dependence in (b), and the Vandermonde
matrix of the points $0,1,2,3$ has determinant $12$ and rank $4$. The coordinates of $p$ are
printed as $(3,-1,-3,3)$, matching the hand computation.

\begin{untukdosen}
Start from the opening table, not from the axioms: ask the students to guess how many columns are
independent, then show that the guess ``four'' is wrong. The axioms come in after that, as an
answer to a need. The subspace test and span in the second meeting, linear independence and basis
in the third, direct sums in the fourth.

Two typical misconceptions. First, testing a subspace only by ``does it contain zero'' and then
stopping; the first quadrant contains zero and is closed under addition, but fails on negative
multiples. Second, taking basis to mean standard basis, so that $(5,7)$ and $(4,3)$ are judged not
to be a basis of $\mathbb{R}^{2}$; stress that being a basis is a property of the set.

Class activity: each group writes three vectors in $\mathbb{R}^{3}$ and decides independence or
dependence by elimination; another group changes one entry and asks for a new decision. Students
are ready to move on when they can explain why three vectors in $\mathbb{R}^{2}$ are always
dependent.
\end{untukdosen}

\section{Summary}
\tingkat{basic}

\begin{ringkasan}
\begin{itemize}
  \item Vector space: a set with addition and scalar multiplication satisfying the eight
        axioms; its members may be vectors, matrices, polynomials, or functions.
  \item A subspace is checked with just three things: it contains $\mathbf{0}$ and is closed
        under both operations.
  \item $\operatorname{span}(S)$ is the smallest subspace containing $S$; $\mathbf{b}$ lies in
        the column space of $A$ exactly when $A\mathbf{x}=\mathbf{b}$ is consistent.
  \item Linear dependence: one member lies in the span of the preceding members; three vectors
        in $\mathbb{R}^{2}$ are always dependent.
  \item The exchange theorem traps $n+1$ linearly independent vectors in a space spanned by $n$
        vectors, so that all bases have the same size and $\dim V$ is defined.
  \item $\dim(U+W)=\dim U+\dim W-\dim(U\cap W)$; $U+W$ is direct exactly when
        $U\cap W=\{\mathbf{0}\}$; every finite-dimensional subspace has a complement.
\end{itemize}
\end{ringkasan}

\section{Exercises}
\tingkat{basic}

\begin{latihan}
\begin{enumerate}
  \item Decide whether $W=\{(x,y,z)\in\mathbb{R}^{3}:x-2y+z=0\ \text{and}\ x+y=0\}$
        is a subspace of $\mathbb{R}^{3}$. If so, find a basis and its dimension.
  \item Prove that the set of all $2\times2$ matrices whose main diagonal entries are zero
        is a subspace of $M_{2\times2}(\mathbb{R})$, then find a basis and its dimension.
  \item Show that $S=\{1+x,\ 1+x^{2},\ 2+x+x^{2}\}$ is linearly dependent in $P_{2}$.
        Find $\dim\operatorname{span}(S)$, then extend a largest linearly independent
        subset of $S$ to a basis of $P_{2}$.
  \item Prove: if $\{\mathbf{v}_{1},\mathbf{v}_{2},\mathbf{v}_{3}\}$ is linearly independent in
        $V$, then $\{\mathbf{v}_{1}+\mathbf{v}_{2},\ \mathbf{v}_{2}+\mathbf{v}_{3},\
        \mathbf{v}_{3}+\mathbf{v}_{1}\}$ is also linearly independent.
  \item Compute $\dim(U\cap W)$ and $\dim(U+W)$ for the two subspaces
        \begin{align*}
          U &= \operatorname{span}\{(1,0,0),(0,1,0)\}, \\
          W &= \operatorname{span}\{(1,1,0),(0,0,1)\}
        \end{align*}
        in $\mathbb{R}^{3}$. Check the dimension-of-a-sum formula, then compute them with numpy
        as well.
  \item For the $3\times4$ matrix with rows $(1,2,3,4)$, $(2,4,6,9)$, and $(0,1,1,1)$,
        find the rank, one basis of the column space, and one basis of the null space.
        Check with the function \texttt{matrix\_rank} from numpy and
        \texttt{null\_space} from scipy.
  \item Assess this claim: \emph{if $U_{1},U_{2},U_{3}$ are subspaces of $V$ and every pairwise
        intersection equals $\{\mathbf{0}\}$, then
        $U_{1}+U_{2}+U_{3}=U_{1}\oplus U_{2}\oplus U_{3}$.} State true, false, or not enough
        evidence, and include a reason or a counterexample.
  \item Prove that if $\dim V=n$ and $U$ is a subspace with $\dim U=n-1$, then there is a
        vector $\mathbf{v}\in V$ such that $V=U\oplus\operatorname{span}\{\mathbf{v}\}$.
  \item Find the dimension of the space of all polynomials of degree at most $4$ that vanish at
        $x=0$ and $x=1$ simultaneously, then give one basis of it. Check your answer by
        computing the rank of the $2\times5$ evaluation matrix at those two points.
  \item Given $B=\{1,\ 1+x,\ 1+x+x^{2},\ 1+x+x^{2}+x^{3}\}$ a basis of $P_{3}$. Find the
        coordinates of $q=4+2x-5x^{2}+x^{3}$ with respect to $B$, then check by solving the
        linear system in numpy.
\end{enumerate}

\noindent\textbf{Hints for the three hardest items.} Item 4: suppose
$\sum c_{i}(\dots)=\mathbf{0}$, collect the coefficient of each $\mathbf{v}_{i}$, then use the
independence of $\{\mathbf{v}_{1},\mathbf{v}_{2},\mathbf{v}_{3}\}$. Item 8: take a basis of $U$,
extend it by one vector, then use consequence (d). Item 9: a polynomial of degree at most
$4$ that vanishes at $0$ and $1$ is divisible by $x(x-1)$; count the free parameters.
\end{latihan}

\begin{untukpeneliti}
Three things are genuinely proved in this chapter, and all three hold over any field: the
uniqueness of the size of a basis through the exchange theorem (Steinitz's result), the formula
$\dim(U+W)=\dim U+\dim W-\dim(U\cap W)$, and the existence of a complement for every
finite-dimensional subspace. The popular claim that ``rank equals the number of informative
features'' is meaningful only if ``informative'' means linearly independent, and on real data that
condition is never satisfied exactly.

Primary sources: Anton, \emph{Elementary Linear Algebra} 11th edition, Chapter 4 sections
4.1--4.5; Axler, \emph{Linear Algebra Done Right} 4th edition, sections 1B, 1C, 2A--2C; Hoffman and
Kunze, \emph{Linear Algebra} 2nd edition, Chapter 2.

Open question: how far does $\dim(U\cap W)$ shift as a result of perturbing the entries? Limits of
validity: all results assume the vector space axioms; for infinite-dimensional spaces, the
existence of a basis requires Zorn's lemma. A rank computed by a computer depends on a threshold.
\end{untukpeneliti}

\peta{This chapter follows Anton, \emph{Elementary Linear Algebra} 11th edition, Chapter 4
(General Vector Spaces) section 4.1 (Real Vector Spaces) for the axioms and examples, 4.2
(Subspaces) for the subspace test and span, 4.3 (Linear Independence) for independence and the
dependence lemma, and 4.5 (Dimension) for the exchange theorem and its consequences; other
editions number these differently. A more abstract counterpart is in Axler, \emph{Linear Algebra
Done Right} 4th edition, Chapter 1 sections 1B and 1C (including sums and direct sums of
subspaces) and Chapter 2 sections 2A, 2B, and 2C, where the exchange theorem appears as a
comparison of the length of a linearly independent list with that of a spanning list. The
treatment of coordinates, linear transformations, and rank-nullity is not in this chapter but in
Chapter 7.}


\chapter{Basis, Dimension, and Linear Transformations}

An anomaly detector does not read raw network traffic. It summarizes it first. Four
counters, say failed connections, outbound packets, session length, and the number of
unique ports, are compressed into three features by the linear map
$F\colon\mathbb{R}^{4}\to\mathbb{R}^{3}$.
The decision threshold only looks at $F(\mathbf{x})$.

The decisive question: is there $\boldsymbol{\delta}\neq\mathbf{0}$ with
$F(\mathbf{x}+\boldsymbol{\delta})=F(\mathbf{x})$? Because $F$ is linear, that condition is
equivalent to
$F(\boldsymbol{\delta})=\mathbf{0}$. The answer is a subspace of dimension
$4-\operatorname{rank}F$; if $F$ has rank three, there is a whole line of blind directions,
and an attacker can shift all four counters along that line without shifting a single one of
the features.

To answer it we need three things Chapter 6 has not yet given: the list of coordinates with
respect to a basis, the rule for moving that list when the basis changes, and the translation
of a linear map into a matrix so that the kernel can be computed by elimination. Chapter 6
already established that a basis exists and that its size is unique; its axioms are not
repeated here.

\begin{untukanak}
Imagine a bookcase with many shelves. Every book has a fixed place. If we rearrange the
bookcase, the books do not change, but each book's address does. To say where a book is, we
must name the arrangement currently in use. A friend who uses another arrangement will name a
different address for the same book,
and the two of you are both right.

There is also a special shelf behind the bookcase. A book moved there cannot be seen from the
front. As long as books only move within that shelf, the list of books in front does not
change at all.

Have you ever moved something to a place nobody inspects? We are going to find out why the
list of addresses can change while the contents of
the bookcase stay the same, why that hidden shelf always has a definite size, and why some
changes do not shift the summary at all.
\end{untukanak}

\section{Coordinates and transition matrices}
\tingkat{intermediate}

Vectors in $\mathbb{R}^{3}$ are already lists of three numbers; vectors in $P_{2}$ or
$M_{2\times2}(\mathbb{R})$ are not. So that both can be computed in the same way, we
need a translation that turns every vector into a list of numbers. That translation depends
on the basis, and that is where the difficulty of this section begins.

Uniqueness we already have: the uniqueness lemma for coordinates in Chapter 6 states that the
expression of a vector in a basis is always unique. Therefore the following definition is
legitimate without choosing among several expressions.

\kamus{coordinate vector}{koordinat}{the address of a book in the bookcase: its contents stay
fixed, its address depends on how the shelves are arranged}

\begin{definisi}[Coordinate vector]
Let $B=(\mathbf{b}_{1},\dots,\mathbf{b}_{n})$ be an ordered basis of the vector space
$V$
over $\mathbb{R}$. For every $\mathbf{v}\in V$ there are unique scalars $a_{1},\dots,a_{n}$
with $\mathbf{v}=a_{1}\mathbf{b}_{1}+\dots+a_{n}\mathbf{b}_{n}$. The column
\[
  [\mathbf{v}]_{B}=(a_{1},\dots,a_{n})^{\top}
\]
is called the \emph{coordinate vector} of $\mathbf{v}$ with respect to $B$.\index{coordinate}
\end{definisi}

The word ``ordered'' is not decoration. If two basis vectors are exchanged, two coordinate
entries are exchanged with them while the vector stays put. A basis set alone does not
determine the coordinate list;
its order does as well.

One warning: $[\mathbf{v}]_{B}$ is not $\mathbf{v}$. Those numbers live in
$\mathbb{R}^{n}$, whereas $\mathbf{v}$ lives in $V$. For $M_{2\times2}(\mathbb{R})$
the difference is palpable: the vector holds four numbers, and its coordinates are four
numbers too, but arranged as a column that depends on the basis.

The first important property of this translation: it respects both operations of the vector
space.

\begin{proposisi}
For every basis $B$ of a space $V$ of dimension $n$, the map
$\varphi_{B}\colon V\to\mathbb{R}^{n}$ with $\varphi_{B}(\mathbf{v})=[\mathbf{v}]_{B}$
is linear, injective, and surjective.
\end{proposisi}

\begin{proof}
Take $\mathbf{v}=\sum_{i}a_{i}\mathbf{b}_{i}$ and $\mathbf{w}=\sum_{i}c_{i}\mathbf{b}_{i}$.
Then $\mathbf{v}+\mathbf{w}=\sum_{i}(a_{i}+c_{i})\mathbf{b}_{i}$. Uniqueness of coordinates
forces $[\mathbf{v}+\mathbf{w}]_{B}=(a_{i}+c_{i})=[\mathbf{v}]_{B}+[\mathbf{w}]_{B}$.
In the same way, $[c\mathbf{v}]_{B}=c[\mathbf{v}]_{B}$.

For injectivity, suppose $[\mathbf{v}]_{B}=\mathbf{0}$. Then all its coefficients are zero,
so $\mathbf{v}=\sum_{i}0\cdot\mathbf{b}_{i}=\mathbf{0}$. For surjectivity, take
any $(a_{1},\dots,a_{n})\in\mathbb{R}^{n}$ and form
$\mathbf{v}=\sum_{i}a_{i}\mathbf{b}_{i}$; its coordinates are exactly that list.
\end{proof}

The map $\varphi_{B}$ turns a question about $V$ into a question about columns in
$\mathbb{R}^{n}$, then turns it back; that transfer depends on $B$. Coordinates
are written as columns so that matrices act from the left: coordinates in, multiplied by the
matrix, coordinates out. With row coordinates, every formula would have to be transposed.

\begin{bacakeras}
Two voices stand in front of the class.

``How long is this table?''
``Three.''
``Three what?''
``Three rulers. Or three spans. Or three paces. Choose the measuring instrument first.''

It is the same with a vector space. A vector has no address until we choose a basis. The basis
is
the ruler: it supplies a scale, a direction, and an order. Change the ruler and you change
the numbers. The object being measured does not change with them.

One ruler, one number. Two rulers, two numbers. Without a ruler, no numbers
at all.

Be careful with a bad ruler: two of its scale marks are almost parallel. Numbers still come
out,
but the numbers wobble, and two people measuring patiently will obtain different
numbers.

Repeat together: a basis is not the object, a basis is the instrument. What changes is the
address, what stays is the content.
\end{bacakeras}

Now two bases are used on the same space. Every vector has two coordinate lists,
and we want a formula that turns the first list into the second. The matrix that
performs that task is built from the coordinates of the old basis vectors in the new basis.

\begin{definisi}[Transition matrix]
Let $B=(\mathbf{b}_{1},\dots,\mathbf{b}_{n})$ and
$C=(\mathbf{c}_{1},\dots,\mathbf{c}_{n})$
be two bases of $V$. The \emph{transition matrix} from $B$ to $C$ is the $n\times n$ matrix
\[
P_{C\leftarrow
B}=\Bigl(\,[\mathbf{b}_{1}]_{C}\;\Big|\;\cdots\;\Big|\;[\mathbf{b}_{n}]_{C}\,\Bigr),
\]
that is, the matrix whose $j$-th column is $[\mathbf{b}_{j}]_{C}$.\index{transition matrix}
\end{definisi}

What still needs proof: multiplying the old coordinate list by this matrix produces the new
coordinate
list.

\begin{teorema}
For every $\mathbf{v}\in V$ we have $[\mathbf{v}]_{C}=P_{C\leftarrow B}[\mathbf{v}]_{B}$.
\end{teorema}

\begin{proof}
Write $\mathbf{v}=\sum_{j}a_{j}\mathbf{b}_{j}$, so that
$[\mathbf{v}]_{B}=(a_{1},\dots,a_{n})^{\top}$.
Linearity of $\varphi_{C}$ gives
$[\mathbf{v}]_{C}=\sum_{j}a_{j}[\mathbf{b}_{j}]_{C}$. That right-hand side is exactly the
matrix
product $P_{C\leftarrow B}$ with the column $(a_{j})$, by the definition of matrix-vector
multiplication.
\end{proof}

A corollary that makes the matrix usable in both directions.

\begin{akibat}
The transition matrix is always invertible, and $P_{B\leftarrow C}=P_{C\leftarrow B}^{-1}$.
Moreover
$P_{C\leftarrow B}=P_{C\leftarrow D}\,P_{D\leftarrow B}$ for three bases $B,C,D$ of $V$.
\end{akibat}

\begin{proof}
The previous theorem gives $P_{B\leftarrow C}P_{C\leftarrow B}[\mathbf{v}]_{B}=[\mathbf{v}]_{B}$
for all $\mathbf{v}$. Because $\varphi_{B}$ is surjective, that product equals the identity
matrix. The same argument applies in the opposite order, so the two matrices are mutually
inverse. For the chain rule, apply the theorem twice:
$[\mathbf{v}]_{C}=P_{C\leftarrow D}[\mathbf{v}]_{D}=P_{C\leftarrow D}P_{D\leftarrow B}[\mathbf{v}]_{B}$,
then compare with $[\mathbf{v}]_{C}=P_{C\leftarrow B}[\mathbf{v}]_{B}$ at $\mathbf{v}=\mathbf{b}_{j}$.
\end{proof}

Note that $P_{C\leftarrow B}$ and $P_{B\leftarrow C}$ are not transposes of each other; they are mutual inverses, and that is a different thing. The direction of the subscript is read from right to left, like the composition of maps: take the $j$-th basis vector of $B$, then ask for its address according to $C$. When $B$ is the standard basis, the columns of $P_{B\leftarrow C}$ are the vectors of $C$ themselves.

Take $B=(\mathbf{e}_{1},\mathbf{e}_{2},\mathbf{e}_{3})$
and $C=(\mathbf{b}'_{1},\mathbf{b}'_{2},\mathbf{b}'_{3})$ with
\[
  \mathbf{b}'_{1}=(1,1,0),\qquad \mathbf{b}'_{2}=(1,0,1),\qquad \mathbf{b}'_{3}=(0,1,1).
\]
Because $B$ is standard, the matrix $P_{B\leftarrow C}$ contains those three vectors as columns:
\[
  P_{B\leftarrow C}=\begin{pmatrix}1&1&0\\ 1&0&1\\ 0&1&1\end{pmatrix},
  \qquad \det P_{B\leftarrow C}=-2\neq0 .
\]
Its inverse is
\[
  P_{C\leftarrow B}=\frac{1}{2}\begin{pmatrix}1&1&-1\\ 1&-1&1\\ -1&1&1\end{pmatrix}.
\]
For $\mathbf{v}=(3,8,11)$ we obtain
\[
  [\mathbf{v}]_{C}=P_{C\leftarrow B}[\mathbf{v}]_{B}
  =\frac{1}{2}\begin{pmatrix}3+8-11\\ 3-8+11\\ -3+8+11\end{pmatrix}
  =\begin{pmatrix}0\\ 3\\ 8\end{pmatrix}.
\]
The check is immediate: $3\mathbf{b}'_{2}+8\mathbf{b}'_{3}=(3,8,11)$. Conversely
$\mathbf{w}=(1,0,0)$ has coordinates $(1/2,\,1/2,\,-1/2)$; a fractional entry is natural,
because
coordinates are nothing but an address.

\section{Linear transformations, kernel, and range}
\tingkat{basic}

Chapter 6 discussed spaces and subspaces, but not maps between them. Yet the main work of
linear algebra is moving vectors: rotating, projecting, differentiating, summarizing.

A rule that respects addition and scalar multiplication is called a linear transformation:
adding first and then mapping gives the same result as mapping first and then adding.

\begin{definisi}[Linear transformation]
A map $T\colon V\to W$ between two vector spaces over $\mathbb{R}$ is called a
\emph{linear transformation} if $T(\mathbf{u}+\mathbf{v})=T(\mathbf{u})+T(\mathbf{v})$ and
$T(c\mathbf{v})=cT(\mathbf{v})$ for all $\mathbf{u},\mathbf{v}\in V$ and $c\in\mathbb{R}$.
\end{definisi}

The two conditions are equivalent to a single condition,
$T(c\mathbf{u}+\mathbf{v})=cT(\mathbf{u})+T(\mathbf{v})$. The first consequence,
$T(\mathbf{0})=\mathbf{0}$, is used to reject linearity quickly: the map
$T(\mathbf{x})=\mathbf{x}+\mathbf{a}$ with $\mathbf{a}\neq\mathbf{0}$ fails at once.

Conversely $T(\mathbf{0})=\mathbf{0}$ is not enough. The map $T(x,y)=(x^{2},y)$ satisfies it,
but
$T(2,0)=(4,0)$ whereas $2T(1,0)=(2,0)$. Squaring breaks linearity.\index{linear
transformation}

A $90^{\circ}$ rotation, $R(x,y)=(-y,x)$, is linear: $R(3,1)=(-1,3)$. The derivative $D\colon P_{3}\to P_{2}$, $D(p)=p'$, is linear because $(p+q)'=p'+q'$. The map $T(x,y)=(x+y,\,2x+2y)$ is linear too, but all of $\mathbb{R}^{2}$ falls onto a single line; it is linear and simultaneously destroys information.

Take
$A=\begin{pmatrix}2&1\\1&2\end{pmatrix}$ and $T(\mathbf{x})=A\mathbf{x}$. The vector $(1,0)$ goes
to $(2,1)$ and $(0,1)$ to $(1,2)$, so the unit square becomes a parallelogram with
vertices
$(0,0)$, $(2,1)$, $(3,3)$, $(1,2)$ of area $3$, since $|2\cdot2-1\cdot1|=3$.
Figure~\ref{fig:kisi-07} shows the grid before and after the map: lines stay lines, the grid
stays a grid, but distances and areas are no longer the same.

\begin{figure}[htbp]
\centering
\begin{tikzpicture}[scale=0.38,>=stealth,font=\scriptsize]
  \foreach \k in {0,1,2,3}{
    \draw[layergrey!45] (\k,0)--(\k,3);
    \draw[layergrey!45] (0,\k)--(3,\k);
    \draw[layerteal,thick] (2*\k,\k)--(2*\k+3,\k+6);
    \draw[layerteal,thick] (\k,2*\k)--(\k+6,2*\k+3);}
  \fill[accent!15] (0,0)--(1,0)--(1,1)--(0,1)--cycle;
  \draw[accent,thick] (0,0)--(1,0)--(1,1)--(0,1)--cycle;
  \fill[layerorange!25] (0,0)--(2,1)--(3,3)--(1,2)--cycle;
  \draw[layerorange,very thick] (0,0)--(2,1)--(3,3)--(1,2)--cycle;
  \draw[->,ink,thick] (-0.6,0)--(9.8,0) node[below right]{$x$};
  \draw[->,ink,thick] (0,-0.6)--(0,9.8) node[left]{$y$};
  \foreach \x/\y in {2/1,1/2,3/3}{
    \fill[ink] (\x,\y) circle (1.8pt);
    \node[ink,anchor=north west] at (\x+0.05,\y-0.05) {$(\x,\y)$};}
  \node[layergrey,anchor=north west] at (0,-1.0) {thin: original grid};
  \node[layerteal,anchor=north west] at (4.0,-1.0) {thick: grid after mapping};
\end{tikzpicture}
\caption{The unit grid and its image under $T(\mathbf{x})=A\mathbf{x}$ with
$A=\begin{pmatrix}2&1\\1&2\end{pmatrix}$. One cell becomes a parallelogram with
vertices $(0,0)$, $(2,1)$, $(3,3)$, and $(1,2)$, so its area grows from $1$ to $3$.
The old grid is drawn thin, the new grid thick; both remain grids.}
\label{fig:kisi-07}
\end{figure}

\begin{tangan}
Take a sheet of graph paper, a pencil, and a ruler. Draw a horizontal axis and a vertical
axis,
then mark their intersection as the center. Draw a small house: a square five cells wide, a
triangular
roof, a door one cell wide.

Now translate. Move every point of the house three cells to the right and two cells up.
Do not guess; count. A point that was one cell right and one cell up
is now
four cells right and three cells up. Redraw the house. Its shape and
size are the same, only its place has moved.

Now rotate. Turn the paper a quarter turn clockwise, then redraw the same
house. The length of each side stays, the direction of each side changes.

What should be visible: under translation, all directions are preserved but the center moves with the figure; under rotation, the center stays and all directions change. Because a translation moves the center, a translation is
not
a linear transformation. A rotation is.
\end{tangan}

The good news: to know a linear map, it is enough to know its values on the basis vectors.

\begin{proposisi}
Let $B=(\mathbf{b}_{1},\dots,\mathbf{b}_{n})$ be a basis of $V$ and let
$\mathbf{w}_{1},\dots,\mathbf{w}_{n}\in W$ be arbitrary. There is exactly one linear map
$T\colon V\to W$ with $T(\mathbf{b}_{j})=\mathbf{w}_{j}$ for every $j$.
\end{proposisi}

\begin{proof}
For existence, define
$T\bigl(\sum_{j}a_{j}\mathbf{b}_{j}\bigr)=\sum_{j}a_{j}\mathbf{w}_{j}$.
This definition does not depend on the expression, because coordinates with respect to a basis
are unique. For linearity, use the linearity of coordinates from the proposition in the
previous section:
$T(\mathbf{u}+\mathbf{v})=\sum_{j}(a_{j}+c_{j})\mathbf{w}_{j}=T(\mathbf{u})+T(\mathbf{v})$,
and similarly for scalar multiplication.

For uniqueness, suppose $S$ is linear and $S(\mathbf{b}_{j})=\mathbf{w}_{j}$. For
$\mathbf{v}=\sum_{j}a_{j}\mathbf{b}_{j}$ we have
$S(\mathbf{v})=\sum_{j}a_{j}S(\mathbf{b}_{j})=\sum_{j}a_{j}\mathbf{w}_{j}=T(\mathbf{v})$.
\end{proof}

The practical corollary: $n$ vectors in $W$ determine one linear map from an $n$-dimensional
$V$.
When $W=\mathbb{R}^{m}$, that data is nothing other than an $m\times n$ matrix.

Every linear map carries two sets worth inspecting: the vectors mapped to zero,
and the values that are actually attained.

\kamus{kernel or null space}{inti}{the list of changes that do not shift the
summary at all}

\begin{definisi}[Kernel, range, rank, and nullity]
For a linear $T\colon V\to W$, the \emph{kernel} is
$\ker T=\{\mathbf{v}\in V: T(\mathbf{v})=\mathbf{0}\}$ and the \emph{range} is
$\operatorname{im}T=\{T(\mathbf{v}):\mathbf{v}\in V\}$. When $V$ is finite-dimensional,
$\operatorname{nullity}(T)=\dim\ker T$ and $\operatorname{rank}(T)=\dim\operatorname{im}T$.
\end{definisi}

Both are subspaces, not merely sets.\index{kernel}\index{range}

\begin{lemma}
$\ker T$ is a subspace of $V$ and $\operatorname{im}T$ is a subspace of $W$.
\end{lemma}

\begin{proof}
Both contain $\mathbf{0}$ because $T(\mathbf{0})=\mathbf{0}$. If
$\mathbf{u},\mathbf{v}\in\ker T$
and $c\in\mathbb{R}$, then $T(\mathbf{u}+\mathbf{v})=\mathbf{0}$ and
$T(c\mathbf{u})=\mathbf{0}$,
so $\ker T$ is closed. If $\mathbf{u}=T(\mathbf{x})$ and $\mathbf{v}=T(\mathbf{y})$, then
$\mathbf{u}+\mathbf{v}=T(\mathbf{x}+\mathbf{y})$ and $c\mathbf{u}=T(c\mathbf{x})$ belong to
$\operatorname{im}T$.
\end{proof}

The kernel measures the failure of injectivity, and that is all that needs checking.

\begin{proposisi}
$T$ is injective if and only if $\ker T=\{\mathbf{0}\}$.
\end{proposisi}

\begin{proof}
If $T$ is injective, only $\mathbf{0}$ maps to $\mathbf{0}$. Conversely, suppose
$\ker T=\{\mathbf{0}\}$ and $T(\mathbf{u})=T(\mathbf{v})$. Then $T(\mathbf{u}-\mathbf{v})=\mathbf{0}$,
so $\mathbf{u}-\mathbf{v}\in\ker T$, hence $\mathbf{u}=\mathbf{v}$.
\end{proof}

Rank counts the directions that survive, nullity counts the directions that vanish. Together
they divide the dimension of the source space.

\kamus{rank}{rank}{the number of independent directions that still survive after the map}

A small example shows both fates. Take
$B=\begin{pmatrix}1&2\\2&4\end{pmatrix}$. Then $B\mathbf{e}_{1}=(1,2)$ and
$B\mathbf{e}_{2}=(2,4)$, two collinear vectors, so $\operatorname{rank}(B)=1$. The direction $(1,2)$
survives and stretches: $B(1,2)=(5,10)$. The direction $(2,-1)$ vanishes: $B(2,-1)=(0,0)$.
Figure~\ref{fig:arah} draws both, and rank-nullity closes the count: $2=1+1$.

\begin{figure}[htbp]
\centering
\begin{tikzpicture}[scale=0.60,>=stealth,font=\scriptsize]
  \draw[layerjade!60,thin] (-0.7,-1.4)--(2.8,5.6);
  \node[layerjade,anchor=west,rotate=63] at (2.0,4.3) {line in direction $(1,2)$};
  \draw[->,ink,thick] (-1.0,0)--(3.3,0) node[below right]{$x$};
  \draw[->,ink,thick] (0,-2.0)--(0,5.0) node[left]{$y$};
  \draw[->,layergrey,thick] (0,0)--(1,0) node[below]{$\mathbf{e}_{1}$};
  \draw[->,layergrey,thick] (0,0)--(0,1) node[left]{$\mathbf{e}_{2}$};
  \draw[->,layerorange,very thick] (0,0)--(1,2) node[above left]{$B\mathbf{e}_{1}=(1,2)$};
  \draw[->,layerbrown,very thick] (0,0)--(2,4) node[right]{$B\mathbf{e}_{2}=(2,4)$};
  \draw[->,layerplum,thick,dashed] (0,0)--(1.8,-0.9) node[below right]{$(2,-1)$};
  \node[layerplum,anchor=north] at (1.5,-1.6) {collapses to $\mathbf{0}$};
  \fill[ink] (0,0) circle (1.6pt) node[below left,ink]{$\mathbf{0}$};
\end{tikzpicture}
\caption{Two directions under $B=\begin{pmatrix}1&2\\2&4\end{pmatrix}$. The direction $(1,2)$
survives and stretches into $(5,10)$; the direction $(2,-1)$ collapses to $\mathbf{0}$. The
images of the two basis vectors,
$B\mathbf{e}_{1}=(1,2)$ and $B\mathbf{e}_{2}=(2,4)$, lie on a single line, so
$\operatorname{rank}(B)=1$ and $\operatorname{nullity}(B)=1$.}
\label{fig:arah}
\end{figure}

Now the central theorem of this chapter. Its content: no dimension is lost along the way.

\begin{teorema}[Rank-nullity]
If $V$ is finite-dimensional and $T\colon V\to W$ is linear, then
\[
  \dim V=\dim\ker T+\dim\operatorname{im}T .
\]
\end{teorema}

\begin{proof}
Let $\dim V=n$ and $\dim\ker T=k$. Take a basis
$(\mathbf{u}_{1},\dots,\mathbf{u}_{k})$ of $\ker T$. By the basis extension corollary in Chapter 6,
that basis can be completed to a basis of $V$ by adding $n-k$ vectors
$\mathbf{w}_{1},\dots,\mathbf{w}_{n-k}$.

First, the set $\{T(\mathbf{w}_{1}),\dots,T(\mathbf{w}_{n-k})\}$ spans
$\operatorname{im}T$. For an arbitrary $T(\mathbf{v})$, write
$\mathbf{v}=\sum_{i}\alpha_{i}\mathbf{u}_{i}+\sum_{j}\beta_{j}\mathbf{w}_{j}$. The terms
$\mathbf{u}_{i}$ vanish because $T(\mathbf{u}_{i})=\mathbf{0}$, so
$T(\mathbf{v})=\sum_{j}\beta_{j}T(\mathbf{w}_{j})$.

Second, that set is linearly independent. Suppose $\sum_{j}\beta_{j}T(\mathbf{w}_{j})=\mathbf{0}$.
Then $T\bigl(\sum_{j}\beta_{j}\mathbf{w}_{j}\bigr)=\mathbf{0}$, so
$\sum_{j}\beta_{j}\mathbf{w}_{j}\in\ker T$, and there are scalars $\alpha_{i}$ with
$\sum_{j}\beta_{j}\mathbf{w}_{j}=\sum_{i}\alpha_{i}\mathbf{u}_{i}$. Move all terms to
one side:
\[
  \sum_{j}\beta_{j}\mathbf{w}_{j}-\sum_{i}\alpha_{i}\mathbf{u}_{i}=\mathbf{0}.
\]
That relation lies inside a basis of $V$, so all its coefficients are zero; in particular every
$\beta_{j}=0$. Hence $\{T(\mathbf{w}_{j})\}$ is a basis of $\operatorname{im}T$, and
$\operatorname{rank}(T)=n-k$.
\end{proof}

Figure~\ref{fig:batang} arranges that theorem as two bars. For a space of dimension
$4$ the total length is always $4$: the identity gives $4+0$, the projection that discards one
coordinate
gives $3+1$, and the feature map $G$ later gives $2+2$. Only the position of the divider
shifts.

\begin{figure}[htbp]
\centering
\begin{tikzpicture}[x=1cm,y=1cm,font=\scriptsize,>=stealth]
  \fill[layerjade!45] (0,2.4) rectangle (4,3.1);
  \draw[layergrey] (0,2.4) rectangle (4,3.1);
  \node at (2,2.75) {rank $4$};
  \node[layerjade,anchor=east] at (-0.12,2.75) {identity};
  \node[layerjade,anchor=west] at (4.12,2.75) {nullity $0$};
  \fill[layerteal!45] (0,1.2) rectangle (3,1.9);
  \fill[layerorange!45] (3,1.2) rectangle (4,1.9);
  \draw[layergrey] (0,1.2) rectangle (4,1.9);
  \draw[layergrey] (3,1.2)--(3,1.9);
  \node at (1.5,1.55) {rank $3$};
  \node at (3.5,1.55) {$1$};
  \node[layerteal,anchor=east] at (-0.12,1.55) {projection};
  \node[layerorange,anchor=west] at (4.12,1.55) {nullity $1$};
  \fill[layerteal!45] (0,0) rectangle (2,0.7);
  \fill[layerorange!45] (2,0) rectangle (4,0.7);
  \draw[layergrey] (0,0) rectangle (4,0.7);
  \draw[layergrey] (2,0)--(2,0.7);
  \node at (1,0.35) {rank $2$};
  \node at (3,0.35) {$2$};
  \node[layerteal,anchor=east] at (-0.12,0.35) {feature map $G$};
  \node[layerorange,anchor=west] at (4.12,0.35) {nullity $2$};
  \draw[layergrey,dashed] (4,-0.25)--(4,3.4);
  \node[layergrey,anchor=south] at (4,3.4) {$\dim V=4$};
  \fill[layerteal!45] (0,-0.95) rectangle (0.35,-0.6);
  \draw[layergrey] (0,-0.95) rectangle (0.35,-0.6);
  \node[anchor=west] at (0.45,-0.78) {directions that survive};
  \fill[layerorange!45] (3.3,-0.95) rectangle (3.65,-0.6);
  \draw[layergrey] (3.3,-0.95) rectangle (3.65,-0.6);
  \node[anchor=west] at (3.75,-0.78) {directions that vanish};
\end{tikzpicture}
\caption{Rank-nullity as two bars whose sum is fixed. For $\dim V=4$: the identity map
gives $4+0$, the projection that discards one coordinate gives $3+1$, and the feature map $G$
gives
$2+2$. Only the position of the divider shifts; the total length stays $4$.}
\label{fig:batang}
\end{figure}

Two corollaries that are used most often.

\begin{akibat}
Let $\dim V=\dim W$ be finite and let $T\colon V\to W$ be linear. Then $T$ is injective if and
only if $T$ is surjective. Moreover
$\operatorname{rank}(T)\le\min(\dim V,\dim W)$.
\end{akibat}

\begin{proof}
By rank-nullity, $\dim\ker T=n-\operatorname{rank}T$ with $n=\dim V$. Injectivity holds exactly
when $\dim\ker T=0$, that is $\operatorname{rank}T=n=\dim W$, which means
$\operatorname{im}T=W$ because a subspace of full dimension fills the whole space. The
reverse implication is obtained by reversing the order of the argument. The upper bound on the rank is clear from
$\operatorname{im}T\subseteq W$ and from rank-nullity.
\end{proof}

As a numerical illustration, take $T\colon\mathbb{R}^{3}\to\mathbb{R}^{2}$ with
$T(x,y,z)=(x-y,\,y-z)$. Because $T(t,t,t)=(0,0)$, its kernel is the line $\{(t,t,t)\}$ and its nullity is
$1$.
The images $T(1,0,0)=(1,0)$ and $T(0,0,1)=(0,-1)$ span $\mathbb{R}^{2}$, so the rank is $2$;
rank-nullity: $3=1+2$.

For machine learning, that corollary is harsh. A feature map $\mathbb{R}^{n}\to\mathbb{R}^{m}$
with $n>m$ cannot possibly be injective, and its kernel has dimension at least $n-m$. Every dimensional
compression stores blind directions; the question is not whether such a direction exists, but
whether an attacker can find it.

The standard procedure for a basis of the kernel: bring the matrix to reduced row echelon form, mark the columns
without a pivot as free variables, then in turn give each free variable the value one and the other
free variables zero. The pivot variables are read off from the rows. The result is one vector per
column without a pivot, linearly independent because each has a one in a position no other
one has.

For the range, the procedure is different and is often confused with the kernel procedure above. The following
four steps are worth memorizing.

\begin{resep}
\begin{enumerate}
  \item Assemble the matrix $A$, then reduce it to reduced row echelon form by elimination.
  \item Mark the columns that contain a pivot. A column without a pivot is not a candidate basis member.
  \item Go back to the original matrix $A$, not to its echelon form. Take the original columns whose positions
        coincide with the pivot columns.
  \item Check: the number of selected vectors equals the number of nonzero rows, and those vectors are
        linearly independent.
\end{enumerate}
A warning about order: step three must not take columns from the echelon form. Row operations
change the columns, whereas the column space does not change. Taking columns from the echelon matrix
gives a wrong answer even when the count is right.
\end{resep}

The kernel also answers a more general question: once a value has been attained, all its
preimages are neatly arranged. Here $T^{-1}(\mathbf{w})$ denotes the set of preimages, not the inverse
map.

\begin{proposisi}
Let $T\colon V\to W$ be linear and $T(\mathbf{v}_{0})=\mathbf{w}$. Then
$T^{-1}(\mathbf{w})=\{\mathbf{v}_{0}+\mathbf{k}:\mathbf{k}\in\ker T\}$.
\end{proposisi}

\begin{proof}
If $\mathbf{k}\in\ker T$, then $T(\mathbf{v}_{0}+\mathbf{k})=\mathbf{w}+\mathbf{0}=\mathbf{w}$,
so every $\mathbf{v}_{0}+\mathbf{k}$ belongs to the preimage of $\mathbf{w}$. Conversely, if
$T(\mathbf{v})=\mathbf{w}$, then $T(\mathbf{v}-\mathbf{v}_{0})=\mathbf{0}$, so
$\mathbf{v}-\mathbf{v}_{0}\in\ker T$ and $\mathbf{v}=\mathbf{v}_{0}+(\mathbf{v}-\mathbf{v}_{0})$.
\end{proof}

The set of inputs that produce the same features is a shifted copy of the kernel: when the kernel
has dimension one, the family is a line, namely the blind direction of the chapter opening; when it has dimension
$k$,
an attacker has $k$ directions.

\section{Representation matrices and composition}
\tingkat{intermediate}

Now the two earlier tools are combined: coordinates turn a vector into a column, and a linear
map is determined by its values on a basis. The combination produces a matrix, where all the
computation takes place.

The representation matrix stores the image of each basis vector as one column, written
in the coordinates of the target basis.

\begin{definisi}[Representation matrix]
Let $B=(\mathbf{b}_{1},\dots,\mathbf{b}_{n})$ be a basis of $V$,
$C=(\mathbf{c}_{1},\dots,\mathbf{c}_{m})$ a basis of $W$, and $T\colon V\to W$ linear.
The \emph{representation matrix} of $T$ with respect to $B$ and $C$ is the $m\times n$ matrix
\[
[T]_{C\leftarrow
B}=\Bigl(\,[T(\mathbf{b}_{1})]_{C}\;\Big|\;\cdots\;\Big|\;[T(\mathbf{b}_{n})]_{C}\,\Bigr).
\]
\end{definisi}

Those columns are obtained by solving systems, not by guessing: for each $j$,
solve $\sum_{i}x_{i}\mathbf{c}_{i}=T(\mathbf{b}_{j})$, then write the solution as
column
$j$. When $C$ is the standard basis, that solving need not be carried out, because coordinates
with respect to the standard basis are the vector itself.

The following claim explains why the matrix is useful: it works in the expected way.

\begin{teorema}
For every $\mathbf{v}\in V$ we have
$[T(\mathbf{v})]_{C}=[T]_{C\leftarrow B}\,[\mathbf{v}]_{B}$.
\end{teorema}

\begin{proof}
Write $\mathbf{v}=\sum_{j}a_{j}\mathbf{b}_{j}$, so $[\mathbf{v}]_{B}=(a_{j})$. Linearity of
$T$ gives $T(\mathbf{v})=\sum_{j}a_{j}T(\mathbf{b}_{j})$. Taking coordinates is also
linear,
so $[T(\mathbf{v})]_{C}=\sum_{j}a_{j}[T(\mathbf{b}_{j})]_{C}$, which is exactly the matrix
product with the column $(a_{j})$.
\end{proof}

Two immediate corollaries. A linear map and its matrix determine each other, because the values on a
basis already suffice. And the rank does not change under that translation.

\begin{proposisi}
$\operatorname{rank}(T)$ equals the rank of the matrix $[T]_{C\leftarrow B}$, and $\ker T$
corresponds to the null space of that matrix. The rank of a matrix here means the dimension of its column
space, as already used in Chapter 6.
\end{proposisi}

\begin{proof}
Column $j$ of the matrix is the coordinate vector of $T(\mathbf{b}_{j})$, and
$\{T(\mathbf{b}_{j})\}$ spans $\operatorname{im}T$. Taking coordinates preserves
linear combinations, so the dimension of the range equals the dimension of the column space. The previous theorem
shows that $T(\mathbf{v})=\mathbf{0}$ exactly when $[T(\mathbf{v})]_{C}=[T]_{C\leftarrow B}[\mathbf{v}]_{B}=\mathbf{0}$,
that is, when $[\mathbf{v}]_{B}$ lies in the null space of that matrix.
\end{proof}

Composition of maps corresponds to matrix multiplication. That is no coincidence; it is
the reason matrix multiplication is defined the way it is.

\begin{teorema}
Let $T\colon V\to W$ and $S\colon W\to U$ be linear, with bases $B,C,E$ on the three spaces.
Then $[S\circ T]_{E\leftarrow B}=[S]_{E\leftarrow C}\,[T]_{C\leftarrow B}$.
\end{teorema}

\begin{proof}
For every $\mathbf{v}\in V$,
\[
  [(S\circ T)(\mathbf{v})]_{E}
  =[S]_{E\leftarrow C}[T(\mathbf{v})]_{C}
  =[S]_{E\leftarrow C}[T]_{C\leftarrow B}[\mathbf{v}]_{B}.
\]
So the matrix $[S]_{E\leftarrow C}[T]_{C\leftarrow B}$ satisfies the definition of the representation matrix
for $S\circ T$. The representation matrix is unique: if two matrices give the same result on
all $\mathbf{v}$, their results agree in particular at $\mathbf{v}=\mathbf{b}_{j}$, so
their columns agree.
\end{proof}

Composition also provides a test for invertibility through matrices.

\begin{proposisi}
Let $B$ and $C$ be bases of a space $V$ of dimension $n$, and let $T\colon V\to V$ be linear. Then
$T$ is invertible if and only if $[T]_{C\leftarrow B}$ is invertible, and in that case
$[T^{-1}]_{B\leftarrow C}=([T]_{C\leftarrow B})^{-1}$.
\end{proposisi}

\begin{proof}
If $T$ is invertible, then $T^{-1}\circ T=I$ and $T\circ T^{-1}=I$. The composition theorem gives
$[T^{-1}]_{B\leftarrow C}[T]_{C\leftarrow B}=I$ and
$[T]_{C\leftarrow B}[T^{-1}]_{B\leftarrow C}=I$, so the two matrices are mutually inverse.

Conversely, suppose $A=[T]_{C\leftarrow B}$ is invertible. If $T(\mathbf{v})=\mathbf{0}$, then
$A[\mathbf{v}]_{B}=\mathbf{0}$, so $[\mathbf{v}]_{B}=\mathbf{0}$ and
$\mathbf{v}=\mathbf{0}$.
Hence $T$ is injective. Because the source and target spaces both have dimension $n$, the
rank-nullity corollary makes $T$ surjective, so $T$ is invertible. The equation for the matrix of
$T^{-1}$
then follows from the first part.
\end{proof}

With that, a question about the inverse of an operator becomes a question about the inverse of a
matrix; the complete answer is in Chapter 8.

The special case used most often is an operator on a single space, $T\colon V\to V$,
with two different bases. Here the two matrices are linked by the transition matrix.

\begin{teorema}[Change of basis]
Let $T\colon V\to V$ be linear, $B$ and $B'$ two bases of $V$, and $P=P_{B\leftarrow B'}$.
Then
\[
  [T]_{B'}=P^{-1}[T]_{B}P .
\]
\end{teorema}

\begin{proof}
Write $T=I\circ T\circ I$ with $I$ the identity map. By the composition theorem,
\[
  [T]_{B'\leftarrow B'}=[I]_{B'\leftarrow B}\,[T]_{B\leftarrow B}\,[I]_{B\leftarrow B'}.
\]
The identity matrix with respect to two bases is the transition matrix, so
$[I]_{B'\leftarrow B}=P^{-1}$ and $[I]_{B\leftarrow B'}=P$.
\end{proof}\index{representation matrix}

Two square matrices linked in that way describe the same operator from
different viewpoints, so they are given a name of their own.

\begin{definisi}[Similar matrices]
Matrices $A$ and $B$ of size $n\times n$ are called \emph{similar} if there is an invertible
matrix $P$ with $B=P^{-1}AP$.\index{similarity}
\end{definisi}

Similarity is therefore another name for the same map seen from different bases.

\begin{proposisi}
Two $n\times n$ matrices are similar if and only if they are the representation matrices of
one and the same linear operator with respect to two bases.
\end{proposisi}

\begin{proof}
If $B=P^{-1}AP$, take the standard basis $B_{0}$ and the basis $B'$ whose columns are
the columns of $P$. Then $P=P_{B_{0}\leftarrow B'}$, and the change of basis theorem gives
$B=[T]_{B'}$ for the operator $T(\mathbf{x})=A\mathbf{x}$. Conversely, the same theorem
gives
that relation for every pair of bases.
\end{proof}

Rank is among the quantities that do not change: the rank of $[T]_{B'}$ equals the rank of $T$,
and likewise the rank of $[T]_{B}$. Determinant and trace are also invariant, but both are only
defined in Chapter 8 and discussed in Chapter 9.

A small example for the derivative operator: $D\colon P_{2}\to P_{2}$, $D(p)=p'$, with basis
$B=(1,x,x^{2})$. Because $D(1)=0$, $D(x)=1$, and $D(x^{2})=2x$,
\[
  [D]_{B}=\begin{pmatrix}0&1&0\\ 0&0&2\\ 0&0&0\end{pmatrix}.
\]
For $p=x^{2}+3x$ the coordinates are $(0,3,1)^{\top}$, and the matrix gives $(3,\,2,\,0)^{\top}$,
namely $3+2x$, exactly the derivative of $p$.

\section{Isomorphisms, dimensional invariance, and spaces of linear maps}
\tingkat{advanced}

Two different spaces can behave the same way. $P_{2}$ and $\mathbb{R}^{3}$ contain objects
of different kinds, but both have three independent directions and addition rules of the same
form.

A correspondence that preserves the whole structure is called an isomorphism: bijective,
linear, and therefore preserving bases, dimension, and kernels.

\kamus{isomorphism}{isomorfisme}{two floor plans of a house whose drawings differ but whose
number of rooms and connections between rooms are the same}

\begin{definisi}[Isomorphism]
A linear map $T\colon V\to W$ is called an \emph{isomorphism} if $T$ is bijective. Two spaces are called
\emph{isomorphic} if there is an isomorphism between them, written $V\cong W$.
\end{definisi}

One lemma before the main theorem: an isomorphism carries bases to bases.

\begin{lemma}
If $T\colon V\to W$ is an isomorphism and $B$ is a basis of $V$, then $T(B)$ is a basis of $W$.
\end{lemma}

\begin{proof}
Let $B=(\mathbf{b}_{1},\dots,\mathbf{b}_{n})$. For linear independence, suppose
$\sum_{j}c_{j}T(\mathbf{b}_{j})=\mathbf{0}$. Because $T$ is linear, the vector
$\sum_{j}c_{j}\mathbf{b}_{j}$ is mapped to $\mathbf{0}$, so that vector lies in the kernel. The kernel
is only $\{\mathbf{0}\}$ because $T$ is injective, so all $c_{j}=0$.
For spanning, take $\mathbf{w}\in W$. There is $\mathbf{v}\in V$ with $T(\mathbf{v})=\mathbf{w}$
because $T$ is surjective, and $\mathbf{v}=\sum_{j}a_{j}\mathbf{b}_{j}$ because $B$ spans.
Then
$\mathbf{w}=\sum_{j}a_{j}T(\mathbf{b}_{j})$.
\end{proof}

With that lemma, equality of dimension turns out to be the only obstacle.

\begin{teorema}
Two finite-dimensional spaces over $\mathbb{R}$ are isomorphic if and only if their dimensions are equal.
\end{teorema}

\begin{proof}
If $V\cong W$ via $T$, the previous lemma turns a basis of $V$ into a basis of $W$ with
the same number of members, so $\dim V=\dim W$.

Conversely, suppose $\dim V=\dim W=n$. Choose a basis $B$ of $V$ and $C$ of $W$. The coordinate
maps $\varphi_{B}\colon V\to\mathbb{R}^{n}$ and $\varphi_{C}\colon W\to\mathbb{R}^{n}$ are
bijective and linear. Then $\varphi_{C}^{-1}\circ\varphi_{B}$ is a bijective linear map
from
$V$ to $W$, because the composition of bijective linear maps is also bijective and linear.
\end{proof}

That conclusion is felt at once in the spaces we know.

\begin{akibat}
Every space of dimension $n$ is isomorphic to $\mathbb{R}^{n}$. In particular
$P_{n-1}\cong\mathbb{R}^{n}$ and $M_{2\times2}(\mathbb{R})\cong\mathbb{R}^{4}$.
\end{akibat}

An isomorphism preserves everything defined from the operations: linear independence, span,
kernel, range, rank, and nullity. Two isomorphic spaces cannot be distinguished
by the tools of this chapter; dimension becomes the only number that needs computing in order to recognize
a finite-dimensional space.

There is a trap in interpretation. The isomorphism above depends on the choice of
basis; there is no canonical isomorphism between $V$ and $\mathbb{R}^{n}$. Hence ``$V$ equals
$\mathbb{R}^{n}$'' is false, while ``$V$ is isomorphic to $\mathbb{R}^{n}$'' is
true.
Polynomials of degree less than three and triples of real numbers are objects of different kinds,
but the two are isomorphic because every polynomial is determined by its three coefficients.

Now the collection of all linear maps between two spaces; that set is itself a vector space.

\begin{teorema}
The set $\mathcal{L}(V,W)$ of all linear maps $V\to W$, with the operations
$(S+T)(\mathbf{v})=S(\mathbf{v})+T(\mathbf{v})$ and $(cT)(\mathbf{v})=cT(\mathbf{v})$,
is
a vector space. If $\dim V=n$ and $\dim W=m$, then $\dim\mathcal{L}(V,W)=mn$. The mapping
$T\mapsto[T]_{C\leftarrow B}$ is an isomorphism $\mathcal{L}(V,W)\to M_{m\times n}(\mathbb{R})$.
\end{teorema}

\begin{proof}
$\mathcal{L}(V,W)$ sits inside the space of all functions $V\to W$, which is a vector space
by the example in Chapter 6. The zero map is linear. The sum of two linear maps is linear, and a scalar multiple of a
linear
map is also linear; so $\mathcal{L}(V,W)$ is a subspace, and therefore a vector space.

For the isomorphism claim, let $\Phi(T)=[T]_{C\leftarrow B}$. Linearity of $\Phi$: column
$j$ of $\Phi(S+cT)$ is $[S(\mathbf{b}_{j})+cT(\mathbf{b}_{j})]_{C}$, which equals
$[S(\mathbf{b}_{j})]_{C}+c[T(\mathbf{b}_{j})]_{C}$ because coordinates are linear. Hence
$\Phi(S+cT)=\Phi(S)+c\Phi(T)$.

Injectivity: if $\Phi(T)$ is the zero matrix, then all $[T(\mathbf{b}_{j})]_{C}=\mathbf{0}$, so
$T(\mathbf{b}_{j})=\mathbf{0}$ for all $j$, and $T=\mathbf{0}$ because a linear map is
determined by its values on a basis.

Surjectivity: given a matrix $A=(a_{ij})$ of size $m\times n$. Define
$T(\mathbf{b}_{j})=\sum_{i}a_{ij}\mathbf{c}_{i}$ and extend linearly; by the
proposition
on values on a basis, that extension exists and is unique. Column $j$ of $\Phi(T)$ is
$(a_{ij})_{i}$, so $\Phi(T)=A$.

Because $\Phi$ is an isomorphism and $\dim M_{m\times n}(\mathbb{R})=mn$, we get
$\dim\mathcal{L}(V,W)=mn$.
\end{proof}

That count gives an immediate corollary.

\begin{akibat}
$\dim\mathcal{L}(V,V)=n^{2}$ if $\dim V=n$. For example,
$\dim\mathcal{L}(P_{2},P_{2})=9$ and $\dim\mathcal{L}(P_{2},\mathbb{R}^{2})=6$.
\end{akibat}

That last result explains why a matrix is more than a computational device. The space of operators
and the space of matrices have the same dimension and are linked by an isomorphism that respects composition,
so every statement about operators has a counterpart about matrices. What is not
inherited
is simplicity of form: the same operator can be complicated in one basis and simple
in
another; the search for the best basis is the program of Chapter 9.

\begin{tantangan}
Find one linear map $T\colon\mathbb{R}^{3}\to\mathbb{R}^{3}$ that maps the unit square
in the $xy$ plane to another unit square of area still $1$, yet
$\operatorname{rank}(T)=2$. Explain why an area preserved in one plane does not
force
full rank.

Hint 1: examine one plane, not the whole space. After that, ask where the third axis
goes.

Hint 2: place a $90^{\circ}$ rotation on the first two coordinates and send the third coordinate to zero,
so that $T(\mathbf{e}_{3})=\mathbf{0}$. Compute the rank from the images of $\mathbf{e}_{1}$ and
$\mathbf{e}_{2}$.

If you get stuck, return to the section Linear transformations, kernel, and range, especially the definition of
rank and the rank-nullity corollary.
\end{tantangan}

\begin{untukanak}
Two classes at school can differ in size, yet have the same number of chairs. If the
number is
the same, every pupil in the first class can be paired with one chair in the second, with no
empty
chair and no pupil getting two chairs. As far as counting chairs goes, the two classes are
equivalent.

What is not equivalent is the pupils' names and faces. The chair pairing does not say
who sits where; it only pairs. Different people can make different pairings,
and all of them are legitimate.

In this chapter we will learn when two spaces that look different may be regarded as the same for
computational purposes, and when that assumption is misleading. The equivalence always depends
on
our choice, not on something already present in the spaces.

If two classes have the same number of chairs, may the headmaster simply swap the class names
just
like that? The answer is no, and the reason is in this chapter.
\end{untukanak}

\section{A worked example}
\tingkat{intermediate}

Consider the matrix
\[
  A=\begin{pmatrix}1&2&1\\ 2&5&3\\ 3&7&4\end{pmatrix},
  \qquad \text{third row}=\text{first row}+\text{second row}.
\]
We determine its rank, its nullity, and both bases, then examine the same operator in another
basis.\index{rank}\index{nullity}

Row elimination:
\[
  \begin{pmatrix}1&2&1\\ 2&5&3\\ 3&7&4\end{pmatrix}
  \xrightarrow[R_{3}\leftarrow R_{3}-3R_{1}]{R_{2}\leftarrow R_{2}-2R_{1}}
  \begin{pmatrix}1&2&1\\ 0&1&1\\ 0&1&1\end{pmatrix}
  \xrightarrow{R_{3}\leftarrow R_{3}-R_{2}}
  \begin{pmatrix}1&2&1\\ 0&1&1\\ 0&0&0\end{pmatrix}.
\]
Then $R_{1}\leftarrow R_{1}-2R_{2}$ gives the reduced echelon form
\[
  \begin{pmatrix}1&0&-1\\ 0&1&1\\ 0&0&0\end{pmatrix}.
\]
There are two nonzero rows, so $\operatorname{rank}A=2$ and $\det A=0$; the pivot columns are the
first and the second.

Kernel: from $x_{1}-x_{3}=0$ and $x_{2}+x_{3}=0$ we obtain $x_{1}=x_{3}$ and
$x_{2}=-x_{3}$.
Hence
\[
  \ker A=\operatorname{span}\{(1,-1,1)^{\top}\},\qquad \operatorname{nullity}(A)=1 .
\]
Rank-nullity is satisfied: $3=2+1$.

Range: the columns of $A$ are $\mathbf{c}_{1}=(1,2,3)^{\top}$,
$\mathbf{c}_{2}=(2,5,7)^{\top}$, and $\mathbf{c}_{3}=(1,3,4)^{\top}$. The first and the
second
pivot columns give the basis $\{\mathbf{c}_{1},\mathbf{c}_{2}\}$. The third column is not independent:
$\mathbf{c}_{2}-\mathbf{c}_{1}=\mathbf{c}_{3}$, and that relation is exactly the kernel vector
$(1,-1,1)$: $\mathbf{c}_{1}-\mathbf{c}_{2}+\mathbf{c}_{3}=\mathbf{0}$.

Now consider $T(\mathbf{x})=A\mathbf{x}$ and the basis
$B'=(\mathbf{b}'_{1},\mathbf{b}'_{2},\mathbf{b}'_{3})$ with
$\mathbf{b}'_{1}=(1,1,0)$, $\mathbf{b}'_{2}=(1,0,1)$, $\mathbf{b}'_{3}=(0,1,1)$. Its
transition matrix from the standard basis $B$ is
\[
  P=P_{B\leftarrow B'}=\begin{pmatrix}1&1&0\\ 1&0&1\\ 0&1&1\end{pmatrix},
  \qquad
  P^{-1}=\frac{1}{2}\begin{pmatrix}1&1&-1\\ 1&-1&1\\ -1&1&1\end{pmatrix}.
\]
The images of the new basis vectors are computed first:
\[
  A\mathbf{b}'_{1}=\begin{pmatrix}3\\7\\10\end{pmatrix},\quad
  A\mathbf{b}'_{2}=\begin{pmatrix}2\\5\\7\end{pmatrix},\quad
  A\mathbf{b}'_{3}=\begin{pmatrix}3\\8\\11\end{pmatrix}.
\]
Hence
\[
  AP=\begin{pmatrix}3&2&3\\ 7&5&8\\ 10&7&11\end{pmatrix},
  \qquad
  [T]_{B'}=P^{-1}AP=\begin{pmatrix}0&0&0\\ 3&2&3\\ 7&5&8\end{pmatrix}.
\]
Check one entry: first row, first column $\tfrac{1}{2}(3+7-10)=0$; second row,
third column $\tfrac{1}{2}(3-8+11)=3$. The rest follows the same pattern.

A first row that is entirely zero is no accident: all three images lie inside
$\operatorname{span}(\mathbf{b}'_{2},\mathbf{b}'_{3})$, for instance
$A\mathbf{b}'_{3}=(3,8,11)=3\mathbf{b}'_{2}+8\mathbf{b}'_{3}$. Hence the first coordinate of
every image is zero, and the range of $T$ is contained in the coordinate plane of the first
coordinate.

The pattern is general. A zero row signals a range contained in a coordinate subspace
of lower dimension; a zero column signals a basis vector that goes into the kernel. Both symptoms
are visible only in the chosen basis, whereas rank and nullity do not change.

A complete check:
\[
  \operatorname{rank}[T]_{B'}=2,\qquad \operatorname{trace}[T]_{B'}=0+2+8=10=\operatorname{trace}A,
  \qquad \det[T]_{B'}=0=\det A .
\]
The equality of traces is not an arithmetic accident. The $j$-th diagonal entry is the coefficient of
$\mathbf{b}'_{j}$ in the image $T(\mathbf{b}'_{j})$, that is, how much $T$
preserves
that direction itself. The sum does not depend on the directions chosen, only on the operator,
so the trace serves as a cheap fingerprint: different traces mean the matrices are not similar, equal
traces
guarantee nothing.

For $\mathbf{v}=(3,8,11)$ we have $[\mathbf{v}]_{B'}=(0,3,8)^{\top}$ and
\[
  [T]_{B'}[\mathbf{v}]_{B'}=\begin{pmatrix}0\\30\\79\end{pmatrix},
\]
whereas $T(\mathbf{v})=A\mathbf{v}=(30,79,109)^{\top}$ and its coordinates are
$(0,30,79)^{\top}$. The two agree.

The feature map $G\colon\mathbb{R}^{4}\to\mathbb{R}^{3}$
is given by the matrix
\[
  G=\begin{pmatrix}1&2&0&1\\ 0&1&1&1\\ 1&3&1&2\end{pmatrix},
\]
whose third row is the sum of the first two rows. Elimination gives the reduced echelon
form
\[
  \begin{pmatrix}1&0&-2&-1\\ 0&1&1&1\\ 0&0&0&0\end{pmatrix},
\]
so $\operatorname{rank}G=2$ and $\operatorname{nullity}(G)=4-2=2$. From
$x_{1}=2x_{2}+x_{4}$ and $x_{3}=-x_{2}-x_{4}$:
\[
  \mathbf{k}_{1}=(-2,1,-1,0)^{\top},\qquad \mathbf{k}_{2}=(-1,0,-1,1)^{\top}.
\]
Both satisfy $G\mathbf{k}_{i}=\mathbf{0}$ and are linearly independent, so the changes that the
three features do not see form a two-dimensional plane, not a single line: an attacker can
shift the input in any direction of that plane without moving the score.

\kertas{Repeat the computation with the basis $B''=(\mathbf{b}'_{2},\mathbf{b}'_{3},\mathbf{b}'_{1})$,
that is, with $B'$ permuted. Predict how the zero row and the other entries move, then compute
$[T]_{B''}$ via $Q^{-1}[T]_{B'}Q$ with $Q$ the corresponding permutation matrix. Check that
rank
and trace do not change, and explain why $\det A=0$ can be concluded directly from ``the third
row is the sum of the first two rows''.}

All the numbers above were checked with exact arithmetic (SymPy); the numpy and
scipy checks are in the next section, including one number that is not exactly zero because of machine
rounding.

\section{Computational lab}
\tingkat{advanced}

\begin{lab}
\begin{verbatim}
import numpy as np
from scipy.linalg import null_space

A = np.array([[1., 2., 1.],
              [2., 5., 3.],
              [3., 7., 4.]])
print("rank A  :", np.linalg.matrix_rank(A))
print("det A   :", np.linalg.det(A))
K = null_space(A)
print("dim ker :", K.shape[1])
print("basis ker:\n", np.round(K, 4))
print("A @ K   :\n", np.round(A @ K, 12))

P = np.array([[1., 1., 0.],
              [1., 0., 1.],
              [0., 1., 1.]])
B = np.linalg.inv(P) @ A @ P
print("[T]_B'  :\n", np.round(B, 10))
print("difference :", np.linalg.norm(B - np.linalg.inv(P) @ A @ P))
print("trace   :", round(np.trace(B), 10), round(np.trace(A), 10))
print("rank B  :", np.linalg.matrix_rank(B))
print("cond P  :", np.linalg.cond(P))

v = np.array([3., 8., 11.])
print("[v]_B'  :", np.round(np.linalg.solve(P, v), 4))
print("[Av]_B' :", np.round(np.linalg.solve(P, A @ v), 4))
print("B[v]_B' :", np.round(B @ np.linalg.solve(P, v), 4))

x0 = np.array([1., 0., 0.])
d = 0.5 * K[:, 0]
print("A@d     :", np.round(A @ d, 12))
print("shift   :", np.linalg.norm(A @ (x0 + d) - A @ x0))
\end{verbatim}
Actual output:
\begin{verbatim}
rank A  : 2
det A   : 6.66133814775094e-16
dim ker : 1
basis ker:
 [[ 0.5774]
 [-0.5774]
 [ 0.5774]]
A @ K   :
 [[-0.]
 [ 0.]
 [ 0.]]
[T]_B'  :
 [[0. 0. 0.]
 [3. 2. 3.]
 [7. 5. 8.]]
difference : 0.0
trace   : 10.0 10.0
rank B  : 2
cond P  : 2.0000000000000004
[v]_B'  : [0. 3. 8.]
[Av]_B' : [ 0. 30. 79.]
B[v]_B' : [ 0. 30. 79.]
A@d     : [-0.  0.  0.]
shift   : 6.280369834735101e-16
\end{verbatim}
\end{lab}

Four things are worth reading from that output.

First, $\det A$ is not zero but $6.66\times10^{-16}$. This matrix is singular
exactly, but binary fraction arithmetic knows no exact zero.
\texttt{matrix\_rank} compares the largest and smallest singular values, not a simple threshold,
so the rank is still reported as $2$; Chapter 22 discusses why such a threshold is necessary.

The choice of singular values rather than the determinant: the determinant measures volume, and a volume can be
very
small without the rank changing. What determines the rank is the ratio of the smallest singular value
to the
largest.

Second, \texttt{null\_space} returns a unit vector with entries $0.5774$,
$-0.5774$, and $0.5774$, that is $(1,-1,1)^{\top}$ divided by $\sqrt{3}$, and
$1/\sqrt{3}=0.57735$ agrees. The direction is the same as in the hand computation; its length is normalized
because
scipy chooses an orthonormal basis. A kernel basis is not unique, its dimension is.

Third, the $[T]_{B'}$ reported by numpy matches the hand computation exactly, and the traces
of both matrices agree, $10$. The difference of $P^{-1}AP$ computed twice is $0.0$ because
$P$
contains small integers. The condition number of $P$ is $2.0000000000000004$, quite
small;
a bad basis inflates the error, and that is discussed in Chapter 10.

Fourth, the kernel direction is genuinely blind to $T$. Pushing $\mathbf{x}_{0}=(1,0,0)$ by
half the kernel vector changes $\mathbf{x}_{0}$ visibly, yet $A\mathbf{x}$ shifts only
$6.28\times10^{-16}$, that is, zero at machine precision.

A final warning concerns the choice of basis. A nearly dependent basis makes coordinates
explode, and the reported numbers lose their meaning.

\begin{lab}
\begin{verbatim}
Pb = np.array([[1., 1., 0.],
               [1., 1.001, 0.],
               [0., 0., 1.]])
print("cond Pb:", np.linalg.cond(Pb))
print("coord   :", np.linalg.inv(Pb) @ v)
print("error   :", np.linalg.norm(np.linalg.inv(Pb) @ (Pb @ v) - v))
\end{verbatim}
Actual output:
\begin{verbatim}
cond Pb: 4002.00075012484
coord   : [-4997.  5000.    11.]
error   : 4.823324033076557e-13
\end{verbatim}
\end{lab}

The first two basis vectors differ only in the third digit, and the condition number of $P_{b}$
reaches $4002$. The coordinates $(3,8,11)$ explode into $(-4997,\,5000,\,11)$. The reconstruction
error
is still $4.8\times10^{-13}$, but the relative error in the coordinates grows by a factor of
the
condition number; with 16 digits of precision, about 12 remain. In a worse basis that remainder
is
used up and the reported rank cannot be trusted.

\begin{untukdosen}
A teaching order that works in practice: coordinates first, linear maps second. Students who
know
coordinates before the abstract definition are readier to accept the representation matrix, because
the matrix appears as a list of coordinates, not as a new object. After rank-nullity, move to
change of basis.

Two typical misconceptions. First, equating a vector with its coordinate list, then getting confused
when two bases give two lists for the same vector. Second, the direction of the transition matrix
backwards. A quick check: the columns must contain the addresses of the basis vectors, and the subscript is read
right to left.

Blackboard activity: give one operator on $\mathbb{R}^{2}$ and two bases. Ask half the
class to compute $[T]_{B'}$ via the column definition, half via $P^{-1}AP$, then compare
on
the board; the noise that arises is usually exactly about the direction of $P$.

A sign of readiness for the next step: students can explain why a zero row appears in the new
matrix, without recomputing. If they connect it with a range
contained in a coordinate subspace, they are ready for Chapter 9.
\end{untukdosen}

\section{Summary}
\tingkat{basic}

\begin{ringkasan}
\begin{itemize}
  \item The coordinates $[\mathbf{v}]_{B}$ are the unique list of coefficients with respect to an ordered basis;
        the map $\mathbf{v}\mapsto[\mathbf{v}]_{B}$ is linear and bijective, but depends on the
        basis.
  \item The transition matrix contains the coordinates of the old basis in the new basis:
        $[\mathbf{v}]_{C}=P_{C\leftarrow B}[\mathbf{v}]_{B}$, with
        $P_{B\leftarrow C}=P_{C\leftarrow B}^{-1}$.
  \item A linear map is determined by its values on a basis; $T$ is injective exactly when
        $\ker T=\{\mathbf{0}\}$.
  \item Rank-nullity: $\dim V=\dim\ker T+\dim\operatorname{im}T$; in equal dimensions,
        injective is equivalent to surjective.
  \item The representation matrix satisfies
$[T(\mathbf{v})]_{C}=[T]_{C\leftarrow B}[\mathbf{v}]_{B}$, composition becomes matrix
        multiplication, and $[T]_{B'}=P^{-1}[T]_{B}P$.
  \item $\dim\mathcal{L}(V,W)=mn$ if $\dim V=n$, $\dim W=m$; two finite-dimensional spaces
        are isomorphic exactly when their dimensions are equal.
\end{itemize}
\end{ringkasan}

\section{Exercises}
\tingkat{intermediate}

\begin{latihan}
\begin{enumerate}
  \item With $B'=(\mathbf{b}'_{1},\mathbf{b}'_{2},\mathbf{b}'_{3})$ as in the worked
example, determine $[\mathbf{w}]_{B'}$ for $\mathbf{w}=(2,-1,3)$, then check the result
        by multiplying back.

\item Prove that every isomorphism $V\to\mathbb{R}^{n}$ has the form $\varphi_{B}$ for
        some basis $B$, and explain why an isomorphism to $\mathbb{R}^{n}$ is not unique.
        (Hint: take the preimages of the standard basis, then show that this set is a basis.)

  \item For $A=\begin{pmatrix}0&1&0\\ 0&0&2\\ 0&0&0\end{pmatrix}$, determine the rank, the nullity,
        a basis of the kernel, and a basis of the range. Explain the connection with the derivative operator on
        $P_{2}$.

\item Prove $\operatorname{rank}(S\circ T)\le\min(\operatorname{rank}S,\operatorname{rank}T)$
        for linear maps that are defined. Use
        $\operatorname{im}(S\circ T)\subseteq\operatorname{im}S$ and compare the kernels.
        (Hint: for the second bound, compare $\ker T$ with $\ker(S\circ T)$ via
        rank-nullity.)

\item Assess this claim: \emph{two square matrices of the same size that have the same rank
are necessarily similar.} Decide whether it is true, false, or not supported by enough evidence, with reasons and a counterexample if needed.

\item Let $T\colon\mathbb{R}^{3}\to\mathbb{R}^{2}$ with
        $T(x,y,z)=(x+2y,\;y+z)$. Compute $[T]_{C\leftarrow B}$ for
$B=(\,(1,0,0),(1,1,0),(1,1,1)\,)$ and $C=(\,(1,1),(1,0)\,)$, then check the formula at
        $\mathbf{v}=(1,2,3)$.
        (Hint: $[T]_{C\leftarrow B}=C^{-1}AB$; compute $AB$ first, then solve
        $C\mathbf{x}=\mathbf{y}$ for each of its columns.)

  \item Prove: if $U$ is a subspace of a finite-dimensional $V$ and $\dim U=\dim V$, then $U=V$.

  \item Prove that if $A$ and $B$ are similar then $\operatorname{nullity}(A)
        =\operatorname{nullity}(B)$. Explain why that statement is equivalent to
        ``the kernel does not depend on the basis''.

  \item The feature map $F\colon\mathbb{R}^{4}\to\mathbb{R}^{3}$ is given by the matrix
        \[
          F=\begin{pmatrix}1&0&2&1\\ 0&1&1&3\\ 1&1&3&4\end{pmatrix}.
        \]
        Determine the rank and nullity of $F$, give one nonzero kernel vector, and explain why
        that vector is a direction of change that does not shift a single feature.
\end{enumerate}
\end{latihan}

\begin{untukpeneliti}
What is proved in this chapter are statements about finite dimension over the field $\mathbb{R}$:
rank-nullity, uniqueness of the representation matrix, the equivalence of isomorphism with equality of
dimension,
and $\dim\mathcal{L}(V,W)=mn$. All of them are theorems, not conventions. The claim ``a matrix is its transformation'' is merely popular; without a choice of basis that claim is meaningless, and the zero-row
example
above shows that the form of a matrix depends on the basis.

Primary sources: Axler, \emph{Linear Algebra Done Right} 4th edition, sections 3A to 3D,
contain
spaces of linear maps, kernel and range, matrices, isomorphisms, and change of basis;
rank-nullity
there is called the fundamental theorem for linear maps. Hoffman and Kunze, \emph{Linear
Algebra}
2nd edition, Chapter 3, is denser but equivalent.

An honest open question. The classification of operators up to similarity is not given
here; Chapter 9 stops at diagonalization. Numerically, the rank decision depends on a
threshold and
condition numbers; this chapter uses library output without discussing its stability, a debt paid
in Chapter 22. Limit of validity: in infinite-dimensional spaces the whole argument collapses,
rank-nullity
fails, and isomorphism is no longer equivalent to equality of dimension.
\end{untukpeneliti}

\peta{This chapter follows Anton, \emph{Elementary Linear Algebra} 11th edition, sections 4.4, 4.6,
4.8, and 8.1 to 8.5, with the abstract counterpart in Axler, \emph{Linear Algebra Done Right}
4th
edition, Chapter 3, and Hoffman and Kunze, \emph{Linear Algebra} 2nd edition, Chapter 3.}

\chapter{Linear Systems, Determinants, and Inverses}

A threat scorer adds three features with unknown weights,
$s=w_{1}f_{1}+w_{2}f_{2}+w_{3}f_{3}$. Calibration on three labeled examples gives three
linear equations at once:
\[
  2w_{1}+w_{2}-w_{3}=8,\qquad
  -3w_{1}-w_{2}+2w_{3}=-11,\qquad
  -2w_{1}+w_{2}+2w_{3}=-3 .
\]
The answer is $w=(2,3,-1)$, and nothing about it is surprising.

Then the coefficient matrix is replaced by the result of another measurement,
\[
  S=\begin{pmatrix}1&2&3\\ 4&5&6\\ 7&8&9.001\end{pmatrix}.
\]
The second row minus the first gives $(3,3,3)$, while the third row minus the second gives
$(3,3,3.001)$: the three rows are nearly dependent. The weights $w=(1,1,1)$ produce the
right-hand side $(6,\,15,\,24.001)$, so the system is consistent and the answer is exactly
$w=(1,1,1)$. Raise the third measurement by $0.001$, and the answer moves to
$(2,-1,2)$. A relative change in the data of $3.46\times10^{-5}$ produces a relative change
in the answer of $1.41$.

There are four questions here, not one: does the system have an answer, how many answers does
it have, which number decides that without solving anything, and how far does the answer shift
when the data shift. Chapter 7 already supplied rank and nullity for the first. Three are still
missing: a procedure that is written down and can be checked by someone else, one number that
decides invertibility and at the same time yields a formula for the answer, and a measure of
sensitivity; those three are the content of this chapter.

The determinant is usually introduced as the cofactor formula, and then its properties are
proved one by one. We reverse that: definition first, through three properties, uniqueness
proved, the cofactor formula appearing as a consequence. The reason is not style: the cofactor
formula hides the fact that the determinant is the only normalized alternating function, and
that fact is what makes the proofs in this chapter short.

\begin{untukanak}
Three friends weigh three boxes on three different scales. Each friend writes down three
numbers, and from those numbers we want to know the weight of each box. As long as the three
scales measure different things, the answer is easy: every new number supplies one new piece
of information.

Here is the trouble. The second scale turns out to nearly copy the first one. The numbers are
not exactly the same, only slightly different, perhaps because of dust on the pan. The third
scale also nearly copies the other two. Now the three pieces of information have nearly the
same content. We feel we hold three pieces of information, when only one or two of them are
really new.

This is where the strangeness begins. If one number changes by a thousandth, the box weights
we compute can change by almost half. Not because the arithmetic is wrong, but because the
information really is nearly identical. There is one number that can tell us in advance
whether this will happen, before we compute anything at all.

Have you ever asked the same question of three people, received three answers with similar
content, and still felt unsure? In this chapter we will learn why similar answers actually
weaken the information, and how to recognize that from the numbers.
\end{untukanak}

\section{Gaussian elimination and three fates of a system}
\tingkat{basic}

A linear system is written as a single unit, not as a list of separate equations. The way to
do it is to attach the right-hand side column to the coefficient matrix. The solution is then
carried out on that arrangement, and every step replaces one system by another system with the
same solution set.

\kamus{elimination}{eliminasi}{Replacing a question by a simpler question without changing the
answer, like simplifying a fraction without changing its value.}

\begin{definisi}[Augmented matrix and elementary row operations]
For the system $A\mathbf{x}=\mathbf{b}$ with $A$ of size $m\times n$, the \emph{augmented
matrix} is the matrix $\bigl(A\mid\mathbf{b}\bigr)$ of size $m\times(n+1)$. The three
\emph{elementary row operations} are: interchanging two rows; multiplying one row by a
nonzero scalar; adding a multiple of one row to another row. Two matrices are called
\emph{row equivalent} if one can be reached from the other by those operations.\index{elimination}
\end{definisi}

Those operations are chosen because each of them can be undone. Interchanging two rows undoes
itself, multiplying by $c\neq0$ is undone by multiplying by $1/c$, and adding $c$ times row $i$
to row $j$ is undone by subtracting $c$ times row $i$. No information is therefore lost.

\begin{teorema}
If two linear systems have row equivalent augmented matrices, the two systems have the same
solution set.
\end{teorema}

\begin{proof}
It is enough to check one operation, because repeated operations follow. Let $S$ be the
original system and $S'$ the system after one operation. Every equation of $S'$ is a linear
combination of the equations of $S$, so every answer of $S$ also answers $S'$. Conversely,
because the operation can be undone by an operation of the same kind, every equation of $S$ is
a linear combination of the equations of $S'$, so every answer of $S'$ also answers $S$. The two
solution sets therefore contain each other.
\end{proof}

The form we aim at is the staircase form. Every nonzero row starts further to the right than
the nonzero row above it, and all entries below that first position are zero. The column where
the first number of a nonzero row sits is called the pivot column, and that number itself is
called a pivot.

\begin{definisi}[Row echelon form]
A matrix is in \emph{row echelon form} if all zero rows lie at the bottom, and the first
nonzero entry of every nonzero row is located further right than the first nonzero entry of the
row above it. The form is called \emph{reduced} if every pivot equals $1$ and is the only
nonzero entry in its column.\index{echelon}\index{pivot}
\end{definisi}

The reduced form is unique. This fact is used every time we speak of ``the rank computed by
elimination'', and its proof is not as short as it looks.

\begin{teorema}[Uniqueness of the reduced echelon form]
Every matrix is row equivalent to exactly one matrix in reduced row echelon form.
\end{teorema}

\begin{proof}
The proof is skipped here. The statement is proved in full in Hoffman and Kunze,
\emph{Linear Algebra} 2nd edition, section 1.4, by showing that the solution space of
$A\mathbf{x}=\mathbf{0}$ determines the position of the pivots uniquely. What we use in this
chapter is only the consequence: the number of pivots and their positions do not depend on the
order of operations chosen.
\end{proof}

With the reduced form, every question about the number of answers can be settled by counting
pivots. These are the three fates of a system.

\begin{teorema}[Three cases]
Suppose $\bigl(A\mid\mathbf{b}\bigr)$ is row equivalent to the reduced echelon form
$\bigl(R\mid\mathbf{c}\bigr)$, with $A$ of size $m\times n$.
\begin{enumerate}
  \item The system is \emph{inconsistent} exactly when there is a pivot in column $(n+1)$.
  \item If it is consistent, the system has a \emph{unique answer} exactly when every column
        of $A$ contains a pivot.
  \item If it is consistent, the system has \emph{many answers} exactly when some column of $A$
        has no pivot, and the number of free parameters equals the number of columns without a
        pivot.
\end{enumerate}
\end{teorema}

\begin{proof}
A pivot in the last column means there is a row of the form $0=\text{nonzero constant}$, which
cannot possibly be satisfied. If there is no such pivot, every nonzero row gives one equation
expressing a pivot variable in terms of the other variables, so the system is consistent.

If every column contains a pivot, the number of pivots equals $n$, so there is no free variable
and every pivot variable is uniquely determined. If some column has no pivot, the variable in
that column is free: it may take any value, and every choice determines the pivot variables
uniquely. The answer is therefore infinite in number, and the number of free parameters is
exactly the number of columns without a pivot.
\end{proof}

Figure~\ref{fig:tiga-nasib} places those three fates in a single plane. The line $x+y=2$ does
not change; only the second line changes. Left panel: $x-y=0$ cuts it at $(1,1)$, two pivots,
a unique answer. Middle panel: $x+y=0$ is parallel, and the last column of the augmented matrix
contains a pivot. Right panel: $2x+2y=4$ is the same line written with different numbers, one
pivot and one free parameter.

\begin{figure}[htbp]
\centering
\begin{tikzpicture}[scale=0.66,>=Stealth]
  \begin{scope}
    \draw[layergrey!28,step=1] (-2.2,-2.2) grid (2.2,2.2);
    \draw[->,ink] (-2.3,0)--(2.4,0) node[right,font=\scriptsize]{$x$};
    \draw[->,ink] (0,-2.3)--(0,2.4) node[above,font=\scriptsize]{$y$};
    \draw[very thick,layerindigo] (-1.2,3.2)--(3.2,-1.2);
    \draw[very thick,layerteal] (-1.6,-1.6)--(2.2,2.2);
    \fill[accent] (1,1) circle (2.4pt);
    \node[font=\scriptsize,below right] at (1,1) {$(1,1)$};
    \node[font=\scriptsize] at (0,-2.9) {(a) intersecting};
  \end{scope}
  \begin{scope}[xshift=6.4cm]
    \draw[layergrey!28,step=1] (-2.2,-2.2) grid (2.2,2.2);
    \draw[->,ink] (-2.3,0)--(2.4,0) node[right,font=\scriptsize]{$x$};
    \draw[->,ink] (0,-2.3)--(0,2.4) node[above,font=\scriptsize]{$y$};
    \draw[very thick,layerindigo] (-1.2,3.2)--(3.2,-1.2);
    \draw[very thick,accent,dashed] (-2.2,2.2)--(2.2,-2.2);
    \node[font=\scriptsize] at (0,-2.9) {(b) parallel};
  \end{scope}
  \begin{scope}[xshift=12.8cm]
    \draw[layergrey!28,step=1] (-2.2,-2.2) grid (2.2,2.2);
    \draw[->,ink] (-2.3,0)--(2.4,0) node[right,font=\scriptsize]{$x$};
    \draw[->,ink] (0,-2.3)--(0,2.4) node[above,font=\scriptsize]{$y$};
    \draw[very thick,layerindigo] (-1.2,3.2)--(3.2,-1.2);
    \draw[line width=2.6pt,accent,densely dotted] (-1.2,3.2)--(3.2,-1.2);
    \node[font=\scriptsize] at (0,-2.9) {(c) coincident};
  \end{scope}
\end{tikzpicture}
\caption{Three fates of a $2\times2$ system with the first line $x+y=2$ held fixed. Left: the
second line $x-y=0$ cuts it at $(1,1)$, two pivots, a unique answer. Middle: the second line
$x+y=0$ is parallel, the right-hand side column contains a pivot, there is no answer. Right: the
second line $2x+2y=4$ coincides with the first line, one pivot, infinitely many answers.}
\label{fig:tiga-nasib}
\end{figure}

\begin{bacakeras}
Two voices. The first voice names the situation, the second voice names the consequence. Read
them alternately, then swap roles.

One system. Two lines. Three fates.

Intersecting: one point, one answer.

Parallel: they never meet, there is no answer.

Coincident: one line disguised as two, infinitely many answers.

Who decides, before we compute anything? The pivots decide.

Two pivots, two columns: a unique answer.

A pivot in the last column: some row reads zero equals nonzero.

A column without a pivot: one free parameter, and the answers line up in one direction.

Repeat the last line together: count the pivots first, then talk about answers.
\end{bacakeras}

The matrix from Chapter 7 displays two fates at once,
\[
  A=\begin{pmatrix}1&2&1\\ 2&5&3\\ 3&7&4\end{pmatrix},
  \qquad \text{third row}=\text{first row}+\text{second row}.
\]
For $\mathbf{b}=(4,9,13)$ elimination gives
\[
  \left(\begin{array}{ccc|c}1&2&1&4\\ 2&5&3&9\\ 3&7&4&13\end{array}\right)
  \longrightarrow
  \left(\begin{array}{ccc|c}1&0&-1&2\\ 0&1&1&1\\ 0&0&0&0\end{array}\right).
\]
There are two pivots and three columns, so there is one free variable. The answer is
$(2,1,0)+t\,(1,-1,1)$ for $t\in\mathbb{R}$, and the check is immediate:
$A(2,1,0)^{\top}=(4,9,13)^{\top}$ while $A(1,-1,1)^{\top}=\mathbf{0}$, so the direction of that
shift really is the blind direction of the matrix from Chapter 7.

For $\mathbf{b}=(4,9,14)$ the last row becomes $0=1$, and the system is inconsistent. The
coefficient matrix is exactly the same; one number on the right-hand side moves the fate of the
system from ``many answers'' to ``no answer''. The question of consistency therefore cannot be
answered from $A$ alone.

The case that appears most often in practice is $m=n$. There the three fates rest on a single
number, and that number is what we look for in the next section.

\begin{akibat}
Suppose $A$ is of size $n\times n$. The system $A\mathbf{x}=\mathbf{b}$ has a unique answer for
every $\mathbf{b}\in\mathbb{R}^{n}$ exactly when $\operatorname{rank}A=n$.
\end{akibat}

\begin{proof}
If $\operatorname{rank}A=n$, the reduced echelon form of $A$ is $I_{n}$, so elimination on
$\bigl(A\mid\mathbf{b}\bigr)$ always ends with a pivot in every column of $A$ and no pivot in
the last column; the previous theorem gives a unique answer. Conversely, if
$\operatorname{rank}A<n$, some column has no pivot, and for a consistent $\mathbf{b}$ the answer
is not unique. Since the rank is less than $n$, the column space of $A$ does not fill
$\mathbb{R}^{n}$, so some $\mathbf{b}$ lies outside it; for that $\mathbf{b}$ the system is
inconsistent.
\end{proof}

\section{The determinant as an alternating function}
\tingkat{intermediate}

We want one number that answers the question of the preceding corollary without solving any
system. That number must be zero exactly when the columns of $A$ are not linearly independent,
must change sign when two columns are interchanged, and must equal one for the identity matrix.
Those three conditions turn out to determine the number uniquely.

\kamus{determinant}{determinan}{The volume factor: how many times space is stretched or
flattened by the matrix.}

\begin{definisi}[Determinant]
A function $D$ from $n\times n$ matrices to $\mathbb{R}$ is called a \emph{determinant} if
\begin{enumerate}
  \item \emph{multilinear in each row}: for scalars $\alpha,\beta$ and rows
        $\mathbf{u},\mathbf{v}$,
        $D(\dots,\alpha\mathbf{u}+\beta\mathbf{v},\dots)
         =\alpha D(\dots,\mathbf{u},\dots)+\beta D(\dots,\mathbf{v},\dots)$;
  \item \emph{alternating}: $D(\dots,\mathbf{u},\dots,\mathbf{u},\dots)=0$ if two rows are
        equal;
  \item \emph{normalized}: $D(I_{n})=1$.
\end{enumerate}
The value $D(A)$ is written $\det A$.\index{determinant}
\end{definisi}

Multilinearity means the determinant is linear in one row while the other rows are held fixed.
The alternating property is the strongest condition, and its consequences are long.

\begin{lemma}
Interchanging two rows changes the sign of the determinant. If one row is a scalar multiple of
another row, or if there is a zero row, the determinant is zero.
\end{lemma}

\begin{proof}
Suppose two adjacent rows both contain $\mathbf{u}+\mathbf{v}$. Since the two rows are equal,
the value is zero. Multilinearity splits that zero into four terms,
\[
  0=D(\dots,\mathbf{u},\mathbf{u},\dots)+D(\dots,\mathbf{u},\mathbf{v},\dots)
   +D(\dots,\mathbf{v},\mathbf{u},\dots)+D(\dots,\mathbf{v},\mathbf{v},\dots),
\]
and the first and fourth terms vanish. Interchanging two adjacent rows therefore reverses the
sign:
\[
  D(\dots,\mathbf{v},\mathbf{u},\dots)=-D(\dots,\mathbf{u},\mathbf{v},\dots).
\]
An interchange of rows that are not adjacent is achieved by several adjacent interchanges, and
every step reverses the sign.

For a multiple, $D(\dots,c\mathbf{u},\dots,\mathbf{u},\dots)
=c\,D(\dots,\mathbf{u},\dots,\mathbf{u},\dots)=0$ by multilinearity and the alternating
property. A zero row is the case $c=0$ of the same form.
\end{proof}

Now uniqueness. The formula that appears in the middle of the proof is the Leibniz formula, and
that formula will later yield the cofactor expansion as a consequence.

\begin{teorema}[Uniqueness and the Leibniz formula]
There is only one function satisfying the three conditions of the definition, and its value is
\[
  \det A=\sum_{\sigma}\operatorname{sgn}(\sigma)\,
  a_{\sigma(1)1}a_{\sigma(2)2}\cdots a_{\sigma(n)n},
\]
where the sum is taken over all permutations $\sigma$ of $\{1,\dots,n\}$ and
$\operatorname{sgn}(\sigma)=(-1)^{k}$ if $\sigma$ is reached by $k$ interchanges.
\end{teorema}

\begin{proof}
Write row $i$ as $\mathbf{a}_{i}=\sum_{j}a_{ij}\mathbf{e}_{j}$. Multilinearity allows every row
to be pulled out into a sum,
\[
  D(A)=\sum_{j_{1},\dots,j_{n}}a_{1j_{1}}a_{2j_{2}}\cdots a_{nj_{n}}\,
  D(\mathbf{e}_{j_{1}},\dots,\mathbf{e}_{j_{n}}).
\]
Terms with $j_{p}=j_{q}$ for $p\neq q$ vanish because $D$ is alternating, so only
$(j_{1},\dots,j_{n})$ that form a permutation $\sigma$ survive. By the previous lemma,
$D(\mathbf{e}_{\sigma(1)},\dots,\mathbf{e}_{\sigma(n)})
=\operatorname{sgn}(\sigma)D(I_{n})=\operatorname{sgn}(\sigma)$.

For uniqueness, suppose $D$ satisfies the three conditions. The same computation gives
$D(A)=\bigl(\sum_{\sigma}\operatorname{sgn}(\sigma)a_{\sigma(1)1}\cdots
a_{\sigma(n)n}\bigr)D(I_{n})=\det A$. Every function satisfying the three conditions therefore
equals that formula.
\end{proof}

The Leibniz formula is correct but impractical: it sums $n!$ terms, and for $n=12$ there are
$479\,001\,600$ of them. What people use is another form of the same formula.

\begin{akibat}[Cofactor expansion]
For every fixed row $i$,
\[
  \det A=\sum_{j=1}^{n}(-1)^{i+j}a_{ij}M_{ij},
\]
where $M_{ij}$ is the determinant of the $(n-1)\times(n-1)$ submatrix obtained by deleting row
$i$ and column $j$. The number $C_{ij}=(-1)^{i+j}M_{ij}$ is called the \emph{cofactor} of the
entry $a_{ij}$. The same formula holds if the expansion is carried out along a column.
\end{akibat}

\begin{proof}
In the Leibniz formula, group the terms by the value of $j$ at position row $i$. Every term
containing $a_{ij}$ has the form $a_{ij}$ times a sign times the product of $n-1$ other
entries, one from each row other than $i$ and one from each column other than $j$. That product
is exactly the Leibniz term for the minor $M_{ij}$, and its sign is $(-1)^{i+j}$ because moving
row $i$ to the last position requires $n-i$ interchanges and moving column $j$ to the last
position requires $n-j$ interchanges, so the sign is $(-1)^{2n-i-j}=(-1)^{i+j}$. Summing over
$j$ gives the formula as written.

For expansion along a column, we use $\det A=\det A^{\top}$. That equality itself follows from
the Leibniz formula: every term for $A^{\top}$ has the form
$a_{\sigma(1)1}\cdots a_{\sigma(n)n}$ with the index order reversed, and the map
$\sigma\mapsto\sigma^{-1}$ preserves the sign because the number of interchanges is the same.
\end{proof}

\begin{contoh}
For
$A=\begin{pmatrix}2&1&-1\\ -3&-1&2\\ -2&1&2\end{pmatrix}$ expansion along the first row gives the
cofactors
\[
  C_{11}=\begin{vmatrix}-1&2\\ 1&2\end{vmatrix}=-4,\qquad
  C_{12}=-\begin{vmatrix}-3&2\\ -2&2\end{vmatrix}=2,\qquad
  C_{13}=\begin{vmatrix}-3&-1\\ -2&1\end{vmatrix}=-5,
\]
so that $\det A=2(-4)+1(2)+(-1)(-5)=-8+2+5=-1$. The alternating sign is not decoration: without
the sign, the result would be $-8+2-5=-11$, and every conclusion afterwards would be wrong.
\end{contoh}

The determinant is useful for product matrices because it is multiplicative.

\begin{teorema}
For $n\times n$ matrices, $\det(AB)=\det A\,\det B$.
\end{teorema}

\begin{proof}
Fix $A$ and define $f(B)=\det(AB)$. Row $i$ of $AB$ is a linear combination of the rows of $B$
with coefficients from row $i$ of $A$, so $f$ is multilinear in the rows of $B$. If two rows of
$B$ are equal, two rows of $AB$ are also equal, so $f$ is alternating. By uniqueness, every
multilinear alternating function $f$ satisfies $f(B)=\det(B)\,f(I_{n})$. Here $f(I_{n})=\det A$,
so $f(B)=\det B\,\det A$.
\end{proof}

\begin{akibat}
$A$ is invertible if and only if $\det A\neq0$. If it is invertible, $\det(A^{-1})=1/\det A$.
Moreover $\det(cA)=c^{n}\det A$, and the determinant of a triangular matrix is the product of
its diagonal entries.
\end{akibat}

\begin{proof}
If $A$ is invertible, $\det A\,\det(A^{-1})=\det I_{n}=1$, so $\det A\neq0$ and its inverse is
$1/\det A$. If $A$ is not invertible, its columns are linearly dependent by Chapter 7, so one
column is a linear combination of the others. Multilinearity splits $\det A$ into a sum of
determinants with two equal columns, so all of them are zero. The formula $\det(cA)=c^{n}\det A$
is obtained by pulling a factor $c$ out of every row. For a triangular matrix, cofactor
expansion along a row that borders only one diagonal entry, repeated $n$ times, leaves the
product of the diagonal.
\end{proof}

That corollary makes elimination and the determinant meet. Elimination brings $A$ to an upper
triangular matrix $U$ by interchanges and additions of row multiples. Additions do not change
the determinant; interchanges reverse its sign. Hence $\det A=(-1)^{s}\det U$ with $s$ the
number of interchanges, and $\det U$ is the product of the pivots: the determinant comes free
after elimination, and that is how libraries compute it.

The other side of the determinant, just as important, is area. The columns of
$B=\begin{pmatrix}3&1\\ 1&2\end{pmatrix}$ are two arrows, $(3,1)$ and $(1,2)$. The unit square
whose sides are carried by those two arrows turns into a parallelogram, and its area is five,
exactly $\det B=3\cdot2-1\cdot1=5$. Figure~\ref{fig:luas-determinan} shows that mapping. If the
two columns are collinear, the parallelogram flattens into a single segment, its area is zero,
and the determinant is zero as well.

\begin{figure}[htbp]
\centering
\begin{tikzpicture}[scale=0.95,>=Stealth]
  \draw[layergrey!28,step=1] (-0.4,-0.4) grid (4.4,3.4);
  \draw[->,ink] (-0.5,0)--(4.6,0) node[right,font=\scriptsize]{$x$};
  \draw[->,ink] (0,-0.5)--(0,3.5) node[above,font=\scriptsize]{$y$};
  \fill[layerteal!16] (0,0)--(3,1)--(4,3)--(1,2)--cycle;
  \draw[thick,layerteal] (0,0)--(3,1)--(4,3)--(1,2)--cycle;
  \fill[layerorange!22] (0,0)--(1,0)--(1,1)--(0,1)--cycle;
  \draw[thick,layerorange] (0,0)--(1,0)--(1,1)--(0,1)--cycle;
  \draw[very thick,accent,->] (0,0)--(3,1)
    node[below right,font=\scriptsize]{$(3,1)$};
  \draw[very thick,accent,->] (0,0)--(1,2)
    node[above left,font=\scriptsize]{$(1,2)$};
  \foreach \t in {1,2,3,4} {\node[font=\tiny,below] at (\t,-0.02) {\t};}
  \foreach \t in {1,2,3} {\node[font=\tiny,left] at (-0.02,\t) {\t};}
  \node[font=\tiny,fill=white,inner sep=0.5pt] at (0.5,0.5) {1};
  \node[font=\scriptsize,fill=white,inner sep=1.5pt] at (2,1.5) {area $5$};
\end{tikzpicture}
\caption{The determinant as area. The unit square (orange, area $1$) is mapped by
$B=\begin{pmatrix}3&1\\ 1&2\end{pmatrix}$ to the parallelogram (teal) spanned by
$(3,1)$ and $(1,2)$. Its area is $5$, the same as $\det B=3\cdot2-1\cdot1=5$: the determinant
multiplies area by five, and its sign records whether the direction of rotation is reversed.}
\label{fig:luas-determinan}
\end{figure}

\begin{tangan}
Materials: graph paper, scissors, ruler, glue.
\begin{enumerate}
  \item Draw a rectangle three squares by one square. Its area is three squares.
  \item Cut out the rectangle. From the left end, cut off a right triangle with base two squares
        and height one square, then glue that piece onto the right end.
  \item The new shape is a slanted parallelogram. Its area is still three squares, because no
        paper was thrown away and none was added.
  \item Repeat with a slant of one square, then two, then three. Each time compute base times
        height, not the length of the slanted side.
\end{enumerate}
What the reader should see: the shape changes, the area does not. A slant only shifts, it does
not stretch, so the determinant does not grow.
\end{tangan}

\begin{untukanak}
Take a rectangular sheet of paper, then slant it by pushing its top edge sideways without
shortening or lengthening anything. Its height stays the same, its width stays the same, its
area stays the same, even though the shape is no longer a rectangle. Only the direction of
rotation changes: trace the bottom edge and then the right edge starting from the lower left
corner, and the slanted sheet still turns the same way; swap the two edges, and the direction
reverses like a left shoe on a right foot.

In this chapter we will learn why a single number can report both area and orientation, and why
that number becomes zero exactly when the sheet is folded into a line. In your opinion, if the
sheet is folded in half, is there still a way to compute its area with the same number?
\end{untukanak}

\section{Cramer's rule and the adjugate}
\tingkat{intermediate}

The determinant answers whether a unique answer exists. The next question: can that answer be
written as a formula, rather than as the result of steps? For small systems it can, and the
formula is tidy.

\kamus{adjugate}{adjoin}{The cofactor map arranged as a transpose, used to write the inverse as
a single formula.}

\begin{definisi}[Adjugate]
For $A$ of size $n\times n$ with cofactors $C_{ij}$, the \emph{cofactor matrix} is the matrix
with entry $C_{ij}$ at position $(i,j)$, and the \emph{adjugate} $\operatorname{adj}A$ is its
transpose, that is, the matrix with entry $C_{ji}$ at position $(i,j)$.
\index{cofactor}\index{adjugate}
\end{definisi}

Note the reversal of the indices: the cofactor $C_{ij}$ comes from deleting row $i$ and column
$j$, but is placed at position $(j,i)$. This is the most frequent source of error when
computing an adjugate by hand.

\begin{teorema}
For every $n\times n$ matrix, $A\operatorname{adj}A
=(\det A)I_{n}=\operatorname{adj}A\,A$.
\end{teorema}

\begin{proof}
Entry $(i,j)$ of $A\operatorname{adj}A$ is $\sum_{k}a_{ik}C_{jk}$. If $i=j$, that sum is exactly
the cofactor expansion along row $i$, namely $\det A$. If $i\neq j$, that sum is the expansion
along row $j$ of the matrix obtained from $A$ by replacing row $j$ with row $i$. That matrix has
two equal rows, so its determinant is zero. Hence $A\operatorname{adj}A=(\det A)I_{n}$. The
second equality is obtained by the same argument on columns, or from
$\operatorname{adj}(A^{\top})=(\operatorname{adj}A)^{\top}$.
\end{proof}

\begin{akibat}
If $\det A\neq0$, then $A^{-1}=\dfrac{1}{\det A}\operatorname{adj}A$.
\end{akibat}

\begin{proof}
From the previous theorem, $A\bigl(\operatorname{adj}A/\det A\bigr)=I_{n}$ and
$\bigl(\operatorname{adj}A/\det A\bigr)A=I_{n}$. The inverse of a matrix is unique, so that
matrix is $A^{-1}$.
\end{proof}

That formula holds for all $n$, but the price is high: $\operatorname{adj}A$ contains $n^{2}$
cofactors, and every cofactor is an $(n-1)\times(n-1)$ determinant.

With that inverse formula, the answer of a system can be written as a ratio of two determinants.

\begin{teorema}[Cramer's rule]
Suppose $A$ is of size $n\times n$ with $\det A\neq0$, and $\mathbf{b}\in\mathbb{R}^{n}$. The
unique answer of $A\mathbf{x}=\mathbf{b}$ is
\[
  x_{j}=\frac{\det A_{j}}{\det A},\qquad j=1,\dots,n,
\]
where $A_{j}$ is the matrix obtained from $A$ by replacing column $j$ with $\mathbf{b}$.
\index{Cramer}
\end{teorema}

\begin{proof}
Since $\det A\neq0$, the previous corollary gives
$\mathbf{x}=A^{-1}\mathbf{b}=(\operatorname{adj}A)\,\mathbf{b}/\det A$. Entry $j$ of that product
is $\sum_{i}b_{i}C_{ji}/\det A$. The sum $\sum_{i}b_{i}C_{ji}$ is the cofactor expansion along
column $j$ of the matrix whose column $j$ has been replaced by $\mathbf{b}$, namely $A_{j}$.
Hence $x_{j}=\det A_{j}/\det A$.
\end{proof}

\begin{contoh}
Back to the opening system, $A\mathbf{x}=\mathbf{b}$ with $\mathbf{b}=(8,-11,-3)^{\top}$. We
already computed $\det A=-1$, and
\[
  \det A_{1}=\begin{vmatrix}8&1&-1\\ -11&-1&2\\ -3&1&2\end{vmatrix}=-2,\qquad
  \det A_{2}=-3,\qquad \det A_{3}=1 .
\]
So $x_{1}=(-2)/(-1)=2$, $x_{2}=(-3)/(-1)=3$, and $x_{3}=1/(-1)=-1$. Check:
$A(2,3,-1)^{\top}=(4+3+1,\,-6-3-2,\,-4+3-2)^{\top}=(8,-11,-3)^{\top}$.

The adjugate of that matrix, with the cofactors of its first row already computed above, is
\[
  \operatorname{adj}A=\begin{pmatrix}-4&-3&1\\ 2&2&-1\\ -5&-4&1\end{pmatrix},
  \qquad
  A^{-1}=\frac{1}{-1}\operatorname{adj}A
  =\begin{pmatrix}4&3&-1\\ -2&-2&1\\ 5&4&-1\end{pmatrix}.
\]
The first row of $A^{-1}$ times $\mathbf{b}$ gives $4(8)+3(-11)-1(-3)=32-33+3=2$, matching
$x_{1}$.
\end{contoh}

Cramer's rule is beautiful and rarely used for computing. Its value lies elsewhere: the formula
shows that the answer depends smoothly on $\mathbf{b}$ and on the entries of $A$, and that
dependence is what we measure through cost and sensitivity.

\section{Elementary matrices, inverses, and equivalent conditions}
\tingkat{advanced}

Row elimination can be seen as matrix multiplication. That view turns the question ``is $A$
invertible'' into a question about a list of factors, and that list of factors yields an
algorithm for the inverse.

\kamus{elementary matrix}{matriks elementer}{The identity matrix with one of its steps already
carried out, so that it records a single row operation.}

\begin{definisi}[Elementary matrix]
An \emph{elementary matrix} is a matrix obtained from $I_{n}$ by one elementary row operation.
There are three kinds, matching the three operations: interchanging two rows, multiplying one
row by a nonzero scalar, and adding a multiple of one row to another row.
\end{definisi}

\begin{lemma}
If $E$ is an $m\times m$ elementary matrix and $A$ an $m\times n$ matrix, then $EA$ is the
result of performing the same row operation on $A$. Every elementary matrix is invertible, and
its inverse is an elementary matrix of the same kind.
\end{lemma}

\begin{proof}
Row $i$ of $EA$ is a linear combination of the rows of $A$ with coefficients from row $i$ of
$E$. For $E$ of the third kind, that row is $\mathbf{a}_{i}+c\mathbf{a}_{j}$, exactly as in the
row operation; for the other two kinds the check is equally direct. The inverse is obtained by
undoing the operation: interchanging rows undoes itself, and
$\operatorname{diag}(1,\dots,c,\dots,1)^{-1}=\operatorname{diag}(1,\dots,1/c,\dots,1)$.
\end{proof}

Figure~\ref{fig:tiga-gerakan} runs the three operations on the system $x+2y=4$, $3x+4y=10$,
which intersect at $(2,1)$. Interchanging rows only interchanges names; multiplying the first
row by $2$ writes the same line with numbers twice as large; adding the second row to the first
gives $4x+6y=14$, which still passes through $(2,1)$. The intersection point does not move, but
the determinant of the coefficient matrix changes sign, doubles, then stays fixed. The solution
set is preserved; the number that measures area is not.

\begin{figure}[htbp]
\centering
\begin{tikzpicture}[scale=0.58,>=Stealth]
  \begin{scope}
    \draw[layergrey!28,step=1] (-0.4,-0.4) grid (4.4,3.2);
    \draw[->,ink] (-0.5,0)--(4.6,0) node[right,font=\scriptsize]{$x$};
    \draw[->,ink] (0,-0.5)--(0,3.4) node[above,font=\scriptsize]{$y$};
    \draw[very thick,layerindigo] (0,2)--(4,0) node[pos=0.15,above,font=\scriptsize]{$R_{1}$};
    \draw[very thick,layerteal] (0,2.5)--(3.33,0) node[pos=0.85,above left,font=\scriptsize]{$R_{2}$};
    \fill[accent] (2,1) circle (2.6pt);
    \node[font=\scriptsize,below right] at (2,1) {$(2,1)$};
    \node[font=\scriptsize] at (2,-1.1) {(i) interchange rows};
  \end{scope}
  \begin{scope}[xshift=5.6cm]
    \draw[layergrey!28,step=1] (-0.4,-0.4) grid (4.4,3.2);
    \draw[->,ink] (-0.5,0)--(4.6,0) node[right,font=\scriptsize]{$x$};
    \draw[->,ink] (0,-0.5)--(0,3.4) node[above,font=\scriptsize]{$y$};
    \draw[very thick,layerindigo] (0,2)--(4,0) node[pos=0.15,above,font=\scriptsize]{$R_{1}$};
    \draw[very thick,layerindigo,densely dashed] (0,2)--(4,0);
    \draw[very thick,layerteal] (0,2.5)--(3.33,0) node[pos=0.85,above left,font=\scriptsize]{$R_{2}$};
    \fill[accent] (2,1) circle (2.6pt);
    \node[font=\scriptsize,below right] at (2,1) {$(2,1)$};
    \node[font=\scriptsize] at (2,-1.1) {(ii) multiply a row};
  \end{scope}
  \begin{scope}[xshift=11.2cm]
    \draw[layergrey!28,step=1] (-0.4,-0.4) grid (4.4,3.2);
    \draw[->,ink] (-0.5,0)--(4.6,0) node[right,font=\scriptsize]{$x$};
    \draw[->,ink] (0,-0.5)--(0,3.4) node[above,font=\scriptsize]{$y$};
    \draw[very thick,layergrey,dash pattern=on 3pt off 2pt] (0,2)--(4,0);
    \draw[very thick,layerteal] (0,2.5)--(3.33,0)
      node[pos=0.8,below right,font=\scriptsize]{$R_{2}$};
    \draw[very thick,layerplum] (0,2.33)--(3.5,0)
      node[pos=0.2,above,font=\scriptsize]{$R_{1}+R_{2}$};
    \fill[accent] (2,1) circle (2.6pt);
    \node[font=\scriptsize,below right] at (2,1) {$(2,1)$};
    \node[font=\scriptsize] at (2,-1.1) {(iii) add a multiple};
  \end{scope}
\end{tikzpicture}
\caption{The three elementary row operations on the system $x+2y=4$, $3x+4y=10$. Interchanging
rows only swaps the names $R_{1}$ and $R_{2}$; multiplying $R_{1}$ by $2$ writes the same line;
adding $R_{2}$ to $R_{1}$ rotates the first line into $4x+6y=14$. In all three panels, both
lines pass through the same intersection point, $(2,1)$.}
\label{fig:tiga-gerakan}
\end{figure}

\begin{teorema}
An $n\times n$ matrix $A$ is invertible if and only if $A$ is a product of elementary matrices.
\end{teorema}

\begin{proof}
A product of invertible matrices is also invertible, so one direction is clear.

For the other direction, suppose $A$ is invertible and bring $A$ to reduced echelon form by row
operations, which is the same as multiplying from the left by elementary matrices
$E_{1},\dots,E_{k}$. By the theorem of three cases, $A$ invertible means every column of $A$
contains a pivot, so the reduced form of $A$ is $I_{n}$, since $A$ is square with $n$ pivots.
Hence $E_{k}\cdots E_{1}A=I_{n}$, and therefore $A=E_{1}^{-1}\cdots E_{k}^{-1}$, a product of
elementary matrices.
\end{proof}

That proof is also an algorithm: performing the same operations on $I_{n}$ gives $A^{-1}$, so
$\bigl(A\mid I_{n}\bigr)\to\bigl(I_{n}\mid A^{-1}\bigr)$ produces the inverse without a single
determinant, far cheaper than the adjugate.

We collect what has been proved into one list. This is the list people use every day, and its
content is another face of a single statement.

\begin{teorema}[Equivalent conditions]
For an $n\times n$ matrix $A$, the following statements are equivalent.
\begin{enumerate}
  \item $A$ is invertible.
  \item $A\mathbf{x}=\mathbf{0}$ has only the trivial answer $\mathbf{x}=\mathbf{0}$.
  \item The reduced echelon form of $A$ is $I_{n}$.
  \item $A$ is a product of elementary matrices.
  \item $A\mathbf{x}=\mathbf{b}$ is consistent for every $\mathbf{b}\in\mathbb{R}^{n}$.
  \item $A\mathbf{x}=\mathbf{b}$ has exactly one answer for every $\mathbf{b}\in\mathbb{R}^{n}$.
  \item $\det A\neq0$.
\end{enumerate}
\end{teorema}

\begin{proof}
We chain the implications $(1)\Rightarrow(2)\Rightarrow(3)\Rightarrow(4)\Rightarrow(1)$,
then $(1)\Rightarrow(5)\Rightarrow(6)\Rightarrow(1)$ and $(1)\Leftrightarrow(7)$.

$(1)\Rightarrow(2)$: multiply $A\mathbf{x}=\mathbf{0}$ from the left by $A^{-1}$.
$(2)\Rightarrow(3)$: the kernel of $A$ is only $\{\mathbf{0}\}$, which means nullity zero, so by
rank-nullity in Chapter 7, $\operatorname{rank}A=n$; hence there are $n$ pivots and the reduced
form is $I_{n}$. $(3)\Rightarrow(4)$: the elimination steps give $E_{k}\cdots E_{1}A=I_{n}$.
$(4)\Rightarrow(1)$: a product of invertible matrices is also invertible.
$(1)\Rightarrow(5)$: take $\mathbf{x}=A^{-1}\mathbf{b}$. $(5)\Rightarrow(6)$: suppose some
$\mathbf{b}$ has two different answers $\mathbf{x}\neq\mathbf{y}$. Then
$\mathbf{x}-\mathbf{y}\neq\mathbf{0}$ and $A(\mathbf{x}-\mathbf{y})=\mathbf{0}$, so the columns
of $A$ are linearly dependent and $\operatorname{rank}A<n$. The column space of $A$ therefore
has dimension less than $n$, so some $\mathbf{b}'$ is out of reach, contradicting (5).
$(6)\Rightarrow(1)$: if $A$ is not invertible, its columns are dependent, so some
$\mathbf{x}\neq\mathbf{0}$ satisfies $A\mathbf{x}=\mathbf{0}$, and the answer for
$\mathbf{b}=\mathbf{0}$ is not unique. Finally $(1)\Leftrightarrow(7)$ is the corollary of the
previous section.
\end{proof}

That list explains why four questions that look different, about solutions, echelon form,
factorization, and the determinant, are in fact one question; only the measuring tool changes.
But the list is silent about cost and sensitivity, and those two things decide whether the
answer obtained can be trusted.

\begin{tantangan}
The following system looks ordinary: three equations for three numbers, no conspicuous
duplicate row, so the answer seems to be unique.
\[
  x+y+z=3,\qquad x+2y+3z=6,\qquad 2x+3y+4z=9 .
\]
Find all its answers. How many answers are there, and why is counting equations misleading here?

Hint 1: add the first two equations, then compare the result with the third equation.

Hint 2: the third equation adds no information at all, so only two pieces of information are
independent. Once $z$ is chosen, $y$ and then $x$ are determined; write the answer as one family
with one parameter.

If you get stuck, go back to the section on three fates and count the pivots.
\end{tantangan}

\subsection*{Cost and sensitivity}

Elimination on an $n\times n$ matrix uses about $n^{3}/3$ multiplications: for $n=1000$ about
$3.33\times10^{8}$, which is still fast. Cramer's rule with cofactors uses $n\cdot n!$
multiplications, about $5.75\times10^{9}$ for $n=12$ and above $10^{19}$ for $n=20$; that is
why it is used only for small $n$ or for theoretical purposes.

The subtler matter is sensitivity. Suppose $A$ is invertible and $\mathbf{b}\neq\mathbf{0}$, and
suppose $\mathbf{b}$ changes to $\mathbf{b}+\delta\mathbf{b}$ so that the answer changes to
$\mathbf{x}+\delta\mathbf{x}$.

\begin{teorema}
With the notation above, for any subordinate matrix norm,
\[
  \frac{\lVert\delta\mathbf{x}\rVert}{\lVert\mathbf{x}\rVert}
  \le\kappa(A)\,
  \frac{\lVert\delta\mathbf{b}\rVert}{\lVert\mathbf{b}\rVert},
  \qquad \kappa(A)=\lVert A\rVert\,\lVert A^{-1}\rVert .
\]
\end{teorema}

\begin{proof}
From $A\,\delta\mathbf{x}=\delta\mathbf{b}$ we get $\delta\mathbf{x}=A^{-1}\delta\mathbf{b}$, so
$\lVert\delta\mathbf{x}\rVert\le\lVert A^{-1}\rVert\lVert\delta\mathbf{b}\rVert$. From
$\mathbf{b}=A\mathbf{x}$ we get $\lVert\mathbf{b}\rVert\le\lVert A\rVert\lVert\mathbf{x}\rVert$,
that is, $1/\lVert\mathbf{x}\rVert\le\lVert A\rVert/\lVert\mathbf{b}\rVert$. Multiplying the two
inequalities gives the required result.
\end{proof}

The number $\kappa(A)$ is called the \emph{condition number}. It does not depend on scale:
multiplying $A$ by $c$ does not change it, because $\lVert cA\rVert\lVert (cA)^{-1}\rVert
=|c|\lVert A\rVert\cdot|c|^{-1}\lVert A^{-1}\rVert$. The determinant has no such property, and
this is why we do not use the determinant as a measure of sensitivity. A small determinant only
means a small volume; it can be small because the whole matrix was scaled down, while the
sensitivity does not change.

For the opening system, $\kappa(S)\approx1.01\times10^{5}$. With machine precision around
$10^{-16}$, the digits that can be trusted in the answer are about $10^{-11}$ relative: enough
for this example, exhausted for a worse matrix. Chapter 22 discusses thresholds and how to
recognize a matrix that can no longer be trusted.
\index{condition}

\begin{untukanak}
Imagine a long plank used as a base for measuring the height of a hill. The plank is almost
horizontal, only slightly tilted. If the lower end is lifted by a finger, the other end rises
far more, because the plank rests on a very small slope. A finger in one place becomes a
hand-span in another.

Now imagine the plank is made very steep. Lift its lower end by a finger, and its upper end
rises only a tiny bit. The same work produces two very different effects, only because the base
is different.

Our computations work like that plank. There are arrangements of numbers in which a small change
in the input produces a small change in the answer; there are also arrangements in which a
change the size of a speck of dust produces a completely different answer. What decides is not
how large the numbers are, but the shape of the arrangement.

In this chapter we will learn why one number can predict that before we compute, and why that
number does not change when all the numbers are enlarged or reduced by the same ratio.

If you weigh on a scale whose base is tilted, is the result wrong, or merely untrustworthy?
\end{untukanak}

\section{A worked example}
\tingkat{intermediate}

We work one system in two ways, compute one $4\times4$ determinant, then measure the sensitivity
of two matrices.

\begin{resep}
\begin{enumerate}
  \item Write the augmented matrix $\bigl(A\mid\mathbf{b}\bigr)$: coefficient columns first, the
        right-hand side column last.
  \item Work on the leftmost column. Choose one nonzero entry as the pivot; if the top entry is
        zero, interchange rows first.
  \item Zero out all entries below the pivot: for each row $j$ below it,
        $R_{j}\leftarrow R_{j}-\frac{a_{ji}}{a_{ii}}R_{i}$.
  \item Shift one column to the right and repeat; stop when there are no nonzero rows left, then
        read the fate of the system from the pivot positions.
\end{enumerate}
Warning about the order: the pivots are zeroed from top to bottom, column by column, before back
substitution begins. Reversing that order destroys the pivots already obtained and forces the
work to be repeated.
\end{resep}

The system is $A\mathbf{x}=\mathbf{b}$ with
\[
  A=\begin{pmatrix}2&1&-1\\ -3&-1&2\\ -2&1&2\end{pmatrix},
  \qquad \mathbf{b}=\begin{pmatrix}8\\ -11\\ -3\end{pmatrix}.
\]

First way, elimination, with the steps written in full:
\[
  \left(\begin{array}{ccc|c}2&1&-1&8\\ -3&-1&2&-11\\ -2&1&2&-3\end{array}\right)
  \xrightarrow[R_{3}\leftarrow R_{3}+R_{1}]{R_{2}\leftarrow R_{2}+\tfrac{3}{2}R_{1}}
  \left(\begin{array}{ccc|c}2&1&-1&8\\ 0&\tfrac{1}{2}&\tfrac{1}{2}&1\\
  0&2&1&5\end{array}\right),
\]
then $R_{3}\leftarrow R_{3}-4R_{2}$ gives
\[
  \left(\begin{array}{ccc|c}2&1&-1&8\\ 0&\tfrac{1}{2}&\tfrac{1}{2}&1\\
  0&0&-1&1\end{array}\right).
\]
Back substitution: from the third row $-x_{3}=1$, so $x_{3}=-1$. The second row
$\tfrac{1}{2}x_{2}+\tfrac{1}{2}(-1)=1$ gives $x_{2}=3$. The first row $2x_{1}+3+1=8$ gives
$x_{1}=2$. The result is $\mathbf{x}=(2,3,-1)$. The pivots that appear are $2$, $\tfrac{1}{2}$,
and $-1$, so the product of the pivots is $\det A=2\cdot\tfrac{1}{2}\cdot(-1)=-1$, matching the
cofactor expansion above. Here there was no row interchange, so the sign did not change.

The second way, Cramer's rule, we already ran on the same system in the previous section. The
three replacement determinants give $\det A_{1}=-2$, $\det A_{2}=-3$, and $\det A_{3}=1$, so
$\mathbf{x}=(-2,-3,1)/(-1)=(2,3,-1)$. Two roads, one answer; a check like this is cheap and
should always be carried out.

Now the $4\times4$ determinant. Take
\[
  M=\begin{pmatrix}2&0&1&3\\ 1&4&0&2\\ 3&1&2&1\\ 0&2&1&5\end{pmatrix}.
\]
Expansion along the first row uses four minors, and one of them we compute in full:
\[
  M_{11}=\begin{vmatrix}4&0&2\\ 1&2&1\\ 2&1&5\end{vmatrix}
  =4(10-1)-0(5-2)+2(1-4)=36-6=30 .
\]
In a similar way $M_{12}=15$, $M_{13}=-45$, and $M_{14}=-15$. Hence
\[
  \det M=2(30)-0(15)+1(-45)-3(-15)=60-45+45=60 .
\]
The number of terms summed is $4!=24$, and for $5\times5$ it is already $120$: the growth is
factorial. This is why libraries do not use this method.

Now sensitivity. The $4\times4$ Hilbert matrix is defined by $h_{ij}=1/(i+j-1)$. Its determinant
is exactly $1/6\,048\,000\approx1.65\times10^{-7}$, very small, so one easily concludes that
this matrix is dangerous. Its condition number is $15\,513.74$. The opening matrix had
$\det S=-0.003$ and $\kappa(S)\approx1.01\times10^{5}$, so its determinant is eighteen
thousand times larger, while its sensitivity is about six and a half times worse. The determinant
is not a measure of sensitivity; the condition number is.

The jump test for $S$ uses $\mathbf{x}=(1,1,1)$, so that $\mathbf{b}=(6,\,15,\,24.001)$.
Changing the third component to $24.002$ changes the answer to $(2,-1,2)$. The relative change
in the data is $3.46\times10^{-5}$, the relative change in the answer is
$\sqrt{2}\approx1.414$, and the ratio is $40\,916$: below $\kappa(S)$ as the theorem promises,
but only by a factor of two and a half.

\kertas{Which is more sensitive, $M$ or the $4\times4$ Hilbert matrix? Predict by comparing their
determinants, then compute the condition number of both with \texttt{numpy.linalg.cond}. Compare
also with $10^{-3}I_{4}$, whose determinant is far smaller, then explain why the rule ``small
determinant means sensitive'' fails. Finally, multiply $M$ by $10^{-4}$: its condition number
does not change, while its determinant shrinks by a factor of $10^{-16}$.}

\section{Computational lab}
\tingkat{advanced}

\begin{lab}
\begin{verbatim}
import numpy as np, math

def det_kofaktor(A):
    n = A.shape[0]
    if n == 1:
        return A[0, 0]
    return sum((-1) ** j * A[0, j] *
               det_kofaktor(np.delete(np.delete(A, 0, 0), j, 1))
               for j in range(n))

A = np.array([[2., 1., -1.], [-3., -1., 2.], [-2., 1., 2.]])
b = np.array([8., -11., -3.])
x = np.linalg.solve(A, b)
print("det A    :", np.linalg.det(A))
print("det cof  :", det_kofaktor(A))
print("rank A   :", np.linalg.matrix_rank(A))
print("x        :", x)
print("residual :", np.linalg.norm(A @ x - b))
for j in range(3):
    Aj = A.copy(); Aj[:, j] = b
    print("det A%d   : %.1f  -> x%d = %.6f" % (j + 1, np.linalg.det(Aj),
          j + 1, np.linalg.det(Aj) / np.linalg.det(A)))

M = np.array([[2., 0., 1., 3.], [1., 4., 0., 2.],
              [3., 1., 2., 1.], [0., 2., 1., 5.]])
print("det M    :", det_kofaktor(M), np.linalg.det(M))

H = np.array([[1. / (i + j + 1) for j in range(4)] for i in range(4)])
print("det H4   :", np.linalg.det(H))
print("cond H4  : %.2f" % np.linalg.cond(H))

S = np.array([[1., 2., 3.], [4., 5., 6.], [7., 8., 9.001]])
bs = S @ np.array([1., 1., 1.])
print("det S    :", np.linalg.det(S), " cond S: %.1f" % np.linalg.cond(S))
print("x        :", np.linalg.solve(S, bs))
bs2 = bs.copy(); bs2[2] += 0.001
print("x'       :", np.linalg.solve(S, bs2))
print("cost     :", math.factorial(12) * 12, round(1000 ** 3 / 3))
\end{verbatim}
Actual output:
\begin{verbatim}
det A    : -1.0
det cof  : -1.0
rank A   : 3
x        : [ 2.  3. -1.]
residual : 1.9860273225978185e-15
det A1   : -2.0  -> x1 = 2.000000
det A2   : -3.0  -> x2 = 3.000000
det A3   : 1.0  -> x3 = -1.000000
det M    : 60.0 60.000000000000036
det H4   : 1.6534391534390412e-07
cond H4  : 15513.74
det S    : -0.0030000000000036655  cond S: 101072.3
x        : [1. 1. 1.]
x'       : [ 2. -1.  2.]
cost     : 5748019200 333333333
\end{verbatim}
\end{lab}

Four things must be read from that output.

First, $x=(2,3,-1)$ appears with residual $1.99\times10^{-15}$, which is zero to machine
precision. The residual is computed on purpose: a library answer is never exact, and the residual
is the cheapest way to check whether that answer satisfies its system. For a bad matrix, a small
residual does not guarantee a small error; that is the topic of Chapter 22.

Second, the determinant from the cofactor expansion and the one from the library are both $-1$,
and for $M$ both are $60$. The library does not use cofactors; it uses a factorization and
multiplies the pivots. The difference of $3.6\times10^{-14}$ on $\det M$ is rounding, not a
difference of definition.

Third, $\det H_{4}$ is reported as $1.65\times10^{-7}$ while its exact value is
$1/6\,048\,000$, so the machine number is correct to seven digits. For higher orders, $1/\det$
passes $10^{12}$ quickly, and that is where the determinant stops being useful as a marker.

Fourth, the jump of the answer for $S$ shows up intact: from $(1,1,1)$ to $(2,-1,2)$ only
because one number changed by a thousandth. The number reported for $\det S$ is not exactly
$-0.003$ either; the difference $3.7\times10^{-15}$ appears because $9.001$ is not
represented exactly in binary fractions.

Finally, cost: $5.748\times10^{9}$ multiplications for Cramer's rule at $n=12$ against
$3.33\times10^{8}$ for elimination at $n=1000$. A far larger matrix is handled with far less
work; that is the difference between a cube and a factorial.

\begin{untukdosen}
A teaching order that has worked: elimination first, with the three cases shown on one fixed
coefficient matrix and three different right-hand sides. Students who see consistency change
because of one number in $\mathbf{b}$ usually accept at once that the question about $A$ and the
question about $\mathbf{b}$ must be separated. Only after that does the determinant enter, as the
answer to the question about $A$ alone.

Two typical misconceptions. First, ``a small determinant means the system is unstable''. The
matrix $10^{-6}I$ has determinant $10^{-18}$ and condition number one; conversely $H_{4}$ has
determinant $1.65\times10^{-7}$ and condition $1.5\times10^{4}$. This mistake appears every
year, and the way to dismantle it is to ask students to multiply one matrix by a scalar and then
compare the two numbers. Second, ``Cramer's rule is the way to solve a system''. Theoretically
true; in cost, false. Give one $12\times12$ system and ask them to estimate the cost of both
ways before computing.

A useful class activity: hand out three systems with the same coefficient matrix and right-hand
sides $(4,9,13)$, $(4,9,14)$, and $(0,0,0)$. Ask each group to determine the fate of its system
from the echelon form alone, then to conclude what can and cannot be concluded from $A$ by
itself.

A sign of readiness for advanced work: students can list the seven equivalent conditions
without memorizing the order, and can explain why the condition about the determinant belongs in
the same list as the conditions about pivots.
\end{untukdosen}

\section{Summary}
\tingkat{basic}

\begin{ringkasan}
\begin{itemize}
  \item Elementary row operations do not change the solution set, because every operation can be
        undone. Because the reduced form is unique, the pivot positions determine three fates:
        inconsistent if a pivot appears in the last column, unique if every coefficient column
        has a pivot, many if some column has no pivot.
  \item The determinant is the only multilinear, alternating, normalized function. The Leibniz
        formula and the cofactor expansion are consequences of it, not its definition.
  \item $\det(AB)=\det A\det B$, so $A$ is invertible exactly when $\det A\neq0$; and $\det A$
        equals the product of the pivots after elimination, times the sign of the interchanges.
  \item The adjugate gives $A\operatorname{adj}A=(\det A)I$ and
        $A^{-1}=\operatorname{adj}A/\det A$; Cramer's rule writes $x_{j}=\det A_{j}/\det A$.
  \item Seven conditions for invertibility are equivalent: trivial kernel, reduced form $I_{n}$,
        product of elementary matrices, consistent and unique for all $\mathbf{b}$, and
        $\det A\neq0$.
  \item The cost of elimination is proportional to $n^{3}/3$, the cost of cofactors to
        $n\cdot n!$. Sensitivity is measured by $\kappa(A)=\lVert A\rVert\lVert A^{-1}\rVert$,
        not by $\det A$.
\end{itemize}
\end{ringkasan}

\section{Exercises}
\tingkat{intermediate}

\begin{latihan}
\begin{enumerate}
  \item Solve by elimination, then check with Cramer's rule:
        \[
          \begin{pmatrix}1&2&3\\ 2&5&3\\ 1&0&8\end{pmatrix}\mathbf{x}
          =\begin{pmatrix}9\\ 8\\ 26\end{pmatrix}.
        \]
        Write down the pivots that appear, record how many times rows were interchanged, then
        conclude $\det A$ from the product of the pivots.

  \item Compute $\det M$ for
        $M=\begin{pmatrix}0&1&2&3\\ 1&0&3&2\\ 2&3&0&1\\ 3&2&1&0\end{pmatrix}$
        by cofactor expansion along one row, then compute it with
        \texttt{numpy.linalg.det}. The two results are not exactly equal: the first is exact, the
        second $5.3\times10^{-15}$. Explain why. Afterwards find a nonzero vector
        $\mathbf{v}$ with $M\mathbf{v}=\mathbf{0}$ and conclude $\det M=0$ without any
        expansion.

  \item Prove that if two columns of $A$ are equal, then $\det A=0$, without using
        $\det A=\det A^{\top}$.
        (Hint: use the Leibniz formula, group the terms by pairs of permutations that swap on
        those two columns, and show that their signs are opposite.)

  \item Prove that $\det(\operatorname{adj}A)=(\det A)^{n-1}$ for $n\ge2$.
        (Hint: take determinants in $A\operatorname{adj}A=(\det A)I$ and use
        $\det(cI)=c^{n}$.)

  \item Evaluate this claim: \emph{if $\det A$ is very small, then the answer of $A\mathbf{x}
        =\mathbf{b}$ is very sensitive to changes in $\mathbf{b}$}. Decide true, false, or not
        yet enough evidence, and include a counterexample or an additional reason that makes it
        true.

  \item For $A=\begin{pmatrix}1&1\\ 1&1+\varepsilon\end{pmatrix}$ with
        $\varepsilon>0$, compute $\kappa(A)$ in the $1$-norm as a function of
        $\varepsilon$. Determine the value of $\varepsilon$ that makes $\kappa(A)>10^{6}$, then
        check it with \texttt{numpy.linalg.cond}.

  \item Prove: if $A$ is invertible and $B$ is obtained from $A$ by adding a multiple of one row
        to another row, then $\det B=\det A$ and $B^{-1}$ can be expressed through $A^{-1}$ and
        one elementary matrix.

  \item Suppose $D=\operatorname{diag}(d_{1},\dots,d_{n})$ and $A$ is of size $n\times n$.
        Show that $\det(A+AD)=\det(A)\prod_{i=1}^{n}(1+d_{i})$, then determine the condition on
        $d_{1},\dots,d_{n}$ for $A+AD$ to be invertible. Check your answer with a $3\times3$
        matrix in Python.

  \item Choose a $3\times3$ matrix with small integer entries and $\det A=1$, compute $A^{-1}$
        through the adjugate and through \texttt{numpy.linalg.inv}, then compare the two. Repeat
        for a matrix with $\det A=7$ and explain why the entries of $A^{-1}$ become fractions.
\end{enumerate}

Hints for the three hardest items. Item 4: take determinants of both sides, remember
$\det(cI)=c^{n}$, and note that $\det A$ may be zero, so the conclusion about
$\operatorname{adj}A$ must be read as a matrix, not as an inverse. Item 6: in the $1$-norm,
$\lVert A\rVert_{1}$ is the largest column sum, and for this matrix the second column is the
largest; compute $\lVert A^{-1}\rVert_{1}$ from the $2\times2$ formula and note that the answer
is proportional to $1/\varepsilon$. Item 8: write $A+AD$ as $A(I+D)$ before pulling out the
determinant, and remember that $I+D$ is diagonal, so its determinant is the product of its
diagonal entries.
\end{latihan}

\begin{untukpeneliti}
What is really proved in this chapter, for the field $\mathbb{R}$ and finite dimension: the
uniqueness of the determinant as a normalized alternating $n$-linear function together with the
permutation formula, in Hoffman and Kunze, \emph{Linear Algebra} 2nd edition, sections 5.2 and
5.3, which hold more generally over a commutative ring with identity (section 5.1);
multiplicativity and the equivalence of invertibility with $\det\neq0$, in Axler,
\emph{Linear Algebra Done Right} 4th edition, section 9C, statements 9.49 and 9.50, with a
basis-free determinant of an operator; and the list of seven equivalent conditions, in Anton
11th edition, Theorem 2.3.8.

What is popularized: ``the determinant is volume''. What is proved is that $|\det A|$ equals the
volume of the parallelepiped spanned by its columns; without the absolute value, the claim is
false for a negative determinant. Also wrong numerically: Cramer's rule as a way to solve
systems above $n=11$, where the cost is factorial.

An open question. The sensitivity bound in the theorem above is normwise and worst-case, often
too pessimistic compared with a componentwise analysis. More fundamentally: the maximum growth
factor of Gaussian elimination with partial pivoting is not known for general $n$. The proved
bound is $2^{n-1}$, while numerical experiments give far less, and that gap is open.

Limits of validity. The argument collapses when the field is not a field, because modules do not
always have a basis and elimination cannot always divide. For very large $n$, the determinant
overflows or underflows long before the rank is lost, so libraries report the logarithm of the
determinant. The sensitivity statements also depend on the norm, and no choice is canonical.
\end{untukpeneliti}

\peta{This chapter follows Anton, \emph{Elementary Linear Algebra} 11th edition, sections 1.1
through 1.6 (linear systems, Gauss-Jordan elimination, inverses, elementary matrices) and 2.1
through 2.3 (cofactors, the determinant through row reduction, Cramer's rule, including Theorem
2.3.8 on equivalent conditions). The axiomatic treatment follows Hoffman and Kunze,
\emph{Linear Algebra} 2nd edition, sections 5.2 and 5.3, with section 1.4 for the uniqueness of
the reduced echelon form. The basis-free view is in Axler, \emph{Linear Algebra Done Right} 4th
edition, section 9C. For elimination as factorization and the view $A=CR$, see Strang,
\emph{Introduction to Linear Algebra} 6th edition, Chapter 2 and the note ``Elimination and
Factorization''. The numerical side, namely condition numbers, rounding error, and cost, is
treated in Chapter 22 of this book.}

\chapter{Eigenvalues and Diagonalization}

Every minute, a network monitor splits all incoming requests into three buckets: host A, host B,
and host C. From old logs it builds a transition rule. Of the requests heading to A, 80 percent
stay at A, 10 percent move to B, and 10 percent move to C. From B, 30 percent move to A, 60
percent stay, and 10 percent move to C. From C, 20 percent move to A, 30 percent to B, and 50
percent stay. If the share of requests in the three buckets at minute $k$ is written as the
column $\mathbf{x}_{k}$, that rule takes the form
$\mathbf{x}_{k+1}=M\mathbf{x}_{k}$ with
\[
  M=\begin{pmatrix}0.8&0.3&0.2\\ 0.1&0.6&0.3\\ 0.1&0.1&0.5\end{pmatrix}.
\]

Starting from $\mathbf{x}_{0}=(1,0,0)$, that is, all requests initially go to A, the monitor
computes ten steps and obtains
$\mathbf{x}_{10}=(0.5672,\;0.2661,\;0.1666)$. Thirty steps later the numbers are
$(0.5667,\;0.2667,\;0.1667)$. The three entries still sum to one, and on the second reading the
numbers stop moving at four decimal places.

The question worth asking is not what the result is after forty minutes; repeated multiplication
has already answered that for one starting point. The more useful question is: is there a vector
that $M$ only enlarges without changing its direction? Here that direction is
$(17/30,\;4/15,\;1/6)=(0.5667,\;0.2667,\;0.1667)$, exactly the numbers that appear after forty
steps, and the multiplier is $1$.

Chapter 8 supplied the determinant: $\det M=1/5$, so $M$ is invertible and every system
$M\mathbf{y}=\mathbf{b}$ has a unique solution. That answers whether the process can be reversed,
not which way it is heading, and not how fast. The number that governs the rate here is $1/2$: the
largest eigenvalue after $1$. Since $(1/2)^{40}\approx9.1\times10^{-13}$, the initial error stays
behind the first six digits, and that is why the second reading looks still.

This chapter prepares the tools for reading those two numbers straight from the matrix: how to
compute them, the condition under which the preserved directions are numerous enough to form a
basis, what changes when the basis is changed, and then three uses: matrix powers, linear
differential systems, and Markov chains.

\begin{untukanak}
Take a sheet of rubber marked with a square grid, then pull it from two sides. The squares turn
into diamonds, and almost every line becomes slanted. But one line stays straight and stays
parallel to itself: the line along the pull. Its length changes, its direction does not. Sometimes
there is also another line that does not slant but reverses its direction, like an arrow turned
half a circle.

If we are told only how many times each of those special lines stretches, we already know the
whole pull, without having to draw the grid again. That is a very economical method: three or four
numbers replace thousands of points on the sheet.

The trouble is that not every pull is that tidy. Try twisting the sheet slightly while pulling.
Now no line is really straight any longer, and the economical method fails completely. Some pulls
can be described entirely by a few stretch numbers, others cannot.

The question: from the picture of the pull alone, can we guess whether there are special lines
that do not slant, and how many? In this chapter we will learn why some pulls have such lines and
some do not, why some pulls can be replaced entirely by just a few numbers, and why a computing
engine sometimes guesses wrong on pulls that are almost a twist.
\end{untukanak}

\section{Directions that do not rotate}
\tingkat{basic}

A square matrix acts on vectors: it rotates, stretches, or flattens. Almost every vector comes out
with a direction different from the one it went in with. A small minority do not. A vector that
comes out in the same direction it went in, only with a different length, marks a direction that
the matrix preserves. It is that kind of direction that we name first, before any computing tool
exists.

\kamus{eigenvalue}{nilai eigen}{The multiplier along a direction that is not rotated: the direction
stays the same, its length is multiplied by that number.}

\begin{definisi}[Eigenvalues and eigenvectors]
Let $A$ be an $n\times n$ matrix with real entries. A scalar $\lambda$ is called an
\emph{eigenvalue} of $A$ if there is a nonzero vector $\mathbf{v}\in\mathbb{R}^{n}$ with
$A\mathbf{v}=\lambda\mathbf{v}$. The vector $\mathbf{v}$ is called an \emph{eigenvector} for
$\lambda$.\index{eigenvalue}\index{eigenvector}
\end{definisi}

The condition $\mathbf{v}\neq\mathbf{0}$ is not decoration. Without it, $\mathbf{0}$ satisfies
$A\mathbf{0}=\lambda\mathbf{0}$ for every $\lambda$, so the term loses its meaning. What we are
looking for is precisely a direction, and the zero point has no direction.

The equation $A\mathbf{v}=\lambda\mathbf{v}$ can be rewritten as
$(A-\lambda I)\mathbf{v}=\mathbf{0}$. So $\lambda$ is an eigenvalue exactly when the matrix
$A-\lambda I$ has a nontrivial kernel, that is, exactly when $A-\lambda I$ is singular. A question
about eigenvalues can therefore always be translated into a question about rank deficiency, and
Chapter 8 already supplied the tool for that.

Before computing, one important warning. Not every real matrix has real eigenvalues. The rotation
matrix
\[
  R=\begin{pmatrix}0&-1\\ 1&0\end{pmatrix}
\]
sends $(1,0)$ to $(0,1)$. Every real vector rotates ninety degrees, so no line is preserved;
$R^{2}=-I$ turns half a circle after two steps. The absence of real eigenvalues here is not a
failure of computation but a geometric statement.

To find eigenvalues without guessing, we use the determinant: the eigenvalues are the numbers that
make the matrix lose its inverse once that number is subtracted from the diagonal.

\kamus{characteristic polynomial}{polinomial karakteristik}{The list of numbers that make the
matrix singular, written as a polynomial in one variable.}

\begin{definisi}[Characteristic polynomial]
For an $n\times n$ matrix, the \emph{characteristic polynomial} is defined as
$p_{A}(\lambda)=\det(\lambda I-A)$.\index{characteristic polynomial}
\end{definisi}

We use $\lambda I-A$, not $A-\lambda I$. The two differ by a factor $(-1)^{n}$, so their roots are
the same, but the first choice makes the polynomial monic, and that makes the other coefficients
easier to read, as the next corollary shows.

\begin{teorema}
A scalar $\lambda$ is an eigenvalue of $A$ if and only if $p_{A}(\lambda)=0$.
\end{teorema}

\begin{proof}
Suppose $A\mathbf{v}=\lambda\mathbf{v}$ with $\mathbf{v}\neq\mathbf{0}$. Then
$(A-\lambda I)\mathbf{v}=\mathbf{0}$, so $A-\lambda I$ has a nontrivial kernel and is therefore
singular. By Chapter 8, singularity is equivalent to a zero determinant, so
$\det(A-\lambda I)=0$, which differs by a factor $(-1)^{n}$ from $\det(\lambda I-A)=0$.

Conversely, if $\det(\lambda I-A)=0$ then $\lambda I-A$ is singular. Its kernel is nontrivial, so
there is $\mathbf{v}\neq\mathbf{0}$ with $(\lambda I-A)\mathbf{v}=\mathbf{0}$, that is,
$A\mathbf{v}=\lambda\mathbf{v}$.
\end{proof}

The polynomial $p_{A}$ has degree $n$ and is monic, so it has at most $n$ distinct roots. Over
$\mathbb{C}$ it always factors into $n$ linear factors, so every complex square matrix has exactly
$n$ eigenvalues when counted with multiplicity. Over $\mathbb{R}$ that can fail:
$p_{R}(\lambda)=\lambda^{2}+1$ has no real root, but has a pair of complex roots $\pm i$.

The top two coefficients of the characteristic polynomial carry two quantities we already know.

\begin{akibat}
With eigenvalues counted together with their multiplicities,
$\operatorname{tr}(A)=\sum_{i}\lambda_{i}$ and $\det(A)=\prod_{i}\lambda_{i}$.
\index{trace}
\end{akibat}

\begin{proof}
The coefficient of $\lambda^{n-1}$ in $\det(\lambda I-A)$ is $-\operatorname{tr}(A)$. That term is
produced only by the identity permutation in the Leibniz formula, because every other permutation
loses at least two diagonal entries. The constant term is
$p_{A}(0)=\det(-A)=(-1)^{n}\det(A)$. Comparing those two coefficients with the form
$\prod_{i}(\lambda-\lambda_{i})=\lambda^{n}-(\sum_{i}\lambda_{i})\lambda^{n-1}+\dots+(-1)^{n}
\prod_{i}\lambda_{i}$ gives both identities.
\end{proof}

That corollary gives a cheap check: after computing eigenvalues with great effort, their sum and
their product must match the trace and the determinant. If they do not, there is an arithmetic
mistake.

\begin{bacakeras}
Two voices, taking turns.

First voice: what does a matrix do?\\
Second voice: it rotates, stretches, flattens.\\
First voice: to whom?\\
Second voice: to almost every vector.\\
First voice: and then some do not rotate?\\
Second voice: some do not. One direction. Sometimes two. Never more than the number of columns.\\
First voice: does the length change?\\
Second voice: it changes. The direction stays.\\
First voice: by how much?\\
Second voice: that is the number. That number is the eigenvalue.\\
First voice: and the direction?\\
Second voice: that is the eigenvector. One direction, one number.\\
First voice: what if the direction flips?\\
Second voice: the number is negative. The line stays the same, it only faces backwards.\\
First voice: what if the number is zero?\\
Second voice: that direction disappears, and the matrix loses its inverse.

Together: a direction that does not change stores a number that changes its length.

Repeat once in a normal voice: a direction that does not change stores a number that changes its
length.
\end{bacakeras}

\begin{resep}
For $A=\begin{pmatrix}a&b\\c&d\end{pmatrix}$ whose entries are already known:
\begin{enumerate}
  \item Subtract $\lambda$ from both diagonal entries, leaving $A-\lambda I$.
  \item Compute $\det(A-\lambda I)$, or use $\lambda^{2}-(a+d)\lambda+(ad-bc)$ directly.
  \item Solve that quadratic equation; check that the sum of the roots is $a+d$ and the product is
        $ad-bc$.
  \item For each root $\lambda$, solve $(A-\lambda I)\mathbf{v}=\mathbf{0}$ and take one nonzero
        solution.
\end{enumerate}
Steps 3 and 4 must not be swapped: an eigenvector only has meaning after $\lambda$ is known, and
the column order of $P$ must follow the diagonal order of $D$.
\end{resep}

Now one example by hand. Take
\[
  A=\begin{pmatrix}4&2\\ 1&3\end{pmatrix},\qquad
  p_{A}(\lambda)=\lambda^{2}-7\lambda+10=(\lambda-5)(\lambda-2).
\]
The trace is $7=5+2$ and the determinant is $10=5\cdot2$, which matches. For $\lambda=5$ we solve
$(A-5I)\mathbf{v}=\mathbf{0}$, that is $-v_{1}+2v_{2}=0$ and $v_{1}-2v_{2}=0$; both rows say the
same thing, so $\mathbf{v}=(2,1)$. For $\lambda=2$ we obtain $2v_{1}+2v_{2}=0$, so
$\mathbf{v}=(1,-1)$. Direct check:
\[
  A\begin{pmatrix}2\\ 1\end{pmatrix}=\begin{pmatrix}10\\ 5\end{pmatrix}=5
  \begin{pmatrix}2\\ 1\end{pmatrix},\qquad
  A\begin{pmatrix}1\\ -1\end{pmatrix}=\begin{pmatrix}2\\ -2\end{pmatrix}=2
  \begin{pmatrix}1\\ -1\end{pmatrix}.
\]
Both of those checks can be seen in the plane. Figure~\ref{fig:panah-eigen} draws both eigen
directions together with their images, and one vector that is not an eigenvector as a comparison.

\begin{figure}[htbp]
\centering
\begin{tikzpicture}[x=0.4cm,y=0.4cm,>=latex]
  \draw[->,layergrey!60] (-0.8,0) -- (11.4,0) node[right,font=\scriptsize] {$x$};
  \draw[->,layergrey!60] (0,-2.9) -- (0,6.4) node[above,font=\scriptsize] {$y$};
  \draw[->,thick,layergrey] (0,0) -- (1,0);
  \draw[->,thick,dashed,layergrey] (0,0) -- (4,1)
        node[right,font=\scriptsize] {$(4,1)$};
  \draw[->,very thick,layerteal] (0,0) -- (2,1)
        node[above left,font=\scriptsize] {$(2,1)$};
  \fill[layerteal] (2,1) circle (2pt);
  \draw[->,very thick,accent] (0,0) -- (10,5)
        node[above,font=\scriptsize] {$(10,5)=5\,(2,1)$};
  \draw[->,very thick,layerindigo] (0,0) -- (1,-1)
        node[below left,font=\scriptsize] {$(1,-1)$};
  \fill[layerindigo] (1,-1) circle (2pt);
  \draw[->,very thick,layerindigo] (0,0) -- (2,-2)
        node[below right,font=\scriptsize] {$(2,-2)=2\,(1,-1)$};
\end{tikzpicture}
\caption{Eigen directions of the matrix $A=\begin{pmatrix}4&2\\1&3\end{pmatrix}$. The vector
$(2,1)$ comes out as $(10,5)$, that is, $5$ times as long without rotating; the vector $(1,-1)$
comes out as $(2,-2)$, that is, $2$ times as long. The vector $(1,0)$ comes out as $(4,1)$ and
rotates, so it is not an eigenvector.}
\label{fig:panah-eigen}
\end{figure}

An eigenvector is determined only up to a multiplier: $(4,2)$ is also an eigenvector for $5$. What
is unique is the direction, not the length. We will use that freedom later to tidy up the matrix.

For a $2\times2$ matrix the roots can be written down directly. Since
$p_{A}(\lambda)=\lambda^{2}-\operatorname{tr}(A)\lambda+\det(A)$, the quadratic formula gives
\[
  \lambda_{1,2}=\frac{\operatorname{tr}(A)\pm\sqrt{\operatorname{tr}(A)^{2}-4\det(A)}}{2},
\]
so the sign of $\Delta=\operatorname{tr}(A)^{2}-4\det(A)$ determines the whole spectrum. For $A$
above, $\Delta=49-40=9$, and the roots are $5$ and $2$. If $\Delta=0$ the two roots coincide, and
if $\Delta<0$ there is no real root: the rotation matrix above gives $\Delta=0-4=-4$.

\begin{figure}[htbp]
\centering
\begin{tikzpicture}[x=0.85cm,y=0.85cm,>=latex,font=\scriptsize]
  \draw[->,layergrey!60] (-1.2,0.35) -- (6.6,0.35);
  \foreach \x in {0,1,2,3,4,5,6} \draw[layergrey!50] (\x,0.26) -- (\x,0.44);
  \fill[layerteal] (2,0.35) circle (2.2pt);
  \fill[layerteal] (5,0.35) circle (2.2pt);
  \node[above=2pt] at (2,0.35) {$2$};
  \node[above=2pt] at (5,0.35) {$5$};
  \node[left,align=right,text width=3.0cm] at (-1.3,0.35) {$\Delta=9$: two real roots};

  \draw[->,layergrey!60] (-1.2,-1.25) -- (6.6,-1.25);
  \foreach \x in {0,1,2,3,4,5,6} \draw[layergrey!50] (\x,-1.34) -- (\x,-1.16);
  \fill[layerorange] (1,-1.25) circle (2.2pt);
  \node[above=2pt] at (1,-1.2) {$1$};
  \node[left,align=right,text width=3.0cm] at (-1.3,-1.25) {$\Delta=0$: one double root};

  \draw[->,layergrey!60] (-1.2,-2.85) -- (6.6,-2.85);
  \foreach \x in {0,1,2,3,4,5,6} \draw[layergrey!50] (\x,-2.94) -- (\x,-2.76);
  \draw[layergrey!80] (0,-3.55) -- (0,-2.15);
  \draw[accent,fill=white,thick] (0,-2.3) circle (2.4pt);
  \draw[accent,fill=white,thick] (0,-3.4) circle (2.4pt);
  \node[left=2pt,accent] at (0,-2.3) {$i$};
  \node[left=2pt,accent] at (0,-3.4) {$-i$};
  \node[left,align=right,text width=3.0cm] at (-1.3,-2.85) {$\Delta=-4$: no real roots};
\end{tikzpicture}
\caption{The spectrum of a $2\times2$ matrix on the number line, separated by the sign of
$\Delta=\operatorname{tr}(A)^{2}-4\det(A)$. Top row: $A$ with $\Delta=49-40=9$ gives the roots
$2$ and $5$. Middle row: a shear matrix with $\Delta=0$ gives the root $1$ counted twice. Bottom
row: the rotation matrix with $\Delta=0-4=-4$ does not meet the real line, and both of its roots,
$\pm i$, lie off that line.}
\label{fig:spektrum}
\end{figure}

Figure~\ref{fig:spektrum} arranges those three possibilities on a single number line: two points,
one point counted twice, or no point at all. The last is not a shortcoming of the picture but a
statement that the real line contains no preserved direction.

\begin{tangan}
Materials: one rubber band, two pencils, a sheet of paper, a ruler.

Place the two pencils parallel, then stretch the rubber band around both so that a rectangle
forms. On the paper, draw three arrows from the lower left corner of that rectangle, each at an
angle of 30, 45, and 60 degrees to the bottom side. Label all three.

Slide the right pencil twice as far, so that the rubber stretches. Compare each arrow with its
original length. The arrow parallel to the pull side keeps its direction; the other two slant, and
the one that slants most is the one farthest from the pull side.

Now rotate the whole frame ninety degrees and repeat. The surviving direction rotates along with
the rubber, because what is special is the direction relative to the rubber, not relative to the
table.

What should be visible: one direction survives, two directions rotate, and the stretch factor
along the surviving direction is the multiplier we call the eigenvalue.
\end{tangan}

\section{Eigen spaces and two multiplicities}
\tingkat{intermediate}

One eigenvalue usually comes with more than one eigenvector. The set of all eigenvectors for one
eigenvalue, plus the zero vector, is not a random collection: it is a subspace. Its dimension will
later decide whether the matrix can be simplified.

\begin{definisi}[Eigen space]
For an eigenvalue $\lambda$ of $A$, the \emph{eigen space} is
$E_{\lambda}=\operatorname{null}(A-\lambda I)=\{\mathbf{v}:A\mathbf{v}=\lambda\mathbf{v}\}$.
\index{eigen space}
\end{definisi}

\kamus{eigen space}{ruang eigen}{All directions that the matrix does not rotate, gathered into a
single subspace.}

An eigen space is always a subspace because it is the kernel of a matrix (Chapter 7). Its
dimension is at least one exactly when $\lambda$ really is an eigenvalue.

Eigen spaces for different eigenvalues intersect only at the zero point.

\begin{proposisi}
If $\lambda\neq\mu$, then $E_{\lambda}\cap E_{\mu}=\{\mathbf{0}\}$.
\end{proposisi}

\begin{proof}
Take $\mathbf{v}$ in that intersection. Then $A\mathbf{v}=\lambda\mathbf{v}$ and
$A\mathbf{v}=\mu\mathbf{v}$, so $(\lambda-\mu)\mathbf{v}=\mathbf{0}$. Since
$\lambda-\mu\neq0$, we must have $\mathbf{v}=\mathbf{0}$.
\end{proof}

That property extends to any finite number of eigenvectors, by induction. The result is the lemma
that forms the backbone of this chapter.

\begin{lemma}
Eigenvectors corresponding to distinct eigenvalues are linearly independent of one another.
\end{lemma}

\begin{proof}
We prove this by induction on the number of vectors. For one vector the statement is true because
an eigenvector is nonzero. Suppose the statement holds for $m-1$ vectors, and let
\[
  c_{1}\mathbf{v}_{1}+\dots+c_{m}\mathbf{v}_{m}=\mathbf{0}
\]
with $\lambda_{1},\dots,\lambda_{m}$ distinct. Multiply that equation by $A$, and use
$A\mathbf{v}_{i}=\lambda_{i}\mathbf{v}_{i}$:
\[
  c_{1}\lambda_{1}\mathbf{v}_{1}+\dots+c_{m}\lambda_{m}\mathbf{v}_{m}=\mathbf{0}.
\]
Subtract $\lambda_{m}$ times the first equation from the second:
\[
  c_{1}(\lambda_{1}-\lambda_{m})\mathbf{v}_{1}+\dots
  +c_{m-1}(\lambda_{m-1}-\lambda_{m})\mathbf{v}_{m-1}=\mathbf{0}.
\]
By the induction hypothesis, $\mathbf{v}_{1},\dots,\mathbf{v}_{m-1}$ are independent, so
$c_{i}(\lambda_{i}-\lambda_{m})=0$ for $i<m$. Since $\lambda_{i}\neq\lambda_{m}$, all $c_{i}=0$
for $i<m$. The original equation is left with $c_{m}\mathbf{v}_{m}=\mathbf{0}$, so $c_{m}=0$.
\end{proof}

That lemma gives a tidy upper bound: an $n\times n$ matrix cannot have more than $n$ independent
eigenvectors. When is that bound attained?

To answer, we need two ways of counting how often one eigenvalue occurs: from the polynomial and
from the dimension of the eigen space. The two often differ, and it is the difference that decides.

\kamus{multiplicity}{multiplisitas}{How many times an eigenvalue occurs, counted in two ways that
are not always equal.}

\begin{definisi}[Algebraic and geometric multiplicity]
Let $\lambda$ be an eigenvalue of $A$. The \emph{algebraic multiplicity} of $\lambda$ is the power
of the factor $(\lambda'-\lambda)$ in $p_{A}(\lambda')$. The \emph{geometric multiplicity} of
$\lambda$ is $\dim E_{\lambda}=\operatorname{nullity}(A-\lambda I)$.\index{multiplicity}
\end{definisi}

Geometric multiplicity is easy to compute: subtract $\lambda$ from the diagonal, then compute the
nullity of the resulting matrix. Algebraic multiplicity demands factoring a polynomial. The
relation between the two always runs in one direction.

\begin{lemma}
For every eigenvalue $\lambda$ we have
$1\le\dim E_{\lambda}\le$ the algebraic multiplicity of $\lambda$.
\end{lemma}

\begin{proof}
The left inequality is true because $\lambda$ is an eigenvalue, so $E_{\lambda}$ contains a
nonzero vector.

For the right inequality, write $k=\dim E_{\lambda}$ and choose a basis
$\mathbf{u}_{1},\dots,\mathbf{u}_{k}$ of $E_{\lambda}$. Extend it to a basis
$\mathbf{u}_{1},\dots,\mathbf{u}_{n}$ of $\mathbb{R}^{n}$, which is possible by Chapter 7. In that
basis, column $j$ of the matrix $A$ for $j\le k$ is $A\mathbf{u}_{j}=\lambda\mathbf{u}_{j}$, that
is, $\lambda$ in the diagonal position and zero everywhere else. So the matrix $A$ in that basis
has the form
\[
  \begin{pmatrix}\lambda I_{k}&C\\ 0&D\end{pmatrix}
\]
with $C$ of size $k\times(n-k)$ and $D$ of size $(n-k)\times(n-k)$. The determinant of
$\lambda' I-A$ in that basis, either through cofactor expansion along the first $k$ columns or
through the determinant formula for a block triangular matrix, equals
$(\lambda'-\lambda)^{k}\det(\lambda' I-D)$. So $(\lambda'-\lambda)^{k}$ divides $p_{A}$, hence the
algebraic multiplicity of $\lambda$ is at least $k$.
\end{proof}

Now look at both possibilities in a single matrix. Take
\[
  A=\begin{pmatrix}2&1&1\\ 1&2&1\\ 1&1&2\end{pmatrix},\qquad
  p_{A}(\lambda)=(\lambda-1)^{2}(\lambda-4).
\]
The eigenvalues are $1$ and $4$, with algebraic multiplicities $2$ and $1$. The trace is
$6=4+1+1$ and the determinant is $4=4\cdot1\cdot1$, which matches. For $\lambda=4$, the equation
$(A-4I)\mathbf{v}=\mathbf{0}$ gives $v_{1}=v_{2}=v_{3}$, so $E_{4}$ is spanned by $(1,1,1)$. For
$\lambda=1$, the matrix $A-I$ is the matrix all of whose entries are $1$, so its kernel is the set
of vectors whose entries sum to zero. That set is two-dimensional, with basis $(1,-1,0)$ and
$(1,0,-1)$. So the geometric multiplicity of $1$ is $2$, equal to its algebraic multiplicity.

A repeated eigenvalue is not automatically a problem; what becomes a problem is an eigen space
that is too small, and that happens with the following matrix.

\section{When a matrix can be diagonalized}
\tingkat{intermediate}

If we may choose the basis, we want the matrix as simple as possible: diagonal, so that every
basis vector is merely scaled and not mixed with the others. Is there always such a basis?

\begin{definisi}[Diagonalizable]
An $n\times n$ matrix is called \emph{diagonalizable} if there is an invertible matrix $P$ and a
diagonal matrix $D$ with $P^{-1}AP=D$, or equivalently $A=PDP^{-1}$.
\end{definisi}

\kamus{diagonalization}{diagonalisasi}{Changing the measuring instrument until the whole work of
the matrix is a few multiplications of numbers.}

The columns of $P$ will contain eigenvectors, and $D$ contains the eigenvalues on its diagonal. A
question about diagonalization is therefore really a question about having enough eigenvectors.

\begin{teorema}
An $n\times n$ matrix is diagonalizable if and only if it has $n$ linearly independent
eigenvectors.\index{diagonalization}
\end{teorema}

\begin{proof}
Suppose $P^{-1}AP=D$ with $D=\operatorname{diag}(d_{1},\dots,d_{n})$. Multiply from the left by
$P$: $AP=PD$. Column $j$ of both sides gives
$A\mathbf{p}_{j}=d_{j}\mathbf{p}_{j}$, with $\mathbf{p}_{j}$ the $j$-th column of $P$. Since $P$ is
invertible, each of its columns is nonzero, so $\mathbf{p}_{j}$ is an eigenvector. The columns of
$P$ are also linearly independent because $P$ is invertible. So $A$ has $n$ independent
eigenvectors.

Conversely, suppose $\mathbf{p}_{1},\dots,\mathbf{p}_{n}$ are independent eigenvectors with
$A\mathbf{p}_{j}=\lambda_{j}\mathbf{p}_{j}$. Assemble $P$ from those vectors as columns and
$D=\operatorname{diag}(\lambda_{1},\dots,\lambda_{n})$. Then $AP=PD$, and since $P$ is invertible
we obtain $P^{-1}AP=D$.
\end{proof}

The corollary most often used in hand computation: enough distinct eigenvalues is already
sufficient.

\begin{akibat}
If an $n\times n$ matrix has $n$ distinct eigenvalues, it is diagonalizable.
\end{akibat}

\begin{proof}
Take one eigenvector for each eigenvalue. By the earlier lemma, $n$ eigenvectors with $n$ distinct
eigenvalues are independent of one another. The theorem above then gives a diagonalization.
\end{proof}

The converse does not hold: the $3\times3$ example above had only two distinct eigenvalues, yet is
still diagonalizable because the eigen space for $\lambda=1$ is two-dimensional. The exact
condition is not the number of eigenvalues but the total dimension of the eigen spaces.

\begin{teorema}[Multiplicity condition]
An $n\times n$ matrix is diagonalizable if and only if for each of its eigenvalues the algebraic
multiplicity equals the geometric multiplicity. That condition is equivalent to
$\sum_{\lambda}\dim E_{\lambda}=n$.
\end{teorema}

\begin{proof}
First, suppose the total dimension of all eigen spaces equals $n$. For each $\lambda$ choose a
basis of $E_{\lambda}$, then combine all those bases into a list $\mathbf{v}_{1},\dots,
\mathbf{v}_{n}$. That list is linearly independent: if some linear combination equals zero, group
the terms according to their eigenvalues to get $\mathbf{u}_{1}+\dots+\mathbf{u}_{m}=\mathbf{0}$
with $\mathbf{u}_{k}\in E_{\lambda_{k}}$. By the lemma on distinct eigenvectors, every
$\mathbf{u}_{k}=\mathbf{0}$; since each $\mathbf{u}_{k}$ is a combination of the basis of
$E_{\lambda_{k}}$, all its coefficients are zero. That list also contains $n$ vectors in a space of
dimension $n$, so it is a basis of eigenvectors, and the previous theorem closes the argument.

Second, suppose $A$ is diagonalizable, so there is a basis of eigenvectors
$\mathbf{p}_{1},\dots,\mathbf{p}_{n}$. In that basis the matrix $A$ is diagonal, so
$p_{A}(\lambda')=\prod_{j}(\lambda'-\lambda_{j})$ and the algebraic multiplicity of each
eigenvalue equals the number of $j$ with $\lambda_{j}=\lambda$. That number of $j$ does not exceed
$\dim E_{\lambda}$, because those vectors lie inside $E_{\lambda}$ and are independent. So the
algebraic multiplicity does not exceed the geometric one. The earlier lemma gives the opposite
direction, so the two are equal. Finally, $\sum_{\lambda}\dim E_{\lambda}=n$ because the basis of
eigenvectors fills the whole space.
\end{proof}

Now an example that fails. Take
\[
  S=\begin{pmatrix}1&1\\ 0&1\end{pmatrix},\qquad p_{S}(\lambda)=(\lambda-1)^{2}.
\]
The only eigenvalue is $1$, with algebraic multiplicity $2$. But
$S-I=\begin{pmatrix}0&1\\ 0&0\end{pmatrix}$ has rank one, so $\dim E_{1}=2-1=1$. Its geometric
multiplicity is one short, and by the theorem above $S$ is not diagonalizable.

Geometrically $S$ is a shear: it preserves the direction $(1,0)$ and slides $(0,1)$ to $(1,1)$. No
basis makes this map a mere scaling, because a scaling never slides something that was previously
parallel. That failure is a property of the map, not a choice of basis.

\begin{untukanak}
Have you ever tried to tidy something that turned out not to be tidyable? Take a stack of cards,
then slide the top part slightly sideways. The picture on each card stays the same, but every
point shifts by its distance from the base. There is no way to rotate that stack so that the shift
becomes a mere lengthening or shortening of something. A shift remains a shift.

There are other objects that are more accommodating. A rubber sheet pulled straight can be
described only by how many times each direction stretches. If we know that number for a few
important directions, we know the whole pull, including for points we never measured.

The question: how do we tell the two apart from the picture alone, without trying every
possibility? In this chapter we will learn why a shear always stores one direction that cannot be
tidied up, why some pulls have enough special directions, and why that difference decides whether
the tenth power from a computing engine can still be trusted.
\end{untukanak}

A shear is easy to recognize because of the triangle. The same failure appears in matrices that do
not look triangular. Take
\[
  B=\begin{pmatrix}-1&4&0\\ -1&3&0\\ 1&-2&2\end{pmatrix},\qquad
  p_{B}(\lambda)=(\lambda-1)^{2}(\lambda-2).
\]
For $\lambda=1$, the matrix $B-I$ has rank two, so $\dim E_{1}=3-2=1$, while its algebraic
multiplicity is $2$. The only eigenvector for $1$ is $(2,1,0)$. For $\lambda=2$ its eigen space is
spanned by $(0,0,1)$ and is one-dimensional, matching its algebraic multiplicity. The total
dimension of all eigen spaces is only $2<3$, so $B$ is not diagonalizable. A matrix of this kind is
called \emph{defective}: it is short of eigenvectors because its eigen spaces are narrow.
\index{defective}

\kamus{defective}{defective}{A matrix whose eigenvalues repeat but whose eigen spaces are too
narrow to form a basis.}

One note before we go on. Whether a matrix is diagonalizable depends on the field of scalars: $R$
is not diagonalizable over $\mathbb{R}$, but over $\mathbb{C}$ it has two distinct eigenvalues,
$\pm i$, so it is diagonalizable. A statement without a field qualifier is therefore incomplete.
For the rest of this chapter the field is $\mathbb{R}$ unless stated otherwise.

\section{Similar matrices and invariants}
\tingkat{advanced}

Diagonalization is a change of basis that produces the simplest possible form. But not every
matrix can reach that form, and the next question is natural: what stays the same across all
matrices connected by a change of basis?\index{invariant}

\begin{definisi}[Similar matrices]
Two $n\times n$ matrices $A$ and $B$ are called \emph{similar} if there is an invertible matrix
$P$ with $B=P^{-1}AP$.\index{similar}
\end{definisi}

\kamus{similar}{serupa}{Two reports about the same map from two different measuring spaces.}

Similarity is an equivalence relation: every matrix is similar to itself, the relation is
symmetric because $(P^{-1})^{-1}BP^{-1}=A$, and transitive because
$(PQ)^{-1}A(PQ)=Q^{-1}(P^{-1}AP)Q$. The classes divide up neatly, and diagonalization means asking
whether a class contains a diagonal matrix.

\begin{proposisi}
If $A$ and $B$ are similar, then $p_{A}=p_{B}$. Consequently $A$ and $B$ have the same eigenvalues
together with their algebraic multiplicities, and also the same trace, determinant, rank, and
nullity.
\end{proposisi}

\begin{proof}
Suppose $B=P^{-1}AP$. Then
\[
  \lambda I-B=\lambda I-P^{-1}AP=P^{-1}(\lambda I-A)P,
\]
so
$\det(\lambda I-B)=\det(P^{-1})\det(\lambda I-A)\det(P)=\det(\lambda I-A)$ because
$\det(P^{-1})\det(P)=1$. So the characteristic polynomials are equal, and all data derived from
their coefficients are equal as well. Equality of rank and nullity follows from $B-\lambda I=
P^{-1}(A-\lambda I)P$: multiplication by an invertible matrix does not change the rank, and it is
that rank which determines the nullity through rank-nullity.
\end{proof}

The converse does not hold, and the counterexample is simple: $I_{2}$ and the shear $S$ have the
same characteristic polynomial, $(\lambda-1)^{2}$, but are not similar. If $P^{-1}I_{2}P=S$, the
left side equals $I_{2}$ for every $P$. So the characteristic polynomial is not a complete
invariant; it only filters.

What changes and what does not can be separated neatly. The eigenvalues do not change. The
eigenvectors do change, and they change in a way that is easy to remember.

\begin{proposisi}
If $B=P^{-1}AP$ and $A\mathbf{v}=\lambda\mathbf{v}$, then
$B(P^{-1}\mathbf{v})=\lambda(P^{-1}\mathbf{v})$.
\end{proposisi}

\begin{proof}
Compute $B(P^{-1}\mathbf{v})=P^{-1}APP^{-1}\mathbf{v}=P^{-1}A\mathbf{v}
=P^{-1}(\lambda\mathbf{v})=\lambda(P^{-1}\mathbf{v})$.
\end{proof}

So eigenvectors move along with their coordinates, while the multipliers stay put. That makes
sense: a preserved direction is a property of the map, while its list of coordinates depends on
the basis. That is why eigenvalues are used as the fingerprint of a matrix, and eigenvectors are
not.

On the other hand, the freedom to choose $P$ explains why two diagonalizations can look different
although both are correct: the columns of $P$ may be multiplied by a nonzero scalar or permuted as
long as $D$ is permuted accordingly. What must be the same is the set of eigenvalues together with
their multiplicities.

\section{A worked example}
\tingkat{intermediate}

Every number in this section is computed by hand first, then checked with Python in the lab
section. We start from the $3\times3$ matrix that already appeared in the discussion of eigen
spaces,
\[
  A=\begin{pmatrix}2&1&1\\ 1&2&1\\ 1&1&2\end{pmatrix}.
\]
Its characteristic polynomial is computed from $\lambda I-A$, whose diagonal entries are all
$\lambda-2$ and whose other entries are all $-1$. Writing $t=\lambda-2$ and expanding along the
first row,
\[
  p_{A}(\lambda)=t\,(t^{2}-1)-(-1)\,(-t-1)+(-1)\,(1+t)=t^{3}-3t-2 .
\]
That form factors because $t=-1$ is a root: $t^{3}-3t-2=(t+1)^{2}(t-2)$. Back to $\lambda$,
\[
  p_{A}(\lambda)=(\lambda-1)^{2}(\lambda-4)=\lambda^{3}-6\lambda^{2}+9\lambda-4 .
\]
The trace is $6$ and the determinant is $4$, matching $4+1+1$ and $4\cdot1\cdot1$.

Eigenvectors for $\lambda=4$: the matrix $A-4I$ has diagonal entries $-2$ and other entries $1$.
Its rows say $-2v_{1}+v_{2}+v_{3}=0$ and so on; all three equations hold exactly when
$v_{1}=v_{2}=v_{3}$. So $E_{4}$ is spanned by $(1,1,1)$.

Eigenvectors for $\lambda=1$: the matrix $A-I$ contains only $1$s, so its kernel is the set of
vectors with $v_{1}+v_{2}+v_{3}=0$. A basis is two vectors, for instance $(1,-1,0)$ and
$(1,0,-1)$. Hence
\[
  P=\begin{pmatrix}1&1&1\\ 1&-1&0\\ 1&0&-1\end{pmatrix},\qquad
  P^{-1}=\frac{1}{3}\begin{pmatrix}1&1&1\\ 1&-2&1\\ 1&1&-2\end{pmatrix},\qquad
  D=\operatorname{diag}(4,1,1).
\]
The determinant of $P$ is $3$, and that is the denominator in $P^{-1}$. We check one entry by
hand: the entry in the first row and first column of $PDP^{-1}$ is
$\tfrac{1}{3}(4\cdot1\cdot1+1\cdot1\cdot1+1\cdot1\cdot1)=2$, matching the entry of $A$. A full
check with SymPy gives $PDP^{-1}-A=0$ exactly.

Now the tenth power.\index{matrix power} Multiplying nine times gives no insight at all;
diagonalization provides it. Since $A=I+J$ with $J$ the matrix all of whose entries are $1$, and
$J^{2}=3J$, every power of $J$ returns as a multiple of $J$:
\[
  J^{2}=3J,\quad J^{3}=9J,\quad\dots,\quad J^{k}=3^{k-1}J .
\]
Then
\[
  A^{k}=(I+J)^{k}=I+\Bigl(\sum_{j=1}^{k}\binom{k}{j}3^{j-1}\Bigr)J
  =I+\frac{4^{k}-1}{3}\,J,
\]
with the last step using the geometric series formula
$\sum_{j=1}^{k}\binom{k}{j}3^{j-1}=\bigl((1+3)^{k}-1\bigr)/3$. For $k=10$ we have
$4^{10}=1048576$ and $(1048576-1)/3=349525$, so that
\[
  A^{10}=I+349525\,J=
  \begin{pmatrix}349526&349525&349525\\ 349525&349526&349525\\ 349525&349525&349526\end{pmatrix}.
\]
Every entry can be checked: $1+349525$ on the diagonal, $349525$ off it. Direct multiplication
with Python gives the same matrix, with largest difference $1.05\times10^{-9}$ because of
floating-point arithmetic.

The form of that formula explains what is not visible in the final result. The term $4^{k}$
dominates the direction $(1,1,1)$; the constant term $1$ survives in the direction whose entries
sum to zero. The eigenvalues $4$ and $1$ are two growth rates, and a matrix power is merely a
mixture of the two. If the matrix is not diagonalizable, that mixture changes shape: a term
$k\lambda^{k}$ appears or, for $\lambda=1$, a pure factor $k$.

The defective example above shows that. For the $B$ that failed to be diagonalizable,
\[
  B^{k}=
  \begin{pmatrix}1-2k&4k&0\\ -k&1+2k&0\\ 2^{k}-1&-2(2^{k}-1)&2^{k}\end{pmatrix},
\]
which can be proved by induction and holds for $k=1$. At $k=10$ that formula gives
\[
  B^{10}=\begin{pmatrix}-19&40&0\\ -10&21&0\\ 1023&-2046&1024\end{pmatrix},
\]
the same as repeated multiplication. The top-left entry, $-19=1-2\cdot10$, grows linearly, not
exponentially: the eigenvalue $1$ controls that direction, and the shortage of eigenvectors adds
one factor of $k$.

Now the first application: a process that settles. For the Markov chain matrix $M$ from the
opening,
\index{Markov chain}
\[
  p_{M}(\lambda)=\lambda^{3}-\tfrac{19}{10}\lambda^{2}+\tfrac{11}{10}\lambda-\tfrac15
  =(\lambda-1)\bigl(\lambda-\tfrac12\bigr)\bigl(\lambda-\tfrac25\bigr).
\]
The trace is $19/10$ and the determinant is $1/5$, matching $1+\tfrac12+\tfrac25$ and
$1\cdot\tfrac12\cdot\tfrac25$. The eigenvector for $\lambda=1$ is found from $M\mathbf{v}=
\mathbf{v}$, that is $(M-I)\mathbf{v}=\mathbf{0}$, together with the condition that the entries sum
to one. The solution gives $\boldsymbol{\pi}=(17/30,\;4/15,\;1/6)$. One coordinate checked by hand:
\[
  \tfrac{8}{10}\cdot\tfrac{17}{30}+\tfrac{3}{10}\cdot\tfrac{8}{30}
  +\tfrac{2}{10}\cdot\tfrac{5}{30}
  =\tfrac{13.6+2.4+1.0}{30}=\tfrac{17}{30}.
\]
So $\boldsymbol{\pi}$ really is a fixed point.

Since $M$ has three distinct eigenvalues, it is diagonalizable. Every starting point therefore
decomposes as
\[
  \mathbf{x}_{k}=\boldsymbol{\pi}+c_{2}\bigl(\tfrac12\bigr)^{k}\mathbf{v}_{2}
  +c_{3}\bigl(\tfrac25\bigr)^{k}\mathbf{v}_{3},
\]
with $c_{2},c_{3}$ depending on $\mathbf{x}_{0}$. The second and third terms shrink, and the rate
is governed by $1/2$. For $\mathbf{x}_{0}=(1,0,0)$ we obtain $c_{2}\mathbf{v}_{2}$ of order
$10^{-1}$, so the error after ten steps is of order
$10^{-1}\cdot(1/2)^{10}\approx10^{-4}$, and after forty steps of order
$(1/2)^{40}\approx9.1\times10^{-13}$. The computed result agrees: $\mathbf{x}_{10}$ differs from
$\boldsymbol{\pi}$ in the fourth digit, while $\mathbf{x}_{40}$ already agrees to six digits.

Second application: differential systems. If $\mathbf{x}(t)$ satisfies
$\mathbf{x}'=A\mathbf{x}$ with $A$ above, the solution is $\mathbf{x}(t)=e^{At}\mathbf{x}(0)$.
Since $A=I+J$ and $J^{2}=3J$,
\[
  e^{At}=e^{t}e^{Jt}=e^{t}\Bigl(I+\frac{e^{3t}-1}{3}J\Bigr)
  =e^{t}I+\frac{e^{4t}-e^{t}}{3}J .
\]
At $t=0$ that formula gives $I$. For $t=0.5$ its off-diagonal entries are
$(e^{2}-e^{0.5})/3=1.9134$, and Python gives the same number to four digits. The component along
$(1,1,1)$ grows like $e^{4t}$; the component in the plane where the entries sum to zero grows like
$e^{t}$, so the solution quickly becomes parallel to $(1,1,1)$. A positive real eigenvalue means
growth; a purely imaginary eigenvalue such as $\pm i$ means oscillation without growth.

Third application: power iteration. Multiplying repeatedly and then normalizing the result,
$\mathbf{x}_{k+1}=A\mathbf{x}_{k}/\lVert A\mathbf{x}_{k}\rVert$, makes the direction turn toward
the dominant eigenvector, because $\lambda_{1}^{k}$ defeats $\lambda_{2}^{k}$ when
$|\lambda_{1}|>|\lambda_{2}|$. For $A$ above starting from $(1,0)$, the angle runs $0^\circ$,
$14.04^\circ$, $21.25^\circ$, $24.39^\circ$, $25.69^\circ$, $26.21^\circ$, and heads toward
$\arctan(1/2)=26.57^\circ$; the Rayleigh quotient $\mathbf{x}^{\top}A\mathbf{x}/\mathbf{x}^{\top}
\mathbf{x}$ runs $4.00$, $4.65$, $4.88$, $4.96$, $4.98$, $4.99$, heading toward $5$.
Figure~\ref{fig:iterasi-pangkat} shows that turning.

\begin{figure}[htbp]
\centering
\begin{tikzpicture}[scale=3.2,>=latex]
  \draw[layergrey!35] (0,0) circle (1);
  \draw[->,layergrey!60] (-1.25,0) -- (1.55,0) node[right,font=\scriptsize] {$x$};
  \draw[->,layergrey!60] (0,-0.35) -- (0,1.2) node[above,font=\scriptsize] {$y$};
  \draw[->,semithick,layergrey] (0,0) -- (0:1);
  \draw[->,thick,layerindigo!45] (0,0) -- (14.04:1);
  \draw[->,thick,layerindigo!70] (0,0) -- (21.25:1);
  \draw[->,very thick,layerindigo] (0,0) -- (24.39:1);
  \foreach \a in {0,14.04,21.25,24.39} \fill[layerindigo!70] (\a:1) circle (1.3pt);
  \draw[->,very thick,dashed,accent] (0,0) -- (26.565:1.42);
  \node[right,font=\scriptsize,accent,align=left] at (26.565:1.42)
        {$(2,1)$\\[1pt]$26.57^\circ$};
\end{tikzpicture}
\caption{Power iteration on $A=\begin{pmatrix}4&2\\1&3\end{pmatrix}$ starting from $(1,0)$. The
solid arrows are the direction after $k=0,1,2,3$ steps, with angles $0^\circ$, $14.04^\circ$,
$21.25^\circ$, and $24.39^\circ$. The dashed arrow is the dominant eigen direction $(2,1)$ at
$26.57^\circ$, which is approached but never passed; the last two arrows coincide because the
angles are getting tighter, and that is precisely the point of the picture.}
\label{fig:iterasi-pangkat}
\end{figure}

\kertas{For $J$ the matrix all of whose entries are $1$, take $B=J-I$. Determine the eigenvalues of
$B$ first, then predict $B^{10}$ using the pattern $I+\frac{\mu^{10}-(-1)^{10}}{3}J$ with $\mu$ the
nonzero eigenvalue of $B$. Check that prediction with Python, then explain why $B$ is
diagonalizable although it has a repeated eigenvalue. As an additional exercise, compute $e^{Bt}$
and determine the limit $t\to\infty$ of $e^{-2t}e^{Bt}$.}

\section{Computational lab}
\tingkat{advanced}

\begin{lab}
\begin{verbatim}
import numpy as np
from scipy.linalg import expm

np.set_printoptions(precision=4, suppress=True)

A = np.array([[2., 1., 1.],
              [1., 2., 1.],
              [1., 1., 2.]])
lam, V = np.linalg.eig(A)
print("eigenvalues A :", np.round(np.sort(lam.real), 10))
for i in range(3):
    sisa = A @ V[:, i] - lam[i] * V[:, i]
    print("  residual Av-lv :", round(float(np.linalg.norm(sisa)), 16))
print("A^10 via diagonalization :")
print(np.round((V @ np.diag(lam ** 10) @ np.linalg.inv(V)).real).astype(int))
print("A^10 via ordinary power :")
print(np.linalg.matrix_power(A, 10).astype(int))
print("max difference :", np.abs((V @ np.diag(lam ** 10)
                              @ np.linalg.inv(V)).real
                             - np.linalg.matrix_power(A, 10)).max())
print("cond V       :", round(float(np.linalg.cond(V)), 4))

A2 = np.array([[4., 2.],
               [1., 3.]])
x = np.array([1., 0.])
print("power iteration on A2:")
for k in range(6):
    sudut = np.degrees(np.arctan2(x[1], x[0]))
    rq = float(x @ A2 @ x / (x @ x))
    print("  k=%d angle=%7.2f rayleigh=%.4f" % (k, sudut, rq))
    x = A2 @ x / np.linalg.norm(A2 @ x)
print("limit angle  :", round(float(np.degrees(np.arctan2(1., 2.))), 2))

B = np.array([[-1., 4., 0.],
              [-1., 3., 0.],
              [1., -2., 2.]])
print("eigen B      :", np.round(np.sort(np.linalg.eigvals(B).real), 10))
print("rank(B - I)  :", np.linalg.matrix_rank(B - np.eye(3)))
print("B^10         :")
print(np.linalg.matrix_power(B, 10).astype(int))

M = np.array([[0.8, 0.3, 0.2],
              [0.1, 0.6, 0.3],
              [0.1, 0.1, 0.5]])
print("column sums  :", M.sum(axis=0))
print("eigen M      :", np.round(np.sort(np.linalg.eigvals(M).real), 10))
x = np.array([1., 0., 0.])
for _ in range(10):
    x = M @ x
print("x after 10   :", np.round(x, 6))
for _ in range(30):
    x = M @ x
print("x after 40   :", np.round(x, 6))

C = np.array([[1., 1000.], [0., 2.]])
Cp = C.copy()
Cp[1, 0] = 1e-6
print("eigen C      :", np.round(np.sort(np.linalg.eigvals(C).real), 6))
print("eigen C+1e-6 :", np.round(np.sort(np.linalg.eigvals(Cp).real), 6))
print("cond V(C)    :", round(float(np.linalg.cond(np.linalg.eig(C)[1])), 1))
print("expm(0.5A)   :", round(float(expm(0.5 * A)[0, 1]), 4))
print("formula      :", round(float((np.exp(2) - np.exp(0.5)) / 3), 4))
\end{verbatim}
Actual output:
\begin{verbatim}
eigenvalues A : [1. 1. 4.]
  residual Av-lv : 8e-16
  residual Av-lv : 1.1e-15
  residual Av-lv : 3e-16
A^10 via diagonalization :
[[349526 349525 349525]
 [349525 349526 349525]
 [349525 349525 349526]]
A^10 via ordinary power :
[[349526 349525 349525]
 [349525 349526 349525]
 [349525 349525 349526]]
max difference : 1.0477378964424133e-09
cond V       : 1.5227
power iteration on A2:
  k=0 angle=   0.00 rayleigh=4.0000
  k=1 angle=  14.04 rayleigh=4.6471
  k=2 angle=  21.25 rayleigh=4.8820
  k=3 angle=  24.39 rayleigh=4.9578
  k=4 angle=  25.69 rayleigh=4.9840
  k=5 angle=  26.21 rayleigh=4.9938
limit angle  : 26.57
eigen B      : [1. 1. 2.]
rank(B - I)  : 2
B^10         :
[[  -19    40     0]
 [  -10    21     0]
 [ 1023 -2046  1024]]
column sums  : [1. 1. 1.]
eigen M      : [0.4 0.5 1. ]
x after 10   : [0.5672 0.2661 0.1666]
x after 40   : [0.5667 0.2667 0.1667]
eigen C      : [1. 2.]
eigen C+1e-6 : [0.999 2.001]
cond V(C)    : 2000.0
expm(0.5A)   : 1.9134
formula      : 1.9134
\end{verbatim}
\end{lab}

The first line of that output lists the eigenvalues in increasing order, $1,1,4$. The next three
lines are the residuals $A\mathbf{v}-\lambda\mathbf{v}$ for the three eigenvectors that the library
returned. A residual of order $10^{-16}$ means that the equation is satisfied to the limit of
machine precision; no residual is exactly zero because eigenvectors are normalized and computed
iteratively. The library eigenvectors are also not $(1,1,1)$ and $(1,-1,0)$ but normalized
multiples: the direction is the same, the length differs, and diagonalization cares only about
direction.

The two $A^{10}$ blocks give identical integer entries, with largest difference
$1.05\times10^{-9}$. That difference is not an algebraic error: entries of size
$3.5\times10^{5}$ are computed through $V$, $V^{-1}$, and powers of real numbers, so rounding
errors of order $10^{-16}$ grow along with them.

The value \texttt{cond V} of $1.5227$ is close to one, and that is good news: $A$ is symmetric, so
its eigenvectors can be chosen mutually perpendicular and $P$ is nearly orthogonal.
Diagonalization for such matrices is far more stable.

The power iteration block repeats the computation in Figure~\ref{fig:iterasi-pangkat}: the angle
rises from $0$ to $26.21$ degrees, staying below the limit $26.57$, while the Rayleigh quotient
rises from $4.0000$ to $4.9938$ toward $5$.

The $B$ block shows the other side: its eigenvalues are $1,1,2$ and \texttt{rank(B - I)} equals
$2$, so the nullity is $1$. The geometric multiplicity for $\lambda=1$ is only one, while the
algebraic one is two; the matrix is defective. The $B^{10}$ block shows the consequence: the
entries $-19$ and $1024$ sit side by side, the first growing linearly and the second
exponentially.

The $M$ block verifies that every column sums to one, so that $\mathbf{x}_{k}$ stays stochastic.
Its eigenvalues are $0.4$, $0.5$, and $1$; after forty steps the numbers stop moving at six
digits, matching the rate $(1/2)^{k}$ derived by hand.

The last block gives one warning: the matrix $C$ has eigenvalues $1$ and $2$, and if one zero
entry below the diagonal is replaced by $10^{-6}$, its eigenvalues shift to $0.999$ and $2.001$; a
shift of $10^{-3}$, a thousand times larger than the perturbation. That number can be computed
exactly: the roots of $\lambda^{2}-3\lambda+1.999$ are $1.5\mp\sqrt{2510}/100$, that is
$0.999000998$ and $2.000999002$. The shift is proportional to $10^{-6}\cdot1000/(2-1)$: the
perturbation times the large entry divided by the gap between the two eigenvalues. The condition
number $2000$ captures the same amplification. Eigenvalues are still continuous in the entries,
but they need not be Lipschitz continuous with constant one.

\begin{tantangan}
Find a $2\times2$ matrix whose two eigenvalues are real and distinct, but whose two eigenvectors
are not perpendicular. Show that your choice is correct by computing the dot product, then explain
why a symmetric matrix cannot be used (the answer waits in Chapter 10).

Hint 1: two vectors are perpendicular exactly when their dot product is zero. Start from a simple
upper triangular matrix, then compute both eigenvectors without rounding.

Hint 2: the matrix $\begin{pmatrix}2&1\\0&3\end{pmatrix}$ is enough. The eigenvector for $3$ is
$(1,1)$ and the one for $2$ is $(1,0)$; their dot product is $1$, not zero. Compare with the
symmetric matrix $A$ in the lab, which gives $\texttt{cond V}\approx1.5$.

If you get stuck, go back to the section Similar matrices and invariants, then to the paragraph on
the condition number in Computational lab.
\end{tantangan}

\begin{untukdosen}
A teaching order that has worked: eigenvalues first, as preserved directions, then the
characteristic polynomial as a computational tool. If the polynomial comes first, students see
roots of a polynomial and forget that what is being sought is a direction. After eigen spaces,
introduce geometric multiplicity through nullity, then the condition for diagonalization. The
shear example should come before the condition theorem, so that the theorem answers a question
that is already felt.

Two typical misconceptions. First, students assume that a repeated eigenvalue always means not
diagonalizable. The matrix $(2,1,1;1,2,1;1,1,2)$ in the example section refutes that with
geometric multiplicity two. Second, they write eigenvectors as rows of $P$, then obtain a
$P^{-1}AP$ that is not diagonal and blame the formula. The rule to repeat: eigenvectors are
columns, and the column order of $P$ must be the same as the diagonal order of $D$.

Class activity: split the class into two groups for the matrix $(2,1,1;1,2,1;1,1,2)$. The first
group computes $A^{10}$ through diagonalization, the second through repeated powers or fast
repeated multiplication. Compare the results on the board, then ask the class to explain why the
two roads must meet. The noise usually arises over the denominator $3$ in $P^{-1}$.

Sign of readiness for what follows: students can say, without computing, why the shear matrix is
not diagonalizable, and can guess the form of $S^{k}$ for a $2\times2$ shear. If they mention the
appearance of a factor $k$, they are ready for Chapter 10.
\end{untukdosen}

\section{Summary}
\tingkat{basic}

\begin{ringkasan}
\begin{itemize}
  \item $\lambda$ is an eigenvalue of $A$ if there is $\mathbf{v}\neq\mathbf{0}$ with
        $A\mathbf{v}=\lambda\mathbf{v}$. Eigenvalues are the roots of
        $p_{A}(\lambda)=\det(\lambda I-A)$, and $E_{\lambda}=\operatorname{null}(A-\lambda I)$ is
        the corresponding subspace.
  \item The degree of $p_{A}$ is $n$, so there are at most $n$ distinct eigenvalues. Over
        $\mathbb{C}$ there are always $n$ when counted with multiplicity; over $\mathbb{R}$ there
        need not be, as with the rotation matrix.
  \item $\operatorname{tr}(A)=\sum_{i}\lambda_{i}$ and $\det(A)=\prod_{i}\lambda_{i}$, both
        counted with multiplicity. This is a cheap check for every hand computation.
  \item For every eigenvalue, $1\le\dim E_{\lambda}\le$ the algebraic multiplicity. Eigenvectors
        with distinct eigenvalues are always linearly independent.
  \item $A$ is diagonalizable exactly when there are $n$ independent eigenvectors, and exactly
        when the algebraic multiplicity equals the geometric one for every eigenvalue. Otherwise
        $A$ is called defective, and $A^{k}$ in general contains a factor $k$, not only
        $\lambda^{k}$.
  \item Similar matrices have the same characteristic polynomial, eigenvalues, trace, and
        determinant; the converse does not hold. Eigenvectors change along with the basis.
\end{itemize}
\end{ringkasan}

\section{Exercises}
\tingkat{intermediate}

\begin{latihan}
\begin{enumerate}
  \item Determine the eigenvalues and eigenvectors of
        $A_{1}=\begin{pmatrix}3&1\\ 2&2\end{pmatrix}$, then assemble $P$ and $D$ and check
        $P^{-1}A_{1}P=D$ by multiplying back.

  \item Compute $G^{10}$ for $G=\begin{pmatrix}4&2\\ 1&3\end{pmatrix}$ through diagonalization,
        then check the result with the NumPy function \texttt{matrix\_power}. Write down the
        general form of $G^{k}$ as well.

  \item Prove that if $\lambda$ is an eigenvalue of $A$ then $\lambda^{2}$ is an eigenvalue of
        $A^{2}$, and in general $\lambda^{k}$ is an eigenvalue of $A^{k}$ for every integer
        $k\ge1$. Explain why the converse fails: give an $A$ with $A^{2}=-I$, then compare the
        eigenvalues of $A^{2}$ with the real eigenvalues of $A$.
        (Hint: use the rotation matrix from the start of the chapter; it has no real eigenvalue.)

  \item Prove that $A$ is invertible if and only if $0$ is not an eigenvalue of $A$, and that in
        that case the eigenvalues of $A^{-1}$ are $1/\lambda$ for every eigenvalue $\lambda$ of
        $A$.
        (Hint: start from $A\mathbf{v}=\lambda\mathbf{v}$, multiply on the left by $A^{-1}$, and
        make sure that $\lambda\neq0$ so that the division is legitimate.)

  \item Assess this claim: \emph{two square matrices of the same size that have the same
        characteristic polynomial are necessarily similar.} Decide whether it is true, false, or
        not supported by enough evidence, and include a reason or a counterexample.

  \item Prove that $A$ and $A^{\top}$ have the same characteristic polynomial, using
        $\det(\lambda I-A)=\det((\lambda I-A)^{\top})$ and the properties of the determinant from
        Chapter 8. Explain why this means that the transpose does not change the eigenvalues but
        usually changes the eigenvectors.

  \item The matrix $K=\begin{pmatrix}2&1&0\\ 0&2&0\\ 0&0&1\end{pmatrix}$ has eigenvalues $2,2,1$.
        Determine the rank of $K-2I$, then decide whether $K$ is diagonalizable. Justify your
        answer with the multiplicity condition.

  \item A two-state Markov chain has matrix
        $T=\begin{pmatrix}0.9&0.4\\ 0.1&0.6\end{pmatrix}$. Determine the eigenvalues of $T$, its
        fixed point, and the number of steps needed for the initial error to shrink below
        $10^{-6}$. Check your answer with Python.
        (Hint: once the eigenvalues are known, the error shrinks like $|\lambda_{2}|^{k}$; solve
        $|\lambda_{2}|^{k}<10^{-6}$ with logarithms.)

  \item Prove: if $A$ is diagonalizable with $A=PDP^{-1}$, then
        $\operatorname{tr}(A^{k})=\sum_{i}\lambda_{i}^{k}$ for every $k\ge1$, using
        $\operatorname{tr}(XY)=\operatorname{tr}(YX)$. For the defective $B$ above, compute
        $\operatorname{tr}(B^{10})$ from the printed matrix, then compare it with
        $1^{10}+1^{10}+2^{10}$. Explain why that equality still holds although $B$ is defective,
        and at which point the proof above used diagonalization.
        (Hint: the trace does not change under a change of basis; use the form of $B^{k}$ already
        derived.)
\end{enumerate}
\end{latihan}

\begin{untukpeneliti}
What is really proved in this chapter are algebraic statements over any field and in finite
dimension: the equivalence of eigenvalues with roots of the characteristic polynomial, the
independence of eigenvectors for distinct eigenvalues, the inequality
$1\le\dim E_{\lambda}\le$ the algebraic multiplicity, and the equivalence of diagonalization with
the two conditions used here. The standard references are Axler, \emph{Linear Algebra Done Right}
4th edition, sections 5A and 5D (5.55, the four equivalent conditions, and 5.54); a matrix version
is in Anton, \emph{Elementary Linear Algebra} 11th edition, sections 5.1 and 5.2.

Two things are more often popularized than proved. First, the claim that eigenvalues ``measure
importance'', the basis of the PageRank explanation: the principal eigenvector of a stochastic
matrix gives steady weights only if the chain is irreducible and aperiodic, and that is the content
of the Perron-Frobenius theorem, which is not proved here; a full proof is in Strang,
\emph{Introduction to Linear Algebra} 6th edition, Appendix 8. Second, the claim that the
characteristic polynomial is a practical way to compute eigenvalues: for $n\ge5$ there is no
general root formula, and libraries use the QR algorithm on matrices, not polynomial factoring.
This chapter also stops at the diagonalizable case; Jordan form and the minimal polynomial are not
discussed, and are in Hoffman and Kunze, \emph{Linear Algebra} 2nd edition, sections 6.3 and 7.3.

Open questions: how far eigenvalues shift when the entries are perturbed, whose quantitative answer
requires the Bauer-Fike theorem or pseudospectra; and whether defectiveness itself is numerically
stable. Limits of validity: finite dimension, real or complex field, with no claim about the
stability of algorithms.
\end{untukpeneliti}

\peta{This chapter follows Anton, \emph{Elementary Linear Algebra} 11th edition, Chapter 5:
section 5.1 (Eigenvalues and Eigenvectors), 5.2 (Diagonalization), 5.3 (Complex Vector Spaces),
5.4 (Application: Differential Equations), and 5.5 (Application: Dynamical Systems and Markov
Chains). A more abstract counterpart is in Axler, \emph{Linear Algebra Done Right} 4th edition,
Chapter 5: 5A (eigenvalues and eigen spaces), 5B (existence of eigenvalues over complex vector
spaces), and 5D (diagonalizable operators, including 5.55). For Markov chains and the
Perron-Frobenius theorem, see Strang, \emph{Introduction to Linear Algebra} 6th edition, sections
6.1, 6.2, 6.5, and Appendix 8. Canonical forms for defective matrices are in Hoffman and Kunze,
\emph{Linear Algebra} 2nd edition, Chapter 6 (Elementary Canonical Forms), section 6.3 for the
minimal polynomial, and Chapter 7 (The Rational and Jordan Forms), section 7.3 for the Jordan
form.}

\chapter{Orthogonality, SVD, and Least Squares}

Five readings from one temperature probe: $100$, $101$, $102$, $103$, and $104$ degrees.
The recorded readings are $302.0$, $305.1$, $307.9$, $311.2$, and
$313.8$ in that order. The calibration model is a single straight line $y=a+bT$, so there
are five equations for two unknowns:
\[
  a+100b=302.0,\qquad a+101b=305.1,\qquad \dots,\qquad a+104b=313.8 .
\]
No pair $(a,b)$ satisfies all five; the five points are never
collinear, and the residual of each equation has a different sign. There is still one pair that
makes the most sense. In what sense is that pair best, and how far does it
move if one reading changes by a tenth?

The same question arrives from the opposite side. A detector computes a score
$s=\mathbf{w}^{\top}\mathbf{x}$ from three counters, and an attacker wants to shift $s$
by exactly $\Delta$ with as small a change in $\mathbf{x}$ as possible, because a large change
is easy to see. There are three unknowns and one condition, so the solution set is too large
rather than empty: among all the changes that give the same shift, which one is smallest?

Those two problems are one problem. In the first case the set of candidate answers is empty and
we look for the nearest one; in the second it is too large and we look for the
shortest. Both are settled by orthogonal projection. Chapter 9 supplied eigenvalues, but
only for square matrices with enough eigenvectors.
The calibration matrix above is $5\times2$, and the detector matrix is $2\times3$.
Forcing both to become square through $\mathbf{A}^{\top}\mathbf{A}$ does work,
but that road squares the sensitivity to error. The sixth section shows the
price in numbers.

\begin{untukanak}
Have you ever stood under a lamp and seen your shadow on the floor? That shadow is
not you: it is shorter, flattened, and it cannot show your face. But it is the best
the floor can do: of all the points on the floor, your shadow is the one
closest to the top of your head. If the lamp is moved, the shadow moves with it, even though you
have not moved.

Now picture two rulers that differ by one hundredth of a centimeter. If you
measure a table with both and then subtract the results, the leftover number comes almost
entirely from reading mistakes. A small difference drowns easily.

In this chapter we will learn why that ``closest'' shadow can be computed
cleanly, why among the many roads to the same destination there is one that is shortest, and
why two nearly twin measuring rules make our answer wobble.
\end{untukanak}

\section{Inner product, norm, and the Cauchy-Schwarz inequality}
\tingkat{basic}

Length and angle are two words we have used since school without ever
defining them. For vectors in $\mathbb{R}^{n}$ both come from a single operation
that eats two vectors and returns one number; that operation is what gives a
vector space a size, because the axioms of a vector space in Chapter 6 say nothing about distance.

The rule: multiply component by component, then add. The result is zero exactly when the
two vectors are at right angles, so this operation is used in place of angle even
in a space of a hundred dimensions.

\kamus{inner product}{hasilkali dalam}{Two lists of numbers folded into a single number that
measures how much they point the same way.}

\begin{definisi}[Inner product and norm]
The \emph{inner product} on a real vector space $V$ is a function
$\langle\cdot,\cdot\rangle\colon V\times V\to\mathbb{R}$ that satisfies three conditions:
\begin{enumerate}
  \item symmetric, $\langle\mathbf{u},\mathbf{v}\rangle=\langle\mathbf{v},\mathbf{u}\rangle$;
  \item linear in the first argument,
        $\langle c\mathbf{u}+\mathbf{w},\mathbf{v}\rangle
         =c\langle\mathbf{u},\mathbf{v}\rangle+\langle\mathbf{w},\mathbf{v}\rangle$;
  \item positive definite, $\langle\mathbf{v},\mathbf{v}\rangle\ge0$ with equality
        only when $\mathbf{v}=\mathbf{0}$.
\end{enumerate}
On $\mathbb{R}^{n}$ the standard inner product is
$\langle\mathbf{u},\mathbf{v}\rangle=\sum_{i}u_{i}v_{i}
=\mathbf{u}^{\top}\mathbf{v}$. The \emph{norm} it induces is
$\|\mathbf{v}\|=\sqrt{\langle\mathbf{v},\mathbf{v}\rangle}$.\index{inner product}
\index{norm}
\end{definisi}

The third condition is what makes $\|\mathbf{v}\|=0$ only for the zero vector, so the norm
deserves the name length. The norm also satisfies $\|c\mathbf{v}\|=|c|\,\|\mathbf{v}\|$ and
$\|\mathbf{u}+\mathbf{v}\|\le\|\mathbf{u}\|+\|\mathbf{v}\|$; the second is not automatic,
but follows from the next inequality.

In everyday language: the inner product of two vectors is never larger than the
product of their lengths. Two vectors can be very long, but if their directions are at right
angles the product stays small.

\begin{teorema}[Cauchy-Schwarz]
For all $\mathbf{u},\mathbf{v}\in V$,
$|\langle\mathbf{u},\mathbf{v}\rangle|\le\|\mathbf{u}\|\,\|\mathbf{v}\|$, with
equality exactly when one vector is a scalar multiple of the other.
\index{Cauchy-Schwarz}
\end{teorema}

\begin{proof}
If $\mathbf{v}=\mathbf{0}$, both sides are zero and there is nothing to prove.
Suppose $\mathbf{v}\ne\mathbf{0}$ and $t\in\mathbb{R}$ is arbitrary. Positive definiteness
gives
\[
  0\le\|\mathbf{u}-t\mathbf{v}\|^{2}
   =\|\mathbf{u}\|^{2}-2t\langle\mathbf{u},\mathbf{v}\rangle+t^{2}\|\mathbf{v}\|^{2}.
\]
The right-hand side is a quadratic polynomial in $t$ that is never negative. The coefficient of
$t^{2}$ is positive, so its discriminant cannot be positive:
\[
  4\langle\mathbf{u},\mathbf{v}\rangle^{2}-4\|\mathbf{u}\|^{2}\|\mathbf{v}\|^{2}\le0 .
\]
After dividing by four and taking square roots, the required inequality follows. Equality
holds exactly when that polynomial has a double root, that is, when
$\mathbf{u}-t\mathbf{v}=\mathbf{0}$ for some $t$.
\end{proof}

That proof used no coordinates, so the inequality holds in every inner-product
space, including function spaces.

\begin{akibat}[Triangle inequality]
$\|\mathbf{u}+\mathbf{v}\|\le\|\mathbf{u}\|+\|\mathbf{v}\|$.
\end{akibat}

\begin{proof}
Expand $\|\mathbf{u}+\mathbf{v}\|^{2}
=\|\mathbf{u}\|^{2}+2\langle\mathbf{u},\mathbf{v}\rangle+\|\mathbf{v}\|^{2}$.
Cauchy-Schwarz bounds the middle term by $2\|\mathbf{u}\|\|\mathbf{v}\|$, so
\[
  \|\mathbf{u}+\mathbf{v}\|^{2}\le\|\mathbf{u}\|^{2}+2\|\mathbf{u}\|\|\mathbf{v}\|
  +\|\mathbf{v}\|^{2}=(\|\mathbf{u}\|+\|\mathbf{v}\|)^{2}.
\]
Both sides are nonnegative, so the square root may be taken.
\end{proof}

Now we have length. Angle follows by definition: two vectors are perpendicular when their
inner product is zero. That makes sense because Cauchy-Schwarz guarantees that
$\langle\mathbf{u},\mathbf{v}\rangle/(\|\mathbf{u}\|\|\mathbf{v}\|)$ lies between $-1$
and $1$, so it may be read as a cosine.

\kamus{orthogonal}{ortogonal}{Two directions that contribute nothing to each other; the shadow
of one direction on the other has length zero.}

\begin{definisi}[Orthogonal and orthonormal]
$\mathbf{u}$ and $\mathbf{v}$ are \emph{orthogonal} when
$\langle\mathbf{u},\mathbf{v}\rangle=0$, written $\mathbf{u}\perp\mathbf{v}$. A set of
vectors is called \emph{orthogonal} when every two distinct members are orthogonal, and
\emph{orthonormal} when all its members also have norm one.\index{orthogonality}
\end{definisi}

One small consequence follows directly from expanding the norm: if $\mathbf{u}\perp\mathbf{v}$
then $\|\mathbf{u}+\mathbf{v}\|^{2}=\|\mathbf{u}\|^{2}+\|\mathbf{v}\|^{2}$. The
Pythagorean theorem is a special case of this definition, not the other way round.

\section{Orthogonal sets, projection, and Gram-Schmidt}
\tingkat{intermediate}

For an arbitrary basis we must solve a linear system; for an orthonormal basis we need not,
because each coefficient is read off with one inner product. One fact makes that
legitimate: an orthogonal set with no zero vector is always linearly independent.

\begin{lemma}
If $\mathbf{v}_{1},\dots,\mathbf{v}_{k}$ are orthogonal and none of them is the zero vector,
then that set is linearly independent.
\end{lemma}

\begin{proof}
Suppose $\sum_{i}c_{i}\mathbf{v}_{i}=\mathbf{0}$. Take the inner product with
$\mathbf{v}_{j}$. All terms with $i\ne j$ vanish because of orthogonality, so only
$c_{j}\langle\mathbf{v}_{j},\mathbf{v}_{j}\rangle=0$ is left. Since $\mathbf{v}_{j}\ne\mathbf{0}$,
its norm is nonzero, hence $c_{j}=0$. This holds for every $j$.
\end{proof}

Now the statement about coefficients, in the form most often used.

\begin{proposisi}
Suppose $(\mathbf{q}_{1},\dots,\mathbf{q}_{n})$ is an orthonormal basis of $V$. Then for
every $\mathbf{v}\in V$,
\[
  \mathbf{v}=\sum_{i=1}^{n}\langle\mathbf{v},\mathbf{q}_{i}\rangle\,\mathbf{q}_{i},
  \qquad
  \|\mathbf{v}\|^{2}=\sum_{i=1}^{n}\langle\mathbf{v},\mathbf{q}_{i}\rangle^{2}.
\]
\end{proposisi}

\begin{proof}
Write $\mathbf{v}=\sum_{i}a_{i}\mathbf{q}_{i}$; those coefficients exist and are unique because
this is a basis. The inner product with $\mathbf{q}_{j}$ gives
$\langle\mathbf{v},\mathbf{q}_{j}\rangle=\sum_{i}a_{i}\langle\mathbf{q}_{i},\mathbf{q}_{j}\rangle
=a_{j}$, since $\langle\mathbf{q}_{i},\mathbf{q}_{j}\rangle$ equals one when $i=j$
and zero otherwise. For the norm formula, use orthogonality on the expansion
$\|\mathbf{v}\|^{2}=\langle\sum_{i}a_{i}\mathbf{q}_{i},\sum_{j}a_{j}\mathbf{q}_{j}\rangle$;
the cross terms all vanish.
\end{proof}

The second equality is called Parseval's identity: in an orthonormal basis the total length is
the square root of the sum of squared coordinates; in a skew basis that property fails.

The next question comes from the opening problem: $\mathbf{b}$ lies outside the subspace
$W$, and we want the point of $W$ closest to it. That point is called the projection,
and the way to find it is to add up the shadow of $\mathbf{b}$ on each orthonormal direction of $W$.

\kamus{projection}{proyeksi}{The shadow of a stick on the floor: the point on the floor that is
closest to the tip of the stick.}

\begin{definisi}[Orthogonal projection]
Suppose $W\subseteq\mathbb{R}^{m}$ is a subspace with orthonormal basis
$(\mathbf{q}_{1},\dots,\mathbf{q}_{k})$. The \emph{orthogonal projection} of $\mathbf{b}$ onto $W$
is
\[
  \operatorname{proj}_{W}\mathbf{b}=\sum_{i=1}^{k}\langle\mathbf{b},\mathbf{q}_{i}\rangle\mathbf{q}_{i},
  \qquad
  \mathbf{r}=\mathbf{b}-\operatorname{proj}_{W}\mathbf{b}.
\]
The vector $\mathbf{r}$ is called the \emph{residual}.\index{projection}
\end{definisi}

That residual is orthogonal to all of $W$, not merely to the basis vectors: because
$\langle\mathbf{r},\mathbf{q}_{j}\rangle=\langle\mathbf{b},\mathbf{q}_{j}\rangle
-\langle\mathbf{b},\mathbf{q}_{j}\rangle=0$ and every vector of $W$ is a combination of
the $\mathbf{q}_{j}$.

What makes the projection useful: it is not just one point in $W$, it is the closest
point. That is a claim about every other point in $W$, so it needs a proof.

\begin{bacakeras}
Read aloud, two voices taking turns.

\textbf{Voice one:} There is a point hanging above the floor, and the floor is asked to produce its stand-in.

\textbf{Voice two:} The floor cannot rise. It can only drop a shadow, and that shadow is the floor point closest to the point above.

\textbf{Voice one:} Then where does the part that never reaches the floor go?

\textbf{Voice two:} That part stands upright. It has no floor to stand on, so no path on the floor can shorten it.

Repeat this line until it sits on the tongue: \emph{what falls to the floor is what is closest, what stands upright is what cannot be brought back down.}
\end{bacakeras}

\begin{teorema}[Best approximation]
For every $\mathbf{w}\in W$,
$\|\mathbf{b}-\operatorname{proj}_{W}\mathbf{b}\|\le\|\mathbf{b}-\mathbf{w}\|$, with
equality only when $\mathbf{w}=\operatorname{proj}_{W}\mathbf{b}$.
\end{teorema}

\begin{proof}
Write $\mathbf{p}=\operatorname{proj}_{W}\mathbf{b}$. For arbitrary $\mathbf{w}\in W$,
\[
  \mathbf{b}-\mathbf{w}=(\mathbf{b}-\mathbf{p})+(\mathbf{p}-\mathbf{w}).
\]
The two terms are orthogonal: the first is orthogonal to all of $W$, and
$\mathbf{p}-\mathbf{w}\in W$. Therefore the squared norm splits:
\[
  \|\mathbf{b}-\mathbf{w}\|^{2}=\|\mathbf{b}-\mathbf{p}\|^{2}+\|\mathbf{p}-\mathbf{w}\|^{2}.
\]
The second term is nonnegative, so the right-hand side is smallest when $\mathbf{w}=\mathbf{p}$.
\end{proof}

The geometric reading is in Figure~\ref{fig:proyeksi-sisa}.

\begin{figure}[htbp]
\centering
\begin{tikzpicture}[scale=1.18]
  \draw[->,layergrey] (-0.5,0)--(3.4,0) node[right,font=\small]{$x_{1}$};
  \draw[->,layergrey] (0,-0.5)--(0,3.9) node[above,font=\small]{$x_{2}$};
  \draw[very thick,accent] (-0.42,-0.21)--(2.75,1.375)
    node[right,font=\small]{$W$};
  \draw[densely dotted,thick,layergrey] (0.2,3.4)--(2.6,1.3);
  \fill[layergrey] (2.6,1.3) circle (1.2pt);
  \node[font=\small,layergrey,below right=-2pt] at (2.6,1.3) {$\mathbf{w}$};
  \node[font=\small,layergrey,above right=-3pt] at (1.05,2.62)
    {$\|\mathbf{b}-\mathbf{w}\|$};
  \draw[->,very thick,layerteal] (0,0)--(0.2,3.4)
    node[left,font=\small]{$\mathbf{b}$};
  \draw[->,thick,layerindigo] (0,0)--(1.52,0.76);
  \node[font=\small,layerindigo,below right=-2pt] at (1.34,0.6) {$\mathbf{p}$};
  \draw[-{Latex},very thick,layerplum] (1.52,0.76)--(0.2,3.4);
  \node[font=\small,layerplum,right=2pt] at (0.95,2.15) {$\mathbf{r}$};
  \draw[thick,layergrey] (1.717,0.858)--(1.619,1.055)--(1.422,0.957);
  \fill[layerplum] (1.52,0.76) circle (1.4pt);
\end{tikzpicture}
\caption{Orthogonal projection of $\mathbf{b}=(0.2;\,3.4)$ onto the line $W$
spanned by $\mathbf{q}=(2,1)/\sqrt{5}$. Its perpendicular foot is
$\mathbf{p}=(1.52;\,0.76)$, marked as a right angle with $W$, and the residual
$\mathbf{r}=\mathbf{b}-\mathbf{p}$ is perpendicular to all of $W$. The point
$\mathbf{w}=(2.6;\,1.3)$, which also lies in $W$, gives distance
$\|\mathbf{b}-\mathbf{w}\|\approx3.19$, longer than
$\|\mathbf{b}-\mathbf{p}\|\approx2.95$.}
\label{fig:proyeksi-sisa}
\end{figure}

\begin{tangan}
Materials: one pencil, one desk lamp, one sheet of paper, and a ruler.

Place the paper on the table, then move the lamp to the far end of the room until its light is almost parallel. A lamp that is too close lengthens the shadow, and what is measured is then no longer an orthogonal projection.

Hold the pencil upright with its tip touching the paper, and mark that point. Tilt the pencil to an angle of sixty degrees with the paper, mark the end of the shadow, then repeat at forty-five and thirty degrees. Measure the three shadows with the ruler.

What to look for: the shadow is longest when the pencil almost lies on the paper, and shortest when the pencil is upright. The length of the shadow does not care how high the tip of the pencil is above the paper; it records only the horizontal part. That upright part is the residual the floor cannot reach.
\end{tangan}

If the columns of $\mathbf{Q}$ are $\mathbf{q}_{1},\dots,\mathbf{q}_{k}$, the projection formula becomes a
single matrix product: $\mathbf{p}=\mathbf{Q}\mathbf{Q}^{\top}\mathbf{b}$. The matrix
$\mathbf{P}=\mathbf{Q}\mathbf{Q}^{\top}$ satisfies $\mathbf{P}^{2}=\mathbf{P}$ and
$\mathbf{P}^{\top}=\mathbf{P}$: projecting twice is the same as once. Its trace is $k$,
the dimension of $W$; for the example in the seventh section the trace is $2$.

Now we know how to project, but we do not yet have a way to obtain an orthonormal basis.
The Gram-Schmidt algorithm builds one vector at a time: normalize the first vector, then strip
from each following vector the part parallel to those already present. Each step
uses only the projection we have just proved.

\begin{teorema}[Gram-Schmidt and the QR factorization]
Suppose $\mathbf{a}_{1},\dots,\mathbf{a}_{n}$ are linearly independent in $\mathbb{R}^{m}$.
Define in order
\[
  \mathbf{v}_{1}=\mathbf{a}_{1},\qquad
  \mathbf{v}_{j}=\mathbf{a}_{j}-\sum_{i<j}\langle\mathbf{a}_{j},\mathbf{q}_{i}\rangle\mathbf{q}_{i},
  \qquad
  \mathbf{q}_{j}=\frac{\mathbf{v}_{j}}{\|\mathbf{v}_{j}\|}
\]
for $j=2,\dots,n$. Then $\mathbf{q}_{1},\dots,\mathbf{q}_{n}$ are orthonormal and
$\operatorname{span}(\mathbf{q}_{1},\dots,\mathbf{q}_{j})
=\operatorname{span}(\mathbf{a}_{1},\dots,\mathbf{a}_{j})$ for every $j$. Consequently
$\mathbf{A}=\mathbf{Q}\mathbf{R}$ with $\mathbf{Q}$ having orthonormal columns and
$\mathbf{R}$ upper triangular with positive diagonal.\index{Gram-Schmidt}\index{QR}
\end{teorema}

\begin{proof}
We prove this by induction on $j$. For $j=1$, $\mathbf{q}_{1}$ is
$\mathbf{a}_{1}$ normalized, so the claim holds. Assume the claim holds up to $j-1$.
The vector $\mathbf{v}_{j}$ is $\mathbf{a}_{j}$ minus its combination with
$\mathbf{q}_{1},\dots,\mathbf{q}_{j-1}$, so $\mathbf{v}_{j}$ still lies in
$\operatorname{span}(\mathbf{a}_{1},\dots,\mathbf{a}_{j})$. For every $i<j$,
\[
  \langle\mathbf{v}_{j},\mathbf{q}_{i}\rangle
  =\langle\mathbf{a}_{j},\mathbf{q}_{i}\rangle-\langle\mathbf{a}_{j},\mathbf{q}_{i}\rangle=0,
\]
because $\mathbf{q}_{i}$ is orthonormal. So $\mathbf{v}_{j}$ is orthogonal to all the
previous vectors. That vector is not zero: if it were, $\mathbf{a}_{j}$ would become a
combination of $\mathbf{a}_{1},\dots,\mathbf{a}_{j-1}$, contradicting linear independence.
Therefore $\mathbf{q}_{j}$ is defined and the set stays orthonormal as it grows.

For the factorization, write each $\mathbf{a}_{j}$ as
$\|\mathbf{v}_{j}\|\mathbf{q}_{j}+\sum_{i<j}\langle\mathbf{a}_{j},\mathbf{q}_{i}\rangle\mathbf{q}_{i}$.
Those coefficients fill column $j$ of $\mathbf{R}$; the entries above the diagonal are
$\langle\mathbf{a}_{j},\mathbf{q}_{i}\rangle$ and the diagonal entries are $\|\mathbf{v}_{j}\|>0$.
\end{proof}

One warning: classical Gram-Schmidt computes each coefficient against the original vector
$\mathbf{a}_{j}$, not against the vector already reduced. If the columns of $\mathbf{A}$ are
nearly dependent, the result stops being orthonormal; subtracting from $\mathbf{v}_{j}$ first
repairs it at almost no cost. The eighth section shows the numbers: on a
matrix with $\kappa\approx6.8\times10^{6}$, the classical version gives
$\|\mathbf{Q}^{\top}\mathbf{Q}-\mathbf{I}\|=2.05\times10^{-2}$, the modified version
$3.40\times10^{-10}$, Householder $1.19\times10^{-15}$; the first number means two
columns that should be at right angles differ by an angle of about one degree.

\section{Least squares and the normal equations}
\tingkat{intermediate}

Back to the five probe readings. The system $\mathbf{A}\mathbf{x}=\mathbf{b}$ has no solution,
but we can demand that the residual $\mathbf{b}-\mathbf{A}\mathbf{x}$ be as small as possible:
demanding a zero residual amounts to demanding a solution, while demanding minimum length always
has an answer. The measure used is the sum of squared differences, not the sum of absolute values,
because the sum of squares is born from the inner product, which puts the whole projection
machinery at our disposal; the sum of absolute values does not, and its answer may fail to be unique.

\kamus{least squares}{kuadrat terkecil}{When no line passes through all the points,
choose the line whose sum of squared errors is smallest.}

\begin{definisi}[Least-squares solution]
For $\mathbf{A}\in\mathbb{R}^{m\times n}$ and $\mathbf{b}\in\mathbb{R}^{m}$,
a vector $\hat{\mathbf{x}}\in\mathbb{R}^{n}$ is called a \emph{least-squares solution}
if $\|\mathbf{A}\hat{\mathbf{x}}-\mathbf{b}\|\le\|\mathbf{A}\mathbf{x}-\mathbf{b}\|$
for every $\mathbf{x}\in\mathbb{R}^{n}$.\index{least squares}
\end{definisi}

That definition does not mention derivatives. What is used is geometry: $\mathbf{A}\mathbf{x}$ always
lies in $\operatorname{col}(\mathbf{A})$, so the smallest residual is the point of
$\operatorname{col}(\mathbf{A})$ closest to $\mathbf{b}$; that is exactly the projection
problem from before.

\begin{teorema}[Normal equations]
The vector $\hat{\mathbf{x}}$ minimizes $\|\mathbf{A}\mathbf{x}-\mathbf{b}\|$ if and
only if
\[
  \mathbf{A}^{\top}\mathbf{A}\hat{\mathbf{x}}=\mathbf{A}^{\top}\mathbf{b}.
\]
\end{teorema}

\begin{proof}
Take arbitrary $\mathbf{x}$ and arbitrary $\mathbf{h}\in\mathbb{R}^{n}$, then write
$\mathbf{x}+\mathbf{h}$ as a neighbor of $\mathbf{x}$. Then
\[
  \|\mathbf{A}(\mathbf{x}+\mathbf{h})-\mathbf{b}\|^{2}
  =\|\mathbf{A}\mathbf{x}-\mathbf{b}\|^{2}
   +2\mathbf{h}^{\top}\mathbf{A}^{\top}(\mathbf{A}\mathbf{x}-\mathbf{b})
   +\|\mathbf{A}\mathbf{h}\|^{2}.
\]

First, suppose $\mathbf{A}^{\top}(\mathbf{A}\mathbf{x}-\mathbf{b})=\mathbf{0}$.
The middle term vanishes for every $\mathbf{h}$, and the last term is nonnegative, so
$\|\mathbf{A}(\mathbf{x}+\mathbf{h})-\mathbf{b}\|^{2}\ge\|\mathbf{A}\mathbf{x}-\mathbf{b}\|^{2}$.
Hence $\mathbf{x}$ minimizes.

Second, suppose $\mathbf{A}^{\top}(\mathbf{A}\mathbf{x}-\mathbf{b})=\mathbf{z}\ne\mathbf{0}$.
Choose $\mathbf{h}=-\epsilon\mathbf{z}$ with small $\epsilon>0$. The middle term becomes
$-2\epsilon\|\mathbf{z}\|^{2}$, while the last term grows by
$\epsilon^{2}\|\mathbf{A}\mathbf{z}\|^{2}$. For $\epsilon$ small enough, the negative part
beats the positive part, so $\mathbf{x}$ is not a minimizer. Hence the condition is necessary.
\end{proof}

The system $\mathbf{A}^{\top}\mathbf{A}\mathbf{x}=\mathbf{A}^{\top}\mathbf{b}$ is called the
\emph{normal equations}. The condition
$\mathbf{A}^{\top}(\mathbf{A}\hat{\mathbf{x}}-\mathbf{b})=\mathbf{0}$ says that the residual
$\mathbf{r}=\mathbf{b}-\mathbf{A}\hat{\mathbf{x}}$ is orthogonal to every column of
$\mathbf{A}$, hence to all of $\operatorname{col}(\mathbf{A})$. In other words
$\mathbf{A}\hat{\mathbf{x}}$ is the projection of $\mathbf{b}$ onto the column space.\index{normal equations}

From here one practical question follows: when is the answer unique? The matrix
$\mathbf{A}^{\top}\mathbf{A}$ is always square and symmetric, but not necessarily
invertible.

\begin{akibat}
If the columns of $\mathbf{A}$ are linearly independent, then $\mathbf{A}^{\top}\mathbf{A}$ is
invertible and the least-squares solution is unique, namely
$\hat{\mathbf{x}}=(\mathbf{A}^{\top}\mathbf{A})^{-1}\mathbf{A}^{\top}\mathbf{b}$.
\end{akibat}

\begin{proof}
Suppose $\mathbf{A}^{\top}\mathbf{A}\mathbf{z}=\mathbf{0}$. Multiply from the left by
$\mathbf{z}^{\top}$, which gives $\|\mathbf{A}\mathbf{z}\|^{2}
=\mathbf{z}^{\top}\mathbf{A}^{\top}\mathbf{A}\mathbf{z}=0$, hence
$\mathbf{A}\mathbf{z}=\mathbf{0}$. Linear independence of the columns forces $\mathbf{z}=\mathbf{0}$.
So the kernel of $\mathbf{A}^{\top}\mathbf{A}$ contains only the zero vector, and a square
matrix with trivial kernel is invertible. Uniqueness of $\hat{\mathbf{x}}$ then follows from
the normal equations theorem.
\end{proof}

\begin{resep}
\begin{enumerate}
  \item Build $\mathbf{A}$ and $\mathbf{b}$: one row of $\mathbf{A}$ per equation, with the columns being the functions used.
  \item Form $\mathbf{A}^{\top}\mathbf{A}$, which is $n\times n$, and $\mathbf{A}^{\top}\mathbf{b}$, which is $n\times1$.
  \item Solve $\mathbf{A}^{\top}\mathbf{A}\hat{\mathbf{x}}=\mathbf{A}^{\top}\mathbf{b}$; check $\det(\mathbf{A}^{\top}\mathbf{A})\ne0$ first, since zero means the answer is not unique.
  \item Compute $\mathbf{r}=\mathbf{b}-\mathbf{A}\hat{\mathbf{x}}$, check $\mathbf{A}^{\top}\mathbf{r}=\mathbf{0}$, then report $\|\mathbf{r}\|$.
\end{enumerate}
Warning about order: do not invert anything before the third step. $\mathbf{A}$ is not square, so $\mathbf{A}^{-1}$ does not exist, and turning $\mathbf{A}^{\top}\mathbf{b}$ into $\mathbf{b}^{\top}\mathbf{A}$ gives a $1\times n$ row vector, not a system of equations.
\end{resep}

Conversely, if the columns of $\mathbf{A}$ are dependent, $\mathbf{A}^{\top}\mathbf{A}$ is singular
and there are infinitely many solutions; the fifth section chooses one of them. The normal
equations turn an $m\times n$ problem into an $n\times n$ one, a large saving for
$m\gg n$, and the price shows up in the sixth section.

\section{Singular value decomposition}
\tingkat{advanced}

Chapter 9 turned square matrices into diagonal form by choosing one basis.
The condition is heavy: the matrix must have enough eigenvectors, and for a real matrix its
eigenvalues must be real. Our calibration matrix is $5\times2$, so it is not even
square, and there is no such thing as an eigenvalue for it.

The way out: allow two different bases, one in the domain and one in the codomain,
both orthonormal. With that relaxation every matrix is diagonal; that is the
singular value decomposition. The numbers on the diagonal are the lengths of the axes of the ellipse
produced from the unit ball, always real and nonnegative even when the matrix is not
symmetric.

\kamus{singular value}{nilai singular}{How far a direction is stretched by the
matrix; zero means that direction is flattened until it disappears.}

\begin{definisi}[Singular values]
Suppose $\mathbf{A}\in\mathbb{R}^{m\times n}$. The \emph{singular values} of $\mathbf{A}$
are the square roots of the eigenvalues of $\mathbf{A}^{\top}\mathbf{A}$, sorted in decreasing order:
$\sigma_{1}\ge\sigma_{2}\ge\dots\ge\sigma_{n}\ge0$ with
$\sigma_{i}=\sqrt{\lambda_{i}(\mathbf{A}^{\top}\mathbf{A})}$.\index{singular value}
\end{definisi}

That definition is legitimate if $\mathbf{A}^{\top}\mathbf{A}$ has $n$ real eigenvalues with
orthonormal eigenvectors; for a symmetric matrix that is true, and we state it as a
lemma, with one step honestly labelled as not done.

\begin{lemma}[Spectral theorem for symmetric matrices]
If $\mathbf{S}\in\mathbb{R}^{n\times n}$ is symmetric, then there is an orthonormal basis
$\mathbf{v}_{1},\dots,\mathbf{v}_{n}$ of $\mathbb{R}^{n}$ and real numbers
$\lambda_{1},\dots,\lambda_{n}$ with $\mathbf{S}\mathbf{v}_{i}=\lambda_{i}\mathbf{v}_{i}$.
\end{lemma}

\begin{proof}[Sketch]
The function $\mathbf{x}\mapsto\mathbf{x}^{\top}\mathbf{S}\mathbf{x}$ is continuous on the closed
and bounded unit sphere, so it attains a maximum at some point
$\mathbf{v}_{1}$ of norm one. By the Lagrange multiplier method, that point satisfies
$\mathbf{S}\mathbf{v}_{1}=\lambda_{1}\mathbf{v}_{1}$; the value $\lambda_{1}$ is that
maximum. Next, $\mathbf{v}_{1}^{\perp}$ is also invariant under $\mathbf{S}$, because
$\langle\mathbf{S}\mathbf{w},\mathbf{v}_{1}\rangle
=\langle\mathbf{w},\mathbf{S}\mathbf{v}_{1}\rangle
=\lambda_{1}\langle\mathbf{w},\mathbf{v}_{1}\rangle=0$ for $\mathbf{w}\perp\mathbf{v}_{1}$.
The restriction of $\mathbf{S}$ to $\mathbf{v}_{1}^{\perp}$ is symmetric and of dimension $n-1$, so
induction supplies the rest of the orthonormal basis.

The step not carried out here is the proof that a continuous function on a closed and bounded set
attains a maximum. That is the extreme value theorem, proved in Chapter 11,
and a full proof of this lemma is in Axler, \emph{Linear Algebra Done Right} 4th edition,
section 7B, or in Hoffman and Kunze's chapter on symmetric operators.
\end{proof}

With that lemma, $\mathbf{A}^{\top}\mathbf{A}$ is symmetric and
$\mathbf{x}^{\top}\mathbf{A}^{\top}\mathbf{A}\mathbf{x}=\|\mathbf{A}\mathbf{x}\|^{2}\ge0$,
so its eigenvalues are nonnegative and their roots are real. Now the main theorem.

\begin{teorema}[Singular value decomposition]
Every $\mathbf{A}\in\mathbb{R}^{m\times n}$ of rank $r$ can be written as
\[
  \mathbf{A}=\mathbf{U}\boldsymbol{\Sigma}\mathbf{V}^{\top},
\]
with $\mathbf{U}\in\mathbb{R}^{m\times m}$ and $\mathbf{V}\in\mathbb{R}^{n\times n}$
matrices with orthonormal columns, and $\boldsymbol{\Sigma}\in\mathbb{R}^{m\times n}$
diagonal with entries $\sigma_{1}\ge\dots\ge\sigma_{r}>0$ and the rest zero.
\index{SVD}
\end{teorema}

\begin{proof}
Because $\mathbf{A}^{\top}\mathbf{A}$ is symmetric, the spectral lemma gives an orthonormal basis
$\mathbf{v}_{1},\dots,\mathbf{v}_{n}$ with
$\mathbf{A}^{\top}\mathbf{A}\mathbf{v}_{i}=\sigma_{i}^{2}\mathbf{v}_{i}$ and
$\sigma_{1}\ge\dots\ge\sigma_{n}\ge0$. The nonzero eigenvalues number $r$;
call their indices $1,\dots,r$.

For $i\le r$ define $\mathbf{u}_{i}=\mathbf{A}\mathbf{v}_{i}/\sigma_{i}$, which
has norm one because $\|\mathbf{A}\mathbf{v}_{i}\|^{2}
=\mathbf{v}_{i}^{\top}\mathbf{A}^{\top}\mathbf{A}\mathbf{v}_{i}=\sigma_{i}^{2}$.
For $i,j\le r$,
\[
  \langle\mathbf{u}_{i},\mathbf{u}_{j}\rangle
  =\frac{1}{\sigma_{i}\sigma_{j}}\mathbf{v}_{i}^{\top}\mathbf{A}^{\top}\mathbf{A}\mathbf{v}_{j}
  =\frac{\sigma_{j}^{2}}{\sigma_{i}\sigma_{j}}\langle\mathbf{v}_{i},\mathbf{v}_{j}\rangle
  =\delta_{ij},
\]
so $\mathbf{u}_{1},\dots,\mathbf{u}_{r}$ are orthonormal. For $i>r$ we have
$\|\mathbf{A}\mathbf{v}_{i}\|^{2}=\sigma_{i}^{2}=0$, hence $\mathbf{A}\mathbf{v}_{i}=\mathbf{0}$.

Those vectors do not necessarily fill $\mathbb{R}^{m}$; complete them to an orthonormal basis with
Gram-Schmidt on any basis of their orthogonal complement. Define
$\mathbf{B}=\sum_{i\le r}\sigma_{i}\mathbf{u}_{i}\mathbf{v}_{i}^{\top}$. For $j\le r$
we have $\mathbf{B}\mathbf{v}_{j}=\sigma_{j}\mathbf{u}_{j}=\mathbf{A}\mathbf{v}_{j}$, and
for $j>r$ both are zero. The two matrices agree on a basis, so $\mathbf{B}=\mathbf{A}$.
\end{proof}

The geometric reading: the unit ball is mapped by $\mathbf{A}$ to an ellipsoid. The columns
$\mathbf{v}_{i}$ point along the axes of the ball, $\mathbf{u}_{i}$ along the axes of the ellipsoid, and $\sigma_{i}$
is the length of that axis: $\mathbf{V}^{\top}$ rotates, $\boldsymbol{\Sigma}$ stretches,
$\mathbf{U}$ rotates again.

Figure~\ref{fig:svd-lingkaran-elips} carries out that reading on
$\mathbf{A}=\begin{pmatrix}0&3\\4&1\end{pmatrix}$: the axis $\mathbf{v}_{1}$, tilted
$26.57^\circ$, ends up as the axis $\mathbf{u}_{1}$ of the ellipse, tilted
$71.57^\circ$.

\begin{figure}[htbp]
\centering
\begin{tikzpicture}[scale=1.55]
  \draw[->,gray] (-1.3,0)--(1.3,0);
  \draw[->,gray] (0,-1.3)--(0,1.3);
  \draw[thick,accent] (0,0) circle (1);
  \draw[->,very thick,blue!55!black] (0,0)--(0.894,0.447)
    node[above right=-3pt] {$\mathbf{v}_{1}$};
  \draw[->,very thick,blue!55!black] (0,0)--(-0.447,0.894)
    node[above left=-3pt] {$\mathbf{v}_{2}$};
  \node[below,font=\small] at (0,-1.5) {unit circle};
\end{tikzpicture}
\hfill
\begin{tikzpicture}[scale=0.68]
  \draw[->,gray] (-3.5,0)--(3.5,0);
  \draw[->,gray] (0,-4.6)--(0,4.6);
  \draw[thick,accent] plot[domain=0:360,samples=160]
    ({3*sin(\x)},{4*cos(\x)+sin(\x)});
  \draw[->,very thick,blue!55!black] (0,0)--(1.342,4.025)
    node[right=3pt] {$\sigma_{1}\mathbf{u}_{1}$};
  \draw[->,very thick,blue!55!black] (0,0)--(2.683,-0.894)
    node[below right=-3pt] {$\sigma_{2}\mathbf{u}_{2}$};
  \node[below,font=\small] at (0,-5.1) {its image};
\end{tikzpicture}
\caption{The action of $\mathbf{A}=\begin{pmatrix}0&3\\4&1\end{pmatrix}$ on the unit circle.
Left: two orthonormal axes of the circle, $\mathbf{v}_{1}=(2,1)/\sqrt{5}$ and
$\mathbf{v}_{2}=(-1,2)/\sqrt{5}$. Right: the image is an ellipse with semi-axes
$\sigma_{1}=3\sqrt{2}\approx4.243$ (numpy prints $4.242641$) along
$\mathbf{u}_{1}=(1,3)/\sqrt{10}$, and $\sigma_{2}=2\sqrt{2}\approx2.828$ (numpy
$2.828427$) along $\mathbf{u}_{2}=(3,-1)/\sqrt{10}$.}
\label{fig:svd-lingkaran-elips}
\end{figure}

The sum form $\mathbf{A}=\sum_{i}\sigma_{i}\mathbf{u}_{i}\mathbf{v}_{i}^{\top}$
is the most honest way to read it: the matrix is broken into $r$ layers of rank
one, ordered from the largest contribution. Truncating after $k$ terms gives the best
rank-$k$ approximation in the Frobenius norm; that is the Eckart-Young theorem, whose proof
is not carried out here and can be read in Golub and Van Loan, section 2.4.

Figure~\ref{fig:potret-singular} displays the singular value portrait of the same matrix:
truncating after the first term kills $\sigma_{2}$, and the remaining error is exactly
$\sigma_{2}=2.828427$.

\begin{figure}[htbp]
\centering
\begin{tikzpicture}
\begin{axis}[
  width=9.6cm, height=5.4cm,
  ybar, bar width=15pt,
  xmin=0.35, xmax=2.65, ymin=0, ymax=6.4,
  ylabel={singular value},
  xtick={1,2}, xticklabels={$\sigma_{1}$,$\sigma_{2}$},
  tick label style={font=\small},
  label style={font=\small},
  axis lines=left,
  legend style={font=\small,draw=none,at={(0.03,0.97)},anchor=north west},
  legend cell align=left,
]
\addplot[fill=accent!55,draw=accent,bar shift=-8.5pt]
  coordinates {(1,4.243) (2,2.828)};
\addlegendentry{full $\mathbf{A}$}
\addplot[fill=blue!25,draw=blue!55!black,bar shift=8.5pt]
  coordinates {(1,4.243) (2,0)};
\addlegendentry{$k=1$}
\draw[<->,gray,thick] (axis cs:2.30,0)--(axis cs:2.30,2.828)
  node[midway,right=2pt,font=\scriptsize] {$\sigma_{2}$};
\draw[gray,densely dotted] (axis cs:2.03,2.828)--(axis cs:2.30,2.828);
\end{axis}
\end{tikzpicture}
\caption{Singular value portrait of $\mathbf{A}=\begin{pmatrix}0&3\\4&1\end{pmatrix}$.
Red bars: full $\mathbf{A}$, $\sigma_{1}=4.243$ and $\sigma_{2}=2.828$. Blue
bars: the rank-one approximation $\sigma_{1}\mathbf{u}_{1}\mathbf{v}_{1}^{\top}$, which
leaves zero in the second slot. The difference of their tips, namely $\sigma_{2}=2.828$,
is the Frobenius error of that approximation.}
\label{fig:potret-singular}
\end{figure}

Several consequences follow directly, and all of them are used in the next section.

\begin{akibat}
For $\mathbf{A}\in\mathbb{R}^{m\times n}$ with singular values
$\sigma_{1}\ge\dots\ge\sigma_{r}>0$:
\begin{enumerate}
  \item $\operatorname{rank}\mathbf{A}=r$, the number of nonzero singular values;
  \item $\|\mathbf{A}\mathbf{x}\|\le\sigma_{1}\|\mathbf{x}\|$ for every $\mathbf{x}$,
        and that bound is attained;
  \item $\|\mathbf{A}\|_{F}^{2}=\sum_{i}\sigma_{i}^{2}$ and
        $\|\mathbf{A}\|_{2}=\sigma_{1}$;
  \item if $\mathbf{A}$ is square and invertible,
        $|\det\mathbf{A}|=\prod_{i}\sigma_{i}$ and
        $\kappa(\mathbf{A})=\sigma_{1}/\sigma_{n}$.
\end{enumerate}
\end{akibat}

\begin{proof}
First item: $\mathbf{A}\mathbf{x}=\sum_{i}\sigma_{i}\mathbf{u}_{i}
(\mathbf{v}_{i}^{\top}\mathbf{x})$, so the range of $\mathbf{A}$ is spanned by
$\mathbf{u}_{1},\dots,\mathbf{u}_{r}$, which are linearly independent. Second item:
$\|\mathbf{A}\mathbf{x}\|^{2}=\sum_{i}\sigma_{i}^{2}(\mathbf{v}_{i}^{\top}\mathbf{x})^{2}
\le\sigma_{1}^{2}\|\mathbf{x}\|^{2}$, with equality at $\mathbf{x}=\mathbf{v}_{1}$.
Third item: the cross terms vanish, so the squared Frobenius norm is the sum of
$\sigma_{i}^{2}$. Fourth item: $|\det\mathbf{A}|=|\det\mathbf{U}\det\boldsymbol{\Sigma}
\det\mathbf{V}^{\top}|=\prod_{i}\sigma_{i}$, since the determinant of an orthogonal matrix is $\pm1$.
\end{proof}

The number $\kappa$ in the fourth item is the \emph{condition number}, the ratio of the largest
singular value to the smallest; for a matrix of rank less than $n$ the ratio is infinite.
A large $\kappa$ means one direction is almost flat, so a small error in the data becomes a
large error in the answer; the sixth section uses it to compare two ways of
solving least squares.\index{condition number}

One warning: singular values are not eigenvalues. The matrix
$\begin{pmatrix}0&1\\0&0\end{pmatrix}$ has eigenvalues $0$ and $0$, while its singular
values are $1$ and $0$. The two agree only when the matrix is symmetric and positive
semidefinite.

\section{Pseudoinverse and minimum-norm solutions}
\tingkat{advanced}

Now we return to the attacker's question. The score is one number, the counters are three,
and the set of changes that give the same shift contains infinitely many answers. All those
answers are correct; what separates them is the size of the change, and the attacker wants
the smallest.

\begin{untukanak}
Imagine you have to go to the shop at the end of the street. There are many roads: through the
market, across the field, or around by the bridge. All of them reach the same shop. If all
you count is ``arrive or not'', every road is equally good. But if you count
footsteps, only one road is shortest, and that road is usually the
straightest.

Remember the two rulers that differ by one hundredth of a centimeter? A difference that small
drowns easily in reading mistakes. That is why two roads that are the same mathematically can
give very different answers when the numbers are computed by a machine.

In this chapter we will learn why among many roads there is one that is shortest, why that
shortest road is always perpendicular to every other road that does not change the
destination, and why a machine sometimes chooses a road that is not the most direct.
\end{untukanak}

The tool that answers this is the generalized inverse. $\mathbf{A}$ is not necessarily square, so it
does not necessarily have an inverse. The singular value decomposition gives a way to invert ``as much
as can be inverted'': invert every $\sigma_{i}$ that is nonzero, and let a direction that is
already lost stay lost.

\kamus{pseudoinverse}{pseudoinvers}{An inverse that still works when some directions are already
flat: what can be inverted is inverted, what is lost is left lost.}

\begin{definisi}[Pseudoinverse]
Suppose $\mathbf{A}=\mathbf{U}\boldsymbol{\Sigma}\mathbf{V}^{\top}$ has rank $r$.
The \emph{pseudoinverse} $\mathbf{A}^{+}\in\mathbb{R}^{n\times m}$ is defined as
\[
  \mathbf{A}^{+}=\mathbf{V}\boldsymbol{\Sigma}^{+}\mathbf{U}^{\top},
  \qquad
  \boldsymbol{\Sigma}^{+}=\operatorname{diag}(1/\sigma_{1},\dots,1/\sigma_{r},0,\dots,0),
\]
that is, the matrix $\boldsymbol{\Sigma}$ transposed and then with every nonzero entry
inverted.\index{pseudoinverse}
\end{definisi}

$\mathbf{A}^{+}$ always exists, even for the zero matrix. If $\mathbf{A}$ is square and
invertible, all $\sigma_{i}>0$ and $\mathbf{A}^{+}=\mathbf{A}^{-1}$; this definition
extends the inverse, it does not replace it.

The vector $\mathbf{A}^{+}\mathbf{b}$ answers two demands at once: the residual is smallest, and
among those with the smallest residual, its length is shortest.

\begin{teorema}
For every $\mathbf{b}\in\mathbb{R}^{m}$, the vector
$\mathbf{x}^{+}=\mathbf{A}^{+}\mathbf{b}$ satisfies
$\|\mathbf{A}\mathbf{x}^{+}-\mathbf{b}\|\le\|\mathbf{A}\mathbf{x}-\mathbf{b}\|$ for
all $\mathbf{x}$, and among all vectors satisfying that inequality,
$\mathbf{x}^{+}$ has the smallest norm.
\end{teorema}

\begin{proof}
Write $\mathbf{A}=\mathbf{U}\boldsymbol{\Sigma}\mathbf{V}^{\top}$ and let
$\mathbf{y}=\mathbf{V}^{\top}\mathbf{x}$. Because $\mathbf{V}$ is square with orthonormal
columns, $\|\mathbf{x}\|=\|\mathbf{y}\|$ and the map $\mathbf{x}\mapsto\mathbf{y}$ is
bijective. Expand $\mathbf{b}$ along the columns of $\mathbf{U}$:
$\mathbf{b}=\sum_{i=1}^{r}b_{i}\mathbf{u}_{i}+\mathbf{b}_{\perp}$ with
$b_{i}=\langle\mathbf{u}_{i},\mathbf{b}\rangle$ and $\mathbf{b}_{\perp}$ orthogonal
to all $\mathbf{u}_{i}$. Then
\[
  \mathbf{A}\mathbf{x}-\mathbf{b}
  =\sum_{i=1}^{r}(\sigma_{i}y_{i}-b_{i})\mathbf{u}_{i}+\sum_{i=r+1}^{m}0\cdot\mathbf{u}_{i}
   -\mathbf{b}_{\perp},
\]
and the three groups of terms are mutually orthogonal. Therefore
\[
  \|\mathbf{A}\mathbf{x}-\mathbf{b}\|^{2}
  =\sum_{i=1}^{r}(\sigma_{i}y_{i}-b_{i})^{2}+\|\mathbf{b}_{\perp}\|^{2}.
\]
The last term does not depend on $\mathbf{x}$, so the minimum is attained with
$y_{i}=b_{i}/\sigma_{i}$ for $i\le r$. The coordinates $y_{r+1},\dots,y_{n}$ appear in
no formula: whatever their values, the squared residual is the same. But
$\|\mathbf{y}\|^{2}=\sum_{i\le r}(b_{i}/\sigma_{i})^{2}+\sum_{i>r}y_{i}^{2}$, so the
minimum norm is attained when $y_{i}=0$ for $i>r$. The resulting vector is
$\mathbf{x}^{+}=\sum_{i\le r}(b_{i}/\sigma_{i})\mathbf{v}_{i}
=\mathbf{V}\boldsymbol{\Sigma}^{+}\mathbf{U}^{\top}\mathbf{b}$.
\end{proof}

That proof contains the answer to the attacker's question. On an underdetermined system, $n>r$, and
the coordinates $y_{r+1},\dots,y_{n}$ do not affect the residual at all; those coordinates
lie along $\mathbf{v}_{r+1},\dots,\mathbf{v}_{n}$, exactly the vectors that span
the kernel of $\mathbf{A}$. So the set of all solutions is $\mathbf{x}^{+}+\ker\mathbf{A}$, an
affine space of dimension $n-r$, and the pseudoinverse picks the point with zero kernel component:
the only one orthogonal to the kernel, that is, in the row space of $\mathbf{A}$.
Every other solution is $\mathbf{x}^{+}$ plus a nonzero vector in the kernel, and that
addition increases the length by Pythagoras:
$\|\mathbf{x}^{+}+\mathbf{k}\|^{2}=\|\mathbf{x}^{+}\|^{2}+\|\mathbf{k}\|^{2}$.

For the two extreme cases there are closed formulas.

\begin{akibat}
If $\mathbf{A}$ has full column rank, $\mathbf{A}^{+}
=(\mathbf{A}^{\top}\mathbf{A})^{-1}\mathbf{A}^{\top}$. If $\mathbf{A}$ has full
row rank, $\mathbf{A}^{+}=\mathbf{A}^{\top}(\mathbf{A}\mathbf{A}^{\top})^{-1}$.
\end{akibat}

\begin{proof}
For the first case we repeat the argument with the SVD. Take
$\mathbf{A}=\mathbf{U}\boldsymbol{\Sigma}\mathbf{V}^{\top}$ with $\mathbf{V}$ square and
orthonormal; then
\[
  (\mathbf{A}^{\top}\mathbf{A})^{-1}\mathbf{A}^{\top}
  =\mathbf{V}\boldsymbol{\Sigma}^{-2}\mathbf{V}^{\top}\mathbf{V}\boldsymbol{\Sigma}\mathbf{U}^{\top}
  =\mathbf{V}\boldsymbol{\Sigma}^{-1}\mathbf{U}^{\top}=\mathbf{A}^{+}.
\]
The second case is obtained by applying the first to $\mathbf{A}^{\top}$ and then
transposing: $(\mathbf{A}^{\top})^{+}=(\mathbf{A}\mathbf{A}^{\top})^{-1}\mathbf{A}$,
and its transpose is $\mathbf{A}^{\top}(\mathbf{A}\mathbf{A}^{\top})^{-1}$. Since
$\mathbf{A}^{+}=(\mathbf{A}^{\top})^{+}{}^{\top}$ by definition, the result is correct.
\end{proof}

That second consequence is what answers the attacker's problem. The $2\times3$ detector matrix has full
row rank, so the minimum change is
$\mathbf{A}^{\top}(\mathbf{A}\mathbf{A}^{\top})^{-1}\boldsymbol{\delta}$, and that formula
needs only the inverse of a $2\times2$ matrix. Note that
$\mathbf{A}\mathbf{A}^{\top}$ is invertible precisely because its rows are linearly independent, not
its columns.

\begin{tantangan}
Take
\[
  \mathbf{A}=\begin{pmatrix}1&0\\0&1\\1&1\end{pmatrix},\qquad \mathbf{b}=(1,1,3).
\]
The system $\mathbf{A}\mathbf{x}=\mathbf{b}$ has no solution: the first two rows force
$\mathbf{x}=(1,1)$, while the third row demands $x_{1}+x_{2}=3$.

There is still one least-squares solution, because the columns are linearly independent. Determine $\hat{\mathbf{x}}$, then find all values of $c$ such that for $\mathbf{b}=(1,1,c)$ the length of the least-squares residual is exactly $1$. There are two answers, and both are symmetric about $c=2$.

Hint 1: the residual is always parallel to one fixed direction; find that direction first.

Hint 2: the distance from $\mathbf{b}$ to the column plane is $|c-2|/\sqrt{3}$, so only one equation remains to be solved.

If you get stuck, go back to the section Least squares and the normal equations.
\end{tantangan}

\section{Why SVD is safer than the normal equations}
\tingkat{advanced}

The normal equations and the pseudoinverse give the same answer mathematically. Numerically
the two differ widely, and the reason is one line: the normal equations work on
$\mathbf{A}^{\top}\mathbf{A}$, whose singular values are $\sigma_{i}^{2}$, not
$\sigma_{i}$. Since the condition number is a ratio of singular values, it is squared as well.

\begin{akibat}
If $\mathbf{A}$ has full column rank, $\kappa(\mathbf{A}^{\top}\mathbf{A})
=\kappa(\mathbf{A})^{2}$.
\end{akibat}

\begin{proof}
From $\mathbf{A}=\mathbf{U}\boldsymbol{\Sigma}\mathbf{V}^{\top}$ we get
$\mathbf{A}^{\top}\mathbf{A}=\mathbf{V}\boldsymbol{\Sigma}^{2}\mathbf{V}^{\top}$, with
$\mathbf{V}$ orthogonal. So the singular values of $\mathbf{A}^{\top}\mathbf{A}$ are
$\sigma_{i}^{2}$, and the ratio of largest to smallest becomes
$\sigma_{1}^{2}/\sigma_{n}^{2}=(\sigma_{1}/\sigma_{n})^{2}$.
\end{proof}

It means this. In binary fractions with about sixteen digits, the relative error of the
answer is roughly $\kappa$ times the relative error of the data: at $\kappa=10^{7}$ nine
digits remain, at $\kappa^{2}=10^{14}$ two remain, and at $\kappa=10^{8}$ every digit is
gone.

The numbers in the following table were computed with numpy and scipy, against a 50-digit reference
from mpmath. The first row is the probe calibration from the chapter opening; the second row is a
degree-nine polynomial fit on twelve points, a problem that looks
harmless.

\begin{center}
\begin{tabular}{lrrr}
\hline
case & $\kappa(\mathbf{A})$ & relative error normal equations & relative error SVD \\
\hline
calibration $5\times2$ & $7.36\times10^{3}$ & $9.10\times10^{-11}$ & $4.72\times10^{-13}$ \\
polynomial $12\times10$ & $6.84\times10^{6}$ & $2.87\times10^{-4}$ & $4.90\times10^{-12}$ \\
\hline
\end{tabular}
\end{center}

Both rows follow the same pattern. In the first row, $\kappa=7.36\times10^{3}$ and
$\kappa^{2}=5.42\times10^{7}$; the normal equations differ from the
reference by $9.10\times10^{-11}$, the SVD by $4.72\times10^{-13}$, a ratio of $193$. In the second row,
$\kappa(\mathbf{A})=6.84\times10^{6}$ and $\kappa(\mathbf{A}^{\top}\mathbf{A})
=4.68\times10^{13}$; the normal equations differ by $2.87\times10^{-4}$, the SVD
by $4.90\times10^{-12}$, a ratio of $5.9\times10^{7}$. That difference grows together with
$\kappa$, far faster for the normal equations; both are still below the bound
$\kappa^{2}\epsilon$ because that bound is a worst case, not a prediction. What matters is the
order, not the absolute numbers.

The squared residuals of the two roads are nearly the same, $4.47019\times10^{-6}$ versus
$4.47018\times10^{-6}$, even though the coefficients differ widely: a small residual does not prove the
coefficients are right.

The honest conclusion: SVD does not remove the sensitivity, it only removes one source of
magnification, namely the squaring, so the error grows like $\kappa$ and not like
$\kappa^{2}$. When $\kappa$ approaches $10^{16}$, no double precision
saves you; what remains is regularization, that is, deliberately discarding directions with
small $\sigma_{i}$, equivalent to approximating $\mathbf{A}$ by a matrix of lower
rank. Chapter 26 discusses how to choose the threshold.

\section{Worked examples}
\tingkat{intermediate}

The four computations below use simple fractions, so they can be checked again by
hand without a machine.

\emph{First, projection onto a plane.} Take
$\mathbf{a}_{1}=(1,1,0)$, $\mathbf{a}_{2}=(1,0,1)$, and $\mathbf{b}=(3,1,3)$ in
$\mathbb{R}^{3}$. The subspace $W=\operatorname{span}(\mathbf{a}_{1},\mathbf{a}_{2})$
has dimension two. Gram-Schmidt gives
\[
  \mathbf{q}_{1}=\frac{1}{\sqrt{2}}(1,1,0),\qquad
  \mathbf{u}_{2}=\mathbf{a}_{2}-\frac{1}{\sqrt{2}}\mathbf{q}_{1}
  =\Bigl(\tfrac12,-\tfrac12,1\Bigr),\qquad
  \mathbf{q}_{2}=\frac{1}{\sqrt{6}}(1,-1,2),
\]
because $\|\mathbf{u}_{2}\|=\sqrt{3/2}=\sqrt{6}/2$. The projection coefficients are
$\langle\mathbf{b},\mathbf{q}_{1}\rangle=4/\sqrt{2}=2\sqrt{2}$ and
$\langle\mathbf{b},\mathbf{q}_{2}\rangle=8/\sqrt{6}$. So
\[
  \mathbf{p}=2\sqrt{2}\,\mathbf{q}_{1}+\frac{8}{\sqrt{6}}\,\mathbf{q}_{2}
  =(2,2,0)+\frac{4}{3}(1,-1,2)
  =\Bigl(\tfrac{10}{3},\tfrac23,\tfrac83\Bigr).
\]
The residual is $\mathbf{r}=\mathbf{b}-\mathbf{p}=(-\tfrac13,\tfrac13,\tfrac13)$, with
$\|\mathbf{r}\|=1/\sqrt{3}=0.577350$. Check: $\mathbf{a}_{1}^{\top}\mathbf{r}=0$
and $\mathbf{a}_{2}^{\top}\mathbf{r}=0$, so the residual is indeed perpendicular to both spanning vectors.
The distance from $\mathbf{b}$ to the plane is $0.577$; no other point in $W$ is closer.

\emph{Second, least squares $5\times2$.} For the chapter's opening data,
$\mathbf{A}$ has columns $(1,1,1,1,1)^{\top}$ and $(100,101,102,103,104)^{\top}$, so
\[
  \mathbf{A}^{\top}\mathbf{A}=
  \begin{pmatrix}5&510\\ 510&52030\end{pmatrix},
  \qquad
  \mathbf{A}^{\top}\mathbf{y}=\begin{pmatrix}1540\\ 157109.7\end{pmatrix},
  \qquad
  \det(\mathbf{A}^{\top}\mathbf{A})=5\cdot52030-510^{2}=50 .
\]
Cramer's rule gives
\begin{align*}
  a&=\frac{1540\cdot52030-510\cdot157109.7}{50}=\frac{253}{50}=5.06,\\
  b&=\frac{5\cdot157109.7-510\cdot1540}{50}=\frac{148.5}{50}=2.97 .
\end{align*}
The predicted readings become $302.06$, $305.03$, $308.00$, $310.97$, and $313.94$,
whereas the actual data are $302.0$, $305.1$, $307.9$, $311.2$, $313.8$. The residual is
\[
  \mathbf{r}=(-0.06,\;0.07,\;-0.10,\;0.23,\;-0.14),
  \qquad
  \|\mathbf{r}\|^{2}=0.0910,\qquad \|\mathbf{r}\|=0.301662 .
\]
Two checks must pass if the computation is right: $\sum_{i}r_{i}=0$ and
$\sum_{i}t_{i}r_{i}=0$. The first because $\mathbf{r}$ is perpendicular to the column of ones, the
second because it is perpendicular to the column $t$. Both are indeed zero here.

\emph{Third, SVD of a $3\times2$ matrix.} Take
$\mathbf{A}=\begin{pmatrix}1&1\\ 1&1\\ 1&-1\end{pmatrix}$. Then
\[
  \mathbf{A}^{\top}\mathbf{A}=\begin{pmatrix}3&1\\ 1&3\end{pmatrix},
  \qquad
  \lambda^{2}-6\lambda+8=0\;\Longrightarrow\;\lambda=4,\,2,
\]
so $\sigma_{1}=2$ and $\sigma_{2}=\sqrt{2}=1.41421$, with
$\kappa(\mathbf{A})=\sqrt{2}$. Its right singular vectors are
$\mathbf{v}_{1}=\frac{1}{\sqrt{2}}(1,1)$ and
$\mathbf{v}_{2}=\frac{1}{\sqrt{2}}(1,-1)$; the left singular vectors are computed from
$\mathbf{u}_{i}=\mathbf{A}\mathbf{v}_{i}/\sigma_{i}$, namely
$\mathbf{u}_{1}=\frac{1}{\sqrt{2}}(1,1,0)$ and $\mathbf{u}_{2}=(0,0,1)$. Check one entry:
$(2\,\mathbf{u}_{1}\mathbf{v}_{1}^{\top}+\sqrt{2}\,\mathbf{u}_{2}\mathbf{v}_{2}^{\top})_{32}
=-\sqrt{2}/\sqrt{2}=-1$, which matches $\mathbf{A}$.

\emph{Fourth, a rank-one matrix.} Take
$\mathbf{B}=\begin{pmatrix}1&2\\ 2&4\\ 3&6\end{pmatrix}$, whose second column is twice
the first column. Then $\mathbf{B}^{\top}\mathbf{B}=14\begin{pmatrix}1&2\\ 2&4\end{pmatrix}$
with eigenvalues $70$ and $0$, so $\sigma_{1}=\sqrt{70}=8.36660$ and
$\sigma_{2}=0$. Its rank is one, and $\kappa$ is infinite. For $\mathbf{b}=(1,0,0)$,
the projection onto the column space is $\frac{1}{14}(1,2,3)$, so $\mathbf{B}\mathbf{x}$
must equal that vector. Since $\mathbf{B}(s,2s)=5s(1,2,3)$, we need
$5s=1/14$, that is, $s=1/70$. The minimum-norm solution is
\[
  \mathbf{x}^{+}=\Bigl(\tfrac{1}{70},\tfrac{2}{70}\Bigr)=(0.0142857,\;0.0285714),
  \qquad
  \|\mathbf{x}^{+}\|=\frac{\sqrt{5}}{70}=0.031944 .
\]
The least-squares residual:
\[
  \mathbf{r}=\mathbf{b}-\mathbf{B}\mathbf{x}^{+}
  =\Bigl(\tfrac{13}{14},-\tfrac17,-\tfrac{3}{14}\Bigr),
  \qquad
  \|\mathbf{r}\|=\sqrt{13/14}=0.963624 .
\]
That residual cannot be zero: $\mathbf{b}$ lies outside the column space. Other solutions still give the
same residual, for instance $\mathbf{x}^{+}+(-\tfrac27,\tfrac17)$, but their norm is
larger.

\kertas{Repeat the second computation with the last reading $314.3$ and not $313.8$.
Predict first whether $a$ goes up or down, then compute. Note that the point
with the largest $t$ shifts the line more than a point in the middle; explain this through the
columns of $\mathbf{A}$ and the formula $(\mathbf{A}^{\top}\mathbf{A})^{-1}\mathbf{A}^{\top}$.
Compute the new $\|\mathbf{r}\|$ as well and compare it with $0.301662$.}

\section{Computational lab}
\tingkat{advanced}

The two blocks below were run with numpy 2.5.2, scipy 1.18.0, Python 3.14.4; the output is
copied as it stands.

\begin{lab}
\begin{verbatim}
import numpy as np
from scipy.linalg import lstsq, qr

np.set_printoptions(precision=6, suppress=True, linewidth=72)

# 1. projection onto a plane
a1 = np.array([1., 1., 0.])
a2 = np.array([1., 0., 1.])
b = np.array([3., 1., 3.])
Q, _ = qr(np.column_stack([a1, a2]), mode="economic")
P = Q @ Q.T
print("Q =\n", Q)
print("P b =", P @ b, " residual =", b - P @ b)
nr = np.linalg.norm(b - P @ b)
print("residual norm =", nr, " 1/sqrt3 =", 1/np.sqrt(3))
print("P^2 = P:", np.allclose(P @ P, P), " P^T = P:", np.allclose(P.T, P))

# 2. least squares 5x2
t = np.array([100., 101., 102., 103., 104.])
A = np.column_stack([np.ones(5), t])
y = np.array([302.0, 305.1, 307.9, 311.2, 313.8])
print("\nA^T A =\n", A.T @ A, " det =", np.linalg.det(A.T @ A))
print("A^T y =", A.T @ y)
xn = np.linalg.solve(A.T @ A, A.T @ y)
xs = lstsq(A, y)[0]
print("normal equations:", xn)
print("lstsq           :", xs)
print("solution gap    :", np.linalg.norm(xn - xs))
cA = np.linalg.cond(A)
print("cond A =", cA, " cond A^T A =", np.linalg.cond(A.T @ A))
r = y - A @ xn
print("residual =", r, " norm =", np.linalg.norm(r))
print("A^T r =", A.T @ r)

# 3. SVD of a 3x2 matrix
S = np.array([[1., 1.], [1., 1.], [1., -1.]])
U, s, Vt = np.linalg.svd(S, full_matrices=False)
print("\nsingular:", s)
print("rank", np.linalg.matrix_rank(S), " cond", np.linalg.cond(S))
print("U =\n", U, "\nV^T =\n", Vt)
print("reconstruction error:", np.linalg.norm(U @ np.diag(s) @ Vt - S))
\end{verbatim}
Actual output:
\begin{verbatim}
Q =
 [[-0.707107  0.408248]
 [-0.707107 -0.408248]
 [-0.        0.816497]]
P b = [3.333333 0.666667 2.666667]  residual = [-0.333333  0.333333  0.333333]
residual norm = 0.5773502691896251  1/sqrt3 = 0.5773502691896258
P^2 = P: True  P^T = P: True

A^T A =
 [[    5.   510.]
 [  510. 52030.]]  det = 49.999999999983494
A^T y = [  1540.  157109.7]
normal equations: [5.06 2.97]
lstsq           : [5.06 2.97]
solution gap    : 5.315035090945416e-10
cond A = 7358.860135917719  cond A^T A = 54152822.48728771
residual = [-0.06  0.07 -0.1   0.23 -0.14]  norm = 0.30166206257998296
A^T r = [0. 0.]

singular: [2.       1.414214]
rank 2  cond 1.414213562373095
U =
 [[-0.707107 -0.      ]
 [-0.707107 -0.      ]
 [-0.        1.      ]] 
V^T =
 [[-0.707107 -0.707107]
 [ 0.707107 -0.707107]]
reconstruction error: 6.568167990716596e-16
\end{verbatim}
\end{lab}

Three things need reading. First, $\mathbf{Q}$ from scipy differs in sign from the hand result:
its columns are $-(1,1,0)/\sqrt{2}$ and $-(1,-1,2)/\sqrt{6}$. The sign of the columns in QR is not unique,
and $\mathbf{Q}\mathbf{Q}^{\top}$ does not change because of it; the entries of $\mathbf{P}$ stay
$(10/3,2/3,8/3)$, with residual $(-0.333333,$ $0.333333,$ $0.333333)$.

Second, the value of $\det(\mathbf{A}^{\top}\mathbf{A})$ is reported as $49.999999999983494$ and not
$50$; the difference $1.7\times10^{-11}$ is purely rounding. What matters is
$\kappa(\mathbf{A}^{\top}\mathbf{A})=5.415\times10^{7}$, the square of
$\kappa(\mathbf{A})=7.359\times10^{3}$.

Third, $\mathbf{A}^{\top}\mathbf{r}$ is reported as $(0,0)$ to six digits and the difference between
the two solutions is $5.32\times10^{-10}$: they only differ in the ninth digit, and that
is no guarantee.

\begin{lab}
\begin{verbatim}
import numpy as np
import mpmath as mp
from scipy.linalg import lstsq

mp.mp.dps = 50
np.set_printoptions(precision=4, suppress=True, linewidth=72)

# 4. rank one: nonzero residual, minimum-norm solution
B = np.array([[1., 2.], [2., 4.], [3., 6.]])
b2 = np.array([1., 0., 0.])
sB = np.linalg.svd(B, compute_uv=False)
print("singular B:", sB, " rank", np.linalg.matrix_rank(B))
xb = np.linalg.pinv(B) @ b2
print("x+ =", xb, " B x+ =", B @ xb)
nr = np.linalg.norm(B @ xb - b2)
print("residual norm =", nr, " sqrt(13/14) =", np.sqrt(13/14))

# 5. underdetermined system: minimum norm
C = np.array([[1., 1., 1.], [0., 1., 1.]])
b3 = np.array([3., 1.])
xp = np.linalg.pinv(C) @ b3
print("\nx+ =", xp, " C x+ =", C @ xp, " norm =", np.linalg.norm(xp))
print("norm (2,1,0) =", np.linalg.norm([2., 1., 0.]))

# 6. least squares: normal equations vs lstsq vs 50-digit reference
x = np.linspace(0, 1, 12)
V = np.vander(x, 10, increasing=True)
b = V @ (0.1 * np.arange(1, 11)) + 0.001 * np.sin(np.arange(12))
Vm, bm = mp.matrix(V.tolist()), mp.matrix(b.tolist())
ref = np.array([float(v) for v in mp.lu_solve(Vm.T * Vm, Vm.T * bm)])
xne = np.linalg.solve(V.T @ V, V.T @ b)
xls = lstsq(V, b)[0]
rel = lambda u: np.linalg.norm(u - ref) / np.linalg.norm(ref)
print("\ncond V     :", np.linalg.cond(V))
print("cond V^T V :", np.linalg.cond(V.T @ V))
print("relative error normal equations:", rel(xne))
print("relative error lstsq           :", rel(xls))
print("residual normal equations:", np.linalg.norm(V @ xne - b))
print("residual lstsq           :", np.linalg.norm(V @ xls - b))

# 7. classical vs modified Gram-Schmidt vs Householder
def cgs(M):
    m, n = M.shape
    Q = np.zeros((m, n)); R = np.zeros((n, n))
    for j in range(n):
        v = M[:, j].astype(float).copy()
        for i in range(j):
            R[i, j] = Q[:, i] @ M[:, j]
            v -= R[i, j] * Q[:, i]
        R[j, j] = np.linalg.norm(v)
        Q[:, j] = v / R[j, j]
    return Q

def mgs(M):
    m, n = M.shape
    Q = np.zeros((m, n)); R = np.zeros((n, n)); W = M.astype(float).copy()
    for j in range(n):
        for i in range(j):
            R[i, j] = Q[:, i] @ W[:, j]
            W[:, j] -= R[i, j] * Q[:, i]
        R[j, j] = np.linalg.norm(W[:, j])
        Q[:, j] = W[:, j] / R[j, j]
    return Q

ort = lambda Q: np.linalg.norm(Q.T @ Q - np.eye(Q.shape[1]))
Qh, _ = np.linalg.qr(V)
print("\n||Q^T Q - I||  classical  :", ort(cgs(V)))
print("||Q^T Q - I||  modified   :", ort(mgs(V)))
print("||Q^T Q - I||  Householder:", ort(Qh))
\end{verbatim}
Actual output:
\begin{verbatim}
singular B: [8.3666 0.    ]  rank 1
x+ = [0.0143 0.0286]  B x+ = [0.0714 0.1429 0.2143]
residual norm = 0.9636241116594315  sqrt(13/14) = 0.9636241116594315

x+ = [2.  0.5 0.5]  C x+ = [3. 1.]  norm = 2.121320343559642
norm (2,1,0) = 2.23606797749979

cond V     : 6841129.486331141
cond V^T V : 46788641042301.25
relative error normal equations: 0.0002869502788738492
relative error lstsq           : 4.899523677448492e-12
residual normal equations: 4.4701872339775605e-06
residual lstsq           : 4.47017846708843e-06

|||Q^T Q - I||  classical  : 0.020506481578936916
|||Q^T Q - I||  modified   : 3.397900311236548e-10
|||Q^T Q - I||  Householder: 1.1933579412809552e-15
\end{verbatim}
\end{lab}

First reading: the least-squares residual of the rank-one matrix is exactly
$\sqrt{13/14}$; numpy gives the same number to sixteen digits, so the fractional formula
in the seventh section is not an approximation.

Second reading: $\mathbf{C}\mathbf{x}=(3,1)$ gives
$\mathbf{x}^{+}=(2,\,0.5,\,0.5)$ with norm $2.12132$, whereas $(2,1,0)$ gives
$2.23607$; the difference of squared norms is $0.5$, exactly $2t^{2}$ at $t=0.5$ for
$\mathbf{k}=(0,1,-1)$. The pseudoinverse chooses $t=0$.

Third and fourth readings: on the $12\times10$ Vandermonde, the normal equations differ by
$2.87\times10^{-4}$ from the 50-digit reference and \texttt{lstsq} by $4.90\times10^{-12}$,
even though the residuals are practically the same ($4.47019\times10^{-6}$ versus $4.47018\times10^{-6}$).
The coefficients are wrong, the residual is right. For Gram-Schmidt,
classical gives $2.05\times10^{-2}$, modified $3.40\times10^{-10}$, Householder
$1.19\times10^{-15}$.

\begin{untukdosen}
The order that has worked: projection first, the definition of the inner product afterwards.
Students who see the shadow first accept $\mathbf{Q}\mathbf{Q}^{\top}$ as the
answer to the question ``which point is closest'', not as a magic formula; after that
Gram-Schmidt becomes a need, not a memorized fact.

Two typical misconceptions. First, $\mathbf{A}^{\top}\mathbf{A}$ is assumed always invertible
because ``squaring a matrix never gives zero''; use the corollary in the third section, then ask for the
kernel of $\mathbf{A}^{\top}\mathbf{A}$ for a rank-one $2\times2$ matrix. Second, the SVD
is assumed to give eigenvalues; the matrix $\begin{pmatrix}0&1\\0&0\end{pmatrix}$ is enough to
break that.

Class activity: give five clearly non-collinear points, ask half the class
to project $\mathbf{b}$ by hand and half to use \texttt{lstsq}; compare the
residual vectors, not the coefficients.

A sign of readiness to go further: a student can explain why adding one equation does not
always improve the answer, and name one situation in which the pseudoinverse is better
than the inverse.
\end{untukdosen}

\section{Summary}

\begin{ringkasan}
\begin{itemize}
  \item The inner product supplies length and angle; Cauchy-Schwarz supports the triangle
        inequality. In an orthonormal basis
        $\mathbf{v}=\sum_{i}\langle\mathbf{v},\mathbf{q}_{i}\rangle\mathbf{q}_{i}$ and
        $\|\mathbf{v}\|^{2}=\sum_{i}\langle\mathbf{v},\mathbf{q}_{i}\rangle^{2}$.
  \item The projection $\mathbf{p}=\mathbf{Q}\mathbf{Q}^{\top}\mathbf{b}$ is the closest point
        in the subspace, and the residual is perpendicular to the whole subspace. Gram-Schmidt builds an
        orthonormal basis and gives $\mathbf{A}=\mathbf{Q}\mathbf{R}$.
  \item Least squares satisfies $\mathbf{A}^{\top}\mathbf{A}\hat{\mathbf{x}}
        =\mathbf{A}^{\top}\mathbf{b}$; if the columns of $\mathbf{A}$ are linearly independent, the answer is unique.
  \item Every matrix has $\mathbf{A}=\mathbf{U}\boldsymbol{\Sigma}\mathbf{V}^{\top}$:
        the rank is the number of $\sigma_{i}>0$, $\|\mathbf{A}\|_{2}=\sigma_{1}$, and
        $\kappa(\mathbf{A})=\sigma_{1}/\sigma_{n}$.
  \item The pseudoinverse $\mathbf{A}^{+}=\mathbf{V}\boldsymbol{\Sigma}^{+}\mathbf{U}^{\top}$
        gives the smallest residual and at the same time the smallest norm, and on an underdetermined system
        it chooses the answer perpendicular to the kernel.
  \item The normal equations work at $\kappa^{2}$, the SVD at $\kappa$; that is why a badly
        conditioned problem loses digits twice as fast through $\mathbf{A}^{\top}\mathbf{A}$.
\end{itemize}
\end{ringkasan}

\section{Exercises}

\begin{latihan}
\begin{enumerate}
  \item Project $\mathbf{b}=(1,2,2)$ onto $W=\operatorname{span}\{(1,0,0),(0,1,1)\}$
        by hand, check with \texttt{numpy.linalg.qr}, then compute the distance from
        $\mathbf{b}$ to $W$.

  \item Prove $\|\mathbf{b}\|^{2}=\|\operatorname{proj}_{W}\mathbf{b}\|^{2}
        +\|\mathbf{r}\|^{2}$ for arbitrary $W$ and $\mathbf{b}$, then explain why
        that equality fails when $\mathbf{p}$ is not the orthogonal projection. (Hint: take a
        skew basis in $\mathbb{R}^{2}$ for the second part.)

  \item If the columns of $\mathbf{A}$ are orthonormal, prove $\mathbf{A}^{\top}\mathbf{A}=\mathbf{I}$,
        $\hat{\mathbf{x}}=\mathbf{A}^{\top}\mathbf{b}$, and $\kappa(\mathbf{A})=1$.

  \item Evaluate this claim: \emph{if $\mathbf{A}^{\top}\mathbf{A}$ is invertible then
        the rows of $\mathbf{A}$ are linearly independent.} Decide whether it is true, false, or not enough
        evidence to decide, and include a counterexample if needed.

  \item Compute the SVD of $\mathbf{A}=\begin{pmatrix}1&1\\ 1&1\end{pmatrix}$ by hand: singular
        values, $\mathbf{U}$, $\mathbf{V}$, rank. Determine $\mathbf{A}^{+}\mathbf{b}$
        for $\mathbf{b}=(1,3)$, compute the residual, check with
        \texttt{numpy.linalg.pinv}.

  \item Prove that the least-squares residual is always perpendicular to $\operatorname{col}(\mathbf{A})$
        and conversely, then state where linear independence of the columns is used.
        (Hint: use the splitting of the norm as in the proof of the normal equations; the
        ``necessary'' argument chooses the direction $\mathbf{h}$.)

  \item Run Gram-Schmidt on $\{(1,1,1),(1,1,0),(1,0,0)\}$ by hand, compare
        with \texttt{numpy.linalg.qr}, compute $\|\mathbf{Q}^{\top}\mathbf{Q}-\mathbf{I}\|$
        for both, and explain why their signs may differ.

  \item Suppose $m<n$ and the rows of $\mathbf{A}$ are linearly independent. Prove that
        $\mathbf{A}\mathbf{A}^{\top}$ is invertible, and that
        \[ \mathbf{A}^{\top}(\mathbf{A}\mathbf{A}^{\top})^{-1}\mathbf{b} \]
        is the solution perpendicular to the kernel of $\mathbf{A}$. (Hint: imitate the proof of the
        corollary for the full-column case with $\mathbf{A}^{\top}$ as the matrix.)

  \item With the calibration data in the seventh section, change the last reading to
        $314.3$. Recompute $a$ and $b$ through Cramer's rule, compare with
        \texttt{scipy.linalg.lstsq}, and explain why a point at the end shifts the line
        more than a point in the middle.

  \item For $\mathbf{B}=\begin{pmatrix}1&2\\ 2&4\\ 3&6\end{pmatrix}$ and
        $\mathbf{b}=(1,0,0)$, write all least-squares solutions in the form
        $\mathbf{x}^{+}+t\mathbf{k}$, determine $\mathbf{k}$, then show that
        $\|\mathbf{x}^{+}+t\mathbf{k}\|^{2}$ is minimal at $t=0$.
\end{enumerate}
\end{latihan}

\begin{untukpeneliti}
What is proved in this chapter: Cauchy-Schwarz and the triangle inequality for an arbitrary
inner-product space; uniqueness of coefficients in an orthonormal basis; the best-approximation property
of the projection; the correctness of Gram-Schmidt; the equivalence of least squares with the normal
equations; the existence of the SVD; and the pseudoinverse theorem. What is borrowed is only the spectral
lemma, whose proof uses the extreme value theorem on the unit sphere (Chapter 11); here it is only sketched.

Two claims are often popularized. First, the claim that ``the SVD is always stable''. What is true:
the SVD algorithm is backward stable, meaning the computed answer is the exact answer for
a matrix that differs slightly from $\mathbf{A}$; the distance to the answer for $\mathbf{A}$
itself can still grow by a factor of $\kappa$. Second, the claim that truncating small singular values
``removes noise''. Truncation removes a direction, and whether that direction is noise
or signal cannot be decided from $\mathbf{A}$ alone.

Primary sources: Axler 4th edition, sections 6A through 6C and 7E through 7F; Anton 11th edition,
sections 6.1 through 6.4 and 9.4 and 9.5; for the numerical side, Trefethen and Bau,
\emph{Numerical Linear Algebra}, and Golub and Van Loan, \emph{Matrix Computations}
4th edition, sections 2.4 and 5.3.

Open questions. The threshold at which a singular value counts as zero has no parameter-free
answer; the choice depends on the scale of the data and on the error that can be accepted (Chapter 22).
For large matrices, randomized algorithms only guarantee anything probabilistically. The limit of this chapter:
real matrices of finite dimension; for complex matrices the transpose is replaced by the conjugate transpose,
and for infinite-dimensional spaces the SVD is replaced by the decomposition of compact operators, whose
singular values can accumulate at zero.
\end{untukpeneliti}

\peta{This chapter follows Anton 11th edition, sections 6.1 (Inner Products), 6.2 (Angle and
Orthogonality), 6.3 (Gram-Schmidt; QR-Decomposition), 6.4 (Best Approximation; Least
Squares), 9.4 (Singular Value Decomposition), and 9.5 (Data Compression Using SVD).
The abstract counterpart is Axler 4th edition, Chapter 6 sections 6A through 6C and Chapter 7 sections 7E
through 7F, with the spectral theorem in 7B. Strang 6th edition gives the geometric road,
but the copy in the source directory contains only the cover and the preface. The numerical side:
Trefethen and Bau; Golub and Van Loan 2.4, 5.3.}


% ---------------------------- part III ----------------------------
\part{Analysis and Calculus}

This part turns the words ``close'', ``smooth'', and ``shrinking'' into statements that can
be checked with numbers. The prerequisites are Chapters~3--5 and~7. After six chapters you
can prove continuity, differentiate and integrate, exchange a limit with an integral, and
solve first-order differential equations and two-dimensional systems.
\chapter{Continuity and Uniform Continuity}

A program samples $f(x)=1/x$ on a grid of points spaced $10^{-4}$ inside
$(0,1]$. Near $x=10^{-4}$ two neighboring points, namely $10^{-4}$ and
$2\times 10^{-4}$, give outputs $10000$ and $5000$; the difference is $5000$. Near
$x=1$ two neighboring points give $1$ and $1.0001$; the difference is about $10^{-4}$.
The same function, the same grid spacing, two differences seven powers apart.

The question that arises is not ``is $1/x$ continuous'', since the answer is yes at
every point of $(0,1]$. The question is a different one: how much grid spacing must be
used so that \emph{every} output on that interval stays inside a single tolerance? For
tolerance $10^{-2}$, the point $x=1$ demands only a spacing below
$9.901\times 10^{-3}$, whereas the point $x=10^{-4}$ demands a spacing below
$1.000\times 10^{-10}$. The ratio of the two demands is $9.9\times 10^{7}$. No single
spacing serves both, and that is not a defect of the program.

Chapter 4 supplied the limit of a sequence, Chapter 3 supplied the supremum, and neither
is enough: the limit of a function compares $f$ with one fixed number, so it cannot
compare two outputs from two nearby inputs. What is needed is one threshold for the
whole set, together with a way to prove that the threshold exists. This chapter builds
that tool, from the $\varepsilon$-$\delta$ definition to compact sets and the class of
Lipschitz functions, then closes with a function that is continuous at one point only.

\begin{untukanak}
Imagine a fire escape whose steps grow taller toward the top. At the bottom, one stride
the width of a foot raises us a little. At the top, the same stride raises us a hand's
width, then a forearm, then almost the height of a body. Our older sibling stands beside
us and asks one thing: how wide is a safe stride for the whole staircase? The honest
answer for a staircase like that does not exist. The width that is safe below is too
wide above, and the width that is safe above makes us crawl below.

Now the question that makes one curious. Is every staircase like that? There are
staircases whose steps are constant, and staircases whose steps change but never pass a
certain bound. For the second kind, one step width works from bottom to top. The
difference is not in the height of the staircase, but in how steeply its slope changes.

The same thing happens with the graph of a function. There are graphs that can be
controlled by one distance rule over their whole domain, and graphs that demand a
different rule at every point. In this chapter we will know why that difference can be
proved with certainty, and why a domain that is closed and bounded always belongs to the
first kind.
\end{untukanak}

\section{Limits of functions and the delta threshold}
\tingkat{basic}

The word ``approaches'' cannot be used as a definition, because everyone may mean
something different by it. What can be used is a wager: however much precision is
demanded, the output will fall inside that precision as long as the input is close
enough. In everyday language: we install a narrow window around $L$, and demand that $f$
enter that window for all $x$ close enough to $c$.

\kamus{functional limit}{limit fungsi}{a promise that the output can be forced into a
narrow window as long as the input is close enough to the door}

Before the definition, one technical term is needed. The point $c$ is called a
\emph{limit point}\index{limit point} of a set $A$ if every open interval around $c$
contains a point of $A$ other than $c$. Without that condition, the definition of limit
becomes vacuous at an isolated point.

\begin{definisi}[Limit of a function]
\index{limit of a function}
Let $A\subseteq\mathbb{R}$, $c$ a limit point of $A$, and $f\colon A\to\mathbb{R}$.
We write $\lim_{x\to c}f(x)=L$ if for every $\varepsilon>0$ there exists
$\delta>0$ such that for every $x\in A$ with $0<|x-c|<\delta$ we have
$|f(x)-L|<\varepsilon$.
\end{definisi}

Two things are easily missed. First, the condition $0<|x-c|$ means the value of $f$ at
$c$ is not counted; the function may be undefined at $c$ and the limit still exists.
Second, the order of the quantifiers: $\delta$ may depend on $\varepsilon$, but not on
$x$. In the section on uniform continuity, it is precisely that freedom which is at
issue.

This definition is a checking tool, not a recipe for finding $\delta$; $\delta$ is found
by working backward from $|f(x)-L|$. For
\[
f(x)=\frac{x^{2}-1}{x-1}\quad (x\neq 1),\qquad L=2,
\]
algebra gives $f(x)=x+1$, hence $|f(x)-2|=|x-1|$. So $\delta=\varepsilon$
works. For $\varepsilon=10^{-2}$ take $\delta=10^{-2}$: at $x=1+10^{-2}$ the
difference is exactly $10^{-2}$, and that does \emph{not} meet the strict condition; at
$x=1+5\times10^{-3}$ the difference is $5\times10^{-3}$. Therefore $\delta=\varepsilon/2$ is the safe
choice, and $\delta=\varepsilon$ is the cheapest choice whose bound is never
attained from inside.

One thing that may reasonably be demanded of a limit is its uniqueness: a function cannot
stop at two different places at once.

\begin{teorema}[Uniqueness of the limit]
\index{limit of a function}
If $\lim_{x\to c}f(x)=L_{1}$ and $\lim_{x\to c}f(x)=L_{2}$, then $L_{1}=L_{2}$.
\end{teorema}

\begin{proof}
Suppose $L_{1}\neq L_{2}$ and set $\varepsilon=|L_{1}-L_{2}|/2>0$. There is
$\delta_{1}>0$ such that $0<|x-c|<\delta_{1}$ gives $|f(x)-L_{1}|<\varepsilon$, and
there is $\delta_{2}>0$ such that $0<|x-c|<\delta_{2}$ gives $|f(x)-L_{2}|<\varepsilon$.
Take $x$ with $0<|x-c|<\min(\delta_{1},\delta_{2})$; such a point exists because
$c$ is a limit point. Then
\[
|L_{1}-L_{2}|\le|L_{1}-f(x)|+|f(x)-L_{2}|<2\varepsilon=|L_{1}-L_{2}|,
\]
an impossibility. Hence $L_{1}=L_{2}$.
\end{proof}

We also use one bridge that is not proved here: the limit of a function exists and
equals $L$ exactly when every sequence $x_{n}\to c$ with $x_{n}\neq c$ satisfies
$f(x_{n})\to L$. Its name is the sequence criterion; its proof is in Abbott,
\emph{Understanding Analysis}, Chapter 4, section Functional Limits, and in Rudin,
\emph{Principles of Mathematical Analysis}, Chapter 4, section Continuous Functions.
That bridge is what allows the sequences of Chapter 4 to be used for dismantling limits
that do not exist. For $g(x)=x/|x|$, the sequence $1/n$ gives $g\to 1$ while $-1/n$ gives
$g\to -1$; two different limits, so $\lim_{x\to 0}g(x)$ does not exist. That pattern is
used again for uniform continuity, with two sequences whose distance tends to zero.

\section{Continuity and its algebraic properties}
\tingkat{intermediate}

A limit compares $f$ with a number $L$ that stands on its own. Continuity asks for
something stricter: that number must be the value of the function at its point. In
everyday language: the graph does not jump at that point, and its value there really is
the value approached from the left and from the right.

\kamus{continuous}{kontinu}{a graph that can be drawn without lifting the pen at that
point}

\begin{definisi}[Continuity at a single point]
\index{continuous}
Let $f\colon A\to\mathbb{R}$ and $c\in A$. The function $f$ is \emph{continuous at $c$}
if for every $\varepsilon>0$ there exists $\delta>0$ such that for every $x\in A$
with $|x-c|<\delta$ we have $|f(x)-f(c)|<\varepsilon$. The function is called continuous
on $A$ if it is continuous at every point of $A$.
\end{definisi}

Note that $x=c$ is now allowed, and what is demanded is not merely that the limit exists,
but that the limit equals the value of the function. Continuity is therefore local: it
speaks only about the behavior of $f$ around one point.

For $f(x)=x^{2}$ at $c=2$ and $\varepsilon>0$ we have $|x^{2}-4|=|x-2|\,|x+2|$, and
if $|x-2|<1$ then $|x+2|<5$. So $\delta=\min(1,\varepsilon/5)$ works, whereas the
largest $\delta$ that is still valid is the positive root of $\delta(4+\delta)=\varepsilon$,
namely $2.49984376953\times10^{-4}$ for $\varepsilon=10^{-3}$. The gap between the two numbers
is the price of safety paid for not wanting to compute a root.
Figure~\ref{fig:epsdelta} shows the two thresholds for $\varepsilon=1$.

\begin{figure}[htbp]
\centering
\begin{tikzpicture}
\begin{axis}[width=0.82\linewidth,
  axis lines=left,xmin=1.5,xmax=2.6,ymin=2.6,ymax=5.8,
  xlabel={$x$},ylabel={$f(x)=x^{2}$},
  xtick={1.7639,2,2.2361},xticklabels={$2-\delta$,$2$,$2+\delta$},
  ytick={3,4,5},yticklabels={$4-\varepsilon$,$4$,$4+\varepsilon$},
  tick label style={font=\scriptsize}]
\fill[blue!8,draw=blue!45] (1.5,3) rectangle (2.6,5);
\fill[orange!14,draw=orange!60] (1.7639,2.6) rectangle (2.2361,5.8);
\draw[densely dashed,black!60] (1.7,2.6) rectangle (2.3,5.8);
\addplot[thick,domain=1.5:2.6] {x^2};
\draw[fill=white] (2.3,5.29) circle (2pt);
\node[font=\scriptsize,anchor=east] at (2.32,5.5) {$\delta=0.3$};
\end{axis}
\end{tikzpicture}
\caption{The $\varepsilon$ band and the $\delta$ band for $f(x)=x^{2}$ at $a=2$ with
$\varepsilon=1$. The $\delta$ bound uses the formula $\delta=\sqrt{4+\varepsilon}-2$
in this section, so $\delta=\sqrt{5}-2\approx0.236$; the dashed band uses
$\delta=0.3$, and at $x=2.3$ the value $f(2.3)=5.29$ has already left the band.}
\label{fig:epsdelta}
\end{figure}

\begin{proposisi}
\index{continuous}
Let $c$ be a limit point of $A$. The function $f$ is continuous at $c$ if and only if
$\lim_{x\to c}f(x)=f(c)$.
\end{proposisi}

\begin{proof}
If $f$ is continuous at $c$ and $\varepsilon>0$, take $\delta$ from the definition of
continuity. For $0<|x-c|<\delta$ we have $|f(x)-f(c)|<\varepsilon$, and that is exactly
the definition of limit with $L=f(c)$. Conversely, if the limit is $f(c)$, take $\delta$
from the definition of limit; the condition $0<|x-c|$ can be dropped because for $x=c$
the inequality becomes $0<\varepsilon$, which is true.
\end{proof}

Continuity is closed under the operations usually used. In everyday language: adding,
multiplying, and dividing functions that do not jump produces functions that also do not
jump, as long as the divisor is not zero.

\begin{proposisi}[Algebraic properties and composition]
\index{continuous}
If $f$ and $g$ are continuous at $c$, then $f+g$, $fg$, and $cf$ are continuous at $c$;
if $g(c)\neq 0$, then $f/g$ is also continuous at $c$. If $f$ is continuous at $c$ and
$g$ is continuous at $f(c)$, then $g\circ f$ is continuous at $c$.
\end{proposisi}

\begin{proof}
The algebraic part follows the same limit rules as for sequences, with
$x_{n}\to c$ replaced by $x\to c$; since $c$ is a limit point, the sequence criterion
carries every statement about a limit of a function over to a statement about a limit of
a sequence, and the algebraic rules for sequences were proved in Chapter 4. For
composition, take $\varepsilon>0$. Since $g$ is continuous at $f(c)$, there is
$\eta>0$ such that $|y-f(c)|<\eta$ implies $|g(y)-g(f(c))|<\varepsilon$. Since $f$ is
continuous at $c$, for that $\eta$ there is $\delta>0$ such that $|x-c|<\delta$ implies
$|f(x)-f(c)|<\eta$. That chain of implications gives $|g(f(x))-g(f(c))|<\varepsilon$
whenever $|x-c|<\delta$.
\end{proof}

The pattern of the composition proof is worth remembering: two $\varepsilon$'s are used
in succession, and $\eta$ plays the role of $\varepsilon$ for the inner function. As a
consequence, polynomials and rational functions are continuous on their whole domain.
What is not guaranteed is a function defined by cases, and from there counterexamples
appear at the end of the chapter.

The crudest form of failure is a jump. The function in
Figure~\ref{fig:lompat} equals $x$ to the left of $1$, equals $2$ at the point $1$,
and equals $x+1$ to its right. At $x=1$ its value is $2$, whereas the values from the
left approach $1$. That difference of $1$ cannot be closed by any $\delta$, however small,
because for every $\delta$ there is still a point on the left with $|f(x)-f(1)|$ close
to $1$.

\begin{figure}[htbp]
\centering
\begin{tikzpicture}
\begin{axis}[width=0.70\linewidth,height=5.0cm,
  axis lines=left,xmin=0,xmax=2.2,ymin=0,ymax=3.2,
  xlabel={$x$},ylabel={$f(x)$},xtick={1},xticklabels={$1$},
  ytick={1,2,3},tick label style={font=\scriptsize},
  label style={font=\scriptsize}]
\fill[layerplum!12,draw=layerplum!60] (0,1.5) rectangle (2.2,2.5);
\addplot[densely dashed,layerplum] coordinates {(0.9,0) (0.9,3.2)};
\addplot[densely dashed,layerplum] coordinates {(1.1,0) (1.1,3.2)};
\addplot[thick,ink,domain=0:1,samples=2] {x};
\addplot[thick,ink,domain=1:2.2,samples=2] {x+1};
\addplot[only marks,mark=*,mark size=1.6pt,accent] coordinates {(0.95,0.95)};
\addplot[only marks,mark=o,mark size=2.6pt,ink,mark options={fill=white,solid}]
  coordinates {(1,1)};
\addplot[only marks,mark=*,mark size=1.7pt,ink] coordinates {(1,2)};
\draw[accent,->,>=stealth,thin] (0.60,1.30) -- (0.93,0.99);
\node[font=\scriptsize,layerplum,anchor=west] at (1.13,2.75) {$\delta$};
\node[font=\scriptsize,accent,anchor=west] at (0.10,1.22) {$f(0.95)=0.95$};
\node[font=\scriptsize,ink,anchor=north east] at (1.02,1.88) {$f(1)=2$};
\end{axis}
\end{tikzpicture}
\caption{A jump at $x=1$: $f(x)=x$ for $x<1$, $f(1)=2$, $f(x)=x+1$ for $x>1$.
The band $\varepsilon=\tfrac12$ around $f(1)=2$ is the horizontal strip between $1.5$ and
$2.5$, while the two dashed lines mark $|x-1|<0.1$. The point $x=0.95$ is already
inside that $\delta$ band, but its value $0.95$ is still $0.55$ below the lower edge of
the $\varepsilon$ band. The empty circle at $(1,1)$ is the value approached from the left,
not the value of the function. Shrinking $\delta$ does not help: a point on the left
closer to $1$ gives a value closer to $1$, not to $2$.}
\label{fig:lompat}
\end{figure}

\section{The intermediate value theorem}
\tingkat{intermediate}

In everyday language: to move from the lower room to the upper room, a climber has to
pass through every height in between. The graph of a continuous function cannot jump over
a value.

\begin{teorema}[Intermediate value]
\index{intermediate value theorem}
Let $f$ be continuous on $[a,b]$ and $f(a)<0<f(b)$. Then there is $c\in(a,b)$ with
$f(c)=0$. More generally, for every $v$ between $f(a)$ and $f(b)$ there is $c\in[a,b]$
with $f(c)=v$.
\end{teorema}

\begin{proof}
It is enough to prove the case $f(a)<0<f(b)$, because the general case is obtained by
applying that theorem to $f-v$. Form $S=\{x\in[a,b] : f(x)<0\}$. That set is not empty
because $a\in S$ and it is bounded above by $b$, so by the supremum property of Chapter 3,
$c=\sup S$ exists and $a\le c\le b$.

First, $f(c)\le 0$. Since $c$ is a supremum, there is a sequence $x_{n}\in S$ with
$x_{n}\to c$, and $f(x_{n})<0$ for every $n$. Continuity gives $f(x_{n})\to f(c)$, and a
limit preserves weak inequalities, so $f(c)\le 0$.

Second, $f(c)\ge 0$. Suppose $f(c)<0$. Since $f(b)>0$, the point $c$ is not $b$, so
$c<b$. Continuity at $c$ gives $\delta>0$ such that $|x-c|<\delta$ implies
$|f(x)-f(c)|<-f(c)$, in particular $f(x)<0$. For $\delta<2(b-c)$ the point $c+\delta/2$
lies in $[a,b]$ and satisfies $f(c+\delta/2)<0$, so it is a member of $S$; yet
$c+\delta/2>c=\sup S$. That impossibility gives $f(c)\ge 0$.

The two inequalities together give $f(c)=0$, and since $f(a)<0<f(b)$, that point lies in
$(a,b)$.
\end{proof}

Note that the proof gives no way of finding $c$. It only shows that the set of negative
points has a roof, and that the roof must be a root. The gap between ``exists'' and
``can be computed'' is what makes the intermediate value theorem used as a guarantee of
existence, not as an algorithm.

\begin{untukanak}
A climber starts in a valley and stops at a summit. They may rest, may walk backward for
a while, may stand still for a long time. What they do not do is jump. Therefore, if
there is a river in the valley and no river at the summit, there must be one point on
that path where they are for the first time exactly at the height of the river.

What makes one curious: how many such points are there? There may be one, there may be
three, there may be infinitely many, and the climber's story does not say which. It only
guarantees at least one. Mathematicians use that guarantee for something else: to prove
that an equation certainly has an answer, even though the answer cannot always be written
with a simple formula.

The next question is harder. What if the path is not connected, for example if there is a
chasm in the middle? The guarantee vanishes at once, and that is no prank: there are
functions whose value jumps from negative to positive without ever touching zero.
In this chapter we will know why the property ``does not jump'' is already enough, and
why that property has to be stated first before the theorem is used.
\end{untukanak}

Two immediate consequences are worth recording. First, the image of an interval under a
continuous function is an interval: $f([a,b])$ contains every value between $f(a)$ and
$f(b)$. Second, every continuous function $f\colon[0,1]\to[0,1]$ has a fixed point.
Consider $h(x)=f(x)-x$; at $x=0$ we have $h(0)=f(0)\ge 0$ and at $x=1$ we have
$h(1)=f(1)-1\le 0$. If either value is zero, that point is a fixed point. If not, the
intermediate value theorem gives $c$ with $h(c)=0$, that is $f(c)=c$.

The numbers in a concrete case can be checked. The equation $x^{3}+x-1=0$ has a root in
$(0,1)$ because $f(0)=-1<0<1=f(1)$. After twenty-one bisections, the root is trapped
between $0.682327747$ and $0.682328224$, with interval width $2^{-21}=4.768\times10^{-7}$;
the true value is $0.6823278038280193$. Bisection is not a proof, but the numbers can be
compared with the theoretical bound.

\section{The extreme value theorem}
\tingkat{intermediate}

In everyday language: on a road that is closed and unbroken, there is certainly a highest
point and a lowest point. Before proving it, we need one control: such a function cannot
run off to infinity.

\begin{lemma}
\index{extreme value theorem}
If $f$ is continuous on $[a,b]$, then $f$ is bounded on $[a,b]$.
\end{lemma}

\begin{proof}
Suppose $f$ is unbounded. Then for every $n$ there is $x_{n}\in[a,b]$ with
$|f(x_{n})|>n$. The sequence $(x_{n})$ is bounded, so by the
Bolzano--Weierstrass theorem of Chapter 4 there is a subsequence $x_{n_{k}}\to x$. Since
$a\le x_{n_{k}}\le b$ for all $k$, a limit preserves inequalities and
$x\in[a,b]$. Continuity at $x$ gives $f(x_{n_{k}})\to f(x)$, so that sequence is
bounded. Yet $|f(x_{n_{k}})|>n_{k}\to\infty$. That impossibility closes the
supposition.
\end{proof}

Now the main statement.

\begin{teorema}[Extreme value]
\index{extreme value theorem}
If $f$ is continuous on $[a,b]$, then there are $c,d\in[a,b]$ such that
$f(c)=\sup f([a,b])$ and $f(d)=\inf f([a,b])$. In other words, $f$ attains its
maximum and its minimum.
\end{teorema}

\begin{proof}
By the lemma, the set $f([a,b])$ is bounded and not empty, so $M=\sup f([a,b])$
exists and is finite. For every $n$ there is $x_{n}\in[a,b]$ with $f(x_{n})>M-1/n$;
such a sequence exists because $M$ is the least upper bound. By
Bolzano--Weierstrass there is a subsequence $x_{n_{k}}\to c$, and since
$a\le x_{n_{k}}\le b$ we get $c\in[a,b]$. Continuity at $c$ gives
$f(x_{n_{k}})\to f(c)$, while the trapping above gives $f(x_{n_{k}})\to M$.
Uniqueness of the limit gives $f(c)=M$. For the minimum, apply that result to $-f$.
\end{proof}

Both proofs use the same bridge: a bounded sequence inside a closed interval has a
subsequence that tends to a point inside that interval. Closedness is what ensures the
limit does not escape, and that is what is lost on an open interval.

The counterexample is simple. The function $f(x)=x$ is continuous on $(0,1)$ but attains
neither a maximum nor a minimum, because $0$ and $1$ are not members of its domain. The
function $g(x)=1/x$ is continuous on $(0,1]$ but unbounded there: $g(1/n)=n$. So
closedness and boundedness of the domain are not decoration; they are the entire reason
the theorem works.

As a consequence, for $f$ continuous on $[a,b]$ with $m=\inf f$ and $M=\sup f$, we have
$f([a,b])=[m,M]$. That equality combines two theorems: the first says $m$
and $M$ are attained, the second says all values in between are used.

\section{Uniform continuity}
\tingkat{advanced}

In the definition of continuity, $\delta$ may depend on $\varepsilon$ \emph{and}
on the point $c$. For $f(x)=1/x$ on $(0,1]$, every point has its own $\delta$,
and that $\delta$ shrinks without bound as $c$ approaches zero. Locally
everything is fine, but there is never one $\delta$ for the whole interval, and that is
what the grid in the chapter opening needs.

\kamus{uniform continuity}{kekontinuan seragam}{one step width that works from the foot of
the staircase to its top, not a different width on every step}

\begin{definisi}[Uniform continuity]
\index{uniform continuity}
A function $f\colon A\to\mathbb{R}$ is called \emph{uniformly continuous} on $A$ if for
every $\varepsilon>0$ there exists $\delta>0$ such that for all $x,y\in A$ with
$|x-y|<\delta$ we have $|f(x)-f(y)|<\varepsilon$.
\end{definisi}

Compare the two orders of quantifiers side by side:
\[
\text{continuous: } \ \forall c\in A\ \forall\varepsilon>0\ \exists\delta>0\
\forall x\in A,
\qquad
\text{uniform: } \ \forall\varepsilon>0\ \exists\delta>0\ \forall x,y\in A .
\]
In the uniform version, $\delta$ appears before $x$ and $y$ are introduced, so it may
not exploit knowledge of any point. That is the whole content of the difference, and
because of it uniformity is always relative to the set $A$: the same function can be
uniform on one set and not on another.

\begin{bacakeras}
Read aloud, two voices taking turns.

\textbf{Voice one:} Close at one point. However narrow the window installed, as long as
the input is close enough, the output is inside the window.

\textbf{Voice two:} Close at all points. It is not a question of which door. One window
width is installed first, then demanded to hold in the whole corridor, from the first
door to the last.

\textbf{Voice one:} If the corridor narrows at the end, that width has to be cut again
and again.

\textbf{Voice two:} And the cutting must not stop at one number. Stopping means one width
for the whole corridor; shrinking forever means a promise in installments.

Repeat this line until it feels right on the tongue: \emph{one point close enough, all
points one width.}
\end{bacakeras}

The statement that uniformity is stronger than continuity is proved directly from its
definition: uniformity gives one $\delta$ for all pairs of points, and
continuity only asks for a $\delta$ for pairs one of whose members is a fixed
point.

\begin{proposisi}
\index{uniform continuity}
If $f$ is uniformly continuous on $A$, then $f$ is continuous at every point of $A$.
\end{proposisi}

\begin{proof}
Take $c\in A$ and $\varepsilon>0$. Since $f$ is uniformly continuous, there is $\delta>0$
such that $|x-y|<\delta$ implies $|f(x)-f(y)|<\varepsilon$ for all $x,y\in A$.
Choosing $y=c$, the condition $|x-c|<\delta$ gives $|f(x)-f(c)|<\varepsilon$. That is
exactly the definition of continuity at $c$.
\end{proof}

The second way, used repeatedly in the rest of this chapter, works by contraposition: if
$f$ fails to be continuous at some point, two sequences whose distance tends to zero can
be built so that their outputs stay apart, and that refutes uniformity. The proof for
$1/x$ below uses that pattern in full.

Now an example showing that uniformity depends on the domain, not on the formula of the
function.

\begin{proposisi}
\index{uniform continuity}
$f(x)=x^{2}$ is uniformly continuous on $[0,1]$, with $\delta=\varepsilon/2$.
\end{proposisi}

\begin{proof}
For $x,y\in[0,1]$ we have $|x+y|\le 2$, hence
\[
|x^{2}-y^{2}|=|x-y|\,|x+y|\le 2|x-y| .
\]
If $|x-y|<\varepsilon/2$, then $|x^{2}-y^{2}|<\varepsilon$.
\end{proof}

That uniformity does not depend on the two endpoints: on $(0,1)$ the same proof still
works, because the bound $|x+y|\le 2$ uses no endpoint at all. The number
$\varepsilon/2$ is safe but not the best. For $|x-y|=\delta$ the largest difference
occurs at the end of the interval and equals $\delta(2-\delta)$, so $\delta$ is valid
provided $\delta\le 1-\sqrt{1-\varepsilon}$. For $\varepsilon=10^{-2}$ that bound is
$5.012563\times10^{-3}$, and with $\delta=\varepsilon/2$ the largest difference is
$9.975\times10^{-3}$.

\begin{proposisi}
\index{uniform continuity}
$f(x)=x^{2}$ is not uniformly continuous on $\mathbb{R}$.
\end{proposisi}

\begin{proof}
Take $\varepsilon=1$. Suppose there is $\delta>0$ that works. Choose a natural number
$n>1/\delta$ and set $x_{n}=n$, $y_{n}=n+1/n$. Then $|x_{n}-y_{n}|=1/n<\delta$,
but
\[
|x_{n}^{2}-y_{n}^{2}|=\left|n^{2}-\left(n+\tfrac1n\right)^{2}\right|
=2+\frac{1}{n^{2}}>2>1=\varepsilon .
\]
So no $\delta$ works for $\varepsilon=1$.
\end{proof}

The number $2$ in that computation is the point: the difference of values does not
shrink, it approaches $2$, even though the difference of inputs tends to zero. What fails
is the rate of change of $x^{2}$, which grows without bound, not its continuity. With
$\delta=10^{-4}$ and the pair $x=10^{4}$, $y=10^{4}+5\times10^{-5}$, the difference of
outputs is $1.0000000149$; with $\delta=10^{-2}$ and $x=100$, the difference is
$1.000025$. That number barely moves when $\delta$ shrinks ten thousand
times, and it is precisely that immobility which proves the failure of uniformity.

An equally important example is $1/x$, because here pointwise continuity never fails.

\begin{proposisi}
\index{uniform continuity}
$f(x)=1/x$ is not uniformly continuous on $(0,1]$, but is uniformly continuous on
$[1,\infty)$ with $\delta=\varepsilon$.
\end{proposisi}

\begin{proof}[Proof of the first part, with two sequences]
Take $x_{n}=1/n$ and $y_{n}=1/(n+1)$. Then
\[
|x_{n}-y_{n}|=\frac{1}{n(n+1)}\longrightarrow 0,
\qquad\text{but}\qquad
\left|\frac{1}{x_{n}}-\frac{1}{y_{n}}\right|=|n-(n+1)|=1
\]
for every $n$. If $f$ were uniformly continuous, take $\delta$ for $\varepsilon=1$;
for large $n$ we have $|x_{n}-y_{n}|<\delta$, so the difference of outputs would have to
be less than $1$. That contradicts the result above.
\end{proof}

\begin{proof}[Proof of the second part]
For $x,y\ge 1$ we have $xy\ge 1$, hence
\[
\left|\frac{1}{x}-\frac{1}{y}\right|=\frac{|x-y|}{xy}\le|x-y| .
\]
So $\delta=\varepsilon$ works, and the uniformity comes from one fact only: the
product $xy$ is never smaller than $1$.
\end{proof}

The pattern is clear. On $[1,\infty)$ the denominator $xy$ is kept away from zero by the
lower bound of the domain; on $(0,1]$ that denominator can be made as small as we like,
and that is where uniformity is lost. Uniformity is a property of the pair of function and
set, not of the function alone: $1/x$ appears in both propositions with two different
answers.

The largest threshold $\delta(c)$ at each point shows the difference as a picture.
For $c\in(0,1]$ and $\varepsilon>0$, the condition $1/(c-\delta)-1/c\le\varepsilon$ is
equivalent to $\delta\le\varepsilon c^{2}/(1+\varepsilon c)$, so the largest threshold at $c$
is $\delta(c)=\varepsilon c^{2}/(1+\varepsilon c)$, and
Figure~\ref{fig:delta-mengecil} plots it together with one fixed value
$\delta=\varepsilon/2$ that works for $x^{2}$ on $[0,1]$. The first curve falls
along a straight line with slope two on the log-log scale, from $9.901\times10^{-3}$ at $c=1$
to $1.000\times10^{-10}$ at $c=10^{-4}$; the second line is horizontal. Uniformity on
$[0,1]$ therefore reads as a horizontal line, and its failure on $(0,1]$
as a curve that never touches the floor.

\begin{figure}[htbp]
\centering
\begin{tikzpicture}
\begin{axis}[width=0.80\linewidth,height=5.8cm,
  xmode=log,ymode=log,xmin=8e-5,xmax=1.4,ymin=8e-11,ymax=4e-2,
  xlabel={point $c$},ylabel={threshold $\delta(c)$},
  grid=major,grid style={black!10},
  tick label style={font=\scriptsize},label style={font=\scriptsize},
  legend style={font=\scriptsize,at={(0.97,0.03)},anchor=south east,
    draw=black!15,fill=white,fill opacity=0.9,text opacity=1,row sep=1pt}]
\addplot[thick,accent,domain=1e-4:1,samples=80] {1e-2*x^2/(1+1e-2*x)};
\addlegendentry{$1/x$ on $(0,1]$, $\varepsilon=10^{-2}$}
\addplot[thick,layerteal,densely dashed,domain=1e-4:1] {5e-3};
\addlegendentry{$x^{2}$ on $[0,1]$, $\delta=\varepsilon/2$}
\addplot[only marks,mark=*,mark size=1.6pt,accent]
  coordinates {(1,9.901e-3) (1e-2,9.999e-7) (1e-4,1.0e-10)};
\node[font=\scriptsize,accent,anchor=south west] at (1.15e-4,1.7e-10)
  {$1.000\times10^{-10}$};
\node[font=\scriptsize,accent,anchor=east] at (0.92,1.05e-2) {$9.901\times10^{-3}$};
\end{axis}
\end{tikzpicture}
\caption{The largest threshold $\delta$ at each point $c$ for $1/x$ on $(0,1]$ with
$\varepsilon=10^{-2}$, namely $\delta(c)=\varepsilon c^{2}/(1+\varepsilon c)$, on a
log-log scale. The value falls from $9.901\times10^{-3}$ at $c=1$ to
$9.999\times10^{-7}$ at $c=10^{-2}$ and $1.000\times10^{-10}$ at $c=10^{-4}$,
so the graph is a straight line with slope two. The dashed line
$\delta=5\times10^{-3}$ is the safe choice $\varepsilon/2$ from the proof for $x^{2}$
on $[0,1]$: one and the same number works on that whole interval.}
\label{fig:delta-mengecil}
\end{figure}

\begin{tangan}
Materials: a corridor two meters long, chalk or paper tape, and a ruler.

Mark the floor at the point $0$, then walk with a stride that is continually halved:
half a meter, a quarter of a meter, an eighth of a meter, and so on. Mark every
footfall. After six or seven steps, stop and look at the marks from the end of the
corridor.

What needs to be seen: the marks cluster near the end, and the remaining distance shrinks
without ever becoming zero. That rule never finishes the journey, although every step is
valid according to its own rule.

Repeat with another rule: a fixed stride of half a meter all the way to the end. The
journey is finished in four steps, and one rule holds from the first mark to the last.
That is the difference between a rule that shrinks forever and one rule for the whole
corridor.
\end{tangan}

\section{Compactness, Lipschitz, and the Heine--Cantor theorem}
\tingkat{advanced}

In everyday language: on a domain that is closed and does not stretch to infinity,
ordinary continuity is already enough to give one threshold for all points. To state it,
we need one term about the domain.

\kamus{compact}{kompak}{a box that is closed tight and does not stick out without end}

\begin{definisi}[Compact]
\index{compact}
A set $K\subseteq\mathbb{R}$ is called \emph{compact} if every family of open sets whose
union contains $K$ has a finite subfamily whose union also contains $K$.
\end{definisi}

That definition is not practical to use directly. What saves us is its equivalence with
two properties that are easy to check: for sets in $\mathbb{R}$, compact is equivalent to
closed and bounded. That equivalence is known as the Heine--Borel theorem; we skip its
proof, and it is in Rudin, \emph{Principles of Mathematical Analysis}, Chapter 2, section
Compact Sets.

\begin{teorema}[Heine--Cantor]
\index{heine-cantor}
If $f$ is continuous on a compact set $K$, then $f$ is uniformly continuous on $K$.
\end{teorema}

\begin{proof}
Suppose $f$ is not uniformly continuous. Then there is $\varepsilon>0$ such that for every $n$
there are $x_{n},y_{n}\in K$ with $|x_{n}-y_{n}|<1/n$ but
$|f(x_{n})-f(y_{n})|\ge\varepsilon$. Since $K$ is bounded, $(x_{n})$ is bounded, and
by Bolzano--Weierstrass there is a subsequence $x_{n_{k}}\to x$. Since $K$ is closed,
$x\in K$. From $|x_{n_{k}}-y_{n_{k}}|<1/n_{k}$ we get $y_{n_{k}}\to x$ as well, and
$y_{n_{k}}\in K$. Continuity of $f$ at $x$ gives $f(x_{n_{k}})\to f(x)$ and
$f(y_{n_{k}})\to f(x)$, so $|f(x_{n_{k}})-f(y_{n_{k}})|\to 0$. That contradicts
$|f(x_{n_{k}})-f(y_{n_{k}})|\ge\varepsilon$. So the supposition is false.
\end{proof}

That proof uses the two properties of compactness separately: boundedness to trigger
Bolzano--Weierstrass, and closedness to ensure the limit point is still inside
$K$. Both are needed. On $(0,1)$, the function $1/x$ is continuous but not uniform, and
its domain is not closed. On $\mathbb{R}$, the function $x^{2}$ is continuous but not
uniform, and its domain is not bounded.

The following class of functions gives a shortcut that is often faster than checking the
definition.

\begin{definisi}[Lipschitz]
\index{lipschitz}
A function $f\colon A\to\mathbb{R}$ is called \emph{Lipschitz} if there is a constant $L\ge 0$
such that $|f(x)-f(y)|\le L|x-y|$ for all $x,y\in A$.
\end{definisi}

\begin{proposisi}
\index{lipschitz}
If $f$ is Lipschitz on $A$, then $f$ is uniformly continuous on $A$.
\end{proposisi}

\begin{proof}
Take $\varepsilon>0$. If $L=0$, the function is constant and any choice of $\delta$ works.
If $L>0$, take $\delta=\varepsilon/L$. For $|x-y|<\delta$ we have
$|f(x)-f(y)|\le L|x-y|<\varepsilon$.
\end{proof}

The converse does not hold, and the example shows that uniformity is weaker than
Lipschitz. The function $f(x)=\sqrt{x}$ is continuous on $[0,1]$, hence uniformly
continuous by the Heine--Cantor theorem. But
\[
\frac{|\sqrt{x}-\sqrt{0}|}{|x-0|}=\frac{1}{\sqrt{x}}\longrightarrow\infty
\quad (x\to 0^{+}),
\]
so no $L$ satisfies the Lipschitz definition. Thus Lipschitz is a sufficient condition,
not a necessary one. Conversely $x^{2}$ on $[0,1]$ is Lipschitz with $L=2$
and $1/x$ on $[1,\infty)$ with $L=1$, exactly the numbers in the preceding proofs.

\begin{tantangan}
The two examples of failure of uniformity in this chapter, $x^{2}$ on $\mathbb{R}$ and
$1/x$ on $(0,1]$, are both unbounded on their domains. Find a function that is
\emph{bounded} on $(0,1)$, continuous at every point, but not uniformly continuous
there, then prove it.

Hint 1: failure does not require values that grow; values that oscillate are enough,
with an oscillation that gets denser as $x$ approaches the end of the interval.

Hint 2: try $f(x)=\sin(1/x)$ with $x_{n}=1/(2\pi n+\pi/2)$ and
$y_{n}=1/(2\pi n-\pi/2)$. The function values are $1$ and $-1$, while
$|x_{n}-y_{n}|=\pi/((2\pi n)^{2}-(\pi/2)^{2})$ tends to zero.

If you are stuck, go back to the section Uniform continuity, the proof for $1/x$ with two
sequences.
\end{tantangan}

\section{A function that is continuous at one point only}
\tingkat{advanced}

Continuity is not built into all functions; it can fail on a very large set of points. The
cheapest example is built from the density of the rationals in Chapter 2.

\begin{contoh}
\index{dirichlet function}
Define $f\colon\mathbb{R}\to\mathbb{R}$ by $f(x)=x$ if $x$ is rational and
$f(x)=0$ if $x$ is irrational. Then $f$ is continuous at $0$ and not continuous at any
other point.
\end{contoh}

\begin{proof}
At $x=0$ we have $|f(x)-f(0)|=|f(x)|\le|x|$, because $f(x)$ always equals either $x$ or
$0$. So $\delta=\varepsilon$ works and $f$ is continuous at $0$.

Now take $c\neq 0$. Since $\mathbb{Q}$ and $\mathbb{R}\setminus\mathbb{Q}$
are both dense, there is a rational sequence $r_{n}\to c$ and an irrational sequence
$s_{n}\to c$. Then $f(r_{n})=r_{n}\to c$ while $f(s_{n})=0\to 0$. Two sequence limits
differ, so by the sequence criterion $\lim_{x\to c}f(x)$ does not exist, and $f$ is not
continuous at $c$.
\end{proof}

A close relative of that example is the indicator of the rationals, which equals $1$ at
rationals and $0$ at irrationals; that function is continuous at no point at all. The next
example is more interesting, because its set of continuity points is in fact large; it is
known as a modification of the Dirichlet function.

\begin{contoh}[Thomae's function, a Dirichlet modification]
\index{thomae function}
Define $t\colon[0,1]\to\mathbb{R}$ by $t(x)=0$ if $x$ is irrational, and
$t(x)=1/q$ if $x=p/q$ in lowest terms with $q\ge 1$. Then $t$ is
continuous at every irrational point and not continuous at every rational point.
\end{contoh}

\begin{proof}
Take $a$ irrational and $\varepsilon>0$. Choose $N$ with $1/N<\varepsilon$. The rational
numbers in $[0,1]$ with denominator $q\le N$ are finite in number, so among them there is
one closest to $a$; call its distance $\eta>0$. Since $a$ is irrational, that distance is
not zero. For $|x-a|<\eta$ there are two possibilities. If $x$ is irrational, then
$|t(x)-t(a)|=0<\varepsilon$. If $x=p/q$ with $q>N$, then
$|t(x)-t(a)|=1/q<1/N<\varepsilon$. So $t$ is continuous at $a$.

Now take $a=p/q$ rational. An irrational sequence $s_{n}\to a$ gives
$t(s_{n})=0\to 0$, while $t(a)=1/q>0$. The limit along that sequence is not equal to the
value of the function, so $t$ is not continuous at $a$.
\end{proof}

Both examples show that the set of discontinuity points can be the whole line, the set of
rationals, or a single point. The full classification, which states that this set is
always a countable union of closed sets, is in Rudin, Chapter 4, section
Discontinuities. Continuity that fails at many points is not an oddity but a structured
symptom.

\section{A worked example: a delta that depends on the point}
\tingkat{intermediate}

We return to $f(x)=x^{2}$ at $x_{0}=2$ with $\varepsilon=10^{-3}$, this time
down to the last digit.

\begin{resep}
Three steps for writing an $\varepsilon$-$\delta$ proof of $|f(x)-f(c)|<\varepsilon$:
\begin{enumerate}
  \item Work backward. Write $|f(x)-f(c)|$ as $|x-c|$ times a remainder, then find an
        upper bound for that remainder on the domain in use.
  \item Bound $|x-c|<1$ first (or $<r$) so that the remainder is bounded. Call
        its upper bound $M$, then set $\delta=\min(1,\varepsilon/M)$.
  \item Write the proof forward. Take $\varepsilon>0$, choose that $\delta$, suppose
        $|x-c|<\delta$, then derive $|f(x)-f(c)|<\varepsilon$ line by line.
\end{enumerate}
Steps 1 and 3 must not be swapped. The sentences of the proof are read from $\varepsilon$
to $\delta$, but $\delta$ is found in the opposite direction; writing directly from
$\varepsilon$ usually produces a $\delta$ that still contains $x$, and a choice like
that is not a valid $\delta$.
\end{resep}

The condition that must be satisfied is
\[
|x^{2}-4|=|x-2|\,|x+2|<\varepsilon
\]
for all $x$ with $|x-2|<\delta$. If $\delta<1$, then $|x+2|<4+\delta$,
so it is enough to demand $\delta(4+\delta)=\varepsilon$. That quadratic equation
has the positive root
\[
\delta=\sqrt{4+\varepsilon}-2 .
\]
For $\varepsilon=10^{-3}$ that formula gives $\delta=2.49984376953\times10^{-4}$.
A check at three points:
\[
\begin{array}{lll}
|x-2|=2.497343925761\times10^{-4} & \Rightarrow & |x^{2}-4|=9.989999375719\times10^{-4},\\
|x-2|=2.499843769530\times10^{-4} & \Rightarrow & |x^{2}-4|=1.000000000000\times10^{-3},\\
|x-2|=2.502343613300\times10^{-4} & \Rightarrow & |x^{2}-4|=1.001000062555\times10^{-3}.
\end{array}
\]
The second row shows that the largest $\delta$ produces a difference of exactly
$\varepsilon$; the third row, that going slightly past it already violates the condition.
The popular choice $\delta=\varepsilon/4=2.5\times10^{-4}$ turns out to \emph{fail}: the
difference is $1.0000625\times10^{-3}$, larger than $\varepsilon$ by $6.25\times10^{-8}$.
A failure that small is easily missed when checked with three-digit truncation.

Now notice how $\delta$ depends on the point. For
$\varepsilon=10^{-2}$ and $f(x)=x^{2}$, the positive root of $\delta(2x_{0}+\delta)=\varepsilon$
is $\delta=\sqrt{x_{0}^{2}+\varepsilon}-x_{0}$.

\begin{center}
\small
\begin{tabular}{rrrr}
\hline
$x_{0}$ & largest $\delta$ & $\delta/x_{0}$ & $\varepsilon/(2x_{0})$ \\
\hline
$1$ & $4.987562\times10^{-3}$ & $4.987562\times10^{-3}$ & $5.000000\times10^{-3}$ \\
$2$ & $2.498439\times10^{-3}$ & $1.249220\times10^{-3}$ & $2.500000\times10^{-3}$ \\
$10$ & $4.999875\times10^{-4}$ & $4.999875\times10^{-5}$ & $5.000000\times10^{-4}$ \\
$100$ & $4.999999\times10^{-5}$ & $4.999999\times10^{-7}$ & $5.000000\times10^{-5}$ \\
$1000$ & $5.000000\times10^{-6}$ & $5.000000\times10^{-9}$ & $5.000000\times10^{-6}$ \\
$10^{4}$ & $5.000002\times10^{-7}$ & $5.000002\times10^{-11}$ & $5.000000\times10^{-7}$ \\
\hline
\end{tabular}
\end{center}

The third column shrinks as fast as $1/x_{0}^{2}$, and for large $x_{0}$ the value of
$\delta$ is practically the same as $\varepsilon/(2x_{0})$. So there is no positive lower
bound for $\delta$ as $x_{0}$ grows: the rate of change of $x^{2}$ at $x_{0}$ is $2x_{0}$,
and that rate is what sets the threshold, not rounding.

The failure of uniformity of $x^{2}$ on $\mathbb{R}$ was proved in the section Uniform
continuity and is repeated by machine in the Computational lab, part 3, so the numbers are
not repeated here. What has not yet appeared as a picture is the dependence on the domain,
and that is shown by $1/x$. On $[1,\infty)$ with $\delta=\varepsilon$ the largest
difference occurs at $x=1$ and equals $\varepsilon/(1+\varepsilon)$, namely
$9.900990\times10^{-3}$ for $\varepsilon=10^{-2}$. On $(0,1]$ the pair $x=1/n$,
$y=1/(n+1)$ gives an output difference of exactly $1$ while the input distance shrinks:

\begin{center}
\small
\begin{tabular}{rrr}
\hline
$n$ & $|x-y|$ & $|1/x-1/y|$ \\
\hline
$10$ & $9.091\times10^{-3}$ & $1$ \\
$100$ & $9.901\times10^{-5}$ & $1$ \\
$1000$ & $9.990\times10^{-7}$ & $1$ \\
$10^{4}$ & $9.999\times10^{-9}$ & $1$ \\
\hline
\end{tabular}
\end{center}

\kertas{For $f(x)=x^{2}$ and $\varepsilon=10^{-2}$, compute the largest $\delta$ at
$x_{0}=1000$ from the formula $\sqrt{x_{0}^{2}+\varepsilon}-x_{0}$, then compare with
$\varepsilon/(2x_{0})$. Check that the difference at $x=x_{0}+\delta$ is exactly
$\varepsilon$, and explain why that formula cannot be used to prove uniformity on
$\mathbb{R}$.}

\section{Computational lab: measuring uniformity}
\tingkat{advanced}

The eight checks below turn the previous sections into numbers that can be rechecked:
explicit $\delta$ and its bound, its dependence on the point, $x^{2}$ and $1/x$ on two
domains, the largest threshold at each point, and bisection for the intermediate value
theorem.

\begin{lab}
\begin{verbatim}
import math

# 1. optimal delta for f(x)=x^2 at x0=2, eps=1e-3
def delta_sq(x0, eps):
    return math.sqrt(x0*x0 + eps) - x0

eps = 1e-3
d = delta_sq(2.0, eps)
print("== 1. explicit delta, x0=2, eps=1e-3 ==")
print("optimal delta = %.12f" % d)
for t in (d*0.999, d, d*1.001):
    print("  x-2=%.12e  |x^2-4|=%.12e" % (t, abs((2+t)**2-4)))
print("  delta=eps/4 -> |x^2-4|=%.12e" % abs((2+eps/4)**2-4))

# 2. delta depending on the point
print("== 2. delta(x0) for eps=1e-2 ==")
print("      x0       delta     delta/x0")
for x0 in (1.0, 2.0, 10.0, 100.0, 1000.0, 1e4):
    d = delta_sq(x0, 1e-2)
    print("%8.0f  %.6e  %.6e" % (x0, d, d/x0))

# 3. failure of uniformity of x^2 on R
print("== 3. x^2 on R, eps=1 ==")
print("   delta   x=1/delta   |x^2-(x+d/2)^2|")
for delta in (1e-1, 1e-2, 1e-3, 1e-4):
    x = 1.0/delta
    print("%8.0e %9.0f   %.10f" % (delta, x, abs(x*x-(x+delta/2)**2)))

# 4. x^2 on [0,1]
print("== 4. x^2 on [0,1], |x-y|=delta ==")
print("   delta    2d-d^2       grid")
for delta in (1e-1, 1e-2, 1e-3):
    n = int(round(1/delta))
    g = max(abs((i*delta/n)**2-(i*delta/n+delta)**2)
            for i in range(0, int(n*(1-delta)/delta)+1))
    print("%8.0e  %.10e %.10e" % (delta, 2*delta-delta*delta, g))

# 5. 1/x on [1,inf) and on (0,1]
print("== 5. 1/x on [1,inf), delta=eps ==")
for eps in (1e-1, 1e-2, 1e-3):
    s = max(abs(1/x-1/(x+eps)) for x in [1+i*eps/500 for i in range(501)])
    print("eps=%.0e  sup=%.10e  eps/(1+eps)=%.10e" % (eps, s, eps/(1+eps)))
print("== 6. 1/x on (0,1]: x=1/n, y=1/(n+1) ==")
print("      n      |x-y|      |1/x-1/y|")
for n in (10, 100, 1000, 10000):
    print("%7d   %.3e       %.1f" % (n, 1/n-1/(n+1), 1.0))

# 7. largest valid delta at c, f=1/x, eps=1e-2
print("== 7. delta(c) for 1/x, eps=1e-2 ==")
print("       c      delta(c)    delta/c")
for c in (1.0, 1e-1, 1e-2, 1e-3, 1e-4):
    dd = 1e-2*c*c/(1+1e-2*c)
    print("%8.0e  %.4e  %.4e" % (c, dd, dd/c))

# 8. numerical IVT for x^3+x-1=0
print("== 8. root of x^3+x-1=0 by bisection ==")
lo, hi = 0.0, 1.0
for k in range(1, 22):
    mid = (lo+hi)/2
    if mid**3+mid-1 < 0:
        lo = mid
    else:
        hi = mid
    if k in (1, 5, 10, 15, 20, 21):
        print("k=%2d  [%.12f, %.12f]  width=%.3e"
              % (k, lo, hi, hi-lo))
print("f(lo)=%.3e  f(hi)=%.3e" % (lo**3+lo-1, hi**3+hi-1))
\end{verbatim}
\end{lab}

\newpage
The output, as it stands:

\begin{verbatim}
== 1. explicit delta, x0=2, eps=1e-3 ==
optimal delta = 0.000249984377
  x-2=2.497343925761e-04  |x^2-4|=9.989999375719e-04
  x-2=2.499843769530e-04  |x^2-4|=1.000000000000e-03
  x-2=2.502343613300e-04  |x^2-4|=1.001000062555e-03
  delta=eps/4 -> |x^2-4|=1.000062499999e-03
== 2. delta(x0) for eps=1e-2 ==
      x0       delta     delta/x0
       1  4.987562e-03  4.987562e-03
       2  2.498439e-03  1.249220e-03
      10  4.999875e-04  4.999875e-05
     100  4.999999e-05  4.999999e-07
    1000  5.000000e-06  5.000000e-09
   10000  5.000002e-07  5.000002e-11
== 3. x^2 on R, eps=1 ==
   delta   x=1/delta   |x^2-(x+d/2)^2|
   1e-01        10   1.0025000000
   1e-02       100   1.0000250000
   1e-03      1000   1.0000002499
   1e-04     10000   1.0000000149
== 4. x^2 on [0,1], |x-y|=delta ==
   delta    2d-d^2       grid
   1e-01  1.9000000000e-01 1.9000000000e-01
   1e-02  1.9900000000e-02 1.9900000000e-02
   1e-03  1.9990000000e-03 1.9990000000e-03
== 5. 1/x on [1,inf), delta=eps ==
eps=1e-01  sup=9.0909090909e-02  eps/(1+eps)=9.0909090909e-02
eps=1e-02  sup=9.9009900990e-03  eps/(1+eps)=9.9009900990e-03
eps=1e-03  sup=9.9900099900e-04  eps/(1+eps)=9.9900099900e-04
== 6. 1/x on (0,1]: x=1/n, y=1/(n+1) ==
      n      |x-y|      |1/x-1/y|
     10   9.091e-03       1.0
    100   9.901e-05       1.0
   1000   9.990e-07       1.0
  10000   9.999e-09       1.0
== 7. delta(c) for 1/x, eps=1e-2 ==
       c      delta(c)    delta/c
   1e+00  9.9010e-03  9.9010e-03
   1e-01  9.9900e-05  9.9900e-04
   1e-02  9.9990e-07  9.9990e-05
   1e-03  9.9999e-09  9.9999e-06
   1e-04  1.0000e-10  1.0000e-06
== 8. root of x^3+x-1=0 by bisection ==
k= 1  [0.500000000000, 1.000000000000]  width=5.000e-01
k= 5  [0.656250000000, 0.687500000000]  width=3.125e-02
k=10  [0.681640625000, 0.682617187500]  width=9.766e-04
k=15  [0.682312011719, 0.682342529297]  width=3.052e-05
k=20  [0.682327270508, 0.682328224182]  width=9.537e-07
k=21  [0.682327747345, 0.682328224182]  width=4.768e-07
f(lo)=-1.354e-07  f(hi)=1.007e-06
\end{verbatim}

There are four things to read from that output, and all of them were predicted before the
program ran. Part 1 confirms that the optimal $\delta$ produces a difference of exactly
$10^{-3}$, whereas $\delta=\varepsilon/4$ produces $1.0000625\times10^{-3}$:
that popular choice does fail, not merely look untidy. Part 2 confirms that
the column $\delta/x_{0}$ falls as fast as $1/x_{0}^{2}$ up to $x_{0}=1000$, then stops
following the pattern at $x_{0}=10^{4}$. At the last point the true value is
$4.99999999987\times10^{-7}$, while float64 prints $5.000002\times10^{-7}$; that is
cancellation error in $\sqrt{x_{0}^{2}+\varepsilon}-x_{0}$, not a property of the function, and
that is why such a formula is not used at large scale.

Parts 3 through 7 show two faces of the same function. On $\mathbb{R}$ the difference
for $x^{2}$ stays above $1$ for every $\delta$ tried; on $[0,1]$ the largest
difference is $2\delta-\delta^{2}$, and the grid prints the same number to ten
digits because the largest difference occurs at $x=1-\delta$, a point that is in the grid
list. On $[1,\infty)$ the difference for $1/x$ is exactly $\varepsilon/(1+\varepsilon)$,
whereas on $(0,1]$ it stays $1.0$ while the input distance shrinks to
$9.999\times10^{-9}$. Part 8 closes with interval width $2^{-21}$ after twenty-one
bisections, and the root of $x^{3}+x-1$ is trapped inside it.

The program does not prove a single one of those things. It checks claims that have already
been proved, and in two places it shows the limits of itself: the failure of
$\delta=\varepsilon/4$ only becomes visible at the eighth digit, and cancellation error
only becomes visible at $x_{0}=10^{4}$.

\begin{untukdosen}
The teaching order I suggest: start from the grid in the chapter opening, not from the
definition. Write $10000$ and $5000$ on the board, ask what grid spacing is needed at
that point, then compare with the answer at $x=1$; from that gap the
$\varepsilon$-$\delta$ definition appears as the answer to a question that has already
been felt, not as a sentence to memorize. After that continuity, the algebraic properties,
the intermediate value theorem, the extreme value theorem, then uniform continuity and
Heine--Cantor.

Two misconceptions appear most often. First, students believe $\delta$ may depend on $x$
inside the definition of continuity, then ``prove'' uniformity by choosing a different
$\delta$ at every point; the symptom shows when they cannot name one $\delta$ for two
distant points. Second, they think the intermediate value theorem gives a way of
computing a root, whereas it only guarantees existence.

A blackboard activity that works: ask two students to play challenger and responder for
$f(x)=x^{2}$ at $x_{0}=2$ and at $x_{0}=1000$ with the same $\varepsilon$,
then have the class guess at which point the responder runs out of numbers. The sign that
the class is ready to move on: they can name one explicit $\delta$ for $1/x$
on $[1,\infty)$ and explain why the same $\delta$ fails on $(0,1]$.
\end{untukdosen}

\section{Summary}
\tingkat{basic}

\begin{ringkasan}
\begin{itemize}
  \item $\lim_{x\to c}f(x)=L$ means that for every $\varepsilon>0$ there is $\delta>0$
        such that $0<|x-c|<\delta$ gives $|f(x)-L|<\varepsilon$. The value at $c$
        itself is not counted, and the limit is unique.
  \item $f$ is continuous at $c$ if the limit exists and equals $f(c)$. Continuity
        is local, closed under algebraic operations, and preserved under
        composition.
  \item On a closed bounded interval, a continuous function is bounded, attains its
        maximum and minimum, and takes all values in between. All three can
        fail on an open interval.
  \item $f$ is uniformly continuous if one $\delta$ works for all pairs of points.
        Uniformity is stronger than continuity, and uniformity is always relative
        to its domain.
  \item $x^{2}$ is uniform on $[0,1]$ but not on $\mathbb{R}$; $1/x$ is uniform
        on $[1,\infty)$ but not on $(0,1]$. A function continuous on a compact
        set is always uniform, and a Lipschitz function is always uniform.
  \item Continuity can fail on a large set of points: the rational indicator fails
        everywhere, $x$ at rationals and $0$ at irrationals fails at every point
        except zero.
\end{itemize}
\end{ringkasan}

\section{Exercises}
\tingkat{intermediate}

\begin{latihan}
\begin{enumerate}
  \item Prove from the definition that $\lim_{x\to 3}(2x-1)=5$, and give the smallest
        explicit $\delta$ as a function of $\varepsilon$. Explain why
        $\delta=\varepsilon/2$ is the largest valid choice, whereas
        $\delta=\varepsilon/3$ is valid but smaller than necessary.
  \item Prove that $f(x)=x^{2}$ is continuous at every point $c\in\mathbb{R}$ with
        $\delta=\min(1,\varepsilon/(2|c|+1))$, then check that formula at $c=2$ and
        $c=100$ for $\varepsilon=10^{-2}$.
  \item Assess this claim: ``Because $1/x$ is continuous at every point of $(0,1]$, that
        function is uniformly continuous on $(0,1]$.'' True, false, or not enough
        evidence? Name the part that is correct and the part that is wrong.
  \item Prove that if $f$ and $g$ are uniformly continuous on $A$, then $f+g$ is uniformly
        continuous on $A$. Show by example that $fg$ need not be uniformly
        continuous on an unbounded $A$.
  \item Prove that $f(x)=\sqrt{x}$ is uniformly continuous on $[0,\infty)$ by
        using $|\sqrt{x}-\sqrt{y}|\le\sqrt{|x-y|}$. Explain why the proof does not
        require the domain to be closed.
  \item Compute the largest $\delta$ for $f(x)=x^{3}$ at $x_{0}=1$ with
        $\varepsilon=10^{-2}$ by solving $(1+\delta)^{3}-1=\varepsilon$, then
        compare with $\delta=\varepsilon/3$. Check both values in Python.
  \item For $f(x)=x^{2}$ on $[0,1]$ and $\varepsilon=10^{-2}$, compute the largest
        difference for $\delta=\varepsilon/2$ and for $\delta=1-\sqrt{1-\varepsilon}$.
        Show that the first choice is safe but not optimal, and that the second
        choice produces a difference of exactly $\varepsilon$.
  \item Prove that if $f$ is continuous on $[a,b]$ and $f(x)\neq 0$ for all
        $x\in[a,b]$, then there is $m>0$ with $|f(x)|\ge m$ for all $x$. Use the
        extreme value theorem, and explain why the conclusion fails for
        $f(x)=x$ on $(0,1)$.
  \item Define $f(x)=x$ if $x$ is rational and $f(x)=-x$ if $x$ is irrational.
        Prove that $f$ is continuous only at $0$, then show that $|f|$ is continuous at
        every point. Explain why the second result does not contradict the first.
\end{enumerate}

Hints for the three hardest items. Item 4: for the second part, take $f(x)=g(x)=x$
on $\mathbb{R}$ and use the pair $x_{n}=n$, $y_{n}=n+1/n$. Item 5: square
$|\sqrt{x}-\sqrt{y}|\le\sqrt{|x-y|}$ to see where it comes from. Item 9: $|f|$ is none
other than $x\mapsto|x|$, which is continuous everywhere; for $c\neq0$ use rational
and irrational sequences tending to $c$.
\end{latihan}

\begin{untukpeneliti}
What is proved in this chapter: uniformity is stronger than continuity; continuity
on a compact set gives uniformity; continuity on a closed bounded interval
gives boundedness, extreme values, and the intermediate value property. What is popular but
not proved is the impression that the intermediate value theorem gives an algorithm: it
only gives existence, and the rate $2^{-n}$ comes from the procedure, not from the
theorem.

Primary sources: Rudin, \emph{Principles of Mathematical Analysis}, Chapter 4, section
Continuity and Compactness, which proves Heine--Cantor in the framework of general metric
spaces; and Abbott, \emph{Understanding Analysis}, Chapter 4, sections Continuous Functions on
Compact Sets and Sets of Discontinuity. The proof of Heine--Borel is deliberately skipped
here and is in Rudin, Chapter 2, section Compact Sets.

Limit of validity: the whole chapter works in $\mathbb{R}$ with the usual distance. In a
metric space, the definition of uniformity stays the same word for word, but the
Heine--Cantor theorem demands compactness, not merely closedness and boundedness; in an
infinite-dimensional space those two properties separate. In a complete metric space,
uniformity on a dense set is enough to extend a function uniquely, and that is the door
into the extension theorems not discussed here.

A question still open for the reader: how fast the modulus of continuity
$\omega(\delta)=\sup_{|x-y|<\delta}|f(x)-f(y)|$ may fall. For a function continuous
on a closed interval, $\omega(\delta)\to 0$ is guaranteed, but the rate of its decrease is
not: there are continuous functions whose modulus falls more slowly than every power
$\delta^{\alpha}$. The Lipschitz guarantee that is often assumed on adversarial attack
bounds in Chapter 29 is therefore sufficient, not necessary.
\end{untukpeneliti}

\peta{Chapter 11 in Rudin, \emph{Principles of Mathematical Analysis}: Chapter 4 (Continuous
Functions, Continuity and Compactness, Continuity and Connectedness, Discontinuities),
with the Heine--Borel theorem taken from Chapter 2, section Compact Sets. In Abbott,
\emph{Understanding Analysis}: Chapter 4 (Discussion: Examples of Dirichlet and Thomae,
Functional Limits, Combinations of Continuous Functions, Continuous Functions on
Compact Sets, The Intermediate Value Theorem, Sets of Discontinuity). The
$\varepsilon$-$\delta$ definition and the sequence criterion are built from the sequences in Chapter 4 of this
book, whereas the supremum used to prove the intermediate value theorem comes from Chapter 3.
The derivative in Chapter 12 continues this chapter by using the difference quotient as a limit,
and Chapter 27 uses uniformity as one of the guarantees that is often assumed
without being checked.}

\chapter{Derivatives: One and Several Variables}

A position logger on a bicycle writes one point every second. To obtain the velocity, we
divide the difference in distance by the difference in time. Two readings one second apart
and four meters apart give four meters per second, even though that is only an average
velocity. If the separation is one hundredth of a second, the difference in distance is a few
centimeters while the instrument noise is several meters; the quotient jumps to hundreds of
meters per second, even with a negative sign. The same division, two different answers.

The derivative resolves that tension by moving the question to the limit, where noise does not
enter the count. But the limit cannot be measured: what is measurable is only the quotient at
a finite separation, and it carries two errors at once, truncation because the separation is
not zero and rounding because the machine stores only so many digits. The derivative is the
exact mathematical object; the quotient is its approximation, with an error that can be
computed and controlled.

The second difficulty appears when there is more than one input. A credit scoring model
receives twelve inputs, and shifting one input gives one number for the rate of change.
Shifting two inputs at once is not the same as adding those two numbers, unless the function
behaves well; later we shall see an example that fails, and that function is still
continuous. The correct answer is not a list of twelve numbers, but a single linear map on the
shift vector: that is the true derivative, and its matrix is the Jacobian. The continuity of
Chapter 11 does not say how fast something changes.

\begin{untukanak}
Have you ever stood on a curved roof? If you step forward, you climb. If you step sideways,
you can stay at the same height. And yet you are standing at the same point, on the same roof.
So the question ``how steep is this place'' is not complete before we name the direction of
the step.

There is another problem. If you look at only two points on the roof, and then divide the
difference in height by the difference in distance, you get an average slope, not the slope at
a single point. To get the slope exactly under your feet, the two points must be brought
closer and closer. The closer they are the more exact it is, but if they are too close your
measurement numbers fall apart because the measuring instrument is not perfect.

In this chapter we shall learn why the slope at a single point needs a direction, why the
steepest direction always exists, and why a computing tool whose step is too tight is
precisely wrong.
\end{untukanak}

\section{The derivative as a limit and its two meanings}
\tingkat{basic}

The instantaneous rate of change is the number approached by the quotient of two differences
as the difference in the input shrinks. The following definition merely writes that sentence
down carefully, because the quotient itself is undefined when its denominator is zero.

\kamus{derivative}{turunan}{The slope of the line that hugs the graph most tightly at a single
point, which is also the instantaneous rate of change at that point.}

\begin{definisi}[The derivative at a single point]
Suppose $f$ is defined on an open interval containing $c$. The derivative of $f$ at $c$ is
\[
  f'(c)=\lim_{h\to0}\frac{f(c+h)-f(c)}{h},
\]
provided that limit exists and is finite. In this case $f$ is called
\emph{differentiable} at $c$. Other notation used in this book: $Df(c)$ and
$\dfrac{df}{dx}(c)$.\index{derivative}
\end{definisi}

That quotient is defined for every $h\neq0$, so there is no division by zero; what may fail
is the limit. The failure is real: for $f(x)=|x|$ at $c=0$ the quotient is
$1$ for $h>0$ and $-1$ for $h<0$, so the two-sided limit does not exist.

Physically, the derivative is the instantaneous rate of change: $s'(t)$ is velocity when
$s(t)$ is position, $Q'(t)$ is current when $Q(t)$ is the charge that flows, and $C'(q)$ is
marginal cost when $C(q)$ is the cost of producing $q$ units.

Geometrically, the quotient is the slope of the secant from $(c,f(c))$ to
$(c+h,f(c+h))$. As $h\to0$, that secant rotates about the first point toward a single
position; the line in the limiting position is called the tangent line, and its slope is
$f'(c)$. The equation of the tangent line at $c$ is
\[
  y=f(c)+f'(c)(x-c).
\]
The tangent line is not ``the line that touches the graph exactly once'', for with
$f(x)=x^{3}$ at $c=0$ the tangent line $y=0$ crosses the graph at the origin. What is correct:
its difference from the graph shrinks faster than the distance does.

That rotation can be followed numerically. For $f(x)=x^{2}$ at $c=2$, the secant with $h=1$
has slope $5$, with $h=0.5$ the slope is $4.5$, and the tangent line has slope $4$.
Figure~\ref{fig:sekan-singgung} draws all three: from $5$ to $4.5$, settling at $4$.

\begin{figure}[htbp]
\centering
\begin{tikzpicture}
\begin{axis}[
  width=\linewidth, height=6.4cm,
  axis lines=left,
  xlabel={$x$}, ylabel={$y$},
  xmin=1.55, xmax=3.25, ymin=2.2, ymax=10.2,
  xtick={2,3}, ytick={4,9},
  tick label style={font=\footnotesize},
  label style={font=\footnotesize},
  legend style={font=\scriptsize, draw=none, fill=none,
                at={(0.03,0.97)}, anchor=north west},
  clip=false,
]
\addplot[ink, thick, domain=1.6:3.15, samples=90] {x^2};
\addlegendentry{$f(x)=x^{2}$}
\addplot[accent, thick, domain=1.95:3.18] {5*(x-2)+4};
\addlegendentry{secant $h=1$, slope $5$}
\addplot[layerteal, thick, dashed, domain=1.95:3.0] {4.5*(x-2)+4};
\addlegendentry{secant $h=0.5$, slope $4.5$}
\addplot[layerindigo, thick, domain=1.68:3.05] {4*(x-2)+4};
\addlegendentry{tangent line, slope $4$}
\addplot[only marks, mark=*, mark size=1.5pt, ink]
  coordinates {(2,4) (2.5,6.25) (3,9)};
\node[font=\scriptsize, anchor=north west, text=ink] at (axis cs:2.04,3.75) {$(2,4)$};
\end{axis}
\end{tikzpicture}
\caption{Secants rotating into the tangent line of $f(x)=x^{2}$ at $(2,4)$.
Their slopes drop from $5$ (for $h=1$) to $4.5$ (for $h=0.5$) and settle at $4$
as $h\to0$.}
\label{fig:sekan-singgung}
\end{figure}

\begin{bacakeras}
Two voices take turns; the answer gets shorter every time.

--- How fast is it moving? --- From which point to which point?

--- From the first second to the second second. --- Three meters.

--- From the first second to second one and a half. --- Two and a half meters.

--- From the first second to second one and one tenth. --- Two and one fifth.

--- From the first second to second one and one hundredth. --- Almost two.

--- From the first second to the first second. --- I cannot. The denominator is zero.

Stop there. What changes is not the answer, but the distance between the two points being
compared. The distance shrinks, the secant rotates, the number settles: a value that is never
reached yet always approached, that is the derivative.

Repeat together: the distance shrinks, the line rotates, the number settles.
\end{bacakeras}

Let us compute one derivative from the definition. For $f(x)=x^{2}$ at $c=2$,
\[
  \frac{f(2+h)-f(2)}{h}=\frac{(4+4h+h^{2})-4}{h}=\frac{h(4+h)}{h}=4+h .
\]
Since $h\neq0$, that division is legitimate, and the limit is $4$. For $h=0.01$ the quotient
is $4.01$; the difference $0.01$ is not an arithmetic mistake but the order-$h$ term that
disappears only in the limit. For $f(x)=1/x$ at $c=2$,
\[
  \frac{1/(2+h)-1/2}{h}
  =\frac{2-(2+h)}{2h(2+h)}=\frac{-1}{2(2+h)}\longrightarrow-\frac14 .
\]
In both computations the algebra removes the factor $h$ from the denominator, so that the
form $0/0$ becomes a function continuous at $h=0$.

That quotient form is important, but it does not survive when the variables become many. The
form that survives is the linear approximation.

\begin{definisi}[Derivative as the coefficient of a linear approximation]
A function $f$ is differentiable at $c$ with derivative $L$ if
\[
  f(c+h)=f(c)+Lh+r(h),
  \qquad\text{with}\qquad
  \lim_{h\to0}\frac{r(h)}{h}=0 .
\]
\end{definisi}

The remainder $r(h)$ is smaller than $h$; that is the content of the definition. For one
variable this form feels excessive, but only this form survives up to the Jacobian section.

\begin{proposisi}
The two definitions of the derivative above are equivalent, and the number $L$ in the second
definition is unique.
\end{proposisi}

\begin{proof}
If the limit of the quotient equals $L$, set $r(h)=f(c+h)-f(c)-Lh$. Then $r(h)/h$ is exactly
the quotient minus $L$, hence tends to $0$. Conversely, if $r$ with that property exists, the
quotient equals $L+r(h)/h$, which tends to $L$.

For uniqueness, suppose $L_{1}$ and $L_{2}$ both satisfy it. Then
$(L_{1}-L_{2})h=r_{2}(h)-r_{1}(h)$, and after dividing by $h$ we obtain
$L_{1}-L_{2}=r_{2}(h)/h-r_{1}(h)/h$. The right-hand side tends to $0$, whereas the left-hand
side does not depend on $h$. Hence $L_{1}=L_{2}$.
\end{proof}

The first property to hold on to: differentiability is stronger than continuity.

\begin{proposisi}
If $f$ is differentiable at $c$, then $f$ is continuous at $c$.
\end{proposisi}

\begin{proof}
For $h\neq0$ we have
\[
  f(c+h)-f(c)=h\cdot\frac{f(c+h)-f(c)}{h}.
\]
The second factor tends to $f'(c)$ and the first factor tends to $0$, so the product tends to
$0$. In other words $\lim_{h\to0}f(c+h)=f(c)$.
\end{proof}

The converse does not hold. The function $|x|$ is continuous at $0$ but not differentiable
there, because the two one-sided quotients have opposite signs. The function $x\sin(1/x)$,
equal to $0$ at $x=0$, is also continuous but not differentiable there: its quotient is
$\sin(1/x)$, which oscillates without approaching a single value.

\begin{tantangan}
Suppose $f(0)=0$ and, for $x\neq0$,
\[
  f(x)=|x|\bigl(2+\sin(1/x)\bigr).
\]
Show that $f$ is continuous at $0$ but has no derivative there. The two properties do not
contradict each other: continuity demands that the values of the function approach $f(0)$,
the derivative demands that the difference quotient approach a single number.

Hint 1: for continuity, one inequality that traps $|f(x)|$ between $0$ and something tending
to $0$ is enough.

Hint 2: for the second, compute $f(h)/h$ separately for $h>0$ and $h<0$. For $h>0$ the
quotient is $2+\sin(1/h)$, which oscillates between $1$ and $3$; for $h<0$ it oscillates
between $-3$ and $-1$. The two sets of values do not overlap.

If you get stuck, go back to the definition of the derivative at the beginning of this section
and compute the quotient for $h>0$ and then $h<0$ for the function $|x|$.
\end{tantangan}

The following computation rules free us from the definition every time.

\begin{teorema}[Basic rules]
Suppose $f$ and $g$ are differentiable at $c$. Then
\begin{enumerate}
  \item $(f+g)'(c)=f'(c)+g'(c)$ and $(af)'(c)=a f'(c)$ for $a\in\mathbb{R}$;
  \item $(fg)'(c)=f'(c)g(c)+f(c)g'(c)$;
  \item $\left(\dfrac{f}{g}\right)'(c)=\dfrac{g(c)f'(c)-f(c)g'(c)}{g(c)^{2}}$ if
        $g(c)\neq0$;
  \item if $n$ is an integer, the derivative of $x^{n}$ at $x\neq0$ is $nx^{n-1}$.
\end{enumerate}
\end{teorema}

\begin{proof}
Part (1) is immediate from the limit of a sum. For (2), write
\[
  f(x)g(x)-f(c)g(c)=f(x)\bigl[g(x)-g(c)\bigr]+g(c)\bigl[f(x)-f(c)\bigr].
\]
After dividing by $x-c$ and letting $x\to c$, the factor $f(x)$ tends to $f(c)$ because $f$ is
continuous at $c$ (the previous proposition), and both quotients tend to their derivatives.

For (3), it suffices to show that $1/g$ is differentiable with derivative $-g'(c)/g(c)^{2}$,
then apply (2) to $f\cdot(1/g)$. The quotient for $1/g$ is
\[
  \frac{1/g(x)-1/g(c)}{x-c}
  =-\frac{1}{g(x)g(c)}\cdot\frac{g(x)-g(c)}{x-c},
\]
which tends to $-g'(c)/g(c)^{2}$ because $g$ is continuous and $g(c)\neq0$. Part (4) follows
from (2) by induction on $n\ge0$, then from (3) for $n<0$.
\end{proof}
\section{The chain rule}
\tingkat{intermediate}

If one quantity depends on a second quantity, and the second quantity depends on time, then
the rate of change of the first quantity is the product of the two rates. That sentence sounds
like everyday arithmetic, and that is indeed the content of the chain rule.

Figure~\ref{fig:roda-gigi} uses two gears: the inner wheel turns three times per turn of the
middle wheel, and the middle wheel twice per turn of the outer wheel. One turn of the outer
wheel turns the inner wheel $3\times2=6$ times; adding the two ratios gives $5$, a number that
appears in no turn at all.

\begin{figure}[htbp]
\centering
\begin{tikzpicture}[scale=0.95]
% inner wheel, ratio 3
\draw[ink, thick] (0,0) circle (1.05);
\foreach \a in {0,20,...,340} { \draw[ink, thick] (\a:1.05) -- (\a:1.17); }
\fill[ink] (0,0) circle (1.5pt);
\draw[->, layerindigo, thick] (0,0) ++(150:0.55) arc (150:70:0.55);
% outer wheel, ratio 2
\begin{scope}[xshift=2.9cm]
  \draw[ink, thick] (0,0) circle (1.8);
  \foreach \a in {10,30,...,350} { \draw[ink, thick] (\a:1.8) -- (\a:1.95); }
  \fill[ink] (0,0) circle (1.5pt);
  \draw[->, layerbrown, thick] (0,0) ++(-20:0.75) arc (-20:-100:0.75);
\end{scope}
\node[font=\footnotesize, text=ink, align=center] at (0,-2.35) {inner wheel:\\3 turns};
\node[font=\footnotesize, text=ink, align=center] at (2.9,-2.35) {outer wheel:\\2 turns};
\node[font=\footnotesize, text=layerplum] at (1.45,2.15) {$3\times2=6$};
\end{tikzpicture}
\caption{The chain rule as two gears. The inner wheel turns $3$ times for one turn of the
outer wheel, and the outer wheel $2$ times for one turn of its output; the total ratio is
$3\times2=6$, not $3+2=5$.}
\label{fig:roda-gigi}
\end{figure}

\kamus{chain rule}{aturan rantai}{The rate at the end of the chain equals the product of the
rates at every link.}

\begin{teorema}[Chain rule]
Suppose $g$ is differentiable at $c$ and $f$ is differentiable at $g(c)$. Then $f\circ g$
is differentiable at $c$ and
\[
  (f\circ g)'(c)=f'\bigl(g(c)\bigr)\,g'(c).
\]
\end{teorema}

\begin{proof}
Write $y=g(c)$. Since $f$ is differentiable at $y$, by the linear approximation definition
there is a function $v$ with $v(s)\to0$ as $s\to y$ and
\[
  f(s)-f(y)=(s-y)\bigl[f'(y)+v(s)\bigr].
\]
In the same way, since $g$ is differentiable at $c$, there is $u$ with $u(t)\to0$ as
$t\to c$ and
\[
  g(t)-g(c)=(t-c)\bigl[g'(c)+u(t)\bigr].
\]
Substitute $s=g(t)$ in the first equation, then use the second:
\[
  f\bigl(g(t)\bigr)-f\bigl(g(c)\bigr)
  =(t-c)\bigl[g'(c)+u(t)\bigr]\bigl[f'(y)+v(g(t))\bigr].
\]
For $t\neq c$, divide by $t-c$. Since $g$ is continuous at $c$, $g(t)\to y$, so that
$v(g(t))\to0$, and the quotient tends to $g'(c)f'(y)$.
\end{proof}

The naive proof that divides by $g(t)-g(c)$ fails when $g(t)=g(c)$ for $t$ near $c$ but
$t\neq c$; a locally constant function, for instance, makes that denominator zero at every
point. The buffer $v$ above avoids that division altogether.

Short example: for $h(x)=(x^{2}+1)^{5}$, take $g(x)=x^{2}+1$ and $f(u)=u^{5}$. Then
$h'(x)=5(x^{2}+1)^{4}\cdot2x$ and $h'(1)=5\cdot16\cdot2=160$; direct expansion is possible,
but unpleasant.

\begin{resep}
\begin{enumerate}
  \item \textbf{Separate the two layers.} Point out which is inside and which is outside. For
        $k(x)=\sin(x^{3}+2x)$: the inner layer is $g(x)=x^{3}+2x$, the outer layer
        $f(u)=\sin u$.
  \item \textbf{Differentiate the outer layer, leave its content intact.} $f'(u)=\cos u$
        becomes $\cos(x^{3}+2x)$; do not simplify the content yet.
  \item \textbf{Multiply by the derivative of the inner layer.} $g'(x)=3x^{2}+2$, so that
        $k'(x)=(3x^{2}+2)\cos(x^{3}+2x)$.
  \item \textbf{Test one point.} At $x=0$ the formula gives $2$, and the central difference
        with $h=10^{-3}$ gives $1.9999997$.
\end{enumerate}
Warning about order: the outer layer is evaluated at $g(x)$, not at $x$. Swapping the place of
evaluation gives $\cos(x)\,(3x^{2}+2)$, whose value at $x=1$ is $2.7015$, whereas the true
value is $-4.9499$.
\end{resep}

Rate example: the area of a circle $A=\pi r^{2}$ with $\dot r=0.5$ cm per second. At
$r=10$ cm,
\[
  \frac{dA}{dt}=2\pi r\frac{dr}{dt}=2\pi\cdot10\cdot0.5=10\pi\approx31.4159
  \ \text{cm}^{2}\ \text{per second}.
\]
The number $10\pi$ is exact; at $r=20$ cm the same rate gives $20\pi$: the answer depends on
$r$.

\section{The mean value theorem and its consequences}
\tingkat{intermediate}

Between any two points on a smooth graph there is always a point whose slope equals the slope
of the secant joining those two points. That statement looks like geometric fantasy, yet it is
the bridge between the derivative and the function.

The key to its proof is the special case in which the two ends have equal value.

\begin{teorema}[Rolle]
If $f$ is continuous on $[a,b]$, differentiable on $(a,b)$, and $f(a)=f(b)$, then there is
$c\in(a,b)$ with $f'(c)=0$.\index{mean value theorem}
\end{teorema}

\begin{proof}
If $f$ is constant, every $c\in(a,b)$ works. If not, $f$ attains a maximum or a minimum that
differs from the value at the ends; since $f$ is continuous on a closed bounded interval,
that value is attained at some point, and that point cannot be an endpoint, so $c\in(a,b)$
(extreme value theorem, Chapter 11). At an interior extreme point, the quotient
$(f(c+h)-f(c))/h$ is $\le0$ for small $h>0$ and $\ge0$ for small $h<0$. The two inequalities
together force its two-sided limit, namely $f'(c)$, to be $0$.
\end{proof}

\begin{teorema}[Mean value]
If $f$ is continuous on $[a,b]$ and differentiable on $(a,b)$, then there is $c\in(a,b)$
with
\[
  f'(c)=\frac{f(b)-f(a)}{b-a}.
\]
\end{teorema}

\begin{proof}
Set $s=(f(b)-f(a))/(b-a)$ and $g(x)=f(x)-f(a)-s(x-a)$. Then $g$ is continuous on $[a,b]$,
differentiable on $(a,b)$, and $g(a)=g(b)=0$. By Rolle's theorem there is $c\in(a,b)$ with
$g'(c)=0$, that is, $f'(c)=s$.
\end{proof}

Numerical example: $f(x)=x^{3}$ on $[0,3]$ has secant slope $27/3=9$, and $3c^{2}=9$ gives
$c=\sqrt{3}\approx1.7321$; there is no guarantee that $c$ sits in the middle.

\begin{tangan}
Prepare a thin 30 cm board, two books, a ruler, and a pencil. Prop up the left end of the
board to a height of 4 cm, then press the middle slightly so that it curves upward.

Hold the ruler upright at 5 cm and at 25 cm from the raised end; the difference in height
divided by 20 cm gives one slope number. Repeat with 5 cm and 10 cm: the number changes,
because what is being measured is the average slope over different intervals.

Now lay the ruler flat at the 15 cm point and rotate it until only one point touches. Lift the
other end slowly: the ruler rotates about that point, exactly like the secant in
Figure~\ref{fig:sekan-singgung}. The slope of the ruler that touches a single point is the
slope at that point, and that is where the mean value theorem promises that such a point
exists.
\end{tangan}

\begin{akibat}[Monotonicity test]
Suppose $f$ is continuous on $[a,b]$ and differentiable on $(a,b)$. If $f'(x)\ge0$ for all
$x\in(a,b)$, then $f$ is nondecreasing on $[a,b]$. If $f'(x)>0$ for all $x$, then $f$ is
strictly increasing. If $f'(x)=0$ for all $x$, then $f$ is constant.
\end{akibat}

\begin{proof}
Take $a\le x<y\le b$. The mean value theorem on $[x,y]$ gives $c\in(x,y)$ with
$f(y)-f(x)=f'(c)(y-x)$. Since $y-x>0$, the sign of the difference in values follows the sign
of $f'(c)$. If $f'$ is zero everywhere, the difference is zero for every pair, so $f$ is
constant.
\end{proof}

The direction of the logic deserves attention: this corollary turns information about the
derivative into information about the function. The converse does not hold: a strictly
increasing function may have zero derivative at some points, such as $x^{3}$ at $0$. Nor does
the conclusion run backward from a single point: $f'(c)=0$ does not mean that $c$ is an
extreme point.

\begin{akibat}[Lipschitz continuity]
If $f$ is differentiable on an interval $I$ and $|f'(x)|\le M$ for all $x\in I$, then
$|f(x)-f(y)|\le M|x-y|$ for all $x,y\in I$.
\end{akibat}

\begin{proof}
For $x\neq y$, the mean value theorem on the interval between them gives $c$
with $f(x)-f(y)=f'(c)(x-y)$. Take the absolute value of both sides and use $|f'(c)|\le M$.
\end{proof}

Therefore a function with bounded derivative is automatically uniformly continuous (take
$\delta=\varepsilon/M$), but the converse does not hold: $\sqrt{x}$ is uniformly continuous on
$[0,1]$ with an unbounded derivative near zero.

\section{Taylor polynomials and their remainder}
\tingkat{advanced}

The tangent line is a polynomial of degree one that clings to the graph at a single point. If
we want it to cling longer, we add terms of degree two, three, and so on. The question is not
whether the additions improve the approximation, but how large the remainder is.

\begin{definisi}[Taylor polynomial]
Suppose $f$ has derivatives up to order $n$ at $a$. The Taylor polynomial of order $n$ of $f$
at $a$ is
\[
  P_{n}(x)=\sum_{k=0}^{n}\frac{f^{(k)}(a)}{k!}(x-a)^{k}.
\]
\end{definisi}

The requirement: $P_{n}$ and $f$ agree in value and in their first $n$ derivatives at $a$;
since $n+1$ coefficients determine one polynomial of degree $n$, $P_{n}$ is unique.

\begin{teorema}[Taylor with Lagrange remainder]
Suppose $f$ has derivatives up to order $n-1$ that are continuous on $[a,b]$ and a derivative
of order $n$ on $(a,b)$. For every $\beta\in(a,b]$ there is $\xi\in(a,\beta)$ with
\[
  f(\beta)=P_{n-1}(\beta)+\frac{f^{(n)}(\xi)}{n!}(\beta-a)^{n}.
\]
\end{teorema}

\begin{proof}
Define $M$ by $f(\beta)=P_{n-1}(\beta)+M(\beta-a)^{n}$, then set
\[
  g(t)=f(t)-P_{n-1}(t)-M(t-a)^{n},\qquad a\le t\le b .
\]
$P_{n-1}$ is built from the derivatives of $f$ at $a$, so $g(a)=\cdots=g^{(n-1)}(a)=0$, and
the choice of $M$ gives $g(\beta)=0$. Rolle gives $x_{1}\in(a,\beta)$ with $g'(x_{1})=0$;
since $g'(a)=0$, Rolle applied to $g'$ on $[a,x_{1}]$ gives $x_{2}$ with $g''(x_{2})=0$.
Repeat through $n$ steps: there is $x_{n}\in(a,\beta)$ with $g^{(n)}(x_{n})=0$. The derivative
of order $n$ of $P_{n-1}$ is zero and that of $(t-a)^{n}$ is $n!$, so
$0=f^{(n)}(x_{n})-n!M$, that is, $n!M=f^{(n)}(\xi)$ with $\xi=x_{n}$.
\end{proof}

For $n=1$ this theorem is exactly the mean value theorem; what is new is the ability to
compute errors of higher order.

\begin{contoh}
Approximate $e$ by $P_{3}(1)$ from $f(x)=e^{x}$ at $a=0$, then compare with the remainder
bound.
\end{contoh}

Here $f^{(k)}(0)=1$ for every $k$, so that
\[
  P_{3}(1)=1+1+\frac{1}{2}+\frac{1}{6}=\frac{8}{3}=2.6666\ldots
\]
The true value is $e=2.71828\ldots$, so the error is $0.0516152$. The Lagrange remainder
gives $|e-P_{3}(1)|=e^{\xi}/24$ for some $\xi\in(0,1)$, and since $e^{\xi}\le e<3$, the bound
is $3/24=0.125$. That bound is correct, but loose by a factor of about two and a half. Using
order five, $P_{5}(1)=2.71666\ldots$ with error $0.0016152$, whereas the bound is
$3/720=0.0041667$.

\begin{untukanak}
Imagine that you are guessing your friend's route from a few clues only. First clue: where he
is now. Your guess is a single point, and that point does not move. Second clue: which way he
is stepping. Your guess is a straight line, already closer, but a real route usually curves.
Third clue: how fast the direction changes. Your guess is a piece of a parabola, which clings
longer.

The more clues, the longer your guess clings, but never exactly. There is always a remainder.
That remainder can be bounded without knowing the true route, only from how large the change
is at the next level. In this chapter we shall learn why a straight line is the guess of order
one, why a parabola is of order two, and why a remainder bound is more useful than an answer
that claims to be exact.
\end{untukanak}
\section{Partial derivatives, the gradient, and directional derivatives}
\tingkat{intermediate}

When there is more than one input, ``rate of change'' is not necessarily a single number. It
depends on the direction in which we step. The simplest way to obtain a single number is to
freeze every coordinate except one, and then use the definition we already have.

\begin{definisi}[Partial derivative]
Suppose $f\colon U\to\mathbb{R}$ is defined on an open set
$U\subseteq\mathbb{R}^{n}$ and $a\in U$. The partial derivative of $f$ with respect to $x_{j}$
at $a$ is
\[
  \frac{\partial f}{\partial x_{j}}(a)
  =\lim_{h\to0}\frac{f(a+h\,e_{j})-f(a)}{h},
\]
with $e_{j}$ the unit vector in the direction of $x_{j}$, provided the limit exists.
\index{partial derivative}
\end{definisi}

Geometrically, the partial derivative with respect to $x$ is the slope of the slice of the
graph by the plane $y$ constant, hence an ordinary derivative of a function of one variable.

\kamus{gradient}{gradien}{An arrow pointing in the steepest direction on a hill, with length
equal to the steepness in that direction.}

\begin{definisi}[Gradient]
If all partial derivatives of $f$ at $a$ exist, the vector
\[
  \nabla f(a)=\left(\frac{\partial f}{\partial x_{1}}(a),\dots,
  \frac{\partial f}{\partial x_{n}}(a)\right)
\]
is called the gradient of $f$ at $a$.\index{gradient}
\end{definisi}

For $f(x,y)=x^{2}y^{3}$ we have $\partial f/\partial x=2xy^{3}$ and
$\partial f/\partial y=3x^{2}y^{2}$, so that $\nabla f(2,-1)=(-4,12)$.

\begin{definisi}[Directional derivative]
For a unit vector $u\in\mathbb{R}^{n}$, the directional derivative of $f$ at $a$ along $u$ is
\[
  D_{u}f(a)=\lim_{t\to0}\frac{f(a+tu)-f(a)}{t},
\]
provided the limit exists.\index{directional derivative}
\end{definisi}

The partial derivatives are the special case $u=e_{j}$. For other directions $D_{u}f$ cannot
be predicted from the list of partials, except in one well-behaved situation.

\begin{proposisi}
If $f$ is differentiable at $a$ (in the sense of the linear approximation of the next section),
then for every unit vector $u$ we have $D_{u}f(a)=\nabla f(a)\cdot u$.
\end{proposisi}

\begin{proof}
With $v=tu$ in the linear approximation definition,
$f(a+tu)=f(a)+t\,\nabla f(a)\cdot u+r(tu)$, so that
\[
  \frac{f(a+tu)-f(a)}{t}=\nabla f(a)\cdot u+\frac{r(tu)}{t}.
\]
Since $\|r(tu)\|/\|tu\|\to0$ and $\|tu\|=|t|$, the last term tends to $0$.
\end{proof}

\begin{proposisi}[Direction of steepest ascent]
If $\nabla f(a)\neq0$, then among all unit vectors $u$ the value $D_{u}f(a)$ is largest at
$u=\nabla f(a)/\|\nabla f(a)\|$, and the largest value is $\|\nabla f(a)\|$.
\end{proposisi}

\begin{proof}
By the Cauchy--Schwarz inequality (Chapter 10),
\[
  |D_{u}f(a)|=|\nabla f(a)\cdot u|\le\|\nabla f(a)\|\,\|u\|=\|\nabla f(a)\|.
\]
Equality holds exactly when $u$ is parallel to $\nabla f(a)$. Since $u$ is a unit vector, only
two signs are possible along that direction, and the positive sign gives the largest value.
\end{proof}

For $f(x,y)=x^{2}y^{3}$ at $(2,-1)$, the length of the gradient is $\sqrt{160}\approx12.6491$,
so the direction of steepest ascent is $(-4,12)/\sqrt{160}\approx(-0.3162,\,0.9487)$ with
rate $12.6491$ per unit of distance. The $x$ direction gives $-4$, the $y$ direction gives
$12$; both are smaller.

If a curve $t\mapsto\gamma(t)$ stays on the level set $f=\text{constant}$, then
$f(\gamma(t))$ is constant and its derivative is zero; the chain rule gives
$\nabla f(\gamma(t))\cdot\gamma'(t)=0$. Hence the gradient is perpendicular to the level
curve, and the direction of steepest ascent is perpendicular to the contour lines on a map.

For example $f(x,y)=x^{2}+2y^{2}$ at $(2,-1)$, a point with $f=6$, so it sits on the level
ellipse $x^{2}+2y^{2}=6$. Its gradient is $(2x,4y)=(4,-4)$, of length
$4\sqrt{2}\approx5.6569$, whereas the tangent direction of the ellipse there is $(1,1)$ and
$(4,-4)\cdot(1,1)=0$. Figure~\ref{fig:kontur-gradien} draws all three: the level ellipse, the
gradient arrow, and the tangent line perpendicular to it. The gradient arrow does not run
along the contour; it pierces the contour in the direction in which the value of $f$
increases.

\begin{figure}[htbp]
\centering
\begin{tikzpicture}[scale=0.92]
\draw[->, layergrey] (-3.0,0) -- (3.4,0) node[right, font=\scriptsize] {$x$};
\draw[->, layergrey] (0,-2.45) -- (0,2.5) node[above, font=\scriptsize] {$y$};
\draw[ink] (0,0) ellipse [x radius=1.414, y radius=1.000];
\draw[ink] (0,0) ellipse [x radius=2.000, y radius=1.414];
\draw[ink, thick] (0,0) ellipse [x radius=2.449, y radius=1.732];
\node[font=\scriptsize, anchor=west, text=ink] at (0.28,-0.98) {$c=2$};
\node[font=\scriptsize, anchor=west, text=ink] at (0.28,-1.40) {$c=4$};
\node[font=\scriptsize, anchor=west, text=ink] at (0.28,-1.98) {$c=6$};
\draw[accent, thick, dashed] (1.30,-1.70) -- (2.72,-0.28);
\draw[->, very thick, layerindigo] (2,-1) -- (2.707,-1.707);
\fill[ink] (2,-1) circle (1.7pt);
\node[font=\scriptsize, anchor=north west, text=ink] at (2.02,-1.05) {$(2,-1)$};
\node[font=\scriptsize, anchor=west, text=layerindigo] at (2.72,-1.86) {$\nabla f=(4,-4)$};
\node[font=\scriptsize, anchor=east, text=accent] at (1.32,-1.76) {direction $(1,1)$};
\end{tikzpicture}
\caption{Level curves $x^{2}+2y^{2}=c$ for $c=2,4,6$ and the gradient at $(2,-1)$. The
gradient arrow $(4,-4)$ is perpendicular to the tangent direction $(1,1)$ of the level curve
$c=6$, because $(4,-4)\cdot(1,1)=0$.}
\label{fig:kontur-gradien}
\end{figure}

\section{Differentiability, the Jacobian, and the multivariable chain rule}
\tingkat{advanced}

The true derivative is not a list of partials, but a single linear map that imitates the
function near a point; the list of partials can mislead.

\begin{definisi}[Differentiability]
Suppose $U\subseteq\mathbb{R}^{n}$ is open and $F\colon U\to\mathbb{R}^{m}$. $F$ is
differentiable at $a\in U$ if there is a linear map $L\colon\mathbb{R}^{n}\to\mathbb{R}^{m}$
with
\[
  F(a+v)=F(a)+Lv+r(v),
  \qquad
  \lim_{v\to0}\frac{\|r(v)\|}{\|v\|}=0 .
\]
The map $L$ is called the derivative of $F$ at $a$ and is written $F'(a)$.
\index{differentiability}
\end{definisi}

\begin{proposisi}
The map $L$ in that definition is unique if it exists.
\end{proposisi}

\begin{proof}
If $L_{1}$ and $L_{2}$ both work, then $(L_{1}-L_{2})v=r_{2}(v)-r_{1}(v)$. Take
$v=tw$ with $w$ fixed and $t>0$, divide by $t$, then let $t\to0$: this gives
$(L_{1}-L_{2})w=0$. Since this holds for every $w$, the two maps are equal.
\end{proof}

\begin{teorema}
If $F$ is differentiable at $a$, then all its partial derivatives at $a$ exist, and the matrix
of $L$ with respect to the standard bases is
\[
  J_{F}(a)=\begin{pmatrix}
    \dfrac{\partial F_{1}}{\partial x_{1}}(a) & \cdots &
    \dfrac{\partial F_{1}}{\partial x_{n}}(a)\\
    \vdots & \ddots & \vdots\\
    \dfrac{\partial F_{m}}{\partial x_{1}}(a) & \cdots &
    \dfrac{\partial F_{m}}{\partial x_{n}}(a)
  \end{pmatrix}.
\]
\end{teorema}

\begin{proof}
Take $v=te_{j}$ with $t\neq0$. The definition gives
$F(a+te_{j})-F(a)=tLe_{j}+r(te_{j})$, and after dividing by $t$ and letting $t\to0$, the $j$-th
column of $L$ appears as the limit of a quotient, namely the vector of partial derivatives
with respect to $x_{j}$.
\end{proof}

\kamus{Jacobian matrix}{Jacobian}{A table of all the slopes of a function of several variables:
each row is one output, each column is one input.}

That $m\times n$ matrix is called the Jacobian matrix of $F$ at $a$.\index{jacobian} For $m=1$
it becomes a single row, namely the transpose of the gradient. Example:
$F(x,y)=(x^{2}y,\;x+y^{2})$ at $(1,2)$ gives
\[
  J_{F}(1,2)=\begin{pmatrix}2xy & x^{2}\\ 1 & 2y\end{pmatrix}_{(1,2)}
  =\begin{pmatrix}4&1\\ 1&4\end{pmatrix},
\]
so that $F(1.01,\,1.98)\approx(2,5)+J_{F}(1,2)(0.01,-0.02)=(2.02,\,4.93)$,
whereas the exact value is $(2.019798,\,4.9304)$; the approximation is off by
$2\times10^{-4}$ in the first component.

Existence of the partials does not guarantee differentiability; the next two examples show two
levels of failure.

\begin{contoh}[Partials exist, function discontinuous]
Let $f(x,y)=\dfrac{xy}{x^{2}+y^{2}}$ for $(x,y)\neq(0,0)$ and $f(0,0)=0$. At the origin
$f(h,0)=0$, so both partial derivatives vanish there. But along the line $y=kx$ the value of
the function is $k/(1+k^{2})$: for $k=1$ the value is $1/2$, for $k=0$ the value is $0$. The
limit does not exist, and $f$ is not continuous at the origin.
\end{contoh}

\begin{contoh}[Continuous, partials exist, not differentiable]
Let $g(x,y)=\dfrac{xy}{\sqrt{x^{2}+y^{2}}}$ for $(x,y)\neq(0,0)$ and $g(0,0)=0$.
Since $2|xy|\le x^{2}+y^{2}$, we have $|g(x,y)|\le\frac12\sqrt{x^{2}+y^{2}}$, so $g$ is
continuous at the origin, and both of its partials vanish there. If $g$ were differentiable,
its linear map would have to be zero, so the remainder $r(v)=g(v)$ would have to satisfy
$\|r(v)\|/\|v\|\to0$. Take $v=t(1,1)/\sqrt{2}$ with $t>0$: then $g(v)=(t^{2}/2)/t=t/2$, so
$g(v)/\|v\|=1/2$ for every $t$. That ratio tends to $1/2$, not to $0$; hence $g$ is not
differentiable at the origin.
\end{contoh}

The consequence matters: in the second example $\nabla g(0,0)=0$ yet the diagonal directional
derivative does not exist, because its quotient is $\pm1/2$ depending on the sign of $t$. Thus
$D_{u}f=\nabla f\cdot u$ requires differentiability.

A condition that is easy to check and sufficient for differentiability is continuity of the
partials.

\begin{teorema}
If the partial derivatives of $F$ exist and are continuous near $a$, then $F$ is
differentiable at $a$.
\end{teorema}

\begin{proof}
It suffices to prove this for $m=1$ and $n=2$; the other cases follow component by component
with increasing indices. Write $a=(p,q)$ and $v=(h,k)$. Then
\[
  f(p+h,q+k)-f(p,q)
  =\bigl[f(p+h,q+k)-f(p,q+k)\bigr]+\bigl[f(p,q+k)-f(p,q)\bigr].
\]
By the mean value theorem in one variable, the first term equals
$h\,\partial_{x}f(\xi,q+k)$ for some $\xi$ between $p$ and $p+h$, and the second term equals
$k\,\partial_{y}f(p,\eta)$ for some $\eta$ between $q$ and $q+k$. Therefore
\[
  f(a+v)-f(a)-\nabla f(a)\cdot v
  =h\bigl[\partial_{x}f(\xi,q+k)-\partial_{x}f(a)\bigr]
  +k\bigl[\partial_{y}f(p,\eta)-\partial_{y}f(a)\bigr].
\]
Divide by $\|v\|=\sqrt{h^{2}+k^{2}}$: since $|h|/\|v\|\le1$ and $|k|/\|v\|\le1$, and both
differences in brackets tend to $0$ (this is where continuity of the partials is used), the
right-hand side divided by $\|v\|$ tends to $0$.
\end{proof}

Now composition.

\begin{teorema}[Multivariable chain rule]
Suppose $H\colon\mathbb{R}^{n}\to\mathbb{R}^{k}$ is differentiable at $a$ and
$G\colon\mathbb{R}^{k}\to\mathbb{R}^{m}$ is differentiable at $H(a)$. Then $G\circ H$
is differentiable at $a$ and
\[
  (G\circ H)'(a)=G'(H(a))\,H'(a),
\]
that is, the composition of two linear maps; in matrices,
$J_{G\circ H}(a)=J_{G}(H(a))\,J_{H}(a)$.
\end{teorema}

\begin{proof}
Write $A=H'(a)$, $B=G'(y_{0})$ with $y_{0}=H(a)$, and
\[
  H(a+v)=H(a)+Av+r(v),\qquad G(y_{0}+w)=G(y_{0})+Bw+q(w),
\]
with $\|r(v)\|\le\varepsilon(v)\|v\|$ and $\|q(w)\|\le\eta(w)\|w\|$, where
$\varepsilon(v)\to0$ and $\eta(w)\to0$. For $w=H(a+v)-H(a)=Av+r(v)$ we have
$\|w\|\le(\|A\|+\varepsilon(v))\|v\|$, so $w\to0$ as $v\to0$. Then
\[
  G\bigl(H(a+v)\bigr)-G(H(a))-BAv
  =B\bigl(H(a+v)-H(a)-Av\bigr)+q(w)
  =B\,r(v)+q(w).
\]
Divide by $\|v\|$:
\[
  \frac{\|B\,r(v)+q(w)\|}{\|v\|}
  \le\|B\|\,\varepsilon(v)+\bigl(\|A\|+\varepsilon(v)\bigr)\eta(w),
\]
and the right-hand side tends to $0$. Hence $BA$ satisfies the definition of the derivative of
$G\circ H$ at $a$.
\end{proof}

The matrix form explains why the order of multiplication may not be reversed: $J_{H}$ is
$k\times n$ and $J_{G}$ is $m\times k$, so only one order can be multiplied. For scalar
functions the result is $\frac{d}{dt}f(\gamma(t))=\nabla f(\gamma(t))\cdot
\gamma'(t)$.

Finally, second derivatives of a function of several variables are stored in a single matrix.

\kamus{Hessian matrix}{Hessian}{A table of curvature: whether the surface bends upward,
downward, or in both directions at once.}

\begin{definisi}[Hessian matrix]
If all second partial derivatives of $f\colon U\to\mathbb{R}$ exist at $a$, the matrix
\[
  H_{f}(a)=\left(\frac{\partial^{2}f}{\partial x_{i}\,\partial x_{j}}(a)\right)_{i,j=1}^{n}
\]
is called the Hessian matrix of $f$ at $a$.\index{hessian}
\end{definisi}

For a function whose second partials are continuous, $H_{f}(a)$ is symmetric (Clairaut's
theorem, proved in Stewart in the section on partial derivatives). Without continuity, that
conclusion can fail.

The Hessian is the derivative of the gradient. If the second partials are continuous near a
critical point $a$ ($\nabla f(a)=0$), then
\[
  f(a+v)=f(a)+\tfrac12\,v^{\top}H_{f}(a)\,v+o\bigl(\|v\|^{2}\bigr),
\]
so that the sign of the quadratic form $v\mapsto v^{\top}Hv$ determines the type of the
critical point.

\begin{proposisi}[Second derivative test]
Suppose $f$ has continuous second partial derivatives near a critical point $a$
($\nabla f(a)=0$).
\begin{enumerate}
  \item If $H_{f}(a)$ is positive definite, $a$ is a strict local minimum.
  \item If $H_{f}(a)$ is negative definite, $a$ is a strict local maximum.
  \item If $H_{f}(a)$ has eigenvalues of different signs, $a$ is a saddle point.
  \item If $H_{f}(a)$ is semidefinite but singular, this test gives no conclusion.
\end{enumerate}
\end{proposisi}

The first three parts follow from the quadratic expansion above: along a direction with
$v^{\top}Hv>0$ the function increases, along a direction with $v^{\top}Hv<0$ it decreases. The
complete proof is in Stewart in the section on maximum and minimum values. The fourth part is
genuinely empty: $x^{4}+y^{4}$ and $x^{4}-y^{4}$ both have $H=0$ at the origin, but the first
is a minimum and the second a saddle. For two variables, $AC-B^{2}>0$ with $A>0$ means
positive definite, $AC-B^{2}<0$ means indefinite.
\section{Worked examples}
\tingkat{intermediate}

The following three computations use the same numbers as the code in the lab section.

First, finite differences for $f(x)=\sin x$ at $x_{0}=1$; the exact value is $\cos 1$, whose
digits appear in the code output in the lab section. The two formulas:
\[
  \text{forward: }\frac{f(x_{0}+h)-f(x_{0})}{h},\qquad
  \text{central: }\frac{f(x_{0}+h)-f(x_{0}-h)}{2h}.
\]

\begin{center}
\begin{tabular}{rrrrr}
\hline
$h$ & forward difference & relative error & central difference & relative error\\
\hline
$10^{-1}$ & $0.497363752535$ & $7.947\times10^{-2}$ & $0.539402252170$ & $1.666\times10^{-3}$\\
$10^{-2}$ & $0.536085981012$ & $7.804\times10^{-3}$ & $0.540293300875$ & $1.667\times10^{-5}$\\
$10^{-4}$ & $0.540260231419$ & $7.787\times10^{-5}$ & $0.540302304968$ & $1.667\times10^{-9}$\\
$10^{-6}$ & $0.540301885121$ & $7.787\times10^{-7}$ & $0.540302305896$ & $5.130\times10^{-11}$\\
$10^{-8}$ & $0.540302302898$ & $5.497\times10^{-9}$ & $0.540302308449$ & $4.777\times10^{-9}$\\
$10^{-10}$ & $0.540302247387$ & $1.082\times10^{-7}$ & $0.540302247387$ & $1.082\times10^{-7}$\\
\hline
\end{tabular}
\end{center}

The pattern of the errors is no accident. For the forward difference, the Taylor expansion
gives
\[
  \frac{f(x_{0}+h)-f(x_{0})}{h}=f'(x_{0})+\frac{h}{2}f''(x_{0})+\mathcal{O}(h^{2}),
\]
so its relative error is of order $h$ and drops by a factor of ten each time $h$ is divided by
ten: $7.947\times10^{-2}$, $7.804\times10^{-3}$, $7.787\times10^{-5}$. The numbers can be
matched: $\frac{h}{2}f''(1)=-\frac{1}{2}\cdot10^{-2}\sin 1=-4.207\times10^{-3}$, and
$4.207\times10^{-3}/0.5403=7.787\times10^{-3}$. The central difference gives
\[
  \frac{f(x_{0}+h)-f(x_{0}-h)}{2h}=f'(x_{0})+\frac{h^{2}}{6}f'''(x_{0})+\mathcal{O}(h^{4}),
\]
so its error is of order $h^{2}$ and drops by a factor of a hundred: $1.666\times10^{-3}$,
$1.667\times10^{-5}$, $1.667\times10^{-9}$. Here $\frac{10^{-4}}{6}\cos 1
=9.005\times10^{-6}$, and $9.005\times10^{-6}/0.5403$ again gives $1.667\times10^{-5}$.

At $h=10^{-8}$ and $10^{-10}$ the central difference deteriorates again: the truncation error
is already smaller than the rounding error, which grows like $\varepsilon/h$, with
$\varepsilon\approx2.2\times10^{-16}$. The best entry in the table, $h=10^{-6}$ with error
$5\times10^{-11}$, has an order consistent with $\varepsilon^{1/3}\approx6\times10^{-6}$.

Second, the gradient of $f(x,y)=x^{2}y^{3}$ at $(2,-1)$ with analytic value $(-4,12)$. The
central difference gives an error of $3.333\times10^{-5}$ in the $y$ component for
$h=10^{-2}$, dropping to $3.333\times10^{-9}$ for $h=10^{-4}$: order $h^{2}$. The $x$
component is exact to machine error, because $f$ is quadratic in $x$.

Third, the chain rule on a composition of different dimensions. Let
\[
  H(x,y)=\bigl(x^{2}+y,\;xy,\;y^{3}\bigr),\qquad
  G(u,v,w)=\bigl(uv+w,\;u^{2}-vw\bigr),
\]
and consider $F=G\circ H$ at $a=(1,2)$. Here $H(a)=(3,2,8)$, and
\[
  J_{H}(1,2)=\begin{pmatrix}2&1\\ 2&1\\ 0&12\end{pmatrix}_{3\times2},
  \qquad
  J_{G}(3,2,8)=\begin{pmatrix}2&3&1\\ 6&-8&-2\end{pmatrix}_{2\times3}.
\]
Their product is $2\times2$:
\[
  J_{G}(3,2,8)\,J_{H}(1,2)=
  \begin{pmatrix}2&3&1\\ 6&-8&-2\end{pmatrix}
  \begin{pmatrix}2&1\\ 2&1\\ 0&12\end{pmatrix}
  =\begin{pmatrix}10&17\\ -4&-26\end{pmatrix}.
\]
\kertas{For the composition above, compute $J_{F}$ at $(2,1)$ in two ways: the product
$J_{G}(H(2,1))J_{H}(2,1)$ and the direct derivative of the explicit form. Then check the
direction $v=(1,3)$ by comparing $J_{F}(2,1)v$ against $\bigl(F(a+sv)-F(a)\bigr)/s$ for
$s=10^{-2}$ and $s=10^{-4}$. Why does the error drop only by a factor of $10$, not $100$?}

\section{Computational lab}
\tingkat{advanced}

The following code is run as is with numpy 2.5.2 on Python 3: the first part repeats the three
computations above, then one experiment that cannot be done by hand, finite differences under
measurement noise.

\begin{lab}
\begin{verbatim}
import numpy as np

# 1. finite differences for f(x) = sin x at x0 = 1
x0 = 1.0
exact = float(np.cos(x0))
print("cos(1) =", exact)
for h in [1e-1, 1e-2, 1e-4, 1e-6, 1e-8, 1e-10]:
    fwd = (np.sin(x0 + h) - np.sin(x0)) / h
    cen = (np.sin(x0 + h) - np.sin(x0 - h)) / (2*h)
    em = abs(fwd - exact) / abs(exact)
    ep = abs(cen - exact) / abs(exact)
    print(f"{h:8.0e} fwd={fwd:.12f} {em:.3e} "
          f"cen={cen:.12f} {ep:.3e}")
\end{verbatim}
Actual output:
\begin{verbatim}
cos(1) = 0.5403023058681398
   1e-01 fwd=0.497363752535 7.947e-02 cen=0.539402252170 1.666e-03
   1e-02 fwd=0.536085981012 7.804e-03 cen=0.540293300875 1.667e-05
   1e-04 fwd=0.540260231419 7.787e-05 cen=0.540302304968 1.667e-09
   1e-06 fwd=0.540301885121 7.787e-07 cen=0.540302305896 5.130e-11
   1e-08 fwd=0.540302302898 5.497e-09 cen=0.540302308449 4.777e-09
   1e-10 fwd=0.540302247387 1.082e-07 cen=0.540302247387 1.082e-07
\end{verbatim}
\end{lab}

The last line deserves a second reading: at $h=10^{-10}$ both formulas give identical digits,
because they are both computing the same rounding error, not the derivative.

\begin{lab}
\begin{verbatim}
def f(p):
    return p[0]**2 * p[1]**3

def grad_fd(fun, a, h):
    g = np.zeros_like(a)
    for j in range(a.size):
        e = np.zeros_like(a)
        e[j] = h
        g[j] = (fun(a + e) - fun(a - e)) / (2*h)
    return g

a = np.array([2.0, -1.0])
g_an = np.array([2*a[0]*a[1]**3, 3*a[0]**2*a[1]**2])
print("gradient =", g_an)
for h in [1e-2, 1e-4, 1e-6]:
    g = grad_fd(f, a, h)
    print(f"{h:8.0e} {g[0]:.12f} {g[1]:.12f} "
          f"{abs(g[0]+4)/4:.3e} {abs(g[1]-12)/12:.3e}")
\end{verbatim}
Actual output:
\begin{verbatim}
gradient = [-4. 12.]
   1e-02 -4.000000000000 12.000400000000 1.577e-14 3.333e-05
   1e-04 -4.000000000004 12.000000039998 1.000e-12 3.333e-09
   1e-06 -4.000000000115 11.999999999679 2.876e-11 2.676e-11
\end{verbatim}
\end{lab}

\begin{lab}
\begin{verbatim}
def hh(p):
    x, y = p
    return np.array([x*x + y, x*y, y**3])

def gg(q):
    u, v, w = q
    return np.array([u*v + w, u*u - v*w])

def F(p):
    return gg(hh(p))

def jac(fun, p, step):
    n = p.size
    J = np.zeros((fun(p).size, n))
    for j in range(n):
        e = np.zeros(n)
        e[j] = step
        J[:, j] = (fun(p + e) - fun(p - e)) / (2*step)
    return J

a = np.array([1.0, 2.0])
q = hh(a)
Jh = np.array([[2*a[0], 1.0], [a[1], a[0]], [0.0, 3*a[1]**2]])
Jg = np.array([[q[1], q[0], 1.0], [2*q[0], -q[2], -q[1]]])
J_chain = Jg @ Jh
J_num = jac(F, a, 1e-6)
print("h(1,2) =", q)
print("Jg @ Jh =\n", J_chain)
print("max difference =", f"{np.max(np.abs(J_num - J_chain)):.3e}")

v = np.array([0.7, -1.3])
d_chain = J_chain @ v
d_num = (F(a + 1e-6*v) - F(a - 1e-6*v)) / 2e-6
print("J v =", d_chain)
print("finite difference =", d_num)
print("relative error =",
      f"{np.max(np.abs(d_num - d_chain)/np.abs(d_chain)):.3e}")
\end{verbatim}
Actual output:
\begin{verbatim}
h(1,2) = [3. 2. 8.]
Jg @ Jh =
 [[ 10.  17.]
 [ -4. -26.]]
max difference = 1.858e-09
J v = [-15.1  31. ]
finite difference = [-15.1  31. ]
relative error = 9.860e-11
\end{verbatim}
\end{lab}

The difference $1.858\times10^{-9}$ between the chain-rule Jacobian and the finite-difference
Jacobian is the truncation error of the central difference, not an algebraic mistake; the
directional check gives $9.86\times10^{-11}$, because both use the same point.

\begin{lab}
\begin{verbatim}
rng = np.random.default_rng(7)
sigma = 1e-3
for h in [1e-1, 1e-2, 1e-3, 1e-4, 1e-5]:
    err = []
    for _ in range(2000):
        s1 = np.sin(x0 + h) + sigma*rng.standard_normal()
        s0 = np.sin(x0) + sigma*rng.standard_normal()
        err.append(abs((s1 - s0)/h - exact))
    print(f"{h:8.0e} mean error = {np.mean(err):.3e}")
\end{verbatim}
Actual output:
\begin{verbatim}
   1e-01 mean error = 4.203e-02
   1e-02 mean error = 1.112e-01
   1e-03 mean error = 1.160e+00
   1e-04 mean error = 1.117e+01
   1e-05 mean error = 1.088e+02
\end{verbatim}
\end{lab}

The last experiment is the heart of this chapter in numbers. With noise
$\sigma=10^{-3}$, the quotient carries two errors: a truncation bias of $\approx0.42h$ and
noise that amplifies like $\sqrt{2}\sigma/h$. At $h=10^{-1}$ the bias dominates ($0.042$),
at $h=10^{-4}$ the noise dominates ($11.2$), and the two balance when
$0.42h=\sqrt{2}\cdot10^{-3}/h$, that is $h\approx0.058$: with an instrument that bad, the
best step is about six hundredths, not one hundred-millionth.

That is the honest answer to the opening question. The derivative is a limit, but its best
approximation is not the closest possible limit. The best is a compromise between two errors
that move in opposite directions, and that compromise can be computed before a single number
is measured.
\begin{untukdosen}
The order that works: the quotient definition, then the linear approximation form, then the
computation rules. That form feels excessive for one variable, but only it survives up to the
Jacobian; if it is introduced later, students memorize the chain rule as matrix multiplication
without knowing why the order is what it is.

Two typical misconceptions. First, ``the partial derivatives exist means the function is
smooth''; the example $xy/\sqrt{x^{2}+y^{2}}$ on the board breaks it, provided the facts are
computed together: partials zero, continuous, yet the diagonal direction gives $1/2$. Second,
the direction of multiplication in the chain rule: students write $J_{H}J_{G}$ because the
order in which the functions are written is the reverse of the order in which they are
applied; the way to check is the size of the matrices.

Blackboard activity: ask the class to compute both partials at the origin for the two
counterexamples, then to decide whether that function has a tangent plane. Sign of readiness
to go further: students explain why $D_{u}f=\nabla f\cdot u$ fails in the second example, and
use the word ``remainder'', not ``formula''.
\end{untukdosen}

\section{Summary}
\tingkat{basic}

\begin{ringkasan}
\begin{itemize}
  \item The derivative $f'(c)$ is the limit of a difference quotient, and also the unique
        coefficient of the linear approximation $f(c+h)=f(c)+Lh+r(h)$ with $r(h)/h\to0$.
        Differentiability guarantees continuity, not the other way around.
  \item The chain rule $(f\circ g)'=f'(g(c))g'(c)$ is proved through the linear approximation,
        and the proof survives into many dimensions: $J_{G\circ H}(a)=J_{G}(H(a))J_{H}(a)$.
  \item The mean value theorem links the sign of the derivative to the behavior of the
        function: $f'\ge0$ means nondecreasing, $|f'|\le M$ means Lipschitz; it fails for
        vector functions.
  \item Taylor with Lagrange remainder uses an unknown $\xi$, so only an upper bound for
        $|f^{(n)}|$ is used.
  \item For several variables, the derivative is a linear map with a Jacobian matrix. Partials
        may exist without continuity, and may exist together with continuity without
        differentiability.
  \item The gradient points in the direction of steepest ascent with rate $\|\nabla f\|$,
        perpendicular to the contours; a definite Hessian determines the type of a critical
        point.
  \item Finite differences are of order $h$ (forward) or $h^{2}$ (central), and the rounding
        error grows like $\varepsilon/h$. The smallest step is not the best step.
\end{itemize}
\end{ringkasan}

\section{Exercises}
\tingkat{intermediate}

\begin{latihan}
\begin{enumerate}
  \item Compute $f'(2)$ from the definition for $f(x)=x^{3}$ and $f(x)=\sqrt{x}$, until the
        factor $h$ disappears from the denominator.

  \item Prove that if $f'(c)>0$, there is $\delta>0$ such that $f(c+h)>f(c)$ for every
        $h\in(0,\delta)$, and explain why this does not mean that $f$ is increasing on some
        interval.

  \item Prove the quotient rule by differentiating $1/g$ first, then using the product rule.
        State at which step the continuity of $g$ is used.

  \item Assess this claim: \emph{if $f$ has all its partial derivatives at a point, then $f$
        is continuous at that point.} State whether it is true, false, or not yet supported by
        proof, and include a counterexample if needed.

  \item For $f(x,y)=x^{3}y-xy^{3}$, compute the gradient at $(2,1)$, determine the direction
        of steepest ascent, and compute its maximum rate. Check with a central difference
        $h=10^{-4}$ and compare the relative error with the order $h^{2}$.

  \item Prove that if $f'=0$ everywhere on an interval, then $f$ is constant. Explain why that
        proof fails on a disconnected set, and give an example.

  \item Suppose $H(x,y)=(e^{x}y,\;x+y^{2})$ and $G(u,v)=(u^{2}v,\;u-v)$. Compute
        $J_{G\circ H}$ at $(0,1)$ in two ways: the matrix product and the direct derivative.
        (Hint: $H(0,1)=(1,1)$.)

  \item Find all critical points of $f(x,y)=x^{3}+y^{3}-3xy$, then classify them with the
        Hessian; explain why the sign of the determinant is not enough when the determinant is
        zero.

  \item Estimate $\sqrt{1.04}$ with the Taylor polynomial of order two at $a=1$, then bound
        its error with the Lagrange remainder. Compare that bound with the actual error: is it tight? (Hint:
        $f''(x)=-\frac14x^{-3/2}$ and $f'''$ is positive on that interval.)
\end{enumerate}
\end{latihan}

\begin{untukpeneliti}
What is proved in this chapter are the classical statements for one and several variables: the
linear approximation and the uniqueness of its map, the chain rule, the mean value theorem and
its consequences, Taylor with Lagrange remainder, the equivalence of differentiability with the
Jacobian, and the sufficient condition of continuous partials. All of them are theorems, not
conventions, and they hold in finite-dimensional spaces over $\mathbb{R}$.

What is often popularized is the claim ``the gradient is the derivative''. Without
differentiability that claim is false, and $xy/\sqrt{x^{2}+y^{2}}$ is a counterexample that
stays continuous. The claim ``the Hessian determines the type of a critical point'' holds only
when the Hessian is nonsingular; the singular case requires higher derivatives or direct
analysis.

Primary sources. Rudin, \emph{Principles of Mathematical Analysis}, the chapters
\emph{Differentiation} and \emph{Functions of Several Variables}, contains all the statements
above. Spivak, \emph{Calculus on Manifolds}, the chapter \emph{Differentiation}, takes the same
route and is shorter. Stewart, \emph{Calculus: Early Transcendentals} ninth edition, sections
14.3 to 14.7, supplies the computational counterparts, including Clairaut's theorem.

Limits of validity: all arguments assume that the domain is open and that the derivatives used
exist. For functions that are differentiable only almost everywhere, the monotonicity and
Taylor conclusions collapse, and that is territory of measure theory which this book does not
cover. For vector functions, the mean value theorem does not hold; what replaces it is the mean
value inequality, which requires an operator norm. The best order of a finite difference
depends on the scale of the function, the noise, and the arithmetic; this chapter has only
shown the order empirically.
\end{untukpeneliti}

\peta{This chapter follows Rudin, \emph{Principles of Mathematical Analysis}, the chapters
\emph{Differentiation} and \emph{Functions of Several Variables}; its computational version is
Stewart, \emph{Calculus: Early Transcendentals} ninth edition, sections 14.3 to 14.7; its
abstract counterpart is Spivak, \emph{Calculus on Manifolds}, the chapter
\emph{Differentiation}. Continuity and the extreme value theorem are used from Chapter 11, the
Cauchy--Schwarz inequality from Chapter 10, and the Jacobian norm returns in Chapter 23 and
Chapter 28.}

% Judul penuh meluber 62,75 pt di kepala halaman; hanya mark header yang
% dipendekkan (sekali pakai, lihat \judulpendek di preamble-en.tex).
\judulpendek{Riemann Integrals}
\chapter{Riemann Integrals and the Fundamental Theorem of Calculus}

A program computes the area under $y=x^2$ on $[0,1]$ with 1000 slices of equal
width. If the height of each rectangle is taken at the left endpoint, the result is $0.3328335$; if at
the right endpoint, $0.3338335$. The difference is $0.001$, and there is no reason to prefer either one.
The correct value is $1/3$, which does not appear in either computation. Increasing the slices to a
million does not help: both answers still bracket the value, they never catch it.

What is needed is not a finer partition, but a way to read two approximations at
once. If the difference can be made smaller than any positive number,
then there is one number trapped between them, and that is the integral. This entire
chapter rests on that simple idea.

Chapter 12 taught us to take change apart into rates: from position to velocity,
from a function to its derivative. What is still missing is the way home. We know that the velocity at
second $t$ is $v(t)$, but we do not yet know how much distance was covered in one
minute. The intuition is clear: distance is the sum of velocity times small slices of time.
What is not clear is how to add up infinitely many such terms correctly.
This chapter answers that, and the answer is surprising: the way home is the antiderivative.

\begin{untukanak}
Imagine we want to know the area of a lake whose shape is irregular. There is no
formula for such a shape. All we have is transparent graph paper. We
lay that paper over the drawing of the lake, then count how many squares lie entirely
inside the water. The result is too small. We count again how many squares are
touched by water, even slightly. The result is too large. The true area lies
between those two numbers.

Now cut the squares into ten times smaller ones. The lower number rises, the upper number
falls, and the two grow closer. If the squares are very small, the two numbers are almost equal,
and we can state the area of the lake to any accuracy we like.

There is one condition that is easy to overlook. If the lake has very thin coral islands
scattered everywhere, the squares that touch them can keep the upper estimate far
from the lower estimate. Mathematicians call that condition integrability. In this chapter we
shall learn why two bracketing estimates can be trusted, and when they fail.
\end{untukanak}

\section{Partitions, Upper Sums, and Lower Sums}
\tingkat{basic}

We want to give meaning to the area under a graph without assuming that the area already
has a meaning: enclose the region with rectangles from inside and from outside, then
shrink the difference.

\begin{definisi}[Partitions and upper-lower sums]\index{partition}
Let $[a,b]$ be a closed interval. A \emph{partition} $\mathcal{P}$ of $[a,b]$
is a finite set of points
\[
a=x_0<x_1<\cdots<x_n=b .
\]
The width of the $i$-th slice is $\Delta x_i=x_i-x_{i-1}$, and the \emph{mesh} of that partition
is the widest slice width, $\|\mathcal{P}\|=\max_i\Delta x_i$.

Let $f$ be bounded on $[a,b]$. On the $i$-th slice write
\[
M_i=\sup\{f(x):x_{i-1}\le x\le x_i\},\qquad
m_i=\inf\{f(x):x_{i-1}\le x\le x_i\}.
\]
The \emph{upper sum} and the \emph{lower sum} of the partition $\mathcal{P}$ are
\begin{equation}
U(\mathcal{P},f)=\sum_{i=1}^{n}M_i\,\Delta x_i,
\qquad
L(\mathcal{P},f)=\sum_{i=1}^{n}m_i\,\Delta x_i .
\end{equation}
\end{definisi}

We use supremum and infimum, not the largest and smallest values, because $f$ need not
attain its extremes on a slice: the function equal to $x$ for $x<1$ and to $0$
at $x=1$ is bounded on $[0,1]$, but its largest value on the last slice does
not exist, whereas its supremum does, namely $1$.

The upper sum is always at least equal to the lower sum, because $M_i\ge m_i$ for
every $i$. For $f(x)=x^2$ on $[0,1]$ with four slices of equal width, the lower
sum is
\[
L=\frac{1}{4}\Bigl(0+\frac{1}{16}+\frac{4}{16}+\frac{9}{16}\Bigr)=\frac{7}{32},
\]
and the upper sum is $\frac{15}{32}$. Both bracket the value sought, with a difference of
$\frac14$ that is still large. If $f$ takes negative values, these sums compute signed
area, and the minus sign is part of the definition, not an error.

Figure~\ref{fig:rect-halus} shows two values of $n$ at once, with the right endpoint
as height. At $n=4$ the staircase is coarse and $U_4=0.46875$; at $n=16$ the serration is
fine and $U_{16}=0.3652343750$. The distance from the staircase to the line $1/3$ shrinks from $0.1354$
to $0.0319$.

\begin{figure}[htbp]
\centering
\begin{tikzpicture}
\begin{axis}[width=0.47\linewidth,height=4.6cm,
  axis lines=left,xmin=0,xmax=1.02,ymin=0,ymax=1.12,
  xtick={0,0.5,1},ytick={0,0.5,1},
  tick label style={font=\scriptsize},label style={font=\scriptsize},
  xlabel={$x$},
  title={\scriptsize $n=4$: $U_4=0.46875$}]
\pgfplotsinvokeforeach{1,2,3,4}{
  \draw[fill=layerteal!18,draw=layerteal!65,line width=0.4pt]
    (axis cs:{(#1-1)/4},0) rectangle (axis cs:{#1/4},{#1*#1/16});
}
\draw[densely dashed,layergrey,line width=0.6pt]
  (axis cs:0,0.33333) -- (axis cs:1,0.33333);
\addplot[thick,accent,domain=0:1,samples=80] {x^2};
\node[font=\scriptsize,text=layergrey,anchor=west] at (axis cs:0.06,0.40) {$1/3$};
\end{axis}
\end{tikzpicture}\hfill
\begin{tikzpicture}
\begin{axis}[width=0.47\linewidth,height=4.6cm,
  axis lines=left,xmin=0,xmax=1.02,ymin=0,ymax=1.12,
  xtick={0,0.5,1},ytick={0,0.5,1},
  tick label style={font=\scriptsize},label style={font=\scriptsize},
  xlabel={$x$},
  title={\scriptsize $n=16$: $U_{16}=0.3652343750$}]
\pgfplotsinvokeforeach{1,...,16}{
  \draw[fill=layerteal!18,draw=layerteal!65,line width=0.3pt]
    (axis cs:{(#1-1)/16},0) rectangle (axis cs:{#1/16},{#1*#1/256});
}
\draw[densely dashed,layergrey,line width=0.6pt]
  (axis cs:0,0.33333) -- (axis cs:1,0.33333);
\addplot[thick,accent,domain=0:1,samples=80] {x^2};
\node[font=\scriptsize,text=layergrey,anchor=west] at (axis cs:0.06,0.40) {$1/3$};
\end{axis}
\end{tikzpicture}
\caption{Rectangular approximation of $f(x)=x^{2}$ on $[0,1]$ with right-endpoint heights.
At $n=4$ the upper sum is $U_4=0.46875$, at $n=16$ the upper sum is $U_{16}=0.3652343750$,
and both are still above the true value $1/3=0.333\ldots$ (dashed line).
The difference $U_n-L_n$ is exactly $1/n$, so going from $n=4$ to $n=16$ cuts the distance from the
staircase to that line down to a quarter of it.}
\label{fig:rect-halus}
\end{figure}


There are two facts that will be used again and again. The first: refining a partition
does not make the estimate worse.

\begin{lemma}[Refinement]\index{refinement}
If $\mathcal{P}^{*}$ is a refinement of $\mathcal{P}$, meaning that
$\mathcal{P}\subseteq\mathcal{P}^{*}$, then
\[
L(\mathcal{P},f)\le L(\mathcal{P}^{*},f)\le U(\mathcal{P}^{*},f)\le U(\mathcal{P},f).
\]
\end{lemma}

\begin{proof}
It suffices to prove this for a refinement that adds one point $c$ inside the slice
$[x_{i-1},x_i]$, since a general refinement is obtained by repeating that step.
Let $m'$ and $m''$ be the infima of $f$ on $[x_{i-1},c]$ and on $[c,x_i]$. Both
intervals are contained in $[x_{i-1},x_i]$, so that $m'\ge m_i$ and $m''\ge m_i$.
Consequently
\[
m'(c-x_{i-1})+m''(x_i-c)\ \ge\ m_i(c-x_{i-1})+m_i(x_i-c)=m_i\,\Delta x_i ,
\]
so the lower sum does not decrease. The same argument with suprema gives $M'\le M_i$
and $M''\le M_i$, so that the upper sum does not increase.
\end{proof}

The second fact is stronger: two different partitions can still be compared, as long as what
is compared is a lower sum with an upper sum.

\begin{lemma}
For any two partitions $\mathcal{P}_1$ and $\mathcal{P}_2$ of $[a,b]$ we have
\[
L(\mathcal{P}_1,f)\le U(\mathcal{P}_2,f).
\]
\end{lemma}

\begin{proof}
Take $\mathcal{P}$ the common refinement, that is, the union of the points of the two partitions. By
the previous lemma, $L(\mathcal{P}_1,f)\le L(\mathcal{P},f)$ and
$U(\mathcal{P},f)\le U(\mathcal{P}_2,f)$. Since $L(\mathcal{P},f)\le U(\mathcal{P},f)$
always holds, that chain of inequalities gives the result asked for.
\end{proof}

That lemma ensures that every lower sum is a lower bound for every upper
sum. Hence the set of all lower sums is bounded above, and the set of all upper sums
is bounded below; the following two numbers always exist.

\begin{definisi}[Riemann integrable]\index{Riemann integrable}
For $f$ bounded on $[a,b]$, the \emph{lower integral} and the \emph{upper integral}
are defined as
\[
\underline{I}(f)=\sup_{\mathcal{P}}L(\mathcal{P},f),
\qquad
\overline{I}(f)=\inf_{\mathcal{P}}U(\mathcal{P},f),
\]
with the supremum and infimum taken over all partitions of $[a,b]$. Since
$\underline{I}(f)\le\overline{I}(f)$, the function $f$ is called \emph{Riemann integrable}
on $[a,b]$ when those two numbers are equal. Their common value is written
$\int_a^b f$ or $\int_a^b f(x)\,dx$.
\end{definisi}

Notice what is absent from that definition: there is no sample point and no
limit over $n$. The measure of success of an approximation is the difference between its two brackets,
not its closeness to a guess.

\kamus{partition}{partisi}{A way of cutting an interval into several slices, like
dividing a bar of chocolate with a knife, not measuring exactly how long each slice is.}

\kamus{upper sum}{jumlah atas}{An estimate that is deliberately inflated, like
weighing goods together with their box.}

\begin{bacakeras}
Two voices. Call the first \emph{Cutter}, the second \emph{Adder}.

Cutter: I have a strip of ribbon. I cut it into four.
Adder: I add up those four pieces. The result is too small.
Cutter: I cut again. So, eight.
Adder: The result goes up a little. Still too small.
Cutter: Sixteen.
Adder: Up again. Closer, not there yet.
Cutter: A hundred.
Adder: Very close.
Cutter: A thousand.
Adder: Closer still, still not there.

Cutter: When do we get there?
Adder: Never, if you only add slices. Stop chasing the number,
start chasing the bracket.
Cutter: What bracket?
Adder: Thin slices give a lower bound. Fat slices give an upper
bound. If those two bounds meet at one point, that point is the answer.
Cutter: So my slices are never enough?
Adder: Never. Slices only bring the two bounds close. What answers
is not the last slice, but the meeting of the two bounds.

Now repeat one line: cut thin, bracket tight, bounds meet.
\end{bacakeras}

\section{The Integrability Criterion}
\tingkat{intermediate}

The definition just given speaks of suprema and infima over all partitions, and there are infinitely
many partitions. The following criterion replaces them with a single inequality that
can be checked on one partition alone.

\begin{teorema}[Integrability criterion]\index{integrability criterion}
A bounded function $f$ is Riemann integrable on $[a,b]$ if and only if for
every $\varepsilon>0$ there is a partition $\mathcal{P}$ with
\[
U(\mathcal{P},f)-L(\mathcal{P},f)<\varepsilon .
\]
\end{teorema}

\begin{proof}
For every partition,
$L(\mathcal{P},f)\le\underline{I}(f)\le\overline{I}(f)\le U(\mathcal{P},f)$; the first inequality
comes from the definition of supremum, the last from the definition of infimum.

$(\Rightarrow)$ Suppose $f$ is integrable with value $I$, and take
$\varepsilon>0$. By the definition of supremum, there is $\mathcal{P}_1$ with
$L(\mathcal{P}_1,f)>I-\varepsilon/2$. By the definition of infimum, there is $\mathcal{P}_2$
with $U(\mathcal{P}_2,f)<I+\varepsilon/2$. Take $\mathcal{P}$ the common refinement
of both. The refinement lemma gives
\[
U(\mathcal{P},f)-L(\mathcal{P},f)
<\Bigl(I+\tfrac{\varepsilon}{2}\Bigr)-\Bigl(I-\tfrac{\varepsilon}{2}\Bigr)
=\varepsilon .
\]

$(\Leftarrow)$ Suppose that condition is met. For every $\varepsilon>0$ there is
$\mathcal{P}$ with $U(\mathcal{P},f)-L(\mathcal{P},f)<\varepsilon$, so that
\[
0\le\overline{I}(f)-\underline{I}(f)\le U(\mathcal{P},f)-L(\mathcal{P},f)<\varepsilon .
\]
The number $\overline{I}(f)-\underline{I}(f)$ is nonnegative and smaller than every
$\varepsilon>0$, so that number is zero. The two integrals are equal, and $f$ is integrable.
\end{proof}

The criterion turns a question about infinitely many partitions into a single
inequality. A direct consequence: if the difference is small on one partition, that difference
stays small on every refinement of it, since the refinement lemma can only
narrow it.

Before looking for classes of functions that pass the criterion, we record the properties of the integral
that will be used throughout. All of them are consequences of the definition of upper and lower sums, not new
assumptions.

\begin{proposisi}[Basic properties of the integral]\index{integral, basic properties}
Let $f$ and $g$ be Riemann integrable on $[a,b]$. Then:
(a) $f+g$ is integrable and $\int_a^b(f+g)=\int_a^b f+\int_a^b g$;
(b) if $f\le g$ on $[a,b]$ then $\int_a^b f\le\int_a^b g$;
(c) if $a<c<b$ then $f$ is integrable on $[a,c]$ and on $[c,b]$, and
$\int_a^b f=\int_a^c f+\int_c^b f$;
(d) $|f|$ is integrable and $\bigl|\int_a^b f\bigr|\le\int_a^b|f|$;
(e) for every constant $k$, $\int_a^b k\,dx=k(b-a)$.
\end{proposisi}

\begin{proof}
(a) On every slice, $\sup(f+g)\le\sup f+\sup g$ and $\inf(f+g)\ge\inf f+\inf g$.
Hence $U(\mathcal{P},f+g)\le U(\mathcal{P},f)+U(\mathcal{P},g)$ and
$L(\mathcal{P},f+g)\ge L(\mathcal{P},f)+L(\mathcal{P},g)$. Take $\varepsilon>0$ and
choose $\mathcal{P}$ that makes $U-L$ less than $\varepsilon/2$ for $f$ and $g$
at once. Then
\[
U(\mathcal{P},f+g)-L(\mathcal{P},f+g)
\le\bigl[U(\mathcal{P},f)-L(\mathcal{P},f)\bigr]
+\bigl[U(\mathcal{P},g)-L(\mathcal{P},g)\bigr]<\varepsilon ,
\]
so that $f+g$ is integrable. Its value is obtained by bracketing from both sides, since
$\varepsilon$ is arbitrary.

(b) If $f\le g$, then $M_i(f)\le M_i(g)$ on every slice, so that
$U(\mathcal{P},f)\le U(\mathcal{P},g)$ for all $\mathcal{P}$. Take the infimum of both
sides.

(c) If $\mathcal{P}$ contains $c$, the sum splits into the parts $[a,c]$ and
$[c,b]$, with $L(\mathcal{P},f)=L(\mathcal{P}_1,f)+L(\mathcal{P}_2,f)$ and the same
for $U$. The supremum gives
$\underline{I}(f)\le\underline{I}_{[a,c]}(f)+\underline{I}_{[c,b]}(f)$, the infimum
gives the opposite direction. For a partition without $c$, add $c$ first; the
refinement lemma preserves the inequality.

(d) Since $-|f|\le f\le|f|$, parts (b) and (a) give
$-\int|f|\le\int f\le\int|f|$. The integrability of $|f|$ follows from
$\bigl||f(x)|-|f(y)|\bigr|\le|f(x)-f(y)|$, which gives
$M_i(|f|)-m_i(|f|)\le M_i(f)-m_i(f)$.

(e) On every slice, $M_i=m_i=k$, so $U=L=k(b-a)$ for all partitions.
\end{proof}

Part (c) is the reason we may cut an interval at any point; without it the accumulation
function could not be defined cleanly.

\section{Classes of Integrable Functions}
\tingkat{intermediate}

The criterion in the previous section is a test. Now we look for classes of functions that
pass that test without having to find a special partition for each case. The two most
useful classes are continuous functions and monotone functions.

\begin{akibat}[Continuity guarantees integrability]
If $f$ is continuous on $[a,b]$, then $f$ is Riemann integrable on $[a,b]$.
\end{akibat}

\begin{proof}
Take $\varepsilon>0$. Since $[a,b]$ is compact, $f$ is uniformly continuous there (Chapter 11):
there is $\delta>0$ such that $|f(x)-f(t)|<\varepsilon/(b-a)$ when $|x-t|<\delta$. Choose
a partition with mesh less than $\delta$. On every slice, for all $x,t$ in
that slice we have $|f(x)-f(t)|<\varepsilon/(b-a)$, so that
$M_i-m_i\le\varepsilon/(b-a)$. Then
\[
U(\mathcal{P},f)-L(\mathcal{P},f)
=\sum_{i=1}^{n}(M_i-m_i)\,\Delta x_i
\le\frac{\varepsilon}{b-a}\sum_{i=1}^{n}\Delta x_i=\varepsilon .
\]
The integrability criterion closes the proof.
\end{proof}

Note the role of uniform continuity: one $\delta$ must work for all slices
at once, whereas ordinary continuity only gives a $\delta$ that depends on the
point. This is where Chapter 11 is used, not as decoration.

The second class does not require continuity: a monotone function may jump, as long as
its jumps are finite in number.

\begin{teorema}[Monotonicity guarantees integrability]
If $f$ is monotone on $[a,b]$, then $f$ is Riemann integrable on $[a,b]$.
\end{teorema}

\begin{proof}
Suppose $f$ is increasing. Divide $[a,b]$ into $n$ slices of equal width, $\Delta x=(b-a)/n$.
On the $i$-th slice, the smallest value is at the left endpoint and the largest value at the right
endpoint, so that $M_i-m_i=f(x_i)-f(x_{i-1})$. Then
\[
U(\mathcal{P},f)-L(\mathcal{P},f)
=\Delta x\sum_{i=1}^{n}\bigl(f(x_i)-f(x_{i-1})\bigr)
=\frac{(b-a)\bigl(f(b)-f(a)\bigr)}{n},
\]
because the sum in the middle telescopes. Choose $n>(b-a)(f(b)-f(a))/\varepsilon$ and the criterion
is satisfied. If $f$ is decreasing, apply this result to $-f$ and use linearity.
\end{proof}

That proof uses no $\delta$: all it uses is that the sum of all the jumps does not
exceed $f(b)-f(a)$, so a monotone function is easy to integrate even though its graph is full of
corners.

Those two classes can be combined: a function that is continuous except at finitely many points
is still integrable (Rudin, Theorem 6.10, not repeated here).

Now the other side: bounded functions that fail. The cleanest example is the
Dirichlet function, equal to $1$ at rational numbers and to $0$ at irrational numbers.

\begin{contoh}[The Dirichlet function is not integrable]\index{Dirichlet function}
Let $D(x)=1$ if $x$ is rational and $D(x)=0$ if $x$ is irrational, on $[0,1]$.
Every slice contains both (Chapter 3, density of $\mathbb{Q}$), so that $M_i=1$ and
$m_i=0$ for every $i$, and
\[
U(\mathcal{P},D)=\sum_{i=1}^{n}1\cdot\Delta x_i=1,
\qquad
L(\mathcal{P},D)=0
\]
for \emph{every} partition. The difference is always $1$, never smaller than
$\varepsilon=1/2$, so by the criterion $D$ is not integrable even though it is bounded.
\end{contoh}

That failure is not because $D$ is broken at many points, but because
its jumps are spread everywhere. The following function shows that the number of
discontinuity points is not an obstacle.

\begin{teorema}[Thomae's function]\index{Thomae function}
Define $T$ on $[0,1]$ by $T(x)=1/q$ if $x=p/q$ in lowest
terms with $q\ge1$, and $T(x)=0$ if $x$ is irrational. Then $T$ is integrable in the
Riemann sense and $\int_0^1 T=0$.
\end{teorema}

\begin{proof}
Take $\varepsilon>0$. The set $S=\{x:T(x)\ge\varepsilon\}$ consists of fractions with
denominator $q\le1/\varepsilon$; for each $q$ there are only $q+1$ such fractions in
$[0,1]$, so that $S$ is finite. Cover $S$ with finitely many open intervals
whose total length is less than $\varepsilon$.

Form a partition that contains all endpoints of those covering intervals. The slices that contain $S$
have total length less than $\varepsilon$ and there $M_i-m_i\le1$, so
their contribution is less than $\varepsilon$. The other slices satisfy $M_i\le\varepsilon$ and
$m_i\ge0$, so their contribution is at most $\varepsilon$. Hence
$U(\mathcal{P},T)-L(\mathcal{P},T)<2\varepsilon$, and the criterion closes the proof. The value is
zero because every slice contains an irrational number, so that $L(\mathcal{P},T)=0$.
\end{proof}

So $T$ is broken at every rational number, which is dense in $[0,1]$, yet $T$ is still
integrable. What decides is not the number of discontinuity points but the amount of space
they occupy: Lebesgue's criterion states that a bounded function is Riemann integrable
if and only if its discontinuity points have measure zero (Rudin, Theorem 11.33(b); Abbott,
Chapter 7). We quote it without proof because measure demands measure theory.

\begin{tantangan}
Define $h$ on $[0,1]$ by $h(x)=1$ if $x=p/q$ in lowest terms
and the denominator $q$ is odd, and $h(x)=0$ for all other $x$, including every irrational
number. Is $h$ Riemann integrable on $[0,1]$? If so, what is the value of
the integral? If not, prove it using the integrability criterion, not by
counting points.

Hint 1: do not count the number of points where $h$ equals one. The number is
infinite, and the criterion does not ask for that count.

Hint 2: take any slice $[x_{i-1},x_i]$ with $x_{i-1}<x_i$. That slice
contains fractions with odd denominator, so $M_i=1$, and it contains irrational numbers, so
$m_i=0$. Thus $U-L=1$ for every partition.

If you get stuck, go back to the section Classes of Integrable Functions, the Dirichlet example.
\end{tantangan}

\kamus{Riemann integrability}{keterintegralan}{The property of a function whose lower and
upper estimates can be brought together until their difference vanishes.}

\begin{untukanak}
Have you ever looked at a car's speedometer and then asked how much distance has been
covered? Distance is not displayed there, only instantaneous speed. But if we
record the speed every second, we can add it up: one second at
ten meters per second means 10 meters, the next second 12 meters, and
so on. The sum of all those slices is the distance.

The problem is that the speed changes in the middle of a second. A car can speed up and then slow down
within the same second. If we use the speed at the start of the second, the distance we
compute is too small. If we use the speed at the end of the second, the result is too
large. The right value lies between them.

If we record ten times per second, the two estimates grow closer; a hundred times,
closer still. Mathematicians ask further: if the recording can be as dense as we
like, do the two estimates head toward one and the same number? The answer is not always yes,
and in this chapter we shall learn when it is yes and why.
\end{untukanak}

\section{The Integral as an Accumulation Function}
\tingkat{intermediate}

Up to here the integral is only a number attached to an interval. We can
view it as a function by letting the upper endpoint move: the integral becomes
a report on how much has accumulated up to that point.

\begin{definisi}[Accumulation function]\index{accumulation function}
Let $f$ be integrable on $[a,b]$. For $x\in[a,b]$ define
\[
F(x)=\int_a^x f(t)\,dt .
\]
For $x<a$ we use the convention $\int_a^x f=-\int_x^a f$, and for $x=a$ the value is
zero. The function $F$ is called the \emph{accumulation function} of $f$ on $[a,b]$.
\end{definisi}

With that convention the additive property holds for any ordering of points, so that
\begin{equation}
F(y)-F(x)=\int_x^{y} f(t)\,dt
\end{equation}
for all $x,y\in[a,b]$. This is the form used in all the proofs that follow:
the increase of savings between two times is the integral of the income rate in between.

The accumulation function is always better than the function it comes from: it is continuous even when $f$
is not, and differentiable at every point where $f$ is continuous.

\begin{proposisi}[Accumulation is always continuous]
If $f$ is bounded on $[a,b]$ with $|f|\le M$ and integrable, then the $F$ that
is defined above satisfies
\[
|F(y)-F(x)|\le M\,|y-x|
\]
for all $x,y\in[a,b]$. In particular $F$ is continuous on $[a,b]$.
\end{proposisi}

\begin{proof}
For $x<y$, monotonicity and part (d) of the basic properties proposition give
\[
|F(y)-F(x)|=\Bigl|\int_x^{y}f\Bigr|\le\int_x^{y}|f|\le M(y-x).
\]
For $y<x$ the same result with the opposite sign. Hence $F$ satisfies a
Lipschitz condition with constant $M$, and every Lipschitz function is continuous.
\end{proof}

Now the first direction of the derivative-integral relationship. In everyday words: the rate of
growth of savings at a given moment equals the income rate at that moment.

\begin{teorema}[Fundamental theorem of calculus, first part]\index{fundamental theorem of calculus}
Let $f$ be integrable on $[a,b]$ and continuous at the point $x_0\in(a,b)$. Then the
accumulation function $F$ is differentiable at $x_0$ and
\[
F'(x_0)=f(x_0).
\]
\end{teorema}

\begin{proof}
Take $h\ne0$ with $x_0+h\in[a,b]$. From the increment form,
\[
\frac{F(x_0+h)-F(x_0)}{h}-f(x_0)
=\frac{1}{h}\int_{x_0}^{x_0+h}\bigl(f(t)-f(x_0)\bigr)\,dt ,
\]
because the integral of the constant function $f(x_0)$ over an interval of length $|h|$ equals
$f(x_0)h$.

Take $\varepsilon>0$. Since $f$ is continuous at $x_0$, there is $\delta>0$ such that
$|t-x_0|<\delta$ implies $|f(t)-f(x_0)|<\varepsilon$. For $0<|h|<\delta$ the whole
interval of integration lies inside $(x_0-\delta,x_0+\delta)$, so that, by monotonicity
and part (d) of the basic properties proposition,
\[
\left|\frac{1}{h}\int_{x_0}^{x_0+h}\bigl(f(t)-f(x_0)\bigr)dt\right|
\le\frac{1}{|h|}\left|\int_{x_0}^{x_0+h}\bigl|f(t)-f(x_0)\bigr|\,dt\right|
\le\frac{1}{|h|}\,\varepsilon\,|h|=\varepsilon .
\]
Since $\varepsilon$ is arbitrary, the difference quotient tends to $f(x_0)$.
\end{proof}

That theorem demands continuity at only one point. Where $f$ jumps, the derivative of
$F$ may fail to exist, and that is not a weakness of the proof: for $f(t)=\operatorname{sgn} t$
on $[-1,1]$ we have $F(x)=\int_0^x\operatorname{sgn} t\,dt=|x|$, which is not
differentiable at $0$, exactly where the sign of $f$ changes.

One important consequence: every continuous function has an antiderivative, and its antiderivative can be
written as an integral. Chapter 12 derived the differentiation rules but did not answer whether
every continuous function is the derivative of something. Now that question is answered,
and the answer is constructive.

\begin{akibat}[An antiderivative always exists]
If $f$ is continuous on $[a,b]$, then $F(x)=\int_a^x f$ is differentiable on $[a,b]$
with $F'=f$. Hence $F$ is an antiderivative of $f$ on $[a,b]$.
\end{akibat}

\begin{proof}
The continuity of $f$ gives integrability (the previous corollary) and continuity at every
point, so that the first part of the theorem applies at every $x_0\in(a,b)$. At the endpoints
one-sided derivatives are used.
\end{proof}

The chain rule gives one further flexibility: if the upper limit depends on $x$,
\[
\frac{d}{dx}\int_a^{g(x)}f(t)\,dt=f\bigl(g(x)\bigr)\,g'(x),
\]
provided $g$ is differentiable. For example,
$\frac{d}{dx}\int_0^{x^2}e^{-t^2}dt=2x\,e^{-x^4}$.

The two directions are easy to confuse, so it is worth drawing them once.
Figure~\ref{fig:tdk-dua-arah} places both in one diagram: the upper arrow
accumulates thin slices, the lower arrow differentiates. What the theorem guarantees is that
differentiating the result of accumulating returns the original function, not a tidy algebraic inverse,
since an antiderivative is fixed only up to a constant.

\begin{figure}[htbp]
\centering
\begin{tikzpicture}[>=Stealth,every node/.style={font=\footnotesize},
  kotak/.style={draw,rounded corners=3pt,align=center,inner sep=5pt,
    minimum height=10mm,text width=3.3cm}]
\node[kotak,fill=layerindigo!8,draw=layerindigo] (f) at (0,0) {function $f$};
\node[kotak,fill=layerbrown!8,draw=layerbrown] (F) at (7.6,0)
  {accumulation function $F$};
\draw[->,line width=0.9pt,draw=layerindigo] ([yshift=2mm]f.north) to[bend left=24]
  node[above,text=layerindigo,align=center]
  {definite integral: accumulate thin slices} ([yshift=2mm]F.north);
\draw[->,line width=0.9pt,draw=layerbrown] ([yshift=-2mm]F.south) to[bend left=24]
  node[below,text=layerbrown,align=center]
  {derivative: $F'=f$, this is part 1} ([yshift=-2mm]f.south);
\node[draw=layergrey,dashed,rounded corners=3pt,align=center,inner sep=4pt]
  at (3.8,-2.7)
  {part 2 closes the loop: $\int_a^b F'(x)\,dx=F(b)-F(a)$};
\end{tikzpicture}
\caption{The two directions of the fundamental theorem of calculus as back-and-forth arrows. Upper arrow:
accumulating thin slices produces $F(x)=\int_a^x f$. Lower arrow: differentiating $F$
returns $f$, and that is the first part. The dashed box below is the second
part, which uses an antiderivative to compute integrals. For $f(x)=x^2$ on
$[0,1]$, $F(x)=x^3/3$, and the second part gives $F(1)-F(0)=1/3$, the same number as
the bracketing in Figure~\ref{fig:rect-halus}.}
\label{fig:tdk-dua-arah}
\end{figure}

\section{The Fundamental Theorem of Calculus, Second Part}
\tingkat{advanced}

The first part runs from the integral to the derivative. The second part runs the other way:
if we already hold an antiderivative, the integral can be computed with one subtraction.
In everyday words: total change equals the final value minus the initial value.

\begin{teorema}[Fundamental theorem of calculus, second part]
Let $F$ be continuous on $[a,b]$, differentiable on $(a,b)$, and let $F'$ be
Riemann integrable on $[a,b]$. Then
\[
\int_a^b F'(x)\,dx=F(b)-F(a).
\]
\end{teorema}

\begin{proof}
Take a partition $\mathcal{P}=\{x_0,\dots,x_n\}$ of $[a,b]$. On every subinterval, $F$ is
continuous on $[x_{i-1},x_i]$ and differentiable inside it, so the mean value
theorem (Chapter 12) gives $t_i\in(x_{i-1},x_i)$ with
\[
F(x_i)-F(x_{i-1})=F'(t_i)\,\Delta x_i .
\]
Since $m_i\le F'(t_i)\le M_i$, summing over $i$ gives
\[
L(\mathcal{P},F')\le\sum_{i=1}^{n}\bigl(F(x_i)-F(x_{i-1})\bigr)\le U(\mathcal{P},F').
\]
The sum in the middle telescopes and equals $F(b)-F(a)$, so that number is trapped
between the lower sum and the upper sum of every partition. Since $F'$ is integrable, both
brackets head toward the same value, namely $\int_a^b F'$.
\end{proof}

\begin{resep}
\begin{enumerate}
\item Find $F$ with $F'=f$, then check by differentiating $F$. If the derivative does not
      return to $f$, the candidate is wrong and the following steps are wasted.
\item Write $F$ at the upper limit, write $F$ at the lower limit, then subtract:
      $F(b)-F(a)$. Do not insert $+C$; the constant cancels.
\item Fix the sign if the interval is reversed: $\int_b^a f=-\int_a^b f$, and
      $\int_a^a f=0$.
\item Test with one coarse slice, a trapezoid or \texttt{quad}. If the result is far off,
      the error is in step 1 or 2, not in step 3.
\end{enumerate}
Warning about order: checking the antiderivative must happen before subtracting. Two values of an
arbitrary function still give a tidy and plausible difference, so that a mistake in
step 1 is not visible from the shape of the answer.
\end{resep}

Likewise the two parts do not replace each other: the first part differentiates the
accumulation function without requiring an antiderivative, whereas the second part requires an
antiderivative and integrability of its derivative.

The second part explains why the result is called net change, not area. For
$F(x)=\frac{x^4}{4}-\frac{x^2}{2}$ on $[-1,1]$, $F'(x)=x^3-x$, so that
\[
\int_{-1}^{1}(x^3-x)\,dx=F(1)-F(-1)=\Bigl(\frac14-\frac12\Bigr)-\Bigl(\frac14-\frac12\Bigr)=0 .
\]
The region above the axis has area $1/4$, the region below the axis also $1/4$, so the total area is
$1/2$ whereas the integral is zero because the signs are opposite.
Figure~\ref{fig:luas-bertanda} shows both.

\begin{figure}[htbp]
\centering
\begin{tikzpicture}
\begin{axis}[width=0.74\linewidth,height=5.4cm,
  axis lines=middle,xmin=-1.3,xmax=1.3,ymin=-0.6,ymax=0.6,
  xtick={-1,0,1},ytick=\empty,
  tick label style={font=\scriptsize},
  xlabel={$x$},ylabel={$y=x^3-x$},
  xlabel style={font=\scriptsize},ylabel style={font=\scriptsize,rotate=-90}]
\addplot[fill=layerjade!22,draw=layerjade!70,domain=-1:0,samples=60]
  {x^3-x} \closedcycle;
\addplot[fill=layerplum!22,draw=layerplum!70,domain=0:1,samples=60]
  {x^3-x} \closedcycle;
\addplot[thick,ink,domain=-1.2:1.2,samples=120] {x^3-x};
\node[font=\scriptsize,text=layerjade!55!black] at (axis cs:-0.52,0.36)
  {$+\frac14$};
\node[font=\scriptsize,text=layerplum!55!black] at (axis cs:0.52,-0.36)
  {$-\frac14$};
\end{axis}
\end{tikzpicture}
\caption{Signed area for $F'(x)=x^3-x$ on $[-1,1]$. The region above the $x$-axis
contributes $+1/4$, the region below the axis contributes $-1/4$, so that
$\int_{-1}^{1}(x^3-x)\,dx=F(1)-F(-1)=0$. The geometric area is still $1/2$; what is zero is
the signed sum, and that is what the second part of the theorem computes.}
\label{fig:luas-bertanda}
\end{figure}

The integrability hypothesis in the second part is not excessive. There are
differentiable functions whose derivative is bounded but not Riemann integrable; the classical
example was found by Volterra in the nineteenth century, and we do not
prove it here. The second part of the theorem is not an unconditional statement.

\section{Substitution and Integration by Parts}
\tingkat{advanced}

The two most frequently used rules for computing integrals are not new axioms: both are
differentiation rules read backwards, and their proofs use only the fundamental theorem of calculus.
The first one reverses the chain rule.

\begin{proposisi}[Substitution]\index{substitution}
Let $\varphi$ be differentiable with $\varphi'$ continuous on $[a,b]$, and let $f$ be
continuous on an interval that contains $\varphi([a,b])$. Then
\[
\int_a^b f\bigl(\varphi(t)\bigr)\,\varphi'(t)\,dt
=\int_{\varphi(a)}^{\varphi(b)}f(u)\,du .
\]
\end{proposisi}

\begin{proof}
Since $f$ is continuous, it has an antiderivative $G(u)=\int_c^u f$, with $c$ inside the
range of $\varphi$; by the first part of the theorem, $G'=f$. The chain rule (Chapter 12)
gives
\[
\frac{d}{dt}G\bigl(\varphi(t)\bigr)=G'\bigl(\varphi(t)\bigr)\varphi'(t)
=f\bigl(\varphi(t)\bigr)\varphi'(t).
\]
The right-hand side is continuous, hence integrable. The second part of the theorem applied to the function
$G\circ\varphi$ gives
\[
\int_a^b f\bigl(\varphi(t)\bigr)\varphi'(t)\,dt
=G\bigl(\varphi(b)\bigr)-G\bigl(\varphi(a)\bigr)
=\int_{\varphi(a)}^{\varphi(b)}f(u)\,du .
\]
\end{proof}

The second rule reverses the product rule.

\begin{proposisi}[Integration by parts]\index{integration by parts}
Let $F$ and $G$ be differentiable on $[a,b]$ with $F'$ and $G'$ continuous. Then
\[
\int_a^b F(x)\,G'(x)\,dx=F(b)G(b)-F(a)G(a)-\int_a^b F'(x)\,G(x)\,dx .
\]
\end{proposisi}

\begin{proof}
Let $H=FG$. The product rule gives $H'=F'G+FG'$, and the right-hand side is continuous
so that $H'$ is integrable. The second part of the theorem applied to $H$ gives
\[
F(b)G(b)-F(a)G(a)=\int_a^b H'=\int_a^b F'G+\int_a^b FG' .
\]
Move the first term to the left-hand side.
\end{proof}

\kamus{antiderivative}{antiturunan}{A function whose derivative is the function we are
integrating, like the way home on a road: it is useful only when the way out is already
known.}

As a numerical check, take $\int_0^1 x\,e^{x^2}dx$. The substitution $u=x^2$,
$du=2x\,dx$ gives
\[
\int_0^1 x\,e^{x^2}dx=\frac12\int_0^1 e^{u}du=\frac{e-1}{2}=0.8591409142295225 .
\]
\texttt{scipy.integrate.quad} gives $0.8591409142295227$, a difference of $2\times10^{-16}$
from machine rounding. For the second rule, $F(x)=x$ and $G'(x)=e^x$ give
$\int_0^1 xe^xdx=e-(e-1)=1$, and \texttt{quad} returns $1.0$ with estimated
error $1.1\times10^{-14}$.

\section{The Mean Value Theorem for Integrals}
\tingkat{intermediate}

An average is usually computed by adding and then dividing by the number of data. For a function
that changes continuously, the sum is replaced by an integral and the divisor by the length of the interval: the average
is the height of the rectangle whose area equals the area under the curve.

\begin{definisi}[Average value]
For $f$ integrable on $[a,b]$, the \emph{average value} of $f$ on that interval
is
\[
\bar{f}=\frac{1}{b-a}\int_a^b f(x)\,dx .
\]
\end{definisi}

\begin{teorema}[Mean value for integrals]\index{mean value of an integral}
If $f$ is continuous on $[a,b]$, then there is $c\in[a,b]$ with
\[
f(c)=\frac{1}{b-a}\int_a^b f(x)\,dx .
\]
\end{teorema}

\begin{proof}
Since $f$ is continuous on a closed bounded interval, it attains its smallest value $m$ and
its largest value $M$ there (Chapter 11). Monotonicity gives
\[
m(b-a)\le\int_a^b f\le M(b-a),
\]
so that $\bar{f}\in[m,M]$. By the intermediate value theorem (Chapter 11), the continuous function $f$
takes every value between $m$ and $M$, so there is $c\in[a,b]$ with
$f(c)=\bar{f}$.
\end{proof}

As a number, for $f(x)=x^2$ on $[0,1]$ the average is $1/3$, and the point that the
theorem reaches is $c=1/\sqrt3=0.5773502691896258$, since
$c^2=0.3333333333333334$. That point is not the midpoint of the interval: the theorem guarantees
existence, not location.

\section{Worked Examples}
\tingkat{intermediate}

Now we work through three things: the upper and lower sums of $f(x)=x^2$ on $[0,1]$,
a comparison of five quadrature rules, and one integral without an elementary antiderivative.
Fractions are computed with \texttt{fractions.Fraction}, without rounding.

Start from $n$ slices of equal width with $x_i=i/n$ and $\Delta x=1/n$, so that
\[
U_n=\sum_{i=1}^{n}\Bigl(\frac{i}{n}\Bigr)^2\frac1n=\frac{1}{n^3}\sum_{i=1}^{n}i^2
=\frac{(n+1)(2n+1)}{6n^2}=\frac13+\frac{1}{2n}+\frac{1}{6n^2},
\]
using the sum-of-squares formula from Chapter 1. In the same way
$L_n=\frac13-\frac{1}{2n}+\frac{1}{6n^2}$, so that
\begin{equation}
U_n-L_n=\frac1n .
\end{equation}
That result is not an approximation. For every $n$ the difference is exactly $1/n$; since $1/n$ can
be made smaller than any $\varepsilon$ by taking $n>1/\varepsilon$, the integrability
criterion is satisfied. Hence the following table is at the same time a proof that
$\int_0^1x^2dx=1/3$: the two sums bracket $1/3$, and their distance tends to zero.

\begin{center}\small
\begin{tabular}{rrrr}
\hline
$n$ & $U_n$ & $L_n$ & $U_n-L_n$ \\
\hline
1 & $1.0000000000$ & $0.0000000000$ & $1$ \\
4 & $0.4687500000$ & $0.2187500000$ & $1/4$ \\
16 & $0.3652343750$ & $0.3027343750$ & $1/16$ \\
64 & $0.3411865234$ & $0.3255615234$ & $1/64$ \\
256 & $0.3352890015$ & $0.3313827515$ & $1/256$ \\
1024 & $0.3338217735$ & $0.3328452110$ & $1/1024$ \\
\hline
\end{tabular}
\end{center}

At $n=1024$ the difference is $9.77\times10^{-4}$ and the distance from $U_n$ to $1/3$ is half of that.
Each time $n$ is doubled the difference is exactly halved, and that is why upper and lower
sums are not an efficient way to reach high accuracy.

\begin{tangan}
Materials: one sheet of graph paper, scissors, a ruler, tape, a pencil.

Steps. (1) On the graph paper draw a curve that rises faster and faster: one square from the
left edge has height one square, two squares have height four squares, three squares have height
nine squares. (2) Cut the region under the curve into vertical strips one square wide.
(3) Lay those strips side by side on the table, from shortest to longest, with level
bases. (4) Look at their upper edge: what appears is a jagged staircase, not a smooth line.
(5) Repeat with strips half a square wide, then a quarter square wide.

What we should see: the serration on the upper edge grows finer as the strips grow thinner,
whereas the total area does not change, since cutting and rearranging neither adds nor removes
paper. Hence the area of that curved region equals the area of the jagged figure, and that jagged
figure increasingly resembles one rectangle. The height of that rectangle is what we shall later call
the average value of the function on that interval.
\end{tangan}

Now compare five quadrature rules on the same integral: left endpoint, right
endpoint, midpoint $x_i=(i-\frac12)/n$, trapezoid, and Simpson with weights
$1,4,2,4,\dots,4,1$ divided by $3n$. The error is computed against the exact value $1/3$.

\begin{center}\small
\begin{tabular}{rrrrrr}
\hline
$n$ & left & right & midpoint & trapezoid & Simpson \\
\hline
10 & $4.833\times10^{-2}$ & $5.167\times10^{-2}$ & $8.333\times10^{-4}$
   & $1.667\times10^{-3}$ & $1.1\times10^{-16}$ \\
100 & $4.983\times10^{-3}$ & $5.017\times10^{-3}$ & $8.333\times10^{-6}$
   & $1.667\times10^{-5}$ & $0$ \\
1000 & $4.998\times10^{-4}$ & $5.002\times10^{-4}$ & $8.333\times10^{-8}$
   & $1.667\times10^{-7}$ & $0$ \\
\hline
\end{tabular}
\end{center}

Three patterns stand out. The error of the left and right rules drops by a factor of ten each time $n$ grows
by a factor of ten: order $\mathcal{O}(1/n)$, about $1/(2n)$ because $f$ is increasing. The error of the midpoint
and trapezoid rules drops by a factor of a hundred, order $\mathcal{O}(1/n^2)$, with values
$1/(12n^2)$ and $1/(6n^2)$; the two have opposite signs and the trapezoid is twice as large,
so that $(2M+T)/3$ cancels the term. Simpson is exact to machine rounding, since the
rule is exact for polynomials of degree at most three.

The third integral has no elementary antiderivative:
\[
\int_0^1 e^{-x^2}dx .
\]
That absence is a consequence of Liouville's theorem, which we quote without proof. The reference
is the error function $\frac{\sqrt{\pi}}{2}\operatorname{erf}(1)=0.746824132812427$; \texttt{quad} gives
the same number with estimated error $8.3\times10^{-15}$. Simpson's rule gives an error of
$3.563\times10^{-4}$ at $n=2$, $1.988\times10^{-6}$ at $n=8$, and
$7.795\times10^{-9}$ at $n=32$.

\kertas{For $f(x)=x^2$ on $[0,1]$, derive the formula
$U_n-L_n=1/n$ yourself from the sum-of-squares formula, then compute $U_{10}$ and $L_{10}$ as
fractions. Check that their average $(U_{10}+L_{10})/2$ equals
$1/3+1/(6\cdot100)$, not $1/3$. Explain why that average is still not the best
estimate, even though its error is far smaller than the error of $U_{10}$ itself.}

\section{Improper Integrals and Their Relation to Series}
\tingkat{advanced}

The interval of integration has so far always been bounded and the function bounded. Many real
computations are not like that: we want the area of a region that extends without end, or the integral
of a function that blows up at the endpoint of the interval. The method is not to widen the Riemann
definition, but to write that integral as a limit of ordinary integrals.

\begin{definisi}[Improper integral]\index{improper integral}
If $f$ is Riemann integrable on $[a,R]$ for every $R>a$, define
\[
\int_a^\infty f=\lim_{R\to\infty}\int_a^R f ,
\]
provided that limit exists and is finite. If $f$ is integrable on $[c,b]$ for every
$c\in(a,b]$ but is unbounded near $a$, define
\[
\int_a^b f=\lim_{c\to a^{+}}\int_c^b f .
\]
In both cases, the integral is called \emph{convergent} if the limit exists, and
\emph{divergent} if not.
\end{definisi}

The most important example is the family $f(x)=x^{-p}$ on $[1,\infty)$. For $p\ne1$,
\[
\int_1^R x^{-p}dx=\frac{R^{1-p}-1}{1-p},
\]
so that the limit exists only when $p>1$, with value $1/(p-1)$; for $p=1$ the result is
$\ln R$, which tends to infinity. Hence $\int_1^\infty x^{-2}dx=1$, and the partial
integrals head toward that value at rate $1/R$: $0.9$, $0.99$, $0.999$ for
$R=10,100,1000$.

The link with the series of Chapter 5 is no coincidence: for $f$ positive and decreasing, the areas of
rectangles of height $f(n)$ bracket the area under the curve, so that
\[
\sum_{n=N+1}^{\infty}f(n)\ \le\ \int_N^\infty f(x)\,dx\ \le\ \sum_{n=N}^{\infty}f(n) .
\]
Hence the integral test holds: $\sum f(n)$ converges if and only if
$\int_1^\infty f$ converges. For $f(x)=1/x^2$ both converge, with
$\sum_{n=1}^{\infty}1/n^2=\pi^2/6=1.6449340668482264$. The left inequality also gives
an error bound, and the larger $N$ is the closer the error comes to that bound:

\begin{center}\small
\begin{tabular}{rrrr}
\hline
$N$ & error of the partial sum & $1/N$ \\
\hline
10 & $9.516634\times10^{-2}$ & $0.1$ \\
100 & $9.950167\times10^{-3}$ & $0.01$ \\
1000 & $9.995002\times10^{-4}$ & $0.001$ \\
\hline
\end{tabular}
\end{center}

It must be stressed: an improper integral is not a Riemann integral, but its limit. Hence
substitution and integration by parts apply only when the limit exists, and for functions
that change sign their convergence need not be absolute.

\section{Computational Lab}
\tingkat{advanced}

The following code is run with Python 3, \texttt{scipy} 1.18.0, and the module
\texttt{fractions} for exact arithmetic. All partitions are of equal width on $[0,1]$.

\begin{lab}
\begin{verbatim}
import math
from fractions import Fraction as F
from scipy.integrate import quad

def both_sums(n):
    U = sum(F(i * i, n ** 3) for i in range(1, n + 1))
    L = sum(F(i * i, n ** 3) for i in range(n))
    return U, L

print("n        U_n          L_n          U_n - L_n")
for n in (1, 2, 4, 8, 16, 32, 64, 128, 256, 512, 1024):
    U, L = both_sums(n)
    print(f"{n:5d}  {float(U):.10f}  {float(L):.10f}  {float(U - L):.2e}")

def rules(f, n):
    h = 1 / n
    left = sum(f(i * h) for i in range(n)) * h
    right = sum(f(i * h) for i in range(1, n + 1)) * h
    mid = sum(f((i + 0.5) * h) for i in range(n)) * h
    trap = (f(0) / 2 + sum(f(i * h) for i in range(1, n)) + f(1) / 2) * h
    simp = (f(0) + f(1) + sum((4 if i % 2 else 2) * f(i * h)
                              for i in range(1, n))) * h / 3
    return left, right, mid, trap, simp

f = lambda x: x * x
print()
print("error |approximation - 1/3|, integral of x^2 on [0,1]")
print("   n      left     right       mid  trapezoid    Simpson")
for n in (10, 100, 1000):
    err = [f"{abs(v - 1 / 3):.3e}" for v in rules(f, n)]
    print(f"{n:4d}  " + "  ".join(err))

g = lambda x: math.exp(-x * x)
v, e = quad(g, 0, 1)
ref = math.sqrt(math.pi) / 2 * math.erf(1)
print()
print(f"quad exp(-x^2) = {v:.15f}   estimated error {e:.1e}")
print(f"erf reference  = {ref:.15f}   difference {v - ref:+.1e}")
print(f"quad x^2       = {quad(f, 0, 1)[0]:.17f}")
print()
print("error |Simpson - reference| for exp(-x^2)")
for n in (2, 4, 8, 16, 32):
    print(f"  n={n:3d}  {abs(rules(g, n)[4] - ref):.3e}")
print()
print("error |trapezoid - reference| for exp(-x^2)")
for n in (10, 100, 1000):
    print(f"  n={n:4d}  {abs(rules(g, n)[3] - ref):.3e}")
\end{verbatim}
\end{lab}

The output obtained, as it stands:

\begin{verbatim}
n        U_n          L_n          U_n - L_n
    1  1.0000000000  0.0000000000  1.00e+00
    2  0.6250000000  0.1250000000  5.00e-01
    4  0.4687500000  0.2187500000  2.50e-01
    8  0.3984375000  0.2734375000  1.25e-01
   16  0.3652343750  0.3027343750  6.25e-02
   32  0.3491210938  0.3178710938  3.12e-02
   64  0.3411865234  0.3255615234  1.56e-02
  128  0.3372497559  0.3294372559  7.81e-03
  256  0.3352890015  0.3313827515  3.91e-03
  512  0.3343105316  0.3323574066  1.95e-03
 1024  0.3338217735  0.3328452110  9.77e-04

error |approximation - 1/3|, integral of x^2 on [0,1]
   n      left     right       mid  trapezoid    Simpson
  10  4.833e-02  5.167e-02  8.333e-04  1.667e-03  1.110e-16
 100  4.983e-03  5.017e-03  8.333e-06  1.667e-05  0.000e+00
1000  4.998e-04  5.002e-04  8.333e-08  1.667e-07  0.000e+00

quad exp(-x^2) = 0.746824132812427   estimated error 8.3e-15
erf reference  = 0.746824132812427   difference +1.1e-16
quad x^2       = 0.33333333333333337

error |Simpson - reference| for exp(-x^2)
  n=  2  3.563e-04
  n=  4  3.125e-05
  n=  8  1.988e-06
  n= 16  1.246e-07
  n= 32  7.795e-09

error |trapezoid - reference| for exp(-x^2)
  n=  10  6.133e-04
  n= 100  6.131e-06
  n=1000  6.131e-08
\end{verbatim}

Four things stand out. The difference halves each time $n$ doubles, consistent with $1/1024$
in the last row. The left and right errors drop by a factor of ten, midpoint and
trapezoid by a factor of a hundred. The Simpson column contains $1.110\times10^{-16}$ and then zero, a residue
of rounding, and Simpson's error for $e^{-x^2}$ drops by about a factor of sixteen
each time $n$ is doubled.

That code does not prove a single theorem. It checks that the formulas of this chapter give
the numbers stated. The proof of integrability of $x^2$ is in the computation $U_n-L_n=1/n$
in the previous section: the output of a program is data, not an argument.

\begin{untukdosen}
The teaching order I suggest: start from the two different answers in the opening, then
the definition of a partition and compute $U$, $L$ for $x^2$ with four slices on the board. After
that the refinement lemma, the criterion, and only then the classes of functions. Do not postpone the fundamental
theorem of calculus: students need to see the hard work of the criterion bear fruit as a computational tool.

Two misconceptions that appear often. First, the integral is taken to be the same as the antiderivative,
so that $F'=f$ is taken to hold by definition, whereas it is a theorem. Second, a function
with many discontinuity points is taken to be certainly not integrable, because the Dirichlet
example is too often used as the only example. Thomae's function removes
that assumption in a single page.

Blackboard activity: ask students to guess $\int_0^1x^2dx$ only from $U_n$ and
$L_n$ for $n=1,2,4$, then to say how many $n$ are needed for the error to be below
$10^{-3}$ without computing the integral. The answer $n>1000$ comes out of $1/n$, and that moment
explains why an antiderivative is more valuable than quadrature. Sign of readiness
to move on: students can explain the difference in direction of the two parts of the fundamental theorem of calculus,
and can name one integrable function that is not continuous.
\end{untukdosen}

\section{Summary}
\tingkat{basic}

\begin{ringkasan}
\begin{itemize}
\item The Riemann integral is the only number trapped between all
      lower sums and all upper sums; when the two brackets meet, that value is called
      the integral.
\item The integrability criterion replaces the examination of infinitely many partitions
      with one inequality $U(\mathcal{P},f)-L(\mathcal{P},f)<\varepsilon$.
\item Continuous functions and monotone functions are integrable; the Dirichlet function is not, because
      every slice of it contains rational and irrational numbers, whereas Thomae's function
      is integrable even though it is broken at every rational number.
\item The accumulation function $F(x)=\int_a^xf$ is always continuous and its derivative is $f$ at every
      point where $f$ is continuous, so every continuous function has an antiderivative.
\item The second part states $\int_a^bF'=F(b)-F(a)$; its proof is the mean value theorem
      plus telescoping, with the condition that $F'$ be integrable.
\item Substitution and integration by parts are the chain rule and the product rule
      read backwards. An improper integral is a limit of Riemann integrals, and the integral
      test links it to the series of Chapter 5.
\end{itemize}
\end{ringkasan}

\section{Exercises}
\tingkat{intermediate}

\begin{latihan}
\begin{enumerate}
  \item Compute $U_n$ and $L_n$ for $f(x)=x$ on $[0,1]$ with a uniform partition,
        express both as fractions in $n$, then determine $U_n-L_n$. Conclude
        that $f$ is integrable and compute the integral.
  \item Prove that $f$ integrable on $[a,b]$ implies $|f|$ integrable
        and $\bigl|\int_a^bf\bigr|\le\int_a^b|f|$. Show that the converse is false with
        one example.
  \item Assess this claim: ``Every bounded function on $[a,b]$ is integrable, since its
        upper sum and lower sum are always finite.'' True, false, or not yet
        enough evidence? Name which part is supported by proof and which part is
        not.
  \item Prove that an increasing function on $[a,b]$ is integrable by showing
        $U(\mathcal{P},f)-L(\mathcal{P},f)\le(b-a)\bigl(f(b)-f(a)\bigr)/n$ for
        a uniform partition into $n$ parts. Explain why that bound does not hold for
        a function that is not monotone.
  \item Compute $\int_0^1(3x^2-2x+1)\,dx$ with the fundamental theorem of calculus, then check
        the result using Simpson's rule with $n=100$ and the function
        \texttt{scipy.integrate.quad}. Report all three numbers and their differences.
  \item Compute $\int_1^{10}x^{-2}dx$ and $\int_1^{\infty}x^{-2}dx$. Report the difference
        between them, then explain why that difference equals $1/10$ without
        recomputing.
  \item Prove that if $f$ is continuous, $f\ge0$ on $[a,b]$, and $\int_a^bf=0$, then
        $f(x)=0$ for every $x\in[a,b]$. Also show that the statement fails
        if $f$ is only integrable.
  \item Compute $\int_0^{\pi/2}\sin x\cos x\,dx$ in two ways: the substitution
        $u=\sin x$, and the identity $\sin2x=2\sin x\cos x$. Compare the results with
        the numerical value from \texttt{quad}.
  \item Show that the Dirichlet function on $[0,1]$ satisfies
        $U(\mathcal{P},D)-L(\mathcal{P},D)=1$ for every partition, then explain
        why that fact is enough to conclude that it is not integrable.
  \item Estimate $\sum_{n=101}^{\infty}1/n^2$ with an integral bound and compare with
        the true value $9.950167\times10^{-3}$; what is the relative error?
\end{enumerate}

\noindent\textbf{Hints.} Item 3: the claim contains two different statements, about
the boundedness of the sums and about the equality of the two integrals. Item 7: use
continuity at one point to make the integral small on a narrow interval around
that point. Item 10: use $\sum_{n=N+1}^\infty1/n^2\le\int_N^\infty x^{-2}dx=1/N$
with $N=100$.
\end{latihan}

\begin{untukpeneliti}
What is genuinely proved in this chapter: the integrability criterion with both of its directions,
the integrability of continuous functions (using uniform continuity from Chapter 11), of monotone
functions, of Thomae's function, both parts of the fundamental theorem of calculus, the mean value theorem
for integrals, and the status of substitution and integration by parts as consequences, not
axioms.

What is quoted without proof and should be called so in class: Lebesgue's criterion
(Rudin, Theorem 11.33(b); Abbott, Chapter 7), Volterra's example of a bounded derivative that
is not integrable, and the absence of an elementary antiderivative for $e^{-x^2}$ (Liouville's
theorem).

Primary sources: Rudin, \emph{Principles of Mathematical Analysis}, Chapter 6, the section
Definition and Existence of the Integral (Theorems 6.4, 6.6, 6.8, 6.9, 6.10) and the section
Integration and Differentiation (6.19--6.22). Rudin sets it up for the
Riemann-Stieltjes integral; the Riemann integral is the case $\alpha(x)=x$. Abbott,
\emph{Understanding Analysis}, Chapter 7, is slower but more readable; Stewart only
for the computational rules.

Open question: how sharp the quadrature error bounds are when only $|f''|$ is
known, and how to estimate the error without any derivative at all; the standard
answers are adaptive quadrature and the difference of two levels of refinement. Limits of validity:
bounded functions on a closed bounded interval, improper integrals as limits, without
measure theory, and vector-valued integrals are not touched.
\end{untukpeneliti}

\peta{Chapter 13 follows Rudin, Chapter 6: the section Definition and Existence of the Integral
(partitions, upper and lower sums, the refinement lemma as Theorem 6.4, the criterion as
6.6, continuous and monotone functions as 6.8 and 6.9) and the section Integration and
Differentiation (6.19 change of variables, 6.21 the fundamental theorem of calculus, 6.22 integration by
parts). In Abbott, Chapter 7, from The Definition of the Riemann Integral to
Lebesgue's Criterion. In Stewart: The Substitution Rule and Average Value of a
Function (Chapter 5), Integration by Parts and Improper Integrals (Chapter 7), and 11.3
The Integral Test, the bridge to Chapter 5 of this book.}


\chapter{Series of Functions and the Exchange of Limits}

A small numerical library computes $\ln(1+x)$ with the series
$\sum_{k\ge 1}(-1)^{k+1}x^{k}/k$. Its author tested the function at $x=0.5$
with twelve terms and compared it against a standard library: the error was
$6.4\times 10^{-6}$. The test passed, the documentation states that the series is
accurate to six digits, and the code went into production. The first user called it with
$x=0.99$. Twelve terms give $0.6529204447$, whereas
$\ln 1.99=0.6881346387$. The error is $3.5\times10^{-2}$, about five thousand times
larger, for the same function with the same number of terms.

That damage is not a typo. For every $x$ in $(-1,1)$ the series converges
to $\ln(1+x)$; what changes is only how many terms are needed: $16$
terms at $x=0.5$, $89$ at $x=0.9$, $721$ at $x=0.99$, $5248$ at $x=0.999$,
all of them for an error below $10^{-6}$. No single number of terms suffices
for all $x$, because $x$ may lie as close as we please to $1$.

Compare with a closed subinterval $[-r,r]$, $0<r<1$. There one estimate
holds simultaneously for all $x$: the remainder of an alternating series is bounded by
$r^{N+1}/(N+1)$, and that bound does not mention $x$. For $r=0.5$ and $N=12$
the bound is $9.4\times10^{-6}$, the actual error $6.4\times10^{-6}$. This is the
difference that this entire chapter is about.

Chapter 5 dealt with series of numbers, and each $x$ gives a different series of
numbers. Chapters 11 and 13 define continuity and the integral point by point,
so neither of them has a way to move a limit past an integral sign without
extra conditions. The tool still missing is a distance between two functions; once that
distance exists, the whole family of series becomes a single sequence in a metric space.

\begin{untukanak}
Imagine a class with thirty students, and the teacher writes a single number
on the board: one hundred. Each student's task is to guess that number, and after every
guess the teacher announces whether the guess is getting closer. The first student is very
fast: within five guesses he is off by less than one. The second student is
slow: he needs twenty guesses. The third is very slow: he needs a hundred.
But all three, given enough time, end up off by less than one.

Now the teacher changes the question. Not ``does every student eventually get
close'', but ``is there a moment when all students are already close at the same
time''. For that class the answer is no: the slowest student is always
lagging, and one can always add a student who is slower still.

The same question arises for the formulas a computer uses. A formula
may be correct for every value put into it, yet the number of terms
needed can blow up without bound as the value approaches the edge. If a
program takes ``correct for every value'' to mean ``one formula suffices
for all values'', that program will be far off at the most troublesome values, and the
error will not show up in the test its author chose.

In this chapter we shall see why those two questions are not a matter of words,
but a matter of whether two operations may be exchanged in order: limit with
integral, limit with derivative, and summation with both.
\end{untukanak}

\section{Pointwise and uniform convergence}
\tingkat{basic}

Take a family of functions $f_1,f_2,\ldots$ all defined on a single
set $E\subseteq\mathbb{R}$. For one point $x$ that we hold fixed,
the values $f_1(x),f_2(x),\ldots$ are nothing but a sequence of numbers, and Chapter 4
already said what it means for that sequence to converge. Hence the first definition
almost writes itself.

\begin{definisi}
Suppose $f_n$ and $f$ are defined on $E\subseteq\mathbb{R}$.
\begin{enumerate}
\item We say $f_n\to f$ \emph{pointwise}\index{pointwise convergence} on $E$
      if for every $x\in E$ and every
      $\varepsilon>0$ there is $N$ such that $n\ge N$ implies
      $|f_n(x)-f(x)|<\varepsilon$. Here $N$ may depend on $x$ and
      on $\varepsilon$.
\item We say $f_n\to f$ \emph{uniformly}\index{uniform convergence} on
      $E$ if for every $\varepsilon>0$ there is $N$ such that $n\ge N$
      implies $|f_n(x)-f(x)|<\varepsilon$ for \emph{all} $x\in E$
      at once. Here $N$ depends only on $\varepsilon$.
\item If $h$ is bounded on $E$, the supremum norm\index{supremum norm} of $h$
      is $\|h\|_\infty=\sup_{x\in E}|h(x)|$.
\end{enumerate}
\end{definisi}

\kamus{uniform convergence}{konvergensi seragam}{all points must already be close
at the same time, not one by one on their own schedules.}

One quantifier changes place, and the whole content of this chapter hangs on
that move. In the first definition the order is ``for every point there is a
threshold''. In the second the order is ``there is one threshold for all points''.
The first sentence is true for the series in the chapter opening, the second is false, and
the difference between them is the difference between $16$ terms and $5248$ terms.

\begin{bacakeras}
Two voices take turns.

First voice: ``Does this series get there?''
Second voice: ``It gets there. Slowly, but it gets there.''
First voice: ``Does it get there for this point?''
Second voice: ``For this point. For the point next door, another time, at its own pace.''

That is the first promise. Every point may be slow, as long as it has one tempo
of its own. A point far from the edge arrives within a dozen steps;
a point pressed against the edge needs thousands.

Now the first voice asks in other words: ``Is there one tempo for
all points at once?''

The second voice stays silent longer. ``On a closed interval that stops before the edge,
yes. On the whole open interval, no. And what is missing there is not a failure of
counting; it merely demands a tempo that moves together with the point.''

Repeat this line until it sits on the tongue: all points may be slow, as long as one
tempo. If the tempo must differ for each point, its name is not uniform.
\end{bacakeras}

\begin{contoh}
Suppose $f_n(x)=x^n$ on $E=[0,1]$. For $0\le x<1$ we have $x^n\to0$,
whereas $x^n=1$ for all $n$ if $x=1$. So $f_n$ converges pointwise
to
\[
  f(x)=\begin{cases}0, & 0\le x<1,\\ 1, & x=1.\end{cases}
\]
That convergence is not uniform, since $\sup_{x\in[0,1]}|x^n-f(x)|=1$ for
every $n$: at $x=1$ the difference is exactly $1$, and for $x$ close enough to $1$
from the left $x^n$ approaches $1$ while $f(x)=0$.
\end{contoh}

That example shows two things. First, the pointwise limit of continuous functions
need not be continuous: every $x^n$ is continuous, its limit jumps at $x=1$. Second,
uniformity depends on the set, not only on the formula of the function.

That dependence can be measured. On $[0,r]$ with $r<1$ the convergence is
uniform, since $\sup_{[0,r]}|x^n|=r^n\to0$. What is interesting is when the right
endpoint also moves with $n$. On $[0,1-1/n]$ the supremum is
$(1-1/n)^n\to1/e=0.3678794412$; so shrinking by $1/n$ does not yet
rescue uniformity. On $[0,1-1/\sqrt{n}]$ the supremum is
$(1-1/\sqrt n)^n$, which collapses like $e^{-\sqrt n}$: $2.2\times10^{-2}$ at
$n=10$, $2.7\times10^{-5}$ at $n=100$, $2.2\times10^{-44}$ at $n=10^4$. The rate
of shrinking, not the formula of the function, is what decides.

The diagonal test yields sharper additional information. Take the points
$x_n=1-1/n$ chasing the break point. Then $f_n(x_n)=(1-1/n)^n\to1/e$,
while $f(x_n)=0$ for all $n$. So not only does the supremum fail to
shrink: the values at well-chosen points also refuse to approach the
pointwise limit. This pattern is used repeatedly in this chapter. To
refute uniformity, find one sequence of points that keeps the difference
large.

Two simple families separate the two symptoms. In the first, $g_n$ is a
triangle of height $1$ with base $1/n$ and apex at $1-1/n$: the apex
moves right, the height never drops, and the area is $1/(2n)$. A point
$x<1$ is touched by the apex only while $1/(2(1-x))<n<3/(2(1-x))$, so at $x=0.9$
only nine values of $n$ give a positive value, at $x=0.99$ ninety-nine,
and once the last apex has passed the value is zero forever. The limit is zero
at every point, whereas $\|g_n\|_\infty=1$. In the second, $h_n$ has height $1/n$
with base $1$ in the same place: the apex shrinks instead of moving, and
$\|h_n\|_\infty=1/n\to0$ (Figure~\ref{fig:puncak}).

\begin{figure}[htbp]
\centering
\begin{tikzpicture}[x=5.6cm,y=3.1cm,font=\small]
  \draw[->,layergrey] (0,0)--(1.15,0) node[right,font=\scriptsize]{$x$};
  \draw[->,layergrey] (0,0)--(0,1.32) node[above,font=\scriptsize]{$y$};
  \node[below,font=\scriptsize,layergrey] at (1,0) {$1$};
  \foreach \n/\c in {16/layerindigo!15,8/layerindigo!30,4/layerindigo!50,2/layerindigo!75}{
    \fill[\c] ({1-1/\n-1/(2*\n)},0) -- ({1-1/\n},1) -- ({1-1/\n+1/(2*\n)},0) -- cycle;
  }
  \draw[dashed,layergrey] (0,1)--(1.06,1);
  \node[above,font=\scriptsize,layergrey] at (0.72,1) {height $1$};
  \begin{scope}[xshift=7.0cm]
    \draw[->,layergrey] (0,0)--(1.15,0) node[right,font=\scriptsize]{$x$};
    \draw[->,layergrey] (0,0)--(0,1.32) node[above,font=\scriptsize]{$y$};
    \node[below,font=\scriptsize,layergrey] at (1,0) {$1$};
    \foreach \n/\c in {8/layerteal!18,4/layerteal!32,2/layerteal!52,1/layerteal!72}{
      \fill[\c] (0,0) -- (0.5,{1/\n}) -- (1,0) -- cycle;
    }
    \node[right,font=\scriptsize,layergrey] at (0.52,0.5) {$n=2$};
    \node[right,font=\scriptsize,layergrey] at (0.52,0.25) {$n=4$};
    \node[right,font=\scriptsize,layergrey] at (0.52,0.125) {$n=8$};
  \end{scope}
\end{tikzpicture}
\caption{Apexes that move and apexes that shrink. Left: $g_n$ has constant height
$1$, base $1/n$, and apex at $1-1/n$; the height does not drop,
so $\|g_n\|_\infty=1$ even though $g_n(x)\to0$ at every point. Right: $h_n$ has
height $1/n$ with base $1$ in the same place; its whole graph fits inside the
band of height $1/n$, so it converges uniformly.}
\label{fig:puncak}
\end{figure}

\begin{lemma}
$f_n\to f$ uniformly on $E$ if and only if $\|f_n-f\|_\infty\to0$.
\end{lemma}

\begin{proof}
Uniform convergence means: for every $\varepsilon>0$ there is $N$ such that
for $n\ge N$ we have $|f_n(x)-f(x)|\le\varepsilon$ for all $x\in E$,
that is, $\sup_x|f_n(x)-f(x)|\le\varepsilon$. That is precisely the statement that the
sequence of real numbers $\|f_n-f\|_\infty$ tends to zero. The only technical
difference is that the definition uses $<\varepsilon$ whereas the norm uses
$\le\varepsilon$; the two are equivalent because $\varepsilon$ is arbitrary.
\end{proof}

That lemma turns a question about functions into a question about a single
sequence of numbers: to check uniformity one supremum suffices, to reject it
it suffices to show that this supremum does not tend to zero, as with $x^n$.

Uniformity has an exact picture, namely a band. Take $\varepsilon=0.1$. If
the whole graph of $f_n$ for $n\ge N$ lies between $f-\varepsilon$ and
$f+\varepsilon$, then that single $N$ works for all points. What is checked is
not one point, but whether any part of the graph escapes the band
of width $0.2$. For $g_n$ that band is always pierced by the apex; the
hoped-for situation is visible in Figure~\ref{fig:pita}.

\begin{figure}[htbp]
\centering
\begin{tikzpicture}[x=2.55cm,y=4.0cm,font=\small]
  \draw[->,layergrey] (0,0)--(4.5,0) node[right,font=\scriptsize]{$x$};
  \draw[->,layergrey] (0,0)--(0,1.18) node[above,font=\scriptsize]{$y$};
  \fill[layergrey!12] plot[domain=0:4,samples=90] (\x,{0.55+0.22*sin(2*\x r)+0.1})
    -- plot[domain=4:0,samples=90] (\x,{0.55+0.22*sin(2*\x r)-0.1}) -- cycle;
  \draw[dashed,layergrey] plot[domain=0:4,samples=90] (\x,{0.55+0.22*sin(2*\x r)+0.1});
  \draw[dashed,layergrey] plot[domain=0:4,samples=90] (\x,{0.55+0.22*sin(2*\x r)-0.1});
  \draw[thick,accent] plot[domain=0:4,samples=90] (\x,{0.55+0.22*sin(2*\x r)});
  \draw[thick,layerindigo] plot[domain=0:4,samples=180]
    (\x,{0.55+0.22*sin(2*\x r)+0.06*cos(9*\x r)});
  \draw[<->,layergrey] (3.78,{0.55+0.22*sin(2*3.78 r)-0.1})
    -- (3.78,{0.55+0.22*sin(2*3.78 r)+0.1});
  \node[right,font=\scriptsize,layergrey] at (3.8,{0.55+0.22*sin(2*3.78 r)}) {$0.2$};
  \node[accent,font=\scriptsize,above left] at (0.42,{0.55+0.22*sin(2*0.42 r)}) {$f$};
  \node[layerindigo,font=\scriptsize,below right] at (0.95,{0.55+0.22*sin(2*0.95 r)+0.06*cos(9*0.95 r)}) {$f_n$};
\end{tikzpicture}
\caption{The uniform band of width $0.2$ around $f$, that is between $f-0.1$
and $f+0.1$. The wavy graph of $f_n$ still lies entirely inside the
band, so $\|f_n-f\|_\infty\le0.1$ and one and the same $N$ works for all
points at once.}
\label{fig:pita}
\end{figure}

\begin{tangan}
Take a two-meter rope, tape, and two people. Stretch the rope straight on the
floor, then tape five marks at equal spacing. The first person holds the left end
and lifts it by just $5$ cm. Watch: only that end rises, the other four marks
still touch the floor. That is pointwise convergence.

Now both people lift both ends together on the same count: $5$ cm, then $3$ cm,
then $1$ cm, then half a centimeter. Watch the whole length of the rope. No point is
left on the floor, and the greatest distance from rope to floor is the same at all
points, because both ends rise together. That is uniform convergence.

Finally, the first person lifts his end while walking toward the middle, while
the second lowers hers. The rope stays tilted, and there is always some part
$5$ cm above the floor: every point eventually comes down, but the greatest distance
never shrinks. That is $g_n$ in Figure~\ref{fig:puncak}.
\end{tangan}

\begin{lemma}[Uniform Cauchy criterion]
The sequence $f_n$ converges uniformly on $E$ if and only if for every
$\varepsilon>0$ there is $N$ such that $n,m\ge N$ implies
$|f_n(x)-f_m(x)|<\varepsilon$ for all $x\in E$. This last condition is called the
uniform Cauchy criterion\index{uniform Cauchy criterion}.
\end{lemma}

\begin{proof}
$(\Rightarrow)$ If $f_n\to f$ uniformly, take $N$ from the definition of convergence
for $\varepsilon/2$. For $n,m\ge N$ and every $x$,
\[
  |f_n(x)-f_m(x)|\le|f_n(x)-f(x)|+|f(x)-f_m(x)|<\varepsilon .
\]
$(\Leftarrow)$ For every fixed $x$, the sequence $f_n(x)$ is a Cauchy sequence
in $\mathbb{R}$, so it converges to a number that we call $f(x)$. This
defines a function $f$ point by point. To obtain
uniformity, fix $n\ge N$ in the Cauchy condition and let
$m\to\infty$. Since $f_m(x)\to f(x)$, we get $|f_n(x)-f(x)|\le\varepsilon$
for all $x\in E$ at once, and the $N$ used does not depend on $x$.
\end{proof}

The proof of the second direction uses only the completeness of $\mathbb{R}$: the space of
functions is complete with respect to the supremum norm, and that completeness is what lets a
series of functions be summed without fear that the limit leaves the class of functions we
use.

\begin{catatan}
Uniform convergence is a property of a pair: a sequence and a set. The sentence
``$x^n$ converges uniformly'' is incomplete without naming the set. On
$[0,1]$ the answer is no, on $[0,0.99]$ yes: the supremum norm of the error is
$0.99^n$, namely $0.37$ at $n=100$ and $4.3\times10^{-5}$ at $n=1000$.
Therefore every claim of uniformity in this chapter names its set.
\end{catatan}

\section{The Weierstrass test: uniformity from a single series of numbers}
\tingkat{intermediate}

A series of functions is a series whose terms are functions. Its partial sums
$S_N=\sum_{n=1}^{N}f_n$ are functions, so a series of functions is nothing but a
particular sequence of functions, and all the definitions of the previous section apply
as soon as we replace $f_n$ by $S_N$.

\begin{definisi}
The series $\sum_{n\ge1}f_n$ \emph{converges pointwise} on $E$ if for
every $x\in E$ the series of numbers $\sum_n f_n(x)$ converges. That series
\emph{converges uniformly}\index{series of functions} on $E$ if its sequence of partial
sums converges uniformly on $E$.
\end{definisi}

The following test is the cheapest way to obtain uniformity, and almost every
example in this book uses it. The idea: if each term is bounded by a number
whose series converges, then the tail of the series of functions is bounded by the tail of
that series of numbers, and the numerical tail knows nothing about $x$.

\begin{teorema}[Weierstrass $M$-test]\index{Weierstrass test}
Suppose $f_n$ is defined on $E$ and there are numbers $M_n\ge0$ with
$|f_n(x)|\le M_n$ for all $x\in E$ and all $n$. If $\sum M_n$
converges, then $\sum f_n$ converges uniformly on $E$. If $E$ is a closed
bounded interval and every $f_n$ is continuous, then the sum is continuous on $E$.
\end{teorema}

\begin{proof}
Write $S_N=\sum_{i=1}^{N}f_i$. For $n>m$ and every $x\in E$,
\[
  |S_n(x)-S_m(x)|\le\sum_{i=m+1}^{n}|f_i(x)|\le\sum_{i=m+1}^{n}M_i .
\]
Since $\sum M_n$ converges, the tail $\sum_{i=m+1}^{n}M_i$ can be made smaller
than $\varepsilon$ provided $m,n$ are large enough, and that choice does not
mention $x$. The uniform Cauchy criterion closes the proof of uniformity. For
continuity, each $S_N$ is continuous as a finite sum of continuous functions, and the
uniform limit of continuous functions is continuous; that last statement is proved in
the next section.
\end{proof}

\begin{contoh}
On $E=[-1,1]$ the series $\sum_{n\ge1}x^n/n^2$ converges uniformly, since
$|x^n/n^2|\le1/n^2$ and $\sum1/n^2=\pi^2/6$ converges. Its sum is continuous on
$[-1,1]$, and its value at $x=1$ is $1.6449340668$. The sum of two million terms
gives $1.6449335668$; the difference is $5.0\times10^{-7}$, proportional to the tail
bound $1/N=5\times10^{-7}$. At $x=-1$ the series converges to $-\pi^2/6$, since
the terms alternate in sign with decreasing absolute value tending to zero.
\end{contoh}

\begin{contoh}
On $E=\mathbb{R}$ the series $\sum_{n\ge0}2^{-n}\sin(2^n x)$ converges uniformly,
since $|2^{-n}\sin(2^n x)|\le2^{-n}$ and $\sum2^{-n}=1$. Every term is continuous,
so the sum is continuous at every point. What is surprising is that the sum is not
differentiable at any point: the classical result of Weierstrass (1872), extended
by Hardy up to ratio $1$, which we quote without proof. Uniformity guarantees
continuity, but it guarantees nothing about derivatives.
\end{contoh}

\begin{catatan}
The $M$-test gives a sufficient condition, not a necessary one, and its failure is not a proof
of anything. The series $\sum_{n\ge1}(-1)^n(1-x)x^n$ on $[0,1]$ converges uniformly:
the remainder after the $N$-th term is bounded by $(1-x)x^{N+1}$, whose maximum on $[0,1]$
is $\tfrac{1}{N+2}\left(\tfrac{N+1}{N+2}\right)^{N+1}\to0$, namely
$3.20\times10^{-2}$ at $N=10$ and $3.67\times10^{-4}$ at $N=1000$. But
$\sup_{[0,1]}|f_n|=n^n/(n+1)^{n+1}$, proportional to $1/(en)$, so that
$\sum_n\sup|f_n|$ diverges: $0.88$ up to $N=10$ and $2.53$ up to $N=1000$.
There is no convergent comparison series, so the $M$-test does not apply even though
the uniformity is genuine.
\end{catatan}

\begin{untukanak}
There are two ways to work through a list of chores. The first: finish all the chores
one by one, then add up the total. The second: keep a running total while
working, then finish the rest. For ordinary addition both ways give the
same result, and we are used to swapping them without thinking.

But some tasks do not like being swapped. Imagine you measure the length of a
pole's shadow every hour for a day, compute the average, then
look for the hour when that average is longest. Reverse the order: first find the
longest hour on each day, then average that hour over the month. The two
results can differ widely, not because the arithmetic is wrong, but because
``longest'' and ``average'' genuinely do not commute.

The same thing happens with the formulas used to add up infinitely
many terms, and then integrate or differentiate them. Sometimes the order may be
swapped and the result is exactly the same; sometimes not, and the difference is large. In this chapter
we shall know the warning signs before the result goes wrong.
\end{untukanak}

\section{When limits may be exchanged}
\tingkat{intermediate}

The three operations we habitually use on functions are taking a limit at
a point, integrating, and differentiating. Each of the three is itself a limit:
continuity at $x_0$ means $\lim_{x\to x_0}f(x)=f(x_0)$, the integral is the
limit of Riemann sums, the derivative is the limit of a difference quotient. So moving the
limit of a sequence past one of those operations means exchanging two limits, and
that exchange is not free: its price is one hypothesis, namely
uniformity\index{exchange of limits}.

\begin{teorema}
If $f_n$ is continuous on $E$ and $f_n\to f$ uniformly on $E$, then $f$ is continuous
on $E$.
\end{teorema}

\begin{proof}
Take $x_0\in E$ and $\varepsilon>0$. Since the convergence is uniform, there is $n$
with $\|f_n-f\|_\infty<\varepsilon/3$. Since $f_n$ is continuous at $x_0$, there is
$\delta>0$ such that $|x-x_0|<\delta$ implies $|f_n(x)-f_n(x_0)|<\varepsilon/3$.
Then for $x\in E$ with $|x-x_0|<\delta$,
\[
  |f(x)-f(x_0)|\le|f(x)-f_n(x)|+|f_n(x)-f_n(x_0)|+|f_n(x_0)-f(x_0)|
  <\varepsilon .
\]
\end{proof}

Note the logical order: $n$ is chosen first, then $\delta$, and
$\delta$ may depend on $n$. What is not allowed is choosing a $\delta$
that must work for all $n$ at once.

\begin{tantangan}
The series $\sum_{n\ge0}\bigl(x^{2n}-x^{2n+1}\bigr)$ has continuous terms on $[0,1]$
and converges at every one of its points. Determine its sum, then show that this sum
is not continuous at $x=1$. Why does this result not contradict the theorem above?
A correct answer must name one number that may depend on $\varepsilon$ but must not
depend on $x$.

Hint 1: the partial sums have a closed form, and that form contains
$x^{2N+2}$.

Hint 2: for $x<1$ the sum is $1/(1+x)$, whereas at $x=1$ every term
equals zero; so the value jumps from $1/2$ to $0$.

If you get stuck, go back to the section on the diagonal test and ask what happens at
$x_N=1-1/N$.
\end{tantangan}

\begin{teorema}[Term-by-term integration]
If $f_n$ is Riemann integrable on $[a,b]$ and $f_n\to f$ uniformly on
$[a,b]$, then $f$ is integrable on $[a,b]$ and
\[
  \int_a^b f=\lim_{n\to\infty}\int_a^b f_n .
\]
\end{teorema}

\begin{proof}
For integrability, take $\varepsilon>0$ and choose $n$ with
$\|f_n-f\|_\infty<\varepsilon$. Since $f_n$ is integrable, there is a partition
$\mathcal{P}$ with $U(\mathcal{P},f_n)-L(\mathcal{P},f_n)<\varepsilon(b-a)$.
On each subinterval, $M_i(f)\le M_i(f_n)+\varepsilon$ and
$m_i(f)\ge m_i(f_n)-\varepsilon$, so that
$U(\mathcal{P},f)\le U(\mathcal{P},f_n)+\varepsilon(b-a)$ and
$L(\mathcal{P},f)\ge L(\mathcal{P},f_n)-\varepsilon(b-a)$. Those three inequalities
give $U(\mathcal{P},f)-L(\mathcal{P},f)<3\varepsilon(b-a)$, and the integrability
criterion of Chapter 13 closes the argument. For the value of the integral, it suffices
to note that
\[
  \left|\int_a^b f_n-\int_a^b f\right|
  =\left|\int_a^b (f_n-f)\right|
  \le\int_a^b|f_n-f|\le(b-a)\,\|f_n-f\|_\infty ,
\]
and the right-hand side tends to zero.
\end{proof}

\begin{akibat}[Integration of series]
If $\sum_{n\ge1}f_n$ converges uniformly on $[a,b]$ and every $f_n$ is
integrable there, then
\[
  \int_a^b\sum_{n=1}^{\infty}f_n=\sum_{n=1}^{\infty}\int_a^b f_n .
\]
\end{akibat}

This corollary is applied to the partial sums, with uniformity usually
checked through the $M$-test. For example, the series $\sum_{n\ge1}x^n/n^2$ converges uniformly
on $[0,1]$, so that
\[
  \int_0^1\sum_{n=1}^{\infty}\frac{x^n}{n^2}\,dx
  =\sum_{n=1}^{\infty}\frac{1}{n^2(n+1)}=\frac{\pi^2}{6}-1=0.6449340668 .
\]
That last number is also obtained by summing the first two million terms of the series on
the right-hand side, which gives $0.6449340668$ at machine precision.

\begin{contoh}[Failure without uniformity]
Suppose $f_n(x)=nx(1-x^2)^n$ on $[0,1]$. For every $x\in(0,1]$ the factor
$(1-x^2)^n$ decreases geometrically, so $n(1-x^2)^n\to0$ and $f_n(x)\to0$;
at $x=0$ every $f_n$ equals $0$. So the pointwise limit is the zero function.
Its integral can be computed exactly by the substitution $u=1-x^2$:
\[
  \int_0^1 nx(1-x^2)^n\,dx=\frac{n}{2}\int_0^1 u^n\,du=\frac{n}{2(n+1)} .
\]
That sequence tends to $1/2$, not to $0$. Numerical quadrature gives
$0.4545454545$ for $n=10$ and $0.4950495050$ for $n=100$, matching
$n/(2(n+1))$ to the last digit. So $\lim_n\int_0^1f_n=1/2$ whereas
$\int_0^1\lim_nf_n=0$.
\end{contoh}

That failure has a geometric shape. The function $f_n$ is zero at both ends
and peaks at $x=1/\sqrt{2n+1}$, with height $1.3396685$ at $n=10$ and
$4.2834707$ at $n=100$: the peak grows taller and narrower, so that its
area stays near $1/2$ while its height at every point tends to zero.

\begin{teorema}[Term-by-term differentiation]
Suppose $f_n$ is differentiable on $[a,b]$, the sequence $f_n'$ converges uniformly
on $[a,b]$, and there is one point $x_0\in[a,b]$ where $f_n(x_0)$ converges.
Then $f_n$ converges uniformly to some $f$ that is differentiable on $[a,b]$,
and
\[
  f'(x)=\lim_{n\to\infty}f_n'(x)
\]
for every $x\in[a,b]$.
\end{teorema}

\begin{proof}[Sketch of proof]
The heart of it is the mean value theorem applied to the difference of two terms: for $m,n$ and every
$x,y\in[a,b]$,
\[
  \bigl|(f_n-f_m)(x)-(f_n-f_m)(y)\bigr|\le\|f_n'-f_m'\|_\infty\,|x-y| .
\]
Since $f_n'$ is uniformly Cauchy, the right-hand side can be made small simultaneously for all
pairs $x,y$ once $m,n$ are large. Together with convergence at one point $x_0$,
that inequality makes $f_n$ uniformly Cauchy, so $f_n\to f$ uniformly. For the
derivative, divide by $|x-y|$, let $y\to x$, then exchange the two limits that
appear; the same inequality guarantees that the exchange is legitimate. The technical details are in
Rudin, Theorem 7.17.
\end{proof}

The condition ``there is one point where $f_n(x_0)$ converges'' cannot be dropped: $f_n$
and $f_n+1$ have the same sequence of derivatives, but only one of them
converges at $x_0$. Conversely, uniformity of $f_n$ alone is also not enough to
conclude anything about the derivative.

\begin{contoh}[Derivatives that do not converge]
Suppose $f_n(x)=\sin(nx)/n$ on $[0,\pi]$. Since $\|f_n\|_\infty=1/n$,
the sequence converges uniformly to the zero function; its value is $0.1$ at $n=10$ and
$0.01$ at $n=100$. But $f_n'(x)=\cos(nx)$: at $x=0$ that sequence equals
$1$ for all $n$, whereas at $x=\pi/2$ its values are $0,-1,0,1,0,-1,\ldots$
according to the parity of $n$, which does not converge. So $f_n'$ fails to converge pointwise
anywhere except at special points, let alone uniformly. Differentiation
does not forgive small high-frequency oscillations: $f_n$ shrinks to zero,
but its slope does not.
\end{contoh}

Figure~\ref{fig:turunan} shows three of those derivative curves on $[0,\pi]$. At
$x=\pi/2$ they take the values $1$, $0$, and $-1$: three different values of one
sequence at one point, and that is already enough to say the limit does not exist.

\begin{figure}[htbp]
\centering
\begin{tikzpicture}[x=3.4cm,y=1.5cm,font=\small]
  \draw[->,layergrey] (0,0)--(3.55,0) node[right,font=\scriptsize]{$x$};
  \draw[->,layergrey] (0,-1.28)--(0,1.28) node[above,font=\scriptsize]{$y$};
  \draw[dashed,layergrey] (1.571,-1.25)--(1.571,1.25);
  \node[below,font=\scriptsize,layergrey] at (1.571,-1.25) {$\pi/2$};
  \draw[thick,layerindigo] plot[domain=0:3.1416,samples=160] (\x,{cos(4*\x r)});
  \draw[thick,layerteal] plot[domain=0:3.1416,samples=160] (\x,{cos(5*\x r)});
  \draw[thick,layerorange] plot[domain=0:3.1416,samples=160] (\x,{cos(6*\x r)});
  \fill[ink] (1.571,1) circle (1.5pt);
  \fill[ink] (1.571,0) circle (1.5pt) node[right,font=\scriptsize] {$0$};
  \fill[ink] (1.571,-1) circle (1.5pt) node[right,font=\scriptsize] {$-1$};
  \node[right,font=\scriptsize] at (1.585,1.06) {$1$};
\end{tikzpicture}
\caption{The derivatives $f_n'=\cos(nx)$ on $[0,\pi]$ for $n=4$ (indigo), $n=5$
(teal), and $n=6$ (orange). At $x=\pi/2$ they take the values $1$, $0$, and $-1$: three
different values of one sequence at one point, so the limit of the derivatives does not exist,
even though $f_n=\sin(nx)/n$ converges uniformly to zero with amplitude $1/4$,
$1/5$, and $1/6$.}
\label{fig:turunan}
\end{figure}

\begin{catatan}
In machine learning, the risk function is often written as an expectation over the
data distribution, $\mathbb{E}_{(x,y)\sim D}[\ell(f(x),y)]$, and then differentiated
with respect to the model parameters. If the distribution is finite discrete, as on
the training set, that expectation is a finite sum and exchanging the derivative with
the sum is always legitimate, with no conditions. If it is continuous, the expectation is an integral, and
moving the derivative inside it demands a domination condition more delicate
than uniformity; the tool is the dominated convergence theorem of measure
theory, which is not discussed here. ``Differentiate inside the integral'' is an operation
with a license, not an operation that is always correct.
\end{catatan}

\section{Power series and the radius of convergence}
\tingkat{advanced}

A power series is a series of functions of the most rigid shape: the $n$-th term is
a fixed coefficient times $x^n$. That rigidity bears fruit. The set on which the series
converges is always an interval symmetric about its center, so that all the
information about where the series is useful is contained in one number.

\begin{definisi}
A series of the form $\sum_{n\ge0}a_nx^n$ with $a_n\in\mathbb{R}$ is called a power
series\index{power series} centered at $0$. The number
\[
  R=\sup\Bigl\{r\ge0:\ \sum_{n\ge0}a_nr^n\ \text{converges}\Bigr\}
\]
is called the radius of convergence\index{radius of convergence} of that series.
\end{definisi}

\kamus{power series}{deret pangkat}{a series whose terms are fixed coefficients times
rising powers of one variable, like the same recipe for every value.}

\kamus{radius of convergence}{jari-jari konvergensi}{the furthest distance from the center
at which that recipe still produces a finite number.}

\begin{teorema}
Suppose $\sum a_nx^n$ is a power series with radius $R$. Then the series
converges absolutely for every $|x|<R$ and diverges for every $|x|>R$.
Moreover
\[
  \frac{1}{R}=\limsup_{n\to\infty}|a_n|^{1/n},
\]
with the conventions $1/0=\infty$ and $1/\infty=0$.
\end{teorema}

\begin{proof}
Fix $x\ne0$ and write $|a_nx^n|^{1/n}=|x|\cdot|a_n|^{1/n}$. By the root test
of Chapter 5, the series $\sum a_nx^n$ converges absolutely if
$\limsup_n|a_nx^n|^{1/n}<1$, that is if $|x|<1/\limsup_n|a_n|^{1/n}$, and
diverges if that limit superior is greater than $1$, since its terms do not
tend to zero. Since $x$ is arbitrary, the boundary between the two regimes is exactly
$R=1/\limsup_n|a_n|^{1/n}$, and the set of convergence lies inside $[-R,R]$.
\end{proof}

At $|x|=R$ the root test gives no verdict, and indeed there is no general
verdict: $\sum x^n/n^2$ converges at both ends, $\sum x^n/n$ only at $x=-1$, and
$\sum x^n$ diverges at both.

\begin{akibat}
If $a_n\ne0$ for $n$ large enough and $|a_{n+1}/a_n|\to L$, then $R=1/L$
with the same conventions as above.
\end{akibat}

The reason is the chain of inequalities
\[
  \liminf|a_{n+1}/a_n|\le\liminf|a_n|^{1/n}\le\limsup|a_n|^{1/n}
  \le\limsup|a_{n+1}/a_n|,
\]
which is proved as an exercise in Chapter 5. If the ratio has a limit, all four of those
numbers are equal, so the ratio formula and the root formula give the same $R$. In
hand computation, the ratio formula is almost always faster.

Now the part that makes power series useful: they may be differentiated and
integrated term by term, and the result holds on the whole interval of
convergence, not only on a small piece of it.

\begin{teorema}[Term-by-term differentiation and integration]
Suppose $\sum a_nx^n$ has radius $R>0$. Then for every $x$ with
$|x|<R$,
\[
  \frac{d}{dx}\sum_{n\ge0}a_nx^n=\sum_{n\ge1}n\,a_nx^{n-1},
  \qquad
  \int_0^x\sum_{n\ge0}a_nt^n\,dt=\sum_{n\ge0}\frac{a_n}{n+1}x^{n+1},
\]
and both new series also have radius $R$.
\end{teorema}

\begin{proof}
Take $r$ with $0<r<R$ and choose $\rho$ with $r<\rho<R$. The series
$\sum|a_n|\rho^n$ converges, so there is $C$ with $|a_n|\rho^n\le C$ for all
$n$. For $|x|\le r$,
\[
  |n\,a_nx^{n-1}|\le n\,\frac{C}{\rho}\left(\frac{r}{\rho}\right)^{n-1},
\]
and the series $\sum_n n(r/\rho)^{n-1}$ converges since $r/\rho<1$. The $M$-test
therefore gives uniformity of the differentiated series on $[-r,r]$, and the term-by-term
differentiation theorem gives equality there. Every point $|x|<R$ lies in some
$[-r,r]$ with $r<R$, so the equality holds on all of $(-R,R)$. For the
integral, the same route is used on $\sum|a_n|r^{n-1}$, which also converges.
The radius does not change because the coefficients of the differentiated series are $na_n$, and
\[
  \limsup_n(n|a_n|)^{1/n}
  =\limsup_n|a_n|^{1/n}\cdot\lim_n n^{1/n}
  =\limsup_n|a_n|^{1/n},
\]
since $n^{1/n}\to1$. The same argument applies to the coefficients $a_n/(n+1)$,
because $(n+1)^{-1/n}\to1$.
\end{proof}

That proof uses the whole content of the chapter: the $M$-test for uniformity, the differentiation
theorem to move the limit, and the radius formula to make sure the new series does not
shrink. Uniformity is claimed only on $[-r,r]$, $r<R$; in the opening example,
$\ln(1+x)$ is uniform on every $[-r,r]$ with $r<1$ but not on $(-1,1)$.

\begin{contoh}
For $|x|<1$, the geometric series gives $\sum_{n\ge0}x^n=1/(1-x)$. Differentiating
term by term,
\[
  \sum_{n\ge1}n\,x^{n-1}=\frac{1}{(1-x)^2},
\]
and at $x=1/2$ the right-hand side is $4$. The sum of the first two hundred terms of the series on the
left gives $4.000000$ at machine precision. Integrating term by term from
$0$ to $x$ gives
\[
  \sum_{n\ge1}\frac{(-1)^{n+1}x^n}{n}=\ln(1+x),\qquad |x|<1,
\]
which is the series from the chapter opening. Both results are not numerical coincidences: both are
consequences of the theorem above, because $1/2$ and the $x$ used lie inside
the radius.
\end{contoh}

With that theorem, the opening question of the chapter is answered. The series for $\ln(1+x)$ converges
uniformly on every $[-r,r]$ with $r<1$, since $|(-1)^{n+1}x^n/n|\le r^n/n$
and $\sum r^n/n$ converges. On $(-1,1)$ uniformity fails, and the price shows
in the number of terms: $16$ at $x=0.5$, $5248$ at $x=0.999$. At the endpoint
$x=1$ the series still converges, since it becomes an alternating harmonic series. That is
not a consequence of the radius but of Abel's theorem: if $\sum a_nR^n$ converges,
then $\sum a_nx^n$ converges uniformly on $[0,R]$. We quote it; the proof is in
Rudin, chapter ``Some Special Functions'', Theorem 8.2.

One practical warning: uniform bounds are designed for the worst case. For the
series $\ln(1+x)$ with $N=12$ and $r=0.5$, the bound $r^{13}/13=9.4\times10^{-6}$
is only three times the actual error at $x=0.5$, but at $x=0.1$ it is loose by about
ten thousand times.

\begin{resep}
\begin{enumerate}
\item Write down the difference, $M_N=\sup_{x\in E}|f_N(x)-f(x)|$, or estimate it with
      the value at the worst point.
\item Find the point $x_N$ that maximizes that difference; usually a peak, or
      a point chasing the edge of the interval.
\item Compute $M_N$ as an expression containing only $N$, then check
      $M_N\to0$.
\end{enumerate}
A warning about order: find $x_N$ first for each $N$, then compute its
supremum. Reversing that order, that is choosing $N$ first and then looking for one $x$ that
represents all points, yields a wrong verdict of uniformity for $g_n$ in
Figure~\ref{fig:puncak}, whose apex actually shifts together with $N$.
\end{resep}

\section{Worked numerical examples}
\tingkat{intermediate}

The three groups of numbers below are recomputed in the computational lab. The first
measures the failure of uniformity of $x^n$ on $[0,1]$.

\begin{center}
\small
\begin{tabular}{rrrrr}
\hline
$n$ & $\sup_{[0,1]}x^n$ & $(1-\tfrac1n)^n$ & $(1-\tfrac1{\sqrt n})^n$ & $0.9^{\,n}$\\
\hline
$10$    & $1$ & $0.3486784401$ & $2.234\times10^{-2}$ & $0.3486784401$\\
$100$   & $1$ & $0.3660323413$ & $2.656\times10^{-5}$ & $2.656\times10^{-5}$\\
$1000$  & $1$ & $0.3676954248$ & $1.108\times10^{-14}$ & $1.748\times10^{-46}$\\
$10000$ & $1$ & $0.3678610464$ & $2.249\times10^{-44}$ & $<10^{-450}$\\
\hline
\end{tabular}
\end{center}

The second and fifth columns are two different rates. At the fixed point $x=0.9$
the error is $0.9^n$, falling by a factor of $0.9$ per term: $132$ terms are needed to
reach $10^{-6}$, whereas at $x=0.5$ $20$ terms suffice and at $x=0.99$ it takes
$1375$. Every point has its own exponential rate, with base $x$. The second
column does not move: the uniform error on $[0,1]$ stays $1$ for all $n$, so
its rate of decrease is zero. What is being compared is therefore not two numbers,
but something that tends to zero and something that does not.

The second group of numbers is the error of approximating $e^{x}$ by a Taylor series at
$x=1$. The remainder of the series is bounded by $\frac{e}{(N+1)!}$ for $x\in[0,1]$, since
the $(N+1)$-st derivative of $e^x$ is $e^x\le e$.

\begin{center}
\small
\begin{tabular}{rrrr}
\hline
$N$ & $S_N(1)=\sum_{k=0}^{N}\frac{1}{k!}$ & $e-S_N(1)$ & bound $\frac{e}{(N+1)!}$\\
\hline
$5$  & $2.716666666667$ & $1.615\times10^{-3}$ & $3.775\times10^{-3}$\\
$10$ & $2.718281801146$ & $2.731\times10^{-8}$ & $6.810\times10^{-8}$\\
$15$ & $2.718281828459$ & $5.063\times10^{-14}$ & $1.299\times10^{-13}$\\
$20$ & $2.718281828459$ & $0$ (below precision) & $5.320\times10^{-20}$\\
\hline
\end{tabular}
\end{center}

The bound is always correct and larger than the actual error, by a factor of
about $2.3$ to $2.6$: the remainder is roughly
$\frac{1}{(N+1)!}\left(1+\frac{1}{N+2}+\cdots\right)$, and $e$ appears because
its terms are summed at $x=1$. At $N=20$ the actual error is no longer measurable,
since $5.3\times10^{-20}$ is far below one unit in the last place of a quantity of
size $2.718$, namely about $4.4\times10^{-16}$.

The third group of numbers is three radii of convergence, computed in two
different ways at term $1000$.

\begin{center}
\small
\begin{tabular}{llrr}
\hline
series & $R$ from the formula & $|a_{1000}/a_{1001}|$ & $|a_{1000}|^{1/1000}$\\
\hline
$\sum x^n/n^2$          & $1$              & $1.0020010000$ & $0.9862794856$\\
$\sum \frac{n!}{n^n}x^n$ & $e=2.7182818285$ & $2.7169239322$ & $0.3694916635$\\
$\sum \frac{2^n}{n}x^n$  & $0.5$            & $0.5005000000$ & $1.9862320968$\\
\hline
\end{tabular}
\end{center}

The third column estimates $R$ from the ratio of two consecutive terms, the fourth from
the reciprocal of the root estimate. For the second series, $|a_n/a_{n+1}|=(1+1/n)^n$ rises to
$e$; at $n=1000$ its value is $2.7169239$, off by $0.0014$. For the first series,
$(1+1/n)^2$ decreases to $1$ from above, so the estimate is $1.0020$, and that error
$0.0020$ is a sign that we are only at term $1000$. The fourth column approaches
its limit more slowly, since $|a_n|^{1/n}$ approaches $1$ at rate $|\ln a_n|/n$; at
$n=1000$ the difference is $0.0016$, proportional to $\ln(2\pi n)/(2n)$ times $1/e$.

\kertas{Summary of numbers. The uniform error of $x^n$ on $[0,1]$ is $1$ for every
$n$, while the error at a fixed point $x$ is $x^n$: $132$ terms for $10^{-6}$
at $x=0.9$, $1375$ terms at $x=0.99$. The series for $e^x$ at $x=1$ reaches an error of
$2.7\times10^{-8}$ with $10$ terms, below the bound $6.8\times10^{-8}$.
The radii of the three series are $1$, $e$, and $0.5$; the integral
$\int_0^1 nx(1-x^2)^n dx=n/(2(n+1))$ equals $0.4545$ at $n=10$, while
the limit of the functions is zero.}

\section{Computational lab}
\tingkat{advanced}

The following code runs as it stands on Python 3.14.4 with NumPy 2.5.2 and
SciPy 1.18.0. It recomputes the three groups of numbers above, plus
quadrature for the counterexample on integration.

\begin{lab}
\begin{verbatim}
import math
import numpy as np
from scipy.integrate import quad

print("n      sup[0,1]   (1-1/n)^n   (1-1/sqrt n)^n     0.9^n")
for n in (10, 100, 1000, 10000):
    a = (1 - 1/n)**n
    b = (1 - 1/math.sqrt(n))**n
    print(f"{n:6d} {1.0:9.6f} {a:11.6f} {b:17.3e} {0.9**n:10.3e}")

print()
print("terms of ln(1+x) for error < 1e-6")
for x in (0.5, 0.9, 0.99, 0.999):
    N = 1
    while x**N / N > 1e-6:
        N += 1
    s = sum(((-1)**(k+1)) * x**k / k for k in range(1, N + 1))
    print(f"x={x:6.3f}  N={N:5d}  err={abs(math.log(1+x) - s):.3e}")

print()
print("N      S_N              e - S_N        e/(N+1)!")
for N in (5, 10, 15, 20):
    s = math.fsum(1/math.factorial(k) for k in range(N + 1))
    b = math.e / math.factorial(N + 1)
    print(f"{N:3d} {s:.12f} {math.e - s:14.3e} {b:13.3e}")

print()
print("radius from the ratio of terms 1000 and 1001 (log-space)")
la = lambda n: -2 * math.log(n)
lb = lambda n: math.lgamma(n + 1) - n * math.log(n)
lc = lambda n: n * math.log(2) - math.log(n)
for name, g in (("a", la), ("b", lb), ("c", lc)):
    R = math.exp(g(1000) - g(1001))
    H = math.exp(g(1000) / 1000)
    print(f"series {name}: R = {R!r}   |a_n|^(1/n) = {H!r}")
print("e =", repr(math.e))

print()
print("f_n(x) = n x (1-x^2)^n on [0,1]")
for n in (10, 100):
    val, _ = quad(lambda x, n=n: n * x * (1 - x*x)**n, 0, 1)
    print(f"n={n:4d}  integral={val!r}  n/(2(n+1))={n/(2*(n+1))!r}")
\end{verbatim}
\end{lab}

The output obtained:

\begin{verbatim}
n      sup[0,1]   (1-1/n)^n   (1-1/sqrt n)^n     0.9^n
    10  1.000000    0.348678         2.234e-02  3.487e-01
   100  1.000000    0.366032         2.656e-05  2.656e-05
  1000  1.000000    0.367695         1.108e-14  1.748e-46
 10000  1.000000    0.367861         2.249e-44  0.000e+00

terms of ln(1+x) for error < 1e-6
x= 0.500  N=   16  err=3.049e-07
x= 0.900  N=   89  err=4.478e-07
x= 0.990  N=  721  err=4.915e-07
x= 0.999  N= 5248  err=4.993e-07

N      S_N              e - S_N        e/(N+1)!
  5 2.716666666667      1.615e-03     3.775e-03
 10 2.718281801146      2.731e-08     6.810e-08
 15 2.718281828459      5.063e-14     1.299e-13
 20 2.718281828459      0.000e+00     5.320e-20

radius from the ratio of terms 1000 and 1001 (log-space)
series a: R = 1.0020010000000001   |a_n|^(1/n) = 0.9862794856312105
series b: R = 2.7169239322350776   |a_n|^(1/n) = 0.36949166347196
series c: R = 0.5004999999999973   |a_n|^(1/n) = 1.9862320968418674
e = 2.718281828459045

f_n(x) = n x (1-x^2)^n on [0,1]
n=  10  integral=0.4545454545454544  n/(2(n+1))=0.45454545454545453
n= 100  integral=0.4950495049504951  n/(2(n+1))=0.49504950495049505
\end{verbatim}

Four things must be read from this output. First, the second column stays at $1.000000$
for all four values of $n$, while the fifth column is already $3.5\times10^{-1}$ at
$n=10$: the failure of uniformity shows up as a number. Second, the last entry of the fifth
column prints $0.000e+00$ not because the error is zero, but because
$0.9^{10000}=10^{-457.6}$ lies below the smallest normal \texttt{float64} bound: underflow.

Third, the fourth column gives numerical evidence for two limits already
proved: $(1-1/n)^n\to1/e=0.3678794$ and $(1-1/\sqrt n)^n\to0$, which collapses
far faster.

Fourth, SciPy's quadrature gives $0.4545454545454544$ for $n=10$, whereas the
formula $n/(2(n+1))$ gives $0.45454545454545453$; a difference of one unit in the last place
of rounding. Without a closed form, we would have no way to make sure the quadrature is not
measuring something else.

One thing that does not show in the output: every bound in this chapter is an upper bound.
The row $N$ that the code reports is chosen by checking $x^N/N$, so that $N$ is
guaranteed to suffice, not to be minimal; the actual errors are $3.0\times10^{-7}$ to
$5.0\times10^{-7}$, below $10^{-6}$.

\begin{untukdosen}
The teaching order I suggest: start from the failure of the $\ln(1+x)$ library in
the opening, then the pointwise and uniform definitions with $x^n$ as the single
example. Do not rush to the $M$-test. Spend one full session making sure
the students can say precisely why $x^n$ fails to be uniform
and why $(1-1/\sqrt n)^n$ succeeds. After that the $M$-test, then the three exchange
theorems, each with its own counterexample.

Two misconceptions appear most often. First, students think ``every
$f_n$ continuous and $f_n\to f$'' already suffices to conclude that $f$ is continuous; they
forget that what is needed is uniformity, and $x^n$ is the cure. Second,
students think uniformity is inherited downward: if $f_n\to f$ uniformly then
$f_n'\to f'$. The example $\sin(nx)/n$ breaks that in one line; it is best
written on the board before the differentiation theorem appears.

A board activity that works: write four unnamed sequences of functions,
namely $x^n$, $x^n/n$, $nx(1-x^2)^n$, and $\sin(nx)/n$, then ask the class
to decide for each whether it converges pointwise, whether uniformly, and
whether its integral and derivative may be exchanged. The answers form a $4\times4$ table
that can be filled in twenty minutes. Students are ready to go on when they can
explain why a power series is uniform only on $[-r,r]$.
\end{untukdosen}

\section{Summary}
\tingkat{basic}

\begin{ringkasan}
\begin{itemize}
\item Pointwise convergence gives every point its own threshold;
      uniform convergence demands one threshold for all points at once,
      and that difference is equivalent to $\|f_n-f\|_\infty\to0$.
\item The Weierstrass $M$-test concludes uniformity from a single comparison
      series of numbers. It is a sufficient condition, not a necessary one; its failure proves
      nothing.
\item The uniform limit of continuous functions is continuous, and its integral may be exchanged with
      the limit. Differentiation demands uniformity of the sequence of derivatives plus convergence
      at one point.
\item Those three hypotheses cannot be removed: $x^n$, $nx(1-x^2)^n$, and
      $\sin(nx)/n$ each break one of them.
\item A power series converges absolutely inside its radius and diverges outside;
      inside the radius it may be differentiated and integrated term by term
      without changing its radius.
\item Uniformity is always local: $\ln(1+x)$ is uniform on every $[-r,r]$, $r<1$,
      but not on $(-1,1)$, and the price shows in the number of terms, from
      $16$ at $x=0.5$ to $5248$ at $x=0.999$.
\end{itemize}
\end{ringkasan}

\section{Exercises}
\tingkat{intermediate}

\begin{latihan}
\begin{enumerate}
\item Compute in Python, for $n=10,100,1000$, the value
      $\sup_{x\in[0,1]}|x^n|$ and the value of $x^n$ at $x=0.9$. Explain why
      the first column is flat and the second is not.

\item Prove: if $f_n\to f$ uniformly on $E$ and every $f_n$ is bounded on
      $E$, then $f$ is bounded on $E$ and the sequence $\|f_n\|_\infty$ is bounded.
      Show by example that pointwise convergence is not enough.

\item Suppose $f_n(x)=x/(1+nx^2)$ on $\mathbb{R}$. Show that
      $f_n\to0$ pointwise, determine $\|f_n\|_\infty$, and prove that
      the convergence is uniform. What is the smallest $n$ such that
      $\|f_n\|_\infty<0.01$?

\item Judge this claim: true, false, or not enough evidence. ``Since $x^n\to0$
      for every $x\in[0,1)$, the sequence $f_n(x)=x^n$ converges uniformly on
      $[0,1)$.'' State exactly where that reasoning breaks, and give one
      sequence of points that proves it.

\item Compute the error $|e-S_N(1)|$ for $N=5,10,15,20$ and compare it with
      the bound $e/(N+1)!$. What is the smallest $N$ for which that bound is below
      $10^{-16}$? Can the actual error still be measured at that $N$?

\item Prove that the series $\sum_{n\ge0}x^n(1-x)$ converges pointwise
      on $[0,1]$ but not uniformly, by computing its sum first
      and then checking its continuity.

\item Compute the radii of convergence of $\sum x^n/n^2$, $\sum(n!/n^n)x^n$, and
      $\sum 2^nx^n/n$ from the ratio of terms $1000$ and $1001$, then compare
      with the exact values. Which series has the slowest estimate, and
      why?

\item Prove that if $f_n\to f$ uniformly on $E$ and $g$ is bounded on $E$,
      then $f_ng\to fg$ uniformly on $E$. Show by example that the
      boundedness of $g$ cannot be removed.

\item For $f_n(x)=nx(1-x^2)^n$ on $[0,1]$, compute the peak point, the peak
      height, and $\int_0^1f_n$ for $n=10,100$. Compare the peak height
      with $\sqrt{n/(2e)}$ and explain why the area does not tend to zero.

\item Differentiate $\sum_{n\ge0}x^n=1/(1-x)$ term by term and use the result
      to compute $\sum_{n\ge1}n/2^n$. How many terms are needed for the
      partial sum to lie within $10^{-10}$ of the exact value?
\end{enumerate}

\noindent\emph{Hints for the three hardest items.} Item 4: the uniform property
demands one $N$ for all $x$, so it suffices to show that for every $n$
there is an $x$ with $x^n\ge\tfrac12$; take $x=2^{-1/n}$. Item 6: sum the geometric
series first; the result is a function that differs from its limit at one point,
and the uniform limit of continuous functions must be continuous. Item 8: write
$|f_ng-fg|=|f_n-f|\,|g|$ and use $\|g\|_\infty$; for the counterexample
take $g(x)=1/x$ on $(0,1)$ and an $f_n$ that decreases uniformly but not
fast enough.
\end{latihan}

\begin{untukpeneliti}
What is proved with complete proof in this chapter: the equivalence of uniformity with the
supremum norm, the uniform Cauchy criterion, the Weierstrass $M$-test, the three exchange
theorems together with their counterexamples, the radius formula through the limit superior, and
the invariance of the radius under term-by-term differentiation and integration. Two results are
quoted without proof: the non-differentiability of the Weierstrass function (1872, extended
by Hardy 1916) and Abel's theorem. The primary sources are Rudin, \emph{Principles of
Mathematical Analysis}, chapters ``Sequences and Series of Functions'' (7.1--7.17) and
``Some Special Functions'' (8.1--8.2); Abbott, \emph{Understanding Analysis}, chapter
``Sequences and Series of Functions''; Stewart, \emph{Calculus}, the section ``Power
Series''.

What is often popularized beyond its due is the idea that uniformity removes all
problems. It removes three particular problems, on the interval where it holds, and
no more. For Fourier series uniformity is in fact rare, and the exchange of
integral remains legitimate for other reasons. In neural-network approximation, what is used is
convergence on a finite set.

Open question: if $f_n$ is known only through its values on a finite grid, how large
is the error of estimating $\|f_n\|_\infty$ by the maximum on the grid? Without
regularity assumptions there is no general answer, and for narrow peaks a grid that looks
dense can miss them entirely. Limits of validity: all functions here are
real-valued on subsets of $\mathbb{R}$, with 64-bit arithmetic in
Python 3.14.4, NumPy 2.5.2, and SciPy 1.18.0.
\end{untukpeneliti}

\peta{Chapter 14 in Rudin, \emph{Principles of Mathematical Analysis}, chapters
``Sequences and Series of Functions'' (7.1--7.17) and ``Some Special Functions''
(8.1--8.2); in Abbott, \emph{Understanding Analysis}, chapter ``Sequences and
Series of Functions''; in Stewart, \emph{Calculus}, the sections ``Power Series'' and
``Taylor and Maclaurin Series''. Its basis is the second half of the chapter \emph{Integral, Series
of Functions, and Exchange of Limits} of the first edition of this book. This chapter uses the definitions of
limit and the Cauchy criterion from Chapter 4, the root test of Chapter 5, uniform continuity of Chapter 11,
and the integrability criterion of Chapter 13, all of them referred to, not repeated.}

\chapter{Multiple Integrals, Change of Variables, and Vector Field Theorems}
\chaptermark{Multiple Integrals and Vector Field Theorems}

A cooling plate is a disk of radius one meter. Its thickness is
uniform, but its density falls off from the center: $e^{-(x^{2}+y^{2})}$ kilograms
per square meter. Its mass is the sum of the density over the whole disk: a numerical library reports
$1.9858653038$ kilograms after $11907$ function evaluations.

A colleague writes a different version of the same computation. He moves to
polar coordinates and writes
\[
\int_{0}^{2\pi}\!\!\int_{0}^{1} e^{-r^{2}}\,dr\,d\theta ,
\]
and gets $4.6924344183$ kilograms. His program runs without a warning.
That number is $2.36$ times larger than the first, and not one of the
programs complains. The mistake is one letter: the factor $r$ is missing, so the area
of every small cell is counted too large.

Damage of that size also shows up on the other side of this chapter. Compute the work of
$F(x,y)=(-y,x)/(x^{2}+y^{2})$ around the unit circle: the result is
$6.2831853072$, even though the curl of that field is zero at every point where it
is defined. If zero curl meant zero closed work, one of those two computations
would lie. What lies is not the computation but the hypothesis that was left
unchecked: the field is not defined at the origin, and the origin lies inside
the circle.

Chapter 13 built the integral on an interval, and Chapter 12 gave the Jacobian as the
matrix of derivatives. Neither is enough yet: there is still no way to add up over a region
of dimension greater than one, no rule for what happens to the area of a cell
when the coordinates are changed, and no meaning for the sign on a boundary. This chapter
fills in all three, then shows that Green, Gauss, and Stokes are in fact one
statement.

\begin{untukanak}
Picture a tailor with a shirt pattern on squared paper. The pattern is square
and regular; the cloth to be cut curves to follow the body. If that square pattern
is pasted on as it is, the shirt is tight on one side and loose on the other,
even though the number of cells stays the same.

An experienced tailor does something else: every small cell will stretch
or shrink by a certain ratio as it is carried onto the cloth. Cells near the waist stretch
a little; cells near the shoulder shrink. He records that ratio, multiplies the area
of the cell by it, and only then adds up.

The same work appears with curved shapes: many objects, areas, or weights.
If the stretch ratio is forgotten, the total is still a tidy number, and
precisely because it is tidy it is hard to distrust. In this chapter we shall learn why
that ratio has to be there, and how to compute it for round shapes.
\end{untukanak}

\section{From rectangles to general regions}
\tingkat{basic}

We begin with the simplest object: a rectangle, and a
number we want to add up over its whole surface. The method is the same as in
Chapter 13, only the pieces are now cells rather than segments.

\begin{definisi}[The double integral]\index{double integral}
Let $R=[a,b]\times[c,d]$ be a closed rectangle. A
\emph{partition} $\mathcal{P}$ of $R$ is a division of $[a,b]$ into
intervals and of $[c,d]$ into intervals; the result is $mn$ small
rectangles $R_{ij}$. The \emph{mesh} of the partition is the longest diagonal among
those cells, and is written $\|\mathcal{P}\|$.

For $f$ bounded on $R$, write $M_{ij}=\sup_{R_{ij}}f$ and
$m_{ij}=\inf_{R_{ij}}f$, then
\[
U(\mathcal{P},f)=\sum_{i,j}M_{ij}\,\mathrm{luas}(R_{ij}),
\qquad
L(\mathcal{P},f)=\sum_{i,j}m_{ij}\,\mathrm{luas}(R_{ij}).
\]
The function $f$ is called \emph{integrable} on $R$ when the supremum of all $L$ equals
the infimum of all $U$; that common value is $\iint_{R}f\,dA$.
\end{definisi}

This definition says nothing about the order of integration. If $f$ is continuous, the upper
and lower sums meet because the largest gap of $f$ on a single cell
shrinks with the mesh: on a compact set, a continuous function is uniformly continuous
(Chapter 11).

The next question is practical. To add up over cells means computing
something two-dimensional, whereas our hands can only compute one-variable
integrals. The bridge: slice along one direction first, then add up
the results of those slices.

\begin{teorema}[Fubini]\index{Fubini's theorem}
If $f$ is continuous on $R=[a,b]\times[c,d]$, then
\[
\iint_{R} f\,dA=\int_{a}^{b}\!\!\int_{c}^{d} f(x,y)\,dy\,dx
=\int_{c}^{d}\!\!\int_{a}^{b} f(x,y)\,dx\,dy .
\]
\end{teorema}

\begin{proof}
A full proof demands matching the two-dimensional upper and lower sums with the one-dimensional
integral uniformly, and it is in Spivak, the chapter \emph{Integration},
Theorem 3-10. The idea runs like this.

For $f\ge 0$, the double integral is the volume of the solid under the graph of $f$.
Slice that solid with a plane of constant $x$. The slice is a region under the curve
$y\mapsto f(x,y)$, so its area is
\[
A(x)=\int_{c}^{d}f(x,y)\,dy .
\]
The volume of the solid is the integral of the area of the slice:
\[
\iint_{R}f\,dA=\int_{a}^{b}A(x)\,dx=\int_{a}^{b}\!\!\int_{c}^{d}f(x,y)\,dy\,dx .
\]
Because $f$ is continuous on a compact set, $A$ is continuous. For $f$ of any
sign, split $f=f^{+}-f^{-}$; the other direction comes from slices of constant $y$.
\end{proof}

What matters in this theorem is not the formula but the permission it grants: the order may
be swapped. That permission lapses for functions that are not integrable, and that is
discussed later.

The solid sliced in that proof appears in Figure~\ref{fig:kolom-tipis}:
each column contributes $f(x,y)\,\Delta A$.

\begin{bacakeras}
--- How much sand is in this pile?
--- Add up its thin columns. Stand on the north side, slice the pile into
strips, measure each strip, then add.
--- What do you get?
--- Eighteen buckets.
--- Now stand on the east side, slice in the other direction, measure, add.
--- What do you get?
--- Eighteen buckets as well.
--- The north-south strips and the east-west strips are not the same at all. Why
is the answer the same?
--- Because the pile is one. We are not adding two piles; we are slicing
the same object with two knives in different directions.
--- So the direction of the slice may be swapped?
--- As long as the object is measurable. If the pile has a hole in the middle, the two
directions still agree; what fails later is the hole, not the order.

Say it once more, in an ordinary voice: the direction of the slice changes, the volume stays.
\end{bacakeras}

\begin{figure}[htbp]
\centering
\begin{tikzpicture}[scale=1.55,font=\small]
  \fill[layergrey!10] (0,0) -- (6.1,0) -- (7.1,0.90) -- (1.0,0.90) -- cycle;
  \draw[layergrey,thick] (0,0) -- (6.1,0) -- (7.1,0.90) -- (1.0,0.90) -- cycle;
  \fill[layerteal!12] (0.91,0.82) rectangle (1.43,1.38);
  \draw[layerteal!70,line width=0.4pt] (0.91,0.82) rectangle (1.43,1.38);
  \fill[layerteal!30] (0.91,1.38) -- (1.08,1.47) -- (1.60,1.47) -- (1.43,1.38) -- cycle;
  \draw[layerteal!70,line width=0.4pt] (0.91,1.38) -- (1.08,1.47) -- (1.60,1.47);
  \fill[layerteal!12] (1.67,0.82) rectangle (2.19,1.60);
  \draw[layerteal!70,line width=0.4pt] (1.67,0.82) rectangle (2.19,1.60);
  \fill[layerteal!30] (1.67,1.60) -- (1.83,1.69) -- (2.35,1.69) -- (2.19,1.60) -- cycle;
  \draw[layerteal!70,line width=0.4pt] (1.67,1.60) -- (1.83,1.69) -- (2.35,1.69);
  \fill[layerteal!12] (2.42,0.82) rectangle (2.93,1.97);
  \draw[layerteal!70,line width=0.4pt] (2.42,0.82) rectangle (2.93,1.97);
  \fill[layerteal!30] (2.42,1.97) -- (2.58,2.06) -- (3.10,2.06) -- (2.93,1.97) -- cycle;
  \draw[layerteal!70,line width=0.4pt] (2.42,1.97) -- (2.58,2.06) -- (3.10,2.06);
  \fill[layerteal!12] (3.17,0.82) rectangle (3.68,2.30);
  \draw[layerteal!70,line width=0.4pt] (3.17,0.82) rectangle (3.68,2.30);
  \fill[layerteal!30] (3.17,2.30) -- (3.33,2.39) -- (3.85,2.39) -- (3.68,2.30) -- cycle;
  \draw[layerteal!70,line width=0.4pt] (3.17,2.30) -- (3.33,2.39) -- (3.85,2.39);
  \fill[layerteal!12] (3.92,0.82) rectangle (4.43,2.32);
  \draw[layerteal!70,line width=0.4pt] (3.92,0.82) rectangle (4.43,2.32);
  \fill[layerteal!30] (3.92,2.32) -- (4.08,2.41) -- (4.60,2.41) -- (4.43,2.32) -- cycle;
  \draw[layerteal!70,line width=0.4pt] (3.92,2.32) -- (4.08,2.41) -- (4.60,2.41);
  \fill[layerteal!12] (4.67,0.82) rectangle (5.18,2.00);
  \draw[layerteal!70,line width=0.4pt] (4.67,0.82) rectangle (5.18,2.00);
  \fill[layerteal!30] (4.67,2.00) -- (4.83,2.09) -- (5.35,2.09) -- (5.18,2.00) -- cycle;
  \draw[layerteal!70,line width=0.4pt] (4.67,2.00) -- (4.83,2.09) -- (5.35,2.09);
  \fill[layerteal!12] (5.42,0.82) rectangle (5.93,1.62);
  \draw[layerteal!70,line width=0.4pt] (5.42,0.82) rectangle (5.93,1.62);
  \fill[layerteal!30] (5.42,1.62) -- (5.58,1.71) -- (6.10,1.71) -- (5.93,1.62) -- cycle;
  \draw[layerteal!70,line width=0.4pt] (5.42,1.62) -- (5.58,1.71) -- (6.10,1.71);
  \fill[layerteal!12] (6.17,0.82) rectangle (6.68,1.39);
  \draw[layerteal!70,line width=0.4pt] (6.17,0.82) rectangle (6.68,1.39);
  \fill[layerteal!30] (6.17,1.39) -- (6.33,1.48) -- (6.85,1.48) -- (6.68,1.39) -- cycle;
  \draw[layerteal!70,line width=0.4pt] (6.17,1.39) -- (6.33,1.48) -- (6.85,1.48);
  \fill[layerteal!12] (0.59,0.50) rectangle (1.11,1.07);
  \draw[layerteal!70,line width=0.4pt] (0.59,0.50) rectangle (1.11,1.07);
  \fill[layerteal!30] (0.59,1.07) -- (0.76,1.16) -- (1.28,1.16) -- (1.11,1.07) -- cycle;
  \draw[layerteal!70,line width=0.4pt] (0.59,1.07) -- (0.76,1.16) -- (1.28,1.16);
  \fill[layerteal!12] (1.34,0.50) rectangle (1.86,1.33);
  \draw[layerteal!70,line width=0.4pt] (1.34,0.50) rectangle (1.86,1.33);
  \fill[layerteal!30] (1.34,1.33) -- (1.51,1.42) -- (2.03,1.42) -- (1.86,1.33) -- cycle;
  \draw[layerteal!70,line width=0.4pt] (1.34,1.33) -- (1.51,1.42) -- (2.03,1.42);
  \fill[layerteal!12] (2.09,0.50) rectangle (2.61,1.75);
  \draw[layerteal!70,line width=0.4pt] (2.09,0.50) rectangle (2.61,1.75);
  \fill[layerteal!30] (2.09,1.75) -- (2.26,1.84) -- (2.78,1.84) -- (2.61,1.75) -- cycle;
  \draw[layerteal!70,line width=0.4pt] (2.09,1.75) -- (2.26,1.84) -- (2.78,1.84);
  \fill[layerteal!12] (2.84,0.50) rectangle (3.36,2.14);
  \draw[layerteal!70,line width=0.4pt] (2.84,0.50) rectangle (3.36,2.14);
  \fill[layerteal!30] (2.84,2.14) -- (3.01,2.23) -- (3.53,2.23) -- (3.36,2.14) -- cycle;
  \draw[layerteal!70,line width=0.4pt] (2.84,2.14) -- (3.01,2.23) -- (3.53,2.23);
  \fill[layerteal!12] (3.59,0.50) rectangle (4.11,2.15);
  \draw[layerteal!70,line width=0.4pt] (3.59,0.50) rectangle (4.11,2.15);
  \fill[layerteal!30] (3.59,2.15) -- (3.76,2.24) -- (4.28,2.24) -- (4.11,2.15) -- cycle;
  \draw[layerteal!70,line width=0.4pt] (3.59,2.15) -- (3.76,2.24) -- (4.28,2.24);
  \fill[layerteal!12] (4.34,0.50) rectangle (4.86,1.79);
  \draw[layerteal!70,line width=0.4pt] (4.34,0.50) rectangle (4.86,1.79);
  \fill[layerteal!30] (4.34,1.79) -- (4.51,1.88) -- (5.03,1.88) -- (4.86,1.79) -- cycle;
  \draw[layerteal!70,line width=0.4pt] (4.34,1.79) -- (4.51,1.88) -- (5.03,1.88);
  \fill[layerteal!12] (5.09,0.50) rectangle (5.61,1.35);
  \draw[layerteal!70,line width=0.4pt] (5.09,0.50) rectangle (5.61,1.35);
  \fill[layerteal!30] (5.09,1.35) -- (5.26,1.44) -- (5.78,1.44) -- (5.61,1.35) -- cycle;
  \draw[layerteal!70,line width=0.4pt] (5.09,1.35) -- (5.26,1.44) -- (5.78,1.44);
  \fill[layerteal!12] (5.84,0.50) rectangle (6.36,1.08);
  \draw[layerteal!70,line width=0.4pt] (5.84,0.50) rectangle (6.36,1.08);
  \fill[layerteal!30] (5.84,1.08) -- (6.01,1.17) -- (6.53,1.17) -- (6.36,1.08) -- cycle;
  \draw[layerteal!70,line width=0.4pt] (5.84,1.08) -- (6.01,1.17) -- (6.53,1.17);
  \fill[layerteal!12] (0.26,0.17) rectangle (0.78,0.73);
  \draw[layerteal!70,line width=0.4pt] (0.26,0.17) rectangle (0.78,0.73);
  \fill[layerteal!30] (0.26,0.73) -- (0.43,0.82) -- (0.95,0.82) -- (0.78,0.73) -- cycle;
  \draw[layerteal!70,line width=0.4pt] (0.26,0.73) -- (0.43,0.82) -- (0.95,0.82);
  \fill[layerteal!12] (1.02,0.17) rectangle (1.54,0.95);
  \draw[layerteal!70,line width=0.4pt] (1.02,0.17) rectangle (1.54,0.95);
  \fill[layerteal!30] (1.02,0.95) -- (1.19,1.04) -- (1.71,1.04) -- (1.54,0.95) -- cycle;
  \draw[layerteal!70,line width=0.4pt] (1.02,0.95) -- (1.19,1.04) -- (1.71,1.04);
  \fill[layerteal!12] (1.76,0.17) rectangle (2.29,1.32);
  \draw[layerteal!70,line width=0.4pt] (1.76,0.17) rectangle (2.29,1.32);
  \fill[layerteal!30] (1.76,1.32) -- (1.93,1.41) -- (2.46,1.41) -- (2.29,1.32) -- cycle;
  \draw[layerteal!70,line width=0.4pt] (1.76,1.32) -- (1.93,1.41) -- (2.46,1.41);
  \fill[layerteal!12] (2.51,0.17) rectangle (3.04,1.65);
  \draw[layerteal!70,line width=0.4pt] (2.51,0.17) rectangle (3.04,1.65);
  \fill[layerteal!30] (2.51,1.65) -- (2.69,1.74) -- (3.21,1.74) -- (3.04,1.65) -- cycle;
  \draw[layerteal!70,line width=0.4pt] (2.51,1.65) -- (2.69,1.74) -- (3.21,1.74);
  \fill[layerteal!12] (3.26,0.17) rectangle (3.79,1.67);
  \draw[layerteal!70,line width=0.4pt] (3.26,0.17) rectangle (3.79,1.67);
  \fill[layerteal!30] (3.26,1.67) -- (3.44,1.76) -- (3.96,1.76) -- (3.79,1.67) -- cycle;
  \draw[layerteal!70,line width=0.4pt] (3.26,1.67) -- (3.44,1.76) -- (3.96,1.76);
  \fill[layerteal!12] (4.01,0.17) rectangle (4.53,1.35);
  \draw[layerteal!70,line width=0.4pt] (4.01,0.17) rectangle (4.53,1.35);
  \fill[layerteal!30] (4.01,1.35) -- (4.18,1.44) -- (4.70,1.44) -- (4.53,1.35) -- cycle;
  \draw[layerteal!70,line width=0.4pt] (4.01,1.35) -- (4.18,1.44) -- (4.70,1.44);
  \fill[layerteal!12] (4.76,0.17) rectangle (5.28,0.97);
  \draw[layerteal!70,line width=0.4pt] (4.76,0.17) rectangle (5.28,0.97);
  \fill[layerteal!30] (4.76,0.97) -- (4.93,1.06) -- (5.45,1.06) -- (5.28,0.97) -- cycle;
  \draw[layerteal!70,line width=0.4pt] (4.76,0.97) -- (4.93,1.06) -- (5.45,1.06);
  \fill[layerteal!12] (5.51,0.17) rectangle (6.03,0.74);
  \draw[layerteal!70,line width=0.4pt] (5.51,0.17) rectangle (6.03,0.74);
  \fill[layerteal!30] (5.51,0.74) -- (5.68,0.83) -- (6.20,0.83) -- (6.03,0.74) -- cycle;
  \draw[layerteal!70,line width=0.4pt] (5.51,0.74) -- (5.68,0.83) -- (6.20,0.83);
  \fill[accent!42] (2.69,0.17) rectangle (2.86,1.65);
  \draw[accent,line width=1pt] (2.69,0.17) rectangle (2.86,1.65);
  \fill[accent!55] (2.69,1.65) -- (2.78,1.70) -- (2.96,1.70) -- (2.86,1.65) -- cycle;
  \draw[accent,line width=1pt] (2.69,1.65) -- (2.78,1.70) -- (2.96,1.70);
  \draw[accent,densely dashed] (2.86,0.16) -- (3.42,-0.26);
  \node[accent,anchor=north west,align=left]
    at (3.50,-0.30) {one thin column: base $\Delta A$,\\ height $f(x,y)$, volume $\approx f(x,y)\,\Delta A$};
  \node[layergrey,anchor=north west] at (-0.05,-0.08) {base $R$};
  \node[layergrey,anchor=north] at (5.95,-0.02) {$x$};
  \node[layergrey,anchor=west] at (7.15,0.55) {$y$};
\end{tikzpicture}
\caption{The solid under the graph of $f$ over a rectangle $R$, cut into
thin columns. Each column is almost a box with base $\Delta A$ and height
$f(x,y)$, so its volume is $f(x,y)\,\Delta A$; adding up all the columns gives
$\iint_R f\,dA$. The column marked in red is one term in that sum.
Slicing along the $x$ direction and along the $y$ direction gives different columns, but
the total volume is the same.}
\label{fig:kolom-tipis}
\end{figure}

The regions that appear in problems are rarely rectangles. We handle them
with one extra step: extend the function by zero outside the region.

\begin{lemma}[A region between two graphs]\index{type I region}
Let $g_{1},g_{2}\colon[a,b]\to\mathbb{R}$ be continuous with $g_{1}\le
g_{2}$, and
\[
D=\{(x,y): a\le x\le b,\ g_{1}(x)\le y\le g_{2}(x)\}.
\]
If $f$ is continuous on $D$, then
\[
\iint_{D}f\,dA=\int_{a}^{b}\!\!\int_{g_{1}(x)}^{g_{2}(x)}f(x,y)\,dy\,dx .
\]
\end{lemma}

\begin{proof}
Take a rectangle $R$ that contains $D$, and extend $f$ by the value $0$ on
$R\setminus D$. The boundary of $D$ can be covered by finitely many rectangles of
total area less than $\varepsilon$: on the graph, uniform continuity on a
compact interval gives a partition whose oscillation of $g$ is small on each piece.
A set that can be covered that way does not change the value of the integral, so
\[
\iint_{D}f\,dA=\iint_{R} f\cdot\mathbf{1}_{D}\,dA .
\]
By Fubini, the integral on the right is computed by integrating $x$ first;
for fixed $x$, the integrand $\mathbf{1}_{D}f$ equals $f$ only on
$g_{1}(x)\le y\le g_{2}(x)$.
\end{proof}

A type I region is enough for almost every problem in this chapter; a boundary that crosses a
vertical line more than twice is split and then added up. The most important case is the disk:
in Cartesian coordinates it demands two pieces, in polar only one.

\begin{untukanak}
Now take a closed room with several stoves inside it.
Each stove gives off heat. The question: how much heat leaves that room
every second, added up over all the walls, the floor, and the roof?

The first way to compute it is hard: walk around every wall, measure the heat
flow at each point, then add up. The second way is much easier: count
how much heat all the stoves inside produce. If no heat escapes
through other gaps, the two numbers must agree: the amount inside equals the
amount through the skin.

That sentence sounds too good to be true. The agreement holds only if
the skin of the room is really closed and the direction ``out'' is fixed unambiguously. If one small
hole is not counted, or if we mislabel which side is outside and
which is inside, the numbers miss even though all our additions are correct.

In this chapter we shall learn why the amount inside can be read off from the skin,
for any curved shape, and how one small hole can ruin
the whole conclusion.
\end{untukanak}

\section{Change of variables and the Jacobian determinant}
\tingkat{intermediate}

Changing coordinates means changing cells: a square cell in the $(u,v)$ plane
becomes a curved cell in the $(x,y)$ plane, and its area changes. We need
one number at each point.

\begin{definisi}[Change of variables]
\index{change of variables}\index{Jacobian} Let $D^{*}$ be an open region in
$\mathbb{R}^{2}$ and $T\colon D^{*}\to\mathbb{R}^{2}$,
$T(u,v)=(x(u,v),y(u,v))$, a function with continuous partial derivatives. $T$ is called a
\emph{valid change of variables} if $T$ is one-to-one on the interior of
$D^{*}$ and its Jacobian determinant
\[
\det J_{T}(u,v)=\frac{\partial x}{\partial u}\frac{\partial y}{\partial v}
-\frac{\partial x}{\partial v}\frac{\partial y}{\partial u}
\]
is nonzero there. The matrix $J_{T}$ itself we already met in Chapter 12.
\end{definisi}

In everyday words: it is the number that tells how many times a small cell
grows when it is carried by $T$. A value of $3$ means the cell becomes three times larger in area, $1/2$ means it shrinks
by half, and zero means the cell is flattened into a line and its information is lost.

Figure~\ref{fig:juring-jacobian} shows one cell and its image:
area $0.04$ becomes $0.064$, and the number $1.6$ is $|{\det J}|$ at the center
of the cell.

\begin{figure}[htbp]
\centering
\begin{tikzpicture}[font=\footnotesize]
  \begin{scope}[x=1.9cm,y=1.9cm]
    \fill[layergrey!7] (1,0) rectangle (2.4,1.2);
    \draw[layergrey!40,very thin] (1,0) grid[xstep=0.2,ystep=0.2] (2.4,1.2);
    \draw[->,ink] (0.88,0) -- (2.62,0) node[right] {$u$};
    \draw[->,ink] (1,-0.08) -- (1,1.42) node[above] {$v$};
    \fill[layerindigo!30] (1.5,0.6) rectangle (1.7,0.8);
    \draw[layerindigo,thick] (1.5,0.6) rectangle (1.7,0.8);
    \node[layerindigo!85!black,anchor=north,align=center,fill=white,inner sep=1.2pt]
      at (1.60,0.55) {cell $0.2\times0.2$\\ area $0.04$};
    \node[layergrey,anchor=north] at (1.7,-0.10) {the $(u,v)$ plane};
  \end{scope}
  \draw[-{Latex[length=2.4mm]},ink,thick] (4.62,1.14) -- (5.85,1.14);
  \node[ink,above] at (5.23,1.18) {$T$};
  \node[ink,below,font=\scriptsize] at (5.23,1.06) {$\times\,1.6$};
  \begin{scope}[xshift=5.9cm,x=1.9cm,y=1.9cm]
    \foreach \r in {1,1.2,1.4,1.6,1.8,2,2.2,2.4}{
      \draw[layergrey!40,very thin] (\r,0) arc (0:69:\r);}
    \foreach \a in {0,11.5,23,34.5,46,57.5,69}{
      \draw[layergrey!40,very thin] (0,0) -- (\a:2.4);}
    \draw[->,ink] (-0.08,0) -- (2.62,0) node[right] {$x$};
    \draw[->,ink] (0,-0.08) -- (0,2.52) node[above] {$y$};
    \fill[accent!30] (34.4:1.5) arc (34.4:45.8:1.5) -- (45.8:1.7)
      arc (45.8:34.4:1.7) -- cycle;
    \draw[accent,thick] (34.4:1.5) arc (34.4:45.8:1.5) -- (45.8:1.7)
      arc (45.8:34.4:1.7) -- cycle;
    \node[accent,anchor=west,align=left] at (1.62,1.42)
      {curved sector\\ area $0.064=1.6\times0.04$};
    \node[layergrey,anchor=north] at (1.6,-0.10) {the $(x,y)$ plane};
  \end{scope}
\end{tikzpicture}
\caption{One small cell in the $(u,v)$ plane and its image in the $(x,y)$ plane under
$T(u,v)=(u\cos v,\;u\sin v)$. The cell $0.2\times0.2$ of area $0.04$ becomes a sector of
area $0.064$, that is $1.6$ times as large, exactly the value $|{\det J}|=u$ at the center of the cell
($u=1.6$). Note that its sides curve as well: what changes is not only the size
but also the shape, and one number at each point is enough to record the change.}
\label{fig:juring-jacobian}
\end{figure}

\begin{tangan}
Take a square sheet of paper $12$ cm by $12$ cm, a ruler, a pencil, scissors, and
tape. Draw $1$ cm cells over the whole sheet; there are $144$ cells. Cut the paper
into twelve vertical strips, each $1$ cm by $12$ cm. On every
strip, mark the point $3$ cm to the right on the top edge, draw a line from the bottom left
corner to that point, then cut along the line. Each strip is now two slanted
pieces. Arrange all the pieces into one tight parallelogram, then
tape them together.

Watch what happens: the outer frame is $15$ cm by $12$ cm, wider by $3$ cm
than the original sheet, but the parallelogram inside still has base
$12$ cm and height $12$ cm. Measure both with the ruler, perpendicular, not
along the slanted side: the result is $144$ square centimeters, the same as the area of the sheet
before it was cut. If your number is different, the height is almost certainly what went wrong.
\end{tangan}

\begin{teorema}[Change of variables in a double integral]
\index{change of variables theorem} Let $T$ be a valid change of variables
from $D^{*}$ to $D=T(D^{*})$, and $f$ continuous on $D$. Then
\[
\iint_{D}f(x,y)\,dx\,dy
=\iint_{D^{*}} f\bigl(T(u,v)\bigr)\,\bigl|\det J_{T}(u,v)\bigr|\,du\,dv .
\]
\end{teorema}

\begin{proof}
We prove the idea; the full proof is in Spivak, the chapter \emph{Integration},
the section on change of variables, and in Rudin, the chapter \emph{Integration of Differential
Forms}.

Partition $D^{*}$ into small cells. On a cell around $(u_{0},v_{0})$,
the map $T$ is almost linear: $T(u_{0}+h,v_{0}+k)\approx
T(u_{0},v_{0})+J_{T}(u_{0},v_{0})(h,k)$. That cell is mapped onto a
parallelogram with sides $h\,\partial T/\partial u$ and $k\,\partial T/\partial v$,
whose area is computed with a determinant (Chapter 7):
\[
\mathrm{luas}
=\Bigl|h\,k\,\det\Bigl(\frac{\partial T}{\partial u},
\frac{\partial T}{\partial v}\Bigr)\Bigr|
=|h\,k|\,|\det J_{T}(u_{0},v_{0})| .
\]
Each cell contributes $f(T(u_{0},v_{0}))$ times its original area times
$|\det J_{T}|$, and the sum is the Riemann sum for the integral on the right-hand
side; because the derivative of $T$ is continuous on a compact set, the error of the linear
approximation shrinks uniformly with the mesh.
\end{proof}

The absolute value on the determinant is not decoration: a determinant can be negative, an area cannot.
The negative sign stores an orientation; in the double integral that orientation is discarded, but
in the line and surface integrals later that sign decides the answer.

Now the determinants for the two coordinate systems used most often,
beginning with polar.

\begin{proposisi}[Polar coordinates]\index{polar coordinates}
With $x=r\cos\theta$, $y=r\sin\theta$, $r>0$, we have
\[
\det J=\begin{vmatrix}\cos\theta & -r\sin\theta\\ \sin\theta & r\cos\theta
\end{vmatrix}=r\cos^{2}\theta+r\sin^{2}\theta=r ,
\]
so that
$\iint_{D}f\,dx\,dy=\iint_{D^{*}}f(r\cos\theta,r\sin\theta)\,r\,dr\,d\theta$.
\end{proposisi}

That factor $r$ is the one missing from the program in the opening of the chapter. At $r=1$
it is $1$, but at $r=0.1$ it is $0.1$: a cell near the center
contributes ten times less than a guess that ignores it, and
the total error can be more than a factor of two.

\begin{proposisi}[Cylindrical and spherical coordinates]
\index{cylindrical coordinates}\index{spherical coordinates} For cylindrical coordinates
$x=r\cos\theta$, $y=r\sin\theta$, $z=z$, the Jacobian determinant is $r$. For
spherical coordinates
\[
x=\rho\sin\varphi\cos\theta,\quad
y=\rho\sin\varphi\sin\theta,\quad
z=\rho\cos\varphi ,
\]
with $\rho\ge0$, $0\le\varphi\le\pi$, $0\le\theta<2\pi$, the Jacobian
determinant is $\rho^{2}\sin\varphi$.
\end{proposisi}

\begin{proof}
For the cylinder, the third column $(0,0,1)$ leaves the determinant as the $2\times2$
polar block determinant, namely $r$.

For the sphere, compute the $3\times3$ determinant by cofactor expansion. With
$s_{\varphi}=\sin\varphi$, $c_{\varphi}=\cos\varphi$,
$s_{\theta}=\sin\theta$, $c_{\theta}=\cos\theta$, the matrix is
\[
J=\begin{pmatrix}
s_{\varphi}c_{\theta} & \rho c_{\varphi}c_{\theta} & -\rho s_{\varphi}s_{\theta}\\
s_{\varphi}s_{\theta} & \rho c_{\varphi}s_{\theta} & \rho s_{\varphi}c_{\theta}\\
c_{\varphi} & -\rho s_{\varphi} & 0
\end{pmatrix}.
\]
The entry in the third row and third column is zero, so expansion along the third column leaves two
terms:
\[
\det J=\rho^{2}s_{\varphi}s_{\theta}^{2}+\rho^{2}s_{\varphi}c_{\theta}^{2}
=\rho^{2}s_{\varphi}\bigl(s_{\theta}^{2}+c_{\theta}^{2}\bigr)=\rho^{2}s_{\varphi}.
\]
\end{proof}

The factor $\rho^{2}\sin\varphi$ is zero at the poles $\varphi=0$ and $\varphi=\pi$,
because there the whole sphere of radius $\rho$ collapses to a single point and the map stops being
one-to-one; that is why the poles are excluded from $D^{*}$. It is largest at the
equator, where a small change in $\varphi$ carries a point farthest.

\begin{resep}
\begin{enumerate}
\item Write $T$ explicitly: $x=x(u,v)$, $y=y(u,v)$, and check that $T$
  is one-to-one on the interior of $D^{*}$.
\item Compute the four partial derivatives, assemble the $2\times2$ matrix, then compute
  $\det J(u,v)$. Keep its sign; the absolute value is put in later.
\item Translate the boundary: draw $D$ in the $(x,y)$ plane, ask which values of $(u,v)$
  land there, then write $D^{*}$.
\item Replace every $x,y$ in $f$ by $x(u,v),y(u,v)$, multiply by
  $|\det J(u,v)|$, then close with $du\,dv$ in the order of the limits.
\end{enumerate}
A warning about order: steps 2 and 3 may not be swapped. The determinant must already
be in terms of $(u,v)$ before the limits are put in. If the limits are written first, that factor is left in $x$ and $y$ and
the integral cannot be computed.
\end{resep}

\kamus{Jacobian determinant}{determinan Jacobian}{The number that says how many times a small cell stretches when the coordinates are changed.}

\kamus{double integral}{integral lipat dua}{Adding up values on small cells of a surface, then narrowing the cells.}

\kamus{orientation}{orientasi}{An agreement about which direction of rotation counts as positive.}

\section{Line integrals and conservative fields}
\tingkat{intermediate}

Now we stop adding up over a surface and start adding up
along a path. Picture a small object dragged along a curved path by a force whose
direction and size change at every point. The work of that force is the sum of the component of the force along the path times the piece of
length of the path.

\begin{definisi}[Line integral]\index{line integral}
Let $\gamma\colon[a,b]\to\mathbb{R}^{2}$ be a curve with piecewise continuous
derivative, and $F=(P,Q)$ a continuous vector field on its trace.
The \emph{line integral} of $F$ along $\gamma$ is
\[
\int_{\gamma}F\cdot d\mathbf{r}
=\int_{a}^{b}\Bigl(P(\gamma(t))\,\gamma_{1}'(t)
+Q(\gamma(t))\,\gamma_{2}'(t)\Bigr)\,dt .
\]
If $\gamma$ is closed, meaning $\gamma(a)=\gamma(b)$, that integral is called
the \emph{circulation} of $F$ around $\gamma$.
\end{definisi}

That formula does not choose a starting point, only a direction. Changing the parameter to a
different speed does not change the value, because substitution (Chapter 13) absorbs
the speed factor; reversing the direction changes the sign. So a line integral is oriented,
and that is its first difference from the double integral.

Figure~\ref{fig:sirkulasi} shows such a rotating field, with one
closed curve for measuring how much rotation is caught.

\begin{figure}[htbp]
\centering
\begin{tikzpicture}[scale=1.45,font=\footnotesize]
  \draw[->,layergrey] (-2.85,0) -- (2.85,0) node[right] {$x$};
  \draw[->,layergrey] (0,-2.85) -- (0,2.85) node[above] {$y$};
  \foreach \x in {-2.4,-1.6,-0.8,0,0.8,1.6,2.4}{
    \foreach \y in {-2.4,-1.6,-0.8,0,0.8,1.6,2.4}{
      \draw[-{Latex[length=1.2mm]},layerteal!75]
        (\x,\y) -- ({\x-0.30*\y},{\y+0.30*\x});}}
  \draw[accent,very thick] (0,0) circle (1.75);
  \draw[-{Latex[length=2.8mm]},accent,very thick]
    (1.126,-1.341) -- (1.340,-1.126);
  \node[accent,anchor=north west,align=left] at (0.28,-1.86)
    {closed curve, positive direction};
  \node[ink,align=center,fill=white,fill opacity=0.85,text opacity=1,inner sep=2pt] at (0,0.95)
    {$\displaystyle\oint_{\gamma}F\cdot d\mathbf{r}=2\pi$};
  \node[ink,font=\scriptsize,fill=white,fill opacity=0.85,text opacity=1,inner sep=1.5pt] at (0,0.42)
    {whereas $\operatorname{curl}F=2$ at every point};
\end{tikzpicture}
\caption{The field $F(x,y)=(-y,x)$ with an arrow at every grid point. The arrows turn
around the origin, and their length grows with the distance of the point from the center. The circulation
around a circle of radius $1$ is $2\pi$, equal to the area $\pi$ times
$\operatorname{curl}F=2$: this is Green for this field. The direction of the arrows on the curve is the positive
direction, counterclockwise.}
\label{fig:sirkulasi}
\end{figure}

If a force field comes from a height function, the work depends only
on the starting point and the ending point. This is the curve version of the fundamental theorem of calculus,
and its name follows from that.

\begin{teorema}[The fundamental theorem for line integrals]
\index{conservative field} Let $f\colon U\to\mathbb{R}$ have continuous partial
derivatives on an open set $U\subseteq\mathbb{R}^{2}$, and let $\gamma$
be a curve from $A$ to $B$ inside $U$. Then
\[
\int_{\gamma}\nabla f\cdot d\mathbf{r}=f(B)-f(A).
\]
\end{teorema}

\begin{proof}
Write $g(t)=f(\gamma(t))$. By the multivariable chain rule (Chapter 12),
\[
g'(t)=\nabla f(\gamma(t))\cdot\gamma'(t),
\]
provided $\gamma$ is differentiable at $t$. So the integrand in the line integral
is the derivative of $g$, and the fundamental theorem of calculus (Chapter 13) gives
$\int_{a}^{b}g'(t)\,dt=g(b)-g(a)=f(B)-f(A)$. For a curve that is only
piecewise differentiable, add up over each piece; the terms at the joining
points cancel each other.
\end{proof}

A field of the form $\nabla f$ is called \emph{conservative}, and $f$ is its potential.
A direct consequence of the theorem: for a conservative field, the circulation around
a closed path is always zero.

The converse is also true.

\begin{teorema}[Converse]
Let $F$ be continuous on a connected open set
$U\subseteq\mathbb{R}^{2}$, and let the line integral of $F$ be independent of the
path inside $U$. Then $F=\nabla f$ for some $f$, namely
\[
f(\mathbf{x})=\int_{\mathbf{x}_{0}}^{\mathbf{x}}F\cdot d\mathbf{r},
\]
with $\mathbf{x}_{0}\in U$ fixed.
\end{teorema}

\begin{proof}
Take $\mathbf{x}\in U$ and $h$ small enough that the segment from $\mathbf{x}$ to
$\mathbf{x}+h\mathbf{e}_{1}$ stays inside $U$. Because the path does not
matter, we may travel from $\mathbf{x}_{0}$ to $\mathbf{x}$ and then
add that short segment, so that
\[
f(\mathbf{x}+h\mathbf{e}_{1})-f(\mathbf{x})
=\int_{0}^{1}F_{1}(\mathbf{x}+th\mathbf{e}_{1})\,h\,dt .
\]
Divide by $h$ and use continuity: the average value of the integrand tends to
$F_{1}(\mathbf{x})$ as $h\to0$, so $\partial f/\partial x_{1}=F_{1}$,
and the same argument gives $\partial f/\partial x_{2}=F_{2}$. Both partial
derivatives are continuous, hence $F=\nabla f$.
\end{proof}

Given $F$, how do we know it is conservative without looking for a potential? In
the plane there is a quick test: $F=(P,Q)$ has a potential only if $\partial P/\partial
y=\partial Q/\partial x$, because both sides are second derivatives of $f$ and Clairaut
(Chapter 12) makes them equal. The other direction demands a domain without holes and is
proved through Green later. The field from the opening of the chapter is the counterexample: its curl is zero,
but it is not the gradient of anything.

\begin{contoh}[Two ways to compute work]
\index{potential} Take $F(x,y)=(2xy,\ x^{2}+3y^{2})$ and the path
$\gamma(t)=(t,2t^{2})$ for $0\le t\le1$, from $(0,0)$ to $(1,2)$.

First way, directly. Here $dx=dt$ and $dy=4t\,dt$, so
\[
\int_{\gamma}F\cdot d\mathbf{r}
=\int_{0}^{1}\Bigl(2t\cdot 2t^{2}+(t^{2}+12t^{4})\cdot 4t\Bigr)dt
=\int_{0}^{1}\bigl(8t^{3}+48t^{5}\bigr)dt=2+8=10 .
\]
Second way, through the potential. First check that $\partial P/\partial
y=2x=\partial Q/\partial x$, then guess $f(x,y)=x^{2}y+y^{3}$. Its derivatives
are indeed $\partial f/\partial x=2xy$ and $\partial f/\partial y=x^{2}+3y^{2}$.
Therefore
\[
\int_{\gamma}F\cdot d\mathbf{r}=f(1,2)-f(0,0)=(2+8)-0=10 .
\]
Both ways give $10$; adaptive quadrature differs by $1.78\times10^{-15}$.
\end{contoh}

\section{Green's theorem}
\tingkat{intermediate}

The idea in the box about the room and the stoves now becomes a theorem: what happens
inside a region can be read off from its edge.

Before that, one agreement on direction.

\begin{definisi}[Positive orientation]
\index{positive orientation} A simple closed curve $\gamma$ is oriented
\emph{positively} if the region it bounds always lies on the left
as we traverse it. For a circle, the positive direction is counterclockwise.
\end{definisi}

\begin{teorema}[Green]\index{Green's theorem}
Let $D$ be a closed bounded region in $\mathbb{R}^{2}$ whose boundary
$\partial D$ consists of finitely many simple closed curves that are piecewise
smooth and positively oriented. Let $P,Q$ have continuous partial
derivatives on an open set containing $D$. Then
\[
\oint_{\partial D}P\,dx+Q\,dy=\iint_{D}\Bigl(\frac{\partial Q}{\partial x}
-\frac{\partial P}{\partial y}\Bigr)dA .
\]
\end{teorema}

\begin{proof}
It is enough to prove it for regions that are both type I and type II, because a general
region is split into finitely many pieces of that kind and the integrals are
added up; the adjacent sides of the pieces are traversed in opposite directions and therefore
cancel each other.

Suppose $D$ is type I with $g_{1},g_{2}$ smooth. On the vertical sides $dx=0$,
so their contribution is zero and
\[
\oint_{\partial D}P\,dx=\int_{a}^{b}P(x,g_{1}(x))\,dx+\int_{b}^{a}P(x,g_{2}(x))\,dx
=\int_{a}^{b}\bigl[P(x,g_{1})-P(x,g_{2})\bigr]dx .
\]
By the fundamental theorem of calculus,
$P(x,g_{1})-P(x,g_{2})=-\int_{g_{1}(x)}^{g_{2}(x)}\partial P/\partial y\,dy$,
so that
\[
\oint_{\partial D}P\,dx=-\int_{a}^{b}\!\!\int_{g_{1}(x)}^{g_{2}(x)}
\frac{\partial P}{\partial y}\,dy\,dx=-\iint_{D}\frac{\partial P}{\partial y}\,dA .
\]
In the same way, on a $D$ of type II, where $dy=0$ on the flat sides,
$\oint_{\partial D}Q\,dy=\iint_{D}\partial Q/\partial x\,dA$.
\end{proof}

\begin{contoh}[Green on the unit square]
Take $P=x^{2}-y$, $Q=x+2y$, and $D=[0,1]\times[0,1]$. Here $\partial
Q/\partial x-\partial P/\partial y=1-(-1)=2$, so the right-hand side of Green is
$2$ times the area of the square, namely $2$.

We compute the left-hand side on the four sides, counterclockwise:
\begin{align*}
\text{bottom }(y=0,\,dy=0):&\quad \int_{0}^{1}x^{2}\,dx=\tfrac13 ,\\
\text{right }(x=1,\,dx=0):&\quad \int_{0}^{1}(1+2y)\,dy=2 ,\\
\text{top }(y=1,\,dy=0):&\quad \int_{1}^{0}(x^{2}-1)\,dx=\tfrac23 ,\\
\text{left }(x=0,\,dx=0):&\quad \int_{1}^{0}2y\,dy=-1 .
\end{align*}
The sum is $\tfrac13+2+\tfrac23-1=2$, the same as the right-hand side. Adaptive quadrature
for the four sides gives $2.000000000000000$, and the double integral also gives
$2.000000000000000$.
\end{contoh}

Now the sign error we asked for. Repeat the same computation clockwise. Every side is traversed
backwards, so every term changes sign and
the sum becomes $-2$; the right-hand side of Green does not change because the region is the same.
The answer that should be $2$ becomes $-2$: not slightly off, but wrong in
sign entirely, with a difference of $4$ that cannot be blamed on
rounding.

So orientation is not a polite convention: Green's formula is correct only for one
direction of rotation, and that direction must be checked before computing.

With Green in hand, the converse direction of the curl test from the previous section
can be proved.

\begin{akibat}[The curl test on a domain without holes]\index{curl test}
Let $U\subseteq\mathbb{R}^{2}$ be open and \emph{simply connected},
meaning that every closed curve in $U$ can be contracted to a single point inside
$U$. Let $P,Q$ have continuous partial derivatives on $U$ and $\partial
P/\partial y=\partial Q/\partial x$ throughout $U$. Then $F=(P,Q)$ is conservative
on $U$.
\end{akibat}

\begin{proof}
Take a simple closed curve $\gamma$ in $U$ and the region $D$ that
it bounds. Because $U$ is simply connected, $D$ stays inside $U$, so
Green applies and
\begin{align*}
\oint_{\gamma}P\,dx+Q\,dy
&=\iint_{D}\Bigl(\frac{\partial Q}{\partial x}-\frac{\partial P}{\partial y}\Bigr)dA\\
&=0 .
\end{align*}
A closed curve that crosses itself is split into simple arcs,
and each piece gives zero. All circulations are zero, so the line
integral is independent of path, and the converse theorem above supplies a potential.
\end{proof}

That hypothesis cannot be dropped. On a region with a hole, Green demands that the inner boundary
also be counted; that is where the field from the opening of the chapter escapes the curl test.

\section{Gauss's and Stokes' theorems}
\tingkat{advanced}

Green talks about a region in the plane and its edge. In three-dimensional space
there are two versions, and both use the divergence and the curl from Chapter 12.

\begin{definisi}[Divergence and curl]
\index{divergence}\index{curl} For a field $F=(F_{1},F_{2},F_{3})$ with
continuous partial derivatives,
\[
\operatorname{div}F=\frac{\partial F_{1}}{\partial x}
+\frac{\partial F_{2}}{\partial y}+\frac{\partial F_{3}}{\partial z},
\qquad
\operatorname{curl}F=\Bigl(\frac{\partial F_{3}}{\partial y}
-\frac{\partial F_{2}}{\partial z},\
\frac{\partial F_{1}}{\partial z}-\frac{\partial F_{3}}{\partial x},\
\frac{\partial F_{2}}{\partial x}-\frac{\partial F_{1}}{\partial y}\Bigr).
\]
\end{definisi}

The divergence measures how much field is born at a point: positive means a
source, negative a sink. The curl measures the rotation of the field around that point.

\begin{teorema}[Gauss, or the divergence theorem]\index{Gauss's theorem}
Let $E\subseteq\mathbb{R}^{3}$ be a closed bounded region whose boundary
$\partial E$ is a piecewise smooth surface, and let $\mathbf{n}$ be the unit
normal vector pointing outward. If $F$ has continuous partial derivatives
on an open set containing $E$, then
\[
\iiint_{E}\operatorname{div}F\,dV=\iint_{\partial E}F\cdot\mathbf{n}\,dA .
\]
\end{teorema}

\begin{proof}
We prove it for a box first.

Take $E=[a,b]\times[c,d]\times[e,f]$. For the third term,
\[
\iiint_{E}\frac{\partial F_{3}}{\partial z}\,dV
=\iint_{[a,b]\times[c,d]}\bigl[F_{3}(x,y,f)-F_{3}(x,y,e)\bigr]dx\,dy ,
\]
by Fubini and the fundamental theorem of calculus in the $z$ direction. That right-hand side is exactly
the flux through the top face (normal $+\mathbf{k}$) and the bottom face (normal
$-\mathbf{k}$), because there $\mathbf{n}\,dA=\pm\mathbf{k}\,dx\,dy$. The other
terms give the flux through the two other pairs of faces.

For a general region, cut it into small boxes. An interior face appears twice
with opposite normals and cancels, so all that remains are the faces
that touch the outer boundary; the limit as the mesh goes to zero gives the statement for a curved
boundary. The full proof is in Spivak, Theorem 5-8.
\end{proof}

\begin{contoh}[Gauss on the unit ball]
Take $F(x,y,z)=(x^{3},y^{3},z^{3})$ and $E$ the unit ball. Here
$\operatorname{div}F=3(x^{2}+y^{2}+z^{2})=3\rho^{2}$, so with the
Jacobian determinant of the sphere
\[
\iiint_{E}3\rho^{2}\,dV=\int_{0}^{2\pi}\!\!\int_{0}^{\pi}\!\!\int_{0}^{1}
3\rho^{4}\sin\varphi\,d\rho\,d\varphi\,d\theta
=3\cdot2\pi\cdot2\cdot\frac15=\frac{12\pi}{5}.
\]
The surface side: on the unit ball $\mathbf{n}=(x,y,z)$, so
$F\cdot\mathbf{n}=x^{4}+y^{4}+z^{4}$ and
\[
\iint_{\partial E}(x^{4}+y^{4}+z^{4})\,dA=3\iint_{\partial E}z^{4}\,dA
=3\int_{0}^{2\pi}\!\!\int_{0}^{\pi}\cos^{4}\varphi\sin\varphi\,d\varphi\,d\theta
=3\cdot2\pi\cdot\frac25=\frac{12\pi}{5},
\]
with the factor $3$ from symmetry. The value is $7.5398223686155$, and quadrature on
the surface gives the same number to the last digit.
\end{contoh}

The second theorem talks about an open surface whose boundary is a curve.
Its formula requires the right-hand rule.

\begin{teorema}[Stokes]\index{Stokes' theorem}
Let $S$ be an oriented piecewise smooth surface with normal
$\mathbf{n}$, and let $\partial S$ be its boundary curve with the orientation given by the
right-hand rule. If $F$ has continuous partial derivatives on an open set
containing $S$, then
\[
\iint_{S}(\operatorname{curl}F)\cdot\mathbf{n}\,dA
=\oint_{\partial S}F\cdot d\mathbf{r}.
\]
\end{teorema}

\begin{proof}
Parametrize $S$ by $\mathbf{r}(u,v)$ over a region $D$ in the $(u,v)$ plane,
with the matching orientation. The chain rule turns the integrand into
derivatives of the components of $F$ with respect to $u$ and $v$, and after rearrangement the surface
integral becomes
\[
\iint_{D}\Bigl(\frac{\partial \tilde{Q}}{\partial u}
-\frac{\partial \tilde{P}}{\partial v}\Bigr)du\,dv ,
\]
for two functions $\tilde{P},\tilde{Q}$ built from $F$ and
$\mathbf{r}$. That is exactly the right-hand side of Green, so Green gives the line integral
on $\partial D$, which is nothing other than the line integral of $F$ on $\partial S$.
The full algebra is in Spivak, Theorem 5-9.
\end{proof}

If $S$ is flat in the $xy$ plane, Stokes becomes Green; if $S$ is only an
interval, it becomes the fundamental theorem of calculus.

\begin{contoh}[Stokes on a hemisphere]
Take $F=(z,x,y)$, whose curl is $(1,1,1)$, and $S$ the upper half of the unit ball
with outward normal, so that its boundary is the unit circle in the $xy$ plane
counterclockwise.

The surface side:
\[
\iint_{S}(1,1,1)\cdot\mathbf{n}\,dA=\iint_{S}(x+y+z)\,dA=\iint_{S}z\,dA
=\int_{0}^{2\pi}\!\!\int_{0}^{\pi/2}\cos\varphi\sin\varphi\,d\varphi\,d\theta
=2\pi\cdot\frac12=\pi ,
\]
because the $x$ and $y$ terms cancel each other by the symmetry of the hemisphere. The boundary side, with
$\mathbf{r}(t)=(\cos t,\sin t,0)$:
\[
\oint_{\partial S}F\cdot d\mathbf{r}=\int_{0}^{2\pi}\cos^{2}t\,dt=\pi .
\]
Both are $3.1415926535898$; quadrature gives the same number on the two sides.
\end{contoh}

Now the sign error in Stokes. Turn the normal of the surface inward:
the surface side becomes $-\pi$, the boundary side does not change as long as the direction of rotation
stays the same, and the equality fails with a difference of $2\pi$. The right-hand rule is what
binds the two sides; reversing both together restores the equality.

\section{Differential forms: one sentence for all of them}
\tingkat{advanced}

We have now seen three theorems with three faces. There is a language that
writes them as one statement, and that language is not mere tidiness: it
explains why the sign of orientation appears everywhere.

\begin{definisi}[Differential forms]
\index{differential form} A \emph{$k$-form} on an open set
$U\subseteq\mathbb{R}^{n}$ is a sum of terms of the form
$a(x)\,dx_{i_{1}}\wedge\cdots\wedge dx_{i_{k}}$, with $a$ a smooth function.
The wedge product $\wedge$ is antisymmetric:
\[
dx\wedge dy=-dy\wedge dx,\qquad dx\wedge dx=0 .
\]
The exterior derivative $d$ acts on the coefficients and adds one wedge:
\[
d(a\,dx_{I})=da\wedge dx_{I},\qquad da=\sum_{j}\frac{\partial a}{\partial x_{j}}dx_{j}.
\]
\end{definisi}

In everyday words: a differential form is an object that is ready to be
integrated over a surface of a given dimension, and $d$ raises the dimension by one
level. Antisymmetry makes reversing the orientation reverse the sign automatically;
the sign difficulties from before are built into that one rule.

\begin{teorema}[General Stokes]
Let $M$ be an oriented manifold of dimension $k$ with boundary, and let $\omega$
be a $(k-1)$-form with smooth coefficients. Then
\[
\int_{\partial M}\omega=\int_{M}d\omega ,
\]
with the orientation of $\partial M$ induced by the orientation of $M$.
\end{teorema}

The complete proof is skipped: Spivak, \emph{Integration on Chains}, Theorem 4-13
for chains, and Theorem 5-5 in \emph{Integration of Manifolds} for
manifolds; also Rudin, the chapter \emph{Integration of Differential Forms}. There
Green, Gauss, and classical Stokes are derived as special cases.

Its translation in three-dimensional space:

\begin{center}
\begin{tabular}{lll}
\hline
Language of forms & Vector field & Theorem \\
\hline
0-form $f$ & scalar function & $df$ gives the gradient \\
1-form $P\,dx+Q\,dy+R\,dz$ & $F=(P,Q,R)$ & $d\omega$ gives $\operatorname{curl}F$ \\
2-form $P\,dy\wedge dz+\cdots$ & $F=(P,Q,R)$ & $d\eta$ gives $\operatorname{div}F$ \\
$\int_{\partial M}\omega=\int_{M}d\omega$ & & Green, Gauss, Stokes, FTC \\
\hline
\end{tabular}
\end{center}

One small consequence: because $d(d\omega)=0$, the exterior derivative taken twice is always zero.
Translated, $\operatorname{curl}\nabla f=0$ and
$\operatorname{div}\operatorname{curl}F=0$. A symbolic check for a general function
and a general field gives zero in all six components, so those two identities that look
like coincidences are one statement.

\section{When these theorems fail}
\tingkat{advanced}

All the theorems in this chapter are conditional equalities. This section collects four
ways those equalities break, with numbers.

The first failure is a singularity. Take $F(x,y)=(-y,x)/(x^{2}+y^{2})$
on the plane without the origin. Compute its curl:
\[
\frac{\partial Q}{\partial x}-\frac{\partial P}{\partial y}
=\frac{(x^{2}+y^{2})-2x^{2}}{(x^{2}+y^{2})^{2}}
+\frac{(x^{2}+y^{2})-2y^{2}}{(x^{2}+y^{2})^{2}}=0 .
\]
A finite difference with $h=10^{-5}$ at $(1,0)$ gives $-2.01\times10^{-10}$, zero up to
the truncation error. But the circulation
around the unit circle is $2\pi=6.2831853072$, because there
$F\cdot\gamma'=1$ at every point. So zero curl does not mean zero circulation.

What is violated is the hypothesis of Green: $P$ and $Q$ must have continuous partial
derivatives on an open set containing $D$, whereas both are not
defined at the origin inside $D$, and $D$ is not simply connected. On the annulus $1\le r\le2$ Green still applies, with two boundaries
of opposite direction, and the integral is $2\pi-2\pi=0$; that zero is correct because there
the field has the potential $\theta$, the polar angle.

The second failure is swapping the order without integrability. Take
\[
f(x,y)=\frac{x^{2}-y^{2}}{(x^{2}+y^{2})^{2}}
\]
on the square $(0,1)\times(0,1)$. For fixed $y>0$,
\[
\int_{0}^{1}f(x,y)\,dx=\Bigl[-\frac{x}{x^{2}+y^{2}}\Bigr]_{x=0}^{x=1}
=-\frac{1}{1+y^{2}},
\]
so the outer integral gives $-\pi/4=-0.7853981634$; in the reverse order
the result is $+\pi/4$. Both iterated integrals exist but they differ, so by
Fubini the double integral does not exist. In polar coordinates
$f=\cos2\theta/r^{2}$, and on the unit quarter disk
\[
\iint|f|\,dA\ge\int_{0}^{\pi/2}|\cos2\theta|\,d\theta
\int_{0}^{1}\frac{dr}{r}=\infty ,
\]
because $\int_{0}^{1}dr/r$ diverges at $r=0$. Fubini is not wrong: its hypothesis
is not satisfied.

The third failure is an unmanageable boundary. Gauss demands that $E$ be closed
and bounded with a piecewise smooth boundary. On an unbounded set,
the boundary can run off to infinity and the surface integral stops meaning anything;
Spivak places this as an exercise after Theorem 5-5, which fails without
compactness. Finitely many corners are still safe, because the normal jumps only at
finitely many points; a boundary with endlessly jagged teeth, a coastline for instance, is not.

The fourth failure is a surface that cannot be oriented. On a Möbius
strip, walking around the strip brings the normal back to the opposite
side, so $\mathbf{n}$ in Stokes' formula is not defined.
The theorem is not wrong; its statement is simply incomplete for objects of that kind.

\begin{tantangan}
Take $T(u,v)=(u^{2},v)$ on $D^{*}=(-2,2)\times(0,1)$. This map seems to
widen things: the point $u=2$ lands at $x=4$, and its derivative $2u$ grows without
bound as $u$ moves away from zero. Compute the area of $T(D^{*})$ from the picture, then
compute it with the change-of-variables formula. The two numbers differ by exactly a factor
$2$. Explain why, and name the line on which the Jacobian determinant is zero.

Hint 1: check whether $T$ is one-to-one. There are two values of $u$ that land at
$x=1$.

Hint 2: the line $u=0$ is folded onto a single point, so every image point
is counted twice by the integral; $|\det J|=|2u|$ does not know that the fold
is only one sheet.

If you are stuck, go back to the definition of a valid change of variables and check the
one-to-one condition.
\end{tantangan}

One closing warning, which is not in the textbooks. Numerical quadrature never
checks the hypotheses. It gives $6.2831853072$ with seven correct digits
for a field that violates the hypothesis of Green, and no error message. A tidy number is not proof that the hypothesis holds; it is only proof that
the arithmetic ran.

\section{Worked examples}
\tingkat{intermediate}

Now three computations that form the backbone of this chapter. All three can be checked
with Python, and all three have gone wrong in the author's own hands.

\textbf{One integral, two coordinate systems.}
\nopagebreak[4]

Compute
\[
I=\iint_{D}e^{-(x^{2}+y^{2})}\,dA,\qquad D=\{(x,y):x^{2}+y^{2}\le1\}.
\]
In polar coordinates, $D$ becomes $0\le r\le1$, $0\le\theta\le2\pi$, and
$dA=r\,dr\,d\theta$. The integral separates into two factors:
\[
I=\int_{0}^{2\pi}\!\!\int_{0}^{1}e^{-r^{2}}r\,dr\,d\theta
=2\pi\cdot\Bigl[-\frac12e^{-r^{2}}\Bigr]_{0}^{1}
=2\pi\cdot\frac{1-e^{-1}}{2}=\pi\bigl(1-e^{-1}\bigr).
\]
Its value is $1.9858653038$. The antiderivative with respect to $r$ exists precisely because of the factor
$r$: without that factor, the antiderivative involves the error function.

In Cartesian coordinates, the region demands two pieces:
\[
I=\int_{-1}^{1}\!\!\int_{-\sqrt{1-x^{2}}}^{\sqrt{1-x^{2}}}
e^{-(x^{2}+y^{2})}\,dy\,dx .
\]
The inner integral has no elementary antiderivative, so we have to use
quadrature. The results of the two forms, at the same tolerance:
\begin{align*}
\text{cartesian}&: 1.9858653037988705 \quad (11907\ \text{evaluations}),\\
\text{polar}&: 1.9858653037988712 \quad (441).
\end{align*}
The errors are $8.9\times10^{-16}$ and $2.2\times10^{-16}$; the difference in accuracy
does not matter, but the difference in cost is a factor of $27$.

The midpoint rule on the same uniform Cartesian grid lowers the error more
slowly: $1.141\times10^{-3}$ on $50$ cells per side, and
$4.667\times10^{-5}$ on $400$. Each time the number of cells is doubled, the error
drops by about a factor of $2.9$, not $4$, because the height of the columns near the edge
behaves like $\sqrt{1-x^{2}}$, so the error is of order $h^{3/2}$.

\textbf{Volume of a ball with the Jacobian determinant.}

The volume of a ball of radius $R$ is
\[
V=\int_{0}^{2\pi}\!\!\int_{0}^{\pi}\!\!\int_{0}^{R}
\rho^{2}\sin\varphi\,d\rho\,d\varphi\,d\theta
=2\pi\cdot2\cdot\frac{R^{3}}{3}=\frac{4\pi R^{3}}{3}.
\]
For $R=1$ its value is $4.1887902048$. Three-dimensional quadrature on the ball gives
\[
4.188790204786397 \quad(\text{difference }6.2\times10^{-15}),
\]
and a Monte Carlo simulation with $10^{6}$ points in the cube $[-1,1]^{3}$ gives
$4.192824\pm0.004$.

Compare that with dropping the Jacobian determinant. Without the factor
$\rho^{2}\sin\varphi$,
\[
\int_{0}^{2\pi}\!\!\int_{0}^{\pi}\!\!\int_{0}^{1}1\,d\rho\,d\varphi\,d\theta
=2\pi^{2}=19.7392088022,
\]
which is $3\pi/2$ times too large, and the result is still a number that looks
plausible for a volume.



The following table collects every number in this chapter.

\begin{center}
{\small
\begin{tabular}{llll}
\hline
Quantity & Exact & Numerical & Error \\
\hline
$\iint_{D}e^{-(x^{2}+y^{2})}dA$, polar & $1.985865$ & $1.985865$ & $2.2\times10^{-16}$ \\
$\iint_{D}e^{-(x^{2}+y^{2})}dA$, Cartesian & $1.985865$ & $1.985865$ & $8.9\times10^{-16}$ \\

$\int_{\gamma}F\cdot d\mathbf{r}$ & $10$ & $10.000000$ & $1.8\times10^{-15}$ \\
$\oint_{\partial D}$ Green, square & $2$ & $2.000000$ & $<10^{-15}$ \\
$\iint_{S}(\operatorname{curl}F)\cdot\mathbf{n}\,dA$ & $3.141593$ & $3.141593$ & $<10^{-13}$ \\
\hline
\end{tabular}}
\end{center}

\kertas{
Test yourself in three steps. (1) Recompute $I$ by hand using the
antiderivative $-\tfrac12e^{-r^{2}}$, then compare with $1.9858653038$;
a difference in the last digit is a rounding error, not a formula error. (2)
Run \texttt{scipy.integrate.dblquad} on both forms and count how many
times the function is called; the numbers $11907$ and $441$ will reappear, provided
the tolerance is the same for both. (3) Change $D$ to the disk of radius $2$
and check that the exact answer becomes $\pi(1-e^{-4})$, namely $3.0841$.
If the third step gives a different number, what is wrong is almost certainly the limit on $r$,
not the formula. }

\section{Computational lab}
\tingkat{advanced}

The following code checks five things at once: the cost of the two coordinate systems, Green,
Gauss, Stokes, and a field with a singularity. The numbers in the previous table come
from this program.

\begin{lab}
\begin{verbatim}
import numpy as np
from scipy import integrate

# 1. One double integral, two coordinate systems
def fc(y, x):
    return np.exp(-(x*x + y*y))

cnt = [0]
def fc_count(y, x):
    cnt[0] += 1
    return fc(y, x)

cart = integrate.dblquad(fc_count, -1, 1,
                         lambda x: -np.sqrt(1-x*x),
                         lambda x: np.sqrt(1-x*x),
                         epsabs=1e-12, epsrel=1e-12)
ncart = cnt[0]
cnt[0] = 0

def fp(r, th):
    cnt[0] += 1
    return np.exp(-r*r)*r

pol = integrate.dblquad(fp, 0, 2*np.pi, 0, 1,
                        epsabs=1e-12, epsrel=1e-12)
value = np.pi*(1 - np.exp(-1))
print("cartesian %.16f error %.2e eval %d"
      % (cart[0], abs(cart[0]-value), ncart))
print("polar     %.16f error %.2e eval %d"
      % (pol[0], abs(pol[0]-value), cnt[0]))
print("exact     %.16f" % value)

# 2. Green: (x^2-y)dx + (x+2y)dy on the unit square
P = lambda x, y: x*x - y
Q = lambda x, y: x + 2*y
edges = [integrate.quad(lambda t: P(t, 0.0), 0, 1)[0],
         integrate.quad(lambda t: Q(1.0, t), 0, 1)[0],
         integrate.quad(lambda t: P(t, 1.0), 1, 0)[0],
         integrate.quad(lambda t: Q(0.0, t), 1, 0)[0]]
dbl = integrate.dblquad(lambda y, x: 2.0, 0, 1, 0, 1)[0]
print("Green edges %.15f  double %.15f" % (sum(edges), dbl))

# 3. Gauss: F=(x^3,y^3,z^3) on the unit ball
def flux(phi, th):
    sp, cp = np.sin(phi), np.cos(phi)
    x, y, z = sp*np.cos(th), sp*np.sin(th), cp
    return (x**4 + y**4 + z**4)*sp

F = integrate.dblquad(flux, 0, 2*np.pi, 0, np.pi,
                      epsabs=1e-11, epsrel=1e-11)[0]
print("Gauss flux %.13f  12pi/5 %.13f" % (F, 12*np.pi/5))

# 4. Stokes: F=(z,x,y) on the upper hemisphere of the unit ball
def curl_n(phi, th):
    sp, cp = np.sin(phi), np.cos(phi)
    x, y, z = sp*np.cos(th), sp*np.sin(th), cp
    return (x + y + z)*sp

S = integrate.dblquad(curl_n, 0, 2*np.pi, 0, np.pi/2)[0]
L = integrate.quad(lambda t: np.cos(t)**2, 0, 2*np.pi)[0]
print("Stokes surface %.13f  boundary %.13f" % (S, L))

# 5. A field with a singularity: F=(-y,x)/(x^2+y^2)
def circle(t):
    x, y = np.cos(t), np.sin(t)
    return (-y*(-np.sin(t)) + x*np.cos(t))/(x*x + y*y)

def Pg(x, y): return -y/(x*x + y*y)
def Qg(x, y): return x/(x*x + y*y)

h = 1e-5
curl_01 = ((Qg(1+h, 0) - Qg(1-h, 0))/(2*h)
           - (Pg(1, h) - Pg(1, -h))/(2*h))
C = integrate.quad(circle, 0, 2*np.pi)[0]
print("circle %.13f  2pi %.13f" % (C, 2*np.pi))
print("curl(1,0) %.2e" % curl_01)
\end{verbatim}
\end{lab}

The actual output, as it came out:

\begin{verbatim}
cartesian 1.9858653037988705 error 8.88e-16 eval 11907
polar     1.9858653037988712 error 2.22e-16 eval 441
exact     1.9858653037988714
Green edges 2.000000000000000  double 2.000000000000000
Gauss flux 7.5398223686155  12pi/5 7.5398223686155
Stokes surface 3.1415926535898  boundary 3.1415926535898
circle 6.2831853071796  2pi 6.2831853071796
curl(1,0) -2.01e-10
\end{verbatim}

Four things deserve a reading. First, the two opening lines agree to the fifteenth digit at a cost that differs
by a factor of $27$: the right coordinates do not change the answer, only
its price. Second, Green matches to the last digit, but that is
only proof that the implementation is consistent with its formula. Third, Gauss and Stokes are
computed from different sides and still meet; that is where verification has teeth.
Fourth, quadrature gives $6.2831853071796$ for a field that violates the hypothesis of
Green, without a single error message.

\begin{untukdosen}
The teaching order I use: the density of the plate and the program that forgets the factor
$r$, then Fubini, then the Jacobian determinant as the area of a stretching cell,
then line integrals and their fundamental theorem, and only then Green. Green is best proved
in full for type I regions at the board, because that proof uses only two tools
the students already have. Gauss and Stokes are enough to prove for a box and for a
parametrized surface.

Two misconceptions come up most often. First, students think
$\partial P/\partial y=\partial Q/\partial x$ is enough for conservative, because
they have only tested fields on a domain without holes. The field
$(-y,x)/(x^{2}+y^{2})$ has to appear at the board, not on the reading list. Second,
they treat the Jacobian determinant as a notational formality and write
$r$ or $\rho^{2}\sin\varphi$ as something to memorize without checking the limits.
Compare $4\pi/3$ with $2\pi^{2}$; that ratio of $4.712$ times is more impressive
than a hundred warnings.

One classroom activity: four groups, each computing one side of
Green, Gauss, and Stokes for the same field, then matching numbers. The group
with the sign error is found out within five minutes, and that is better than being told.
The sign of readiness for what comes next: students can explain why orientation
changes the sign without looking at the formula.
\end{untukdosen}

\section{Summary}
\begin{ringkasan}
\begin{itemize}
\item The double integral is defined through upper and lower sums over cells;
Fubini allows the order of integration to be swapped when the function is integrable.
\item Changing coordinates multiplies the area of a cell by $|\det J|$; polar gives
$r$, spherical gives $\rho^{2}\sin\varphi$. Without that factor the result is tidy but
wrong.
\item For a conservative field, the line integral depends only on the endpoints.
\item Green, Gauss, and Stokes state one thing: the integral of a derivative over a region
equals the integral of the field over its boundary, and orientation decides the sign.
\item All those theorems are conditional: singularities, unmanageable boundaries, surfaces that cannot
be oriented, and non-integrable functions make them fail without a warning.
\end{itemize}
\end{ringkasan}

\section{Exercises}
\begin{latihan}
\begin{enumerate}
\item Compute $\iint_{D}(x^{2}+y^{2})\,dA$ on the unit disk in Cartesian and polar coordinates, check it with \texttt{dblquad}, and state the number of evaluations for each method.

\item Prove that if $T(u,v)=A(u,v)$ is linear and $\det A\ne0$, the
change-of-variables formula becomes $|\det A|\iint_{D}f\circ T\,dA$.

\item Repeat the Green example on the square in the clockwise direction: show that the result is
$-2$ and state which side contributes which sign.

\item Compute the volume of the ellipsoid $x^{2}/a^{2}+y^{2}/b^{2}+z^{2}/c^{2}\le1$ with
the change of variables $x=au$, $y=bv$, $z=cw$, then check the formula on
$a=b=c=1$.

\item Assess this claim: true, false, or not enough evidence. ``Because
$\operatorname{curl}F=0$ at every point where $F$ is defined, the field
$F=(-y,x)/(x^{2}+y^{2})$ is conservative on
$\mathbb{R}^{2}\setminus\{(0,0)\}$.'' Include a reason and one supporting
number.

\item Prove that the line integral is independent of path if and only if the
circulation of every closed curve is zero, by using two paths from $A$ to
$B$ that together form one closed curve.

\item Compute both sides of Stokes for $F=(-y,x,z)$ on the paraboloid
$z=1-x^{2}-y^{2}$, $z\ge0$, with the normal pointing upward, then check
with quadrature.

\item Compute both sides of Gauss for $F=(x,y,z)/(x^{2}+y^{2}+z^{2})^{3/2}$ on the unit
ball. Explain why the result does not contradict the theorem, and
which hypothesis is violated. (Hint: compute $\operatorname{div}F$ first.)

\item Show that the Jacobian determinant of spherical coordinates is zero at $\varphi=0$ and
$\varphi=\pi$, then explain why the map may still be used on
$0<\varphi<\pi$. (Hint: check where the map stops being one-to-one.)
\end{enumerate}
\end{latihan}

\begin{untukpeneliti}
What is really proved in this chapter, with a written proof: the lemma on a region
between two graphs, the fundamental theorem for line integrals, its converse through the integral
potential, Green for type I and type II regions, the consequence for the curl test on a
simply connected domain, and the sketch of the proof of Gauss through boxes and of Stokes
through a parametrization. What I quote without proof: Fubini (Spivak, the chapter
\emph{Integration}, Theorem 3-10), the full change of variables (Spivak's same
chapter; Rudin, the chapter \emph{Integration of Differential Forms}), Gauss (Spivak,
Theorem 5-8), and Stokes (Spivak, Theorem 5-9); general Stokes is Theorem
4-13 and 5-5.

What is overpopularized is the phrase ``zero curl means
conservative''. What is correct: zero curl means closed form, and a closed form is
exact only when the first cohomology group of the domain is zero. For $\mathbb{R}^{3}$
that holds, for $\mathbb{R}^{2}\setminus\{0\}$ it does not, and the difference is
measured by the integral on the circle, namely $2\pi$ here.

The open question that is genuinely open for me: can quadrature give
a certificate that the hypotheses of a theorem are satisfied, and not merely an estimate of the value?
The output in the lab shows that a numerical answer correct to fifteen
digits can come from a field that violates the hypotheses, with no warning sign. For
a field that is known only through samples on a grid, there is no general way
to be sure that there is no singularity between the sample points.

Limits of validity: all functions here are real-valued and piecewise
smooth; all numbers come from Python 3.14.4, NumPy 2.5.2, SciPy 1.18.0
with a tolerance of $10^{-12}$ unless stated otherwise, and spherical surfaces
are parametrized with $\mathbf{n}$ pointing outward.
\end{untukpeneliti}

\peta{Chapter 15 follows Stewart, \emph{Calculus: Early Transcendentals}, the chapter ``Multiple Integrals'' (15.1--15.3 and 15.6--15.9) and the chapter ``Vector Calculus'' (16.2
line integrals, 16.3 the fundamental theorem, 16.4 Green, 16.5 curl
and divergence, 16.8 Stokes, 16.9 Gauss); on Spivak, \emph{Calculus on
Manifolds}, the chapter ``Integration'' (Fubini, Theorem 3-10, change of variables),
the chapter ``Integration on Chains'' and the chapter ``Integration of Manifolds'' (Theorems
5-5, 5-8, 5-9); on Rudin, \emph{Principles of Mathematical Analysis}, the chapters
``Functions of Several Variables'' and ``Integration of Differential Forms''; on Active Calculus, chapters 11 and 12 for a more gradual order.
The basis is the chapter \emph{Integral, Series of Functions, and Exchange of Limits} (the integral part)
and the chapter \emph{Continuity and Derivatives} (partial derivatives, Jacobian) of the first edition
of this book. This chapter uses the Riemann integral and the fundamental theorem of calculus
from Chapter 13, uniform continuity and compact sets from Chapter 11, determinants
from Chapter 7, and gradients and Jacobian matrices from Chapter 12, all of them referred to
and not repeated.}

\chapter{Differential Equations and Dynamical Systems}

The pendulum clock in the living room swings with an amplitude of sixty degrees, and the
length of its arm is one meter. The textbook gives the period
\[
T=2\pi\sqrt{L/g}=2.006409\ \text{seconds},
\]
and that formula is indeed correct for a very small swing. Measured at an amplitude of
sixty degrees, the same pendulum needs $2.153242$ seconds for one full swing: $7.32$
percent longer. A clock set with the small-angle formula falls behind by $5892$ seconds,
about ninety-eight minutes, every day.

The formula is not a miscalculation. It is the exact answer for a different equation.
The actual swing satisfies
\[
\frac{d^{2}\theta}{dt^{2}}=-\frac{g}{L}\sin\theta ,
\]
whereas the small-period formula comes from replacing $\sin\theta$ with $\theta$. At
$\theta=60^{\circ}$ that replacement turns $0.866025$ into $1.047198$, a difference of
seventeen percent at the farthest point, and the period error left behind is $7.32$
percent.

Chapter 12 taught us to differentiate, Chapter 14 series of functions, Chapter 15 vector
fields. Not one of them answers the question here: if the rate of change is known, how do we
find the function, is it unique, and what happens in the long run. This chapter fills those
four gaps: how to solve, conditions for existence and uniqueness, the phase plane, and the
stability test.

\begin{untukanak}
Imagine the water tank in a bathroom. The tap above runs at a steady rate, the drain below
runs out, and the higher the water, the faster it leaves. If we open the tap and then walk
away, tomorrow's water level cannot be computed from the size of the tap alone: what
determines it is the water level now. High water pushes out faster, the rise slows down,
and finally it stops when the water coming in equals the water going out.

Other things work the same way. Hot coffee cools faster while it is still very hot. Germs
multiply faster when there are many of them, then slow down when the food runs out. What we
know is the rule of change, not the final result.

Now the question that makes this work hard. Two tanks the same, taps the same, drains the
same, but one starts empty and one starts half full: do they end at the same level?
Sometimes yes, sometimes no. And there is a wicked drain, partly clogged and then suddenly
open: two different stories can come out of the same initial state.

In this chapter we shall learn why the rule of change sometimes suffices to determine the
whole future, and sometimes does not.
\end{untukanak}

\section{A rate that depends on itself}
\tingkat{basic}

A differential equation sets the rate of change of a function equal to something that
contains the function itself. To solve it means to find a function that satisfies that
sentence on an entire interval. The derivatives of Chapter 12 were computed at a single
point; the solution here is a function on an interval, and that difference is the one most
often missed.

\begin{definisi}[Ordinary differential equation]\index{differential equation}
The equation
\[
F\bigl(t,y,y',y'',\dots,y^{(n)}\bigr)=0
\]
is called an \emph{ordinary differential equation} (ODE) of order $n$ for a function $y$ of
the variable $t$ if $F$ contains the $n$-th derivative and contains no higher derivative.
The ODE is called \emph{linear} if $F$ has the form
\[
a_{n}(t)\,y^{(n)}+\cdots+a_{1}(t)\,y'+a_{0}(t)\,y=b(t),
\]
and is called \emph{autonomous} if $F$ does not contain $t$ explicitly.
\end{definisi}

 All the examples in
this chapter are of first or second order, because those two already contain every important
idea.

\begin{definisi}[Solution and initial value problem]\index{initial value problem}
A function $y$ is called a \emph{solution} of the ODE on an interval $I$ if $y$ has
derivatives up to order $n$ on $I$ and satisfies the equation at every point of $I$. If the
condition $y(t_{0})=y_{0}$ is added, the combination of the two is called an \emph{initial
value problem}, and $y_{0}$ is called the initial value.
\end{definisi}

One example to make sure the definition works. Take the population model
\[
y'=y(1-y),\qquad y(0)=\tfrac12 ,
\]
with $y$ denoting the fraction of the population relative to the carrying capacity. The
candidate solution is
\[
y(t)=\frac{1}{1+e^{-t}} .
\]
Its derivative is $y'(t)=e^{-t}/(1+e^{-t})^{2}$, and the right-hand side
$y(1-y)=e^{-t}/(1+e^{-t})^{2}$ is exactly the same. The initial value $y(0)=1/2$ matches as
well. So the function is a solution, and that check uses only the quotient rule of
Chapter 12.

What deserves attention: $y(t)=1$ and $y(t)=0$ are also solutions, both of them constant,
and both appear when we ask when the right-hand side is zero. So one equation can have many
solutions; the initial value selects one of them. The collection of all solutions is called
the \emph{general solution}, and it contains one constant for each order.

The interval on which a solution lives is part of the answer. Take
\[
y'=y^{2},\qquad y(0)=1 .
\]
Separation of variables gives $y(t)=1/(1-t)$, and that formula makes sense only on $t<1$.
At $t=0.9$ its value is $10$; at $t=0.99$ its value is $100$; at $t=0.999$ its value is
$1000$. There is no finite number at $t=1$, and no solution continues that function past it.
The solution \emph{blows up} in finite time, and a program that reports a value at $t=2$ does
not find a solution; it finds rounding.

\kamus{differential equation}{persamaan diferensial}{A sentence that ties the rate of change of a quantity to the size of that quantity itself.}

\kamus{general solution}{solusi umum}{The family of all functions satisfying the equation, usually containing one free constant per order.}

\begin{bacakeras}
--- If I let this ball go on the slope, it will move downhill. How fast?
--- It depends on where the ball is now. Near the top, slowly. In the middle of the slope,
quickly.
--- So its speed is not a single fixed number?
--- No. It is a rule: at each point, the speed is determined by that point.
--- If the rule already knows everything, is the path already determined?
--- Not yet. The rule gives a direction at every point; the path is born from the starting
point you choose.
--- Two balls, two starting points, one rule?
--- One rule, two paths.
--- If two paths meet at the same point, how do we know which one came from where?
--- We cannot. That is why this chapter demands one starting point for one path.

Repeat it once, in a normal voice: same rule, different starting point, different path.
\end{bacakeras}

\section{Two families that can be solved by hand}
\tingkat{basic}

Not every differential equation has a solution formula, but two large families do, and both
appear in almost every first-level application. The first family separates variables; the
second is the first-order linear equation.

\subsection{Separation of variables}

If the equation can be written $y'=f(t)g(y)$, we move all $y$ to one side and all $t$ to the
other side:
\[
\frac{dy}{g(y)}=f(t)\,dt,\qquad \int\frac{dy}{g(y)}=\int f(t)\,dt .
\]
That step is not a notational trick but the chain rule run backwards: if $H$ is an
antiderivative of $1/g$, then $\frac{d}{dt}H(y(t))=H'(y)y'=y'/g(y)$, so the original
equation becomes $\frac{d}{dt}H(y(t))=f(t)$ and Chapter 13 gives $H(y)=\int f\,dt+C$. The
result is usually implicit, and that is a legitimate answer.

The logistic model above is used throughout the chapter. With $g(y)=y(1-y)$ and partial
fractions,
\[
\int\frac{dy}{y(1-y)}=\int\Bigl(\frac{1}{y}+\frac{1}{1-y}\Bigr)dy
=\ln\Bigl|\frac{y}{1-y}\Bigr| ,
\]
so $\ln|y/(1-y)|=t+C$ and
\[
y(t)=\frac{1}{1+Ae^{-t}},\qquad A=\frac{1-y_{0}}{y_{0}} .
\]
For $y_{0}=1/2$ we obtain $A=1$, which is the formula we checked a moment ago. Note that $A$
is determined by the initial value, and that $A=0$ gives the constant solution $y=1$ which
the division above does not cover: when separating variables we divide by $g(y_{0})$, and
that division is valid only if $g(y_{0})\ne0$.

\subsection{Linear equations and the integrating factor}

The standard form of a first-order linear equation is
\[
y'+p(t)\,y=q(t).
\]
The left-hand side is almost the derivative of a product, only one factor short. The missing
factor is exactly what the following theorem looks for.

\begin{teorema}[Integrating factor]\index{integrating factor}
Suppose $p$ and $q$ are continuous on an open interval $I$ containing $t_{0}$, and suppose
\[
\mu(t)=\exp\Bigl(\int_{t_{0}}^{t}p(s)\,ds\Bigr).
\]
Then the initial value problem $y'+p(t)y=q(t)$, $y(t_{0})=y_{0}$ has exactly one solution on
$I$, namely
\[
y(t)=\frac{1}{\mu(t)}\Bigl(y_{0}+\int_{t_{0}}^{t}\mu(s)\,q(s)\,ds\Bigr).
\]
\end{teorema}

\begin{proof}
Since $p$ is continuous, $\mu$ is defined, positive, and differentiable on $I$, with
$\mu'=p\mu$ by the fundamental theorem of calculus (Chapter 13). Multiply the equation by
$\mu$:
\[
\mu y'+\mu p\,y=\mu q
\quad\Longleftrightarrow\quad
(\mu y)'=\mu q ,
\]
because the product rule gives $(\mu y)'=\mu'y+\mu y'=\mu(py+y')$. Integrate both sides from
$t_{0}$ to $t$ and use the fundamental theorem of calculus:
\[
\mu(t)y(t)-\mu(t_{0})y(t_{0})=\int_{t_{0}}^{t}\mu(s)q(s)\,ds .
\]
Since $\mu(t_{0})=1$, division by $\mu(t)>0$ gives the formula requested. Uniqueness: if
$y_{1},y_{2}$ are two solutions, their difference $w=y_{1}-y_{2}$ satisfies $w'+pw=0$ with
$w(t_{0})=0$; multiplying by $\mu$ gives $(\mu w)'=0$, so $\mu w$ is constant, and since its
value is zero at $t_{0}$ it is zero everywhere.
\end{proof}

That proof guesses nothing: $\mu$ is found by demanding $\mu'=p\mu$, an equation easier than
the original. This pattern repeats: one problem is solved by turning it into a simpler
problem.

\begin{contoh}[One linear equation, two checks]
Solve $y'+2y=4t$ with $y(0)=1$. Here $p=2$ and $q=4t$, so $\mu=e^{2t}$ and
\[
(e^{2t}y)'=4t\,e^{2t}.
\]
Integration by parts gives
\[
\int 4t\,e^{2t}\,dt=\Bigl(2t-1\Bigr)e^{2t}+C ,
\]
because the derivative of $(2t-1)e^{2t}$ is $2e^{2t}+2(2t-1)e^{2t}=(4t)e^{2t}$. Thus
\[
y(t)=2t-1+Ce^{-2t},
\]
and the condition $y(0)=1$ gives $C=2$, that is $y(t)=2t-1+2e^{-2t}$.

Check: substitution gives $y'+2y=4t$, and $y(1)=1+2e^{-2}=1.2706705665$; an adaptive
integrator with tolerance $10^{-12}$ gives $1.270670566473$.
\end{contoh}

\subsection{The RC circuit: the simplest real example}

A voltage source $V_{s}$ is connected in series with a resistor $R$ and a capacitor $C$.
Kirchhoff's law of voltages gives $V_{s}=Ri+V_{C}$, and the law of the capacitor gives
$i=C\,dV_{C}/dt$. The two together give
\[
RC\,\frac{dV_{C}}{dt}+V_{C}=V_{s},
\qquad\text{or}\qquad
\frac{dV_{C}}{dt}+\frac{1}{RC}V_{C}=\frac{V_{s}}{RC}.
\]
This is a linear equation with $p=1/(RC)$ constant. By the theorem above, with
$V_{C}(0)=0$,
\[
V_{C}(t)=V_{s}\bigl(1-e^{-t/\tau}\bigr),\qquad \tau=RC .
\]
The number $\tau$ is called the time constant. For $R=1\,\mathrm{k}\Omega$ and
$C=1\,\mu\mathrm{F}$, $\tau=10^{-3}$ seconds. With $V_{s}=5$ volts, the voltage at
$t=\tau$ is $3.160603$ volts and at $t=5\tau$ it is $4.966310$ volts, that is $99.33$
percent of $5$ volts. To reach $90$ percent requires $t=\tau\ln10=2.302585$ milliseconds.

\begin{resep}
\begin{enumerate}
\item Write the equation in standard form $y'+p(t)y=q(t)$: divide by the coefficient of
  $y'$ so that the coefficient is exactly one.
\item Compute $\mu(t)=\exp\bigl(\int p(t)\,dt\bigr)$, with no constant of integration.
\item Multiply the whole equation by $\mu$, then write the left-hand side as $(\mu y)'$. If
  the left-hand side refuses to become a single derivative, step 1 is wrong.
\item Integrate both sides with respect to $t$, then divide by $\mu$.
\item Use the initial value to determine the constant, and write down the interval of
  validity.
\end{enumerate}
Warning: step 1 must not be skipped, and step 3 must not be done before step 2. Computing
$\mu$ from coefficients that have not yet been normalized gives a wrong integrating factor,
and the error only becomes visible when the solution is tested.
\end{resep}

\kamus{integrating factor}{faktor integrasi}{A multiplier that turns the left-hand side of a linear equation into a single derivative of a product.}

\section{Existence, uniqueness, and Picard iteration}
\tingkat{intermediate}

The example $y'=y^{2}$ shows that a solution can stop existing. Given $f$ and one initial
point, does the solution exist and is it only one? The answer depends on one property of $f$
with respect to $y$.

\begin{definisi}[Lipschitz condition]\index{Lipschitz condition}
A function $f(t,y)$ satisfies the Lipschitz condition with respect to $y$ on a region
$D\subseteq\mathbb{R}^{2}$ if there is a constant $L$ such that
\[
|f(t,y_{1})-f(t,y_{2})|\le L\,|y_{1}-y_{2}|
\]
for every $(t,y_{1}),(t,y_{2})\in D$. The smallest possible constant is called the Lipschitz
constant.
\end{definisi}

Its everyday sentence: $f$ must not change too abruptly when $y$ shifts a little. If $f$ has
a partial derivative with respect to $y$ that is continuous and bounded by $L$ on $D$, that
condition holds automatically, because the mean value theorem (Chapter 12) gives
$f(t,y_{1})-f(t,y_{2})=\partial_{y}f(t,\xi)(y_{1}-y_{2})$.

\begin{teorema}[Picard--Lindel\"of]\index{Picard-Lindelof theorem}
Suppose $f$ is continuous on the closed rectangle
$D=[t_{0}-a,t_{0}+a]\times[y_{0}-b,y_{0}+b]$ and satisfies the Lipschitz condition with
respect to $y$ on $D$ with constant $L$. Suppose $M=\max_{D}|f|$ and
$h=\min(a,\,b/M)$. Then the initial value problem $y'=f(t,y)$, $y(t_{0})=y_{0}$ has exactly
one solution on the interval $|t-t_{0}|\le h$.
\end{teorema}

The full proof uses the Banach fixed point theorem on the space of continuous functions with
the maximum norm, which is complete because uniform continuity is discussed in Chapter 11
and Chapter 14. I give the sketch, because the sketch contains the whole idea and can be
turned into an algorithm.

Write the equation as an integral equation, which is equivalent because both sides are
continuous:
\[
y(t)=y_{0}+\int_{t_{0}}^{t}f\bigl(s,y(s)\bigr)\,ds .
\]
Define the map $T$ by that formula, starting from the constant $y_{0}(t)=y_{0}$, then
iterate:
\[
y_{n+1}(t)=y_{0}+\int_{t_{0}}^{t}f\bigl(s,y_{n}(s)\bigr)\,ds .
\]
Every $y_{n}$ lies inside $D$ as long as $|t-t_{0}|\le h$, because
$|y_{n+1}-y_{0}|\le M|t-t_{0}|\le Mh\le b$. The difference of two successive iterations
satisfies
\[
|y_{n+1}(t)-y_{n}(t)|\le L\Bigl|\int_{t_{0}}^{t}|y_{n}(s)-y_{n-1}(s)|\,ds\Bigr|
\le \frac{(L|t-t_{0}|)^{n}}{n!}\,\|y_{1}-y_{0}\| ,
\]
and induction on $n$ gives that bound. Since $\sum_{n}(Lh)^{n}/n!$ converges
(ratio test, Chapter 5), the series $y_{0}+\sum(y_{n+1}-y_{n})$ converges uniformly, and
its limit is the solution. Uniqueness comes from the same inequality: two different solutions
would satisfy $|y_{1}-y_{2}|\le L|t-t_{0}|\,\|y_{1}-y_{2}\|$, which for $Lh<1$ is possible
only if the difference is zero; the case $Lh\ge1$ is handled by shrinking $h$, and that is
why the theorem promises only a short interval.

Run that iteration on $y'=y$, $y(0)=1$. Starting from $y_{0}=1$:
\[
y_{1}=1+t,\quad y_{2}=1+t+\frac{t^{2}}{2},\quad
y_{3}=1+t+\frac{t^{2}}{2}+\frac{t^{3}}{6},
\]
so $y_{n}$ is the partial sum of the exponential series from Chapter 5. At $t=1$, the eighth
iteration gives $2.718278769841$, whereas $e=2.718281828459$; the error is
$3.06\times10^{-6}$.

The Lipschitz condition must not be thrown away. Take
\[
y'=\sqrt{|y|},\qquad y(0)=0 .
\]
The right-hand side is continuous, so a solution exists; but it does not satisfy the
Lipschitz condition near $y=0$, because its slope is unbounded there. And indeed there are
two solutions: the zero function, $y(t)\equiv0$, and
\[
y(t)=\frac{t^{2}}{4}\quad (t\ge0).
\]
The check is one line: $y'=t/2$ and $\sqrt{|y|}=\sqrt{t^{2}/4}=t/2$ for
$t\ge0$, whereas $y(0)=0$. Two different histories from one and the same initial state,
exactly like the drain that opens suddenly in the opening box of the chapter.

Lipschitz is a \emph{sufficient} condition, not a necessary one. Existence without
uniqueness is guaranteed by Peano's theorem, which uses compactness (Chapter 11) and
Arzel\`a--Ascoli; see Coddington and Levinson, chapters 1--2.

\kamus{uniqueness}{ketunggalan}{The guarantee that one initial state gives birth to only one history, with no branch.}
\section{Systems, the phase plane, and oscillators}
\tingkat{intermediate}

One second-order equation is not a new animal; it is two first-order equations running
together. By naming $v=\theta'$, the pendulum equation becomes
\[
\theta'=v,\qquad v'=-\frac{g}{L}\sin\theta .
\]
This pair is autonomous: its right-hand side does not contain $t$, so the initial state
determines the whole history (provided the Lipschitz condition holds) and we may forget time
for a moment in order to ask what happens on the plane $(\theta,v)$.

\begin{definisi}[Autonomous systems and the phase plane]\index{phase plane}
The system
\[
\mathbf{x}'=\mathbf{f}(\mathbf{x}),\qquad \mathbf{x}\in\mathbb{R}^{n},
\]
is called \emph{autonomous}. For $n=2$, the plane on which the point $\mathbf{x}$ lives is
called the \emph{phase plane}; the curve $\{\mathbf{x}(t)\}$ is called a
\emph{trajectory}, and the picture of all the arrows $\mathbf{f}(\mathbf{x})$ is called the
\emph{direction field}. A point $\mathbf{x}^{*}$ with $\mathbf{f}(\mathbf{x}^{*})=0$ is
called an \emph{equilibrium point} or a fixed point.
\end{definisi}

The horizontal axis of the phase plane represents position, the vertical axis represents
velocity, and one point means ``here, with this speed''. Time is not drawn; it is only the
parameter that walks along the curve. That is its advantage: the whole history, for all
initial points at once, appears as a single picture.

The equilibrium points of the pendulum are easy to find. $\theta'=0$ demands $v=0$, and
$v'=0$ demands $\sin\theta=0$, so $\theta=k\pi$. At $\theta=0$ the pendulum hangs straight
down; at $\theta=\pi$ it stands straight up. The first is stable in the everyday sense, the
second is not: the slightest touch brings it down.

\begin{untukanak}
Go to the park, find a swing, and sit down. Ask a friend to push you gently, only a little.
Count ten swings by counting inside your head: one, two, three, up to ten. Now ask for a push
that is much higher, until the swing rises almost as high as the upper bar, then count ten
swings again.

What you will see: the high swing takes longer for one full round, even though you feel as
though the high swing is faster. That feeling cheats because the distance traveled is much
longer, while the time for one round is almost the same.

Now look at the top of the swing. When you are at the highest point, there is a short moment
when you almost stop and hang. At the lowest point, you move fastest. A swing never stops in
the middle, and is never fastest at the ends. Those two things always change places.

In this chapter we shall learn why a high swing is slower, and why a pendulum standing
straight up is almost impossible to keep there.
\end{untukanak}

Now the exact period formula. Multiply $v'=-(g/L)\sin\theta$ by $v=\theta'$
and integrate:
\[
\frac{d}{dt}\Bigl(\frac{v^{2}}{2}+\frac{g}{L}(1-\cos\theta)\Bigr)=0 .
\]
That bracket is the energy per unit mass, and it is constant along the trajectory.
With $E$ denoting its value,
\[
\frac{v^{2}}{2}+\frac{g}{L}\bigl(1-\cos\theta\bigr)=E ,
\]
so $v=\pm\sqrt{2E-2(g/L)(1-\cos\theta)}$. The trajectories in the phase plane are those
energy curves; in Figure~\ref{fig:bidang-fase-pendulum} they are drawn for the unit
$g/L=1$.

\begin{figure}[htbp]
\centering
\begin{tikzpicture}[x=1.05cm,y=1.0cm,font=\footnotesize]
  \draw[->,layergrey] (-4.15,0) -- (4.15,0) node[right] {$\theta$};
  \draw[->,layergrey] (0,-2.95) -- (0,2.95) node[above] {$v$};
  \foreach \tick/\lab in {-3.1416/{$-\pi$},-1.5708/{$-\pi/2$},
                          1.5708/{$\pi/2$},3.1416/{$\pi$}}{
    \draw[layergrey!55] (\tick,0.07) -- (\tick,-0.07);
    \node[layergrey,font=\scriptsize,below] at (\tick,-0.07) {\lab};}
  \node[layergrey,font=\scriptsize,below left] at (-0.06,0.02) {$0$};
  \foreach \th in {-3.5,-2.6,-1.7,-0.85,0.85,1.7,2.6,3.5}{
    \foreach \w in {-2.4,-1.6,-0.8,0,0.8,1.6,2.4}{
      \pgfmathsetmacro{\nn}{sqrt(\w*\w+sin(\th r)*sin(\th r))}
      \pgfmathsetmacro{\dx}{0.32*\w/\nn}
      \pgfmathsetmacro{\dy}{-0.32*sin(\th r)/\nn}
      \draw[-{Latex[length=0.9mm]},layerteal!60]
        (\th,\w) -- ({\th+\dx},{\w+\dy});}}
  \foreach \En/\tm in {0.4/0.9273,1.0/1.5708,1.6/2.2143}{
    \draw[layerindigo,thick]
      plot[domain=-\tm:\tm,samples=121,variable=\t]
        ({\t},{sqrt(max(0,2*(\En-1+cos(\t r))))});
    \draw[layerindigo,thick]
      plot[domain=-\tm:\tm,samples=121,variable=\t]
        ({\t},{-sqrt(max(0,2*(\En-1+cos(\t r))))});}
  \draw[accent,very thick]
    plot[domain=-3.1416:3.1416,samples=161,variable=\t]
      ({\t},{2*cos(\t r/2)});
  \draw[accent,very thick]
    plot[domain=-3.1416:3.1416,samples=161,variable=\t]
      ({\t},{-2*cos(\t r/2)});
  \draw[layerplum,thick,densely dashed]
    plot[domain=-3.6:3.6,samples=161,variable=\t]
      ({\t},{sqrt(2*(1.6+cos(\t r)))});
  \fill[ink] (0,0) circle (1.6pt);
  \node[ink,anchor=north east,font=\scriptsize] at (-0.12,-0.52) {center};
  \fill[accent] (3.1416,0) circle (1.6pt);
  \fill[accent] (-3.1416,0) circle (1.6pt);
  \node[accent,anchor=north west,font=\scriptsize,align=left]
    at (2.35,-0.18) {saddle};
  \node[accent,anchor=west,font=\scriptsize,align=left]
    at (2.30,1.30) {separatrix $E=2$};
  \node[layerplum,anchor=west,font=\scriptsize,align=left]
    at (0.55,2.72) {open trajectory $E=2.6$};
\end{tikzpicture}
\caption{The phase plane of the pendulum in units $g/L=1$, with the horizontal axis the angle
$\theta$ and the vertical axis the velocity $v$. The small arrows show the direction field
$(\theta',v')=(v,-\sin\theta)$. The dark blue curves are three closed trajectories with
energy $E=0.4$, $1.0$, and $1.6$; the larger $E$, the wider the turn. The red curve is
the separatrix $E=2$, the trajectory that reaches exactly $\theta=\pm\pi$ with zero velocity;
it separates closed swings from open rotations. The dashed curve is the open trajectory
$E=2.6$, where the pendulum rotates fully without ever stopping. The point at $\theta=0$ is
a center, the points at $\theta=\pm\pi$ are saddles.}
\label{fig:bidang-fase-pendulum}
\end{figure}

That figure contains the entire qualitative behavior of the pendulum without a single
equation being solved; all that is needed is the sign of $\sin\theta$ and the symmetry of the
energy.

For small amplitude, $\sin\theta\approx\theta$ and the system becomes the harmonic
oscillator
\[
\theta'=v,\qquad v'=-\omega_{0}^{2}\theta,\qquad \omega_{0}=\sqrt{g/L},
\]
with solution $\theta(t)=\theta_{0}\cos(\omega_{0}t)$ and period
$2\pi/\omega_{0}=2.006409$ seconds for $L=1$ meter. Its trajectories are the ellipses
$v^{2}+\omega_{0}^{2}\theta^{2}=\text{constant}$, and those are the three small rings near
the center in the figure above.

The exact period comes from energy conservation, not from linearization.
With $E=(g/L)(1-\cos\theta_{0})$ for a pendulum released from rest at
$\theta_{0}$, the time to fall from $\theta_{0}$ to zero is
\[
\frac{T}{4}=\sqrt{\frac{L}{2g}}\int_{0}^{\theta_{0}}
\frac{d\theta}{\sqrt{\cos\theta-\cos\theta_{0}}}
=\sqrt{\frac{L}{g}}\,K(k),\qquad k=\sin\frac{\theta_{0}}{2},
\]
with $K$ the complete elliptic integral of the first kind. For $\theta_{0}=60^{\circ}$
we obtain $k=0.5$ and $K$ at the parameter $m=k^{2}=0.25$ has the value
$1.6857503548$, so $T=2.153242$ seconds. Quadrature on that integral gives the same
number to within $9\times10^{-13}$; two different roads, one result. For $\theta_{0}=5^{\circ}$
the ratio is $1.00047617$, that is an error of less than one part in two thousand: that is
where the small formula deserves to be used, and where its limit of validity lies.

\begin{tangan}
Take a string one meter long, a metal nut as the weight, and a stopwatch on the phone. Tie
the nut to the end of the string, stick the other end to the door frame with tape, and
measure the length of the string from the suspension point to the center of the nut: try to
make it exactly one meter. Pull the nut aside only about five centimeters, release it, and
record the time for ten full swings. Divide by ten; the result is about two seconds.

Now pull much higher, roughly half a meter above the lowest point, and repeat. Compare the
two numbers. What you will see: the second number is larger, but not twice as large. The
difference is about seven percent, and that agrees with the computation above if the angle is
about sixty degrees.

Repeat with a string a quarter of a meter long: the period becomes half, because the period
depends on the square root of the length.
\end{tangan}

\section{Stability: eigenvalues, linearization, and Lyapunov}
\tingkat{advanced}

The question most often asked about a system is not ``what is its value at $t=7$'' but ``does it
come back after being disturbed''. To answer it we do not need the solution; the derivatives
around the equilibrium point are enough.

\begin{definisi}[Stability]\index{stability}
Let $\mathbf{x}^{*}$ be an equilibrium point of $\mathbf{x}'=\mathbf{f}(\mathbf{x})$.
That point is called \emph{stable} if for every $\varepsilon>0$ there is
$\delta>0$ such that $\|\mathbf{x}(0)-\mathbf{x}^{*}\|<\delta$ guarantees
$\|\mathbf{x}(t)-\mathbf{x}^{*}\|<\varepsilon$ for every $t\ge0$. The point is
called \emph{asymptotically stable} if it is stable and there is $\delta>0$ such that
$\mathbf{x}(t)\to\mathbf{x}^{*}$ as $t\to\infty$. If it is not stable, it is
called \emph{unstable}.
\end{definisi}

\begin{teorema}[Linearization]\index{linearization}
Suppose $\mathbf{f}$ has continuous partial derivatives near the equilibrium point
$\mathbf{x}^{*}$ and suppose $A=J_{\mathbf{f}}(\mathbf{x}^{*})$ is the Jacobian matrix from
Chapter 12. If no eigenvalue of $A$ has real part zero (a \emph{hyperbolic} equilibrium
point), then the qualitative behavior of the nonlinear system near $\mathbf{x}^{*}$ is the
same as that of the linear system $\mathbf{u}'=A\mathbf{u}$: the type of stability is the
same, and nearby trajectories can be mapped one-to-one and continuously onto linear
trajectories.
\end{teorema}

This is the Hartman--Grobman theorem; its proof is long and uses the fixed point theorem on a
space of functions, so I cite it: Hartman, chapter 9, Perko, section 2.8. The practical
consequence: local stability is read off from the eigenvalues.

For a two-dimensional system, every possibility is read from the trace
$\tau=\operatorname{tr}A$ and the determinant $\Delta=\det A$; the eigenvalues are the roots
of $\lambda^{2}-\tau\lambda+\Delta=0$. If $\Delta<0$, the roots are real and of
opposite sign: a \emph{saddle}. If $\Delta>0$ and $\tau^{2}\ge4\Delta$,
both are real and of the same sign: a \emph{node}, stable if $\tau<0$. If
$\tau^{2}<4\Delta$, the roots are conjugates with real part $\tau/2$:
a \emph{spiral}, stable if $\tau<0$, and a center if $\tau=0$.

Example: the matrix $A=\begin{pmatrix}0&1\\-2&-3\end{pmatrix}$ has $\tau=-3$
and $\Delta=2$, so its roots are $\lambda^{2}+3\lambda+2=0$, namely $-1$ and
$-2$: a stable node. The matrix $B=\begin{pmatrix}0&1\\-2&0\end{pmatrix}$ has
$\tau=0$ and $\Delta=2>0$, its roots are $\pm i\sqrt{2}$: a center.

Figure~\ref{fig:potret-linear} shows the two portraits that differ most in character: the
saddle, where almost every initial point moves away, and the stable spiral, where everything
spirals inward.

\begin{figure}[htbp]
\centering
\begin{tikzpicture}[font=\footnotesize]
  \begin{scope}[x=1.35cm,y=1.35cm]
    \begin{scope}
      \clip (-1.62,-1.52) rectangle (1.62,1.52);
      \foreach \c in {0.05,0.30,0.70,1.20}{
        \draw[layerteal,thick] plot[domain=-1.55:1.55,samples=90,variable=\s]
          ({sqrt(\c+\s*\s)},{\s});
        \draw[layerteal,thick] plot[domain=-1.55:1.55,samples=90,variable=\s]
          ({-sqrt(\c+\s*\s)},{\s});}
      \draw[layerindigo,densely dashed] (-1.5,-1.5) -- (1.5,1.5);
      \draw[layerindigo,densely dashed] (-1.5,1.5) -- (1.5,-1.5);
    \end{scope}
    \draw[->,layergrey] (-1.75,0) -- (1.75,0) node[right] {$x$};
    \draw[->,layergrey] (0,-1.55) -- (0,1.55) node[above] {$y$};
    \fill[ink] (0,0) circle (1.5pt);
    \node[layerindigo,font=\scriptsize,anchor=south east] at (-1.52,1.30)
      {eigen directions};
    \node[layergrey,anchor=north] at (0,-1.62)
      {(a) saddle: $\lambda=1,-1$};
  \end{scope}
  \begin{scope}[xshift=6.6cm,x=1.35cm,y=1.35cm]
    \draw[->,layergrey] (-1.75,0) -- (1.75,0) node[right] {$x$};
    \draw[->,layergrey] (0,-1.55) -- (0,1.55) node[above] {$y$};
    \foreach \r in {0.35,0.7,1.1,1.5}{
      \draw[accent,thick]
        plot[domain=0:12.6,samples=200,variable=\t]
          ({\r*exp(-0.3*\t)*cos(\t r)},{\r*exp(-0.3*\t)*(-sin(\t r))});}
    \fill[ink] (0,0) circle (1.5pt);
    \node[layergrey,anchor=north] at (0,-1.62)
      {(b) stable spiral: $\lambda=-0.3\pm i$};
  \end{scope}
\end{tikzpicture}
\caption{Two linear phase portraits of different types. (a) The system
$x'=y$, $y'=x$: trajectories satisfy $x^{2}-y^{2}=\text{constant}$, hyperbolas
with asymptotes $y=\pm x$, and the dashed lines are the eigen directions, one in
and one out, so the origin is a saddle and is unstable. (b) The system $x'=-0.3x+y$, $y'=-x-0.3y$:
the eigenvalues are $-0.3\pm i$, so trajectories spiral while shrinking by the
factor $e^{-0.3t}$ and the origin is a stable spiral. The four trajectories in (b)
differ only in their initial radius, because the system is linear.}
\label{fig:potret-linear}
\end{figure}

One warning that is expensive to learn. If $\tau=0$ and $\Delta>0$, linearization gives a
center, but the nonlinear system at that point is not necessarily a center. Take
\[
x'=-y+x(x^{2}+y^{2}),\qquad y'=x+y(x^{2}+y^{2}).
\]
Its Jacobian at the origin is a rotation matrix, its eigenvalues are $\pm i$:
linearization says center. But in polar coordinates $r^{2}=x^{2}+y^{2}$,
the cross multiplication gives
\[
\frac{1}{2}\frac{d}{dt}r^{2}=xx'+yy'=r^{2}(x^{2}+y^{2})=r^{4},
\]
so $r'=r^{3}>0$ for $r>0$: trajectories spiral outward, and the origin is not
stable. Linearization is not wrong; it is only insufficient when there is an eigenvalue on
the imaginary axis. That is where the theorem above demands a hyperbolic equilibrium point.

For the non-hyperbolic case, a second tool is used: a decreasing energy function, which we
compose ourselves.

\begin{definisi}[Lyapunov function]\index{Lyapunov function}
A function $V$ that is continuous and has continuous derivatives on a neighborhood $U$ of
$\mathbf{x}^{*}$ is called a \emph{Lyapunov function} if $V(\mathbf{x}^{*})=0$ and
$V(\mathbf{x})>0$ for $\mathbf{x}\in U\setminus\{\mathbf{x}^{*}\}$. Its derivative along a
trajectory is written
\[
\dot V(\mathbf{x})=\nabla V(\mathbf{x})\cdot\mathbf{f}(\mathbf{x}).
\]
\end{definisi}

That artificial hill is $V$, and $\dot V$ is the steepness of its descent.

\begin{teorema}[Lyapunov test]
Suppose $V$ is a Lyapunov function near the equilibrium point $\mathbf{x}^{*}$.
\begin{enumerate}
\item If $\dot V\le0$ on $U$, then $\mathbf{x}^{*}$ is stable.
\item If $\dot V<0$ on $U\setminus\{\mathbf{x}^{*}\}$, then
  $\mathbf{x}^{*}$ is asymptotically stable.
\end{enumerate}
\end{teorema}

\begin{proof}
For the first statement, take $\varepsilon>0$ small enough that the ball
$\|\mathbf{x}-\mathbf{x}^{*}\|\le\varepsilon$ stays inside $U$. On that ball
$V$ is positive and the ball is compact, so
$m=\min_{\|\mathbf{x}-\mathbf{x}^{*}\|=\varepsilon}V>0$. Choose $\delta>0$
such that $V(\mathbf{x})<m$ for $\|\mathbf{x}-\mathbf{x}^{*}\|<\delta$;
this is possible because $V$ is continuous and $V(\mathbf{x}^{*})=0$. A trajectory starting
in the $\delta$ ball satisfies $V(\mathbf{x}(t))\le V(\mathbf{x}(0))<m$, since $\dot V\le0$,
so it cannot cross the $\varepsilon$ ball; hence it is stable.

For the second statement, $V(\mathbf{x}(t))$ decreases and is bounded below by
zero, so it has a limit $L\ge0$. Suppose $L>0$; the trajectory then stays in the compact set
$\{V\ge L\}$, and there $-\dot V$ has a positive minimum, so $V$ would have to fall through
$L$ in finite time. Contradiction. Hence $L=0$, and because $V$ vanishes only at
$\mathbf{x}^{*}$, the trajectory heads there.
\end{proof}

An example is the damped pendulum, $\theta'=v$, $v'=-(g/L)\sin\theta-bv$ with
$b>0$, and $V=v^{2}/2+(g/L)(1-\cos\theta)$, which is the mechanical energy, positive
except at the equilibrium point. Its derivative is
\[
\dot V=v\Bigl(-\frac{g}{L}\sin\theta-bv\Bigr)+\frac{g}{L}\sin\theta\,v=-bv^{2}\le0 .
\]
So the lower equilibrium point is stable and the energy never increases. This test
does not give asymptotic stability, because $\dot V$ is zero along the whole line $v=0$;
the LaSalle version adds that the only invariant set inside $\{\dot V=0\}$ is the equilibrium
point, so all trajectories head there.

Linearization gives the same answer. At $\theta=0$ the Jacobian is
$\begin{pmatrix}0&1\\-g/L&-b\end{pmatrix}$, with $\tau=-b$ and
$\Delta=g/L>0$; for $b=0.5$ and $g/L=9.80665$ the roots are $-0.25\pm3.121562\,i$,
that is a stable spiral. At $\theta=\pi$ a $+g/L$ appears in the lower right corner,
so $\Delta=-g/L<0$: a saddle with eigenvalues $\pm3.131557$. A pendulum standing
straight up is indeed a saddle, and not as a figure of speech: there is exactly one direction
that keeps it up there, and every disturbance carries it away.

\kamus{asymptotically stable}{stabil asimtotik}{Stays close after a disturbance, then actually arrives back at the equilibrium point.}

\kamus{Lyapunov function}{fungsi Lyapunov}{An artificial hill used as a ruler: as long as the system descends it, the system cannot run far away.}

\section{Euler, Runge--Kutta, and stiffness}
\tingkat{advanced}

Almost no equation can be solved by a formula, so numerical methods are not a supplement but
the main road. The idea goes back to Chapter 12: if we know the value and the derivative at
one point, we can estimate the value a little to its right.

Euler's method does the simplest thing. From $y'=f(t,y)$ and the initial value,
\[
y_{n+1}=y_{n}+h\,f(t_{n},y_{n}),\qquad t_{n+1}=t_{n}+h .
\]
It is exactly the first two terms of the Taylor series around $t_{n}$, so its one-step error
is the term $h^{2}y''(\xi)/2$, of order $h^{2}$. That local error accumulates into a global
error of order $h$: doubling the accuracy demands twice as many steps, not four times.

Numbers for $y'=y$, $y(0)=1$, at $t=1$ with $h=0.1$: Euler gives
$(1.1)^{10}=2.5937424601$, whereas $e=2.7182818285$. The error is $0.1245$, that is
$4.58$ percent, and the error of the first single step is $5.171\times10^{-3}$.

The fourth-order Runge--Kutta method repairs this by taking four samples of the slope at
every step:
\[
y_{n+1}=y_{n}+\frac{h}{6}\bigl(k_{1}+2k_{2}+2k_{3}+k_{4}\bigr),
\]
with $k_{1}=f(t_{n},y_{n})$, $k_{2}=f(t_{n}+h/2,\,y_{n}+hk_{1}/2)$,
$k_{3}=f(t_{n}+h/2,\,y_{n}+hk_{2}/2)$, and $k_{4}=f(t_{n}+h,\,y_{n}+hk_{3})$.
Its local error is of order $h^{5}$ and its global error of order $h^{4}$: half the step
means the error falls by about sixteen times.

The following table comes from the program in the lab section, on the logistic model with
$y(0)=1/2$ and $t=2$; the exact value is $0.8807970780$.

\begin{center}
{\small
\begin{tabular}{lrrr}
\hline
$h$ & Euler error & RK4 error & Euler ratio \\
\hline
$0.400$ & $1.958\times10^{-2}$ & $1.245\times10^{-5}$ & --- \\
$0.200$ & $9.423\times10^{-3}$ & $7.237\times10^{-7}$ & $2.08$ \\
$0.100$ & $4.630\times10^{-3}$ & $4.359\times10^{-8}$ & $2.04$ \\
$0.050$ & $2.296\times10^{-3}$ & $2.674\times10^{-9}$ & $2.02$ \\
$0.025$ & $1.143\times10^{-3}$ & $1.656\times10^{-10}$ & $2.01$ \\
\hline
\end{tabular}}
\end{center}

The last column is the proof of first order: every time $h$ is halved, the Euler error falls
by almost exactly two; for RK4 the ratio is about $16$, in keeping with fourth order.
Figure~\ref{fig:galat-numerik} displays the same numbers on a logarithmic scale on both
axes.

\begin{figure}[htbp]
\centering
\begin{tikzpicture}
\begin{axis}[width=0.80\linewidth,height=5.6cm,
  xlabel={step $h$},ylabel={$|y(2)-y_{h}(2)|$},
  xmode=log,ymode=log,xmin=0.02,xmax=0.5,ymin=1e-11,ymax=1e-1,
  grid=both,minor tick num=1,
  legend style={font=\scriptsize,at={(0.5,-0.22)},anchor=north,
    legend columns=2,draw=none,fill=none},
  tick label style={font=\scriptsize},label style={font=\small},
  axis line style={layergrey},tick style={layergrey}]
  \addplot[layerindigo,thick,mark=*,mark size=1.7pt] coordinates {
    (0.4,0.01958) (0.2,0.009423) (0.1,0.004630)
    (0.05,0.002296) (0.025,0.001143)};
  \addplot[accent,thick,mark=square*,mark size=1.6pt] coordinates {
    (0.4,0.00001245) (0.2,0.0000007237) (0.1,0.00000004359)
    (0.05,0.000000002674) (0.025,0.0000000001656)};
  \legend{Euler (slope $1$), Runge--Kutta 4 (slope $4$)}
\end{axis}
\end{tikzpicture}
\caption{The error at $t=2$ for the logistic model as a function of the step $h$, on a
logarithmic scale on both axes. The Euler line falls with slope one: the error is
proportional to $h$. The Runge--Kutta line falls with slope four:
the error is proportional to $h^{4}$. At $h=0.4$ the two methods are about three
orders apart, and that distance widens when $h$ is made smaller. The numbers in the figure are
the same numbers as in the table in the text.}
\label{fig:galat-numerik}
\end{figure}

Now the property that deceives many people: stiffness. A system
is called \emph{stiff} if inside it there is a process that moves
far faster than the process we are interested in. Its practical signature: explicit methods
demand a very small step in order to remain \emph{stable}, even when the accuracy we ask for
demands only a large step.

The example is small and honest: $y'=-50(y-\cos t)-\sin t$, $y(0)=1$, whose exact
solution is $y=\cos t$. Here $y$ follows $\cos t$ slowly, but every
deviation is dragged back at the rate $50$ per second. Euler is stable only if
$|1-50h|<1$, that is $h<0.04$.

Results at $t=3$, with exact value $\cos3=-0.989992$: Euler with $h=0.02$
gives an error of $1.98\times10^{-4}$, with $h=0.039$ it gives
$7.97\times10^{-4}$, and with $h=0.05$ it gives $-1.839\times10^{7}$. One
step up from $0.039$ to $0.05$ changes the answer from almost right to
meaningless. RK4 holds out longer, because its stability region is wider: at
$h=0.05$ its error is $1.11\times10^{-3}$, but at $h=0.1$ it gives
$-1.32\times10^{31}$; its practical limit is $2.785/50=0.0557$.

The lesson to carry away: error and stability are two different things. The local error on
the first step is small and warns us of nothing; what explodes is
stability, and that is a matter of eigenvalues.

The antidote is implicit methods, the simplest being backward Euler:
\[
y_{n+1}=y_{n}+h\,f(t_{n+1},y_{n+1}),
\]
which demands solving an equation at every step, but is stable for large $h$
on stiff equations; for the example above, backward Euler with $h=0.5$
gives an answer close to $\cos t$. Reference: Hairer and Wanner, chapter IV.

\begin{tantangan}
Linearization gives a center, but a nonlinear system need not be a center. Construct a
two-dimensional system whose origin is \emph{not} a center even though its linearization is a
center, determine the character of that point (stable, asymptotically stable, or unstable),
and include an explicit formula, the eigenvalues of the Jacobian at the origin, and a brief
proof of your claim.

Hint 1: linearization sees only terms of degree one. Add a degree-three term that is
symmetric under rotation, then write the system in polar coordinates.

Hint 2: in polar coordinates, the rotation matrix becomes $\theta'=1$, and the radial
disturbance becomes $r'=cr^{3}$ for some constant $c$. The sign of $c$
decides everything, and $V=r^{2}$ is a Lyapunov function that works for
one of the two signs.

If you get stuck, look again at the example $r'=r^{3}$ in the Stability section.
\end{tantangan}
\section{Worked examples}
\tingkat{intermediate}

The four computations below are the backbone of this chapter; all of them can be checked
again with a calculator or with the program in the lab section, and all of them have gone
wrong in my own hands.

\textbf{Logistic model: one formula, two methods.}
Take $y'=y(1-y)$ with $y(0)=1/2$ and ask for the value at $t=2$. The formula from
the separation of variables section gives $y(2)=1/(1+e^{-2})=0.8807970780$. With
$h=0.1$ there are twenty steps: Euler gives $0.8854273272$, an error of
$4.630\times10^{-3}$ or $0.53$ percent, whereas Runge--Kutta with the same step
gives $0.8807970344$, an error of $4.359\times10^{-8}$. The cost is four times
as much, the accuracy about one hundred thousand times.

Note also the direction of the error. Euler is too high here, because
$y(1-y)$ is concave down on that interval, so a straight line always lies above the curve.
Checking the sign of the error is a cheap way to find out whether our program
works as it should.

\textbf{Variable coefficient: $y'+y/t=t^{2}$, $y(1)=0$.}
Here $p(t)=1/t$, so
\[
\mu(t)=\exp\Bigl(\int\frac{dt}{t}\Bigr)=t
\]
on $t>0$. Multiplying the equation by $t$ gives
\[
t\,y'+y=t^{3}
\quad\Longleftrightarrow\quad
\bigl(t\,y\bigr)'=t^{3},
\]
so $t\,y=t^{4}/4+C$ and
\[
y(t)=\frac{t^{3}}{4}+\frac{C}{t}.
\]
The condition $y(1)=0$ gives $C=-1/4$, so
\[
y(t)=\frac{t^{3}}{4}-\frac{1}{4t}=\frac{t^{4}-1}{4t},
\qquad y(2)=\frac{15}{8}=1.875 .
\]
Check: $y'=3t^{2}/4+1/(4t^{2})$ and $y/t=t^{2}/4-1/(4t^{2})$, their sum is
$t^{2}$, exactly the right-hand side. An adaptive integrator gives $1.875000000000$.

Without $\mu=t$, the left-hand side $y'+y/t$ does not look like the derivative of anything;
after multiplication by $t$, it becomes the derivative of $ty$.

\textbf{How large a swing can be before the small formula misses by one percent.}
The ratio of the exact period to the small period is
\[
\frac{T(\theta_{0})}{T_{0}}
=\frac{2}{\pi}K\Bigl(\sin^{2}\frac{\theta_{0}}{2}\Bigr),
\]
with $K$ the complete elliptic integral. Some of its values:

\begin{center}
{\small
\begin{tabular}{rrrr}
\hline
$\theta_{0}$ & $T/T_{0}$ & error & remark \\
\hline
$5^{\circ}$ & $1.000476$ & $0.048$ percent & pendulum clock \\
$10^{\circ}$ & $1.001907$ & $0.191$ percent & still safe \\
$20^{\circ}$ & $1.007669$ & $0.767$ percent & nearing the limit \\
$22.81^{\circ}$ & $1.010000$ & $1.000$ percent & the one-percent limit \\
$30^{\circ}$ & $1.017409$ & $1.741$ percent & correction needed \\
$60^{\circ}$ & $1.073182$ & $7.318$ percent & clock off by an hour \\
\hline
\end{tabular}}
\end{center}

So the small formula is fit to be used up to about twenty degrees, and
$22.81^{\circ}$ is the one-percent error threshold; the error grows like the square of the
amplitude, not like the amplitude.
Doubling the amplitude from $10^{\circ}$ to $20^{\circ}$ raises the error from
$0.19$ to $0.77$ percent, that is four times, in keeping with
$1+(\theta_{0}^{2}/16)$ as the first two terms of that estimate.

\kertas{
Test yourself in three steps. (1) Work the logistic example again with
$h=0.2$; you should get an Euler error of $9.423\times10^{-3}$ and an RK4 error of
$7.237\times10^{-7}$. (2) Derive the elliptic period formula yourself from energy
conservation, then check that $K$ at the parameter $0.25$ has the value
$1.6857503548$. (3) Find the threshold $h$ for Euler on $y'=-100y$; the answer is
$h<0.02$, and at $h=0.02$ exactly, the method oscillates between $+1$ and $-1$
without either decaying or exploding.}

\section{Computational lab}
\tingkat{advanced}

The following code computes all the numbers above: the Euler and RK4 errors on the logistic
model, the pendulum period by two different roads, and the stiffness threshold. It
uses only NumPy and SciPy.

\begin{lab}
\begin{verbatim}
import numpy as np
from scipy import integrate, special

def euler(f, y0, T, h):
    y, t, n = y0, 0.0, int(round(T / h))
    for _ in range(n):
        y += h * f(t, y); t += h
    return y

def rk4(f, y0, T, h):
    y, t, n = y0, 0.0, int(round(T / h))
    for _ in range(n):
        k1 = f(t, y); k2 = f(t + h/2, y + h*k1/2)
        k3 = f(t + h/2, y + h*k2/2); k4 = f(t + h, y + h*k3)
        y += h * (k1 + 2*k2 + 2*k3 + k4) / 6; t += h
    return y

logi = lambda t, y: y * (1 - y)
ex = lambda t: 1 / (1 + np.exp(-t))
print("logistic y(2) exact %.10f" % ex(2.0))
for h in (0.4, 0.2, 0.1, 0.05, 0.025):
    ee = abs(euler(logi, 0.5, 2, h) - ex(2))
    er = abs(rk4(logi, 0.5, 2, h) - ex(2))
    print("h=%.3f euler %.3e rk4 %.3e" % (h, ee, er))

g, L = 9.80665, 1.0
T0 = 2 * np.pi * np.sqrt(L / g)
th0 = np.radians(60.0)
k = np.sin(th0 / 2)
Tq = 4 * np.sqrt(L / (2*g)) * integrate.quad(
        lambda th: 1/np.sqrt(np.cos(th) - np.cos(th0)), 0, th0, limit=200)[0]
Te = 4 * np.sqrt(L / g) * special.ellipk(k**2)
print("T0 %.6f T60quad %.6f T60elliptic %.6f" % (T0, Tq, Te))
print("ratio %.6f percent %.4f" % (Te / T0, 100 * (Te / T0 - 1)))

stiff = lambda t, y: -50 * (y - np.cos(t)) - np.sin(t)
for h in (0.02, 0.039, 0.05):
    print("stiff euler h=%.3f y(3)=%.3e" % (h, euler(stiff, 1.0, 3, h)))
for h in (0.05, 0.1):
    print("stiff rk4 h=%.3f y(3)=%.3e" % (h, rk4(stiff, 1.0, 3, h)))
print("exact cos(3) %.10f" % np.cos(3))
\end{verbatim}
\end{lab}

The actual output, verbatim:

\begin{verbatim}
logistic y(2) exact 0.8807970780
h=0.400 euler 1.958e-02 rk4 1.245e-05
h=0.200 euler 9.423e-03 rk4 7.237e-07
h=0.100 euler 4.630e-03 rk4 4.359e-08
h=0.050 euler 2.296e-03 rk4 2.674e-09
h=0.025 euler 1.143e-03 rk4 1.656e-10
T0 2.006409 T60quad 2.153242 T60elliptic 2.153242
ratio 1.073182 percent 7.3182
stiff euler h=0.020 y(3)=-9.902e-01
stiff euler h=0.039 y(3)=-9.908e-01
stiff euler h=0.050 y(3)=-1.839e+07
stiff rk4 h=0.050 y(3)=-9.883e-01
stiff rk4 h=0.100 y(3)=-1.317e+31
exact cos(3) -0.9899924966
\end{verbatim}

Three things deserve reading. First, the Euler error ratios $2.08$, $2.04$, $2.02$,
$2.01$ and RK4 heading to $16=2^{4}$: first and fourth order are visible as numbers. Second,
the two roads to the pendulum period meet to six digits, even though quadrature and the
elliptic integral do not use the same formula. Third, the stiff lines: $h=0.039$
is still reasonable, $h=0.05$ is meaningless, and not one of those lines warns
us.

The program does not check the hypotheses of the theorems: whether $f$ satisfies the Lipschitz
condition, whether the equilibrium point is hyperbolic. That remains our job.

\begin{untukdosen}
The teaching order I use: start from the pendulum clock and the $7.32$
percent difference, then the definition of a solution along with the blow-up example
$1/(1-t)$, then separation of variables and the integrating factor. The integrating factor is
best proved in full at the board, because its proof uses only the product rule and the
fundamental theorem of calculus, two tools they have had since Chapter 13. Picard--Lindelof is
enough as a sketch with one iteration carried out; what must remain is the meaning of the
Lipschitz condition, not the fixed point technique. Phase plane before eigenvalues:
draw first, compute later.

Two misconceptions almost always appear. First, students take
``solution'' to mean a closed formula, then conclude that an equation without a formula has no
solution; the example $y'=\sqrt{|y|}$ helps, because there two solutions exist and
no single formula does. Second, they treat a small error as proof of correctness, and are
surprised when Euler with $h=0.05$ gives $-1.8\times10^{7}$.
Show them that output line; do not merely warn them about it.

One blackboard activity: draw the pendulum phase plane empty, then ask
the students to mark where the velocity is zero, where the acceleration is zero, and where the
energy is highest. The class will argue about the region between the two saddles, and that
argument is exactly the content of the theorem. A sign of readiness to go further: students
can guess the type of an equilibrium point from the trace and the determinant without
computing the roots.

\end{untukdosen}

\section{Summary}
\begin{ringkasan}
\begin{itemize}
\item A differential equation sets the rate of change equal to the state; its solution is a
function on an interval, and that interval is part of the answer.
\item Two families are solvable by hand: separation of variables and the first-order linear
equation with integrating factor $\mu=\exp(\int p\,dt)$.
\item Picard--Lindel\"of guarantees a unique solution on a short interval if $f$ is
continuous and satisfies the Lipschitz condition with respect to $y$; without that condition,
one initial state can give birth to two histories.
\item The phase plane replaces time with a picture: equilibrium points, closed trajectories,
separatrices, and open trajectories are read without any formula for the solution.
\item Local stability is read from the eigenvalues of the Jacobian, except when an eigenvalue
lies on the imaginary axis; there linearization does not decide, and a Lyapunov function is
used.
\item Euler is first order globally and RK4 is fourth order; both are explicit and therefore
fail on stiff systems, where stability, not accuracy, sets the size of the step.
\end{itemize}
\end{ringkasan}

\section{Exercises}
\begin{latihan}
\begin{enumerate}
\item Solve $y'=2ty$, $y(0)=1$, then compute $y(1)$ and compare it with
  $e$. Also compute Euler with $h=0.1$ and state its error. (My result:
  $2.33463336$, error $3.84\times10^{-1}$.)
\item Prove that $y'+p(t)y=q(t)$ with $p,q$ continuous has at most
  one solution, without an explicit formula: the difference of two solutions satisfies
  $w'=-p(t)w$ and $w(t_{0})=0$.
\item True, false, or not enough information: ``because linearization at the origin gives
  eigenvalues $\pm i$, the origin is a stable center''. Give one supporting example
  or one counterexample.
\item Prove that $E=v^{2}/2+\omega_{0}^{2}\theta^{2}/2$ is constant for the harmonic
  oscillator, then explain why its trajectories are ellipses and its period is
  $2\pi/\omega_{0}$.
\item An RC circuit with $R=2.2\,\mathrm{k}\Omega$, $C=10\,\mu\mathrm{F}$,
  $V_{s}=12$ volts starts from zero. Compute $\tau$, the voltage at $t=\tau$, and
  the time to reach $6$ volts. ($22$ milliseconds; $7.585447$ volts; $15.249$ milliseconds.)
\item Show that $v=\pm2\cos(\theta/2)$ satisfies $E=2$, then explain
  why that curve joins $(-\pi,0)$ to $(\pi,0)$ and is called a
  separatrix.
\item A pendulum with $L=0.25$ meters is released from $30^{\circ}$. Compute the small
  period and the exact period, then state the percentage difference. ($1.003205$ seconds;
  $1.020669$ seconds; $1.74$ percent.)
\item For $y'=-100y$ with $y(0)=1$, determine the threshold $h$ that keeps explicit Euler
  stable, then explain what happens exactly at that threshold.
  Check it with the program.
\item Prove that the scalar autonomous equation $y'=f(y)$ with $f$ continuous and
  satisfying the Lipschitz condition cannot have a nonconstant periodic solution.
\end{enumerate}

\textbf{Hints for the three hardest items.} Item 3: linearization sees only terms
of degree one; add a symmetric cubic term as in the example $r'=r^{3}$
and compare it with $r'=-r^{3}$. Item 6: use
$\cos\theta=2\cos^{2}(\theta/2)-1$ and check the value of $v$ at $\theta=\pm\pi$.
Item 9: suppose $y$ is periodic and nonconstant; at a maximum $y'=0$, so
$f(y_{\max})=0$, and uniqueness forces that trajectory to be constant.
\end{latihan}

\begin{untukpeneliti}
What is genuinely proved in this chapter, with written proofs: the integrating factor
formula together with its uniqueness, the complete Picard--Lindelof sketch with the error
bounds of the iteration, uniqueness of solutions of the linear equation through $w'=-pw$, the
Lyapunov test for both of its statements, and all the energy identities for the oscillator
and the pendulum. What I cite without proof: the Hartman--Grobman theorem (Hartman,
\emph{Ordinary Differential Equations}, chapter 9; Perko, \emph{Differential
Equations and Dynamical Systems}, section 2.8), Peano's theorem and
Arzel\`a--Ascoli, and the LaSalle invariance principle.

What is overpopularized is the phrase ``the eigenvalues decide everything''. What is true:
eigenvalues determine local behavior only when none of them has real part zero; on the
imaginary axis, degree-three terms can turn a center into a spiral, as in the example
$r'=r^{3}$ here. For high-dimensional systems, local ``asymptotic stability'' does not even
mean safe from large disturbances, because the basin of attraction can be very small.

An open question: how far can qualitative properties be certified numerically? A small error
on a single trajectory says nothing about stability; Euler with $h=0.039$ is right to four
digits, whereas $h=0.05$ gives $10^{7}$. For a system known only through samples of its
field, there is no general way to make sure that no stiff region hides between two sample
points.

Limits of validity: all functions here are real and smooth enough; all numbers are from
Python 3.14.4, NumPy 2.5.2, and SciPy 1.18.0, with quadrature at tolerance
$10^{-13}$ unless stated otherwise; $g=9.80665\ \mathrm{m/s^{2}}$ and friction is
neglected in the period computations.
\end{untukpeneliti}

\peta{Chapter 16 in Boyce and DiPrima, \emph{Elementary Differential Equations
and Boundary Value Problems}, chapters 1 (introduction), 2 (first-order equations), 3
(second-order linear equations), 4 (first-order systems and the phase plane), and 9
(nonlinear equations and stability); in Strogatz, \emph{Nonlinear Dynamics
and Chaos}, chapters 2 (flows on the line), 5 (two-dimensional systems), 6 (phase planes),
7 (oscillators), and 8 (bifurcations); in Arnold, \emph{Ordinary Differential
Equations}, chapters 1--4 for existence, uniqueness, and the classification of equilibrium
points; in Coddington and Levinson, chapters 1--2 for Picard--Lindelof and
Peano; in Hartman, chapter 9, and Perko, section 2.8, for Hartman--Grobman;
in Hairer and Wanner, \emph{Solving Ordinary Differential Equations II}, chapter
IV, for stiffness and implicit methods. The basis for the computations is the chapter
\emph{Derivatives: One and Several Variables} (partial derivatives and the Jacobian), the
chapter \emph{Series of Functions and Exchange of Limits} (uniform limits, which make Picard
iteration converge), and the chapter \emph{Multiple Integrals} (vector fields) in the first
edition of this book; the eigenvalues of Chapter 9 are used directly without being repeated,
and the floating-point error analysis of Chapter 22 is the continuation of the numerical part
of this chapter.}


% ---------------------------- part IV ----------------------------
\part{Probability, Statistics, and Computer Arithmetic}

This part prepares two things that look unrelated: how to measure uncertainty and how a
computer stores numbers. The prerequisites are Chapters~2, 5, and~10. After six chapters you
can compute probabilities, estimate parameters, test claims, fit a regression, and explain
why $0.1+0.2$ does not equal $0.3$.
\chapter{Probability and Conditional Probability}

A clinic advertises a screening test for a rare disease: only one person in a thousand in
that region has it. The test is good. Out of every hundred sick people, ninety-nine come
back positive. Out of every hundred healthy people, only one is wrongly accused of being
positive. The advertisement writes both numbers as ``99 percent accuracy''.

Now the positive result is in our hands. What is the probability that we are really sick?
In a population of $100{,}000$ people, about $100$ have the disease and $99$ of them are
positive. The rest, $99{,}900$ healthy people, contribute $999$ false positives. So there
are $1.098$ positive results and only $99$ that are right: $99/1.098=11/122$,
that is $0.0902$. About one in eleven, not ninety-nine in a hundred.

Both numbers are right and both answer different questions. The advertisement gives the
probability of a positive result if the person is sick. The patient wants to know the
probability that the person is sick if the result is positive. Swapping those two
directions is not an arithmetic mistake; it is a mistake in naming the quantity, and the
consequence is often a factor of ten.

Chapter 2 already gave us a language for events: union, intersection, complement. Chapter 5
gave us a way to sum infinitely many terms. Neither is enough, because a set carries no
weight: $A\cup B$ is the same set whether $A$ is almost certain or almost certainly not.
What we need is three things: a weight function with rules that must not be broken, a way
to count weights, and a rule for updating weights when new information arrives. This
chapter builds the three, then checks them with a seeded simulation.

\begin{untukanak}
Imagine the guard at a school gate checking bags with a device that beeps. The device is
good: out of a hundred bags that really carry a forbidden object, ninety-nine make it
beep. But the device also beeps for one out of a hundred bags whose contents are perfectly
ordinary. Now suppose that at that school only one bag in a thousand really carries a
forbidden object.

The question is simple and surprising. If the device beeps, is that bag certainly
dangerous? It turns out not. Dangerous bags are indeed rare, while ordinary bags are very
many, so the false beeps from ordinary bags far outnumber the true ones. A guard who looks
only at the red light will hold back many innocent children.

The question that makes us curious: what fraction of all those beeps are right? The answer
can be computed without guessing, as long as we know how many bags there are of each kind.
In this chapter we shall learn why a device that is ``ninety-nine percent accurate'' can
turn into a false accusation in nine cases out of ten, and why the direction of the
question determines the answer.
\end{untukanak}

\section{Probability axioms and their consequences}
\tingkat{basic}

Before a single number is computed, we have to agree on what object is being weighted. In
everyday language: probability is the weight we place on each possibility, under the rule
that the total weight of all possibilities is one and no weight is ever negative. That rule
does not explain the philosophical meaning of probability; it fixes the arithmetic that
every computation has to obey.

\kamus{probability}{peluang}{The weight placed on each possibility, with the total weight
of all possibilities equal to one}

\index{probability}
\index{sample space}

\begin{definisi}[Probability space]
Let $\Omega$ be the set of outcomes, $\mathcal{F}$ a collection of events, that is, the
subsets of $\Omega$ that contain $\Omega$ and are closed under complement and countable
union, and $P\colon\mathcal{F}\to\mathbb{R}$. The function $P$ is called a
\emph{probability} if
\begin{enumerate}
  \item $P(A)\ge 0$ for every $A\in\mathcal{F}$;
  \item $P(\Omega)=1$;
  \item if $A_{1},A_{2},\dots$ are pairwise disjoint, then
        $P\bigl(\bigcup_{i\ge 1}A_{i}\bigr)=\sum_{i\ge 1}P(A_{i})$.
\end{enumerate}
\end{definisi}

The third axiom, countable additivity, separates modern probability theory from the
gambling computations of the seventeenth century. Without it, a probability series such as
$\sum_{n\ge 1}2^{-n}=1$ has no foundation, even though that series is exactly what appears
when we ask on which toss the first success occurs. This axiom is also what allows
probability to be used in limits, the subject of Chapter 4 and Chapter 5.

Three immediate consequences are used over and over, so they are worth proving once.

\begin{proposisi}[Consequences of the axioms]
\index{probability axioms}
For events $A$ and $B$ in a probability space we have
$P(A^{c})=1-P(A)$, $P(\varnothing)=0$, and if $A\subseteq B$ then
$P(A)\le P(B)$. Moreover
$P(A\cup B)=P(A)+P(B)-P(A\cap B)$.
\end{proposisi}

\begin{proof}
Since $A$ and $A^{c}$ are disjoint and their union is $\Omega$, the second and third
axioms give $1=P(\Omega)=P(A)+P(A^{c})$. For the second claim, take $A=\varnothing$ in the
first result. For the third, write $B$ as the disjoint union of $A$ and
$B\setminus A$; then $P(B)=P(A)+P(B\setminus A)\ge P(A)$ because probability is not
negative. For the last, split $A\cup B$ into three disjoint parts $A\setminus B$,
$A\cap B$, and $B\setminus A$. In the same way $A$ and $B$ each split into two parts, so
that
\[
P(A\cup B)=P(A\setminus B)+P(A\cap B)+P(B\setminus A)
=P(A)+P(B)-P(A\cap B).\qedhere
\]
\end{proof}

That last property, inclusion--exclusion, is the reason ``or'' may not simply be added.
For example, in a class of $30$ students, $18$ like coffee and $12$ like tea, and $8$ like
both. Then the probability that a randomly chosen student likes coffee or tea is
\[
\frac{18+12-8}{30}=\frac{22}{30}=0.7333,
\]
not $30/30=1$ as the naive sum $18+12=30$ would give. Eight people are counted twice, and
inclusion--exclusion subtracts them once.

For three events the formula extends with alternating signs:
\[
P(A\cup B\cup C)=P(A)+P(B)+P(C)-P(A\cap B)-P(A\cap C)-P(B\cap C)+P(A\cap B\cap C).
\]
The proof applies the two-event version twice, and that is the first exercise of this
chapter. The pattern of signs is not decoration: whatever was subtracted too often has to
be added back.

One warning for Chapter 18. In a finite space where every outcome has positive probability,
$P(A)=0$ means that $A$ is empty; in a continuous space it does not, since a single point
can have probability zero without being impossible. This chapter works in finite spaces.

\begin{contoh}
Two dice are thrown. What is the probability that the sum of the faces is $8$ or that both
faces are equal? With $36$ equally likely outcomes, the first event contains $5$ outcomes,
the second contains $6$, and their intersection is one outcome, namely $(4,4)$. So the
probability is $(5+6-1)/36=10/36=0.2778$. Without subtracting the intersection the answer
is $11/36=0.3056$, too large by $1/36$.
\end{contoh}

\section{Counting and equally likely outcomes}
\tingkat{intermediate}

If all outcomes are equally likely, the probability of an event is just the ratio of the
number of favorable outcomes to the number of all outcomes. In everyday language: divide
the number of ways we want by the number of ways there are. That definition is not a law of
nature; it is an assumption about symmetry, and the assumption has to be checked, not
presupposed.

\index{counting}
\index{combination}

\kamus{combination}{kombinasi}{A choice of several objects without regard to their order,
like taking five cards from one deck}

\begin{proposisi}[Product rule and counting]
\index{product rule}
If $|A|=m$ and $|B|=n$, then $|A\times B|=mn$. The number of sequences of length $k$ from
$n$ distinct objects without repetition is $n!/(n-k)!$, and the number of subsets of size
$k$ is
\[
\binom{n}{k}=\frac{n!}{k!\,(n-k)!}.
\]
In particular $\sum_{k=0}^{n}\binom{n}{k}=2^{n}$.
\end{proposisi}

\begin{proof}
For the product, every pair $(a,b)$ is determined by one choice from $m$ and one choice
from $n$. For sequences there are $n$ choices for the first place, $n-1$ for the second,
and so on up to $n-k+1$. For subsets, every set of size $k$ can be ordered in $k!$ ways, so
the number of sequences divided by $k!$ gives the number of sets. The last identity counts
all subsets in two ways: summing by size, or choosing in or out for every object.
\end{proof}

Neat counting often shortens a long computation. A full house in poker, three cards of one
rank plus a pair of another rank, is a good example: choose the rank of the triple ($13$
ways), choose three of its four suits ($\binom{4}{3}=4$), choose the rank of the pair ($12$
ways), choose two of its four suits ($\binom{4}{2}=6$). The total is
$13\cdot 4\cdot 12\cdot 6=3.744$ out of $\binom{52}{5}=2{,}598{,}960$ hands, so the
probability is $0.001441$, that is one hand in $694.2$. In $694$ hands the expected number
of full houses is $0.9998$, so we wait about seven hundred hands to see the first one.

Counting also explains events that feel miraculous. Take the birthday problem: $k$ people
in one room, every birthday equally likely among $365$ days, and birthdays independent
between people. The probability that no two people share a birthday is
\[
P(\text{no pair})=\frac{365\cdot 364\cdots (365-k+1)}{365^{k}},
\]
so the probability that there is a pair is its complement.

\index{birthday paradox}

\begin{center}
\small
\begin{tabular}{rrr}
\hline
$k$ & $P(\text{no pair})$ & $P(\text{a pair})$ \\
\hline
$10$ & $0.8831$ & $0.1169$ \\
$23$ & $0.4927$ & $0.5073$ \\
$41$ & $0.0968$ & $0.9032$ \\
$50$ & $0.0296$ & $0.9704$ \\
$57$ & $0.0099$ & $0.9901$ \\
\hline
\end{tabular}
\end{center}

At $k=23$ the probability passes one half, and at $k=50$ it is already $0.9704$. The source
of the confusion is the wrong unit: what collides is not people but pairs of people, and
$23$ people form $\binom{23}{2}=253$ pairs. Every pair matches with probability $1/365$, so
the expected number of pairs is $253/365=0.69315$. If collisions are treated as rare and
independent, the probability of no collision approaches $e^{-0.69315}=0.5000$, differing by
$0.0073$ from the exact value $0.4927$. That approximation is not a proof; it explains why
the threshold falls near twenty-three.

The counting above assumes every birthday equally likely. In real populations that
inequality actually raises the probability of a collision; the symmetry assumption here is
an approximation that has to be checked.

Counting cannot answer the clinical question from before. The outcomes ``sick'' and
``healthy'' are not equally likely: their ratio is one to nine hundred ninety-nine. What is
needed is not a way to count outcomes but a way to use the information that someone has
already been declared positive. That is conditional probability.

\section{Conditional probability and the product rule}
\tingkat{intermediate}

New information narrows the world. If we know that $B$ occurred, the relevant world is no
longer $\Omega$ but $B$, and all weights have to be read again inside that narrower world.
In everyday language: conditional probability is the probability recomputed after part of
the possibilities has been crossed out.

\kamus{conditional probability}{peluang bersyarat}{The probability recomputed after we know
that part of the possibilities no longer applies}

\index{conditional probability}

\begin{definisi}[Conditional probability]
If $A$ and $B$ are events with $P(B)>0$, the probability of $A$ given $B$ is
\[
P(A\mid B)=\frac{P(A\cap B)}{P(B)}.
\]
\end{definisi}

Dividing by $P(B)$ is the choice that makes conditional probabilities sum to one over $B$.
If we also demand $P(B\mid B)=1$ and that ratios between events inside $B$ do not change,
that division is the only possible choice; hence this definition is not an arbitrary
convention.

The narrowing can be seen directly in a cross table. Figure~\ref{fig:silang} arranges the
clinical population into four cells. Before any result, the denominator is all $100{,}000$
people and $P(S)=100/100{,}000=0.001$. After a positive result is known, only the $+$
column, which holds $1.098$ people, is still counted, and
$P(S\mid+)=99/1.098=0.0902$. The person does not change; only the set that may still be
counted changes.

    \begin{figure}[htbp]
\centering
\begin{tikzpicture}[font=\small, x=1cm, y=1cm]
  \fill[accent!10] (1.95,1.05) rectangle (3.25,-1.1);
  \draw[layergrey, thick] (1.9,1.05)--(7.3,1.05);
  \draw[layergrey] (1.9,0.35)--(7.3,0.35);
  \draw[layergrey] (1.9,-0.35)--(7.3,-0.35);
  \draw[layergrey, thick] (1.9,-1.1)--(7.3,-1.1);
  \node[anchor=east] at (1.75,1.4) {test result};
  \node at (2.6,1.4) {$+$};
  \node at (4.3,1.4) {$-$};
  \node at (6.5,1.4) {total};
  \node[anchor=east] at (1.75,0.7) {sick $S$};
  \node at (2.6,0.7) {99};
  \node at (4.3,0.7) {1};
  \node at (6.5,0.7) {100};
  \node[anchor=east] at (1.75,0) {healthy $S^{c}$};
  \node at (2.6,0) {999};
  \node at (4.3,0) {98{,}901};
  \node at (6.5,0) {99{,}900};
  \node[anchor=east] at (1.75,-0.7) {total};
  \node at (2.6,-0.7) {1.098};
  \node at (4.3,-0.7) {98{,}902};
  \node at (6.5,-0.7) {100{,}000};
  \draw[->, accent, thick] (2.6,-1.65)--(2.6,-1.15);
  \node[anchor=north, align=center, text width=10.4cm, font=\footnotesize] at (4.6,-1.95)
    {once $+$ is known, only this column is counted:
     $99/1.098=11/122=0.0902$};
\end{tikzpicture}
\caption{Cross table of a screening test on $100{,}000$ people. Rows are the true state,
columns are the test result. Once a positive result is known, the sample space narrows
from the whole table to the $+$ column, which holds $1.098$ people, and $99$ of them
are really sick.}
\label{fig:silang}
\end{figure}    

\begin{proposisi}[Conditional probability is a probability]
\index{conditional probability}
For fixed $B$ with $P(B)>0$, the function $A\mapsto P(A\mid B)$ satisfies the three axioms
of probability on $\mathcal{F}$.
\end{proposisi}

\begin{proof}
Nonnegativity comes from $P(A\cap B)\ge 0$. Normality comes from $P(\Omega\mid B)=1$,
because $P(B)/P(B)=1$. For additivity, if $A_{1},A_{2},\dots$ are pairwise disjoint then
$A_{i}\cap B$ are also pairwise disjoint, so that
\[
P\Bigl(\bigcup_{i}A_{i}\mid B\Bigr)
=\frac{P\bigl(\bigcup_{i}(A_{i}\cap B)\bigr)}{P(B)}
=\frac{\sum_{i}P(A_{i}\cap B)}{P(B)}
=\sum_{i}P(A_{i}\mid B).\qedhere
\]
\end{proof}

Hence all consequences of the axioms hold for conditional probability: complement,
monotonicity, and inclusion--exclusion remain valid when every $P$ is replaced by
$P(\cdot\mid B)$. What does not hold is swapping the direction of the condition.

From the definition the product rule is born at once
\[
P(A\cap B)=P(A\mid B)\,P(B)=P(B\mid A)\,P(A),
\]
which extends to
$P(A_{1}\cap\cdots\cap A_{n})=P(A_{1})P(A_{2}\mid A_{1})\cdots
P(A_{n}\mid A_{1}\cap\cdots\cap A_{n-1})$. This rule reads a staged experiment as a
product of probabilities at each stage.

A single deck of $52$ cards gives an example of a direction that is easy to swap. Let $A$
be the event that the first card is a heart and $B$ the event that the second card is red.
Then
\[
P(A\cap B)=\frac{13\cdot 25}{52\cdot 51}=\frac{25}{204},
\]
so that
\[
P(A\mid B)=\frac{25/204}{1/2}=\frac{25}{102}=0.2451,
\qquad
P(B\mid A)=\frac{25/204}{1/4}=\frac{25}{51}=0.4902 .
\]
Two numbers that differ by a factor of two, from the same data. The mistake is not in the
arithmetic but in the placement of the events to the left and right of the vertical bar.

\begin{bacakeras}
Two voices. The first does not know yet, the second already knows. Read them in turn.

First voice: ``My world is as wide as a playing field. Every door is possible. I guess with
my eyes closed.''

Second voice: ``My world is one room. One door is already open, and I count only what is
inside it.''

First voice: ``If I know something, does my probability change?''

Second voice: ``It is not your probability that changes. It is the world you count that
changes. The world shrinks after we learn something: weight that was spread out now
gathers in the places that are still possible. The denominator narrows, and the number that
comes out sounds new even though the arithmetic has not changed since the beginning.''

First voice: ``So it is the denominator that moves?''

Second voice: ``The denominator moves, and what is left to be recounted is the remainder
that may still take part.''

Repeat together, three times, more and more slowly: ``The world shrinks, the probability
does not betray us. Only what may still be counted changes.''
\end{bacakeras}

The second tool splits the population into several groups that do not overlap. If
$B_{1},\dots,B_{k}$ are pairwise disjoint and their union is $\Omega$, then
\[
P(A)=\sum_{i=1}^{k}P(A\mid B_{i})\,P(B_{i}).
\]

\begin{proposisi}[Law of total probability]
\index{total probability}
If $B_{1},\dots,B_{k}$ are pairwise disjoint, their union is $\Omega$, and every
$P(B_{i})>0$, then $P(A)=\sum_{i}P(A\mid B_{i})P(B_{i})$ for every event $A$.
\end{proposisi}

\begin{proof}
The events $A\cap B_{1},\dots,A\cap B_{k}$ are pairwise disjoint and their union is $A$.
So $P(A)=\sum_{i}P(A\cap B_{i})$, and every term is rewritten with the product rule as
$P(A\mid B_{i})P(B_{i})$.
\end{proof}

This law looks like a formality until it is used on the clinical example: it is what turns
two numbers from an advertisement into one useful number, and that is the content of the
next section.

Before that, one example showing that ordinary sentences and mathematical events do not
always correspond. A family has two children, each child a boy or a girl with equal
probability, and the sexes of the two children are independent. If we know that
\emph{at least one} child is a boy, the probability that both are boys is $1/3=0.3333$: of
the four equally likely arrangements three remain, and only one fits. If what we know is
that \emph{the older child} is a boy, the probability is $1/2=0.5$. Two similar sentences
give different answers, because the first excludes one arrangement and the second two.

\section{Bayes' rule and the interpretation of test results}
\tingkat{intermediate}

Bayes' rule reverses the direction of the condition: from the probability of a result given
sickness it computes the probability of sickness given the result. In everyday language: of
all the ways a result can appear, what fraction comes through the path we are asking about.

\index{Bayes' rule}
\index{prevalence}

\begin{teorema}[Bayes' rule]
If $B_{1},\dots,B_{k}$ are pairwise disjoint with union $\Omega$ and $P(B_{i})>0$, then
for every event $A$ with $P(A)>0$,
\[
P(B_{j}\mid A)=\frac{P(A\mid B_{j})\,P(B_{j})}
{\sum_{i=1}^{k}P(A\mid B_{i})\,P(B_{i})}.
\]
\end{teorema}

\begin{proof}
By definition, $P(B_{j}\mid A)=P(A\cap B_{j})/P(A)$. The numerator is
$P(A\mid B_{j})P(B_{j})$ by the product rule, and the denominator is $P(A)$ by the law of
total probability. So the formula contains no new assumption; it is the definition of
conditional probability written twice, in two different directions.
\end{proof}

Hence there is no ``Bayesian method'' standing on top of the axioms; what there is, is
bookkeeping. The numerator is the path we are asking about, the denominator is the sum of
all paths that produce the same result.

\begin{resep}
\begin{enumerate}
  \item First name two things: the hypothesis $H$ we are asking about and the evidence $E$
        that has already been seen. Write the question as $P(H\mid E)$, not $P(E\mid H)$.
  \item Write the numerator as a single path: $P(E\mid H)\,P(H)$. That is the weight of the
        world in which $H$ is true.
  \item Write the denominator as the sum of all paths that produce $E$, including the path
        through $H^{c}$: $P(E\mid H)P(H)+P(E\mid H^{c})P(H^{c})$.
  \item Divide, then check for sense: the result lies between $0$ and $1$, and is smaller
        than $P(E\mid H)$ when $H$ is rare.
\end{enumerate}
A warning about order: numerator and denominator both use $P(E\mid\cdot)$. Writing
$P(\cdot\mid E)$ in the denominator is the mistake that happens most often, and the size of
the number does not tell us that the direction has been swapped.
\end{resep}

The tree in Figure~\ref{fig:pohon} writes out every road to a positive result. Each leaf is
the product of the probabilities on its branch, and the two leaves marked $+$ sum to
$P(+)=0.01098$. The Bayes numerator is just the first leaf, while the denominator is the
sum of those two leaves.

    \begin{figure}[htbp]
\centering
\begin{tikzpicture}[font=\small, every node/.style={inner sep=2pt}]
  \node (r) at (0,0) {start};
  \node (s) at (2.2,1.7) {$S$};
  \node (sc) at (2.2,-1.7) {$S^{c}$};
  \node (sp) at (4.6,2.5) {$+$};
  \node (sm) at (4.6,0.9) {$-$};
  \node (scp) at (4.6,-0.9) {$+$};
  \node (scm) at (4.6,-2.5) {$-$};
  \draw[->] (r)--(s) node[midway, above left=-2pt] {$0.001$};
  \draw[->] (r)--(sc) node[midway, below left=-2pt] {$0.999$};
  \draw[->, accent, thick] (s)--(sp) node[midway, above left=-2pt] {$0.99$};
  \draw[->, layergrey] (s)--(sm) node[midway, below left=-2pt] {$0.01$};
  \draw[->, accent, thick] (sc)--(scp) node[midway, above left=-2pt] {$0.01$};
  \draw[->, layergrey] (sc)--(scm) node[midway, below left=-2pt] {$0.99$};
  \node[anchor=west, font=\footnotesize] at (5.0,2.5) {$0.001\cdot 0.99=0.00099$};
  \node[anchor=west, font=\footnotesize] at (5.0,0.9) {$0.001\cdot 0.01=0.00001$};
  \node[anchor=west, font=\footnotesize] at (5.0,-0.9) {$0.999\cdot 0.01=0.00999$};
  \node[anchor=west, font=\footnotesize] at (5.0,-2.5) {$0.999\cdot 0.99=0.98901$};
  \node[anchor=north, align=center, text width=10.6cm, font=\footnotesize] at (4.4,-3.0)
    {$P(+)=0.00099+0.00999=0.01098$, \quad $P(S\mid+)=0.00099/0.01098=0.0902$};
\end{tikzpicture}
\caption{Two-stage tree of a screening test. The first branch separates sick from healthy,
the second separates positive results from negative ones; the product at each leaf is the
probability of that path, and the two positive leaves sum to $P(+)=0.01098$.}
\label{fig:pohon}
\end{figure}    

Back to the clinic. Write $S$ for the event of being sick and $+$ for a positive result.
From the advertisement and the prevalence of the disease we have $P(S)=0.001$,
$P(+\mid S)=0.99$, and $P(+\mid S^{c})=0.01$. The law of total probability gives
\[
P(+)=0.99\cdot 0.001+0.01\cdot 0.999=0.00099+0.00999=0.01098,
\]
so that
\[
P(S\mid +)=\frac{0.00099}{0.01098}=\frac{11}{122}=0.0902 .
\]
The false positive term, $0.00999$, is ten times larger than the true detection term,
$0.00099$; their ratio is $999/99=10.09$. Not because the test is bad: healthy people are
$999$ times more numerous, $99{,}900$ against $100$.

A change of prevalence changes the answer far more strongly than a change in the quality of
the test. With sensitivity and specificity fixed at $0.99$, the following table shows
$P(S\mid +)$ for several prevalences.

\begin{center}
\small
\begin{tabular}{rrrr}
\hline
$P(S)$ & $P(S\mid +)$ & $1/P(S\mid +)$ & reading \\
\hline
$0.0001$ & $0.0098$ & $102.0$ & one in a hundred and two \\
$0.001$ & $0.0902$ & $11.09$ & one in eleven \\
$0.010$ & $0.5000$ & $2.00$ & one in two \\
$0.100$ & $0.9167$ & $1.09$ & eleven in twelve \\
\hline
\end{tabular}
\end{center}

At a prevalence of $1$ percent the same test gives $P(S\mid+)=1/2$ exactly, not $0.99$.
At a prevalence of $10$ percent, $P(S\mid+)=11/12=0.9167$. These numbers explain a health
policy that is often criticized: mass screening for a rare disease produces many false
positives, and for that reason a positive result is usually followed by a second test
rather than by treatment.

The second test uses the same rule, with the prior already updated. If the second test is
independent of the first and of the same quality, then
\[
P(S\mid +,+)=\frac{0.99\cdot (11/122)}{0.99\cdot (11/122)+0.01\cdot (111/122)}
=\frac{10.89}{10.89+1.11}=\frac{363}{400}=0.9075 .
\]
Note that $0.9075$ is neither $0.99$ nor $0.9999$. Two positive tests raise the probability
from $0.0902$ to $0.9075$, almost ten times, but they do not yet give certainty. What
changes is only the prevalence inside the group that is already positive; the quality of
the test is unchanged. With this quality of $0.99$, the posterior exceeds $0.99$ only when
the initial prevalence is already past $1/2$.

One technical point: $P(S\mid+)=0.0902$ applies to a population with prevalence $1/1000$,
not to a patient who arrives with symptoms. Symptoms change the prevalence first, and that
is why physicians ask about history before interpreting a laboratory result.

The last example of this section releases Bayes from the clinical test. Three bags hold ten
marbles each: bag $A$ one red, $B$ five red, $C$ nine red. We choose a bag at random, draw
one marble, and the one that comes out is red. Before the draw the three bags each have
probability $1/3$; afterwards they have $0.0667$, $0.3333$, and $0.6$, as in
Figure~\ref{fig:bayes-batang}. No weight is lost; weight merely moves to the bag that more
easily produces that evidence.

    \begin{figure}[htbp]
\centering
\begin{tikzpicture}[font=\small, x=1cm, y=1cm]
  \draw[->, layergrey] (0,0)--(7.4,0);
  \draw[->, layergrey] (0,0)--(0,4.3);
  \foreach \y/\lab in {0/0, 1.2/0.2, 2.4/0.4, 3.6/0.6} {
    \draw[layergrey] (0,\y)--(-0.08,\y);
    \node[anchor=east, font=\footnotesize] at (-0.12,\y) {\lab};
  }
  \fill[layerindigo] (0.65,0) rectangle (1.20,2.0);
  \fill[layerteal] (1.40,0) rectangle (1.95,2.0);
  \fill[accent] (2.15,0) rectangle (2.70,2.0);
  \fill[layerindigo] (4.35,0) rectangle (4.90,0.4);
  \fill[layerteal] (5.10,0) rectangle (5.65,2.0);
  \fill[accent] (5.85,0) rectangle (6.40,3.6);
  \draw[dashed, layergrey] (0,2.0)--(6.9,2.0);
  \node[font=\footnotesize, anchor=south] at (0.92,2.02) {$1/3$};
  \node[font=\footnotesize, anchor=south] at (1.67,2.02) {$1/3$};
  \node[font=\footnotesize, anchor=south] at (2.42,2.02) {$1/3$};
  \node[font=\footnotesize, anchor=south] at (4.62,0.44) {$0.0667$};
  \node[font=\footnotesize, anchor=south] at (5.37,2.02) {$0.3333$};
  \node[font=\footnotesize, anchor=south] at (6.12,3.62) {$0.6$};
  \node[font=\footnotesize, anchor=north] at (0.92,-0.08) {$A$};
  \node[font=\footnotesize, anchor=north] at (1.67,-0.08) {$B$};
  \node[font=\footnotesize, anchor=north] at (2.42,-0.08) {$C$};
  \node[font=\footnotesize, anchor=north] at (4.62,-0.08) {$A$};
  \node[font=\footnotesize, anchor=north] at (5.37,-0.08) {$B$};
  \node[font=\footnotesize, anchor=north] at (6.12,-0.08) {$C$};
  \node[font=\footnotesize] at (1.67,-0.85) {before};
  \node[font=\footnotesize] at (5.37,-0.85) {after};
\end{tikzpicture}
\caption{Bayes update on three bags. Before any evidence the three bags share the weight
$1/3$; after one red marble is seen, the weights move to $0.0667$, $0.3333$, and $0.6$. The
bag that most easily produces that evidence rises the highest.}
\label{fig:bayes-batang}
\end{figure}    

\section{Independence, and why disjoint is not independent}
\tingkat{advanced}

Two events are called disjoint if they cannot occur together. Two events are called
independent if knowing one does not change the probability of the other. The two sentences
sound similar, and swapping them is the most expensive mistake in this chapter. In everyday
language: disjoint means never together; independent means no news from one to the other.

\kamus{independent}{independen}{Events that send each other no news: knowing one does not
change the probability of the other}

\begin{untukanak}
There are two things that are often swapped, and the difference can be seen without numbers.
At one school, two students are never in class at the same time: one studies in the
morning, one studies in the afternoon. If we know that Andi has come in, we know that Budi
has not. They are never together, and precisely for that reason news about one gives news
about the other.

Now two other students. Both attend the same class, but they sit far apart and work on
different exams. Their grades send each other no news: knowing Sari's grade does not help
us guess Tono's.

The first group is called never together. The second group is called no news to each other.
Both sound like ``no relation'', even though the first is in fact the strongest relation:
news about one determines the other. In this chapter we shall learn why two things that are
never together cannot possibly send each other news, and why swapping the two makes the
computation miss by more than a little.
\end{untukanak}

\index{independence}
\index{disjoint events}

\begin{definisi}[Independence of two events]
Events $A$ and $B$ are \emph{independent} if
\[
P(A\cap B)=P(A)\,P(B).
\]
\end{definisi}

This definition is symmetric, and when $P(B)>0$ it is equivalent to $P(A\mid B)=P(A)$. Note
that the definition does not mention cause and effect: two events can be independent
without one causing the other, and can be dependent without one causing the other.

\begin{proposisi}[Disjoint is not independent]
\index{disjoint events}
If $A$ and $B$ are disjoint and $P(A)>0$ and $P(B)>0$, then $A$ and $B$ are not
independent. More precisely $P(A\cap B)=0<P(A)P(B)$, and $P(A\mid B)=0<P(A)$.
\end{proposisi}

\begin{proof}
Disjoint means $A\cap B=\varnothing$, so $P(A\cap B)=P(\varnothing)=0$. Since
$P(A)P(B)>0$, the equality in the definition of independence fails. For the conditional
statement, $P(A\mid B)=0/P(B)=0$, which is smaller than $P(A)>0$.
\end{proof}

The mistake moves in two directions, and both directions can be computed.

Take one die and two pairs of events with the same marginals but a different relation. First
pair: $A=\{1,2\}$ and $B=\{3,4\}$, disjoint, with $P(A)=P(B)=1/3$. Second pair:
$C=\{1,2,3\}$ and $D=\{1,4\}$, with $P(C)=1/2$ and $P(D)=1/3$. For the first pair,
\[
P(A)P(B)=\frac{1}{9}=0.1111,\qquad P(A\cap B)=0,\qquad P(A\mid B)=0,
\]
so the gap in the definition of independence is $0.1111$ and the gap in conditional
probability is $0.3333$, because $P(A)=1/3$. For the second pair,
\[
P(C)P(D)=\frac{1}{6}=0.1667=P(C\cap D),\qquad P(C\mid D)=\frac{1}{2}=P(C),
\]
so both gaps are zero.

What matters is not the number $1/6$ but the fact that two pairs can have the same marginals
with a completely opposite relation. The pair $E=\{1,2,3\}$ and $F=\{4,6\}$ also has
marginals $1/2$ and $1/3$ but is disjoint. Marginals do not determine the relation; only the
intersection determines it, so every claim of independence is an additional claim about the
world.

For three events, pairwise independence is not enough. The events $A,B,C$ are called
mutually independent if $P(A\cap B)=P(A)P(B)$, $P(A\cap C)=P(A)P(C)$,
$P(B\cap C)=P(B)P(C)$, and $P(A\cap B\cap C)=P(A)P(B)P(C)$.

\begin{contoh}[Pairwise but not mutually independent]
\index{pairwise independence}
Two coins are tossed. Let $A$ be heads on the first coin, $B$ heads on the second coin, and
$C$ both coins the same; each has probability $1/2$. Every
pair intersects in one outcome, so $P(A\cap B)=P(A\cap C)=P(B\cap C)=1/4$, which equals
$1/2\cdot 1/2$. But $A\cap B\cap C$ occurs only when both are heads,
so its probability is $1/4$, whereas $P(A)P(B)P(C)=1/8$. So the three are pairwise
independent and not mutually independent; the gap is $0.125$.
\end{contoh}

A final warning: independence does not survive conditioning. If $A$ and $B$ are independent
with probability $1/2$, then given $A\cup B$, $P(A\mid A\cup B)=2/3$
while $P(A\cap B\mid A\cup B)=1/3\neq 4/9$. Combining two independent events and then
conditioning on their union actually creates dependence, so the sentence ``$A$ and $B$ are
independent'' is always relative to the probability space in use.

\begin{tangan}
Materials: twelve cards, six pairs of equal rank (ace through six), each pair one red and
one black.

Step one, make the pairs. Lay the twelve cards face up in two rows; put each pair of equal
rank side by side, red on the left, black on the right. Check with your eyes: rank sends no
news about color, since every rank has a red and a black. Those two properties send each
other no news.

Step two, filter. Shuffle, cover them all, then turn over one card: say the four of hearts.
Give one hint, ``this card's mate is black'', then set aside every red card that is still
face down. Of the eleven face-down cards, only six are black.

What you must see: the denominator changes from eleven to six, while the cards do not change
at all. The probability of guessing its mate rises from one in eleven to one in six. Two
properties that sent each other no news now send each other news, and that is what one hint
does.
\end{tangan}

\section{Two paradoxes that demand a computation}
\tingkat{advanced}

Paradoxes in probability are almost always a collision between intuition and definition, and
the collision is settled only by computing. The two examples below show different things:
the answer depends on the rules of the game, and on the unit that is counted.

\index{Monty Hall}

Monty Hall. Three closed doors; one holds a car, two hold goats. We choose one door. The
host, who knows what is behind all the doors, opens one other door with a goat, then offers
the chance to switch to the door that is still closed. Does switching raise the
probability?

Compute with the law of total probability over the three possible locations of the car,
each with probability $1/3$, and suppose that we first choose door $1$. If the car is
behind door $1$, switching loses. If the car is behind door $2$, the host must open door
$3$, so switching wins. If the car is behind door $3$, he opens door $2$, and switching
wins. So switching wins in two of the three cases:
\[
P(\text{win by switching})=\frac{2}{3},\qquad
P(\text{win by staying})=\frac{1}{3}.
\]
The gap is $1/3$, not zero.

The argument ``two doors left, so $1/2$'' is wrong because it ignores that the host does not
choose a door freely. He is forbidden to open the door with the car and forbidden to open
the door we chose, so his action carries information about the door he does not open.

That information disappears when the rule changes. Suppose the host opens one of the other
two doors at random, without looking at what is behind it, and the game is cancelled if the
car is revealed. Given that the game runs as far as the offer,
\[
P(\text{car behind door }1\mid \text{goat revealed})
=\frac{1/3}{1/3+1/6+1/6}=\frac{1}{2},
\]
because if the car is behind door $2$ or $3$, he has only probability $1/2$ of showing a
goat. Under that rule switching does not help: ``switching has probability $2/3$'' is not a
property of the doors but a property of the protocol.

\index{birthday paradox}

Birthdays. The previous section already computed $0.5073$ for $23$ people and $0.9704$ for
$50$ people. The paradox is not in the numbers but in the unit: we compare it with the
probability that one particular person matches us, whereas what is computed is $253$
pairs at once. The same mistake appears in hash tables and in DNA match tests.

The prosecutor's fallacy is the same directional mistake in a criminal case. Suppose the
probability that a feature appears by coincidence is one in a million, and the feature is
found in a suspect taken from a database of a million people with no other evidence. Then
$P(\text{match}\mid\text{innocent})=10^{-6}$, but
\[
P(\text{guilty}\mid\text{match})
=\frac{10^{-6}\cdot 1}{10^{-6}\cdot 1+(1-10^{-6})\cdot 10^{-6}}=0.5000003,
\]
that is practically one half, not near certainty. The two numbers $10^{-6}$ are different
quantities: the first is the probability of the feature by coincidence, the second is the
prevalence of the suspect inside the database. Swapping them is another face of the same
mistake, this time in a courtroom.

\begin{tantangan}
Three cards on a table. The first card is red on both sides, the second is black on both
sides, the third is red on one side and black on the other. We shuffle the three face down,
take one, and lay it on the table. The visible side is red. What is the probability that
the side underneath is also red?

Almost everyone answers one half: they say only two cards have a red side, so there are two
possibilities. The correct answer is not one half.

Hint 1: do not count cards, count sides. How many red sides are there in total, and how many
of them sit next to a red side?

Hint 2: the visible red side can come from three different red sides, and two of them are on
the red-red card.

If you are stuck, go back to the section on conditional probability and rewrite the sample
space as six sides, not three cards.
\end{tantangan}

\section{Worked computations}
\tingkat{intermediate}

The three computations below are carried out with exact fractions, so that no rounding hides
a step. All of them have been checked with Python, and the numbers that appear are the same
numbers used in the previous sections.

First, two-step Bayes. With $P(S)=1/1000$, $P(+\mid S)=99/100$, and
$P(+\mid S^{c})=1/100$, the law of total probability gives
\[
P(+)=\frac{99}{100}\cdot\frac{1}{1000}+\frac{1}{100}\cdot\frac{999}{1000}
=\frac{99+999}{100{,}000}=\frac{1098}{100{,}000}.
\]
The numerator is two separate groups: $99$ sick people who are detected and $999$ healthy
people who are wrongly accused. Hence
\[
P(S\mid +)=\frac{99/100{,}000}{1098/100{,}000}=\frac{99}{1098}=\frac{11}{122}=0.090164 .
\]
A second test of the same quality uses $11/122$ as the new prevalence and gives $0.9075$,
just as in the computation of the previous section. The denominator is $1200$, not a number
in the millions: the factor $1/122$ appears in both terms and cancels, a sign that what
changed is only the prevalence, not the quality of the test.

Second, disjoint versus independent, with matched marginals.

\begin{center}
\small
\begin{tabular}{lrrrr}
\hline
pair & $P(\text{first})$ & $P(\text{second})$ & $P(\text{intersection})$ & $P(\text{first}\mid\text{second})$ \\
\hline
$C=\{1,2,3\}$, $D=\{1,4\}$ & $0.5$ & $0.3333$ & $0.1667$ & $0.5$ \\
$E=\{1,2,3\}$, $F=\{4,6\}$ & $0.5$ & $0.3333$ & $0$ & $0$ \\
\hline
\end{tabular}
\end{center}

In the first row $P(C)P(D)=1/6=0.1667$ equals the intersection, so the gap is $0$. In the
second row $P(E)P(F)=1/6$ while the intersection is $0$, so the gap is $0.1667$; the
conditional probability is $0$ while $P(E)=0.5$. Anyone who takes ``cannot happen together''
to mean ``does not influence each other'' will report $0.1667$ where $0$ belongs.

Third, Monty Hall with the protocol stated. If the host always opens a goat door, the
probability of winning by switching is $2/3=0.6667$ and by staying $1/3=0.3333$. If he opens
at random and the game is cancelled when the car is revealed, the game runs with probability
$2/3$, and given that, both strategies are worth $1/2=0.5$. Two protocols, one set of doors,
two answers.

\kertas{For a test with prevalence $1/2000$, sensitivity $0.95$, and specificity
$0.90$, compute $P(S\mid+)$ with the $100{,}000$ table as above, then compute it again with
specificity $0.99$. Write both answers as fractions and as decimals with four digits, and
explain which part changes the most: the numerator or the denominator. After that compute
$P(S\mid +,+)$ for a second test of the same quality, and check whether the result is still
below $0.95$.}

\section{Computational lab: Monte Carlo as a verification tool}
\tingkat{advanced}

All the numbers in this chapter are derived, not measured. Simulation does not replace a
derivation, but it catches the mistakes that happen most often: a swapped direction of the
condition, a wrong denominator, and an independence assumption used without grounds.
Generate the same experiment many times, compute the ratio of successes, then compare it
with the exact value. By the law of large numbers (Chapter 19), that ratio tends to the true
probability; what remains is only how fast, measured by the standard error
$\sqrt{p(1-p)/m}$ with $m$ the number of trials that produce the event.

\index{Monte Carlo}
\index{standard error}

A fixed seed is used so that the output can be checked again word for word. A fixed seed
does not make the result more correct; it only makes it repeatable.

\begin{lab}
\begin{verbatim}
import numpy as np

rng = np.random.default_rng(20261024)

print("== 1. diagnostic test: prevalence 0.1%, sens 99%, spec 99% ==")
n = 1_000_000
sick = rng.random(n) < 0.001
positive = np.where(sick, rng.random(n) < 0.99, rng.random(n) < 0.01)
m = int(positive.sum())
ppv = sick[positive].mean()
err = np.sqrt(ppv * (1 - ppv) / m)
exact = 0.99 * 0.001 / (0.99 * 0.001 + 0.01 * 0.999)
print("sick=%d  positive=%d  true_positive=%d"
      % (sick.sum(), m, int((sick & positive).sum())))
print("PPV empirical=%.6f  exact=%.6f  std error=%.6f  z=%.3f"
      % (ppv, exact, err, (ppv - exact) / err))

print("== 2. Monty Hall, 10^5 games ==")
N = 100_000
car = rng.integers(0, 3, N)
pick = rng.integers(0, 3, N)
opened = np.where(pick == car,
                  (pick + 1 + rng.integers(0, 2, N)) % 3,
                  3 - pick - car)
switch = 3 - pick - opened
win_stay = float((pick == car).mean())
win_switch = float((switch == car).mean())
se_switch = np.sqrt(win_switch * (1 - win_switch) / N)
print("win without switch=%.4f  win switch=%.4f  std error=%.5f"
      % (win_stay, win_switch, se_switch))
print("exact 1/3=%.6f  2/3=%.6f" % (1/3, 2/3))

print("== 3. host opens a door at random ==")
open_random = (pick + 1 + rng.integers(0, 2, N)) % 3
valid = open_random != car
switch_random = 3 - pick - open_random
print("valid games=%d of %d (%.4f)" % (valid.sum(), N, valid.mean()))
print("win switch (valid)=%.4f  without switch (valid)=%.4f"
      % (float((switch_random[valid] == car[valid]).mean()),
         float((pick[valid] == car[valid]).mean())))

print("== 4. birthday probability, 20,000 rooms ==")
for k in (23, 50):
    days = rng.integers(0, 365, size=(20000, k))
    days.sort(axis=1)
    dup = (np.diff(days, axis=1) == 0).any(axis=1)
    q = 1.0
    for i in range(k):
        q *= (365 - i) / 365
    f = float(dup.mean())
    se = np.sqrt(f * (1 - f) / 20000)
    print("k=%2d  empirical=%.4f  exact=%.4f  std error=%.5f  z=%+.2f"
          % (k, f, 1 - q, se, (f - (1 - q)) / se))

print("== 5. disjoint versus independent, 10^6 die throws ==")
N = 1_000_000
die = rng.integers(1, 7, N)
A = np.isin(die, [1, 2])
B = np.isin(die, [3, 4])
C = np.isin(die, [1, 2, 3])
D = np.isin(die, [1, 4])
print("disjoint:  P(A)=%.4f P(B)=%.4f P(AnB)=%.4f P(A|B)=%.4f gap=%.4f"
      % (A.mean(), B.mean(), (A & B).mean(), (A & B).sum() / B.sum(),
         A.mean() * B.mean() - (A & B).sum() / B.sum()))
print("independent:  P(C)=%.4f P(D)=%.4f P(CnD)=%.4f P(C|D)=%.4f gap=%.4f"
      % (C.mean(), D.mean(), (C & D).mean(), (C & D).sum() / D.sum(),
         abs(C.mean() - (C & D).sum() / D.sum())))
print("exact disjoint: 1/3, 1/3, 0, 0, 1/9=%.6f" % (1 / 9))
print("exact independent: 1/2, 1/3, 1/6, 1/2, 0")
\end{verbatim}    
\end{lab}

The actual output, exactly as it came out:

\begin{verbatim}
== 1. diagnostic test: prevalence 0.1%, sens 99%, spec 99% ==
sick=963  positive=11094  true_positive=952
PPV empirical=0.085812  exact=0.090164  std error=0.002659  z=-1.637
== 2. Monty Hall, 10^5 games ==
win without switch=0.3325  win switch=0.6675  std error=0.00149
exact 1/3=0.333333  2/3=0.666667
== 3. host opens a door at random ==
valid games=66472 of 100000 (0.6647)
win switch (valid)=0.4998  without switch (valid)=0.5002
== 4. birthday probability, 20,000 rooms ==
k=23  empirical=0.5019  exact=0.5073  std error=0.00354  z=-1.54
k=50  empirical=0.9701  exact=0.9704  std error=0.00120  z=-0.19
== 5. disjoint versus independent, 10^6 die throws ==
disjoint:  P(A)=0.3333 P(B)=0.3335 P(AnB)=0.0000 P(A|B)=0.0000 gap=0.1112
independent:  P(C)=0.5001 P(D)=0.3334 P(CnD)=0.1667 P(C|D)=0.5000 gap=0.0001
exact disjoint: 1/3, 1/3, 0, 0, 1/9=0.111111
exact independent: 1/2, 1/3, 1/6, 1/2, 0
\end{verbatim}    

Four things deserve reading from that output, and two of them actually show the limits of
simulation.

First, the first line of part 1 reports $963$ sick people out of a million, whereas the
expectation is $1{,}000$ with standard deviation $\sqrt{10^{6}(0.001)(0.999)}=31.6$. So
$963$ lies $1.17$ standard deviations below the expectation: a normal swing, not a
programming oversight. The expected number of positive results becomes
$963\cdot 0.99+999{,}037\cdot 0.01=10{,}943.7$,
and the observed $11{,}094$ is $1.4$ standard deviations away. The empirical conditional
probability $0.0858$ differs from the exact $0.0902$ by $0.0044$, that is $1.64$ standard
errors. That gap is far smaller than the distance from $0.09$ to $1$, so the interpretation
does not change.

Second, Monty Hall gives $0.6675$ and $0.3325$ with standard error $0.0015$, while the
expected gap between the strategies is $1/3$. Agreement to three digits is not luck; the
most useful form of verification is indeed not confirming the third digit but confirming
that the gap we claim exists and has the right sign.

Third, parts 3 and 5 show that the protocol and the relations between events enter the
numbers. A host who opens doors at random lets only $66.472$ of
$100{,}000$ games run as far as the offer, and among those switching wins
$0.4998$ while staying wins $0.5002$; that gap of $0.0004$ is a quarter of a standard error.
The disjoint pair gives $P(A)P(B)-P(A\cap B)=0.1112$, matching $1/9$ to
three digits, while the independent pair gives $0.0001$, that is zero within random swing.
The number $0.1112$ is the size of the mistake made by anyone who thinks two disjoint events
must be independent.

What the program does not prove: not one of the five parts proves a formula.
The program only checks formulas that have already been proved, with a sample far larger
than the reach of the hand; Chapter 19 explains why that is legitimate.

\begin{untukdosen}
The teaching order I suggest: start from the $100{,}000$ table, not from
the axioms. Have the students fill in four cells: sick and positive, sick and negative,
healthy and positive, healthy and negative. The numbers $99$, $1$, $999$, and $98{,}901$
appear on their own, and from there $99/1.098$ becomes the conclusion of the class, not
the result of a lecture. After that the axioms, their consequences, counting, conditional
probability, Bayes, then independence. Monty Hall is best placed last, after the students
can write the host's protocol as a conditional event.

Two misconceptions appear most often. First, ``disjoint means independent''.
The symptom shows up in the arithmetic: students add the probabilities of two events that
should be multiplied, or the other way round, and then cannot explain why the result exceeds
one. Second, ``sensitivity equals the probability of being sick given a positive result''.
The symptom is the answer $0.99$ for the clinical problem, without using the prevalence
already given in the problem.

A blackboard activity that works: two students play the host with different written
protocols, one always opening a goat door and one opening at random; the class records
twenty games and then compares them with $2/3$ and $1/2$. Students are ready
to go further when they can give one example of two independent events and one example of
disjoint events with the same marginals, and compute a conditional probability from a
two-way table without reversing the direction.
\end{untukdosen}

\section{Summary}

\begin{ringkasan}
\begin{itemize}
  \item Probability is a function on events that is nonnegative, takes the value one on the
        whole space, and is countably additive on disjoint events; complement,
        monotonicity, and inclusion--exclusion follow from that.
  \item When outcomes are equally likely, probability is a ratio of counts, with the product
        rule and $\binom{n}{k}$; the assumption of equal likelihood has to be checked.
  \item Conditional probability $P(A\mid B)=P(A\cap B)/P(B)$ recomputes inside $B$
        and satisfies the three axioms; the product rule and the law of total probability
        are its consequences.
  \item Bayes' rule is not a new assumption but a definition read in two directions.
        Prevalence determines the answer more strongly than the quality of the test: with
        sensitivity and specificity $0.99$, a prevalence of $0.1$ percent gives
        $P(S\mid+)=0.0902$.
  \item Disjoint and independent are different; when both probabilities are positive,
        disjoint events cannot be independent, and only the intersection determines the
        relation.
  \item The answer to a paradox depends on the protocol and the unit of counting: Monty Hall
        $2/3$ under a fixed protocol and $1/2$ under a random one; birthdays $0.5073$ for
        $23$ people because what is counted is $253$ pairs. A seeded simulation
        checks, it does not prove.
\end{itemize}
\end{ringkasan}

\section{Exercises}

\begin{latihan}
\begin{enumerate}
  \item Prove the inclusion--exclusion formula for three events by applying the two-event
        version to $A\cup B$ and $C$. Show at which step the term
        $P(A\cap B\cap C)$ appears, and why its sign is positive.
  \item Assess this claim: ``Because $P(A\mid B)=0$, $A$ and $B$ are disjoint.'' True,
        false, or not enough evidence? Give an example with $P(A\mid B)=0$ but
        $A\cap B\neq\varnothing$, or explain why that is impossible.
  \item A test has prevalence $2$ percent, sensitivity $0.95$, specificity $0.90$.
        Compute $P(S\mid+)$ from a table of $10.000$ people, then repeat with
        specificity $0.99$; state both as fractions.
  \item Prove that if $A$ and $B$ are independent, then $A$ and $B^{c}$ are also
        independent. Use $P(A)=P(A\cap B)+P(A\cap B^{c})$ and explain why
        the same result holds for $A^{c}$ and $B^{c}$.
  \item Compute the probability that a pair with the same birthday exists for $30$ people,
        compare it with $1-e^{-\binom{30}{2}/365}$, then explain the source of the gap.
  \item Prove the law of total probability for $k=2$ directly from the axioms, without using
        the law for general $k$: write $A=(A\cap B)\cup(A\cap B^{c})$ and check
        that the two parts are disjoint.
  \item Four doors, one car, three goats. A host who knows the contents opens two
        goat doors among the three other doors and offers a switch to the one door
        that remains. Compute the probability of winning by switching, then check it with a
        simulation of $10^{4}$ games.
  \item Assess this claim: ``Two independent tests with sensitivity and specificity $0.99$
        give certainty $0.9999$ when both are positive, as long as the prevalence is $0.1$
        percent.'' Compute the correct number and state the assumptions.
  \item Is pairwise independence for three events enough to conclude
        $P(A\cap B\cap C)=P(A)P(B)P(C)$? Refute it with an example from two coin tosses,
        and compute both values.
\end{enumerate}

Hints for the three hardest items. Item 2: in a finite space where every outcome has
positive probability, $P(A\mid B)=0$ forces $A\cap B=\varnothing$, so for a die throw
the claim is true; a counterexample has to come from a continuous space, for instance
$A=\{1/2\}$ and $B=(0,1)$ with uniform probability, where $P(A)=0$ even though $A$ contains
a point. Item 7: at first the car is behind the chosen door with probability $1/4$; if not,
 the only door the host does not open must hold the car, so the probability of switching is
$3/4$. Item 9: take $A$ heads on the first coin, $B$ heads on the second coin, and $C$ both
coins the same; all three pairs give $1/4$ while the intersection is $1/4$, not $1/8$.
\end{latihan}

\begin{untukpeneliti}
What is genuinely proved in this chapter: the three axioms together with their consequences;
Bayes' rule as a theorem derived from the definition of conditional probability; the
impossibility of two disjoint events with positive probability also being independent; and
the dependence of the Monty Hall answer on the protocol, an analysis stated sharply by
Morgan, Chaganty, Dahiya, and Doviak, \emph{The American Statistician} 1991, in the article
``Let's Make a Deal: The Player's Dilemma''. What is popular but not proved here
is the claim that the natural frequency format improves prevalence reasoning in humans;
that is an empirical claim about readers, not a theorem, and its evidence lies in the
psychology of decision making literature.
\end{untukpeneliti}

Primary sources: Blitzstein and Hwang, \emph{Introduction to Probability}, Chapter 1 (Why
probability, Naive definition of probability, How to count, Story proofs, Non-naive
definition of probability) and Chapter 2 (The importance of thinking conditionally,
Definition and intuition, Bayes' rule and the law of total probability, Conditional
probabilities are probabilities, Independence of events, Coherency of Bayes' rule);
as well as Wasserman, \emph{All of Statistics}, Chapter 1 in the section Probability for the
formulation of the axioms and Chapter 2 in the section Independent Events.

Limits of validity: the whole chapter works in finite spaces where every outcome has
positive probability, where $P(A)=0$ is equivalent to $A=\varnothing$. In uncountable
spaces, conditioning on an event of probability zero is not defined by the ratio formula
and requires a regular conditional distribution; the Borel paradox shows that the answer
can depend on how the limit is taken. Independence is not stable either: it can disappear
or appear under conditioning, and in graphical models it appears as selection bias.

An open question: how far must the prevalence be known before Bayes' rule is used
honestly, and what can be said when the prevalence is known only from a small sample. The
formal answer is Bayesian inference with a prior, which brings in the problem of choosing
the prior; the frequentist answer is the hypothesis test of Chapter 20, which does not use
a prior but answers a different question. This chapter does not choose either.

\peta{Chapter 17 in Blitzstein and Hwang, \emph{Introduction to Probability}: Chapter 1
(Probability and counting, including How to count and the birthday problem example) for the
axioms, counting, and combinatorial probability; Chapter 2 (Conditional probability, Bayes'
rule and the law of total probability, Independence of events) for conditional probability,
Bayes, and independence. In Wasserman, \emph{All of Statistics}: Chapter 1 in the section
Probability and Chapter 2 in the section Independent Events. Sets and complements are used
from Chapter 2 of this book; the series underlying countable additivity from Chapter 5; the
law of large numbers that supports the simulation is proved in Chapter 18. Random variables
in Chapter 17 move the weight from events to numbers, while the hypothesis tests in Chapter
19 use Bayes' rule and the law of total probability for a different question.}

\chapter{Random Variables and Distributions}

A messaging service receives an average of three requests per second at peak hours.
That number is printed on the dashboard, and it does not answer the pressing question: how
often does a single second hold eight requests or more, so that the queue overflows? Two systems can have exactly the same average and behave
very differently: one almost never overflows, the other several times a day.
The average alone is not enough, and not because the counting is insufficiently careful.

Chapter 17 gave us a language for events: the probability of one outcome, the probability of a union,
conditional probability. That language stops working as soon as what we measure is a number.
``The number of requests in one second'' is not one event but a family of
infinitely many events, one for each possible value. For
counts, the probability of each value can still be added up one by one. For a quantity that
is continuous, that way collapses: a waiting time can be $0.317$ seconds or
$0.3170001$ seconds, and there is no way to state ``the probability of exactly $0.317$ seconds''
without assigning weight zero to every value.

So two things are needed: a way to move probability onto the number line, and a way
to replace summation by integration when the values are continuous.
Once both exist, the average, the spread, and the shape of the tail can be computed.

\begin{untukanak}
Imagine a bakery. Every hour, the seller writes down how many loaves were
sold: sometimes 3, sometimes 11, sometimes none at all. Those records are
whole numbers, and cannot take the value one half. If we want to know the probability
of selling exactly 7 loaves, we can compute it: count how many hours sold 7 loaves,
divide by the number of hours recorded. All the probabilities can be added one by one,
because the options are few and can be listed.

Now weigh one loaf on a digital scale. The result is not a whole number:
250.3 grams, or 249.8 grams. However many decimal places we take, there is always
another possibility closer still. Hence the question ``what is the probability that the weight is
exactly 250.0000 grams'' is useless, whereas ``what is the probability that the weight lies between 249
and 251 grams'' is very useful. Probability here belongs not to a single point but to
an interval.

In this chapter we will learn why those two ways differ, and why both can still
give a sensible average.
\end{untukanak}

\section{From Outcomes to Numbers}
\tingkat{basic}

A measuring instrument does not record events; it records numbers. The bridge is a function
that takes each outcome of the experiment and assigns one real number. The name
``random variable'' is misleading: it is not a variable, and it is not random. The randomness
comes from the probability on the outcome space, whereas the function is ordinary.

\begin{definisi}[Random variable]\index{random variable}
Let $\Omega$ be the outcome space and $\mathcal{F}$ a collection of events. The mapping
$X:\Omega\to\mathbb{R}$ is called a \emph{random variable} if for every $x\in\mathbb{R}$
the set
\[
\{X\le x\}=\{\omega\in\Omega : X(\omega)\le x\}
\]
is an event, that is, a member of $\mathcal{F}$.
\end{definisi}

That last condition guarantees that $P(X\le x)$ has a meaning. For almost every function that
appears in practice it holds automatically, and it fails only for set indicators that
are not measurable.

As an example, toss three fair coins and let $X$ be the number of heads that appear.
Its values are $0,1,2,3$ with probabilities $1/8,3/8,3/8,1/8$. Note that $X$ maps
eight outcomes onto four numbers, so several outcomes share the same value. That is
why the probability of a value is the probability of a union of events.

We need one tool that answers ``how far have the values moved already'' for all
thresholds at once, because the threshold is chosen by the reader, not by the author of the book.

\begin{definisi}[Cumulative distribution function]\index{cumulative distribution function}
The cumulative distribution function of a random variable $X$ is
$F_X(x)=P(X\le x)$ for $x\in\mathbb{R}$.
\end{definisi}

\begin{proposisi}[Properties of the distribution function]
For every random variable $X$ the following hold:
\begin{enumerate}
  \item $F_X$ is nondecreasing: if $a\le b$ then $F_X(a)\le F_X(b)$;
  \item $\lim_{x\to-\infty}F_X(x)=0$ and $\lim_{x\to+\infty}F_X(x)=1$;
  \item $F_X$ is right-continuous, that is, $F_X(x)=F_X(x^+)$ for every $x$;
  \item $P(a<X\le b)=F_X(b)-F_X(a)$ for every $a<b$.
\end{enumerate}
\end{proposisi}

\begin{proof}
For (1), the event $\{X\le a\}$ is contained in $\{X\le b\}$ when $a\le b$, and
probability is monotone with respect to inclusion.

For (2), let $A_n=\{X\le -n\}$. These events are decreasing, and
$\bigcap_n A_n=\varnothing$ because every outcome yields a finite value. Continuity of
probability from above, which follows from countable additivity, gives $F_X(-n)\to0$. The
opposite direction uses $B_n=\{X\le n\}$, which increase to $\Omega$, so that
$F_X(n)\to1$.

For (3), take a decreasing sequence $y_n\downarrow x$ and $C_n=\{X\le y_n\}$. Then
$\bigcap_n C_n=\{X\le x\}$, hence $F_X(y_n)\to F_X(x)$ by continuity from above.
Since the sequence is arbitrary, the right-hand limit equals the value at $x$.

For (4), the event $\{X\le b\}$ splits into $\{X\le a\}$ and $\{a<X\le b\}$,
so that $F_X(b)=F_X(a)+P(a<X\le b)$.
\end{proof}

For two tosses of a fair coin with $X$ the number of heads,
\[
F_X(x)=\begin{cases}
0, & x<0,\\
1/4, & 0\le x<1,\\
3/4, & 1\le x<2,\\
1, & x\ge 2.
\end{cases}
\]
This function is defined on the whole number line, even though $X$ takes only the values $0,1,2$.
That is why $F_X(1.4)=3/4$: the question is not ``is $X$ equal to $1.4$'',
but ``is $X$ at most $1.4$''. Every jump of the graph is the probability
mass at a single value. Figure~\ref{fig:tangga-cdf} draws that staircase: three horizontal
steps, three vertical jumps, and the height of each jump equals the probability
mass at the point of the jump.

\begin{figure}[htbp]
\centering
\begin{tikzpicture}[font=\footnotesize,>=Stealth,x=1.75cm,y=4.2cm]
  \foreach \y/\lab in {0.25/{1/4},0.75/{3/4},1/{1}}{
    \draw[layergrey!35,dashed] (-1.35,\y) -- (3.55,\y);
    \node[left,font=\scriptsize] at (0,\y) {$\lab$};
  }
  \draw[->,layergrey] (-1.35,0) -- (3.6,0) node[below] {$x$};
  \draw[->,layergrey] (0,-0.1) -- (0,1.18) node[left] {$F_X(x)$};
  \draw[very thick,layerindigo] (-1.25,0) -- (0,0);
  \draw[very thick,layerindigo] (0,0.25) -- (1,0.25);
  \draw[very thick,layerindigo] (1,0.75) -- (2,0.75);
  \draw[very thick,layerindigo] (2,1) -- (3.5,1);
  \draw[very thick,layerplum] (0,0) -- (0,0.25);
  \draw[very thick,layerplum] (1,0.25) -- (1,0.75);
  \draw[very thick,layerplum] (2,0.75) -- (2,1);
  \fill[layerindigo] (0,0.25) circle (1.8pt);
  \fill[layerindigo] (1,0.75) circle (1.8pt);
  \fill[layerindigo] (2,1) circle (1.8pt);
  \draw[layerindigo,fill=white,thick] (1,0.25) circle (1.8pt);
  \draw[layerindigo,fill=white,thick] (2,0.75) circle (1.8pt);
  \node[below,font=\scriptsize] at (0,-0.02) {$0$};
  \node[below,font=\scriptsize] at (1,-0.02) {$1$};
  \node[below,font=\scriptsize] at (2,-0.02) {$2$};
  \node[right,font=\scriptsize,text=layerplum] at (1.06,0.5) {jump $1/2$};
  \node[right,font=\scriptsize,text=layerplum] at (0.06,0.12) {jump $1/4$};
\end{tikzpicture}
\caption{The cumulative distribution function of the number of heads in two tosses of a coin. The graph is
a flat staircase between the possible values, then jumps exactly at those values. The jumps of
$1/4$ at $x=0$, $1/2$ at $x=1$, and $1/4$ at $x=2$ are the probability masses $p_X(0)$,
$p_X(1)$, and $p_X(2)$.
The filled dot at the left end of every step marks that $F_X$ is right-continuous, so that
$F_X(1)=3/4$ and not $1/4$.}
\label{fig:tangga-cdf}
\end{figure}

\kamus{random variable}{variabel acak}{a function that attaches one number to each
outcome of the experiment, like a clerk who records one number for every event}

\kamus{cumulative distribution function}{fungsi distribusi kumulatif}{a rising
staircase: its height at the point $x$ states how much probability has already
accumulated up to that threshold}

\section{Mass Functions and Density Functions}
\tingkat{basic}

There are two ways of weighting the values of a random variable. For a variable whose values can
be listed, the weight attaches to points; for one whose values fill an interval, the weight attaches
to length. The most common mistake in this topic is to use the first way to
read the outcome of the second.

\begin{definisi}[Discrete variables and probability mass function]\index{probability mass function}
$X$ is called \emph{discrete} if its set of values is countable, say
$\{x_1,x_2,\dots\}$. The \emph{probability mass function} of $X$ is $p_X(x)=P(X=x)$.
This function satisfies $p_X(x)\ge0$, $p_X(x)\le1$, and $\sum_i p_X(x_i)=1$, as well as
$F_X(x)=\sum_{x_i\le x}p_X(x_i)$.
\end{definisi}

For $X$ binomial with $n=2$ and $p=1/2$, the probability mass is $1/4,1/2,1/4$ at
$x=0,1,2$ and zero elsewhere. Probability mass is never greater than one,
because it is a probability after all.

\begin{definisi}[Continuous variables and probability density function]\index{density function}
$X$ is called \emph{continuous} if there exists a function $f_X\ge0$ such that
\[
P(a\le X\le b)=\int_a^b f_X(x)\,dx
\]
for every $a\le b$, with $\int_{-\infty}^{\infty}f_X(x)\,dx=1$. The function $f_X$
is called the \emph{probability density function}.
\end{definisi}

Density is not probability, and the difference is not a matter of terminology. Take $X$ with
density $f_X(x)=5$ on $[0,1/5]$ and zero outside it. Then $f_X(0.1)=5$, a number
that is impossible for a probability, whereas
$P(X\le0.1)=\int_0^{0.1}5\,dx=0.5$. Density may also be unbounded:
$f_X(x)=\tfrac23x^{-1/3}$ on $(0,1)$ integrates to one even though it blows up near zero.
The units differ: probability per unit of length, just as population density is people
per square kilometer. Figure~\ref{fig:batang-dan-kurva} puts
the two ways side by side: the bars on the left are added by height, the triangle on the
right by area, and both total one.

\begin{figure}[htbp]
\centering
\begin{tikzpicture}[font=\footnotesize,>=Stealth]
  \begin{scope}[x=1.05cm,y=5.4cm]
    \draw[->,layergrey] (0,0) -- (5.4,0) node[right] {$x$};
    \draw[->,layergrey] (0,0) -- (0,0.62) node[left] {$p_X(x)$};
    \foreach \x/\p in {1/0.1,2/0.2,3/0.3,4/0.4}{
      \fill[layerorange!40] (\x-0.3,0) rectangle (\x+0.3,\p);
      \draw[layerbrown,thick] (\x-0.3,0) rectangle (\x+0.3,\p);
      \node[below,font=\scriptsize] at (\x,0) {$\x$};
    }
    \node[font=\scriptsize,text=layerbrown,align=left,anchor=west] at (0.45,0.44)
      {four bars,\\ total height $=1$};
  \end{scope}
  \begin{scope}[xshift=7.7cm,x=1.05cm,y=5.4cm]
    \draw[->,layergrey] (0,0) -- (5.4,0) node[right] {$x$};
    \draw[->,layergrey] (0,0) -- (0,0.62) node[left] {$f_X(x)$};
    \fill[layerteal!25] (0.4,0) -- (2.4,0.5) -- (4.4,0) -- cycle;
    \draw[layerteal,thick] (0.4,0) -- (2.4,0.5) -- (4.4,0) -- cycle;
    \draw[layergrey,dashed] (2.4,0) -- (2.4,0.5);
    \node[above,font=\scriptsize] at (2.4,0.5) {$0.5$};
    \node[below,font=\scriptsize] at (0.4,0) {$0$};
    \node[below,font=\scriptsize] at (4.4,0) {$4$};
    \node[font=\scriptsize,text=layerteal,align=center] at (2.4,0.14)
      {area $=\tfrac12\cdot4\cdot0.5=1$};
  \end{scope}
\end{tikzpicture}
\caption{Two ways of weighting that both total one. Left: four bars
of height $0.1$, $0.2$, $0.3$, $0.4$ at the values $1,2,3,4$, added by
height. Right: a triangle-shaped density with peak $0.5$ on $[0,4]$,
added by area, and $\tfrac12\cdot4\cdot0.5=1$. The height $0.5$ in the right-hand panel
is not a probability of anything; what has probability is the area of the slice beneath it.}
\label{fig:batang-dan-kurva}
\end{figure}

\begin{bacakeras}
Read aloud in pairs, one person taking the Height voice and one the Weight voice. Start
slowly, then speed up.

Height: I am the height of the curve. I tower at one point, and everyone points at me.
Weight: I am the weight under the curve. I pile up along the interval, and everyone weighs me.
Height: My number may be five. May be ten. May be unbounded.
Weight: My number never exceeds one, because I am the whole of probability.
Height: Without me, you do not know where the most frequent value lies.
Weight: Without me, you do not know whether that peak contains anything.
Together: Height points to the place. Weight weighs the content.

Repeat together three times, more slowly each time: probability is weight, not height.
\end{bacakeras}

\begin{teorema}\index{density}
If $X$ is continuous with density $f_X$, then $F_X(x)=\int_{-\infty}^{x}f_X(t)\,dt$,
and $f_X(x)=F_X'(x)$ at every point $x$ where $F_X$ is differentiable.
\end{teorema}

\begin{proof}
The first equation is the definition of density with $a=-\infty$ and $b=x$. The second
equation is the first part of the fundamental theorem of calculus: the derivative of the accumulation
function equals $f_X$ at every point where $f_X$ is continuous.
\end{proof}

The question that most often causes confusion: if $X$ is continuous then
$P(X=x)=0$ for every $x$, because $\int_x^x f_X=0$. Does that mean that the event $X=x$
is impossible? No. The set $\{\omega:X(\omega)=x\}$ still exists and can still contain
outcomes; what is zero is its probability, not its existence.

The difficulty appears when we try to add up all the points. The interval $[0,1]$ is an uncountable
union of points of probability zero, yet $P(0\le X\le1)=1$. There is no
contradiction: the third axiom of Chapter 17 guarantees additivity only for
\emph{countable} unions, and the integral measures length, it does not count points.

There is a practical layer. NumPy's built-in generator forms values
as multiples of $2^{-53}$, so there are only $2^{53}$ possible outputs on $[0,1)$, and
the probability that one particular value appears in $10^7$ draws is about
$10^7/2^{53}=1.1\times10^{-9}$. On a computer ``probability zero'' is really
$2^{-53}\approx1.1\times10^{-16}$; the zero we write is the idealization.

\begin{contoh}
Let $X$ be uniform on $(0,1)$, with $f_X(x)=1$ for $0<x<1$. Then
\[
P(0.49<X<0.51)=\int_{0.49}^{0.51}1\,dx=0.02,
\qquad P(X=0.5)=0 .
\]
The density around $0.5$ is as large as the density around $0.01$; what distinguishes them
is only the length of the interval. That is why the right question for a continuous variable
is always an interval.
\end{contoh}

\begin{tangan}
Materials: a tape measure, scissors, card stock, and ten blank cards.

Measure the heights of ten friends to the nearest centimeter, then write one height on each
card. Make five boxes on the floor from strips of card stock five centimeters wide:
$150$--$154$, $155$--$159$, up to $170$--$174$. Put each card in the box that
matches, stacking them upward. What we must see: those stacks have a shape, and the number of
cards in all the boxes is still ten.

Cut the card stock again into boxes ten centimeters wide, arrange the cards again,
and see what changes. The stacks become almost half as high, but not a single card
is lost. The width of the box changes the height of the stack; the number of cards does not
change. That is the difference between height and weight, and that is why a density must be multiplied
by the width of the interval before it may be called a probability.
\end{tangan}

\begin{untukanak}
Imagine a football field full of people. We want to know where the people are
most numerous. We cannot answer that by asking ``how many people are exactly at
this point''. The answer is always zero, because a point has no area. What can
be answered is ``how many people are there per square meter''. That number is called density,
and it differs from place to place: dense at the entrance, sparse in the corner of the field.

High density does not mean that there are many people at one point there. High
density means that around that place, within a small region, people are packed together. If we
multiply the density by the area of the region, only then does the number of people come out. The same
holds for probability: density times the length of the interval gives probability, and density
itself is not probability.

This is why two curves of very different shape can give the same total,
namely one. What matters is not the height of the curve at one point, but the area beneath it.
In this chapter we will learn why a question about points must be replaced by a question
about intervals.
\end{untukanak}

\section{Expectation and Variance}
\tingkat{intermediate}

If the whole distribution has to be compressed into a single number, the number chosen
is its balance point: the location of the center of mass when probability is treated as weight.

\begin{definisi}[Expectation]\index{expectation}
The expectation of a random variable $X$ is
\[
\mathbb{E}[X]=\sum_i x_i\,p_X(x_i)
\qquad\text{or}\qquad
\mathbb{E}[X]=\int_{-\infty}^{\infty}x\,f_X(x)\,dx,
\]
depending on whether $X$ is discrete or continuous. The expectation exists only if
$\mathbb{E}[|X|]<\infty$; otherwise the expectation of $X$ is said not to exist.
\end{definisi}

That absolute-value condition is not excessive caution. For the Cauchy distribution with
density $f(x)=1/(\pi(1+x^2))$, $\int|x|f(x)dx$ diverges, so that $\mathbb{E}[X]$
does not exist even though the distribution is symmetric about zero; the sample mean does not
concentrate anywhere, however much data we add. Chapter 19 uses this example
as a boundary of the law of large numbers.

As short computations, for $X$ Bernoulli with $P(X=1)=p$,
$\mathbb{E}[X]=0\cdot(1-p)+1\cdot p=p$. For $X$ uniform on $(0,1)$,
$\mathbb{E}[X]=\int_0^1x\,dx=1/2$.

What is often needed is not the expectation of $X$ but the expectation of a function of $X$, for example
$\mathbb{E}[X^2]$. The direct route is to find the density of $X^2$ first, then
integrate. There is a much shorter route.

\begin{teorema}[Expectation of a function of a random variable]
For a function $r$ and a random variable $X$,
\[
\mathbb{E}[r(X)]=\sum_i r(x_i)\,p_X(x_i)
\qquad\text{or}\qquad
\mathbb{E}[r(X)]=\int_{-\infty}^{\infty}r(x)\,f_X(x)\,dx .
\]
\end{teorema}

\begin{proof}
For the discrete case, group the values according to $y=r(x)$:
\[
\sum_i r(x_i)p_X(x_i)=\sum_y y\sum_{i:\,r(x_i)=y}p_X(x_i)=\sum_y y\,P(r(X)=y),
\]
and the last expression is the definition of $\mathbb{E}[r(X)]$ because $r(X)$ is discrete with
probability mass $P(r(X)=y)$. For the continuous case, the proof is the change-of-variables
formula for integrals in Chapter 15; the same grouping applies with integrals
replacing sums.
\end{proof}

\begin{proposisi}[Linearity of expectation]
For random variables $X,Y$ and constants $a,b$,
$\mathbb{E}[aX+bY]=a\,\mathbb{E}[X]+b\,\mathbb{E}[Y]$.
\end{proposisi}

\begin{proof}
For the discrete case with joint mass $p(x,y)$,
\[
\mathbb{E}[aX+bY]=\sum_{x,y}(ax+by)p(x,y)
=a\sum_x x\sum_y p(x,y)+b\sum_y y\sum_x p(x,y)=a\,\mathbb{E}[X]+b\,\mathbb{E}[Y].
\]
The continuous case is identical with integrals replacing sums.
\end{proof}

What is not used here is independence: linearity holds for any variables,
including strongly dependent ones. Hence it is safe for sums of indicators, but it may not
be used for products: $\mathbb{E}[XY]=\mathbb{E}[X]\mathbb{E}[Y]$ requires
independence.

Expectation gives the center, but it does not give the width; the most widely used measure of width
is the average squared distance to the center.

\begin{definisi}[Variance]\index{variance}
The variance of $X$ with $\mu=\mathbb{E}[X]$ is
$\operatorname{Var}(X)=\mathbb{E}\bigl[(X-\mu)^2\bigr]$, and its standard deviation is
$\sigma_X=\sqrt{\operatorname{Var}(X)}$.
\end{definisi}

\begin{proposisi}
If the variance exists, then $\operatorname{Var}(X)=\mathbb{E}[X^2]-\mu^2$ and
$\operatorname{Var}(aX+b)=a^2\operatorname{Var}(X)$.
\end{proposisi}

\begin{proof}
Expand the square and use linearity:
\[
\mathbb{E}[(X-\mu)^2]=\mathbb{E}[X^2]-2\mu\,\mathbb{E}[X]+\mu^2
=\mathbb{E}[X^2]-\mu^2 .
\]
For the second property, $\mathbb{E}[aX+b]=a\mu+b$, so that
$\operatorname{Var}(aX+b)=\mathbb{E}[(aX+b-a\mu-b)^2]=a^2\mathbb{E}[(X-\mu)^2]$.
\end{proof}

For $X$ Bernoulli, $\mathbb{E}[X^2]=p$ because $X^2=X$ for the values $0$ and $1$,
so that $\operatorname{Var}(X)=p-p^2=p(1-p)$. For $X$ uniform on $(0,1)$,
$\mathbb{E}[X^2]=1/3$, so that $\operatorname{Var}(X)=1/3-1/4=1/12\approx0.083333$.
Variance has squared units, whereas the standard deviation has the same units as $X$. One
costly warning: $\mathbb{E}[X^2]\ne\mathbb{E}[X]^2$, and the difference is precisely
the variance; for the uniform variable above, $0.25$ against $0.333333$.

\begin{resep}
Materials: a table of values $x_i$ and probabilities $p_X(x_i)$.
\begin{enumerate}
  \item Add up all the $p_X(x_i)$. If the result is not one, the table is not yet valid.
  \item Multiply each $x_i$ by $p_X(x_i)$, then add: that is $\mathbb{E}[X]$.
  \item Square each $x_i$ first, multiply by $p_X(x_i)$, then add:
        that is $\mathbb{E}[X^2]$.
  \item Subtract the square of the result of step 2 from the result of step 3, namely
        $\operatorname{Var}(X)=\mathbb{E}[X^2]-(\mathbb{E}[X])^2$; its square root
        is the standard deviation.
\end{enumerate}
Example: for $x=0,1,2,3$ with $p=0.1;\,0.2;\,0.3;\,0.4$, step 2 gives
$2$, step 3 gives $5$, and step 4 gives $1$.

Warning about the order: square $x_i$ before multiplying by the probability, and subtract
$\mu^2$ after $\mathbb{E}[X^2]$ has been fully added up. Swapping that order produces
a number that measures no width at all.
\end{resep}

\kamus{expectation}{harapan}{the balance point of a loaded ruler: one number that
marks where the values pile up}

\kamus{variance}{variansi}{a measure of how far typical values stray from the balance
point, computed from squared distances so that plus and minus signs do not cancel each other}

\section{Four Discrete Distributions}
\tingkat{intermediate}

Four families of discrete distributions suffice for most practical work: one
binary trial, the number of binary trials, the waiting time until the first success, and
rare events in a time interval.

\paragraph{Bernoulli.} $X$ Bernoulli with parameter $p\in[0,1]$ takes the value $1$ with
probability $p$ and $0$ with probability $1-p$, so that $p_X(x)=p^x(1-p)^{1-x}$ for
$x\in\{0,1\}$. From those two points, $\mathbb{E}[X]=p$ and
$\operatorname{Var}(X)=p(1-p)$; the variance is largest at $p=1/2$, namely $1/4$, exactly
when the two possibilities are balanced.

\paragraph{Binomial.} If $n$ independent Bernoulli trials with success probability $p$
are run, the number of successes $X$ is binomial with
\[
p_X(k)=\binom{n}{k}p^k(1-p)^{n-k},\qquad k=0,1,\dots,n .
\]
Its mean and variance are derived most easily by splitting $X$ into indicators,
not from that formula.

\begin{teorema}\index{binomial distribution}
If $X$ is binomial with parameters $n$ and $p$, then
$\mathbb{E}[X]=np$ and $\operatorname{Var}(X)=np(1-p)$.
\end{teorema}

\begin{proof}
Write $X=\sum_{i=1}^{n}X_i$ with $X_i=1$ if trial $i$ succeeds and $0$ if it
fails. Each $X_i$ is Bernoulli with parameter $p$, hence $\mathbb{E}[X_i]=p$ and
$\operatorname{Var}(X_i)=p(1-p)$. Linearity gives
$\mathbb{E}[X]=\sum_i\mathbb{E}[X_i]=np$ without using independence at all.

For the variance, independence is needed for the first time. Because the trials are independent, the pairs
$X_i,X_j$ with $i\ne j$ are independent too, so that
$\operatorname{Cov}(X_i,X_j)=p^2-p^2=0$. Hence
\[
\operatorname{Var}(X)=\sum_{i=1}^{n}\operatorname{Var}(X_i)
+2\sum_{i<j}\operatorname{Cov}(X_i,X_j)=np(1-p).
\]
\end{proof}

The variance formula is more fragile than the mean formula: the former requires independence,
the latter does not. For $n=20$ and $p=0.35$, $\mathbb{E}[X]=7$ and
$\operatorname{Var}(X)=4.55=91/20$, with standard deviation $2.133073$; the mass
around the center is $p_X(7)=0.184401$, $p_X(10)=0.068614$, $p_X(12)=0.013564$.
The tail is small but not zero: $P(X\ge12)=0.019579$ and $P(X\le4)=0.118197$, so the
rule ``reject a box if it contains twelve defects or more'' rejects about two percent of
boxes.

\paragraph{Geometric.} If independent Bernoulli trials are repeated until the first success,
the number of trials needed is $X$ geometric with
$p_X(k)=p(1-p)^{k-1}$ for $k=1,2,\dots$ Its total mass is one because of the geometric series:
$\sum_{k\ge1}p(1-p)^{k-1}=p/(1-(1-p))=1$. With $q=1-p$,
\[
\sum_{k\ge1}kq^{k-1}=\frac{d}{dq}\frac{1}{1-q}=\frac{1}{(1-q)^2}=\frac{1}{p^2},
\]
so that $\mathbb{E}[X]=p\cdot p^{-2}=1/p$ and, in a similar way,
$\operatorname{Var}(X)=(1-p)/p^2$ from $\mathbb{E}[X^2]=(2-p)/p^2$. For $p=0.2$:
$\mathbb{E}[X]=5$, $\operatorname{Var}(X)=20$, $P(X=5)=0.081920$ and
$P(X>10)=0.10737418$; a simulation of $10^5$ draws gives mean $4.99176$ with
standard error $0.01414$.

\paragraph{Poisson.} For $X$ Poisson with parameter $\lambda>0$,
\[
p_X(k)=e^{-\lambda}\frac{\lambda^k}{k!},\qquad k=0,1,2,\dots
\]
Its total mass is one because of the exponential series:
$e^{-\lambda}\sum_k\lambda^k/k!=1$. Its mean and variance are equal, and both
come out of a single trick. Because $k/k!=1/(k-1)!$,
\[
\mathbb{E}[X]=\sum_{k\ge1}k\,e^{-\lambda}\frac{\lambda^k}{k!}
=\lambda e^{-\lambda}\sum_{k\ge1}\frac{\lambda^{k-1}}{(k-1)!}=\lambda .
\]
Writing $k(k-1)/k!=1/(k-2)!$ gives $\mathbb{E}[X(X-1)]=\lambda^2$, hence
$\operatorname{Var}(X)=\lambda$. One parameter governs the center and the width at once, and that
restricts the Poisson: data whose variance far exceeds their mean (called
\emph{overdispersion}) do not fit the Poisson model.

For $\lambda=3$: $p_X(0)=0.049787$, $p_X(3)=0.224042$,
$P(X\ge8)=0.011905$. That last number answers the queueing question from the opening of the chapter:
with three requests per second, about $1.2$ percent of seconds contain eight
requests or more.

The Poisson is the limit of the binomial when $n$ is large, $p$ is small, and $np\to\lambda$. For
$n=1000$ and $p=0.003$, the binomial gives $P(0)=0.049563083$ whereas the Poisson
gives $e^{-3}=0.049787068$, a relative difference of $-4.5\times10^{-3}$. The error is small but
not zero and its sign changes with $k$; the Poisson is used as an approximation, not as an
identity.

\section{Three Continuous Distributions}
\tingkat{intermediate}

\paragraph{Uniform.} $X$ is uniform on $(a,b)$ if its density is $1/(b-a)$ inside
that interval and zero outside. Its distribution function rises linearly from $0$ to $1$, and
\[
\mathbb{E}[X]=\frac{a+b}{2},\qquad
\operatorname{Var}(X)=\frac{(b-a)^2}{12}.
\]
The factor $12$ comes from $\int_a^b(x-\tfrac{a+b}{2})^2\frac{dx}{b-a}=(b-a)^2/12$; for
$(0,1)$ the result is $1/12\approx0.083333$, the same as the computation through
$\mathbb{E}[X^2]$ in the previous section.

\paragraph{Exponential.} $X$ exponential with parameter $\beta>0$ has
density $f_X(x)=\beta^{-1}e^{-x/\beta}$ for $x>0$, with
$F_X(x)=1-e^{-x/\beta}$, $\mathbb{E}[X]=\beta$, and $\operatorname{Var}(X)=\beta^2$.
The parameter $\beta$ is the mean waiting time. The property that distinguishes it from every
other continuous distribution is the absence of memory.

\begin{proposisi}[Memoryless]\index{exponential distribution}
For $s,t\ge0$, $P(X>s+t\mid X>s)=P(X>t)$.
\end{proposisi}

\begin{proof}
$P(X>u)=e^{-u/\beta}$, so that
\[
P(X>s+t\mid X>s)=\frac{P(X>s+t)}{P(X>s)}
=\frac{e^{-(s+t)/\beta}}{e^{-s/\beta}}=e^{-t/\beta}=P(X>t).
\]
\end{proof}

A component that has already lasted five years is no more likely to last one more year
than a new component: sensible for a process without aging, absurd for a lamp that wears out.
For $\beta=2$, $P(X>3)=0.22313016$ and a simulation of $10^5$ draws
gives $0.22237$;
conditional on $X>2$, the empirical probability of $X>5$ is $0.22028$ against the theoretical
$0.22313$. The exponential is the continuous twin of the geometric: if arrivals follow a
Poisson process with rate $\lambda$, the gaps between arrivals are exponential with
$\beta=1/\lambda$.

\paragraph{Normal.} $X$ normal with parameters $\mu$ and $\sigma^2>0$ has
density
\[
f_X(x)=\frac{1}{\sigma\sqrt{2\pi}}
\exp\!\left(-\frac{(x-\mu)^2}{2\sigma^2}\right),\qquad x\in\mathbb{R}.
\]
The constant in front is not a choice: it is the only value that makes the area under the
curve equal to one. That the area is one is proved with polar coordinates in Chapter 15:
write $I^2$ as the double integral of $e^{-(x^2+y^2)/2}$, change to polar, and obtain
$I^2=2\pi$.

\begin{proposisi}
If $X\sim N(\mu,\sigma^2)$, then $Z=(X-\mu)/\sigma\sim N(0,1)$.
\end{proposisi}

\begin{proof}
$F_Z(z)=P(X\le\mu+\sigma z)=F_X(\mu+\sigma z)$. Differentiate with respect to $z$ by the chain
rule: $f_Z(z)=\sigma f_X(\mu+\sigma z)=\varphi(z)$, the standard normal density.
\end{proof}

The function $\Phi(z)=P(Z\le z)$ has no elementary form; it is computed numerically or
read from a table. Three numbers worth memorizing: $P(|Z|\le1)=0.682689492137$,
$P(|Z|\le2)=0.954499736104$, $P(|Z|\le3)=0.997300203937$. For a probability that is not
symmetric, subtract two values of $\Phi$: $P(1<Z<2)=0.135905122$. Its quantile,
$\Phi^{-1}(0.975)=1.959964$, appears in every 95 percent confidence interval in
Chapter 20.

The normal is not a more correct distribution than the others. It matters for two
different reasons: many natural quantities approach it, and the sum of many independent variables
tends to approach it whatever their original shape. The second reason is the central limit
theorem, proved in Chapter 19.

\kamus{normal distribution}{distribusi normal}{the symmetric bell curve: almost
all values pile up near the center, and the tails thin out very quickly}

\kamus{exponential distribution}{distribusi eksponensial}{waiting time in a process
without aging: one that has waited a long time is no closer to the next event}

\section{Joint Distributions, Covariance, and Correlation}
\tingkat{advanced}

Two separate distributions do not store information about whether the two quantities
tend to rise together. The tool for that is the joint distribution.

\begin{contoh}
Let $X,Y\in\{0,1\}$ with joint mass $P(0,0)=0.1$, $P(0,1)=0.2$,
$P(1,0)=0.3$, $P(1,1)=0.4$. Adding up rows and columns gives the marginals
$P(X=1)=0.7$ and $P(Y=1)=0.6$. If $X$ and $Y$ were independent, the mass at $(1,1)$
would have to be $0.42$, whereas the actual value is $0.4$; so the two are not independent,
even though the difference is small.
\end{contoh}

In the continuous case, the joint distribution is described by the joint density
$f_{X,Y}\ge0$ with $\iint f_{X,Y}=1$, and the marginals are obtained by
integrating one variable out: $f_X(x)=\int f_{X,Y}(x,y)\,dy$. Independence means
$f_{X,Y}(x,y)=f_X(x)f_Y(y)$ for all $x,y$.

\begin{definisi}[Covariance and correlation]\index{covariance}\index{correlation}
With $\mu_X=\mathbb{E}[X]$ and $\mu_Y=\mathbb{E}[Y]$,
\[
\operatorname{Cov}(X,Y)=\mathbb{E}\bigl[(X-\mu_X)(Y-\mu_Y)\bigr],
\qquad
\rho(X,Y)=\frac{\operatorname{Cov}(X,Y)}{\sigma_X\sigma_Y}.
\]
\end{definisi}

Positive covariance means the two variables tend to deviate in the same direction;
negative means in opposite directions. Correlation is the unitless version of covariance.

\begin{proposisi}
$\operatorname{Cov}(X,Y)=\mathbb{E}[XY]-\mathbb{E}[X]\mathbb{E}[Y]$ and
$-1\le\rho(X,Y)\le1$.
\end{proposisi}

\begin{proof}
Expand the product and use linearity. For the bound on correlation, consider the inner product
$\langle U,V\rangle=\mathbb{E}[UV]$ on centered variables; the Cauchy-Schwarz
inequality of Chapter 10 gives
$|\langle U,V\rangle|^2\le\langle U,U\rangle\langle V,V\rangle$, namely
$\operatorname{Cov}(X,Y)^2\le\sigma_X^2\sigma_Y^2$.
\end{proof}

\begin{teorema}
$\operatorname{Var}(X+Y)=\operatorname{Var}(X)+\operatorname{Var}(Y)
+2\operatorname{Cov}(X,Y)$.
\end{teorema}

\begin{proof}
Write $X+Y-\mu_X-\mu_Y=(X-\mu_X)+(Y-\mu_Y)$, square it, and then use linearity on
the three terms that appear; the cross term is exactly twice the covariance.
\end{proof}

Back to the table above: $\mathbb{E}[X]=7/10$, $\mathbb{E}[Y]=3/5$,
$\mathbb{E}[XY]=2/5$, so that $\operatorname{Cov}(X,Y)=-1/50=-0.02$. With
$\operatorname{Var}(X)=21/100$ and $\operatorname{Var}(Y)=6/25$, the correlation is
$-0.089087$ and $\operatorname{Var}(X+Y)=0.41$.

Zero covariance is often mistaken for independence, whereas the implication runs only one way:
independence implies zero covariance, not the reverse. Take $X$ taking the values $-1,0,1$ with
equal probability and $Y=X^2$; then $\mathbb{E}[X]=0$ and
$\mathbb{E}[XY]=\mathbb{E}[X^3]=0$, so that the correlation is zero, even though $Y$ is determined
completely by $X$. Covariance captures only dependence in the form of a straight line.

\kamus{covariance}{kovarians}{a measure of whether two quantities tend to rise and fall
together, computed from the product of their deviations}

\section{Transformations of Random Variables}
\tingkat{advanced}

This work comes up again and again: prices converted to logarithms, lengths to areas, waiting times
to rates. The question is always the same, namely how the distribution changes when
$Y=r(X)$.

In the discrete case the answer is simple: mass is moved from the old values to the new
values. Let $P(X=-1)=1/4$, $P(X=0)=1/2$, $P(X=1)=1/4$, and $Y=X^2$. Then
$P(Y=0)=1/2$ and $P(Y=1)=P(X=-1)+P(X=1)=1/2$; $Y$ has only two values, fewer
than $X$, because $r(x)=x^2$ is not one-to-one.

In the continuous case the mass is zero at every point, so what is moved is length,
and length changes scale. Three standard steps answer everything: determine
$A_y=\{x:r(x)\le y\}$, compute $F_Y(y)=\int_{A_y}f_X(x)\,dx$, then differentiate
$f_Y=F_Y'$.

\begin{teorema}[Monotone transformation]\index{transformation of random variables}
If $r$ is strictly monotone and differentiable, with inverse $s=r^{-1}$, then
\[
f_Y(y)=f_X\bigl(s(y)\bigr)\left|\frac{ds(y)}{dy}\right| .
\]
\end{teorema}

\begin{proof}
Suppose $r$ is strictly increasing. Then $r(x)\le y$ is equivalent to $x\le s(y)$, so that
$F_Y(y)=F_X(s(y))$, and the chain rule gives $f_Y(y)=f_X(s(y))s'(y)$ with
$s'(y)>0$. If $r$ is strictly decreasing, the inequality reverses: $F_Y(y)=1-F_X(s(y))$, so that
$f_Y(y)=-f_X(s(y))s'(y)=f_X(s(y))|s'(y)|$.
\end{proof}

The factor $|ds/dy|$ is local stretching: if $r$ stretches the scale, the density
decreases, because the same probability is spread over a longer interval. In more than
one dimension that factor becomes the Jacobian determinant, whose general form is in Chapter 15.
Figure~\ref{fig:transformasi-sumbu} shows that stretching for $Y=X^2$: four
intervals of equal length land on intervals that grow longer as $y$ approaches one,
while the density of $Y$ rises near zero, where the intervals shorten.

\begin{figure}[htbp]
\centering
\begin{tikzpicture}[font=\footnotesize,>=Stealth]
  \draw[->,layergrey] (-0.25,0) -- (7.3,0) node[right] {$x$};
  \foreach \px/\lab in {0/0,1.75/0.25,3.5/0.5,5.25/0.75,7/1}{
    \draw[layergrey] (\px,0.07) -- (\px,-0.07);
    \node[above,font=\scriptsize] at (\px,0.07) {$\lab$};
  }
  \node[font=\scriptsize,text=layergrey,anchor=west] at (0.15,0.55)
    {equal-length intervals};
  \draw[->,layergrey] (-0.25,-2.6) -- (7.3,-2.6) node[right] {$y=x^2$};
  \foreach \px/\lab/\anc in {0/0/east,0.44/0.0625/east,1.75/0.25/north,
                             3.94/0.5625/west,7/1/west}{
    \draw[layergrey] (\px,-2.53) -- (\px,-2.67);
    \node[below,font=\scriptsize,anchor=\anc] at (\px,-2.67) {$\lab$};
  }
  \foreach \px/\py in {0/0,1.75/0.44,3.5/1.75,5.25/3.94,7/7}{
    \draw[->,layerindigo] (\px,-0.05) -- (\py,-2.55);
  }
  \node[font=\scriptsize,text=layerindigo,align=center,anchor=west] at (1.0,-1.85)
    {interval shortens,\\ density rises};
  \node[font=\scriptsize,text=layerindigo,align=center,anchor=east] at (6.6,-1.85)
    {interval lengthens,\\ density falls};
  \begin{scope}[xshift=8.7cm,yshift=-2.6cm]
    \draw[->,layergrey] (0,0) -- (5.4,0) node[right] {$y$};
    \draw[->,layergrey] (0,0) -- (0,3.95) node[left] {$f_Y(y)$};
    \draw[layerplum,very thick,domain=0.12:1,samples=100,smooth,variable=\t]
      plot ({5.0*\t},{1.175/sqrt(\t)});
    \draw[layerplum,dashed] (0,0) -- (0,3.8);
    \node[below,font=\scriptsize,anchor=east] at (0,0) {$0$};
    \node[below,font=\scriptsize] at (1.25,0) {$0.25$};
    \node[below,font=\scriptsize,anchor=west] at (5.0,0) {$1$};
    \draw[layergrey,dotted] (5.0,0) -- (5.0,1.175);
    \node[font=\scriptsize,text=layerplum,anchor=west] at (0.25,3.45) {blows up near $0$};
  \end{scope}
\end{tikzpicture}
\caption{The transformation $Y=X^2$ for $X$ uniform on $(0,1)$, drawn on two axes.
The point $x=0.25$ lands at $y=0.0625$, $x=0.5$ at $y=0.25$, and $x=0.75$ at
$y=0.5625$; an interval of length $0.25$ near zero shortens to $0.0625$, whereas the same interval
near one lengthens to $0.4375$. The right-hand panel is the density
$f_Y(y)=1/(2\sqrt{y})$: its value is $2$ at $y=0.0625$, $1$ at $y=0.25$, and $0.5$ at
$y=1$. The density rises exactly where the interval shortens, so that the total area stays
one.}
\label{fig:transformasi-sumbu}
\end{figure}

If $r$ is not monotone, the theorem above does not apply directly, but the three standard
steps still work with $A_y$ a union of intervals.

\begin{contoh}
Let $X$ be uniform on $(-1,3)$, so that $f_X(x)=1/4$ inside it, and $Y=X^2$
taking values in $(0,9)$. For $0<y<1$ the set $A_y=[-\sqrt y,\sqrt y]$ lies entirely
inside $(-1,3)$, so $F_Y(y)=\tfrac14\cdot2\sqrt y=\sqrt y/2$; for $1\le y<9$
that set becomes $[-1,\sqrt y]$, so $F_Y(y)=(\sqrt y+1)/4$. Its derivative is
\[
f_Y(y)=\begin{cases}
\dfrac{1}{4\sqrt y}, & 0<y<1,\\[4pt]
\dfrac{1}{8\sqrt y}, & 1<y<9,\\[4pt]
0, & \text{elsewhere}.
\end{cases}
\]
The two branches add up to one, $\tfrac12$ and $\tfrac12$. The density of $Y$ blows up
near zero: $x^2$ compresses the intervals there.
\end{contoh}

If $X$ is uniform on $(0,1)$ and $Y=X^2$, the result is tidier: $F_Y(y)=\sqrt y$ and
$f_Y(y)=1/(2\sqrt y)$, which is checked numerically in the Lab. If $Y=-\ln X$, then
$F_Y(y)=1-e^{-y}$, exponential with $\beta=1$: taking the logarithm of a uniform number
produces an exponential waiting time, and the converse is what computers use.
If $U$ is uniform and $F$ is the distribution function we want, then $F^{-1}(U)$
has distribution $F$.

\begin{tantangan}
A set of lecture notes proposes the waiting-time density
\[
f(x)=\frac{1}{2\sqrt{x}},\qquad 0<x<4,
\]
on the grounds that it is ``positive, decreasing, and similar in shape to the example above''. From that density
the author computes $P(X\le4)=1$ and then concludes that the distribution is already valid.

Examine that claim. Determine the total area under $f$, repair the density in two different valid
ways, and explain why the repair is not unique.

Hint 1: what decides the matter is not the value of $f$ at a single point, but the whole region
under its curve.

Hint 2: compute $\int_0^4\tfrac12 x^{-1/2}\,dx$; the result is $2$, not $1$. For the
total to be one, the constant in front must be halved to $1/4$, or the interval must be
narrowed to $(0,1)$ with the constant unchanged.

If you get stuck, go back to the section ``Mass Functions and Density Functions''.
\end{tantangan}

\begin{untukanak}
Imagine a lamp on the floor and a ball between the lamp and the wall.
The shadow of the ball on the wall is much larger than the ball, and the farther the ball is from
the lamp, the larger the shadow. Nothing increases in amount: the paint on the ball
stays the same, it is merely spread over a larger area.

Now move the ball a little when it is already far from the lamp. Its shadow moves
a long way. When the ball is still close to the lamp, the same movement shifts the shadow only
a little. That is, one centimeter near the lamp becomes many centimeters on the wall,
whereas one centimeter far from the lamp becomes very many centimeters.

Changing a random variable works like that. Some regions are compressed and some
are stretched, and the resulting density has to adjust so that the total
stays one. If it does not adjust, probability will appear out of thin air or vanish
somewhere.
\end{untukanak}

\section{A Worked Example}
\tingkat{intermediate}

A production line fills boxes containing 20 fuses. Each fuse is defective with
probability $0.35$, independently of the others, and a box is rejected if it contains twelve defects or
more.

With $n=20$ and $p=0.35$, the rejection rule gives $P(X\ge12)=0.019579$ from
the binomial tail, not from a normal approximation. Even the most likely value, $X=7$, has
probability only $18.4$ percent, so that one box does not tell us much.

A simulation of $10^5$ boxes with \texttt{default\_rng(20261024)} gives mean
$7.005260$ and variance $4.552918$. The standard error of the mean is $0.006745$, so the
deviation is $+0.78$ standard errors; the standard error of the variance is about $0.020348$, so the
deviation is $+0.14$. Both are close, and that is to be expected: a simulation does not
prove a formula, it only fails to refute it.

The second check concerns the total probability mass of the Poisson distribution. The sum
$e^{-\lambda}\lambda^k/k!$ up to a sufficiently large $K$ must equal one. It is computed
in log space so that nothing overflows:

\begin{center}\small
\begin{tabular}{rrrr}
\hline
$\lambda$ & $K$ & sum & $1-\text{sum}$ \\
\hline
$0.5$ & 68 & $1.00000000000000000$ & $0$ \\
$1$ & 73 & $1.00000000000000000$ & $0$ \\
$3$ & 83 & $1.00000000000000022$ & $-2.2\times10^{-16}$ \\
$10$ & 107 & $1.00000000000000311$ & $-3.1\times10^{-15}$ \\
$50$ & 194 & $0.99999999999999778$ & $+2.2\times10^{-15}$ \\
$100$ & 280 & $1.00000000000003886$ & $-3.9\times10^{-14}$ \\
\hline
\end{tabular}
\end{center}

The largest deviation appears at $\lambda=100$, namely $3.9\times10^{-14}$, about
175 times the machine epsilon $2.2\times10^{-16}$. That is not a sign that the formula is wrong, but a
sign that adding $281$ terms of very different sizes loses a few of the last
digits; the first fourteen digits are still correct.

The third check concerns the total of the normal density, because the constant
$1/\sqrt{2\pi}$ is easily mistyped. Adaptive quadrature on $[-8,8]$ gives
$0.99999999999999878$, so the shortfall is $1.22\times10^{-15}$. Theory
predicts that shortfall as exactly the tail outside the interval, namely
$2\Phi(-8)=1.24\times10^{-15}$. The two numbers agree to two significant figures: what
is missing is indeed the tail, not an error in the constant. On the whole number line the result is
$0.99999999999999978$ with a printed error estimate of $1.0\times10^{-8}$, a conservative
estimate.

The fourth check, the transformation: the density $1/(2\sqrt y)$ has integral
$0.9999999999999992$, with a remainder of $8\times10^{-16}$ from the end at $0$ where the density
blows up. The simulation gives $F_Y(0.09)=0.29944$, $F_Y(0.25)=0.50003$,
$F_Y(0.49)=0.69944$, $F_Y(0.81)=0.90057$ against the theoretical $0.3$, $0.5$,
$0.7$, $0.9$. The standard error of each number is about $0.00145$ at $p=0.3$ and
$0.00095$ at $p=0.9$; all four deviations are below one standard error.

\kertas{Check by hand. Compute $\binom{20}{7}0.35^7 0.65^{13}$ and
compare it with $0.184401$; explain why $\sum_{k=0}^{20}p_X(k)=1$ follows from the
binomial theorem, not from measurement. Finally: the sample mean $7.005260$
differs from $7$ by $0.00526$. Is that a bug in the program? Answer with the
standard error $0.006745$.}

\section{Computational Lab}
\tingkat{advanced}

The following code is run with Python 3, NumPy 2.5.2, and SciPy 1.18.0. All the numbers
in the table and in the preceding paragraphs come from this run; the seed is fixed at
$20261024$, so that the results can be repeated.

\begin{lab}
\begin{verbatim}
import numpy as np
from math import exp, lgamma, log, sqrt, pi
from scipy import stats
from scipy.integrate import quad

rng = np.random.default_rng(20261024)
R = 100_000

# 1. binomial: exact formula vs Monte Carlo
n, p = 20, 0.35
mu, var = n * p, n * p * (1 - p)
sd = sqrt(var)
x = rng.binomial(n, p, size=R)
se = sd / sqrt(R)
print("Binomial n=20 p=0.35")
print(f"  theory E={mu:.6f} Var={var:.6f} sd={sd:.6f}")
print(f"  MC     mean={x.mean():.6f} Var={x.var(ddof=1):.6f}")
print(f"  SE mean={se:.6f} z={(x.mean()-mu)/se:+.4f}")
print(f"  SE Var  ={var*sqrt(2/(R-1)):.6f}")
print(f"  P(X>=12)={stats.binom.sf(11, n, p):.6f}")

# 2. Poisson pmf sum, computed in log space
print("\nPoisson pmf sum, k=0..K")
for lam in (0.5, 1.0, 3.0, 10.0, 50.0, 100.0):
    K = int(lam + 12 * sqrt(lam) + 60)
    s = sum(exp(-lam + k*log(lam) - lgamma(k+1)) for k in range(K+1))
    print(f"  lam={lam:6.1f} K={K:4d} sum={s:.17f}"
          f" 1-sum={1-s:+.2e}")

# 3. total normal density must be 1
phi = lambda t: exp(-t*t/2) / sqrt(2*pi)
print("\nStandard normal density")
print(f"  phi(0)={phi(0):.15f} 1/sqrt(2pi)={1/sqrt(2*pi):.15f}")
print(f"  quad [-8,8]  ={quad(phi, -8, 8)[0]:.17f}")
print(f"  remainder={1-quad(phi, -8, 8)[0]:.2e}"
      f"  2*Phi(-8)={2*stats.norm.cdf(-8):.2e}")
v, e = quad(phi, -np.inf, np.inf)
print(f"  quad (-inf,inf)={v:.17f} reported error={e:.1e}")

# 4. transformation X~Unif(0,1)
u = rng.uniform(0, 1, R)
y = u * u
print("\nY=X^2, X~Unif(0,1)")
print(f"  integral of f_Y={quad(lambda t: 1/(2*sqrt(t)), 0, 1)[0]:.16f}")
for yy in (0.09, 0.25, 0.49, 0.81):
    print(f"  F_Y({yy:.2f}) MC={np.mean(y<=yy):.5f} theory={sqrt(yy):.5f}")
z = -np.log(u)
print("Y=-ln(X), should be Exp(1)")
print(f"  mean={z.mean():.5f} SE={1/sqrt(R):.5f}"
      f" Var={z.var(ddof=1):.5f}")
print(f"  P(Y>2) MC={np.mean(z>2):.5f} theory={exp(-2):.5f}")

# 5. geometric and exponential
g = rng.geometric(0.2, R)
print("\nGeometric p=0.2: E=5 Var=20")
print(f"  mean={g.mean():.5f} SE={sqrt(20/R):.5f}"
      f" Var={g.var(ddof=1):.5f}")
w = rng.exponential(2.0, R)
print("Exponential beta=2: E=2 Var=4")
print(f"  mean={w.mean():.5f} SE={2/sqrt(R):.5f}"
      f" Var={w.var(ddof=1):.5f}")
print(f"  P(X>3) MC={np.mean(w>3):.5f} theory={exp(-1.5):.5f}")
print(f"  P(X>5|X>2) MC={np.mean(w[w>2]>5):.5f} theory={exp(-1.5):.5f}")
\end{verbatim}
\end{lab}

The output obtained, as it stands:

\begin{verbatim}
Binomial n=20 p=0.35
  theory E=7.000000 Var=4.550000 sd=2.133073
  MC     mean=7.005260 Var=4.552918
  SE mean=0.006745 z=+0.7798
  SE Var  =0.020348
  P(X>=12)=0.019579

Poisson pmf sum, k=0..K
  lam=   0.5 K=  68 sum=1.00000000000000000 1-sum=+0.00e+00
  lam=   1.0 K=  73 sum=1.00000000000000000 1-sum=+0.00e+00
  lam=   3.0 K=  83 sum=1.00000000000000022 1-sum=-2.22e-16
  lam=  10.0 K= 107 sum=1.00000000000000311 1-sum=-3.11e-15
  lam=  50.0 K= 194 sum=0.99999999999999778 1-sum=+2.22e-15
  lam= 100.0 K= 280 sum=1.00000000000003886 1-sum=-3.89e-14

Standard normal density
  phi(0)=0.398942280401433 1/sqrt(2pi)=0.398942280401433
  quad [-8,8]  =0.99999999999999878
  remainder=1.22e-15  2*Phi(-8)=1.24e-15
  quad (-inf,inf)=0.99999999999999978 reported error=1.0e-08

Y=X^2, X~Unif(0,1)
  integral of f_Y=0.9999999999999992
  F_Y(0.09) MC=0.29944 theory=0.30000
  F_Y(0.25) MC=0.50003 theory=0.50000
  F_Y(0.49) MC=0.69944 theory=0.70000
  F_Y(0.81) MC=0.90057 theory=0.90000
Y=-ln(X), should be Exp(1)
  mean=0.99811 SE=0.00316 Var=0.99623
  P(Y>2) MC=0.13478 theory=0.13534

Geometric p=0.2: E=5 Var=20
  mean=4.99176 SE=0.01414 Var=19.92571
Exponential beta=2: E=2 Var=4
  mean=1.99895 SE=0.00632 Var=3.98400
  P(X>3) MC=0.22237 theory=0.22313
  P(X>5|X>2) MC=0.22028 theory=0.22313
\end{verbatim}

Four things can be read from that output. First, every sample mean deviates by less
than one standard error: $z=+0.78$ for the binomial, whereas $0.99811$, $4.99176$,
and $1.99895$ against $1$, $5$, and $2$ for the other three cases. A deviation
of that size has to appear; if all of them were zero, something would be wrong.

Second, the sample variances are slightly below the theoretical value in the long-tailed case,
$19.92571$ against $20$, because a single very large draw dominates the sum
of squares. Third, the error of the Poisson summation grows with $\lambda$, from zero to
$10^{-14}$ as the number of terms increases from $74$ to $281$. Fourth, the resulting
density integrates to one to $16$ digits, and $F_Y$ matches $\sqrt y$ at four test
points; the second is more important, because an integral of one only checks the constant,
whereas agreement at four points checks the shape of the function.

That code does not prove a single theorem in this chapter; it checks that the formulas
written down give the numbers stated. The proof of $\mathbb{E}[X]=np$ is in the indicator
argument, the proof of $f_Y(y)=1/(2\sqrt y)$ in the computation of the distribution function;
the output of a program is data, not an argument.

\begin{untukdosen}
The order I suggest: two systems with the same mean but different tails, then
the definition of a random variable, then straight to the cumulative distribution function. Do one small
table on the board, draw the staircase, ask for $F(1.4)$, and only after that the density.
Resist the temptation to mention the integral before students see why summation
is no longer enough.

The two misconceptions that appear most often. First, $f_X(x)$ is read as a probability,
so that a density of value $5$ feels impossible and $P(X=x)$ is taken to equal
$f_X(x)$. Second, $P(X=x)=0$ is concluded to mean impossible, so that
$P(0\le X\le1)=1$ becomes a paradox. The answer is countable additivity: an uncountable
union is not added up.

Sign of readiness to move on: students can explain the difference between mass and density without
reading their notes, and can compute $\mathbb{E}[X^2]$ without finding the density of $X^2$.
\end{untukdosen}

\section{Summary}
\tingkat{basic}

\begin{ringkasan}
\begin{itemize}
\item A random variable is a function from the outcome space to the real numbers, not a random
      number. The distribution function $F_X(x)=P(X\le x)$ carries the whole information of the distribution:
      nondecreasing, right-continuous, with limits $0$ and $1$.
\item A mass function weights points and its values are at most one; a density function
      weights length, may exceed one, and $P(X=x)=0$ at every point. Probability
      zero is not impossibility, because only countable unions are added up.
\item Expectation is the balance point and is linear without any condition of independence;
      $\operatorname{Var}(X)=\mathbb{E}[X^2]-\mathbb{E}[X]^2$ is the mean squared
      deviation.
\item Binomial $\mathbb{E}=np$, $\operatorname{Var}=np(1-p)$, derived through
      indicators; geometric $\mathbb{E}=1/p$, $\operatorname{Var}=(1-p)/p^2$; Poisson
      $\mathbb{E}=\operatorname{Var}=\lambda$, the binomial limit for rare events.
\item The normal is standardized to $N(0,1)$ and read through $\Phi$; the exponential is
      memoryless, the continuous twin of the geometric.
\item Covariance measures only linear dependence: zero does not mean independent, and a
      monotone transformation changes the density by the factor $|ds/dy|$.
\end{itemize}
\end{ringkasan}

\section{Exercises}
\tingkat{intermediate}

\begin{latihan}
\begin{enumerate}
  \item Let $p_X(-2)=0.1$, $p_X(0)=0.4$, $p_X(3)=c$. Determine $c$, then
        compute $\mathbb{E}[X]$, $\mathbb{E}[X^2]$, and $\operatorname{Var}(X)$.
  \item Prove $P(a<X\le b)=F_X(b)-F_X(a)$, then explain why
        $P(a\le X\le b)$ can differ from $P(a<X\le b)$ for a discrete variable and
        why the two are equal for a continuous variable.
  \item Assess this claim: ``because $P(X=x)=0$ for every $x$, the event $X=x$
        is impossible.'' Name the part that is correct, the part that is wrong, and the role of countable
        additivity.
  \item Derive $\mathbb{E}[X^2]=(2-p)/p^2$ for the geometric by differentiating
        $\sum k^2q^{k-1}$ twice, then conclude $\operatorname{Var}(X)$ and check
        for $p=0.2$.
  \item Let $X\sim N(3,4)$. Compute $P(X>1)$, $P(1<X<5)$, and the 0.2 quantile of
        $X$; compare with \texttt{scipy.stats.norm}.
  \item Compute $\sum_{k=0}^{40}e^{-2.5}2.5^k/k!$ and $P(X\ge8)$ for the Poisson
        with $\lambda=2.5$, then compare with the binomial $n=1000$, $p=0.0025$;
        explain the direction of the difference.
  \item Prove that the exponential is memoryless, then explain why the geometric is also
        memoryless and why the uniform on the integers is not.
  \item For $X$ uniform on $(-1,3)$ and $Y=X^2$, derive $f_Y$, check
        $\int f_Y=1$, then compute $\mathbb{E}[Y]$ through $f_Y$ and through
        $\mathbb{E}[X^2]$.
  \item Construct two variables whose covariance is zero but which are not independent. Compute
        $\operatorname{Cov}$ and $\rho$ exactly, and exhibit a joint probability
        that proves the dependence.
  \item How many repetitions $R$ are needed for the standard error of the mean of a binomial with $n=20$, $p=0.35$
        to be at most $0.005$? And for the standard error of the variance estimate to be at
        most $0.02$?
\end{enumerate}

\noindent\textbf{Hints.} Item 3: separate the statement about probability from the
statement about the existence of a set. Item 4: differentiate $\sum_{k\ge0}q^k=1/(1-q)$
twice with respect to $q$, and remember that $\sum k^2q^{k-1}$ contains $\sum k(k-1)q^{k-1}$.
Item 8: for $0<y<1$ the set $A_y$ has two branches, for $y\ge1$ only one.
\end{latihan}

\begin{untukpeneliti}
What is genuinely proved in this chapter: the properties of the distribution function together with their proofs,
the relation of the density to the distribution function through the fundamental theorem of calculus, the formula for the
expectation of a function of a random variable in the discrete case, linearity of expectation, the variance formula,
the mean and variance of the binomial through indicators, the geometric and Poisson means through
series, memorylessness of the exponential, and the monotone transformation formula.

What is quoted, not proved: that a function satisfying the three properties of a distribution function
always comes from some probability space (Wasserman, Theorem 2.8, whose direction uses the
Caratheodory theorem), that the distribution function determines every probability (Wasserman,
Theorem 2.7), the normalization of the normal through polar coordinates (Chapter 15), and the correlation
bound through Cauchy-Schwarz (Chapter 10).

Primary sources: Wasserman, \emph{All of Statistics}, Sections 2.1--2.4 (including Lemma
2.15 and the warning that a density may exceed one) and 2.11 (Example 2.47 for
$Y=X^2$ on $U(-1,3)$), as well as 3.1--3.4. Blitzstein and Hwang, \emph{Introduction to
Probability}, Chapters 3--5 (LOTUS in 4.5, Poisson in 4.7, exponential in 5.5) and 7.3
for covariance.

Limits of validity: the density is not unique, so every statement about $f_X$
is a statement about the same class of functions almost everywhere; the Jacobian formula
requires monotonicity and differentiability; the Cauchy expectation does not exist, so every
``mean'' of very heavy-tailed data must be suspected first. Open question: how well a density can be estimated
from data without assuming a shape, and whether the tails of real distributions are thin enough to
justify the use of variance.
\end{untukpeneliti}

\peta{Chapter 18 follows Wasserman, \emph{All of Statistics}, Chapter 2 (2.1--2.4, Lemma 2.15,
and 2.11 for transformation) and Chapter 3 (3.1--3.4). In Blitzstein and Hwang,
\emph{Introduction to Probability}, Chapter 3 (3.6--3.7), Chapter 4 (4.2, 4.5, 4.7), Chapter 5
(5.1--5.5), and 7.3. The integrals used are the Riemann integrals of Chapter 13; change of
variables in more than one dimension is in Chapter 15.}

% Judul penuh meluber 14,54-20,02 pt di kepala halaman; hanya mark header yang
% dipendekkan (sekali pakai, lihat \judulpendek di preamble-en.tex).
\judulpendek{Law of Large Numbers and Central Limit Theorem}
\chapter{The Law of Large Numbers and the Central Limit Theorem}

A short program computes the average of $\exp(U)$ over ten thousand random numbers
$U$ uniform on $(0,1)$. The result it prints is $1.722206$. The true value is
$e-1=1.7182818285$, so the error of that run is $+0.003924$. The program is run
again with four times as many numbers: $1.717502$, error $-0.000780$. With ten
times as many: $1.717767$, error $-0.000515$. The pattern is not tidy. Four
times the data cuts the error by more than a factor of five, whereas ten times the
data cuts it by only $7.6$ times. The naive guess "ten times the data, one tenth
of the error" does not hold for a single run. One run of a program carries no rate
at all; what has a rate is the distribution over many runs.

That distribution can be measured. With $n=10^4$ the standard deviation of this
estimator is $0.00492$; with $n=4\cdot10^4$ it becomes $0.00246$; with $n=10^5$
it becomes $0.00156$. Every time $n$ grows fourfold, that deviation is exactly
halved. Here we have two different statements that are often confused. First, the
average really does approach the true value: that is a tendency, and its name is
the law of large numbers. Second, the remaining error has a definite size and
shape: that is fluctuation, and its name is the central limit theorem. This chapter
supplies names, proofs, and limits of validity for both.

Chapter 18 has already prepared the ingredients: expectation, variance, and the
distribution of a sum through convolution. Convolution is useful only when $n$ is
small and the original shape is simple; for most distributions no closed form
exists. What we need is not the exact distribution at one $n$, but a statement that
holds for every large $n$. The practical questions are two: how close may the
average be trusted, and how large an error is reasonable. The answers come from
examining how the sequence of distributions behaves, not from computing one
distribution.

\begin{untukanak}
Imagine a class is asked to guess the number of marbles in a large jar. The
children's guesses miss badly in both directions. One child guesses three hundred,
another guesses two thousand, when the jar holds eight hundred. The teacher then
writes all the guesses on the board and computes their average. The number turns
out to be very close to eight hundred.

Now the teacher adds students. With thirty guessers the average still wobbles. With
three hundred guessers the average is calmer. With three thousand guessers the
average hardly shifts any more. There are two strange things here. First, not one
of the guessers has become better; only the number of guessers changed. Second,
adding guessers from thirty to three hundred helps far more than adding from three
hundred to three thousand. The help keeps growing, but more and more slowly.

The question that makes us curious: if each guess can miss by a great deal, why
can the average be trusted? And why does the help slow down in a pattern we can
predict, rather than at random? In this chapter we will find out why, and at the
same time find out when the habit of trusting the average deceives us.
\end{untukanak}

\section{Three Meanings of Approaching for a Random Variable}
\tingkat{basic}

In everyday language: a sequence of random variables can "approach" in three
different senses, and the loosest sense says nothing at all about the values coming
closer.

A sequence of real numbers has only one sense of approaching, namely that the
distance $|x_n-L|$ becomes small. For random variables that distance is itself a
random variable, so we have to choose which quantity is made small. The first
choice is the probability. This is the choice that is weakest intuitively but most
useful practically.

\kamus{convergence in probability}{konvergensi dalam peluang}{a promise that the
probability of missing by more than a tolerance shrinks to zero, not a promise that
the value comes closer on a single run}

\begin{definisi}[Convergence in probability]
A sequence of random variables $X_1,X_2,\dots$ \emph{converges in probability}
\index{convergence in probability} to a random variable $X$ if for every
$\varepsilon>0$
\[
P\bigl(|X_n-X|>\varepsilon\bigr)\to 0 \qquad (n\to\infty).
\]
\end{definisi}

Note where the quantifier sits. The value of $\varepsilon$ is chosen first, and the
probability may shrink only after that. This definition does not say that
$|X_n-X|$ shrinks for one particular outcome of the experiment; it says that the
event "missing by more than $\varepsilon$" becomes rarer and rarer. On the same
experiment, $X_n$ may still miss by a lot, provided that event is rare.

\kamus{convergence in distribution}{konvergensi dalam distribusi}{what resembles is
the shape of the distribution, not the value itself; two variables can share a shape
without ever coming close}

\begin{definisi}[Convergence in distribution]
Let $F_n$ be the distribution function of $X_n$ and $F$ the distribution function
of $X$. We write $X_n\rightsquigarrow X$ when $F_n(t)\to F(t)$ at every point $t$
at which $F$ is continuous. That limit is called the \emph{limit in
distribution}\index{convergence in
distribution}.
\end{definisi}

The condition "at continuity points" is not a technical whim. Take $X_n$ equal to
$1/n$ with probability one, and $X$ equal to $0$ with probability one. Then
$F_n(0)=0$ for every $n$, whereas $F(0)=1$, so the sequence does not converge at
$t=0$. At every $t\neq0$ the sequence does converge, and $t=0$ is the only point of
discontinuity of $F$. Without that condition, a correct conclusion would be
rejected because of a single point.

The third sense is stronger and more literal: almost every outcome of the
experiment, as one point $\omega$ of the sample space, produces a convergent
sequence of real numbers.

\begin{definisi}[Almost sure convergence]
$X_n$ \emph{converges almost surely}\index{convergence almost surely} to $X$ when
\[
P\bigl(\{\omega : X_n(\omega)\to X(\omega)\}\bigr)=1 .
\]
\end{definisi}

Here the sequence of real numbers $X_1(\omega),X_2(\omega),\dots$ really does
converge in the sense of Chapter 4, and the set of outcomes that fail to converge
may only have probability zero. Therefore almost sure convergence speaks about one
long run, whereas convergence in probability speaks about one $n$ across many runs
at once.

The three senses are not equivalent, but there is a fixed direction of implication.

\begin{teorema}[Relations among the three modes]
For random variables $X_n$ and $X$:
(a) $X_n\to X$ almost surely implies $X_n\to X$ in probability;
(b) $X_n\to X$ in probability implies $X_n\rightsquigarrow X$;
(c) if $X_n\rightsquigarrow c$ and $c$ is a constant, then $X_n\to c$ in
probability.
\end{teorema}

\begin{proof}[Proof]
For (a), fix $\varepsilon>0$ and let
$A_n=\{\sup_{k\ge n}|X_k-X|>\varepsilon\}$. The sets $A_n$ decrease, and every
$\omega$ in $\bigcap_n A_n$ satisfies $|X_k(\omega)-X(\omega)|>\varepsilon$ for
infinitely many $k$, so the sequence does not converge to $X(\omega)$. Hence
$\bigcap_n A_n$ is contained in a set of probability zero. Since probability is
continuous from above on sets of finite probability,
$P(A_n)\to P(\bigcap_n A_n)=0$. The event $\{|X_n-X|>\varepsilon\}$ is contained in
$A_n$, so its probability also tends to zero.

For (b), take $t$ a continuity point of $F$ and $\varepsilon>0$. From
$\{X_n\le t\}\subseteq\{X\le t+\varepsilon\}\cup\{|X_n-X|>\varepsilon\}$ we obtain
\[
F_n(t)\le F(t+\varepsilon)+P\bigl(|X_n-X|>\varepsilon\bigr).
\]
In the same way, from $\{X\le t-\varepsilon\}\subseteq\{X_n\le t\}\cup
\{|X_n-X|>\varepsilon\}$ we obtain
$F(t-\varepsilon)\le F_n(t)+P(|X_n-X|>\varepsilon)$. The two inequalities give
\[
F(t-\varepsilon)-\delta_n\le F_n(t)\le F(t+\varepsilon)+\delta_n,
\qquad \delta_n=P\bigl(|X_n-X|>\varepsilon\bigr).
\]
Since $\delta_n\to0$, the limit inferior and the limit superior of $F_n(t)$ are
trapped between $F(t-\varepsilon)$ and $F(t+\varepsilon)$. Let
$\varepsilon\downarrow0$ and use the continuity of $F$ at $t$; both tend to $F(t)$.

For (c), take $\varepsilon>0$. Since $F$ is the point-mass distribution at $c$, the
points $c-\varepsilon$ and $c+\varepsilon$ are both continuity points of $F$, so
\[
P\bigl(|X_n-c|>\varepsilon\bigr)
\le P(X_n<c-\varepsilon)+P(X_n>c+\varepsilon)
\le F_n(c-\varepsilon)+1-F_n(c+\varepsilon)
\to 0+1-1=0 .
\]
\end{proof}

The three reverse implications fail, and the counterexamples are worth remembering
because each one shows what is actually lost.

\begin{contoh}[In probability but not almost surely]
Split $(0,1)$ into $k$ pieces of equal length, for $k=1,2,3,\dots$, then arrange all
those pieces into one list. Define $X_n(\omega)=1$ if $\omega$ lies in the $n$-th
piece and $0$ otherwise. The $n$-th piece comes from the split into $k_n$ pieces, so
$P(X_n=1)=1/k_n\to0$ and $X_n\to0$ in probability. But for every
$\omega\in(0,1)$, within every block of splits there is exactly one $n$ with
$X_n(\omega)=1$. That value one appears infinitely often, so $X_n(\omega)$ does not
converge, for every $\omega$. Hence there is no almost sure convergence. This
sequence even converges in $L^1$, because $\mathbb{E}|X_n|=1/k_n\to0$, so $L^1$
strength does not close that gap either.
\end{contoh}

\begin{contoh}[In distribution but not in probability]
Take $X\sim N(0,1)$ and set $X_n=-X$ for every $n$. The distribution of $X_n$ is
exactly the distribution of $X$, so $F_n(t)=F(t)$ and $X_n\rightsquigarrow X$
vacuously. But $P(|X_n-X|>\varepsilon)=P(2|X|>\varepsilon)>0$ for every
$\varepsilon$, and that probability does not depend on $n$. Two variables of the
same shape need not be close: $X$ and $-X$ can be far apart on the same run.
\end{contoh}

The lesson of those two examples: convergence in distribution touches the
distributions one at a time, whereas the other two modes touch the relation between
variables on the same sample space. That is why part (c) of the theorem matters
practically: in asymptotic statistics the limit is almost always a constant, and
there the three modes meet.

One further warning has to be recorded from the start, because the mistake is
expensive.

\begin{catatan}
Convergence in probability does not carry convergence of expectations with it.
Suppose $P(X_n=n^2)=1/n$ and $P(X_n=0)=1-1/n$. For every $\varepsilon>0$,
$P(|X_n|>\varepsilon)=1/n\to0$, so $X_n\to0$ in probability. But
$\mathbb{E}[X_n]=n$, so the expectation tends to infinity. An average that
concentrates is no guarantee that its expectation concentrates.
\end{catatan}

\section{Chebyshev's Inequality and the Weak Law of Large Numbers}
\tingkat{intermediate}

In everyday language: the probability of straying far from the mean is bounded by
the variance divided by the square of that stray distance.

To prove the law of large numbers we need only one inequality, and that inequality
needs only one simple observation about expectation.

\begin{lemma}[Markov]
If $Y\ge0$ and $a>0$, then $P(Y\ge a)\le \mathbb{E}[Y]/a$.
\index{Markov's inequality}
\end{lemma}

\begin{proof}[Proof]
Write $Y=Y\mathbf{1}_{\{Y<a\}}+Y\mathbf{1}_{\{Y\ge a\}}$. The first term is
nonnegative. On the event $Y\ge a$ we have $Y\ge a$, so the expectation of the
second term is at least $a\,P(Y\ge a)$. Hence
$\mathbb{E}[Y]\ge 0+a\,P(Y\ge a)$.
\end{proof}

\begin{akibat}[Chebyshev's inequality]
If $\mathbb{E}[X]=\mu$ and $\operatorname{Var}(X)=\sigma^2<\infty$, then for every
$a>0$
\[
P\bigl(|X-\mu|\ge a\bigr)\le\frac{\sigma^2}{a^2}.
\]
\index{Chebyshev's inequality}
\end{akibat}

\begin{proof}[Proof]
Apply Markov's lemma to $Y=(X-\mu)^2$ with threshold $a^2$. The event
$\{Y\ge a^2\}$ is the same as $\{|X-\mu|\ge a\}$, and $\mathbb{E}[Y]=\sigma^2$.
\end{proof}

This inequality is weak, and its weakness is not a hidden defect but its strength.
It knows nothing about the shape of any distribution, only two moments. Therefore it
holds for every distribution with finite variance, and precisely for that reason it
is enough to prove a theorem that holds in general.

\kamus{law of large numbers}{hukum bilangan besar}{a promise that the average of
many observations sticks to the true mean, with a probability of missing that
shrinks}

\begin{teorema}[Weak law of large numbers]
Suppose $X_1,X_2,\dots$ are independent and identically distributed with
$\mathbb{E}[X_i]=\mu$ and $\operatorname{Var}(X_i)=\sigma^2<\infty$. Then for every
$\varepsilon>0$\index{law of large numbers}
\[
P\bigl(|\bar X_n-\mu|\ge\varepsilon\bigr)\to0,
\qquad \bar X_n=\frac1n\sum_{i=1}^n X_i .
\]
\end{teorema}

\begin{proof}[Proof]
Since expectation is linear, $\mathbb{E}[\bar X_n]=\mu$. Since the variables are
independent, the variance of a sum is the sum of the variances, so
$\operatorname{Var}(\bar X_n)=\sigma^2/n$. Chebyshev's inequality with
$a=\varepsilon$ gives
\[
P\bigl(|\bar X_n-\mu|\ge\varepsilon\bigr)\le\frac{\sigma^2}{n\varepsilon^2}.
\]
The right-hand side tends to zero for each fixed $\varepsilon>0$.
\end{proof}

Note what has just been proved and what has not. What is proved: for a tolerance
$\varepsilon$ that we choose ourselves, the probability that the average misses by
more than $\varepsilon$ tends to zero. What is not proved: how fast, and what the
shape of the remaining error is. Those are two different questions.

Figure~\ref{fig:tiga-lebar} draws the content of that statement for dice. Only two
things change from $n=4$ to $n=64$: the width shrinks and the peak rises, while the
center stays at $3.5$. The law of large numbers is not a promise that one run is
more accurate, but that the distribution of possibilities contracts to a single
point.

\begin{figure}[htbp]
\centering
\begin{tikzpicture}
\begin{axis}[
  width=0.95\linewidth, height=5.4cm,
  xmin=1, xmax=6, ymin=0, ymax=2.05,
  xlabel={value of the average $\bar x_n$}, ylabel={density},
  label style={font=\small}, tick label style={font=\footnotesize},
  axis lines=left,
  legend style={font=\footnotesize, draw=layergrey!35, fill=white,
    fill opacity=0.9, text opacity=1, at={(0.015,0.985)},
    anchor=north west, cells={anchor=west}},
  legend cell align=left,
]
\addplot[draw=layerteal!60!black, fill=layerteal!18, fill opacity=0.85,
  line width=0.4pt] coordinates{
  (1,0.0031) (1.25,0.0123) (1.5,0.0309) (1.75,0.0617) (2,0.108) (2.25,0.1728)
  (2.5,0.2469) (2.75,0.321) (3,0.3858) (3.25,0.4321) (3.5,0.4506) (3.75,0.4321)
  (4,0.3858) (4.25,0.321) (4.5,0.2469) (4.75,0.1728) (5,0.108) (5.25,0.0617)
  (5.5,0.0309) (5.75,0.0123) (6,0.0031)};
\addlegendentry{$n=4$}
\addplot[draw=layerindigo!70!black, fill=layerindigo!22, fill opacity=0.8,
  line width=0.4pt] coordinates{
  (1.6875,0) (1.75,0.0001) (1.8125,0.0002) (1.875,0.0004) (1.9375,0.0008) (2,0.0014)
  (2.0625,0.0025) (2.125,0.0043) (2.1875,0.0072) (2.25,0.0117) (2.3125,0.0184) (2.375,0.028)
  (2.4375,0.0416) (2.5,0.0601) (2.5625,0.0847) (2.625,0.1163) (2.6875,0.1558) (2.75,0.2039)
  (2.8125,0.2607) (2.875,0.3258) (2.9375,0.3981) (3,0.4759) (3.0625,0.5567) (3.125,0.6374)
  (3.1875,0.7145) (3.25,0.7843) (3.3125,0.8431) (3.375,0.8878) (3.4375,0.9156) (3.5,0.9251)
  (3.5625,0.9156) (3.625,0.8878) (3.6875,0.8431) (3.75,0.7843) (3.8125,0.7145) (3.875,0.6374)
  (3.9375,0.5567) (4,0.4759) (4.0625,0.3981) (4.125,0.3258) (4.1875,0.2607) (4.25,0.2039)
  (4.3125,0.1558) (4.375,0.1163) (4.4375,0.0847) (4.5,0.0601) (4.5625,0.0416) (4.625,0.028)
  (4.6875,0.0184) (4.75,0.0117) (4.8125,0.0072) (4.875,0.0043) (4.9375,0.0025) (5,0.0014)
  (5.0625,0.0008) (5.125,0.0004) (5.1875,0.0002) (5.25,0.0001) (5.3125,0)};
\addlegendentry{$n=16$}
\addplot[draw=accent!80!black, fill=accent!30, fill opacity=0.8,
  line width=0.5pt] coordinates{
  (2.5312,0) (2.5469,0.0001) (2.5625,0.0001) (2.5781,0.0001) (2.5938,0.0002) (2.6094,0.0003)
  (2.625,0.0004) (2.6406,0.0005) (2.6562,0.0007) (2.6719,0.0009) (2.6875,0.0012) (2.7031,0.0016)
  (2.7188,0.0021) (2.7344,0.0028) (2.75,0.0036) (2.7656,0.0047) (2.7812,0.0061) (2.7969,0.0079)
  (2.8125,0.01) (2.8281,0.0127) (2.8438,0.0161) (2.8594,0.0202) (2.875,0.0252) (2.8906,0.0312)
  (2.9062,0.0385) (2.9219,0.0472) (2.9375,0.0576) (2.9531,0.0698) (2.9688,0.0842) (2.9844,0.1009)
  (3,0.1203) (3.0156,0.1426) (3.0312,0.1681) (3.0469,0.1971) (3.0625,0.2297) (3.0781,0.2663)
  (3.0938,0.307) (3.1094,0.3521) (3.125,0.4015) (3.1406,0.4553) (3.1562,0.5136) (3.1719,0.5762)
  (3.1875,0.6429) (3.2031,0.7134) (3.2188,0.7875) (3.2344,0.8645) (3.25,0.944) (3.2656,1.0252)
  (3.2812,1.1075) (3.2969,1.19) (3.3125,1.2718) (3.3281,1.352) (3.3438,1.4295) (3.3594,1.5035)
  (3.375,1.573) (3.3906,1.6369) (3.4062,1.6943) (3.4219,1.7445) (3.4375,1.7867) (3.4531,1.8202)
  (3.4688,1.8445) (3.4844,1.8592) (3.5,1.8641) (3.5156,1.8592) (3.5312,1.8445) (3.5469,1.8202)
  (3.5625,1.7867) (3.5781,1.7445) (3.5938,1.6943) (3.6094,1.6369) (3.625,1.573) (3.6406,1.5035)
  (3.6562,1.4295) (3.6719,1.352) (3.6875,1.2718) (3.7031,1.19) (3.7188,1.1075) (3.7344,1.0252)
  (3.75,0.944) (3.7656,0.8645) (3.7812,0.7875) (3.7969,0.7134) (3.8125,0.6429) (3.8281,0.5762)
  (3.8438,0.5136) (3.8594,0.4553) (3.875,0.4015) (3.8906,0.3521) (3.9062,0.307) (3.9219,0.2663)
  (3.9375,0.2297) (3.9531,0.1971) (3.9688,0.1681) (3.9844,0.1426) (4,0.1203) (4.0156,0.1009)
  (4.0312,0.0842) (4.0469,0.0698) (4.0625,0.0576) (4.0781,0.0472) (4.0938,0.0385) (4.1094,0.0312)
  (4.125,0.0252) (4.1406,0.0202) (4.1562,0.0161) (4.1719,0.0127) (4.1875,0.01) (4.2031,0.0079)
  (4.2188,0.0061) (4.2344,0.0047) (4.25,0.0036) (4.2656,0.0028) (4.2812,0.0021) (4.2969,0.0016)
  (4.3125,0.0012) (4.3281,0.0009) (4.3438,0.0007) (4.3594,0.0005) (4.375,0.0004) (4.3906,0.0003)
  (4.4062,0.0002) (4.4219,0.0001) (4.4375,0.0001) (4.4531,0.0001) (4.4688,0)};
\addlegendentry{$n=64$}
\draw[layergrey, densely dotted, line width=0.6pt] (axis cs:3.5,0) -- (axis cs:3.5,1.95);
\node[font=\footnotesize, layergrey, anchor=south east] at (axis cs:3.45,1.95) {$3.5$};
\end{axis}
\end{tikzpicture}
\caption{The three distributions of the average of $n$ die rolls, drawn on the same
scale: $n=4$, $16$, and $64$. Their widths are $0.853913$, $0.426956$, and
$0.213478$, so each time $n$ is multiplied by four the width is halved. The peak
rises from $0.4506$ to $0.9251$ and $1.8641$, while the center stays at
$3.5$.}
\label{fig:tiga-lebar}
\end{figure}

The numbers make that difference tangible. For a coin toss with $p=1/2$ and
tolerance $\varepsilon=0.1$, Chebyshev gives
$\sigma^2/(n\varepsilon^2)=0.25/(0.01n)=25/n$. For the probability of missing to
be at most $0.3$, it suffices that $n\ge25/0.3=83.33$, that is $n=84$. At
$n=84$ that bound equals $0.297619$, so Chebyshev guarantees that the probability
of landing between $0.4$ and $0.6$ is at least $0.702381$. The true
probability, computed from the binomial distribution, is $0.937028$, and the
normal approximation gives $0.933197$. The guarantee is correct but too low by
about $23$ percentage points.

The same weakness becomes severe when $\sigma$ is large. For a fair die,
$\operatorname{Var}(X)=35/12=2.916667$ and $\sigma=1.707825$. At $n=100$ with
$\varepsilon=0.1$, Chebyshev gives the bound $2.916667$, a number larger than
one and therefore meaningless. To obtain an error of $0.01$ with $95$ percent
confidence, the normal formula demands at least $112043$ observations, whereas
Chebyshev demands far more still.

This is where the problem that sets up the next section lies. Chebyshev answers "may
the average be trusted" correctly but bluntly. It cannot answer "how close, with
what probability" in a usable way. A sharp answer demands knowledge of the shape of
the error distribution, not only of its size.

\begin{bacakeras}
Two voices, taking turns. The first names the number of tosses; the second names the
distance off that is still reasonable.

--- Ten tosses. The miss is still wide.
--- A hundred tosses. Already narrower.
--- A thousand tosses. Narrower still.

Now quadruple, do not add ten: ten becomes forty, a hundred becomes four hundred, a
thousand becomes four thousand. The miss is half. Not a quarter. Half. Every time
fourfold, exactly half.

Then reverse the direction, and read it more slowly. What is the cost of one more
digit after the decimal point? A hundred times the tosses. A hundred times the work,
one digit. The average is impatient: it demands four times, not two, and it offers
no discount.

The line repeated three times: four times the data, half the error.
\end{bacakeras}

\begin{untukanak}
Imagine a machine that fills bags of rice. Each bag misses the requested weight by a
little, sometimes over and sometimes under, and the size of the miss differs from
bag to bag; weighing one bag tells us nothing about the direction of its miss.

Now the machine fills a hundred bags, and we weigh all of them at once. The misses
upward and downward cancel each other, so the total weight is much closer to a
hundred times the requested weight. If there are ten thousand bags, the total is
almost exact.

The surprising part is that we do not need to know the shape of the distribution of
the miss. Whether the miss has a bell shape, takes only two values, or has any odd
shape at all, the distribution of the total still resembles the same bell, provided
there are enough bags and the miss does not turn into a giant too often. This is
what lets gamblers, bridge builders, and weather researchers use the same formula.

The question that makes us curious: why does the original shape of the miss no
longer matter once there are many bags? And why can one bag with a giant miss ruin
the whole conclusion? In this chapter we will find out why.
\end{untukanak}

\section{The Strong Law: What It Actually Adds}
\tingkat{advanced}

In everyday language: with confidence one, the average along one long run
eventually settles, rather than merely coming close often.

The weak law speaks about each $n$ separately: at a sufficiently large $n$, the
probability of missing is small. It does not forbid the average from missing badly
at $n=10^6$, then calming down, and then missing again at $n=10^9$, provided such
events are rare at each $n$. The strong law closes that gap.

\begin{teorema}[Strong law of large numbers]
Suppose $X_1,X_2,\dots$ are independent and identically distributed with
$\mathbb{E}|X_1|<\infty$ and $\mathbb{E}[X_1]=\mu$. Then
$\bar X_n\to\mu$ almost surely.\index{strong law of large numbers}
\end{teorema}

The difference is subtle but sharp. The sequence $X_n$ of the previous section
converges in probability to zero, yet at each $\omega$ its value is one infinitely
often. The weak law cannot reject it; the strong law does, because it demands
convergence at every $\omega$ outside a set of probability zero.

Note the hypothesis as well. The strong law demands only $\mathbb{E}|X_1|<\infty$,
with no variance. The weak law we proved above demands $\sigma^2<\infty$ only
because its proof uses Chebyshev; that hypothesis is not the natural hypothesis of
the theorem. For the weak law, a necessary and sufficient condition for an iid
sequence is $n\,P(|X_1|>n)\to0$, which is weaker than $\mathbb{E}|X_1|<\infty$. We
leave those details aside, and the reader can find them in Feller, \emph{An
Introduction to Probability Theory and Its Applications}, volume 2, in the chapter
on the law of large numbers.

We do not carry out the proof of the strong law here. Its classical form truncates
the variables at the level $n$, applies Borel--Cantelli to the tails of the
truncation, and then uses the Cauchy criterion for almost sure convergence; Etemadi
gives a proof that avoids the variance assumption and needs only pairwise
independence. Where the proofs live: Blitzstein and Hwang,
\emph{Introduction to Probability}, chapter 10, section 10.2; Wasserman, \emph{All
of Statistics}, appendix to chapter 5, section 5.7.1. What must be remembered from
that statement: the strong law gives no rate. It says that the error is eventually
small and stays small, not how fast.

\begin{tangan}
Take one coin and a sheet of lined paper. Heads counts one, tails counts zero.

Toss forty times, write 1 or 0 in a row. On the second line, write the running
average under each toss: the sum so far divided by the number of tosses so far.

What we should see: the second line swings most wildly over the first five tosses,
then grows calmer. At the fortieth toss the number need not be $0.5$; $0.475$ or
$0.525$ is still normal.

Now repeat that experiment four times and draw the four paths on one graph with
different pencils. All four paths narrow toward $0.5$, and the width of their
spread is halved each time the number of tosses is multiplied by four.

Tape the sheet to the wall. If someone says the results must balance, point to the
fortieth toss, not the tenth.
\end{tangan}

\section{The Central Limit Theorem}
\tingkat{advanced}

In everyday language: the error of the average, once magnified by a factor
$\sqrt n$, always has the same bell shape, whatever the original shape of the data.

The law of large numbers tells us that $\bar X_n$ sticks to $\mu$. It does not tell
us how wide that sticking is. The central limit theorem answers that, and the answer
is surprising because it does not depend on the original distribution.

\kamus{central limit theorem}{teorema limit pusat}{the bell shape that appears in
the error of the average, not in the data, and only after that error is magnified by
the square root of the number of data points}

\begin{teorema}[Central limit, Lindeberg--Levy form]\index{central limit theorem}
Suppose $X_1,X_2,\dots$ are independent and identically distributed with
$\mathbb{E}[X_i]=\mu$ and $0<\operatorname{Var}(X_i)=\sigma^2<\infty$. Then
\[
Z_n=\frac{\sqrt n\,(\bar X_n-\mu)}{\sigma}
\;\rightsquigarrow\; Z\sim N(0,1),
\]
meaning that $P(Z_n\le z)\to\Phi(z)$ for every $z\in\mathbb{R}$.
\end{teorema}

Note where $\sqrt n$ sits. The standard deviation of $\bar X_n$ is $\sigma/\sqrt n$,
not $\sigma/n$. That is why cutting the error in half requires four times the data,
and adding one digit of precision requires a hundred times the data. This is bad
news that cannot be dodged: precision comes at a quadratic price.

Figure~\ref{fig:lonceng-dari-kotak} shows where the bell shape comes from. The left
panel is the distribution of a single die: flat, with no bell at all. The right panel
is the distribution of the average of twelve dice after standardization. The raw
material is not a bell; the result is a bell.

\begin{figure}[htbp]
\centering
\begin{tikzpicture}
\begin{axis}[
  name=satu, width=4.3cm, height=4.5cm,
  ybar, bar width=11pt,
  xmin=0.5, xmax=6.5, ymin=0, ymax=0.21,
  xtick={1,2,3,4,5,6}, ytick={0,0.0833,0.1667},
  yticklabels={$0$,$\frac{1}{12}$,$\frac16$},
  xlabel={die face}, ylabel={probability},
  title={\footnotesize one die},
  label style={font=\footnotesize}, tick label style={font=\scriptsize},
  axis lines=left,
]
\addplot[fill=layerindigo!50, draw=layerindigo!80!black, line width=0.3pt]
  coordinates{
  (1,0.1667) (2,0.1667) (3,0.1667) (4,0.1667) (5,0.1667) (6,0.1667)};
\end{axis}
\begin{axis}[
  name=duabelas, at={(satu.east)}, anchor=west, xshift=7pt,
  width=9.3cm, height=4.5cm,
  ybar, bar width=3.4pt,
  xmin=-5.3, xmax=5.3, ymin=0, ymax=0.45,
  xtick={-4,-2,0,2,4}, ytick={0,0.2,0.4},
  xlabel={score $z$ from the mean}, ylabel={density},
  title={\footnotesize twelve dice, already standardized},
  label style={font=\footnotesize}, tick label style={font=\scriptsize},
  axis lines=left,
]
\addplot[fill=accent!45, draw=accent!70!black, line width=0.2pt]
  coordinates{
  (-4.3948,0) (-4.2258,0.00001) (-4.0567,0.00003) (-3.8877,0.00009) (-3.7187,0.0002) (-3.5496,0.00044)
  (-3.3806,0.00091) (-3.2116,0.00177) (-3.0426,0.00327) (-2.8735,0.00575) (-2.7045,0.00966) (-2.5355,0.01559)
  (-2.3664,0.02418) (-2.1974,0.03614) (-2.0284,0.05215) (-1.8593,0.07276) (-1.6903,0.09826) (-1.5213,0.1286)
  (-1.3522,0.16323) (-1.1832,0.20113) (-1.0142,0.24073) (-0.8452,0.28001) (-0.6761,0.31669) (-0.5071,0.34838)
  (-0.3381,0.37287) (-0.169,0.38835) (0,0.39365) (0.169,0.38835) (0.3381,0.37287) (0.5071,0.34838)
  (0.6761,0.31669) (0.8452,0.28001) (1.0142,0.24073) (1.1832,0.20113) (1.3522,0.16323) (1.5213,0.1286)
  (1.6903,0.09826) (1.8593,0.07276) (2.0284,0.05215) (2.1974,0.03614) (2.3664,0.02418) (2.5355,0.01559)
  (2.7045,0.00966) (2.8735,0.00575) (3.0426,0.00327) (3.2116,0.00177) (3.3806,0.00091) (3.5496,0.00044)
  (3.7187,0.0002) (3.8877,0.00009) (4.0567,0.00003) (4.2258,0.00001) (4.3948,0)};
\addplot[domain=-4.6:4.6, samples=120, layerindigo, line width=0.9pt]
  {0.39894228*exp(-x^2/2)};
\end{axis}
\end{tikzpicture}
\caption{From flat to bell. Left: the distribution of one die, flat over six values.
Right: the distribution of the average of twelve dice after standardization, stacked
with the standard normal curve. The supremum distance between the two is
$0.03384$, less than half the distance at $n=2$, which is $0.08333$. The raw
material is not a bell, the result is a bell.}
\label{fig:lonceng-dari-kotak}
\end{figure}

The following proof uses moment generating functions. It is complete for the case
in which the third moment is finite, and we state that limitation honestly at the
end.

\begin{lemma}[Convergence through moment generating functions]\index{moment generating function}
Suppose $\psi_n(t)=\mathbb{E}[e^{tZ_n}]$ and $\psi(t)=\mathbb{E}[e^{tZ}]$ are finite
on some open interval around $t=0$. If $\psi_n(t)\to\psi(t)$ for every $t$ in that
interval, then $Z_n\rightsquigarrow Z$.
\end{lemma}

That lemma is the continuity theorem for moment generating functions. Its proof uses
the fact that the distribution function is determined by its moment generating
function on an interval around zero, and we do not repeat it here; see Wasserman,
\emph{All of Statistics}, chapter 5, section 5.7.2, Lemma 5.20.

\begin{proof}[Proof of the central limit theorem]
Set $Y_i=(X_i-\mu)/\sigma$, so that $\mathbb{E}[Y_i]=0$,
$\operatorname{Var}(Y_i)=1$, and $Z_n=n^{-1/2}\sum_{i=1}^n Y_i$. Let
$\psi(t)=\mathbb{E}[e^{tY_1}]$ be the moment generating function of $Y_1$, which is
finite on the interval $|t|\le\delta$. Since the variables are independent, the
moment generating function of $Z_n$ is
\[
\xi_n(t)=\Bigl[\psi\bigl(t/\sqrt n\bigr)\Bigr]^n .
\]
Derivatives at zero give $\psi(0)=1$, $\psi'(0)=\mathbb{E}[Y_1]=0$, and
$\psi''(0)=\mathbb{E}[Y_1^2]=1$. The second-order Taylor expansion with third-order
remainder, valid when $\mathbb{E}|Y_1|^3<\infty$, reads
\[
\psi(s)=1+\frac{s^2}{2}+R(s),\qquad |R(s)|\le C|s|^3 \quad (|s|\le\delta).
\]
For $n$ large enough we have $|t|/\sqrt n\le\delta$, so
\[
\xi_n(t)=\Bigl(1+\frac{t^2}{2n}+R\bigl(t/\sqrt n\bigr)\Bigr)^n
=\Bigl(1+\frac{t^2}{2n}+O\bigl(n^{-3/2}\bigr)\Bigr)^n .
\]
Take logarithms. With $\log(1+a)=a-a^2/2+\dots$ and
$a_n=t^2/(2n)+O(n^{-3/2})$,
\[
\log\xi_n(t)=n\Bigl(a_n-\frac{a_n^2}{2}+\dots\Bigr)
=\frac{t^2}{2}+O\bigl(n^{-1/2}\bigr)-\frac{t^4}{8n}+\dots
\longrightarrow \frac{t^2}{2}.
\]
Hence $\xi_n(t)\to e^{t^2/2}$ for every fixed $t$. That last function is the moment
generating function of $N(0,1)$, so the lemma above gives
$Z_n\rightsquigarrow N(0,1)$.
\end{proof}

There is one further hypothesis that slips in: the third-order Taylor remainder
demands $\mathbb{E}|Y_1|^3<\infty$. For distributions with very heavy tails that
hypothesis can fail even though the variance is finite, and the proof above has to
be replaced. The usual route is to use the characteristic function
$\varphi(t)=\mathbb{E}[e^{itX}]$, which always exists because $|e^{itX}|=1$, and then
to apply Levy's continuity theorem. The skeleton of the proof is the same: the
logarithm of the characteristic function of $Z_n$ tends to $-t^2/2$. The version
without the third-moment assumption is in Wasserman, chapter 5, section 5.7.2, and
in Blitzstein and Hwang, chapter 10, section 10.3.

This theorem does not say that the error is small. It says that the error has a bell
shape once it is divided by $\sigma/\sqrt n$. Its size is still given by
$\sigma/\sqrt n$, and $\sigma$ is a quantity we usually do not know. In practice
$\sigma$ is replaced by the sample standard deviation $s$, and that replacement does
not destroy the limit; the statement becomes $\sqrt n(\bar X_n-\mu)/s_n
\rightsquigarrow N(0,1)$. The quantity $\sigma/\sqrt n$ is called the standard error
\index{standard error} (\emph{galat baku}), and it is the basis of the confidence
intervals in Chapter 20.

How fast that approximation improves can also be bounded.

\begin{teorema}[Berry--Esseen inequality]
If $\rho=\mathbb{E}|X_1-\mu|^3<\infty$, then
\[
\sup_z\bigl|P(Z_n\le z)-\Phi(z)\bigr|\le \frac{C\rho}{\sigma^3\sqrt n},
\]
where $C$ is a universal constant; the value $C=0.4748$ suffices for all
distributions.
\end{teorema}

The shape of this inequality matters: the distance shrinks like $1/\sqrt n$, with a
constant that depends on the skewness of the original distribution through
$\rho/\sigma^3$. A strongly skewed distribution will need a far larger $n$ than a
symmetric one. The numbers in the lab section show this quantitatively for the
exponential distribution, whose third moment is $\rho=2.414553$ and whose bound is
therefore $1.1464/\sqrt n$.

\begin{resep}
Three steps to standardize a sample average. The order must not be swapped.
\begin{enumerate}
\item \textbf{Subtract the center.} Compute $\bar x-\mu$. The bell shifts so that
      its peak falls at zero.
\item \textbf{Divide by the width.} Divide the result by the standard error
      $\sigma/\sqrt n$. The shape of the bell does not change; only its scale becomes
      one unit.
\item \textbf{Read the table.} The remaining number $z$ is read off $\Phi(z)$; for
      the right tail use $1-\Phi(z)$.
\end{enumerate}
A warning about order: dividing before subtracting the center shifts the result by
$\mu\sqrt n/\sigma$. For dice with $\mu=3.5$, $\sigma=1.707825$, and $n=12$,
that shift is $7.10$; the number is already outside any normal table.
\end{resep}

\section{The Normal Approximation and the Continuity Correction}
\tingkat{intermediate}

In everyday language: to compute the probability of a sum we may use the normal
table, but when the sum can only take integer values, the threshold is shifted by
one half.

The normal approximation for sums works most simply on discrete distributions, and
there a small correction appears whose effect is large. It is called the continuity
correction\index{continuity correction}. Take $S\sim
\text{Binomial}(100;\,0.5)$, so that $\mu=50$ and $\sigma=\sqrt{100\cdot0.5
\cdot0.5}=5$. The event $\{S\le40\}$ is the same event as $\{S<40.5\}$, because
$S$ only takes integer values. Hence the correct threshold to standardize is
$40.5$, not $40$.

\begin{center}
\begin{tabular}{lrrr}
\hline
Event & Exact & No correction & With correction \\
\hline
$P(S\le40)$ & $0.028444$ & $0.022750$ & $0.028717$ \\
$P(S\le55)$ & $0.864373$ & $0.841345$ & $0.864334$ \\
$P(S=60)$   & $0.010844$ & $0.010798$ & $0.010852$ \\
\hline
\end{tabular}
\end{center}

Read that table line by line. In the first row, the approximation without the
correction misses by $0.005694$, whereas with the correction it misses by
$0.000273$, almost twenty-one times closer. In the third row both are already
close, because near the center of the distribution a width of one unit is already
small compared with $\sigma=5$. The general lesson: the continuity correction helps
most in the tails and when $\sigma$ is small. The size of the shift in $z$ that it
produces is $1/(2\sigma)$, namely $0.1$ in this example.

\kertas{Recompute the first row of this table by hand. $z=(40-50)/5=-2$,
$\Phi(-2)=0.022750$; then $z=(40.5-50)/5=-1.9$, $\Phi(-1.9)=0.028717$.
The exact value is $\sum_{k=0}^{40}\binom{100}{k}2^{-100}=0.028444$. Check also
that the error without the correction $0.005694$ is about $21$ times the error
with the correction $0.000273$.}

This correction is not a new formula, only a rereading of the event. For a discrete
variable with spacing $h$ between values, the threshold is shifted by $h/2$ inward.
For a continuous variable nothing needs shifting, because there are no jumps.

One bad habit deserves mention here. Many books write that the normal approximation
"may be used when $n\ge30$". That is not a theorem but a habit. The numbers in the
lab section show why: for the exponential average, the supremum distance to the
normal at $n=30$ is $0.02429$, and the $95$ percent quantile of $Z_{30}$ is
$1.7419$ whereas the normal gives $1.6449$. That difference of $0.097$ means a
confidence interval that is too narrow by about six percent. What decides is not $n$
alone, but the skewness and kurtosis of the original distribution.

\section{When the Central Limit Theorem Misleads: Heavy Tails}
\tingkat{advanced}

In everyday language: when one observation can occasionally be far larger than all
the others, the average no longer settles, and the bell does not appear.

\kamus{heavy tail}{ekor berat}{a probability of huge values that does not shrink
fast enough, so that one observation can outweigh all the rest}

All the theorems above use one hypothesis that is easy to pass over: finite variance.
The Cauchy distribution\index{Cauchy distribution}, with density
\[
f(x)=\frac{1}{\pi(1+x^2)},\qquad x\in\mathbb{R},
\]
does not have it. Its tail shrinks like $1/x$, which is too slow. The probability of
exceeding ten is $1-(2/\pi)\arctan 10=0.063451$, and the probability of exceeding
one hundred is still $0.006366$. Compare the standard normal, whose probability of
exceeding ten is practically zero. Even $\mathbb{E}|X|$ is not finite: the
expectation truncated at level $M$ equals $(2/\pi)\ln(1+M^2)$, namely $2.9381$ for
$M=10$, $8.7952$ for $M=10^3$, and $17.5905$ for $M=10^6$. That number grows
like a logarithm, so there is no limit.

The consequence is blunt, and not merely "the CLT is less accurate".

\begin{proposisi}
If $X_1,\dots,X_n$ are iid Cauchy, then $\bar X_n$ has a Cauchy distribution for
every $n$.
\end{proposisi}

\begin{proof}[Proof]
The characteristic function of the Cauchy distribution is $\varphi(t)=e^{-|t|}$.
Since the variables are independent, the characteristic function of $\bar X_n$ is
$\bigl[e^{-|t|/n}\bigr]^n=e^{-|t|}$, which is the Cauchy characteristic function
again. Since the characteristic function determines the distribution, $\bar X_n$ has
a Cauchy distribution.
\end{proof}

So the average of $10^5$ Cauchy observations is no more informative than one
observation. The law of large numbers fails, the central limit theorem fails, and
"standard error" means nothing. The numbers in the lab section show all three
symptoms at once. The median of $|\bar X_n|$ stays around one for all $n$
($1.004$, $1.008$, $1.040$, $0.973$), whereas for a distribution with finite
variance it must shrink like $1/\sqrt n$. The probability $|\bar X_n|>10$ stays
around $0.0635$, its theoretical value, for all $n$. The empirical standard
deviation jumps about: $33.04$, $863.61$, $23.87$, $379.53$ for $n=10$ up to
$10^4$. Those jumps are not a bug in the program. The sample standard deviation of
Cauchy data is itself a random variable with no typical value.

Figure~\ref{fig:cauchy-berjalan} shows the difference on one run. The running
average of uniform data sticks to $0.5$ after a few dozen draws. The running
average of Cauchy data jumps from $-0.0866$ to $4.4572$ at the twentieth draw,
because one value $90.79$ happened to arrive, and then wanders again. One giant
value is enough to cancel the entire history before it.

\begin{figure}[htbp]
\centering
\begin{tikzpicture}
\begin{axis}[
  width=0.95\linewidth, height=6.0cm,
  xmode=log, xmin=1, xmax=4000, ymin=-1.1, ymax=4.9,
  xlabel={number of draws $n$ (logarithmic scale)}, ylabel={running average},
  label style={font=\small}, tick label style={font=\footnotesize},
  axis lines=left,
  legend style={font=\footnotesize, draw=layergrey!35, fill=white,
    fill opacity=0.9, text opacity=1, at={(0.015,0.985)},
    anchor=north west, cells={anchor=west}},
  legend cell align=left,
]
\addplot[accent, line width=0.7pt] coordinates{
  (1,-0.253) (2,1.045) (3,0.582) (4,0.564) (5,0.553) (6,1.064)
  (7,-0.579) (8,-0.517) (9,-0.409) (10,-0.325) (11,-0.409) (12,-0.361)
  (13,-0.344) (14,-0.262) (15,0.008) (16,-0.038) (17,0.017) (18,-0.08)
  (19,-0.087) (20,4.457) (21,4.29) (22,4.153) (23,4.046) (24,3.762)
  (25,3.656) (26,3.551) (27,3.425) (28,3.257) (29,3.186) (30,3.083)
  (31,2.934) (32,2.786) (33,2.837) (34,2.743) (35,2.633) (36,2.528)
  (37,2.464) (38,2.381) (39,2.279) (40,2.218) (41,2.157) (44,1.732)
  (47,1.387) (50,1.367) (53,1.191) (57,1.113) (61,1.077) (65,0.942)
  (70,0.796) (75,0.773) (80,0.725) (85,0.492) (91,0.79) (97,0.655)
  (104,0.54) (111,0.478) (119,0.494) (127,0.401) (135,0.981) (145,0.868)
  (155,2.5) (165,2.351) (177,2.207) (189,2.168) (202,1.935) (216,1.854)
  (230,1.776) (246,1.665) (263,1.573) (281,0.848) (300,0.794) (321,0.791)
  (343,0.511) (367,0.922) (392,0.986) (419,0.984) (447,1.03) (478,0.951)
  (511,0.827) (546,0.685) (583,0.37) (623,0.548) (666,0.953) (712,1.021)
  (761,0.945) (813,2.066) (869,2.135) (929,1.855) (992,1.814) (1060,1.649)
  (1133,1.31) (1211,0.557) (1294,0.422) (1383,0.434) (1478,0.466) (1579,0.324)
  (1688,0.406) (1803,0.386) (1927,0.444) (2059,0.403) (2201,0.585) (2352,0.419)
  (2513,0.528) (2686,0.397) (2870,0.356) (3067,0.326) (3278,0.276) (3503,-0.096)
  (3743,-0.159) (4000,-0.203)};
\addlegendentry{Cauchy}
\addplot[layerteal, dashed, line width=1.0pt] coordinates{
  (1,0.221) (2,0.272) (3,0.323) (4,0.472) (5,0.462) (6,0.519)
  (7,0.524) (8,0.573) (9,0.55) (10,0.574) (11,0.605) (12,0.601)
  (13,0.57) (14,0.581) (15,0.543) (16,0.532) (17,0.529) (18,0.528)
  (19,0.532) (20,0.544) (21,0.545) (22,0.547) (23,0.539) (24,0.531)
  (25,0.54) (26,0.555) (27,0.561) (28,0.541) (29,0.545) (30,0.547)
  (31,0.562) (32,0.549) (33,0.533) (34,0.542) (35,0.531) (36,0.535)
  (37,0.537) (38,0.526) (39,0.519) (40,0.516) (41,0.517) (44,0.518)
  (47,0.519) (50,0.503) (53,0.518) (57,0.515) (61,0.525) (65,0.53)
  (70,0.529) (75,0.53) (80,0.535) (85,0.538) (91,0.539) (97,0.538)
  (104,0.541) (111,0.535) (119,0.532) (127,0.533) (135,0.528) (145,0.526)
  (155,0.518) (165,0.523) (177,0.513) (189,0.516) (202,0.513) (216,0.507)
  (230,0.5) (246,0.498) (263,0.495) (281,0.496) (300,0.496) (321,0.495)
  (343,0.499) (367,0.499) (392,0.501) (419,0.503) (447,0.497) (478,0.505)
  (511,0.508) (546,0.509) (583,0.504) (623,0.502) (666,0.497) (712,0.494)
  (761,0.488) (813,0.484) (869,0.485) (929,0.484) (992,0.485) (1060,0.484)
  (1133,0.483) (1211,0.485) (1294,0.487) (1383,0.489) (1478,0.488) (1579,0.491)
  (1688,0.489) (1803,0.489) (1927,0.49) (2059,0.491) (2201,0.49) (2352,0.492)
  (2513,0.491) (2686,0.493) (2870,0.492) (3067,0.49) (3278,0.492) (3503,0.491)
  (3743,0.493) (4000,0.493)};
\addlegendentry{uniform $(0,1)$}
\draw[layergrey, densely dotted, line width=0.6pt]
  (axis cs:1,0.5) -- (axis cs:4000,0.5);
\node[font=\scriptsize, layergrey, anchor=south west] at (axis cs:1.3,0.53) {$0.5$};
\end{axis}
\end{tikzpicture}
\caption{Running averages over 4000 draws. Dashed line: uniform on $(0,1)$, settling
near $0.5$ and ending at $0.4930$. Solid line: Cauchy, jumping from $-0.0866$
to $4.4572$ at the twentieth draw when one value $90.79$ arrived, then wandering
again and ending at $-0.203$. The horizontal axis is logarithmic in width, so that
the early jump is not hidden.}
\label{fig:cauchy-berjalan}
\end{figure}

What replaces the central limit theorem is the family of stable
distributions\index{stable distribution}. For iid variables with tails that shrink
like $x^{-\alpha}$ with $0<\alpha<2$, there are constants $a_n$ such that
\[
\frac{S_n-a_n}{n^{1/\alpha}}\;\rightsquigarrow\;Y_\alpha,
\]
with $Y_\alpha$ the stable distribution of index $\alpha$. For $\alpha=1$ we get
Cauchy and $n^{1/1}=n$, exactly as in the proposition above. For
$\alpha\in(1,2)$ the average still converges to $\mu$, but the error shrinks like
$n^{1-1/\alpha}$, slower than $1/\sqrt n$, and its shape is not a bell. This general
statement is the result of Gnedenko and Kolmogorov; the details are in Feller volume
2, in the chapter on stable distributions.

If $\alpha>2$ but the third moment does not exist, the central limit theorem still
holds because the variance is finite. What is lost is the guarantee of a rate: the
Berry--Esseen bound becomes infinite, so there is no longer any handle on the claim
that the error shrinks like $1/\sqrt n$ with a controlled constant. In practice this
is what bites most often: the distribution is not extreme like Cauchy, but it has
enough tail to make confidence intervals too narrow at realistic sample sizes.

\begin{tantangan}
A distribution with no mean. Take $X$ with
$P(X=2^k)=P(X=-2^k)=2^{-k-1}$ for $k=1,2,3,\dots$

Show two things. First, $\mathbb{E}|X|=\infty$, so there is no $\mu$ that may be
approached. Second, the sample average does not converge in probability to any
constant: the law of large numbers fails, it does not merely slow down.

Hint 1: what decides is not the average of all the terms, but the largest term. Ask
how often one term exceeds the threshold $n$.

Hint 2: compute $n\,P(|X|>n)$ for $n$ between two powers of two, then use the
necessary condition $n\,P(|X_1|>n)\to0$ stated in the section on the strong law.

If you get stuck, go back to the section on heavy tails and to the paragraph on the
necessary and sufficient condition in the section on the strong law.
\end{tantangan}

\begin{untukanak}
In a village there are a hundred people with ordinary incomes, and one person who is
extremely rich. If we compute the average income of that village, the number is
decided almost entirely by that rich person. If the rich person happens to be away,
the average drops drastically. If that person comes back, the average jumps up.

Now add inhabitants. Another hundred people, another thousand; as long as the
newcomers are ordinary people, the average changes little. But if among the newcomers
there is one more giant, the average jumps again. Adding many people does not help,
because what decides is not the number of people but who is the largest.

This is what happens with data whose tails are heavy: one observation can outweigh
all the others. The average becomes unstable, and the bell that is usually used to
compute "how confident" does not appear at all.

The question that makes us curious: how can we know, before computing, whether our
data are of the calm kind or of the kind ruled by giants? In this chapter we will
find out why that question cannot be answered just by looking at the average.
\end{untukanak}

\section{Computational Lab: Five Checks}
\tingkat{advanced}

The five claims of this chapter can be checked with one script: the $1/\sqrt n$
rate in the law of large numbers, the skewness and kurtosis of the exponential
average, the supremum distance to the normal, the Monte Carlo standard error, and
total failure for Cauchy. All the numbers below are real output from an execution
with NumPy 2.5.2 and SciPy 1.18.0, generator \texttt{default\_rng(20261024)}.

\begin{lab}
\begin{verbatim}
import numpy as np
from scipy import stats

rng = np.random.default_rng(20261024)

print("== 1. law of large numbers: dice average ==")
for k in range(1, 6):
    n = 10**k
    m = rng.integers(1, 7, size=(2000, n)).mean(axis=1)
    print("n=%6d  mean=%.4f  rms=%.4f  sd/sqrt n=%.4f"
          % (n, m.mean(), np.sqrt(np.mean((m-3.5)**2)),
             np.sqrt(35/12/n)))

print("== 2. CLT: exponential average ==")
for n in (2, 5, 10, 30, 100, 500):
    z = (rng.exponential(size=(50000, n)).mean(axis=1)
         - 1)*np.sqrt(n)
    ks = stats.kstest(z, "norm")
    print("n=%4d skew=%+.3f (%+.3f) kurt=%+.3f (%+.3f)"
          % (n, stats.skew(z), 2/np.sqrt(n),
             stats.kurtosis(z), 6.0/n))
    print("      KS D=%.4f p=%.3f" % (ks.statistic, ks.pvalue))

print("== 3. sup distance to normal, exact via gamma ==")
rho = 2 + 2*(6/np.e - 2)          # E|X-1|^3 for X ~ Exp(1)
for n in (2, 5, 10, 30, 100, 500):
    zg = np.linspace(-4, 4, 40001)
    Fn = stats.gamma.cdf(n*(1+zg/np.sqrt(n)), a=n)
    D = np.max(np.abs(Fn - stats.norm.cdf(zg)))
    print("n=%4d sup|Fn-Phi|=%.5f  D*sqrt n=%.3f  BE=%.3f"
          % (n, D, D*np.sqrt(n), 0.4748*rho/np.sqrt(n)))

print("== 4. Monte Carlo standard error for E[exp(U)] ==")
for k in (2, 3, 4, 5):
    n = 10**k
    est = np.exp(rng.random(size=(2000, n))).mean(axis=1)
    print("n=%6d  sd=%.5f  sd*sqrt n=%.5f"
          % (n, est.std(), est.std()*np.sqrt(n)))

print("== 5. Cauchy: the average does not concentrate ==")
for k in (1, 2, 3, 4):
    n = 10**k
    m = rng.standard_cauchy(size=(4000, n)).mean(axis=1)
    print("n=%5d  median|m|=%.3f P(|m|>10)=%.4f sd=%.2f"
          % (n, np.median(np.abs(m)), np.mean(np.abs(m) > 10),
             m.std()))
print("theory P(|mean|>10) = %.4f"
      % (1 - 2/np.pi*np.arctan(10)))
\end{verbatim}
\end{lab}

Its output, verbatim:

\begin{verbatim}
== 1. law of large numbers: dice average ==
n=    10  mean=3.5014  rms=0.5440  sd/sqrt n=0.5401
n=   100  mean=3.4992  rms=0.1698  sd/sqrt n=0.1708
n=  1000  mean=3.4986  rms=0.0527  sd/sqrt n=0.0540
n= 10000  mean=3.5001  rms=0.0171  sd/sqrt n=0.0171
n=100000  mean=3.4997  rms=0.0053  sd/sqrt n=0.0054
== 2. CLT: exponential average ==
n=   2 skew=+1.440 (+1.414) kurt=+3.130 (+3.000)
      KS D=0.0929 p=0.000
n=   5 skew=+0.919 (+0.894) kurt=+1.393 (+1.200)
      KS D=0.0602 p=0.000
n=  10 skew=+0.631 (+0.632) kurt=+0.589 (+0.600)
      KS D=0.0480 p=0.000
n=  30 skew=+0.355 (+0.365) kurt=+0.181 (+0.200)
      KS D=0.0219 p=0.000
n= 100 skew=+0.196 (+0.200) kurt=+0.038 (+0.060)
      KS D=0.0167 p=0.000
n= 500 skew=+0.112 (+0.089) kurt=-0.003 (+0.012)
      KS D=0.0062 p=0.043
== 3. sup distance to normal, exact via gamma ==
n=   2 sup|Fn-Phi|=0.09446  D*sqrt n=0.134  BE=0.811
n=   5 sup|Fn-Phi|=0.05963  D*sqrt n=0.133  BE=0.513
n=  10 sup|Fn-Phi|=0.04211  D*sqrt n=0.133  BE=0.363
n=  30 sup|Fn-Phi|=0.02429  D*sqrt n=0.133  BE=0.209
n= 100 sup|Fn-Phi|=0.01330  D*sqrt n=0.133  BE=0.115
n= 500 sup|Fn-Phi|=0.00595  D*sqrt n=0.133  BE=0.051
== 4. Monte Carlo standard error for E[exp(U)] ==
n=   100  sd=0.04928  sd*sqrt n=0.49281
n=  1000  sd=0.01586  sd*sqrt n=0.50156
n= 10000  sd=0.00495  sd*sqrt n=0.49538
n=100000  sd=0.00155  sd*sqrt n=0.49056
== 5. Cauchy: the average does not concentrate ==
n=   10  median|m|=1.004 P(|m|>10)=0.0630 sd=33.04
n=  100  median|m|=1.008 P(|m|>10)=0.0680 sd=863.61
n= 1000  median|m|=1.040 P(|m|>10)=0.0653 sd=23.87
n=10000  median|m|=0.973 P(|m|>10)=0.0580 sd=379.53
theory P(|mean|>10) = 0.0635
\end{verbatim}

Part 1 shows the law of large numbers at work, and shows its price. The rms column
matches $\sigma/\sqrt n$ to two decimal places for $n\ge100$: $0.1698$ against
$0.1708$, $0.0171$ against $0.0171$, $0.0053$ against $0.0054$. But look
at the magnitudes: at $n=10^5$ the average still misses by about $0.005$, which is
$0.15$ percent of $3.5$. The average of $10^5$ tosses is not "practically
exact"; it is exact to two decimal places.

Parts 2 and 3 show two ways of measuring closeness to the normal. Skewness and
kurtosis follow the theoretical prediction neatly from $n=10$ on: $0.631$ against
$2/\sqrt{10}=0.632$ and $0.589$ against $6/10=0.600$. At $n=2$ and $n=5$ the
kurtosis estimator is still noisy, because its deviation depends on the eighth moment
of that distribution, not only on the number of repetitions. Part 3 computes the
supremum distance exactly, without simulation, by using the Gamma distribution for
the exponential average: $0.04211$ at $n=10$ and $0.00595$ at $n=500$. The column
$D\sqrt n$ is nearly constant at $0.133$, direct evidence that the rate is
$1/\sqrt n$. The Berry--Esseen bound $1.1464/\sqrt n$ holds, with the numbers
$0.209$ against $0.02429$ at $n=30$; that bound is about $8.6$ times looser
than reality, which is normal for a bound that is uniform over all distributions.

The Kolmogorov--Smirnov test on fifty thousand draws rejects normality with
$p=0.000$ for all $n$ up to $100$, then gives $p=0.043$ at $n=500$. That is not a
sign that the program is wrong: the normal approximation is never exact for finite
$n$, and the distance merely shrinks until a test that large can no longer see it.
The right question is not "is it normal" but "how large an error can we still
accept", and part 3 answers it with numbers.

Part 4 separates two quantities that are often confused. The standard deviation of
the estimator falls from $0.04928$ to $0.00155$ when $n$ grows a thousandfold,
whereas $\text{sd}\cdot\sqrt n$ is nearly constant at $0.49$. That constant value
is the $\sigma$ of $\exp(U)$, namely $\sqrt{(e^2-1)/2-(e-1)^2}=0.491971$. This is
the $1/\sqrt n$ law in its barest form: for a precision of $0.001$ at $95$ percent
confidence we need $(1.96\cdot0.491971/0.001)^2=929803$ draws.

Part 5 is the counterweight. The median of $|\bar X_n|$ does not shrink; the
probability of exceeding ten does not shrink; and the empirical standard deviation
jumps from $33.04$ to $863.61$ and then down to $23.87$. Theory says
$0.0635$, and all five simulation numbers hover around it. There is no law of large
numbers here, and no central limit theorem. The only way to rescue the analysis is to
change the question: estimate the tail index, not the mean.

\begin{untukdosen}
The teaching order I suggest: start from the simulation table, not from the
definitions. Show the class that the error of one run follows no pattern at all,
whereas the distribution of the error does. After that, the three modes of
convergence, with a counterexample for each direction that fails; do not postpone
those examples, because that is where students understand that "approaching" needs to
be made precise. Only then Markov, Chebyshev, the weak law, the strong law as a
statement, and then the central limit theorem with the moment generating function
proof.

Two misconceptions appear most often. First, students read "converges in
probability" as "the value becomes fixed", and are then surprised when a simulation
shows big jumps at large $n$; the split $(0,1)$ example in this chapter is the
antidote. Second, they read the central limit theorem as "the average becomes
normal", whereas what becomes normal is the error divided by $\sigma/\sqrt n$; as a
result they write $\bar X_n\approx N(\mu,\sigma^2)$ and forget $\sqrt n$.

A blackboard activity that works: ask two students to play the gambler and the
bookmaker. The gambler names a tolerance, the bookmaker names the number of tosses
that guarantees it, and the class computes with Chebyshev and then compares against
the exact numbers from SciPy. The gap is always large, and that is the doorway into
the normal approximation. Sign that the class is ready to move on: they can explain
why Chebyshev is still useful although it is loose, and they can name one
distribution that makes the central limit theorem fail.
\end{untukdosen}

\section{Summary}

\begin{ringkasan}
\begin{itemize}
  \item There are three different modes of convergence: almost sure, in probability,
        and in distribution. Almost sure implies in probability, in probability
        implies in distribution, and the reverse directions fail unless the limit is
        a constant.
  \item Markov's and Chebyshev's inequalities use only moments, which is why they
        hold for every distribution with finite variance, and why they are loose.
  \item Weak law of large numbers: $P(|\bar X_n-\mu|\ge\varepsilon)\le
        \sigma^2/(n\varepsilon^2)\to0$. The strong law adds that
        $\bar X_n\to\mu$ almost surely when $\mathbb{E}|X_1|<\infty$, without
        demanding a variance.
  \item Central limit theorem: $\sqrt n(\bar X_n-\mu)/\sigma\rightsquigarrow N(0,1)$
        when the variance is finite. The rate is $1/\sqrt n$, hence four times the
        data for half the error, and a hundred times the data for one extra digit.
  \item The normal approximation for discrete sums requires a continuity correction
        of half a unit; in the binomial example of this chapter that correction cuts
        the error in the tail by about a factor of twenty.
  \item If the variance is infinite, the central limit theorem fails, it does not
        merely slow down. The average of $n$ Cauchy observations is still Cauchy, and
        what replaces the bell is a stable distribution of index $\alpha<2$.
\end{itemize}
\end{ringkasan}

\section{Exercises}

\begin{latihan}
\begin{enumerate}
  \item Prove: if $X_n\to X$ and $Y_n\to Y$ in probability, then
        $X_n+Y_n\to X+Y$ in probability. Write down the set inequality you use, then
        explain why $\varepsilon/2$ appears.
  \item Prove that Chebyshev's inequality is sharp: for $0<\sigma^2\le a^2$, take $X$
        with $P(X=a)=P(X=-a)=\sigma^2/(2a^2)$ and
        $P(X=0)=1-\sigma^2/a^2$, then compute $\operatorname{Var}(X)$ and
        $P(|X|\ge a)$.
  \item Evaluate this claim: "Since the central limit theorem holds for all
        distributions, the average of thirty observations may always be treated as
        normal." True, false, or not enough evidence? Name the hypothesis that is
        violated or that has not been checked, and use numbers from the lab section.
  \item For a coin with $p=1/2$ and tolerance $\varepsilon=0.05$, find the smallest
        $n$ that Chebyshev guarantees for a probability of missing of at most
        $0.05$. Then find the smallest $n$ that really satisfies that condition, by
        computing the binomial distribution in Python. Compare the two numbers.
  \item Prove that convergence in probability does not guarantee convergence of
        expectations, using $P(X_n=n^2)=1/n$ and $P(X_n=0)=1-1/n$. Compute
        $\mathbb{E}[X_n]$ and $P(|X_n|>\varepsilon)$ for fixed $\varepsilon$.
  \item For $S\sim\text{Binomial}(200;\,0.3)$, approximate $P(S\le50)$ with the
        normal without and with the continuity correction, then compare both with the
        exact value from SciPy. Report both errors.
  \item Generate $4000$ averages of $n$ Cauchy numbers for
        $n=10,10^2,10^3,10^4$, then report the median of $|\bar X_n|$, the
        probability $|\bar X_n|>10$, and the standard deviation. Compare with
        $1-(2/\pi)\arctan10$ and explain why the standard deviation column does not
        shrink.
  \item Prove with characteristic functions that the average of $n$ iid Cauchy
        variables has a Cauchy distribution. Then explain why a report of the form
        "average $\pm$ two standard errors" on such data cannot be interpreted.
  \item How many Monte Carlo draws are needed for the estimator
        $\mathbb{E}[\exp(U)]$ to have a $95$ percent interval half-width of
        $0.0005$? Use $\sigma=0.491971$ from the lab section, then compare with
        the $n$ needed for a half-width of $0.001$. State the ratio and explain
        where it comes from.
  \item Prove that if $X_n\rightsquigarrow c$ with $c$ a constant, then
        $X_n\to c$ in probability, without using part (c) of the theorem on the
        relations among the three modes at the start of the chapter, then show by
        example that the hypothesis "constant" cannot be dropped.
\end{enumerate}

Hints for the three hardest items. Item 1: write
$\{|(X_n+Y_n)-(X+Y)|>\varepsilon\}$ as contained in the union of two events, then use
the union bound. Item 8: the characteristic function of a sum is the product of the
characteristic functions, and the Cauchy characteristic function $e^{-|t|}$ is
special because its form does not change when it is raised to a power. Item 10: write
$P(X_n\le t)\le P(X_n\le c-\varepsilon)+P(|X_n-c|>\varepsilon)$ for $t$ to the left
of $c$, and remember that a point-mass distribution jumps exactly at $c$.
\end{latihan}

\begin{untukpeneliti}
What is really proved in this chapter: the implications almost sure $\Rightarrow$ in
probability $\Rightarrow$ in distribution, the weak law via Chebyshev, and the
central limit theorem via moment generating functions under the assumption
$\mathbb{E}|Y_1|^3<\infty$. What is popular but not proved here: the strong law (the
statement only, Etemadi's proof is skipped), the version of the central limit theorem
without a third moment (it needs characteristic functions and Levy's continuity
theorem), and the whole theory of stable distributions (the Gnedenko--Kolmogorov
statement, without proof).

Primary sources I used: Wasserman, \emph{All of Statistics}, chapter 5, sections 5.2
(types of convergence and the implication diagram), 5.3 (the weak law), 5.4 (the
central limit theorem and Berry--Esseen), plus appendix 5.7.1 and 5.7.2 for the strong
law and the proof of the central limit theorem. Blitzstein and Hwang,
\emph{Introduction to Probability}, chapter 10, sections 10.1 (inequalities), 10.2
(the law of large numbers), 10.3 (the central limit theorem). For the heavy-tail
side: Feller volume 2.

Limits of validity that must be mentioned. First, all the theorems here assume
independence and identical distribution; for dependent data, for instance time series
or iterates of a Markov chain, what holds are martingale or mixing versions, and the
rate can be far slower. Second, the optimal constant in Berry--Esseen is not known;
$0.4748$ suffices, and the best lower bound for that constant is still below
$0.5$. Third, the rate $1/\sqrt n$ is the worst rate that can be guaranteed
uniformly; for particular distributions the true rate is better, and for heavy tails
with $\alpha\in(1,2)$ the true rate is worse. Open questions that are still alive:
how sharp the Berry--Esseen bound is for classes of distributions with bounded
moments, and how to estimate the tail index with valid confidence intervals at small
$n$.
\end{untukpeneliti}

\peta{Chapter 19 in Wasserman, \emph{All of Statistics}: chapter 5 (Convergence of
Random Variables), sections 5.2 Types of Convergence, 5.3 The Law of Large Numbers,
5.4 The Central Limit Theorem, and appendix 5.7.1 Almost Sure and L1 Convergence as
well as 5.7.2 Proof of the Central Limit Theorem. In Blitzstein and Hwang,
\emph{Introduction to Probability}: chapter 10 (Inequalities and limit theorems),
sections 10.1 Inequalities, 10.2 Law of large numbers, 10.3 Central limit theorem.
For stable distributions: Feller, \emph{An Introduction to Probability Theory and Its
Applications}, volume 2. The basic distributions used here are built in Chapter 17,
and the confidence intervals that rest on this chapter are discussed in Chapter 19.}

% Judul penuh meluber 26,51 pt di kepala halaman; hanya mark header yang
% dipendekkan (sekali pakai, lihat \judulpendek di preamble-en.tex).
\judulpendek{Estimation, Confidence Intervals, and Testing}
\chapter{Estimation, Confidence Intervals, and Hypothesis Testing}

A team tests an attack detector on $50$ labeled examples, and the detector catches $33$ of
them. The report written that night reads: recall $66$ percent. A second detector, tested on
$100$ examples, catches $53$: recall $53$ percent. The morning meeting decides to use the
first, because $66$ is larger than $53$, and that decision is finished in one minute.

The arithmetic is correct: $33/50=0.66$. What is wrong is the assumption that this number
points at one number. If those fifty examples are replaced by fifty other examples from the
same collection, the result is no longer $33$; it wobbles, and the wobble is not small. After
two hours of writing the program, the team is in fact holding two sentences of very different
strength. The first sentence, ``on this data we see $66$ percent'', says almost nothing. The
second, ``the true catch rate lies somewhere inside this interval'', says something usable,
provided we know how wide that interval is and what that width means.

Chapter 19 already gave two tools: the law of large numbers, which says the sample mean
concentrates at the true value, and the central limit theorem, which says the error spreads
normally with standard deviation $\sigma/\sqrt n$. Both are statements about many trials at
once, and neither answers the question of the morning meeting: from one trial already
finished, what sentence may we utter about a parameter we have never seen, and how often is
that sentence wrong if the procedure is repeated. This chapter builds the answer, from
estimators to p-values.

\begin{untukanak}
A farmer wants to know how many weeds there are in his field. The whole field is too large to
inspect, so he takes twenty small plots spread around, then counts the weeds in each plot.
From those twenty numbers he guesses the state of the whole field. The trouble is that a guess
is always a single number, whereas the state of a field is never that simple.

Now the part that makes one curious. In those twenty plots he finds no weeds at all. May he go
home and say that his field is clean? Certainly not. It may be that there are really no weeds,
and it may be that he merely happened to choose plots made of bare soil. All he can say is:
if the weeds were plentiful, twenty empty plots would be almost impossible. From that he can
compute an honest upper bound, for instance ``at most three in a hundred plots''.

The same holds when a teacher grades a few students' quizzes and when a computer program is
tested on a few examples.
In this chapter we shall learn why one number is never enough, why the honest answer is always
an interval, and why the interval widens when the evidence is thin.
\end{untukanak}

\section{Estimators, bias, and standard error}
\tingkat{basic}

Every statistical conclusion begins with one small decision: which quantity we compute from
the data. The mean, a proportion, the sample standard deviation, all are functions of the data
and require no knowledge of the parameter being sought. In everyday language: an estimator is
a surveyor who guesses the height of a tree from the length of its shadow, and we need to know
whether that surveyor is systematically off or merely unsteady.

\kamus{estimator}{penaksir}{a surveyor who guesses the height of a tree from the length of
its shadow, rather than from the tree itself}

\begin{definisi}[Statistic and estimator]
\index{estimator}
Let $X_1,\dots,X_n$ be a random sample from a distribution that depends on a parameter
$\theta$. A \emph{statistic} is a function of $X_1,\dots,X_n$ that does not contain
$\theta$. A statistic used to guess $\theta$ is called an \emph{estimator} and written
$\hat\theta$; the value obtained on one sample is called an \emph{estimate}.
\end{definisi}

An estimator is a random variable, because it is a function of random variables. Hence it has
an expectation and a variance, and those two quantities determine its quality. The expectation
answers the question ``does the guess keep missing in one direction'', while the variance
answers ``how far does the guess wobble from sample to sample''.

\begin{definisi}[Bias and unbiased estimator]
\index{unbiased}
The bias of the estimator $\hat\theta$ for $\theta$ is $\operatorname{bias}(\hat\theta)
=\mathbb{E}[\hat\theta]-\theta$. An estimator is called \emph{unbiased} if its bias is zero,
namely $\mathbb{E}[\hat\theta]=\theta$ for every possible value of $\theta$.
\end{definisi}

The sample mean is unbiased for $\mu$, and that was already used in the proof of the law of
large numbers in Chapter 19, so we do not repeat it here. More interesting is the estimator
used for the variance. With divisor $n$, the estimator
$S_n^2=\frac1n\sum(X_i-\bar X_n)^2$ is biased downward; with divisor $n-1$, the estimator
$S^2=\frac1{n-1}\sum(X_i-\bar X_n)^2$ is unbiased. The reason is that $\sum(X_i-\bar X_n)^2$
loses one degree of freedom because $\bar X_n$ is itself computed from the same data, so its
expectation is $(n-1)\sigma^2$ and not $n\sigma^2$. The divisor $n-1$ is a correction for that
loss, not mere tidiness.

The quality of an estimator is usually summarized by one number, the mean squared error. In
everyday language: how large an error we expect, computed in a way that punishes large errors
more heavily.

\begin{definisi}[Mean squared error]
\index{mean squared error}
$\operatorname{MSE}(\hat\theta)=\mathbb{E}[(\hat\theta-\theta)^2]$.
\end{definisi}

\begin{teorema}[MSE decomposition]
For an estimator $\hat\theta$ with finite variance,
\[
\operatorname{MSE}(\hat\theta)=\operatorname{bias}(\hat\theta)^2
+\operatorname{Var}(\hat\theta).
\]
\end{teorema}

\begin{proof}[Proof]
Write $\hat\theta-\theta=(\hat\theta-\mathbb{E}\hat\theta)
+(\mathbb{E}\hat\theta-\theta)$. Square it and take the expectation:
\[
\mathbb{E}[(\hat\theta-\theta)^2]
=\mathbb{E}[(\hat\theta-\mathbb{E}\hat\theta)^2]
+2(\mathbb{E}\hat\theta-\theta)\,\mathbb{E}[\hat\theta-\mathbb{E}\hat\theta]
+(\mathbb{E}\hat\theta-\theta)^2 .
\]
The middle term vanishes because $\mathbb{E}[\hat\theta-\mathbb{E}\hat\theta]=0$, while the
first term is the variance and the third is the bias squared.
\end{proof}

That identity explains why ``unbiased'' is not a championship title. An unbiased estimator may
have a large variance, and a slightly biased estimator may have a far smaller variance so that
its MSE is lower. To estimate $\mu$ from $n$ observations, the statistic $X_1$ is unbiased with
variance $\sigma^2$, whereas $\bar X_n$ is unbiased with variance $\sigma^2/n$; at $n=20$ the
difference is a factor of twenty.

The second example shows a bias that must be computed, because it is used every day for
proportions. If $K\sim\mathrm{Bin}(n,p)$ and $\hat p=K/n$, then
$\hat p(1-\hat p)$ is the natural estimator for $p(1-p)$, that is the variance of one trial.
Its expectation is
\[
\mathbb{E}[\hat p(1-\hat p)]
=\mathbb{E}[\hat p]-\mathbb{E}[\hat p^2]
=p-\Bigl(p^2+\frac{p(1-p)}{n}\Bigr)
=p(1-p)\Bigl(1-\frac1n\Bigr),
\]
so its bias is $-p(1-p)/n$. For $p=0.3$ and $n=20$ the value is $p(1-p)=0.21$
and the bias is $-0.0105$, that is five percent downward. At $n=200$ the bias is only
$-0.00105$. A bias that shrinks like $1/n$ is typical: it disappears at infinity, but it is
present in every finite sample.

The second property demanded of an estimator is that it actually heads for the target as the
data grow. In everyday language: with enough evidence, a wrong guess becomes increasingly
impossible, not merely increasingly rare.

\begin{definisi}[Consistent]
\index{consistent}
A sequence of estimators $\hat\theta_n$ is \emph{consistent} for $\theta$ if for every
$\varepsilon>0$ we have $P(|\hat\theta_n-\theta|\ge\varepsilon)\to0$ as
$n\to\infty$.
\end{definisi}

\begin{proposisi}[MSE going to zero suffices for consistency]
\index{consistent}
If $\operatorname{MSE}(\hat\theta_n)\to0$, then $\hat\theta_n$ is consistent for
$\theta$.
\end{proposisi}

\begin{proof}[Proof]
By Chebyshev's inequality, already proved in Chapter 19,
\[
P(|\hat\theta_n-\theta|\ge\varepsilon)
\le\frac{\mathbb{E}[(\hat\theta_n-\theta)^2]}{\varepsilon^2}
=\frac{\operatorname{MSE}(\hat\theta_n)}{\varepsilon^2}.
\]
For fixed $\varepsilon$ the right-hand side goes to zero, and the inequality closes the
proof.
\end{proof}

Note that Chebyshev uses the MSE, not the variance. The bias term is squared, so a biased
estimator is still consistent provided its bias goes to zero. Conversely, the estimator $X_1$
is unbiased for every $n$ but never uses the rest of the data, so it is not consistent.

The next quantity is used continually until the end of the chapter. The variance of an
estimator says how far the guess wobbles, and its square root has a name of its own.

\kamus{standard error}{galat baku}{the range of wobble of one guess, not the width of the whole
data set}

\begin{definisi}[Standard error]
\index{standard error}
The \emph{standard error} of the estimator $\hat\theta$ is $\operatorname{se}(\hat\theta)
=\sqrt{\operatorname{Var}(\hat\theta)}$. If the formula contains an unknown parameter, the
\emph{estimated} standard error is obtained by replacing that parameter with its estimate.
\end{definisi}

For a proportion, $\operatorname{se}(\hat p)=\sqrt{p(1-p)/n}$, and the estimated version is
$\sqrt{\hat p(1-\hat p)/n}$. For the data $33$ out of $50$, $\hat p=0.66$ and
\[
\sqrt{\frac{0.66\cdot0.34}{50}}=\sqrt{0.004488}=0.066993 .
\]
For twelve measurements of delay time (in milliseconds),
\begin{align*}
&18.2,\ 19.1,\ 17.8,\ 20.3,\ 18.9,\ 19.7,\\
&18.4,\ 21.0,\ 17.5,\ 19.9,\ 18.6,\ 20.1,
\end{align*}
the mean is $19.125$ and the sample standard deviation $1.082191$. The standard error is
$1.082191/\sqrt{12}=0.312401$.
That number does not say that the true mean is $0.312$ away from
$19.125$; it says that guesses like this one, if repeated, wobble on a scale of about
$0.31$. The confidence interval is the way to turn the last sentence into a sentence we can
use.

\section{Confidence intervals for a mean}
\tingkat{intermediate}

A single number without a measure of wobble is easily misused. What is needed is a pair of
numbers that, with high probability, traps the parameter. In everyday language: we do not
shoot at one point, we set a net and promise that the net catches the fish we are after on most
throws.

\kamus{confidence interval}{interval kepercayaan}{a net set with a certain width, and a promise
that the net catches the target on most throws}

\begin{definisi}[Confidence interval and coverage]
\index{confidence interval}
Let $C(X_1,\dots,X_n)$ be an interval computed from the data. $C$ is called a
\emph{confidence interval} of level $1-\alpha$ for $\theta$ if
$P_\theta(\theta\in C)\ge1-\alpha$ for every $\theta$. That probability is called the
\emph{coverage} of the procedure.
\end{definisi}

Note where the probability sits: $P_\theta$ is computed before the data are taken, and what is
random is $C$, not $\theta$. This is where misinterpretation most often occurs, and we return
there once the formula exists.

\begin{teorema}[The z-interval]
\index{confidence interval}
Let $X_1,\dots,X_n$ be independent and distributed $N(\mu,\sigma^2)$ with $\sigma$
known. For $z=z_{1-\alpha/2}$, the interval
\[
C=\Bigl[\bar X_n-z\frac{\sigma}{\sqrt n},\ \bar X_n+z\frac{\sigma}{\sqrt n}\Bigr]
\]
has coverage exactly $1-\alpha$. For other distributions with finite variance, its
coverage tends to $1-\alpha$ as $n\to\infty$.
\end{teorema}

\begin{proof}[Proof]
Let $Z=\sqrt n(\bar X_n-\mu)/\sigma$. For a normal sample, $Z\sim N(0,1)$ exactly;
for other distributions, $Z$ tends to $N(0,1)$ by the central limit theorem in Chapter 19.
The inequality $\mu\in C$ is equivalent to $|Z|\le z$, because
\[
|\bar X_n-\mu|\le z\frac{\sigma}{\sqrt n}
\iff \frac{\sqrt n\,|\bar X_n-\mu|}{\sigma}\le z .
\]
So $P(\mu\in C)=P(|Z|\le z)=1-\alpha$ in the normal case, and tends to $1-\alpha$
in the general case.
\end{proof}

If $\sigma$ is unknown, as in almost all real data, replacing
$\sigma$ by $s$ adds uncertainty and thickens the tails. The result uses the
$t$ distribution with $n-1$ degrees of freedom, published by Gosset in 1908 under the pen
name Student:
\[
C=\Bigl[\bar X_n-t_{n-1,1-\alpha/2}\frac{s}{\sqrt n},\
\bar X_n+t_{n-1,1-\alpha/2}\frac{s}{\sqrt n}\Bigr],
\]
with coverage exactly $1-\alpha$ if the sample is normal. The proof requires the joint
distribution of $\bar X_n$ and $s$ to be independent, and it is in Wasserman, \emph{All of
Statistics}, Chapter 6. For the $n=12$ delay-time data above, $t_{11,0.975}=2.200985$, so
the half-width is $2.200985\cdot0.312401=0.687591$ and the interval is
$(18.4374;\ 19.8126)$. With $z=1.96$ the half-width is $0.612296$: admitting that
$\sigma$ is unknown widens the interval by twelve percent at $n=12$, and by one
percent at $n=100$.

Now the interpretation, and here we must go slowly. Once the data are taken,
$C$ is no longer a random variable: it is the fixed interval $(18.4374;\ 19.8126)$. The
true mean $\mu$ is also a fixed number, only unknown. The sentence
``the probability that $\mu$ lies inside $(18.4374;\ 19.8126)$ equals $0.95$''
is therefore meaningless in this framework: no probability is left, because
nothing is random. What has probability $0.95$ is \emph{the procedure}, and that is a
statement about a trial not yet run.

Figure~\ref{fig:dua-puluh-selang} runs that trial twenty times, with
$\mu=19.125$ and $\sigma=1$ known: nineteen intervals contain $\mu$, one
does not.

\begin{figure}[htbp]
\centering
\begin{tikzpicture}[x=4.15cm,y=0.235cm,font=\scriptsize]
  \draw[->,layergrey] (17.80,0)--(20.22,0);
  \draw[layergrey,densely dashed,thick] (19.125,0)--(19.125,21.4);
  \node[above,font=\scriptsize] at (19.125,21.4) {$\mu=19.125$};
  \foreach \y/\a/\b in {
    1/18.5955/19.7271,
    2/18.5211/19.6527,
    3/18.7441/19.8757,
    4/18.5895/19.7211,
    5/18.4046/19.5362,
    6/18.6636/19.7952,
    7/18.9356/20.0672,
    8/18.8326/19.9642,
    9/18.3561/19.4876,
    10/18.1939/19.3255,
    11/18.3793/19.5109,
    12/18.5711/19.7027,
    14/18.4960/19.6276,
    15/18.1995/19.3311,
    16/18.3478/19.4794,
    17/18.4021/19.5337,
    18/18.4679/19.5995,
    19/18.6780/19.8096,
    20/18.8602/19.9917
  } {
    \draw[very thick,layerindigo] (\a,\y)--(\b,\y);
    \draw[layerindigo] (\a,\y-0.22)--(\a,\y+0.22);
    \draw[layerindigo] (\b,\y-0.22)--(\b,\y+0.22);
  }
  \draw[very thick,accent] (17.8880,13)--(19.0196,13);
  \draw[accent] (17.8880,12.7)--(17.8880,13.3);
  \draw[accent] (19.0196,12.7)--(19.0196,13.3);
  \node[accent,anchor=west,font=\scriptsize] at (17.82,0.62) {one interval misses};
  \foreach \t in {18,19,20} {\draw[layergrey] (\t,-0.2)--(\t,0.2);}
  \foreach \t in {18,19,20} {\node[below,font=\scriptsize] at (\t,-0.45) {\t};}
\end{tikzpicture}
\caption{Twenty $95$ percent confidence intervals for $\mu$, each from one
sample of size $n=12$ with $\sigma=1$ known, so the half-width is
$1.96/\sqrt{12}=0.5658$ and the width is $1.1316$. The vertical dashed line is the
true $\mu=19.125$. Nineteen intervals contain it; the interval in row
13, namely $(17.8880;\,19.0196)$, does not. One failure in twenty is a
reasonable number for a $95$ percent promise, not a sign of a broken procedure.}
\label{fig:dua-puluh-selang}
\end{figure}

\begin{catatan}
Three sentences that are true, and one that is false. True: ``this procedure catches $\mu$
on $95$ percent of repeated trials''. True: ``we are $95$ percent confident that $\mu$
lies between $18.44$ and $19.81$'', provided ``confident'' is read as a summary of
repeated trials. True: ``this procedure is wrong in about five cases out of a hundred''.
False: ``the probability that $\mu$ is in this interval is $0.95$''. The last sentence
pretends to know something about $\mu$ that the data do not provide.
\end{catatan}

The width of an interval shrinks like $1/\sqrt n$, not $1/n$: halving it demands four times
the data, making it a tenth demands a hundred times. This rule can be computed before the
experiment starts, since a half-width $z\sigma/\sqrt n\le d$ is equivalent
to $n\ge(z\sigma/d)^2$; with $\sigma=1$, $z=1.96$, and $d=0.1$, we need
$n\ge385$.

\begin{tangan}
Prepare twenty cards, write the numbers $1$ to $20$ on the cards, one number per card.
Shuffle, draw five cards without looking, add them up, then divide by five. Write the result on
one line of paper. Return the five cards, shuffle again, repeat. Do this thirty times, always
five cards, and arrange those numbers as points on one number line.

What you should see: the points do not pile up at one number. The mean of all the
cards is really $10.5$, and your points scatter around it, some $8$, some
$13$. Now repeat the whole job with ten cards per draw, thirty
times more. The second spread is narrower, but not by half: it narrows by roughly
the square root of two, because the number of cards doubled. Paste the two rows of points side
by side, and see for yourself that doubling the work only trims the wobble by about thirty
percent.
\end{tangan}

\section{Proportions: Wald, Wilson, and the rule of three}
\tingkat{intermediate}

For proportions, the most widely used formula is the simplest version of the previous
theorem: replace $\bar X_n$ by $\hat p$ and $\sigma/\sqrt n$ by the standard error of the
proportion. In everyday language: we set a symmetric interval around the observed proportion,
twice its standard error wide.

\begin{definisi}[Wald interval]
\index{wald interval}
The Wald interval of level $1-\alpha$ for a proportion is
\[
\hat p\pm z\sqrt{\frac{\hat p(1-\hat p)}{n}},\qquad z=z_{1-\alpha/2}.
\]
\end{definisi}

That formula is reasonable as long as $\hat p$ is far from zero and one. The difficulty
appears exactly where the data are scarcest, and its shape is not smooth. If $\hat p=0$, the
formula gives the interval $[0;\,0]$, of zero width: a claim of complete confidence that the
event rate is zero, when all we have is the absence of events in $n$ trials. If $\hat p=1/20$,
the lower limit is $-0.0455$, a number that is not a probability.

The cause is not rounding. Wald uses $\hat p$ for two jobs at once: as the center point and as
the variance estimator. When $\hat p$ approaches zero, the estimated variance shrinks with it,
although the real uncertainty is precisely large on that scale.

The way out is to invert the question: instead of assuming that $\hat p$ spreads normally
around the true $p$, we ask which values of $p$ are still plausible given the observed data.
The criterion is that the normalized distance between $\hat p$ and $p$ does not exceed $z$,
with the variance in the denominator using $p$, not $\hat p$.

\kamus{Wilson interval}{interval wilson}{an interval found by asking which values are still
plausible, rather than by setting a symmetric interval around the guess}

\begin{teorema}[Wilson interval]
\index{wilson interval}
For $0\le k\le n$, $\hat p=k/n$, and $z>0$, the set of $p\in[0,1]$ satisfying
\[
n(\hat p-p)^2\le z^2p(1-p)
\]
is a closed interval with endpoints
\[
\frac{\hat p+\dfrac{z^2}{2n}\ \pm\
z\sqrt{\dfrac{\hat p(1-\hat p)}{n}+\dfrac{z^2}{4n^2}}}{1+\dfrac{z^2}{n}} .
\]
Both endpoints always lie in $[0,1]$, and the interval is never empty.
\end{teorema}

\begin{proof}[Proof]
Rearrange the inequality into $(n+z^2)p^2-(2n\hat p+z^2)p+n\hat p^2\le0$.
Its discriminant is
\[
(2n\hat p+z^2)^2-4(n+z^2)n\hat p^2
=4z^2\Bigl(n\hat p(1-\hat p)+\frac{z^2}{4}\Bigr)\ge0,
\]
so the roots are real and the quadratic formula gives the form in the theorem statement after
numerator and denominator are divided by $2n$.

For the location, let $f(p)=n(\hat p-p)^2-z^2p(1-p)$. The coefficient of $p^2$ is
$n+z^2>0$, so $\{f\le0\}$ is the interval $[r_-,r_+]$. Because
$f(0)=n\hat p^2\ge0$, the point $0$ is not in the interior of that interval, so
$r_-\ge0$; because $f(1)=n(1-\hat p)^2\ge0$, neither is the point $1$, so
$r_+\le1$. Finally $f(\hat p)=-z^2\hat p(1-\hat p)\le0$ places $\hat p$ inside the
interval, so the interval is not empty.
\end{proof}

Now compare the two on the same data. The following table uses $z=1.96$ and
$95$ percent confidence intervals; every number is computed with the formulas above and
rechecked with Python.

\begin{center}
\small
\begin{tabular}{lrrrr}
\hline
data & $\hat p$ & Wald $95$ percent & Wilson $95$ percent & Wilson width \\
\hline
$33/50$  & 0,660 & $(0.5287;\ 0.7913)$ & $(0.5215;\ 0.7756)$ & 0,254 \\
$53/100$ & 0,530 & $(0.4322;\ 0.6278)$ & $(0.4329;\ 0.6249)$ & 0,192 \\
$13/20$  & 0,650 & $(0.4410;\ 0.8590)$ & $(0.4329;\ 0.8188)$ & 0,386 \\
$1/20$   & 0,050 & $(-0.0455;\ 0.1455)$ & $(0.0089;\ 0.2361)$ & 0,227 \\
$0/20$   & 0,000 & $(0;\ 0)$ & $(0;\ 0.1611)$ & 0,161 \\
\hline
\end{tabular}
\end{center}

The first two rows barely differ: when $\hat p$ is near one half, both are legitimate and
agree to two decimal places. The fourth and fifth rows show the separation. At $1/20$ the Wald
interval is \emph{narrower} than Wilson, $0.191$ against $0.227$, and that narrowness is
the sign of an error: its lower limit falls outside $[0,1]$. At $0/20$ Wald gives zero width,
whereas Wilson leaves room up to $0.1611$. The conclusion is not that Wilson is always
better, but that the width of an interval cannot be judged without its coverage.

For the zero-event case there is a special result older than Wilson and easier to remember. In
everyday language: no events at all is not proof that the event rate is zero, and the upper
limit is about three divided by the number of trials.

\kamus{rule of three}{aturan tiga}{if no event occurs in a given number of trials,
the event rate may still be as high as three divided by the number of trials}

\begin{teorema}[Rule of three]
\index{rule of three}
Let $K\sim\mathrm{Bin}(n,p)$ and suppose $K=0$ is observed. The one-sided upper limit at the
$95$ percent level is $p_{\mathrm{hi}}=1-0.05^{1/n}$, and
\[
p_{\mathrm{hi}}\le\frac{\ln 20}{n}<\frac{3}{n}.
\]
\end{teorema}

\begin{proof}[Proof]
The probability of observing zero events is $(1-p)^n$. Setting that probability equal to
$0.05$ gives $(1-p_{\mathrm{hi}})^n=0.05$, that is
$p_{\mathrm{hi}}=1-0.05^{1/n}=1-e^{-c/n}$ with $c=\ln 20=2.9957\ldots$ Because
$1-e^{-x}\le x$ for $x\ge0$, we obtain $p_{\mathrm{hi}}\le c/n<3/n$.
\end{proof}

For $n=100$, the rule of three gives $3/100=0.03$, the exact limit $0.0295130$, and the
two-sided Wilson upper limit at $k=0$ is $z^2/(n+z^2)=0.0369935$. All three numbers answer
the same question with three different promises; the rule of three is the exact one-sided
limit, Wilson a two-sided approximation. What is not allowed is to quote a number without
naming the procedure.

\section{Hypothesis tests and p-values}
\tingkat{intermediate}

Up to here the question took the form ``what is its value''. The next question takes the form
``is its value equal to a particular number''. In everyday language, we file a charge that
there is no effect, then check whether the data are surprising enough to overturn that charge.

\begin{definisi}[Null hypothesis, alternative, and level]
\index{null hypothesis}
The \emph{null hypothesis} $H_0$ is a statement about the parameter to be tested; the
\emph{alternative hypothesis} $H_1$ is the set of values held to be true if $H_0$ is
rejected. The \emph{level} $\alpha$ is the probability of rejecting $H_0$ when $H_0$ is in fact
true. The rejection rule is determined by one test statistic $T$ and one rejection region
$R$: we reject $H_0$ if $T\in R$.
\end{definisi}

The level $\alpha$ is chosen before the data are seen, and therein lies the whole credibility
of the procedure: we are willing to wrongly accuse in a small fraction of trials, and we state
how large that fraction is. The p-value moves the same idea to a single data set.

\begin{definisi}[p-value]
\index{p-value}
For data $x$ and test statistic $T$, the p-value is the probability of obtaining a result at
least as extreme as the observed one if $H_0$ is true,
\[
\text{p-value} p=\sup_{\theta\in H_0}P_\theta\bigl(T\text{ as far as }T(x)\bigr).
\]
We reject $H_0$ at level $\alpha$ if the p-value $\le\alpha$.
\end{definisi}

The phrase ``at least as far'' needs to be fixed for a two-sided test, and the usual
convention is to add the two tails: for a symmetric $z$ statistic, the
p-value $=2(1-\Phi(|z|))$. That definition makes the p-value a random
variable, and its behavior under $H_0$ can be stated exactly.

\begin{teorema}[The p-value is uniform under $H_0$]
\index{p-value}
If $H_0$ is true and $T$ has a continuous distribution with distribution function $F_0$, then
$P_{H_0}(\text{p-value} p\le\alpha)=\alpha$ for every $\alpha\in(0,1)$. Consequently
$P_{H_0}(\text{p-value} p\le0.05)=0.05$: a test with the rule ``reject if the p-value
$\le\alpha$'' has level exactly $\alpha$.
\end{teorema}

\begin{proof}[Proof]
For a right one-sided test, the p-value $=1-F_0(T)$, so
\[
P_{H_0}(1-F_0(T)\le\alpha)
=P_{H_0}\bigl(F_0(T)\ge1-\alpha\bigr)
=P_{H_0}\bigl(T\ge F_0^{-1}(1-\alpha)\bigr)=\alpha,
\]
because $F_0(T)$ is uniformly distributed on $(0,1)$ when $F_0$ is continuous. For a two-sided
test, the p-value depends on $|T|$ and the proof uses the same step on the
distribution of $|T|$.
\end{proof}

That uniformity has a consequence often missed: under $H_0$ the p-value spreads evenly, so a
p-value $=0.04$ appears in four percent of trials whose null hypothesis is true. What
distinguishes two trials is not the single p-value but the power of each.

\begin{bacakeras}
Two voices, taking turns.

``The p-value is point zero four. So the probability that there is no effect is four
percent.''

``No. The p-value counts the data, not the charge. If there were no effect, and
this trial were repeated a hundred times, about four times the data would look this far out.''

``So it is the data that get counted.''

``Yes, the data. The charge is only a place to stand.''

``And if the p-value is large?''

``Then the data are not surprising. The silence of a scale is not proof that the goods are
light.''

Say it once more, slowly: the p-value is the probability of the data, not the probability of
the charge. Repeat that line until it feels natural in the mouth: probability of the data, not
probability of the charge.
\end{bacakeras}


\begin{resep}
\begin{enumerate}
  \item Write $H_0$ and $H_1$ as statements about one parameter, before
        computing anything.
  \item Fix $\alpha$ and its rejection region, then pause. Do not look at
        the data yet.
  \item Compute the test statistic with the standard error under $H_0$, then the p-value.
  \item Compare the p-value with the $\alpha$ already fixed, and report the
        interval together with that decision.
\end{enumerate}
A warning about order: step 2 must be finished before step 3. If $\alpha$ is chosen
after the data are seen, $\alpha$ is no longer a promise but a record of the result, and step 4
loses its meaning. What may not be swapped: step 2 and step 3.
\end{resep}

Now an example with numbers. For the data $33$ out of $50$ and $H_0\colon p=0.5$,
the $z$ statistic uses the standard error under $H_0$, namely $\sqrt{0.25/50}=0.070711$:
\[
z=\frac{0.66-0.5}{0.070711}=2.2627,\qquad
\text{p-value} p=2(1-\Phi(2.2627))=0.0237 .
\]
The exact binomial test on the same rule, $P(X\ge33)+P(X\le17)$ for
$X\sim\mathrm{Bin}(50;0.5)$, gives $0.0328$. Both numbers reject $H_0$ at the five percent
level, but the normal approximation reports a p-value almost thirty percent smaller. At $n=50$
a difference that large is enough to flip the decision when the level is $0.03$; a good
report states which test was used, not only its p-value.

Figure~\ref{fig:dua-ekor} places those two limits on the distribution of $z$ under $H_0$: two
tails of area $0.0250$ each, and the observed statistic $2.2627$ inside the right tail.

\begin{figure}[htbp]
\centering
\begin{tikzpicture}
\begin{axis}[
  width=11.2cm, height=5.4cm,
  xmin=-4, xmax=4, ymin=0, ymax=0.46,
  axis lines=left,
  xtick={-1.96,0,1.96},
  xticklabels={$-1.96$,$0$,$1.96$},
  ytick=\empty,
  tick label style={font=\small},
  xlabel={value of $z$ if $H_0$ is true},
  xlabel style={font=\small},
]
\addplot[fill=accent!25,draw=none,domain=-3.95:-1.96,samples=100]
  {0.39894228*exp(-0.5*x^2)} \closedcycle;
\addplot[fill=accent!25,draw=none,domain=1.96:3.95,samples=100]
  {0.39894228*exp(-0.5*x^2)} \closedcycle;
\addplot[very thick,layerindigo,domain=-3.95:3.95,samples=200]
  {0.39894228*exp(-0.5*x^2)};
\draw[densely dashed,layergrey] (axis cs:-1.96,0)--(axis cs:-1.96,0.0584);
\draw[densely dashed,layergrey] (axis cs:1.96,0)--(axis cs:1.96,0.0584);
\draw[very thick,accent] (axis cs:2.2627,0)--(axis cs:2.2627,0.0308);
\node[font=\scriptsize] (tl) at (axis cs:-2.55,0.135) {$0.025$};
\draw[-{Stealth[length=4pt]},layergrey] (tl)--(axis cs:-2.55,0.012);
\node[accent,font=\scriptsize] (tz) at (axis cs:2.05,0.215) {$z=2.2627$};
\draw[-{Stealth[length=4pt]},accent] (tz)--(axis cs:2.2627,0.055);
\node[font=\scriptsize] (tr) at (axis cs:3.05,0.135) {$0.025$};
\draw[-{Stealth[length=4pt]},layergrey] (tr)--(axis cs:3.05,0.004);
\end{axis}
\end{tikzpicture}
\caption{The distribution of the statistic $z=\sqrt n(\hat p-p_0)/\sqrt{p_0(1-p_0)}$ when
$H_0$ is true, for the data $33$ out of $50$ with $H_0\colon p=0.5$. The two dashed lines at
$\pm1.96$ cut off the two shaded tails, each of area
$1-\Phi(1.96)=0.0250$, so the two-sided level is $0.05$. The thick line at
$2.2627$ is the statistic observed on that data; the right tail from that point has area
$0.0118$, and twice that area is the p-value $0.0237$ that the normal test reports.}
\label{fig:dua-ekor}
\end{figure}

A confidence interval and a hypothesis test are not two different procedures; they are two
ways of reading the same set.

\begin{proposisi}[Equivalence of interval and test]
\index{confidence interval}
For the Wald interval and the two-sided $z$-test that uses the standard error $\hat p$, we have
that $p_0$ is contained in the Wald interval of level $1-\alpha$ if and only if $H_0\colon
p=p_0$ is not rejected at level $\alpha$. More generally, the Wilson interval is the set of
$p_0$ not rejected by the score test at level $\alpha$.
\end{proposisi}

\begin{proof}[Proof]
For the Wald pair, both statements use the same inequality, namely
$|\hat p-p_0|\le z\sqrt{\hat p(1-\hat p)/n}$: the first reads it as
``$p_0$ lies inside the interval'', the second as ``the test statistic does not exceed
$z$''. For Wilson, the set $\{p_0:n(\hat p-p_0)^2\le z^2p_0(1-p_0)\}$ is the
definition of the score test turned around, and the Wilson theorem has already computed its
endpoints.
\end{proof}

\section{Type I error, type II error, and power}
\tingkat{advanced}

Every binary decision can be wrong in two ways, and the two are not symmetric.
In everyday language: a fire alarm can sound when there is no fire, and
can stay silent when there is one. The first mistake sends us out of the house for nothing;
the second is far more costly.

\begin{definisi}[Type I error, type II error, and power]
\index{type i error}
A \emph{type I error} is rejecting $H_0$ when $H_0$ is true; its probability is the level
$\alpha$. A \emph{type II error} is not rejecting $H_0$ when $H_1$ is true;
its probability is written $\beta$. The \emph{power} of a test is $1-\beta$, the probability of
rejecting $H_0$ when the effect sought is really there.
\end{definisi}

Figure~\ref{fig:dua-galat} arranges those four possibilities as a decision box,
with probabilities from the exact computation at $n=20$, $p_0=0.5$, and $p_1=0.75$ that
appears in the computational lab: $\alpha=0.0414$, power $0.6172$, so $\beta=0.3828$.

\begin{figure}[htbp]
\centering
\begin{tikzpicture}[font=\small,
  keputusan/.style={draw=layergrey,thick,rounded corners=2pt,
    minimum width=4.3cm,minimum height=1.6cm,align=center,inner sep=4pt}]
\node[keputusan,fill=layerjade!10] (a) at (0,0)
  {does not reject $H_0$\\ correct, probability $1-\alpha=0.9586$};
\node[keputusan,fill=accent!12] (b) at (5.1,0)
  {rejects $H_0$\\ type I error, probability $\alpha=0.0414$};
\node[keputusan,fill=accent!12] (c) at (0,-2.4)
  {does not reject $H_0$\\ type II error, probability $\beta=0.3828$};
\node[keputusan,fill=layerjade!10] (d) at (5.1,-2.4)
  {rejects $H_0$\\ correct, power $1-\beta=0.6172$};
\node[font=\scriptsize,above=2pt] at (0,1.0) {$H_0$ is true: $p=0.5$};
\node[font=\scriptsize,above=2pt] at (5.1,1.0) {$H_1$ is true: $p=0.75$};
\end{tikzpicture}
\caption{The two kinds of error as four decision boxes. Top row: the actual
state $H_0$ is true; bottom row: $H_1$ is true with $p=0.75$. Left column: the test
does not reject; right column: the test rejects. The boxes in red are the two errors that
cannot be suppressed together: shrinking $\alpha=0.0414$ enlarges
$\beta=0.3828$, and conversely. The numbers are computed from the exact binomial
distribution at $n=20$ with the rejection rule $k\le5$ or $k\ge15$.}
\label{fig:dua-galat}
\end{figure}

Power is not a property of the test alone, but of the test together with the sample size and
the size of the effect. Hence it must be computed for one value $p_1$ deemed important, not in
general. For a two-sided test with $H_0\colon p=0.5$, $n=100$, and the rule to reject when
$|z|>1.96$, the rejection region is $k\le40$ or $k\ge60$. The exact binomial computation
gives size $0.0569$, not $0.05$, and power $0.8750$ at
$p_1=0.65$, whereas the normal approximation gives $0.8508$. The normal approximation,
namely
\[
1-\beta\approx\Phi\Bigl(\frac{(p_1-p_0)\sqrt n}{\sqrt{p_0(1-p_0)}}-z\Bigr),
\]
is close enough for planning, but the two-and-a-half-point difference is real and comes from
the discreteness of the data.

First lesson: the nominal level is seldom attained exactly. At $n=20$ the same rule rejects
when $k\le5$ or $k\ge15$, with exact size $0.0414$; at $n=100$
the size is $0.0569$. A discrete test cannot put a threshold between two whole
numbers, so the actual level hangs on $n$ and $p_0$. Second lesson: at
$n=20$ the power against $p_1=0.75$ is only $0.6172$, so almost four in ten
trials of size twenty fail to find it.

That last number opens the most costly misinterpretation: reading ``p-value
$>0.05$'' as ``no effect''. First, the p-value is
$P(\text{data}\mid H_0)$, not $P(H_0\mid\text{data})$; both are conditional probabilities
with the conditions swapped, and we already met that swap in the diagnostic tests of
Chapter 17. Second, a large p-value is a statement about the data, not about the world: scanty
data cannot surprise anyone. Third, failure to reject can occur precisely because the test is
weak, as at $n=20$ above.

A small example makes all three concrete. Ten successes out of twenty trials against
$H_0\colon p=0.5$ give a two-sided p-value of exactly $1$: that data set sits
precisely in the middle. Is that proof that $p=0.5$? No. The $95$ percent Wilson interval
for the same data is $(0.2993;\ 0.7007)$. So the true rate may be
$0.30$, may be $0.70$, and this data set cannot tell it apart from $0.50$.
A p-value $=1$ only means ``no direction of deviation'', and that is far less than what people
usually read into it.

\begin{tantangan}
Two versions of a detector are tested on \mbox{$100\,000$} examples each. Version A catches
\mbox{$66\,000$}, version B catches \mbox{$65\,000$}. A two-proportion test gives a p-value of about
$0.0000026$, and the team writes in the report: ``A is better, $p<0.001$''. Does that
sentence have the support of the data? Answer with one interval, not with a p-value.

Hint 1: the p-value measures surprise, not the size of the effect. Compute the $95$ percent
interval for the difference of the two proportions, and mind the units.

Hint 2: the standard error of the difference is $0.002126$, so the interval runs from
$0.0058$ to $0.0142$, that is zero point six to one point four percentage points.
Compare it with the practical threshold the team itself uses: three percentage points.

If you get stuck, go back to the section Type I error, type II error, and power.
\end{tantangan}

\begin{untukanak}
A line judge is tasked with raising the flag when the ball goes out of the field. He raises the
flag only when he sees the ball go out himself, clearly enough. If he does not raise the flag,
does that mean the ball was always in? Certainly not. The ball may have gone out, but his view
was blocked by another player, or the ball was very fast.

The way he works is similar to the way a hypothesis test works. He is given one rule that must
be settled before the match: how sure he must be before raising the flag. If the rule is too
loose, he often raises the flag when the ball is in, and the crowd is angry. If the rule is too
strict, he almost never raises the flag, and big fouls slip by. The two mistakes are
different, and he cannot suppress both at once by carefulness alone.

In this chapter we shall learn why ``no evidence'' is not the same as ``no foul'', why the
number of observations matters as much as the rule chosen, and why a judge who never raises
the flag is not a more careful judge.
\end{untukanak}

\section{Worked example: 33 out of 50 and 53 out of 100}
\tingkat{intermediate}

We return to the two detectors in the chapter opening and compute both intervals all the way.
Every number below uses $z=1.96$, so $z^2=3.8416$.

For A, $\hat p=0.66$ and $n=50$. The first term inside the Wilson root is
$0.66\cdot0.34/50=0.2244/50=0.004488$; the second term
$z^2/(4n^2)=3.8416/10000=0.00038416$; their sum is $0.00487216$ and its root
$0.069801$. The center is
\[
\frac{0.66+3.8416/100}{1+3.8416/50}
=\frac{0.698416}{1.076832}=0.648584,
\]
and the half-width is $1.96\cdot0.069801/1.076832=0.127048$. So the Wilson interval
for A is $(0.5215;\ 0.7756)$, of width $0.2541$. The Wald interval for the
same data uses standard error $0.066993$ and half-width
$1.96\cdot0.066993=0.131303$, hence $(0.5287;\ 0.7913)$ of width
$0.2626$.

For B, $\hat p=0.53$ and $n=100$. The first term
$0.53\cdot0.47/100=0.002491$, the second term $3.8416/40000=0.00009604$,
the sum $0.00258704$ and its root $0.050863$. The center
$(0.53+0.019208)/1.038416=0.528890$, the half-width
$1.96\cdot0.050863/1.038416=0.096003$, hence $(0.4329;\ 0.6249)$ of width
$0.1920$. Wald gives $(0.4322;\ 0.6278)$ of width $0.1956$.

Now compare. The intervals A and B overlap on $(0.5215;\ 0.6249)$, an extent of
$0.1034$: about $41$ percent of the length of A and $54$ percent of the length of B. The
normal-based two-proportion test gives $z=1.5175$ and a two-sided p-value of $0.1291$.
Although the raw numbers differ by thirteen points, this data set is not enough to conclude
that A is better; the interval for A is also wider, because it is built from fifty examples.

\kertas{Recompute part A without looking at the text: $0.66\cdot0.34$, divide by $50$,
add $3.8416/10000$, take its root, then divide by $1.076832$ after
multiplying by $1.96$. Check that the result is $0.127048$. Then answer: how many
examples are needed for the width of the interval for A to fall to $0.15$, if $\hat p$
stays $0.66$? Hint: the width is proportional to $1/\sqrt n$, and you already know the width
at $n=50$.}

\section{Computational lab: coverage, type I error, and power}
\tingkat{advanced}

Every claim above can be checked with a computer, and there is one claim that can be checked
only that way: the actual coverage of a procedure. The key is knowing the true $p$, something
that never happens with real data. Hence we work with simulation and with exact binomial
computations.

\begin{lab}
\begin{verbatim}
import numpy as np
from math import sqrt
from scipy.stats import norm, binom

z = norm.ppf(0.975)


def wald(k, n):
    p = k / n
    se = sqrt(p * (1 - p) / n)
    return p - z * se, p + z * se


def wilson(k, n):
    p = k / n
    z2 = z * z
    mid = (p + z2 / (2 * n)) / (1 + z2 / n)
    half = z * sqrt(p * (1 - p) / n + z2 / (4 * n * n)) / (1 + z2 / n)
    return mid - half, mid + half


# 1. width of the 95% interval for several proportions
print("== 1. width of the 95% interval ==")
print("   n    k   phat  Wald width    Wilson width")
for n in (20, 100):
    for p in (0.00, 0.05, 0.25, 0.50, 0.66):
        k = round(p * n)
        a = wald(k, n)
        b = wilson(k, n)
        print("%4d %4d  %5.2f   %9.4f    %9.4f"
              % (n, k, k / n, a[1] - a[0], b[1] - b[0]))

# 2. exact coverage at p = 0.05
print("== 2. exact coverage, p=0.05 ==")
for n in (20, 50, 100):
    for nama, f in (("Wald", wald), ("Wilson", wilson)):
        c = 0.0
        for k in range(n + 1):
            lo, hi = f(k, n)
            if lo <= 0.05 <= hi:
                c += binom.pmf(k, n, 0.05)
        print("n=%3d %-6s coverage=%.4f" % (n, nama, c))

# 3. simulated coverage with a fixed seed
rng = np.random.default_rng(20261024)
R = 20000
print("== 3. simulated coverage, R=%d ==" % R)
for n in (20, 100):
    X = rng.binomial(n, 0.05, R)
    for nama, f in (("Wald", wald), ("Wilson", wilson)):
        ok = 0
        w = 0.0
        for k in X:
            lo, hi = f(int(k), n)
            ok += (lo <= 0.05 <= hi)
            w += hi - lo
        print("n=%3d %-6s coverage=%.4f  width=%.4f" % (n, nama, ok / R, w / R))

# 4. type I error and power for H0: p=0.5
print("== 4. type I error and power, H0 p=0.5 ==")


def tolak(k, n, p0=0.5):
    return abs((k / n - p0) / sqrt(p0 * (1 - p0) / n)) > z


for n in (20, 100):
    rej = [k for k in range(n + 1) if tolak(k, n)]
    lo = max(k for k in rej if k < n / 2)
    hi = min(k for k in rej if k > n / 2)
    print("n=%3d reject if k<=%d or k>=%d" % (n, lo, hi))
    print("      exact size at p=0.5 : %.4f"
          % sum(binom.pmf(k, n, 0.5) for k in rej))
    for p1 in (0.65, 0.75):
        ex = sum(binom.pmf(k, n, p1) for k in rej)
        zc = (p1 - 0.5) / sqrt(0.25 / n)
        ap = norm.cdf(zc - z) + norm.cdf(-zc - z)
        print("      power at p=%.2f: exact %.4f  approx %.4f"
              % (p1, ex, ap))

# 4b. simulated type I error rate when H0 is true
print("== 4b. simulated type I error, R=%d ==" % R)
for n in (20, 100):
    X0 = rng.binomial(n, 0.5, R)
    tolak0 = np.abs((X0 / n - 0.5) / sqrt(0.25 / n)) > z
    print("n=%3d p=0.5: reject=%.4f" % (n, tolak0.mean()))

# 5. t-interval for 12 measurements
x = np.array([18.2, 19.1, 17.8, 20.3, 18.9, 19.7,
              18.4, 21.0, 17.5, 19.9, 18.6, 20.1])
n = len(x)
m = x.mean()
s = x.std(ddof=1)
from scipy.stats import t
tq = t.ppf(0.975, n - 1)
h = tq * s / sqrt(n)
print("== 5. t-interval, n=12 ==")
print("mean=%.6f  s=%.6f  se=%.6f" % (m, s, s / sqrt(n)))
print("t=%.6f  half-width=%.6f  CI=(%.6f, %.6f)" % (tq, h, m - h, m + h))
print("half-width with z=1.96: %.6f" % (z * s / sqrt(n)))
\end{verbatim}
\end{lab}

The actual output, verbatim:

\begin{verbatim}
== 1. width of the 95% interval ==
   n    k   phat  Wald width    Wilson width
  20    0   0.00      0.0000       0.1611
  20    1   0.05      0.1910       0.2272
  20    5   0.25      0.3795       0.3568
  20   10   0.50      0.4383       0.4014
  20   13   0.65      0.4181       0.3860
 100    0   0.00      0.0000       0.0370
 100    5   0.05      0.0854       0.0902
 100   25   0.25      0.1697       0.1676
 100   50   0.50      0.1960       0.1923
 100   66   0.66      0.1857       0.1826
== 2. exact coverage, p=0.05 ==
n= 20 Wald   coverage=0.6389
n= 20 Wilson coverage=0.9245
n= 50 Wald   coverage=0.9199
n= 50 Wilson coverage=0.9622
n=100 Wald   coverage=0.8775
n=100 Wilson coverage=0.9659
== 3. simulated coverage, R=20000 ==
n= 20 Wald   coverage=0.6430  width=0.1472
n= 20 Wilson coverage=0.9224  width=0.2193
n=100 Wald   coverage=0.8783  width=0.0829
n=100 Wilson coverage=0.9667  width=0.0884
== 4. type I error and power, H0 p=0.5 ==
n= 20 reject if k<=5 or k>=15
      exact size at p=0.5 : 0.0414
      power at p=0.65: exact 0.2457  approx 0.2687
      power at p=0.75: exact 0.6172  approx 0.6088
n=100 reject if k<=40 or k>=60
      exact size at p=0.5 : 0.0569
      power at p=0.65: exact 0.8750  approx 0.8508
      power at p=0.75: exact 0.9997  approx 0.9988
== 4b. simulated type I error, R=20000 ==
n= 20 p=0.5: reject=0.0430
n=100 p=0.5: reject=0.0554
== 5. t-interval, n=12 ==
mean=19.125000  s=1.082191  se=0.312401
t=2.200985  half-width=0.687591  CI=(18.437409, 19.812591)
half-width with z=1.96: 0.612296
\end{verbatim}

Four things deserve reading from that output. First, parts 2 and 3 agree to almost
three digits, and that agreement is itself a check, since a simulation with
$R=20000$ repetitions can approach the exact computation only to about
$\sqrt{0.25/20000}=0.0035$. At $p=0.05$ and $n=20$ the difference is $0.0041$ for
Wald and $0.0021$ for Wilson. An honest simulation must report $R$.

Second, the coverage of Wald at $n=20$ and $p=0.05$ is $0.6389$, not $0.95$.
The cause can be computed by hand: the probability of observing zero events is $0.3585$,
and at $k=0$ the Wald interval is $[0;0]$, which does not contain $0.05$; that single point
alone consumes almost thirty-six points of coverage. Wilson fails only when
$k\ge3$, whose probability is $0.0755$, and its failure is not total because its interval
stays wide. That is why it reaches $0.9245$. Note also the width column: the mean
Wald width $0.1472$ is smaller than the Wilson $0.2193$, and precisely therein lies the
point. A narrower interval is not a better interval if it is narrow because it sits in the
wrong place.

Third, part 4 separates the size of a test from the nominal level: $0.0414$ at $n=20$ and
$0.0569$ at $n=100$, both for the rule $|z|>1.96$; part 4b repeats it
by simulation and lands close, $0.0430$ and $0.0554$. Power moves far more
drastically: $0.2457$ at $n=20$ and $0.8750$ at $n=100$ for $p_1=0.65$.
The normal approximation is not always on the safe side: at $n=20$ it reports $0.2687$
while the truth is $0.2457$.

Fourth, part 5 merely repeats the $t$-interval already computed in the text, with a
difference in half-width between $t$ and $z$ of $0.075295$.

The program does not prove any one of all those things. It checks claims already
proved, and in two places it shows its own limits: exact coverage
uses the distribution we assume to be true, so if the model is wrong, the number
$0.9245$ means nothing; and simulation only measures how often
the procedure succeeds under that model, not how right the model is.

\begin{untukdosen}
The teaching order I suggest: start from the morning meeting in the chapter opening, not from
the definition of an estimator. Write $33/50$ and $53/100$ on the board, ask the class to choose, then
ask what happens if those fifty examples are replaced. Out of that unease the
definitions of estimator, bias, and standard error emerge as the answer. The MSE
decomposition is proved at the board with one line of algebra. The interpretation of the
confidence interval comes after that, as material for a lecture of its own, not a footnote.

Two misconceptions appear most often. First, students read ``$95$ percent''
as a conditional probability on data already in hand, then write ``the probability that $\mu$ is in
this interval is $0.95$''; the symptom shows when they cannot explain what is
random about an interval already computed. Second, they take a large p-value as
proof that $H_0$ is true, and will compute the p-value for $10/20$ and then conclude that
$p=0.5$ is established.

A blackboard activity that works: split the class into two groups, one computing the
Wald interval and one the Wilson interval for $1/20$, then have both defend their result
in front of the class. The Wald group will find a negative lower limit and have to
explain it. The sign that the class is ready to go on: they can name one example of a
sample size that makes the p-value large although the effect is large, and can explain
why that happens.
\end{untukdosen}

\section{Summary}
\tingkat{basic}

\begin{ringkasan}
\begin{itemize}
  \item An estimator is a statistic that guesses a parameter, so it is a random variable with
        an expectation and a variance. Its quality is measured by
        $\operatorname{MSE}=\operatorname{bias}^2+\operatorname{Var}$; unbiasedness alone
        guarantees nothing, and MSE going to zero guarantees consistency through
        Chebyshev.
  \item The standard error is the square root of the variance of the estimator, not the width
        of the data. For a proportion its value is $\sqrt{\hat p(1-\hat p)/n}$, shrinking like
        $1/\sqrt n$, so four times the data only cuts it in half.
  \item A $95$ percent confidence interval is a statement about the procedure, not
        about the parameter: the procedure catches the true value on $95$ percent
        of repeated trials. Once the data are taken, the interval is fixed and no longer
        carries a probability.
  \item The Wald interval collapses near $0$ and $1$: at $0/20$ its width is zero, and at
        $p=0.05$ with $n=20$ its coverage is only $0.6389$. The Wilson interval, the
        result of inverting the score test, reaches $0.9245$ in the same setting.
  \item If no event occurs in $n$ trials, the $95$ percent upper limit is about
        $3/n$; for $n=100$ the exact limit is $0.0295130$ and the two-sided Wilson
        limit $0.0369935$.
  \item The p-value is the probability of data at least that extreme if $H_0$ is true, and
        under $H_0$ it is uniformly distributed. A large p-value is not proof that $H_0$ is
        true: at $n=20$ the power against $p=0.75$ is only $0.6172$.
\end{itemize}
\end{ringkasan}

\section{Exercises}
\tingkat{intermediate}

\begin{latihan}
\begin{enumerate}
  \item Assess this claim: ``Because the p-value equals $0.42$, the null hypothesis
        $p=0.5$ is proved true.'' True, false, or not enough evidence? Name
        the part that is right, the part that is wrong, and one additional computation that
        would strengthen your answer.
  \item Prove that $\mathbb{E}[\hat p(1-\hat p)]=p(1-p)(n-1)/n$ for
        $\hat p=K/n$ with $K\sim\mathrm{Bin}(n,p)$, then determine the unbiased estimator
        of $p(1-p)$ obtained from it. Check the result at $p=0.3$ and
        $n=20$.
  \item Prove that if $\operatorname{bias}(\hat\theta_n)\to0$ and
        $\operatorname{Var}(\hat\theta_n)\to0$, then $\hat\theta_n$ is consistent for
        $\theta$. Explain why that condition is weaker than
        $\operatorname{MSE}\to0$ yet sufficient for the same conclusion.
  \item Compute the $95$ percent Wilson and Wald intervals for $7/40$, then compare
        their widths with $13/20$. Explain why the smaller width at $7/40$
        is not a sign of higher precision, and check both intervals with
        Python.
  \item For $n=30$ and $H_0\colon p=0.5$ with a normal-based two-sided test,
        determine the rejection region, compute its exact size, then compute its power at
        $p_1=0.7$. Compare that power with $n=100$ at the same
        $p_1$.
  \item Prove that the Wald interval and the two-sided $z$-test using the standard error
        $\hat p$ really are equivalent, by writing the two inequalities side by side.
        Explain why the equivalence fails when the test uses the standard error under
        $H_0$.
  \item A report writes: ``there is no significant difference, $p=0.31$,
        so the two methods are equally good''. Rewrite that sentence honestly, naming the
        interval, the level, and the power. You may invent the numbers as long as they are
        consistent.
  \item For $n=100$ with not a single event, compute the rule of three, the exact one-sided
        limit, and the Wilson upper limit. Explain why the three numbers differ, then
        state which one may be used for the sentence ``at most $5$ percent''.
  \item Compute the $90$ percent $t$-interval for the twelve delay-time data in this chapter,
        then compare it with the $95\%$ interval. How much data is needed for the width
        of the $95\%$ interval to fall by half, if $s$ does not change?
\end{enumerate}

Hints for the three hardest items. Item 3: write
$\hat\theta_n-\theta=\operatorname{bias}_n+(\hat\theta_n-\mathbb{E}\hat\theta_n)$,
then use the triangle inequality and Chebyshev on the second term. Item 6: the first inequality
is $|\hat p-p_0|\le z\sqrt{\hat p(1-\hat p)/n}$, and a test that uses
$\sqrt{p_0(1-p_0)/n}$ uses a different denominator, so the set of $p_0$ not
rejected is no longer the Wald interval but a more complicated set. Item 5: the rejection
region is $|k-15|>1.96\sqrt{30\cdot0.25}$, so $k\le9$ or $k\ge21$, and
the power is computed by summing the binomial pmf at $p=0.7$ over those two tails.
\end{latihan}

\begin{untukpeneliti}
What is really proved in this chapter: the MSE decomposition; consistency from MSE going to
zero; exact coverage of the $z$-interval under the normal model; the closed form of the Wilson
interval and its location in $[0,1]$; the rule of three as an exact one-sided limit together
with the inequality $p_{\mathrm{hi}}<\ln 20/n$; the uniformity of the p-value under $H_0$ for a
continuous statistic; and the equivalence of interval and test in the pairs mentioned.
What is merely popular, not a theorem: that $95\%$ means a conditional probability on
data already observed; that a large p-value supports $H_0$; and that a smaller interval width
is always better, when a comparison of widths without coverage means nothing. The subnominal
coverage of Wilson at some $(n,p)$ is also a computational fact, not a theorem.

Primary source: Wasserman, \emph{All of Statistics}, Chapter 6 (Models, Statistical
Inference, and Learning) for bias, MSE, and the interpretation of intervals; Chapter 9
(Parametric Inference) for intervals and the plug-in method; Chapter 10 (Hypothesis Testing
and p-values) for p-values, type I and II errors, power, and the relation of tests to
intervals. The Wilson interval comes from E.~B. Wilson, ``Probable Inference, the Law of
Succession, and Statistical Inference'' (1927); the rule of three from B.~D. Jovanovic and
P.~S. Levy, ``A Look at the Rule of Three'' (1997). Neither is discussed by Wasserman, so the
references are cited directly, not through an intermediate textbook.

Limits of validity: the whole chapter works with independent and identically distributed
samples, one parameter, and a known normal or binomial approximation. The $z$-interval is
exact only under the normal model with $\sigma$ known; beyond that it is an asymptotic
statement, and at small $n$ with heavy-tailed distributions its coverage can fall far below
nominal. The Wilson interval is a score approximation, not an exact interval; exact intervals
for proportions are computed from the binomial distribution and are wider. Every reported
coverage is conditional on the binomial model being correct.

Open questions worth a reader's work: how well proportion intervals behave when $n$ is not
fixed but stops adaptively, as in sequential sampling; and how to build intervals for
classification metrics computed from dependent data, for instance when one example appears in
several cross-validation folds. Both are still research material.
\end{untukpeneliti}

\peta{Chapter 20 in Wasserman, \emph{All of Statistics}: Chapter 6 (Models, Statistical
Inference, and Learning: estimators, bias, MSE, standard error, interpretation of intervals),
Chapter 9 (Parametric Inference: confidence intervals and the plug-in method), Chapter 10
(Hypothesis Testing and p-values: p-values, type I and II errors, power, equivalence of tests
and intervals). The Wilson interval and the rule of three are taken from their original
papers, Wilson (1927) and Jovanovic--Levy (1997), because Wasserman discusses neither. The law
of large numbers, the central limit theorem, and Chebyshev's inequality used here are proved
in Chapter 18; conditional probability and the error of swapping the conditions are discussed
in Chapter 16. Chapter 25 continues this chapter for classification metrics, including the
rule of three in the reporting of model evaluation.}

\chapter{Regression and Applied Statistics}

A small café records nine trading days. Two numbers are written every evening: the number of
online orders in tens, and the kitchen's overtime hours. The café owner needs one number for
one question: if orders rise by ten, how many extra overtime hours must she prepare for
tomorrow?

\begin{center}
\small
\begin{tabular}{l*{9}{r}}
\hline
day                & 1 & 2 & 3 & 4 & 5 & 6 & 7 & 8 & 9 \\
orders (tens)      & 1 & 2 & 3 & 4 & 5 & 6 & 7 & 8 & 9 \\
overtime (hours)   & 10 & 12 & 13 & 14 & 16 & 18 & 19 & 20 & 22 \\
\hline
\end{tabular}
\end{center}

A quick answer is available without any tool: from the first day to the last, orders rise by
eight tens and overtime rises by twelve hours, so roughly $1.5$ hours per ten orders. Least
squares, which we will build in this chapter, gives $1.466667$. The two numbers are close, and
precisely because they are close we are tempted to treat the matter as settled. Three
important questions are left unanswered by both. How far does that number move if those nine days
are replaced by nine other days? If tomorrow's orders are eighty, how many hours must be
prepared? And if one day turns out to be misrecorded, how much damage does that do?

Chapter 10 already gave the picture: least squares is the orthogonal projection of the
observation vector onto the column space of the design matrix, and the normal equations say
that the residual is perpendicular to that space. What is missing there is uncertainty.
Chapter 20 already gave confidence intervals and hypothesis tests, but for a single parameter
from observations that are independent and share one distribution. Here the parameter sought
is a slope, a ratio of two sums that are both random, and the observations are no longer on
equal footing: a point at the end of the range carries more influence than a point in the
middle. This chapter joins the two and stops at an honest place: how far a conclusion may be
carried.

\begin{untukanak}
A grandmother sells coconut ice in front of her house. She never records anything neatly, but
her memory is good enough: on a hot day her stock sells out before the afternoon. On a cloudy
day half the ice melts in the cooler. One day her grandchild arrives with an exercise book and
records nine days in a row: roughly how hot it was, and how many cups were sold.

From those nine records the grandchild draws one straight line: the hotter it gets, the more is
sold. The grandmother nods, but then asks, ``if tomorrow is a little hotter, how many cups
should I prepare?'' The grandchild falls silent. The line gives the direction, not the number.
Besides, on one day it rained hard although it was hot, and the sales were dreadful; that point
sits far from the line.

In this chapter we will learn why the best-fitting line is still not a sure forecast, why one
strange day can shift the line further than five ordinary days, and why anyone who sells
numbers must state how wide their uncertainty is.
\end{untukanak}

\section{Nine days and one line}
\tingkat{basic}

Nine pairs of numbers, two unknowns, nine equations. We met a system like this in Chapter 8:
more equations than unknowns, so an exact solution exists only if all nine points happen to be
collinear. No line passes exactly through every point above, and that is not a measurement
failure but the usual state of affairs. All we can look for is the closest line, and
``closest'' has to be defined first.

In everyday language: we look for the one line whose total error is as small as possible, with
the error counted as the perpendicular distance from the point to the line.

\kamus{regression}{regresi}{drawing the one line closest to all the points, rather than a line
that passes through all of them}

\begin{definisi}[Simple linear regression model]
\index{regression}
For observations $i=1,\dots,n$,
\[
y_i = a + b\,x_i + \varepsilon_i ,
\]
with $x_i$ the observed explanatory variable, $y_i$ the explained variable, $a$ the intercept,
$b$ the slope, and $\varepsilon_i$ the unobserved error. The model is called \emph{linear}
because it is linear in the parameters $a$ and $b$, not because the line is straight in $x$.
\end{definisi}

One small distinction governs the whole chapter: $\varepsilon_i$ and the residual are not the
same object. The error $\varepsilon_i$ is the difference between the true value and the
population line that we never see; the residual $e_i=y_i-\hat y_i$ is the difference between
the observed value and the line we fitted from the data itself. A residual can be computed and
drawn, an error can only be assumed. Whenever a book states an assumption about errors and then
checks it through a picture of residuals, there is one jump that must not be forgotten: what is
inspected is a stand-in, not the thing itself.

\section{Least squares as a projection}
\tingkat{intermediate}

Stack $y\in\mathbb{R}^n$ as the observation vector and let $X$ be the $n\times2$ matrix whose
columns are one and $x$. Finding the best line means choosing $\hat y$ in the column space of
$X$ that is closest to $y$: the problem of Chapter 10, with the addition that we need the
explicit form of the coefficients in order to compute their uncertainty.

\begin{definisi}[Least-squares estimator]
\index{least squares}
$\hat\beta=(\hat a,\hat b)$ is the least-squares estimator if it minimizes
\[
\operatorname{RSS}(\beta)=\|y-X\beta\|^2=\sum_{i=1}^{n}\bigl(y_i-a-bx_i\bigr)^2 .
\]
\end{definisi}

\kamus{least squares}{kuadrat terkecil}{adding up all the errors after squaring them, then
choosing the line whose total is smallest}

\begin{teorema}[The normal equations]
\index{normal equations}
$\hat\beta$ minimizes RSS if and only if
\[
X^{\top}(y-X\hat\beta)=0 ,
\qquad\text{that is}\qquad X^{\top}X\hat\beta=X^{\top}y .
\]
If the columns of $X$ are linearly independent, $\hat\beta=(X^{\top}X)^{-1}X^{\top}y$ is
unique.
\end{teorema}

\begin{proof}[Proof]
Write $f(\beta)=\|y-X\beta\|^2=y^{\top}y-2\beta^{\top}X^{\top}y+\beta^{\top}X^{\top}X\beta$.
Its gradient is $-2X^{\top}y+2X^{\top}X\beta$, so the stationary condition is exactly the normal
equations. Since $X^{\top}X$ is positive definite when its columns are linearly independent,
$f$ is strictly convex and that stationary point is the global minimum.

The converse direction is more informative. Suppose $X^{\top}(y-X\hat\beta)=0$. For any
$\beta$,
\[
\|y-X\beta\|^2=\|y-X\hat\beta\|^2+\|X(\hat\beta-\beta)\|^2
+2(\hat\beta-\beta)^{\top}X^{\top}(y-X\hat\beta),
\]
and the cross term vanishes. Hence $\operatorname{RSS}(\beta)\ge\operatorname{RSS}(\hat\beta)$
for every $\beta$, with equality only when $X\beta=X\hat\beta$.
\end{proof}

That proof contains the whole geometric content: the residual is perpendicular to every column
of $X$. Since the first column is the vector of ones and the second is $x$, we get
$\sum_i e_i=0$ and $\sum_i x_ie_i=0$, and those two identities are what support the
decomposition of sums of squares in the next section.

\begin{proposisi}[Hat matrix]
\index{hat matrix}
With $H=X(X^{\top}X)^{-1}X^{\top}$ we have $\hat y=Hy$, $H^{\top}=H$, $H^2=H$,
$\operatorname{tr}H=2$, and $0\le h_{ii}\le1$ with $\sum_i h_{ii}=2$. The value $h_{ii}$ is
called the \emph{leverage} of point $i$.
\end{proposisi}

\begin{proof}[Proof]
Symmetry and idempotence follow directly from the definition:
\[
H^2=X(X^{\top}X)^{-1}X^{\top}X(X^{\top}X)^{-1}X^{\top}=H .
\]
Hence $H$ is an orthogonal projection matrix, so
$h_{ii}=e_i^{\top}He_i=\|He_i\|^2\ge0$ and $h_{ii}\le\|e_i\|^2=1$. Finally
$\operatorname{tr}H=\operatorname{tr}\bigl((X^{\top}X)^{-1}X^{\top}X\bigr)
=\operatorname{tr}I_2=2$.
\end{proof}

Now the numbers, because the whole chapter rests on one computation that can be rechecked. For
the café data,
\[
X^{\top}X=\begin{pmatrix}9&45\\45&285\end{pmatrix},
\qquad X^{\top}y=\begin{pmatrix}144\\808\end{pmatrix},
\qquad \det(X^{\top}X)=9\cdot285-45^2=540 .
\]
Therefore
\[
\hat\beta=\frac1{540}\begin{pmatrix}285&-45\\-45&9\end{pmatrix}
\begin{pmatrix}144\\808\end{pmatrix}
=\frac1{540}\begin{pmatrix}4680\\792\end{pmatrix}
=\begin{pmatrix}8.666667\\1.466667\end{pmatrix}.
\]
The number $808$ is worth checking by hand: $\sum x_iy_i=10+24+39+56+80+108+133+160+198$. So is
$4680$: $285\cdot144-45\cdot808=41040-36360$.

One more theorem is needed before we may speak about the quality of an estimator. That theorem
explains why least squares is used in almost every field, and at the same time why it is not a
sacred choice.

\begin{teorema}[Gauss--Markov]
\index{gauss-markov}
Assume (A1) $y=X\beta+\varepsilon$; (A2) $\mathbb{E}[\varepsilon\mid X]=0$;
(A3) $\operatorname{Var}(\varepsilon\mid X)=\sigma^2I$; (A4) the columns of $X$ are linearly
independent. Then $\hat\beta$ is unbiased with $\operatorname{Var}(\hat\beta)=\sigma^2
(X^{\top}X)^{-1}$, and among all linear unbiased estimators, $\hat\beta$ has the smallest
variance. Such an estimator is called \emph{BLUE} (\emph{best linear unbiased estimator}).
\end{teorema}

\begin{proof}[Proof]
Unbiasedness: $\mathbb{E}[\hat\beta\mid X]=(X^{\top}X)^{-1}X^{\top}X\beta=\beta$, using (A2).
Variance: $\hat\beta-\beta=(X^{\top}X)^{-1}X^{\top}\varepsilon$, so
$\operatorname{Var}(\hat\beta)=(X^{\top}X)^{-1}X^{\top}(\sigma^2I)X(X^{\top}X)^{-1}
=\sigma^2(X^{\top}X)^{-1}$.

The third part is not proved in full here. The sketch: every linear unbiased estimator has the
form $\tilde\beta=Cy$ with $CX=I$, so $\tilde\beta-\beta=C\varepsilon$ and
$\operatorname{Var}(\tilde\beta)=\sigma^2CC^{\top}$. Writing
$C=(X^{\top}X)^{-1}X^{\top}+D$ gives $DX=0$ and
\[
CC^{\top}-(X^{\top}X)^{-1}=DD^{\top},
\]
which is positive semi-definite, so every component of $\tilde\beta$ has variance no smaller
than the matching component of $\hat\beta$. The full proof is in Seber and Lee,
\emph{Linear Regression Analysis}, Chapter 3, and in Greene, \emph{Econometric Analysis},
Chapter 3.
\end{proof}

\section{Coefficients and their units}
\tingkat{basic}

The slope $1.466667$ means nothing until its unit is named. The unit is hours per ten orders:
every additional ten orders raises the forecast overtime by $1.466667$ hours, that is $88$
minutes, or $8.8$ minutes per order. The café owner can use that sentence; the sentence ``the
coefficient is $1.47$'' she cannot.

The intercept $8.666667$ hours is the forecast for $x=0$, a day with no online orders at all.
No such day occurs in the data, and outside the range $1$ to $9$ the line is not guaranteed to
hold. Calling that number ``fixed overtime cost'' is \index{extrapolation}extrapolation wrapped
in a name: we are using the line outside the place where it was fitted. What is safe: the
intercept is the number that makes the line pass through $(\bar x,\bar y)$, because
$\hat a=\bar y-\hat b\bar x$; it is useful for computation, not for interpretation.

A change of units does not change the conclusion. If $x$ is written in orders rather than tens,
the slope is $0.146667$ hours per order and its standard error is $0.004714$, both divided by
ten; $t$ and $R^2$ do not move. Multiplying one variable by a constant $c$ multiplies the
coefficient and its standard error by $c$ or $1/c$, while their ratio stays put, so a report
that quotes $t$ without units cannot be read.

\begin{tangan}
Prepare graph paper, nine coins, and a piece of string as long as the side of the paper. Place
the coins on the grid according to the nine pairs of numbers at the start of the chapter:
horizontally orders, vertically overtime, one square to one unit. Tie the ends of the string to
the two edges of the paper so that the string can be slid as a straight line. Slide and rotate
the string until the sum of the perpendicular distances from all the coins to the string feels
smallest; do not measure, just compare by eye, then lock it with tape.

What should be visible: two people will put the string in slightly different places, and the
difference is largest when you shift the ends of the string. Now lift the leftmost coin by one
square and notice how far the string has to turn. Put it back, then lift the middle coin by one
square: almost nothing changes. That is leverage, and you have just seen it with your hands
before meeting the formula.
\end{tangan}

\section{R squared and its correlation}
\tingkat{intermediate}

Once the line exists, people ask ``how good is it''. The most common answer is one number that
compares the residual with the original spread of $y$. In everyday language: what share of the
wobble in $y$ the line manages to explain, and what share is left over.

\begin{definisi}[$R^2$]
\index{coefficient of determination}
With $\mathrm{SST}=\sum_i(y_i-\bar y)^2$, $\mathrm{SSE}=\sum_i e_i^2$, and
$\mathrm{SSR}=\sum_i(\hat y_i-\bar y)^2$, the coefficient of determination is
\[
R^2=1-\frac{\mathrm{SSE}}{\mathrm{SST}}=\frac{\mathrm{SSR}}{\mathrm{SST}} .
\]
\end{definisi}

\kamus{coefficient of determination}{koefisien determinasi}{the share of the wobble that the
line manages to explain, expressed as a fraction of the whole}

\begin{teorema}[Decomposition of sums of squares]
If $X$ contains a column of ones, then $\mathrm{SST}=\mathrm{SSR}+\mathrm{SSE}$ and
$0\le R^2\le1$.
\end{teorema}

\begin{proof}[Proof]
Write $y-\bar y\mathbf{1}=(\hat y-\bar y\mathbf{1})+e$. Since $\sum e_i=0$, the vector $e$ is
perpendicular to $\bar y\mathbf{1}$, and by the normal equations $e$ is also perpendicular to
$\hat y$. Hence the two terms on the right are perpendicular, and Pythagoras gives
$\|y-\bar y\mathbf{1}\|^2=\|\hat y-\bar y\mathbf{1}\|^2+\|e\|^2$.
\end{proof}

For the café data, $\mathrm{SST}=130$, $\mathrm{SSE}=0.933333$, so
$R^2=1-0.933333/130=0.992821$. In a simple regression, $R^2$ equals the squared correlation:
$r=88/\sqrt{60\cdot130}=0.996404$ and $r^2=0.992821$. That equality is special to one
explanatory variable; as soon as there are two, $R^2$ is no longer the square of any
correlation.

The number $0.992821$ easily makes people stop thinking, so let us tame it. What is explained is
$99.28$ percent of the \emph{wobble} in overtime, and that wobble is measured against its own
mean. The residuals still have variance $s^2=0.933333/7=0.133333$, that is a standard deviation
$s=0.365148$ hours or $21.9$ minutes. A café owner who prepares overtime on the basis of this
line can still miss by more than twenty minutes. $R^2$ is relative; $s$ is absolute, and what
people feel is $s$.

\begin{untukanak}
Imagine two students. The first runs a hundred-meter race and finishes $0.2$ seconds behind the
winner. The second runs a marathon and finishes twenty minutes behind the winner. Both lose, but
not in the same sense. Two tenths of a second is very close; twenty minutes is very far. If the
teacher writes only ``second place'' for both, we lose the whole story.

Numbers of that kind occur in this chapter. There is a number that measures the \emph{share}
guessed correctly, and there is a number that measures the \emph{size} of the error left over.
The first can look almost perfect while the second stays large. Someone who reads only the first
number will prepare too little, then wonder why the forecast misses.
\end{untukanak}

\begin{proposisi}[Adding a variable never lowers $R^2$]
If a column block $X_2$ is added to $X_1$ so that its column space contains the column space of
$X_1$, then $\mathrm{SSE}$ does not increase and $R^2$ does not decrease.
\end{proposisi}

\begin{proof}[Proof]
The set of vectors $X_1\beta_1$ is contained in the set $X_2\beta_2$, so the minimum distance to
$y$ taken over the larger set cannot be larger. Hence $R^2$ cannot fall, even when the added
variable has nothing to do with $y$.
\end{proof}

The practical consequence matters: comparing models by raw $R^2$ always favors the fatter model.
The usual correction is $R^2_{\text{adj}}=1-(1-R^2)(n-1)/(n-p)$ with $p$ the number of
parameters. For the multiple model of the next section, $R^2=0.997951$ falls to $0.997496$ when
$p=3$ and $n=12$; at small $n$ with many variables the correction can reverse the ranking.

\section{Multiple regression and collinearity}
\tingkat{advanced}

So far, one explanatory variable. In reality almost every problem demands more than one, and
that is where its characteristic disease appears: two explanatory variables carrying almost the
same information.

Twelve branches of a food stall are recorded: daily sales $y$ in millions of rupiah, the number
of employees $x_1$, and the seating area $x_2$ in square meters.

\begin{center}
\small
\begin{tabular}{l*{12}{r}}
\hline
branch   & 1 & 2 & 3 & 4 & 5 & 6 & 7 & 8 & 9 & 10 & 11 & 12 \\
$x_1$    & 6 & 7 & 7 & 8 & 8 & 9 & 9 & 10 & 10 & 11 & 12 & 12 \\
$x_2$    & 18 & 22 & 20 & 26 & 24 & 28 & 26 & 32 & 30 & 34 & 36 & 38 \\
$y$      & 11.2 & 12.4 & 12.0 & 14.1 & 13.6 & 15.2 & 14.8 & 17.0 & 16.4 & 18.1 & 19.4 & 19.9 \\
\hline
\end{tabular}
\end{center}

Least squares gives
\[
\hat y=2.442991+0.686916\,x_1+0.239252\,x_2,
\]
with standard errors $0.207583$, $0.135366$, and $0.042101$; the $t$ values are
$11.7687$, $5.0745$, and $5.6828$ on $9$ degrees of freedom; $s=0.143607$ and
$R^2=0.997951$. All of it looks healthy. Then we fit one explanatory variable at a time: $y$ on
$x_1$ gives $1.446019$, $y$ on $x_2$ gives $0.450075$, both nearly twice their multiple
coefficients.

Not because of a miscalculation. Employees and floor area almost determine each other: their
correlation is $r_{12}=0.986801$. Once $x_2$ is in, $x_1$ carries almost no new information, so
its coefficient shrinks to its additional effect only, namely $0.687$ million per employee at
fixed area. The single coefficient $1.446$ mixes the employee effect with the area effect that
usually accompanies it, and the gap between the two is a measure of how far the two variables
overlap.

\begin{definisi}[Variance inflation factor]
\index{variance inflation}
For explanatory variable $j$, let $R_j^2$ be the coefficient of determination of the regression
of $x_j$ on all the other explanatory variables. Its variance inflation factor is
\[
\mathrm{VIF}_j=\frac{1}{1-R_j^2}.
\]
\end{definisi}

\kamus{variance inflation factor}{varians inflasi}{how much the standard error of a coefficient
is blown up by other explanatory variables that carry similar information}

\begin{proposisi}[VIF with two explanatory variables]
\index{collinearity}
If there are only two explanatory variables, $\operatorname{Var}(\hat b_1)=\sigma^2/\bigl(S_{11}
(1-r_{12}^2)\bigr)$, with $S_{11}=\sum_i(x_{i1}-\bar x_1)^2$. Hence the standard error of
$\hat b_1$ is inflated by exactly $\sqrt{\mathrm{VIF}_1}$ compared with the case of mutually
orthogonal explanatory variables.
\end{proposisi}

\begin{proof}[Proof]
For two explanatory variables, inverting the $2\times2$ matrix of the centered block
$X^{\top}X$ gives first element $S_{22}/(S_{11}S_{22}-S_{12}^2)$, and
$S_{12}=r_{12}\sqrt{S_{11}S_{22}}$. Dividing numerator and denominator by $S_{22}$ yields
$\sigma^2/\bigl(S_{11}(1-r_{12}^2)\bigr)$.
\end{proof}

The numbers can be checked: $S_{11}=42.9167$, $r_{12}^2=0.973775$, so $\mathrm{VIF}=38.132232$
and $\sqrt{\mathrm{VIF}}=6.1751$. The standard error of $\hat b_1$ is indeed $0.135366$, while
$\sigma/\sqrt{S_{11}}$ is only $0.021921$; the ratio is $6.1751$, matching the root of the VIF
to four digits.

Collinearity does not damage the forecast; it damages the interpretation. The value
$R^2=0.997951$ shows that the line is still close to the points, and a forecast at $x_1=9$,
$x_2=28$ is still useful. What cannot be trusted is the sentence ``one extra employee raises
sales by $0.687$ million'': its $95$ percent interval $(0.380697;\ 0.993135)$ is wide, and if
the data shift a little that coefficient can change a lot without the line changing much.

\section{Intervals, the $t$ test, and prediction intervals}
\tingkat{advanced}

The next question is the same as in Chapter 20, only the object changes: not a mean, but a
slope. We need one quantity whose distribution is known when the model is correct, then we
invert it into an interval.

In everyday language: we measure how far the observed slope sits from zero, in units of the
slope's own typical wobble.

\kamus{prediction interval}{selang prediksi}{an interval for one day that has not happened yet,
always wider than the interval for the mean because it has to bear the surprise of that day as
well}

\begin{teorema}[The $t$ test and an interval for one coefficient]
Assume (A1) to (A4) and normal errors. With $p$ the number of parameters,
$s^2=\mathrm{SSE}/(n-p)$, and $\operatorname{se}(\hat b_j)=s\sqrt{[(X^{\top}X)^{-1}]_{jj}}$,
we have
\[
\frac{\hat b_j-\beta_j}{\operatorname{se}(\hat b_j)}\sim t_{n-p},
\]
so the $1-\alpha$ interval for $\beta_j$ is
$\hat b_j\pm t_{n-p,1-\alpha/2}\operatorname{se}(\hat b_j)$, and the test of
$H_0\colon\beta_j=0$ rejects at level $\alpha$ when $|t|>t_{n-p,1-\alpha/2}$.
\end{teorema}

\begin{proof}[Proof]
$\hat b_j$ is a linear combination of $y$ with fixed coefficients, so it is normal with
expectation $\beta_j$ and variance $\sigma^2[(X^{\top}X)^{-1}]_{jj}$ by Gauss--Markov. Moreover
$\mathrm{SSE}/\sigma^2\sim\chi^2_{n-p}$ and is independent of $\hat\beta$, because the residual
is perpendicular to the column space while $\hat\beta$ lies entirely inside it. The ratio of a
standard normal variable to the root of a chi-square variable divided by its degrees of freedom
follows a $t$ distribution. The full proof is in Wasserman, \emph{All of Statistics},
Chapter 13.
\end{proof}

For the café data, $\hat b=1.466667$, $s=0.365148$, $S_{xx}=60$, so
$\operatorname{se}(\hat b)=0.047140$. With $t_{7,0.975}=2.364624$,
\[
\hat b\pm t\,\operatorname{se}(\hat b)=1.466667\pm0.111470
=(1.355197;\ 1.578136),
\]
and $t=31.1127$ with a p-value below $10^{-6}$. Rejecting $\beta=0$ here is no achievement: with
residuals as tight as $22$ minutes and an $x$ range this wide, almost any data would reject it.
What deserves reporting is the interval, because it says that the extra overtime per ten orders
lies between $1.36$ and $1.58$ hours.

Now distinguish two questions that are often confused. What is the average overtime on a day
with eighty orders? And what is the overtime on one particular day with eighty orders? The first
answer narrows like $1/\sqrt n$; the second never narrows below the standard deviation of the
residuals, because a single day brings its own surprise.

\begin{proposisi}[Two kinds of forecast variance]
For fixed $x_0$ and a new observation $y_0=x_0^{\top}\beta+\varepsilon_0$ independent of the old
data,
\[
\operatorname{Var}(\hat y_0)=\sigma^2x_0^{\top}(X^{\top}X)^{-1}x_0,
\qquad
\operatorname{Var}(y_0-\hat y_0)=\sigma^2\bigl(1+x_0^{\top}(X^{\top}X)^{-1}x_0\bigr).
\]
\end{proposisi}

\begin{proof}[Proof]
The first already appeared in Gauss--Markov. For the second, write
$y_0-\hat y_0=\varepsilon_0-x_0^{\top}(\hat\beta-\beta)$; the two terms are independent because
$\varepsilon_0$ is independent of the old data, so the variance is the sum of the two
variances.
\end{proof}

For $x_0=8$ in the café data, $x_0^{\top}(X^{\top}X)^{-1}x_0=1/9+9/60=0.261111$. The half-width
of the mean interval is $2.364624\sqrt{0.133333\cdot0.261111}=0.441209$, so
$(19.9588;\ 20.8412)$ with width $0.882$. The half-width of the prediction interval is
$2.364624\sqrt{0.133333\cdot1.261111}=0.969635$, so $(19.4304;\ 21.3696)$ with width $1.939$.
The ratio of widths is $\sqrt{1.261111/0.261111}=2.198$: the prediction interval is more than
twice as wide, and no amount of data can remove that excess. Figure~\ref{fig:dua-selang} shows
both bands over the whole range of $x$.

\begin{figure}[htbp]
\centering
\begin{tikzpicture}
\begin{axis}[
  width=11.4cm, height=6.0cm,
  xmin=0.6, xmax=9.4, ymin=8.2, ymax=25.4,
  axis lines=left,
  xtick={1,2,3,4,5,6,7,8,9},
  tick label style={font=\scriptsize},
  xlabel={online orders (tens)}, ylabel={overtime hours},
  xlabel style={font=\small}, ylabel style={font=\small},
  legend style={font=\scriptsize,draw=layergrey,at={(0.02,0.98)},
    anchor=north west,row sep=0.5pt},
]
\addplot[accent,very thick,domain=0.7:9.3,samples=60]
  {8.666667+1.466667*x+0.863439*sqrt(1.111111+(x-5)^2/60)};
\addplot[accent,very thick,domain=0.7:9.3,samples=60]
  {8.666667+1.466667*x-0.863439*sqrt(1.111111+(x-5)^2/60)};
\addplot[accent,very thick,domain=0.7:9.3,samples=60,forget plot]
  {8.666667+1.466667*x};
\addplot[layerindigo,densely dashed,thick,domain=0.7:9.3,samples=60]
  {8.666667+1.466667*x+0.863439*sqrt(0.111111+(x-5)^2/60)};
\addplot[layerindigo,densely dashed,thick,domain=0.7:9.3,samples=60]
  {8.666667+1.466667*x-0.863439*sqrt(0.111111+(x-5)^2/60)};
\addplot[only marks,mark=*,mark size=1.8pt,layerindigo] coordinates {
  (1,10) (2,12) (3,13) (4,14) (5,16) (6,18) (7,19) (8,20) (9,22)};
\legend{prediction bands, , fitted line, mean bands, }
\end{axis}
\end{tikzpicture}
\caption{The $95$ percent confidence interval for mean overtime (dashed line) and the $95$
percent prediction interval for one new day (thick line), both around the fitted line
$\hat y=8.666667+1.466667x$. At $x=8$ the mean interval
$(19.9588;\ 20.8412)$ has width $0.882$, while the prediction interval
$(19.4304;\ 21.3696)$ has width $1.939$. Both bands widen at the ends of the range, where the
line is less supported by data; the ratio of their widths, $2.198$, does not depend on $n$.}
\label{fig:dua-selang}
\end{figure}

\section{Residual diagnostics}
\tingkat{intermediate}

Every interval above is conditional on the model being right. Hence the first step after fitting
a line is not to compute p-values but to inspect the residuals. In everyday language: we look at
whether what is left over from the line still has a pattern; a residual with a pattern means the
line has the wrong shape, not that chance is at work.

\index{residual}
Four patterns deserve recognition: a curve against $x$, sign of a non-linear relation; a funnel
that widens with $x$, a sign that (A3) fails; one faraway point on its own, a suspect
misrecording; and a pattern by time order, a sign that observations are not independent.

For the café data, the residuals in order are
\[
-0.1333,\ 0.4,\ -0.0667,\ -0.5333,\ 0,\ 0.5333,\ 0.0667,\ -0.4,\ 0.1333 .
\]
There is no evident curvature, and the sum is zero as it should be. What is interesting is their
size: the largest residual is $0.5333$ hours, about $32$ minutes, on the fourth and sixth days,
both near the center of the data. Those points are in fact \emph{not} very influential, because
their leverage is only $0.1278$.

The leverage of point $i$ in a simple regression has a closed form
\[
h_{ii}=\frac1n+\frac{(x_i-\bar x)^2}{S_{xx}},
\]
so on this data $h_{11}=h_{99}=1/9+16/60=0.3778$, while $h_{55}=0.1111$. An end point holds
$3.4$ times the weight of a middle point, and that is not a designer's choice but a consequence
of geometry: a line rotated about the center of the data moves furthest at the ends. Actual
influence combines the size of the residual with leverage, and the usual measure is Cook's
distance
\[
D_i=\frac{e_i^2}{p\,s^2}\cdot\frac{h_{ii}}{(1-h_{ii})^2},
\]
which here is largest, $0.2870$, on the second and eighth days, both with residual $0.4$ and
leverage $0.2611$. Values below one are generally considered not worrying, but that threshold
is a convention, not a law.

One warning that is rarely written: with $n=9$ diagnostics have almost no power. A pattern in
nine points easily arises from chance, and a pattern that is really there is easily invisible.
Diagnostics offer suspicion, not verdict.

\begin{bacakeras}
Two voices.

``The $R$ squared is zero point nine nine. So the model is almost perfect.''

``Perfect for what? The residuals are still twenty-two minutes.''

``But ninety-nine percent of the wobble is already explained.''

``Explained by what? By orders. Not by a kitchen that suddenly slowed down.''

``So a high $R$ squared is not a guarantee.''

``A high $R$ squared only means the line is tight to the points that exist. It knows nothing
about a day that has not come yet.''

Say it once more, slowly: tight to the old points is not a promise about a new day. Repeat it
until it feels natural in the mouth: tight to the old points, not a promise.
\end{bacakeras}

\begin{resep}
\begin{enumerate}
  \item Plot $y$ against $x$ before computing anything, and ask whether a straight line makes
        sense.
  \item Fit the line by least squares, then write the equation complete with the unit of every
        coefficient.
  \item Inspect the residuals against $x$ and against time order; look for curvature, a funnel,
        and isolated points.
  \item Report $s$ together with $R^2$, then an interval for the coefficient that carries the
        main conclusion.
  \item If the aim is to forecast one event, use a prediction interval; if the aim is a mean,
        use a mean interval.
\end{enumerate}
A warning about order: step 1 must not be swapped with step 2. Plotting after fitting the line
makes the eye look for justification, not for faults.
\end{resep}

\section{When the answer is yes or no}
\tingkat{advanced}

So far $y$ has been a number that can be measured. Many important questions are binary: a
customer returns or not, a transaction is fraudulent or not, a patient recovers or not. It is
still tempting to fit an ordinary straight line to $0$ and $1$, then read it as a probability.
Two things go wrong at once, and both can be computed.

The following data: $26$ customers, the distance of their home from the stall $x$ in kilometers,
and $y=1$ if they returned within thirty days. Of the $26$, $15$ returned. A linear regression
on the binary $y$, called the \emph{linear probability model}, gives
$\hat y=1.346154-0.096154x$. At $x=2$ the forecast is $1.153846$: a probability greater than
one, and that is not a rounding error but a defect of form. At $x=20$, still outside the data
range but only seven kilometers outside, the forecast is $-0.576923$.

The second failure is subtler. For binary $y$, the variance of the residual is really $p(1-p)$
and therefore depends on $x$: at $x=2$ its value is $0.007225$, at $x=14$ it is $0.028890$. The
linear probability model assumes one variance for all points, namely $s^2=0.124199$ on this
data, so its standard errors are wrong exactly where it matters most.
Ordinary least squares, which works well when the variance is constant, here relies on an
assumption that plainly does not hold.

\index{logistic regression}
\begin{definisi}[Logistic model]
For $p_i=P(y_i=1\mid x_i)$,
\[
\log\frac{p_i}{1-p_i}=a+b\,x_i,
\qquad\text{that is}\qquad
p_i=\frac{1}{1+e^{-(a+bx_i)}} .
\]
The parameters are estimated by maximum likelihood, not by least squares, so there is no closed
form and the computation is iterative.
\end{definisi}

\kamus{odds ratio}{rasio peluang}{how many times the odds multiply for one additional unit, not
by how many points the probability rises}

For that data, maximum likelihood gives $\hat a=6.315432$ and $\hat b=-0.699945$ with standard
errors $2.383413$ and $0.258594$. The odds ratio is $e^{-0.699945}=0.496613$: each extra
kilometer cuts the odds of returning to about a half, in the sense of a ratio, not a difference
in points. The estimated probabilities are $0.992722$ at $x=2$ and $0.029777$ at $x=14$, both
legitimate as probabilities. The $95$ percent interval for $b$ is
$(-1.233657;\ -0.166233)$; it is entirely negative, so the direction is clear even if the
magnitude is not.

One costly misreading: a logistic coefficient is not the marginal effect on probability. The
value $b=-0.699945$ acts on the log-odds scale, and its effect on probability is largest around
$p=0.5$ and smallest at the ends. A good report states probabilities at several meaningful
values of $x$, not just the coefficient.

\section{Correlation is not causation}
\tingkat{basic}

Now the part most often violated in reports that look tidy. Two groups of employees attend
training of different length, and the number of errors per hundred transactions is recorded.

\begin{center}
\small
\begin{tabular}{l*{6}{r}}
\hline
morning shift & training hours & 1 & 2 & 3 & 4 & 5 \\
              & errors         & 4.0 & 3.6 & 3.2 & 2.9 & 2.5 \\
\hline
night shift   & training hours & 5 & 6 & 7 & 8 & 9 \\
              & errors         & 8.0 & 7.5 & 7.0 & 6.4 & 6.0 \\
\hline
\end{tabular}
\end{center}

Within the morning shift, the slope is $-0.36$ errors per training hour. Within the night shift,
$-0.474286$. Both are negative: longer training goes with fewer errors. But when the two groups
are pooled, the slope is $+0.359036$ with $R^2=0.226121$. The sign flips, and not one of the
numbers is wrong.

The cause is not magic but the addition of two sources:
\[
b_{\text{pooled}}=\frac{\sum_g S_{xy,g}+\sum_g n_g(\bar x_g-\bar x)(\bar y_g-\bar y)}
{\sum_g S_{xx,g}+\sum_g n_g(\bar x_g-\bar x)^2}.
\]
The morning group contributes $S_{xy}=-6.3$ and $S_{xx}=17.5$; the night group $S_{xy}=-8.3$ and
$S_{xx}=17.5$; the between-group term is $6\cdot(-2)(-1.85)+6\cdot(2)(1.85)=44.4$; the
denominator is $83$. Hence
$b_{\text{pooled}}=(-6.3-8.3+44.4)/83=29.8/83=0.359036$. The between-group term $44.4$ defeats
the sum of within-group terms $-14.6$, and that is the whole secret.

The interpretation: training hours are not assigned at random. Night-shift employees, who make
more errors because of working hours and the kinds of transactions, also receive more training,
so what is measured as the effect of training is mostly the difference between shifts.
Figure~\ref{fig:simpson} places everything on one plane.

\begin{figure}[htbp]
\centering
\begin{tikzpicture}
\begin{axis}[
  width=11.4cm, height=6.2cm,
  xmin=0.5, xmax=10.5, ymin=1.6, ymax=8.9,
  axis lines=left,
  xtick={1,2,3,4,5,6,7,8,9,10},
  tick label style={font=\scriptsize},
  xlabel={training hours}, ylabel={errors per 100 transactions},
  xlabel style={font=\small}, ylabel style={font=\small},
  legend style={font=\scriptsize,draw=layergrey,at={(0.98,0.98)},
    anchor=north east,row sep=0.5pt},
]
\addplot[only marks,mark=*,mark size=2pt,layerteal] coordinates {
  (1,4.0) (2,3.6) (3,3.2) (4,2.9) (5,2.5)};
\addplot[only marks,mark=square*,mark size=1.9pt,layerplum] coordinates {
  (6,7.5) (7,7.0) (8,6.4) (9,6.0) (10,5.7)};
\addplot[layerteal,very thick,domain=0.8:5.4,samples=2] {4.326667-0.36*x};
\addplot[layerplum,very thick,domain=5.6:10.2,samples=2]
  {10.32381-0.474286*x};
\addplot[accent,thick,densely dashed,domain=0.8:10.2,samples=2]
  {2.941968+0.359036*x};
\legend{morning, night, pooled}
\end{axis}
\end{tikzpicture}
\caption{Simpson's paradox in the training data. Within each shift the line falls: $-0.36$ errors
per hour for the morning and $-0.4743$ for the night. The pooled line, computed from all twelve
points regardless of shift, instead rises with slope $+0.3590$. The change of sign comes from
the difference between groups, not from one odd point; the between-group term $44.4$ defeats the
within-group term $-14.6$.}
\label{fig:simpson}
\end{figure}

\begin{untukanak}
Two classes take part in a reading contest. Class A reads thin books, many pages, and gets good
marks. Class B reads thick books, fewer pages, and gets poor marks. If all the pupils are pooled,
it looks as if the fewer pages read, the higher the mark. Yet inside each class, the more pages,
the higher the mark.

What is wrong is not the numbers but the pooling. Class B really was given thicker books, and
class B really is weaker at reading. The two facts blend into one deceptive line.

In this chapter we will learn why pooling two groups can reverse the direction of a conclusion,
why the question ``does training reduce errors'' cannot be answered just by plotting all the
points, and why anyone changing policy must know which groups they are really comparing.
\end{untukanak}

\begin{tantangan}
Build a table with two groups, five points each, such that both within-group slopes are positive
but the pooled slope is negative, and the difference between the group means of $y$ is as small
as still possible. How small is that difference, and why can it not be smaller?

Hint 1: the sign of the pooled slope is decided by the between-group term. Make the first group
have small $x$ and small $y$, the second large $x$ and large $y$, then reverse the direction
inside each group.

Hint 2: with five points and within-group slope $+1$ at $S_{xx}=10$, each group contributes
$S_{xy}=+10$ and the between-group term must be less than $-20$. Since
$n_g(\bar x_g-\bar x)(\bar y_g-\bar y)$ is negative only when one of the mean differences has
the opposite sign to the other, the mean difference must pass the point where the two terms
cancel.

If you get stuck, go back to the section Correlation is not causation and recompute
$29.8/83$ step by step.
\end{tantangan}

\section{Honest reporting}
\tingkat{intermediate}

A regression analysis can be done correctly and reported so as to mislead without a single number
being falsified. Four habits close most of that gap.

First, state how many observations and how many parameters, because the reader needs both to
compute the degrees of freedom and to judge whether $s$ can be trusted. On the café data, $n=9$
and $p=2$, so $s$ is built from seven degrees of freedom; hiding that number makes $s$ look
sturdier than it is.

Second, separate the data used to choose the model from the data used to judge it. Choosing the
model form, choosing the variables, and computing the error on the same set always produces
numbers that are too good. At $n=9$ that separation is impossible, and precisely for that reason
the conclusion must be called weak, not reinforced with $R^2=0.992821$.

Third, state how many models were tried. If ten possibilities are tested and the best is
reported, the chance that at least one of them looks meaningful is $1-0.95^{10}=0.401263$ even
though none of the effects is real. That number is not the probability of being wrong, but the
probability of picking something that is a coincidence.

Fourth, report magnitudes with units and intervals, not just stars. The sentence
``$\hat b=0.687$ ($95$ percent: $0.381$ to $0.993$ million per employee)'' can be checked; the
sentence ``significantly affects'' cannot. If the conclusion is about causation, say why:
randomization, or adjustment for the named confounders.

\section{A worked example from end to end}
\tingkat{intermediate}

We repeat the whole chain on the café data, with the intermediate numbers written out, so that
every step can be checked with a calculator.

The means are $\bar x=5$ and $\bar y=16$. The squared deviations are
$S_{xx}=\sum(x_i-5)^2=60$ and $S_{yy}=130$; their cross product is
$S_{xy}=\sum(x_i-5)(y_i-16)=808-9\cdot5\cdot16=88$. Hence $\hat b=88/60=1.466667$ and
$\hat a=16-1.466667\cdot5=8.666667$.

The forecasts at the nine points are $10.1333$; $11.6$; $13.0667$; $14.5333$; $16$; $17.4667$;
$18.9333$; $20.4$; $21.8667$. The residuals are $-0.1333$; $0.4$; $-0.0667$; $-0.5333$; $0$;
$0.5333$; $0.0667$; $-0.4$; $0.1333$, and their sum of squares is $\mathrm{SSE}=0.933333$. With
$n-p=7$, $s^2=0.133333$ and $s=0.365148$ hours.

The standard error of the slope is $s/\sqrt{S_{xx}}=\sqrt{0.133333/60}=\sqrt{0.00222222}=0.047140$.
The statistic is $t=1.466667/0.047140=31.1127$, and $t_{7,0.975}=2.364624$, so the half-width of
the interval is $0.111470$ and the interval is $(1.355197;\ 1.578136)$.

\kertas{Repeat the same computation for one change of data. The ninth day is recorded with $22$
hours of overtime; suppose the correct record is $26$ hours. Recompute $\hat b$, $s$,
$\operatorname{se}(\hat b)$, and its $95$ percent interval, then compare with the table above.
The correct answers: $\hat b=1.733333$, $s=1.306880$, $\operatorname{se}(\hat b)=0.168718$,
interval $(1.3344;\ 2.1323)$. The width is $0.798$ against $0.223$, that is $3.6$ times.
Explain why one point at the end can do that, when it is only one of nine.}

\section{Computational lab}
\tingkat{advanced}

Every number in this chapter can be recomputed in one short file. The following program uses
$X^{\top}X$ directly, so it also checks the algebra of the second section.

\begin{lab}
\begin{verbatim}
import numpy as np
from math import sqrt
from scipy import stats

x = np.arange(1, 10, dtype=float)
y = np.array([10, 12, 13, 14, 16, 18, 19, 20, 22], dtype=float)
n = len(y)
X = np.column_stack([np.ones(n), x])

beta = np.linalg.solve(X.T @ X, X.T @ y)
fit = X @ beta
e = y - fit
s2 = (e ** 2).sum() / (n - 2)
sst = ((y - y.mean()) ** 2).sum()

print("X'X =", (X.T @ X).tolist(), " X'y =", (X.T @ y).tolist())
print("beta = %.6f %.6f" % tuple(beta))
sse = (e ** 2).sum()
print("SSE = %.6f  R2 = %.6f  s = %.6f"
      % (sse, 1 - sse / sst, sqrt(s2)))
print("residuals =", np.round(e, 4).tolist())

cov = s2 * np.linalg.inv(X.T @ X)
se = np.sqrt(np.diag(cov))
tq = stats.t.ppf(0.975, n - 2)
print("se = %.6f %.6f  t = %.4f %.4f"
      % (se[0], se[1], beta[0] / se[0], beta[1] / se[1]))
lo, hi = beta[1] - tq * se[1], beta[1] + tq * se[1]
print("CI95 slope = (%.6f, %.6f)" % (lo, hi))

for x0 in (5.0, 8.0):
    c = np.array([1.0, x0])
    y0 = c @ beta
    v = s2 * (c @ np.linalg.inv(X.T @ X) @ c)
    print("x0=%.0f  yhat=%.4f  mean +/-%.4f  prediction +/-%.4f"
          % (x0, y0, tq * sqrt(v), tq * sqrt(v + s2)))

h = np.einsum("ij,jk,ik->i", X, np.linalg.inv(X.T @ X), X)
cook = e ** 2 / (2 * s2) * h / (1 - h) ** 2
print("h_ii =", np.round(h, 4).tolist())
print("D_Cook =", np.round(cook, 4).tolist())

x1 = np.array([6, 7, 7, 8, 8, 9, 9, 10, 10, 11, 12, 12], dtype=float)
x2 = np.array([18, 22, 20, 26, 24, 28, 26, 32, 30, 34, 36, 38], dtype=float)
yb = np.array([11.2, 12.4, 12.0, 14.1, 13.6, 15.2, 14.8, 17.0, 16.4, 18.1,
               19.4, 19.9])
Xb = np.column_stack([np.ones(12), x1, x2])
bb = np.linalg.solve(Xb.T @ Xb, Xb.T @ yb)
eb = yb - Xb @ bb
s2b = (eb ** 2).sum() / 9
seb = np.sqrt(np.diag(s2b * np.linalg.inv(Xb.T @ Xb)))
r12 = np.corrcoef(x1, x2)[0, 1]
print("multiple beta = %.4f %.4f %.4f  s = %.4f"
      % (bb[0], bb[1], bb[2], sqrt(s2b)))
print("multiple se   = %.4f %.4f %.4f" % tuple(seb))
print("r12 = %.4f  VIF = %.2f  sqrt VIF = %.2f"
      % (r12, 1 / (1 - r12 ** 2), 1 / sqrt(1 - r12 ** 2)))


def ols(xx, yy):
    A = np.column_stack([np.ones(len(xx)), xx])
    return np.linalg.solve(A.T @ A, A.T @ yy)[1]


xp = np.array([1, 2, 3, 4, 5, 6], dtype=float)
yp = np.array([4.0, 3.6, 3.2, 2.9, 2.5, 2.2])
xm = np.array([5, 6, 7, 8, 9, 10], dtype=float)
ym = np.array([8.0, 7.5, 7.0, 6.4, 6.0, 5.7])
print("Simpson: morning %.4f  night %.4f  pooled %.4f"
      % (ols(xp, yp), ols(xm, ym), ols(np.r_[xp, xm], np.r_[yp, ym])))

xd = np.array([2, 2, 3, 3, 4, 4, 5, 5, 6, 6, 7, 7, 8, 8, 9, 9, 10, 10,
               11, 11, 12, 12, 13, 13, 14, 14], dtype=float)
yd = np.array([1, 1, 1, 1, 1, 1, 1, 1, 1, 1, 1, 0, 1, 0, 1, 0, 1, 0, 0,
               1, 0, 0, 0, 0, 0, 0], dtype=float)
A = np.column_stack([np.ones(len(yd)), xd])
lpm = np.linalg.solve(A.T @ A, A.T @ yd)
print("LPM: %.4f %+.4f  prediction at x=2: %.4f  at x=20: %.4f"
      % (lpm[0], lpm[1], lpm @ [1, 2], lpm @ [1, 20]))


def nll(th):
    z = A @ th
    return float(np.sum(np.log1p(np.exp(-z)) + (1 - yd) * z))


g = np.array([0.0, 0.0])
for _ in range(40):
    p = 1 / (1 + np.exp(-(A @ g)))
    grad = A.T @ (p - yd)
    Hm = A.T @ (A * (p * (1 - p))[:, None])
    g = g - np.linalg.solve(Hm, grad)
se_g = np.sqrt(np.diag(np.linalg.inv(Hm)))
p = 1 / (1 + np.exp(-(A @ g)))
print("logit: %.4f %+.4f  se %.4f %.4f" % (g[0], g[1], se_g[0], se_g[1]))
print("deviance = %.4f  p(x=2) = %.4f  p(x=14) = %.4f"
      % (2 * nll(g), p[0], p[-1]))
print("odds ratio per km = %.4f" % np.exp(g[1]))
\end{verbatim}
\end{lab}

The output, exactly as it comes:

\begin{verbatim}
X'X = [[9.0, 45.0], [45.0, 285.0]]  X'y = [144.0, 808.0]
beta = 8.666667 1.466667
SSE = 0.933333  R2 = 0.992821  s = 0.365148
residuals = [-0.1333, 0.4, -0.0667, -0.5333, 0.0, 0.5333, 0.0667, -0.4, 0.1333]
se = 0.265274 0.047140  t = 32.6706 31.1127
CI95 slope = (1.355197, 1.578136)
x0=5  yhat=16.0000  mean +/-0.2878  prediction +/-0.9101
x0=8  yhat=20.4000  mean +/-0.4412  prediction +/-0.9696
h_ii = [0.3778, 0.2611, 0.1778, 0.1278, 0.1111, 0.1278, 0.1778, 0.2611, 0.3778]
D_Cook = [0.0651, 0.287, 0.0044, 0.1792, 0.0, 0.1792, 0.0044, 0.287, 0.0651]
multiple beta = 2.4430 0.6869 0.2393  s = 0.1436
multiple se   = 0.2076 0.1354 0.0421
r12 = 0.9868  VIF = 38.13  sqrt VIF = 6.18
Simpson: morning -0.3600  night -0.4743  pooled 0.3590
LPM: 1.3462 -0.0962  prediction at x=2: 1.1538  at x=20: -0.5769
logit: 6.3154 -0.6999  se 2.3834 0.2586
deviance = 17.8843  p(x=2) = 0.9927  p(x=14) = 0.0298
odds ratio per km = 0.4966
\end{verbatim}

Two things must be read from that output. First, $\mathrm{X}^{\top}\mathrm{X}$ and
$\mathrm{X}^{\top}y$ are printed exactly as computed by hand, and $\hat\beta$ comes out of
solving the system; agreement of this kind is not proof that the theory is true, but it removes
arithmetic slips, which are the most frequent cause. Second, the last three parts show why
interpretation is more fragile than computation: $0.6869$ against $1.4460$ in the single
regression, VIF $38.13$, and a pooled slope whose sign reverses. Not one of the three is a
programming error.

\begin{untukdosen}
The teaching order I suggest: start from the table of nine days, ask the class to guess the slope
in two ways, by the difference of the ends and by ``the closest line''. The two answers differ
only slightly, and that is exactly the entry point: the real question is not the number but its
wobble. Only after the class is restless, move to the normal equations and show that
$\sum x_ie_i=0$ is a geometric statement, not an algebraic trick.

The two misconceptions that appear most often. First, students read $R^2$ as ``the probability
that the forecast is right'', then are surprised when the prediction interval is two hours wide
at $R^2=0.99$; the symptom shows when they cannot explain why $s$ does not shrink along with
$R^2$. Second, they interpret multiple coefficients like single ones, even though on the
twelve-branch data the number moves from $1.446$ to $0.687$ only because another explanatory
variable enters.

A blackboard activity that works: have two groups compute the slope on the training data, one
group using all twelve points, one separating the shifts. Write both results side by side, then
ask each group to defend theirs. The sign that the class is ready to move on: they can name one
situation in which a multiple coefficient changes sign, and can explain why the forecast stays
good.
\end{untukdosen}

\section{Summary}
\tingkat{basic}

\begin{ringkasan}
\begin{itemize}
  \item Regression fits the line that minimizes the sum of squared residuals. By the normal
        equations the residual is perpendicular to every column of $X$; on the café data
        $\hat b=88/60=1.466667$ and $\hat a=8.666667$.
  \item A coefficient is meaningful only together with its unit: $1.466667$ hours per ten orders
        equals $8.8$ minutes per order. Changing the scale of $x$ changes the coefficient and its
        standard error, not $t$ or $R^2$.
  \item $R^2$ is relative, $s$ is absolute. On the café data $R^2=0.992821$ but $s=0.365148$
        hours, that is $21.9$ minutes.
  \item Under (A1)--(A4) Gauss--Markov guarantees that least squares is BLUE with
        $\operatorname{Var}(\hat\beta)=\sigma^2(X^{\top}X)^{-1}$; that guarantee lapses when the
        error variance is not constant or the explanatory variables are collinear.
  \item Collinearity does not damage the forecast but it damages the interpretation: on the
        twelve branches, $r_{12}=0.986801$ gives $\mathrm{VIF}=38.13$ and standard errors $6.18$
        times larger, while $R^2$ stays at $0.997951$.
  \item A prediction interval is always wider than a mean interval; at $x_0=8$ the ratio is
        $2.198$ and does not shrink with $n$. The sign of a slope can reverse when two groups are
        pooled: $-0.36$ and $-0.474286$ within groups, $+0.359036$ once pooled.
\end{itemize}
\end{ringkasan}

\section{Exercises}
\tingkat{intermediate}

\begin{latihan}
\begin{enumerate}
  \item Evaluate this claim: ``Because $R^2=0.99$, this line can forecast a single day with an
        error of at most twenty minutes.'' True, false, or not enough evidence? Name the part
        that is right, the part that is wrong, and the additional number that settles the
        assessment.
  \item Prove that with a column of ones in $X$ we have $\sum_i e_i=0$ and
        $\sum_i x_ie_i=0$, then explain why both fail when the line is forced through the
        origin.
  \item Prove $\operatorname{Var}(\hat b)=\sigma^2/S_{xx}$ from
        $\hat b=\sum_i(x_i-\bar x)y_i/S_{xx}$ without matrix formulas.
  \item Compute $\hat a$, $\hat b$, $s$, $\operatorname{se}(\hat b)$, and the $95$ percent
        interval for $(1,2)$, $(2,2.5)$, $(3,4)$, $(4,3.5)$, $(5,5)$, then check with Python. Why
        is the interval much wider than the interval for the café data?
  \item Compute the $95\%$ prediction interval for $x_0=2$ and $x_0=9$ on the café data.
        Why do the widths differ when $n$ is the same, and why does the interval never
        shrink to zero when $n$ is increased?
  \item Show that adding explanatory variables cannot raise SSE, then explain why
        $R^2_{\text{adj}}$ can fall. Compute $R^2_{\text{adj}}$ for the café data if $p$ rises
        from $2$ to $3$ at the same $R^2$.
  \item For $S_{11}=50$, $S_{22}=200$, $r_{12}=0.9$, and $\sigma^2=2$, compute the VIF and the
        standard error of $\hat b_1$ with and without $x_2$ in the model.
  \item Build a table with two groups with within-group slope $-0.2$ and pooled slope $+0.4$,
        then check the result with the decomposition formula.
  \item Rewrite this sentence so that it is honest: ``the best model of the twelve we tried gives
        $p=0.03$, so the effect is real''.
\end{enumerate}

Hints for the three hardest items. Item 3: write
$\hat b-\beta=\sum_i(x_i-\bar x)\varepsilon_i/S_{xx}$, take the variance, use independence and
$\operatorname{Var}(\varepsilon_i)=\sigma^2$. Item 6: compare the column spaces of the two
models, then remember the divisor $n-p$ that shrinks. Item 7: $\mathrm{VIF}=1/(1-0.81)=5.263158$,
so the standard error is $\sqrt{5.263158}$ times larger.
\end{latihan}

\begin{untukpeneliti}
What is genuinely proved here: the normal equations as a necessary and sufficient condition for
the minimum, together with uniqueness when $X$ has full column rank; the properties of the hat
matrix and $\operatorname{tr}H=p$; unbiasedness and $\operatorname{Var}(\hat\beta)=\sigma^2
(X^{\top}X)^{-1}$; the Gauss--Markov inequality within the class of linear unbiased estimators;
the decomposition $\mathrm{SST}=\mathrm{SSR}+\mathrm{SSE}$ when a column of ones is present; the
non-monotonicity of $R^2$ under added columns; the $t$ distribution for the coefficient ratio
under normal errors; and the variance of the forecast error.

What is merely popular: that $R^2$ measures the quality of a model, that multiple coefficients
always represent effects, and that the thresholds $\mathrm{VIF}>10$ or $D_i>1$ are derived
results. Primary sources: Seber and Lee, \emph{Linear Regression Analysis}, Chapter 3; Greene,
\emph{Econometric Analysis}, Chapter 3; Wasserman, \emph{All of Statistics}, Chapter 13;
Simpson (1951) for the reversal of direction in contingency tables; Hosmer, Lemeshow, and
Sturdivant, \emph{Applied Logistic Regression}, Chapters 1--3.

Limits of validity: every interval is conditional on the model being right, the errors being
independent, and the variance being constant. On grouped data, time series, or repeated
observations, independence fails and the reported standard errors are too small. Open question:
how far post-selection intervals depart from their nominal level, and how influence diagnostics
are constructed when $n$ is far smaller than $p$.
\end{untukpeneliti}

\peta{Chapter 21 in Wasserman, \emph{All of Statistics}: Chapter 13 (Linear and Logistic
Regression: least squares, the $t$ distribution for coefficients, and logistic regression) and
Chapter 6 (Models, Statistical Inference, and Learning: the interpretation of intervals and
standard errors). The geometry of least squares as an orthogonal projection is set out in
Chapter 10 of this book, where the normal equations first appear as a consequence of
orthogonality, and is used again in Chapter 24 for regularization. Confidence intervals,
p-values, and the interpretation of levels are taken from Chapter 20 and not repeated here. The
Gauss--Markov theorem and its conditions are referred to Seber and Lee Chapter 3 and Greene
Chapter 3, because Wasserman does not prove it. Simpson's paradox is referred to Simpson's
(1951) paper and used again in Chapter 25 when discussing data separation and evaluation
leakage, and in Chapter 26 when discussing causal claims that must be proved, not inferred from
correlation.}

\chapter{Floating-Point Arithmetic and Error Analysis}

A program solves $x^{2}+10^{8}x+1=0$ with the quadratic formula we have used since
school, $x=(-b\pm\sqrt{b^{2}-4ac})/(2a)$. For $a=1$, $b=10^{8}$, $c=1$ it reports
the roots $-10^{8}$ and $-7.450580596923828\times10^{-9}$. The first root is correct
to the last digit. The second is wrong by a quarter: the true value is $-10^{-8}$, and
its relative error is $0.255$. The formula is not mistyped. Neither is
the machine broken. The coefficients are exact, the discriminant is exact, and the
subtraction is even exact according to the theorem we shall prove below.

What collapses is the algebraic form. The numbers $b=10^{8}$ and $\sqrt{b^{2}-4ac}$
differ by only one ulp, that is one smallest step the machine can still tell apart at
that scale: the difference is $1.49\times10^{-8}$, or $1.5\times10^{-16}$ times the
size of the numbers themselves. A difference that small means nothing as long as we
regard numbers of size $10^{8}$. As soon as the two are subtracted, the result has
size $10^{-8}$ and that difference becomes the whole answer: what remains is error,
not a number.

Chapters 4 and 5 build limits, continuity, and derivatives over $\mathbb{R}$ whose
every point is exact. Chapter 20 uses the variance formula and the square root under
the same assumption. A computer does not work in $\mathbb{R}$. It works in a finite
set of numbers that has rules of its own: every operation is rounded, addition is not
associative, and two algebraically equivalent routes can give answers far apart. This
chapter builds a way of reading those rules, of measuring the damage before the program
runs, and of choosing a form of the formula that does not collapse.

\begin{untukanak}
Imagine a digital kitchen scale whose display holds only five digits. We weigh a sack
of rice together with its container, then weigh the empty container, then subtract the
two numbers to get the weight of the rice. If container and rice each weigh about
twenty kilograms, the display can only show numbers to the nearest tens of grams. The
difference of two large numbers that are nearly equal is correct as arithmetic, but its
last digits no longer mean anything. Rice that actually weighs eight hundred grams can
show up as seven hundred grams, or even as an absurd number.

There is a second, subtler problem. The same scale cannot weigh a single grain of rice:
the needle does not move at all. If we pour in a million grains one at a time, not one
of them is recorded, even though a million grains have weight. A calculating machine
works in exactly that way: some numbers are too small to change anything, and some
subtractions throw away almost every digit.

The question: if the machine is like that, how can we know when an answer can still be
trusted? In this chapter we shall learn why that error is not a random event but
something whose bound can be computed before a single line of the program runs.
\end{untukanak}

\section{Numbers without a representative}
\tingkat{basic}

Between $0$ and $1$ there are infinitely many real numbers. Inside a computer, all
numbers are stored in 32 or 64 bits, so only finitely many are available. The
compromise adopted from the first machines down to the IEEE-754 standard is binary
scientific notation with fixed field widths: one sign, one exponent, one mantissa.

\kamus{floating-point system}{sistem floating-point}{Binary scientific notation with a
fixed number of digits, like writing $1.01\times2^{7}$ in limited room.}

\begin{definisi}[Floating-point number system]
Let $\beta=2$ be the base, $p$ the precision, and $e_{\min}\le e\le e_{\max}$ the
exponent bounds. A normal number is
\[
  x=\pm\left(1.d_{1}d_{2}\dots d_{p-1}\right)_{2}\times2^{e},
\]
with $d_{i}\in\{0,1\}$. Subnormal numbers use $e=e_{\min}$ with leading digit
$0$.\index{subnormal}\index{floating-point arithmetic} We call the set of all
representable numbers, together with zero, the two infinities, and NaN,
$\mathbb{F}$.
\end{definisi}

The leading digit is always $1$ for a normal number, so that digit is not stored;
that is why the mantissa field has $p-1$ bits. One bit saved doubles the precision:
23 stored bits give 24 significant
bits.



\begin{teorema}[Number of machine numbers]
For every exponent $e$, the number of normal numbers in the interval
$\left[2^{e},2^{e+1}\right)$ is $2^{p-1}$. The set $\mathbb{F}$ is finite, hence
countable.
\end{teorema}

\begin{proof}
On that interval every number has the form $M\cdot2^{e-p+1}$ with $M$ an integer
and $2^{p-1}\le M<2^{p}$. The number of such $M$ is exactly $2^{p}$, minus
$2^{p-1}$, that is $2^{p-1}$. There are finitely many allowed exponents, so
$\mathbb{F}$ is a finite union of finite sets.
\end{proof}

For binary64, $2^{p-1}=2^{52}=4\,503\,599\,627\,370\,496$. That is a lot, but Chapter 2
already showed that $\mathbb{R}$ is uncountable. So almost every real number has no
representative: representation always means rounding, and rounding always means
error. The question is not whether there is error, but how large it is and where it
goes.

\begin{bacakeras}
Two voices, read alternately with a friend. Tap the table every time the word ``piece''
is said.

``One divided by three, how much?'' ``Zero point three three three, on and on, without
end.'' ``Does the machine write it out whole?'' ``No. It stops where there is room.''
``So what is stored?'' ``A piece. Not the value, the piece.''
``And if the pieces are added?'' ``What grows is the pieces, not the value.''
``And if the pieces are subtracted?'' ``The leading digits are used up, dust is what
remains.''

Repeat the three middle lines, faster and faster, with the same tapping: a piece is not
the value, a piece is not the value, a piece is not the value.

The second voice closes quietly: the machine does not lie, it only cuts. We are the ones
who go wrong if we read the piece as the value.

The line the reader must repeat, once, in a clear voice: a computer stores pieces, not
values, so every computed result must be checked, not trusted.
\end{bacakeras}

The IEEE-754 standard fixes two widths that are used everywhere. This table contains all
the numbers; nothing needs memorizing except $p$ and the bias.

\begin{center}
\begin{tabular}{lcc}
\hline
 & binary32 & binary64 \\
\hline
sign / exponent / mantissa bits & $1/8/23$ & $1/11/52$ \\
precision $p$ & $24$ & $53$ \\
exponent bias & $127$ & $1023$ \\
$e_{\min},\ e_{\max}$ & $-126,\ 127$ & $-1022,\ 1023$ \\
$\mathrm{eps}=2^{-p+1}$ & $1.1920929\times10^{-7}$ & $2.2204460\times10^{-16}$ \\
$u=2^{-p}$ & $5.9604645\times10^{-8}$ & $1.1102230\times10^{-16}$ \\
smallest normal & $1.1754944\times10^{-38}$ & $2.2250739\times10^{-308}$ \\
smallest subnormal & $1.4012985\times10^{-45}$ & $4.9406565\times10^{-324}$ \\
largest & $3.4028235\times10^{38}$ & $1.7976931\times10^{308}$ \\
\hline
\end{tabular}
\end{center}

Exponents are stored with a bias so that no separate sign is needed: a stored value of
$123$ means an actual exponent of $123-127=-4$. This trick makes the bit patterns from
zero up to positive infinity increase monotonically, so that two positive numbers can be
compared as integers.

Let us take apart three bit patterns that Python actually produces. In binary32,
$0.1$ prints as \texttt{0 01111011 10011001100110011001101} with
$e=123-127=-4$; $6$ becomes \texttt{0 10000001 10000000000000000000000}, that is
$1.5\times2^{2}$; and $-12.375$ becomes
\texttt{1 10000010 10001100000000000000000}, that is $-1.546875\times2^{3}$.
Two boundary cases complete the picture:

\begin{center}
\begin{tabular}{ll}
\hline
smallest subnormal & \texttt{0 00000000 00000000000000000000001} \\
largest & \texttt{0 11111110 11111111111111111111111} \\
\hline
\end{tabular}
\end{center}

The first has the value $2^{-149}$: all its exponent bits are zero and only the last
mantissa bit is on. The second has the value $3.4028235\times10^{38}$, with stored
exponent $254$ and a full mantissa.

A more readable form for binary64 is provided by \texttt{float.hex()}. Here is the result for several numbers:

\begin{center}
\begin{tabular}{ll}
\hline
$0.1$ & \texttt{0x1.999999999999ap-4} \\
$6.0$ & \texttt{0x1.8000000000000p+2} \\
$10^{-8}$ & \texttt{0x1.5798ee2308c3ap-27} \\
$5\times10^{-324}$ & \texttt{0x0.0000000000001p-1022} \\
\hline
\end{tabular}
\end{center}

That form shows three things. The digits after the hexadecimal point are always
thirteen, and $13\times4=52$, the width of the binary64 mantissa. The last digit of
$0.1$ is \texttt{a}, that is $1010$, a rounding up from the repeating pattern
$1001$. The last row has a zero before the point: the mark of a subnormal number.

Special values complete this system: $\pm\infty$ for results beyond
$e_{\max}$, signed zero, subnormal numbers below $2^{e_{\min}}$, and NaN
(\emph{not a number}) for indeterminate forms such as $0/0$. One property of NaN
deserves remembering because it often surprises: it is not equal to itself, so that
the comparison \texttt{x == x} is false precisely when $x$ is NaN.\index{NaN}

\section{The only bridge: rounding}
\tingkat{basic}

No operation other than rounding connects $\mathbb{R}$ with $\mathbb{F}$. That is why
the whole of error analysis rests on one question: how far is the nearest machine number
from the original number, and how does that distance scale.

IEEE-754 lays down the rule \emph{round to nearest, ties to even}: choose the nearest
machine number, and if two candidates are equally far, choose the one whose last
mantissa digit is even. The second rule is no theoretical luxury. Without that rule,
rounding at a tie point always leans one way, and the error of a million operations
piles up in the same direction instead of canceling out.

Tie points really do occur, and we can create one deliberately. In
binary64, the distance between $1$ and the number after it is $2^{-52}$. Hence
$1+2^{-53}$ falls exactly halfway, and the even rule chooses $1$ because its mantissa
ends in zero. Python gives

\begin{center}
\texttt{1 + 2.0**-53 = 1.0}, whereas \texttt{1 + 3*2.0**-53 = 1.0000000000000004}.
\end{center}

The first case rounds down, the second up, although both are exactly halfway. Only the
parity of the last digit distinguishes them. This is why
$1+10^{-16}=1$ in binary64: not because $10^{-16}$ is ignored carelessly,
but because there is no room to store it.

\kamus{machine epsilon}{epsilon mesin}{The distance between $1$ and the next machine
number: the size of one step at scale one.}\index{machine epsilon}

Two numbers are often confused, and the confusion changes an error bound by a factor of
two. The machine epsilon $\mathrm{eps}=2^{-p+1}$ is the distance between $1$ and the
number after it, namely $2^{-52}$ for binary64. The unit roundoff
$u=2^{-p}=2^{-53}$ is half that distance, and only $u$ appears in error
bounds. Python calls the first \texttt{numpy.finfo(float).eps}; the second we must
compute ourselves.

The distance between numbers is not uniform. Near $1$ it is $2^{-52}$, near
$10^{10}$ it is $1.9\times10^{-6}$. What is uniform is that distance relative
to the size of the number, and that is the content of the following lemma.

\begin{lemma}[Ulp]
If $2^{e}\le|x|<2^{e+1}$, then the distance between $x$ and the next machine number is
$\operatorname{ulp}(x)=2^{e-p+1}$.
\end{lemma}

\begin{proof}
On that interval all normal numbers have the form $M\cdot2^{e-p+1}$ with $M$ an integer
and $2^{p-1}\le M<2^{p}$; the number of such $M$ is $2^{p-1}$ and the gap between two
consecutive $M$ is one unit $2^{e-p+1}$. The next interval uses $e+1$ so that the gap
doubles, and that is the only point where the density changes.
\end{proof}

\kamus{ulp}{unit in the last place}{The smallest step the machine can tell apart at the
scale of a number.}\index{ulp}

The lemma explains two things at once.
$\operatorname{ulp}(1)=2^{-52}=2.220446049250313\times10^{-16}$, whereas
$\operatorname{ulp}(10^{10})=2^{-52}\cdot2^{33}\approx1.9\times10^{-6}$. So
adding $1$ to $10^{16}$ changes nothing in binary64, because
$\operatorname{ulp}(10^{16})=2$: there really is no machine number between the two.

Figure~\ref{fig:titikmesin} shows the consequence of that lemma over the first four
octaves: each octave holds $2^{23}$ points, but its width doubles, so that the
distance between points doubles as well.

\begin{figure}[htbp]
\centering
\begin{tikzpicture}[x=2.1cm,y=1cm,every node/.style={font=\footnotesize}]
  \draw[ink,thick,->] (-0.55,0) -- (4.30,0);
  \draw[ink,thick,fill=white] (-0.42,0) circle (1.7pt);
  \draw[ink,thick] (-0.48,-0.09) -- (-0.36,0.09);
  \foreach \x in {0,0.125,...,0.875} {\draw[ink] (\x,0.055) -- (\x,-0.055);}
  \foreach \x in {1,1.25,...,1.875} {\draw[ink] (\x,0.055) -- (\x,-0.055);}
  \foreach \x in {2,2.5,...,3.0} {\draw[ink] (\x,0.055) -- (\x,-0.055);}
  \foreach \x in {3,4} {\draw[ink] (\x,0.055) -- (\x,-0.055);}
  \foreach \x/\lab in {0/{$2^{0}$},1/{$2^{1}$},2/{$2^{2}$},3/{$2^{3}$},4/{$2^{4}$}} {
    \draw[ink,thick] (\x,0.14) -- (\x,-0.14);
    \node[below=1pt] at (\x,-0.14) {\lab};
  }
  \draw[{Latex}-{Latex},accent] (0.125,-0.62) -- (0.25,-0.62) node[midway,below] {$d$};
  \draw[{Latex}-{Latex},accent] (1.25,-0.62) -- (1.50,-0.62) node[midway,below] {$2d$};
  \draw[{Latex}-{Latex},accent] (2.50,-0.62) -- (3.00,-0.62) node[midway,below] {$4d$};
  \draw[{Latex}-{Latex},accent] (3.00,-0.62) -- (4.00,-0.62) node[midway,below] {$8d$};
  \node[align=center,text width=3.0cm] at (1.95,0.92)
    {each octave holds $2^{23}$ points, twice the width of the one before};
  \node[align=center,text width=3.0cm] at (0.40,-1.32)
    {leftward toward zero: the gaps shrink};
\end{tikzpicture}
\caption{The first four octaves of binary32 machine numbers. Each octave holds $2^{23}$
points, but its width is twice that of the octave before, so the spacing between points
doubles: $d=2^{-23}$ on $[1,2)$, $2d$ on $[2,4)$, $4d$ on $[4,8)$, and $8d$ on
$[8,16)$. The tail does not hold fewer points, only sparser ones; toward zero the
spacing shrinks down to the subnormal limit.}
\label{fig:titikmesin}
\end{figure}

There is one subtlety worth naming before we use the ulp as a unit of error. The largest
rounding error is half an ulp, but expressed as a relative error that figure is not one
number but an interval. For $x$ exactly equal to a power of two, half an ulp is
$2^{-53}$ relative; for $x$ that nearly reaches the next power of two, the value
approaches $2^{-54}$. So the relative error of one rounding can swing by a factor of
two. Goldberg named that swing \emph{wobble}, and that is why the bound $u$ is always
used as a worst case, not as a typical value.

\begin{teorema}[Relative rounding error]
For every $x\in\mathbb{R}$ in the normal range, correct rounding satisfies
$\operatorname{fl}(x)=x(1+\delta)$ with $|\delta|\le u=2^{-p}$.
\end{teorema}

\begin{proof}
Take $2^{e}\le|x|<2^{e+1}$. The nearest machine number is at most half an
ulp away, that is $2^{e-p}$, so that
$|\operatorname{fl}(x)-x|\le2^{e-p}=u\cdot2^{e}\le u|x|$. Dividing by $|x|$
gives $|\delta|\le u$, and the sign of $\delta$ follows the direction of rounding.
\end{proof}

The form $x(1+\delta)$ is the beginning of everything: it turns the question ``how
wrong'' into ``how many factors $(1+\delta)$ pile up''. The operations
$+,-,\times,\div,\sqrt{\phantom{x}}$ in IEEE-754 are required to return the correctly
rounded exact result, so every operation contributes one such factor
and no more. This guarantee is what makes error analysis possible.
\section{Error, condition number, and cancellation}
\tingkat{intermediate}

Now we have a unit. What we do not yet have is a measure of how much the error of a
computation exceeds the error of its input. Two problems can both receive input that is
wrong by $10^{-16}$ and come out with an answer that is wrong by $10^{-16}$ or by
$10^{-1}$, depending on the problem itself.

\kamus{condition number}{bilangan kondisi}{The multiplier showing how many times a
computation amplifies the error of its input.}\index{condition number}

\begin{definisi}[Absolute, relative, and backward error]
Let $\tilde{x}$ be an approximation to $x$. Its absolute error is $|\tilde{x}-x|$, its
relative error is $|\tilde{x}-x|/|x|$ for $x\neq0$. If $\tilde{x}$ is the exact answer
to a problem with data $x+\Delta x$, then $\|\Delta x\|/\|x\|$ is called the backward
error, and an algorithm that always gives a small backward error is called
\emph{backward stable}.\index{relative error}\index{backward stable}
\end{definisi}

The last distinction is the most useful in practice. A backward stable algorithm may be
badly wrong, provided its answer is right for a slightly different problem: the error
belongs to the problem. A large backward error is the mistake of the person writing the
code.

For a function of one variable, that amplification is computed directly with the
derivative: the condition number is the slope of the function at that point, multiplied
by the size of the point, divided by the size of its value.

\begin{teorema}[Condition number]
Let $f$ be differentiable on an interval containing $x$ and $f(x)\neq0$. If
$\tilde{x}=x(1+\delta)$ with $|\delta|\le u$, then
\[
  \frac{|f(\tilde{x})-f(x)|}{|f(x)|}\;\le\;\kappa_{f}(x)\,u+O(u^{2}),
  \qquad \kappa_{f}(x)=\frac{|x\,f'(x)|}{|f(x)|}.
\]
\end{teorema}

\begin{proof}
By the mean value theorem there is $\xi$ between $x$ and $\tilde{x}$ with
$f(\tilde{x})-f(x)=f'(\xi)(\tilde{x}-x)$. Since $|\tilde{x}-x|\le u|x|$ and $f'$ is
continuous at $x$, we have $|f'(\xi)|=|f'(x)|+O(u)$. Hence
$|f(\tilde{x})-f(x)|\le u|x|\,|f'(x)|+O(u^{2}|x|)$. Dividing by $|f(x)|$ gives the
claimed inequality.
\end{proof}

The number $\kappa_{f}$ does not depend on the machine, only on the problem. For
$f(x)=\sqrt{x}$ it is $1/2$: a square root never makes the relative error worse.
For $f(x)=\log x$ it is $1/|\log x|$, which is large when $x$ is near $1$ because
there $\log x$ is almost zero. For $f(x)=x-c$ with $c$ a large constant, it is
$|x|/|x-c|$, which grows without bound as $x$ approaches $c$. That last pattern is
what becomes the source of the whole trouble in this chapter.

Subtracting two numbers that are close together is the most dangerous operation in
machine arithmetic, and the danger is often blamed on the wrong thing. A common
misunderstanding runs: subtraction loses digits, so the operation of subtraction is
itself inaccurate. The truth is subtler and more reassuring: when the two numbers are
close together, the subtraction is in fact exact.

\begin{teorema}[Sterbenz]
If $a,b\in\mathbb{F}$ and $b/2\le a\le2b$ with $b>0$, then $a-b$ is represented exactly
in $\mathbb{F}$.
\end{teorema}

\begin{proof}
Let $2^{e}\le b<2^{e+1}$. Every machine number on that interval is an integer multiple of
$2^{e-p+1}$, so $a$ is as well, and hence $a-b$ is a multiple of $2^{e-p+1}$. From
$b/2\le a\le2b$ we get $0\le a-b\le b/2<2^{e}$. A multiple of $2^{e-p+1}$ that
lies in $[0,2^{e})$ has the form $M\cdot2^{e-p+1}$ with $0\le M<2^{p}$, and
all of them are represented with $p$ digits. Hence the result is exact.
\end{proof}

So when the result of a subtraction looks bad, what is bad is the input. The digits that
``disappear'' had already disappeared before the subtraction was carried out: $a$ and $b$
are each known only to relative error $u$, and an absolute error of size $u|a|$ does not
shrink along with the result when the two numbers cancel each other. Only the size of the
result shrinks, so that the relative error explodes.

\begin{proposisi}[Condition number of subtraction]
Let $x=a-b\neq0$ and let the machine work with $\tilde{a}=a(1+\delta_{a})$,
$\tilde{b}=b(1+\delta_{b})$, $|\delta_{a}|,|\delta_{b}|\le u$. Then
\[
  \frac{|\operatorname{fl}(\tilde{a}-\tilde{b})-x|}{|x|}
  \;\le\;\kappa\,u+u,
  \qquad \kappa=\frac{|a|+|b|}{|a-b|}.
\]
\end{proposisi}

\begin{proof}
Since $\operatorname{fl}(\tilde{a}-\tilde{b})$ is the correctly rounded version of
$\tilde{a}-\tilde{b}$, the additional error is at most $u|\tilde{a}-\tilde{b}|$.
Furthermore
$|\tilde{a}-\tilde{b}-(a-b)|\le|a||\delta_{a}|+|b||\delta_{b}|\le u(|a|+|b|)$,
and dividing by $|x|=|a-b|$ gives $u\kappa$. Adding the two errors with the
triangle inequality completes the proof.
\end{proof}

The number $\kappa$ is never less than one and grows without bound as
$|a-b|$ shrinks. For $a=12345.678$ and $b=12345.671$ we get $a-b=0.007$
and $\kappa\approx3.53\times10^{6}$, so that the relative error of the result can reach
$3.9\times10^{-10}$ even though each input is off by only $10^{-16}$. Since every
factor of ten in $\kappa$ erases one significant digit,
$\log_{10}\kappa\approx6.5$ means six and a half of the sixteen machine digits are gone,
and the result $0.007$ leaves only about ten digits that can be trusted.

One note that saves us: Sterbenz's rule applies to differences, not to
sums. For the sum $a+b$ the condition number is
$(|a|+|b|)/|a+b|$, which approaches one when $a$ and $b$ have the same sign. Adding
large numbers is not dangerous. Subtracting two large numbers that are nearly equal is
dangerous, and that is the only habit that needs watching.

\section{Catastrophic cancellation: two faces}
\tingkat{intermediate}

We already know the mechanism. Now let us give it a name, because two similar kinds of
cancellation are often taken for the same thing although their consequences differ
widely.

\begin{untukanak}
If we want to know the thickness of a sheet of paper, an ordinary ruler will not help:
one sheet is far thinner than a millimeter, so the needle seems not to move. The way is
not to look for a finer ruler but to change what is measured: stack five hundred sheets,
measure the thickness of the stack, then divide by five hundred. The result is far more
trustworthy, even though the ruler is the same ruler.

It turns out that a formula can be treated in the same way. Some formulas work by
subtracting two large numbers that are nearly equal, and that is the same as weighing a
single sheet of paper with a ruler. Other formulas are mathematically exactly the same,
but order their steps so that the small number is measured directly rather than
recovered from a difference. The results can differ widely, even though the paper has
not changed.

The question: if two formulas that are equally correct give different answers, which
one should we trust? In this chapter we shall learn why the wrong formula is not the one
whose arithmetic is mistaken but the one that chooses the step that throws information
away, and how to recognize it before the program runs.
\end{untukanak}

\kamus{catastrophic cancellation}{pembatalan katastrofik}{The loss of significant digits
when two nearly equal numbers are subtracted, so that only error is left.}
\index{catastrophic cancellation}

\begin{definisi}[Catastrophic and benign cancellation]
Let $x$ and $y$ be machine numbers computed from earlier data.
The subtraction $x-y$ is called \emph{benign cancellation} if $x$ and $y$ are exact, so
that the result is exact by Sterbenz's theorem. It is called \emph{catastrophic
cancellation} if $x$ and $y$ themselves already carry rounding error, so that the digits
the two share disappear as well and what is left is dominated by that error.
\end{definisi}

A benign example: $1.0000000000000002-1$ in binary64. Both are machine numbers, the
difference is $2^{-52}$ and is represented exactly. The result is small, but not one
digit disappears: we really do know both inputs to the last bit.

The catastrophic example is the quadratic formula from the opening of the chapter. There
$b^{2}$ and $4ac$ are each the result of a multiplication, so both are rounded. For
$a=1$, $b=10^{8}$, $c=1$ both happen to be exact, and it is precisely that specialness
which makes this example worth examining down to the bit: the collapse does not come
from the multiplications but from the rounding of the square root. Python computes

\begin{center}
\begin{tabular}{ll}
\hline
discriminant $b^{2}-4ac$ & \texttt{9999999999999996.0} \\
$\sqrt{b^{2}-4ac}$ & \texttt{99999999.99999999} \\
$b-\sqrt{b^{2}-4ac}$ & \texttt{1.4901161193847656e-08} \\
$\operatorname{ulp}(b)$ & \texttt{1.4901161193847656e-08} \\
\hline
\end{tabular}
\end{center}

The true value of $b-\sqrt{b^{2}-4ac}$ is $2\times10^{-8}$, that is $1.34$ ulp.
What is computed is exactly $1$ ulp, and half an ulp of error in the square root is
enough for that: the square root is correct to $0.46$ ulp at the scale $10^{8}$, but a
difference that ought to be $1.34$ ulp can only land on $1$ or $2$ ulp, so its
information is down to one binary choice. The final result is $-7.45\times10^{-9}$, that
is $-2^{-27}$, whereas the true root is $-10^{-8}$. Relative error $0.255$.

Note the division of labor. The square root works with relative error
$5\times10^{-17}$, far below $u$. The problem itself is well conditioned: shifting
$b$ by one ulp shifts the small root by one ulp, so
$\kappa\approx|b|/\sqrt{b^{2}-4ac}\approx1$. What is bad is the \emph{form} we have
chosen: the condition number of the subtraction at that step is
\[
  \kappa=\frac{|b|+\sqrt{b^{2}-4ac}}{\left|\sqrt{b^{2}-4ac}-|b|\right|}
  =\frac{2\times10^{8}}{1.49\times10^{-8}}\approx1.34\times10^{16},
\]
so that $\kappa u\approx1.49$: the relative error bound passes one, which means not a
single digit is guaranteed.

Figure~\ref{fig:katastrofik} separates two things that are often mixed up at that step:
the size of the numbers, and how many digits still mean something after the subtraction.

\begin{figure}[htbp]
\centering
\begin{tikzpicture}[every node/.style={font=\footnotesize}]
  \draw[ink,thick,->] (1.2,3.20) -- (10.6,3.20);
  \foreach \dec/\tk in {-9/{$10^{-9}$},-6/{$10^{-6}$},-3/{$10^{-3}$},0/{$10^{0}$},
                        3/{$10^{3}$},6/{$10^{6}$},9/{$10^{9}$}} {
    \draw[ink] ({1.4+(\dec+9)*0.5},3.26) -- ({1.4+(\dec+9)*0.5},3.14);
    \node[below] at ({1.4+(\dec+9)*0.5},3.12) {\tk};
  }
  \fill[accent] ({1.4+17*0.5},3.20) circle (2.4pt);
  \node[above=2pt] at ({1.4+17*0.5},3.24) {$b=10^{8}$};
  \fill[layerteal] ({1.4+1.173*0.5},3.20) circle (2.4pt);
  \node[above=2pt] at ({1.4+1.173*0.5},3.24) {$b-\sqrt{\Delta}$};
  \draw[<->,thick,dashed,layergrey] ({1.4+1.173*0.5},3.72) -- ({1.4+17*0.5},3.72);
  \node[above] at ({1.4+9.1*0.5},3.74) {sixteen orders of magnitude};
  \foreach \i in {0,...,16} {
    \fill[accent!28] ({1.4+\i*0.45},1.60) rectangle ++(0.45,0.40);
    \draw[ink!45] ({1.4+\i*0.45},1.60) rectangle ++(0.45,0.40);
  }
  \foreach \i in {0,...,15} {
    \fill[layergrey!16] ({1.4+\i*0.45},0.85) rectangle ++(0.45,0.40);
    \draw[ink!45] ({1.4+\i*0.45},0.85) rectangle ++(0.45,0.40);
  }
  \fill[accent!70] ({1.4+16*0.45},0.85) rectangle ++(0.45,0.40);
  \draw[ink!45] ({1.4+16*0.45},0.85) rectangle ++(0.45,0.40);
  \node[left] at (1.35,1.80) {$b$};
  \node[left] at (1.35,1.05) {$b-\sqrt{\Delta}$};
  \draw[{Latex}-{Latex},layergrey] (1.40,0.62) -- ({1.4+16*0.45},0.62);
  \node[below] at ({1.4+8*0.45},0.60) {sixteen identical digits, canceling out};
  \node[right] at ({1.4+17*0.45},1.05) {one digit};
\end{tikzpicture}
\caption{Top: $b$ and $\sqrt{b^{2}-4ac}$ on a log scale, sixteen orders of
magnitude away from their difference. Bottom: the boxes are significant digits, seventeen
go in and sixteen of them are the same so that they cancel out. What remains is one
digit, namely the difference $1.49\times10^{-8}$ itself; since half an ulp of error in
the square root is enough to shift the last digit, the small root comes out with a
relative error of $0.255$.}
\label{fig:katastrofik}
\end{figure}

The fix is not more precision but choosing which root to compute first.
The product of the two roots is $c/a$, so once the large root $r_{1}$ is known, the
small root is obtained by a division, not by a subtraction:
\[
  r_{1}=\frac{-b-\sqrt{b^{2}-4ac}}{2a},\qquad r_{2}=\frac{c}{a\,r_{1}}
  =\frac{2c}{-b-\sqrt{b^{2}-4ac}}.
\]
The second form never subtracts two numbers that are close together. For the example
above it gives $-10^{-8}$ exactly, that is the correctly rounded root. The same formula
with the signs arranged according to the sign of $b$ is used in almost every library
that solves quadratic equations, and the reason is stability, not tidiness.

The second face of catastrophic cancellation appears in a place that looks harmless:
the difference $1-\cos x$ for small $x$. For small $x$,
\[
  \frac{1-\cos x}{x^{2}}\longrightarrow\frac{1}{2},
\]
because $\cos x=1-x^{2}/2+x^{4}/24-\cdots$. The formula is right, and it is also the
most natural one to write. But $\cos x$ is stored with relative error $u$, and because
$|\cos x|\le1$, its absolute error is about $u$. After division by $x^{2}$, that error
becomes a relative error of about $2u/x^{2}$. For $x=10^{-8}$ that factor is $2.2$:
every digit is wiped out.

Here are the results of an actual run for ten values of $x$, compared with a reference
computed at 50 digits and with the Taylor series $0.5-x^{2}/24+x^{4}/720$; for
$x\le10^{-3}$ the two comparators agree to the last digit. The ``naive'' column is
$(1-\cos x)/x^{2}$, the ``identity'' column is
$\tfrac12\left(\sin(x/2)/(x/2)\right)^{2}$, which comes from
$1-\cos x=2\sin^{2}(x/2)$.

\begin{center}
\begin{tabular}{rccc}
\hline
$x$ & naive & relative error & identity \\
\hline
$10^{-1}$ & $0.499583472197418$ & $1.1\times10^{-14}$ & $0.499583472197423$ \\
$10^{-2}$ & $0.499995833347366$ & $2.9\times10^{-13}$ & $0.499995833347222$ \\
$10^{-3}$ & $0.499999958325503$ & $1.6\times10^{-11}$ & $0.499999958333335$ \\
$10^{-4}$ & $0.499999996961265$ & $5.2\times10^{-9}$ & $0.499999999583333$ \\
$10^{-5}$ & $0.500000041370185$ & $8.3\times10^{-8}$ & $0.499999999995833$ \\
$10^{-6}$ & $0.500044450291171$ & $8.9\times10^{-5}$ & $0.499999999999958$ \\
$10^{-7}$ & $0.499600361081320$ & $8.0\times10^{-4}$ & $0.500000000000000$ \\
$10^{-8}$ & $0.000000000000000$ & $1.0$ & $0.500000000000000$ \\
$10^{-9}$ & $0.000000000000000$ & $1.0$ & $0.500000000000000$ \\
$10^{-10}$ & $0.000000000000000$ & $1.0$ & $0.500000000000000$ \\
\hline
\end{tabular}
\end{center}

The pattern is typical and worth memorizing. As long as $x^{2}/2$ is far above $u$, the
naive formula is nearly right. When $x^{2}/2$ shrinks to the scale of $u$, the
significant digits run out and the error grows like $2u/x^{2}$: each tenfold decrease of
$x$ adds two digits of loss. At $x=10^{-8}$, $1-\cos x$ is smaller than half an ulp
around $1$, so that $\cos x$ is rounded exactly to $1$ and the result is zero. That
breakdown point is no mystery: $x^{2}/2\approx u$, that is $x\approx\sqrt{2u}\approx1.5\times10^{-8}$.

The last column shows that the problem is well conditioned. The double-angle identity
gives an answer correct to the last digit for every $x$ in the table, because it never
subtracts two numbers that are close together. Error is not a property of the problem.
It is a property of the formula chosen.

\begin{resep}
Four steps for choosing a formula form that is safe from cancellation:
\begin{enumerate}
  \item Mark every subtraction $a-b$ inside the formula, then compute
        $\kappa=(|a|+|b|)/|a-b|$ on the worst data that can reasonably arrive.
  \item Translate $\kappa$ into lost digits: $\kappa u\approx10^{-d}$
        means the last $d$ digits cannot be trusted.
  \item Look for an equivalent form that contains no such difference: multiply by the
        conjugate, use the double-angle identity, factor, or call \texttt{log1p},
        \texttt{expm1}, and \texttt{hypot}.
  \item Repeat the computation at higher precision. The difference between the two
        results should be comparable to $\kappa u$; if it is far larger, the new form
        still hides a subtraction.
\end{enumerate}
A warning about order: step 1 must not be swapped with step 3. Rewriting a formula
before knowing which subtraction is dangerous only moves the problem, as in the one-pass
form of the variance that trades one dangerous subtraction for another.
\end{resep}
\section{Summation: error that grows with $n$}
\tingkat{advanced}

Summation is the operation most often taken for granted, and yet it is where error is
easiest to compute. Before we come to the bound, here is a demonstration that dismantles
the assumption that order does not matter:

\begin{center}
\texttt{(0.1 + 0.2) + 0.3 = 0.6000000000000001}, \quad
\texttt{0.1 + (0.2 + 0.3) = 0.6}.
\end{center}

The difference is $1.11\times10^{-16}$, exactly one ulp at the scale $0.6$. No digit is
written down wrongly: only three numbers, each with no finite binary
representative, summed in a different order. In mathematics, addition is
associative. Inside the machine, it is associative only when no differing rounding
happens to occur.

\begin{teorema}[Error of sequential summation]
Let $s_{1}=x_{1}$ and $s_{k}=\operatorname{fl}(s_{k-1}+x_{k})$ for
$k=2,\dots,n$ with $x_{i}\in\mathbb{F}$. Then
\[
  \left|s_{n}-\sum_{i=1}^{n}x_{i}\right|
  \;\le\;(n-1)\,u\sum_{i=1}^{n}|x_{i}|+O(u^{2}).
\]
\end{teorema}

\begin{proof}
Every step uses one rounding, so $s_{k}=(s_{k-1}+x_{k})(1+\delta_{k})$
with $|\delta_{k}|\le u$. Expanding the result recursively gives
\[
  s_{n}=\sum_{i=1}^{n}x_{i}\prod_{k=i}^{n}(1+\delta_{k}),
\]
where the first term carries no $\delta_{1}$ because $s_{1}=x_{1}$ has not been rounded.
Since $\prod(1+\delta_{k})=1+\theta$ with
$|\theta|\le(1+u)^{n-1}-1\le(n-1)u+O(u^{2})$, the triangle inequality gives the bound
claimed.
\end{proof}

That bound is honest but loose: it assumes every rounding reaches its
worst error and that all of them point the same way, which in practice never happens.
The following table shows the gulf between bound and reality for the harmonic series
$H_{n}=\sum_{i=1}^{n}1/i$ in binary32, with the exact value computed by
\texttt{math.fsum}, which sums exactly.

\begin{center}
\begin{tabular}{rccc}
\hline
$n$ & sequential & exact & bound $(n-1)u\sum|x_{i}|$ \\
\hline
$10^{2}$  & $5.1873779297$ & $5.1873775176$ & $3.1\times10^{-5}$ \\
$10^{3}$  & $7.4854784012$ & $7.4854708606$ & $4.5\times10^{-4}$ \\
$10^{4}$  & $9.7876129150$ & $9.7876060360$ & $5.8\times10^{-3}$ \\
$10^{5}$  & $12.0908508301$ & $12.0901461299$ & $7.2\times10^{-2}$ \\
$10^{6}$  & $14.3573579788$ & $14.3927267229$ & $8.6\times10^{-1}$ \\
\hline
\end{tabular}
\end{center}

The actual error at $n=10^{6}$ is $3.54\times10^{-2}$, about twenty-five
times smaller than the bound, but it has already erased two and a half significant
digits. Nor does the error rise smoothly: it dips at $n=10^{4}$ because roundings in the
same and in opposite directions happened to cancel. What can be relied on is not the
size of the error at one $n$, but its tendency to grow, at a worst-case rate
proportional to $n$.

The cause is not large rounding but an accumulator that moves away from zero.
Once $s$ is around $14$, the term $1/i$ for large $i$ lies below half an
ulp of $s$ and every addition loses all its information. The extreme case is easy to
construct: add $1$ to a million terms of value $10^{-16}$; exactly the result is
$1.0000000001$, whereas sequential summation stays at $1$ for however many terms there
are. What is lost is not one term but the whole contribution.

Figure~\ref{fig:galatjumlah} summarizes the same pattern for the five values of $n$ in the
table above, this time with Kahan summation as the comparison.

\begin{figure}[htbp]
\centering
\begin{tikzpicture}
\begin{axis}[
  width=0.88\linewidth, height=6.0cm,
  xmode=log, ymode=log,
  xlabel={$n$}, ylabel={$|s_{n}-H_{n}|$},
  xmin=70, xmax=1.6e6, ymin=1e-8, ymax=1,
  grid=both, grid style={layergrey!15},
  tick label style={font=\footnotesize},
  label style={font=\footnotesize},
  legend style={font=\footnotesize,draw=layergrey!40,
                at={(0.03,0.97)},anchor=north west},
  legend cell align=left,
]
\addplot[ink,thick,mark=*,mark size=1.7pt] coordinates {
  (100,4.12e-7) (1000,7.54e-6) (10000,6.88e-6) (100000,7.05e-4) (1000000,3.54e-2)};
\addlegendentry{sequential}
\addplot[accent,thick,mark=square*,mark size=1.6pt] coordinates {
  (100,6.48e-8) (1000,8.88e-8) (10000,2.03e-7) (100000,6.51e-8) (1000000,1.75e-7)};
\addlegendentry{Kahan}
\end{axis}
\end{tikzpicture}
\caption{Absolute error of the harmonic sum in binary32 at five sizes $n$, log-log
scale. Sequential summation rises from $4.1\times10^{-7}$ at $n=10^{2}$ to
$3.54\times10^{-2}$ at $n=10^{6}$; Kahan holds at order $10^{-7}$ across the whole
range, so that at $n=10^{6}$ the two are $2\times10^{5}$ times apart. The data, the
precision, and the number of operations are the same; only the order of summation
differs.}
\label{fig:galatjumlah}
\end{figure}

\begin{tangan}
Materials: eleven blank cards, a marker, and one rule of the game. Every time two cards
are added, write the result on a new card with at most four significant figures,
then throw the old cards away.

Write $1.000$ on one card and $0.0001$ on ten other cards. First route: put down the big
card, then add the ten small cards one at a time. Each time
$1.000+0.0001=1.0001$ is cut down to $1.000$, so that the small card disappears
without a trace and the pile stays $1.000$.

Shuffle the cards again, second route: first add up the ten small cards into
$0.0010$, then add the big card. The final pile is $1.001$.

What must be visible: the same eleven cards, the same cutting rule, two
different results in the last digit. Only the order changed, and the machine repeats that
order a million times.
\end{tangan}

Kahan summation attacks that problem directly.\index{Kahan summation}
The idea is simple: every time
a term is lost to rounding, keep the lost remainder in a compensation variable,
then put it back in at the next step.

\begin{quote}
\texttt{s = 0.0; c = 0.0}\\
\texttt{for x in data:}\\
\hspace*{1.5em}\texttt{y = x - c}\\
\hspace*{1.5em}\texttt{t = s + y}\\
\hspace*{1.5em}\texttt{c = (t - s) - y}\\
\hspace*{1.5em}\texttt{s = t}
\end{quote}

The line \texttt{c = (t - s) - y} is the heart of it, and it works precisely because of
Sterbenz: $t-s$ is the difference of two numbers that are close together and hence exact,
and $(t-s)-y$ is therefore exactly the error that was just discarded. The result no
longer piles up error proportional to $n$.

\begin{teorema}[Kahan summation bound]
With Kahan's algorithm above, the computed sum satisfies the inequality
\[
  s=\sum_{i=1}^{n}x_{i}(1+\delta_{i})+O(nu^{2})\sum_{i=1}^{n}|x_{i}|,
  \qquad |\delta_{i}|\le2u .
\]
\end{teorema}

We do not repeat the proof of this theorem here. It is Theorem 8 in Goldberg's
survey, and its proof is in the appendix of that work, taken from Knuth. What is worth
reading from the statement: the relative error of each term is bounded by $2u$, twice the
error of a single rounding, but the factor $n$ is gone from the leading term. For small
$n$ this bound is in fact looser than the naive one, because $2u>u$. Its advantage
appears only when $n$ is large, and that is where its whole use lies.

The real numbers on the harmonic series with $n=10^{6}$: sequential summation in
binary32 gives an error of $3.54\times10^{-2}$, whereas Kahan gives
$1.75\times10^{-7}$, that is $0.18$ ulp from the correct answer. In binary64
the difference is just as sharp: sequential is off by $7.35\times10^{-13}$, or $414$ ulp,
whereas Kahan produces a figure identical to \texttt{math.fsum}.

Two other improvements deserve mention. Pairwise summation, which
\texttt{numpy.sum} uses,\index{pairwise summation} splits the data recursively
so that the error grows like $\log n$; it is nearly free and suits data in
memory. \texttt{math.fsum} uses an exact accumulator and gives a result that is correctly
rounded for the whole list, at a higher cost in time. Summing from the smallest term to
the largest is the cheapest improvement, but it only reduces the error.

\section{Practical guidance}
\tingkat{basic}

Everything we have built can be compressed into a few habits that are cheap to
keep and expensive to neglect.

First, choose precision from the condition number, not from habit.
The relative error of the result is roughly $\kappa u$. With binary64, a computation
with condition $10^{6}$ still leaves ten digits; with binary32, the same computation
leaves only one digit. For $10^{8}$ sequential operations the crude bound is
$nu$, and in binary32 that already passes one, which means the answer is meaningless.

Second, do not test fractional numbers for equality with \texttt{==}. Two paths that
are mathematically identical may produce different bits, like
$0.1+0.2$ above. The right test uses a scaled tolerance, for instance
$|x-y|\le10^{-12}\max(1,|x|,|y|)$, and that tolerance must come from error
analysis, not from experimenting until the test passes.

Third, distrust every formula that contains a difference of two nearly equal quantities.
The standard form of the variance, $(1-\cos x)/x^{2}$, and $\log(1+x)$ for small $x$ can
all be written differently. The standard library already does so: \texttt{log1p},
\texttt{expm1}, and \texttt{hypot} exist precisely because their naive forms are not
stable.

Fourth, for a long sum, or one with terms of very different sizes, use
\texttt{math.fsum} or pairwise summation, and keep the accumulator in higher precision
where available. Sort from the smallest term if there is no other choice.

Fifth, do not conclude truth from agreement between numbers. Two programs that agree
to twelve digits are fairly convincing; ones that agree to three digits only tell
us that both use the same formula. A better test is to repeat the computation at a
different precision: if the result shifts far beyond the prediction
$\kappa u$, something is wrong with the formula, not with the machine.


\pagebreak[4]
\begin{tantangan}
Compute $10^{8}\left(\sqrt{10^{16}+1}-10^{8}\right)$ in binary64. The true value is
$0.5$, but a program that writes the formula down as it stands reports $0.0$.

Explain why the result is zero, then compute the form that gives $0.5$ without raising
the precision at all.

Hint 1: ask what $\operatorname{ulp}(10^{16})$ is, then whether $10^{16}+1$
still has room on the grid of machine numbers.

Hint 2: multiply numerator and denominator by the conjugate
$\sqrt{10^{16}+1}+10^{8}$; the dangerous subtraction becomes a division.

If you are stuck, go back to the section Catastrophic cancellation: two faces.
\end{tantangan}

\section{Worked example}
\tingkat{intermediate}

Every number in this section can be checked with a calculator, \texttt{struct}, and
ordinary fractions. We begin by taking apart one bit pattern, namely $0.1$ in binary32:

\[
  \underbrace{0}_{1}\;\underbrace{01111011}_{8}\;
  \underbrace{10011001100110011001101}_{23}.
\]

The stored exponent is $123$, so after subtracting the bias $127$ we get $e=-4$. The
23-bit integer in the mantissa is $5033165$, so that the significand is
$1+5033165/2^{23}$. Hence
\[
  x=\left(1+\frac{5033165}{2^{23}}\right)2^{-4}=\frac{13421773}{134217728}.
\]
The numerator comes from $2^{23}+5033165=13421773$. Meanwhile the denominator is
$134217728$, and a tenth of it is $13421772.8$, so that the error is exactly
$0.2/134217728$, that is $1.4901161\times10^{-9}$,
and the relative error is $1.49\times10^{-8}$, that is $u_{32}/4$. That figure makes
sense: the number $1/10$ equals $1.1001100110011\dots_{2}\times2^{-4}$, which needs a
periodic mantissa and does not terminate, so that correct rounding cuts it at the 24th
bit. The pattern that is stored
ends in \texttt{\dots 1101}, not \texttt{\dots 1100}, because the remainder behind it is
greater than half an ulp and the rounding goes up.

\kertas{Take apart the pattern \texttt{0 01111100 01000000000000000000000} by hand:
determine the sign, the exponent after subtracting the bias, the value of the mantissa,
then write the number as a fraction and as a decimal. Compute the
$\operatorname{ulp}$ of that number, then decide the result of adding $2^{-27}$ and
$3\cdot2^{-27}$ under the \emph{tie to even} rule. Check your answer with
\texttt{struct.unpack} and \texttt{numpy.float32}. This number was chosen so that its
mantissa stops after two digits, so that all of its arithmetic can be done without a
calculator.}

Now the case $0.1+0.2$ in binary64, which can be computed entirely with
fractions. All three numbers involved have power-of-two denominators:
\[
  \operatorname{fl}(0.1)=\frac{3602879701896397}{2^{55}},\qquad
  \operatorname{fl}(0.2)=\frac{3602879701896397}{2^{54}},
\]
\[
  \operatorname{fl}(0.3)=\frac{5404319552844595}{2^{54}}.
\]
The sum of the first two before rounding is $10808639105689191/2^{55}$,
whereas $\operatorname{fl}(0.3)=10808639105689190/2^{55}$. So the exact sum
lies exactly half an ulp above $\operatorname{fl}(0.3)$: a tie point, not something near
one. The last digit of its mantissa is $1$, that is odd, so that \emph{tie to even}
rounds up: $0.300000000000000044408920985006\dots$

The difference from $\operatorname{fl}(0.3)$ is therefore exactly one ulp,
$2^{-54}=5.551115123125783\times10^{-17}$, not a random error. The subtraction that
follows is exact by Sterbenz, so that the figure the program prints is that figure
precisely.

The same numbers explain why the conclusion depends on the precision: in
binary32, $0.1+0.2$ and $0.3$ produce identical bit patterns because the two
rounding errors happen to cancel each other. So \texttt{0.1+0.2 != 0.3} is not a law of
nature but the consequence of a coincidence at one precision.

Now the square root from before, with ulp arithmetic. The discriminant $b^{2}-4ac$ is
exactly $10^{16}-4$, because $10^{16}$ is still below $2^{54}$ and that number is a
multiple of $4$. Its square root is correct to $0.46$ ulp, so that it lands one ulp
below $b$ because $\operatorname{ulp}(b)=2^{-26}$. The true difference is $2\times10^{-8}=1.34$ ulp,
the computed difference $1$ ulp; that error of $0.34$ ulp at the scale $10^{8}$, divided
by two, becomes the whole error of the small root:
\[
  r_{2}^{\text{naive}}=-\frac{2^{-26}}{2}=-2^{-27}
  =-7.450580596923828\times10^{-9},
  \qquad r_{2}=-10^{-8}.
\]
The stable form $2c/(-b-\sqrt{b^{2}-4ac})$ suffers no dangerous subtraction, and
its result is $-10^{-8}$ exactly, the correctly rounded root.
\section{Computational lab}
\tingkat{advanced}

\begin{lab}
\begin{verbatim}
import math, struct
import numpy as np

def bit32(x):
    return ''.join(f'{b:08b}' for b in struct.pack('>f', float(x)))

print("0.1     :", bit32(0.1))
print("6.0     :", bit32(6.0))
print("-12.375 :", bit32(-12.375))
print("hex     :", (0.1).hex(), (6.0).hex(), (1e-8).hex())
print("eps     :", np.finfo(np.float32).eps, np.finfo(np.float64).eps)
print("ulp     :", math.ulp(1.0), math.ulp(1e10))
print("tie     :", 1 + 2.0**-53, 1 + 3*2.0**-53)

s = (0.1 + 0.2) - 0.3
print("0.1+0.2-0.3 :", repr(s), "=", s / math.ulp(0.3), "ulp")
print("associative :", repr((0.1 + 0.2) + 0.3), repr(0.1 + (0.2 + 0.3)))

n = 10**6
exact = math.fsum(1.0 / i for i in range(1, n + 1))

def jumlah(prec, kahan):
    s, c = prec(0), prec(0)
    for i in range(1, n + 1):
        y = prec(1) / prec(i)
        if kahan:
            y = y - c
            t = s + y
            c = (t - s) - y
            s = t
        else:
            s = s + y
    return s

print("H_n exact   :", f"{exact:.12f}")
for prec, nama in ((np.float32, "float32"), (np.float64, "float64")):
    for kahan in (False, True):
        v = float(jumlah(prec, kahan))
        print(f"{nama} {'Kahan' if kahan else 'naive':7} {v:.12f}"
              f"  error {abs(v - exact):.2e}")

print(f"{'x':>8} {'naive':>19} {'error':>9} {'identity':>19} {'error':>9}")
for k in range(1, 11):
    x = 10.0**-k
    lugu = (1 - math.cos(x)) / x**2
    stabil = 2 * math.sin(x / 2)**2 / x**2
    acuan = 0.5 - x**2/24 + x**4/720 - x**6/40320
    print(f"{x:8.0e} {lugu:19.15f} {abs(lugu-acuan)/acuan:9.1e}"
          f" {stabil:19.15f} {abs(stabil-acuan)/acuan:9.1e}")

a, b, c = 1.0, 1e8, 1.0
d = b*b - 4*a*c
sq = math.sqrt(d)
print("discriminant :", d, " root:", sq)
print("b - root    :", b - sq, " ulp(b):", math.ulp(b))
print("naive root  :", (-b + sq) / (2*a))
print("stable root :", 2*c / (-b - sq))
\end{verbatim}
The actual output:
\begin{verbatim}
0.1     : 00111101110011001100110011001101
6.0     : 01000000110000000000000000000000
-12.375 : 11000001010001100000000000000000
hex     : 0x1.999999999999ap-4 0x1.8000000000000p+2 0x1.5798ee2308c3ap-27
eps     : 1.1920929e-07 2.220446049250313e-16
ulp     : 2.220446049250313e-16 1.9073486328125e-06
tie     : 1.0 1.0000000000000004
0.1+0.2-0.3 : 5.551115123125783e-17 = 1.0 ulp
associative : 0.6000000000000001 0.6
H_n exact   : 14.392726722866
float32 naive   14.357357978821  error 3.54e-02
float32 Kahan   14.392726898193  error 1.75e-07
float64 naive   14.392726722865  error 7.35e-13
float64 Kahan   14.392726722866  error 0.00e+00
       x               naive     error            identity     error
   1e-01   0.499583472197418   5.6e-15   0.499583472197423   5.4e-15
   1e-02   0.499995833347366   2.9e-13   0.499995833347222   1.1e-16
   1e-03   0.499999958325503   1.6e-11   0.499999958333335   0.0e+00
   1e-04   0.499999996961265   5.2e-09   0.499999999583333   1.1e-16
   1e-05   0.500000041370185   8.3e-08   0.499999999995833   1.1e-16
   1e-06   0.500044450291171   8.9e-05   0.499999999999958   0.0e+00
   1e-07   0.499600361081320   8.0e-04   0.500000000000000   1.1e-16
   1e-08   0.000000000000000   1.0e+00   0.500000000000000   0.0e+00
   1e-09   0.000000000000000   1.0e+00   0.500000000000000   0.0e+00
   1e-10   0.000000000000000   1.0e+00   0.500000000000000   0.0e+00
discriminant : 9999999999999996.0  root: 99999999.99999999
b - root    : 1.4901161193847656e-08  ulp(b): 1.4901161193847656e-08
naive root  : -7.450580596923828e-09
stable root : -1e-08
\end{verbatim}
\end{lab}

The first three lines confirm that the table at the beginning of the chapter is not
invented: the bit patterns of $0.1$, $6$, and $-12.375$ come out exactly as in the hand
decomposition, sign bit included. The \texttt{hex} line shows thirteen hexadecimal
digits. The \texttt{tie} line is the plainest demonstration of the even rule: adding
half an ulp to $1$ changes nothing, adding one and a half ulp gives
$1+2^{-51}$. The \texttt{0.1+0.2-0.3} line prints
$5.551115123125783\times10^{-17}$ and declares itself to be $1.0$ ulp; the
\texttt{associative} line gives two different results, differing by
$1.11\times10^{-16}$.

The next four lines are the core of this chapter in its most compact form: at
$n=10^{6}$, binary32 sequential loses $3.54\times10^{-2}$ whereas Kahan loses only
$1.75\times10^{-7}$; binary64 sequential $7.35\times10^{-13}$ whereas Kahan zero, the
same bits as \texttt{math.fsum}. That error of $7.35\times10^{-13}$ is equivalent to
$414$ ulp, so double precision only moves the problem to a smaller scale.

The table for $(1-\cos x)/x^{2}$ shows the pattern already computed on paper. At
$x=10^{-1}$ both methods differ from the reference by about $5\times10^{-15}$ because the
reference itself is truncated at the $x^{6}$ term; from $x=10^{-2}$ on that reference is
far more accurate than the machine, and the error of the naive formula rises from
$2.9\times10^{-13}$ to $1.0$ as $x$ falls to $10^{-8}$, whereas the identity column holds
at $10^{-16}$. The zero at $x=10^{-8}$ is not lazy rounding: $\cos(10^{-8})$ really
equals $1$.

The last block checks the square root: discriminant $10^{16}-4$, root $10^{8}-2^{-26}$,
a difference of one ulp of $b$, namely $2^{-26}$, naive root $-2^{-27}$, and stable root
$10^{-8}$, which is correctly rounded. The last two lines
are $25.5$ percent apart, even though the formulas are algebraically equivalent and the
data are exactly the same.

\begin{untukdosen}
A teaching order that has worked in practice: begin with one failure that can be
recomputed, namely the quadratic formula at the opening of the chapter, and only then go
into representation. If representation comes first, students memorize $1/8/23$ and
$1/11/52$ without knowing what for. After that failure, the order that works is: the
number of machine numbers (so that it is clear rounding cannot be avoided), the rounding
rule, the ulp lemma, the $(1+\delta)$ bound, then the condition number. Summation is
better placed after the condition number, because the bound $(n-1)u\sum|x_{i}|$ only
reads as ``error that grows'' once $\kappa$ is familiar.

Two typical misconceptions. First, students believe that rounding error is random and
therefore cancels out in long sums. The harmonic table contradicts that sharply: the
error moves in one direction, and at $n=10^{6}$ it reaches thirty-five thousand ulp.
Second, they believe that subtraction always damages digits, and then avoid all
subtraction. Sterbenz's theorem shows the opposite: the subtraction of two nearby machine
numbers is in fact exact, and what does the damage is the error already present in the
input. The rule that needs repeating: what is dangerous is not the minus sign, but a
small difference of two numbers that have each already been rounded.

A cheap and effective whiteboard activity: ask the whole class to compute
$1+2^{-53}$ and $1+3\cdot2^{-53}$ with the even rule, then compare with Python.
After that, ask two volunteers to work out
$(0.1+0.2)+0.3$ and $0.1+(0.2+0.3)$ by writing the 52-bit mantissa
as a fraction. The hubbub usually starts when they discover that both are right.

A sign of readiness to move on: students can guess, without running a program, that
$\operatorname{fl}(1+10^{-16})=1$, can explain why the form
$2c/(-b-\sqrt{\Delta})$ is better than $(-b+\sqrt{\Delta})/(2a)$, and can
estimate how many digits are lost from a given $\kappa$. If
they say the words ``backward error'' without being prompted, they are ready for Chapter 25,
where the gradient of a network is computed with a similar rule.
\end{untukdosen}

\section{Summary}
\tingkat{basic}

\begin{ringkasan}
\begin{itemize}
  \item $\mathbb{F}$ is finite: $2^{p-1}$ normal numbers per power of two, so
        almost every real number has no representative and every act of storage means
        rounding.
  \item IEEE rounding is \emph{round to nearest, ties to even}. Consequently
        $\operatorname{fl}(x)=x(1+\delta)$ with $|\delta|\le u=2^{-p}$, that is
        $1.11\times10^{-16}$ for binary64 and $5.96\times10^{-8}$ for
        binary32.
  \item $\operatorname{eps}=2^{-p+1}$ is the distance from $1$ to the number after it;
        $\operatorname{ulp}(x)=2^{e-p+1}$ for $2^{e}\le|x|<2^{e+1}$. The error of one
        rounding is half an ulp, and expressed as a relative error it swings by
        a factor of two.
  \item The condition number $\kappa_{f}=|xf'(x)|/|f(x)|$ measures the amplification of
        input error. For subtraction, $\kappa=(|a|+|b|)/|a-b|$, which grows without
        bound as $|a-b|$ shrinks.
  \item Subtracting two nearby machine numbers is exact (Sterbenz). Catastrophic
        cancellation happens because the inputs already carry error, not because of
        the subtraction.
  \item The naive quadratic formula computes the small root through a subtraction of two
        numbers of size $|b|$; the form $2c/(-b-\sqrt{\Delta})$ avoids it and gives the
        correctly rounded root.
  \item Sequential summation has a worst-case error of $(n-1)u\sum|x_{i}|$ and is not
        associative. Kahan moves the error into a compensation variable so that the
        factor $n$ disappears from the leading term.
\end{itemize}
\end{ringkasan}

\section{Exercises}
\tingkat{intermediate}

\begin{latihan}
\begin{enumerate}
  \item Determine the $\operatorname{ulp}$ of $1$; $1024$; $0.1$; and $3\times10^{8}$
        in binary64, then compare with $|x|\,u$. How many binary64
        numbers lie in $[1,2)$?

  \item Show that $1+2^{-53}=1$ and $1+2^{-52}>1$ in binary64, then determine
        the smallest positive number $x$ with $1+x>1$ in binary32. Express
        your answer as $2^{-k}$ and compare it with the table in the opening section.

  \item Prove that for every normal $x$ we have
        $|\operatorname{fl}(x)-x|\le\tfrac12\operatorname{ulp}(x)$, then explain
        why that inequality cannot be replaced by equality for all $x$.

  \item Assess this claim: \emph{because every IEEE-754 operation is correctly rounded,
        the relative error of the final result of a program that uses only $+,-,\times,\div$ is always
        at most a few $u$.} Decide whether it is true, false, or not supported by
        evidence, and name one example from this chapter that supports your decision.

  \item Compute by hand $(0.1+0.2)+0.3$ and $0.1+(0.2+0.3)$ as
        fractions with power-of-two denominators, then show that the difference is one ulp.
        Check with \texttt{math.fsum} and with \texttt{float.hex()}.
        (Hint: write each number as $m/2^{55}$ or $m/2^{54}$ and
        compute the numerators as integers.)

  \item Repeat the computation of $(1-\cos x)/x^{2}$ for $x=10^{-7}$ with the double-angle
        identity, then compare with $0.5-x^{2}/24+x^{4}/720$. What relative error
        remains, and where does it come from?

  \item For $a=1$, $b=10^{4}$, $c=1$, compute both roots with the naive formula and with
        the stable form in binary64. Determine the condition number of the subtraction, then
        predict how many digits are lost before running the program. Check
        your prediction with Python.
        (Hint: $\sqrt{10^{8}-4}\approx10^{4}-2\times10^{-4}$; compare with
        $\operatorname{ulp}(10^{4})=2^{-39}$.)

  \item Prove that pairwise summation over $n=2^{k}$ terms has a worst-case error
        proportional to $ku\sum|x_{i}|$, not $nu\sum|x_{i}|$, by expanding
        the summation tree and counting the depth of the tree.

  \item Let $x_{i}=10^{-16}$ for $i=1,\dots,10^{6}$ and $x_{0}=1$. Compute
        sequential summation and Kahan summation both by hand (a calculator may be
        used), then explain why the Kahan answer differs even though
        each of its steps is also rounded.
        (Hint: in sequential summation every term lies below half an ulp
        of the accumulator; in Kahan, the compensation $c$ keeps that remainder and adds it
        back at the next step.)

  \item A program computes the variance with
        $\frac1n\sum x_{i}^{2}-\left(\frac1n\sum x_{i}\right)^{2}$ on data that
        are all around $999$. Determine the condition number of that form, predict
        how many digits are lost in binary32, and compare with the
        two-pass form $\frac1n\sum(x_{i}-\bar{x})^{2}$. Explain why the second form is
        better even though it requires two passes over the data.
\end{enumerate}
\end{latihan}

\begin{untukpeneliti}
What is really proved in this chapter are standard statements that do not
depend on any particular machine: the bound $\operatorname{fl}(x)=x(1+\delta)$ with
$|\delta|\le u$, the ulp lemma, Sterbenz's theorem, the error bound for sequential
summation, and the condition-number inequality via the mean value theorem. The primary
reference is Goldberg,
\emph{What Every Computer Scientist Should Know About Floating-Point Arithmetic}
(ACM Computing Surveys 23(1), 1991), sections ``Relative Error and Ulps'',
``Cancellation'', ``Exactly Rounded Operations'', and ``Errors in Summation''. The Kahan
bound is Theorem 8 of that work; its proof is in the appendix, taken from Knuth, and is
not repeated here. For a systematic treatment of stability, the further reference is
Higham, \emph{Accuracy and Stability of Numerical Algorithms}, the chapters on error
analysis and summation.

Three things are more often popularized than proved. First, the claim that
``rounding error is random and cancels out''. That is true only as a statistical
model for uncorrelated rounding, and it fails precisely in the most important case:
accumulation in long sums is systematic, as the harmonic table shows. Second, the claim
that the error can always be reduced with higher precision. For catastrophic
cancellation, raising the precision only moves the breakdown point, it does not remove
it: in binary64 that point lies at
$x\approx1.5\times10^{-8}$, and in binary128 it lies around $10^{-17}$. Third, the claim
that the ulp is an unambiguous unit of error. Goldberg shows that the
ulp may be defined as the distance between numbers or as half that distance, and the
difference between the two definitions appears as a factor of two in every bound. The
figure ``$0.46$ ulp'' in the discussion of the square root in this chapter uses the first
definition.

The limits of validity of the results here must be stated. The whole analysis assumes a model of
correct rounding and no \emph{fused multiply-add} that changes the order of operations,
no rearrangement by the compiler, and no rounding mode other than
\emph{to nearest}. On real hardware, vectorization and FMA instructions can change
results bit for bit; that is not a violation of the standard, but it makes comparisons
between machines non-deterministic. Error analysis also does not touch model error,
data error, or truncation error: this chapter discusses only rounding error.

Honest open questions. First, classical error analysis gives worst-case bounds,
whereas what practitioners need is a tight estimate for the typical case;
the probabilistic treatment of rounding is still an actively researched tool and not a
finished theory. Second, for iterative algorithms and the solution of linear systems,
rounding error interacts with the condition number of the matrix in such a way that
naive bounds are far too pessimistic; backward analysis, which is the central idea of
Higham and Wilkinson, gives a sharper framework but is not discussed here.
Third, practical decisions such as ``is binary32 enough'' depend on the distribution of
the data, and no answer holds in general.
\end{untukpeneliti}

\peta{This chapter follows Goldberg, \emph{What Every Computer Scientist Should Know About
Floating-Point Arithmetic} (ACM Computing Surveys, 1991), sections ``Introduction'',
``Relative Error and Ulps'', ``Cancellation'', ``Exactly Rounded Operations'',
``Formats and Operations'', ``Special Quantities'', and ``Errors in Summation'', with
Theorem 8 for the Kahan summation bound. For the real-analysis side of the limit ideas
and the mean value theorem used in the proof of the condition number, see Abbott,
\emph{Understanding Analysis}, sections 2.2 (Limit of a Sequence), 5.3 (The Mean Value
Theorem), and 6.6 (Taylor Series). The example $0.1+0.2$, the decomposition of bit
patterns, the analysis of $(1-\cos x)/x^{2}$, and all the numbers in the tables and the
lab are computations of our own run with Python 3.14 and NumPy 2.5 and re-verified with
ordinary fractions and 50-digit arithmetic. For an advanced treatment of stability
and backward error, the reference commonly used is Higham, \emph{Accuracy and
Stability of Numerical Algorithms}.}


% ---------------------------- part V ----------------------------
\part{Optimization}

This part turns ``find the best'' into a problem with a stopping condition. The prerequisites
are Chapters~6 and~12. After two chapters you can separate a local minimum from a global one,
use duality to prove bounds, and recognise when a linear problem can be solved in polynomial
time and when it cannot.
\chapter{Convex Optimization and Duality}

A paint mixer works with two pigments. Each pigment has a proportion the machine prefers
most: two parts for the first pigment, two parts for the second. But the recipe demands that
the two amounts total exactly two parts, and the solubility of the first pigment does not
allow more than half a part. We want the mixture as close as possible to the machine's
preference, under two demands that cannot be violated. Without the demands the answer is two
and two; with the sum demand alone, one and one. The solubility demand shifts the answer
again, and it is no longer clear in which direction.

That shape of problem shows up where nothing smells of pigment. Two model weights have best
values from an earlier fit, a calibration condition demands that their sum stay fixed, and
one device limit caps one of them. What are the best weights that still satisfy both, and
what is the price of each demand in units of the objective function? The second question is
the one engineers use to decide which demand is worth relaxing.

Chapter 10 gave least squares, but only without constraints. Chapter 12 gave the gradient,
but the gradient speaks only about neighbors: a point with zero gradient may be a minimum, a
maximum, or a saddle. Chapter 9 gave eigenvalues, which are used here, but eigenvalues do not
answer the most important question: if an algorithm stops at some point, is that the best
answer over the whole feasible region? For a convex function, yes. For a general problem, no,
and the fourth section shows how duality supplies a valid lower bound without examining the
entire feasible region.

\begin{untukanak}
Have you ever put a marble inside a bowl? The marble always stops at the bottom, in the same
place, wherever it is released, because the bottom curves upward like the letter U. Now
imagine the same marble released on a rocky hillside. It stops in the first hollow that is
deep enough, and that hollow is not necessarily the lowest one. From where it stops, you
cannot tell whether another valley is deeper without walking far and looking.

In this chapter we shall learn why a bowl is so friendly and a hillside so deceptive, and why
we can tell the two apart before computing anything. There is a second, subtler idea:
sometimes we do not need to know where the bottom of the valley lies, only the lowest bound
that can possibly be reached. If that bound equals the place where we stand, we are sure that
nothing is better, without examining a single point. Such bounds are sold by a price list for
each demand, and that price list is what makes many optimization problems trustworthy without
being solved to the end.
\end{untukanak}

\section{Convex sets and convex combinations}
\tingkat{basic}

In everyday language the word ``convex'' means bulging outward, like the back of a turtle. In
optimization it means almost the opposite: a set is called convex when it has no inward dent,
so that the line segment between any two of its members stays inside. The bowl above is
convex; the region between the two horns of a crescent moon is not.

There is one rule of the game: mix two points with nonnegative weights that add up to one, and
the mixture must still be a member of the set. Weight one means choosing the first point,
weight zero the second point, weight one half exactly halfway.

\kamus{convex combination}{kombinasi konveks}{A recipe for mixing several points with
proportions that add up to exactly one, like mixing paint: the result always lies among the
ingredients.}

\begin{definisi}[Convex set and convex combination]
A set $C\subseteq\mathbb{R}^{n}$ is called \emph{convex} when for every
$\mathbf{x},\mathbf{y}\in C$ and every $t\in[0,1]$ we have
$t\mathbf{x}+(1-t)\mathbf{y}\in C$. That point is called a \emph{convex combination} of
$\mathbf{x}$ and $\mathbf{y}$. In general, $\sum_{i=1}^{k}\theta_{i}\mathbf{x}_{i}$ is called
a convex combination when $\theta_{i}\ge0$ and $\sum_{i}\theta_{i}=1$.
\emph{The convex hull} $\operatorname{conv}S$ is the set of all finite convex combinations of
points in $S$.\index{convex set}\index{convex combination}
\end{definisi}

The definition with two points suffices for all points: if $C$ contains every mixture of two
members, it contains every finite mixture of its members, by induction on the number of
points.

Four important examples. A \emph{hyperplane}
$\{\mathbf{x}:\mathbf{a}^{\top}\mathbf{x}=b\}$ is convex, because
$\mathbf{a}^{\top}(t\mathbf{x}+(1-t)\mathbf{y})=tb+(1-t)b=b$. A \emph{half-space}
$\{\mathbf{x}:\mathbf{a}^{\top}\mathbf{x}\le b\}$ is convex for the same reason, using the
inequality. A \emph{ball} $\{\mathbf{x}:\|\mathbf{x}-\mathbf{c}\|\le r\}$ is convex because
the norm satisfies the triangle inequality (Chapter 10):
\[
  \|t\mathbf{x}+(1-t)\mathbf{y}-\mathbf{c}\|
  \le t\|\mathbf{x}-\mathbf{c}\|+(1-t)\|\mathbf{y}-\mathbf{c}\|\le tr+(1-t)r=r .
\]
A \emph{polyhedron}, that is a finite intersection of hyperplanes and half-spaces, is convex
for the following reason.

\begin{proposisi}[Intersections preserve convexity]
If $C_{1},\dots,C_{k}$ are convex, then $\bigcap_{i=1}^{k}C_{i}$ is convex.
\end{proposisi}

\begin{proof}
Take $\mathbf{x},\mathbf{y}$ in the intersection and $t\in[0,1]$. For every $i$, both points
are members of $C_{i}$, so $t\mathbf{x}+(1-t)\mathbf{y}\in C_{i}$ because $C_{i}$ is convex.
Hence the mixture lies in every $C_{i}$, that is in the intersection.
\end{proof}

Intersection is generous; union is not. Take $C_{1}=\{(x_{1},0):x_{1}\in[0,1]\}$ and
$C_{2}=\{(x_{1},1):x_{1}\in[0,1]\}$, both convex. The points $(0,0)$ and $(0,1)$ lie in the
union, but their halfway mixture $(0,\tfrac12)$ does not. The same happens with
$\{(x_{1},x_{2}):x_{1}x_{2}\ge1\}$: the points $(-1,-1)$ and $(1,1)$ are both members,
whereas the origin is not, because $0<1$. That set is two separate regions.

The convex hull is the cheapest patch for a hole of that kind: the smallest convex set that
contains $S$. It always exists, because the intersection of all convex sets containing $S$ is
convex by the proposition above and nonempty, since $\mathbb{R}^{n}$ itself is convex.

\section{Convex functions and their two tests}
\tingkat{intermediate}

Convex sets speak about regions; convex functions speak about the shape of the surface above
that region. What is demanded: the value at a mixture is no greater than the mixture of the
values. In the left panel of Figure~\ref{fig:fungsi-konveks} the chord from $(-1,1)$ to
$(1,1)$ lies above the graph of $f(t)=t^{2}$; in the right panel the chord cuts through it.

\kamus{convex function}{fungsi konveks}{A surface that never bends upward like a roof; the
chord between any two of its points always lies above the surface.}

\begin{definisi}[Convex function]
Let $C\subseteq\mathbb{R}^{n}$ be convex. A function $f\colon C\to\mathbb{R}$ is called
\emph{convex} when
\[
  f\bigl(t\mathbf{x}+(1-t)\mathbf{y}\bigr)
  \le t\,f(\mathbf{x})+(1-t)\,f(\mathbf{y})
\]
for all $\mathbf{x},\mathbf{y}\in C$ and $t\in[0,1]$. It is \emph{strictly convex} when that
inequality is strict for $\mathbf{x}\ne\mathbf{y}$ and $t\in(0,1)$, and \emph{concave} when
$-f$ is convex.\index{convex function}
\end{definisi}

\begin{figure}[!ht]
\centering
\begin{tikzpicture}
\begin{axis}[width=6cm,height=4.8cm,xlabel={$t$},ylabel={$f(t)$}]
\addplot[domain=-1.3:1.3,samples=50,thick] {x^2};
\addplot[dashed,thick,red] coordinates {(-1,1) (1,1)};
\addplot[only marks,mark=*] coordinates {(-1,1) (1,1)};
\end{axis}
\end{tikzpicture}\quad
\begin{tikzpicture}
\begin{axis}[width=6cm,height=4.8cm,xlabel={$t$},ylabel={$g(t)$}]
\addplot[domain=-1.3:1.3,samples=80,thick] {x^3-x};
\addplot[dashed,thick,red] coordinates {(-1,0) (1,0)};
\addplot[only marks,mark=*] coordinates {(-1,0) (1,0)};
\end{axis}
\end{tikzpicture}
\caption{The chord between two points of the graph. Left: $f(t)=t^{2}$, the chord from
$(-1,1)$ to $(1,1)$ has value $1$ at $t=0$ whereas the graph has $0$, so the chord lies above
the graph. Right: $g(t)=t^{3}-t$, the chord from $(-1,0)$ to $(1,0)$ crosses the graph at
$t=-1$, $0$, $1$; at $t=-0.5$ the graph has value $0.375$ whereas the chord has $0$.}
\label{fig:fungsi-konveks}
\end{figure}

That definition is equivalent to a geometric statement about the \emph{epigraph}, that is the
set of points above the graph: $f$ is convex exactly when
$\operatorname{epi}f=\{(\mathbf{x},t):f(\mathbf{x})\le t\}$ is convex. The equivalence is
purely algebraic, and it turns a question about functions into a question about sets.

For differentiable functions there are two far faster tests, and both are
necessary-and-sufficient, not merely sufficient.

\begin{teorema}[First-order test]
Let $f$ be differentiable on an open convex set $C$. The function $f$ is convex if and only
if
\[
  f(\mathbf{y})\ge f(\mathbf{x})+\nabla f(\mathbf{x})^{\top}(\mathbf{y}-\mathbf{x})
\]
for all $\mathbf{x},\mathbf{y}\in C$.
\end{teorema}

\begin{proof}
Suppose $f$ is convex, take $\mathbf{x}\ne\mathbf{y}$ and $t\in(0,1)$. From the definition,
\[
  \frac{f\bigl(\mathbf{x}+t(\mathbf{y}-\mathbf{x})\bigr)-f(\mathbf{x})}{t}
  \le f(\mathbf{y})-f(\mathbf{x}).
\]
As $t\downarrow0$ the left side tends to the directional derivative
$\nabla f(\mathbf{x})^{\top}(\mathbf{y}-\mathbf{x})$, so the required inequality holds in the
limit. For the converse direction, suppose that inequality holds and take
$\mathbf{z}=t\mathbf{x}+(1-t)\mathbf{y}$ with $t\in[0,1]$. Apply the inequality at
$\mathbf{z}$ against $\mathbf{x}$ and against $\mathbf{y}$:
\[
  f(\mathbf{x})\ge f(\mathbf{z})+\nabla f(\mathbf{z})^{\top}(\mathbf{x}-\mathbf{z}),
  \qquad
  f(\mathbf{y})\ge f(\mathbf{z})+\nabla f(\mathbf{z})^{\top}(\mathbf{y}-\mathbf{z}).
\]
Multiply the first by $t$, the second by $1-t$, then add. The gradient terms vanish because
$t(\mathbf{x}-\mathbf{z})+(1-t)(\mathbf{y}-\mathbf{z})=\mathbf{0}$, and what remains is
$t\,f(\mathbf{x})+(1-t)f(\mathbf{y})\ge f(\mathbf{z})$.
\end{proof}

The practical reading: the graph lies above every tangent plane. A single tangent plane
already gives a global lower bound, and that is why the gradient at any point becomes an
honest tool rather than a local forecast.

\begin{teorema}[Second-order test]
Let $f$ be twice differentiable on an open convex set $C$. The function $f$ is convex if and
only if the Hessian $\nabla^{2}f(\mathbf{x})$ is positive semidefinite for every
$\mathbf{x}\in C$, that is
$\mathbf{v}^{\top}\nabla^{2}f(\mathbf{x})\mathbf{v}\ge0$ for all $\mathbf{v}$.
\index{Hessian}
\end{teorema}

\begin{proof}
If $f$ is convex, then for every $\mathbf{x}$ and $\mathbf{v}$ the one-variable function
$g(s)=f(\mathbf{x}+s\mathbf{v})$ is convex on an interval around zero, because composition
with an affine map preserves convexity. The first-order test on $g$ shows that $g'$ is
increasing, so $g''(0)\ge0$; and $g''(0)=\mathbf{v}^{\top}\nabla^{2}f(\mathbf{x})
\mathbf{v}$.

Conversely, if the Hessian is positive semidefinite on all of $C$, Taylor's theorem with
integral remainder (Chapter 12) gives
\[
  f(\mathbf{y})=f(\mathbf{x})+\nabla f(\mathbf{x})^{\top}(\mathbf{y}-\mathbf{x})
  +\int_{0}^{1}(1-s)\,(\mathbf{y}-\mathbf{x})^{\top}
   \nabla^{2}f\bigl(\mathbf{x}+s(\mathbf{y}-\mathbf{x})\bigr)(\mathbf{y}-\mathbf{x})\,ds .
\]
The integrand is nonnegative because the Hessian is positive semidefinite and $(1-s)\ge0$, so
the integral term is nonnegative. The first-order inequality holds, and $f$ is convex.
\end{proof}

Two warnings about the second-order test. First, a Hessian that is positive
\emph{definite} suffices for strict convexity, but is not a necessary condition:
$f(x)=x^{4}$ is strictly convex although $f''(0)=0$. Second, a positive semidefinite Hessian
gives convexity only, not strictness; $f(x_{1},x_{2})=x_{1}^{2}$ has Hessian
$\operatorname{diag}(2,0)$ and its minimum is the entire line $x_{1}=0$.

\begin{bacakeras}
Two voices, read alternately with a friend. The first voice asks, the second answers.

``I drop a marble into the bowl. Where does it go?'' ``To the bottom. Always to the same
bottom.'' ``And if I release it from the other side?'' ``To that bottom too.'' ``So there is
only one lowest place?'' ``Only one, and every descending path finds it.''

``Now I release it on a stone slope. Where does it go?'' ``Into the first hollow that is deep
enough.'' ``Is that the deepest one?'' ``Not necessarily. That hollow holds it, it does not
convince it.''

Repeat alternately, faster and faster, with drumming on the table: a bowl has no traps, a
hillside has traps.

The line the reader must repeat, once, in a clear voice: if every descending path ends at the
same point, there are no traps there, and that is what convex means.
\end{bacakeras}

Several operations preserve convexity. The sum of two convex functions stays convex, because
the inequalities are added; a positive multiple does too. The composition
$f(\mathbf{A}\mathbf{x}+\mathbf{b})$ is convex when $f$ is convex, because an affine map
sends mixtures to mixtures. The pointwise maximum is convex as well: if $f=\max_{i}f_{i}$,
the inequality for each $f_{i}$ holds for the maximum. That explains why
$\max\{|x_{1}|,|x_{2}|\}$ has a crease and is still convex.

\begin{resep}
Four steps to check convexity, from the cheapest.
\begin{enumerate}
  \item Check the region first: take two points in $C$ and make sure the segment between them
        stays in $C$. If that fails, stop, because the next tests assume that $C$ is convex.
  \item Test through the definition with a general pair $\mathbf{x},\mathbf{y}$: write
        $f(t\mathbf{x}+(1-t)\mathbf{y})$ and compare it with
        $t\,f(\mathbf{x})+(1-t)f(\mathbf{y})$ for $t\in[0,1]$. One pair that reverses the
        inequality is enough to reject.
  \item Compute the Hessian $\nabla^{2}f(\mathbf{x})$ as a function of $\mathbf{x}$, not as
        one matrix of numbers at one point.
  \item Decide the sign: nonnegative eigenvalues, or all leading principal minors from
        $1\times1$ up to $n\times n$ nonnegative. One negative eigenvalue closes the case.
\end{enumerate}
A warning about the order: step 4 must not be taken before step 1. A Hessian that is positive
definite at one point speaks only about that point, whereas the definition speaks about all of
$C$.
\end{resep}

\section{Standard form and why a local optimum becomes global}
\tingkat{intermediate}

Up to here we have two separate ideas: a region with no holes and a surface that does not bend
upward. Optimization unites the two through a single way of writing. Every optimization
problem can be written in the same way: one function to be made small, a set of inequality
conditions, and a set of equality conditions.

\kamus{standard form}{bentuk baku}{A way of writing every optimization problem with one
pattern: one objective, conditions of the form ``no more than zero'', and conditions of the
form ``equals''.}

\begin{definisi}[Optimization problem in standard form]
An \emph{optimization problem in standard form} is
\[
  \begin{array}{ll}
  \text{minimize} & f_{0}(\mathbf{x})\\
  \text{subject to} & f_{i}(\mathbf{x})\le0,\qquad i=1,\dots,m,\\
                & h_{j}(\mathbf{x})=0,\qquad j=1,\dots,p,
  \end{array}
\]
with variable $\mathbf{x}\in\mathbb{R}^{n}$, objective function $f_{0}$, inequality constraint
functions $f_{i}$, and equality constraint functions $h_{j}$. Its domain is
$D=\bigcap_{i}\operatorname{dom}f_{i}\cap\bigcap_{j}\operatorname{dom}h_{j}$, its feasible
region is $X=\{\mathbf{x}\in D:f_{i}(\mathbf{x})\le0,\ h_{j}(\mathbf{x})=0\}$, and the
\emph{optimal value} is $p^{*}=\inf_{\mathbf{x}\in X}f_{0}(\mathbf{x})$, possibly $-\infty$
when the problem is unbounded below. The problem is called a \emph{convex problem} when
$f_{0},\dots,f_{m}$ are convex and $h_{1},\dots,h_{p}$ are affine.
\index{standard form}
\end{definisi}

That every constraint can be brought to this pattern is easy to check. The condition
$f(\mathbf{x})\ge b$ becomes $b-f(\mathbf{x})\le0$; a linear equality stays an equality;
maximizing $f$ is the same as minimizing $-f$. What is not free is convexity alone: flipping
the sign of an inequality flips the convexity of the function, so the direction of an
inequality must not be changed without checking again.

\begin{akibat}[Convex feasible region]
In a convex problem, the feasible region $X$ is convex.
\end{akibat}

\begin{proof}
The sublevel set $S_{i}=\{\mathbf{x}:f_{i}(\mathbf{x})\le0\}$ is convex for every $i$: if
$f_{i}(\mathbf{x})\le0$ and $f_{i}(\mathbf{y})\le0$, convexity gives
$f_{i}(t\mathbf{x}+(1-t)\mathbf{y})\le t\,f_{i}(\mathbf{x})+(1-t)f_{i}(\mathbf{y})\le0$.
The set $\{\mathbf{x}:h_{j}(\mathbf{x})=0\}$ is a hyperplane because $h_{j}$ is affine, hence
convex. The feasible region is the intersection of all of them, and an intersection of convex
sets is convex by the proposition in the first section.
\end{proof}

That corollary is what makes this whole chapter make sense: two feasible points can always be
joined without leaving the feasible region. The simplest polyhedron is the feasible region of a
two-variable linear program in Figure~\ref{fig:lp-verteks}: a quadrilateral, and its linear
objective is optimal at a vertex.

\begin{figure}[!ht]
\centering
\begin{tikzpicture}
\begin{axis}[width=9cm,height=6cm,xlabel={$x_{1}$},ylabel={$x_{2}$}]
\addplot[forget plot,fill=blue!8,draw=blue!50!black,thick]
  coordinates {(0,0) (4,0) (3,1.5) (0,3)} \closedcycle;
\addplot[domain=0:2.1,densely dotted,thick] {2.5-1.25*x};
\addlegendentry{profit $10$}
\addplot[domain=0:2.9,densely dotted,thick] {3.75-1.25*x};
\addlegendentry{profit $15$}
\addplot[domain=0:4.1,thick,red] {5.25-1.25*x};
\addlegendentry{profit $21$}
\addplot[forget plot,only marks,mark=*] coordinates {(3,1.5)};
\node[font=\scriptsize,below right] at (axis cs:3,1.5) {$(3;\,1.5)$};
\end{axis}
\end{tikzpicture}
\caption{The feasible region of the workshop linear program of Chapter 24:
$6x_{1}+4x_{2}\le24$ (plywood), $x_{1}+2x_{2}\le6$ (hours), $x_{1},x_{2}\ge0$, objective
$5x_{1}+4x_{2}$. The vertices $(0;0)$, $(4;0)$, $(3;1.5)$, $(0;3)$ have values $0$, $20$,
$21$, $12$. The level line $21$ touches only the vertex $(3;1.5)$; the line $20$ still cuts
the feasible region.}
\label{fig:lp-verteks}
\end{figure}


Now the statement promised by the chapter opening, and its proof is short precisely because
the feasible region is convex.

\begin{teorema}[A local minimum is a global minimum]
Let $f_{0}$ be convex on a convex set $X$, and let $\mathbf{x}^{*}\in X$ satisfy
$f_{0}(\mathbf{x}^{*})\le f_{0}(\mathbf{x})$ for every $\mathbf{x}\in X$ close enough to
$\mathbf{x}^{*}$. Then $f_{0}(\mathbf{x}^{*})\le f_{0}(\mathbf{x})$ for every
$\mathbf{x}\in X$. The set of all minimum points is convex as well.
\index{global optimum}
\end{teorema}

\begin{proof}
Suppose there is $\mathbf{y}\in X$ with $f_{0}(\mathbf{y})<f_{0}(\mathbf{x}^{*})$. For
$t\in(0,1]$ the point $\mathbf{z}_{t}=(1-t)\mathbf{x}^{*}+t\mathbf{y}$ lies in $X$ and
concentrates at $\mathbf{x}^{*}$ as $t\to0$, so for $t$ small enough it enters the
neighborhood where $\mathbf{x}^{*}$ is a local minimum. But convexity gives
\[
  f_{0}(\mathbf{z}_{t})\le(1-t)f_{0}(\mathbf{x}^{*})+t\,f_{0}(\mathbf{y})
  <f_{0}(\mathbf{x}^{*}),
\]
contrary to local minimality. For the second part, if $\mathbf{x}^{*},\mathbf{y}^{*}$ are both
minimum with value $p^{*}$, then $f_{0}(\mathbf{z}_{t})\le p^{*}$, whereas $p^{*}$ is the
smallest possible value, so $f_{0}(\mathbf{z}_{t})=p^{*}$ and $\mathbf{z}_{t}$ is a minimum
too.
\end{proof}

\begin{akibat}[Stationary condition without constraints]
If $f$ is convex and differentiable on an open set, then
$\nabla f(\mathbf{x}^{*})=\mathbf{0}$ if and only if $\mathbf{x}^{*}$ is a global minimum.
\end{akibat}

\begin{proof}
If $\nabla f(\mathbf{x}^{*})=\mathbf{0}$, the first-order test gives
$f(\mathbf{y})\ge f(\mathbf{x}^{*})$ for every $\mathbf{y}$. Conversely, if
$\mathbf{x}^{*}$ is a minimum, then for every direction $\mathbf{v}$ and small $t>0$ we have
$f(\mathbf{x}^{*}+t\mathbf{v})\ge f(\mathbf{x}^{*})$, hence the directional derivative
$\nabla f(\mathbf{x}^{*})^{\top}\mathbf{v}\ge0$. Since $\mathbf{v}$ may also be replaced by
$-\mathbf{v}$, it must be that $\nabla f(\mathbf{x}^{*})=\mathbf{0}$.
\end{proof}

\begin{tangan}
Materials: two meters of string, one small ring, a blindfold, and a friend.

Tie the string between two posts of equal height, then stretch it straight. Ask the friend to
close their eyes, to hang the ring at the middle of the string, and to slide it slowly left and
right with one finger. The ring slides to the lowest point, wherever it is released. Ask the
friend to mark the stopping point without opening their eyes; repeat five times from different
starting positions. All the marks must fall on one point.

Now tie a small knot in the string, a little to the left, so that the string forms two
hollows. Repeat the experiment. The ring stops in the left hollow when released near there,
and in the middle when released to the right of the knot: two stopping places, two valleys.

What must be visible: a string without a knot has only one lowest point, a string with a knot
has two, and a finger that stops in the left hollow does not know what happens on the right,
exactly like an algorithm that stops at the first stationary point.
\end{tangan}

Without convexity, that corollary is false. Take $f(x_{1},x_{2})=x_{1}x_{2}$ on the square
$[-1,1]^{2}$. Its gradient $(x_{2},x_{1})$ vanishes at the origin, whereas
$f(\mathbf{0})=0$ while $f(1,-1)=-1$. Its Hessian is
$\begin{pmatrix}0&1\\1&0\end{pmatrix}$ with eigenvalues $1$ and $-1$, so the second-order test
rejects convexity: the origin is a saddle, not a valley. Chapter 24 gives a more deceptive
example: rounding a continuous answer gives a profit of 19, whereas the best integer answer
gives 20.

\section{The Lagrangian and lower bounds from the dual}
\tingkat{advanced}

A constraint always has the form ``something must not be positive'', so we can put a price on
violating it: every time $f_{i}(\mathbf{x})$ is positive we pay a fine $\lambda_{i}$ per
unit, and every deviation $h_{j}(\mathbf{x})$ from zero is paid or received at $\nu_{j}$ per
unit. The resulting function is no larger than the original objective, provided the price is
nonnegative. A bad price gives a loose bound; a good price gives an exact bound.

\begin{definisi}[Lagrangian and dual function]
For a problem in standard form, its \emph{Lagrangian} is
\[
  L(\mathbf{x},\boldsymbol{\lambda},\boldsymbol{\nu})
  =f_{0}(\mathbf{x})+\sum_{i=1}^{m}\lambda_{i}f_{i}(\mathbf{x})
   +\sum_{j=1}^{p}\nu_{j}h_{j}(\mathbf{x}),
\]
with $\boldsymbol{\lambda}\in\mathbb{R}^{m}$ and $\boldsymbol{\nu}\in\mathbb{R}^{p}$ called
\emph{multipliers} or \emph{dual variables}. The \emph{dual function} is its infimum over the
original variable,
\[
  g(\boldsymbol{\lambda},\boldsymbol{\nu})
  =\inf_{\mathbf{x}\in D}L(\mathbf{x},\boldsymbol{\lambda},\boldsymbol{\nu}),
\]
and the \emph{dual problem} is to maximize $g$ under the condition
$\boldsymbol{\lambda}\ge\mathbf{0}$. A pair $(\boldsymbol{\lambda},\boldsymbol{\nu})$ is
called \emph{dual feasible} when $\boldsymbol{\lambda}\ge\mathbf{0}$ and
$g(\boldsymbol{\lambda},\boldsymbol{\nu})>-\infty$. Its optimal value is written $d^{*}$.
\index{Lagrangian}\index{duality}
\end{definisi}

\begin{proposisi}[The dual function is always concave]
$g$ is the infimum of affine functions in $(\boldsymbol{\lambda},\boldsymbol{\nu})$, because
for every fixed $\mathbf{x}$ the function $L(\mathbf{x},\cdot,\cdot)$ is affine. Hence $g$ is
concave, and the dual problem is always a convex problem, whatever the shape of the original
problem.
\end{proposisi}

\begin{proof}
Take two dual points $(\boldsymbol{\lambda},\boldsymbol{\nu})$ and
$(\boldsymbol{\lambda}',\boldsymbol{\nu}')$ and $t\in[0,1]$. For every $\mathbf{x}$,
linearity gives
\[
  L\bigl(\mathbf{x},\,t\boldsymbol{\lambda}+(1-t)\boldsymbol{\lambda}',\,
  t\boldsymbol{\nu}+(1-t)\boldsymbol{\nu}'\bigr)
  =t\,L(\mathbf{x},\boldsymbol{\lambda},\boldsymbol{\nu})
   +(1-t)\,L(\mathbf{x},\boldsymbol{\lambda}',\boldsymbol{\nu}').
\]
Take the infimum of both sides over $\mathbf{x}$; the infimum of a sum is no less than the sum
of the infima, so
$g\ge t\,g(\boldsymbol{\lambda},\boldsymbol{\nu})
+(1-t)g(\boldsymbol{\lambda}',\boldsymbol{\nu}')$, that is $g$ is concave.
\end{proof}

A pleasant corollary: maximizing a concave function under the constraint
$\boldsymbol{\lambda}\ge\mathbf{0}$ is always a convex problem. Every problem, even one full of
holes, has a convex twin; the twin does not always give the same answer, but it always gives a
bound.

\begin{teorema}[Weak duality]
For every dual feasible $(\boldsymbol{\lambda},\boldsymbol{\nu})$ and every primal feasible
$\mathbf{x}$ we have $L(\mathbf{x},\boldsymbol{\lambda},\boldsymbol{\nu})\le
f_{0}(\mathbf{x})$. Hence $g(\boldsymbol{\lambda},\boldsymbol{\nu})\le p^{*}$, that is
$d^{*}\le p^{*}$.
\end{teorema}

\begin{proof}
Take $\mathbf{x}\in X$. Every $f_{i}(\mathbf{x})\le0$ and $\lambda_{i}\ge0$, so
$\lambda_{i}f_{i}(\mathbf{x})\le0$. Every $h_{j}(\mathbf{x})=0$, so
$\nu_{j}h_{j}(\mathbf{x})=0$ whatever $\nu_{j}$ is. Therefore
$L(\mathbf{x},\boldsymbol{\lambda},\boldsymbol{\nu})\le f_{0}(\mathbf{x})$. That inequality
holds for every feasible $\mathbf{x}$, hence also for the infimum over all feasible
$\mathbf{x}$; the infimum over all of $D$ is no larger, so
$g(\boldsymbol{\lambda},\boldsymbol{\nu})\le p^{*}$. The last inequality is taken over all
dual feasible pairs.
\end{proof}

That proof uses neither convexity nor differentiability, so weak duality holds for integer
problems, combinatorial problems, and any problem at all in standard form. That is why dual
lower bounds are used in branch and bound in Chapter 24: the bound is valid even though the
relaxation is not the answer.

The difference $p^{*}-d^{*}$ is called the \emph{duality gap}. Since
$d^{*}\le p^{*}$, the gap is never negative. When the gap is zero, we say that
\emph{strong duality} holds; the next section discusses when that happens.
\index{weak duality}\index{duality gap}

Before that, one corollary used to check computed results. If there is a feasible point
$\mathbf{x}$ and a dual feasible pair $(\boldsymbol{\lambda},\boldsymbol{\nu})$ with
$f_{0}(\mathbf{x})=g(\boldsymbol{\lambda},\boldsymbol{\nu})$, then
$d^{*}\ge g=f_{0}(\mathbf{x})\ge p^{*}\ge d^{*}$, so all the inequalities become equalities:
the point is optimal, the pair is optimal, and the gap is zero. Equality of values is a
certificate of optimality that can be checked without knowing the answer in advance.

\begin{untukanak}
Have you ever haggled in a market? You name one price, the seller names another, and if the
bargaining stops halfway, neither of you knows the true price yet. But you already know one
thing: the seller will not go below a certain number, and you will not go above a certain
number. Those two numbers cage the true price, even though nobody has ever spoken it.

The idea of duality works exactly like that. Instead of looking for the answer, we attach a
price to each demand, then ask: with this price list, what is the worst outcome still possible?
Every sensible price list gives one number that the true answer cannot exceed. If some price
list gives a number equal to the answer we found, the matter is settled.

In this chapter we shall learn why such a price list always exists, why it sometimes fails to
close the gap, and why the gap that remains can be measured by a number.
\end{untukanak}

\section{Strong duality and Slater's condition}
\tingkat{advanced}

Convexity turns a local minimum into a global one, but it does not yet guarantee a zero
duality gap. What is missing is a condition about the \emph{geometry of feasibility}: if the
feasible region only touches the boundary of a constraint without entering it, there is no
sufficiently clean separating plane. That condition is named after Slater.

\begin{teorema}[Slater's condition]
Let the problem be convex, $f_{0},\dots,f_{m}$ differentiable, and
$h_{j}(\mathbf{x})=\mathbf{a}_{j}^{\top}\mathbf{x}-b_{j}$. If there is a feasible point
$\mathbf{x}\in\operatorname{ri}D$ with $f_{i}(\mathbf{x})<0$ for every $i$ and
$h_{j}(\mathbf{x})=0$ for every $j$, then strong duality holds, $d^{*}=p^{*}$, and the optimal
value of the dual is attained. If some of the inequality constraints turn out to be affine,
only the constraints that are not affine need be satisfied strictly.
\index{Slater}
\end{teorema}

A point with $f_{i}(\mathbf{x})<0$ is called \emph{strictly feasible}. The existence of one such
point is enough, and the condition demands nothing about the objective function.

The proof is not carried out here. The route is this: the set
$\mathcal{A}=\{(u,t):\exists\mathbf{x},\ f_{i}(\mathbf{x})\le u_{i},\
h_{j}(\mathbf{x})=0,\ f_{0}(\mathbf{x})\le t\}$ is convex, and the point $(0,p^{*})$ lies on
its boundary; the hyperplane separation theorem (Chapter 10 uses a similar idea for projection)
gives a supporting plane with coefficients $(\boldsymbol{\lambda},\boldsymbol{\nu})$ and value
$d^{*}$. The full proof is in Boyd and Vandenberghe section 5.3.2, and that theorem is what
justifies the step ``solve KKT and then trust the result''.

Slater is sufficient, not necessary, and the difference shows in the following small example.
Take $\min x$ with constraint $x^{2}\le0$. Its feasible region is a single point, $x=0$, so
there is no strictly feasible point and Slater fails. Its optimal value is $p^{*}=0$. The
Lagrangian $L=x+\lambda x^{2}$ gives $g(\lambda)=-1/(4\lambda)$ for $\lambda>0$, because the
minimum occurs at $x=-1/(2\lambda)$, so $\sup_{\lambda\ge0}g(\lambda)=0=d^{*}$. Strong duality
still holds, but that value is only approached, not attained: at every finite $\lambda$,
$g(\lambda)<0$. So what is lost when Slater fails is not always equality of values, but the
guarantee that the dual optimum is attained.

Conversely, Slater guarantees something that strong duality alone does not give: optimal
multipliers exist, so that the KKT equations have a solution to look for.

\section{The KKT conditions: necessary, and sufficient when convex}
\tingkat{advanced}

If the multipliers are already optimal, the optimal primal point must minimize the Lagrangian
without constraints, and there the gradient vanishes. Expanding that condition component by
component gives four groups of equations: the Karush-Kuhn-Tucker conditions.

\kamus{Lagrange multiplier}{pengali Lagrange}{The price per unit of violating a constraint; zero
means that constraint is not pressing, positive means that it is pressing.}

\begin{definisi}[KKT conditions]
A point $(\mathbf{x},\boldsymbol{\lambda},\boldsymbol{\nu})$ satisfies the \emph{KKT
conditions} when
\[
  \begin{array}{ll}
  f_{i}(\mathbf{x})\le0,\quad h_{j}(\mathbf{x})=0, & \text{(primal feasibility)}\\
  \lambda_{i}\ge0, & \text{(dual feasibility)}\\
  \lambda_{i}f_{i}(\mathbf{x})=0,\quad i=1,\dots,m, & \text{(complementary slackness)}\\
  \nabla f_{0}(\mathbf{x})+\sum_{i}\lambda_{i}\nabla f_{i}(\mathbf{x})
  +\sum_{j}\nu_{j}\nabla h_{j}(\mathbf{x})=\mathbf{0}. & \text{(stationarity)}
  \end{array}
\]
\index{KKT}\index{Lagrange multiplier}\index{complementary slackness}
\end{definisi}

The third group causes the most misunderstanding, so let us translate it: a multiplier may be
nonzero only for a constraint that is really pressing. A loose constraint is paid nothing,
because loosening it changes nothing. A zero multiplier on a pressing constraint is also
possible, and that is a sign of degeneracy.

The geometric side of the last group is visible in Figure~\ref{fig:kkt-singgung}.

\begin{figure}[!ht]
\centering
\begin{tikzpicture}[scale=1.42]
  \draw[->,layergrey] (-0.3,0)--(4.4,0) node[right,font=\small]{$x_{1}$};
  \draw[->,layergrey] (0,-0.3)--(0,4.4) node[above,font=\small]{$x_{2}$};
  \fill[layerteal!10] (0,0) circle (2);
  \begin{scope}
    \clip (-0.3,-0.3) rectangle (4.3,4.3);
    \draw[layergrey,densely dashed] (2.6,2.6) circle (1.25);
    \draw[layergrey,densely dashed] (2.6,2.6) circle (2.15);
    \draw[thick,accent] (2.6,2.6) circle (1.677);
  \end{scope}
  \draw[very thick,layerteal] (0,0) circle (2);
  \node[font=\scriptsize,layerteal,fill=white,inner sep=1pt,anchor=west]
    at (-0.4,1.9) {$f_{1}(\mathbf{x})=0$};
  \node[font=\scriptsize,accent,fill=white,inner sep=1pt,anchor=south east]
    at (1.1,2.02) {$f_{0}=2.81$};
  \node[font=\scriptsize,layergrey,fill=white,inner sep=1pt]
    at (2.6,3.9) {smaller value, not feasible};
  \fill[layergrey] (2.6,2.6) circle (1.4pt);
  \node[font=\scriptsize,layergrey,anchor=west] at (2.72,2.6) {$\mathbf{c}$};
  \fill[accent] (1.414,1.414) circle (1.9pt);
  \node[font=\scriptsize,accent,below left=-2pt] at (1.414,1.414)
    {$\mathbf{x}^{*}$};
  \draw[-{Latex},very thick,accent] (1.414,1.414)--(0.8,0.8)
    node[below right=-1pt,font=\scriptsize]{$\nabla f_{0}$};
  \draw[-{Latex},very thick,layerteal] (1.414,1.414)--(2.03,2.03)
    node[above left=-1pt,font=\scriptsize]{$\lambda\nabla f_{1}$};
\end{tikzpicture}
\caption{The contours of the objective $f_{0}(\mathbf{x})=(x_{1}-2.6)^{2}+(x_{2}-2.6)^{2}$,
concentric around the point $\mathbf{c}=(2.6;\,2.6)$, which is the unconstrained minimum and
falls outside the feasible region
$f_{1}(\mathbf{x})=x_{1}^{2}+x_{2}^{2}-4\le0$. The contour $f_{0}=2.81$ only touches the
constraint circle at $\mathbf{x}^{*}=(\sqrt{2};\sqrt{2})$; the deeper dashed contours have
smaller values but lie entirely outside the feasible region, whereas the outer ones still cut
the feasible region with larger values. At that point of tangency $\nabla f_{0}$ and
$\nabla f_{1}$ are parallel and oppositely directed, with a positive multiplier.}
\label{fig:kkt-singgung}
\end{figure}

Now the statement that gives this section its title, and which must be read carefully because
its two directions have different conditions.

\begin{teorema}[KKT is sufficient for optimality]
Let the problem be convex, $f_{0},\dots,f_{m}$ differentiable, and $h_{j}$ affine. If
$(\mathbf{x},\boldsymbol{\lambda},\boldsymbol{\nu})$ satisfies the KKT conditions, then
$\mathbf{x}$ is primal optimal, $(\boldsymbol{\lambda},\boldsymbol{\nu})$ is dual optimal, and
the duality gap is zero.
\end{teorema}

\begin{proof}
Because $\lambda_{i}\ge0$ and $f_{i}$ is convex, every term $\lambda_{i}f_{i}$ is convex;
because $h_{j}$ is affine, every term $\nu_{j}h_{j}$ is affine. So
$L(\cdot,\boldsymbol{\lambda},\boldsymbol{\nu})$ is convex in $\mathbf{x}$. The stationarity
condition says that its gradient vanishes at $\mathbf{x}$, so by the corollary in the third
section that point minimizes $L$ globally:
$g(\boldsymbol{\lambda},\boldsymbol{\nu})=L(\mathbf{x},\boldsymbol{\lambda},
\boldsymbol{\nu})$. Primal feasibility and complementary slackness remove all the constraint
terms:
\[
  g(\boldsymbol{\lambda},\boldsymbol{\nu})
  =f_{0}(\mathbf{x})+\sum_{i}\lambda_{i}f_{i}(\mathbf{x})
   +\sum_{j}\nu_{j}h_{j}(\mathbf{x})=f_{0}(\mathbf{x}).
\]
Because $g\le d^{*}\le p^{*}\le f_{0}(\mathbf{x})$ by weak duality, all those inequalities
become equalities. Hence $\mathbf{x}$ is primal optimal, the pair is dual optimal, and
$p^{*}-d^{*}=0$.
\end{proof}

The opposite direction is different in nature, and the difference is not technical.

\begin{teorema}[KKT is necessary under a qualification condition]
Let the problem be convex and differentiable, satisfy Slater's condition, and let
$\mathbf{x}^{*}$ be primal optimal with $p^{*}$ finite. Then there exists
$(\boldsymbol{\lambda}^{*},\boldsymbol{\nu}^{*})$ such that
$(\mathbf{x}^{*},\boldsymbol{\lambda}^{*},\boldsymbol{\nu}^{*})$ satisfies the KKT conditions.
\end{teorema}

The proof uses the hyperplane separation theorem on the achievable set, and is not repeated
here; see Boyd and Vandenberghe section 5.5.3. What must be remembered is the role of Slater in
it: without a qualification condition, an optimal point may fail KKT altogether, so that
``solve KKT'' is not a legitimate route for finding the answer.

The smallest counterexample is $\min x$ with constraint $x^{2}\le0$. Its optimal point is
$x^{*}=0$ and $\nabla f_{0}=1$, whereas $\nabla f_{1}=2x=0$ at that point. Stationarity demands
$1+\lambda\cdot0=0$, which is impossible. So that optimal point does not satisfy KKT, and what
fails is only the constraint qualification: the feasible region is one point, there is no
direction to move, there is no supporting plane.

The logical conclusion, and it is often quoted wrongly: in a convex problem, KKT is
\emph{sufficient} without extra conditions, but \emph{necessary} only when a constraint
qualification is present. The statement ``$\mathbf{x}$ is optimal if and only if there are
multipliers satisfying KKT'' demands both directions, so Slater is part of it.

One interpretation for remembering the multipliers. Suppose constraint $i$ is relaxed to
$f_{i}(\mathbf{x})\le u_{i}$, and the optimal value becomes $p^{*}(u)$. If $p^{*}$ is
differentiable at $u=0$, then
\[
  \lambda_{i}^{*}=-\frac{\partial p^{*}}{\partial u_{i}}(0).
\]
The multiplier is a shadow price: how much the objective value falls when that constraint is
relaxed by one unit. In the worked example later, relaxing the solubility limit from one half
to $0.6$ lowers the deviation from $2.5$ to $2.32$, whereas the shadow price predicts
$2.5-2(0.1)=2.30$. The difference $0.02$ is a second-order term, and its direction is always
the same: the linear forecast is a lower bound, so the improvement it promises is never
exceeded.

The mechanical interpretation, already used by Lagrange himself: two blocks are connected by a
spring and must not pass through each other or through the wall. Potential energy is the
objective function; at equilibrium, stationarity is the balance of forces, and the multiplier
is the contact force of the wall. A contact force only pushes, never pulls, that is
$\lambda_{i}\ge0$; it appears only when there is contact, that is complementary slackness (Boyd
and Vandenberghe section 5.5.4).

\begin{catatan}[Subgradients]
If $f$ is convex but not differentiable, the gradient is replaced by a \emph{subgradient}:
$\mathbf{s}$ is a subgradient of $f$ at $\mathbf{x}$ when
$f(\mathbf{y})\ge f(\mathbf{x})+\mathbf{s}^{\top}(\mathbf{y}-\mathbf{x})$ for every
$\mathbf{y}$; the set of all subgradients is written $\partial f(\mathbf{x})$. For
$f(x)=|x|$ at $x=0$ that set is the interval $[-1,1]$, because every slope between $-1$ and
$1$ gives a line below the graph. The optimality condition without constraints becomes
$\mathbf{0}\in\partial f(\mathbf{x}^{*})$, and KKT still holds with $\nabla f_{i}$ replaced by
a subgradient (Rockafellar section 23).\index{subgradient}
\end{catatan}

\section{Worked examples}
\tingkat{intermediate}

The following three computations are carried out to the last digit, with simple fractions so
that they can be checked by hand.

\emph{First, a convexity test through eigenvalues.} Take
$f(x_{1},x_{2})=x_{1}^{2}+4x_{1}x_{2}+5x_{2}^{2}$. Its Hessian is constant,
\[
  \nabla^{2}f=\begin{pmatrix}2&4\\4&10\end{pmatrix},
  \qquad
  \operatorname{tr}=12,\quad \det=20-16=4 .
\]
Its eigenvalues are $\lambda_{1,2}=6\pm\sqrt{36-4}=6\pm4\sqrt{2}$, that is $0.343146$ and
$11.656854$. Both are positive, so the Hessian is positive definite and $f$ is strictly convex.
For comparison, $g(x_{1},x_{2})=x_{1}^{2}-x_{2}^{2}$ has Hessian
$\operatorname{diag}(2,-2)$ with eigenvalues $2$ and $-2$: not convex, and its critical point at
the origin is a saddle. The function $h(x_{1},x_{2})=x_{1}^{2}$ has eigenvalues $0$ and $2$:
convex, not strict, and its minimum is the entire line $x_{1}=0$, so the answer is not unique.

\begin{tantangan}
Take $f(x_{1},x_{2})=x_{1}^{2}x_{2}^{2}$. Hold $x_{2}=c$ fixed: the graph is a parabola opening
upward, so every horizontal slice looks convex, and the same holds for every vertical slice. Is
$f$ convex on $\mathbb{R}^{2}$? If not, exhibit two points that prove it, then explain why
checking slice by slice is never enough.

Hint 1: do not move parallel to an axis. Move both coordinates at once, then check the midpoint.

Hint 2: compute $f(1,0)$, $f(0,1)$, and the value of $f$ at the midpoint of that segment; or
compute $\det\nabla^{2}f$ at $(1,1)$, which equals $-12x_{1}^{2}x_{2}^{2}$.

If you get stuck, go back to the section Convex functions and their two tests.
\end{tantangan}

\emph{Second, the paint mixer problem through KKT.} With $x_{1},x_{2}$ the amounts of the two
pigments, the problem is
\[
  \text{minimize } (x_{1}-2)^{2}+(x_{2}-2)^{2}
  \quad\text{subject to}\quad x_{1}+x_{2}=2,\quad x_{1}\le\tfrac12 .
\]
Write $g(\mathbf{x})=x_{1}-\tfrac12\le0$ and $h(\mathbf{x})=x_{1}+x_{2}-2=0$. Stationarity gives
\[
  2(x_{1}-2)+\lambda+\nu=0,\qquad 2(x_{2}-2)+\nu=0 ,
\]
and complementary slackness gives $\lambda(x_{1}-\tfrac12)=0$. If $\lambda=0$, the two
equations give $x_{1}=x_{2}$, and with $x_{1}+x_{2}=2$ we obtain $x_{1}=x_{2}=1$, which
violates $x_{1}\le\tfrac12$. So $\lambda>0$: $x_{1}=\tfrac12$ and $x_{2}=\tfrac32$. The second
equation gives $\nu=2(2-\tfrac32)=1$, the first gives
$\lambda=4-2x_{1}-\nu=4-1-1=2$. Since $\lambda=2\ge0$, that point is globally optimal by the
sufficiency theorem, with
\[
  p^{*}=\bigl(\tfrac12-2\bigr)^{2}+\bigl(\tfrac32-2\bigr)^{2}
  =\tfrac94+\tfrac14=\tfrac52=2.5 .
\]

Checking through the dual, not through trust. The dual function is computed by completing the
square in each variable:
\[
  \min_{x_{1}}\Bigl[(x_{1}-2)^{2}+(\lambda+\nu)x_{1}\Bigr]
  =4-\frac{(\lambda+\nu-4)^{2}}{4},
  \qquad
  \min_{x_{2}}\Bigl[(x_{2}-2)^{2}+\nu x_{2}\Bigr]
  =4-\frac{(\nu-4)^{2}}{4},
\]
so that
\[
  g(\lambda,\nu)=8-\frac{(\lambda+\nu-4)^{2}}{4}-\frac{(\nu-4)^{2}}{4}
  -\frac{\lambda}{2}-2\nu .
\]
At $(\lambda,\nu)=(2,1)$ we get $g=8-\tfrac14-\tfrac94-1-2=\tfrac52=p^{*}$: the gap is zero, as
the strong duality guaranteed by Slater requires, because $(0.4;\,1.6)$ is strictly feasible.
The partial derivatives give $-\tfrac12(\lambda+\nu-4)-\tfrac12(\nu-4)-2=0$ and
$-\tfrac12(\lambda+\nu-4)-\tfrac12=0$; from the second $\lambda+\nu=3$, then $\nu=1$ and
$\lambda=2$. That pair is the only maximizer.

The shadow prices can be checked with numbers. If the solubility limit is relaxed to $0.6$, the
answer is $x_{1}=0.6$, $x_{2}=1.4$, and $p^{*}=1.96+0.36=2.32$. The linear forecast
$p^{*}-\lambda\cdot0.1=2.5-0.2=2.30$ is off by $0.02$, and $0.02=\tfrac12\cdot4\cdot(0.1)^{2}$
exactly, because the second derivative of $p^{*}$ with respect to that limit is $4$. If the
demanded sum is raised to $2.1$ with the limit unchanged, the answer is $x_{1}=\tfrac12$,
$x_{2}=1.6$, and $p^{*}=2.25+0.16=2.41$; the forecast $2.5-\nu\cdot0.1=2.40$ is off by
$0.01=\tfrac12\cdot2\cdot(0.1)^{2}$. Both prices are slopes that can be used to decide, with an
error that can be computed.

\emph{Third, the duality gap on a nonconvex problem.} Take one variable,
\[
  \text{minimize } (x-2)^{2}\quad\text{subject to}\quad x^{3}-x\le0 .
\]
That constraint is equivalent to $x(x-1)(x+1)\le0$, so the feasible region is
$(-\infty,-1]\cup[0,1]$: two separate pieces, not convex. On $[0,1]$ the objective decreases,
so its smallest value is $1$ at $x=1$. On $(-\infty,-1]$ the objective increases to the left
from $x=-1$, which has value $9$. So $p^{*}=1$, attained at $x=1$.

The dual side is very different. The Lagrangian is $L=(x-2)^{2}+\lambda(x^{3}-x)$. For
$\lambda>0$ the cubic term dominates as $x\to-\infty$ and carries a negative sign, so
$L\to-\infty$ and $g(\lambda)=-\infty$: there is no lower bound to collect. Only $\lambda=0$
leaves a finite value, and there $g(0)=\min_{x}(x-2)^{2}=0$ at $x=2$. Hence $d^{*}=0$ and
\[
  p^{*}-d^{*}=1-0=1 ,
\]
a gap of $100\%$ of the optimal value itself. The lower bound from the dual here is not a loose
bound but a bound in the wrong place: $x=2$ is the unconstrained answer, and the constraint is
what determines the answer.

The same example shows why convexity cannot be removed from the sufficiency theorem. Check the
point $x=-1$, which is feasible because $(-1)^{3}-(-1)=0$. There
$\nabla f_{0}=2(x-2)=-6$ and $\nabla g=3x^{2}-1=2$, so stationarity demands $-6+2\lambda=0$, that
is $\lambda=3$; complementary slackness holds because the constraint is active. So
$(x,\lambda)=(-1,3)$ satisfies all four groups of KKT conditions, even though
$f_{0}(-1)=9$ while $p^{*}=1$.

\kertas{Repeat the second computation with the limit $x_{1}\le1$. Before computing, decide
whether that constraint is still pressing, then check whether its multiplier is still $2$.
After that compute the optimal value and compare it with $2.5$. As an extra exercise, find the
value of the limit $c$ at which the constraint stops pressing, and explain why the multiplier
must be zero exactly at that $c$.}

\section{Computational lab}
\tingkat{advanced}

The two blocks below were run with Python 3.14.4, numpy 2.5.2, and scipy 1.18.0. The output is
copied verbatim.

\begin{lab}
\begin{verbatim}
import numpy as np
from scipy.optimize import minimize

np.set_printoptions(precision=8, suppress=True, linewidth=78)

# mixer problem: min (x1-2)^2+(x2-2)^2  s.t. x1+x2=2, x1<=0.5
f  = lambda x: (x[0]-2)**2 + (x[1]-2)**2
df = lambda x: np.array([2*(x[0]-2), 2*(x[1]-2)])
cons = [{"type": "eq",   "fun": lambda x: x[0]+x[1]-2,
         "jac": lambda x: np.array([1., 1.])},
        {"type": "ineq", "fun": lambda x: 0.5-x[0],
         "jac": lambda x: np.array([-1., 0.])}]
r = minimize(f, [0., 0.], jac=df, constraints=cons, method="SLSQP",
             options={"ftol": 1e-15, "maxiter": 500})
xs = r.x
print("x* =", xs, " f(x*) =", f(xs))
print("multipliers (scipy):", r.multipliers)

# multipliers from stationarity, using the active set {eq, ineq}
A = np.column_stack([[1., 0.], [1., 1.]])   # columns: grad g, grad h
lam, nu = np.linalg.solve(A, -df(xs))
print("lambda, nu =", lam, nu)

# KKT residuals
print("h(x*)          =", xs[0]+xs[1]-2)
print("g(x*)          =", 0.5-xs[0])
print("lambda*g(x*)   =", lam*(0.5-xs[0]))
print("||grad L||     =", np.linalg.norm(df(xs) + lam*np.array([1.,0.])
                                         + nu*np.array([1.,1.])))

# dual function and its maximization
g = lambda lam, nu: (8 - (lam+nu-4)**2/4 - (nu-4)**2/4
                     - 0.5*lam - 2*nu)
print("g(2,1)         =", g(2., 1.))
rd = minimize(lambda z: -g(z[0], z[1]), [0.5, 0.5], method="SLSQP",
              bounds=[(0, None), (None, None)], options={"ftol": 1e-15})
print("dual argmax    =", rd.x, " d* =", -rd.fun)
print("gap p*-d*      =", f(xs) + rd.fun)

# Hessian eigenvalue test
for name, H in [("x1^2+4x1x2+5x2^2", np.array([[2.,4.],[4.,10.]])),
                ("x1^2-x2^2",        np.array([[2.,0.],[0.,-2.]])),
                ("x1^2",             np.array([[2.,0.],[0.,0.]]))]:
    print(name, "eig =", np.linalg.eigvalsh(H))
\end{verbatim}
Actual output:
\begin{verbatim}
x* = [0.5 1.5]  f(x*) = 2.500000000000001
multipliers (scipy): [-1.  2.]
lambda, nu = 2.0000000000000027 0.9999999999999987
h(x*)          = 0.0
g(x*)          = 6.106226635438361e-16
lambda*g(x*)   = 1.2212453270876738e-15
||grad L||     = 0.0
g(2,1)         = 2.5
dual argmax    = [1.99999997 0.99999999]  d* = 2.499999999999999
gap p*-d*      = 1.7763568394002505e-15
x1^2+4x1x2+5x2^2 eig = [ 0.34314575 11.65685425]
x1^2-x2^2 eig = [-2.  2.]
x1^2 eig = [0. 2.]
\end{verbatim}
\end{lab}

\begin{lab}
\begin{verbatim}
import numpy as np
from scipy.optimize import minimize

# nonconvex: min (x-2)^2  s.t. x^3 - x <= 0
f3 = lambda x: (x[0]-2)**2
r3 = minimize(f3, [0.5], constraints=[{"type": "ineq",
              "fun": lambda x: -(x[0]**3 - x[0])}],
              method="SLSQP", options={"ftol": 1e-14, "maxiter": 500})
print("x* =", r3.x, " p* =", f3(r3.x))

# Lagrangian for several lambda, evaluated far to the left
for lam in (0.0, 1e-3, 1e-2, 0.1):
    L = lambda x: (x[0]-2)**2 + lam*(x[0]**3 - x[0])
    print(f"lam={lam:<6} L(-1e4)={L(np.array([-1e4])):.3e}",
          f" L(-1e6)={L(np.array([-1e6])):.3e}")

# points that satisfy KKT but are not the answer
for xv in (-1.0, 1.0, 0.0):
    f0p = 2*(xv-2); gp = 3*xv*xv - 1
    lam = -f0p/gp if gp != 0 else float("nan")
    print(f"x={xv:>4}  f0'={f0p:>5}  g'={gp:>4}  lam={lam:>6}"
          f"  f0={ (xv-2)**2 }  feasible={xv**3-xv <= 0}")
\end{verbatim}
Actual output:
\begin{verbatim}
x* = [1.]  p* = 0.999999999995369
lam=0.0    L(-1e4)=1.000e+08  L(-1e6)=1.000e+12
lam=0.001  L(-1e4)=-9.000e+08  L(-1e6)=-9.990e+14
lam=0.01   L(-1e4)=-9.900e+09  L(-1e6)=-9.999e+15
lam=0.1    L(-1e4)=-9.990e+10  L(-1e6)=-1.000e+17
x=-1.0  f0'= -6.0  g'= 2.0  lam=   3.0  f0=9.0  feasible=True
x= 1.0  f0'= -2.0  g'= 2.0  lam=   1.0  f0=1.0  feasible=True
x= 0.0  f0'= -4.0  g'=-1.0  lam=  -4.0  f0=4.0  feasible=True
\end{verbatim}
\end{lab}

Four things must be read. First, the reported $x^{*}$ is $(0.5;\,1.5)$ with value
$2.500000000000001$; the difference of $10^{-15}$ from $2.5$ is binary fraction rounding.
Second, the KKT residuals computed separately are very small: the inequality constraint
$6.1\times10^{-16}$, the product $\lambda g$ about $1.2\times10^{-15}$, and the norm of the
Lagrangian gradient zero, so that point really satisfies the conditions, not almost satisfies
them.

Third, the easiest to trip over, \texttt{scipy} reports the multipliers $(-1,\,2)$, whereas
stationarity gives $\lambda=2.0000000000000027$ and $\nu=0.9999999999999987$, that is
$(\nu,\lambda)=(1,\,2)$ to fourteen digits. The second entry matches; the first has the opposite
sign. The sign convention for multipliers in SQP routines differs between libraries, so the sign
must not be guessed: check again through the stationarity equations. The number $2$ on the
inequality constraint happens to match, and that masks the sign error.

Fourth, $d^{*}=2.499999999999999$ and the gap $p^{*}-d^{*}=1.78\times10^{-15}$: zero to fifteen
digits. The dual maximizer returns $(\lambda,\nu)=(1.99999997,\;0.99999999)$, that is $(2,1)$ to
eight digits, with a difference from the \texttt{ftol} tolerance. Agreement of three routes, KKT
by hand, KKT numerically, and dual maximization, is what makes the number $2.5$ worth trusting.

The second output points the other way. The primal gives $x^{*}=1$ with $p^{*}\approx1$, as in
the hand computation. The Lagrangian with $\lambda=0.001$ is already worth $-9\times10^{8}$ at
$x=-10^{4}$ and $-9.99\times10^{14}$ at $x=-10^{6}$; adding $\lambda$ makes it worse in
proportion to the cube, so no $\lambda>0$ is dual feasible and $d^{*}=0$ from $\lambda=0$ is the
only bound available. Its three critical points: $x=-1$ with $\lambda=3$ satisfies KKT and has
value $9$; $x=1$ is the true answer with value $1$; $x=0$ fails because $\lambda=-4$.

\begin{untukdosen}
The sequence that works: start from the bowl and the marble, then the definition of a convex set,
then straight to the local minimum theorem. The first-order and second-order tests are better
placed after the students have seen why convexity is used; otherwise both feel like two formulas
to memorize.

Two typical misconceptions. First, ``convex means bulging'', so students reject the bowl as
convex; what helps is the line segment between two points inside the bowl. Second, KKT is taken
as necessary and sufficient without conditions, and then students are disappointed when the
example $x^{3}-x\le0$ gives a KKT point that is not the answer. That example is worth writing on
the board early, so that sufficiency appears as a consequence of convexity.

Board activity: give one two-variable QP with one inequality constraint, ask the class to guess
whether the constraint is pressing, then solve KKT by branching on $\lambda=0$ and $\lambda>0$,
then compute its dual function and show that the values agree. Students who can explain why the
$\lambda=0$ branch fails here are ready for duality.

A sign of readiness for the advanced part: a student can name one convex and one nonconvex
problem from their own reading, and state what changes in their optimality guarantees.
\end{untukdosen}

\section{Summary}

\begin{ringkasan}
\begin{itemize}
  \item A convex set contains every segment between two of its members. Intersection preserves
        convexity, union does not. Balls, hyperplanes, half-spaces, and polyhedra are convex;
        the set $\{x_{1}x_{2}\ge1\}$ is not.
  \item A convex function lies below its chords. First-order test: the graph lies above every
        tangent plane. Second-order test: the Hessian is positive semidefinite, necessary and
        sufficient. A positive definite Hessian is sufficient for strict convexity, but not
        necessary.
  \item In standard form with $f_{i}$ convex and $h_{j}$ affine, the feasible region is convex,
        so every local minimum is a global minimum. Without convexity, a stationary point may be
        a saddle, such as $x_{1}x_{2}$ at the origin.
  \item The Lagrangian puts a price on the constraints, and the dual function is its infimum over
        the original variable. The dual function is always concave, so the dual problem is always
        convex, and weak duality $d^{*}\le p^{*}$ holds without any condition at all.
  \item Strong duality $d^{*}=p^{*}$ needs Slater's condition, that is one strictly feasible
        point, which at the same time guarantees that optimal multipliers exist. Slater is
        sufficient, not necessary: for $\min x$
        with $x^{2}\le0$ strong duality holds without the dual optimum being attained.
  \item KKT is sufficient for optimality in a convex problem without extra conditions, but
        necessary only when there is a constraint qualification. The multipliers are shadow
        prices, and the duality gap can be positive, even infinite, in nonconvex problems:
        $\min (x-2)^{2}$
        with $x^{3}-x\le0$ gives $p^{*}=1$ and $d^{*}=0$.
\end{itemize}
\end{ringkasan}

\section{Exercises}

\begin{latihan}
\begin{enumerate}
  \item Decide whether $f=x_{1}^{2}+x_{1}x_{2}+x_{2}^{2}$, $g=x_{1}x_{2}$, and
        $h=2x_{1}^{2}-4x_{1}x_{2}+3x_{2}^{2}$ are convex, concave, or neither.
        State the Hessian and its eigenvalues; for the one that is positive semidefinite but not
        definite, explain the consequence for the uniqueness of the minimum.

  \item Prove that the intersection of two convex sets is convex, then give two convex sets
        whose union is not convex, with two concrete points and one mixture point that fails.

  \item Prove that if $f_{1},f_{2}$ are convex then $\max\{f_{1},f_{2}\}$ is convex.
        Check $f(x_{1},x_{2})=\max\{x_{1}^{2},x_{2}\}$ at $(0,0)$ and $(1,1)$: does
        the gradient exist, and does the subdifferential at the crease contain zero?
        (Hint: use the definition of a subgradient and the first-order test for nonsmooth
        functions.)

  \item Solve by KKT: minimize $x_{1}^{2}+x_{2}^{2}$ subject to $x_{1}+2x_{2}=4$
        and $x_{1}\ge1$. Determine the optimal point, both multipliers, and the optimal value,
        then check against \texttt{scipy.optimize.minimize}. (Hint: solve first without
        $x_{1}\ge1$; if that answer violates it, the constraint is pressing and $x_{1}=1$.)

  \item Compute the dual function of $\min x_{1}^{2}+x_{2}^{2}$ subject to $x_{1}+x_{2}\ge2$,
        determine its maximizer, and compare it with the primal optimal value $2$. Why is the
        gap zero even though there is no equality constraint?

  \item Assess this claim: \emph{strong duality guarantees that the primal optimum is attained.}
        True, false, or not enough evidence? Include a counterexample if needed.

  \item Assess this claim: \emph{complementary slackness states that the multiplier of an
        active constraint is always positive.} True, false, or not enough evidence? Show what
        that condition actually states.

  \item In the paint mixer problem, change the limit to $x_{1}\le1$ and compute again.
        Determine the value of the limit $c$ at which the constraint stops pressing, and show
        that its multiplier is zero at that $c$. (Hint: use the condition that the answer without
        the limit already satisfies $x_{1}\le c$.)

  \item For the nonconvex problem $\min (x-2)^{2}$ with $x^{3}-x\le0$, prove that
        $g(\lambda)=-\infty$ for every $\lambda>0$ by choosing a sequence
        $x_{k}\to-\infty$, then explain why $d^{*}=0$. Compute $p^{*}$ as well by
        examining both pieces of the feasible region.

  \item Prove that if $f$ is convex and differentiable and
        $\nabla f(\mathbf{x}^{*})=\mathbf{0}$, then $\mathbf{x}^{*}$ is a global minimum,
        without using strict convexity. Why does that proof fail for $f(x)=x^{3}$,
        and at which point is its condition violated?
\end{enumerate}
\end{latihan}

\begin{untukpeneliti}
What is really proved in this chapter: the closure of convexity under intersection; the
equivalence of convexity with both derivative tests, in both directions; the convexity of the
feasible region; the theorem that a local minimum is a global one; the concavity of the dual
function; weak duality without any condition at all; the sufficiency of KKT without a constraint
qualification. Two things are borrowed: Slater's theorem and the necessity of KKT, both through
the hyperplane separation theorem, with proofs in Boyd and Vandenberghe sections 5.3.2 and 5.5.3.

What is popularized beyond its proof. ``Convex means easy'': convexity guarantees that every KKT
point is global, but it gives neither a rate nor a guarantee that an algorithm reaches it. ``The
multiplier is always a price'': that interpretation is valid only when $p^{*}$ is differentiable
with respect to the right-hand side of the constraint. ``The gap measures nonconvexity'': the
example $x^{3}-x$ shows that the dual can be weaker than the minimum over the convex hull of the
feasible region.

Primary source: Boyd and Vandenberghe chapters 2 to 5; Nesterov, \emph{Lectures on Convex
Optimization}, chapters 1 to 3; Rockafellar, \emph{Convex Analysis}, sections 2 and 23.

Limits of validity: finite dimension, differentiable functions except in the note on
subgradients, and explicitly stated constraints. For implicit constraints, for instance ones
defined by a simulation, KKT cannot be written down and what remains are derivative-free methods.
Open question: the rate of first-order methods on large convex problems is still under
investigation.
\end{untukpeneliti}

\peta{This chapter follows Boyd and Vandenberghe, \emph{Convex Optimization}: sections 2.1 to
2.3 (convex sets, important examples, operations that preserve them), 3.1 (basic properties of
convex functions, including the first- and second-order tests), 4.1 and 4.2 (optimization
problems and convex problems), 5.1 and 5.2 (the Lagrangian, weak and strong duality, Slater's
condition), 5.5 (KKT and the mechanical interpretation). The viewpoint of the dual certificate is
taken from Bertsimas and Tsitsiklis, \emph{Introduction to Linear Optimization}, chapters 1 and
4; for subgradients, Rockafellar, \emph{Convex Analysis}, sections 2 and 23.}

\chapter{Linear and Integer Programming}

A woodworking shop receives 24 sheets of plywood and 6 hours of assembly-machine time for the
day. One bookshelf uses 6 sheets of plywood and 1 hour of assembly. One bench uses 4 sheets
and 2 hours. A shelf earns a profit of 5 units, a bench 4 units. How many shelves and how many
benches should it make?

The first instinct picks shelves, because the profit per unit is higher. That instinct is
wrong: one shelf consumes as much plywood as one and a half benches. The best answer of the
continuous model is 3 shelves and 1.5 benches, with profit 21, and half a bench cannot be
sold. Rounding down gives 19; 4 shelves without benches gives 20. So the best integer answer
is not the rounding of the best fractional answer, and the difference is an error of reasoning,
not an error of arithmetic.

The tools of the previous chapters are not enough. Chapter 8 taught us to solve
$A\mathbf{x}=\mathbf{b}$, but here there are fewer equations than variables, so there are
infinitely many answers and we must choose the best one. Chapter 23 gives the KKT conditions,
but in this problem the gradient of the objective is constant and never vanishes: the answer
clings to the boundary of the feasible region, not to a stationary point. The guarantee of
Chapter 21 that a local optimum is also global does not show where the optimum is, and it
disappears as soon as the variables are required to be integers.

For that reason this chapter takes three steps: writing the problem as a linear program and
recognizing that its answer lies at a corner of the feasible region; putting a price on every
constraint through duality, so that optimality can be checked without sweeping the whole
region; then facing the requirement of integrality, when the continuous model gives a valid
bound that is not necessarily attained.

\begin{untukanak}
Imagine you have 24 sheets of colored paper and 6 hours to make two kinds of classroom
decoration. The first decoration uses 6 sheets and 1 hour, the second uses 4 sheets and 2
hours. The first decoration is in greater demand, so its price is higher. How many of each
should you make?

If you try things out, there are many mixtures that fit. What is surprising is that the most
profitable mixture turns out to use half a decoration. Half a decoration cannot be sold, and
this is where common sense begins to deceive: if you round down, the result is not the best
one. There is another mixture, entirely whole, that is in fact more profitable, even though it
does not look like the rounding of anything.

The question that makes one curious is this. If there are thousands of possible mixtures, how
do we know that one mixture is the best without trying them all? And if your friend offers one
extra sheet of paper, what is a fair price for that sheet? In this chapter we will learn why
the best answer always lies at a corner of the possible region, why every material has a price
that can be computed, and why a fractional answer must be rounded with care, not with instinct.
\end{untukanak}

\section{From a production plan to standard form}
\tingkat{basic}

A production plan is a pair of numbers: how many shelves and how many benches. A plan is called
feasible when its use of materials does not exceed the supply. For $x_{1}$ shelves and
$x_{2}$ benches, the conditions are
\[
  6x_{1}+4x_{2}\le24,\qquad x_{1}+2x_{2}\le6,\qquad x_{1}\ge0,\quad x_{2}\ge0 .
\]
Each condition cuts the plane with one line and leaves one side; the intersection is a
quadrilateral with four corners. The profit is $5x_{1}+4x_{2}$, and we want that number as
large as possible.

The pattern of this kind of problem is called a linear program. The name is apt: everything
involved is a straight line, with no curves and no powers.

\kamus{linear program}{program linear}{A problem of finding the best value of a quantity
computed by a weighted sum, when all the rules take the shape of straight lines.}

\begin{definisi}[Linear program]
A linear program is the problem of maximizing or minimizing a linear function
$\mathbf{c}^{\top}\mathbf{x}$ over variables $\mathbf{x}\in\mathbb{R}^{n}$ subject to
constraints in the form of linear inequalities or equalities. The set of points satisfying all
the constraints is called the \emph{feasible region}. If the feasible region is empty, the
program is called \emph{infeasible}; if the objective value can be made as large as we please,
the program is called \emph{unbounded}.\index{linear program}
\end{definisi}

Software and textbooks use one uniform written pattern, and that is the pattern used
throughout this chapter.

\begin{definisi}[Standard form]
A linear program in \emph{standard form} is
\[
  \text{maximize } \mathbf{c}^{\top}\mathbf{x}
  \quad\text{with}\quad A\mathbf{x}=\mathbf{b},\qquad \mathbf{x}\ge\mathbf{0},
\]
where $A$ is $m\times n$ of rank $m$, $\mathbf{b}\in\mathbb{R}^{m}$, and
$\mathbf{c}\in\mathbb{R}^{n}$. Every linear program can be brought to this form without
changing its optimal value.\index{standard form}
\end{definisi}

Four translations do all that work. A constraint
$\mathbf{a}^{\top}\mathbf{x}\le b$ becomes an equality by adding a variable that records the
unused remainder; a constraint $\mathbf{a}^{\top}\mathbf{x}\ge b$ by subtracting one. A
variable that may take negative values is replaced by the difference of two nonnegative
variables, and maximizing $\mathbf{c}^{\top}\mathbf{x}$ is the same as minimizing
$-\mathbf{c}^{\top}\mathbf{x}$.

\kamus{slack variable}{variabel kendur}{A number recorded as the unused remainder, so that
``no more than'' turns into ``equal to''.}

For the shop above, the standard form uses $s_{1}$ as the plywood remainder and $s_{2}$ as the
remainder of assembly hours:
\[
  \begin{array}{rcl}
    6x_{1}+4x_{2}+s_{1} &=& 24,\\
    x_{1}+2x_{2}+s_{2} &=& 6,
  \end{array}
  \qquad x_{1},x_{2},s_{1},s_{2}\ge0 .
\]
Here $m=2$ and $n=4$, so $A$ is $2\times4$. The two slack variables are not decoration: later
both appear as columns in the simplex tableau and as a sign of which constraint is still loose.

Standard form also demands $\mathbf{b}\ge\mathbf{0}$, and that can always be arranged by
multiplying one row by $-1$. What cannot be arranged is the rank: rows of $A$ that are
linearly dependent must be discarded first, since otherwise the basis matrix chosen may be
singular.

\section{Geometry: why the answer is at a corner}
\tingkat{intermediate}

The four corners of that quadrilateral are $(0;0)$, $(4;0)$, $(3;1.5)$, and $(0;3)$, with
profits $0$, $20$, $21$, and $12$. The best is at a corner, and that is no coincidence: in a
linear program the best value can always be found at a corner, which is why the simplex method
only travels from corner to corner.

\begin{figure}[!ht]
\centering
\begin{tikzpicture}
\begin{axis}[width=9.3cm,height=6.3cm,
  xlabel={$x_{1}$ (shelves)},ylabel={$x_{2}$ (benches)},
  xmin=-0.35,xmax=5.1,ymin=-0.35,ymax=3.9,
  axis lines=left,xtick={0,1,2,3,4,5},ytick={0,1,2,3},
  tick label style={font=\scriptsize},label style={font=\scriptsize}]
\addplot[forget plot,fill=layerteal!10,draw=layerteal,thick]
  coordinates {(0,0) (4,0) (3,1.5) (0,3)} \closedcycle;
\addplot[forget plot,domain=0:2.4,densely dotted,thick,layerindigo] {3-1.25*x};
\addplot[forget plot,domain=0:3.6,densely dotted,thick,layerindigo] {4.5-1.25*x};
\addplot[forget plot,domain=0.2:4.2,thick,accent] {5.25-1.25*x};
\addplot[forget plot,only marks,mark=*,mark size=1.7pt,draw=ink]
  coordinates {(0,0) (4,0) (0,3)};
\addplot[forget plot,only marks,mark=*,mark size=2.6pt,draw=accent]
  coordinates {(3,1.5)};
\node[font=\scriptsize,anchor=north west] at (axis cs:0,0) {$(0;0)$};
\node[font=\scriptsize,anchor=north east] at (axis cs:4,0) {$(4;0)$};
\node[font=\scriptsize,anchor=south west] at (axis cs:0,3) {$(0;3)$};
\node[font=\scriptsize,anchor=west] at (axis cs:3.05,1.45) {$(3;1.5)$};
\node[font=\scriptsize,anchor=west,layerindigo] at (axis cs:0.15,1.15) {$z=12$};
\node[font=\scriptsize,anchor=west,layerindigo] at (axis cs:0.15,2.15) {$z=18$};
\node[font=\scriptsize,anchor=west,accent] at (axis cs:0.15,3.45) {$z=21$};
\draw[-{Stealth},thick,accent] (axis cs:2.3,2.75) -- (axis cs:3.25,3.35)
  node[font=\scriptsize,above,accent,pos=0.9] {profit rises};
\end{axis}
\end{tikzpicture}
\caption{The feasible region of the shop and three lines of equal profit. The polygon has four
vertices, $(0;0)$, $(4;0)$, $(3;1.5)$, $(0;3)$, with values $0$, $20$, $21$, $12$. The lines
$z=12$ and $z=18$ still cut the feasible region; the line $z=21$ only touches the vertex
$(3;1.5)$, so no feasible point has a higher value.}
\label{fig:ch22-daerah-optimum}
\end{figure}

Figure~\ref{fig:ch22-daerah-optimum} sharpens that coincidence. The three lines there are
lines of equal profit, $5x_{1}+4x_{2}=12$, $18$, and $21$. A higher profit line is farther from
the origin, and shifting it in that direction moves it out of the feasible region; the line
$21$ only touches the vertex $(3;1.5)$.

A corner has a more precise name, and the name does not refer to a picture.

\begin{definisi}[Vertex]
Let $P\subseteq\mathbb{R}^{n}$ be a convex set. A point $\mathbf{x}\in P$ is called a
\emph{vertex} or an \emph{extreme point} if there are no $\mathbf{y},\mathbf{z}\in P$ with
$\mathbf{y}\ne\mathbf{z}$ and $t\in(0,1)$ such that
$\mathbf{x}=t\mathbf{y}+(1-t)\mathbf{z}$. In other words, $\mathbf{x}$ is not the midpoint of
a segment between two other points of $P$.\index{vertex}
\end{definisi}

In the quadrilateral above, the point $(3;1.5)$ is a vertex: every segment inside the
quadrilateral that passes through it must end at that point itself. Conversely $(2;1)$ is not a
vertex, since it is the midpoint between $(1;1)$ and $(3;1)$.

\kamus{vertex}{verteks}{A corner of the feasible region, namely a point that never becomes the
midpoint between two other feasible points.}

The first step of the proof turns the negation into a direction that can be used.

\begin{lemma}[A point that is not a vertex gives a line]
If $\mathbf{x}\in P$ is not a vertex, then there is $\mathbf{d}\ne\mathbf{0}$ such that
$\mathbf{x}+\mathbf{d}\in P$ and $\mathbf{x}-\mathbf{d}\in P$.
\end{lemma}

\begin{proof}
Since $\mathbf{x}$ is not a vertex, there are $\mathbf{y}\ne\mathbf{z}\in P$ and $t\in(0,1)$
with $\mathbf{x}=t\mathbf{y}+(1-t)\mathbf{z}$. Set
$s=\tfrac12\min\{t,1-t\}>0$ and $\mathbf{d}=s(\mathbf{y}-\mathbf{z})$. Then
$\mathbf{x}+\mathbf{d}=(t+s)\mathbf{y}+(1-t-s)\mathbf{z}$ and
$\mathbf{x}-\mathbf{d}=(t-s)\mathbf{y}+(1-t+s)\mathbf{z}$. Since $0<t\pm s<1$, both
points are convex combinations of $\mathbf{y}$ and $\mathbf{z}$, hence members of $P$ because
$P$ is convex. Also $\mathbf{d}\ne\mathbf{0}$, since $\mathbf{y}\ne\mathbf{z}$ and $s>0$.
\end{proof}

Now the statement that makes the simplex method reasonable. In everyday language: if the
feasible region is enclosed and the objective does not grow in any direction, then the highest
summit is at one of the corners, not in the middle of the plane.

\begin{teorema}[The optimum is attained at a vertex]
Let $P=\{\mathbf{x}:A\mathbf{x}\le\mathbf{b}\}$ be nonempty and bounded, and suppose
$\mathbf{c}^{\top}\mathbf{x}$ is bounded above on $P$. Then there is a vertex
$\mathbf{x}^{*}$ of $P$ with $\mathbf{c}^{\top}\mathbf{x}^{*}
=\max_{\mathbf{x}\in P}\mathbf{c}^{\top}\mathbf{x}$.
\index{optimum at a vertex}
\end{teorema}

\begin{proof}
The optimal set $F=\{\mathbf{x}\in P:\mathbf{c}^{\top}\mathbf{x}=p^{*}\}$ is nonempty, closed,
bounded, and convex, so $F$ is also a polytope. Write the constraints of $P$ as
$\mathbf{a}_{i}^{\top}\mathbf{x}\le b_{i}$, $i=1,\dots,k$. Choose $\mathbf{x}^{*}\in F$
that maximizes the number of active constraints, namely the constraints with
$\mathbf{a}_{i}^{\top}\mathbf{x}^{*}=b_{i}$.

Suppose $\mathbf{x}^{*}$ is not a vertex. By the lemma there is $\mathbf{d}\ne\mathbf{0}$
with $\mathbf{x}^{*}\pm\mathbf{d}\in F$. Every constraint active at
$\mathbf{x}^{*}$ is also active at both of those points, so
$\mathbf{a}_{i}^{\top}\mathbf{d}=0$ for all active $i$. Consider the line
$\mathbf{x}^{*}+s\mathbf{d}$ for $s\in\mathbb{R}$. Its intersection with $P$ is a closed
interval $[s_{-},s_{+}]$ that contains $[-1,1]$, and because $P$ is bounded both endpoints are
finite. The endpoint $s_{+}$ is determined by the constraint with $\mathbf{a}_{i}^{\top}\mathbf{d}>0$
that is reached first; that constraint is active at $s_{+}$ but not at $\mathbf{x}^{*}$, so the
new point is feasible and active at more constraints. Its value stays $p^{*}$ because
$\mathbf{c}^{\top}(\mathbf{x}^{*}+s_{+}\mathbf{d})=p^{*}+s_{+}\mathbf{c}^{\top}
\mathbf{d}=p^{*}$, and $\mathbf{c}^{\top}\mathbf{d}=0$ since both points
$\mathbf{x}^{*}\pm\mathbf{d}$ have value $p^{*}$. That new point is a member of $F$ with more
active constraints, contradicting the choice of $\mathbf{x}^{*}$; hence $\mathbf{x}^{*}$ is a
vertex.
\end{proof}

\begin{bacakeras}
Two voices, read alternately with a friend. Tap the table once each time the word ``summit'' is
spoken.

``Where is the best answer?'' ``At the summit.''
``Why not in the middle?'' ``Because in the middle there is a neighbor on two sides.''
``If there is a neighbor, why is that wrong?'' ``Because there is still a direction to move.''
``Move in which direction?'' ``In the direction that climbs.''
``If there is no climbing direction?'' ``Then we are already standing at the summit.''

Repeat the three middle lines, slower and slower, with the same taps: in the middle there is a
neighbor, in the middle one can still move, a summit has no climbing direction.

The line the reader must repeat, once, in a clear voice: the best answer is at the summit, not
in the middle, and a summit is recognized by the absence of a climbing direction.
\end{bacakeras}

That proof uses the boundedness of $P$. If $P$ is unbounded but the objective is bounded above,
the conclusion is still true; the route goes through the decomposition of $P$ into the convex
hull of its vertices plus the cone of recession directions, and the complete proof is in
\texttt{WSPrinted.pdf}, in the section on the structure of polyhedral sets.

The number of vertices is finite, but not always small. In standard form, a vertex corresponds
to a choice of $m$ linearly independent columns of $A$: the chosen columns determine
$\mathbf{x}_{B}=B^{-1}\mathbf{b}$, the other columns are zero, and the point is feasible when
$\mathbf{x}_{B}\ge\mathbf{0}$. There are $\binom{n}{m}$ candidates, which for $n=100$ and
$m=50$ exceeds $10^{29}$; checking them all is not a sensible plan. This is where the simplex
method comes in: it walks from a vertex to a better neighboring vertex.

\begin{tangan}
Materials: a tiled floor, two lengths of string or paper tape, and a handful of buttons.

Stretch two straight strings as two boundaries; one for plywood, one for assembly hours. Stand
on a corner tile and call that tile the origin: each tile to the right adds one shelf, each
tile forward adds one bench.

Walk along the strings and stop only at tiles where two strings meet. Put a button there, and
only there. What must be visible: all the buttons cling to the turns of the fence, not one of
them stands in the middle of the yard. After that, count the profit at each button on your
fingers, then move one button to the most profitable corner.

Shift one string one tile inward, then repeat: the buttons move, the winner is still the
button at the corner.
\end{tangan}

\section{The simplex method: one tableau, two pivots}
\tingkat{intermediate}

The idea is simple. Stand at one vertex, check whether any edge direction raises the profit; if
one does, walk to the vertex at its end. If none does, stop: no neighbor is better, which in a
linear program means no other point is better.

That edge direction is determined by the \emph{reduced cost}: how much the profit increases
when a variable that is now zero is raised by one unit, at the optimal expense of the basis
variables.

\kamus{basic feasible solution}{solusi layak basis}{A point obtained by setting some variables
to zero and then solving for the rest, which happens to still satisfy all the bounds; such a
point is always a vertex.}

For the shop above, start from the vertex $(0;0)$, namely $s_{1}=24$ and $s_{2}=6$. Its first
tableau contains the profit row $z$ and one row for each constraint:
\begin{center}
\begin{tabular}{c|cccc|c}
\hline
basis & $x_{1}$ & $x_{2}$ & $s_{1}$ & $s_{2}$ & right-hand side \\
\hline
$z$   & $-5$ & $-4$ & $0$ & $0$ & $0$ \\
$s_{1}$ & $6$ & $4$ & $1$ & $0$ & $24$ \\
$s_{2}$ & $1$ & $2$ & $0$ & $1$ & $6$ \\
\hline
\end{tabular}
\end{center}
The $z$ row is written as $z-5x_{1}-4x_{2}=0$; a negative coefficient means that raising that
variable raises $z$, so the column with the most negative coefficient enters, here $x_{1}$ with
$-5$.

The ratio test decides who leaves: divide the right-hand side by the positive entries in the
entering column, then take the smallest ratio. $24/6=4$ and $6/1=6$, so $s_{1}$ leaves and the
pivot is $6$. Dividing the pivot row by $6$ and then clearing that column from the other rows
gives
\[
  \begin{array}{rcl}
    x_{1}+\tfrac{2}{3}x_{2}+\tfrac{1}{6}s_{1} &=& 4,\\
    \tfrac{4}{3}x_{2}-\tfrac{1}{6}s_{1}+s_{2} &=& 2,\\
    z-\tfrac{2}{3}x_{2}+\tfrac{5}{6}s_{1} &=& 20 .
  \end{array}
\]
The point is now $(4;0)$ with profit $20$, and a negative coefficient remains: $x_{2}$ with
$-2/3$. The ratio test gives $4/(2/3)=6$ and $2/(4/3)=1.5$, so $s_{2}$ leaves and the pivot is
$4/3$. After one more round,
\[
  \begin{array}{rcl}
    x_{1}+\tfrac{1}{4}s_{1}-\tfrac{1}{2}s_{2} &=& 3,\\
    x_{2}-\tfrac{1}{8}s_{1}+\tfrac{3}{4}s_{2} &=& \tfrac{3}{2},\\
    z+\tfrac{3}{4}s_{1}+\tfrac{1}{2}s_{2} &=& 21 .
  \end{array}
\]
The $z$ row now has no negative coefficient, so no direction raises the profit:
$(3;1.5)$ is optimal with profit $21$. All the arithmetic above was rechecked with exact
fractions, so that no decimal rounding slipped in.

That last row is more than a stopping sign. The coefficients of $s_{1}$ and
$s_{2}$ there, $3/4$ and $1/2$, are the prices of the two materials. Those numbers come from the
basis matrix, $\mathbf{y}^{\top}=\mathbf{c}_{B}^{\top}B^{-1}$, and
$\mathbf{b}^{\top}\mathbf{y}=24\cdot\tfrac34+6\cdot\tfrac12=21$, the same as the profit.
That equality is duality, and the next part explains why it is itself already
a proof of optimality.

Two remarks. The most-negative-coefficient rule does not guarantee termination, since there are
small examples with an endless cycle; Bland's rule, which picks the smallest index, guarantees
termination. Modern software uses the dual simplex or interior-point methods, but reports the
same vertices and the same prices.

\section{Duality: a price for every constraint}
\tingkat{advanced}

Every scarce material has a hidden price. One extra sheet of plywood raises the profit, and the
size of that rise is the price of the sheet to the shop. That price is not written in the
original problem, but it can be computed, and duality is the way to compute it.

Multiply the plywood constraint by $y_{1}\ge0$ and the hours constraint by $y_{2}\ge0$, then
add them. The result is
\[
  (6y_{1}+y_{2})x_{1}+(4y_{1}+2y_{2})x_{2}\le24y_{1}+6y_{2}.
\]
If $6y_{1}+y_{2}\ge5$ and $4y_{1}+2y_{2}\ge4$, the left-hand side is at least
$5x_{1}+4x_{2}$, which is the profit. So $24y_{1}+6y_{2}$ becomes an upper bound on the profit
for every price satisfying those two conditions, and the best bound is made as small as
possible.

\begin{definisi}[Primal and dual pair]
For the linear program
\[
  \text{(P)}\quad \text{maximize } \mathbf{c}^{\top}\mathbf{x}
  \ \text{ with } A\mathbf{x}\le\mathbf{b},\ \mathbf{x}\ge\mathbf{0},
\]
\emph{its dual program} is
\[
  \text{(D)}\quad \text{minimize } \mathbf{b}^{\top}\mathbf{y}
  \ \text{ with } A^{\top}\mathbf{y}\ge\mathbf{c},\ \mathbf{y}\ge\mathbf{0}.
\]
The variables $\mathbf{y}$ are called \emph{dual variables} or \emph{shadow prices}.
\index{duality}\index{shadow price}
\end{definisi}

For the shop above, the dual minimizes $24y_{1}+6y_{2}$ subject to
$6y_{1}+y_{2}\ge5$, $4y_{1}+2y_{2}\ge4$, and $y_{1},y_{2}\ge0$. Its solution is
$y_{1}=3/4$, $y_{2}=1/2$, with value $21$, the same as the best profit.

That the upper bound from the dual never falls below the primal value is the first thing to
hold on to, and its proof is three lines.

\begin{teorema}[Weak duality]
For every feasible $\mathbf{x}$ of (P) and every feasible $\mathbf{y}$ of (D) we have
$\mathbf{c}^{\top}\mathbf{x}\le\mathbf{b}^{\top}\mathbf{y}$.
\index{weak duality}
\end{teorema}

\begin{proof}
Since $\mathbf{y}\ge\mathbf{0}$ and $\mathbf{b}-A\mathbf{x}\ge\mathbf{0}$, we have
$\mathbf{b}^{\top}\mathbf{y}\ge\mathbf{y}^{\top}A\mathbf{x}$. Since
$\mathbf{x}\ge\mathbf{0}$ and $A^{\top}\mathbf{y}-\mathbf{c}\ge\mathbf{0}$, we have
$\mathbf{y}^{\top}A\mathbf{x}\ge\mathbf{c}^{\top}\mathbf{x}$. The two chains together give the
required result.
\end{proof}

The practical consequence is large: a pair of feasible points, one primal and one dual, with
the same value already proves both optimal without inspecting any other point. In the shop
above, $(3;1.5)$ and $(0.75;\,0.5)$ give the value $21$, so $21$ can be neither improved nor
lowered.

What weak duality does not give is a guarantee that the two values always meet. That guarantee
is the content of the next theorem.

\begin{teorema}[Strong duality]
If (P) or (D) has a finite optimal value, then both have optimal values and those values are
equal.
\end{teorema}

The proof is not carried out here. There are two routes: through Farkas' lemma on systems of
linear inequalities, or through the fact that the simplex stops at a tableau with no negative
coefficient, and that tableau at the same time gives a feasible dual $\mathbf{y}$ of the same
value. The complete proof is in \texttt{WSPrinted.pdf}, in the section on linear programming
duality, and in Boyd and Vandenberghe, section 5.2.

From that equality comes the condition used to check a computation quickly. In everyday
language: a constraint that is still loose does not need to be paid for, and a variable that is
used must pay its price exactly.

\begin{akibat}[Complementary slackness]
If $\mathbf{x}$ and $\mathbf{y}$ are both optimal, then
\[
  y_{i}\bigl(b_{i}-\mathbf{a}_{i}^{\top}\mathbf{x}\bigr)=0 \ \text{ for every } i,
  \qquad
  x_{j}\bigl((A^{\top}\mathbf{y})_{j}-c_{j}\bigr)=0 \ \text{ for every } j .
\]
Conversely, if $\mathbf{x}$ is feasible for (P), $\mathbf{y}$ is feasible for (D), and both
groups of equalities hold, then both are optimal.
\index{complementary slackness}
\end{akibat}

\begin{proof}
Since both are optimal, $\mathbf{c}^{\top}\mathbf{x}=\mathbf{b}^{\top}\mathbf{y}$, and the
chain in the proof of weak duality becomes a chain of equalities. Every term summed inside it is
nonnegative, so every term is zero; those terms are exactly
$y_{i}(b_{i}-\mathbf{a}_{i}^{\top}\mathbf{x})$ and
$x_{j}((A^{\top}\mathbf{y})_{j}-c_{j})$. The converse direction: if all terms are zero, the same
chain gives equality of values, and by weak duality that value is optimal for both problems.
\end{proof}

Now the shadow prices themselves. If the right-hand side is perturbed, the optimal value
changes, and the rate of change is the component of $\mathbf{y}$. Adding one sheet of plywood
raises the profit from $21$ to $21.75$; adding one assembly hour raises it to
$21.5$. So a sheet of plywood is worth $0.75$ units of profit, an assembly hour
$0.5$.

That price does not hold forever, and its range of validity can be computed. As long as the set
of active constraints does not change, $\mathbf{x}_{B}=B^{-1}\mathbf{b}$ is linear in
$\mathbf{b}$, so its value is linear too. For the shop above,
\[
  B^{-1}=\begin{pmatrix}\tfrac14&-\tfrac12\\[2pt]-\tfrac18&\tfrac34\end{pmatrix},
  \qquad
  \mathbf{x}_{B}=\begin{pmatrix}\tfrac14b_{1}-\tfrac12b_{2}\\[2pt]
  -\tfrac18b_{1}+\tfrac34b_{2}\end{pmatrix}.
\]
With $b_{2}=6$ fixed, the condition $\mathbf{x}_{B}\ge\mathbf{0}$ gives $b_{1}\in[12,36]$;
outside that interval one basis variable becomes negative and the slope changes. With $b_{1}=24$
fixed, the interval is $b_{2}\in[4,12]$; a shadow price is therefore always
reported together with its interval, not as a bare number.

One warning is useful from the start: shadow prices are an idea of the continuous model. In an
integer model, one extra sheet of plywood may raise the profit by $1$ or not at all, and there is
no single price that summarizes both.

\section{The LP relaxation: a valid bound, a deceptive rounding}
\tingkat{intermediate}

If shelves and benches may be fractional, the problem above is finished: $3$ shelves and
$1.5$ benches. In reality they may not, and the relation between the continuous problem and
the integer problem must be stated with care.

\begin{definisi}[Integer program, relaxation, and integrality gap]
An \emph{integer program} is a linear program in which some or all of the variables are
required to be integer. If all the variables are required to be integer, it is called a pure
integer program. Its \emph{LP relaxation} is the same linear program without the integrality
requirement. For a maximization problem with optimal values $z^{*}_{\mathrm{LP}}$ and
$z^{*}_{\mathrm{IP}}$, the \emph{integrality gap} is $z^{*}_{\mathrm{LP}}-z^{*}_{\mathrm{IP}}$.
\index{integer program}\index{relaxation}\index{integrality gap}
\end{definisi}

\kamus{integrality gap}{celah integralitas}{The distance between the bound that the continuous
model promises and the actual integer answer; that distance measures how far the continuous
model misleads.}

The reason the relaxation is useful is one sentence: the set of feasible integer points is part
of the set of feasible continuous points, so maximizing over the larger set cannot give a
smaller value.

\begin{proposisi}[The relaxation gives a valid bound]
For a maximization problem, $z^{*}_{\mathrm{IP}}\le z^{*}_{\mathrm{LP}}$. For a minimization
problem, the inequality reverses.
\index{relaxation bound}
\end{proposisi}

\begin{proof}
Every integer point that is feasible for the integer program is also feasible for its
relaxation, since the relaxation only removes conditions, it does not add any. So the feasible
region of the integer program is contained in the feasible region of the relaxation. A maximum
over a larger set is not smaller than the maximum over its subset.
\end{proof}

The bound is valid, but not always tight, and the difference can be very misleading when it is
used for rounding. Rounding fails in two ways, and both appear in our small example.

The first way: the rounded result is infeasible. The point $(3;1.5)$ rounded up becomes
$(3;2)$. Its profit is $23$, higher than $21$, but its plywood is
$6\cdot3+4\cdot2=26$ sheets, exceeding $24$: that is not a production plan, only a row of
numbers.

The second way, subtler: the rounded result is feasible but not the best. Rounding down gives
$(3;1)$ with profit $19$, although $(4;0)$ has profit $20$ and is entirely feasible.
Rounding the coordinates one by one in the nearest direction therefore gives a worse answer
than another integer point that is farther from $(3;1.5)$.

Both failures are a consequence of the shape of the feasible region. The set of feasible
integer points is not convex: $(3;1)$ and $(4;0)$ are members, but their midpoint
$(3.5;\,0.5)$ is not, because it is not integer. On a nonconvex set the vertex-optimum theorem
does not apply: no vertex is guaranteed to contain the answer, and no neighbor can be relied
upon.

\begin{figure}[!ht]
\centering
\begin{tikzpicture}
\begin{axis}[width=9.3cm,height=6.3cm,
  xlabel={$x_{1}$ (shelves)},ylabel={$x_{2}$ (benches)},
  xmin=-0.35,xmax=5.1,ymin=-0.35,ymax=3.9,
  axis lines=left,xtick={0,1,2,3,4,5},ytick={0,1,2,3},
  tick label style={font=\scriptsize},label style={font=\scriptsize}]
\addplot[forget plot,fill=layerteal!8,draw=layerteal,thick]
  coordinates {(0,0) (4,0) (3,1.5) (0,3)} \closedcycle;
\addplot[forget plot,only marks,mark=*,mark size=1.5pt,draw=layerindigo]
  coordinates {(0,0) (1,0) (2,0) (3,0) (4,0)
               (0,1) (1,1) (2,1) (3,1)
               (0,2) (1,2) (2,2) (0,3)};
\addplot[forget plot,only marks,mark=square*,mark size=2.4pt,draw=layerjade]
  coordinates {(4,0)};
\addplot[forget plot,only marks,mark=*,mark size=1.5pt,draw=layerplum]
  coordinates {(3,1)};
\addplot[forget plot,only marks,mark=o,mark size=3pt,thick,accent]
  coordinates {(3,1.5)};
\draw[densely dashed,layergrey] (axis cs:3,1) -- (axis cs:3,2.35);
\node[font=\scriptsize,anchor=west,layerjade] at (axis cs:4.1,-0.05) {$(4;0)$, $z=20$};
\node[font=\scriptsize,anchor=east,layerplum] at (axis cs:2.9,0.85) {$(3;1)$, $z=19$};
\node[font=\scriptsize,anchor=west,accent] at (axis cs:3.15,1.6) {$(3;1.5)$, $z=21$};
\node[font=\scriptsize,anchor=north,layergrey] at (axis cs:3,-0.05) {};
\end{axis}
\end{tikzpicture}
\caption{The integer lattice versus the relaxation point. There are $13$ feasible integer points,
marked with filled dots. The relaxation gives $(3;1.5)$ with value $21$, which stands midway
between $(3;1)$ with value $19$ and $(3;2)$, which is infeasible since
$6\cdot3+4\cdot2=26>24$. The best integer answer $(4;0)$ with value $20$ is in fact farther
from the relaxation point.}
\label{fig:ch22-kisi-relaksasi}
\end{figure}

Figure~\ref{fig:ch22-kisi-relaksasi} contains everything: thirteen feasible integer points, and
one relaxation point standing between two points. The one in between is not an integer point,
and the nearest is not necessarily the best.

Rounding also has no single meaning in many dimensions. With two variables, ``nearest integer
point'' can still be pictured; with a hundred variables, the nearest is not necessarily
feasible, and the feasible one is not necessarily the best. The practice is not to round, but to
use the relaxation bound to prune the search.

\begin{tantangan}
Maximize $6x_{1}+5x_{2}$ subject to $2x_{1}+3x_{2}\le12$, $4x_{1}+x_{2}\le13$, and
$x_{1},x_{2}\ge0$ integer. The relaxation gives $(2.7;\,2.2)$ with value $27.2$.
Rounding to the nearest integer, $(3;2)$, gives value $28$: higher than the value of any
relaxation point that is still feasible. Show why $28$ cannot be trusted, then determine the
correct integer answer together with its value.

Hint 1: check feasibility, not only the value. A higher number may come from a point that is
not a plan.

Hint 2: squeeze $x_{1}$. The largest integer that is still feasible is $3$; for
$x_{1}=3$, compute the bound on $x_{2}$ from both constraints and take the smaller one.

If you get stuck, go back to the section The LP relaxation: a valid bound, a deceptive
rounding.
\end{tantangan}

\section{Why integer programs are harder}
\tingkat{intermediate}

Up to here the question is fair: if the relaxation only adds one requirement, why is an integer
program not solved by solving its relaxation and then rounding? The short answer: the difficulty
is not in finding an answer, but in making sure that the answer is the best one.

Three differences explain it.

First, there is no simple certificate of optimality. In a linear program, a pair $\mathbf{x}$
and $\mathbf{y}$ of equal value is a finished proof that can be checked in time proportional to
the size of the data. In an integer program, the prices from the relaxation give only a bound,
not an integer point that attains it: the dual of the relaxation gives a lower bound, not an
answer.

Second, the difference between bound and answer can be large. In the vertex cover of the
complete graph on $n$ vertices, the relaxation allows $x_{v}=1/2$ for all $v$ with value $n/2$,
whereas the integer answer is $n-1$; the ratio $2(n-1)/n$ approaches $2$. The complete
computation is in the computational lab and in Exercise 10.

Third, the computational cost is of a different class, and the number of candidates grows
exponentially. A linear program is solved in time polynomial in the input size; Khachiyan's
result through the ellipsoid method and Karmarkar's through interior-point methods are the two
landmarks. The decision version of the integer program, ``is there a feasible integer point of
value at least $K$'', is NP-complete, and Karp's list of twenty-one problems contains knapsack,
which by itself has $2^{n}$ subsets. No polynomial algorithm is known for integer programming,
and if one existed, many other things would collapse with it.

\begin{untukanak}
Have you ever laid tiles on a floor? You measure the length of the room, then divide by the size
of the tile. The result is rarely a whole number. If you round down, there is a gap at the end
of the floor. If you round up, you cut a tile and throw away the rest. And sometimes the best
answer is neither of the two, but a laying pattern that is altogether different: start from the
middle, or use cut tiles on both sides.

Someone who looks only at the single number that comes out of the division will not find that
pattern. He needs two bounds: at most how many tiles are used, and at least how many. With those
two numbers he stops thinking about patterns that are plainly too wasteful, and concentrates on
the ones that may still win.

In this chapter we will learn why the bound from the continuous model is always valid, why
rounding can give an answer that is feasible but not the best, and how a computer uses that
bound to prune millions of possibilities without inspecting them one by one.
\end{untukanak}

\section{Branch and bound}
\tingkat{advanced}

If the relaxation answer is fractional, we do not round it, but split the problem into two and
ask again with the relaxation on each piece. Every integer answer satisfies one of two mutually
exclusive possibilities, $x_{2}\le1$ or $x_{2}\ge2$, so solving both is the same as solving the
original problem.

The procedure has four steps. Solve the relaxation at one node. If the result is integer, record
it as a candidate answer and do not continue with that node. If it is fractional, choose one
fractional variable, $x_{j}=v$, then create two children: $x_{j}\le\lfloor v\rfloor$ and
$x_{j}\ge\lceil v\rceil$. Before inspecting a node, compare its relaxation bound with the best
candidate; if it is not better, discard that node together with all its descendants. Stop when
no node remains.

That discarding rule is valid, and the reason is one line.

\begin{proposisi}[Pruning discards no answer]
Let $z_{\mathrm{LP}}(N)$ be the relaxation value at node $N$ and $z_{\mathrm{best}}$ the best
candidate already recorded. If $z_{\mathrm{LP}}(N)\le z_{\mathrm{best}}$, then no feasible
integer point at $N$ is better than $z_{\mathrm{best}}$.
\end{proposisi}

\begin{proof}
Every integer point that is feasible at $N$ is also feasible for the relaxation of $N$, so its
value is at most $z_{\mathrm{LP}}(N)$ by the relaxation-bound proposition. If
$z_{\mathrm{LP}}(N)\le z_{\mathrm{best}}$, all those points have value at most
$z_{\mathrm{best}}$, so there is nothing left to inspect.
\end{proof}

The tree for the shop above is small, and can be written entirely as a table.
\begin{center}
\begin{tabular}{llcc}
\hline
node & added constraint & relaxation $\mathbf{x}$ & $z_{\mathrm{LP}}$ \\
\hline
root & -- & $(3;\,1.5)$ & $21$ \\
left & $x_{2}\le1$ & $(10/3;\,1)$ & $20.\overline{6}$ \\
right & $x_{2}\ge2$ & $(2;\,2)$ & $18$ \\
left-left & $x_{2}\le1,\ x_{1}\le3$ & $(3;\,1)$ & $19$ \\
left-right & $x_{2}\le1,\ x_{1}\ge4$ & $(4;\,0)$ & $20$ \\
\hline
\end{tabular}
\end{center}
The root gives $21$ with $x_{2}$ fractional. Branching on $x_{2}$ gives two children: left
$20.\overline{6}$ with $x_{1}=10/3$ still fractional, right $18$ and already integer. The right
node is recorded as a candidate. At the left node, branching on $x_{1}$ gives $(3;1)$ with value
$19$ and $(4;0)$ with value $20$; the second updates the candidate. No node remains, and the
answer is $20$.

\begin{figure}[!ht]
\centering
\begin{tikzpicture}[
  simpul/.style={draw=layergrey,fill=layergrey!6,rounded corners=1.5pt,
    inner sep=3pt,font=\scriptsize,align=center,text width=3.0cm},
  pangkas/.style={draw=accent,fill=accent!7,rounded corners=1.5pt,
    inner sep=3pt,font=\scriptsize,align=center,text width=3.0cm},
  panah/.style={-{Stealth},thick,draw=layergrey}]
\node[simpul] (akar) at (0,0) {root\\ $z_{\mathrm{LP}}=21$\\ $(3;\,1.5)$};
\node[simpul] (kiri) at (-3.6,-1.75) {$x_{2}\le1$\\ $z_{\mathrm{LP}}=20.\overline{6}$\\
  $(10/3;\,1)$};
\node[pangkas] (kanan) at (3.6,-1.75) {$x_{2}\ge2$\\ $z_{\mathrm{LP}}=18$\\ $(2;\,2)$ integer};
\node[pangkas] (kk) at (-5.1,-3.6) {$x_{1}\le3$\\ $z_{\mathrm{LP}}=19$\\ $(3;\,1)$ integer};
\node[simpul] (kn) at (-2.1,-3.6) {$x_{1}\ge4$\\ $z_{\mathrm{LP}}=20$\\ $(4;\,0)$ integer};
\draw[panah] (akar.south) -- (kiri.north);
\draw[panah] (akar.south) -- (kanan.north);
\draw[panah] (kiri.south) -- (kk.north);
\draw[panah] (kiri.south) -- (kn.north);
\node[font=\scriptsize,accent,anchor=north] at (3.6,-2.55) {pruned: $18<20$};
\node[font=\scriptsize,accent,anchor=north] at (-5.1,-4.45) {pruned: $19<20$};
\end{tikzpicture}
\caption{The branch-and-bound tree for the shop: five nodes. The root gives $z=21$
with $x_{2}$ fractional. The branch $x_{2}\ge2$ ends at $(2;\,2)$ with value $18$; the branch
$x_{2}\le1$ branches again into $(3;\,1)$ with value $19$ and $(4;\,0)$ with value $20$.
The two nodes labeled integer whose values are below $20$ have a red frame: their bounds are no
longer better than the best candidate, so both are pruned without further inspection.}
\label{fig:ch22-pohon-percabangan}
\end{figure}

Figure~\ref{fig:ch22-pohon-percabangan} draws the same tree. Two nodes there
have a red frame, namely $(2;\,2)$ with value $18$ and $(3;\,1)$ with value $19$: both are
integer, but their bounds are below the best candidate $20$, so both are discarded.

The right node with value $18$ would not need further inspection if it were
fractional: the bound $18$ is already below the candidate $20$, so its whole branch can be
discarded. On an example this small the saving is not noticeable; on a problem with hundreds of
variables, that rule is what separates finishing in a minute from not finishing at all, since
without good bounds branch and bound degenerates into $2^{n}$ nodes.

Two things must be said honestly. The worst-case time is still exponential, and there is no
guarantee that the tree is small. The choice of branching variable and the order of the nodes
also strongly affect the size of the tree; software uses heuristics, for example the variable
closest to one half. Both explain why an integer program can be far slower than a linear program
of the same size.

\section{Worked example}
\tingkat{intermediate}

Now all the numbers are carried to the end, so that every claim above can be
rechecked by hand or with Python.

\begin{resep}
Four steps solve a two-variable linear program graphically:
\begin{enumerate}
  \item Turn each constraint into a line by replacing the inequality sign with an equals sign,
        then draw the line through two axis intercepts.
  \item Determine the feasible side by testing one point, usually the origin; shade the wrong
        side so that the feasible region stays clean.
  \item Intersect every pair of lines, discard the intersection points that do not satisfy all
        the constraints, and what remains are the vertices.
  \item Compute the objective value at every vertex, then take the largest for maximization or
        the smallest for minimization.
\end{enumerate}
A warning about order: step 3 must not be carried out before step 2. Intersection points
computed earlier often lie outside the feasible region and their values are in fact higher than
the true answer; a nice number that appears before feasibility is checked is hard to throw away.
\end{resep}

\emph{First, the linear program of the shop.} Its feasible region is a quadrilateral with four
vertices, and all four are inspected:
\begin{center}
\begin{tabular}{lcc}
\hline
vertex & active constraints & profit $5x_{1}+4x_{2}$ \\
\hline
$(0;\,0)$ & $x_{1}=0,\ x_{2}=0$ & $0$ \\
$(4;\,0)$ & $6x_{1}+4x_{2}=24,\ x_{2}=0$ & $20$ \\
$(3;\,1.5)$ & $6x_{1}+4x_{2}=24,\ x_{1}+2x_{2}=6$ & $21$ \\
$(0;\,3)$ & $x_{1}+2x_{2}=6,\ x_{1}=0$ & $12$ \\
\hline
\end{tabular}
\end{center}
The point $(3;1.5)$ is computed from the two equations: $x_{1}=6-2x_{2}$ substituted into
$6x_{1}+4x_{2}=24$ gives $36-8x_{2}=24$, so $x_{2}=3/2$, $x_{1}=3$, and the profit is
$15+6=21$.

\emph{Second, the dual and strong duality.} The dual minimizes $24y_{1}+6y_{2}$ subject to
$6y_{1}+y_{2}\ge5$ and $4y_{1}+2y_{2}\ge4$. Since $x_{1},x_{2}>0$, complementary slackness
requires both dual constraints to be active; solving both as equalities gives $y_{1}=3/4$ and
$y_{2}=1/2$. The value is $18+3=21$, the same as the profit. A cross-check:
$6\cdot\tfrac34+\tfrac12=5$ and $4\cdot\tfrac34+2\cdot\tfrac12=4$, exactly the right-hand
sides, so the gap is zero.

\emph{Third, shadow prices and their range of validity.} Perturbing the right-hand side gives
the following table, with $z$ recomputed from scratch in every row.
\begin{center}
\begin{tabular}{lccc}
\hline
perturbation & optimal $\mathbf{x}$ & $z$ & change \\
\hline
$b=(25;\,6)$ & $(3.25;\,1.375)$ & $21.75$ & $+0.75$ \\
$b=(24;\,7)$ & $(2.5;\,2.25)$ & $21.50$ & $+0.50$ \\
$b=(12;\,6)$ & $(0;\,3)$ & $12.00$ & $-9.00$ \\
$b=(24;\,4)$ & $(4;\,0)$ & $20.00$ & $-1.00$ \\
\hline
\end{tabular}
\end{center}
The third row reduces the plywood to $12$, with change
$-9=0.75\cdot(-12)$; the fourth row reduces the hours by $2$, with change
$-1=0.5\cdot(-2)$. Both are still inside the intervals $b_{1}\in[12,36]$ and
$b_{2}\in[4,12]$, so the linearity holds all the way to the end. Outside them the slope changes:
with $b_{1}=40$, the relaxation gives $(0;\,3)$ with value $12$, not
$21+0.75\cdot16=33$.

\emph{Fourth, the integer program of the shop.} The enumeration of the two directions above gave
the answer $20$ at $(4;0)$, whereas its relaxation is $21$ at $(3;1.5)$. The gap is $1$, namely
$1/21=4.76\%$ of the relaxation value. Rounding down gives $19$, so it loses
$1$ unit of profit, or $5\%$ of the true answer, and that happens on a problem with only two
variables.

\emph{Fifth, a knapsack with four items.} Four items with weights $5,4,3,2$ and
values $10,7,5,3$ are loaded into a container of capacity $8$. Their value densities are $2$,
$1.75$, $1.667$, $1.5$, so the greedy rule takes the first item and then $\tfrac34$ of the second:
weight $5+3=8$, value $15.25$. The best integer answer is the first and the third item, weight
$8$ exactly full, value $15$. Its integrality gap is $0.25$, only $1.64\%$ of the bound. But
rounding down throws the second item away entirely and leaves the value $10$, namely
$33.3\%$ below the true answer: a small gap does not mean that rounding is safe, since the two
measure different things.

\kertas{For the shop above, suppose the assembly machine is available for only $5$ hours, not
$6$. Recompute the vertices, determine the optimal continuous profit, then determine the optimal
integer profit with a small enumeration. Before computing, write down your guess: does the
integrality gap grow or shrink? After that, inspect the shadow price of assembly hours at
$b_{2}=5$ and compare it with $0.5$. As an extra, find the $b_{2}$ at which the best integer
answer moves away from $(4;0)$.}

\section{Computational lab}
\tingkat{advanced}

The two blocks below are run with Python 3.14.4, numpy 2.5.2, and scipy 1.18.0. The first block
solves the linear program, its dual, and the perturbations of the right-hand side. The second
block solves the branch-and-bound tree, the knapsack, and the vertex cover of the complete graph
on five vertices. The output is copied as it stands.

\begin{lab}
\begin{verbatim}
import numpy as np
from scipy.optimize import linprog

np.set_printoptions(precision=6, suppress=True, linewidth=78)

c = np.array([5., 4.])              # profit per shelf, per bench
A = np.array([[6., 4.], [1., 2.]])  # plywood, assembly hours
b = np.array([24., 6.])

r = linprog(-c, A_ub=A, b_ub=b, bounds=[(0, None)]*2, method="highs")
print("LP   x* =", r.x, "  z* =", -r.fun)
print("residual b - Ax =", b - A @ r.x)
print("marginals (scipy):", r.ineqlin.marginals)

rd = linprog(b, A_ub=-A.T, b_ub=-c, bounds=[(0, None)]*2, method="highs")
print("dual y* =", rd.x, "  dual value =", rd.fun)
print("duality gap =", rd.fun + r.fun)

print("B^-1 =", np.linalg.inv(A), "  x_B =", np.linalg.inv(A) @ b)

for nama, db in [("b=(25, 6)", (1., 0.)), ("b=(24, 7)", (0., 1.)),
                 ("b=(12, 6)", (-12., 0.)), ("b=(24, 4)", (0., -2.))]:
    rr = linprog(-c, A_ub=A, b_ub=b + np.array(db), bounds=[(0, None)]*2,
                 method="highs")
    print(f"{nama:<10} x = {np.round(rr.x, 6) + 0.0}  z = {-rr.fun:7.4f}"
          f"  delta z = {-rr.fun + r.fun:+.4f}")
\end{verbatim}
The actual output:
\begin{verbatim}
LP   x* = [3.  1.5]   z* = 21.0
residual b - Ax = [0. 0.]
marginals (scipy): [-0.75 -0.5 ]
dual y* = [0.75 0.5 ]   dual value = 21.0
duality gap = 0.0
B^-1 = [[ 0.25  -0.5  ]
 [-0.125  0.75 ]]   x_B = [3.  1.5]
b=(25, 6)  x = [3.25  1.375]  z = 21.7500  delta z = +0.7500
b=(24, 7)  x = [2.5  2.25]  z = 21.5000  delta z = +0.5000
b=(12, 6)  x = [0. 3.]  z = 12.0000  delta z = -9.0000
b=(24, 4)  x = [4. 0.]  z = 20.0000  delta z = -1.0000
\end{verbatim}
\end{lab}

\begin{lab}
\begin{verbatim}
import itertools
import numpy as np
from scipy.optimize import linprog, milp, LinearConstraint, Bounds

c = np.array([5., 4.])
A = np.array([[6., 4.], [1., 2.]])
b = np.array([24., 6.])


def lp_bengkel(tambahan_A=None, tambahan_b=None):
    Aub = A if tambahan_A is None else np.vstack([A, tambahan_A])
    bub = b if tambahan_b is None else np.concatenate([b, tambahan_b])
    r = linprog(-c, A_ub=Aub, b_ub=bub, bounds=[(0, None)]*2, method="highs")
    if r.status != 0:
        return (None, "infeasible")
    return (-r.fun, np.round(r.x, 6) + 0.0)


print("-- branch and bound tree --")
simpul = [("root",        None, None),
          ("x2<=1",       [[0., 1.]],  [1.]),
          ("x2>=2",       [[0., -1.]], [-2.]),
          ("x2<=1,x1<=3", [[0., 1.], [1., 0.]],  [1., 3.]),
          ("x2<=1,x1>=4", [[0., 1.], [-1., 0.]], [1., -4.])]
for nama, ta, tb in simpul:
    z, x = lp_bengkel(None if ta is None else np.array(ta),
                      None if tb is None else np.array(tb))
    print(f"{nama:<14} z_LP = {z:8.4f}  x = {x}")

rb = milp(-c, constraints=LinearConstraint(A, -np.inf, b),
          integrality=[1, 1], bounds=Bounds(0, np.inf))
print("IP optimum   z_IP =", -rb.fun, " x =", np.round(rb.x, 6) + 0.0,
      " gap =", 21.0 + rb.fun)

print("-- knapsack, 4 items, capacity 8 --")
w = np.array([5., 4., 3., 2.])
v = np.array([10., 7., 5., 3.])
for mask in range(16):
    S = [i + 1 for i in range(4) if mask >> i & 1]
    bw, bv = w[[i - 1 for i in S]].sum(), v[[i - 1 for i in S]].sum()
    print(f"{str(S):<12} weight {bw:4.0f}  value {bv:4.0f}"
          f"  {'feasible' if bw <= 8 else 'infeasible'}")
rk = milp(-v, constraints=LinearConstraint(w, -np.inf, 8.),
          integrality=[1]*4, bounds=Bounds(0, 1))
print("enumeration/MILP z_IP =", -rk.fun, " x =", np.round(rk.x, 6) + 0.0)
rl = linprog(-v, A_ub=[w], b_ub=[8.], bounds=[(0, 1)]*4, method="highs")
print("relaxation      z_LP =", -rl.fun, " x =", np.round(rl.x, 6) + 0.0)
print("integrality gap =", -rl.fun + rk.fun)

print("-- vertex cover of K5: LP vs IP --")
n = 5
E = list(itertools.combinations(range(n), 2))
M = np.zeros((len(E), n))
for k, (i, j) in enumerate(E):
    M[k, i] = M[k, j] = 1.
rv = linprog(np.ones(n), A_ub=-M, b_ub=-np.ones(len(E)),
             bounds=[(0, 1)]*n, method="highs")
rc = milp(np.ones(n), constraints=LinearConstraint(M, 1, np.inf),
          integrality=[1]*n, bounds=Bounds(0, 1))
print("LP =", rv.fun, " x =", np.round(rv.x, 6) + 0.0)
print("IP =", rc.fun, " x =", np.round(rc.x, 6) + 0.0)
\end{verbatim}
The actual output:
\begin{verbatim}
-- branch and bound tree --
root           z_LP =  21.0000  x = [3.  1.5]
x2<=1          z_LP =  20.6667  x = [3.333333 1.      ]
x2>=2          z_LP =  18.0000  x = [2. 2.]
x2<=1,x1<=3    z_LP =  19.0000  x = [3. 1.]
x2<=1,x1>=4    z_LP =  20.0000  x = [4. 0.]
IP optimum   z_IP = 20.0  x = [4. 0.]  gap = 1.0
-- knapsack, 4 items, capacity 8 --
[]           weight    0  value    0  feasible
[1]          weight    5  value   10  feasible
[2]          weight    4  value    7  feasible
[1, 2]       weight    9  value   17  infeasible
[3]          weight    3  value    5  feasible
[1, 3]       weight    8  value   15  feasible
[2, 3]       weight    7  value   12  feasible
[1, 2, 3]    weight   12  value   22  infeasible
[4]          weight    2  value    3  feasible
[1, 4]       weight    7  value   13  feasible
[2, 4]       weight    6  value   10  feasible
[1, 2, 4]    weight   11  value   20  infeasible
[3, 4]       weight    5  value    8  feasible
[1, 3, 4]    weight   10  value   18  infeasible
[2, 3, 4]    weight    9  value   15  infeasible
[1, 2, 3, 4] weight   14  value   25  infeasible
enumeration/MILP z_IP = 15.0  x = [1. 0. 1. 0.]
relaxation      z_LP = 15.25  x = [1.   0.75 0.   0.  ]
integrality gap = 0.25
-- vertex cover of K5: LP vs IP --
LP = 2.5  x = [0.5 0.5 0.5 0.5 0.5]
IP = 4.0  x = [1. 1. 1. 1. 0.]
\end{verbatim}
\end{lab}

Three things must be read from the first output. The reported $\mathbf{x}^{*}$,
$(3;\,1.5)$, with $z^{*}$ exactly $21$ and no residual in either constraint, matches
complementary slackness. \texttt{scipy} reports the multipliers $(-0.75,\,-0.5)$, the negative
of the shadow prices, because the routine minimizes $-\mathbf{c}^{\top}\mathbf{x}$; the true
numbers are $0.75$ and $0.5$, and the way to be sure is to recompute through $B^{-1}$. The
duality gap comes out exactly $0$, not $10^{-15}$, because $21$ is represented exactly in
binary; on decimal data a residual of order $10^{-12}$ is normal.

The second output matches the hand computation: five nodes with the same values to four
decimals, the integer answer $20$ at $(4;\,0)$, the knapsack relaxation $15.25$, and its integer
answer $15$ with gap $0.25$. The enumeration of all sixteen subsets confirms that $15$ is the
best: every subset of higher value exceeds the capacity.

The last part of the second output checks the factor-two claim above. For the complete graph on
five vertices, the relaxation is $2.5$ whereas the integer answer is $4$; the ratio is $1.6$,
and for $n$ vertices the ratio $2(n-1)/n$ approaches $2$ from below. The relaxation bound here
can be off by nearly twice.

\begin{untukdosen}
The order that works: start from the shop and the half bench, then standard form, then vertices,
and only then the simplex. Do not start from the tableau. Students who see the tableau first
memorize the pivot step without knowing why that step may stop; those who have already seen that
the answer is at a corner will treat a pivot as a way of moving from corner to corner, and the
optimality test as a way of knowing that no corner is better.

Two typical misconceptions. First, ``the vertices are all the intersections of the constraint
lines''. Many intersections lie outside the feasible region, for example $(4;3)$ in our example,
which cuts both lines but violates both; checking the feasibility of every candidate must be
done on the board. Second, ``the integer answer is the rounding of the fractional answer''. The
shop gives $19$ against $20$, the knapsack gives $10$ against $15$; both are enough to break
that guess.

Class activity: write the initial tableau on the board, ask the class to choose the entering
column and compute the ratio test, then predict the profit before the pivot. After the final
tableau, ask for the profit when the plywood becomes $25$ sheets; $21.75$ opens the way to the
interval $[12,36]$.

Sign of readiness to go on: students can explain why the simplex method does not need to inspect
every vertex, and can name one reason why the relaxation must not be rounded.
\end{untukdosen}

\section{Summary}

\begin{ringkasan}
\begin{itemize}
  \item A linear program demands an objective and constraints of straight-line shape. The
        standard form $\max\mathbf{c}^{\top}\mathbf{x}$ with $A\mathbf{x}=\mathbf{b}$ and
        $\mathbf{x}\ge\mathbf{0}$ covers every case, with slack variables to turn inequalities
        into equalities.
  \item If the feasible region is bounded and the objective is bounded above, the best value is
        attained at a vertex. In standard form, a vertex is the same as a basic feasible
        solution, and their number is $\binom{n}{m}$, so inspecting all of them is not a
        sensible plan.
  \item The simplex method moves from vertex to vertex through columns with negative reduced
        cost and a ratio test, and stops when there is no negative reduced cost; at that moment
        the $z$ row already contains the dual prices.
  \item Weak duality holds without conditions; strong duality states that the two values are
        equal when both are finite; complementary slackness states that a loose constraint has
        price zero.
  \item The shadow price $y_{i}$ is the rate of change of the optimal value with respect to the
        right-hand side, and it holds only as long as the set of active constraints does not
        change; in the shop $0.75$ per sheet and $0.5$ per hour, valid for $b_{1}\in[12,36]$
        and $b_{2}\in[4,12]$.
  \item The LP relaxation gives a valid bound, not an answer. Rounding can give an infeasible
        point, such as $(3;2)$, or a feasible point that is not the best, such as $(3;1)$ with
        value $19$ against $20$.
  \item Branch and bound uses the relaxation bound to discard branches. The shop tree has five
        nodes, its integer answer is $20$ with gap $1$. The worst-case time is still
        exponential, and the decision version is NP-complete.
\end{itemize}
\end{ringkasan}

\section{Exercises}

\begin{latihan}
\begin{enumerate}
  \item A bakery uses $10$ kg of flour and $8$ oven hours. Cake A uses $2$ kg and
        $1$ hour with profit $30$; cake B uses $1$ kg and $2$ hours with profit $20$.
        Write the linear program, draw its feasible region, and compute the profit at every
        vertex. Then determine the best integer answer by enumeration, and compare it
        with the result of rounding down the continuous answer.

  \item Prove that the set $\{\mathbf{x}:A\mathbf{x}\le\mathbf{b}\}$ is convex, then
        explain why that proof does not apply to the set
        $\{\mathbf{x}\in\mathbb{Z}^{n}:A\mathbf{x}\le\mathbf{b}\}$.

  \item Change the profit of the shop to $x_{1}+2x_{2}$. Compute its value at the four vertices,
        state which vertex is optimal, and show that every point of the segment between the two
        optimal vertices is also optimal. Explain why a point in the middle of that segment is
        not a vertex, then connect the result with the shadow price of the hours constraint.

  \item Solve with the simplex tableau: maximize $3x_{1}+2x_{2}$ subject to
        $x_{1}+x_{2}\le4$, $x_{1}\le3$, $x_{2}\le3$, $x\ge0$. Write every intermediate tableau,
        name the entering column and the reason for it, then read the dual prices from the final
        tableau. (Hint: two pivots suffice; compare the final result with the vertex obtained
        from the intersection of $x_{1}+x_{2}=4$ and $x_{1}=3$.)

  \item Write the dual of the program in exercise 1, solve it, and check strong duality
        numerically. Explain the meaning of every shadow price in one sentence.

  \item Assess this claim: \emph{rounding down the answer of the relaxation always gives a
        feasible integer point.} Decide true, false, or not enough evidence, and
        exhibit one counterexample with numbers.

  \item Assess this claim: \emph{a positive shadow price means that the constraint must be
        binding.} Decide true, false, or not enough evidence, then explain the actual relation
        between shadow prices and complementary slackness.
        (Hint: inspect what happens to a constraint with $y_{i}>0$ according to
        $y_{i}(b_{i}-\mathbf{a}_{i}^{\top}\mathbf{x})=0$.)

  \item For the knapsack with weights $(5,4,3,2)$, values $(10,7,5,3)$, and capacity
        $8$, compute its LP relaxation by hand using the density order, then compare it with the
        enumeration result. Compute the integrality gap and state it as a percentage of the
        bound. (Hint: sort by value per unit of weight, then fill the container to the top; the
        last item enters partly.)

  \item Repeat exercise 8 with capacity $7$. Determine the relaxation, the integer answer, and
        the gap, then explain why the gap can change although only one data value changed.

  \item For the vertex cover of the complete graph on $n$ vertices, prove that the relaxation has
        value $n/2$ and the integer answer $n-1$, then compute the ratio of the two for
        $n=4,5,10$. Explain why that result shows that rounding is not a general way to solve
        integer programs.
\end{enumerate}
\end{latihan}

\begin{untukpeneliti}
What is really proved in this chapter: weak duality for a primal-dual pair in standard form,
without any condition; the converse direction of complementary slackness as an optimality test;
optimality at a vertex for a bounded feasible region, with a proof through the active-constraint
argument; and the validity of pruning in branch and bound. What is borrowed: strong duality,
whose usual proof goes through Farkas' lemma or through the termination of the simplex; both are
recorded in \texttt{WSPrinted.pdf}, in the section on linear programming duality, and in Boyd and
Vandenberghe, section 5.2.

What is popularized beyond its proof is quite a lot, and three things must be named.
``The LP relaxation is nearly right'': its worst-case approximation factor can grow without
bound on certain classes of problems, and even for the vertex cover it reaches $2$. ``The shadow
price is the derivative of the optimal value'': that is true only when the optimal value is
differentiable with respect to the right-hand side; at a degenerate point, where two optimal
bases share the same vertex, the derivative does not exist and what remains is only a
subgradient. ``Branch and bound is always fast'': its worst-case time is exponential, and
slowness on a particular instance is not a sign of a bad implementation.

Open questions: how strong the Chvatal-Gomory cutting-plane hierarchy is for particular classes
of polytopes, and whether a superpolynomial lower bound exists for the size of extended
formulations of the traveling salesman polytope; the results of Fiorini and coauthors in 2012
show an exponential bound for that case, while the search for similar bounds for other problems
continues. Limits of validity of this chapter: rational data, exact arithmetic for proofs,
floating-point tolerance in practice; all statements about shadow prices hold for the continuous
model, not for the integer model.

\peta{This chapter follows \emph{Introduction to Linear Optimization} in the file
\texttt{WSPrinted.pdf}: the sections on formulation, the geometry of polyhedral sets (vertices
and extreme points), systems of linear inequalities and duality (complementary slackness,
sensitivity analysis), and the simplex method (basic feasible solutions, reduced costs, the
optimality test). The integer-programming side follows Wolsey, \emph{Integer Programming},
chapter 1 (formulation) and chapter 2 (optimality, relaxation, bounds), which places the
relaxation as the source of bounds and branching as the way to use them. The convex viewpoint in
the second and fourth parts is Boyd and Vandenberghe, \emph{Convex Optimization}, chapters 4 and
5. Relation to other chapters: standard form and the feasible region continue Chapter 21,
solving linear systems uses Chapter 8, and the integrality gap is used again in Chapter 25.}

\end{untukpeneliti}


% ---------------------------- part VI ----------------------------
\part{Machine Learning, Security, and Privacy}

This part closes the book with the question most often avoided: when is a machine that looks
clever actually fragile, and how is its fragility proved with numbers. The prerequisites are
Chapters~10, 12, 18, and~23. After eight chapters you can train and evaluate a model honestly,
attack it, prove limits on attacks, and explain cryptography and differential privacy well
enough to read the original papers.
\chapter{Neural Networks and Backpropagation}

An intrusion prevention system observes four quantities on every connection: the number of
requests per second, the connection duration, the number of rejected packets, and the average
payload size. An analyst trying to separate malicious traffic from ordinary traffic finds a
pattern that is awkward to read. Connections with many requests but no rejected packets are
usually legitimate, because that is what fetching a large file looks like. Connections with many
rejected packets but few requests are also usually legitimate, because that is what a failed scan
looks like. The combination of the two is almost always an attack. Each feature by itself
separates nothing, whereas their combination separates sharply.

A linear threshold cannot draw such a separator, and the first part of this chapter shows that
the obstacle is not the amount of data. What must change is the shape of the function. One hidden
layer is already enough for that pattern. Once the shape has changed, the next question appears,
and it is more serious: how do we tune the thirteen parameters of that small network without
deriving a new formula for every architecture?

The answer is the chain rule of Chapter 12, run backward on the computational graph. This chapter
carries that rule down to the numbers: one forward step with given weights, all thirteen
gradients computed by hand and then checked with NumPy, one gradient-descent step, and the
gradient with respect to the input that later becomes a weapon in Chapter 29. What this chapter
promises is not a clever network, but a differentiation engine that can be checked line by line.

\begin{untukanak}
Imagine an assembly line with several workers. The first worker works on the raw material, the
second works on the output of the first, and so on until the last worker hands the finished item
to an inspector. The inspector weighs the item, then says: too heavy, two hundred grams over.

Now the question: who has to change the way they work? The first worker? The last worker? If the
foreman can change only one setting and then rerun the whole line, he needs two tries per setting.
With thirteen settings, that is twenty-six runs of the line for a single fix.

It turns out there is a far cleverer way. After one run, we can walk back from the inspector to
the last worker, then to the previous worker, all the way to the first. At each stop we compute
one number: how much the final error changes if this person changes their output a little. With a
single backward pass, every worker knows their own share, and nothing has to be run again.

In this chapter we will learn why information can flow backward even though the goods only flow
forward, and why that backward pass gives an answer exactly as accurate as running the line
twenty-six times.
\end{untukanak}

\section{Understanding bridge: from story to symbol}
\tingkat{basic}

The formulas in this chapter are easier to read if the order is reversed: a story without
symbols, then arithmetic on paper, then symbols as shorthand for steps we have already
performed ourselves.

\subsection{Plain analogy: how much, not who}

The assembly-line analogy answered ``who should change''. It did not answer ``how much''.
Picture three taps in a row above one bucket, with a target line on the wall. The third tap
feels closest, but it can only pass on what the second tap sends; the first tap matters
most, yet its effect depends on the two taps between it and the bucket. What we need is not
turning every tap and re-measuring, but one number per tap: how many millimetres does the
water rise if this tap moves a little? One pass, traced backwards from the bucket.

\subsection{Compute by hand first}

\begin{contoh}
One neuron, four numbers: weight $w = 0.5$, bias $b = 0.1$, input $x = 2.0$, correct
answer $y = 1$, learning rate $\eta = 0.1$. Follow the steps in order; every number here
can be checked on an ordinary calculator.

\begin{enumerate}
  \item Raw score: $w x = 0.5 \times 2.0 = 1.0$, plus the bias gives $z = 1.1$.
  \item Probability: $e^{-1.1} = 0.33287$, so $\sigma = 1/(1+0.33287) = 0.7502601$.
  \item Loss: because $y = 1$ only one term survives, $L = -\ln 0.7502601 = 0.2873353$.
  \item Gap between probability and truth: $\sigma - y = 0.7502601 - 1 = -0.2497399$.
  \item Chain rule to the weight: multiply that gap by the input, $-0.2497399 \times 2.0
        = -0.4994798$. The bias is not multiplied by anything, so it keeps $-0.2497399$.
  \item Move: $w' = 0.5 - 0.1(-0.4994798) = 0.5499480$ and $b' = 0.1 - 0.1(-0.2497399)
        = 0.1249740$. The new score is $z' = 1.2248699$, the new probability $0.7729194$,
        the new loss $0.2575805$: down $0.0297549$ from $0.2873353$.
\end{enumerate}

Mind the sign in step 5: the gradient is negative, and because we \emph{subtract} it, the
weight rises. The minus sign in the update rule is not decoration; it reverses the
direction.
\end{contoh}

One check you must not skip: replace the analytic derivative by a finite difference, that
is, evaluate the loss at $w + 0.000001$ and $w - 0.000001$ and divide the difference by
$0.000002$. The answer is $-0.4994798$, equal to seven digits. When the hand result and the
numeric result disagree, the sign is almost always what is wrong.

\subsection{Symbol dictionary}

The symbols below appear in any book on neural networks. Nothing here needs memorising:
reading each one aloud takes a single sentence, and that sentence is what we computed above.

\begin{description}
  \item[$\partial L/\partial w$] read ``del L del w'': if weight $w$ moves a little, how
        much does the loss $L$ change? In our example the answer is $-0.4994798$ (step 5).
  \item[$z = Wx + b$] read ``z equals W times x plus b'': the raw score before it is
        squeezed into a probability. The one-number version of this formula is step 1.
  \item[$\sigma(z)$] read ``sigma of z'': the S-shaped function that keeps the score
        between zero and one. In step 2 it gives $0.7502601$.
  \item[$\nabla_w L$] read ``nabla w L'': every derivative of the loss with respect to
        every weight at once; its arrow answers ``which way does the loss fall fastest''.
        Its entries are numbers like the one in step 5, only many.
\end{description}

\subsection{Why does this matter for security?}

The same derivative that repairs a model also breaks it. An attacker needs one number: if
one input feature moves a little, how much does the loss rise? In our example $x$ has
derivative $-0.1248699$, so shrinking $x$ raises the loss. Moving it by $1.9$ sends $x$ to
$0.1$, the probability drops to $0.5374298$, and the loss rises from $0.2873353$ to
$0.6209570$, a factor of $2.16$. One feature, one direction, one step.

For the defender the same question runs backwards: which feature has the largest
derivative? That feature is the cheapest to fake, so it most needs rounding, clipping, or
sanity checks. Chapter 29 uses this recipe in full, with perturbation size $\epsilon$ and
the formula $x_{\mathrm{adv}} = x + \epsilon \cdot \operatorname{sign}(\nabla_x L)$.

\begin{catatan}
A warning that is easily missed: on real systems an attacker rarely computes the gradient
exactly; it is estimated by changing the input many times. That estimate is more expensive
and noisier, but it still works. A closed model is therefore not a defence.
\end{catatan}

\subsection{Hand computation drills}

Work with a calculator, write on paper, then compare. Answers sit in parentheses.

\begin{catatan}
\begin{enumerate}
  \item Repeat the whole example with $b = 0.3$, everything else unchanged. Compute $z$,
        $\sigma$, $L$, $\partial L/\partial w$, $\partial L/\partial b$.
        ($z = 1.3$; $\sigma = 0.7858350$; $L = 0.2410085$; $-0.4283300$; $-0.2141650$.)
  \item Repeat the original example with $\eta = 0.5$. What are $w'$, $b'$, and the new
        loss? ($w' = 0.7497399$; $b' = 0.2248699$; $L' = 0.1640633$, down $0.1232720$.)
  \item Repeat the original example with $y = 0$, everything else unchanged. What is
        $\partial L/\partial w$ now, and which way does the weight move?
        ($0.7502601$ and $1.5005202$; positive, so the weight falls.)
  \item In the original example, move $x$ by $\epsilon = 0.5$ in the direction that raises
        the loss. What are the new $x$ and the new loss? ($x = 1.5$; $L = 0.3558651$.)
\end{enumerate}
\end{catatan}

\subsection{From paper to code}

The code below runs the same six steps, with variable names that point at their step.
Nothing here is new: every line is one of the lines above.

\begin{lab}
\begin{verbatim}
import numpy as np

# One example, the same numbers as in "Compute by hand first".
x, y = 2.0, 1.0
w, b, eta = 0.5, 0.1, 0.1

# Step 1: z = w*x + b  (raw score)
z = w*x + b
print("Step 1  z        =", round(z, 7))

# Step 2: sigma(z) = 1/(1+e^-z)  (squeeze the score into 0..1)
s = 1/(1+np.exp(-z))
print("Step 2  sigma    =", round(s, 7))

# Step 3: L = -ln(s)  (loss; y=1 leaves a single term)
L = -np.log(s)
print("Step 3  L        =", round(L, 7))

# Step 4: dL/dz = sigma - y  (probability minus reality)
dLdz = s - y
print("Step 4  dL/dz    =", round(dLdz, 7))

# Step 5: chain rule down to weight and bias
dLdw, dLdb = dLdz*x, dLdz
print("Step 5  dL/dw    =", round(dLdw, 7), " dL/db =", round(dLdb, 7))

# Step 6: one gradient-descent step
w2, b2 = w - eta*dLdw, b - eta*dLdb
s2 = 1/(1+np.exp(-(w2*x + b2)))
print("Step 6  w', b'   =", round(w2, 7), round(b2, 7),
      " L' =", round(-np.log(s2), 7))

# Check: analytic gradient against a finite difference.
def loss(ww, bb):
    return -np.log(1/(1+np.exp(-(ww*x + bb))))
h = 1e-6
print("finite difference dL/dw =", round((loss(w+h, b)-loss(w-h, b))/(2*h), 7),
      " dL/db =", round((loss(w, b+h)-loss(w, b-h))/(2*h), 7))

# Gradient with respect to the INPUT: this is what attacks a model (Chapter 29).
dLdx = dLdz*w
xa = x + 1.9*np.sign(dLdx)          # the sign of the gradient sets the direction
sa = 1/(1+np.exp(-(w*xa + b)))
print("dL/dx =", round(dLdx, 7), " x_adv =", round(xa, 4),
      " sigma =", round(sa, 7), " L =", round(-np.log(sa), 7))
\end{verbatim}

The output must match the paper: $1.1$, $0.7502601$, $0.2873353$, $-0.2497399$,
$-0.4994798$, then $0.2575805$. If one number is off, do not patch the code: find the step
that differs on paper.
\end{lab}

\section{Linear thresholds and the perceptron}
\tingkat{basic}

The oldest form of classifier is the linear threshold. For an input
$\mathbf{x}\in\mathbb{R}^{d}$, weights $\mathbf{w}\in\mathbb{R}^{d}$, and bias
$b\in\mathbb{R}$, we compute a single score and decide the class from the sign of that score.

\kamus{perceptron}{perceptron}{A classifier that decides only from which side of a single
separating line or plane a point falls.}

\begin{definisi}[Linear threshold]
A function $f:\mathbb{R}^{d}\to\{-1,+1\}$ is called a linear threshold if there are
$\mathbf{w}$ and $b$ such that $f(\mathbf{x})=\operatorname{sign}(\mathbf{w}^{\top}\mathbf{x}+b)$.
The set $\{\mathbf{x}:\mathbf{w}^{\top}\mathbf{x}+b=0\}$ is the separating hyperplane, and that
hyperplane is the only form of separator this model can learn.\index{perceptron}
\end{definisi}

Perceptron training works by correcting the weights only when there is a mistake. For a label
$y\in\{-1,+1\}$ and an example $\mathbf{x}$, if $y\,\mathbf{w}^{\top}\mathbf{x}\le 0$ then we
update
\[
  \mathbf{w}\leftarrow\mathbf{w}+y\,\mathbf{x}.
\]
Note that there is no learning rate here: the size of the step is determined by the size of the
example, and training stops by itself as soon as all examples are correct. Does it always stop?

\begin{teorema}[Perceptron convergence]
Suppose every example satisfies $\lVert\mathbf{x}_{i}\rVert\le1$ and there is a unit vector
$\mathbf{w}^{*}$ and $\gamma>0$ such that $y_{i}\,\mathbf{w}^{*\top}\mathbf{x}_{i}\ge\gamma$ for
every $i$. Then the perceptron algorithm starting from $\mathbf{w}_{0}=\mathbf{0}$ makes at most
$\lfloor1/\gamma^{2}\rfloor$ updates.\index{perceptron}
\end{teorema}

\begin{proof}
Suppose an update occurs $k$ times, on the examples $\mathbf{x}_{1},\dots,\mathbf{x}_{k}$. Every
update occurs only when that example is misclassified, so
$y_{j}\,\mathbf{w}_{j-1}^{\top}\mathbf{x}_{j}\le0$. Hence
\[
  \lVert\mathbf{w}_{j}\rVert^{2}
  =\lVert\mathbf{w}_{j-1}\rVert^{2}+2y_{j}\mathbf{w}_{j-1}^{\top}\mathbf{x}_{j}
   +\lVert\mathbf{x}_{j}\rVert^{2}
  \le\lVert\mathbf{w}_{j-1}\rVert^{2}+1 .
\]
By induction, $\lVert\mathbf{w}_{k}\rVert^{2}\le k$.

On the other side, the inner product with $\mathbf{w}^{*}$ grows by at least $\gamma$ each
time:
\[
  \mathbf{w}^{*\top}\mathbf{w}_{j}
  =\mathbf{w}^{*\top}\mathbf{w}_{j-1}+y_{j}\mathbf{w}^{*\top}\mathbf{x}_{j}
  \ge\mathbf{w}^{*\top}\mathbf{w}_{j-1}+\gamma ,
\]
so $\mathbf{w}^{*\top}\mathbf{w}_{k}\ge k\gamma$. Since $\mathbf{w}^{*}$ is a unit vector, the
Cauchy-Schwarz inequality gives
\[
  k\gamma\le\mathbf{w}^{*\top}\mathbf{w}_{k}\le\lVert\mathbf{w}_{k}\rVert\le\sqrt{k}.
\]
After dividing by $\sqrt{k}$ and squaring, $k\gamma^{2}\le1$, that is $k\le1/\gamma^{2}$.
\end{proof}

That bound does not depend on the dimension or on the number of examples, and that is its
appeal. But when $\gamma$ does not exist, this theorem says nothing, and the perceptron really
does oscillate forever; it is also not a statement about generalization.

Now a more fundamental bound. Four points in the plane are already enough to show that not every
labeling can be separated by a single hyperplane. The points $(0,0)$ and $(1,1)$ have label $0$,
the points $(1,0)$ and $(0,1)$ have label $1$. This pattern is usually called XOR.

\begin{proposisi}
The four XOR points are not linearly separable.
\end{proposisi}

\begin{proof}
Suppose there are $w_{1},w_{2},b$ that separate them, with label $1$ on the positive side. Label
$0$ at $(0,0)$ gives $b\le0$. Label $1$ at $(1,0)$ gives $w_{1}+b>0$, hence $w_{1}>-b\ge0$. By
the same route, $w_{2}>-b\ge0$. Therefore $w_{1}+w_{2}>-2b$, and
\[
  w_{1}+w_{2}+b>-b\ge0 .
\]
But label $0$ at $(1,1)$ demands $w_{1}+w_{2}+b\le0$. Contradiction.
\end{proof}

This contradiction holds for every choice of $\mathbf{w}$ and $b$, not only for some particular
training outcome. Adding examples, changing the learning rate, or normalizing the features will
not help. What must change is the shape of the function, and the smallest change available is to
insert one layer in between.

\kamus{hidden layer}{lapisan tersembunyi}{A layer of units that the outside world does not see:
its input comes from the previous layer, its output goes to the next layer.}

A hidden layer works in two steps. First, compute
$\mathbf{z}_{1}=\mathbf{W}_{1}\mathbf{x}+\mathbf{b}_{1}$ with
$\mathbf{W}_{1}\in\mathbb{R}^{2\times4}$, then apply
$\operatorname{ReLU}(t)=\max\{t,0\}$ to every component, and write the result $\mathbf{h}$.
Second, compute a linear score from $\mathbf{h}$. For XOR, with
\[
  \mathbf{W}_{1}=\begin{pmatrix}1&1\\1&1\end{pmatrix},\qquad
  \mathbf{b}_{1}=\begin{pmatrix}0\\-1\end{pmatrix},\qquad
  \mathbf{w}_{2}=\begin{pmatrix}1\\-2\end{pmatrix},\qquad b_{2}=0,
\]
the four points give: $(0,0)$ produces $\mathbf{z}_{1}=(0,-1)$ and $\mathbf{h}=(0,0)$ with score
$0$; the points $(1,0)$ and $(0,1)$ both produce $\mathbf{z}_{1}=(1,0)$ and $\mathbf{h}=(1,0)$
with score $1$; the point $(1,1)$ produces $\mathbf{z}_{1}=(2,1)$ and $\mathbf{h}=(2,1)$ with
score $1(2)-2(1)=0$. With a threshold of $0.5$ all four points are classified correctly.

What makes this work is not magic but simple arithmetic. The first unit computes
$\max\{x_{1}+x_{2},0\}$ and the second computes $\max\{x_{1}+x_{2}-1,0\}$. Both are linear, but
their difference is a piecewise function that turns at two places. A linear combination of linear
functions is still linear; a linear combination of piecewise functions is not.

Such a boundary is made of line segments. Each unit in the first layer contributes one straight
line, and the second layer chooses which side is used, so the final boundary is a polygon;
Figure~\ref{fig:dua-lapis} shows its shape.

\begin{figure}[htbp]
\centering
\begin{tikzpicture}[>=latex, font=\small]
  % left panel: two hidden layers
  \foreach \i/\yy in {1/1.5,2/-0.5}
    \node[circle,draw=layergrey,fill=layergrey!10,inner sep=1.7pt] (I\i) at (0.4,\yy) {};
  \foreach \i/\yy in {1/2.4,2/1.65,3/0.9,4/0.15,5/-0.6}
    \node[circle,draw=layerteal,fill=layerteal!12,inner sep=1.7pt] (H1\i) at (2.1,\yy) {};
  \foreach \i/\yy in {1/1.6,2/0.7,3/-0.2}
    \node[circle,draw=layerteal,fill=layerteal!12,inner sep=1.7pt] (H2\i) at (3.7,\yy) {};
  \node[circle,draw=layerindigo,fill=layerindigo!12,inner sep=1.7pt] (O) at (5.1,0.7) {};
  \foreach \i in {1,2} \foreach \j in {1,...,5} \draw[layergrey!45] (I\i) -- (H1\j);
  \foreach \i in {1,...,5} \foreach \j in {1,2,3} \draw[layergrey!45] (H1\i) -- (H2\j);
  \foreach \j in {1,2,3} \draw[layergrey!45] (H2\j) -- (O);
  \node[font=\scriptsize] at (0.4,-1.35) {input};
  \node[font=\scriptsize] at (2.1,-1.35) {five units};
  \node[font=\scriptsize] at (3.7,-1.35) {three units};
  \node[font=\scriptsize] at (5.1,-1.35) {output};
  % panel kanan: bidang input
  \draw[layergrey] (6.7,-0.9) rectangle (14.4,2.9);
  \node[font=\scriptsize] at (14.1,-1.35) {$x_{1}$};
  \node[font=\scriptsize] at (6.4,2.75) {$x_{2}$};
  \fill[layerorange!18] (8.3,0.3) -- (9.7,2.4) -- (12.6,2.55) -- (13.3,1.1) -- (11.2,0.0) -- cycle;
  \draw[layerorange,thick] (8.3,0.3) -- (9.7,2.4) -- (12.6,2.55) -- (13.3,1.1) -- (11.2,0.0) -- cycle;
  \draw[dashed,layergrey!70] (7.50,-0.9) -- (11.43,2.9);
  \draw[dashed,layergrey!70] (6.70,2.23) -- (14.40,2.64);
  \draw[dashed,layergrey!70] (12.43,2.9) -- (14.27,-0.9);
  \draw[dashed,layergrey!70] (14.40,1.68) -- (9.48,-0.9);
  \draw[dashed,layergrey!70] (6.70,0.50) -- (14.40,-0.33);
  \foreach \xx/\yy in {9.9/1.6, 11.4/1.3, 10.6/0.6}
    \draw[accent,thick] (\xx-0.09,\yy-0.09) -- (\xx+0.09,\yy+0.09)
      (\xx-0.09,\yy+0.09) -- (\xx+0.09,\yy-0.09);
  \foreach \xx/\yy in {7.7/1.4, 8.5/-0.3, 12.3/-0.5, 13.6/2.2}
    \draw[layergrey,fill=white] (\xx,\yy) circle (2.1pt);
  \draw[accent,thick] (11.0,-1.30) -- (11.15,-1.15) (11.0,-1.15) -- (11.15,-1.30);
  \node[font=\scriptsize,anchor=west] at (11.2,-1.23) {class 1};
  \draw[layergrey,fill=white] (13.0,-1.23) circle (2.1pt);
  \node[font=\scriptsize,anchor=west] at (13.2,-1.23) {class 0};
\end{tikzpicture}
\caption{A sketch of how two hidden layers turn one line into a fold. Left: each unit in the
first layer pulls one straight line across the input plane, and the second layer combines the
sides of those lines. Right: the decision boundary is a five-sided polygon, drawn thick; the
dashed lines are the extension of each side, namely the line belonging to one unit. Crosses are
class 1 inside the polygon, circles class 0 outside. The figure is a sketch, not a training
result: the number of sides follows the number of units in the first layer.}
\label{fig:dua-lapis}
\end{figure}


\section{Activation functions and two diseases of the inner layers}
\tingkat{intermediate}

Without a nonlinear function between the layers, the whole network collapses into a single
matrix. If ReLU is replaced by the identity, the output is
$\mathbf{W}_{2}\mathbf{W}_{1}\mathbf{x}$, that is, a single matrix product
$\mathbf{W}\mathbf{x}$ with $\mathbf{W}=\mathbf{W}_{2}\mathbf{W}_{1}$. Every hidden layer
therefore needs a nonlinear activation, and the choice of activation determines which disease
will appear when the network is made deep.

\kamus{activation function}{fungsi aktivasi}{The gate that determines how much of a unit's signal
is passed on, and at the same time the one that fixes the shape of its derivative.}

\begin{definisi}[ReLU, sigmoid, and tanh]
For $t\in\mathbb{R}$,
\[
  \operatorname{ReLU}(t)=\max\{t,0\},\qquad
  \sigma(t)=\frac{1}{1+e^{-t}},\qquad
  \tanh(t)=\frac{e^{t}-e^{-t}}{e^{t}+e^{-t}} .
\]
The derivative of ReLU is $1$ for $t>0$ and $0$ for $t<0$; at $t=0$ the derivative does not
exist, and this book uses the value $0$ as a convention. The derivative of tanh satisfies
$\tanh'(t)=1-\tanh^{2}(t)$.\index{activation function}\index{relu}
\end{definisi}

The value $0$ at the break point of ReLU is a convention: the set where the derivative fails is
only $\{0\}$, of measure zero, so it does not affect the derivative almost everywhere. The
practical consequence: a unit whose input is always negative receives zero gradient and will
never change.

Sigmoid has a neater property, and that property is the source of the trouble.

\begin{lemma}[Derivative of sigmoid]
For every $t$ we have $\sigma'(t)=\sigma(t)\bigl(1-\sigma(t)\bigr)$, hence
$0<\sigma'(t)\le\tfrac14$.\index{saturation}
\end{lemma}

\begin{proof}
Write $\sigma(t)=(1+e^{-t})^{-1}$. Then
\[
  \sigma'(t)=-(1+e^{-t})^{-2}\cdot(-e^{-t})
  =\frac{e^{-t}}{(1+e^{-t})^{2}}
  =\frac{1}{1+e^{-t}}\cdot\frac{e^{-t}}{1+e^{-t}}
  =\sigma(t)\bigl(1-\sigma(t)\bigr).
\]
Since $0<\sigma(t)<1$, the product is positive. Writing $s=\sigma(t)$ gives
$s(1-s)=\tfrac14-(s-\tfrac12)^{2}\le\tfrac14$, and equality is attained when $s=\tfrac12$.
\end{proof}

The bound $\tfrac14$ is the one to hold on to. The following table computes $\sigma$,
$\sigma'$, and $\tanh$ at several values; the numbers were produced by NumPy and are written to
six decimals.

\begin{center}
\small
\begin{tabular}{rrrrr}
\hline
$t$ & $\sigma(t)$ & $\sigma'(t)$ & $\tanh(t)$ & $1-\tanh^{2}(t)$ \\
\hline
$0$  & $0.500000$ & $0.250000$ & $0.000000$ & $1.000000$ \\
$1$  & $0.731059$ & $0.196612$ & $0.761594$ & $0.419974$ \\
$2$  & $0.880797$ & $0.104994$ & $0.964028$ & $0.070650$ \\
$5$  & $0.993307$ & $0.006648$ & $0.999909$ & $0.000181$ \\
$10$ & $0.999955$ & $0.000045$ & $0.999999$ & $0.000002$ \\
\hline
\end{tabular}
\end{center}

The fourth and fifth columns show what is called saturation: once $t$ moves away from zero, the
output barely changes with the input, and the derivative tends to zero. At $t=10$ the derivative
is only $4.54\times10^{-5}$, about one twenty-thousandth of its maximum.

The consequence can be computed. If the error propagates backward through $L$ consecutive layers
whose derivative is at most $\tfrac14$, the remaining signal is at most $4^{-L}$: for $L=5$ about
$9.77\times10^{-4}$, and for $L=10$ about $9.54\times10^{-7}$. A layer far from the output
receives a signal that is only one millionth, so it practically does not learn. This symptom is
called the vanishing gradient. ReLU avoids it for a different reason: its derivative is $1$ on
the positive side, so there is no shrinkage, but its derivative is $0$ on the negative side, so
there is unit death. Both diseases are not implementation errors; both are consequences of the
shape of the function.\index{vanishing gradient}

One more statement is often quoted beyond its content: one hidden layer is enough for all
continuous functions.

\begin{teorema}[Universal approximation, statement]
Suppose $\varphi$ is bounded, monotonically increasing, and has different limits at $-\infty$
and $+\infty$. For every compact set $K\subset\mathbb{R}^{d}$, every continuous function
$f:K\to\mathbb{R}$, and every $\varepsilon>0$, there are $N$ and parameters
$\mathbf{w}_{i},b_{i},v_{i}$ such that
\[
  \sup_{\mathbf{x}\in K}\Bigl|f(\mathbf{x})-\sum_{i=1}^{N}v_{i}\,
  \varphi(\mathbf{w}_{i}^{\top}\mathbf{x}+b_{i})\Bigr|<\varepsilon .
\]
\end{teorema}

The proof of this statement is not included here; the original proof is in Cybenko (1989) and a
more general version in Hornik (1991), and the extension to ReLU is in Leshno et al. (1993).
What is important to read carefully is what the theorem does \emph{not} say. It states that such
a set of functions is dense, not that a particular network equals $f$. It does not say how large
$N$ must be for error $\varepsilon$ in dimension $d$; in the worst case $N$ grows exponentially
with $d$ (Goodfellow et al., section 6.4.1). It also does not say that those parameters can be
found by gradient descent, and it says nothing about data not yet seen. The sentence ``a network
can approximate anything'' is therefore true as a statement of density and misleading as a
practical promise.

\section{From score to loss}
\tingkat{intermediate}

The network's output must be compared with the label through a single number that can be
minimized, and the choice of that number determines how much signal reaches the lower layers.
For binary classification with label $y\in\{0,1\}$ and output $a=\sigma(z_{2})$, the loss used
is the binary cross-entropy
\[
  L(a,y)=-\bigl[y\ln a+(1-y)\ln(1-a)\bigr].
\]
Cross-entropy is the negative logarithm of the confidence the model gives to the correct class,
so a confident mistake is punished far more heavily than a hesitant one.

\kamus{cross-entropy}{entropi silang}{A measure of surprise: the more confident the model is in
a wrong answer, the larger the penalty, and the penalty grows without bound.}\index{cross-entropy}

\begin{proposisi}[Gradient of cross-entropy]
For $a=\sigma(z)$ and $L$ the binary cross-entropy we have
$\dfrac{\partial L}{\partial z}=a-y$.
\end{proposisi}

\begin{proof}
Differentiate $L$ with respect to $a$:
\[
  \frac{\partial L}{\partial a}=-\frac{y}{a}+\frac{1-y}{1-a}
  =\frac{a-y}{a(1-a)} .
\]
By the chain rule and the lemma on the derivative of the sigmoid,
$\dfrac{\partial a}{\partial z}=a(1-a)$, so
\[
  \frac{\partial L}{\partial z}=\frac{a-y}{a(1-a)}\cdot a(1-a)=a-y .
\]
Numerator and denominator cancel, and that is the point.
\end{proof}

Compare with the squared loss $L=\tfrac12(a-y)^{2}$, which gives
$\partial L/\partial z=(a-y)\,a(1-a)$. The factor $a(1-a)$ does not cancel; it shrinks exactly
when the model is most wrong with most confidence. For $z=10$ and $y=0$,
$a(1-a)=4.5396\times10^{-5}$, so the signal to the weights is only $4.5394\times10^{-5}$ even
though the error is almost one. With cross-entropy the signal stays $a-y=0.999955$. The ratio of
the two numbers is about $2.2\times10^{4}$; this is why the sigmoid output is paired with
cross-entropy, not with the square.

For $K$ classes, the same pairing holds, and this time we prove it.

\kamus{softmax}{softmax}{Turns a list of scores into a division of probabilities that sums to one,
by comparing every score with the sum of them all.}\index{softmax}

\begin{proposisi}[Softmax and cross-entropy]
Suppose $\mathbf{p}=\operatorname{softmax}(\mathbf{z})$ with
$p_{k}=e^{z_{k}}/\sum_{j}e^{z_{j}}$, and $L=-\sum_{k}y_{k}\ln p_{k}$ with
$\sum_{k}y_{k}=1$. Then $\dfrac{\partial L}{\partial z_{i}}=p_{i}-y_{i}$.
\end{proposisi}

\begin{proof}
The component derivative of softmax is
$\dfrac{\partial p_{k}}{\partial z_{i}}=p_{k}(\delta_{ki}-p_{i})$, with $\delta_{ki}$ the
Kronecker delta. Hence
\[
  \frac{\partial L}{\partial z_{i}}
  =-\sum_{k}\frac{y_{k}}{p_{k}}\cdot p_{k}(\delta_{ki}-p_{i})
  =-\sum_{k}y_{k}\delta_{ki}+p_{i}\sum_{k}y_{k}
  =-y_{i}+p_{i},
\]
because $\sum_{k}y_{k}=1$.
\end{proof}

For two classes, $\operatorname{softmax}((t,0))_{1}=\sigma(t)$, so this formula returns to $a-y$
as before. Note also that adding a constant $c$ to all $z_{k}$ does not change $\mathbf{p}$,
because the factor $e^{c}$ cancels; only differences of scores are meaningful, which is also why
adding a large constant to all logits damages floating-point accuracy (Chapter 22).

\begin{untukanak}
Have you ever played a number-guessing game where, every time you are wrong, you are only told
``too big'' or ``too small'', without being told by how much? If your guess is nine hundred and
the answer is one hundred, the remark ``too big'' is exactly the same as if your guess were one
hundred and one.

A neural network has the same problem, but more complicated. It has thirteen knobs, and every
knob affects the answer through several intermediaries. If we only know that the answer is
wrong, we do not know which knob to turn, in which direction, or by how much.

Fortunately, a wrong answer has a degree: how far off, and in which direction. With one measure
called the loss, we can ask every knob: if you were turned a little, does the loss go up or down,
and how fast? The answer is one number per knob.

In this chapter we will learn why a single loss measure is enough to give thirteen answers at
once, and why those answers can be trusted.
\end{untukanak}

\section{Backpropagation}
\tingkat{advanced}

A network is a composition of functions, so its derivative is computed with the chain rule of
Chapter 12. What makes it cheap is the order of work: from the output backward to the input, not
from the input forward to the output. For every intermediate quantity $u$ in the computational
graph, we assign one number stating how much the loss changes if $u$ changes.

\kamus{adjoint}{adjoint}{The image of how much the final error is affected by one quantity in the
middle of the network; one number per quantity.}

\begin{definisi}[Adjoint]
Suppose $C$ is the loss on one example and $u$ one of the intermediate quantities. The adjoint of
$u$ is $\bar{u}=\partial C/\partial u$.\index{adjoint}
\end{definisi}

\begin{proposisi}[Adjoint recursion]
If $u$ affects $C$ only through $v_{1},\dots,v_{m}$, then
$\bar{u}=\sum_{i=1}^{m}\bar{v}_{i}\,\dfrac{\partial v_{i}}{\partial u}$.
\end{proposisi}

\begin{proof}
This is the chain rule for functions of several variables (Chapter 12) applied to $C$ as a
function of $v_{1},\dots,v_{m}$, with each $v_{i}$ as a function of $u$. Since $u$ does not
affect $C$ in any other way, all other terms in that chain rule are zero.
\end{proof}

That recursion turns the work into a single pass. Starting from $\bar{C}=1$, we visit the
intermediate quantities in the reverse of the forward order. Every quantity receives
contributions from all the quantities that use it, so nothing is missed or counted twice. This
is backpropagation, and it is not a new algorithm: it is only the chain rule run on the graph in
the right order.\index{backpropagation}

Figure~\ref{fig:graf-balik} places both directions on one graph. Full arrows carry the forward
values; dashed arrows carry the adjoints, and the numbers on those arrows are taken from the
worked example in the next section.

\begin{figure}[htbp]
\centering
\begin{tikzpicture}[>=latex, font=\small]
  \node[draw=layerteal,rounded corners=2pt,fill=layerteal!10,inner sep=4pt] (X)  at (0,0)    {$\mathbf{x}$};
  \node[draw=layerteal,rounded corners=2pt,fill=layerteal!10,inner sep=4pt] (Z1) at (2.7,0)  {$\mathbf{z}_{1}$};
  \node[draw=layerteal,rounded corners=2pt,fill=layerteal!10,inner sep=4pt] (H)  at (5.3,0)  {$\mathbf{h}$};
  \node[draw=layerteal,rounded corners=2pt,fill=layerteal!10,inner sep=4pt] (Z2) at (7.5,0)  {$z_{2}$};
  \node[draw=layerteal,rounded corners=2pt,fill=layerteal!10,inner sep=4pt] (A)  at (9.3,0)  {$a$};
  \node[draw=layerteal,rounded corners=2pt,fill=layerteal!10,inner sep=4pt] (L)  at (11.4,0) {$L$};
  \node (W1) at (2.7,1.45) {$\mathbf{W}_{1},\mathbf{b}_{1}$};
  \node (W2) at (7.5,1.45) {$\mathbf{w}_{2},b_{2}$};
  \foreach \a/\b in {X/Z1,Z1/H,H/Z2,Z2/A,A/L} \draw[->,thick,ink] (\a) -- (\b);
  \draw[->,thick,ink] (W1) -- (Z1);
  \draw[->,thick,ink] (W2) -- (Z2);
  \draw[->,dashed,accent] (L.south)  .. controls +(0,-1.0) and +(0,-1.0) .. (Z2.south);
  \draw[->,dashed,accent] (Z2.south) .. controls +(0,-1.0) and +(0,-1.0) .. (H.south);
  \draw[->,dashed,accent] (H.south)  .. controls +(0,-1.0) and +(0,-1.0) .. (Z1.south);
  \draw[->,dashed,accent] (Z1.south) .. controls +(0,-1.0) and +(0,-1.0) .. (X.south);
  \draw[->,dashed,accent] (3.05,0.25) -- (3.05,1.25);
  \draw[->,dashed,accent] (7.85,0.25) -- (7.85,1.25);
  \node[font=\scriptsize] at (9.9,-1.4)  {$\bar{z}_{2}=a-y=-0.2554$};
  \node[font=\scriptsize] at (6.5,-2.0)  {$\bar{\mathbf{h}}=(-0.2043;\,0.1532)$};
  \node[font=\scriptsize] at (4.0,-1.4)  {$\bar{\mathbf{z}}_{1}=\bar{\mathbf{h}}$};
  \node[font=\scriptsize,align=center] at (1.1,-1.95)
    {$\nabla_{\mathbf{x}}L=(-0.0868;\,0.1328;$\\$(-0.0971;\,-0.1430)$};
  \node[font=\scriptsize] at (2.7,2.15) {$\nabla_{\mathbf{W}_{1}}L=\bar{\mathbf{z}}_{1}\mathbf{x}^{\top}$};
  \node[font=\scriptsize] at (9.0,2.15) {$\nabla_{\mathbf{w}_{2}}L=\bar{z}_{2}\mathbf{h}$};
\end{tikzpicture}
\caption{The computational graph of the 4--2--1 network on the worked example. Thick full arrows
are the forward direction: input and weights become $\mathbf{z}_{1}$, the ReLU gate becomes
$\mathbf{h}$, then the score $z_{2}$, the output $a$, and the loss $L$. Dashed arrows are the
backward direction that carries the adjoints: $\bar{z}_{2}=a-y=-0.2554$ enters the output
weights, $\bar{\mathbf{h}}=(-0.2043;\,0.1532)$ enters the gate, and
$\nabla_{\mathbf{x}}L=(-0.0868;\,0.1328;\,-0.0971;\,-0.1430)$ exits at the input. One backward
pass produces all thirteen gradients at once.}
\label{fig:graf-balik}
\end{figure}

\begin{bacakeras}
Two voices. First voice: ``The loss is too large.'' Second voice: ``Who is at fault?'' First
voice: ``Not one person. The fault is at the end, and the fault is divided backward, little by
little.''

Listen once more, slowly. The error is born at the end, in the last number. From there it walks
backward, stops at every junction, and asks: through here, how much is your share? The one
closest to the end answers largest. The one far behind answers smallest, because the answer has
already been filtered twice on the way home.

Second voice: ``So the ones in front are not at fault?'' First voice: ``They are at fault to the
extent of the share that reaches them. No more. That is why we do not accuse, we compute.''

Repeat three times until it sounds right: fault at the end, divided backward, little by little,
along the path.
\end{bacakeras}


Now derive it in full for the 4--2--1 network with ReLU in the hidden layer and sigmoid at the
output. The forward computation is
\[
  \mathbf{z}_{1}=\mathbf{W}_{1}\mathbf{x}+\mathbf{b}_{1},\qquad
  \mathbf{h}=\operatorname{ReLU}(\mathbf{z}_{1}),\qquad
  z_{2}=\mathbf{w}_{2}^{\top}\mathbf{h}+b_{2},\qquad
  a=\sigma(z_{2}),
\]
then $L$ is the binary cross-entropy of $a$ and $y$. Every link in the chain has its own local
derivative, and we already have them all:
\[
  \frac{\partial L}{\partial z_{2}}=a-y\ \text{(previous proposition)},\qquad
  \frac{\partial z_{2}}{\partial h_{i}}=w_{2}[i],
\]
\[
  \frac{\partial h_{i}}{\partial z_{1}[i]}=\mathbb{1}\bigl[z_{1}[i]>0\bigr],\qquad
  \frac{\partial z_{1}[i]}{\partial W_{1}[i,j]}=x_{j}.
\]
Multiplying backward along that chain gives, one by one,
\[
  \bar{z}_{2}=a-y,\qquad
  \nabla_{\mathbf{w}_{2}}L=\bar{z}_{2}\,\mathbf{h},\qquad
  \frac{\partial L}{\partial b_{2}}=\bar{z}_{2},
\]
then the signal that reaches the hidden layer
\[
  \bar{\mathbf{h}}=\bar{z}_{2}\,\mathbf{w}_{2},
  \qquad
  \bar{\mathbf{z}}_{1}=\bar{\mathbf{h}}\odot\mathbb{1}\bigl[\mathbf{z}_{1}>0\bigr],
\]
with $\odot$ the componentwise product. Note that $\bar{\mathbf{z}}_{1}$ is $\bar{\mathbf{h}}$
shaved by the ReLU gate: components whose gate is closed become exactly zero, not small. After
that
\[
  \nabla_{\mathbf{W}_{1}}L=\bar{\mathbf{z}}_{1}\mathbf{x}^{\top},\qquad
  \nabla_{\mathbf{b}_{1}}L=\bar{\mathbf{z}}_{1},\qquad
  \nabla_{\mathbf{x}}L=\mathbf{W}_{1}^{\top}\bar{\mathbf{z}}_{1}.
\]
The outer-product form of $\nabla_{\mathbf{W}_{1}}L$ is no accident: the derivative
$\partial z_{1}[i]/\partial W_{1}[i,j]=x_{j}$ does not depend on $i$, so the whole $i$-th row is
$\bar{z}_{1}[i]$ times the vector $\mathbf{x}$. Every weight receives a gradient whose size is
proportional to the input it multiplies, and that is why large-scale features automatically
receive large steps.

None of those formulas depends on $d=4$ or on two hidden units; the pattern is the same for any
layer.

The cost is comparable: one forward step needs one matrix-vector product per layer, one backward
step one transposed product and one outer product per layer. A centered finite difference for $P$
parameters needs $2P$ forward steps, whereas backpropagation needs about two. For $P=13$ in this
chapter, twenty-six against two; for a million parameters, two million against two, and that is
the whole difference. Nielsen names this observation as the reason backpropagation is regarded
as a great discovery, not merely an algebraic trick.

\begin{tangan}
Take three sheets of paper and a pencil, and label them: INPUT, WEIGHTS, OUTPUT. On the INPUT
sheet write four numbers: $1$; $0.5$; $-1$; $2$. On the WEIGHTS sheet write the two weights to
the output: $0.8$ and $-0.6$. On the OUTPUT sheet write $0.7446$, the probability the network
gives, and remember that the true label is $1$.

Lay the three sheets in a row from left to right: INPUT, WEIGHTS, OUTPUT. Now work backward,
always from right to left. On the OUTPUT sheet, subtract $0.7446$ from $1$ and write $-0.2554$.
Move to the WEIGHTS sheet: multiply $-0.2554$ by $0.8$, write $-0.2043$; multiply by $-0.6$,
write $0.1532$. Move to the INPUT sheet: multiply $-0.2043$ by each number on that sheet, then
write the row: $-0.2043$; $-0.1022$; $0.2043$; $-0.4086$.

What you must see: every time a sheet moves one step to the left, the sign of the number can
flip, and its size follows the contents of the sheet you are holding. If you reverse the
direction, from INPUT to OUTPUT, the order of work is wrong and the result means nothing.
\end{tangan}


\section{Initialization and gradient descent}
\tingkat{intermediate}

Once the gradient exists, the next step is simple: move a little in the direction that lowers the
loss. For all parameters $\boldsymbol{\theta}$,
\[
  \boldsymbol{\theta}\leftarrow\boldsymbol{\theta}-\eta\,\nabla_{\boldsymbol{\theta}}L,
\]
with $\eta$ the learning rate. For $B$ examples and the average loss
$L=\frac{1}{B}\sum_{i}L_{i}$, the gradient is also the average of the per-example gradients; the
backpropagation formula does not change, only how many backward passes. The conditions and
bounds on $\eta$ are in Chapter~23; because the network is nonconvex, the guarantees there are
only a benchmark.

\kamus{stochastic gradient descent}{penurunan gradien stokastik}{Walking down a slope using a
slope estimate from a small part of the data, not from all of the data.}\index{gradient descent}

Three variants deserve mention. SGD takes a small $B$ and updates the parameters every time; the
gradient is noisy, and the spread of that noise shrinks like $1/\sqrt{B}$ because the variance of
the average of $B$ independent variables is $\sigma^{2}/B$. Momentum keeps a velocity,
\[
  \mathbf{v}\leftarrow\beta\mathbf{v}+\mathbf{g},\qquad
  \boldsymbol{\theta}\leftarrow\boldsymbol{\theta}-\eta\mathbf{v},
\]
so that consistent directions accumulate and oscillating directions cancel. Adam normalizes every
coordinate with an estimate of its second moment,
\[
  \mathbf{m}\leftarrow\beta_{1}\mathbf{m}+(1-\beta_{1})\mathbf{g},\qquad
  \mathbf{v}\leftarrow\beta_{2}\mathbf{v}+(1-\beta_{2})\mathbf{g}^{2},
\]
\[
  \hat{\mathbf{m}}=\frac{\mathbf{m}}{1-\beta_{1}^{t}},\qquad
  \hat{\mathbf{v}}=\frac{\mathbf{v}}{1-\beta_{2}^{t}},\qquad
  \boldsymbol{\theta}\leftarrow\boldsymbol{\theta}
  -\eta\,\frac{\hat{\mathbf{m}}}{\sqrt{\hat{\mathbf{v}}}+\epsilon},
\]
with $t$ the step number. The factors $1-\beta_{1}^{t}$ and $1-\beta_{2}^{t}$ are called bias
correction, and their purpose is to remove the tendency of the moments toward zero in the early
steps. On the first step, $\hat{m}_{1}=g$ and $\hat{v}_{1}=g^{2}$, so the update is exactly
$-\eta\,g/(\lvert g\rvert+\epsilon)$: its size is almost exactly $\eta$, however large the
gradient. That property makes Adam scale-insensitive, and at the same time explains its weakness:
because its step does not shrink when the gradient shrinks, Adam is not guaranteed to converge on
all convex problems, and the counterexample can be read in Reddi, Kale, and Kumar (2018).

What remains, and what is most often underestimated, is the starting point.

\begin{proposisi}[Symmetry that never breaks]
If two rows of $\mathbf{W}_{1}$ are identical and $b_{1}[1]=b_{1}[2]$ at the start of training,
then both stay identical after every gradient update, for every data set and every learning
rate.\index{initialization}
\end{proposisi}

\begin{proof}
We prove it by induction on the update step. Suppose the first and second rows of
$\mathbf{W}_{1}$ are the same and $b_{1}[1]=b_{1}[2]$. Then for every input $z_{1}[1]=z_{1}[2]$,
so both ReLU gates make the same decision and $\bar{z}_{1}[1]=\bar{z}_{1}[2]$. Because
$\nabla_{\mathbf{W}_{1}}L=\bar{\mathbf{z}}_{1}\mathbf{x}^{\top}$, those two rows of the gradient
are identical, and so are the two components of $\nabla_{\mathbf{b}_{1}}L$. An update with the
same rate preserves the equality. For a set of examples, the total gradient is the sum of the
per-example gradients, so the equality survives that too.
\end{proof}

The consequence is blunt: eight units initialized uniformly behave like one unit, however long
it is trained.

Zero initialization is worse still. If all of $\mathbf{W}_{1}$ and $\mathbf{b}_{1}$ are zero,
then $\mathbf{z}_{1}=\mathbf{0}$, and because the gate is $\mathbb{1}[z>0]$, all gates are
closed: the gradient with respect to $\mathbf{W}_{1}$ is exactly zero, and the weights stay zero
forever. The only thing that can still learn is $b_{2}$, and the result is a constant prediction
equal to the class proportion. On the lab data, the final loss is $0.690923$, exactly the label
entropy. Eight units, zero benefit.

\begin{tantangan}
Design a 2--2--1 network with ReLU in the hidden layer whose gradient is exactly zero on all six
weights and both biases of the first layer from the first step, while $\mathbf{W}_{1}$ is not
zero, and the network is still wrong on at least one point. Only $b_{2}$ may move. Show why your
choice works.

Write your answer as a table: unit, $z_{1}[i]$ for each example, its gate state, and the gradient
that comes in.

Hint 1: what freezes the weights is not the value of the weights, but the ReLU gate closing for
all examples.

Hint 2: use data in the nonnegative quadrant, then take
$\mathbf{W}_{1}=-\mathbf{1}\mathbf{1}^{\top}$ and $\mathbf{b}_{1}=\mathbf{0}$. Then
$z_{1}[i]=-(x_{1}+x_{2})\le0$, so $\mathbf{h}=\mathbf{0}$,
$\bar{\mathbf{z}}_{1}=\mathbf{0}$, and $\nabla_{\mathbf{W}_{1}}L=\mathbf{0}$ even though
$\mathbf{W}_{1}\neq\mathbf{0}$.

If you get stuck, go back to the section Initialization and gradient descent.
\end{tantangan}


Random initialization is therefore not decoration, but a condition for those units to explore
different functions. The next question: how large. The answer comes from computing variances. If
the components of $\mathbf{x}$ are independent with mean zero and variance $\sigma_{x}^{2}$, and
the weights are independent with variance $\sigma_{w}^{2}$, then
\[
  \operatorname{Var}\bigl(z_{1}[i]\bigr)
  =\sum_{j=1}^{n_{\mathrm{in}}}\operatorname{Var}\bigl(W_{1}[i,j]x_{j}\bigr)
  =n_{\mathrm{in}}\,\sigma_{w}^{2}\sigma_{x}^{2}.
\]
In order for that variance not to depend on $n_{\mathrm{in}}$, we need
$\sigma_{w}^{2}\propto1/n_{\mathrm{in}}$. The He initialization uses
$\sigma_{w}^{2}=2/n_{\mathrm{in}}$, and the number $2$ is not decoration: for symmetric $z$,
ReLU cuts away half the distribution, so $\mathbb{E}[h^{2}]=\operatorname{Var}(z)/2=\sigma_{x}^{2}$,
and that variance survives layer after layer. The Xavier initialization uses
$\sigma_{w}^{2}=2/(n_{\mathrm{in}}+n_{\mathrm{out}})$, a compromise between preserving the forward
signal and the backward signal. That compromise is exact only when
$n_{\mathrm{in}}=n_{\mathrm{out}}$.

For sigmoid layers, an initialization that is too large punishes immediately: on a two-layer
network with standard deviation $1.0$, the initial loss reaches $2.223899$, whereas with standard
deviation $0.5$ it is only $1.220409$. The reverse is more serious: with standard deviation $0.5$
and eight layers, training stops at accuracy $0.5333$. Initialization is not an administrative
detail; it determines whether the signal is still there when it reaches the first layer.

\section{Worked example}
\tingkat{intermediate}

\begin{resep}
\begin{enumerate}
  \item Start from the end: $\bar{z}_{2}=a-y$ for sigmoid with cross-entropy, or
        $\mathbf{p}-\mathbf{y}$ for softmax. Without this number there is nothing to divide.
  \item Output weights: $\nabla_{\mathbf{w}_{2}}L=\bar{z}_{2}\mathbf{h}$ and
        $\partial L/\partial b_{2}=\bar{z}_{2}$.
  \item Signal to the previous layer: $\bar{\mathbf{h}}=\bar{z}_{2}\mathbf{w}_{2}$, then shave
        with the gate
        $\bar{\mathbf{z}}_{1}=\bar{\mathbf{h}}\odot\mathbb{1}[\mathbf{z}_{1}>0]$.
  \item First-layer weights: $\nabla_{\mathbf{W}_{1}}L=\bar{\mathbf{z}}_{1}\mathbf{x}^{\top}$
        and $\nabla_{\mathbf{b}_{1}}L=\bar{\mathbf{z}}_{1}$; if needed, the input gradient
        $\nabla_{\mathbf{x}}L=\mathbf{W}_{1}^{\top}\bar{\mathbf{z}}_{1}$.
\end{enumerate}
A warning about order: step 3 must not come before step 1. $\bar{\mathbf{z}}_{1}$ is computed
from $\bar{z}_{2}$, and the gate uses $\mathbf{z}_{1}$ from the forward computation that has
already been stored. Swapping the two steps kills the entire gradient of the first-layer weights.
\end{resep}


We use one example and compute everything with exact numbers. The input is
$\mathbf{x}=(1.0;\,0.5;\,-1.0;\,2.0)$, the label is $y=1$, and the initial parameters are
\[
  \mathbf{W}_{1}=\begin{pmatrix}
  0.5&-0.5&0.25&1.0\\
  0.1&0.2&-0.3&0.4\end{pmatrix},\quad
  \mathbf{b}_{1}=\begin{pmatrix}0.1\\-0.2\end{pmatrix},\quad
  \mathbf{w}_{2}=\begin{pmatrix}0.8\\-0.6\end{pmatrix},\quad b_{2}=0.05 .
\]
That is all the parameters.

\emph{Forward step.} The first unit receives
$0.5(1)-0.5(0.5)+0.25(-1)+1(2)+0.1=0.5-0.25-0.25+2+0.1=2.1$.
The second unit receives
$0.1(1)+0.2(0.5)-0.3(-1)+0.4(2)-0.2=0.1+0.1+0.3+0.8-0.2=1.1$.
Both are positive, so $\mathbf{h}=(2.1;\,1.1)$ without clipping. The output score is
$z_{2}=0.8(2.1)-0.6(1.1)+0.05=1.68-0.66+0.05=1.07$, so
$a=\sigma(1.07)=0.7445969158$ and the loss is $L=-\ln a=0.2949122596$.

\emph{Backward step.} The output adjoint is $\bar{z}_{2}=a-y=-0.2554030842$. From there
\[
  \nabla_{\mathbf{w}_{2}}L=\bar{z}_{2}\mathbf{h}
  =(-0.5363464768;\,-0.2809433926),\qquad
  \frac{\partial L}{\partial b_{2}}=-0.2554030842 .
\]
The signal to the hidden layer is
$\bar{\mathbf{h}}=\bar{z}_{2}\mathbf{w}_{2}=(-0.2043224673;\,0.1532418505)$; both gates are
open, so $\bar{\mathbf{z}}_{1}=\bar{\mathbf{h}}$. The outer product gives
\[
  \nabla_{\mathbf{W}_{1}}L=
  \begin{pmatrix}
  -0.2043224673&-0.1021612337&0.2043224673&-0.4086449347\\
  0.1532418505&0.0766209253&-0.1532418505&0.3064837010
  \end{pmatrix},
\]
and $\nabla_{\mathbf{b}_{1}}L=(-0.2043224673;\,0.1532418505)$. Two quick checks: the second row
is exactly $-0.75$ times the first, because $-0.6/0.8=-0.75$; and every column is proportional
to $x_{j}$, so the fourth column is twice the first, the second is half of it, and the third is
its negative. The Euclidean norm of the whole gradient is $0.9796964023$.

One gradient-descent step with $\eta=0.5$ updates every parameter according to
$\theta\leftarrow\theta-\eta\,\partial L/\partial\theta$, giving
\[
  \mathbf{W}_{1}'=\begin{pmatrix}
  0.6021612337&-0.4489193832&0.1478387663&1.2043224673\\
  0.0233790747&0.1616895374&-0.2233790747&0.2467581495
  \end{pmatrix},
\]
\[
  \mathbf{b}_{1}'=(0.2021612337;\,-0.2766209253),\qquad
  \mathbf{w}_{2}'=(1.0681732384;\,-0.4595283037),
\]
\[
  b_{2}'=0.1777015421 .
\]
After that update $\mathbf{z}_{1}'=(2.8406689;\,0.5444983)$ and $a'=0.9508189702$, so the loss
falls to $0.0504315919$; the difference is $-0.2444806677$, or about $83\%$ of the initial loss.

The number $83\%$ invites a misreading. What is measured is a single example, with no test data,
on a landscape that is not convex. The valid conclusion is not ``the rate $0.5$ is good'', but
``this gradient direction lowers the loss on this example''.

\kertas{Check one row of the lab table by hand. Take $W_{1}[2,1]$ and compute
$\bar{z}_{1}[2]\,x_{1}$ to get $0.1532418505$. Then compute the centered finite difference for
that parameter yourself: run one forward step with $W_{1}[2,1]+10^{-6}$ and one forward step with
$W_{1}[2,1]-10^{-6}$, divide the difference of the losses by $2\times10^{-6}$, and compare with
the analytic number above. After that, explain why the result is not exactly the same, and why
the sign of the difference can change if $h$ is made smaller, $10^{-12}$.}

\section{Computational lab}
\tingkat{advanced}

The three blocks are run with Python 3.14.4 and NumPy 2.5.2, without machine-learning libraries;
the output is copied as it stands.

The first block checks the hand computation above, then compares every analytic gradient with a
centered finite difference at $h=10^{-6}$.

\begin{lab}
\begin{verbatim}
import numpy as np
np.set_printoptions(precision=10, suppress=True, linewidth=76)

# 4-2-1 network: ReLU hidden, sigmoid output
x  = np.array([1.0, 0.5, -1.0, 2.0])
y  = 1.0
W1 = np.array([[0.5, -0.5, 0.25, 1.0],
               [0.1,  0.2, -0.3,  0.4]])
b1 = np.array([0.1, -0.2])
w2 = np.array([0.8, -0.6])
b2 = 0.05

def forward(x, W1, b1, w2, b2):
    z1 = W1 @ x + b1
    h  = np.maximum(z1, 0.0)
    o  = w2 @ h + b2
    p  = 1.0/(1.0+np.exp(-o))
    return z1, h, o, p, -(y*np.log(p) + (1-y)*np.log(1-p))

def backward(x, W1, b1, w2, b2):
    z1, h, o, p, L = forward(x, W1, b1, w2, b2)
    d   = p - y                    # dL/dz2 for sigmoid + cross-entropy
    dz1 = d*w2*(z1 > 0)            # ReLU gate
    return dict(W1=np.outer(dz1, x), b1=dz1, w2=d*h, b2=d,
                x=W1.T @ dz1), L, p, z1, h, dz1

g, L, p, z1, h, dz1 = backward(x, W1, b1, w2, b2)
print("z1 =", z1, "  h =", h)
print("p = %.12f  L = %.12f" % (p, L))
print("dz1 =", dz1)
print("g_w2 =", g["w2"], " g_b2 = %.12f" % g["b2"])
for i in range(2):
    print("g_W1[%d] =" % i, np.array2string(g["W1"][i], precision=10,
                                            suppress_small=False))
print("g_x =", g["x"])

# centered finite-difference check, one parameter at a time
eps = 1e-6
def loss_at(nama, idx, delta):
    W, B, V, c = W1.copy(), b1.copy(), w2.copy(), b2
    if nama == "W1":  W[idx] += delta
    if nama == "b1":  B[idx] += delta
    if nama == "w2":  V[idx] += delta
    if nama == "b2":  c += delta
    return forward(x, W, B, V, c)[4]

daftar = [("W1", i, j) for i in range(2) for j in range(4)]
daftar += [("b1", i) for i in range(2)] + [("w2", i) for i in range(2)]
print("\nparameter     analytic        numeric     relative error")
for nama, *idx in daftar + [("b2",)]:
    idx = tuple(idx)
    ana = {"W1": g["W1"], "b1": g["b1"], "w2": g["w2"], "b2": np.array(g["b2"])}
    a = ana[nama][idx] if idx else ana[nama]
    n = (loss_at(nama, idx, eps) - loss_at(nama, idx, -eps))/(2*eps)
    rel = abs(a-n)/max(abs(a), abs(n), 1e-300)
    print("%-10s %15.10f %15.10f %12.2e" %
          (nama + "".join("[%d]" % i for i in idx), a, n, rel))

# one gradient-descent step with rate 0.5
eta = 0.5
W1n, b1n = W1 - eta*g["W1"], b1 - eta*g["b1"]
w2n, b2n = w2 - eta*g["w2"], b2 - eta*g["b2"]
z1b, hb, ob, pb, Lb = forward(x, W1n, b1n, w2n, b2n)
print("\none step: p = %.12f  L = %.12f  difference = %.12f"
      % (pb, Lb, L - Lb))
for i in range(2):
    print("W1'[%d] =" % i, np.array2string(W1n[i], precision=10))
print("w2' =", w2n, " b2' = %.12f" % b2n)

# gradient with respect to the input: analytic vs finite difference
E = np.eye(4)
gx = np.array([(forward(x + eps*E[j], W1, b1, w2, b2)[4]
              - forward(x - eps*E[j], W1, b1, w2, b2)[4])/(2*eps)
               for j in range(4)])
print("\ng_x analytic    =", g["x"])
print("g_x finite diff =", gx)
print("relative error   =", ["%.1e" % v for v in np.abs(gx-g["x"])/np.abs(g["x"])])
\end{verbatim}
The actual output:
\begin{verbatim}
z1 = [2.1 1.1]   h = [2.1 1.1]
p = 0.744596915827  L = 0.294912259605
dz1 = [-0.2043224673  0.1532418505]
g_w2 = [-0.5363464768 -0.2809433926]  g_b2 = -0.255403084173
g_W1[0] = [-0.2043224673 -0.1021612337  0.2043224673 -0.4086449347]
g_W1[1] = [ 0.1532418505  0.0766209253 -0.1532418505  0.306483701 ]
g_x = [-0.0868370486  0.1328096038 -0.097053172  -0.1430257271]

parameter     analytic        numeric     relative error
W1[0][0]     -0.2043224673   -0.2043224674     1.17e-10
W1[0][1]     -0.1021612337   -0.1021612337     2.53e-10
W1[0][2]      0.2043224673    0.2043224674     1.17e-10
W1[0][3]     -0.4086449347   -0.4086449347     4.95e-11
W1[1][0]      0.1532418505    0.1532418505     7.22e-11
W1[1][1]      0.0766209253    0.0766209252     8.33e-10
W1[1][2]     -0.1532418505   -0.1532418505     7.22e-11
W1[1][3]      0.3064837010    0.3064837010     1.09e-10
b1[0]        -0.2043224673   -0.2043224674     1.17e-10
b1[1]         0.1532418505    0.1532418506     6.16e-10
w2[0]        -0.5363464768   -0.5363464768     9.80e-11
w2[1]        -0.2809433926   -0.2809433925     2.41e-10
b2           -0.2554030842   -0.2554030842     3.59e-11

one step: p = 0.950818970185  L = 0.050431591881  difference = 0.244480667724
W1'[0] = [ 0.6021612337 -0.4489193832  0.1478387663  1.2043224673]
W1'[1] = [ 0.0233790747  0.1616895374 -0.2233790747  0.2467581495]
w2' = [ 1.0681732384 -0.4595283037]  b2' = 0.177701542086

g_x analytic    = [-0.0868370486  0.1328096038 -0.097053172  -0.1430257271]
g_x finite diff = [-0.0868370486  0.1328096037 -0.097053172  -0.1430257272]
relative error   = ['7.7e-10', '5.0e-10', '1.4e-10', '9.8e-11']
\end{verbatim}
\end{lab}

That table is the proof that was asked for. All thirteen parameters agree to nine or ten digits,
with the largest relative error $8.33\times10^{-10}$ at $W_{1}[1,2]$. That number should not be
zero: the centered finite difference carries a truncation error of order $h^{2}$ and a rounding
error of order $\varepsilon_{\mathrm{mach}}/h$ (Chapter 22), so there is a best $h$ that cannot
make the error zero. For $h=10^{-6}$ an order of $10^{-10}$ is exactly what is expected.
 

The gradient with respect to the input agrees to $7.7\times10^{-10}$ on the first component. That
number means: raising the second feature slightly raises the loss fastest, whereas raising the
fourth feature lowers it fastest. In Chapter 29 this direction is used in reverse, to raise the
loss as fast as possible on inputs the attacker is allowed to change.

The second block trains a 2--8--1 network from scratch on two-dimensional data with an XOR
pattern, then compares it with a linear baseline and with two defective initializations.

\begin{lab}
\begin{verbatim}
import numpy as np

rng = np.random.default_rng(7)
n = 240
X = rng.uniform(-1.0, 1.0, size=(n, 2))
y = (X[:, 0]*X[:, 1] > 0).astype(float)        # label 1 when the signs agree
print("examples =", n, " class 1 =", int(y.sum()))

def sig(o):
    return 1.0/(1.0+np.exp(-o))

def bce(p, t):
    p = np.clip(p, 1e-12, 1-1e-12)
    return float(-np.mean(t*np.log(p) + (1-t)*np.log(1-p)))

def init(mode, rs, nh=8):
    if mode == "nol":
        return np.zeros((nh, 2)), np.zeros(nh), np.zeros(nh), 0.0
    if mode == "seragam":
        return np.full((nh, 2), 0.5), np.zeros(nh), np.full(nh, 0.5), 0.0
    return (rs.normal(0, np.sqrt(2/2), (nh, 2)), np.zeros(nh),
            rs.normal(0, np.sqrt(2/nh), nh), 0.0)          # He

def fwd(X, W1, b1, w2, b2):
    Z1 = X @ W1.T + b1
    H  = np.maximum(Z1, 0.0)
    return Z1, H, sig(H @ w2 + b2)

def grads(X, t, W1, b1, w2, b2):
    Z1, H, P = fwd(X, W1, b1, w2, b2)
    d  = (P - t)/len(t)
    dZ = np.outer(d, w2)*(Z1 > 0)
    return dZ.T @ X, dZ.sum(0), H.T @ d, d.sum()

def latih(mode, langkah=3000, lr=0.5, seed=0):
    W1, b1, w2, b2 = init(mode, np.random.default_rng(seed))
    jejak = []
    for t in range(langkah + 1):
        gW1, gb1, gw2, gb2 = grads(X, y, W1, b1, w2, b2)
        W1, b1 = W1 - lr*gW1, b1 - lr*gb1
        w2, b2 = w2 - lr*gw2, b2 - lr*gb2
        if t % 500 == 0:
            jejak.append((t, bce(fwd(X, W1, b1, w2, b2)[2], y)))
    Z1, H, P = fwd(X, W1, b1, w2, b2)
    mati = int(np.sum(np.all(H <= 0, axis=0)))
    return jejak, float(np.mean((P > 0.5) == (y > 0.5))), bce(P, y), mati

jejak, akur, L, mati = latih("he")
for t, v in jejak:
    print("step %4d   L = %.6f" % (t, v))
print("He      : accuracy = %.4f  dead units = %d/8" % (akur, mati))
for mode in ("nol", "seragam"):
    j, a, l, m = latih(mode)
    print("%-8s: L final = %.6f  accuracy = %.4f  dead units = %d/8"
          % (mode, l, a, m))

w, b = np.zeros(2), 0.0             # garis dasar linear, tanpa lapisan
for t in range(4000):
    p = sig(X @ w + b)
    w -= 0.5*(X.T @ (p - y))/n
    b -= 0.5*np.mean(p - y)
pl = sig(X @ w + b)
print("linear  : L = %.6f  accuracy = %.4f"
      % (bce(pl, y), float(np.mean((pl > 0.5) == (y > 0.5)))))
\end{verbatim}
The actual output:
\begin{verbatim}
examples = 240  class 1 = 112
step    0   L = 0.648815
step  500   L = 0.060245
step 1000   L = 0.036445
step 1500   L = 0.027164
step 2000   L = 0.022039
step 2500   L = 0.018717
step 3000   L = 0.016349
He      : accuracy = 1.0000  dead units = 0/8
nol     : L final = 0.690923  accuracy = 0.5333  dead units = 8/8
seragam : L final = 0.539475  accuracy = 0.7167  dead units = 0/8
linear  : L = 0.686130  accuracy = 0.7208
\end{verbatim}
\end{lab}

Its loss curve decreases at every sampling, from $0.648815$ to $0.016349$. Training accuracy
reaches $1.0000$ with eight units and zero dead units. The linear baseline, which uses the same
optimization and the same data but no hidden layer, stops at $0.7208$. That difference of $0.28$
is the price of a single separating plane, and it agrees with the XOR proposition: this pattern
is not linearly separable.

Figure~\ref{fig:kurva-kerugian} plots those seven points.

\begin{figure}[htbp]
\centering
\begin{tikzpicture}
\begin{axis}[
  width=12.0cm, height=5.6cm,
  xlabel={training step}, ylabel={loss $L$},
  xmin=-80, xmax=3080, ymin=0, ymax=0.75,
  axis lines=left,
  grid=major, grid style={layergrey!15},
  tick label style={font=\scriptsize},
  ytick={0,0.2,0.4,0.6}, yticklabels={$0$,$0.2$,$0.4$,$0.6$},
  label style={font=\small},
  legend style={font=\scriptsize, draw=none, fill=none,
                at={(0.98,0.98)}, anchor=north east},
]
\addplot[thick, accent, mark=*, mark size=1.7pt] coordinates {
  (0,0.648815) (500,0.060245) (1000,0.036445) (1500,0.027164)
  (2000,0.022039) (2500,0.018717) (3000,0.016349) };
\addplot[dashed, layergrey, domain=0:3000, samples=2] {0.693147};
\legend{training loss, $\ln 2=0.693147$}
\end{axis}
\end{tikzpicture}
\caption{The training loss curve of the 2--8--1 network in the second block, from $0.648815$ at
step $0$ to $0.016349$ at step $3000$. The dashed line is $\ln 2=0.693147$, the loss of a
uniform guess; the first point is already below it because it is measured after the first
update.}
\label{fig:kurva-kerugian}
\end{figure}


Zero initialization teaches a harsher lesson. Its final loss is $0.690923$, equal to the label
entropy to six decimals, and its accuracy is $0.5333$, namely the proportion of the majority
class; all eight units are dead. Uniform initialization gives $0.7167$, practically the same as
the linear baseline.

The third block measures two things: how many layers a sigmoid network can still learn through,
and how fast the signal shrinks from the output to the first layer.

\begin{lab}
\begin{verbatim}
import numpy as np

rng = np.random.default_rng(7)
n = 240
X = rng.uniform(-1.0, 1.0, size=(n, 2))
y = (X[:, 0]*X[:, 1] > 0).astype(float)

def sig(o):
    return 1.0/(1.0+np.exp(-o))

def bce(p, t):
    p = np.clip(np.ravel(p), 1e-12, 1-1e-12)
    t = np.ravel(t)
    return float(-np.mean(t*np.log(p) + (1-t)*np.log(1-p)))

def bangun(dalam, std, seed=0):
    rs = np.random.default_rng(seed)
    ukuran = [2] + [16]*(dalam-1) + [1]
    W = [rs.normal(0, std, (ukuran[i+1], ukuran[i])) for i in range(dalam)]
    return W, [np.zeros(ukuran[i+1]) for i in range(dalam)]

def maju(W, b, X):
    A = [X]
    for Wl, bl in zip(W[:-1], b[:-1]):
        A.append(sig(A[-1] @ Wl.T + bl))
    return A, sig(A[-1] @ W[-1].T + b[-1])

def mundur(W, b, X, t):
    A, P = maju(W, b, X)
    d = (P - t[:, None])/len(t)
    gW, gb = [None]*len(W), [None]*len(b)
    gW[-1], gb[-1] = d.T @ A[-1], d.sum(0)
    delta = d @ W[-1]
    for l in range(len(W)-1, 0, -1):
        delta = delta*A[l]*(1 - A[l])
        gW[l-1], gb[l-1] = delta.T @ A[l-1], delta.sum(0)
        if l > 1:
            delta = delta @ W[l-1]
    return gW, gb, bce(P, t), P

print("depth  std   L(start)   L(final)  accuracy  ||gW1||/||gWlast||")
for dalam in (1, 2, 4, 8):
    for std in (0.5, 1.0):
        W, b = bangun(dalam, std)
        gW, gb, L0, P = mundur(W, b, X, y)
        rasio = np.linalg.norm(gW[0])/np.linalg.norm(gW[-1])
        for t in range(2000):
            gW, gb, L, P = mundur(W, b, X, y)
            for l in range(dalam):
                W[l] -= 0.5*gW[l]
                b[l] -= 0.5*gb[l]
        akur = float(np.mean((np.ravel(P) > 0.5) == (y > 0.5)))
        print("%5d %5.1f %9.6f %9.6f %8.4f %14.2e"
              % (dalam, std, L0, bce(P, y), akur, rasio))
\end{verbatim}
The actual output:
\begin{verbatim}
depth  std   L(start)   L(final)  accuracy  ||gW1||/||gWlast||
    1   0.5  0.693190  0.686130   0.7208       1.00e+00
    1   1.0  0.693951  0.686130   0.7208       1.00e+00
    2   0.5  1.220409  0.332913   0.8417       1.66e-02
    2   1.0  2.223899  0.299519   0.8542       2.16e-02
    4   0.5  0.691597  0.147887   0.9458       5.77e-02
    4   1.0  0.701977  0.023491   0.9958       2.05e-01
    8   0.5  0.739920  0.690867   0.5333       1.26e-04
    8   1.0  0.875314  0.002383   1.0000       3.76e-03
\end{verbatim}
\end{lab}

The first three rows of that table are a promise and a warning at once. Depth one gives $0.7208$,
the same as the linear baseline in the previous block; depth two gives $0.8417$ and depth four
$0.9458$. But the last row with standard deviation $0.5$ gives $0.5333$ with loss $0.690867$:
eight sigmoid layers with a small initialization practically do not learn. The gradient-norm
ratio in the last column explains why: $1.26\times10^{-4}$, four orders of magnitude smaller than
at depth four. The same row with standard deviation $1.0$ succeeds perfectly. So the failure is
not about depth alone, but about the initial scale.

Two limits of these lab results must be stated. The data is only 240 self-made two-dimensional
examples, with no test set, so all of it is training accuracy; generalization is the business of
Chapter 26 and Chapter 27. All experiments use a single seed, so variation between seeds is
invisible. What may be concluded is only what was checked directly: backpropagation agrees with
finite differences, the loss decreases, zero and uniform initialization fail in ways that can be
explained, and a deep network with a small initialization loses its signal.

\begin{untukdosen}
The order that has worked: start from one separating line and four XOR points, then show the
2--2--1 construction on the board, and only after that the general definition. The proof that XOR
is impossible is usually accepted easily; what stops students is the backward step, so set aside
the most time for backpropagation, not for activation functions.

Two typical misconceptions. First, ``backpropagation is another algorithm, not the chain rule'',
so that students memorize formulas without being able to derive them for another architecture;
the cure is to ask them to re-derive the 4--2--1 network with three hidden units on the board.
Second, ``a relative error of $10^{-10}$ means the implementation is proved correct''; a small
error only checks agreement at one point, for one set of parameters, and does not check the sign,
the scaling, or the case of a closed gate.

The most helpful board activity: give a network with one hidden unit whose $z_{1}<0$, ask the
class to compute the gradient of the whole row $\mathbf{W}_{1}$, then ask what happens to that
unit after a hundred steps. The answer ``nothing happens, ever'' usually opens the discussion
about unit death and about why a large learning rate is dangerous.

Sign of readiness to go on: students can write the backpropagation for a network with two hidden
layers without looking at their notes, and can explain why zero initialization fails not because
of ``too little data'' but because of symmetry.
\end{untukdosen}

\section{Summary}

\begin{ringkasan}
\begin{itemize}
  \item A linear threshold can only separate sets that are linearly separable, and the four XOR
        points are not one of them. The proof does not depend on the algorithm, so adding data
        does not help. One hidden layer with ReLU is already enough for that pattern.
  \item Without a nonlinear activation, a layered network collapses into a single matrix. Sigmoid
        and tanh shrink the signal because their derivatives are at most $\tfrac14$ and $1$;
        through $L$ sigmoid layers the remaining signal is at most $4^{-L}$, namely
        $9.54\times10^{-7}$ for $L=10$. ReLU does not shrink, but it kills units whose input is
        always negative.
  \item Cross-entropy is paired with sigmoid or softmax so that the output adjoint is exactly
        $a-y$ or $\mathbf{p}-\mathbf{y}$. With the squared loss, the factor $a(1-a)$ is left
        behind and the signal shrinks exactly when the model is most confidently wrong; at $z=10$
        the ratio is about $2.2\times10^{4}$.
  \item Backpropagation is the chain rule of Chapter 12 run on the computational graph in reverse
        order. For 4--2--1, all the gradients are:
        $\bar{z}_{2}=a-y$, $\nabla_{\mathbf{w}_{2}}L=\bar{z}_{2}\mathbf{h}$,
        $\bar{\mathbf{z}}_{1}=\bar{z}_{2}\mathbf{w}_{2}\odot\mathbb{1}[\mathbf{z}_{1}>0]$,
        $\nabla_{\mathbf{W}_{1}}L=\bar{\mathbf{z}}_{1}\mathbf{x}^{\top}$,
        $\nabla_{\mathbf{x}}L=\mathbf{W}_{1}^{\top}\bar{\mathbf{z}}_{1}$. Its cost is about two
        forward steps, not $2P$. On the 4--2--1 example, those thirteen gradients agree with
        centered finite differences to a relative error of $8.33\times10^{-10}$.
  \item Initialization determines whether the network can learn. Zero closes all gates and leaves
        a prediction equal to the class proportion; uniform preserves the symmetry so that the
        effective capacity is one unit; too small shrinks the signal before the first layer; too
        large saturates the sigmoid. He uses $\sigma_{w}^{2}=2/n_{\mathrm{in}}$, Xavier
        $2/(n_{\mathrm{in}}+n_{\mathrm{out}})$.
  \item Gradient descent, SGD, momentum, and Adam differ only in how they use the same gradient.
        Adam's first step has a size of almost exactly $\eta$ however large the gradient, and that
        property explains both its robustness and its failure on some problems.
  
\end{itemize}
\end{ringkasan}

\section{Exercises}

\begin{latihan}
\begin{enumerate}
  \item For data with margin $\gamma=0.2$ and $\lVert\mathbf{x}_{i}\rVert\le1$, what is the upper
        bound on the number of perceptron updates according to the convergence theorem, and what
        is it for $\gamma=0.05$? Explain why that bound does not depend on the number of examples
        and why it is useless when $\gamma$ is small.

  \item Prove that a two-layer network with identity activation in the hidden layer is no stronger
        than a linear threshold, that is, that
        $\mathbf{W}_{2}\operatorname{id}(\mathbf{W}_{1}\mathbf{x}+\mathbf{b}_{1})+\mathbf{b}_{2}$
        can be written as a single linear threshold. Explain the consequence for further identity
        layers.

  

  \item Compute $\nabla_{\mathbf{W}_{1}}L$ for the network in the worked example if
        $\mathbf{x}=(0.5;\,-1.0;\,0.25;\,2.0)$ and $y=0$, with the same parameters. Determine
        which unit's gate is closed, and why that row has zero gradient without meaning that the
        loss is minimal.

  \item Assess this claim: \emph{because the largest relative error is only
        $8.33\times10^{-10}$, the backpropagation implementation in this chapter is proved
        correct.} Decide true, false, or not enough evidence, and name the additional checks
        required.

  \item Assess this claim: \emph{the universal approximation theorem guarantees that a network
        with one hidden layer can reach an arbitrarily small error on test data.} Decide true,
        false, or not enough evidence, and name two things the theorem does not state.

  \item With $\bar{z}_{2}=-0.2554030842$, compute one Adam step for $b_{2}=0.05$ with
        $\eta=0.01$, $\beta_{1}=0.9$, $\beta_{2}=0.999$, $\epsilon=10^{-8}$, at $t=1$. Show that
        the step size is practically $\eta$, then compute the same step if the gradient is $100$
        times larger.

  \item Show that for two ReLU units in $\mathbb{R}^{2}$ the output is a piecewise linear
        function, and that the input plane is divided into at most four regions with a linear
        formula each. How many regions for $k$ ReLU units in $\mathbb{R}^{2}$ in general
        position? (Hint: each unit gives one line $z_{1}[i]=0$, and inside each cell all gates
        stay fixed.)

  \item Compute the gradient with respect to the input for the worked example, then determine the
        perturbation $\boldsymbol{\delta}$ with $\lVert\boldsymbol{\delta}\rVert_{\infty}\le0.05$
        that raises the loss most according to the linearization, namely
        $\boldsymbol{\delta}=0.05\operatorname{sign}(\nabla_{\mathbf{x}}L)$. Compute the loss at
        $\mathbf{x}+\boldsymbol{\delta}$ and compare it with the linear prediction
        $L+\nabla_{\mathbf{x}}L^{\top}\boldsymbol{\delta}$.

  \item Without running a program, predict what happens if all of $\mathbf{b}_{1}$ is initialized
        to $+10$ while $\mathbf{W}_{1}$ stays small and random, using the saturation table; then
        if $\mathbf{b}_{1}=-10$. Check your prediction by running the second lab block after
        changing one line.
\end{enumerate}
\end{latihan}

Hints for the three hardest items. Item 2: write $\mathbf{W}=\mathbf{W}_{2}\mathbf{W}_{1}$ and
$\mathbf{b}=\mathbf{W}_{2}\mathbf{b}_{1}+\mathbf{b}_{2}$, then compare with the definition of a
linear threshold. Item 7: compute the number of regions for $k=1,2,3$ (they are $2,4,7$), then look for the
pattern $1+k+k(k-1)/2$. Item 9: the second lab block uses ReLU, so $\mathbf{b}_{1}=+10$ saturates nothing; what
saturates is $\mathbf{b}_{1}=-10$, when every unit is dead and its derivative zero. In the
third block, which uses the sigmoid, $z\approx10$ makes the derivative of the hidden layer
about $4.5\times10^{-5}$, so the gradient to $\mathbf{W}_{1}$ shrinks proportionally, while
$\mathbf{b}_{1}$ still receives the full gradient.

\begin{untukpeneliti}
What is really proved in this chapter: the impossibility of linear separation for XOR; the bound
$\lfloor1/\gamma^{2}\rfloor$ on the perceptron, with a complete proof; the derivative of the
sigmoid and the bound $\tfrac14$; the equality $\partial L/\partial z=a-y$ for sigmoid with
cross-entropy, and $\mathbf{p}-\mathbf{y}$ for softmax, both with proofs; the adjoint recursion as
the chain rule; and the proposition on initialization symmetry. What is borrowed without proof:
the universal approximation theorem, with the original proofs in Cybenko (1989) and Hornik
(1991), both using the representation theorem for continuous linear functionals. What is
popularized beyond its proof: ``a network can approximate anything'' as a practical promise, when
it is a statement about the density of functions, not about the reachability of parameters by
optimization; and ``Adam is always better'', when the counterexample is in Reddi, Kale, and Kumar
(2018).

The primary readings that shaped this chapter: Nielsen sections 2.4 and 2.8 (the four
backpropagation equations and why it is cheap) and section 3.1 (cross-entropy, softmax);
Goodfellow et al. section 6.4.1 (universal approximation) and 6.5 (backpropagation). For
initialization, Glorot and Bengio (2010) and He et al. (2015); for Adam, Kingma and Ba (2015).

Limits of validity: all computations in this chapter are for one example on a network with one
hidden layer, with the convention that the ReLU derivative is zero at the break point. There is no
statement about generalization here; that is the content of Chapter 26. Open questions:
why wide networks are easier to train than narrow ones at the same error, why a network with far
more parameters than examples still generalizes, and whether there is an initialization that is
optimal for any depth.
\end{untukpeneliti}

\peta{This chapter follows Nielsen, \emph{Neural Networks and Deep Learning}: sections 1.1
(perceptron), 1.2 (sigmoid neuron), 1.5 (gradient descent), 2.1 to 2.5 (matrix notation, the two
assumptions about the loss function, the Hadamard product, the four backpropagation equations
with their proofs), 2.8 (why backpropagation is fast), 3.1 (cross-entropy, softmax), and 3.3
(weight initialization). For the general formulation and for the side Nielsen does not prove,
Goodfellow, Bengio, and Courville, \emph{Deep Learning} is used: section 6.2 (gradient-based
learning, loss functions, and output units), 6.3 (hidden units and the choice of activation),
6.4.1 (universal approximation and depth), 6.5 (backpropagation and computational graphs), 8.3
(basic algorithms: SGD, momentum), 8.4 (parameter initialization strategies), and 8.5 (algorithms
with adaptive learning rates, including Adam). For the gradient with respect to the input as a
basis for attacks, see the discussion of FGSM in Chapter 27.}

\chapter{Generalization, Regularization, and Training}

A security team builds an intrusion detector from two raw quantities: how many
requests per second and the average payload size. From those two quantities they
derive twenty-one features by multiplying and raising both to the fifth power,
then train a logistic classifier on sixty labeled connections. On the training
data, the detector is wrong on zero of the sixty connections. On twenty thousand
new connections it has never seen, it is wrong on thirty-two percent. The same
model, the same weights, two numbers far apart.

That zero is neither a coincidence nor an achievement. With twenty-one features
and sixty examples, a separating plane can always be found that cuts every
training example correctly, including the example whose label is wrong. What the
model learns is not the shape of the correct separator but the position of those
sixty points. As soon as the points change, the plane is in the wrong place.

Chapter 25 taught us to compute gradients and follow them down until the training
error is small; Chapter 20 taught us to read a single number together with its
uncertainty. Neither answers the question that decides matters here: which
quantity must be made small so that the number in the field shrinks as well, and
how to choose training settings without deceiving ourselves.

\begin{untukanak}
A student is given sixty practice problems before an exam. He studies by
memorizing: every problem together with its answer, down to the letter. On the
practice problems he is always right, a hundred out of a hundred, and he goes
home proud. On the exam the problems are different: different numbers, different
sentences, different order. His grade collapses, and the teacher is puzzled,
because this student looked like the most diligent one.

A second student learns more slowly. He does not memorize answers; he looks for
the pattern: when to divide, when to multiply, when an answer makes sense. On the
practice problems he gets several wrong, and he is not proud. On the exam he does
better.

Now the part that makes one curious. Suppose the first student were given two
thousand practice problems instead of sixty. Is memorizing still bad? It turns
out that it is not always: if there are enough problems and they vary enough,
memorizing becomes an expensive way of learning that works. So the problem is not
memorizing but memorizing from few examples. In this chapter we will learn why
the number of examples and the complexity of the learning method must be
balanced, why adding complexity actually makes the exam grade worse, and how to
measure that balance without waiting for the exam.
\end{untukanak}

\section{Two errors and one gap}
\tingkat{basic}

Every trained model has two numbers, and both are computed in exactly the same
way: the proportion of mistakes on a set of examples. The only difference is
which set is used. In everyday language: the grade on a test at home differs from
the grade on a test at school, even though the problems are of the same kind.

\kamus{generalization}{generalisasi}{the ability to work on new problems, not on problems
that have already been memorized}

\begin{definisi}[Training error, population error, and the gap]
\index{generalization}
Suppose the examples $(x_i,y_i)$ are drawn from an unknown distribution $\mathcal{D}$.
For a hypothesis $h$ and 0-1 error, the \emph{population error} is
$R(h)=\mathbb{P}_{(x,y)\sim\mathcal{D}}\bigl(h(x)\neq y\bigr)$, whereas the
\emph{training error} on $n$ examples is
$\hat R_n(h)=\frac1n\sum_{i=1}^n \mathbf{1}\{h(x_i)\neq y_i\}$.
The difference $R(h)-\hat R_n(h)$ is called the \emph{generalization gap} for that $h$.
\end{definisi}

Training error can be computed; population error cannot, because the distribution
is unknown. All we can do is estimate it on another set that is not used for
training, and that estimate itself wobbles like every proportion in Chapter 20. So
every "validation error" number in this chapter is an estimate, not the truth.

Now two opposite diseases. In everyday language: a model that is too simple cannot
capture the pattern, whereas a model that is too complex memorizes coincidence.

\begin{definisi}[Underfitting and overfitting]
\index{underfitting}
\index{overfitting}
A model \emph{underfits} when its training error and its population error are both
high: the model is too rigid for the actual shape. A model \emph{overfits} when its
training error is low but its generalization gap is large: the model adapts itself
to the coincidence in the training data, not to the actual shape.
\end{definisi}

The two diseases are not symmetric in practice. Underfitting is visible: the
training error is high, and adding features helps. Overfitting deceives, because
the only number that is easily available, namely the training error, is precisely
the one that looks perfect. The model in the chapter opening is the purest
example: its largest weight is $406.204$, a sign that its score depends on small
differences between points, not on a plausible scale.

What governs both is \emph{capacity}: how rich the set of functions is that the
model may choose from. For a linear classifier over $d$ features, capacity rises
with $d$; for a neural network, with the number of units, the depth, and the number
of training steps. Capacity is not merely a parameter count: a large model with
strong regularization can behave like a small model.

\begin{bacakeras}
Two voices, alternating.

``A hundred out of a hundred. That means he understands, doesn't it?''

``Not necessarily. A hundred out of a hundred only means there was no new problem.''

``So memorizing is wrong?''

``Memorizing is not wrong. Memorizing from few examples is.''

``Where is the difference?''

``In the problems that are not in the notes. Memorization answers what has already
been seen. Understanding answers what has not been seen.''

``And if there are two thousand problems instead of sixty?''

``Then memorization starts to help. Not because memorization turns into
understanding, but because the stored coincidences are used again more often.''

Say it once more, slowly: memorization answers what is already there, understanding
answers what is not yet there. Repeat that line until it feels natural in the
mouth: already there, not yet there.
\end{bacakeras}

The generalization gap shrinks as examples are added. In the capacity experiment
with a fixed step budget, the gap is $0.1338$ at $n=60$, $0.0274$ at $n=150$,
$0.0193$ at $n=600$, and $0.0097$ at $n=2400$, whereas the training error barely
moves, staying in the range $0.17$ to $0.21$.

\section{The bias-variance decomposition}
\tingkat{intermediate}

A sharper question: where the gap comes from. In everyday language: a shooter can
miss because the sights are shifted to one side, or because his hands shake, and
the two causes call for different remedies.

\kamus{bias-variance}{bias-variansi}{a systematic miss versus shaking; both contribute
error, and pressing one down sometimes raises the other}

\begin{teorema}[Bias-variance decomposition]
\index{bias-variance}
Suppose $y=f(x)+\varepsilon$ with $\mathbb{E}[\varepsilon]=0$ and
$\operatorname{Var}(\varepsilon)=\sigma^2$, and suppose $\hat f$ is trained on a
random training set $D$ that is independent of $\varepsilon$. For every fixed $x$,
\[
\mathbb{E}_{D,\varepsilon}\bigl[(y-\hat f(x))^2\bigr]
=\sigma^2
+\bigl(\mathbb{E}_D\hat f(x)-f(x)\bigr)^2
+\operatorname{Var}_D\bigl(\hat f(x)\bigr).
\]
\end{teorema}

\begin{proof}[Proof]
Write $y-\hat f(x)=\varepsilon+\bigl(f(x)-\mathbb{E}_D\hat f(x)\bigr)
+\bigl(\mathbb{E}_D\hat f(x)-\hat f(x)\bigr)$, namely three terms: noise that cannot
be reduced, bias, and the deviation from the mean. Square it and take the
expectation. The cross terms containing $\varepsilon$ vanish because
$\mathbb{E}[\varepsilon]=0$ and $\varepsilon$ is independent of $D$; the cross term
between bias and deviation vanishes because
$\mathbb{E}_D[\hat f(x)-\mathbb{E}_D\hat f(x)]=0$. What remains is exactly the three
terms of the theorem's statement.
\end{proof}

The first term cannot be reduced by any model; it belongs to the data, is called
the irreducible error, and on the data of the chapter opening its value is about
$0.21$ because the labels are random. The second and third terms can be traded:
adding capacity lowers bias and raises variance, reducing capacity does the
opposite. Regularization is the name for that trade.

The cheapest example that contains both terms at once is the shrinkage estimator.
To measure $\mu$ from $n$ noisy observations with variance $\sigma^2$, people
usually use $\bar X_n$, which is unbiased. We may use $c\bar X_n$ for $c\in(0,1)$,
and that estimator is deliberately biased. Its mean squared error
\[
\operatorname{MSE}(c)=(c-1)^2\mu^2+c^2\frac{\sigma^2}{n},
\]
is minimized at $c^\ast=\mu^2/(\mu^2+\sigma^2/n)$. For $\mu=2$, $\sigma=1$, and
$n=25$ we obtain $c^\ast=0.990099$: the best estimator is in fact slightly shrunk.

\begin{center}
\small
\begin{tabular}{rrrr}
\hline
$c$ & bias$^2$ & variance & MSE \\
\hline
0.70 & 0.360000 & 0.019600 & 0.379600 \\
0.90 & 0.040000 & 0.032400 & 0.072400 \\
0.99 & 0.000400 & 0.039204 & 0.039604 \\
1.00 & 0.000000 & 0.040000 & 0.040000 \\
1.10 & 0.040000 & 0.048400 & 0.088400 \\
1.30 & 0.360000 & 0.067600 & 0.427600 \\
\hline
\end{tabular}
\end{center}

Three things must be read from that table. The gain from shrinking is small: MSE
falls from $0.040000$ to $0.039604$, about one percent. The loss from deviating in
the wrong direction is far larger: $c=1.3$ gives $0.4276$, whereas $c=0.7$ gives
only $0.3796$. And the minimum is flat, so choosing a slightly wrong $c$ is not a
disaster.

The theorem does not give a formula for an $\hat f$ as complicated as a neural
network. Its decomposition holds formally, but neither term can be computed from
data alone; what is observed is only their sum, namely the validation error. For
that reason the rest of this chapter works with observed numbers, not with the true
bias and variance.
\section{L2 and L1 regularization}
\tingkat{intermediate}

The most direct way to press down variance is to forbid the weights from growing
large. In everyday language: we add a fine to the objective, so that the model has
to pay for every large weight, and will pay only if the gain on the training data
is worth it.

\kamus{regularization}{regularisasi}{a fine that the model pays for large weights, so that
it does not memorize the training data}

\begin{definisi}[The regularized objective]
\index{regularization}
For an original objective $J(\mathbf{w})$ and a penalty $\Omega$, the objective being
minimized is $\tilde J(\mathbf{w})=J(\mathbf{w})+\lambda\Omega(\mathbf{w})$ with
$\lambda\ge0$ as a hyperparameter. For the L2 penalty, $\Omega=\frac12\|\mathbf{w}\|_2^2$;
for the L1 penalty, $\Omega=\|\mathbf{w}\|_1=\sum_j|w_j|$.
\end{definisi}

For L2, the gradient of the objective is $\nabla J+\lambda\mathbf{w}$, so one
gradient-descent step with rate $\eta$ has the form
\[
\mathbf{w}\leftarrow\mathbf{w}-\eta\bigl(\nabla J+\lambda\mathbf{w}\bigr)
=(1-\eta\lambda)\mathbf{w}-\eta\nabla J .
\]
The factor $(1-\eta\lambda)$ is the origin of the name \emph{weight decay}: at every
step, before the gradient is used, the weights are multiplied by a number slightly
less than one. If the gradient is zero, the weights decay toward zero geometrically
with half-life $\ln 2/(\eta\lambda)$ steps.

\begin{akibat}[Weight decay is the same as an L2 penalty]
\index{weight decay}
For a rate $\eta$ and an L2 penalty with coefficient $\lambda$, gradient descent on
the regularized objective is identical to gradient descent on the original
objective followed by multiplying the weights by $1-\eta\lambda$. Thus $\lambda$ and
$\eta$ appear only as the product $\eta\lambda$ in the decay step, so doubling the
rate and halving $\lambda$ give the same decay trajectory.
\end{akibat}

The L2 penalty also has a Bayesian reading: it is equivalent to a Gaussian prior on
the weights, so the solution is a MAP estimate, not a maximum-likelihood one. That
explains why L2 believes small weights are more likely, and why it never forces a
single weight to zero.

L1 behaves differently, and the whole difference lies in the shape of its penalty
near zero. In everyday language: a quadratic penalty pulls ever more weakly once the
weight is small, whereas an absolute-value penalty pulls with constant strength
until the weight is exactly zero.

\begin{lemma}[Soft thresholding]
\index{regularization}
For $a\in\mathbb{R}$ and $t>0$, the minimizer of $u\mapsto\frac12(u-a)^2+t|u|$ is
\[
u^\ast=\operatorname{sign}(a)\max\{|a|-t,0\},
\]
so that $u^\ast=0$ exactly when $|a|\le t$.
\end{lemma}

For $u>0$ the function is $\frac12(u-a)^2+tu$ with derivative $u-a+t$, so its
critical point is $u=a-t$, valid only when $a>t$. For $u<0$ the derivative is
$u-a-t$ with critical point $u=a+t$, valid only when $a<-t$. At $u=0$ the
subgradient is the interval $[-a-t,-a+t]$, which contains zero exactly when
$|a|\le t$; in that case zero is the minimizer because the function is convex. The
three cases join into the formula above.

That lemma explains the whole behavior of L1. With a quadratic approximation around
the unpenalized solution $w^\ast$ and curvature $H$, the objective is approximately
$\frac H2(w-w^\ast)^2+\lambda|w|$, so that
$w=\operatorname{sign}(w^\ast)\max\{|w^\ast|-\lambda/H,0\}$: every weight below
$\lambda/H$ becomes exactly zero, not merely small. For L2 the solution is
$w=w^\ast/(1+\lambda/H)$, which is never zero unless $w^\ast$ itself is zero.

\begin{proposisi}[Why L1 produces true zeros]
\index{regularization}
Suppose $J$ is differentiable and convex around zero. The point $\mathbf{w}=\mathbf{0}$
is a minimizer of $J(\mathbf{w})+\lambda\|\mathbf{w}\|_1$ if and only if
$|\partial J/\partial w_j|\le\lambda$ for every $j$. For the L2 penalty, the same
condition is $\partial J/\partial w_j=0$ for every $j$.
\end{proposisi}

\begin{proof}[Proof]
The subgradient of $\|\cdot\|_1$ at zero is the box $[-1,1]^d$, so the condition
$0\in\nabla J(\mathbf{0})+\lambda[-1,1]^d$ is equivalent to $|(\nabla J)_j|\le\lambda$
for all $j$; for a convex function, that condition is necessary and sufficient. For
the L2 penalty, the gradient is $\nabla J+\lambda\mathbf{w}$, which at $\mathbf{0}$
equals $\nabla J$, so the condition is $\nabla J(\mathbf{0})=\mathbf{0}$.
\end{proof}

The practical consequence matters for reading results. L1 gives a flat region of
width $2\lambda$ around zero, so that weights that are "almost useless" are switched
off entirely; L2 only shrinks them. For that reason L1 is used as a feature
selector, and for the same reason the number of zero weights under L1 is not a
delicate measure of importance: it jumps, it does not shift slowly.

Now the numbers. The same logistic classifier, the same data, twenty-one features
on sixty training examples, everything computed with the script in the lab section.

\begin{center}
\small
\begin{tabular}{rrrr}
\hline
$\lambda$ & $\max_j|w_j|$ & training error & validation error \\
\hline
0       & 406.204 & 0.0000 & 0.3262 \\
0.001   &   3.056 & 0.1000 & 0.2927 \\
0.003   &   1.814 & 0.1000 & 0.2688 \\
0.010   &   1.018 & 0.1167 & 0.2475 \\
0.030   &   0.631 & 0.1167 & 0.2433 \\
0.100   &   0.399 & 0.1500 & 0.2406 \\
0.300   &   0.268 & 0.2333 & 0.2452 \\
1.000   &   0.142 & 0.2500 & 0.2695 \\
3.000   &   0.062 & 0.3333 & 0.3205 \\
10.000  &   0.022 & 0.3333 & 0.3395 \\
100.000 &   0.003 & 0.3333 & 0.3395 \\
\hline
\end{tabular}
\end{center}

\begin{figure}[htbp]
\centering
\begin{tikzpicture}
\begin{axis}[
  width=12.2cm, height=5.6cm,
  xmin=-0.4, xmax=10.4, ymin=0, ymax=0.40,
  axis lines=left,
  xtick={0,3,5,7,10},
  xticklabels={$0$,$10^{-2}$,$0.1$,$1$,$100$},
  ytick={0,0.1,0.2,0.3,0.4},
  yticklabels={$0$,$0.10$,$0.20$,$0.30$,$0.40$},
  tick label style={font=\scriptsize},
  xlabel={regularization strength $\lambda$ (capacity decreases to the right)},
  ylabel={error},
  xlabel style={font=\small},
  ylabel style={font=\small},
  legend style={font=\scriptsize,draw=none,at={(0.03,0.97)},anchor=north west},
]
\addplot[very thick,layerindigo,mark=*,mark size=1.5pt] coordinates {
  (0,0.0000) (1,0.1000) (2,0.1000) (3,0.1167) (4,0.1167) (5,0.1500)
  (6,0.2333) (7,0.2500) (8,0.3333) (9,0.3333) (10,0.3333)};
\addlegendentry{training error}
\addplot[very thick,accent,mark=square*,mark size=1.5pt] coordinates {
  (0,0.3262) (1,0.2927) (2,0.2688) (3,0.2475) (4,0.2433) (5,0.2406)
  (6,0.2452) (7,0.2695) (8,0.3205) (9,0.3395) (10,0.3395)};
\addlegendentry{validation error}
\addplot[only marks,mark=o,mark size=3.4pt,accent] coordinates {(5,0.2406)};
\draw[densely dashed,layergrey] (axis cs:5,0)--(axis cs:5,0.2406);
\node[accent,font=\scriptsize,anchor=north] at (axis cs:4.55,0.238)
  {minimum $0.2406$};
\end{axis}
\end{tikzpicture}
\caption{Training error and validation error against the regularization strength
$\lambda$; the direction to the right means capacity decreases, because the weights
are held back ever more tightly. The eleven points are the rows of the previous
table. The training error rises monotonically from $0.0000$ to $0.3333$, whereas the
validation error forms a valley: $0.3262$ at $\lambda=0$, a minimum of $0.2406$ at
$\lambda=0.1$, then $0.3395$ at $\lambda=10$ and above. The bottom of the valley is
wide, from $0.2433$ at $\lambda=0.03$ to $0.2452$ at $\lambda=0.3$.}
\label{fig:lembah-kapasitas}
\end{figure}

Figure~\ref{fig:lembah-kapasitas} plots the third and fourth columns of that table
as a single valley. Read that table as three regions, not eleven rows. Without
regularization the training error is zero and the validation error is $0.3262$, with
weights up to the hundreds because nothing holds them back. In the region of
$\lambda$ between $0.001$ and $0.1$, the training error rises to $0.1500$ while the
validation error falls to $0.2406$. In the region $\lambda\ge3$ both are bad:
$0.3333$ and $0.3395$.

The last row deserves emphasis. With $\lambda=10$ the validation error is $0.3395$,
\emph{worse} than the $0.3262$ of the model without regularization. So regularization
is not a medicine that always helps; it helps at a certain dose and does damage
beyond it. At $\lambda=10$ the model practically always guesses the majority class,
whose error is $0.3419$: its weights are so small that the score is almost entirely
determined by the bias.

The L1 penalty gives a similar picture with a different shape. At $\lambda=0.05$ its
validation error is $0.2334$ with eleven of the twenty-one weights alive; at
$\lambda=0.2$ only one weight remains, with error $0.3395$. The count of live weights
does not fall monotonically: $\lambda=0.005$ leaves eleven, $\lambda=0.01$ leaves
fourteen. That number is stable against a fourfold increase in iterations, so it is
not an optimizer failure; the L1 trajectory for the logistic loss really does not
guarantee a monotone count of zeros, and that is one more reason not to read it as a
measure of importance.

The final comparison: L1 reaches $0.2334$ and L2 $0.2406$. The difference $0.0072$ is
smaller than the wobble between partitions, about $0.02$, so this data is not enough
to conclude that L1 is better; all that can be concluded is that both improve on the
model without regularization.
\begin{untukanak}
Imagine a class that has to guess a friend's body weight. The teacher gives two
methods. Method one: every guess that is too far from the answer is fined an amount
that grows the farther away it is. Method two: every guess that is not exactly the
answer is fined a fixed amount, no matter whether the difference is small or large.

Method one makes all guesses shift slightly toward the middle, but none of them
becomes zero. Method two rounds guesses that are already close straight down to
zero, because the fine is fixed and no longer worth the improvement that might be
gained. That is the difference between two fines that look alike: one shrinks, the
other switches off.

The question: if we have twenty-one features and only sixty examples, what kind of
fine makes the model stop memorizing? The answer is not "a large fine", because a
fine that is too large switches everything off, including what is correct. In this
chapter we will learn why there is a dose, why that dose must be searched for, and
why a dose that is searched for on the wrong data will deceive us neatly.
\end{untukanak}

\section{Early stopping, dropout, and augmentation}
\tingkat{intermediate}

Three other ways of pressing down variance without touching the formula of the
objective. In everyday language: we can stop learning before coincidence begins to
be memorized, we can switch off part of the memory at random, and we can multiply
examples by changing the examples we already have.

Early stopping is the cheapest form of regularization, because it adds not a single
line to the objective: it only stops the trajectory. During training, the training
error almost always falls, but the validation error falls and then rises. The lowest
point of the validation error is the one taken, and the number of steps becomes a
hyperparameter to be searched for.

\begin{center}
\small
\begin{tabular}{rrrr}
\hline
epoch & training error & validation error & $\max_j|w_j|$ \\
\hline
1    & 0.1333 & 0.2481 & 0.15 \\
5    & 0.1333 & 0.2429 & 0.38 \\
10   & 0.1333 & 0.2432 & 0.48 \\
50   & 0.1167 & 0.2454 & 0.77 \\
200  & 0.1167 & 0.2541 & 1.38 \\
800  & 0.1000 & 0.2879 & 2.85 \\
2400 & 0.0833 & 0.3109 & 4.81 \\
\hline
\end{tabular}
\end{center}

\begin{figure}[htbp]
\centering
\begin{tikzpicture}
\begin{axis}[
  width=11.8cm, height=5.6cm,
  xmode=log,
  xmin=1, xmax=3000, ymin=0.05, ymax=0.35,
  axis lines=left,
  xtick={1,10,100,1000},
  xticklabels={$1$,$10$,$100$,$1000$},
  ytick={0.10,0.15,0.20,0.25,0.30},
  yticklabels={$0.10$,$0.15$,$0.20$,$0.25$,$0.30$},
  tick label style={font=\scriptsize},
  xlabel={training epoch (log scale)},
  ylabel={error},
  xlabel style={font=\small},
  ylabel style={font=\small},
  legend style={font=\scriptsize,draw=none,at={(0.03,0.97)},anchor=north west},
]
\addplot[very thick,layerteal,mark=*,mark size=1.5pt] coordinates {
  (1,0.1333) (5,0.1333) (10,0.1333) (50,0.1167) (200,0.1167) (800,0.1000)
  (2400,0.0833)};
\addlegendentry{training error}
\addplot[very thick,accent,mark=square*,mark size=1.5pt] coordinates {
  (1,0.2481) (5,0.2429) (10,0.2432) (50,0.2454) (200,0.2541) (800,0.2879)
  (2400,0.3109)};
\addlegendentry{validation error}
\draw[densely dashed,layergrey] (axis cs:8,0.05)--(axis cs:8,0.2427);
\addplot[only marks,mark=o,mark size=3.4pt,accent] coordinates {(8,0.2427)};
\node[accent,font=\scriptsize,anchor=north west] at (axis cs:8.6,0.2325)
  {epoch 8: $0.2427$};
\end{axis}
\end{tikzpicture}
\caption{Both errors along one training trajectory, horizontal axis on a log scale.
The training error falls without ever rising, from $0.1333$ at epoch 1 to $0.0833$
at epoch 2400. The validation error falls to $0.2427$ at epoch 8 and then rises to
$0.3109$. The dashed line marks the best stopping point: stopping there gives an
error $0.0682$ lower than training all the way to epoch 2400, and the curves begin to
separate just after that point.}
\label{fig:dua-kurva}
\end{figure}

Figure~\ref{fig:dua-kurva} draws those two columns. The training error falls along
the trajectory, from $0.1333$ to $0.0833$, without ever rising, whereas the
validation error falls until epoch 8 with a best value of $0.2427$ and then rises to
$0.3109$ at epoch 2400. The last column explains why: what grows throughout training
is the size of the weights, and that is the variance we are paying for.

Early stopping is not merely a tactic. For a quadratic objective and a small rate,
stopping at step $\tau$ is nearly equivalent to an L2 penalty with coefficient about
$1/(\tau\eta)$; that is an approximation, not an exact identity. For that reason the
number of steps is a capacity hyperparameter, not a matter of "training until
finished".

\begin{resep}
\begin{enumerate}
  \item Split three sets before a single training step is run: train, validation,
        test. The test set is sealed.
  \item Train on the training set only; every few epochs compute the error on the
        validation set, and record the three values (epoch, validation error,
        training error).
  \item Take the epoch with the smallest validation error, then return the weights to
        that point; do not use the weights of the last epoch.
  \item Only now open the test set, compute its error once, and report it together
        with the chosen epoch and the width of the validation curve.
\end{enumerate}
A warning about order: step 1 must be finished before step 2, and step 4 must not be
used to shift the choice made in step 3. If the test set is looked at while the epoch
is being chosen, it is no longer a test set but a second validation set, and its
numbers stop being an honest estimate.
\end{resep}

Dropout takes another route: during training every unit is switched off at random
with probability $1-p$, and the ones that survive are divided by $p$ so that the
expectation does not shift. In everyday language: at every step some members of the
team do not show up, so that no member can be relied on alone.

\begin{lemma}[Expectation and variance of the dropout mask]
\index{dropout}
Suppose $m=B/p$ with $B\sim\mathrm{Bernoulli}(p)$ and $0<p\le1$. Then
$\mathbb{E}[m]=1$ and $\operatorname{Var}(m)=(1-p)/p$. For a unit with inputs
$z_1,\dots,z_k$ and independent masks $m_1,\dots,m_k$, its output
$s=\sum_j m_jw_jz_j$ satisfies $\mathbb{E}[s]=\sum_jw_jz_j$ and
$\operatorname{Var}(s)=\frac{1-p}{p}\sum_j(w_jz_j)^2$.
\end{lemma}

\begin{proof}[Proof]
$\mathbb{E}[m]=p/p=1$ and $\operatorname{Var}(m)=\mathbb{E}[m^2]-1=p/p^2-1=(1-p)/p$.
For $s$, linearity of expectation gives the first term, and independence of the masks
gives the variance of the sum, namely $\sum_j(w_jz_j)^2\operatorname{Var}(m_j)$.
\end{proof}

That variance formula explains two things that are often forgotten. Dropout noise is
proportional to the sum of squared contributions, so units with large weights receive
relatively more noise, and dropout points in the same direction as an L2 penalty. And
that noise disappears entirely at test time, so that the model that is used differs
from the model that was trained; on very small networks that price can exceed the
benefit.

Data augmentation multiplies examples by changing the examples that already exist,
for instance shifting them or adding small noise, provided that the label does not
change. For a linear model, its relation to L2 can be computed exactly.

\begin{proposisi}[Gaussian noise on the input is the same as an L2 penalty]
\index{data augmentation}
Suppose $\hat y=\mathbf{w}^{\top}\mathbf{x}$ and $\boldsymbol{\delta}\sim
N(\mathbf{0},\sigma^2I)$ is added to the input. Then
\[
\mathbb{E}_{\boldsymbol{\delta}}\bigl[(\mathbf{w}^{\top}(\mathbf{x}+\boldsymbol{\delta})-y)^2\bigr]
=(\mathbf{w}^{\top}\mathbf{x}-y)^2+\sigma^2\|\mathbf{w}\|_2^2 .
\]
\end{proposisi}

\begin{proof}[Proof]
Expand: the linear term is $2(\mathbf{w}^{\top}\mathbf{x}-y)\,
\mathbb{E}[\mathbf{w}^{\top}\boldsymbol{\delta}]=0$, and the quadratic term is
$\mathbb{E}[(\mathbf{w}^{\top}\boldsymbol{\delta})^2]
=\mathbf{w}^{\top}\mathbb{E}[\boldsymbol{\delta}\boldsymbol{\delta}^{\top}]\mathbf{w}
=\sigma^2\|\mathbf{w}\|_2^2$.
\end{proof}

So adding small Gaussian noise to the input is the same as adding an L2 penalty with
$\lambda=\sigma^2$: the two are one idea written twice. The numerical check, for
$\mathbf{w}=(1,2)$, $\mathbf{x}=(1,0)$, $y=3$, and $\sigma=0.1$: the original
objective is $4$, the formula predicts $4+0.01\cdot5=4.05$, and a simulation of one
million noises gives $4.049107$. That difference of $0.0009$ is comparable to the
random error of the simulation, not to a mistake in the formula.

A warning for the security field: augmentation is valid only when the label is
genuinely invariant to the change that is made. Adding small noise to traffic
features is allowed, but reversing the order of packets or shifting arrival times can
turn an attack into not-an-attack. Augmentation with a wrong label is not
regularization but a lie that will keep looking good as long as the validation data
is changed in the same way.
\section{SGD, momentum, and Adam}
\tingkat{advanced}

Up to here we have treated "training" as a single box. The contents of that box
determine how fast the error falls and, on a limited step budget, how low an error is
reached. In everyday language: the right direction is not enough; what decides is how
the step is taken and how large it is.

Stochastic gradient descent replaces the full gradient by the gradient over one
minibatch of size $b$. That estimate is unbiased, and its variance shrinks
proportionally to $1/b$: $\operatorname{Var}(\hat{\mathbf{g}}_b)=\sigma_g^2/b$. So a
large minibatch reduces the wobble, but does not change the mean direction. What
changes the direction is momentum and Adam.

\begin{definisi}[Momentum and Adam]
\index{momentum}
\index{adam}
Momentum keeps a running average of the gradient:
$\mathbf{v}_t=\beta\mathbf{v}_{t-1}+\mathbf{g}_t$ and
$\mathbf{w}_t=\mathbf{w}_{t-1}-\eta\mathbf{v}_t$, with $0\le\beta<1$. Adam keeps two
moments, $\mathbf{m}_t=\beta_1\mathbf{m}_{t-1}+(1-\beta_1)\mathbf{g}_t$ and
$\mathbf{v}_t=\beta_2\mathbf{v}_{t-1}+(1-\beta_2)\mathbf{g}_t^2$ (squared per
coordinate), then uses bias correction and the update
$\mathbf{w}_t=\mathbf{w}_{t-1}-\eta\,\hat{\mathbf{m}}_t/(\sqrt{\hat{\mathbf{v}}_t}+\epsilon)$
with $\hat{\mathbf{m}}_t=\mathbf{m}_t/(1-\beta_1^t)$ and
$\hat{\mathbf{v}}_t=\mathbf{v}_t/(1-\beta_2^t)$.
\end{definisi}

That bias correction is not decoration: without it $\mathbf{m}_1=(1-\beta_1)\mathbf{g}_1$
is very small, so the first step is almost zero. Division by $\sqrt{\hat{\mathbf{v}}}$
matters even more for badly scaled problems: it makes the step size in each
coordinate roughly $\eta$, no matter whether the gradient is large or small. Momentum
only smooths the direction over time.

The numbers on the following quadratic show the difference. Take
$f(\mathbf{w})=\frac12\mathbf{w}^{\top}H\mathbf{w}$ with $H=\operatorname{diag}(1,10^4)$
and $\mathbf{w}_0=(1,0)$; the condition number of $H$ is $10^4$, which means the curve
is narrow in one direction and very flat in the other. We count steps until
$f\le10^{-8}f(\mathbf{w}_0)$.

\begin{center}
\small
\begin{tabular}{lrl}
\hline
method & rate & steps to $10^{-8}$ \\
\hline
gradient descent & $10^{-4}$ & 92099 \\
momentum ($\beta=0.9$) & $10^{-4}$ & 9132 \\
Adam & $0.1$ & 82 \\
\hline
\end{tabular}
\end{center}

Gradient descent needs $92099$ steps because its largest rate is limited by the
largest curvature, while the flat direction moves only at that rate. Momentum cuts
one digit, in line with the estimate that acceleration changes the dependence on the
condition number $\kappa$ into about $\sqrt\kappa$. Adam cuts three digits, because
per-coordinate division means the flat direction is no longer punished by the steep
one. That is not a general promise: on well-scaled problems the advantage of Adam
shrinks.

On the logistic problem of this chapter, with minibatch eight and a target training
error of $0.40$, the median of five random seeds gives SGD with rate $0.05$ a hundred
steps, whereas Adam with rate $0.1$ needs only ten. The conclusion must be careful:
the advantage is about one digit and partly comes from the larger rate, not from a
better direction.

The learning rate is usually not constant. Two common schedules are step decay,
namely the rate divided by ten at certain epochs, and cosine decay
$\eta_t=\eta_{\min}+\frac12(\eta_0-\eta_{\min})(1+\cos(\pi t/T))$: large steps at the
start so as to approach quickly, small steps at the end so as not to jump around in
the valley. With $3000$ minibatch steps and $\lambda=0.1$, a constant rate of $0.2$
gives a median training error of $0.3866$, whereas a cosine from $0.5$ to $0.01$ gives
$0.3549$.

\section{Cross validation, hyperparameter selection, and leakage}
\tingkat{advanced}

All the validation numbers in this chapter come from a set that is not used for
training. That is already correct, but it is not enough: as soon as we choose
$\lambda$ with an eye that sees the validation error, that validation error stops
being an unbiased estimate of the population error. In everyday language: if we
choose the exam questions by the grade we get, that grade no longer measures our
ability.

\kamus{cross validation}{validasi silang}{rotating the roles of the data in turn, so that
every example has once been an exam for a model that did not see it}

\begin{definisi}[$k$-fold cross validation]
\index{cross validation}
The training data is split into $k$ folds that are disjoint and of similar size. For
each $j$, the model is trained on $k-1$ folds and tested on fold $j$; the
cross-validation error is the average error over those $k$ folds, and the
hyperparameter that is chosen is the one that minimizes it. For $k=n$ this procedure
is called \emph{leave-one-out}.
\end{definisi}

Cross validation answers one weakness of a single split: at $n=60$ the validation set
contains twelve examples, and the error on twelve examples wobbles with a standard
deviation of about $0.12$. With five folds every example is tested once, but the
wobble does not disappear.

\begin{center}
\small
\begin{tabular}{rrr}
\hline
$\lambda$ & CV error (average of 5 partitions) & spread between partitions \\
\hline
0     & 0.3233 & 0.1245 \\
0.003 & 0.2400 & 0.0253 \\
0.030 & 0.2467 & 0.0217 \\
0.100 & 0.2467 & 0.0274 \\
1     & 0.2967 & 0.0298 \\
\hline
\end{tabular}
\end{center}

The first row shows two things at once: without regularization the CV error is both
the worst and the least stable, with a spread of $0.1245$, four times that of the
other rows. The second to fourth rows are almost indistinguishable: $0.2400$,
$0.2467$, and $0.2467$, with a spread of about $0.025$, and across five different
partitions the winner keeps changing. This CV curve is flat around its minimum, so
that choosing its exact value is an illusion of precision.

The consequence is real: CV selects $\lambda=0.003$ with validation error $0.2688$,
whereas the unselected $\lambda=0.1$ gives $0.2406$; that difference of $0.0282$ is
not far from one standard deviation of the CV curve. The lesson is not that cross
validation is broken, but that with sixty examples, hyperparameter selection carries
an uncertainty as large as the effect being measured. Reporting "the best $\lambda$
is $0.003$" without mentioning the width of the curve is the subtlest lie in
this chapter.

\begin{tangan}
Materials: ten cards, a pencil, paper. Write one two-digit number on each card,
shuffle, then lay out seven cards face up. Build a "rule" in the form of a list:
position one is worth this much, position two is worth that much, through position
seven.

Cover one card with your palm and guess. Because your list contains that card, your
guess is exact. Repeat for all seven positions: seven out of seven correct, and that
is your training score.

Now gather the ten cards, shuffle, lay out seven new cards, and use the same list.
Cover one card, guess, uncover. What must be visible: your list does not contain the
new cards, so your guess is pure luck. Repeat ten times with new layouts and record
how often you are exact; that is your exam score. Finally write a second rule that
uses only one property, for example "the card on the left is not larger than the card
on its right", then compare the two scores.
\end{tangan}

There is a source of error larger than that wobble and invisible in any number:
leakage, namely information from the test data entering the training process. Its
most common form is standardizing features using the mean and standard deviation of
all the data before the split is made.

The numbers, with an artificial covariate shift: if validation is shifted by $0.5$,
standardization with training statistics only gives $0.3765$, whereas pooled
statistics give $0.3361$; at a shift of $1.0$ the numbers are $0.5387$ against
$0.3819$. Pooled statistics always look better, and that is precisely the problem: they
use information about where the test data lies that is not available in the field.
Without a shift the two are equal, $0.2452$ against $0.2463$, a sign that the earlier
difference came from leakage.

The honest protocol is layered: separate train, validation, and test from the start;
compute every quantity, including means, standard deviations, decision thresholds,
and feature selection, inside the training folds only; choose hyperparameters on the
validation set or through cross validation inside the training data; touch the test
set once, then report it as it is.

The relation between capacity and the amount of data can be made quantitative. For a
finite hypothesis class, the generalization gap is controlled by the square root of
the logarithm of the number of hypotheses divided by the number of examples.

\begin{teorema}[A uniform bound for finite classes]
\index{capacity}
Suppose $\mathcal{H}$ is finite with $|\mathcal{H}|=m$, its error is 0-1, and $n$
examples are drawn independently from $\mathcal{D}$. With probability at least
$1-\delta$ over the drawing of the examples, it holds simultaneously for all
$h\in\mathcal{H}$
\[
\bigl|R(h)-\hat R_n(h)\bigr|\le\sqrt{\frac{\ln m+\ln(2/\delta)}{2n}} .
\]
\end{teorema}

\begin{proof}[Proof]
For a fixed $h$, $\hat R_n(h)$ is the average of $n$ Bernoulli variables scaled to
$[0,1]$, so Hoeffding's inequality gives
$\mathbb{P}(|R(h)-\hat R_n(h)|>t)\le2e^{-2nt^2}$. Combine over all $h$ with the union
bound: the probability that some $h$ exceeds $t$ is at most
$2m\,e^{-2nt^2}$. Setting that expression equal to $\delta$ gives
$t=\sqrt{(\ln m+\ln(2/\delta))/(2n)}$, and its complement is the statement of the
theorem.
\end{proof}

For the class that "memorizes", namely a model that gives one label to every binary
pattern over $d$ features, $m=2^d$. With $d=21$ and $\delta=0.05$ the bound is
$\sqrt{(21\ln2+\ln40)/(2n)}$, namely $0.3899$ at $n=60$ and $0.0617$ at
$n=2400$. The measured gap is $0.1338$ and $0.0097$ at the same $n$, so the bound
holds but is loose; that looseness is not a defect, because the bound must hold for
every distribution and every algorithm.

What matters is not the numbers but their shape. The gap shrinks proportionally to
$1/\sqrt n$, whereas the number of hypotheses enters through a logarithm. Quadrupling
the data only halves the gap, and increasing capacity tenfold only raises the bound by
the square root of the logarithmic increment. This is the honest reason why large
models are still useful: as long as their gap is controlled by regularization, rising
capacity does not automatically mean a rising gap.
\section{Worked example: closed-form ridge on four points}
\tingkat{intermediate}

All the formulas above can be checked by hand on a small example: four points with
one feature, $\mathbf{x}=(1,2,3,4)$ and $\mathbf{y}=(2,4,6,9)$, a model without an
intercept $\hat y=wx$, and the objective $\frac12\sum_i(y_i-wx_i)^2+\frac{\lambda}{2}w^2$.
Setting its derivative to zero gives
\[
w(\lambda)=\frac{\sum_ix_iy_i}{\sum_ix_i^2+\lambda}=\frac{64}{30+\lambda},
\]
since $\sum x_i^2=30$ and $\sum x_iy_i=2+8+18+36=64$.

\begin{center}
\small
\begin{tabular}{rrrr}
\hline
$\lambda$ & $w$ & training SSE & $|w\cdot5-10|$ \\
\hline
0   & 2.133333 & 0.466667 & 0.666667 \\
1   & 2.064516 & 0.608741 & 0.322581 \\
2   & 2.000000 & 1.000000 & 0.000000 \\
5   & 1.828571 & 3.253061 & 0.857143 \\
10  & 1.600000 & 9.000000 & 2.000000 \\
30  & 1.066667 & 34.600000 & 4.666667 \\
100 & 0.492308 & 81.255621 & 7.538462 \\
\hline
\end{tabular}
\end{center}

Figure~\ref{fig:lintasan-bobot} plots both weight trajectories against $\lambda$ on
this data: L2 shrinks without end, L1 stops exactly at zero. The last column uses one
test point at $x=5$ with true value $10$, which does not take part in determining
$w$. That is where the curve takes a U shape and its minimum falls exactly at
$\lambda=2$: there $w=2$, so that the prediction at $x=5$ is $10$ and the error is
zero. Without regularization, $w=64/30=2.1333$ is more attractive on the training data
(SSE $0.466667$ against $1.000000$) but predicts $10.667$, off by $0.667$. This is the
bias-variance trade in its barest form: the training SSE doubles, the test error
becomes zero.

\begin{figure}[htbp]
\centering
\begin{tikzpicture}
\begin{axis}[
  width=11.8cm, height=5.4cm,
  xmode=log,
  xmin=0.5, xmax=300, ymin=-0.12, ymax=2.45,
  axis lines=left,
  xtick={1,2,10,64,100},
  xticklabels={$1$,$2$,$10$,$64$,$100$},
  ytick={0,0.5,1.0,1.5,2.0},
  yticklabels={$0$,$0.5$,$1.0$,$1.5$,$2.0$},
  tick label style={font=\scriptsize},
  xlabel={regularization strength $\lambda$ (log scale)},
  ylabel={weight $w$},
  xlabel style={font=\small},
  ylabel style={font=\small},
  legend style={font=\scriptsize,draw=none,at={(0.97,0.97)},anchor=north east},
]
\addplot[domain=0.5:300,samples=200,very thick,layerindigo] {64/(30+x)};
\addlegendentry{L2: $64/(30+\lambda)$}
\addplot[domain=0.5:63.9,samples=200,very thick,layerorange] {(64-x)/30};
\addlegendentry{L1: $\max\{(64-\lambda)/30,\,0\}$}
\addplot[domain=64:300,samples=2,very thick,layerorange] {0};
\draw[densely dashed,layergrey] (axis cs:2,-0.12)--(axis cs:2,2.20);
\draw[densely dashed,layergrey] (axis cs:64,-0.12)--(axis cs:64,2.20);
\addplot[only marks,mark=*,mark size=1.8pt,layerindigo] coordinates {(2,2.0000) (64,0.6809)};
\addplot[only marks,mark=square*,mark size=1.8pt,layerorange] coordinates {(2,2.0667)};
\addplot[only marks,mark=o,mark size=3.4pt,layerorange] coordinates {(64,0)};
\node[font=\scriptsize,anchor=south] at (axis cs:2,2.24) {$\lambda=2$};
\node[font=\scriptsize,anchor=south] at (axis cs:64,2.24) {$\lambda=64$};
\node[font=\scriptsize,anchor=south west] at (axis cs:70,0.06) {L1 exactly zero};
\end{axis}
\end{tikzpicture}
\caption{Weight trajectories against regularization strength on the worked example
with four points, with $\sum x_i^2=30$ and $\sum x_iy_i=64$. The L2 penalty gives
$w=64/(30+\lambda)$, which shrinks but never touches zero. The L1 penalty gives
$w=\max\{(64-\lambda)/30,\,0\}$, which reaches zero exactly at $\lambda=64$. At
$\lambda=2$ the two are nearly the same, $2.0000$ against $2.0667$; at $\lambda=64$ L2
is still $0.6809$ whereas L1 is already $0.0000$. The dashed lines mark the two values
of $\lambda$ that are compared.}
\label{fig:lintasan-bobot}
\end{figure}

Note the right-hand wall. At $\lambda=100$, $w$ is only $0.4923$, the training SSE
explodes to $81.255621$, and the test error is $7.538462$, worse than a model that
learns nothing at all. This example also shows why searching for $\lambda$ on the
training data is misleading: the training SSE curve rises monotonically, so that it
always chooses $\lambda=0$.

\begin{tantangan}
A model reports a training error of $0.0000$ on $60$ examples and a test error of
$0.4100$. A colleague proposes three changes: (i) train longer, (ii) add features up
to degree seven, (iii) raise $\lambda$ from $0$ to $0.1$. For each proposal, predict
the sign of the change in training error and in test error, then state which one
cannot possibly help, with the reason.

Hint 1: first ask in which direction capacity moves under each proposal.

Hint 2: two of the three proposals enlarge capacity, and this model is already at zero
training error, so that enlarging capacity only widens the gap. What helps is what
narrows it, and that must raise the training error.

If you get stuck, go back to the Definition of underfitting and overfitting in the
first section.
\end{tantangan}

\kertas{Check the second column without a calculator: $64/32=2$ exactly, $64/31=2.064516$,
and $64/40=1.6$. Then compute the SSE at $\lambda=2$ yourself: the residuals are
$(0,0,0,1)$, so SSE $=1$. Finally answer: which $\lambda$ makes $w=1.28$, and what is
the prediction at $x=5$ then? Hint: solve $64/(30+\lambda)=1.28$.}
\section{Computational lab}
\tingkat{advanced}

Every number in this chapter comes from a single numpy script with no machine-learning
library. The data is deliberately small so that it can be rechecked within minutes: two
raw features, twenty-one polynomial features up to degree five, sixty training examples,
and twenty thousand validation examples from the same distribution. The labels are not
deterministic, so that there is an irreducible error of about $0.21$ which no model can
go below.

\begin{lab}
\begin{verbatim}
import numpy as np

def data(n, rng):
    x = rng.uniform(-1, 1, size=(n, 2))
    p = 1/(1 + np.exp(-4.0*(0.9 - x[:,0]**2 - x[:,1]**2)))
    return x, (rng.random(n) < p).astype(float)

def poly(x, deg=5):
    return np.column_stack([x[:,0]**i * x[:,1]**j
                            for i in range(deg+1) for j in range(deg+1-i)])

def sig(z):
    return 1/(1 + np.exp(-np.clip(z, -400, 400)))

def soft(t, l):
    return np.sign(t)*np.maximum(np.abs(t) - l, 0.0)

def er(Z, y, w, b):
    return float(np.mean((sig(Z@w + b) > 0.5) != y))

def fit(Z, y, lam, iters=30000, lr=0.05):
    n, d = Z.shape
    w = np.zeros(d); b = 0.0
    mw = np.zeros(d); vw = np.zeros(d); mb = 0.0; vb = 0.0
    for t in range(1, iters+1):
        p = sig(Z@w + b); g = (p - y)/n
        gw = Z.T@g + lam*w; gb = g.sum()
        mw = 0.9*mw + 0.1*gw; vw = 0.999*vw + 0.001*gw*gw
        mb = 0.9*mb + 0.1*gb; vb = 0.999*vb + 0.001*gb*gb
        w -= lr*(mw/(1-0.9**t))/(np.sqrt(vw/(1-0.999**t)) + 1e-8)
        b -= lr*(mb/(1-0.9**t))/(np.sqrt(vb/(1-0.999**t)) + 1e-8)
    return w, b

def fit_l1(Z, y, lam, iters=150000):
    n, d = Z.shape
    w = np.zeros(d); b = 0.0
    st = 0.5/(0.25*np.linalg.eigvalsh(Z.T@Z/n).max())
    for _ in range(iters):
        p = sig(Z@w + b); g = (p - y)/n
        w = soft(w - st*(Z.T@g + lam*w), st*lam); b -= st*g.sum()
    return w, b
\end{verbatim}
\end{lab}

The first part computes the error for eleven values of $\lambda$; the second part repeats
that with the L1 penalty through proximal gradient descent, namely a gradient step
followed by the soft threshold from the lemma above; the third part stops early; the
fourth and fifth parts compare the optimizers; the sixth to eighth parts examine cross
validation, leakage, and capacity. The actual output, cut only where it repeats:

\begin{verbatim}
d=21  n_train=60  n_validation=20000
==1 L2: lam  |w|max  train  validation
  0.0000 406.204  0.0000  0.3262
  0.0010   3.056  0.1000  0.2927
  0.0030   1.814  0.1000  0.2688
  0.0100   1.018  0.1167  0.2475
  0.0300   0.631  0.1167  0.2433
  0.1000   0.399  0.1500  0.2406
  0.3000   0.268  0.2333  0.2452
  1.0000   0.142  0.2500  0.2695
  3.0000   0.062  0.3333  0.3205
 10.0000   0.022  0.3333  0.3395
100.0000   0.003  0.3333  0.3395
==2 L1: lam  zeros/total  train  validation
  0.0020 16/21  0.1167  0.2630
  0.0050 11/21  0.1167  0.2455
  0.0100 14/21  0.1167  0.2412
  0.0200 14/21  0.1500  0.2369
  0.0500 11/21  0.1833  0.2334
  0.1000  7/21  0.2667  0.2492
  0.2000  1/21  0.3333  0.3395
==3 early stopping: epoch  train  validation  |w|max
ep=   1  0.1333  0.2481  0.15
ep=   5  0.1333  0.2429  0.38
ep=  10  0.1333  0.2432  0.48
ep=  50  0.1167  0.2454  0.77
ep= 200  0.1167  0.2541  1.38
ep= 800  0.1000  0.2879  2.85
ep=2400  0.0833  0.3109  4.81
best 0.2427 at epoch 8
==4 quadratic diag(1,1e4): steps to error 1e-8
GD=1e-4  92099 | momentum  9132 | Adam=0.1  82
==5 logistic minibatch: steps to training error 0.40 (median of 5 seeds)
sgd  lr=0.05 [100, 100, 90, 90, 100] median=100
sgd  lr=0.2  [30, 50, 40, 40, 20] median=40
adam lr=0.03 [30, 20, 20, 20, 30] median=20
adam lr=0.1  [20, 10, 20, 10, 10] median=10
==6 CV 5 partitions x 5 folds: lam  CV error (mean +- sd between partitions)
 0.000  0.3233 +- 0.1245
 0.003  0.2400 +- 0.0253
 0.030  0.2467 +- 0.0217
 0.100  0.2467 +- 0.0274
 1.000  0.2967 +- 0.0298
==7 standardization: shift  training stats  pooled stats
0.0  0.2452  0.2463
0.5  0.3765  0.3361
1.0  0.5387  0.3819
==8 capacity vs data (d=21, lam=0, nested sets)
n=   60 train=0.1667 validation=0.3005 diff=+0.1338 bound=0.3899
n=  150 train=0.2133 validation=0.2407 diff=+0.0274 bound=0.2466
n=  600 train=0.2067 validation=0.2260 diff=+0.0193 bound=0.1233
n= 2400 train=0.2058 validation=0.2155 diff=+0.0097 bound=0.0617
\end{verbatim}

Three things need to be read from that output, and all three disappoint ordinary
expectations.

First, the training error is nearly useless as a guide: it moves in the opposite
direction to the validation error. The validation curve is also not symmetric. From
$\lambda=0$ to $0.1$ its error falls by $0.086$; from $0.1$ to $10$ it rises by $0.099$
until it passes its starting point. When in doubt between two values, the smaller one
is safer.

Second, selection through cross validation is not automatically correct. CV chooses
$\lambda=0.003$, whereas validation chooses $0.1$; the difference between their CV errors
is $0.0006$, far below the spread between partitions of $0.025$, even though on the
validation data the two differ by $0.0282$.

Third, the theoretical bound is loose but its direction is right. The capacity
experiment here trains the model with a budget of sixty thousand steps on nested sets,
so that its training error is not zero. At $n=60$ the bound is $0.3899$ whereas the gap
is $0.1338$; at $n=2400$ the bound is $0.0617$ whereas the gap is $0.0097$. The ratio is
about three, so that the bound is useful for comparing settings, not for predicting
numbers.

\begin{untukdosen}
The teaching order I suggest: start from the opening numbers of the chapter, sixty
examples with training error zero and validation error $0.3262$, then ask whether that
model "learned". The answer of the class usually splits in two, and that split is the
best lecture material there is. After that the formal definitions, and let the
bias-variance decomposition be derived on the board in five lines. L1 regularization is
better introduced through the soft-threshold lemma than through a picture of the L1
ball, because students can check that lemma themselves.

The two misconceptions that appear most often. Students treat a low training error as
proof of a good model, then are surprised when its validation error is high; another
symptom is that they suggest "training longer" as the cure, when in the early-stopping
table that is precisely what does the damage. And they treat regularization as always
helpful, then install a large $\lambda$ with the conviction that it is "safer". The row
$\lambda=10$ with error $0.3395$ is the cure for both.

A blackboard activity that works: ask the class to guess, before it is computed, where
the minimum of the validation curve lies, then compute just three rows together; the
guess usually misses toward a $\lambda$ that is too small. The sign that the class is
ready to go on: they can explain why choosing $\lambda$ with the validation data makes
that number optimistic, and they can name one leakage that is invisible in the numbers.
\end{untukdosen}

\section{Summary}
\tingkat{basic}

\begin{ringkasan}
\begin{itemize}
  \item Training error and population error are computed in the same way on different
        sets; a training error of $0.0000$ with a validation error of $0.3262$ signals
        memorization.
  \item The population error decomposes into irreducible noise, squared bias, and
        variance; regularization trades bias for variance, and that trade pays only as
        long as the fall in variance is larger than the rise in bias.
  \item The L2 penalty adds $\lambda\mathbf{w}$ to the gradient, equivalent to
        multiplying the weights by $1-\eta\lambda$ at every step; L1 uses soft
        thresholding so that weights below $\lambda/H$ become exactly zero. Both have a
        dose: $\lambda=0.1$ gives $0.2406$, whereas $\lambda=10$ gives $0.3395$, worse
        than without regularization.
  \item Early stopping is regularization through the number of steps: $0.2427$ at
        epoch 8 against $0.3109$ at epoch 2400. Dropout adds noise with variance
        $(1-p)/p$ per unit, and Gaussian-noise augmentation with variance $\sigma^2$ is
        equivalent to an L2 penalty of scale $\sigma^2$.
  \item Optimizers differ on badly scaled problems: on a quadratic with condition
        number $10^4$, gradient descent needs $92099$ steps, momentum $9132$, Adam
        $82$; a cosine schedule lowers the training error from $0.3866$ to $0.3549$.
  \item Cross validation reduces the wobble without removing it: the CV curve is flat
        with a spread of $0.025$, and its choice misses by $0.0282$. Standardization
        leakage makes the numbers look better ($0.3361$ against $0.3765$) without being
        really better, and the generalization gap shrinks proportionally to $1/\sqrt n$:
        from $0.1338$ at $n=60$ to $0.0097$ at $n=2400$.
\end{itemize}
\end{ringkasan}

\section{Exercises}
\tingkat{intermediate}

\begin{latihan}
\begin{enumerate}
  \item Assess this claim: "Because the training error of our model is zero, the model
        is already good enough." True, false, or not enough evidence? Name the part that
        is mistaken and one additional number that you need.
  \item Prove the bias-variance decomposition for $y=f(x)+\varepsilon$ with
        $\operatorname{Var}(\varepsilon)=\sigma^2$: write the three terms, show that both
        cross terms vanish, and state where the independence of $D$ from $\varepsilon$
        is used.
  \item Derive $w(\lambda)=64/(30+\lambda)$ for the data in the worked example, then show
        that $w$ decreases monotonically in $\lambda$ but is never zero. Compare it with
        the L1 penalty on the same one-dimensional problem, and determine the $\lambda$
        that makes $w=0$ exactly.
  \item Compute the validation error for $c=0.95$ and $c=1.05$ in the shrinkage example
        with $\mu=2$, $\sigma=1$, $n=25$, then compare them with $c=1$. Explain why two
        deviations of equal size do not give losses of equal size.
  \item A report writes: "We chose the best $\lambda$ through cross validation, so that
        cross-validation number is an estimate of the error in the field." Fix that
        sentence, and name two additional procedures that are needed.
  \item Prove the soft-threshold lemma for $a<0$ in full, then explain why the result is
        symmetric in the sign of $a$. Hint: treat the cases $u<0$ and $u=0$ separately.
  \item Compute the number of weight-decay steps needed for a weight to become a quarter
        of its initial value if $\eta=0.1$ and $\lambda=0.02$, using the half-life
        $\ln2/(\eta\lambda)$. Check with a numpy simulation for a zero gradient.
  \item On the data of this chapter, $\lambda=0.1$ gives a training error of $0.1500$ and
        a validation error of $0.2406$. Compute the gap, compare it with the gap at
        $\lambda=0$ and $\lambda=10$, then explain why a small gap is not enough to
        conclude that a model is better.
  \item Derive the uniform bound $\sqrt{(\ln m+\ln(2/\delta))/(2n)}$ from Hoeffding's
        inequality and the union bound, then compute it for $d=10$, $n=200$, and
        $\delta=0.01$. How many times must the data be doubled for the bound to halve?
\end{enumerate}

Hints for the three hardest items. Item 3: its one-dimensional objective is
$\frac12\sum_i(y_i-wx_i)^2+\lambda|w|$; at $w=0$ its subgradient contains zero when
$\lambda\ge|\sum_ix_iy_i|=64$, so that $w=0$ exactly at $\lambda=64$. Item 6: at $u=0$
the subgradient is the interval $[-a-t,\,-a+t]$, which contains zero when $|a|\le t$.
Item 9: with $m=2^{10}$ and $\delta=0.01$ the bound is
$\sqrt{(10\ln2+\ln200)/400}=0.1749$; halving it demands four times the data.
\end{latihan}

\begin{untukpeneliti}
What is proved in this chapter: the bias-variance decomposition together with its
conditions; the equivalence of the weight-decay step with an L2 penalty; the closed form
of the minimizer of $\frac H2(w-w^\ast)^2+\lambda|w|$ and the subgradient condition that
makes L1 produce true zeros; the equivalence of Gaussian-noise augmentation with an L2
penalty for linear models; the formulas for the expectation and variance of the dropout
mask; and the uniform bound for finite hypothesis classes. What is only popular: that
regularization always helps; that Adam is always faster; that early stopping "is the
same as" L2 (it is an approximation for a quadratic objective with a small rate); that
the count of zeros under L1 measures feature importance.

Primary sources: Goodfellow, Bengio, and Courville, \emph{Deep Learning}, Chapter 5
(capacity, bias-variance), Chapter 7 (L1 and L2 penalties, augmentation, noise, early
stopping, dropout), Chapter 8 (SGD, momentum, Adam, learning-rate schedules,
hyperparameter selection). Wasserman, \emph{All of Statistics}, Chapter 13 and Chapter 20.
The uniform bound is a direct consequence of Hoeffding's inequality and the union bound.

Limits of validity: every number comes from one synthetic set with a single $60/20000$
split, so that each validation error carries an uncertainty of order $0.01$ to $0.02$.
For that reason the comparison of L1 and L2 here is not a recommendation. The uniform
bound is only for finite classes and 0-1 error.

Open questions: how large the leakage of batch normalization is, since it computes
statistics from the minibatch and so mixes information between examples; how to choose
$\lambda$ when the cost of mistakes is asymmetric; and whether a flat CV curve can be
turned into a valid interval, not merely a spread between partitions.
\end{untukpeneliti}

\peta{Chapter 26 in Goodfellow, Bengio, and Courville, \emph{Deep Learning}: Chapter 5
(capacity, overfitting and underfitting, the bias-variance decomposition), Chapter 7 (L1
and L2 penalties, weight decay, data augmentation, noise as regularization, early
stopping, dropout), Chapter 8 (SGD, momentum, Adam, learning-rate schedules,
hyperparameter selection). In Wasserman, \emph{All of Statistics}: Chapter 13 (model
selection and cross validation), Chapter 20 (bias-variance in smoothing). Backpropagation
and network training are covered in Chapter 23; honest reading of evaluation metrics,
including leakage and data splitting, continues in Chapter 25; the convex optimization
underlying penalties and subgradients is in Chapter 21.}

\chapter{Honest Model Evaluation}

A short report lands on a security analyst's desk. It holds one table: a new intrusion detector,
tested on network traffic captured from a campus network over a month, reaches an accuracy of
$99$ percent. That number is framed in bold type and closed by a single sentence:
``ready to deploy''.

The data contains one attack for every thousand connections. On the same collection, a three-line
program that always answers ``clean'' obtains an accuracy of $99.9$ percent: it reads no feature at
all and never raises a false accusation. The detector trained for hours loses, and its report
contains not one number that leaks that defeat.

What is wrong is the summary: a single number stands in for four numbers with different roles, on
data whose composition is never stated. Chapter 26 explained why test error differs from training
error; this chapter closes three gaps: which metrics are honest about what they ignore, how much
those numbers wobble, and whether the test data is truly unseen by the procedure that chose the
model.

\begin{untukanak}
Imagine a metal detector at a school gate. Every morning it beeps. Yesterday it beeped a thousand
times, and the guard searched a thousand bags. Only sixty of those bags really carried a forbidden
object. The other nine hundred and forty students had to open their bags for no reason.

Now an inspector from the district office arrives and asks: what percentage of the device's answers
were right? The guard answers with pride: ninety-nine percent. He is not lying; out of a hundred
thousand bags, the device was wrong on only a thousand. But that answer does not answer the
question that upsets the students: they are upset not because the device is often wrong, but
because almost every time it beeps, it is wrong.

The two questions sound alike and their answers are far apart: one asks about all the bags, the
other about the bags that were beeped. In this chapter we will learn why one device can excel at
the first question and collapse on the second, why a number of a hundred percent is sometimes not
good news, and why an honest report always carries more than one number.
\end{untukanak}

\section{Four numbers that are not symmetric}
\tingkat{basic}

A binary classifier, however sophisticated its interior, leaves the same trace: four numbers. What
separates one evaluation report from another is whether all four are reported. In everyday
language: every decision has two ways to be right and two ways to be wrong, and closing off one of
them is the same as hiding.

\kamus{confusion matrix}{matriks konfusi}{a four-column ledger that records every decision, right
or wrong, together with the direction of the error}

\begin{definisi}[Confusion matrix]
\index{confusion matrix}
Suppose every example has a true label $y\in\{0,1\}$ and a prediction
$\hat y\in\{0,1\}$, with $1$ meaning positive (attack, disease, harmful message).
The four numbers
\begin{align}
\mathrm{TP}&=\#\{y=1,\ \hat y=1\}, & \mathrm{FP}&=\#\{y=0,\ \hat y=1\},
\nonumber\\
\mathrm{FN}&=\#\{y=1,\ \hat y=0\}, & \mathrm{TN}&=\#\{y=0,\ \hat y=0\}
\nonumber
\end{align}
together with their arrangement are called the \emph{confusion matrix}. The total number of
examples is $n=\mathrm{TP}+\mathrm{FP}+\mathrm{FN}+\mathrm{TN}$.
\end{definisi}

From those four numbers we derive metrics that answer different questions. Precision asks: when the
model says positive, how likely is it to be right. Recall asks: of all the positive cases, what
fraction was caught.

\begin{definisi}[Precision, recall, specificity]
\index{precision}
\index{recall}
With $\mathrm{TP}+\mathrm{FP}>0$ and $\mathrm{TP}+\mathrm{FN}>0$,
\[
\operatorname{prec}=\frac{\mathrm{TP}}{\mathrm{TP}+\mathrm{FP}},\qquad
\operatorname{rec}=\frac{\mathrm{TP}}{\mathrm{TP}+\mathrm{FN}},\qquad
\operatorname{spec}=\frac{\mathrm{TN}}{\mathrm{TN}+\mathrm{FP}} .
\]
Accuracy is $(\mathrm{TP}+\mathrm{TN})/n$, and $F_1$ is
$2\,\mathrm{TP}/(2\,\mathrm{TP}+\mathrm{FP}+\mathrm{FN})$.
\end{definisi}

Let us use one worked example: a detector is tested on $1000$ labeled connections, $100$ of them
attacks. It flags $62$ attacks correctly, misses $38$, and accuses $17$ clean connections of being
attacks.

\begin{center}
\small
\begin{tabular}{l|rr|r}
\hline
 & predicted positive & predicted negative & total \\
\hline
actually positive & $\mathrm{TP}=62$ & $\mathrm{FN}=38$ & 100 \\
actually negative & $\mathrm{FP}=17$ & $\mathrm{TN}=883$ & 900 \\
\hline
total & 79 & 921 & 1000 \\
\hline
\end{tabular}
\end{center}

Its precision is $62/79=0.784810$: about four out of five alarms point to an attack. Its recall is
$62/100=0.620000$: almost four out of ten attacks go unseen. Its specificity is
$883/900=0.981111$. Its accuracy is $945/1000=0.945000$, and $F_1=124/179=0.692737$. All five
numbers are true about the same model, and not one of them may be used on its own. A check with
\texttt{sklearn.metrics} gives the same five numbers to six decimal places.

Figure~\ref{fig:konfusi-jujur} puts the four numbers in their places, together with the row and
column totals; every number in the rest of the chapter must be traceable to one of the cells of
that figure.

\begin{figure}[htbp]
\centering
\begin{tikzpicture}[every node/.style={font=\footnotesize}]
  % two correct cells (green) and two wrong cells (red)
  \fill[layerjade!14] (4.40,2.00) rectangle (6.70,2.90);
  \fill[layerjade!14] (6.70,1.10) rectangle (9.00,2.00);
  \fill[accent!8]     (6.70,2.00) rectangle (9.00,2.90);
  \fill[accent!8]     (4.40,1.10) rectangle (6.70,2.00);
  \draw[ink,thick] (4.40,0.55) rectangle (10.30,2.90);
  \draw[ink] (6.70,0.55) -- (6.70,2.90);
  \draw[ink] (9.00,0.55) -- (9.00,2.90);
  \draw[ink] (4.40,2.00) -- (10.30,2.00);
  \draw[ink] (4.40,1.10) -- (10.30,1.10);
  \draw[ink!35] (4.40,0.55) -- (10.30,0.55);
  % cell contents
  \node at (5.55,2.45) {$\mathrm{TP}=62$};
  \node at (7.85,2.45) {$\mathrm{FN}=38$};
  \node at (5.55,1.55) {$\mathrm{FP}=17$};
  \node at (7.85,1.55) {$\mathrm{TN}=883$};
  % column headers
  \node at (5.55,3.14) {predicted positive};
  \node at (7.85,3.14) {predicted negative};
  \node at (9.65,3.14) {total};
  % row headers
  \node[anchor=east] at (4.24,2.45) {actually positive $(100)$};
  \node[anchor=east] at (4.24,1.55) {actually negative $(900)$};
  \node[anchor=east] at (4.24,0.82) {total};
  % totals
  \node at (5.55,0.82) {$79$};
  \node at (7.85,0.82) {$921$};
  \node at (9.65,2.45) {$100$};
  \node at (9.65,1.55) {$900$};
  \node at (9.65,0.82) {$1000$};
  \node[anchor=north,font=\scriptsize,align=center,text width=11.4cm] at (7.00,0.30)
    {accuracy $=(62+883)/1000=0.945$;\quad precision $=62/79$;\quad recall $=62/100$};
\end{tikzpicture}
\caption{The confusion matrix of the worked example on $1000$ connections: $62$ attacks caught,
$38$ missed, $17$ false alarms, and $883$ clean connections correctly rejected. The rows sum to
$100$ positive and $900$ negative, the columns to $79$ predicted positive and $921$ predicted
negative, and everything to $1000$. The two green cells are the two ways to be right, the two red
cells the two ways to be wrong. Accuracy $0.945$ adds the two green cells and therefore cannot
distinguish the $62$ that were caught from the $38$ that were missed.}
\label{fig:konfusi-jujur}
\end{figure}

The roles of the four numbers are not symmetric. Precision answers the operator's question: when
the alarm sounds, how likely is a real attack. Recall answers the policy designer's question: of
all attacks, what fraction is caught. Both depend on the threshold, so a metric without a
threshold is not a complete statement: the same model gives precision $0.25$ with recall $0.45$,
or precision $0.013$ with recall $0.86$.

$F_1$ tries to compress precision and recall into a single number, and that is exactly where it is
easiest to misuse.

\begin{proposisi}
\index{f1}
$F_1$ is the harmonic mean of precision $P$ and recall $R$, that is $F_1=2PR/(P+R)$,
so that $F_1\le(P+R)/2$ with equality holding exactly when $P=R$.
\end{proposisi}

\begin{proof}[Proof]
Write $a=\mathrm{TP}$, $b=\mathrm{FP}$, $c=\mathrm{FN}$. Then $P=a/(a+b)$,
$R=a/(a+c)$, and
\[
\frac{2PR}{P+R}
=\frac{2a^2\big/\bigl((a+b)(a+c)\bigr)}{a(2a+b+c)\big/\bigl((a+b)(a+c)\bigr)}
=\frac{2a}{2a+b+c}=F_1 .
\]
The harmonic--arithmetic inequality $4PR\le(P+R)^2$ is equivalent to $(P-R)^2\ge0$, and equality
occurs only when $P=R$.
\end{proof}

The harmonic property pulls the value down toward the smaller component: precision $0.9$ with
recall $0.2$ gives $F_1\approx0.327$, not $0.55$; in our example $F_1=0.6927$, slightly below the
midpoint $0.7024$.

Three defects of the summary must be named one by one, because the remedies differ. Precision is
blind to attacks that slip through: a model that flags every connection obtains precision $0.1$,
namely its own prevalence, but recall $1$. $F_1$ ignores $\mathrm{TN}$ entirely: two models with
the same precision and recall but ten times as many false alarms report an identical $F_1$, even
though their operating costs are far apart. $F_1$ also weighs both components as equal, even though
one missed attack is usually far more expensive than one false alarm.

\begin{tantangan}
A report carries only three numbers: ``accuracy $0.9720$ on $5000$ test connections; precision
$0.000$; model ready to deploy''. Determine the whole confusion matrix that is consistent with
that report, explain why its recall is zero and not some other number, then find the largest
prevalence that is still possible.

Hint 1: accuracy only adds the two diagonal cells, so $5000(1-0.9720)=140$ errors are free to be
split between the two remaining cells.

Hint 2: precision $0.000$ means $\mathrm{TP}=0$, so $\mathrm{FN}$ equals the number of positive
examples and $\mathrm{FP}$ is the rest. But precision with $\mathrm{TP}=\mathrm{FP}=0$ is $0/0$,
not $0$; the figure $0.000$ in that report therefore demands at least one false alarm.

If you get stuck, go back to the section Deceptive accuracy on unbalanced data in this chapter.
\end{tantangan}

\section{Deceptive accuracy on unbalanced data}
\tingkat{basic}

Accuracy is the most often reported and the easiest to deceive: it counts all errors at the same
price on a collection whose composition is never stated. In everyday language: accuracy is a report
card computed as if the two groups were balanced.

\index{balanced accuracy}
\begin{definisi}[Balanced accuracy]
Let $\pi$ be the prevalence of the positive class. \emph{Balanced accuracy} is
$(\operatorname{rec}+\operatorname{spec})/2$, that is the average of the correct rates on the two
classes. This quantity does not change when the numbers of positive and negative examples are
changed without changing those two rates.
\end{definisi}

Ordinary accuracy has no such property, and the consequence can be written exactly.

\begin{proposisi}
A classifier that always answers negative obtains accuracy $1-\pi$. Hence a model whose accuracy
is below $1-\pi$ loses to a baseline that learns nothing, and a report that states an accuracy
without stating $\pi$ cannot be judged.
\end{proposisi}

\begin{proof}[Proof]
Of $n$ examples, $\pi n$ are labeled positive and all of them are wrong, while $(1-\pi)n$ are
labeled negative and all of them are right. The accuracy is $(1-\pi)n/n=1-\pi$. Any other model
with smaller accuracy therefore makes more errors than that baseline.
\end{proof}

\begin{bacakeras}
Two voices, read alternately with a friend. Tap the table every time the word ``zero'' is spoken.

``A hundred cards. How many red?'' ``Ten.'' ``How many white?'' ``Ninety.'' ``What does the machine
answer?'' ``White. All white.'' ``How many does it get right?'' ``Ninety.'' ``So its accuracy?''
``Ninety percent.'' ``And the red ones?'' ``Zero. Not a single one.'' ``Was that number lying?''
``No. It answered only one question.'' ``Which question?'' ``How many of everything there is did I
guess right.'' ``Which question went unanswered?'' ``Of the truly dangerous ones, how many did I
find.''

Repeat the three middle lines, faster and faster, with the same taps: what is counted is
everything, what is lost goes unmentioned, what is lost goes unmentioned.

The second voice closes softly: the small class drowns easily inside the average of the large
class, and that average never names who drowned.

The line the reader must repeat, once, in a clear voice: accuracy adds up what is right, not what
matters, so on an unbalanced class ninety percent can mean zero.
\end{bacakeras}

In our worked example $\pi=0.1$, so the baseline is $0.9$ and the model with $0.945$ really is
better. Now change only the prevalence, without changing the quality of the model, and watch what
happens to the same number.

\begin{contoh}[One attack for every thousand connections]
Take $100\,000$ connections with $100$ attacks. Scores on clean connections are normally
distributed with mean $0$ and standard deviation $1$; scores on attacks with mean $2.6$ and
standard deviation $1$. The threshold is set at $2.3263$, the $0.99$ quantile of the negative score
distribution, so by design about one percent of clean connections will be flagged. One run with a
fixed seed gives:
\[
\mathrm{TP}=62,\quad \mathrm{FP}=1013,\quad \mathrm{FN}=38,\quad \mathrm{TN}=98\,887 .
\]
Its accuracy is $(62+98\,887)/100\,000=0.989490$. The baseline is
$99\,900/100\,000=0.999000$. Its precision is $62/1075=0.057674$, its recall
$62/100=0.620000$, and its balanced accuracy $0.804930$.
\end{contoh}

That number $0.9895$, rounded, becomes ``$99$ percent''. It is not merely uninformative: it is
lower than the baseline by $0.009510$, nearly a full percentage point. The model makes ten times
as many errors as the program that always answers ``clean'', $1051$ against $100$, and all of them
are false alarms, $1013$ of them, or about $16.3$ per attack caught.

Here rounding adds damage: $0.9895$ and $0.9990$ both become ``$99$ percent'', even though they
are separated by the whole difference that matters. A good report therefore writes three or four
decimal places, states $n$, and states the prevalence; without all three, the reader does not know
whether $99$ percent is an achievement or a failure.

Should accuracy then not be used at all? It is legitimate when the two errors cost the same, when
the prevalence is near one half, and when the threshold is fixed. Outside that it must be
accompanied by precision and recall, or replaced by balanced accuracy: $0.8049$ is more honest than
$0.9895$, because the minority class is caught only $62$ percent of the time.

\begin{tangan}
Materials: $100$ cards, $90$ white-faced and $10$ red-faced, a marker, and one sheet of tally paper
with two columns.

Write ``clean'' on the ninety white cards and ``attack'' on the ten red ones. Shuffle, then deal
them into ten piles. First rule: every pile is answered ``clean'' without turning over a single
card. Record the correct answers, then turn all the cards over: ninety right, ten red ones through.
Write on the board: accuracy ninety percent, recall zero percent.

Now the second rule. Put a dot in the corner of six red cards, and without telling the player, put
the same dot in the corner of four white cards. Answer ``attack'' only when that dot is visible.
Count again: six red caught, four red through, four white accused, eighty-six white right.

What must become visible: accuracy moves only from ninety to ninety-two percent, while recall jumps
from zero to sixty percent. The number that changes most is precisely the number missing from the
first report.
\end{tangan}

\section{The ROC curve, AUC, and prevalence}
\tingkat{intermediate}

Every threshold gives one point on two different rates; choosing a threshold without looking at all
the possibilities is guessing. In everyday language: we sweep the threshold from loose to strict,
then draw the trail the model leaves behind.

\kamus{receiver operating characteristic}{kurva roc}{the trail a model leaves when its threshold is
swept from loose to strict, drawn on two rates that do not depend on prevalence}

\begin{definisi}[ROC curve and AUC]
\index{roc curve}
\index{auc}
For a threshold $t$, the true positive rate and the false positive rate are
\[
\mathrm{TPR}(t)=P(S>t\mid y=1),\qquad \mathrm{FPR}(t)=P(S>t\mid y=0),
\]
with $S$ the model score. The ROC curve is the set of points $(\mathrm{FPR}(t),\mathrm{TPR}(t))$
for all $t$. The area under that curve is called the AUC, and for finite data it is computed as the
trapezoid area between the points actually observed.
\end{definisi}

The AUC has a far simpler meaning than its name, and that meaning is what made it popular.

\begin{proposisi}
For data without tied scores, the AUC equals the probability that the score of one randomly drawn
positive example exceeds the score of one randomly drawn negative example. When there are tied
scores, each tied pair counts as one half.
\end{proposisi}

\begin{proof}[Proof]
Sort all scores from large to small and sweep the threshold downward. Each time the threshold
passes one positive example, $\mathrm{TPR}$ rises by $1/n_+$ at the height of $\mathrm{FPR}$ at
that moment, namely the fraction of negative examples already passed. Each time the threshold
passes a negative example, $\mathrm{FPR}$ rises at constant $\mathrm{TPR}$ and adds no area. The
total area is therefore
\[
\mathrm{AUC}=\frac1{n_+}\sum_{i:\ y_i=1}\mathrm{FPR}_i
=\frac{1}{n_+n_-}\#\{(i,j):\ y_i=1,\ y_j=0,\ S_i>S_j\},
\]
which is the probability claimed, and that last term is what the Mann--Whitney statistic computes.
Tied scores add no area, so each such pair contributes one half.
\end{proof}

The AUC does not see prevalence: if every negative example is duplicated ten times, the ordering of
scores does not change and the AUC stays. How to compute it on paper is shown in the Worked
examples section.

Now the part that makes the ROC curve dangerous: sweep the threshold in a one-in-a-thousand
simulation.

\begin{center}
\small
\begin{tabular}{rrrrr}
\hline
threshold & recall & FPR & precision & accuracy \\
\hline
3.0000 & 0.450 & 0.00132 & 0.254 & 0.99813 \\
2.8000 & 0.510 & 0.00263 & 0.162 & 0.99688 \\
2.6000 & 0.550 & 0.00467 & 0.105 & 0.99488 \\
2.3263 & 0.620 & 0.01014 & 0.058 & 0.98949 \\
2.0000 & 0.740 & 0.02269 & 0.032 & 0.97707 \\
1.5000 & 0.860 & 0.06729 & 0.013 & 0.93264 \\
\hline
\end{tabular}
\end{center}

On the ROC curve those rows move from $(0.001;0.45)$ to $(0.067;0.86)$: a trail that rises steeply
near the vertical axis, exactly what every practitioner wants to see. Its AUC is $0.966448$. But
the precision column says something else. Even at the strictest threshold, where only $132$ clean
connections are flagged, the precision is $0.254$: three out of four alarms are still wrong. When
recall is raised to $0.86$, the precision is $0.013$.

Figure~\ref{fig:presisi-recall-ambang} draws the six rows of that table on one threshold axis:
recall climbs while precision falls, and no point gives both at once.

\begin{figure}[htbp]
\centering
\begin{tikzpicture}
\begin{axis}[
  width=0.78\linewidth, height=6.2cm,
  x dir=reverse, xmin=1.45, xmax=3.05,
  ymin=0.35, ymax=0.95, ytick={0.4,0.5,0.6,0.7,0.8,0.9},
  xlabel={threshold $t$ (to the right: the threshold is loosened)}, ylabel={recall},
  axis y line*=left, axis x line*=bottom,
  tick label style={font=\footnotesize}, label style={font=\footnotesize},
  grid=major, grid style={layergrey!15},
]
\addplot[accent,thick,mark=*,mark size=1.8pt] coordinates {
  (3.0,0.450) (2.8,0.510) (2.6,0.550) (2.3263,0.620) (2.0,0.740) (1.5,0.860)};
\node[accent,font=\footnotesize,anchor=south west] at (axis cs:2.62,0.665) {recall};
\end{axis}
\begin{axis}[
  width=0.78\linewidth, height=6.2cm,
  x dir=reverse, xmin=1.45, xmax=3.05,
  ymin=0, ymax=0.30, ytick={0,0.1,0.2,0.3},
  axis y line*=right, axis x line*=none, ylabel={precision (right scale)},
  tick label style={font=\footnotesize}, label style={font=\footnotesize},
]
\addplot[layerteal,thick,mark=square*,mark size=1.6pt] coordinates {
  (3.0,0.254) (2.8,0.162) (2.6,0.105) (2.3263,0.058) (2.0,0.032) (1.5,0.013)};
\node[layerteal,font=\footnotesize,anchor=north east] at (axis cs:2.78,0.232) {precision};
\end{axis}
\end{tikzpicture}
\caption{Precision and recall at the six thresholds of the previous table, drawn on one threshold
axis that is turned so that to the right the threshold is loosened. Recall climbs from $0.450$ at
$t=3.0$ to $0.860$ at $t=1.5$, while precision falls from $0.254$ to $0.013$: the two curves move in
opposite directions and there is not one point that gives both at once. The precision scale is on
the right, $0$ to $0.30$; at $t=2.3263$, the threshold used in the simulation, precision $0.058$ with
recall $0.620$.}
\label{fig:presisi-recall-ambang}
\end{figure}

The two readings do not contradict each other: they measure different things.

\begin{proposisi}
The ROC curve and the AUC depend only on the conditional score distributions $S\mid y=1$ and
$S\mid y=0$, so they do not change when the prevalence $\pi$ is changed. Precision at a fixed
threshold does change, and follows
\[
\operatorname{prec}(t)=\frac{\pi\,\mathrm{TPR}(t)}
{\pi\,\mathrm{TPR}(t)+(1-\pi)\,\mathrm{FPR}(t)} .
\]
\end{proposisi}

\begin{proof}[Proof]
By definition, $\mathrm{TPR}$ and $\mathrm{FPR}$ are probabilities conditional on $y$, so neither
contains $\pi$; a curve assembled from the two contains no $\pi$ either. For precision, the law of
total probability gives $P(\hat y=1)=\pi\,\mathrm{TPR}+(1-\pi)\mathrm{FPR}$, while
$P(y=1,\hat y=1)=\pi\,\mathrm{TPR}$. Bayes' rule from Chapter 17 gives the form above.
\end{proof}

The consequence can be measured. If prevalence is changed to one half without changing the
conditional score distributions, the AUC stays at $0.967061$ (the difference $0.0006$ is only
Monte Carlo error), while \emph{average precision} jumps from $0.287636$ to $0.967544$. At a
prevalence of one in a thousand the baseline is $\pi=0.001$, so $0.288$ is still far better than
guessing, but far closer to operational reality than $0.966$.

The rule of thumb is simple: when the positive class is rare, report a precision--recall curve, not
only the ROC curve, and state the prevalence. The ROC curve remains useful for comparing two models
on the same population because it is not polluted by composition; it misleads as soon as it is used
to conclude usefulness on another population or at a particular threshold. The AUC is also blind to
calibration: a model whose scores are twice too large has exactly the same AUC, because the
ordering does not change. Those two defects carry us to the next section.

\section{Calibration: does the number mean a probability}
\tingkat{intermediate}

An operator who receives an alarm needs to know one thing that neither precision nor the AUC
answers: when the model writes $0.8$, are eight of ten such cases really positive. In everyday
language: calibration checks whether the weather forecast that says ``seventy percent rain'' is
really right on seven of ten such days.

\kamus{calibration}{kalibrasi}{the agreement between the number a model promises and what actually
happens in the group sharing that same promise}

\begin{definisi}[Perfect calibration]
\index{calibration}
A score $S\in[0,1]$ is called perfectly calibrated when $P(y=1\mid S=s)=s$ for every possible $s$.
On finite data, calibration is checked by grouping scores into intervals, then comparing the mean
score in that interval with the observed positive rate. The absolute difference, weighted by the
size of the interval, is called the \emph{expected calibration error}.
\end{definisi}

A worked example with $1000$ predictions shows the shape of the check.

\begin{center}
\small
\begin{tabular}{lrrrr}
\hline
score interval & count & positive & mean score & positive rate \\
\hline
$[0.5;0.6)$ & 40 & 21 & 0.55 & 0.525 \\
$[0.6;0.7)$ & 60 & 36 & 0.65 & 0.600 \\
$[0.7;0.8)$ & 120 & 90 & 0.75 & 0.750 \\
$[0.8;0.9)$ & 200 & 176 & 0.85 & 0.880 \\
$[0.9;1.0]$ & 580 & 551 & 0.95 & 0.950 \\
\hline
\end{tabular}
\end{center}

The largest difference is in the second interval, $0.05$, and the overall calibration error is
\[
\tfrac{40}{1000}(0.025)+\tfrac{60}{1000}(0.050)+\tfrac{120}{1000}(0)
+\tfrac{200}{1000}(0.030)+\tfrac{580}{1000}(0)=0.010 .
\]
The direction of the error differs between intervals, but its size is small: one percentage point.
The cross-check is free: the counts in the five intervals add to $1000$, and the positives to $874$.

Calibration is measured by squared error, and there is a neat identity there.

\begin{proposisi}[Brier for calibrated scores]
\index{brier}
Let $S$ be perfectly calibrated and $B=\mathbb{E}[(S-y)^2]$ be the Brier score. Then
\[
B=\mathbb{E}\bigl[S(1-S)\bigr].
\]
In particular the constant score $S=\pi$ is perfectly calibrated with $B=\pi(1-\pi)$, and that
constant score is the least sharp calibrated forecaster.
\end{proposisi}

\begin{proof}[Proof]
Conditional on $S=s$, the variable $y$ takes the value $1$ with probability $s$. Then
\[
\mathbb{E}\bigl[(s-y)^2\mid S=s\bigr]
=s(s-1)^2+(1-s)s^2=s(1-s)\bigl[(1-s)+s\bigr]=s(1-s).
\]
Taking the expectation over $S$ gives the first form. For constant $S=\pi$, the calibration
condition demands $P(y=1)=\pi$, which is indeed true, and the formula gives $\pi(1-\pi)$.
\end{proof}

That identity explains two things: calibration alone is not enough, since the constant score $\pi$
is perfectly calibrated but distinguishes no examples; and for rare events the Brier baseline is
very low, so a model that looks good on accuracy can lose badly on Brier. In our simulation
$\pi(1-\pi)=0.000999$, while the hard labels of the model give a Brier equal to its error rate,
$1051/100\,000=0.010510$. The Brier skill score, $1-B/B_{\mathrm{garis\ dasar}}$, is $-9.52$: the
model is ten times worse than forecasting the constant prevalence.

That negative number is no oddity: hard labels throw away the information about uncertainty, and
for rare events that information is the most valuable of all. A model that outputs $0.06$ is saying
something true, and a report that carries only accuracy throws that sentence away.

\section{Uncertainty of metrics}
\tingkat{advanced}

All the numbers above are computed from a limited number of examples, so all of them wobble when
the data is replaced. In everyday language: a metric without an interval is a single shot with no
indication of how wide the target is.

Chapter 20 already derived the Wilson interval and proved its properties, so we do not repeat it
here. What must be remembered is its form for a proportion $k/n$:
\[
\frac{\hat p+\dfrac{z^2}{2n}\pm
z\sqrt{\dfrac{\hat p(1-\hat p)}{n}+\dfrac{z^2}{4n^2}}}{1+\dfrac{z^2}{n}},
\qquad z=1.96 .
\]

\begin{untukanak}
Two goalkeepers practice stopping penalties. The first stops nine out of ten shots. The second
stops nine hundred out of a thousand shots. The newspaper prints the same headline for both: ninety
percent.

But we feel that the two sentences are not equally strong, and that feeling is right. The first
keeper has faced only ten shots. Perhaps he really is excellent, perhaps he was merely lucky ten
times in a row. The second keeper has faced a thousand shots, and luck over that long a run is
almost impossible. The same number, a different strength.

The question that makes one curious: if we may write only one sentence about the first keeper, may
``ninety percent'' still be written? Yes, provided the sentence is completed. In this chapter we
will learn why the same number means ``between fifty-nine and ninety-eight percent'' on ten shots,
but only ``between eighty-eight and ninety-two percent'' on a thousand shots, and why multiplying
the shots a hundredfold narrows the answer only to a tenth.

\end{untukanak}

The following table applies that formula to accuracy, with $z=1.96$.

\begin{center}
\small
\begin{tabular}{lrrr}
\hline
data & $\hat p$ & Wilson interval & width \\
\hline
$9/10$ & 0.900 & $(0.5959;\ 0.9821)$ & 0.3863 \\
$900/1000$ & 0.900 & $(0.8798;\ 0.9171)$ & 0.0372 \\
$10/10$ & 1.000 & $(0.7225;\ 1.0000)$ & 0.2775 \\
$0/1000$ & 0.000 & $(0.0000;\ 0.0038)$ & 0.0038 \\
\hline
\end{tabular}
\end{center}

The first two rows answer the two goalkeepers. The point estimate is the same, $0.9$, but the first
interval is $0.3863$ wide and the second $0.0372$: ten times narrower, not a hundred times.

The third row is the case that appears most often: not a single error in $10$ trials, and then
``accuracy $100$ percent'' is written. The honest Wilson interval, $(0.7225;1)$, is $0.2775$ wide:
the error rate may still be $27.75$ percent. At the other end, zero events in $1000$ trials gives an
upper bound of $0.003827$, while the rule of three from Chapter 20 gives $0.003$; either is
acceptable, as long as the procedure is named.

Comparing two models by checking whether their intervals overlap is a wrong habit: both are tested
on the same data, and that dependence sharpens the comparison rather than blurring it.

\begin{proposisi}[The difference of two paired models]
Let $b$ be the number of examples that model A gets right and model B gets wrong, and $c$ the
reverse, on the same $n$ examples. Their accuracy difference is $d=(b-c)/n$ with variance
\[
\operatorname{Var}(d)=\frac{1}{n}\Bigl(p_b+p_c-(p_b-p_c)^2\Bigr),
\]
so that its estimated standard error is $\sqrt{b+c-(b-c)^2/n}\,/\,n$.
\end{proposisi}

\begin{proof}[Proof]
Every example falls in one of four cells: right for both, wrong for both, right only for A, right
only for B. Since the examples are independent, $(b,c)$ is multinomial with probabilities $p_b,p_c$.
Then $\operatorname{Var}(b-c)=n\bigl(p_b+p_c-(p_b-p_c)^2\bigr)$, and dividing by $n^2$ gives the
formula above.
\end{proof}

At $b=37$, $c=15$, and $n=1000$, the difference is $0.022$ with standard error $0.007177$, so the
$95$ percent interval is $(0.007932;0.036068)$, entirely above zero, and the exact McNemar test
gives a p-value of $0.003185$. The naive treatment, which assumes the two models independent, gives
the interval $(-0.004296;0.048296)$, which contains zero. The two models can therefore be declared
different, provided the comparison is paired and the test uses the difference, not two separate
intervals.

Finally, the size of the test determines how small a difference can be seen. To detect a difference
$\delta$ at level $\alpha$ and power $1-\beta$,
\[
n\approx\frac{(z_{1-\alpha/2}+z_{1-\beta})^2\,\bigl(p_1(1-p_1)+p_0(1-p_0)\bigr)}
{\delta^2}.
\]
At $p_0=p_1=0.9$, level $0.05$, and power $0.8$, a difference of one point demands
$n\approx14\,128$ examples per model, a difference of two points demands $3532$. A test set of a
thousand examples cannot separate two models one point apart.

\section{Data leakage and nested cross-validation}
\tingkat{advanced}

All the computations above assume one thing that must not be violated: the test number is measured
on data the model has never seen, and never seen either by the procedure that chose the model. In
everyday language: an exam score means something only if the questions were never leaked, including
by the teacher who wrote them.

\kamus{data leakage}{kebocoran data}{exam answers that slip accidentally into the study material, so
that the score goes up without any skill}

\begin{definisi}[Data leakage]
\index{data leakage}
Leakage occurs when information that would not be available at the moment the model is used enters
training, model selection, or reporting, so that the evaluation number measures the ability to
memorize, not the ability to generalize.
\end{definisi}

Its forms are various, and the most frequent ones are invisible. Computing a preprocessing mean or
median from all the data, test data included, is leakage: that statistic should come from training
data only. Selecting the $50$ best features by correlation with the label before separating
training from test is also leakage, since the feature filter has already seen the test labels.
Tuning hyperparameters on the test set and reporting the best number is the leakage most often
committed knowingly. On grouped data, a random split places the same group on both sides; on time
series, it leaks the future into the past. The hardest to detect: a feature that is a consequence
of the label, for example a column ``action taken'' filled in after the incident is known.

Feature-selection leakage can be measured, and its size is startling. Take $200$ randomly labeled
examples and $2000$ features of normal noise, with no signal at all. Select the $20$ features with
the largest absolute correlation with the label, then cross-validate with five folds using a nearest
centroid classifier. If the selection is done once on all the data before the folds are formed, the
accuracy is $0.7800$; if the selection is repeated inside each fold using only its training rows,
$0.4600$. On three other seeds the numbers are $0.775$ and $0.415$, $0.820$ and $0.545$: the leaky
version is always high, the honest one always near $0.5$. That number $0.78$ is an illusion from
beginning to end, and not one line of its code is mistyped.

The protocol is one sentence: separate first, then compute anything. Split the data into training,
validation, and test; compute preprocessing statistics from training only; tune hyperparameters on
validation; touch test once, at the end, for one number. When cross-validation is used, the
preprocessing lives inside the fold. For grouped data, the folds follow the groups; for time data,
they follow the time order.

Figure~\ref{fig:pemisahan-tiga} draws that order as three boxes and one forbidden arrow.

\begin{figure}[htbp]
\centering
\begin{tikzpicture}[every node/.style={font=\footnotesize},
  blok/.style={draw=ink,thick,rounded corners=2pt,align=center,
               minimum width=3.5cm,minimum height=1.5cm,inner sep=3pt}]
  \node[blok,fill=layerteal!12]   (latih) at (0.0,0) {\textbf{train} $60\%$\\ compute statistics};
  \node[blok,fill=layerorange!12] (val)   at (4.8,0) {\textbf{validation} $20\%$\\ choose hyperparameters};
  \node[blok,fill=layerindigo!12] (uji)   at (9.6,0) {\textbf{test} $20\%$\\ touch once};
  \draw[-{Latex},ink,thick] (latih.east) -- (val.west)
    node[midway,above,font=\scriptsize] {retrain};
  \draw[-{Latex},ink,thick] (val.east) -- (uji.west)
    node[midway,above,font=\scriptsize] {locked before inspection};
  \draw[-{Latex},accent,thick,dashed] (uji.south) .. controls (9.6,-1.85) and (0.0,-1.85)
    .. (latih.south);
  \draw[accent,very thick] (4.62,-1.42) -- (4.98,-1.06);
  \draw[accent,very thick] (4.62,-1.06) -- (4.98,-1.42);
  \node[accent,anchor=west,font=\scriptsize] at (5.12,-1.24) {forbidden};
  \node[anchor=north,align=center,text width=11.6cm] at (4.8,-1.98)
    {test statistics, feature filters, and hyperparameter tuning must not cross
     back into the training box};
\end{tikzpicture}
\caption{Three separate data sets, with proportions $60/20/20$; on $1000$ connections that means
$600$, $200$, and $200$. Preprocessing statistics are computed from the training box,
hyperparameters are chosen in the validation box, and the test box is touched once at the end. The
dashed arrow from test back to training is leakage: that single arrow is enough to raise accuracy
from $0.4600$ to $0.7800$ on data that is pure noise.}
\label{fig:pemisahan-tiga}
\end{figure}

What remains is one subtle problem. When cross-validation is used both to choose hyperparameters and
to report quality, the reported number is the best of many trials on the same data, so it is too
optimistic. The way out is two layers of folds.

\kamus{nested cross-validation}{validasi silang bersarang}{two layers of folds, the outer layer to
judge and the inner layer to choose, like two separate examination rooms}

\begin{definisi}[Nested cross-validation]
\index{cross-validation}
\index{nested cross-validation}
Let $\Lambda$ be a set of hyperparameters. The data is split into $K$ outer folds. For each outer
fold $k$: train on the remaining data with every $\lambda\in\Lambda$ using $J$ inner folds, choose
the best $\lambda_k$, retrain on the whole training part of fold $k$ with $\lambda_k$, then measure
on fold $k$. The reported number is the average of those $K$ numbers. Judging and selection
therefore use different data.
\end{definisi}

Its cost is $K$ times higher, and that is why this protocol is often skipped; without the outer
layer, not one number in the report can be read as an estimate of quality on new data.

\section{Why one test is not enough}
\tingkat{intermediate}

The sentence ``we tested the model on $1000$ examples and obtained an accuracy of $0.900$'' sounds
conclusive, yet it carries three uncertainties. In everyday language: one weighing does not tell
you whether the scale wobbles, whether the goods changed, or whether we picked the most pleasant
result.

The first is statistical and already computed: at $n=1000$ the $95$ percent interval is $0.037$
wide, and separating two models one point apart demands $14\,128$ examples. The conclusion ``A is
better'' from a difference that small is a conclusion about a coin toss.

The second is selective, and its name is \emph{winner's curse}. Training $20$ variants of identical
quality and then reporting the best means reporting the maximum of $20$ wobbles: in a simulation
with true accuracy $0.90$ and $n=1000$, the mean of each variant is $0.899987$ while the mean of
their maximum is $0.917337$, so $1.7$ points of that number are luck. When the variants are
dependent, the wobble is smaller, but its direction is always upward.

The third is procedural: a test set used repeatedly stops being a test set. If each comparison has
a $5$ percent chance of deceiving, the chance that at least one of $20$ comparisons deceives is
$1-0.95^{20}=0.6415$, larger than the chance of being right.

The remedies match: report intervals, separate the data used for selection from the data used for
reporting, and lock the test set before the experiment begins. If the test has already been used
repeatedly, the next number is called an exploratory number.

\section{A reporting checklist}
\tingkat{basic}

These seven lines separate a report that can be checked from one that merely convinces. In everyday
language: a receipt does not make the goods arrive, but without a receipt nothing can be traced.

\begin{enumerate}
  \item The size of each class and its prevalence, not only the total $n$.
  \item The four numbers of the confusion matrix, then precision and recall at the threshold
        actually used, together with the value of that threshold.
  \item A confidence interval for every main metric, with the method named; for small $n$ use
        Wilson, not Wald.
  \item A precision--recall curve when the positive class is rare, together with average precision
        and its baseline $\pi$.
  \item Calibration: the score interval bounds, the number of examples per interval, and the
        calibration error.
  \item The separation protocol: what was computed from which data, how many times the test data
        was touched, and how many variants were tried; the best number without that denominator is
        a selected number, not a quality number.
  \item The spread of results across seeds or across folds, together with the code and the random
        seed.
\end{enumerate}

What must not be there: accuracy without prevalence, AUC without a curve, and the phrase ``ready to
deploy''.

\begin{resep}
Four steps for reporting evaluation results that do not mislead:
\begin{enumerate}
  \item Separate the data first: training, validation, test, with their proportions and the
        splitting key, written down before a single computation is run.
  \item Compute preprocessing and select features from training only, tune hyperparameters on
        validation only, then save test for a single touch.
  \item Report all four cells of the confusion matrix together with $n$, $\pi$, and the threshold
        used, then precision, recall, and a confidence interval for every metric.
  \item State how many variants were tried and the spread of results across seeds or across folds,
        so that the reader knows how many degrees of freedom have already been spent.
\end{enumerate}
A warning about order: step 1 must not come after step 2 or step 3. Statistics computed before the
data is separated have already moved test information into training, and after that no number can
repair it; the whole report then measures memory.
\end{resep}

\section{Worked examples}
\tingkat{intermediate}

The three computations below can be repeated on paper.

First, the confusion matrix from the author's own working notes, on $35\,069$ test connections:
$\mathrm{TP}=23\,279$, $\mathrm{FP}=997$, $\mathrm{FN}=590$, $\mathrm{TN}=10\,203$.
\begin{align*}
\operatorname{prec}&=\frac{23\,279}{24\,276}=0.958931, &
\operatorname{rec}&=\frac{23\,279}{23\,869}=0.975282,\\
F_1&=\frac{46\,558}{48\,145}=0.967037, &
\mathrm{acc}&=\frac{33\,482}{35\,069}=0.954746,
\end{align*}
and its balanced accuracy is $0.943132$. One thing in those notes is inconsistent: one sentence
states $11\,200$ attacks and $23\,869$ normal, while the matrix places $23\,869$ in the positive
row. The arithmetic is right, the class labels are swapped, and the prevalence must be read from
the matrix, not from the sentence.

Second, the AUC of a small data set. Positives scoring $0.92$, $0.81$, $0.55$ against negatives
scoring $0.70$, $0.41$, $0.35$, $0.22$ give $12$ pairs; $0.92$ and $0.81$ beat all four, $0.55$
beats three, so $11/12=0.916667$. If one positive score equals one negative score, that pair counts
as one half and the result is $0.875$; \texttt{sklearn} gives the same number when $y$ and the
scores are fed in as vectors.

Third, the effect of prevalence on precision when sensitivity $0.90$ and specificity $0.95$ are held
fixed. With $\operatorname{prec}=\pi(0.9)/\bigl(\pi(0.9)+(1-\pi)(0.05)\bigr)$:

\begin{center}
\small
\begin{tabular}{rrrr}
\hline
$\pi$ & $\operatorname{prec}$ & true alarms & false alarms \\
\hline
0.001 & 0.017699 & 1 & 56 \\
0.010 & 0.153846 & 1 & 5.5 \\
0.100 & 0.666667 & 2 & 1 \\
0.500 & 0.947368 & 18 & 1 \\
\hline
\end{tabular}
\end{center}

The last column is the reciprocal of precision minus one, namely the number of false alarms for
every true alarm. At a prevalence of one in a thousand, one true alarm arrives with $55.5$ false
alarms; at a prevalence of one half, with $0.056$. The same model, the same quality, and a
completely different experience for the operator.

\kertas{Take the matrix of $1000$ predictions from earlier and recompute its five metrics without
looking at the text. Then answer: how many $\mathrm{FP}$ must be removed, without changing
$\mathrm{TP}$ and $\mathrm{FN}$, for precision to become $0.900$? Hint: precision $62/(62+x)=0.9$
gives $x=6.89$, so $\mathrm{FP}$ must fall from $17$ to at most $6$. Check that the accuracy after
that change is $0.956$.}

\section{Computational lab: one number, four ways to read it}
\tingkat{advanced}

Every number in this chapter is computed by one short program, which also shows two things a
formula cannot show: how far leakage raises a number, and how far the winner among many variants is
flattered. The code is run as it stands with \texttt{numpy} and \texttt{scipy}.

\begin{lab}
\begin{verbatim}
import numpy as np
from math import sqrt
from scipy.stats import norm

z = norm.ppf(0.975)


def wilson(k, n):
    p, z2 = k / n, z * z
    mid = (p + z2 / (2 * n)) / (1 + z2 / n)
    half = z * sqrt(p * (1 - p) / n + z2 / (4 * n * n)) / (1 + z2 / n)
    return mid - half, mid + half


def roc_pr(spos, sneg):
    s = np.concatenate([spos, sneg])
    y = np.concatenate([np.ones(len(spos)), np.zeros(len(sneg))])
    o = np.argsort(-s)
    s, y = s[o], y[o]
    tp, fp = np.cumsum(y), np.cumsum(1 - y)
    tpr, fpr = tp / len(spos), fp / len(sneg)
    prec = tp / (tp + fp)
    auc = np.sum(np.diff(np.concatenate([[0], fpr])) *
                 (tpr + np.concatenate([[0], tpr[:-1]])) / 2)
    ap = np.sum(np.diff(np.concatenate([[0], tpr])) * prec)
    return auc, ap


rng = np.random.default_rng(20261025)
Nneg, Npos = 99900, 100
sneg = rng.normal(0.0, 1.0, Nneg)
spos = rng.normal(2.6, 1.0, Npos)
t = 2.3263
tp = int((spos > t).sum()); fp = int((sneg > t).sum())
fn, tn = Npos - tp, Nneg - fp
n = Nneg + Npos
print("TP=%d FP=%d FN=%d TN=%d" % (tp, fp, fn, tn))
print("accuracy %.6f  baseline %.6f" % ((tp + tn) / n, Nneg / n))
print("precision %.6f  recall %.6f" % (tp / (tp + fp), tp / Npos))
print("balanced accuracy %.6f" % ((tp / Npos + tn / Nneg) / 2))
print("brier hard labels %.6f  brier constant %.6f"
      % ((fp + fn) / n, 0.001 * 0.999))
a1, p1 = roc_pr(spos, sneg)
print("pi=0.001: AUC=%.6f  AP=%.6f" % (a1, p1))
a2, p2 = roc_pr(rng.normal(2.6, 1.0, 50000), rng.normal(0.0, 1.0, 50000))
print("pi=0.500: AUC=%.6f  AP=%.6f" % (a2, p2))
for k, m in ((9, 10), (900, 1000), (10, 10), (0, 1000)):
    lo, hi = wilson(k, m)
    print("%4d/%-4d (%.6f, %.6f) width %.6f" % (k, m, lo, hi, hi - lo))
b, c = 37, 15
d = (b - c) / 1000
se = sqrt(b + c - (b - c) ** 2 / 1000) / 1000
print("difference %.6f  se %.6f  CI (%.6f, %.6f)"
      % (d, se, d - z * se, d + z * se))
rt = np.random.default_rng(11)
a = rt.binomial(1000, 0.90, (20000, 20)) / 1000
print("20 variants: mean %.6f  mean max %.6f"
      % (a.mean(), a.max(axis=1).mean()))

rl = np.random.default_rng(4)
nn, dd, kk = 200, 2000, 20
X = rl.normal(size=(nn, dd))
y = rl.permutation(np.repeat([0, 1], nn // 2))
corr = np.array([abs(np.corrcoef(X[:, j], y)[0, 1]) for j in range(dd)])
sel_all = np.argsort(-corr)[:kk]


def cv(pick):
    idx = rl.permutation(nn)
    ok = 0
    for f in np.array_split(idx, 5):
        tr = np.setdiff1d(idx, f)
        cols = pick(tr)
        m1 = X[np.ix_(tr[y[tr] == 1], cols)].mean(0)
        m0 = X[np.ix_(tr[y[tr] == 0], cols)].mean(0)
        d1 = ((X[np.ix_(f, cols)] - m1) ** 2).sum(1)
        d0 = ((X[np.ix_(f, cols)] - m0) ** 2).sum(1)
        ok += (np.where(d1 < d0, 1, 0) == y[f]).sum()
    return ok / nn


def top_in(tr):
    c = np.array([abs(np.corrcoef(X[tr, j], y[tr])[0, 1]) for j in range(dd)])
    return np.argsort(-c)[:kk]


print("feature selection: leaky %.4f  honest %.4f"
      % (cv(lambda tr: sel_all), cv(top_in)))
\end{verbatim}
\end{lab}

The actual output, as it stands:

\begin{verbatim}
TP=62 FP=1013 FN=38 TN=98887
accuracy 0.989490  baseline 0.999000
precision 0.057674  recall 0.620000
balanced accuracy 0.804930
brier hard labels 0.010510  brier constant 0.000999
pi=0.001: AUC=0.966448  AP=0.287636
pi=0.500: AUC=0.967061  AP=0.967544
   9/10   (0.595850, 0.982124) width 0.386274
 900/1000 (0.879848, 0.917091) width 0.037243
  10/10   (0.722467, 1.000000) width 0.277533
   0/1000 (0.000000, 0.003827) width 0.003827
difference 0.022000  se 0.007177  CI (0.007932, 0.036068)
20 variants: mean 0.899987  mean max 0.917330
feature selection: leaky 0.7800  honest 0.4600
\end{verbatim}

Four things must be read from that output. First, the two numbers on the second line must appear
side by side: $0.989490$ and $0.999000$. The direction of the difference decides the whole
interpretation: the model under test is below the baseline.

Second, the AUC barely moves, $0.966448$ against $0.967061$, while AP moves from $0.287636$ to
$0.967544$, a factor of $3.4$: numerical evidence for the previous proposition.

Third, the width of the interval shrinks proportionally to $1/\sqrt n$, from $0.386274$ at $n=10$ to
$0.037243$ at $n=1000$, and $0/1000$ gives an upper bound of $0.003827$.

Fourth, the last two lines carry two diseases that are invisible in a single number: the paired
difference $0.022$ with interval $(0.007932;0.036068)$ does not contain zero, yet
$(0.8798;0.9171)$ and $(0.8563;0.8969)$ overlap widely; and the winner among $20$ identical
variants is reported as $0.917337$ although the true quality is $0.899987$.

The program proves not one proposition in this chapter; it checks claims that have already been
proved. The simulation uses distributions we assume, and honest cross-validation on noise data
gives $0.4600$, not $0.5000$: even a number without signal wobbles.

\begin{untukdosen}
The teaching order I suggest: start from the vendor report in the chapter opening, write $99$
percent and $99.9$ percent on the board, then ask the class to decide which is better. Almost the
whole class picks the detector, and out of that unease the confusion matrix appears as an answer,
not as a definition to be memorized. Once precision and recall exist, then $F_1$ and balanced
accuracy can enter. The ROC curve is best drawn by hand for the six points of the threshold table,
because there the students see for themselves that a point which looks good on the ROC curve can
have precision $0.013$.

Two typical misconceptions. First, the AUC is taken to be a context-free measure of quality, so the
threshold is chosen at random; the symptom shows when students cannot explain the precision at the
threshold in use. Second, cross-validation is taken to be always honest, even though feature
selection or standardization outside the fold is enough to raise the number by $30$ points on data
that is pure noise.

Board activity: split the class into two groups, one computing precision for $\pi=0.001$ and one
for $\pi=0.5$ with the same sensitivity and specificity, then ask both to defend their results. The
first group finds $0.018$ and must explain why the same model feels different. Sign that the class
is ready to go on: they can name one situation in which accuracy $0.99$ is worse than not using a
model at all.
\end{untukdosen}

\section{Summary}
\tingkat{basic}

\begin{ringkasan}
\begin{itemize}
  \item The four numbers of the confusion matrix play unsymmetric roles: precision answers the
        operator's question, recall the policy designer's question, and both depend on the
        threshold; $F_1$ is their harmonic mean and ignores $\mathrm{TN}$.
  \item The baseline accuracy is $1-\pi$: at $\pi=0.001$ a model with $0.989490$ loses to a program
        that always answers negative ($0.999000$), while balanced accuracy $0.804930$ is more honest
        because it does not depend on composition.
  \item The AUC equals the probability that a random positive score exceeds a random negative
        score, so it depends only on the conditional score distributions; it does not change when
        prevalence changes, whereas precision and average precision do, from $0.288$ to $0.968$.
  \item Calibration checks whether a score of $0.8$ means eight out of ten. For calibrated scores,
        Brier equals $\mathbb{E}[S(1-S)]$; the constant score $\pi$ is perfectly calibrated but not
        sharp. Hard labels here give $0.010510$ against $0.000999$.
  \item The width of the interval shrinks proportionally to $1/\sqrt n$: $0.3863$ at $n=10$ and
        $0.0372$ at $n=1000$. To compare two models, use the paired difference: $b=37$, $c=15$
        gives $(0.007932;0.036068)$, whereas the unpaired treatment gives an interval that contains
        zero.
  \item One test is not enough, and leakage raises numbers falsely: $0.7800$ when features are
        selected outside the fold, $0.4600$ when inside; the winner among $20$ identical variants is
        flattered by $1.7$ points; and $20$ comparisons at the $5$ percent level give a chance of
        $0.6415$ that one of them deceives.
\end{itemize}
\end{ringkasan}

\section{Exercises}
\tingkat{intermediate}

\begin{latihan}
\begin{enumerate}
  \item Assess this claim: ``Model A reaches an accuracy of $99.1$ percent and model B $98.7$
        percent on the same test set, so A is better.'' True, false, or not enough evidence? Name
        the part that is wrong and two additional numbers that would strengthen it.
  \item Prove that the Brier score for hard $0/1$ labels equals the error rate,
        $(\mathrm{FP}+\mathrm{FN})/n$. Why does that equality fail for scores that are
        probabilities?
  \item Prove that a constant score $S=s$ minimizes the Brier score exactly when $s=\pi$, with
        minimum value $\pi(1-\pi)$.
  \item Given $\mathrm{TP}=180$, $\mathrm{FP}=45$, $\mathrm{FN}=20$, $\mathrm{TN}=1755$. Compute
        precision, recall, $F_1$, accuracy, specificity, and balanced accuracy to six decimal
        places, then check with \texttt{sklearn.metrics}.
  \item Compute the $95$ percent Wilson interval for $95/100$ and $950/1000$. Compare the two
        widths, then explain why that comparison must not be used to conclude which model is
        better.
  \item Compute the AUC from five positives scoring $0.95$, $0.88$, $0.70$, $0.60$, $0.40$ and five
        negatives scoring $0.75$, $0.65$, $0.50$, $0.30$, $0.10$. Write down the number of correct
        pairs, then check with \texttt{roc\_auc\_score}.
  \item A model has sensitivity $0.80$ and specificity $0.99$. Compute its precision at
        $\pi=0.002$ and $\pi=0.2$, then determine the minimum prevalence for precision to be at
        least $0.5$.
  \item You have $6000$ records from $300$ hosts, about $20$ records each, and $40$ hyperparameter
        settings. Design an evaluation protocol without leakage: state the data split, where
        preprocessing is computed, where hyperparameters are chosen, and how many times test is
        touched.
  \item Assess this claim: ``Because the AUC of this model is $0.95$, the model is useful at any
        threshold we choose.'' True, false, or not enough evidence? Include one supporting
        computation.
\end{enumerate}

Hints for the three hardest items. Item 2: write $(s-y)^2$ for $s\in\{0,1\}$; its value is $1$
exactly on an error. Item 3: differentiate $\mathbb{E}[(s-y)^2]=s^2-2s\pi+\pi$ with respect to
$s$, find its minimum, then use the Brier proposition in this chapter. Item 7: its precision is
$0.8\pi/\bigl(0.8\pi+0.01(1-\pi)\bigr)$; solve
$0.8\pi\ge0.5\bigl(0.8\pi+0.01(1-\pi)\bigr)$ for $\pi$.
\end{latihan}

\begin{untukpeneliti}
What is proved in this chapter: the accuracy of a constant classifier is $1-\pi$; $F_1$ is the
harmonic mean of precision and recall with $F_1\le(P+R)/2$; the AUC equals the probability of a
correct ranking together with its Mann--Whitney form; ROC and AUC are invariant to prevalence while
precision is not; for calibrated scores Brier equals $\mathbb{E}[S(1-S)]$; and the variance formula
for the difference of two paired proportions.

What is popular but not a theorem: that ``AUC $0.9$ means a good model''; that the threshold $0.5$
is neutral; that cross-validation is always honest; and that overlapping intervals mean there is no
difference. The last three are mathematically wrong, not merely incomplete. \emph{Average
precision} is also not an unbiased estimate of the area under the population precision--recall
curve: it is a summary of the observed curve, and its value depends on the sample prevalence.

Primary sources: Wasserman, \emph{All of Statistics}, Chapter 6, Chapters 9--10, Chapter 13 in the
Model Selection part, and Chapter 22 in the part ``Assessing Error Rates and Choosing a Good
Classifier''. Goodfellow, Bengio, and Courville, \emph{Deep Learning}, Chapter 11 in the part
``Performance Metrics'', give the one-in-a-million disease example and the notion of
\emph{coverage}. Athalye, Carlini, and Wagner (2018) show that weak evaluation produces false
security claims: seven of nine defenses without a certificate at ICLR 2018 relied on obfuscated
gradients, and six fell once their evaluation was fixed.

Limits of validity: all simulations use assumed normal distributions, so they measure the behavior
of the procedure, not the truth of the model. The Wilson interval is a score approximation, and
McNemar assumes the discordant pairs are symmetric under the null hypothesis. The leakage result
depends on $n$, $d$, and the number of features selected: $0.78$ is not a constant of nature but
one point on a curve that rises with $d/n$.

Open questions: honest intervals for metrics from mutually dependent folds; calibration tests
without the assumption of uniform score intervals; and reporting uncertainty when the test set is
chosen adaptively.
\end{untukpeneliti}

\peta{Chapter 27 of Wasserman, \emph{All of Statistics}: Chapter 6 (estimators, bias, standard
error), Chapters 9--10 (intervals and tests), Chapter 13 in the Model Selection part, and Chapter 22
in the part ``Assessing Error Rates and Choosing a Good Classifier''. The Wilson interval and the
rule of three are referred from Chapter 20 of this book and from Wilson (1927) and Jovanovic--Levy
(1997). The example of deceptive accuracy follows Goodfellow, Bengio, and Courville, \emph{Deep
Learning}, Chapter 11; the limits of adversarial-defense evaluation follow Athalye, Carlini, and
Wagner (2018), discussed in Chapters 27 and 28. Bayes' rule for precision is proved in Chapter 16;
data splitting in Chapter 24.}

\chapter{Formal Methods with Lean 4}

A file named \texttt{Sertifikat.lean} arrived together with a single sentence of introduction:
``everything has already been proved by machine''. It contained forty one-line theorems, almost all
of them ending with \texttt{by native\_decide}, and at the bottom there was one summary constant
that mentioned the threshold $10^{-3}$. Not one of its lines can we inspect with our eyes in a
reasonable amount of time: what is being compared are sixteen-digit integers that come from
turning binary fractions into exact form. That introductory sentence is true in a narrow sense and
misleading in an important sense. The machine does check every line, but it does not check whether
what is being checked is what was meant.

Testing and proving promise different things, and the difference is not a question of degree of
confidence. Testing takes a number of inputs, runs the program, and then reports what happened on
those inputs. The conclusion always has the form ``for $10^{4}$ test samples, nothing failed''. A
formal proof promises more: one statement quantified over all inputs, together with an object that
a third party can check without trusting its author. That promise is kept only when three
conditions hold at once, namely that the statement is written correctly, that the object being
checked really is a proof of that statement, and that the foundation of the checking has not been
misused. This chapter takes those three conditions apart on one concrete system, Lean 4.

The tools of the previous chapter are not enough for this work. Chapter 27 taught us how to report
what is actually measured: keep the training data apart from the test data, state the distribution,
report the variance across seeds. That is a discipline about \emph{empirical claims}. Here we face a
claim of a different shape: not an average over samples, but a statement quantified as
$\forall x,\ P(x)$. For a claim of that kind no confidence interval helps, because a single
counterexample is enough to bring it down, and the absence of a counterexample is not a proof. What
is needed is an object that can be rechecked by someone who does not believe us.

\begin{untukanak}
Imagine a very meticulous chess referee. He never asks whether the game is interesting, whether
both players play well, or whether the rules of chess are in fact the rules that ought to be used.
His task is only one: every time a move appears, he checks whether that move is legal according to
the rule book. If all the moves are legal, he signs the board and says ``this game is valid''. He
does not care who wins.

Now imagine someone hands over a board full of legal moves, then says ``this proves that I can beat
anyone''. The referee will answer honestly: what I checked is only that the moves are legal, not
that your sentence is true. If the sentence is wrong, a legal board does not help at all, because
what is wrong is not the moves but the question.

Have you ever given a correct answer to the wrong problem? Or done your homework neatly when what
was asked for was a different page? In this chapter we will find out why a computer can check a
proof down to the smallest step, why that checking does not in the least guarantee that the
question is right, and how people write the right question in a language that a machine can check.
\end{untukanak}

\section{Informal proof and formal proof}

\tingkat{intermediate}

A proof in a mathematics book is a text that convinces a trained reader. It contains jumps that are
agreed on in common: ``clearly'', ``without loss of generality'', ``analogous for the other cases''.
Those jumps are not cheating; they are a compression accepted by the community. The trouble is that
such compression cannot be checked by a machine, and cannot be separated from the knowledge of its
reader.

A formal proof is another version of the same object. Its statement is written in a language with a
strict grammar and strict types, and its proof is a \emph{term}, namely an expression whose type is
that statement. The checking is carried out by the kernel of the system, a small piece of code that
knows only a few rules and knows nothing at all about mathematics. Because the kernel is small, its
checking can be trusted far more strongly than any large program around the system.

\kamus{formal verification}{verifikasi formal}{Step-by-step checking by a clerk who is never bored
and never trusts the author.}

\begin{definisi}[Formal proof]
Let $P$ be a statement written in the formal language of the system. A \emph{formal proof} of $P$ is
a term $p$ such that $p:P$ can be derived by the kernel of the system from definitions, theorems
that have already been checked, and axioms declared open. The system accepts $P$ if and only if a
term of that kind is found and passes the type check.\\index{formal verification}
\end{definisi}

That definition contains three words that are easy to pass over. ``Can be derived'' means that
there is a mechanical procedure, not that there is a person who agrees. ``Checked'' means that the
outcome is binary: accepted or rejected, with no degrees of confidence. ``Declared open'' means that
the axiom is recorded, not hidden; we will use that recording later as our main instrument.

The difference that is most often misunderstood is the difference between \emph{truth} and
\emph{provability}. Lean does not know whether a statement is true in the real world. It knows only
whether there is a term whose type is that statement. If the statement is false, that term will
never be found, but the absence of a term is not a proof that the statement is false; it may simply
be that we have not found it yet.

\begin{bacakeras}
First voice: ``I have checked it, everything is correct.''
Second voice: ``Show me how.''
First voice: ``I read it again three times.''
Second voice: ``That is a habit, not a method.''
First voice: ``I am sure.''
Second voice: ``Certainty is not proof. Give me the steps, one by one.''

That is what a conversation with a machine is like. It does not ask whether we are good people,
whether we have had breakfast, or whether we graduated with high marks. It asks one thing,
patiently, every time, with no angry tone: show me the steps. If the steps are there, the door
opens. If the steps are not there, the door stays shut, as tightly as before, even if we repeat the
request a hundred times. The machine does not judge our intentions; it only checks the form.

Say it together, slowly, then repeat it once more: the machine does not believe, it checks.
\end{bacakeras}

\section{Propositions as types, proofs as terms}

\tingkat{intermediate}

In Lean there is no gulf between a statement and its proof. Both are terms, checked by the same
rules. Every sentence that can take the value true or false is given the type $\mathtt{Prop}$; a
proof of that sentence is a term whose type is that sentence. The statement ``$P$ is proved''
therefore does not mean that $P$ is true in some model, but that there is a term of type $P$ that
passes the check.

\kamus{proposition}{proposisi}{A sentence that can be judged true or false, like an exam question
whose answer must be proved, not merely chosen.}

\begin{definisi}[Inhabitant of a type]
A type $T$ is called \emph{inhabited} when there is a term $t$ with $t:T$; the notation $t:T$ is
read ``$t$ has type $T$''. If $P:\mathtt{Prop}$, then $P$ is called a \emph{proposition} and $t:P$
is called a \emph{proof term} for $P$.\\index{proof term} Two terms are called \emph{definitionally
equal} when both reduce to an identical form by symbolic computation.\\index{proposition-type}
\end{definisi}

This dictionary, namely the correspondence between logic and types, is known as
Curry--Howard.\\index{Curry-Howard} Implication becomes a function type, because a proof of
$P\Rightarrow Q$ must turn an arbitrary proof of $P$ into a proof of $Q$, and that is the work of a
function. The universal quantifier becomes a dependent function, because a proof of
$\forall x,\ P(x)$ must produce a proof of $P(a)$ for every $a$ that is given. Conjunction becomes
a pair, because a proof of $P\wedge Q$ contains both of its proofs at once. The last row of the
table below is the most important: $\bot$ has no inhabitant.

\begin{center}
\begin{tabular}{lll}
\hline
Logic & Type in Lean & How to build the term \\
\hline
$P\Rightarrow Q$   & $P\to Q$              & \texttt{fun h => ...} \\
$P\wedge Q$        & \texttt{And P Q}      & \texttt{And.intro} \\
$P\vee Q$          & \texttt{Or P Q}       & \texttt{Or.inl}, \texttt{Or.inr} \\
$\top$             & \texttt{True}         & \texttt{True.intro} \\
$\bot$             & \texttt{False}        & no inhabitant \\
$\forall x,\ P(x)$ & $\forall x,\ P\,x$    & \texttt{fun x => ...} \\
$\exists x,\ P(x)$ & \texttt{Exists}       & \texttt{Exists.intro} \\
\hline
\end{tabular}
\end{center}

\kamus{proof term}{term bukti}{A building frame whose shape follows the blueprint of its statement
exactly; if the shape does not fit, the building is rejected.}

The following terms are the smallest useful proofs: the identity function as a proof of the identity
implication, and the composition of two functions as a proof of the transitivity of implication.
Both are checked without using any knowledge at all about the contents of $P$, $Q$, and $R$.

\begin{verbatim}
-- proof of the identity implication, term style
def bukti_id {P : Prop} : P -> P := fun h => h

-- composition of two implications
def bukti_komposisi {P Q R : Prop} (h1 : P -> Q) (h2 : Q -> R) : P -> R :=
  fun hp => h2 (h1 hp)

#check bukti_id
#check bukti_komposisi
#check Prop
\end{verbatim}

The output of the last command, as it stands, is \texttt{Prop : Type}: the set of propositions is
itself a type. The two lines before it print the types of the two definitions. Note that
\texttt{\#check} does not prove anything; it only asks for the type, and that is also what the
kernel of the system does when it checks.

\begin{proposisi}
For every $P:\mathtt{Prop}$ we have $P\to P$. For $P,Q,R:\mathtt{Prop}$ we have
$(P\to Q)\to(Q\to R)\to P\to R$.
\end{proposisi}

\begin{proof}
The term $\lambda h.\,h$ has type $P\to P$: in the context $h:P$, the variable $h$ itself is a
proof of $P$, so abstraction over $h$ produces a function from $P$ to $P$. For the second statement,
let $h_{1}:P\to Q$, $h_{2}:Q\to R$, and $h_{p}:P$. The application $h_{1}\,h_{p}$ has type $Q$, then
the application $h_{2}\,(h_{1}\,h_{p})$ has type $R$. Three abstractions in succession over $h_{p}$,
$h_{2}$, and $h_{1}$ give a term of type $(P\to Q)\to(Q\to R)\to P\to R$.
\end{proof}

The tree in Figure~\ref{fig:pohon-tipe} writes that second proof as a derivation tree: three leaves,
two application stems, one root. The rightmost leaf is a hypothesis that is given, not a result of
derivation.

\begin{figure}[htbp]
\centering
\begin{tikzpicture}[font=\small,
  daun/.style={draw=layergrey,fill=layergrey!6,rounded corners=2pt,inner sep=4pt},
  hip/.style={draw=layergrey,dashed,fill=layergrey!3,rounded corners=2pt,inner sep=4pt},
  batang/.style={thick,layergrey}]
  \node[daun] (h1) at (-3.7,2.7) {\texttt{h1 : P -> Q}};
  \node[daun] (hp) at (-0.6,2.7) {\texttt{hp : P}};
  \node[daun] (ap1) at (-2.2,1.35) {\texttt{h1 hp : Q}};
  \node[hip]  (h2) at (2.7,1.35) {\texttt{h2 : Q -> R}};
  \node[daun] (akar) at (0,0) {\texttt{h2 (h1 hp) : R}};
  \draw[batang] (h1.south) -- (ap1.north west);
  \draw[batang] (hp.south) -- (ap1.north east);
  \draw[batang] (ap1.south) -- (akar.north west);
  \draw[batang] (h2.south) -- (akar.north east);
\end{tikzpicture}
\caption{The type tree for the term \texttt{h2 (h1 hp)}: three leaves and two application stems, with
a root of type $R$. The dashed boxes are hypotheses that are given, not results of derivation.}
\label{fig:pohon-tipe}
\end{figure}

There is one property that must be mentioned because it explains why proofs are used as
certificates rather than as data. The Lean kernel treats any two proofs of the same proposition as
definitionally equal. This property is called \emph{proof irrelevance}, and its consequence is that
a proof carries no information other than ``this statement is true''. For that reason a proof can be
thrown away once it has been checked, and what remains is only the status of the statement.

One warning is installed from the start. The Curry--Howard correspondence holds in full for
constructive logic. As soon as we use the law of excluded middle, $\mathtt{Classical.em}:P\vee\neg
P$, we add an axiom, and a proof that uses it stops being computational: it remains valid, but it is
no longer read as a program that produces a witness. Many theorems of the large libraries carry the
axiom of choice in their list of axioms, and the last part of this chapter shows how to ask Lean for
that list directly.

\section{Quantifiers: introduction and elimination rules}

\tingkat{intermediate}

An interesting statement is almost always quantified.\\index{quantifier} ``For every natural
number'' and ``there is a natural number'' are not decorations of language; both change the kind of
object that must be supplied. In the type view, that difference becomes the difference between a
function and a pair.

\kamus{quantifier}{kuantor}{The words ``every'' and ``there is''; the first demands a recipe for all
cases, the second demands one example together with its proof.}

The rule for the universal quantifier is entirely the same as the rule for a function. To prove
$\forall x,\ P(x)$, we take an arbitrary $x$ and prove $P(x)$ for that $x$; that is abstraction. To
use the hypothesis $\forall x,\ P(x)$, we apply that hypothesis to an object; that is application.
The rule for the existential quantifier is the rule for a pair: to prove $\exists x,\ P(x)$ we
supply one witness $a$ together with a proof of $P(a)$; to use an existential hypothesis we unpack
that pair into the witness and its proof.

\begin{center}
\begin{tabular}{lll}
\hline
Form of the statement & How to prove it & How to use it \\
\hline
$\forall x,\ P(x)$      & \texttt{intro x}              & apply to an object \\
$\exists x,\ P(x)$      & \texttt{exact Exists.intro a h} & \texttt{cases} on the hypothesis \\
$P\wedge Q$             & \texttt{constructor}          & \texttt{h.1}, \texttt{h.2} \\
$P\vee Q$               & \texttt{left}, \texttt{right} & \texttt{cases} on the hypothesis \\
$\neg P$                & \texttt{intro h} then reach $\bot$ & apply to a proof of $P$ \\
\hline
\end{tabular}
\end{center}

\begin{tangan}
Prepare ten blank cards and one marker. On five cards write introduction rules: ``take an
arbitrary'', ``supply a witness'', ``pair two proofs'', ``choose left'', ``choose right''. On the
other five cards write use rules: ``apply to an object'', ``unpack the pair'', ``take the first
component'', ``split into cases'', ``use it on the opponent's proof''.

Lay one goal card at the end of the table, for example ``every natural number is zero or a
successor''. Build a chain toward the goal card: the first card must be an introduction rule whose
shape fits the goal, not a use rule. Every following card may only work on the part that is still
open. A card that has no remainder to work on is pushed to the edge of the table; in Lean that card
will never be used.

What you must see: the chain gets shorter toward the goal, and there are never two cards working on
the same remainder. If that happens, your chain branches, and the machine demands two proofs, not
one.
\end{tangan}

The following example shows both at work on the natural numbers. The first statement is a tautology
that requires no reasoning at all; the second statement says that every natural number is zero or
the successor of something, and its proof requires case splitting on the shape of the number.

\begin{verbatim}
-- universal quantifier: take an arbitrary n, then compute
theorem nol_kanan : forall n : Nat, n + 0 = n := by
  intro n
  rfl

-- existential quantifier: supply a witness, then prove its property
theorem ada_tiga : exists n : Nat, n = 3 := by
  exact Exists.intro 3 rfl

-- every natural number is zero or a successor
theorem nol_atau_suksesor (n : Nat) : n = 0 \/ exists m : Nat, n = m + 1 := by
  cases n with
  | zero => left; rfl
  | succ k => right; exact Exists.intro k rfl
\end{verbatim}

The three theorems are checked without any additional axiom, and that can be confirmed with the
command \texttt{\#print axioms}, which we will use many times. Note also that
\texttt{Exists.intro k rfl} uses \texttt{rfl} as a proof of equality: in the case
\texttt{succ k}, the goal is $k+1=k+1$, which is true by definition.

\begin{catatan}
The reverse direction does not hold automatically. From $\exists x,\ P(x)$ we may not conclude
$\forall x,\ P(x)$, and Lean will reject every attempt in that direction. The mistake of swapping
the order of quantifiers is the most frequent specification error of all, and it slips through
testing without complaint because all test samples satisfy it. The part on limits at the end of this
chapter discusses its consequences.
\end{catatan}

\section{Tactics: turning the goal into a term}

\tingkat{basic}

Writing proof terms directly quickly becomes unreadable, because a deep term has to be read from the
inside out. That is why Lean provides a tactic mode.\\index{tactic} After the word \texttt{by}, we
do not write a term, but a sequence of commands that change the \emph{goal state}. The goal state is
printed as a list of hypotheses, then a separator line, then the target.

\kamus{tactic}{taktik}{A step-by-step cooking instruction, not a complete recipe that has to be
memorized all at once.}

\begin{verbatim}
-- the goal state is printed after every tactic; the line "|-" reads "therefore"
theorem imp_trans {P Q R : Prop} (h1 : P -> Q) (h2 : Q -> R) : P -> R := by
  intro hp
  trace_state          -- P Q R : Prop, h1 : P -> Q, h2 : Q -> R, hp : P, |- R
  apply h2
  trace_state          -- the target becomes |- Q
  exact h1 hp
\end{verbatim}

The output of \texttt{trace\_state} on this test machine is the following two blocks. Lean prints
the arrow as a Unicode character; here everything is rewritten in ASCII so as to be consistent with
the whole chapter.

\begin{verbatim}
P Q R : Prop
h1 : P -> Q
h2 : Q -> R
hp : P
|- R

P Q R : Prop
h1 : P -> Q
h2 : Q -> R
hp : P
|- Q
\end{verbatim}

The tactic \texttt{intro} takes a target of implication shape and moves its antecedent into the
hypotheses; that is exactly function abstraction. The tactic \texttt{exact} closes the target with a
term whose type fits exactly. The tactic \texttt{apply} matches the conclusion of a lemma against
the target, then leaves the remaining arguments as new goals. These three tactics are already enough
for all the pure-logic proofs in this chapter, and all three are surface forms of the same rule.

Figure~\ref{fig:dua-keadaan} compresses those two states into two boxes. What changes is only the
line below the bar: the target shrinks from $R$ to $Q$. The list of hypotheses does not grow by a
single line, because \texttt{apply} does not create a new proof; it only lowers the demand.

\begin{figure}[htbp]
\centering
\begin{tikzpicture}[font=\small,
  kotak/.style={draw=layergrey,fill=layergrey!6,rounded corners=3pt,align=left,inner sep=6pt,
                text width=3.9cm},
  panah/.style={-{Stealth[length=2.2mm]},thick,layerteal}]
  \node[kotak] (awal) at (0,0) {\textbf{state 1}\\[3pt]
    \texttt{P Q R : Prop}\\
    \texttt{h1 : P -> Q}\\
    \texttt{h2 : Q -> R}\\
    \texttt{hp : P}\\[3pt]
    \rule{3.6cm}{0.4pt}\\
    $\vdash$ \texttt{R}};
  \node[kotak] (akhir) at (6.7,0) {\textbf{state 2}\\[3pt]
    \texttt{P Q R : Prop}\\
    \texttt{h1 : P -> Q}\\
    \texttt{h2 : Q -> R}\\
    \texttt{hp : P}\\[3pt]
    \rule{3.6cm}{0.4pt}\\
    $\vdash$ \texttt{Q}};
  \draw[panah] (awal.east) -- node[above,font=\scriptsize]{\texttt{apply h2}} (akhir.west);
\end{tikzpicture}
\caption{The two goal states in the proof of \texttt{imp\_trans}: the hypotheses stay at four lines,
the target shrinks from $R$ to $Q$ after one tactic. The second box is still open; it is closed by
\texttt{exact h1 hp}.}
\label{fig:dua-keadaan}
\end{figure}

\begin{center}
\begin{tabular}{lll}
\hline
Tactic & What it does & When to use it \\
\hline
\texttt{intro}     & antecedent becomes a hypothesis   & target of shape $P\to Q$ \\
\texttt{exact}     & closes the target with a term    & the term is already available \\
\texttt{apply}     & lemma conclusion matches target & the lemma ends the same way \\
\texttt{rw}        & rewrites with an equation & there is a hypothesis $a=b$ \\
\texttt{simp}      & simplifies repeatedly      & many trivial steps \\
\texttt{cases}     & splits the target by constructor & case analysis is needed \\
\texttt{induction} & base case and inductive step     & recursive target over \texttt{Nat} \\
\hline
\end{tabular}
\end{center}

The tactic \texttt{rw} rewrites the target using an equation. It works in the direction that is
given, and that direction matters: rewriting in the wrong direction will not find the pattern that
is being looked for. Note the two attempts below, one of which succeeds and one of which is
rejected.

\begin{verbatim}
theorem geser (a b : Nat) (h : a = b) : a + 1 = b + 1 := by
  rw [h]        -- the target becomes b + 1 = b + 1, then it is closed by definition
  rfl

-- commutativity of addition does NOT hold by definition, so rfl fails
theorem tukar_salah (n : Nat) : n + 1 = 1 + n := by
  rfl           -- error: The rfl tactic failed

-- the correct version: rewrite with the commutativity lemma
theorem tukar_benar (n : Nat) : n + 1 = 1 + n := by
  rw [Nat.add_comm]
\end{verbatim}

The failure in the middle is a lesson, not an accident. Addition on the natural numbers is defined
by recursion on the second argument, so $n+1$ reduces to $\mathtt{succ}\,n$ whereas $1+n$ reduces to
$\mathtt{succ}\,(0+n)$, and the two stop at different forms. What we call the ``commutative
property'' at school is a theorem, not a definition, and Lean refuses to treat it as if it were a
definition. This is the most noticeable difference between reading and writing formal proofs: every
thing that we usually skip over has to be named.

The tactic \texttt{cases} splits the target according to the constructors of the type. On the
natural numbers it gives two cases, zero and successor. The tactic \texttt{simp} closes a goal by
rewriting with a collection of lemmas that are marked; it is useful but also dangerous, because its
output depends on the collection of lemmas that is currently active. A proof that begins with
\texttt{simp} without arguments can stop working when the library changes, and that is why in this
chapter \texttt{simp} is always mentioned together with the list of lemmas it uses.

One warning about tactics that are not here. Names such as \texttt{linarith}, \texttt{ring},
\texttt{norm\_num}, and \texttt{induction'} come from the Mathlib library, not from the Lean kernel.
All the code fragments in this chapter were run with Lean 4.10.0 without Mathlib, so those tactics
are not available and are not used; in their place we use \texttt{omega} (which is in the kernel)
and \texttt{simp} with an explicit list of lemmas. This difference is not a matter of style: the
library adds thousands of lemmas and several tactics, and every addition also determines what can be
closed in one line.

\begin{untukanak}
Take a pile of small stones. You arrange one stone into the smallest square. Then you add three
stones around it, and the result is a larger square. Add five more stones, and the result is a
larger square still. Add seven, nine, eleven, and so on, always increasing by two, and every time
the result is still a tidy square.

Now the question: will this go on being so if we add up to a hundred times, up to a million times?
You cannot try all of them. Nor do you need to. If you can show that the pattern of adding stones
always closes one corner and one side of the square neatly, then one small explanation holds for
every size, including sizes that you will never be able to arrange by hand.

The lazy computer does not want to try a million cases, and the overconfident computer does not want
to check anything. In this chapter we will find out why there is a third way: explain the pattern
once, then ask the machine to check that explanation down to its last letter.
\end{untukanak}

\section{A complete inductive proof: the sum of odd numbers}

\tingkat{advanced}

We are going to prove one arithmetic statement that can be checked entirely by machine, namely that
the sum of the first $n$ odd numbers equals $n^{2}$. This statement was chosen for three reasons:
the formula can be computed with Python without Lean, the proof requires induction, and the
algebraic step is small enough that every tactic can be accounted for.

\kamus{induction}{induksi}{Explaining a pattern with only two things: the smallest case, and how to
climb one step from any case.}

Before writing code, we state the statement and its proof in ordinary language. The function is
defined recursively: the sum of the first zero odd numbers is zero, and the sum of the first $n+1$
odd numbers is the sum of the first $n$ odd numbers plus the odd number number $n+1$, namely
$2n+1$.

\begin{definisi}[Sum of odd numbers]
The function $S:\mathbb{N}\to\mathbb{N}$ is defined by $S(0)=0$ and $S(n+1)=S(n)+(2n+1)$ for every
$n\in\mathbb{N}$.\\index{induction}
\end{definisi}

The number $2n+1$ is not an arbitrary choice: the $n$th odd number really does have that form, and
$2n+1$ is also the difference between $(n+1)^{2}$ and $n^{2}$. That last equality is the heart of the
proof.

\begin{lemma}
For every $k\in\mathbb{N}$ we have $(k+1)^{2}=k^{2}+2k+1$.
\end{lemma}

\begin{proof}
$(k+1)^{2}=(k+1)(k+1)=k(k+1)+(k+1)=k^{2}+k+k+1=k^{2}+2k+1$. Each step uses distributivity and
associativity of addition on the natural numbers.
\end{proof}

\begin{teorema}
For every $n\in\mathbb{N}$ we have $S(n)=n^{2}$.
\end{teorema}

\begin{proof}
By induction on $n$. For $n=0$, both sides have the value $0$. Suppose the statement is true for
$k$, namely $S(k)=k^{2}$. Then $S(k+1)=S(k)+(2k+1)=k^{2}+2k+1=(k+1)^{2}$, with the lemma above at
the last step. So the statement is true for $k+1$. By the principle of induction, the statement is
true for every $n\in\mathbb{N}$.
\end{proof}

The proof above is a complete informal proof, in the sense that no step is skipped. Now we write it
for the machine. Note that the recursive definition and the theorem are almost word for word what we
have just written, whereas the proof turns into a sequence of tactics.

\begin{verbatim}
-- recursive definition: two equations, without ":=" on the first line
def oddSum : Nat -> Nat
  | 0     => 0
  | n + 1 => oddSum n + (2 * n + 1)

theorem oddSum_sq (n : Nat) : oddSum n = n * n := by
  induction n with
  | zero => rfl
  | succ k ih =>
    rw [oddSum]        -- unfold the definition: oddSum (k+1) = oddSum k + (2*k+1)
    rw [ih]            -- use the induction hypothesis: oddSum k = k*k
    simp [Nat.succ_eq_add_one, Nat.mul_succ, Nat.add_mul, Nat.one_mul]
    omega              -- close the linear remainder
\end{verbatim}

The line \texttt{induction n with} generates two goals. The first goal is labelled \texttt{zero} and
reads $S(0)=0\cdot 0$, and it is closed by \texttt{rfl} because both sides reduce to zero. The
second goal is labelled \texttt{succ k} and carries one extra hypothesis named \texttt{ih}, namely
$S(k)=k\cdot k$. The name \texttt{ih} is chosen by us; its role is what matters.

The line \texttt{rw [oddSum]} unfolds the definition on the left-hand side. After that line the goal
reads $S(k)+(2k+1)=(k+1)(k+1)$. The line \texttt{rw [ih]} replaces $S(k)$ with $k\cdot k$, so the
goal becomes $k\cdot k+(2k+1)=(k+1)(k+1)$.

At this point we stand exactly where the lemma above does its work, but the machine does not yet know
that $(k+1)(k+1)$ can be written as $k\cdot k+2k+1$. \texttt{simp} with four explicit lemmas changes
the shape of the right-hand side into $k\cdot k+k+(k+1)$, which is a shape that is entirely linear in
$k$. The last line, \texttt{omega}, closes a linear goal of that kind.

It must be stressed that \texttt{omega} is not magic and is not a hidden axiom either. It is a
decision procedure for linear arithmetic over the integers, and it accepts only goals all of whose
operations are linear. If we call it before \texttt{simp}, the goal still contains both $k\cdot k$
and $(k+1)(k+1)$. On this test machine, the temporary experiment file named \texttt{odd.lean}
produces a rejection:

\begin{verbatim}
odd.lean:11:4: error: omega could not prove the goal:
a possible counterexample may satisfy the constraints
  a - b + 2*c <= -2
  ...
\end{verbatim}

That message is honest and useful: it names a possible counterexample, and precisely because of that
we know that what fails is the shape of the goal, not the truth of the statement. The order of
tactics is therefore not a matter of taste; it changes the kind of goal that the decision procedure
faces.

\begin{catatan}
If \texttt{simp} above is run without \texttt{omega}, the remaining goal has the shape
$k\cdot k+(1+2k)=k+(k+(k\cdot k+1))$. Mathematically that is true, but it is no longer linear, and
the simplifier does not close it. On this test machine that experiment ends with
\texttt{unsolved goals}. This is the most common shape of a failed formal proof: not that the
statement is false, but that the shape of the goal does not fit the tool in use.
\end{catatan}

\begin{akibat}
For every $n\in\mathbb{N}$, the sum of the first $n$ odd numbers is a perfect square.
\end{akibat}

That corollary is only a repetition of the theorem in other words, and that is precisely its use: it
restates the result in the language used in this part, so that a reader who meets the formula
elsewhere can recognize it.

\section{Structures, typeclasses, and what is really guaranteed}

\tingkat{advanced}

A certificate is rarely a single loose theorem. Usually it is a collection of data together with
claims about that data: thresholds, tables, file names, and several statements. Lean provides two
tools for grouping things of this kind, namely \emph{structure} and \emph{typeclass}.

\kamus{typeclass}{type class}{A list of capabilities that a type must have in order to be usable in
a certain place, like the licensing requirement for driving a particular kind of vehicle.}

\begin{definisi}[Structure and typeclass]
A \emph{structure} is a type whose fields are named, so that a value of it is built by naming the
fields. A \emph{typeclass} is a structure that is used to look for an \emph{instance}
automatically: when a definition requires a certain capability from a type, Lean searches for a
matching instance in the environment that is currently active.\\index{typeclass}
\end{definisi}

The practical difference is simple. A structure is used when the data itself is what matters, for
example a certificate with a claim field and a threshold field. A typeclass is used when what
matters is the capability: something that can be used as a proof, something whose norm can be
computed, something that can be compared. In both, the checking remains equally strict: every field
must be filled, and every instance must pass the type check.

\begin{verbatim}
-- structure: data with named fields
structure Sertifikat where
  klaim  : String
  ambang : Nat

-- typeclass: a capability that is looked up automatically
class Terperiksa (a : Type) where
  jumlah_teorema : a -> Nat

instance : Terperiksa Sertifikat where
  jumlah_teorema _ := 1

def contoh : Sertifikat := { klaim := "oddSum", ambang := 100 }

#eval (Terperiksa.jumlah_teorema contoh)   -- prints 1
\end{verbatim}

Now the part that is most often skipped in popular talk about Lean, namely the question of what the
acceptance of a theorem actually guarantees. Four things are guaranteed, and four other things are
not.

Figure~\ref{fig:alur-periksa} puts those three stages on one line: the statement is written, the
term is submitted, the kernel decides. There are only two outputs; there is no third box that says
``perhaps''.

\begin{figure}[htbp]
\centering
\begin{tikzpicture}[font=\small,
  tahap/.style={draw=layergrey,fill=layergrey!6,rounded corners=3pt,align=center,inner sep=5pt,
                text width=2.9cm,minimum height=1.7cm},
  panah/.style={-{Stealth[length=2.2mm]},thick,layergrey},
  terima/.style={draw=layerjade,fill=layerjade!8,rounded corners=3pt,align=center,inner sep=5pt,
                 text width=2.9cm},
  tolak/.style={draw=accent,fill=accent!6,rounded corners=3pt,align=center,inner sep=5pt,
                text width=2.9cm}]
  \node[tahap] (s1) at (0,0) {\textbf{1. statement}\\$P$ has type \texttt{Prop}};
  \node[tahap] (s2) at (3.9,0) {\textbf{2. term}\\$p$ is submitted for $P$};
  \node[tahap] (s3) at (7.8,0) {\textbf{3. kernel}\\check whether $p:P$};
  \draw[panah] (s1) -- (s2);
  \draw[panah] (s2) -- (s3);
  \node[terima] (ya) at (5.4,-2.9) {\textbf{accepted}\\there is a term of type $P$};
  \node[tolak] (tidak) at (9.9,-2.9) {\textbf{rejected}\\there is no term};
  \draw[panah] (s3.south west) -- node[above left,font=\scriptsize]{matches} (ya.north);
  \draw[panah] (s3.south east) -- node[above right,font=\scriptsize]{does not match} (tidak.north);
\end{tikzpicture}
\caption{The checking flow: three stages (statement, term, kernel) and two outputs, accepted or
rejected. There is no third output; a file that stops in the middle is reported as a goal that is
not yet closed, not as ``almost correct''.}
\label{fig:alur-periksa}
\end{figure}

What is guaranteed: (i) there is a term whose type is exactly the statement that was written; (ii)
that term is checked by the kernel, not by the tactic that produced it; (iii) the checking uses only
definitions, other theorems that have already been checked, and recorded axioms; (iv) the list of
axioms that are used can be asked for at any time with \texttt{\#print axioms}. Point (iv) is the
most useful instrument in this whole chapter, because it turns ``just believe me'' into ``check it
yourself''.

What is not guaranteed: (i) that the statement says what we mean; (ii) that the definitions model
objects in the real world; (iii) that the kernel, the parser, and the hardware are free of errors;
(iv) that the tactics in use do not add hidden assumptions. Point (iv) is not a theoretical flaw: it
can be checked, and the result is often surprising.

\begin{contoh}[Two ways to prove the same thing, two lists of axioms]
Consider two proofs of the statement $S(10)=100$, both of them correct. The first uses
\texttt{decide}, which runs the computation inside the kernel. The second uses
\texttt{native\_decide}\\index{native-decide}, which compiles the computation into machine code and
runs it. The two statements are the same, but their lists of axioms differ, and that difference is
visible from the command \texttt{\#print axioms}.
\end{contoh}

\begin{verbatim}
def oddSum : Nat -> Nat
  | 0     => 0
  | n + 1 => oddSum n + (2 * n + 1)

theorem cek_native : oddSum 10 = 100 := by native_decide
theorem cek_kernel : oddSum 10 = 100 := by decide

#print axioms cek_native
#print axioms cek_kernel
\end{verbatim}

The output on this test machine, as it stands:

\begin{verbatim}
'cek_native' depends on axioms: [Lean.ofReduceBool]
'cek_kernel' does not depend on any axioms
\end{verbatim}

The first line is the heart of the whole discussion about \texttt{native\_decide}. The computation
is not carried out by the kernel, but by compiled code, and Lean records that trust as an axiom
named \texttt{Lean.ofReduceBool}. In practice this is not a disaster: the Lean compiler is used to
compile Lean itself every day. In principle it is a large difference, because the boundary of trust
moves from the small kernel to a long chain of compiler, linker, and processor. A certificate built
on top of \texttt{native\_decide} is therefore weaker than a certificate built on top of
\texttt{decide}, even though both print ``accepted'' without complaint.

In the induction theorem above, the list of axioms is not empty either. The command
\texttt{\#print axioms} for \texttt{oddSum\_sq} prints \texttt{[propext, Quot.sound]}. Both are
standard axioms used for classical reasoning and for quotients, and both appear because
\texttt{omega} proves the goal by contradiction inside the integers. This is not an error; it is a
fact that ought to be mentioned when a certificate is used as the basis of a claim.

\section{Limits: rejecting wrong proofs, \texttt{sorry}, and axioms}

\tingkat{advanced}

The ability to reject is as important as the ability to accept. The following file deliberately
contains errors of different kinds, and the output is Lean's answer to each of them.

\begin{verbatim}
-- File: Salah.lean
theorem salah_arit : 2 + 2 = 5 := by decide

theorem salah_rfl (n : Nat) : n + 1 = 1 + n := by rfl

axiom magic : 1 = 2
theorem from_magic : 1 = 2 := magic
#print axioms from_magic

theorem belum_selesai (n : Nat) : n + 0 = n := by sorry
\end{verbatim}

The actual output, cut at the long part:

\begin{verbatim}
Salah.lean:2:37: error: tactic 'decide' proved that the proposition
  2 + 2 = 5
is false
Salah.lean:4:50: error: The rfl tactic failed. Possible reasons:
- The goal is not a reflexive relation ...
- The arguments of the relation are not equal.
...
n : Nat
|- n + 1 = 1 + n
'from_magic' depends on axioms: [magic]
Salah.lean:10:8: warning: declaration uses 'sorry'
\end{verbatim}

Three kinds of rejection appear here, and all three differ in nature. The first is the strongest
rejection: \texttt{decide} does not merely fail to close the goal, it proves that the proposition is
false. For a finite statement, failure and refutation are the same thing. The second is a subtler
rejection: \texttt{rfl} refuses because the two sides are not equal by definition, even though the
statement is true. Here Lean does not blame the statement, it blames the tool in use; the message
even suggests an explicit reflexivity lemma. The third is not a rejection at all.

The line \texttt{axiom magic : 1 = 2} is accepted without complaint. Lean does not check whether the
content of an axiom is reasonable; it only records it.\\index{axiom} After that,
\texttt{from\_magic} is accepted as well, and the whole of arithmetic can be taken apart with
\texttt{False.elim} from there. The only danger sign is the \texttt{\#print axioms} line, which
mentions \texttt{[magic]} calmly, without a warning and without the color red. This is the practical
reason why inspecting the list of axioms is a compulsory habit: the system will not reproach us.

\kamus{axiom}{aksioma}{A rule that we ask to be accepted without proof, like an opening assumption
in a debate; legitimate, but it has to be named.}

The last line of that file uses \texttt{sorry}, and its output is only the warning
\texttt{declaration uses 'sorry'}.\\index{sorry} The file is still processed, the theorem still
``exists'', and anyone who reads the output carefully will see the word \texttt{warning}, not
\texttt{error}. In a large library, one \texttt{sorry} left behind among thousands of theorems is
almost invisible. Lean provides the answer, and the answer is the same as the answer for axioms:
\texttt{\#print axioms} prints \texttt{[sorryAx]}. Lean's mathematics library, namely Mathlib, uses
a stricter policy: a file that contains \texttt{sorry} is rejected by the automatic check before it
is accepted. That policy deserves to be copied by anyone who publishes certificates.

\kamus{sorry}{sorry}{A note saying ``I will do it later'' that makes the work look finished; the
machine accepts it, but records it as a debt.}

Now the most important limit, and the one least often mentioned. Formal verification guarantees that
the term fits the statement. It does not guarantee that the statement fits the intention. If the
specification is too weak, the proof can be perfect and useless.\\index{specification} The following
example shows this in its most naked form: the first statement says ``there is a small perturbation
that does not change the output'', the second statement says ``every small perturbation does not
change the output''. The first can be proved for \emph{all} functions, including the worst function
imaginable, because the witness $\delta=0$ is always available.

\begin{verbatim}
def identitas (n : Nat) : Nat := n

-- weak specification: THERE IS a delta that does not change the output (for all f)
theorem ada_delta (f : Nat -> Nat) (x : Nat) :
    exists delta : Nat, f (x + delta) = f x := by
  exact Exists.intro 0 rfl

-- the strong version is false for f = identity, so it cannot be proved
theorem ada_penyangkal :
    exists x : Nat, exists delta : Nat,
      Not (identitas (x + delta) = identitas x) := by
  exact Exists.intro 0 (Exists.intro 1 (by decide))
\end{verbatim}

Both theorems pass without additional axioms, and that is exactly the problem. The first theorem says
nothing about $f$; it only says that $\delta=0$ is a perturbation that changes nothing, and that
holds for every function, including a classifier that is wrong one hundred percent of the time. A
certificate of this kind passes every check and protects nothing. The second theorem shows the other
side: the statement we actually want, namely the universally quantified version, is false for the
identity function, so no proof can possibly be found.

\begin{tantangan}
A certificate contains one line that looks impossible to be wrong: for every natural number $a$ and
$b$ we have $(a-b)+b=a$. On paper the reason is simple: subtract, add again, and you are back where
you started. Write that line in Lean and close it with \texttt{omega}; the checker rejects it, and
the rejection is not because the tactic is not clever enough.

Two questions. Why is that statement not true for all $a$ and $b$? And what additional condition
makes it true while keeping it checkable by machine?

Hint 1: do not trust the parentheses; ask what happens when $b$ is larger than $a$. Hint 2:
subtraction on the natural numbers stops at zero, so $0-1=0$, hence the left-hand side becomes
$0-1+1=1$ whereas the right-hand side becomes $0$.

If you get stuck, go back to the section ``Tactics: turning the goal into a term'', where $n+1$ and
$1+n$ stop at different forms.
\end{tantangan}

Before reading further, try to answer on your own. The output of this test machine, as it stands, is
the comparison for your answer: the second file contains the same statement with the hypothesis
$b\le a$.

\begin{verbatim}
$ lean t.lean
t.lean:1:49: error: omega could not prove the goal:
a possible counterexample may satisfy the constraints
  d >= 0
  c >= 0
  c - d >= 1
where
 c := b
 d := a
$ lean t2.lean
'potong_benar' depends on axioms: [propext, Quot.sound]
\end{verbatim}

The first output mentions the constraint $b-a\ge 1$, that is a counterexample with $b>a$; at $a=0$
and $b=1$ the left-hand side has the value $1$, the right-hand side $0$. The hypothesis $b\le a$
makes the same line pass, with the two standard axioms that we already know from \texttt{omega}. The
shape of the repair is not an additional tactic, but a condition.

\kertas{The claim that may be made from this chapter: the Lean code fragments that are shown have
been run with Lean 4.10.0 without Mathlib on the machine of writing, and their output is copied as
it stands; Unicode output is reprinted in ASCII and cuts are marked with three dots. What was not
run: the Mathlib tactics (\texttt{linarith}, \texttt{ring}, \texttt{norm\_num},
\texttt{induction'}) and the Mathlib library itself, because it is not installed. What is not proved
here and is not claimed: any certificate about neural networks, language models, or robustness to
attack. Chapter 28 does not contain a single proof of that kind.}

These limits have a direct connection with Chapter 27 and Chapter 29. An honest evaluation of a model
demands that we state the distribution in which the numbers are measured; formal verification
demands that we state the proposition that is proved. Both are the same disease in a different form:
a claim that does not state its conditions. The certificate ``this model is proved safe'' without a
formal statement behind it is as useless as ``99 percent accuracy'' without a note on the test data.
In Chapter 29 and Chapter 30, when we discuss attacks and defenses, the first question that must be
put to every robustness certificate is: how many quantifiers does it have, and in what space do those
quantifiers run?

\section{Worked example}

\tingkat{intermediate}

The numbers in this part are computed, not estimated. The first column is $n$, the second column is
the $n$th odd number that is added at that step, the third column is $S(n)$, and the fourth column is
$n^{2}$.

\begin{center}
\begin{tabular}{rrrr}
\hline
$n$ & term $n$ & $S(n)$ & $n^{2}$ \\
\hline
0  & --  & 0     & 0 \\
1  & 1   & 1     & 1 \\
2  & 3   & 4     & 4 \\
5  & 9   & 25    & 25 \\
10 & 19  & 100   & 100 \\
20 & 39  & 400   & 400 \\
100 & 199 & 10000 & 10000 \\
\hline
\end{tabular}
\end{center}

The value $S(100)=10000$ is worth checking by hand once, because it shows that the formula is not a
coincidence for small numbers: the sum of a hundred odd numbers, from $1$ to $199$, equals
$100\cdot 100$. A check with Python over all values $n<5000$ gives \texttt{True}, and a check of the
algebraic identity $(k+1)^{2}=k^{2}+2k+1$ for all $k<100000$ also gives \texttt{True}. Neither check
is a proof, but both rule out the possibility that the formula being proved has a typo.

There is one more number that is more informative than the table above, namely the size of the proof
after elaboration. The tactic script for \texttt{oddSum\_sq} is eight lines long. When Lean prints
the term that results from elaborating it with \texttt{\#print}, the result is 10,783 characters in
149 lines. That means one line of tactics expands on average into about nineteen lines of term. No
human reads that term; what reads it is the kernel, and that is the division of labor that is the
heart of the whole method.

The last number concerns the limit of direct checking. In the default setting, the recursion-depth
limit has the value 1000. With that limit, the tactic \texttt{decide} succeeds in proving
$S(n)=n^{2}$ only for $n\le 122$. The first experiment that fails is $n=123$, and the message that
appears mentions that the maximum recursion depth has been reached. If the limit is raised to 100000,
the value $n=2000$ is proved successfully in 0.48 seconds. At $n=20000$ the check fails because the
recursion is too deep for the stack space. So direct checking is not an alternative that scales,
whereas an inductive proof does not care about the size of $n$.

\section{Computational lab}

\tingkat{advanced}

This lab runs two things: the Lean file that is the same as the one discussed above, and a
cross-check with Python. The Lean file is written without \texttt{import}, so it uses only the
kernel.

\begin{lab}
\begin{verbatim}
-- BilanganGanjil.lean : Lean 4.10.0, without Mathlib
def oddSum : Nat -> Nat
  | 0     => 0
  | n + 1 => oddSum n + (2 * n + 1)

theorem oddSum_sq (n : Nat) : oddSum n = n * n := by
  induction n with
  | zero => rfl
  | succ k ih =>
    rw [oddSum]
    rw [ih]
    simp [Nat.succ_eq_add_one, Nat.mul_succ, Nat.add_mul, Nat.one_mul]
    omega

#eval oddSum 0
#eval oddSum 5
#eval oddSum 10
#eval oddSum 100
#print axioms oddSum_sq
\end{verbatim}

\begin{verbatim}
def odd_sum(n):
    return sum(2 * i + 1 for i in range(n))

print([odd_sum(n) for n in range(11)])
print(odd_sum(20), odd_sum(100))
print(all(odd_sum(n) == n * n for n in range(5000)))
print(all((k + 1) * (k + 1) == k * k + 2 * k + 1 for k in range(100000)))
\end{verbatim}
\end{lab}

The commands and the output, as they stand:

\begin{verbatim}
$ lean --version
Lean (version 4.10.0, x86_64-unknown-linux-gnu, commit c375e19f6b65, Release)
$ time lean BilanganGanjil.lean
0
25
100
10000
'oddSum_sq' depends on axioms: [propext, Quot.sound]
time 0.14 s
$ python3 BilanganGanjil.py
[0, 1, 4, 9, 16, 25, 36, 49, 64, 81, 100]
400 10000
True
True
\end{verbatim}

The first four lines of the Lean output are the results of \texttt{\#eval}, and all four match the
table in the previous part. The fifth line is the list of axioms, and its contents must be read
carefully: \texttt{propext} and \texttt{Quot.sound} are not an error, but a consequence of the way
\texttt{omega} works. The Python output confirms the same thing from the other direction: the list of
the first eleven values is a list of perfect squares, $S(100)=10000$, and both exhaustive checks
print \texttt{True}. The time 0.14 seconds is the time of this test machine, and a number of that
kind depends on the hardware; what is meaningful is not the number, but the fact that the check
finishes before we have time to read the file.

A reader who wants to repeat this lab needs two things. First, Lean 4, which can be installed through
\texttt{elan}, the official version manager of the Lean project; after installation, the command
\texttt{lean --version} must print a version. Second, the file above is saved as
\texttt{BilanganGanjil.lean} and run with \texttt{lean BilanganGanjil.lean} from the same directory;
no \texttt{lake} project is needed because there is no import. To use \texttt{linarith} or
\texttt{ring}, the reader needs a project with Mathlib, and from that moment \texttt{lean} alone is
not enough; the command becomes \texttt{lake env lean}. This difference is often the first source of
confusion for a beginner, and mentioning it is better than letting it be discovered through an
\texttt{unknown tactic} message.

\begin{resep}
\begin{enumerate}
  \item Create a plain text file named \texttt{Uji.lean}; that extension is what makes the command
        \texttt{lean} recognize it as a source file.
  \item Write all \texttt{import}s on the topmost lines, one per line, before any declaration.
  \item Write the statement first, complete with the type of each variable, then the proof after
        \texttt{:= by}.
  \item Run \texttt{lean Uji.lean} from the directory of that file and read the output from the first
        line; the first error is usually the real one.
  \item Ask for \texttt{\#print axioms name\_of\_theorem} before calling the result proved.
\end{enumerate}
A warning about order: \texttt{import} must not be preceded by a declaration. One \texttt{def}
slipped above it makes Lean reject the whole file with the message
\texttt{invalid 'import' command}, and the file is not checked at all. The correct order is not the
author's taste; it is a condition for the file to be checkable.
\end{resep}

\begin{untukdosen}
The teaching order that works in class: start from the failure of \texttt{rfl} on $n+1=1+n$, not
from the definition. Students who see that something ``obvious'' is rejected will ask, and that
question is the doorway into Curry--Howard. Only after that come the definition of a proposition as a
type, then the first three tactics, then induction. Leave one full session for \texttt{\#print
axioms}; that is where students change from users into checkers.

Two misconceptions appear most often. First, ``Lean proves that my program is correct''. What is
true: Lean checks that a term has a certain type of statement. If the statement does not mention the
important thing, for example a quantifier over perturbations, the certificate is empty. Second,
``\texttt{simp} is a proof''. \texttt{simp} is a procedure whose output depends on the active
collection of lemmas, so a proof that depends on it can stop working when the library changes; that
is why we train students to write the list of lemmas explicitly.

A blackboard activity that works: ask two groups to write a specification for ``a classifier that is
robust to perturbation''. One group will almost always write an existential quantifier. Put both
formulations on the board, then prove the weak one for a constant function. The class will see for
itself that the sentence that is proved and the sentence that is meant are two different sentences.

A sign of readiness to go on: students can explain why $n+1=1+n$ requires a lemma, can write an
inductive proof without copying, and can name one axiom that their own theorem uses.
\end{untukdosen}

\section{Summary}

\begin{ringkasan}
\begin{itemize}
  \item An informal proof convinces a trained reader; a formal proof is a term whose type is its
        statement, checked by a small kernel that knows no mathematics.
  \item According to Curry--Howard, implication is a function type, conjunction is a pair, the
        universal quantifier is a dependent function, and $\bot$ is uninhabited.
  \item Tactics change the goal state; \texttt{intro}, \texttt{exact}, \texttt{apply},
        \texttt{rw}, \texttt{cases}, and \texttt{induction} are enough for basic arithmetic.
  \item Induction proves a statement for all $n$ with two goals, not by checking all $n$;
        \texttt{decide} by contrast stops scaling very quickly.
  \item \texttt{\#print axioms} is the main instrument of inspection. The tactic
        \texttt{native\_decide} leaves behind the axiom \texttt{Lean.ofReduceBool}. The tactic
        \texttt{sorry} leaves behind \texttt{sorryAx}. A homemade axiom leaves behind its own name.
  \item Verification guarantees that the term fits the statement, not that the statement fits the
        intention; a weak specification produces an empty certificate.
\end{itemize}
\end{ringkasan}

\section{Exercises}

\begin{latihan}
\begin{enumerate}
  \item Assess this claim: \emph{if a Lean file is accepted without an error message, then all the
        statements inside it are true}. Decide true, false, or not enough evidence, and name one
        command that can be used to check that claim for a particular file.

  \item Prove in Lean, without Mathlib, that for every $n$ we have $2\cdot S(n)=n\cdot(2n)$ with $S$
        as in this chapter. State first the auxiliary lemma that you need.

  \item Prove that the sum of the first $n$ even numbers, namely $0+2+\dots+(2n-2)$, equals
        $n(n-1)$. Write the informal proof, then write the Lean version and name the tactic that is
        used at every step.

  \item Compute with Python: the sum of the first 37 odd numbers, the sum of the first 250 odd
        numbers, and check that both are perfect squares. Include the code and its output.

  \item Compute with Python how many cases must be checked if one wants to prove $S(n)=n^{2}$ by
        direct checking for all $n\le 1000$, then compare it with the number of goals that an
        inductive proof generates.

  \item Explain why \texttt{rfl} closes the goal $n+0=n$ but does not close $0+n=n$. The hint lies
        in the way addition on the natural numbers is defined recursively; name which argument is
        reduced first.

  \item Write a weak specification and a strong specification for the claim ``this classifier does
        not change under a perturbation of size at most one''. Prove the weak version in Lean for an
        arbitrary function, then explain why the strong version cannot be proved for the identity
        function.

  \item Inspect with \texttt{\#print axioms} three theorems that you wrote yourself: one using
        \texttt{decide}, one using \texttt{omega}, one using \texttt{sorry}. Report the three lists
        of axioms and explain the origin of every name.

  \item A certificate contains \texttt{theorem aman : exists delta : Nat, f (x + delta) = f x} with
        \texttt{delta} as a free variable in the context. Explain two errors that may be present in
        that statement, and write the version that you consider correct.

  \item Prove in Lean that $\neg(P\wedge\neg P)$ for every proposition $P$. Hint: start with
        \texttt{intro h}, then take the components with \texttt{h.1} and \texttt{h.2}, and use the
        second component on the first one.
\end{enumerate}

Hints for the three hardest items. Item 2: the auxiliary lemma that is needed has the shape
$2\,(k\cdot k+(2k+1))=(k+1)\cdot(2(k+1))$, and after \texttt{rw} on the definition and on the
induction hypothesis, the remaining goal is linear so that \texttt{omega} is enough. Item 7: in the
weak version, the witness $\delta=0$ makes both sides equal by definition, so \texttt{rfl} closes it
immediately; in the strong version, the identity function gives a counterexample with $\delta=1$.
Item 9: note that the free \texttt{delta} can be chosen by the reader of the statement, not by the
attacker, so that statement asserts no robustness at all; the second error is the absence of a
quantifier over $x$.
\end{latihan}

\begin{untukpeneliti}
What is really proved in this chapter is only small arithmetic statements: $S(n)=n^{2}$ for all
natural numbers, a few propositional tautologies, and two statements about specifications that are
too weak. All of them are checked by the Lean 4.10.0 kernel without Mathlib, and their lists of
axioms are reported. What is popularized far beyond that: the phrase ``formally verified'' is often
used for ``checked by an automatic tool'' or even for ``tested exhaustively''. Only the first
produces an object that a third party can recheck.

Primary sources. Avigad, de Moura, Kong, and Ullrich, \emph{Theorem Proving in Lean 4}, the sections
``Propositions and Proofs'' (propositions as types, proof irrelevance, the axiom warning),
``Tactics'' (the goal state), and ``Induction and Recursion'' (\texttt{induction},
\texttt{sorryAx}). Avigad and Massot, \emph{Mathematics in Lean}, Chapters 2, 3.1, 5.2, and 8; the
Indonesian edition uses the same section numbering.

Open questions. This chapter does not contain a formal proof about neural networks. An interesting
robustness specification has the shape
$\forall x\,\forall\delta,\ \lVert\delta\rVert\le\varepsilon\to f(x+\delta)=f(x)$, that is a
statement over a continuous space with a function with hundreds of thousands of parameters.
Formalizing it demands real arithmetic that is decidable at that scale, and the Lean libraries today
do not provide it for trained models. The claim ``we verified the robustness of the model'' in a
paper usually refers to verification with interval abstraction or with SMT, not to a Lean proof; both
are legitimate, but their guarantees differ and must not be called by the same word.

Limits of validity. All the claims in this chapter hold for Lean 4.10.0 with its default kernel, and
the details of the output can change between versions. What does not change is the structure of the
argument: statement, term, kernel, list of axioms.
\end{untukpeneliti}

\peta{This chapter follows Avigad, de Moura, Kong, and Ullrich, \emph{Theorem Proving in Lean 4}, the
sections ``Propositions and Proofs'' (propositions as types, Curry--Howard, proof irrelevance, the
axiom warning), ``Quantifiers and Equality'', ``Tactics'' (the goal state, \texttt{intro},
\texttt{apply}, \texttt{exact}), and ``Induction and Recursion'' (\texttt{induction}, \texttt{sorry}
and \texttt{sorryAx}). The counterpart in Avigad and Massot, \emph{Mathematics in Lean}, is in
Chapter 2 (Basics, in particular 2.3 Using Theorems and Lemmas and 2.4 More Examples Using apply and
rw), Chapter 3.1 (Implications and Universal Quantifiers), Chapter 5.2 (Induction and Recursion), and
Chapter 8 (Hierarchies) for typeclasses; the Indonesian edition keeps the same section numbering. For
the connection with evaluation practice, see Chapter 27 of this book; for attacks that make
robustness specifications nontrivial, see Chapter 27 and Chapter 28.}

\chapter[Adversarial Machine Learning]{Adversarial Machine Learning: Attacks and Defenses}

In 2014 Goodfellow, Shlens, and Szegedy showed one picture of a panda. The GoogLeNet
network called it a panda with confidence $57.7$ percent. They added a perturbation
that the eye cannot see, of size $0.007$ on the scale $[0,1]$, which is about two
units out of 255 gray levels. The same network then called that picture a gibbon with
confidence $99.3$ percent. Nothing changed for a human. What changed was only the
direction of the perturbation: every pixel was shifted a little, and the direction of
that shift was chosen from the sign of the derivative of the error with respect to the
pixel.

Numbers like that are not an oddity of one model. In this chapter we compute for
ourselves a small case that can be rechecked. A two-layer network with 64 inputs
reaches accuracy $97.99$ percent on 597 test images of digits. With a perturbation
budget $0.15$ per pixel, an iterative attack of 40 steps lowers that accuracy to
$10.72$ percent; at the same budget, a random perturbation without a direction only
lowers it to $95.48$ percent. So the direction that the gradient chooses is not
decoration: it buys the attacker about eighty points of accuracy.

Chapter 25 taught us to compute the derivative of the error with respect to the
weights, and Chapter 27 taught us to measure accuracy honestly. Neither is enough.
Accuracy is an average over one data distribution; security is a question about the
worst point inside one allowed set. A high average and a bad worst point can live
together, and that is what makes an ordinary evaluation report misleading. This
chapter builds the vocabulary for that worst-case question, computes its first attack,
then separates the defenses that raise the standard from those that only hide the
gradient.

\begin{untukanak}
Imagine a building guard who recognizes people by small habits: how many buttons of
the shirt are fastened, the bag in the left hand or the right, and hat or no hat. So
far that method has worked. Everyone entitled to enter really does differ in those
three things. Then someone comes who knows exactly how the guard sees. He undoes one
button, moves the bag to the other hand, and puts on a hat. His face is the same, his
name is the same, he is the same person. The guard still lets him in, because
according to his notes a person with those three marks is the cleaning staff.

The question that makes this story important is not ``how do we fool the guard'',
but ``how do we know that the guard can be fooled''. If we test only with random
people in random clothes, the guard will pass almost always. The error appears only
when someone deliberately chooses those three marks. Testing with chance is not the
same as testing with malice.

A computer that recognizes images works like that guard, only with thousands of small
features at once. In this chapter we will find out why a feature that looks trivial
can be chosen deliberately, why a change that the eye cannot see can flip the answer,
and how to tell a defense that is truly strong from a defense that only leaves the
attacker confused.
\end{untukanak}

\section{Threat model: what the attacker may do}
\tingkat{basic}

The sentence ``the attacker changes the input a little'' cannot be tested before we
agree on the meaning of ``a little''. Every robustness claim is always relative to
that choice, so that choice is part of the claim, not a footnote.

\kamus{threat model}{model ancaman}{The complete list of what the attacker may do:
how large the perturbation is, what it knows, and what its goal is.}
\index{threat model}

\begin{definisi}[Adversarial example and threat model]
Let $f:\mathbb{R}^{d}\to\{1,\dots,K\}$ be a classifier and
$(\mathbf{x},y)$ a correctly classified example. An
\emph{adversarial example} is $\mathbf{x}'=\mathbf{x}+\boldsymbol{\delta}$ with
$f(\mathbf{x}')\neq y$ but $\boldsymbol{\delta}$ small according to an agreed
measure. The threat model is determined by four things:
\begin{enumerate}
  \item the set of allowed perturbations
        $S(\mathbf{x})=\{\mathbf{x}+\boldsymbol{\delta}:
        \lVert\boldsymbol{\delta}\rVert\le\varepsilon\}$, where
        $\lVert\cdot\rVert$ is usually the $\ell_\infty$ norm and $\varepsilon>0$
        is the \emph{budget};
  \item the knowledge of the attacker: full access to $f$ and its gradient
        (\emph{white box}) or only to the output of $f$ (\emph{black box});
  \item the goal: any misclassification (\emph{untargeted}) or a particular
        class (\emph{targeted});
  \item other constraints, for example the perturbation must still be a valid
        object.
\end{enumerate}
The attack succeeds when there is a $\boldsymbol{\delta}\in S(\mathbf{x})$ that
changes the label.
\end{definisi}

The $\ell_\infty$ norm is chosen because it bounds every coordinate separately:
$\lVert\boldsymbol{\delta}\rVert_\infty\le\varepsilon$ means every pixel shifts by
at most $\varepsilon$. For an 8-bit image with values $0$ to $255$, the budget
$\varepsilon=1/255$ means every pixel may change by at most one level.
The $\ell_2$ norm bounds the total length of the vector and allows one pixel to
change a lot as long as the others barely change. The two answer different
questions, and a robustness claim that does not name its norm cannot be compared
with another claim.
\index{infinity norm}
\index{perturbation budget}

The most honest measure of robustness is the distance to the nearest
misclassification point,
\[
  r(\mathbf{x})=\inf\{\lVert\boldsymbol{\delta}\rVert:
  f(\mathbf{x}+\boldsymbol{\delta})\neq y\}.
\]
This quantity does not depend on the algorithm of any attacker. Unfortunately $r$
is hard to compute: finding it means maximizing a nonconvex function on a large
set. For a network with ReLU units, $r$ can be computed exactly by turning the
problem into a mixed-integer linear program, and its cost grows quickly with the
number of units. That is why almost every report in this field names what a
particular attacker managed to achieve, not the value of $r$. The gap between the
two is the source of most conclusions that are too optimistic.

One fact deserves to be recorded from the start. On tabular data, the
$\ell_\infty$ ball contains impossible points: negative ages, non-integer packet
counts, combinations of TCP flags that never appear in a real network. Examples
like that fool the model but cannot be realized, so their success tells us nothing
about the real risk. A budget that is too loose produces attacks that are
spectacular and meaningless.

\section{Why the sign of the gradient is enough}
\tingkat{intermediate}

Now we attack. Let the error of the model be expressed by a function
$L(\mathbf{x},y)$, for example cross-entropy. We want to raise that error as much
as possible while keeping every coordinate inside the box. For a linear function,
the answer can be written in closed form.

\kamus{fast gradient sign method}{metode tanda gradien}{A one-step attack that
shifts every coordinate by the budget, with the direction following the sign of the
derivative of the error.}
\index{FGSM}

\begin{lemma}[Maximum of a linear form on the $\ell_\infty$ ball]
For $\mathbf{w}\in\mathbb{R}^{d}$ and $\varepsilon>0$,
\[
  \max_{\lVert\boldsymbol{\delta}\rVert_\infty\le\varepsilon}
  \mathbf{w}^{\top}\boldsymbol{\delta}
  =\varepsilon\lVert\mathbf{w}\rVert_1,
\]
and the maximum is attained at $\delta_j=\varepsilon\,\mathrm{sign}(w_j)$ for every
coordinate with $w_j\neq0$.
\end{lemma}

\begin{proof}
Write $\mathbf{w}^{\top}\boldsymbol{\delta}=\sum_j w_j\delta_j$. Because
$|\delta_j|\le\varepsilon$,
\[
  \sum_j w_j\delta_j\le\sum_j|w_j|\,|\delta_j|\le\varepsilon\sum_j|w_j|.
\]
If $\delta_j=\varepsilon\,\mathrm{sign}(w_j)$, then $w_j\delta_j=\varepsilon|w_j|$
for every $j$, so that bound is attained. Coordinates with $w_j=0$ do not
affect the value, so their choice is free.
\end{proof}

For a smooth function, we use the first-order approximation of $L$ around
$\mathbf{x}$ together with the lemma above. The result is the recipe that
Goodfellow and colleagues call the \emph{fast gradient sign method}:
\[
  \mathbf{x}^{\mathrm{adv}}
  =\mathbf{x}+\varepsilon\,\mathrm{sign}\bigl(\nabla_{\mathbf{x}}L(\mathbf{x},y)\bigr).
\]
In everyday language it goes like this. The gradient tells us in which direction
the error rises fastest, but we may only walk a distance $\varepsilon$; only the
sign of each component is used, while the size is fixed by the budget. That is why
this method is called the sign attack. For a linear function, that step is exactly
optimal; Figure~\ref{fig:batas} shows the result on a model with two features.

\begin{figure}[htbp]
\centering
{\scriptsize
\begin{tikzpicture}
\node[at={(0.000cm,0.000cm)},inner sep=0pt,minimum size=0pt] (o) {};
\node[at={(8.400cm,0.000cm)},inner sep=0pt,minimum size=0pt] (ex) {};
\node[at={(0.000cm,8.400cm)},inner sep=0pt,minimum size=0pt] (ey) {};
\draw[-latex,black!75] (o) -- (ex);
\draw[-latex,black!75] (o) -- (ey);
\node[anchor=north,at={(4.200cm,-1.120cm)}]{upper-half ink};
\node[rotate=90,anchor=south,at={(-1.15cm,5.10cm)}]{lower-half ink};
\node[anchor=north,at={(2.380cm,-0.280cm)}]{0.20};
\node[anchor=north,at={(6.580cm,-0.280cm)}]{0.35};
\node[anchor=east,at={(-0.224cm,2.660cm)}]{0.35};
\node[anchor=east,at={(-0.224cm,6.860cm)}]{0.50};
\node[at={(1.439cm,0.000cm)},inner sep=0pt,minimum size=0pt] (b0) {};
\node[at={(8.400cm,7.224cm)},inner sep=0pt,minimum size=0pt] (b1) {};
\draw[black,line width=1.2pt] (b0) -- (b1);
\node[circle,fill=black!55,inner sep=1.35pt,at={(2.358cm,4.399cm)}]{};
\node[circle,fill=black!55,inner sep=1.35pt,at={(3.014cm,5.602cm)}]{};
\node[circle,fill=black!55,inner sep=1.35pt,at={(2.796cm,2.376cm)}]{};
\node[circle,fill=black!55,inner sep=1.35pt,at={(3.178cm,5.274cm)}]{};
\node[circle,fill=black!55,inner sep=1.35pt,at={(0.444cm,3.524cm)}]{};
\node[circle,fill=black!55,inner sep=1.35pt,at={(1.975cm,4.290cm)}]{};
\node[circle,draw=black!55,line width=0.5pt,inner sep=1.15pt,at={(5.585cm,2.649cm)}]{};
\node[circle,draw=black!55,line width=0.5pt,inner sep=1.15pt,at={(5.749cm,2.868cm)}]{};
\node[circle,draw=black!55,line width=0.5pt,inner sep=1.15pt,at={(5.147cm,1.391cm)}]{};
\node[circle,draw=black!55,line width=0.5pt,inner sep=1.15pt,at={(3.780cm,0.844cm)}]{};
\node[circle,draw=black!55,line width=0.5pt,inner sep=1.15pt,at={(6.678cm,3.032cm)}]{};
\node[circle,draw=black!55,line width=0.5pt,inner sep=1.15pt,at={(5.803cm,2.868cm)}]{};
\node[circle,draw=black,line width=1pt,inner sep=2.8pt,at={(2.303cm,7.079cm)}]{};
\node[circle,fill=black,inner sep=0.9pt,at={(2.303cm,7.079cm)}]{};
\node[rectangle,fill=accent,inner sep=2.3pt,at={(5.663cm,3.719cm)}]{};
\node[at={(2.303cm,7.079cm)},inner sep=0pt,minimum size=0pt] (x0) {};
\node[at={(5.663cm,3.719cm)},inner sep=0pt,minimum size=0pt] (xa) {};
\draw[-latex,accent,line width=1.1pt] (x0) -- (xa);
\node[anchor=south east,at={(2.135cm,7.443cm)}]{clean};
\node[anchor=north west,at={(5.943cm,3.551cm)}]{adversarial};
\node[anchor=east,at={(2.660cm,5.600cm)}]{class 6};
\node[anchor=south west,at={(5.684cm,1.260cm)}]{class 0};
\node[anchor=south east,at={(5.475cm,4.647cm)}]{boundary};
\end{tikzpicture}}
\caption{Decision boundary of a model with two features on the digits $0$ against
$6$ ($236$ training data, $123$ test, accuracy $98.37$ percent). Thick circle: a
clean example $(0.197;0.508)$ with score $+1.936$. Red square: the result of FGSM
with $\varepsilon=0.12$ at $(0.317;0.388)$ with score $-0.208$. Its arrow is the
perturbation $(0.12;-0.12)$, $\lVert\boldsymbol{\delta}\rVert_\infty=0.12$
and $\lVert\boldsymbol{\delta}\rVert_2=0.1697$.}
\label{fig:batas}
\end{figure}

\begin{teorema}[Linearization and the remainder]
Let $L:\mathbb{R}^{d}\to\mathbb{R}$ be differentiable and let its gradient be
$M$-Lipschitz on the segment between $\mathbf{x}$ and $\mathbf{x}+\boldsymbol{\delta}$.
For every $\lVert\boldsymbol{\delta}\rVert_\infty\le\varepsilon$,
\[
  \left|L(\mathbf{x}+\boldsymbol{\delta})-L(\mathbf{x})
  -\nabla L(\mathbf{x})^{\top}\boldsymbol{\delta}\right|
  \le\frac{M}{2}\lVert\boldsymbol{\delta}\rVert_2^2
  \le\frac{Md}{2}\varepsilon^2 .
\]
\end{teorema}

\begin{proof}
With the integral form of the mean value theorem,
\[
  L(\mathbf{x}+\boldsymbol{\delta})-L(\mathbf{x})
  =\int_0^1\nabla L(\mathbf{x}+t\boldsymbol{\delta})^{\top}\boldsymbol{\delta}\,dt .
\]
Subtract $\nabla L(\mathbf{x})^{\top}\boldsymbol{\delta}
=\int_0^1\nabla L(\mathbf{x})^{\top}\boldsymbol{\delta}\,dt$, then use the
Cauchy--Schwarz inequality and the Lipschitz property:
\[
  \left|\int_0^1\bigl(\nabla L(\mathbf{x}+t\boldsymbol{\delta})
  -\nabla L(\mathbf{x})\bigr)^{\top}\boldsymbol{\delta}\,dt\right|
  \le\int_0^1 Mt\lVert\boldsymbol{\delta}\rVert_2^2\,dt
  =\frac{M}{2}\lVert\boldsymbol{\delta}\rVert_2^2 .
\]
Because $\lVert\boldsymbol{\delta}\rVert_2^2\le d\,\varepsilon^2$, the second
inequality follows.
\end{proof}

For a linear model, that remainder is zero and FGSM is not an approximation but the
exact answer. Goodfellow and colleagues use the case of logistic regression to show
this. There the sign of the gradient of the error with respect to the input is
$-\mathrm{sign}(\mathbf{w})$ times the label, so that
$\mathbf{w}^{\top}\mathrm{sign}(\mathbf{w})=\lVert\mathbf{w}\rVert_1$ and the margin
of the model shrinks by exactly $\varepsilon\lVert\mathbf{w}\rVert_1$. That is why
adversarial training on logistic regression can be written as a penalty
$\zeta\bigl(y(\varepsilon\lVert\mathbf{w}\rVert_1-\mathbf{w}^{\top}\mathbf{x}-b)\bigr)$,
that is a penalty resembling an $\ell_1$ regularizer that enters the margin, not the
cost; this exact form is used in the worked example later.

What needs to be stated honestly: the theorem above justifies the \emph{direction} of
the step, not the \emph{distance} of the step. The remainder $Md\varepsilon^{2}/2$
depends on the Lipschitz constant of the gradient, and for a ReLU network that
constant can be very large. In our experiment, the predicted rise in the error from
the linear term at $\varepsilon=0.15$ is $0.0008$ nats, whereas the rise that
actually occurs is $4.862$ nats. So linearization does not explain the size of the
damage; it only explains why the chosen direction makes sense.

The other side of the story is stronger instead. A perturbation of the same size but
with random signs only lowers accuracy from $97.99$ to $95.48$ percent, whereas FGSM
lowers it to $14.91$ percent. So what is valuable is the pattern of signs, even
though the error at the starting point is almost flat, so that the length of the
gradient and its linear forecast are small.
\index{input gradient}

So the sign of the gradient is enough not because its linear term is large, but
because the sign is more resistant to saturation of the error. The number $0.0008$
against $4.862$ is a warning that linearity holds only as an argument about
direction.

\begin{bacakeras}
Two voices, read alternately with a friend. Tap the table every time the word ``a
little'' is said.

``Shift it a little.'' ``How little?'' ``A fingertip. No more.''
``Does the answer change?'' ``It changes completely. Not doubtful, not off target: changed.''
``What do our eyes see?'' ``Nothing. Exactly the same picture.''
``And if it is shifted that far in another direction?'' ``The answer becomes correct again.''
``So what is doing the work?'' ``The direction. Not the size.''
``Where does it learn the direction?'' ``It reads the slope, asks every pixel: up or
down? Then all the answers are used at once.''
``And if it misreads the slope?'' ``Nothing happens. The picture stays, the answer stays.''

Repeat three times, faster and faster: what works is the direction, not the size.
\end{bacakeras}

\section{Iterative attacks and the min-max game}
\tingkat{intermediate}

One step is a compromise. If the gradient may be computed many times, we can climb
further: a small step in the direction of the sign of the gradient, then pull back
into the box.

\begin{definisi}[Iterative attack]
Let $\Pi$ be the projection onto $S(\mathbf{x})$. The \emph{projected gradient
descent} attack (PGD) computes
\[
  \mathbf{x}_{t+1}=\Pi\Bigl(\mathbf{x}_t+\alpha\,
  \mathrm{sign}\bigl(\nabla_{\mathbf{x}}L(\mathbf{x}_t,y)\bigr)\Bigr),
  \qquad \mathbf{x}_0=\mathbf{x},
\]
with $\alpha>0$ the step size. If $\mathbf{x}_0$ is taken at random inside
$S(\mathbf{x})$, the attack is called PGD with a \emph{random start}.
\end{definisi}

\begin{resep}
Four steps to run a bounded iterative attack:
\begin{enumerate}
  \item Fix the norm, $\varepsilon$, $\alpha$, and the number of steps; record the
        clean accuracy before a single perturbation is made.
  \item Start from $\mathbf{x}$ or from a random point inside the ball, compute the
        gradient of the error with respect to the input, then take only its sign.
  \item Step: $\mathbf{x}_{t+1}=\mathbf{x}_t+\alpha\,\mathrm{sign}(\mathbf{g})$,
        project onto the $\varepsilon$ ball, then onto $[0,1]$.
  \item Repeat, then report the accuracy at several numbers of steps and several
        $\varepsilon$; compare with FGSM, which uses $\alpha=\varepsilon$.
\end{enumerate}
Warning about order: the projection is always with respect to the original
$\mathbf{x}$, not to the previous step point, and always after the gradient step.
If the reference is swapped, the ball moves along and the perturbation exceeds the
reported budget.
\end{resep}

The projection onto the $\ell_\infty$ box cut to $[0,1]^d$ has a very
simple form: clip every coordinate, one by one.

\begin{proposisi}[Projection onto the box]
For $S=\{\mathbf{z}:\lVert\mathbf{z}-\mathbf{x}\rVert_\infty\le\varepsilon\}
\cap[0,1]^{d}$,
\[
  \Pi(\mathbf{z})_j=\min\bigl\{\max\{z_j,\,x_j-\varepsilon,\,0\},\,
  x_j+\varepsilon,\,1\bigr\}
\]
for every $j$.
\end{proposisi}

\begin{proof}
The objective function $\lVert\mathbf{y}-\mathbf{z}\rVert_2^2$ separates per
coordinate, and the allowed set also separates: every $y_j$ must lie in the interval
$[\max\{0,x_j-\varepsilon\},\min\{1,x_j+\varepsilon\}]$. For one coordinate,
minimizing $(y_j-z_j)^2$ on a closed interval means taking the nearest point, that
is $z_j$ if it is inside the interval, and the nearest end of the interval if it is
outside. That is exactly the clipping formula above.
\end{proof}

With a projection of this kind, PGD adds only two things compared with FGSM:
repeated small steps and the pull back into the box. The increase in cost is
proportional to the number of steps, because every step requires one backward pass.
Figure~\ref{fig:langkah} shows three such steps on the same model.

\begin{figure}[htbp]
\centering
{\scriptsize
\begin{tikzpicture}
\node[at={(0.000cm,0.000cm)},inner sep=0pt,minimum size=0pt] (p) {};
\node[at={(8.400cm,0.000cm)},inner sep=0pt,minimum size=0pt] (px) {};
\node[at={(0.000cm,8.400cm)},inner sep=0pt,minimum size=0pt] (py) {};
\draw[-latex,black!75] (p) -- (px);
\draw[-latex,black!75] (p) -- (py);
\node[anchor=north,at={(4.200cm,-1.344cm)}]{upper-half ink};
\node[rotate=90,anchor=south,at={(-1.05cm,4.20cm)}]{lower-half ink};
\node[anchor=north,at={(1.848cm,-0.336cm)}]{0.20};
\node[anchor=north,at={(6.888cm,-0.336cm)}]{0.35};
\node[at={(2.661cm,0.000cm)},inner sep=0pt,minimum size=0pt] (c0) {};
\node[at={(8.400cm,5.956cm)},inner sep=0pt,minimum size=0pt] (c1) {};
\draw[black,line width=1.2pt] (c0) -- (c1);
\node[at={(1.756cm,6.479cm)},inner sep=0pt,minimum size=0pt] (q0) {};
\node[circle,fill=white,draw=accent,line width=0.9pt,inner sep=2.0pt,at={(1.756cm,6.479cm)}]{};
\node[at={(3.100cm,5.135cm)},inner sep=0pt,minimum size=0pt] (q1) {};
\node[circle,fill=white,draw=accent,line width=0.9pt,inner sep=2.0pt,at={(3.100cm,5.135cm)}]{};
\node[at={(4.444cm,3.791cm)},inner sep=0pt,minimum size=0pt] (q2) {};
\node[circle,fill=white,draw=accent,line width=0.9pt,inner sep=2.0pt,at={(4.444cm,3.791cm)}]{};
\node[at={(5.788cm,2.447cm)},inner sep=0pt,minimum size=0pt] (q3) {};
\node[circle,fill=white,draw=accent,line width=0.9pt,inner sep=2.0pt,at={(5.788cm,2.447cm)}]{};
\draw[-latex,accent,line width=1.0pt,shorten >=5pt,shorten <=5pt] (q0) -- (q1);
\draw[-latex,accent,line width=1.0pt,shorten >=5pt,shorten <=5pt] (q1) -- (q2);
\draw[-latex,accent,line width=1.0pt,shorten >=5pt,shorten <=5pt] (q2) -- (q3);
\node[anchor=south west,at={(1.891cm,6.613cm)}]{$\mathbf{x}_0$};
\node[anchor=south west,at={(3.235cm,5.269cm)}]{$\mathbf{x}_1$};
\node[anchor=south west,at={(4.579cm,3.925cm)}]{$\mathbf{x}_2$};
\node[anchor=south west,at={(5.923cm,2.581cm)}]{$\mathbf{x}_3$};
\node[anchor=south east,at={(5.616cm,3.444cm)}]{boundary};
\end{tikzpicture}}
\caption{Three sign-of-gradient steps from $\mathbf{x}_0$, each step
$\alpha=0.04$. The $\lVert\boldsymbol{\delta}\rVert_\infty$ norm: $0.04$,
$0.08$, $0.12$; the $\lVert\boldsymbol{\delta}\rVert_2$ norm: $0.0566$,
$0.1131$, $0.1697$. The score of class $6$: $+1.936$, $+1.221$, $+0.506$,
$-0.208$; the boundary is crossed at the third step.}
\label{fig:langkah}
\end{figure}

There is one warning that is often missed. FGSM is PGD with one step and
$\alpha=\varepsilon$, not PGD with one step and $\alpha=\varepsilon/10$.
In our experiment, one step with $\alpha=\varepsilon/10$ only lowers
accuracy to $94.64$ percent, whereas FGSM lowers it to $14.91$
percent. So ``one step'' is not the name of one algorithm; it only means something
once its step size is named. A fair comparison uses the same step size; the
following table is computed at $\varepsilon=0.15$ with a random start.

\begin{center}
\begin{tabular}{rr}
\hline
steps & accuracy (\%) \\
\hline
1  & $94.64$ \\
2  & $88.78$ \\
5  & $66.67$ \\
10 & $24.79$ \\
20 & $11.22$ \\
40 & $10.72$ \\
\hline
\end{tabular}
\end{center}

That curve falls sharply up to ten steps, then is almost flat; adding steps after
that point rarely changes the conclusion, and that is why PGD is used as the default
attacker. Madry and colleagues read it as a sign that the inner problem is easy
enough to solve with first-order methods, even though in general it is nonconcave.

\begin{tangan}
Materials: ten cards, a marker, and a ruler. On each card write one number between
zero and one, two digits after the decimal point. Make two cards whose numbers differ
by only 0.01, for example 0.62 and 0.63.

Extra rule: when two numbers differ by less than 0.02, the higher card in the pile
wins. Place the two nearly tied cards in the top two positions, then call the topmost
card the answer of the model.

Now shift the card 0.62 up by one millimeter, until it passes the card 0.63. No
number changes, no card is swapped. Look at the topmost card: the answer changes,
because the tie-breaking rule now chooses the other card.

What must be visible: the answer changes without a single number changing, because
only the difference decides, and a thin difference can be flipped by a shift that the
eye cannot see.
\end{tangan}

From the point of view of training, attack and defense are two sides of one problem.
If the attacker maximizes the error inside the box, then the careful trainer
minimizes that worst error:
\[
  \min_{\boldsymbol{\theta}}\;
  \mathbb{E}_{(\mathbf{x},y)\sim\mathcal{D}}
  \Bigl[\max_{\boldsymbol{\delta}\in S(\mathbf{x})}
  L(\boldsymbol{\theta},\mathbf{x}+\boldsymbol{\delta},y)\Bigr].
\]
This formulation is a saddle problem: a maximum inside, a minimum outside. The inner
term is the problem of the attacker, the outer term is the problem of ordinary
training. This is \emph{adversarial training}. If the inner maximum is approximated
with FGSM, we get cheap training; if it is approximated with PGD of several steps, we
get training that is more expensive and usually more robust.
\index{adversarial training}\index{PGD}

What that formulation can and cannot promise must be separated. If the saddle value
is really attained and its value is zero, no perturbation inside the box can fool the
model, so we have a guarantee. But that saddle problem is not concave inside and not
convex outside, so any algorithm only gives an approximation: a small saddle value is
not necessarily attained, and the value that is obtained is not necessarily the
saddle value. That is why adversarial training produces a model that is \emph{more
resistant to the attacker used during training}, not a model that is proved resistant
to all attackers.

Our numbers show both sides. Model $D$, trained with FGSM at
$\varepsilon=0.15$, loses $1.51$ points of clean accuracy: from $97.99$ to
$96.48$ percent. In exchange, at $\varepsilon=0.15$ its accuracy under
FGSM is $66.67$ percent and under PGD of 40 steps $62.65$ percent, whereas
the ordinary model is left with $14.91$ and $10.72$ percent. The trade is real, but
$62.65$ percent is also not security; it only means the attacker must work
harder.

\section{Transfer, black box, and targeted attacks}
\tingkat{advanced}

Up to here our attacker knows everything. In the real world it can often only send
an input and look at the answer. What saves us when there is no gradient is not the
secrecy of the model, but the fact that the direction of attack tends to carry over
between models.

\kamus{transferability}{transferabilitas}{The property of an adversarial example that
is built for one model but also fools another model that was trained separately.}
\index{transferability}

This phenomenon was reported by Szegedy and colleagues in 2014: an example that fools
one network often fools another network with a different architecture and different
training data. Goodfellow and colleagues record its quantitative version on MNIST: an
example built from an old model gives an error of $19.6$ percent on a model trained
adversarially, whereas an example from the new model gives an error of $40.9$ percent
on the old model. The same direction, a different strength.

In our experiment, model $A$ (64 hidden units) and model $B$ (128 units) are trained
on the same training data but with different initialization. At
$\varepsilon=0.15$, PGD of 40 steps built on $A$ lowers the accuracy of $A$ itself to
$10.72$ percent, and lowers the accuracy of $B$ to $23.79$ percent. FGSM built on
$A$ gives $14.91$ on $A$ and $25.96$ on $B$. So the attack transfers, but
it weakens: about half of its strength is lost when it crosses models.

The most economical explanation is still the one of this chapter: at a small scale
the two models learn functions that are almost linear, and a direction that raises
the error on one linear function also raises it on another linear function whose
normal is close by. That is why transfer works best for one-step attacks. Our numbers
are consistent: the gap between white box and transfer is smaller for FGSM ($11$
points) than for PGD of 40 steps ($13$ points).

For a black-box attacker, transfer offers a cheap path: train a substitute model,
attack that model, then send the result to the victim model. Another route is to
estimate the gradient from differences of output over many queries, which is accurate
but expensive and easy to detect. The black box is weaker because its information is
less, not because its attacker is less clever.

The goal of the attacker also changes the numbers drastically. An untargeted attack
only has to make the label change; a targeted attack has to make the label become a
particular class. At $\varepsilon=0.15$, our PGD of 40 steps reaches the target on
$33.00$ percent of the examples, whereas the untargeted attack succeeds on
$89.06$ percent of the examples (accuracy is left at $10.72$ percent). With
$\varepsilon=0.30$, the target is reached on $99.50$ percent of the
examples. That is why the number ``the attack succeeds 90 percent'' without a note on
whether it is targeted cannot be compared with another number. For a defense
evaluator, a targeted test is harder for the attacker, whereas an untargeted test is
harder for the defender.

\section{Defenses: what works and what only appears to work}
\tingkat{advanced}

The defender has three families of answers. The first changes the data: add
adversarial examples to the training. The second changes the model: make its surface
smoother, for example with quantized inputs or smoothed outputs. The third adds tools
outside the model: detect strange inputs, or refuse answers that are not confident.
Only the first family has a clear basis, namely the saddle formulation above; the
second and third families often give good numbers for the wrong reason.

The most frequently repeated mistake has a name.

\begin{definisi}[Gradient obfuscation]
A defense performs \emph{gradient obfuscation} when it makes the gradient that the
attacker uses useless: missing, pointing the wrong way, or random. Such a defense
lowers the success of gradient-based attacks without raising the distance to the
decision boundary.
\end{definisi}

In everyday language: that guard does not become more observant, he starts using a
form that the attacker cannot read. Athalye, Carlini, and Wagner distinguish three
ways obfuscation happens. First, a \emph{shattered} gradient: an operation that is
nondifferentiable or numerically unstable makes the gradient absent or wrong. Second,
a \emph{random} gradient: a defense that adds randomness makes one sample of the
gradient not the true gradient. Third, a gradient that \emph{explodes or vanishes}: a
defense that evaluates the network many times in sequence, as in repeated
quantization, behaves like a very deep network when it is differentiated.
\index{obfuscated gradient}

Why is obfuscation not a defense? Because it changes nothing that a persistent
attacker sees. If the gradient is useless, the attacker uses a substitute gradient
from a smoother model; if randomness is the obstacle, he averages the gradient over
many samples; if the gradient explodes, he uses a finite-difference approximation.

We build the small example ourselves. Take model $A$ and put input quantization in
front of it: every pixel is rounded to a multiple of $1/8$. The derivative of rounding
is zero almost everywhere, so an attacker who differentiates through that path gets a
zero vector and does not move. The result looks convincing: clean accuracy
$97.82$ percent, and after ``PGD of 40 steps'' its accuracy is $96.98$ percent. Almost
no damage. But if the attacker uses the gradient of the model without quantization at
the point that has already been quantized, its accuracy falls to $13.23$ percent, only
a little higher than $10.72$ percent on the model without quantization. That defense
does not close the road; it only closes the most naive road. The shape of the picture
is in Figure~\ref{fig:kuantisasi-gradien}.

\begin{figure}[htbp]
\centering
\begin{tikzpicture}[scale=1.05,font=\scriptsize]
  \fill[layerteal!8] (0,0.5) -- (4,2.9) -- (4,4) -- (0,4) -- cycle;
  \fill[accent!6] (0,0) -- (4,0) -- (4,2.9) -- (0,0.5) -- cycle;
  \draw[very thick,ink] (0,0.5)--(4,2.9);
  \draw[->,layergrey] (-0.2,0)--(4.35,0) node[right]{$x_{1}$};
  \draw[->,layergrey] (0,-0.2)--(0,4.35) node[above]{$x_{2}$};
  \node[font=\scriptsize\bfseries] at (2,4.68) {original model};
  \node[layerteal] at (0.78,3.55) {class $+1$};
  \node[accent] at (3.05,0.8) {class $-1$};
  \node[layergrey,rotate=31] at (2.75,2.95) {boundary};
  \draw[-{Latex},very thick,accent] (1.2,2.6)--(1.92,1.4);
  \node[accent,above left=-2pt] at (1.56,2.05) {$\boldsymbol{\delta}$};
  \fill[ink] (1.2,2.6) circle (1.6pt);
  \node[ink,above left=-2pt] at (1.2,2.6) {$\mathbf{x}$};
  \draw[accent,thick,fill=accent!35] (1.92,1.4) ++(135:0.07) -- ++(-45:0.14)
    -- ++(-135:0.14) -- ++(45:0.14) -- cycle;
  \node[accent,below right=-2pt] at (1.92,1.4) {$\mathbf{x}+\boldsymbol{\delta}$};
\end{tikzpicture}
\hfill
\begin{tikzpicture}[scale=1.05,font=\scriptsize]
  \fill[layerteal!8] (0,0.5) -- (4,2.9) -- (4,4) -- (0,4) -- cycle;
  \fill[accent!6] (0,0) -- (4,0) -- (4,2.9) -- (0,0.5) -- cycle;
  \foreach \g in {0.5,1,1.5,2,2.5,3,3.5}
    {\draw[layergrey!30,very thin] (\g,0)--(\g,4);
     \draw[layergrey!30,very thin] (0,\g)--(4,\g);}
  \draw[very thick,ink] (0,0.5)--(4,2.9);
  \draw[->,layergrey] (-0.2,0)--(4.35,0) node[right]{$x_{1}$};
  \draw[->,layergrey] (0,-0.2)--(0,4.35) node[above]{$x_{2}$};
  \node[font=\scriptsize\bfseries] at (2,4.68) {with rounding to $1/8$};
  \fill[ink] (1.2,2.6) circle (1.6pt);
  \node[ink,above left=-2pt] at (1.2,2.6) {$\mathbf{x}$};
  \draw[layerplum,thick] (1.2,2.6) circle (0.13);
  \node[layerplum,anchor=west] at (1.4,2.78) {$\nabla$ path $=\mathbf{0}$};
  \draw[-{Latex},very thick,dashed,layerplum] (1.2,2.6)--(1.92,1.4);
  \node[layerplum,anchor=east] at (1.44,1.72) {$\nabla$ substitute};
  \draw[layerplum,thick,fill=layerplum!35] (1.92,1.4) ++(135:0.07)
    -- ++(-45:0.14) -- ++(-135:0.14) -- ++(45:0.14) -- cycle;
\end{tikzpicture}
\caption{Input quantization on the same model. Left: the true gradient gives a
direction $\boldsymbol{\delta}$ that takes $\mathbf{x}$ across to the other side of
the boundary, so that its label changes. Right: rounding every feature to a multiple
of $1/8$ makes the derivative of the path zero at almost all points, so the reported
gradient arrow stops; the decision boundary itself does not shift, and the gradient
of the model without rounding still takes it across. The square marks the point after
the perturbation.}
\label{fig:kuantisasi-gradien}
\end{figure}

Athalye and colleagues examined nine white-box defenses at ICLR 2018. Seven of them
depend on gradient obfuscation, and adapted attacks succeeded in breaking six
completely and one partially, under the threat model each of them claimed. This is not
a small error in one paper; it is a recurring pattern. That is why there is a list of
symptoms that can be used before trusting a robustness number. If a one-step attack is
stronger than an iterative attack, something is wrong: an iterative attack should
always be stronger in the white box. If a black-box attack is stronger than a white-box
one, something is wrong: the black box is a smaller set of information. If an attack
without a bound on the perturbation does not reach one hundred percent success,
something is wrong. If a random search finds an adversarial example while the gradient
attack does not, something is wrong. If enlarging the budget does not raise success,
something is wrong.

\begin{tantangan}
A service refuses to answer when the absolute value of its score margin is below 0.50,
and answers only when it is confident enough. In the linear example of the previous
part, the clean margin is 1.500 and $\lVert\mathbf{d}\rVert_1=6$. An attack at
$\varepsilon=0.25$ pushes the margin to zero, exactly at the threshold of reversal, so
it is refused, and the service looks safe. Determine the smallest budget that keeps
the attacker served, that is the one that switches the refusal off on the wrong side.
Is that rule a defense or obfuscation?

Hint 1: the refusal looks only at the size of the margin, not at its sign.

Hint 2: solve $1.500-6\varepsilon=-0.500$.

If you get stuck, go back to the section Worked example: a linear model you can compute
by hand.
\end{tantangan}

What does a correct evaluation look like? It is not one number, but one protocol. Fix
the threat model first: norm, $\varepsilon$, white box or black box, targeted or
not. Use the strongest attacker you can afford, and adapt it to the defense: if the
defense quantizes, attack through a substitute gradient; if it adds randomness,
average over several samples; if it filters the input, attack that filter as well.
Report the accuracy at several $\varepsilon$, not at one point, and check that the
curve falls monotonically. Report more than one random seed and include transfer
attacks from another model. Report the cost: the number of backward passes per
example, the time, and the hardware. If you want to talk about guarantees, use
certification, not testing; that is another path that is not discussed here.

\begin{untukanak}
Think about a multiple-choice exam that is graded by machine. The machine reads the
answer from the position of the bubble that is filled in. If we test the machine only
with answers that are filled in neatly, the machine will always be right. Then comes a
student who knows how the machine reads: he fills one other bubble a little more
thickly, and the machine reads it as a different answer. The paper is still the same
paper, the writing is still the same writing.

Now the interesting part. One school tries to fix that machine by scanning the paper
twice and taking the answer that appears most often. The same student just fills two
bubbles, and the machine can still be fooled. Another school adds a ruler so that the
position of the bubbles shifts a little; the patient student just tries a few
positions. All those additions make the cheating more troublesome, but they do not
make the machine able to tell an honest answer.

In this chapter we will find out why making things harder for the attacker is not the
same as closing the hole, and how to test a fix so that we are not fooled by a fix
that only appears to work.
\end{untukanak}

\section{Worked example: a linear model you can compute by hand}
\tingkat{intermediate}

Now we leave the networks and return to a model that can be computed without a
computer: a linear classifier with two classes and four features. Every number in this
part can be checked with a calculator.

Let class $A$ use the weights $\mathbf{w}_A=(1;-2;0.5;1.5)$ and class $B$
use $\mathbf{w}_B=(0;1;0.5;-0.5)$, without a bias. The example that is examined is
$\mathbf{x}=(1;0.5;1;1)$. The scores of the two classes are
\[
  s_A=\mathbf{w}_A^{\top}\mathbf{x}=1-1+0.5+1.5=2.0,
  \qquad
  s_B=\mathbf{w}_B^{\top}\mathbf{x}=0+0.5+0.5-0.5=0.5 .
\]
Its margin is $m=s_A-s_B=1.5$, so the model is confident but not excessively so. The
difference of the weights is $\mathbf{d}=\mathbf{w}_A-\mathbf{w}_B=(1;-3;0;2)$, with
$\lVert\mathbf{d}\rVert_1=1+3+0+2=6$ and
$\lVert\mathbf{d}\rVert_2=\sqrt{1+9+0+4}=\sqrt{14}\approx3.7417$.

An untargeted attack has to lower the margin. By the lemma, the largest decrease
is attained by $\boldsymbol{\delta}=-\varepsilon\,\mathrm{sign}(\mathbf{d})
=\varepsilon(-1;1;0;-1)$, and the size of the decrease is
$\varepsilon\lVert\mathbf{d}\rVert_1=6\varepsilon$. So the margin after the attack
is $1.5-6\varepsilon$, and the label changes when
\[
  \varepsilon>\varepsilon^{*}=\frac{1.5}{6}=0.25 .
\]
At $\varepsilon=0.25$ the margin is exactly zero, so the two classes tie and
$\arg\max$ chooses the first class; a clear reversal happens only a little above
$0.25$. The following table contains the whole course of it, and every row can be
checked by multiplying $\mathbf{d}$ by $\mathbf{x}+\boldsymbol{\delta}$.

\begin{center}
\begin{tabular}{rccc}
\hline
$\varepsilon$ & $s_A$ & $s_B$ & margin \\
\hline
$0$    & $2.000$ & $0.500$ & $+1.500$ \\
$0.05$ & $1.775$ & $0.575$ & $+1.200$ \\
$0.10$ & $1.550$ & $0.650$ & $+0.900$ \\
$0.15$ & $1.325$ & $0.725$ & $+0.600$ \\
$0.20$ & $1.100$ & $0.800$ & $+0.300$ \\
$0.25$ & $0.875$ & $0.875$ & $0.000$ \\
$0.30$ & $0.650$ & $0.950$ & $-0.300$ \\
\hline
\end{tabular}
\end{center}

If the attacker is bounded by the $\ell_2$ norm instead of $\ell_\infty$, the answer
changes. The maximum of $\mathbf{d}^{\top}\boldsymbol{\delta}$ on the $\ell_2$ ball of
radius $\varepsilon$ is $\varepsilon\lVert\mathbf{d}\rVert_2$, attained at
$\boldsymbol{\delta}=\varepsilon\mathbf{d}/\lVert\mathbf{d}\rVert_2$. So
$\varepsilon^{*}=1.5/\sqrt{14}\approx0.4009$. The $\ell_2$ budget that gives damage of
the same size as $\varepsilon=0.25$ on $\ell_\infty$ is
$0.4009$; both move the example to the same point, namely the point of intersection
$\mathbf{x}-\varepsilon\,\mathrm{sign}(\mathbf{d})$. This is a small example of a
general rule: the number $\varepsilon$ cannot be compared across norms.

Now the part that is often asked: is PGD stronger on this model? Every PGD step adds a
decrease of the margin of $\alpha\lVert\mathbf{d}\rVert_1
=6\alpha$, but the whole trajectory stays inside the ball, so the total decrease never
exceeds $6\varepsilon$. For a linear model, the worst point inside the ball really is
the point that FGSM chooses, so iteration cannot beat it. Hand computation shows it:
with $\varepsilon=0.10$ and 40 steps of
$\alpha=0.01$, the final margin $0.900$ is exactly the same as the result of FGSM at
$\varepsilon=0.10$.

\kertas{Repeat the example above for $\mathbf{w}_A=(2;1;-1;0.5)$,
$\mathbf{w}_B=(0;1;1;-0.5)$, and $\mathbf{x}=(1;1;0.5;1)$. Compute $s_A$, $s_B$,
$\mathbf{d}$, $\lVert\mathbf{d}\rVert_1$, and $\varepsilon^{*}$ for the
$\ell_\infty$ norm. Then compute $\varepsilon^{*}$ for the $\ell_2$ norm and explain
why the two numbers differ. Check with a three-line Python snippet. After that decide,
without running a program, whether PGD of 20 steps with
$\alpha=\varepsilon/10$ can change the label when $\varepsilon$ is only half of
$\varepsilon^{*}$.}

\section{Computational lab}
\tingkat{advanced}

This lab trains three small networks on handwritten digit data provided by the
\texttt{scikit-learn} library: 1797 examples, 64 features, values $0$ to $16$ which we
divide by $16$ so that they fall in $[0,1]$. Twelve hundred examples are used for
training and 597 for test. Model $A$ has one hidden layer of 64 units, model $B$ 128
units, and model $D$ is the same as $A$ but trained with FGSM at $\varepsilon=0.15$
(the reason for using $D$ rather than $A$ is that the comparison of clean and robust
accuracy should come from models of the same size).

All the gradients are computed by hand with NumPy, without an automatic differentiation
library, so that every step can be checked. The function \texttt{grad\_x} is the
backward pass of Chapter 25 that stops at the input, not at the weights.

\begin{lab}
\begin{verbatim}
import numpy as np, time
from sklearn.datasets import load_digits

X, y = load_digits(return_X_y=True)
X = X / 16.0                        # 0..16 -> 0..1
perm = np.random.default_rng(0).permutation(len(X))
tr, te = perm[:1200], perm[1200:]
Xtr, ytr, Xte, yte = X[tr], y[tr], X[te], y[te]
K = 10
Ytr, Yte = np.eye(K)[ytr], np.eye(K)[yte]

def init(dims, seed):
    r = np.random.default_rng(seed); W, b = [], []
    for i in range(len(dims) - 1):
        W.append(r.normal(0, np.sqrt(2.0 / dims[i]),
                          (dims[i], dims[i + 1])))
        b.append(np.zeros(dims[i + 1]))
    return W, b

def fwd(W, b, X):
    acts, pre, h = [X], [], X
    for i in range(len(W) - 1):
        z = h @ W[i] + b[i]; pre.append(z)
        h = np.maximum(z, 0); acts.append(h)
    z = h @ W[-1] + b[-1]; pre.append(z)
    z = z - z.max(1, keepdims=True); p = np.exp(z)
    return p / p.sum(1, keepdims=True), acts, pre

def loss_grad(W, b, X, Y):
    p, acts, pre = fwd(W, b, X); n = len(X)
    loss = -np.mean(np.sum(Y * np.log(p + 1e-12), 1))
    d = p - Y; gW = [None] * len(W); gb = [None] * len(W)
    for i in range(len(W) - 1, -1, -1):
        gW[i] = acts[i].T @ d / n + 1e-4 * W[i]; gb[i] = d.mean(0)
        if i > 0:
            d = (d @ W[i].T) * (pre[i - 1] > 0)
    return loss, gW, gb

def grad_x(W, b, X, Y):             # gradient of the error wrt the input
    p, acts, pre = fwd(W, b, X)
    d = (p - Y) / len(X)
    for i in range(len(W) - 1, 0, -1):
        d = (d @ W[i].T) * (pre[i - 1] > 0)
    return d @ W[0].T

def train(dims, seed, Xa, Ya, epochs=350, lr=0.01, adv_eps=None):
    W, b = init(dims, seed)
    mW = [np.zeros_like(w) for w in W]
    vW = [np.zeros_like(w) for w in W]
    mb = [np.zeros_like(v) for v in b]
    vb = [np.zeros_like(v) for v in b]
    for t in range(1, epochs + 1):
        if adv_eps is None:
            Xb = Xa
        else:
            g = grad_x(W, b, Xa, Ya)
            Xb = np.clip(Xa + adv_eps * np.sign(g), 0, 1)
        _, gW, gb = loss_grad(W, b, Xb, Ya)
        for i in range(len(W)):
            mW[i] = 0.9 * mW[i] + 0.1 * gW[i]
            vW[i] = 0.999 * vW[i] + 0.001 * gW[i] ** 2
            mb[i] = 0.9 * mb[i] + 0.1 * gb[i]
            vb[i] = 0.999 * vb[i] + 0.001 * gb[i] ** 2
            mh = mW[i] / (1 - 0.9 ** t)
            vh = vW[i] / (1 - 0.999 ** t)
            W[i] -= lr * mh / (np.sqrt(vh) + 1e-8)
            mh = mb[i] / (1 - 0.9 ** t)
            vh = vb[i] / (1 - 0.999 ** t)
            b[i] -= lr * mh / (np.sqrt(vh) + 1e-8)
    return W, b

def acc(W, b, X, y):
    return float(np.mean(fwd(W, b, X)[0].argmax(1) == y))

def fgsm(W, b, X, y, eps):
    g = grad_x(W, b, X, np.eye(K)[y])
    return np.clip(X + eps * np.sign(g), 0, 1)

def pgd(W, b, X, y, eps, steps, Ylab=None, arah=1):
    alpha = eps / 10
    r = np.random.default_rng(1)
    xa = np.clip(X + r.uniform(-eps, eps, X.shape), 0, 1)
    Y = np.eye(K)[y] if Ylab is None else Ylab
    for _ in range(steps):
        g = arah * np.sign(grad_x(W, b, xa, Y))
        xa = np.clip(np.clip(xa + alpha * g, X - eps, X + eps), 0, 1)
    return xa
\end{verbatim}
\end{lab}

The rest is calls and printing. We train $A$, $B$, and $D$, then sweep
$\varepsilon$ from $0.02$ to $0.30$.

\begin{verbatim}
WA, bA = train([64, 64, 10], 1, Xtr, Ytr)
WB, bB = train([64, 128, 10], 2, Xtr, Ytr)
WD, bD = train([64, 64, 10], 3, Xtr, Ytr, epochs=120, adv_eps=0.15)

print(f"clean test accuracy: A {acc(WA,bA,Xte,yte):.4f}"
      f"  B {acc(WB,bB,Xte,yte):.4f}  D {acc(WD,bD,Xte,yte):.4f}")
print(f"{'eps':>5} {'A clean':>9} {'A FGSM':>7} {'A PGD40':>8}"
      f" {'D FGSM':>7} {'D PGD40':>8}")
for eps in [0.02, 0.05, 0.10, 0.15, 0.20, 0.30]:
    xf = fgsm(WA, bA, Xte, yte, eps)
    xp = pgd(WA, bA, Xte, yte, eps, 40)
    xd = fgsm(WD, bD, Xte, yte, eps)
    xq = pgd(WD, bD, Xte, yte, eps, 40)
    print(f"{eps:5.2f} {acc(WA,bA,Xte,yte):9.4f} {acc(WA,bA,xf,yte):7.4f}"
          f" {acc(WA,bA,xp,yte):8.4f} {acc(WD,bD,xd,yte):7.4f}"
          f" {acc(WD,bD,xq,yte):8.4f}")

r = np.random.default_rng(7)
rs = np.sign(r.normal(size=Xte.shape))          # control: random signs
print("random signs eps=0.15:",
      f"{acc(WA,bA,np.clip(Xte+0.15*rs,0,1),yte):.4f}")
print("PGD steps, eps=0.15, alpha=eps/10, random start")
for k in [1, 2, 5, 10, 20, 40]:
    print(f"  {k:2d} steps:",
          f"{acc(WA,bA,pgd(WA,bA,Xte,yte,0.15,k),yte):.4f}")
print("  FGSM (1 step, alpha=eps):",
      f"{acc(WA,bA,fgsm(WA,bA,Xte,yte,0.15),yte):.4f}")
\end{verbatim}

The actual output, as it stands, for the block above:

\begin{verbatim}
clean test accuracy: A 0.9799  B 0.9849  D 0.9648
  eps  A clean  A FGSM  A PGD40  D FGSM  D PGD40
 0.02    0.9799  0.9548   0.9548  0.9497   0.9497
 0.05    0.9799  0.8476   0.8425  0.9162   0.9129
 0.10    0.9799  0.4992   0.4590  0.8241   0.8174
 0.15    0.9799  0.1491   0.1072  0.6667   0.6265
 0.20    0.9799  0.0184   0.0050  0.4858   0.3719
 0.30    0.9799  0.0000   0.0000  0.0687   0.0201
random signs eps=0.15: 0.9548
PGD steps, eps=0.15, alpha=eps/10, random start
   1 steps: 0.9464
   2 steps: 0.8878
   5 steps: 0.6667
  10 steps: 0.2479
  20 steps: 0.1122
  40 steps: 0.1072
  FGSM (1 step, alpha=eps): 0.1491
\end{verbatim}

The first reading of that table: clean accuracy $97.99$ percent, and at
$\varepsilon=0.15$ only $14.91$ percent is left.

The second reading: the gap between FGSM and PGD widens in the middle, from zero at
$0.02$ to $4.19$ points at $0.15$. At a small budget one step already touches the
wall of the box; at a large budget several steps are needed to find a worse
direction.

The third reading, and the most important one for this chapter: the row of random
signs. If a perturbation of size $\varepsilon=0.15$ is given with random signs on
every pixel, accuracy only falls to $95.48$ percent. FGSM with the same size lowers it
to $14.91$ percent. So the choice of signs by the gradient is worth about eighty
points of accuracy: it is the direction that works, not the size.

The fourth reading: at $\varepsilon=0.30$, $A$ is already at zero percent while $D$ is
still at $6.87$ percent under FGSM and $2.01$ percent under PGD. Adversarial training
does not make the model immune; it shifts the curve to the right.

The next block measures two things that are not visible in the table: whether the
gradient we compute really is the gradient, and how much it costs.

\begin{verbatim}
i0 = int(np.argmax(np.abs(grad_x(WA, bA, Xte, Yte)).sum(1)))
x0, y0 = Xte[i0:i0 + 1], Yte[i0:i0 + 1]
g0 = grad_x(WA, bA, x0, y0)
L1 = lambda z: loss_grad(WA, bA, z, y0)[0]
print(f"finite difference at sample i={i0}, h=1e-6")
for j in np.argsort(-np.abs(g0[0]))[:3]:
    e = np.zeros_like(x0); e[0, j] = 1e-6
    fd = (L1(x0 + e) - L1(x0 - e)) / 2e-6
    print(f"  j={j:2d} analytic {g0[0,j]:+.6e} finite-diff {fd:+.6e}"
          f" error {abs(fd-g0[0,j])/abs(g0[0,j]):.0e}")
g_all = grad_x(WA, bA, Xte, Yte)
lin = 0.15 * np.abs(g_all).sum(1).mean()
lc = loss_grad(WA, bA, Xte, Yte)[0]
la = loss_grad(WA, bA, fgsm(WA, bA, Xte, yte, 0.15), Yte)[0]
print(f"mean |g|_1 {np.abs(g_all).sum(1).mean():.4f};"
      f" linear prediction {lin:.4f} nats; actual {la-lc:.3f} nats")

t0 = time.perf_counter(); fgsm(WA, bA, Xte, yte, 0.15)
t1 = time.perf_counter(); pgd(WA, bA, Xte, yte, 0.15, 40)
t2 = time.perf_counter()
print(f"time for 597 examples: FGSM {t1-t0:.4f} s  PGD40 {t2-t1:.4f} s")
\end{verbatim}

Its output:

\begin{verbatim}
finite difference at sample i=145, h=1e-6
  j=53 analytic -6.391833e+00 finite-diff -6.391833e+00 error 2e-11
  j=38 analytic +6.211518e+00 finite-diff +6.211518e+00 error 1e-10
  j=13 analytic -6.175321e+00 finite-diff -6.175321e+00 error 8e-11
mean |g|_1 0.0056; linear prediction 0.0008 nats; actual 4.862 nats
time for 597 examples: FGSM 0.0006 s  PGD40 0.0302 s
\end{verbatim}

The first three lines confirm that our backward pass is correct: the analytic gradient
and the finite difference with $h=10^{-6}$ agree to about ten digits, with relative
errors $10^{-10}$ and $10^{-11}$. For a network that uses ReLU, agreement that smooth
is reasonable: at the scale $10^{-6}$ no unit changes sign, so the function really is
linear around that point. If these numbers differed a lot, what would be wrong is the
backward-pass code, not the attack.

The fourth line is the surprising part. The average $\lVert\nabla L\rVert_1$ is only
$0.0056$, so the linear prediction for $\varepsilon=0.15$ is $0.0008$
nats. The actual rise in the error is $4.862$ nats. So the linear term
underestimates the damage by about four thousand times. This is not a sign that the
gradient is wrong (the three lines above prove the opposite), but a sign that $L$ is
strongly curved at the distance $0.15$: ReLU units change sign, and the first-order
approximation loses its meaning. What remains useful is the sign, as the row of random
signs showed.

The last line is about time: for 597 examples, FGSM needs $0.6$ milliseconds and
PGD of 40 steps $30$ milliseconds, about fifty times more expensive. This number
depends on the machine and on vectorization, so what is worth remembering is its
order of magnitude, not the exact number.

Finally, a demonstration of gradient obfuscation. We add input quantization to
multiples of $1/8$ in front of model $A$, then attack in two ways.

\begin{verbatim}
def q(z): return np.round(z * 8) / 8
accq = lambda W, b, Z, yy: float(np.mean(fwd(W, b, q(Z))[0]
                                        .argmax(1) == yy))
def pgd_naif(W, b, Z, yy, eps, steps):
    r = np.random.default_rng(1)
    z = np.clip(Z + r.uniform(-eps, eps, Z.shape), 0, 1)
    for _ in range(steps):          # gradient = zero, so it stands still
        z = np.clip(np.clip(z, Z - eps, Z + eps), 0, 1)
    return z
def pgd_bpda(W, b, Z, yy, eps, steps):
    r = np.random.default_rng(1)
    z = np.clip(Z + r.uniform(-eps, eps, Z.shape), 0, 1)
    for _ in range(steps):
        g = np.sign(grad_x(W, b, q(z), np.eye(K)[yy]))
        z = np.clip(np.clip(z + eps / 10 * g, Z - eps, Z + eps), 0, 1)
    return z
print("inputs quantized to multiples of 1/8")
print(f"  clean {accq(WA,bA,Xte,yte):.4f}"
      f"  PGD40 naive {accq(WA,bA,pgd_naif(WA,bA,Xte,yte,0.15,40),yte):.4f}"
      f"  PGD40 BPDA"
      f" {accq(WA,bA,pgd_bpda(WA,bA,Xte,yte,0.15,40),yte):.4f}")
\end{verbatim}

Its output:

\begin{verbatim}
inputs quantized to multiples of 1/8
  clean 0.9782  PGD40 naive 0.9698  PGD40 BPDA 0.1323
\end{verbatim}

This is the most concise form of a false defense. Its clean number is only $0.17$
points below the original model, and the naive attack does not move at all: if that
were all that was reported, this defense would look almost perfect. The substitute
gradient brings it down to $13.23$ percent, almost the same as $10.72$ percent on the
model without quantization. What changes above all is the readability of its gradient.

\section{Cost, limitations, and how to report}
\tingkat{basic}

The computational budget determines which attack is honestly used. FGSM requires one
backward pass per example. PGD of $k$ steps requires $k$ passes. In our experiment, the
cost of PGD of 40 steps is about fifty times that of FGSM, for a result that, at small
$\varepsilon$, is almost the same. For the assessment of a defense, that cost is
reasonable: a robustness number measured with a weak attack is meaningless. For
training, that cost multiplies again: every epoch demands a new attack.

Adversarial training also has an accuracy price. Model $D$ loses $1.51$ points of
clean accuracy in exchange for a rise of about fifty points at $\varepsilon=0.15$.
Madry and colleagues report that larger models give better robustness, which shows
that robustness is not a free property of the architecture, but capacity that has to
be paid for.

There are three limitations before the numbers in this chapter are used.

First, scale. Our data has only 64 features and 10 classes, with 597 test examples. On
large images, a meaningful budget is much smaller than $0.15$; on ImageNet,
Goodfellow and colleagues used $0.007$ and already flipped the classification. The
number $\varepsilon$ cannot be moved between datasets, because its meaning depends on
the scale of the features. On our data, the pixels are already quantized to multiples
of $1/16$, so that $\varepsilon=0.05$ means less than one quantization step per
pixel.

Second, the coverage of attacks. What we run is only two families: FGSM and PGD with
$\ell_\infty$. A model that resists both is not necessarily resistant to an
$\ell_2$ attack, to an attack that changes a single pixel, or to an attack that sticks
an object into the real world. Every robustness claim holds only for the family of
attacks that was tested, and that is why the evaluation protocol in the previous part
demands several kinds of attack at once.

Third, there is no guarantee. All the numbers here are results of testing, not a
proved lower bound. Model $D$, which reaches $62.65$ percent under PGD of 40 steps,
may still be broken by a better attack; we do not know, because what was not found is
not proof that it does not exist. For a stronger statement, certification is needed:
for example with randomized smoothing or with formal verification on a small network.
Both give a guarantee on a certain radius, at the price of far higher accuracy and
computation cost, and both are outside the scope of this chapter.

An honest way of reporting contains five things: the complete threat model (norm,
$\varepsilon$, the knowledge of the attacker, targeted or not), the attacker used
together with the number and size of its steps, the clean accuracy and the accuracy at
several $\varepsilon$, the number of seeds and the spread of the results, and the
computational cost. Without that, the number ``90 percent robust'' cannot be checked
and cannot be compared. Chapter 27 has already taught us that metrics can deceive;
here the metric deceives in a subtler way, because its denominator is right and only
its numerator is chosen.

\begin{untukdosen}
The teaching order that has worked: start from the table of accuracy against
$\varepsilon$, not from the definition. Students see $97.99$ and $14.91$ on the same
model, then ask why. Only after that come the definition of the threat model, the
lemma on the maximum of a linear form, and FGSM. Only after FGSM runs do we enter PGD
and the saddle formulation. The reverse order, that is definitions first and then
attacks, makes students memorize notation without ever seeing the damage.

Two typical misconceptions. First, students think an adversarial attack is the same as
adding noise, so they conclude that a model resistant to noise will be resistant to
attacks. The row of random signs in the lab refutes that with one number: noise of
size $\varepsilon=0.15$ only lowers accuracy to $95.48$ percent, whereas FGSM with the
same size lowers it to $14.91$ percent. Second, they think one accuracy number under
an attack is already evidence of robustness. The quantization demonstration in the lab
shows the opposite: $96.98$ percent looks safe until someone uses a substitute
gradient and gets $13.23$ percent.

A cheap blackboard activity: write the four-feature linear example from the worked
example part, then ask the class to compute $\varepsilon^{*}$ for $\ell_\infty$ and
$\ell_2$. After that ask them to predict, without computing, whether adding PGD steps
changes the result on that linear model. The correct answer, together with its reason,
forces them to understand that iteration is useful only because $f$ is not linear.

A sign of readiness to go on: students can explain why only the sign of the gradient
is used, can compute $\varepsilon^{*}$ for a linear model, can name one symptom of
gradient obfuscation, and can design one attack adapted to a defense that quantizes
the input. If they demand numbers at several $\varepsilon$ before believing, they are
ready to enter Chapter 30.
\end{untukdosen}

\section{Summary}

\begin{ringkasan}
\begin{itemize}
  \item A robustness claim is always relative to the threat model: the norm, the
        budget $\varepsilon$, the knowledge of the attacker, and whether the goal is
        targeted or not. Without that, two numbers cannot be compared.
  \item FGSM shifts every coordinate by $\varepsilon$ following the sign of the
        gradient of the error. On a linear model it is exact and optimal, because
        $\max_{\lVert\boldsymbol{\delta}\rVert_\infty\le\varepsilon}
        \mathbf{w}^{\top}\boldsymbol{\delta}=\varepsilon\lVert\mathbf{w}\rVert_1$.
  \item Linearization justifies the direction, not the distance. In our experiment the
        linear prediction is $0.0008$ nats, whereas the actual rise is $4.862$ nats;
        what remains useful is the pattern of signs.
  \item PGD adds iteration and projection onto the box. At $\varepsilon=0.15$
        accuracy falls from $94.64$ percent (1 small step) to $10.72$ percent
        (40 steps), and is flat after about twenty steps.
  \item Adversarial training is an approximation of the saddle problem
        $\min_\theta\mathbb{E}\max_\delta L$. It raises robustness at the price of
        clean accuracy: our model loses $1.51$ points to rise from $10.72$
        to $62.65$ percent.
  \item Gradient obfuscation lowers the success of attacks without raising the distance
        to the decision boundary. Our example: $96.98$ percent under the naive
        attack, $13.23$ percent under the attack with a substitute gradient.
  \item An honest evaluation uses the strongest attack that can be afforded, adapted
        to the defense, at several $\varepsilon$ and several seeds, with the cost
        reported.
\end{itemize}
\end{ringkasan}

\section{Exercises}

\begin{latihan}
\begin{enumerate}
  \item Compute $\varepsilon^{*}$ for the linear model with
        $\mathbf{w}_A=(1;-1;2)$, $\mathbf{w}_B=(0;1;1)$ and
        $\mathbf{x}=(1;2;1)$ under the $\ell_\infty$ norm, then under the $\ell_2$ norm.
        Check both answers with Python. (Hint: compute
        $\mathbf{d}=\mathbf{w}_A-\mathbf{w}_B$ first, then compare
        $\lVert\mathbf{d}\rVert_1$ with $\lVert\mathbf{d}\rVert_2$.)

  \item Prove that $\max_{\lVert\boldsymbol{\delta}\rVert_2\le\varepsilon}
        \mathbf{w}^{\top}\boldsymbol{\delta}=\varepsilon\lVert\mathbf{w}\rVert_2$,
        and determine all $\boldsymbol{\delta}$ that attain the maximum. Explain
        why the answer has a different form from the $\ell_\infty$ case.

  \item Prove that on a linear classifier, PGD with any number of
        steps never produces an error larger than FGSM
        with the same $\varepsilon$. (Hint: the error is a linear function of
        $\boldsymbol{\delta}$ on a convex set; the maximum of a linear function on a
        compact convex set is attained at an extreme point. Show that an extreme point
        of the $\ell_\infty$ ball has the form $\pm\varepsilon$ in every coordinate,
        then compare with the FGSM point.)

  \item Assess this claim: \emph{because a model reaches accuracy $98$ percent on
        test data and stays at $95$ percent under random perturbation of size
        $\varepsilon$, that model is resistant to adversarial attacks.} Decide true,
        false, or not enough evidence, and name one number from this chapter that
        supports your decision.

  \item From the lab table, compute the success rate of the untargeted attack for
        model $A$ at $\varepsilon=0.05$, $0.10$, and $0.15$, then for
        model $D$ at the same budgets. State the results in percent and
        explain why the gap between $A$ and $D$ narrows at large $\varepsilon$.

  \item Show that the projection onto
        $\{\mathbf{y}:\lVert\mathbf{y}-\mathbf{x}\rVert_\infty\le\varepsilon\}
        \cap[0,1]^{d}$ may be done by clipping first to the $\varepsilon$ box and then
        to $[0,1]$, or the other way round, and the result is the same.
        (Hint: both sets are products of intervals, so it is enough to check
        one coordinate.)

  \item On model $D$, the clean accuracy is $96.48$ percent and the accuracy under PGD
        of 40 steps at $\varepsilon=0.15$ is $62.65$ percent. On model $A$,
        the numbers are $97.99$ and $10.72$. Compute the price of robustness in
        points of clean accuracy per point of robustness gained, then explain why that
        ratio is not a complete measure.

  \item A defense rounds the probability output to $0$ or $1$
        before the decision is taken, so that the gradient with respect to the input
        becomes zero. Design one evaluation protocol to test its robustness claim, and
        name which symptom from the list in this chapter is most easily seen if that
        defense only obfuscates the gradient.

  \item Given two models with clean accuracy $97.99$ and $98.49$ percent.
        At $\varepsilon=0.15$, PGD of 40 steps built on the first model lowers
        the accuracy of the first model to $10.72$ percent and of the second model to
        $23.79$ percent. Compute the gap between the success of the white-box attack
        and of the transfer attack, then explain why that gap is not evidence that the
        second model is safer. (Hint: an attack built specifically for the second model
        is also available; see row B in the lab output.)

  \item Explain why a targeted attack is harder than an untargeted
        attack, then use the numbers $33.00$ percent at $\varepsilon=0.15$ and
        $99.50$ percent at $\varepsilon=0.30$ to show that the difficulty
        depends on the budget, not only on the kind of attack.
\end{enumerate}
\end{latihan}

\begin{untukpeneliti}
What is proved in this chapter is little. The lemma on the maximum of a linear form is
a one-line proof. The linearization inequality with remainder
$M\lVert\boldsymbol{\delta}\rVert_2^2/2$ is the mean value theorem in integral form,
valid for every differentiable function with a Lipschitz gradient. The exactness of
FGSM on logistic regression is a result of Goodfellow, Shlens, and Szegedy (ICLR 2015),
the section ``Adversarial Training of Linear Models versus Weight Decay''. The saddle
formulation and the claim that PGD can be relied on come from Madry, Makelov, Schmidt,
Tsipras, and Vladu (ICLR 2018), sections 2 and 3; the latter is an empirical claim,
not a theorem, and its own authors call the inner problem nonconcave.

What is more often popularized than proved is the claim that linearity explains
adversarial examples. As an argument about direction it is useful; as an explanation of
the size of the damage it fails in our experiment: a prediction of $0.0008$ nats
against an actual $4.862$ nats, with a gradient that has been verified to $10^{-10}$.
What works is the pattern of signs, and that is why FGSM beats random signs
($14.91$ against $95.48$ percent). A second claim that is exaggerated: adversarial
training gives robustness to the attacker and the budget used during training, at the
price of clean accuracy, without a guarantee for other attacks.

Limits of validity: 64 features, 597 test examples, one hidden layer, one seed, two
families of $\ell_\infty$ attacks, without certification. The numbers here demonstrate
a mechanism, not a ranking of robustness. Open questions: how large a radius can be
certified without erasing the benefit in accuracy; why transfer works that well;
and what threat model makes sense for real systems such as intrusion detectors, which
is the subject of Chapter 30.
\end{untukpeneliti}

\peta{This chapter follows Goodfellow, Shlens, and Szegedy, \emph{Explaining and
Harnessing Adversarial Examples} (ICLR 2015), sections 3 (``The Linear Explanation of
Adversarial Examples''), 4 (``The Fast Gradient Sign Method''), and 5 (``Adversarial
Training of Linear Models versus Weight Decay'') for the derivation of FGSM, its
exactness on logistic regression, and the MNIST and CIFAR-10 numbers. The saddle
formulation, the PGD algorithm, and the MNIST and CIFAR-10 robustness numbers come from
Madry, Makelov, Schmidt, Tsipras, and Vladu, \emph{Towards Deep Learning Models
Resistant to Adversarial Attacks} (ICLR 2018), sections 2, 3, and 4, including
Algorithm 1. The discussion of gradient obfuscation follows Athalye, Carlini, and
Wagner, \emph{Obfuscated Gradients Give a False Sense of Security} (ICML 2018),
sections 2 to 4, including the three kinds of obfuscation and the five symptoms used
here. The phenomenon of transferability was first reported by Szegedy and colleagues,
\emph{Intriguing Properties of Neural Networks} (ICLR 2014). All the numbers in the
tables, the lab, and the worked example are our own computation with Python 3.14,
NumPy 2.5, and the \texttt{load\_digits} dataset from scikit-learn 1.9: Adam training
for 350 epochs, a train-test split of 1200/597 with seed 0, and input gradients
verified with finite differences.}

% The full chapter title is too wide for the running head, so the header mark is
% shortened once (see \judulpendek in preamble-en.tex). The chapter page and the
% table of contents keep the complete title. The macro must come *before* \chapter:
% \chaptermark consumes it, so placing it after leaks the short title to the next chapter.
\judulpendek{The Limits of Attack}
\chapter{The Limits of Attack: Obfuscated Gradients and NIDS Evasion}

A group of students added one line of rounding to the model input, then reported that their own
attack failed completely: the success rate fell from seventeen percent to zero point zero two
percent. That number went into the executive summary, then into the slides, then into the abstract.
Two weeks later an examiner asked for one thing: prove that the gradient the attacker used still
points in the right direction. It did not. Rounding is flat almost everywhere, so its derivative is
zero, so the attacker takes a zero step, so the attack that failed is really an attack that never
happened. When the examiner replaced the derivative with a reasonable approximation, the success
rate returned to seventeen percent within minutes.

This chapter is about that mistake. The first mistake is simple: rounding erases the gradient, and
erasing the gradient means blinding the attacker, not making the model strong. The second mistake is
more interesting: there are times when obfuscating the gradient really does raise the cost of
attack, so part of that robustness is real. Telling the two apart demands experiment, not rhetoric.

The second half of the chapter carries the same problem to a much noisier place: network intrusion
detection systems, or NIDS. There packets must still arrive, protocols must remain valid, decisions
are made in milliseconds, and the model is retrained every day. All of that makes an evasion attack
that is beautiful on paper impractical in the field. But ``impractical'' is not ``impossible''.

\section{Gradients that vanish, and the sense of safety born from them}
\tingkat{basic}

Chapter 29 taught two attacks that rest on one object: the gradient of the loss with respect to the
input. The one-step attack takes the direction of the sign of the gradient and steps by
$\varepsilon$,

\begin{equation}
x' = x + \varepsilon\,\operatorname{sign}\bigl(\nabla_x L(x,y)\bigr),
\end{equation}

whereas the many-step attack repeats a small step and projects the result back onto the
$\ell_\infty$ ball of radius $\varepsilon$,

\begin{equation}
x_{t+1} = \Pi_{B(x,\varepsilon)}\Bigl(x_t - \alpha\,\operatorname{sign}
\bigl(\nabla_x L(x_t,y)\bigr)\Bigr).
\end{equation}

Both work for one reason that already appeared in the chapter on linearization: for a small
perturbation the loss can be approximated by a linear function, so the direction that raises the loss
fastest is the direction of the gradient. If that gradient is correct, the attack works at the cost
of one backward pass. If that gradient is wrong, the attacker goes blind. Those two sentences are not
equivalent to ``if the gradient is wrong, the model is safe''. A blind attacker can still guess, and
a guess without a gradient, for example random noise on packets, can still succeed in some cases.
What is lost is only the cheapest road, not every road.

\index{obfuscated gradient}
The term \emph{obfuscated gradient} names the situation in which the gradient an attacker sees does
not give enough information about the direction that actually raises the loss, not because the model
is strong, but because of the way the model is wrapped. Athalye, Carlini, Wagner and colleagues gave
this symptom a sharper name: a false sense of security. In their review of nine defenses from a
machine learning conference, seven depended on an obfuscated gradient; six were broken completely and
one partially, at an attack cost roughly the same as an ordinary attack.

\begin{untukanak}
Imagine a teacher who does not want her students to cheat, so she seals the notebook shut with duct
tape. The honest student does have trouble. But the creative student does not stop: he notices that
the answer to question one is always larger than the answer to question two, then guesses the whole
sheet from that pattern. The tape does not erase the answers; it only blocks one road to them.

A model that rounds its input works like that tape. The numbers that enter the model are cut to a
fixed multiple, for instance two decimal places, so that small changes are invisible altogether. A
program that tries to find the ``raise it a little at a time'' direction will find that nothing
changes, then give up and report that the model is immune. That report is wrong, and not because the
program is stupid. The program simply asked in a way that does not fit.

Now the part that matters most to remember. The same model can still be wrong if the numbers are
changed enough to cross a cut boundary. So the right question is not ``does a small attack succeed'',
but ``how large a change is needed''. This chapter teaches how to ask that question, and how to
recognize a report that never asks it.
\end{untukanak}

Why does rounding erase the gradient? Because rounding is a step function. Take
$q>0$ and define $Q_q(t)=q\,\mathrm{round}(t/q)$, the nearest multiple of $q$. On
the interval $kq-q/2 \le t < kq+q/2$ this function is constant at $kq$, so its derivative is zero
inside the interval and undefined at the boundary points. The set of boundary points has measure
zero, so for any smooth function $f$, the composition $f\circ Q_q$ has zero gradient almost
everywhere.

\begin{proposisi}
Let $Q_q\colon\mathbb{R}^d\to\mathbb{R}^d$ act coordinate by coordinate as above, and let
$f\colon\mathbb{R}^d\to\mathbb{R}$ be continuously differentiable. Then
$\nabla (f\circ Q_q)(x)=0$ for almost every $x\in\mathbb{R}^d$.
\end{proposisi}

\begin{proof}
Partition $\mathbb{R}^d$ into half-open boxes of side $q$, bounded by the
hyperplanes $x_j=k_jq\pm q/2$. In the interior of every box, $Q_q$ is
constant, so $f\circ Q_q$ is also constant and its gradient is zero. The complement of the union of
all those interiors is the union of the boundary hyperplanes, which has measure zero. So the
gradient is zero almost everywhere.
\end{proof}

\begin{akibat}
A one-step attack that uses the exact gradient of $f\circ Q_q$ produces
$x'=x$ for every $\varepsilon>0$. The measurable success rate is therefore
exactly zero, and the reported ``accuracy under attack'' is precisely the clean accuracy.
\end{akibat}

This corollary is the heart of the whole chapter. The statement is not about the strength of the
model; it is about the arithmetic of the attacker, and the zero in the results table does not measure
robustness, but the absence of a step. If the ``accuracy under attack'' column is identical to the
clean accuracy down to two decimal places, you are not looking at a defense that worked, but at an
attacker that stopped.

\begin{bacakeras}
Two voices. The first carries a map, the second carries a flashlight.

--- Look. I have made the map I drew blurry. Every uphill road I erased, there are no arrows, no
numbers.
--- Then the mountain is gone.
--- No. The mountain is still there.
--- So what is gone is only the picture?
--- What is gone is only the picture.
--- Someone who can smell the soil can still climb.
--- Still can.
--- Someone who throws a pebble and listens to the echo can also still climb.
--- Also still can.
--- So your map is not a wall.
--- Your map is not a wall. Your map is only a map that refuses to answer.

Repeat the last line three times, slower and slower: a map that refuses to answer is not a wall. Not
a wall. Not a wall.
\end{bacakeras}

Athalye and colleagues divided the mechanism into three families. The first family, shattered
gradients, appears when part of the model is non-differentiable or differentiable but useless, such
as rounding, median filtering, and compression. The second family, stochastic gradients, appears when
the model randomizes something on every call, so that the gradient of one call does not hold for the
next. The third family, vanishing or exploding gradients, appears when the defense suppresses the
gradient during training, so that the gradient becomes too small or too large to use for a step.

All three share one feature: they make the measurement look better without making the model better,
and because the measurement is usually the only evidence that is reported, this error spreads fast.
This chapter does not say that all such defenses are useless; some raise the cost of attack, and that
cost is real. What is wrong is concluding robustness from a broken measurement.

\section{Seven ways to obfuscate the gradient}
\tingkat{intermediate}

The table below summarizes the seven techniques that appear most often, their mechanism family, and
the symptom that ought to make you suspicious. The numbers in the last column are not a diagnosis,
but a trigger for a question.

\begin{center}
\small
\begin{tabular}{@{}p{0.6cm}p{4.6cm}p{2.1cm}p{6.4cm}@{}}
\hline
No & Technique & Family & Visible symptom \\
\hline
1 & Input or output quantization & shattered & zero gradient; attack accuracy equals clean accuracy \\
2 & Non-differentiable preprocessing: median, blur, compression & shattered & zero gradient, or spikes with no pattern \\
3 & Sample filtering and rejection & shattered & discrete decisions; no gradient \\
4 & Random transforms: resize, pad, shift & stochastic & a different gradient on every call \\
5 & Random classifiers and random ensembles & stochastic & one sample points the wrong way \\
6 & Gradient suppressed during training & vanishing & small gradient despite a large loss \\
7 & Multi-step optimization defenses & vanishing & exploding gradient; attack does not converge \\
\hline
\end{tabular}
\end{center}

\index{quantization}
The first technique, quantization, has already been proved above, and its effect is exact: not
approximately zero, but zero. The second technique works the same way but more dirtily. Median
filtering replaces every value with the median of its neighborhood, so the function stays flat in
most places, but the boundaries between its pieces are no longer evenly spaced. Compression such as
JPEG discards high-frequency components; the derivative exists but is unstable, and a finite
difference gradient can change sign when the step is shifted. The symptom is characteristic: the
gradient is not zero, but large-scale garbage.

The third technique is the crudest form of obfuscation. The model refuses to answer for inputs it
considers strange, for example those whose distance to the data center is too large, then returns the
label ``unknown''. For the attacker this is not a bad gradient but no gradient at all: the function
is an indicator function, whose derivative is zero everywhere except on the decision surface. A
defense of this kind also rejects legitimate new data, and legitimate data is the reason the system
was installed.

The fourth and fifth techniques are the stochastic family. In the fourth, the input is randomly
transformed before it enters the model: the image is enlarged, shifted by one pixel, or some of its
features are dropped. In the fifth, what is randomized is the model: one ensemble is chosen at random
on every call, or the prediction is taken from the average of many models. In both, the gradient of
one call is the gradient of one realization, whereas what the attacker faces is the distribution of
realizations. A bad gradient sample is not due to bias but to variance: the sign of a direction that
keeps changing is not a good direction.

\index{random filtering}
The sixth technique attacks from the inside. Training deliberately makes the
gradient small around the data points, for example by penalizing the gradient norm or
by using saturating activations, so that the loss surface is very flat around the examples
and steep elsewhere. An attacker that steps by $\varepsilon$ in a flat region
goes nowhere, and the gradient really is small, but not because the model is right.
The seventh technique is its relative: the input is first cleaned by a multi-step
optimization procedure before being classified, and the gradient of that pipeline easily
explodes.

Figure~\ref{fig:permukaan} shows the shape of the surface that quantization
produces on two coordinates. Every flat plateau there is one step. Inside
one plateau, the loss does not change at all, so an attacker that takes a small
step stops where it is --- the three short arrows in that figure. An attacker that is allowed
to step far enough can cross to the next plateau, and the dashed arrow
shows that crossing.

\begin{figure}[!htbp]
\centering
\begin{tikzpicture}
\begin{axis}[
  width=0.80\linewidth, height=7.2cm,
  view={52}{24},
  xlabel={$u$}, ylabel={$v$}, zlabel={$L$},
  domain=-0.35:0.35, y domain=-0.35:0.35,
  samples=27, samples y=27,
  zmin=0, zmax=0.52,
  tick label style={font=\scriptsize},
  label style={font=\small},
  colormap={peta}{color=(layerteal!10) color=(layerteal!60)},
  colorbar=false, axis line style={layergrey}, tick style={layergrey},
]
\addplot3[surf, shader=faceted, faceted color=layergrey!40, line width=0.15pt]
  {0.1*round(10*((x-0.15)^2+(y-0.10)^2))};
\draw[-latex, accent, line width=1pt] (axis cs:0.02,0.02,0) -- (axis cs:0.10,0.06,0);
\draw[-latex, accent, line width=1pt] (axis cs:0.02,0.02,0) -- (axis cs:-0.02,0.08,0);
\draw[-latex, accent, line width=1pt] (axis cs:0.02,0.02,0) -- (axis cs:0.04,-0.05,0);
\node[circle, fill=accent, inner sep=1.5pt] at (axis cs:0.02,0.02,0) {};
\draw[-latex, layerteal, line width=1.2pt, dashed]
  (axis cs:0.02,0.02,0) -- (axis cs:0.30,0.30,0.10);
\node[font=\footnotesize, accent, anchor=north west] at (axis cs:0.09,0.05,0)
  {$\nabla L=0$};
\node[font=\footnotesize, layerteal, anchor=south west] at (axis cs:0.28,0.28,0.12)
  {BPDA};
\end{axis}
\end{tikzpicture}
\caption{A loss surface made flat by quantization:
$L(u,v)=\frac{1}{10}\operatorname{round}(10\,s)$ with
$s=(u-0.15)^2+(v-0.10)^2$.
Every plateau of height $0.1$ is one step: inside it $\nabla L=0$,
so the three short arrows from the red point do not change the height at all.
The dashed arrow crosses to the next step, and only an attacker willing to step
as far as the distance between boundaries reaches it.}
\label{fig:permukaan}
\end{figure}

\begin{tangan}
Materials: one long ruler, one coin, one friend, one table.

Prop up one end of the ruler with a book so that its surface is tilted, then place the coin
in the middle. Close your eyes and hold the coin with one finger. Your friend taps the
table slowly, irregularly, and your task is to guess in which direction the ruler rises
--- only from the vibration that reaches your finger.

Repeat ten times and record right or wrong. Then ask your friend to tap harder,
and repeat ten times. Then ask him to stop tapping altogether, and repeat again.

What you should see: harder tapping does not make your guesses right. Tapping that stops
makes your guesses pure guesses, about half right, no better than tossing
a coin. The vibration tells you that something is moving; it does not tell you the direction.
That is the difference between a bad gradient and a gradient that does not exist.
\end{tangan}

Those seven techniques produce seven symptoms. The first five are the list of Athalye and
colleagues; the last two we add from the same typology.

\begin{enumerate}
  \item \textbf{The attack is unbounded.} The report allows a perturbation of any size,
        then concludes robustness from an attack that ignores the bound.
  \item \textbf{The attack uses a single random sample.} Against a defense that randomizes,
        the attacker draws one realization then reports its failure. What is measured is
        one throw of the die, not the probability of winning.
  \item \textbf{The attack is only one step.} A one-step attack is easily stopped by
        an uneven loss surface, and the robustness it implies disappears
        as soon as the step is repeated.
  \item \textbf{Increasing the distortion does not increase success.} If the budget is enlarged
        but the success does not move, what is broken is usually the gradient.
  \item \textbf{The gradient is reported as zero or NaN.} This is the most naked symptom: a gradient
        that is entirely zero on a non-constant model signals a non-differentiable
        piece, not a flat surface.
  \item \textbf{Accuracy under attack equals clean accuracy.} Our addition.
        Equality down to several decimal places at several different $\varepsilon$
        is the mark of an attack that does not move.
  \item \textbf{There is no adaptive attack.} Our addition. The defense is tested only
        with attacks designed before that defense existed, so it has not been tested
        by an attacker who knows that he is being tested.
\end{enumerate}

\kamus{obfuscated gradient}{gradien dikaburkan}{the gradient given to the attacker
does not point in the right direction, usually because there is a non-differentiable or
random piece inside the pipeline, not because the model is strong}

\kamus{quantization}{kuantisasi}{replacing every value with the nearest multiple of
a fixed distance; the result is a step function whose derivative is zero almost everywhere}

\kamus{adaptive attack}{serangan adaptif}{an attack designed with knowledge of
the defense being tested, including knowing that its gradient may be broken}

That list of symptoms is useful for rejecting weak claims but not for accepting strong ones: a defense
can pass all seven and still be fragile, and conversely it can trigger one symptom and still raise the
cost of attack.


\section{Breaking through: BPDA, EOT, and reparameterization}
\tingkat{advanced}

If the gradient that is available is wrong, the attacker has three options. Replace the wrong
gradient with the gradient of a function that approximates it; average the gradient over the
randomness that makes it change; or move the work to another space, where the defense obfuscates
nothing. All three have names, and all three work for different reasons.

\index{bpda}
The first option is called \emph{backward pass differentiable approximation}, abbreviated
BPDA. The idea is as simple as it sounds: on the forward pass use the actual function,
including the rounding or the filtering; on the backward pass skip that piece as if
it were the identity. If the wrapped model is $h(x)=f\bigl(g(x)\bigr)$,
the attacker uses

\begin{equation}
\widehat{\nabla}_x h(x) \;=\; \nabla f\bigl(g(x)\bigr),
\end{equation}

that is, the gradient of $f$ evaluated at the point after wrapping, not the gradient of the
composition. For rounding, this means the attacker stops asking ``what is the derivative of
rounding'' and starts asking ``if the numbers are already rounded, in which direction does $f$
rise fastest''.

Why is this approximation reasonable? Because the error is bounded by how far $g$
shifts the point, not by how strange the shape of $g$ is.

\begin{proposisi}
Let $f$ be continuously differentiable with a Lipschitz gradient, that is,
$\lVert \nabla f(u)-\nabla f(v)\rVert \le L\lVert u-v\rVert$ for all $u,v$, and
let $g$ satisfy $\lVert g(x)-x\rVert\le\delta$. Then
\begin{equation}
\bigl\lVert \nabla f\bigl(g(x)\bigr)-\nabla f(x)\bigr\rVert \;\le\; L\,\delta .
\end{equation}
\end{proposisi}

\begin{proof}
Write $v=g(x)-x$ and use the mean value theorem in integral form,
\begin{equation}
\nabla f(x+v)-\nabla f(x) = \int_0^1 H_f(x+tv)\,v\,dt,
\end{equation}
with $H_f$ the Hessian matrix. Because the spectral norm of $H_f$ is bounded by $L$,
\begin{equation}
\bigl\lVert \nabla f(x+v)-\nabla f(x)\bigr\rVert
\le \int_0^1 \lVert H_f(x+tv)\rVert\,\lVert v\rVert\,dt
\le L\lVert v\rVert \le L\delta .
\end{equation}
\end{proof}

This proposition explains when BPDA can be trusted. If rounding only shifts the
input by $q/2$ and $f$ is smooth enough at that scale, the gradient error is small and
the attack step still moves in the right direction. If $q$ is large so that the shift is
comparable to the scale of the curvature of $f$, the bound $L\delta$ becomes large and nothing is
guaranteed. This proposition does not say that $\widehat{\nabla}$ is the gradient of $h$; that
gradient does not exist. It only says that the direction used is close to the correct direction for
$f$ at the point that has already been wrapped.

The second option is \emph{expectation over transformation}, abbreviated EOT. Against
a stochastic defense, the attacker does not face one function but a distribution of functions.
What makes sense to maximize is not the loss on one realization, but the expected
loss over realizations,

\begin{equation}
\overline{L}(x) \;=\; \mathbb{E}_{T\sim\mathcal{T}}\; L\bigl(f_T(x),y\bigr),
\end{equation}

with $T$ denoting the randomness the defense uses, for example a feature mask or a
transformation parameter.

\index{eot}
\begin{lemma}
If $\mathcal{T}$ is finite and $\nabla_x L\bigl(f_T(x),y\bigr)$ is uniformly bounded
in $T$, then
\begin{equation}
\nabla_x \overline{L}(x) \;=\; \mathbb{E}_{T}\;\nabla_x L\bigl(f_T(x),y\bigr).
\end{equation}
Moreover, the average of $m$ independent samples,
$\frac{1}{m}\sum_{i=1}^{m}\nabla_x L\bigl(f_{T_i}(x),y\bigr)$, is an unbiased estimator
whose variance shrinks proportionally to $1/m$.
\end{lemma}

\begin{proof}
Because $\mathcal{T}$ is finite, the expectation is a finite sum, and the derivative may
enter the sum because the derivative is a linear operator. The sample average has the
same expectation as each of its terms, so it is unbiased, and its variance is the variance
of one term divided by $m$ because the samples are independent.
\end{proof}

This lemma answers an objection that comes up often: if the defense is random, is not any
attack just luck? No, as long as the attacker may call the model
repeatedly. Randomness raises the cost, because the attacker needs $m$ calls per step,
but it does not change the direction being pursued; what changes is only the precision of the
estimate. This is the first example of a pattern that recurs throughout the chapter: a defense that
raises the cost is real, whereas one that only makes measurement harder is illusory.

The third option is reparameterization. Instead of attacking $x$, the attacker attacks
the variable that produces $x$ through a transformation that is not obfuscated. If
the defense wraps the model with a function $g$ that has an inverse, the attacker works in the
space $z$ with $x=g^{-1}(z)$ and optimizes $z$. This is useful when $g$ is a
rotation or another linear transformation that does not destroy information, and it is not
useful when $g$ destroys information: no $z$ can recover numbers
that have already been thrown away.

In practice, the three options are used together, plus two habits. First: start
from several random initial points inside the $\varepsilon$ ball, because a loss surface that
is not convex makes a single initial point easy to trap. Second: do not stop at one
value of $\varepsilon$. Robustness measured at a single budget is robustness that is not
measured.

\begin{resep}
Four steps for testing a defense claim, in this order.
\begin{enumerate}
  \item \textbf{Write the budget first.} Fix $\varepsilon$ or the byte limit,
        the number of queries, and the time. These numbers must not change after the first
        result comes out.
  \item \textbf{Attack by transfer, without queries.} Train a substitute model on the data
        you have, build adversarial examples there, then send them to the victim model.
        This step measures whether the defense merely blocks one model.
  \item \textbf{Spend the remaining queries adaptively.} Only now use BPDA or
        EOT on the victim, and record the queries per step, not just the final result.
  \item \textbf{Report both.} Transfer success, adaptive success, and the cost of each,
        always together with clean accuracy.
\end{enumerate}
A warning about order: step 2 must be finished before step 3. If you measure
transfer after seeing the adaptive results, you already know which examples succeed,
and your transfer numbers leak from the future.
\end{resep}

\kamus{BPDA}{backward pass differentiable approximation}{a technique that replaces the gradient
of a non-differentiable piece with the gradient of the function at the point after that piece,
as if the piece were the identity}

\kamus{EOT}{expectation over transformation}{an attack that maximizes the expected
loss over the randomness of the defense, usually by averaging gradients over several
realizations}

\begin{untukdosen}
This part is usually taught in the wrong order, namely the definition of BPDA first.
A better order is the reverse: show a model whose accuracy does not move at all
under attack, ask the students to guess why, then show that its gradient is
zero. Almost all of them will guess ``the model is strong'', and the silence that follows is a
teaching moment in itself.

Two mistakes need to be anticipated. First, concluding that BPDA ``removes''
the defense; what is true is that BPDA restores access to an approximately correct direction, and
its usefulness depends on how large $\delta$ is relative to the scale of the model curvature.
Second, concluding that stochastic defenses are useless because EOT exists; what is true is that
a stochastic defense raises the per-step cost linearly in $m$, so that the defense changes the
attacker's budget, not the direction of the attack. The exercises at the end of the chapter
force students to compute that cost, not merely name it.
\end{untukdosen}

\section{Worked example: one point on a toy model}
\tingkat{intermediate}

\begin{untukanak}
A shopkeeper has an odd habit: every price is rounded to the nearest rupiah
before it is recorded. If a customer quietly shifts a price by one cent,
the shopkeeper's record does not change at all. That customer can conclude that there is no
point in shifting prices, then stop trying. The conclusion is wrong, but
the reasoning is sensible: what he measured really did not move.

Now imagine another customer. She does not ask whether the record changes because of a one-cent
shift; she asks in which direction the record rises fastest, then shifts the
price far enough that the rounding changes. The shopkeeper's record changes, and
she succeeds without breaking any rule. What separates the two is not cleverness,
but the question that was asked.

A model that uses rounding works like that shopkeeper. A program that only tries
small shifts will stop and report that the model is safe; a program that understands
rounding will step further. So when you read a report that a model
cannot be defeated, ask first: did the attacker stop because it lost, or because it
only tried one way?
\end{untukanak}

Before trusting the table, let us compute one point by hand. The model is made as
simple as possible: eight features with all weights equal to one, so the score is the
sum of the eight features. The data come from two normal distributions centered at $-0.5$ and
$+0.5$, with clean accuracy around $92$ percent. The defense is rounding to
multiples of $q=0.05$ on every feature, plus random filtering with drop probability $r$
that multiplies the rest by $1/(1-r)$. The attacker budget is $\varepsilon=0.25$ in the
$\ell_\infty$ norm.

Take one example that is correctly labeled $+1$,

\begin{equation}
x_0 = (0.30;\ 0.20;\ 0.15;\ 0.10;\ 0.05;\ 0.00;\ -0.10;\ -0.08).
\end{equation}

The sum is $0.62$, so the model answers $+1$ with a thin margin. After
rounding, the vector that enters the classifier is
$(0.30;0.20;0.15;0.10;0.05;0.00;-0.10;-0.10)$ and the score is $0.60$.
Note that one coordinate shifts from $-0.08$ to $-0.10$: that is the work of
rounding.

Now compare two attackers. The naive attacker computes the derivative of the actual
function, that is, the derivative of the composition that contains the rounding. Because rounding
is flat around $x_0$, the derivative is zero in all eight coordinates, so its step is
$x_0 + \varepsilon\cdot 0 = x_0$. The score stays $0.60$ and the label stays $+1$.
That attacker's report will read: attack failed, robust accuracy one hundred
percent of clean accuracy.

The second attacker uses BPDA. It does not ask about the derivative of rounding; it
evaluates the gradient of the classifier at the point that has already been rounded, that is, the
vector $(1,1,\dots,1)$, then steps against the label direction by $\varepsilon$ in each
coordinate,

\begin{equation}
x_1 = x_0 - 0.25\cdot(1,1,\dots,1)
= (0.05;\ -0.05;\ -0.10;\ -0.15;\ -0.20;\ -0.25;\ -0.35;\ -0.33).
\end{equation}

After rounding, that vector becomes
\[
(0.05;\ -0.05;\ -0.10;\ -0.15;\ -0.20;\ -0.25;\ -0.35;\ -0.35)
\]
with score $-1.40$. The label flips to $-1$. One step, the cost of one
gradient computation, and the model that was just reported immune turns out to be wrong with a margin
that is not thin.

This example shows something sharper than ``BPDA works''. The second attacker
does not use information that the first does not have: both know the
architecture, the weights, and $q$. What differs is one decision, namely whether to respect the
non-differentiable piece or skip over it. So the ``robustness'' reported by the first attacker
is not a property of the model, but a property of that attacker's notation.

\begin{contoh}
On the same model, with $q=0.05$, $r=0.5$, and a fixed seed, a finite
difference gradient with step $h=10^{-5}$ is zero in $15\,998$ of $16\,000$ entries;
the two remaining entries are exactly $q/h=5\,000$, because their points happen to sit at a step
boundary. The median finite difference gradient norm is zero, whereas the norm of the
BPDA gradient for one mask that keeps four features is $4$. But the direction of a single
mask is still far from the direction of the weights, with a cosine between $0.35$ and $0.94$ over
$200$ masks. The average gradient over $200$ masks gives a vector of norm $2.738$ and
cosine $0.9982$ with respect to the weight vector.
\end{contoh}

Those numbers illustrate three things. First, a shattered gradient is not merely small; it can
also be large and wrong, and two entries that spike will dominate every step
so that the naive attack fails in a way that looks random. Second, BPDA restores the direction
almost perfectly: cosine $0.9982$ means an angle of about three degrees. Third,
the norm of the BPDA gradient for a single mask, namely $4$, differs from its expected norm, $2.738$;
that difference in scale does not hurt an attack that uses the sign, but it does hurt one that uses
the magnitude of the gradient, and that is why EOT is still needed.

\kertas{Recompute the example above without looking at the text. Sum $x_0$, round to
multiples of $0.05$, then sum again; check that the result is $0.60$. Subtract
$0.25$ from each coordinate, round again, sum, and check that the result is
$-1.40$. After that, answer: what is the smallest value of $q$ that makes this one-step BPDA
attack fail to flip the label, and why does that value depend on the initial margin
$0.60$ and not on $\varepsilon$ alone?}

\section{Computational lab: accuracy at every level of obfuscation}
\tingkat{advanced}

\begin{figure}[!htb]
\centering
\begin{tikzpicture}
\begin{axis}[
  width=0.86\linewidth, height=5.6cm,
  xlabel={budget $\varepsilon$ in the $\ell_\infty$ norm},
  ylabel={attack success (\%)},
  xmin=0.03, xmax=0.37, ymin=-0.6, ymax=30,
  xtick={0.05,0.10,0.15,0.20,0.25,0.30,0.35},
  xticklabels={0.05,0.10,0.15,0.20,0.25,0.30,0.35},
  ytick={0,5,10,15,20,25,30},
  tick label style={font=\footnotesize},
  label style={font=\small},
  legend style={font=\scriptsize, at={(0.02,0.98)}, anchor=north west,
                draw=layergrey!35, fill=white, fill opacity=0.92, text opacity=1},
  grid=major, grid style={layergrey!15},
  axis line style={layergrey},
]
\addplot[color=layerindigo, mark=*, mark size=1.5pt, line width=1pt]
  coordinates {(0.05,2.38) (0.10,5.27) (0.15,8.66) (0.20,12.84)
               (0.25,17.40) (0.30,22.32) (0.35,27.71)};
\addlegendentry{no rounding, BPDA attacker}
\addplot[color=accent, mark=square*, mark size=1.4pt, line width=1pt, dashed]
  coordinates {(0.05,2.34) (0.10,5.33) (0.15,8.72) (0.20,12.92)
               (0.25,17.48) (0.30,22.46) (0.35,27.79)};
\addlegendentry{rounding $q=0.05$, BPDA attacker}
\addplot[color=layergrey, mark=triangle*, mark size=1.6pt, line width=0.9pt,
         dotted]
  coordinates {(0.05,0.00) (0.10,0.01) (0.15,0.01) (0.20,0.01)
               (0.25,0.01) (0.30,0.02) (0.35,0.02)};
\addlegendentry{rounding $q=0.05$, naive attacker}
\end{axis}
\end{tikzpicture}
\caption{Attack success against $\varepsilon$ on the eight-feature toy model,
$2\,000$ examples, average of five seeds. At $\varepsilon=0.25$ the three curves give
$17.40$, $17.48$, and $0.01$ percent; at $\varepsilon=0.35$ they give
$27.71$, $27.79$, and $0.02$ percent. The first two curves coincide within the
seed-to-seed deviation, which is $0.35$ to $1.08$ percentage points.}
\label{fig:kurva-epsilon}
\end{figure}

Now we repeat that experiment on two thousand examples, five seeds, and several levels of
obfuscation, with \texttt{numpy} alone. What matters is not the code itself,
but one decision inside it: the function \texttt{fgsm} computes the gradient by finite
differences on the function that has already been wrapped, whereas \texttt{pgd} uses the gradient
of the classifier at the point after wrapping. The results, swept over seven values of
$\varepsilon$, are in Figure~\ref{fig:kurva-epsilon}.

\begin{lab}
\begin{verbatim}
import numpy as np

D, N, EPS = 8, 2000, 0.25
W = np.ones(D)

def data(rng, n, mu=0.5):
    y = 2 * rng.integers(0, 2, n) - 1
    return rng.normal(mu * y[:, None], 1.0, (n, D)), y

def quant(x, q):
    return x if q == 0 else np.round(x / q) * q

def mask(shape, r, rng):
    if r == 0:
        return np.ones(shape)
    return (rng.random(shape) >= r) / (1 - r)

def logit(x, q, m):
    return (quant(x, q) * m) @ W

def fgsm(X, y, q, r, rng):
    m, h = mask(X.shape, r, rng), 1e-5
    G = np.zeros_like(X)
    for j in range(D):
        e = np.zeros(D)
        e[j] = h
        G[:, j] = (logit(X + e, q, m) - logit(X - e, q, m)) / (2 * h)
    return X - EPS * y[:, None] * np.sign(G)

def pgd(X, y, q, r, rng, m_eot=1, steps=40):
    Xa, a = X.copy(), EPS / 8
    for _ in range(steps):
        G = sum(W * mask(X.shape, r, rng) for _ in range(m_eot))
        Xa = np.clip(Xa - a * y[:, None] * np.sign(G / m_eot),
                     X - EPS, X + EPS)
    return Xa

def run(seed, q, r):
    rng = np.random.default_rng(seed)
    X, y = data(rng, N)
    def pred(Z):
        return np.sign(logit(Z, q, mask(Z.shape, r, rng)))
    ok = pred(X) == y
    def succ(Z):
        return 100 * np.mean(pred(Z)[ok] != y[ok])
    return (100 * ok.mean(), succ(fgsm(X, y, q, r, rng)),
            succ(pgd(X, y, q, r, rng)), succ(pgd(X, y, q, r, rng, 20)))

print("  q     r    clean  naive   bpda1   eot")
for q, r in ((0.00, 0.0), (0.01, 0.0), (0.05, 0.0),
             (0.05, 0.5), (0.20, 0.0)):
    M = np.array([run(s, q, r) for s in (1, 2, 3, 4, 5)]).mean(0)
    print("%5.2f %5.2f  %6.2f  %5.2f  %5.2f  %5.2f"
          % (q, r, M[0], M[1], M[2], M[3]))
\end{verbatim}

The output of that program, as it stands:

\begin{verbatim}
  q     r    clean  naive   bpda1   eot
 0.00  0.00   92.02  17.40  17.40  17.40
 0.01  0.00   92.00   0.02  17.40  17.40
 0.05  0.00   91.97   0.01  17.48  17.48
 0.05  0.50   82.87  12.16  25.48  25.67
 0.20  0.00   91.77   0.00  17.65  17.65
\end{verbatim}
\end{lab}

Read that table line by line. The first line is the baseline without obfuscation:
clean accuracy $92.02$ percent, one-step attack success $17.40$ percent.
That number is small because $\varepsilon=0.25$ really is small compared with the spread of the
data, not because the attack is weak.

The second line is the heart of this chapter. With $q=0.01$, clean accuracy barely changes,
$92.00$ percent, but the naive success falls from $17.40$ to $0.02$ percent,
about $870$ times smaller. The attacker that uses BPDA gets $17.40$ percent,
exactly the baseline. That defense changes nothing that an attacker who computes correctly can
observe; it only changes what an attacker who computes incorrectly sees.

The third line repeats the same thing at $q=0.05$: naive $0.01$ percent, BPDA
$17.48$ percent, clean accuracy $91.97$ percent. Raising $q$ fivefold does not
change the BPDA success, but it makes the naive report look better. That is the
fourth symptom from the list in section two, in numerical form.

The fourth line adds random filtering with $r=0.5$, and two things happen at once.
Clean accuracy falls from $91.97$ to $82.87$ percent, because discarding half the features and
multiplying the rest by two raises the variance of the score. BPDA success rises to $25.48$
percent, above the baseline, not because the defense is worse but because the examples
that are still correct after obfuscation have a smaller margin.

The fifth line shows that coarse quantization does not save the model. With
$q=0.20$, comparable to $\varepsilon$, BPDA success is $17.65$ percent, slightly
above the baseline: a point that sits near a step boundary needs only a small push to
cross.

The variance across seeds is small but not zero: at $q=0.05$ with $r=0$ the standard deviation is
about $0.6$ percentage points, and on the line with random filtering about $1.4$ points. The standard
error of a single proportion at $2\,000$ examples and a rate of $0.17$ is about $0.8$ points, so
both numbers are consistent with ordinary sampling.

There are four things this experiment does not measure. First, the model is linear, so the BPDA
gradient is almost perfect and the cosine $0.9982$ is too good for a real model. Second, the
attacker knows $q$ and $r$. Third, only one value of $\varepsilon$ was tested.
Fourth, there is no search over attack hyperparameters. All of that makes the numbers above
a lower bound on attack success, not an upper bound.


The third shortcoming can be closed, and the result is in Figure~\ref{fig:kurva-epsilon}.
The same program is run at seven values of $\varepsilon$ from $0.05$ to
$0.35$, five seeds, without random filtering. The curve of the model without rounding and the curve
of the model with rounding $q=0.05$ almost coincide: at $\varepsilon=0.25$ the two
give $17.40$ and $17.48$ percent, and at $\varepsilon=0.35$ the two give
$27.71$ and $27.79$ percent. The curve of the naive attacker stays below $0.02$ percent
across that whole range. This defense therefore does not shift the curve and does not
break it; it merely adds a third curve that is flat.


\section{Case study: NIDS, and why evasion remains impractical}
\tingkat{advanced}

Up to here we have been dealing with feature vectors that may be changed at will. A network
intrusion detection system, or NIDS, is a different world.
\index{nids}
A NIDS reads a packet stream, extracts features such as flow duration,
packet count, byte count, TCP flags, and inter-arrival times, then decides
whether that flow is an attack. The decision must be made fast, often in milliseconds, for
thousands of flows per second.

Figure~\ref{fig:rantai} puts that chain on one line. What matters is not
the boxes, but the red mark between flow reconstruction and feature extraction.
That is where the attacker works: he does not touch the model, he touches the packets, and the
packets are what determine the feature values that enter the model. Everything to the right of that
mark sees only the consequence; everything to the left is a cause that he can control, within the
protocol limits discussed below.

\begin{figure}[!htbp]
\centering
\begin{tikzpicture}[font=\footnotesize,
  tahap/.style={draw=layergrey!70, fill=layergrey!6, rounded corners=1.5pt,
    align=center, inner sep=3pt, minimum height=1.1cm, text width=1.95cm},
  panah/.style={-latex, layergrey, line width=0.8pt}]
\node[tahap] (n1) at (0,0) {packets\\on the wire};
\node[tahap] (n2) at (2.50,0) {flow\\reconstruction};
\node[tahap] (n3) at (5.00,0) {feature\\extraction};
\node[tahap] (n4) at (7.50,0) {feature\\scaling};
\node[tahap] (n5) at (10.00,0) {NIDS\\model};
\node[tahap] (n6) at (12.50,0) {decision\\threshold};
\draw[panah] (n1) -- (n2);
\draw[panah] (n2) -- (n3);
\draw[panah] (n3) -- (n4);
\draw[panah] (n4) -- (n5);
\draw[panah] (n5) -- (n6);
\node[circle, fill=accent, inner sep=2pt] (sisip) at (3.75,0) {};
\draw[accent, line width=0.8pt] (sisip) -- ++(0,0.62);
\node[accent, anchor=south, font=\scriptsize] at (3.75,0.66)
  {insertion point: packet timing and size};
\node[layergrey, anchor=north, font=\scriptsize, align=center] at (12.5,-0.62)
  {alert or block};
\draw[-latex, layergrey, line width=0.6pt] (n6.south) -- (12.5,-0.5);
\end{tikzpicture}
\caption{The NIDS processing chain, from the wire to the decision. The red mark is one
insertion point that an attacker can use: features such as packet count, byte count, and
inter-arrival times are all determined to the left of that mark, whereas the model on
the right only reads the consequence. A defense that adds obfuscation to the right of the mark
does not touch the cause, and that is the limit discussed in this section.}
\label{fig:rantai}
\end{figure}

\kamus{NIDS}{network intrusion detection system}{a system that monitors network
traffic and decides whether a packet flow is an attack}

\kamus{problem space}{ruang masalah}{the space of real objects that an attacker may change,
for example packets and timing, as distinct from the feature space in which the model works}

The first difficulty is the distance between the two spaces. An attack is usually formulated in
feature space, whereas the attacker sends packets: every feature change must be translated
into a packet change that is still valid under the protocol and still carries out its malicious
function, and that translation does not always exist. Some features cannot even be changed
by the sender, because the packet and byte counts in the reverse direction are determined by the
receiving system.

The second difficulty is the correlation between features. Duration is tied to packet count and
inter-arrival times, and TCP flags are tied to the transport protocol, so an
attacker cannot raise the duration while lowering both inter-arrival times;
discrete features such as protocol number and packet length add further constraints.

The third difficulty is time. An attack that requires ten thousand queries means real
traffic that passes and is logged, and a trace is the first thing an analyst looks at.
A real cost is a real defense; the absence of a gradient is not.

Numbers from the literature make that distance concrete. On the CSE-CIC-IDS2018 data with fifty
filtered features, three models were trained and tested on twenty percent of the data, then
attacked with FGSM, PGD, and BIM. Before the attack all three were almost perfect: accuracy
$0.997$ for ANN, $0.998$ for SVM, and $0.997$ for CNN. On the day of the attack,
with full knowledge, FGSM dropped ANN to $0.756$ with F1 $0.677$, CNN to
$0.842$ with F1 $0.863$, and SVM to $0.150$ with F1 $0.000$. PGD was slightly
stronger on CNN, $0.584$ with F1 $0.689$, and BIM gave $0.688$ with F1
$0.751$.

Those numbers look like a victory for the attacker, and on a static model they are;
but the same experiment continues with a model retrained every day.
With knowledge of a model that is one day old, SVM returns to $0.998$, ANN to
$0.938$, and CNN to $0.957$ with F1 $0.959$, and on the second day the numbers are between
$0.966$ and $0.998$. That retraining does not use adversarial examples, so
the recovery is not the result of an adversarial defense but a consequence of the fact that the
gradient computed on yesterday's model no longer points in the right direction.

This is the NIDS version of the whole chapter: deliberate gradient obfuscation is fragile, because
an attacker can break it with BPDA, whereas obfuscation that comes from model change lasts longer,
since the attacker must update his knowledge, and updating knowledge means gaining access, and
access leaves a trace. The difference is not in the technique, but in who pays the cost.

\index{transferability}
The authors of that review built an attack tree and a taxonomy of practicality problems, then
concluded that only three paths are viable. The first path is blind perturbation:
adding random bytes to UDP packets, or to TCP packets with the PSH flag, while
recomputing the checksum so that the protocol stays valid. This path needs no knowledge of the
model and no queries, so it is immune to retraining, but it has no
direction. The second path is querying through a side channel, which demands continuous network
access and depends on indicators that change with network load and
hardware. The third path is a substitute model trained on data captured from one's own
botnet, which works through transferability but demands training data that is still
fresh.

All three paths are real, and the authors stress that the distance between research and
practice is wide and that the threat model assumptions in many papers are too generous
to the attacker.

\index{evasion}
The authors' own work, which we will call a working note and not a published
result, moves on the other side of the same problem. On one network flow dataset,
a model reaches F1 $0.967037$ from $46\,558$ out of $48\,145$ test examples.
The question is not ``does the attack succeed'', but ``how many features must be
changed for an example to cross the decision boundary while remaining valid under
the protocol''. With a mixed-integer linear program, the answer depends sharply
on the number of features that may be changed: at seven features, not one of fifty
examples could be moved across; at twenty features, thirty-three of fifty could.
The ratio of distance to the boundary between successful and control examples is about $3.28$
for the linear model and $2.13$ for the neural network.

\index{rule of three}
Two more numbers show the caution involved. Of twenty thousand candidates tested for
protocol consistency, none passed; with the rule of three, the upper bound at $95$ percent
for that rate is about $0.019$ percent. On the model with an enlarged safety margin,
the condition takes the form $\varepsilon+40\delta$ with
$\delta=1.245\times10^{-7}$, so that the effective budget is $0.150004981$ and not
$0.15$.

The limits of that working note must be stated together with its results. It uses one dataset,
fifty test examples, and two model families. It verifies the absence of an attack that
meets the conditions, not the discovery of an attack. The number $0$ out of $20\,000$ is a statement
about the candidate generator that was used, not about all possible packets, and
the result says nothing about dynamic NIDSs.

Now the conclusion that must be stated clearly. The toy model in section five does not
prove anything about real NIDSs; it shows that rounding that erases the
gradient produces a misleading report, and that the mechanism can be broken
at a reasonable cost. It does not measure the cost of an attack in packet space. What can be carried
away from this chapter is a habit: every time a defense reports zero success,
ask about its gradient, then ask who pays the cost.

\begin{tantangan}
A paper reports a defense that lowers attack success from $96$
percent to $3$ percent, with clean accuracy that does not change at all. Their attacker
uses one-step FGSM, $\varepsilon=8/255$, the exact gradient, one initial point,
and one seed. The question is: is that robustness real, and what is the cheapest experiment
that most quickly shows that the $3$ percent figure is an artifact?

First hint: look at what is not being paid for. A defense that genuinely
raises the cost almost always pays something, either accuracy, or latency, or
a number of parameters.

Second hint: if the gradient is computed through a non-differentiable piece, the
$3$ percent figure is not attack success, but the residue of examples that happen to lie
near the decision boundary. Repeat with forty-step PGD, BPDA, and three initial points;
then sweep $\varepsilon$ and report the curve.

If you get stuck, go back to section three.
\end{tantangan}

\begin{untukpeneliti}
Three things in this section are often misquoted. First, the categories shattered, stochastic, and
vanishing are the taxonomy of the obfuscated gradients paper; the list of seven techniques in
section two is our pedagogical decomposition. In that paper there are three categories, nine
defenses tested, seven of which depend on an obfuscated gradient, and
six broken completely plus one broken partially.

Second, BPDA is not a general solution. The error bound $L\delta$ holds only when $f$ is smooth
and $\delta$ is small relative to the scale of the curvature; for rounding with $q$ comparable to
$\varepsilon$ that guarantee is empty. Conversely, EOT is not a cure for an attack that points in
the wrong direction: averaging wrong directions still produces a wrong direction.

Third, the NIDS results we quote come from feature-space experiments with full
knowledge and daily retraining, not from packet-space attacks that were actually run. They
support a narrow claim: retraining restores accuracy against attacks that use
the gradient of an old model. They do not support the claim that NIDSs are immune to an attacker who
may change packets.
\end{untukpeneliti}

\section{A checklist before writing that an attack succeeded}
\tingkat{basic}

This section is a summary that can be used directly. Every item is a question
that must be answerable before a claim about an attack or a defense enters the
abstract.

\begin{enumerate}
  \item \textbf{The threat model is written in full.} What does the attacker know: architecture,
        weights, training data, probabilities, or only labels? How many queries may he
        make, and for how long?
  \item \textbf{The attack is bounded by a reasonable budget.} State $\varepsilon$
        and its norm, or the byte and packet limits if the work is in problem space. An attack
        without a bound measures something that nobody faces.
  \item \textbf{The attack is adaptive.}
        \index{adaptive attack}
        The attacker must know the defense being tested, including the possibility that
        its gradient is broken. A defense tested only by attacks from an earlier
        paper has not been tested.
  \item \textbf{The gradient is inspected before it is trusted.} Report the gradient norm, the fraction
        of entries that are zero, and their correlation with the correct direction. A zero gradient on a
        non-constant model is a finding, not a result.
  \item \textbf{Success is reported together with clean accuracy.} A defense that
        lowers clean accuracy from $92$ to $60$ percent must not be compared
        directly with one that keeps $92$ percent.
  \item \textbf{Several seeds, with mean and deviation.} One seed is one
        throw of the die; a difference between two lines that is smaller than their
        deviation is not a difference.
  \item \textbf{A curve, not a single point.} Sweep the value of $\varepsilon$, the number of steps,
        the number of EOT samples, and the number of queries. A good defense shifts the curve;
        an illusory one shifts a single point.
  \item \textbf{The attacker's compute budget is stated:} how many times the model is called,
        how many hours, how many GPUs? An attack that succeeds after a million queries is not
        the same attack as one that succeeds after a hundred.
  \item \textbf{Code and weights are available.} A claim that cannot be repeated cannot be
        checked. If release is impossible, the procedure and hyperparameters must be detailed
        enough to reconstruct.
  \item \textbf{Negative results are reported as negative results.} A defense that fails
        is information; a report that succeeds because its attack was stopped before
        time is an expensive mistake.
\end{enumerate}

\kamus{rule of three}{aturan tiga}{if no event occurs in $n$ independent trials,
the upper bound at $95$ percent for the probability of that event is about $3/n$}

\begin{ringkasan}
\begin{itemize}
  \item An obfuscated gradient blinds the attacker, it does not make the model strong.
        The seven techniques produce the same symptom: the measurement looks better
        without anything getting better.
  \item For a step function such as rounding, the gradient of the composition is zero at almost every
        point, so a one-step attack produces $x'=x$ and ``accuracy under
        attack'' is precisely the clean accuracy.
  \item BPDA replaces the gradient of a non-differentiable piece with the gradient of the
        classifier at the point after the piece. Its error is bounded by $L\delta$: it makes
        sense when the shift is small relative to the scale of the curvature, and it gives no guarantee
        when it is not.
  \item EOT maximizes the expected loss over the randomness of the defense; the average of $m$
        samples is unbiased with variance proportional to $1/m$. Randomness raises the per-step
        cost, it does not change the direction.
  \item On the eight-feature toy model with $\varepsilon=0.25$, rounding
        $q=0.01$ drops naive success from $17.40$ to $0.02$ percent,
        whereas BPDA stays at $17.40$ percent; $q=0.20$ does not hold BPDA back either.
  \item On NIDS, FGSM drops SVM accuracy from $0.998$ to $0.150$, but
        daily retraining restores it to $0.998$. The durable defenses there
        come from the cost of access and protocol constraints, not from gradient obfuscation.
\end{itemize}
\end{ringkasan}

\section{Exercises}
\tingkat{intermediate}

\begin{latihan}
\begin{enumerate}
  \item A paper writes: ``Our defense lowers attack success from
        $98$ to $0$ percent at $\varepsilon=8/255$.'' State three checks
        before accepting that sentence, and what to expect from each check if
        the defense only obfuscates the gradient.
  \item Derive the BPDA error bound for $f(x)=w\cdot x+b$, show that the error is zero, then
        explain why that result does not carry over to a two-layer model with ReLU activations.
  \item On the toy model of section four, determine the smallest value of $q$ that makes one
        BPDA step fail to flip the label of $x_0$, express it as a function of the initial margin and
        $\varepsilon$, then check it with the program.
  \item You have $1\,000$ queries. The defense randomizes a feature mask with $r=0.5$ and
        refuses to answer if more than half the features are dropped. Compute how many PGD steps
        you can run with $m=20$ and with $m=1$, then explain what you gain and what you pay.
  \item Repeat the experiment of section five with a two-layer model with ReLU activations. Report
        clean accuracy, naive success, and BPDA success for three values of $q$,
        then explain why the cosine is no longer close to one.
  \item A NIDS reads fifty features, seven of which are binary flags and three
        protocol numbers. Compute how many features an attacker can actually change, state
        your assumptions, and compare with a $32\times32$ image that has $3\,072$ values.
  \item Explain why daily retraining restores accuracy against attacks that use
        gradients but not against blind perturbation. Connect this with the EOT lemma and
        the definition of an adaptive attack.
  \item Design an experiment that separates two hypotheses: (a) the defense raises the cost of
        attack, (b) the defense only corrupts the measurement. State the metrics, the number of seeds,
        the shape of the curve, and what results would make you accept (a).
\end{enumerate}
\end{latihan}

\peta{Chapter 30 of Athalye et al., \emph{Obfuscated Gradients Give a False Sense of
Security} (2018): section 3 for the three families of obfuscated gradients, section 4 for
BPDA and adaptive attacks, section 6 for the checklist. In Madry et al., \emph{Towards
Deep Learning Models Resistant to Adversarial Attacks} (2018): section 3 for PGD and
random initial points. In elShehaby and Matrawy, \emph{Evasion Adversarial Attacks Remain
Impractical Against ML-based Network Intrusion Detection Systems}, for the attack tree,
the practicality taxonomy, and the retraining experiment. In Chapter 20 of this book for the rule of
three, and in Chapter 27 for FGSM and PGD.}

\chapter[Cryptography: Modular to Quantum]{Cryptography: Modular Integers to the Quantum Threat}

A student submits an assignment containing a cipher of his own making. Each letter is replaced by
its number, that number is multiplied by three, and the remainder after division by twenty-six is
used as the new letter. To open the message, he multiplies again by nine. When asked where the
number nine came from, his answer is nearly right: ``because there is no division here, I multiply
by the inverse.'' The examiner then replaces the multiplier three with the multiplier four. The
assignment is never finished. With the multiplier four, the first letter and the fourteenth letter
produce the same symbol, so two different messages produce an identical cipher.

The question that arises from that failure is the question of this whole chapter: when may a number
be divided inside remainder arithmetic? The answer is short: exactly when that number shares no
factor with its modulus. That short answer, pressed hard, yields the three things every browser
uses: key exchange, public keys, and digital signatures.

The previous three chapters are used explicitly. Chapter 2 gives the language for bijections and
partitions; here a bijection proves that multiplication by a unit only shuffles the order, and
partitions count the order. Chapter 3 gives the completeness property of the real numbers, and that
property does not cross over, because $\mathbb{Z}_n$ has no order compatible with multiplication;
what replaces it is finiteness, since on a finite set an injective map is automatically surjective.
Chapter 28 explains why ``secure'' cannot be a theorem without conditions.

One warning applies to every page here. All the moduli are small, chosen so that they can be checked
by hand and by Python in a matter of seconds. The largest modulus, $1\,050\,589$, is factored by a
six-line script. Not one cipher here is safe to use.

\begin{untukanak}
Imagine the wall clock in a classroom. Its hand moves from the number one to twelve, then back to
one. That clock never shows the number thirteen or twenty. If we turn the hand five steps, then
five steps again, then five steps again, we do not need to know how much five plus five plus five
is; we only look at where the hand stops.

Now imagine a combination lock with thirteen positions. We may turn it as far as we like, and only
the final position is recorded. If we turn three steps, how many times must we turn three steps
again to get back to where we started? The answer is nine, and that is not a coincidence, but
because three and thirteen share no factor. Try the same thing with four steps on a twelve-position
lock: however many times we turn, we will never land on all positions, because four and twelve
share a factor.

In this chapter we will understand why that small difference decides whether a message can be opened
again or is lost forever, and why two people who have never met can agree on one secret number in
front of a third person who hears everything.
\end{untukanak}

\section{Residue classes: arithmetic that goes around}
\tingkat{basic}

We begin by fixing the notion of ``remainder'' already used in the first paragraph.
\index{modular arithmetic}
Two integers are called congruent modulo $n$ if their difference is divisible by $n$.

\begin{definisi}
Let $n\ge1$ be an integer. For $a,b\in\mathbb{Z}$ we write $a\equiv b\pmod{n}$ if $n$ divides
$a-b$. The set of equivalence classes under this relation is called $\mathbb{Z}_n$, and the class
containing $a$ is written $[a]_n$.
\end{definisi}

We owe two things: that this really is an equivalence relation, and that the remainder is unique.

\begin{lemma}
Congruence modulo $n$ is an equivalence relation on $\mathbb{Z}$, and every class $[a]_n$ contains
exactly one representative in $\{0,1,\dots,n-1\}$, so that $|\mathbb{Z}_n|=n$.
\end{lemma}

\begin{proof}
Reflexive: $n$ divides $0$. Symmetric: if $n\mid(a-b)$ then $n\mid(b-a)$. Transitive: if
$n\mid(a-b)$ and $n\mid(b-c)$ then $n\mid(a-c)$ because $(a-c)=(a-b)+(b-c)$. For uniqueness, take
$a$ and consider $S=\{a-qn\ge0: q\in\mathbb{Z}\}$. That set is nonempty by the Archimedean property
(Chapter 3), so it has a least element $r$. If $r\ge n$ then $r-n\in S$ is smaller than $r$, a
contradiction; so $0\le r<n$. If $r'$ also satisfies $0\le r'<n$, then $r-r'$ is a multiple of $n$
whose value lies between $-(n-1)$ and $n-1$, hence must be zero.
\end{proof}

Operations on $\mathbb{Z}_n$ are defined through their representatives, and a definition of this
kind is dangerous: if the result depends on the representative, the operation is not a function. The
next lemma closes that hole with one line of algebra.

\begin{lemma}
If $a\equiv a'$ and $b\equiv b'$ modulo $n$, then $a+b\equiv a'+b'$ and $ab\equiv a'b'$ modulo $n$.
\end{lemma}

\begin{proof}
Write $a'=a+kn$ and $b'=b+ln$. Then $a'+b'=(a+b)+(k+l)n$ and $a'b'=ab+(al+bk+kln)n$.
\end{proof}

Addition and multiplication are therefore well defined, and both inherit associativity and
distributivity from $\mathbb{Z}$. What is lost is order: there is no way to put a sign ``$<$'' on
$\mathbb{Z}_n$ such that $a<b$ implies $a+c<b+c$ and $ac<bc$ for positive $c$. The supremum tool of
Chapter 3 stops at the entrance, and what works inside is counting.

Division need not exist. The class $[3]_{26}$ has an inverse, namely $[9]_{26}$, since
$3\cdot9=27\equiv1$. The class $[2]_{26}$ has none: every multiple of two is even, whereas $27$ is
odd, so $2b\equiv1$ has no solution. Consequently multiplication by two modulo twenty-six is not a
bijection, because $[0]$ and $[13]$ both go to $[0]$.

\kamus{residue class}{kelas sisa}{all the numbers that stop at the same clock position, like three
in the morning and three in the afternoon}

\begin{definisi}
A class $[a]_n$ is called a unit if there is $[b]_n$ with $[a]_n[b]_n=[1]_n$. The set of all units
modulo $n$ is written $U(n)$, and its number of members is written $\varphi(n)$.
\end{definisi}

For $n=26$ there are twelve units, namely the classes of $1,3,5,7,9,11,15,17,19,21,23,25$. The
program in the lab counts them again; that number reappears in part three under a different name.

\section{Common factors, Euclid, and when division exists}
\tingkat{basic}

\index{greatest common divisor}
The question ``when does $[a]_n$ have an inverse'' turns out to be a question about the greatest
common factor, which has an algorithm more than two thousand years old.

\begin{definisi}
For $a,b\in\mathbb{Z}$ not both zero, $\gcd(a,b)$ is the largest natural number that divides $a$ and
$b$.
\end{definisi}

The key to Euclid's algorithm is one equality of sets: the common divisors of $a$ and $b$ are exactly
the common divisors of $b$ and $a-qb$.

\begin{lemma}
For every $q\in\mathbb{Z}$, $\gcd(a,b)=\gcd(b,a-qb)$.
\end{lemma}

\begin{proof}
If $d$ divides $a$ and $b$, then $d$ divides $a-qb$. Conversely, if $d$ divides $b$ and $a-qb$, then
$d$ divides $(a-qb)+qb=a$. So the two pairs have the same set of common divisors, hence the same
greatest one.
\end{proof}

Replacing $(a,b)$ by $(b,a\bmod b)$ therefore does not change the answer but shrinks $b$, and the
process stops because natural numbers do not shrink without bound.

\begin{teorema}[Bezout]
For $a,b$ not both zero, there are $x,y\in\mathbb{Z}$ with $ax+by=\gcd(a,b)$.
\end{teorema}

\begin{proof}
Consider $S=\{ax+by>0: x,y\in\mathbb{Z}\}$. This set is nonempty because it contains $|a|$ or $|b|$,
and every nonempty subset of the natural numbers has a least element; this is the discrete form of
the infimum property in Chapter 3. Let $d=ax_0+by_0$ be that least element. Divide $a$ by $d$:
$a=qd+r$ with $0\le r<d$. Then $r=a(1-qx_0)+b(-qy_0)$, so $r\in S$ if $r>0$. Since $d$ is least,
$r>0$ is impossible; so $r=0$ and $d\mid a$, and in the same way $d\mid b$. Hence $d\le\gcd(a,b)$.
Conversely $\gcd(a,b)$ divides every expression $ax+by$, in particular $d$, so that
$\gcd(a,b)\le d$.
\end{proof}

\begin{teorema}
The class $[a]_n$ is a unit if and only if $\gcd(a,n)=1$.
\end{teorema}

\begin{proof}
If $[a]_n$ is a unit, there is $b$ with $ab\equiv1$, that is $ab-1=kn$. Every common divisor of $a$
and $n$ divides $ab$ and $kn$, hence divides $1$; so $\gcd(a,n)=1$. Conversely, if $\gcd(a,n)=1$,
Bezout's theorem gives $ax+ny=1$. Read modulo $n$, that equality says $ax\equiv1$, so $[x]_n$ is the
inverse of $[a]_n$.
\end{proof}

That proof also gives a recipe: run Euclid, then build the combination backward. For $a=1071$ and
$b=462$ the steps are three divisions,
\[
1071=2\cdot462+147,\qquad 462=3\cdot147+21,\qquad 147=7\cdot21+0,
\]
so that $\gcd=21$. Building backward gives $21=1071\cdot(-3)+462\cdot7$, and since the common factor
is $21$ rather than $1$, the class $[462]_{1071}$ is not a unit.

How long does this algorithm run? Every two steps the remainder at least halves, so the number of
divisions is proportional to the number of digits. The sharp bound was given by Lam\'e in 1844: no
more than $\log_\phi(b)+1$ divisions with $\phi=(1+\sqrt5)/2$, and the pairs that attain it are
consecutive Fibonacci pairs. The pair $(F_{47},F_{46})=(2\,971\,215\,073,\ 1\,836\,311\,903)$ needs
$45$ divisions, whereas $\log_\phi(F_{46})=44.33$. For two thousand random thirty-two-bit pairs, the
average is $18.75$ with maximum $32$.

\kamus{modular inverse}{invers modular}{the number that, multiplied by the original number, gives
remainder one, like the ``undo'' button on a cash register}

\section{Fermat's little theorem and Euler}
\tingkat{intermediate}

This chapter needs one fact about powers, and that fact is born from counting: multiply every unit by
one fixed unit; the result is the same list in a different order.

\begin{lemma}
If $\gcd(a,n)=1$, the map $[x]_n\mapsto[ax]_n$ is a bijection from $U(n)$ to $U(n)$.
\end{lemma}

\begin{proof}
The result lies in $U(n)$ because the product of two units is a unit: if $xx'\equiv1$ and
$yy'\equiv1$ then $(xy)(x'y')\equiv1$. The map is injective: if $ax\equiv ay$, multiply both sides by
the inverse of $a$, which exists by the unit theorem in part two. The set is finite, and on a finite
set an injective map from the set to itself is automatically surjective (Chapter 2). So it is a
bijection.
\end{proof}

\begin{teorema}[Euler]
If $\gcd(a,n)=1$, then $a^{\varphi(n)}\equiv1\pmod n$.
\end{teorema}

\begin{proof}
Let $u_1,\dots,u_k$ be all the units modulo $n$, with $k=\varphi(n)$, and $P=u_1\cdots u_k$. By the
lemma above, the list $au_1,\dots,au_k$ contains every unit exactly once, so
\[
a^kP\equiv(au_1)(au_2)\cdots(au_k)\equiv u_1u_2\cdots u_k=P\pmod n .
\]
Since $P$ is itself a unit, it has an inverse, and both sides may be multiplied by it. The result is
$a^k\equiv1$.
\end{proof}

\begin{akibat}[Fermat's little]
If $p$ is prime and $p\nmid a$, then $a^{p-1}\equiv1\pmod p$. For all $a$ we have
$a^p\equiv a\pmod p$.
\end{akibat}

The second form follows from the first: if $p\mid a$ both sides are zero, if not, multiply
$a^{p-1}\equiv1$ by $a$. The name ``little'' distinguishes it from Fermat's last theorem, which is
far harder.

The number twelve from part one now has two names: $\varphi(26)=12$, and $3^{12}\equiv1\pmod{26}$;
the program checks that $5^{12}$ and $7^{12}$ are also $\equiv1$. Large powers can therefore be
trimmed: for $7^{222}\bmod13$ write $222=12\cdot18+6$, so that $7^{222}\equiv7^6\equiv9\cdot10\equiv12$
because $7^2\equiv10$ and $7^4\equiv9$.

The converse does not hold, and that matters for the honesty of this chapter. The number
$561=3\cdot11\cdot17$ is not prime, yet $a^{560}\equiv1\pmod{561}$ for every $a$ that shares no
factor with it; the program checks all $a$ without finding a counterexample. Numbers of this kind are
called Carmichael, and for that reason practical primality tests use Miller--Rabin.

\kamus{Euler totient function}{fungsi phi}{the number of positions on a combination lock that can
reach every other position, that is the positions that share no factor with the number of positions}

\section{Cyclic groups, order, and primitive roots}
\tingkat{intermediate}

If we multiply one unit repeatedly, the result goes around and finally returns to one. The useful
question is how long that circuit is.

\begin{definisi}
For $\gcd(a,n)=1$, the order of $a$ modulo $n$ is the least natural number $k$ with
$a^k\equiv1\pmod n$, written $\operatorname{ord}_n(a)$.
\end{definisi}

An order always exists: the sequence $a^0,a^1,\dots$ has only $\varphi(n)$ distinct values, so two of
them are equal and $a^{i-j}\equiv1$ for some $i>j$. That bound can be sharpened: the order divides
$\varphi(n)$.

\begin{lemma}
If $\gcd(a,n)=1$, then $\operatorname{ord}_n(a)$ divides $\varphi(n)$.
\end{lemma}

\begin{proof}
Let $k=\operatorname{ord}_n(a)$ and $H=\{a^0,\dots,a^{k-1}\}$, whose members are all distinct, so
$|H|=k$. For every unit $x$, the set $xH=\{xa^0,\dots,xa^{k-1}\}$ also has $k$ members, because
multiplication by a unit is a bijection. The sets $xH$ form a partition of $U(n)$: two of them are
equal or disjoint, since $xH$ and $yH$ intersect if and only if $y^{-1}x\in H$ (Chapter 2). So
$\varphi(n)$ is $k$ times the number of parts.
\end{proof}

\begin{definisi}
If $\operatorname{ord}_n(g)=\varphi(n)$, then $g$ is called a primitive root modulo $n$. For a unit
$h$, the number $k$ with $g^k\equiv h$ is called the discrete logarithm of $h$ to the base $g$.
\end{definisi}

Modulo $13$, $2,6,7,11$ have full order $12$ and are therefore primitive roots, whereas $3$ and $9$
have order $3$, $5$ and $8$ have order $4$, $4$ and $10$ have order $6$, and $12$ has order $2$. All
those numbers divide $12$.

For $2$, the circuit visits all twelve nonzero classes before returning: $2,4,8,3,6,12$, then
$11,9,5,10,7,1$. Figure~\ref{fig:jam} shows it as a clock with thirteen positions. For $3$, the
circuit is only $3,9,1$ and then repeats.

\begin{figure}[htbp]
\centering
\begin{tikzpicture}
\foreach \k in {0,...,12}
  \node[circle,draw,inner sep=1pt,minimum size=6mm,font=\scriptsize] (n\k)
    at ({90-360*\k/13}:2.1cm) {\k};
\node[font=\footnotesize,align=center] at (0,0) {multiply\\ by 2};
\foreach \a/\b in {1/2,2/4,4/8,8/3,3/6,6/12,12/11,11/9,9/5,5/10,10/7,7/1}
  \draw[-{Latex[length=1.5mm]},semithick] (n\a) to[bend left=18] (n\b);
\end{tikzpicture}
\caption{Multiplication by $2$ modulo $13$ goes around through all twelve nonzero classes before
returning to $1$. The length of that circuit is the order, and full order means $2$ is a primitive
root modulo $13$.}
\label{fig:jam}
\end{figure}

The existence of a primitive root for every prime $p$ is Gauss's theorem, and we skip its proof. What
must be carried away: for a prime $p$, the units are a single circuit of $p-1$ positions, and every
unit is a power of one primitive root. If $h=g^k$, knowing $k$ means knowing $h$, but not
conversely; that asymmetry is the raw material of public-key cryptography.

\kamus{primitive root}{akar primitif}{one starting position that, if turned continuously, passes
through every other position before returning}

\section{Diffie-Hellman key exchange}
\tingkat{intermediate}

\index{diffie-hellman}
Two people want to agree on one secret number, and the only channel they have is overheard by a third
party. In 1976 Diffie and Hellman showed that this is possible, and the mechanism uses only two
formulas.

They first agree on a prime $p$ and a primitive root $g$ modulo $p$; both may be announced. Alice
chooses $a$ secretly and sends $A=g^a\bmod p$. Bob chooses $b$ secretly and sends $B=g^b\bmod p$.
Alice computes $K=B^a\bmod p$, Bob computes $K=A^b\bmod p$. Both obtain the same number, because
\[
B^a\equiv(g^b)^a\equiv g^{ab}\equiv(g^a)^b\equiv A^b\pmod p .
\]
What the eavesdropper sees is only $p,g,A,B$. To obtain $K$ he needs $g^{ab}$, and the easiest known
way is to find $a$ or $b$, that is to solve the discrete logarithm.

\begin{bacakeras}
Read it as a dialogue for two voices, alternating, one line per breath.

First voice: ``Agree first, announce two.''
Second voice: ``Keep your own, never say it.''
First voice: ``Raise, take the remainder, send the result.''
Second voice: ``Receive the delivery, raise it again.''
First voice: ``The last remainder, that is the key.''
Together: ``Two directions, one number, never crossed.''

Repeat three times, faster and faster. The line that makes the tongue stumble is the step not yet
understood, not the line that needs to be memorized harder. Once it flows, close the book and name
four things in order: what is announced, what is kept, what is sent, what is computed. Naming the
secret number first cancels the whole game.

The line repeated three times aloud: ``what is sent is the power, what is kept is the exponent.''
\end{bacakeras}

The concrete numbers, all checked by the program in the lab part: $p=10\,007$, $g=5$ has order
$10\,006$ and is therefore a primitive root, $a=1234$, $b=5678$. Then $A=5^{1234}\bmod p=4873$ and
$B=5^{5678}\bmod p=3189$. Alice computes $3189^{1234}\bmod p=5674$; Bob computes
$4873^{5678}\bmod p=5674$. Their shared key is $5674$, and that number never crosses the channel.
Figure~\ref{fig:dh} shows the flow together with what the eavesdropper hears.

\begin{figure}[htbp]
\centering
\begin{tikzpicture}[font=\footnotesize,
  kotak/.style={draw,rounded corners=2pt,align=left,inner sep=4pt,text width=4.4cm}]
\node[kotak] (a) {\textbf{Alice}\\ $a=1234$ secret\\ $A=5^{1234}\bmod p=4873$};
\node[kotak,right=3.4cm of a] (b) {\textbf{Bob}\\ $b=5678$ secret\\
  $B=5^{5678}\bmod p=3189$};
\node[kotak,below=1.3cm of $(a)!0.5!(b)$, text width=9.2cm, fill=black!4]
  (e) {\textbf{Eavesdropper} hears $p=10\,007$, $g=5$, $A=4873$, $B=3189$, then intends
  to compute $K=A^{b}$ or $K=B^{a}$; for that it must find $a$ or $b$.};
\draw[-{Latex}] ([yshift=3mm]a.east) -- node[above,font=\scriptsize]{$A$}
  ([yshift=3mm]b.west);
\draw[-{Latex}] ([yshift=-3mm]b.west) -- node[below,font=\scriptsize]{$B$}
  ([yshift=-3mm]a.east);
\draw[-{Latex},dashed] (e.north -| a.south) -- (a.south);
\draw[-{Latex},dashed] (e.north -| b.south) -- (b.south);
\end{tikzpicture}
\caption{The flow of Diffie-Hellman key exchange on the numerical example in the text. The values $A$
and $B$ cross the open channel; what does not cross are the exponents $a$ and $b$, and it is exactly
those two that the eavesdropper needs.}
\label{fig:dh}
\end{figure}

That equality only states that the protocol is consistent, not that it is secure. All that can be
proved is a conditional statement: if the discrete logarithm modulo $p$ is hard, then computing
$g^{ab}$ from $A$ and $B$ is also hard. That condition is not a theorem; there are groups in which
the discrete logarithm is easy, so security is a property of particular groups.

\begin{tangan}
Materials: yellow watercolor, one secret color per person (red and blue), four plastic cups, one
spoon, two sheets of white paper.

First step, in front of the table: pour one spoonful of yellow paint into each cup. The yellow plays
the role of the prime and the base, and both are announced. Second step, behind your back: each
person adds two spoonfuls of their secret color, stirs, then hands over the mixed cup. Third step:
each adds two spoonfuls of their own secret color to the cup just received, stirs, then puts one drop
on the paper.

What must be visible: the two drops are the same color, even though neither of the two people ever
handed over their secret color. The first mixed cup held by the third person cannot be separated
again; he sees only two different dull colors.
\end{tangan}

Two parameter traps deserve mention. First, $g$ must generate a large subgroup of prime order: the
modulus $10\,007$ is a safe prime with $p-1=2\cdot5003$, and the program records that
$5^{5003}\bmod p=10\,006$ has order only $2$. If Alice uses that base, their shared key is only $1$
or $10\,006$, and the eavesdropper need not compute any logarithm at all. Second, this protocol
names no one's identity, so an eavesdropper can stand in the middle and read everything.
Figure~\ref{fig:dh} depicts only a passive eavesdropper.

About size: the modulus $10\,007$ can be broken in less than one second by the index calculus
algorithm, and we use it because the powers fit on one page. Real practice uses $2048$ bits or more,
and even that is a time limit, not a guarantee.

\begin{untukanak}
Imagine two people who want to agree on one secret color, even though all their conversation is
overheard. They first agree on a base color, say yellow. Alice chooses her secret color, mixes it with
yellow, then sends the mixture. Bob does the same with his own secret color. Now each mixes the color
received with their own secret color. Both get the same color, and the eavesdropper sees only two
mixtures that cannot be separated again.

That separation is what is considered hard: mixing colors is easy, separating a mixture without one
of its ingredients is considered not. Cryptography replaces colors with numbers and ``considered not
easy'' with a mathematical sentence.

In this chapter we will find out how numbers replace colors, and why the security depends on one
``considered'' sentence that nobody has ever proved.
\end{untukanak}

\section{RSA: from key to message and back}
\tingkat{intermediate}

\index{rsa}
Diffie-Hellman agrees on a key, but does not send a message. A year later Rivest, Shamir, and
Adleman showed that the same property can be used for two different keys: one to lock, one to open.

Choose two large primes $p$ and $q$, compute $n=pq$ and $\varphi(n)=(p-1)(q-1)$. Choose $e$ with
$\gcd(e,\varphi(n))=1$; the unit theorem gives its inverse $d$ modulo $\varphi(n)$. The public key is
$(n,e)$, the secret key $d$. For a message $0\le m<n$, the cipher is $c=m^e\bmod n$ and its opening is
$m=c^d\bmod n$.

Why does opening return $m$? Because $ed=1+k\varphi(n)$, so that
$c^d=m^{ed}=m\,(m^{\varphi(n)})^k$. If $\gcd(m,n)=1$, Euler's theorem makes that bracket one. The
remaining case explains why $n$ must be the product of two distinct primes.

\begin{teorema}
Let $n=pq$ with $p,q$ distinct primes, $ed\equiv1\pmod{(p-1)(q-1)}$, and $0\le m<n$. Then
$m^{ed}\equiv m\pmod n$.
\end{teorema}

\begin{proof}
If $p\nmid m$, Fermat's little theorem gives $m^{p-1}\equiv1\pmod p$, so that
$m^{ed}=m\,(m^{p-1})^{k(q-1)}\equiv m\pmod p$. If $p\mid m$, both sides are congruent to zero modulo
$p$. So the congruence holds modulo $p$ for every $m$, and for the same reason modulo $q$. Since $p$
and $q$ both divide $m^{ed}-m$ and $\gcd(p,q)=1$, their product also divides that number: if
$m^{ed}-m=pt$ and $q\mid pt$ with $q$ prime and $q\nmid p$, then $q\mid t$, so that $pq\mid m^{ed}-m$.
\end{proof}

The numerical example is small and entirely checked by the program. Take $p=1019$, $q=1031$, then
$n=1\,050\,589$ and $\varphi(n)=1018\cdot1030=1\,048\,540$. Choose $e=65\,537$; Euclid gives
$d=60\,493$, and indeed $65\,537\cdot60\,493=3\,964\,529\,741=1+3781\cdot1\,048\,540$. The message
$m=424\,242$ produces the cipher $c=552\,241$, and $c^d\bmod n$ returns $424\,242$. That modulus is
only twenty-one bits and is factored by a six-line script.

Plain RSA, the formula above with no addition of any kind, is not secure; its three reasons can be
shown with the same modulus. First, it is deterministic: the same message always produces the same
cipher, so that repeated messages can be recognized. Second, it is multiplicative:
$E(m_1)E(m_2)\equiv E(m_1m_2)$, and the program shows
$E(1234)\cdot E(567)\bmod n=551\,173$ opened to $699\,678$, exactly $1234\cdot567$. Third, small
messages are not protected: with $e=3$ and $m=42$, the cipher is $74\,088$, coincidentally equal to
$42^3$ because the result has not yet passed $n$, so that an ordinary cube root recovers the message.

All three are closed by padding: the message receives a random addition before exponentiation, so
that the result is full length, always different, and no longer multiplicative. RSA-OAEP, the
standard used today, has a security proof against chosen-ciphertext attacks in the random oracle
model under the assumption that RSA holds. In practice RSA wraps a $256$-bit symmetric key and then
the data is encrypted with AES; this pattern is called hybrid encryption, because RSA is slow: its
time is proportional to the cube of the modulus length.

\kamus{padding}{bantalan}{extra contents put into a package before shipping, like bubble wrap that
makes every package a different shape and size}

\section{Hash functions, MACs, and signatures}
\tingkat{intermediate}

\index{hash function}
Encryption hides the contents. Sometimes what is needed is not hiding but summarizing: one short
fingerprint that changes when the contents change.

\begin{definisi}
A cryptographic hash function maps a message of any length to an output of fixed size. It must be
preimage resistant (from $h$, it is hard to find $m$ with $H(m)=h$), second-preimage resistant (from
$m$, it is hard to find $m'\neq m$ with $H(m')=H(m)$), and collision resistant (it is hard to find a
distinct pair $m,m'$ with the same hash).
\end{definisi}

Those three properties are not equivalent; the weakest is the one most often misunderstood. Collision
resistance is limited by birthdays: if the output has $N$ values, a collision appears after about
$1.177\sqrt N$ trials, not after $N$; for $128$ bits that number is $2.17\cdot10^{19}$ and for $64$
bits $5.06\cdot10^{9}$. The program in the lab tests that prediction on SHA-256 truncated to $16$
bits: an average of $362.7$ samples against a prediction of $320.8$, with a spread from $100$ to
$743$. So a hash with an $n$-bit output gives only $n/2$ bits of collision resistance.

A MAC adds a key to the hash: two parties who share a key can check that a message came from the
other party and did not change, because without the key there is no way to compute the correct tag.
The limit: a MAC is useful only between two parties who both hold the key, so that a third party
cannot be a witness.

A signature uses a key pair in the opposite direction: the signer uses the secret key, the verifier
uses the public key. In RSA, the signature on $m$ is $s=H(m)^d\bmod n$, and verification compares
$s^e\bmod n$ with $H(m)$. The need for $H$ is not a matter of convenience: without a hash, RSA
signatures are multiplicative like its encryption, so that from the signatures of two documents
someone can assemble a valid signature for their product without holding the secret key. The hash
breaks that property, because $H(m_1m_2)$ is unrelated to $H(m_1)H(m_2)$. The security requirement is
resistance to existential forgery under chosen-message attack, and once again that is an assumption.

\kamus{digital signature}{tanda tangan digital}{a wax seal on a document that anyone can check but
only the holder of the stamp can make}

\section{What the security model actually assumes}
\tingkat{advanced}

Every sentence with the word ``hard'' in this chapter is a promise about what the known algorithms
cannot do, not a statement about what is impossible.

\begin{definisi}
A function $f$ is called one-way if there is a polynomial algorithm to compute $f(x)$ from $x$, but
for random $y$ in the range, every polynomial algorithm fails to find $x$ with $f(x)=y$ except with
negligible probability.
\end{definisi}

The direction is depicted in Figure~\ref{fig:tangga-satu-arah}.

\begin{figure}[htbp]
\centering
\begin{tikzpicture}[font=\small,
  kotak/.style={draw,rounded corners=2pt,align=center,inner sep=5pt,
    minimum height=0.95cm}]
  \node[kotak,draw=layerbrown] (pq) at (0,0) {$p$, $q$ random primes};
  \node[kotak,draw=ink] (n) at (4.0,0) {$n=p\,q$};
  \node[kotak,draw=layergrey,dashed] (s) at (8.2,0) {attacker};
  \draw[-{Latex},very thick,layerjade] (pq) -- (n);
  \node[font=\scriptsize,layerjade,above] at (2.0,0.06)
    {multiply: easy, polynomial};
  \draw[-{Latex},thick,layergrey] (n) -- (s);
  \node[font=\scriptsize,layergrey,above] at (6.1,0.06) {sees $n$};
  \draw[layerplum,thick,dashed] (n.south) -- (4.0,-1.35);
  \draw[layerplum,thick,dashed] (0,-1.35) -- (pq.south);
  \draw[-{Latex},layerplum,thick,dashed] (4.0,-1.35) -- (0,-1.35);
  \draw[layerplum,very thick] (1.86,-1.12)--(1.86,-1.58);
  \draw[layerplum,very thick] (2.04,-1.12)--(2.04,-1.58);
  \node[font=\scriptsize,layerplum,below] at (2.0,-1.62)
    {factor: hard, no known polynomial algorithm};
\end{tikzpicture}
\caption{Two directions that are not equally expensive. From $p$ and $q$ to $n$ one multiplication
suffices, whereas from $n$ back to $p$ and $q$ is a problem that to this day has no known polynomial
algorithm; that gap is what RSA and Diffie-Hellman use. The bar on the return arrow marks a direction
assumed closed, not a direction already proved impossible.}
\label{fig:tangga-satu-arah}
\end{figure}

That definition does not name a single function, and that is not an oversight. To this day there is
no proof that one-way functions exist; their existence would imply $P\neq NP$, which is also still
open. Public-key cryptography therefore stands on a conjecture, and the chain ends in an assumption
that does not turn into a theorem by being used for a long time.

\index{factoring assumption}
The factoring assumption says: for $n$ the product of two random primes of size $k$ bits, there is no
polynomial algorithm in $k$ that factors it. The best known algorithm, the number field sieve, runs
in time $L_n[1/3,1.923]$, which grows subexponentially. The number RSA-250, $250$ digits long or
about $829$ bits, was factored in 2020 at a cost of about $2\,700$ core-years.

The strength equivalences recorded by NIST in SP 800-57 are useful as a rough map.

\begin{center}
\small
\begin{tabular}{@{}lll@{}}
\toprule
Strength (bits) & RSA or DH modulus (bits) & Symmetric cipher \\
\midrule
112 & 2048 & 3DES three keys \\
128 & 3072 & AES-128 \\
192 & 7680 & AES-192 \\
256 & 15360 & AES-256 \\
\bottomrule
\end{tabular}
\end{center}

The numbers in the middle column are cost estimates for the best algorithm at the time the table was
compiled, so a table of this kind has an expiry date.

Formal security models close the gap in the rhetoric with a game: the challenger generates keys,
answers encryption queries, then chooses two messages and sends the encryption of one of them; the
scheme is secure if the attacker cannot guess which one far above the chance of one half. Security is
proved as a reduction: an attacker who wins is turned into an algorithm that factors $n$, so that all
the doubt moves to one place that is narrow but does not disappear.

\section{Shor, and what collapses with it}
\tingkat{advanced}

\index{Shor's algorithm}
The whole building above depends on one difficulty, and in 1994 Peter Shor showed that a quantum
computer solves it in polynomial time.

The idea splits in two, and the first part is entirely classical. To factor $n$, choose $a$ with
$\gcd(a,n)=1$ and compute $r=\operatorname{ord}_n(a)$. If $r$ is odd, repeat with another $a$. If $r$
is even and $a^{r/2}\not\equiv-1\pmod n$, then $\gcd(a^{r/2}-1,n)$ is a proper factor of $n$: because
$a^r\equiv1$, the number $a^{r/2}$ is a square root of one that is not $\pm1$, and a root of that kind
appears only when the modulus splits. The rest of the work is finding the order, and this is the
quantum part.

The classical part alone is not enough. For $n=1019\cdot1031$ we tried $2\,996$ random units and only
$319$, about $10.6$ percent, satisfied the condition and so gave a factor; success therefore comes
from repeating with different $a$. What is missing without a quantum computer is a fast way to find
$r$: no known classical algorithm runs in polynomial time in $\log n$.

A quantum computer finds $r$ with the quantum Fourier transform, and that is where the speedup comes
from. The consequence is not only that RSA collapses: Shor's algorithm also solves the discrete
logarithm, both over $\mathbb{Z}_p$ and over elliptic curves, so that Diffie-Hellman, DSA, ECDH, and
ECDSA collapse too, including the elliptic curves with the shortest moduli because they are immune to
index calculus.

What does not collapse is symmetric ciphers and hash functions. Grover gives a quadratic speedup in
key search, so that AES-128 drops from $2^{128}$ to about $2^{64}$ quantum operations, whereas
AES-256 stays at about $2^{128}$; the collision resistance of hashes is practically unchanged. The
answer for symmetric data is therefore simple: use keys twice as long.

Two things that are often mixed together in the news must be distinguished. Shor proves that the
factorization can be computed in polynomial time on a quantum computer without error; building the
machine is another technical matter. The estimate of Gidney and Ekera (2021) mentions about twenty
million physical qubits and eight hours for a $2048$-bit modulus, almost all of it for error
correction. That number has dropped by several orders of magnitude in one decade, but no machine has
come close. What applies now is the pattern ``harvest now, open later'': today's messages can be
stored and opened when the machine exists.

\section{Post-quantum cryptography: lattices and Kyber}
\tingkat{advanced}

\index{post-quantum cryptography}
If factorization falls, what is needed is another problem that does not fall with it. The strongest
candidates come from lattices, that is sets of points generated by integer combinations of a few
vectors.

\index{lattice}
\emph{Learning with errors} sets up the problem like this. Given a random $m\times n$ matrix $A$ over
$\mathbb{Z}_q$ and a vector $b=As+e\bmod q$, with $s$ secret and $e$ a small error, find $s$. Without
$e$ this is only a system of linear equations solved by Gaussian elimination (Chapter 8); with $e$
small but not zero, the system is almost consistent and finding $s$ is considered hard. Regev's
reduction (2005) adds to its appeal: attacking the random scheme is as hard as solving the worst case
of the lattice problem, not an average case chosen by the attacker.

Kyber, which NIST standardized in August 2024 as ML-KEM in FIPS 203, is a module variant of that idea
with $n=256$, $k=3$, and $q=3329$. Its public key is $1\,184$ bytes and its ciphertext $1\,088$
bytes, about $37$ and $34$ times the $32$-byte symmetric key it wraps, and that space cost is real.
Post-quantum signatures have a similar story: ML-DSA (Dilithium) is also lattice-based, whereas
SLH-DSA (SPHINCS+) takes the conservative hash-based road.

An honesty that is obligatory here: the lattice assumptions are much younger than the factoring
assumption. Factorization has been attacked since Gauss; LWE is only twenty years old, and the
history of cryptography is full of schemes that looked strong and then broke after ten years of
analysis. The recommended transition is therefore not replacement but hybrid: the session key is
formed from Kyber and from Diffie-Hellman exchange at once.

\section{When output becomes an oracle: padding and models}
\tingkat{advanced}

\index{padding oracle}
There is one family of attacks that does not attack the mathematics, but the way the program answers.
Bleichenbacher in 1998 showed that RSA with PKCS\#1 v1.5 padding can be opened completely if the
server answers differently for ``bad padding'' and ``good padding but wrong contents''. Every answer
leaks one bit about the message, and with a few thousand queries the message is read. The mathematics
of RSA is not violated; what is violated is the assumption that the attacker only sees the
ciphertext.

The threat pattern is the same as the one discussed in Chapter 29 and Chapter 30. A machine learning
model that returns probabilities, or error messages that differ, is an oracle; an attacker who may
ask many times uses those small differences to infer what he should not know, including whether an
example is in the training data. Chapter 30 adds a warning that applies in both worlds: a measurement
that seems to improve is not necessarily improving because the defense works.

The defensive lesson is therefore the same. Output must be made uniform: one error message for all
failures, response times that do not depend on secrets, refusals that do not distinguish the cause.
For models, this means limiting the number of queries and disguising scores. All of that reduces
leakage; it does not remove it.

\begin{tantangan}
A library uses AES-GCM with the same key $K$ for two messages, but its nonce counter is reset so that
the same $96$-bit nonce is used twice. The attacker holds both ciphers, both tags, and the contents of
the first message, namely a login page that is fixed. Show how he obtains the authentication key $H$,
then why $H$ is enough to forge a tag for a message of his own making without knowing $K$.

Hint 1: the same nonce gives the same stream key, so that the difference of the two ciphers removes
that key.

Hint 2: the tag is a polynomial in $H$ with known coefficients; two tags with the same nonce give two
equations for one $H$, and their difference becomes a polynomial whose roots can be computed.

If you get stuck, go back to the section ``Hash functions, MACs, and signatures'' and reread the
paragraph about MACs.
\end{tantangan}

\section{Worked example}
\tingkat{intermediate}

We compute one inverse by hand, because the backward steps are where mistakes happen most often. Take
$\varphi(n)=1\,048\,540$ from the toy modulus in part six and $e=65\,537$. Forward Euclid gives
\[
1\,048\,540=15\cdot65\,537+65\,485,\qquad 65\,537=1\cdot65\,485+52,
\]
\[
65\,485=1259\cdot52+17,\qquad 52=3\cdot17+1,\qquad 17=17\cdot1+0 .
\]
The last nonzero remainder is $1$, so that $\gcd=1$ and the inverse exists. Build backward starting
from the remainder one:
\[
1=52-3\cdot17=52-3(65\,485-1259\cdot52)=3778\cdot52-3\cdot65\,485 .
\]
Continue with $52=65\,537-65\,485$ and $65\,485=1\,048\,540-15\cdot65\,537$:
\[
1=3778\cdot65\,537-3781\cdot65\,485
=3778\cdot65\,537-3781(1\,048\,540-15\cdot65\,537).
\]
Group: $1=(3778+15\cdot3781)\cdot65\,537-3781\cdot1\,048\,540
=60\,493\cdot65\,537-3781\cdot1\,048\,540$. So $d=60\,493$, and the negative sign disappears
because it is read modulo $\varphi(n)$. The check is one line: $65\,537\cdot60\,493
=3\,964\,529\,741$, whereas $3781\cdot1\,048\,540=3\,964\,529\,740$; the difference is one.

Two other computations in this chapter can be checked in the same way: the trimming of
$7^{222}\bmod13$ in part three, which uses $222=12\cdot18+6$ and gives $12$, and the small key
exchange $p=23$, $g=5$, $a=4$, $b=3$ which gives $A=4$, $B=10$, and key $18$ from both directions.

\begin{resep}
Four steps to find the inverse of $a$ modulo $n$.
\begin{enumerate}
  \item Write $(a,n)$. Run Euclid: replace $(a,n)$ by $(n,a\bmod n)$ until the remainder is zero;
        record each line $n_0=q\,n_1+r$.
  \item The last nonzero remainder is $\gcd$. If it is not $1$, stop; the inverse does not exist.
  \item From the last line with remainder $1$, write $1=r-q\,s$, then replace $s$ and $r$ by their
        forms from the line above, one line up for each step.
  \item Collect the coefficient of $a$; that is the inverse. If it is negative, add $n$ until it lies
        in $0,\dots,n-1$, then check $a\cdot x\bmod n=1$.
\end{enumerate}
Warning about order: build backward from the smallest remainder upward, not from the first line.
Swapping the roles of $a$ and $n$ also changes the answer: the inverse of $a$ modulo $n$ is not the
inverse of $n$ modulo $a$.
\end{resep}

\kertas{Repeat the first computation without looking at the text. Write five lines of forward Euclid
for $65\,537$ and $1\,048\,540$, then build backward until you obtain $d$. Check by multiplying $e$
and $d$ and dividing by $\varphi(n)$: the remainder must be one. After that, answer: why can the same
steps not be used to open RSA without the secret key, even though this whole computation needs only
five divisions?}

\section{Computational lab}
\tingkat{advanced}

The following program uses only the standard library and runs in under one second. It checks four
claims in the body of the chapter: the units modulo $26$, the statistics of Euclid's steps, the
equality of the Diffie-Hellman keys, and the looseness of plain RSA. The last part tests the birthday
prediction on SHA-256 truncated to two bytes.

\begin{lab}
\begin{verbatim}
import hashlib, math, random

def egcd(a, b):
    if b == 0:
        return a, 1, 0
    g, x, y = egcd(b, a % b)
    return g, y, x - (a // b) * y

def invers(a, n):
    g, x, _ = egcd(a, n)
    return None if g != 1 else x % n

def langkah(a, b):
    k = 0
    while b:
        a, b = b, a % b
        k += 1
    return k

def orde(a, n):
    k, y = 1, a % n
    while y != 1:
        y = y * a % n
        k += 1
    return k

unit = [a for a in range(26) if invers(a, 26) is not None]
print("units mod 26:", len(unit), unit)
print("inverse 3 mod 26 =", invers(3, 26), " inverse 2 mod 26 =", invers(2, 26))

g0, x0, y0 = egcd(1071, 462)
print("gcd(1071,462) =", g0, " x,y =", x0, y0, " steps =", langkah(1071, 462))

p, g, a, b = 10007, 5, 1234, 5678
A, B = pow(g, a, p), pow(g, b, p)
print("DH: A =", A, " B =", B, " K =", pow(B, a, p), pow(A, b, p))
print("order(5) =", orde(5, p), " order(5^5003) =", orde(pow(5, 5003, p), p))

random.seed(29)
L = [langkah(random.getrandbits(32), random.getrandbits(32))
     for _ in range(20000)]
print("euclid 32-bit: mean %.3f, max %d, min %d" % (sum(L) / len(L), max(L), min(L)))

p, q, e = 1019, 1031, 3
n, phi = p * q, (p - 1) * (q - 1)
d = invers(e, phi)
print("RSA: n =", n, " bits =", n.bit_length(), " phi =", phi, " d =", d)
for m in (424242, 42):
    c = pow(m, e, n)
    print("  m =", m, " c =", c, " D(c) =", pow(c, d, n), " m^3 =", m ** 3)
kali = pow(1234, e, n) * pow(567, e, n) % n
print("E(1234)E(567) mod n =", kali, " D(kali) =", pow(kali, d, n),
      " 1234*567 mod n =", 1234 * 567 % n)

def tabrakan():
    terlihat, k = {}, 0
    while True:
        m = random.getrandbits(32)
        h = hashlib.sha256(m.to_bytes(4, "big")).digest()[:2]
        k += 1
        if h in terlihat:
            return k
        terlihat[h] = m

T = [tabrakan() for _ in range(20)]
print("16-bit collisions: mean %.1f, min %d, max %d" % (sum(T) / 20, min(T), max(T)))
print("prediction sqrt(pi/2 * 2^16) =", round(math.sqrt(math.pi / 2 * 2**16), 1))
\end{verbatim}

The output, as it stands:

\begin{verbatim}
units mod 26: 12 [1, 3, 5, 7, 9, 11, 15, 17, 19, 21, 23, 25]
inverse 3 mod 26 = 9  inverse 2 mod 26 = None
gcd(1071,462) = 21  x,y = -3 7  steps = 3
DH: A = 4873  B = 3189  K = 5674 5674
order(5) = 10006  order(5^5003) = 2
euclid 32-bit: mean 18.750, max 32, min 6
RSA: n = 1050589  bits = 21  phi = 1048540  d = 699027
  m = 424242  c = 618000  D(c) = 424242  m^3 = 76355615883580488
  m = 42  c = 74088  D(c) = 42  m^3 = 74088
E(1234)E(567) mod n = 551173  D(kali) = 699678  1234*567 mod n = 699678
16-bit collisions: mean 362.7, min 100, max 743
prediction sqrt(pi/2 * 2^16) = 320.8
\end{verbatim}
\end{lab}

The first line confirms part one: twelve units modulo $26$, with $1$ and $25$ their own inverses and
the remaining ten paired; that number $12$ is the same as $\varphi(26)$.

The second and third lines show two faces of Euclid: three divisions for $1071$ and $462$ with Bezout
coefficients $(-3,7)$, and an average of $18.75$ divisions with maximum $32$ for two thousand random
$32$-bit pairs. The fourth and fifth lines are the promised key exchange: $4873$ and $3189$ on the
channel, $5674$ at both ends, whereas $5^{5003}$ has order $2$ and is therefore the trap of part five.

The three RSA lines show two opposite things. For $m=424\,242$ everything works: the cipher
$618\,000$ returns to $424\,242$. For $m=42$ with $e=3$, the cipher is $74\,088$, exactly $42^3$, so
that no key needs to be found, only a cube root; the next line shows the multiplicative property: the
product of two ciphers opens the product message.

The last four lines test the birthday prediction: an average of $362.7$ samples against a prediction
of $320.8$, with the wide spread expected from an exponential distribution.

\begin{untukdosen}
The teaching order I suggest: start from the student's cipher in the opening, then the
well-definedness lemma and the unit theorem in the first meeting; leave Diffie-Hellman and RSA for the
second meeting, and put Shor and post-quantum as reading for the third. Do not teach Euclid before
asking ``when does division exist''.

Two misconceptions almost always appear. First, students think $\varphi(n)=n-1$ for all $n$, then are
startled when $\varphi(26)=12$; cure it by asking them to list the units modulo $26$ on the board
before any formula is mentioned. Second, they think RSA is secure because $n$ is hard to factor,
without mentioning the size of $n$ and without mentioning that this difficulty is an assumption; ask
them to write the security claim in the form ``if ... then ...''.

A blackboard activity: two groups compute $3^{12}\bmod26$, one by repeated multiplication and one by
Fermat's little theorem, then compare time and results. A sign of readiness: they can explain without
notes why $[2]_{26}$ has no inverse.
\end{untukdosen}

\section{Summary}
\begin{ringkasan}
\begin{itemize}
    \item Arithmetic modulo $n$ works on residue classes, and division exists only for classes that
        share no factor with $n$; Euclid gives the inverse in time proportional to the number of
        digits.
    \item The shuffling lemma, namely that multiplication by a unit is only a shuffling of the order,
        gives Euler's theorem in five lines and Fermat's little theorem as a corollary. The order
        divides $\varphi(n)$, and a primitive root is an element of full order.
    \item Diffie-Hellman agrees on a key through the equality $g^{ab}$ from two directions, with
        security conditional on the difficulty of the discrete logarithm; without authentication it
        falls to the man-in-the-middle attack.
    \item RSA reverses the direction: the public key locks, the secret key opens, and the proof uses
        Fermat's little theorem on both prime factors. Plain RSA is deterministic, multiplicative, and
        does not protect small messages; padding closes all three.
    \item A hash function gives a fingerprint, a MAC gives authenticity between two parties who share
        a key, a signature gives evidence that anyone can check; collision resistance is only half the
        output length, and the numbers $362.7$ against $320.8$ in the lab show that the prediction
        works.
    \item Every guarantee in this chapter is conditional. Shor brings down RSA and the whole of
        discrete logarithms if a large-scale quantum computer exists; symmetric ciphers and hashes
        only need to double the key length, and lattice schemes such as ML-KEM are the possible
        replacements.
\end{itemize}
\end{ringkasan}

\section{Exercises}
\begin{latihan}
\begin{enumerate}
  \item Assess this claim: ``Because $2^{100}\equiv1\pmod{101}$, the number $101$ is proved prime.''
        Decide true, false, or not enough evidence, and name what that congruence does prove.
  \item Prove that the product of two units modulo $n$ is also a unit, then conclude that $U(n)$ is
        closed under multiplication; show also that if $a$ is not a unit, there is $b\neq0$ with
        $ab\equiv0\pmod n$.
  \item Prove that for prime $p$, $a^2\equiv1\pmod p$ is satisfied only by $a\equiv1$ or
        $a\equiv-1$; explain why the same fails for modulus $8$, and name all square roots of one
        modulo $8$.
  \item Prove $\varphi(pq)=\varphi(p)\varphi(q)$ for distinct primes $p,q$ by counting the classes
        $0\le a<pq$ that share a factor with $pq$, then check it on $p=1019$ and $q=1031$.
  \item Compute $17^{-1}\bmod3120$ with backward Euclid, then check it with Python. Repeat for
        $65\,537^{-1}\bmod1\,048\,540$ and compare the number of divisions of the two.
  \item Compute $5^{100}\bmod101$, $3^{200}\bmod101$, and $2^{1000}\bmod101$ using Fermat's little
        theorem, without computing the full power. Check all three with \texttt{pow(a,k,m)}.
  \item On modulus $23$ with base $5$, Alice uses $a=7$ and Bob $b=15$. Compute $A$, $B$, and the
        shared key, then compute the key obtained by an active eavesdropper who replaces $A$ with
        $11$ in front of Bob, and explain why two different keys appear.
  \item Explain why RSA with $e=3$ and a message $m<n^{1/3}$ can be opened without the secret key,
        how long the longest message affected on the modulus $1\,050\,589$ is, and what padding does
        to prevent it.
  \item Assess this claim: ``ML-KEM is already proved secure, so we can stop using RSA now.'' Name
        two things that must be true for the first sentence to make sense, and one reason why the
        second sentence remains wrong even if the first is true.
  \item A system uses a hash with a $48$-bit output. How many samples are needed for a collision
        probability of one in two, and how many if the output is truncated to $32$ bits?
\end{enumerate}

Hints for the three hardest items. Item 3: write $a^2-1=(a-1)(a+1)$ and use Euclid's lemma; for
modulus $8$, find the roots by trying $a=0,\dots,7$. Item 7: compute $11^7\bmod23$ as the key that
Bob sees, then compare it with $11^{15}\bmod23$ computed by Alice. Item 9: distinguish ``proved
secure'' from ``secure against known attacks''.
\end{latihan}

\begin{untukpeneliti}
What is really proved in this chapter is its mathematical part: Bezout's theorem, the unit theorem,
Euler's theorem, the order lemma, and the RSA theorem all have complete proofs and do not depend on
any conjecture. What is only popularized is the word ``secure''. No theorem states that RSA-2048 is
hard; what exists are cost estimates for the best known algorithms, and those estimates change
whenever cryptanalysis makes progress.

The primary sources: Diffie and Hellman (1976), sections 3--4; Rivest, Shamir, and Adleman (1978),
sections 2--4; Shor (1997), section 5 for the reduction of factorization to order finding; Regev
(2005), section 3 for the LWE reduction; Bleichenbacher (1998); Gidney and Ekera (2021); FIPS 203 and
204.

Limits of validity that must be mentioned: security reductions usually produce loose bounds, the
attack models do not contain side channels, and padding attacks show that the difference of one error
message is enough to bring down a scheme that is proved secure in the model. Open questions: how tight
the reduction for ML-KEM is, whether there is a faster quantum algorithm for lattice problems, and how
to prove the security of schemes used together with secret data without leaking access patterns.
\end{untukpeneliti}

\peta{Chapter 31 follows Diffie and Hellman (1976), section 3 for key exchange; Rivest, Shamir, and
Adleman (1978), sections 2--4 for RSA and signatures; Shor (1997), sections 5--6 for the reduction of
factorization to order finding; Regev (2005), section 3 for the LWE reduction; Bleichenbacher (1998);
Gidney and Ekera (2021); FIPS 203 and 204. On Chapter 2 for bijections and partitions, Chapter 3 for
the Archimedean property and the discrete form of the infimum property, Chapter 8 for Gaussian
elimination, Chapter 26 for the reason a security claim cannot be a theorem without conditions, and
Chapters 29--30 for the oracle threat pattern.}

% Running head would be 68.9 pt too wide with the full title; the chapter page and
% the table of contents keep the complete title.
\judulpendek{Differential Privacy and QML}
\chapter{Differential Privacy and the Mathematical Foundations of QML}

A dental clinic collects the records of ten thousand patients, trains a model to
predict who is at high risk, and then publishes that model as a file of weights. The
following month a student takes one name from the patient list that leaked through
another channel, feeds that person's medical record into the model, and reads the
output score. The score is far higher than the scores of other people with similar
medical records who are not patients of that clinic. The student did not break into
the database and did not use anyone's password. He only asked the model, again and
again, and the model answered too confidently for someone it had seen during training.

A parallel to that incident sits at the other end of this book. A paper reports
``quantum advantage'' because its circuit prepares a state with $2^{50}$ amplitudes,
a number that no computer can store. The number is correct. What the paper does not
say: the amplitudes cannot be read, and the measurement result is only a sample.

The two stories have a defect of the same shape: the claim runs faster than the
mathematics that supports it. The first half of the chapter builds the definition
that makes the sentence ``this model is private'' testable, namely \emph{differential
privacy}, and then shows why that definition takes the form of a comparison of two
worlds rather than a prohibition on remembering. The second half builds the foundation
for assessing the sentence ``quantum models are superior'': complex amplitudes,
unitary gates, the Born rule, and the cost of simulation.

\index{differential privacy}
\index{quantum advantage}

\kamus{differential privacy}{privasi diferensial}{a promise that the result of a
procedure barely changes, in the sense of probability, when one person enters or leaves
the training data}

\begin{untukanak}
Imagine an ice cream shop that records the average amount spent by all customers that
day, and then sticks that number on the door. If Sari is the only customer who bought a
large birthday cake, the average that is posted goes up too, and a neighbor who knows
that it is Sari's birthday can guess that Sari is on that day's list of customers. The
number is not a name, but the number still tells a story about one person.

Now imagine a more careful shop. Before posting the number, it adds a small random
jitter, say thirty cents up or down, with equal probability. Whether Sari is included or
not, the number that is posted is nearly the same, because the jitter is much larger
than the influence of one person. Other customers can still learn that day's average
spending accurately enough, but there is no longer any way to determine whether Sari is
on the list.

The question that makes one curious: if the jitter is too small, the secret leaks; if
it is too large, the number is useless. Where is the honest boundary between them, and
how can we prove that the boundary is right? In this chapter we will learn why a privacy
promise has to take the form of a comparison of two nearly identical worlds, and not a
promise that the computer will forget.
\end{untukanak}

\section{What leaks from a model trained on personal data}
\tingkat{basic}

A model does not store a copy of the training data. It stores parameters, and those
parameters are chosen so as to fit the data. Precisely because of that choice the data
leave a trace, and the trace can be read in three ways of different difficulty.

The first way, \emph{membership inference}. The attacker has one medical record and
wants to know whether that record is in the training data. He does not need to open the
model; it is enough to send the record and read the output. A model trained on that
record usually answers with higher confidence, because the parameters have already been
adjusted to include it. Shokri, Stronati, Song, and Shmatikov showed in 2017 that this
difference in confidence is enough to be separated by a second classifier trained on
``shadow models'', namely copies of the model trained on data with and without a
particular example. The attack works on a model that can only be queried, with no
weights and no gradients.

The second way, \emph{attribute inference}. The attacker does not ask about membership,
but about a property that is not in the data: whether this person smokes, whether she is
pregnant. A model trained on a particular population can reveal correlations that should
not be used to decide anything about one person.

The third way, \emph{reconstruction}: the attacker recovers part of the contents of the
data. Dinur and Nissim showed in 2003 that a sequence of aggregate queries with small
noise still leaks most of the contents of a database when the queries are numerous and
interdependent. Zhu, Liu, and Han showed in 2019 that gradients in distributed training
can be used to recover training images almost exactly.

Those three channels explain why the sentence ``we publish only parameters, not data''
is not enough. What has to be promised is not ``nothing was copied'', but ``the output of
this procedure barely depends on one person''. To state that mathematically we need two
definitions of distance: the distance between two datasets that differ in one person,
and the distance between two output distributions.

\begin{definisi}[Neighboring datasets]
Two datasets $D$ and $D'$ over the same domain are called \emph{neighboring} when the
two are identical except in one record: $D'$ is obtained from $D$ by adding or removing
one record, or by replacing one record with another. The set of all such pairs is written
$\mathcal{N}$.
\end{definisi}

\index{membership inference}
The choice ``add or remove'' is called unbounded adjacency, ``replace'' is called bounded
adjacency, and the two give different sensitivity numbers. A privacy claim that does not
state the choice is incomplete.

\section{The definition of differential privacy, and why not an average}
\tingkat{intermediate}

``Nearly the same'' is measured with a ratio of probabilities, not with a difference,
because the ratio does not change when the probabilities are very small, and it is
precisely very small probabilities that attackers most often use to infer something.

\begin{definisi}[$(\varepsilon,\delta)$ differential privacy]
\index{differential privacy}
Let $\varepsilon\ge 0$ and $\delta\in[0,1]$. A randomized procedure $M$ that maps
datasets to an output space $\mathcal{Y}$ satisfies $(\varepsilon,\delta)$ differential
privacy when for every neighboring pair $(D,D')\in\mathcal{N}$ and every measurable set
$S\subseteq\mathcal{Y}$ we have
\begin{equation}
P\bigl[M(D)\in S\bigr] \;\le\; e^{\varepsilon}\,P\bigl[M(D')\in S\bigr] + \delta .
\end{equation}
If $\delta=0$, the procedure satisfies pure $\varepsilon$ differential privacy.
\end{definisi}

\begin{figure}[htbp]
\centering
\begin{tikzpicture}
\begin{axis}[width=6.2cm,height=4.1cm,xmin=-2.6,xmax=3.6,ymin=0,ymax=4.6,
  domain=-2.6:3.6,samples=200,axis lines=left,
  xtick={-2,-1,0,1,2,3},ytick={0,2,4},
  tick label style={font=\scriptsize},label style={font=\scriptsize},
  title={$\varepsilon=8$},title style={font=\scriptsize},
  xlabel={output},ylabel={density},
  legend style={font=\scriptsize,draw=none,at={(0.03,0.97)},anchor=north west},
  legend cell align=left]
  \addplot[thick,accent,fill=accent!12] {4*exp(-8*abs(x))} \closedcycle;
  \addlegendentry{data $D$}
  \addplot[thick,layerteal,densely dashed] {4*exp(-8*abs(x-1))};
  \addlegendentry{data $D'$}
\end{axis}
\end{tikzpicture}
\hfill
\begin{tikzpicture}
\begin{axis}[width=6.2cm,height=4.1cm,xmin=-2.6,xmax=3.6,ymin=0,ymax=0.142,
  domain=-2.6:3.6,samples=200,axis lines=left,
  xtick={-2,-1,0,1,2,3},ytick={0,0.05,0.1},yticklabels={$0$,$0.05$,$0.1$},
  tick label style={font=\scriptsize},label style={font=\scriptsize},
  title={$\varepsilon=0.25$},title style={font=\scriptsize},
  xlabel={output},
  legend style={font=\scriptsize,draw=none,at={(0.03,0.97)},anchor=north west},
  legend cell align=left]
  \addplot[thick,accent,fill=accent!12] {0.125*exp(-0.25*abs(x))} \closedcycle;
  \addlegendentry{data $D$}
  \addplot[thick,layerteal,densely dashed] {0.125*exp(-0.25*abs(x-1))};
  \addlegendentry{data $D'$}
\end{axis}
\end{tikzpicture}
\caption{The output distribution of the Laplace mechanism at sensitivity $1$ for two
neighboring databases that differ in one record, with a different vertical scale in the
two panels. At $\varepsilon=8$ the two curves are far apart, so one output is already
enough to guess which data were used. At $\varepsilon=0.25$ the two almost coincide: the
ratio of their ordinates at every point is at most $e^{0.25}\approx1.284$, so one output
hardly favors either dataset.}
\label{fig:dua-sebaran}
\end{figure}

A new reader usually asks two things: why a comparison of probabilities rather than an
average, and why there is a $\delta$ at all if $\varepsilon$ already measures the
leakage.

The answer to the first question is the core of the first half of the chapter. Note the
order of the quantifiers: \emph{for every} $S$, \emph{for every} neighboring pair. That
promise has to hold for any test $S$, including a test that the attacker picks after
seeing the mechanism. A definition in the form of an average puts the quantifier inside:
``there is a small expectation''. The two forms are not equivalent;
Figure~\ref{fig:dua-sebaran} shows the difference at two values of $\varepsilon$.

\begin{proposisi}
Let $M$ satisfy $(\varepsilon,0)$ differential privacy. Then for every attacker
algorithm $A$ whose output is only $0$ or $1$, and every neighboring pair $(D,D')$,
\begin{equation}
P\bigl[A\bigl(M(D)\bigr)=1\bigr] \;\le\; e^{\varepsilon}\,
P\bigl[A\bigl(M(D')\bigr)=1\bigr] .
\end{equation}
\end{proposisi}

\begin{proof}
The set $\{y : A(y)=1\}$ is measurable, so it may be used as $S$ in the definition.
Direct substitution gives the result.
\end{proof}

That proposition looks trivial, and its use is precisely to close the door that is often
used to defend a weak mechanism. If the promise holds only for the \emph{average} output,
an attacker can average the results of many repetitions so that the random components
cancel one another and what remains is a signal about one person. With a per-event
promise, every repetition pays its own privacy budget.

\begin{bacakeras}
Read alternately, two voices. The first voice asks, the second voice answers.

``If one person leaves the list, does your answer change?''\\
``Slightly.''\\
``How slightly?''\\
``Slightly enough that nobody can be certain.''

Repeat that conversation three times, faster and faster, until it feels like a
recitation. Then the first voice asks again, more slowly:

``If I repeat that question a hundred times, is a hundred times slight still slight?''

Stay silent for one breath. Answer together: something slight repeated a hundred times
is no longer slight. Then say together, once: one person must not be able to change the
conclusion. Repeat that line once more, more slowly, and then once again while pointing
at one empty chair in the room.
\end{bacakeras}

\begin{contoh}[Randomized response]
\index{randomized response}
A respondent is asked whether he has ever used an illegal drug. The procedure: toss a
coin; if side A comes up, answer honestly; if side B comes up, toss again and answer
``yes'' if heads, ``no'' if tails. The probability of answering ``yes'' if he really uses
the drug is $1/2+1/2\cdot 1/2=3/4$, and if he does not use it the probability is
$1/2\cdot 1/2=1/4$. The ratio of the two is $3$, so $\varepsilon=\ln 3=1.0986$. No answer
proves anything about one person, and that is what the definition promises.
\end{contoh}

\begin{tangan}
Materials: one coin per person, paper, pencil. Ten friends each write one yes/no secret
on a slip of paper and fold it.

The rule is one. Toss the coin: heads, answer honestly; tails, toss again, heads means
``yes'' and tails means ``no''. Each person says only one word, and the recorder writes
only two groups of tally marks.

Now ask one person to change his secret, then repeat the whole round. Look at the tally
marks: they shift by at most one, and nobody can point to anyone. Conclusions about the
group survive. A conclusion about one person never appears.
\end{tangan}

Now $\delta$, namely the failure probability: with probability at most $\delta$ the ratio
promise may be violated entirely. Hence $\delta$ is not a tolerance that can be chosen at
will. If $\delta$ is comparable to $1/n$ for $n$ records, an attacker can target one
person and let the rest serve as camouflage; that is why $\delta$ must be much smaller
than $1/n$. The number $\varepsilon$ is not a probability either: $\varepsilon=1$ means a
probability ratio of at most $e\approx 2.718$, whereas $\varepsilon=8$ means a ratio up to
$2\,981$.

\section{Global sensitivity and the Laplace mechanism}
\tingkat{intermediate}

The only thing that stops us from adding noise arbitrarily is the question: how far can
one record shift the answer?

\begin{definisi}[Global sensitivity]
\index{global sensitivity}
For $f$ that maps datasets to $\mathbb{R}^d$, its global sensitivity is
\begin{equation}
\Delta f \;=\; \sup_{D\sim D'} \bigl\| f(D)-f(D') \bigr\|_{1},
\end{equation}
with the supremum taken over all neighboring pairs.
\end{definisi}

A count query, $f(D)=|D|$, has $\Delta f=1$. A sum query over values in $[0,R]$ has
$\Delta f=R$. An average query over $n$ records has $\Delta f=R/n$, because one person
shifts the average by only $R/n$; here it becomes visible that an average is safer than a
sum, but only as long as $n$ is known and fixed. Sensitivity also depends on the choice
of adjacency: if $D'$ may replace, the count query has $\Delta f=2$.

\begin{teorema}[Laplace mechanism]
\index{Laplace mechanism}
Let $f$ have global sensitivity $\Delta f$, and let $Y$ be distributed according to a
Laplace distribution with scale $b=\Delta f/\varepsilon$, that is with density
$p_b(y)=\frac{1}{2b}e^{-|y|/b}$. Then $M(D)=f(D)+Y$ satisfies $(\varepsilon,0)$
differential privacy.
\end{teorema}

\begin{proof}
Take $S\subseteq\mathbb{R}$ measurable and $D\sim D'$. It is enough to show that for
every $y$,
\begin{equation}
\frac{p_b\bigl(y-f(D)\bigr)}{p_b\bigl(y-f(D')\bigr)} \;\le\; e^{\varepsilon}.
\end{equation}
With $a=f(D)$ and $a'=f(D')$, the left-hand side is $\exp\bigl((|y-a'|-|y-a|)/b\bigr)$.
The triangle inequality gives $|y-a'|-|y-a|\le|a-a'|$, so the left-hand side is at most
$\exp(|a-a'|/b)\le\exp(\Delta f/b)=e^{\varepsilon}$, because $\Delta f$ is the supremum.
Hence $p_b(y-a)\le e^{\varepsilon}p_b(y-a')$ for every $y$; integrating both sides over
$S$ gives $P[M(D)\in S]\le e^{\varepsilon}P[M(D')\in S]$, and because there is no
remainder term, $\delta=0$.
\end{proof}

A numerical example goes like this. For a count query with $\varepsilon=0.5$, the scale
is $b=1/0.5=2$; the probability that the noise exceeds $5$ in absolute value is
$e^{-5/2}=0.0821$, so about one answer in twelve misses by more than five people.
Conversely, with probability $0.95$ the error does not exceed
$b\ln 20=2\times 2.9957=5.991$. Here the trade-off that recurs throughout the chapter
becomes visible: shrinking $\varepsilon$ enlarges $b$, and enlarging $b$ worsens
precision linearly.

\begin{resep}
\begin{enumerate}
  \item Fix $f$ and clip every value to the range $[0,R]$ before anything is computed.
  \item Compute $\Delta f$ from the worst case of one record on the clipped data.
  \item Choose $\varepsilon>0$, compute $b=\Delta f/\varepsilon$, add Laplace noise of
        scale $b$ to the answer.
  \item Record $\varepsilon$ in the budget book; the remaining budget is what has not
        been spent.
\end{enumerate}
A warning about order: clipping after the noise has been added changes the output
distribution and cancels the guarantee, because $\Delta f$ was computed from the data
before clipping.
\end{resep}

\kamus{global sensitivity}{sensitivitas global}{how far one person can shift the answer
of a query in the worst case; this is what determines the amount of noise that is needed}

Sometimes Laplace noise does not fit, for instance because the output is
high-dimensional. There Gaussian noise is used: for $f$ with $\ell_2$ sensitivity
$\Delta_2 f$, the mechanism $M(D)=f(D)+N(0,\sigma^{2}\Delta_2 f^{2}I)$ satisfies
$(\varepsilon,\delta)$ when $\sigma\ge\sqrt{2\ln(1.25/\delta)}/\varepsilon$. The proof is
skipped here; it compares two shifted normals with a concentration inequality, and the
full version is in Dwork and Roth (2014), appendix A. That condition is sufficient, not
necessary: for small $\sigma$ it is far looser than the true bound, and the tight bound
is given by the analytic Gaussian mechanism of Balle and Wang (2018).

\section{Composition: the privacy budget and why it runs out}
\tingkat{advanced}

A single mechanism is rarely used once. A model is trained with hundreds of steps, each
leaking a little, and what is needed is the total.

\begin{teorema}[Basic composition]
\index{composition}
Let $M_1,\dots,M_T$ be chosen adaptively, that is $M_t$ may depend on the outputs
$M_1,\dots,M_{t-1}$, and let every $M_t$ satisfy $(\varepsilon_t,\delta_t)$ for all
possible inputs. Then the joint output $(M_1,\dots,M_T)$ satisfies
$\bigl(\sum_t\varepsilon_t,\sum_t\delta_t\bigr)$.
\end{teorema}

\begin{proof}
It is enough to take $T=2$; the rest is induction. By the law of total probability,
\begin{align}
P\bigl[(M_1,M_2)\in S_1\times S_2\bigr]
&= \int_{S_1} P\bigl[M_1(D)\in dy_1\bigr]\,
   P\bigl[M_2(D,y_1)\in S_2 \mid M_1(D)=y_1\bigr] \notag\\
&\le \int_{S_1} e^{\varepsilon_1}P\bigl[M_1(D')\in dy_1\bigr]\,
   \bigl(e^{\varepsilon_2}P[M_2(D',y_1)\in S_2]+\delta_2\bigr) + \delta_1 \notag\\
&\le e^{\varepsilon_1+\varepsilon_2}P\bigl[(M_1,M_2)(D')\in S_1\times S_2\bigr]
   + \delta_1 + \delta_2 .
\end{align}
The first step uses the guarantee for $M_2$ at fixed $y_1$, the second step uses the
guarantee for $M_1$ and $P[M_1(D')\in S_1]\le 1$.
\end{proof}

Basic composition adds up, and that addition is what makes the budget run out quickly: a
hundred steps with $\varepsilon=0.1$ per step already spend $\varepsilon=10$. There is a
sharper bound when the leakage per step is small.

\begin{tantangan}
An office publishes the average salary of its team with Laplace noise of scale
$b=\Delta f/\varepsilon$, with $\varepsilon=1$ and $\Delta f=1$. It says, ``one answer
already satisfies $\varepsilon=1$, so one person is safe.'' How many times may that
answer be published before the average of those answers is sharp enough to guess the
salary of one person, and why?

Hint 1: random noise has mean zero; watch what happens to the average of the noise when
the answer is repeated.

Hint 2: the standard error of the average of $k$ answers is $b\sqrt{2}/\sqrt{k}$. With
$b=1$ and $k=100$ its value is $0.1414$, whereas one person can shift the answer by $1$.

If you get stuck, go back to the section on Composition.
\end{tantangan}

\begin{teorema}[Advanced composition]
For every $\varepsilon'>0$, the composition of $T$ procedures that each satisfy
$(\varepsilon,\delta)$ satisfies $(\varepsilon_{\mathrm{tot}},\delta_{\mathrm{tot}})$
with
\begin{equation}
\varepsilon_{\mathrm{tot}} = \sqrt{2T\ln(1/\delta')}\;\varepsilon
+ T\,\varepsilon\bigl(e^{\varepsilon}-1\bigr),
\qquad
\delta_{\mathrm{tot}} = T\delta + \delta' .
\end{equation}
\end{teorema}

The proof is skipped: it is obtained by bounding the moments of the probability ratio, or
with a moments \emph{accountant}. The references are Dwork and Roth section 3.5, and
Kairouz, Oh, and Viswanath (2015) for optimal composition. Its first term grows like
$\sqrt{T}$, not like $T$.

However, advanced composition is not always better, and this is often misquoted. The
second term, $T\varepsilon(e^{\varepsilon}-1)$, blows up when $\varepsilon$ is not small.
In the experiment in section eleven, with $T=200$ and $\delta=10^{-5}$, basic composition
gives $\varepsilon=36.5$ at $\sigma=32$ whereas advanced composition gives $20.6$; at
$\sigma=2$ the two give $583.7$ and $11\,302.8$. So advanced composition wins when
$\varepsilon$ per step is small and loses when it is large.

One more result makes the budget go further: when several procedures work on disjoint
parts of the data, their leakage is not added but the largest is taken, because one
record is in only one part. Splitting the data among several users, each with its own
budget, is far more economical than one user who runs all the queries.

\index{privacy budget}

\section{DP-SGD: clip, add noise, and the price that is paid}
\tingkat{advanced}

Ordinary training satisfies no definition at all, because one example can shift the
gradient by any amount: there is no upper bound on $\Delta f$, so the noise that is
needed is infinite. The repair consists of two steps that support each other.

The first step, clipping: every per-example gradient is clipped so that its norm does not
exceed $C$,
\begin{equation}
\operatorname{clip}_C(g) \;=\; g\cdot\min\Bigl(1,\;\frac{C}{\|g\|_2}\Bigr).
\end{equation}
After clipping, one example contributes at most $C$ to the sum, so the sensitivity of the
average over a batch of size $B$ is $C/B$. The second step, Gaussian noise with standard
deviation $\sigma C/B$ is added to that average,
\begin{equation}
w_{t+1} = w_t - \eta\Bigl(\frac{1}{B}\sum_{i\in\mathcal{B}_t}
\operatorname{clip}_C\bigl(\nabla\ell_i(w_t)\bigr) + \frac{\sigma C}{B}\,\xi_t\Bigr),
\qquad \xi_t\sim N(0,I).
\end{equation}
This algorithm is what is called DP-SGD, introduced by Abadi and colleagues in 2016
together with the moments \emph{accountant} that makes its budget computable for hundreds
of thousands of steps.

\begin{teorema}
Every DP-SGD step satisfies $(\varepsilon_0,\delta)$ with
\begin{equation}
\varepsilon_0=\sqrt{2\ln(1.25/\delta)}\,\big/\,\sigma ,
\end{equation}
and after $T$ steps the whole procedure satisfies $(T\varepsilon_0,T\delta)$.
\end{teorema}

\begin{proof}
The $\ell_2$ sensitivity of the average of clipped gradients is $C/B$, because replacing
one example shifts the average by at most $C/B$. Apply the Gaussian mechanism from
section three, then basic composition from section four.
\end{proof}

\index{dp-sgd}
The numbers from our experiment explain why the noise does not immediately erase the
learning. With $n=200$, $C=1$, and $\sigma=32$, the standard deviation of the noise per
coordinate is $32/200=0.16$, whereas the norm of the average gradient at the starting
point is $0.83$: the noise is about one fifth of the signal. But the noise changes sign
at random, so its contribution to the final position grows like $\sqrt{T}$, whereas the
signal grows like $T$. After two hundred steps the signal wins, and the model still
reaches $93.71$ percent.

What has to be paid is also visible. At $\sigma=128$ the accuracy falls to $85.51$
percent with a spread of $6.16$ across seeds; at $\sigma=512$ the accuracy is $70.39$
percent with a spread of $17.26$. That growing spread is not an oddity: at large noise
the result of training depends on the noise draws that happen to occur, so a single model
is no longer a sample that represents the procedure. Reporting one seed at large $\sigma$
would be misleading.

Four limits have to be mentioned together with that result. First, clipping introduces
bias: examples with large gradients, usually the hard or rare ones, are clipped harder, so
the model pays them less attention. Second, DP-SGD demands more steps, and that
computational cost is real. Third, the privacy guarantee is uniform but its utility cost
is not: Bagdasaryan, Poursaeed, and Shmatikov showed in 2019 that minority groups can
suffer a much larger drop in accuracy at the same budget. Fourth, and most important:
differential privacy guarantees something about \emph{individuals}, not about the
\emph{truth of the model}. A model with $\varepsilon=0.1$ can still be completely wrong,
can still be discriminatory, and can still be used for unfair decisions. Privacy is a
promise about how much one person can be learned from the output; it is not in the least a
promise that the output is any good.

\kamus{privacy budget}{anggaran privasi}{the total amount of $\varepsilon$ that a system
may spend; every query or training step spends some of it, and once it is exhausted there
is no guarantee left}

\section{Why amplitudes use complex numbers}
\tingkat{basic}

Now the second half, starting from a question that is often taken as technical although
its answer is conceptual. Why is a quantum state represented by complex numbers, and not
by real numbers?

The short answer: because two paths to the same result have to be able to cancel each
other, and cancellation of that kind demands a sign that can rotate, not merely positive
and negative. With real numbers an amplitude has only two signs, and interference can
only add or subtract. With complex numbers an amplitude has a phase that rotates
continuously, so two paths can meet at any angle and the result varies smoothly.

The complex number $z=a+bi$ with $i^2=-1$ has modulus $|z|=\sqrt{a^2+b^2}$, conjugate
$\bar z=a-bi$, and phase $\theta$ in the notation
$z=|z|e^{i\theta}=|z|(\cos\theta+i\sin\theta)$. Modulus and phase play different roles and
cannot be exchanged: the modulus determines the size of the probability, the phase
determines the sign of the interference. Multiplying two complex numbers multiplies the
moduli and adds the phases, and that property makes the phase useful as a coordinate of
the state.

\kamus{probability amplitude}{amplitudo}{a complex number whose modulus squared is the
probability; different from a probability because amplitudes can cancel one another}

The state space of one qubit is $\mathbb{C}^2$, and the state space of $n$ qubits is
$\mathbb{C}^{2^n}$ with the inner product
$\langle\phi|\psi\rangle=\sum_k\overline{\phi_k}\psi_k$. In bra-ket notation,
$|\psi\rangle$ is a column vector and $\langle\psi|$ is its conjugate transpose. A state
is always normalized, $\langle\psi|\psi\rangle=1$, because the sum of all probabilities
has to be one.

\begin{definisi}[Finite-dimensional Hilbert space]
\index{Hilbert space}
A finite-dimensional Hilbert space is a finite-dimensional complex vector space equipped
with a complete inner product. In this book the only example that is used is
$\mathbb{C}^{d}$ with the standard inner product, usually $d=2^{n}$.
\end{definisi}

Two properties of this notation have to be held from the start. First, the global phase
has no physical meaning: $|\psi\rangle$ and $e^{i\theta}|\psi\rangle$ give the same
probability for every measurement outcome, because $|e^{i\theta}a_k|^2=|a_k|^2$. Second,
the \emph{relative} phase has full meaning: changing the phase of one component alone
changes the result of interference, and that is what quantum algorithms use. The
difference between the invisible global phase and the visible relative phase becomes
clear in the worked example in section nine.

\begin{proposisi}
If $|\psi\rangle$ and $|\phi\rangle$ are both of norm one, then
$|\langle\phi|\psi\rangle|\le 1$.
\end{proposisi}

\begin{proof}
The Cauchy-Schwarz inequality gives
$|\langle\phi|\psi\rangle|^2\le\langle\phi|\phi\rangle\langle\psi|\psi\rangle=1$, then
take the square root.
\end{proof}

\section{Gates as unitary matrices, and why they must be unitary}
\tingkat{intermediate}

A change of state due to a gate is a matrix multiplication,
$|\psi'\rangle=U|\psi\rangle$. Not every matrix may be used. The condition is one: the
total probability has to stay one for every initial state, so
$\langle U\psi|U\psi\rangle=\langle\psi|\psi\rangle$ for all $\psi$.

\begin{definisi}[Unitary matrix]
\index{unitary gate}
A square matrix $U$ is called unitary when $U^{\dagger}U=UU^{\dagger}=I$, with
$U^{\dagger}$ the conjugate transpose of $U$.
\end{definisi}

\begin{teorema}
For a square matrix $U$ over $\mathbb{C}^{d}$, the following statements are equivalent:
(i) $U$ is unitary; (ii) $\langle U\phi|U\psi\rangle=\langle\phi|\psi\rangle$ for all
$\phi,\psi$; (iii) $\|U\psi\|=\|\psi\|$ for all $\psi$.
\end{teorema}

\begin{proof}
(i) $\Rightarrow$ (ii): $\langle U\phi|U\psi\rangle=\phi^{\dagger}U^{\dagger}U\psi
=\phi^{\dagger}\psi$. (ii) $\Rightarrow$ (iii): take $\phi=\psi$. (iii)
$\Rightarrow$ (i): from $\langle U\psi|U\psi\rangle=\langle\psi|\psi\rangle$ for all
$\psi$, use the polarization identity to conclude
$\langle\phi|U^{\dagger}U\psi\rangle=\langle\phi|\psi\rangle$ for all $\phi,\psi$;
therefore $U^{\dagger}U=I$, and for a square matrix in finite dimension a left inverse
means a right inverse.
\end{proof}

That theorem explains three things. First, a gate must be unitary because probability must
be conserved; no other reason is needed. Second, every gate has an inverse $U^{\dagger}$
that is also unitary, so a circuit can be reversed and no information is lost while the
gates act. Third, the only non-unitary operation is measurement, and that is where
information is lost. That is why quantum computation is not computation that erases
information, but computation that postpones erasure until the end.

\begin{proposisi}[No cloning]
There is no unitary $U$ on two qubits such that
$U\bigl(|\psi\rangle\otimes|0\rangle\bigr)=|\psi\rangle\otimes|\psi\rangle$ for all
$|\psi\rangle$.
\end{proposisi}

\begin{proof}
Suppose there is one, and take $|\psi\rangle,|\phi\rangle$ of norm one. Because $U$ is
unitary, the inner products of the two sides are equal:
$\langle\phi|\psi\rangle=\langle\phi|\psi\rangle^2$. The equation $x=x^2$ has only the
solutions $0$ and $1$, so every pair of states has to be orthogonal or identical; that is
impossible for more than two states.
\end{proof}

\kamus{unitary}{uniter}{a property of a matrix that preserves the length of a vector; in
everyday language, a gate that neither adds nor removes total probability}

\section{Fourier as a change of basis, and its connection to convolution}
\tingkat{advanced}

The Fourier transform is a change of basis, and that is the most useful way to see it
here. In a space of dimension $N$, the standard basis is replaced by the basis
$|f_k\rangle$ whose components are $(|f_k\rangle)_j=\omega^{jk}/\sqrt{N}$ with
$\omega=e^{2\pi i/N}$. The matrix is $F_{jk}=\omega^{jk}/\sqrt{N}$.

\begin{proposisi}
The Fourier matrix $F$ is unitary.
\end{proposisi}

\begin{proof}
Compute $(F^{\dagger}F)_{jk}=\frac{1}{N}\sum_{l=0}^{N-1}\omega^{-lj}\omega^{lk}
=\frac{1}{N}\sum_{l=0}^{N-1}\omega^{l(k-j)}$. If $k=j$ every term has the value one, so
the result is $1$. If $k\neq j$, this is a geometric series with ratio
$\omega^{k-j}\neq1$ and its sum is $\frac{1-\omega^{N(k-j)}}{1-\omega^{k-j}}=0$ because
$\omega^{N}=1$. Hence $F^{\dagger}F=I$.
\end{proof}

Because it is unitary, $F$ may legitimately be used as a gate, and its $n$-qubit version
is called the quantum Fourier transform. What makes it interesting is not only that it is
a legitimate gate, but that it turns convolution into pointwise multiplication.

\begin{teorema}[Convolution becomes multiplication]
\index{convolution}
For $x,y\in\mathbb{C}^{N}$, define the cyclic convolution
$(x*y)_k=\sum_{j=0}^{N-1}x_j\,y_{(k-j)\bmod N}$. Then
\begin{equation}
F(x*y) \;=\; \sqrt{N}\,\bigl(Fx\bigr)\odot\bigl(Fy\bigr),
\end{equation}
with $\odot$ the componentwise product.
\end{teorema}

\begin{proof}
The $k$-th component of the left-hand side is
$\frac{1}{\sqrt{N}}\sum_m\omega^{km}\sum_j x_j y_{m-j}$. Exchange the order of summation
and write $r=m-j$:
\[
\frac{1}{\sqrt{N}}\sum_j x_j\omega^{kj}\sum_r\omega^{kr}y_r
=\sqrt{N}\,(Fx)_k\,(Fy)_k .
\]
The exchange of the order of summation is legitimate because both sums are finite.
\end{proof}

That exchange of two finite sums looks trivial, but the infinite version of the same step
is the content of Chapter 14: a Fourier series may be differentiated term by term only
when its convergence is uniform enough. The connection with Chapter 15 is more
conceptual: a change of basis is a change of variable, and a unitary matrix is a change of
variable whose Jacobian determinant has modulus one, so volumes and probabilities do not
change.

The connection with quantum computation lies in its cost. A discrete Fourier transform
over $N=2^{n}$ points is computed classically in $\mathcal{O}(N\log N)$ operations with
the FFT. Its quantum version is built from $\mathcal{O}(n^{2})$ gates. The difference is
real and is the basis of Shor's algorithm. But there is a condition that is often left
out: once the transform is finished, we cannot read the $N$ amplitudes, and all we can do
is take samples. Speed in preparing a state does not automatically become speed on the
task one wants to finish; it has to be paired with a procedure that extracts the answer
from only a few samples. Many claims of quantum advantage in machine learning have no
procedure of that kind.

\section{The Born rule, and why the output is always probabilistic}
\tingkat{intermediate}

Up to here everything is deterministic: given a state and a gate, the next state is
unique. Probability enters at one step that cannot be avoided, namely measurement.

\begin{definisi}[Born rule]
\index{Born rule}
If the state $|\psi\rangle=\sum_k a_k|k\rangle$ is measured in the orthonormal basis
$\{|k\rangle\}$, the probability of obtaining outcome $k$ is
\begin{equation}
p(k) \;=\; \bigl|\langle k|\psi\rangle\bigr|^{2} \;=\; |a_k|^{2},
\end{equation}
and after the outcome $k$ is obtained, the state becomes $|k\rangle$.
\end{definisi}

Why must it be so, and not merely our ignorance of the state that is ``really'' there?
The strongest answer is not philosophical but experimental: Bell's inequality, tested by
Aspect, Grangier, and Roger in 1982, shows that the correlations between two separated
measurements violate a bound that every theory storing hidden definite values has to obey.
So probabilism here is not a way of saying ``we do not know yet''; it is the structure of
the theory. The practical consequence: one measurement gives one outcome, not the whole
state.

\kamus{Born rule}{aturan Born}{the rule that turns an amplitude into a probability by
squaring its modulus; the only place where probability enters the theory}

The consequence for machine learning is often underestimated. To estimate $p(k)$ with
standard error $\gamma$ one needs about $p(1-p)/\gamma^{2}$ shots. In our experiment, the
state $R_y(60^\circ)|0\rangle$ has true probabilities $0.75$ and $0.25$; with a thousand
shots the recorded counts were $727$ and $273$, the estimate $0.727$, differing by
$0.023$ from the true value. Its theoretical standard error is
$\sqrt{0.75\cdot 0.25/1000}=0.0137$, so that difference is about $1.7$ standard errors, an
ordinary event. Three digits after the decimal point need about a million shots, and every
shot means one state preparation and one measurement. This is where the factor $\sqrt{N}$
appears, and this is where many ``quadratic speeds'' turn into ordinary speeds once the
cost of readout is counted.

\section{Worked example: one and two qubits by hand}
\tingkat{intermediate}

\begin{untukanak}
Two children toss coins. The first child tosses one coin, then closes his hand and asks:
heads or tails? The result is one, and nobody can know anything before that coin is
opened. The second child prepares two coins in a strange way: he tells the coins to
``consult'' first, so that the two always agree, without ever agreeing in advance whether
they will both be heads or both be tails. When the first coin is opened and shows heads,
the second coin is certainly heads. When the first coin shows tails, the second coin is
certainly tails.

What makes it strange: there is no way to make two ordinary coins behave like that. If
both were fixed from the start, whoever made the coins would know both. In coins that
consult there is no record of that kind, and real experiments show that no hidden record
explains the outcome.

That last example is what is called a Bell state, and this chapter uses it to show why
``zero'' and ``very small'' are two different things.
\end{untukanak}

All the numbers in this section are rechecked with NumPy in section eleven, and the reader
can check them with an ordinary calculator.

The Hadamard gate on one qubit is
$H=\frac{1}{\sqrt{2}}\begin{pmatrix}1&1\\1&-1\end{pmatrix}$. On $|0\rangle=(1,0)$ it gives
\[
H|0\rangle=\frac{1}{\sqrt{2}}(1,1)=(0.7071;\;0.7071),
\]
so the probabilities of the outcomes $0$ and $1$ are $0.5$ each.

Now the circuit $H$, then $Z=\operatorname{diag}(1,-1)$, then $H$ again. The first step
gives $(0.7071;0.7071)$; the second step flips the sign of the second component to
$(0.7071;-0.7071)$; the third step multiplies by $H$:
\[
H\begin{pmatrix}0.7071\\-0.7071\end{pmatrix}
=\frac{1}{2}\begin{pmatrix}0.7071-0.7071\\0.7071+0.7071\end{pmatrix}
=\begin{pmatrix}0\\1\end{pmatrix}.
\]
The result is $|1\rangle$ with probability one. The amplitude of the $|0\rangle$ component
becomes zero not because something was thrown away, but because two paths meet with
opposite signs. This is interference, and this is why the relative phase matters: if $Z$
is replaced by the identity, the result is $|0\rangle$ again.

The third example uses a rotation. The gate
$R_y(\theta)=\begin{pmatrix}\cos(\theta/2)&-\sin(\theta/2)\\
\sin(\theta/2)&\cos(\theta/2)\end{pmatrix}$ rotates the state in the real plane. With
$\theta=60^\circ$ and initial state $|0\rangle$,
\[
R_y(60^\circ)|0\rangle=(\cos 30^\circ;\;\sin 30^\circ)=(0.8660;\;0.5),
\]
so the probabilities are $\cos^{2}30^\circ=0.75$ and $\sin^{2}30^\circ=0.25$. Note that
$|a_0|^2+|a_1|^2=0.75+0.25=1$.

The fourth example is two qubits. Start from $|00\rangle$, apply $H$ to the first qubit,
then CNOT with the first qubit as control and the second as target. After $H$ the state is
$\frac{1}{\sqrt{2}}(|00\rangle+|10\rangle)$. CNOT flips the second qubit only when the
first qubit is $1$, so
\[
|\Phi^{+}\rangle=\frac{1}{\sqrt{2}}\bigl(|00\rangle+|11\rangle\bigr)
=(0.7071;\;0;\;0;\;0.7071).
\]
The probabilities of $00$ and $11$ are $0.5$ each; the probabilities of $01$ and $10$ are
zero. The two amplitudes in the middle are zero because the two paths that lead there
cancel one another, not because both are impossible. Consequently, if both qubits are
measured, the outcomes always agree: $00$ or $11$, never a mixture.

\kertas{Check the examples above without looking at the text. (a) Compute $HZH|0\rangle$
by writing out every matrix multiplication, then explain why the result is $|1\rangle$ and
not $|0\rangle$. (b) Replace $R_y(60^\circ)$ with $R_y(90^\circ)$, compute both
probabilities, then explain why two different states can give the same probabilities.
(c) For the Bell state above, compute the probability of obtaining $01$ and explain why
that number is zero, and not ``very small''.}

\section{Why classical simulation explodes, and what that means for claims of advantage}
\tingkat{advanced}

The state of $n$ qubits is a vector in $\mathbb{C}^{2^{n}}$, so it has $2^{n}$ complex
amplitudes. For $n=10$ that is $1\,024$ numbers, no larger than a small image. For $n=50$
it is $2^{50}=1.1259\times10^{15}$ numbers; with $16$ bytes per amplitude, what is needed
is $2^{54}$ bytes, that is exactly $16$ pebibytes. It is a limit that doubles every time
one qubit is added.

The experiment in section eleven measures the consequence. Preparing a state with one
Hadamard gate on every qubit takes $0.2$ milliseconds for $n=10$, $2.3$ for $n=14$, $81.4$
for $n=18$, and $1\,357$ milliseconds for $n=22$. That time column is the noisiest in the
whole chapter: two runs on the same machine can differ by almost a factor of two. What
does not change is the pattern, namely a straight line on a logarithmic scale against
$2^{n}$.

What may and may not be concluded? What may: a general quantum state cannot be stored
classically for large $n$, so an algorithm that genuinely uses all $2^{n}$ amplitudes
cannot be imitated directly. What may not: that every quantum circuit is faster. There are
three reasons.

The first reason, the state cannot be read. One measurement gives one outcome, whereas a
classical simulation that stores the whole vector can compute the probability of every
outcome. So more information is available in the simulation than in the hardware, and the
inability to copy a state is not evidence of advantage.

The second reason, the circuits that can be simulated include many circuits that are used
in practice: small depth, limited entanglement, product states, and Clifford gates have
efficient classical imitations. Goguen, Clarke, Martens, and Selman showed that family of
circuits can be computed classically in polynomial time, and that family is wide enough to
include many variational ansatzes.

The third reason, advantage can be removed by data structure. Ewin Tang showed in 2019
that a quantum recommendation algorithm that was thought to give an exponential speedup
has a classical counterpart with a polynomial speedup, because the structure of its data
matrix can be exploited without quantum mechanics. Huang, Kueng, and Preskill (2021) give
a strict separation, but only for a particular family of data distributions.

\kamus{barren plateau}{plateau tandus}{a region of the loss landscape of a variational
circuit where the gradients shrink exponentially in the number of qubits, so that
training does not move}

An honest statement about QML in 2026 therefore takes this form. There is evidence of
separation for particular tasks in oracle models, and there are theoretical claims of
speedup for learning tasks under particular distributional assumptions. There is no
established evidence of practical advantage for machine learning tasks on classical data,
in the sense of an experiment repeated on real hardware that beats the best classical
method at a meaningful data size. Cerezo and colleagues (2021, 2022) call the
\emph{barren plateau}, namely gradients that vanish exponentially in deep circuits, an
unresolved obstacle rather than a technical detail.

\section{Computational lab: privacy budget, accuracy, and the size of a state}
\tingkat{advanced}

The following program runs the two experiments above with no library other than NumPy:
DP-SGD on two-class logistic regression with five seeds per noise level, then several
quantum states and the cost of simulating them. All the numbers in this chapter come from
the output below, except those marked as hand computations.

\begin{lab}
\begin{verbatim}
import numpy as np
import time

D, N, T, C, DELTA, LR = 8, 200, 200, 1.0, 1e-5, 0.5


def data(rng, n, mu=0.6):
    y = 2 * rng.integers(0, 2, n) - 1
    return rng.normal(mu * y[:, None], 1.0, (n, D)), y


def sig(t):
    return 1.0 / (1.0 + np.exp(-t))


def clip(G):
    nrm = np.linalg.norm(G, axis=1, keepdims=True)
    return G * np.minimum(1.0, C / np.maximum(nrm, 1e-9))


def dp_sgd(X, y, sigma, seed):
    rng = np.random.default_rng(seed)
    w = np.zeros(D)
    for _ in range(T):
        G = -(1.0 - sig(y * (X @ w)))[:, None] * (y[:, None] * X)
        w -= LR * (clip(G).mean(0) + rng.normal(0.0, sigma * C / len(X), D))
    return w


def eps_basic(s):
    return T * np.sqrt(2 * np.log(1.25 * T / DELTA)) / s


def eps_adv(s):
    d = DELTA / 2
    e = np.sqrt(2 * np.log(1.25 * T / d)) / s
    return np.sqrt(2 * T * np.log(1 / d)) * e + T * e * (np.exp(e) - 1)


rng = np.random.default_rng(11)
Xn, yn = data(rng, N)
Xe, ye = data(rng, 3000)
print("sigma  accuracy     eps_basic  eps_adv")
A0 = [100 * np.mean(np.sign(Xe @ dp_sgd(Xn, yn, 1e-9, k)) == ye) for k in range(5)]
print("%5s  %5.2f+-%4.2f  %8s  %9s" % ("0", np.mean(A0), np.std(A0), "-", "-"))
for s in (2.0, 8.0, 32.0, 128.0, 512.0):
    A = [100 * np.mean(np.sign(Xe @ dp_sgd(Xn, yn, s, k)) == ye) for k in range(5)]
    print("%5.0f  %5.2f+-%4.2f  %8.1f  %9.1f"
          % (s, np.mean(A), np.std(A), eps_basic(s), eps_adv(s)))

I2 = np.eye(2)
Z = np.diag([1.0, -1.0])
S = np.diag([1.0, 1.0j])
H = np.array([[1.0, 1.0], [1.0, -1.0]]) / np.sqrt(2)
CNOT = np.array([[1, 0, 0, 0], [0, 1, 0, 0], [0, 0, 0, 1], [0, 0, 1, 0]], float)


def ry(t):
    c, s = np.cos(t / 2), np.sin(t / 2)
    return np.array([[c, -s], [s, c]])


def kron1(U, n, q):
    out = np.array([[1.0]])
    for j in range(n):
        out = np.kron(out, U if j == q else I2)
    return out


s0 = np.array([1.0, 0.0])
print()
print("H|0>        ", np.round(H @ s0, 4))
print("HZH|0>      ", np.round(H @ Z @ (H @ s0), 4))
print("HSH|0>      ", np.round(H @ S @ (H @ s0), 4))
print("Ry(60)|0>   ", np.round(ry(np.pi / 3) @ s0, 4),
      " probability", np.round((ry(np.pi / 3) @ s0) ** 2, 4))
psi = CNOT @ (kron1(H, 2, 0) @ np.kron(s0, s0))
print("Bell        ", np.round(psi, 4), " probability", np.round(psi ** 2, 4))
p1 = (ry(np.pi / 3) @ s0) ** 2
sh = np.bincount(np.random.default_rng(4).choice(2, 1000, p=p1), minlength=2)
print("1000 shots", sh, " estimate", np.round(sh / 1000, 3))
print()
print("  n   amplitudes    memory (GB)   time (ms)")
for n in (10, 14, 18, 22):
    best = 1e9
    for _ in range(3):
        v = np.zeros(2 ** n, dtype=complex)
        v[0] = 1.0
        t0 = time.perf_counter()
        for q in range(n):
            v = v.reshape([2] * n)
            v = np.moveaxis(np.tensordot(H, v, axes=([1], [q])), 0, q).reshape(-1)
        best = min(best, 1000 * (time.perf_counter() - t0))
    print("%3d  %11.3e  %12.3e  %10.1f" % (n, 2.0 ** n, 2.0 ** n * 16 / 2 ** 30, best))
\end{verbatim}

The output of that program, as it stands:

\begin{verbatim}
sigma  accuracy     eps_basic  eps_adv
    0  95.20+-0.00         -          -
    2  95.13+-0.05     583.7    11302.8
    8  95.03+-0.13     145.9      216.5
   32  93.71+-0.81      36.5       20.6
  128  85.51+-6.16       9.1        3.7
  512  70.39+-17.26       2.3        0.8

H|0>         [0.7071 0.7071]
HZH|0>       [0. 1.]
HSH|0>       [0.5+0.5j 0.5-0.5j]
Ry(60)|0>    [0.866 0.5  ]  probability [0.75 0.25]
Bell         [0.7071 0.     0.     0.7071]  probability [0.5 0.  0.  0.5]
1000 shots [727 273]  estimate [0.727 0.273]

  n   amplitudes    memory (GB)   time (ms)
 10    1.024e+03     1.526e-05         0.2
 14    1.638e+04     2.441e-04         2.3
 18    2.621e+05     3.906e-03        81.4
 22    4.194e+06     6.250e-02      1356.6
\end{verbatim}
\end{lab}

The zero row is the baseline: $95.20$ percent with zero spread. The row $\sigma=2$ gives
$\varepsilon=583.7$, a guarantee that is practically empty because the ratio $e^{583}$
constrains nothing; the row $\sigma=8$ lowers $\varepsilon$ to $145.9$ with no measurable
loss of accuracy, so a little noise is almost free. The price only begins to be felt at
$\sigma=32$: $\varepsilon$ drops to $36.5$ or $20.6$ at a cost of $1.49$ points, then
$85.51$ percent at $\sigma=128$, and $70.39$ percent at $\sigma=512$ with a spread across
seeds of $17.26$ points, so that a single model no longer represents the procedure.

The second block verifies the hand computations of section nine. $HZH|0\rangle$ gives
$(0;1)$, $R_y(60^\circ)$ gives $(0.8660;0.5)$, and the Bell state appears as
$(0.7071;0;0;0.7071)$ with two exact zeros. The row $HSH|0\rangle$ gives
$(0.5+0.5i;\;0.5-0.5i)$, and that pair is used in Figure~\ref{fig:fase}.

The third block measures the cost of simulation. Memory grows by exactly a factor of two
per qubit, from $1.5\times10^{-5}$ GB at $n=10$ to $6.25\times10^{-2}$ GB at $n=22$. Time
grows faster than memory over this range for architectural reasons, not mathematical ones;
on a machine with a larger cache the jump shifts to a higher $n$, but the slope stays.

\begin{figure}
\centering
\begin{tikzpicture}
\begin{axis}[
  width=0.86\linewidth, height=6.0cm,
  xlabel={$\varepsilon$ (basic composition, log scale)},
  ylabel={test accuracy (\%)},
  xmode=log, log basis x=10,
  xmin=1, xmax=1200, ymin=62, ymax=100,
  xtick={1,10,100,1000},
  ytick={65,70,75,80,85,90,95,100},
  grid=major, grid style={black!12},
  tick label style={font=\footnotesize},
  label style={font=\small},
  legend style={font=\footnotesize, at={(0.03,0.03)}, anchor=south west},
]
\addplot[only marks, mark=*, mark size=2.2pt, black] coordinates {
  (2.3,70.39) (9.1,85.51) (36.5,93.71) (145.9,95.03) (583.7,95.13)
};
\addlegendentry{DP-SGD, mean of 5 seeds}
\addplot[black, dashed, domain=1:1200, samples=2] {95.20};
\addlegendentry{no noise: $95.20$}
\node[font=\scriptsize, anchor=south west] at (axis cs:2.3,70.39) {$\sigma{=}512$};
\node[font=\scriptsize, anchor=south west] at (axis cs:9.1,85.51) {$\sigma{=}128$};
\node[font=\scriptsize, anchor=south west] at (axis cs:36.5,93.71) {$\sigma{=}32$};
\node[font=\scriptsize, anchor=north east] at (axis cs:145.9,95.03) {$\sigma{=}8$};
\end{axis}
\end{tikzpicture}
\caption{Test accuracy against the privacy budget in the Computational lab experiment
($n=200$, $d=8$, $T=200$, $\delta=10^{-5}$, five seeds). Every point is the result of a
real execution, paired with the \texttt{eps\_basic} column in the output. The dashed line
is the accuracy without noise. The horizontal axis is logarithmic: lowering
$\varepsilon$ from $36.5$ to $2.3$ costs $23$ points of accuracy, whereas going from
$583.7$ to $145.9$ is almost free.}
\label{fig:tradeoff}
\end{figure}

\begin{figure}
\centering
\begin{tikzpicture}[scale=1.15]
\draw[->, black!60] (-1.25,0) -- (1.45,0) node[right, font=\footnotesize]
  {$\operatorname{Re}$};
\draw[->, black!60] (0,-1.1) -- (0,1.25) node[above, font=\footnotesize]
  {$\operatorname{Im}$};
\draw[black!35, dashed] (0,0) circle (0.7071);
\draw[black!55, thin] (0.7071,0) -- (0.7071,0.7071) -- (0,0.7071);
\draw[black!55, thin] (0.5,0) -- (0.5,0.5) -- (0,0.5);
\draw[->, very thick, black] (0,0) -- (0.7071,0.7071)
  node[above left, font=\footnotesize] {$a_0$};
\draw[->, very thick, black] (0,0) -- (-0.7071,0.7071);
\draw[->, thick, black!55] (0,0) -- (0.5,0.5);
\draw[->, thick, black!55] (0,0) -- (0.5,-0.5) node[below right, font=\footnotesize]
  {$a_1$};
\node[font=\scriptsize, anchor=west] at (0.75,0.16) {$\tfrac{1}{\sqrt2}=0.7071$};
\node[font=\scriptsize, anchor=north] at (0.5,-0.55) {$0.5-0.5i$};
\node[font=\scriptsize, anchor=south] at (0.5,0.56) {$0.5+0.5i$};
\draw[black!35] (0.09,0) arc (0:45:0.42);
\node[font=\scriptsize] at (0.30,0.16) {$45^\circ$};
\end{tikzpicture}
\caption{Amplitudes as arrows in the complex plane. The thick arrows are the state
$H|0\rangle$ with $a_0=a_1=1/\sqrt2$; the thin arrows are $HSH|0\rangle$ with
$S=\operatorname{diag}(1,i)$, namely $0.5+0.5i$ and $0.5-0.5i$. The modulus stays
$0.7071$, so the probability of both outcomes stays $1/2$, but the phase shifts by
$45^\circ$ in opposite directions. A measurement does not distinguish the two states; the
next gate can. The arrow coordinates are taken from the output of the Computational lab,
not drawn as an illustration.}
\label{fig:fase}
\end{figure}

Figure~\ref{fig:tradeoff} is the quantitative shape of the privacy-utility trade-off: the
curve is not linear, and the most expensive part is at the left end, in the region that is
most often praised in abstracts. Figure~\ref{fig:fase} is the geometric shape of the
sentence ``the modulus determines the probability, the phase determines the interference''.

\begin{untukdosen}
The teaching order I recommend: start from Figure~\ref{fig:tradeoff}, not from the
definition. Show that the curve has an expensive left end, then ask why. Students usually
guess ``because the noise is large'', which is true but not yet enough; what has to come
out is that the privacy budget is summed along the steps, so shrinking $\varepsilon$ per
step is paid twice, in noise and in the number of steps.

Two misconceptions almost always appear. First, students read ``$(\varepsilon,\delta)$-DP''
as ``this model is safe''. What is true is that it is a promise about the influence of one
record on the output distribution, and the model can still be completely wrong. Second,
students think the quantum phase is a decoration of notation because ``probability depends
only on the modulus''. The example $HZH|0\rangle$ breaks that in three lines.

A blackboard activity that worked in our class: ask two students to compute
$HZH|0\rangle$ and $HHH|0\rangle$ separately, then compare. Both use three gates, they
differ in only one matrix in the middle, and the results are $|1\rangle$ against
$|0\rangle$. The class is ready to go on when a student can explain why $HSH|0\rangle$
gives the same probabilities as $H|0\rangle$ but a different state.
\end{untukdosen}

\section{Summary}

\begin{ringkasan}
\begin{itemize}
  \item A model trained on personal data leaks information through three channels:
        membership, attributes, and reconstruction, and the third works from gradients.
  \item $(\varepsilon,\delta)$ differential privacy compares two output distributions for
        data that differ in one person, with the quantifier ``for every output set''. That
        quantifier is what makes the promise hold against an attacker, not merely against
        an average.
  \item The Laplace mechanism with scale $b=\Delta f/\varepsilon$ satisfies
        $(\varepsilon,0)$; its proof only compares the ratio of two densities. The Gaussian
        mechanism uses $\sigma\ge\sqrt{2\ln(1.25/\delta)}/\varepsilon$ and gives
        $\delta>0$.
  \item Basic composition adds up $\varepsilon$; advanced composition gives the better
        $\sqrt{T}$ term only when $\varepsilon$ per step is small: at $\sigma=32$ the
        results are $36.5$ versus $20.6$, whereas at $\sigma=2$ they are $583.7$ versus
        $11\,302.8$.
  \item DP-SGD clips per-example gradients at norm $C$ and then adds noise $\sigma C/B$.
        Accuracy falls from $95.20$ percent without noise to $93.71$ percent at
        $\varepsilon=36.5$ and $70.39$ percent at $\varepsilon=2.3$.
  \item Differential privacy guarantees something about individuals, not the truth of the
        model; it does not remove bias, and its utility cost can be uneven across groups.
  \item A quantum state lives in $\mathbb{C}^{2^{n}}$; a gate must be unitary because
        probability must be conserved, and measurement is the only non-unitary step. The
        output is therefore always a sample.
  \item Classical simulation needs $2^{n}$ amplitudes: $n=50$ means $16$ pebibytes. That
        fact limits imitation, but it does not prove advantage, and so far there is no
        established evidence of practical advantage for QML.
\end{itemize}
\end{ringkasan}

\section{Exercises}
\tingkat{intermediate}

\begin{latihan}
\begin{enumerate}
  \item Assess this claim: ``Our model satisfies differential privacy with
        $\varepsilon=0.5$, so no personal information can leak from it.'' State true,
        false, or not enough evidence, and name two things that must hold for the claim to
        apply.
  \item Prove that the Laplace mechanism with scale $b=\Delta f/\varepsilon$ satisfies
        $(\varepsilon,0)$ when sensitivity is measured with the $\ell_1$ norm on an output
        of dimension $d>1$. Show where the one-dimensional argument has to be replaced.
  \item Prove parallel composition: if $M_1$ works on $D_1$ and $M_2$ on $D_2$ with
        $D_1\cap D_2=\varnothing$, and each is $(\varepsilon_i,\delta_i)$, then their
        output satisfies $(\max_i\varepsilon_i,\max_i\delta_i)$.
  \item Compute $\varepsilon$ for randomized response when the probability of answering
        honestly is $p$ and the rest is answered at random with a fair coin. State
        $\varepsilon$ as a function of $p$, then determine the $p$ that gives
        $\varepsilon=2$.
  \item For a count query with $\varepsilon=0.5$ on $n=10\,000$, compute the probability
        that its absolute error exceeds $10$ people, then compare it with the same
        probability when the query is an average with $R=1$. Check both numbers by
        simulation.
  \item Compute $H\,S\,H|0\rangle$ with $S=\operatorname{diag}(1,i)$ step by step, then
        compare its probabilities with those of $H|0\rangle$. Explain why the measurement
        results are the same although the states differ, and which gate distinguishes them.
  \item For the Bell state $|\Phi^{+}\rangle$, compute the probability of the outcome $01$
        when the second qubit is measured first and the result is $0$; repeat for the
        result $1$. Explain why both answers show perfect correlation.
  \item Show that $\begin{pmatrix}1&0\\0&e^{i\varphi}\end{pmatrix}$ is unitary for every
        $\varphi$, then explain why that gate does not change probabilities in the
        computational basis but can change the result of the next gate.
  \item Compute the number of amplitudes and the memory requirement for $n=40$ and $n=60$
        qubits with $16$ bytes per amplitude. State them in sensible units, then compare
        them with the storage capacity that you have.
  \item A paper writes: ``Our circuit prepares a state with $2^{60}$ amplitudes, so it
        beats classical simulation for machine learning.'' Name three questions that must
        be answered before that sentence is accepted, and for each question state what
        kind of result would make you accept it.
\end{enumerate}

\noindent\textbf{Hints for the three hardest items.} Item 2: write the ratio of densities
as a product of per-coordinate ratios, then use $\sum_j|a_j-a'_j|\le\Delta f$. Item 3: at
every step only one of the two mechanisms can change. Item 10: whether all the amplitudes
are used, whether the output can be extracted with a number of measurements that is not
exponential, and whether there is a classical comparison on the same task and data.
\end{latihan}

\begin{untukpeneliti}
Three things in the first half are often quoted without their conditions. First, the
$(\varepsilon,\delta)$ definition here is the definition of Dwork, McSherry, Nissim, and
Smith (2006) in the form of the book by Dwork and Roth; the condition $\delta\ll 1/n$
comes from the camouflage argument for one individual, not from analytical convenience.
Second, the Gaussian bound $\sigma\ge\sqrt{2\ln(1.25/\delta)}/\varepsilon$ is loose for
small $\sigma$; the tight bound is given by the analytic Gaussian mechanism of Balle and
Wang (2018), and the difference can reach a factor of two at $\varepsilon<1$. Third, the
moments \emph{accountant} of Abadi and colleagues (2016) uses subsampling, so an
$\varepsilon$ that is reported without the subsampling rate and the number of steps cannot
be compared with other numbers.

In the second half, what is really proved is a conditional statement: complexity
separations for particular families of tasks, and the impossibility of classical
simulation for general states. What is popular but unproved is the claim that variational
circuits give an advantage on classical data. Tang (2019) and Huang, Kueng, Preskill
(2021) limit that claim from two directions: data structure can remove the speedup, and
the separations that exist depend on distributional assumptions. McClean and colleagues
(2018) show gradients that vanish exponentially in deep circuits, so variational training
has an unresolved obstacle.

The most important question is not ``is there a speedup'', but ``what kind of speedup
survives once the cost of readout, state preparation, and error correction is counted''.
Until there is an answer measured on hardware with real data, every claim of practical
advantage for QML should be read as a hypothesis.
\end{untukpeneliti}

\peta{The privacy half follows Dwork and Roth, \emph{The Algorithmic Foundations of
Differential Privacy} (chapters 2--3 and appendix A), Abadi et al., \emph{Deep Learning
with Differential Privacy} (2016) for DP-SGD, Shokri et al. (2017) for membership
inference, Dinur and Nissim (2003) for reconstruction from aggregate queries, and Zhu,
Liu, Han (2019) for reconstruction from gradients. The quantum half follows Nielsen and
Chuang, \emph{Quantum Computation and Quantum Information} (chapters 2 and 4), Cerezo et
al. (2021, 2022) for the limits of variational algorithms, Tang (2019) for
dequantization, and Huang, Kueng, Preskill (2021) for data-based separation. The
prerequisites are in Chapter 10 (SVD and orthogonality), Chapter 16 (conditional
probability), and Chapter 17 (random variables), whereas the connection to Fourier and
convolution is in Chapter 14 and Chapter 15.}


% ---------------------------- appendices ----------------------------
\appendix
\appendix
\renewcommand{\appendixname}{Appendix}
\chapter{Table of Notation}
\label{app:notasi}

The following notation is used throughout the book. When a chapter uses a symbol that is
not listed here, that symbol is introduced where it is used. The number at the end of each
entry refers to the chapter that uses it.

\section*{Sets, logic, and numbers (Chapters 1--5, 22)}
\begin{description}[style=nextline,itemsep=0.25em]
  \item[$\mathbb{N}$, $\mathbb{Z}$, $\mathbb{Q}$, $\mathbb{R}$, $\mathbb{C}$]
        the natural, integer, rational, real, and complex numbers.
  \item[$\forall$, $\exists$, $\neg$] for all, there exists, and not.
  \item[$\Rightarrow$, $\Leftrightarrow$] implies, and if and only if.
  \item[$\in$, $\notin$, $\subseteq$, $\cup$, $\cap$, $\setminus$] membership,
        not an element of, subset, union, intersection, difference (Chapter 2).
  \item[$\mathcal{P}(S)$, $|S|$, $S\times T$] power set, cardinality,
        Cartesian product (Chapter 2).
  \item[$\sup S$, $\inf S$] least upper bound and greatest lower bound (Chapter 3).
  \item[$a_n\to a$] the sequence $(a_n)$ converges to $a$ (Chapter 4).
  \item[$\sum_{n=1}^{\infty} a_n$] series and its sum (Chapter 5).
  \item[$\mathrm{fl}(x)$, $\mathrm{ulp}$] the rounded value of a real number and the
        smallest unit that can be represented (Chapter 22).
\end{description}

\section*{Linear algebra (Chapters 6--10)}
\begin{description}[style=nextline,itemsep=0.25em]
  \item[$\mathbf{x}$, $\mathbf{x}^{\top}$, $\|\mathbf{x}\|$] column vector,
        transpose, and Euclidean norm.
  \item[$\operatorname{span} S$, $\dim V$] span of a set of vectors and dimension of
        a vector space.
  \item[$\mathcal{L}(V,W)$] the space of all linear transformations from $V$ to $W$ (Chapter 7).
  \item[$\ker T$, $\operatorname{im} T$] kernel and image of a transformation (Chapter 7).
  \item[$\det A$, $\operatorname{tr} A$, $A^{-1}$] determinant, trace, and inverse of
        a matrix (Chapter 8).
  \item[$\lambda$, $\mathbf{v}$] eigenvalue and eigenvector (Chapter 9).
  \item[$\sigma_i$, $A=U\Sigma V^{\top}$] singular values and the singular value
        decomposition (Chapter 10).
  \item[$A^{+}$, $\kappa(A)$] pseudoinverse and condition number (Chapters 10, 20).
\end{description}

\section*{Analysis and calculus (Chapters 11--16)}
\begin{description}[style=nextline,itemsep=0.25em]
  \item[$\varepsilon$, $\delta$] small positive numbers in the definition of limit,
        continuity, and uniform continuity.
  \item[$f'(x)$, $\nabla f$] ordinary derivative and gradient.
  \item[$J_f$, $H_f$] Jacobian and Hessian matrices (Chapter 12).
  \item[$\int_a^b f$] Riemann integral on $[a,b]$ (Chapter 13).
  \item[$\sum_{\mathcal{P}} M_i \Delta x_i$] upper sum for the partition
        $\mathcal{P}$ (Chapter 13).
  \item[$f_n \rightrightarrows f$] sequence of functions converging uniformly (Chapter 14).
  \item[$\dfrac{\partial(x,y)}{\partial(u,v)}$] Jacobian determinant of a change of
        variables (Chapter 15).
  \item[$x'=f(t,x)$, $\dot{x}=F(x)$] ordinary differential equation and autonomous
        system, where $'$ and $\dot{\ }$ both denote the derivative with respect to
        time (Chapter 16).
  \item[$\varphi_t(x_0)$] flow: the state at time $t$ when starting from $x_0$
        (Chapter 16).
  \item[$x^{*}$, $J_F(x^{*})$] equilibrium point and the Jacobian matrix at that point;
        the sign of the real part of its eigenvalues determines stability (Chapter 16).
\end{description}

\section*{Probability and statistics (Chapters 17--21, 26)}
\begin{description}[style=nextline,itemsep=0.25em]
  \item[$P(A)$, $P(A\mid B)$] probability of the event $A$ and conditional probability.
  \item[{$\mathbb{E}[X]$, $\operatorname{Var}X$}] expectation and variance.
  \item[$\operatorname{Cov}(X,Y)$] covariance of two random variables (Chapter 18).
  \item[$\bar{X}_n$, $\mathrm{se}$] sample mean and its standard error.
  \item[$\mathrm{CI}_{1-\alpha}$] confidence interval with level
        $1-\alpha$ (Chapter 20).
  \item[$\mathrm{TP}$, $\mathrm{FP}$, $\mathrm{FN}$, $\mathrm{TN}$] the four cells of
        a confusion matrix (Chapter 27).
  \item[$\hat{\beta}$, $\hat{y}=X\hat{\beta}$] least squares regression
        coefficients and predicted values (Chapter 21).
  \item[$e_i=y_i-\hat{y}_i$, $R^2$] residual and the proportion of variance explained
        by the model (Chapter 21).
  \item[$\mathrm{se}(\hat{\beta}_j)$, $\mathrm{VIF}_j$] standard error of a coefficient and
        variance inflation factor (Chapter 21).
\end{description}

\section*{Optimization, machine learning, and verification (Chapters 23--30)}
\begin{description}[style=nextline,itemsep=0.25em]
  \item[$\min f_0(x)$ with $f_i(x)\le 0$] standard form of an optimization problem (Chapter 23).
  \item[$L(x,\lambda,\nu)$] Lagrangian; $\lambda,\nu$ multipliers (Chapter 23).
  \item[$c^{\top}x$, $Ax\le b$, $x\in\mathbb{Z}^n$] LP form and integer program
        (Chapter 24).
  \item[$\sigma(\cdot)$, $\mathrm{softmax}$] activation function and exponential
        normalization (Chapter 25).
  \item[$\nabla_{\mathbf{x}}\ell$] gradient of the loss function with respect to the input;
        the fuel for attacks (Chapters 25, 29).
  \item[$L_2$, $L_1$, $\lambda_{\text{reg}}$] kinds of regularization and their strength
        (Chapter 26).
  \item[$\|\delta\|_{\infty}\le\varepsilon$] perturbation bound in adversarial
        attacks (Chapters 29, 30).
  \item[$\vdash$, \texttt{by}, \texttt{sorry}] provability, the tactic block opener,
        and deferred proof in Lean~4 (Chapter 28).
\end{description}

\chapter{Glossary}
\label{app:glosarium}

The terms below are used with fixed technical meanings throughout the book. Indonesian
equivalents are given in parentheses, because the Indonesian edition of this book uses those
terms and much of the local literature follows it.

\begin{description}[style=nextline,itemsep=0.35em]
  \item[Absolute convergence] a series that still converges when all its terms are forced positive (\emph{konvergensi mutlak}).
  \item[Absolute value] the distance of a number from zero on the number line, regardless of its sign (\emph{nilai mutlak}).
  \item[Accumulation point] a point every neighborhood of which contains a member of the set other than itself (\emph{titik akumulasi}).
  \item[Activation function] a nonlinear mapping on the output of a unit; it decides how much signal is passed on and the shape of its derivative (\emph{fungsi aktivasi}).
  \item[Adaptive attack] an attack designed with knowledge of the defense being tested, including knowledge that its gradient is manipulated (\emph{serangan adaptif}).
  \item[Adjoint] in the backward pass, the number that says how strongly the final loss is affected by one quantity in the middle of the network; one adjoint per quantity (\emph{adjoint}).
  \item[Adjugate] the cofactor matrix arranged in transpose; used to write the inverse as a single formula, $A^{-1}=\operatorname{adj}(A)/\det A$ (\emph{adjoin}).
  \item[Adversarial example] a small disturbance to the input, deliberately made to change the output of a model (\emph{serangan adversarial}).
  \item[Adversarial learning] the study of examples deliberately designed to fool a model, and of how to defend against them (\emph{pembelajaran adversarial}).
  \item[Antiderivative] a function whose derivative is the function being integrated; the way home from a derivative (\emph{antiturunan}).
  \item[Asymptotically stable] stays close after being disturbed, and then actually returns to the equilibrium point (\emph{stabil asimtotik}).
  \item[Axiom] a statement accepted without proof as a starting point; the whole content of a chapter rests on it (\emph{aksioma}).
  \item[Backpropagation] an algorithm that computes the gradient of every parameter by applying the chain rule from the output back to the input (\emph{turunan balik}).
  \item[Backward gradient] see backpropagation (\emph{gradien terbalik}).
  \item[Backward pass differentiable approximation] a stand-in for the gradient on a piece that is not differentiable, using the gradient of the function at the point just after it, so that an attack can still be computed (\emph{BPDA}).
  \item[Barren plateau] a region of the loss landscape of a variational circuit where the gradient shrinks exponentially with the number of qubits (\emph{plateau tandus}).
  \item[Basic feasible solution] a point obtained by setting some variables to zero and then solving for the rest, and still satisfying every constraint (\emph{solusi layak basis}).
  \item[Basis] a set of vectors that is linearly independent and spans the whole space (\emph{basis}).
  \item[Bias-variance tradeoff] the split of model error into error from being too simple and error from being too sensitive to the data (\emph{bias-variansi}).
  \item[Bijection] a perfect pairing of two sets: injective and surjective at once, with nothing left over on either side (\emph{bijeksi}).
  \item[Bootstrap] estimating the distribution of a statistic by resampling from the data itself, many times, with replacement (\emph{bootstrap}).
  \item[Born rule] the rule that turns an amplitude into a probability by squaring its modulus (\emph{aturan Born}).
  \item[Calibration] the agreement between the numbers a model promises and the actual frequency in the group promised the same (\emph{kalibrasi}).
  \item[Catastrophic cancellation] the loss of significant digits when two nearly equal numbers are subtracted, so that only error is left (\emph{pembatalan katastrofik}).
  \item[Cauchy sequence] a sequence whose terms crowd closer and closer together, without needing to name the value they approach (\emph{barisan Cauchy}).
  \item[Central limit theorem] the average of many independent random variables approaches the normal distribution once it is standardized (\emph{Teorema Limit Pusat}).
  \item[Chain rule] the rate of change of a composed function equals the product of the rates at each link in the chain (\emph{aturan rantai}).
  \item[Characteristic polynomial] $\det(A-\lambda I)$; its roots are the eigenvalues of $A$ (\emph{polinomial karakteristik}).
  \item[Coefficient of determination] the share of variation explained by the model, stated as a fraction of the whole; written $R^2$ (\emph{koefisien determinasi}).
  \item[Combination] a choice of several members without regard to their order; there are $\binom{n}{k}$ of them (\emph{kombinasi}).
  \item[Compact] a closed and bounded set; every sequence in it has a subsequence that converges inside that set (\emph{kompak}).
  \item[Completeness axiom] the promise that every nonempty set bounded above in $\mathbb{R}$ has a supremum; this is what leaves the number line with no holes (\emph{aksioma kelengkapan}).
  \item[Condition number] a measure of how much the relative error of the output can grow compared with the relative error of the input (\emph{bilangan kondisi}).
  \item[Conditional probability] the probability of an event given that another event is known to have occurred; written $P(A\mid B)$ (\emph{peluang bersyarat}).
  \item[Confidence interval] a range of values that, if the sampling procedure were repeated, contains the true parameter with a given probability (\emph{interval kepercayaan}).
  \item[Confusion matrix] a four-cell table that compares predictions with reality (\emph{matriks konfusi}).
  \item[Continuous] a function whose graph can be drawn without lifting the pen at that point, that is, $\lim_{x\to c}f(x)=f(c)$ (\emph{kontinu}).
  \item[Convergence in distribution] its distribution function tends to the limit distribution function at every point of continuity; what comes to resemble is the shape of the distribution, not its value (\emph{konvergensi dalam distribusi}).
  \item[Convergence in probability] $P(|X_n-X|>\varepsilon)\to 0$ for every $\varepsilon>0$ (\emph{konvergensi dalam peluang}).
  \item[Convex combination] a weighted mixture of vectors with nonnegative weights that sum to one (\emph{kombinasi konveks}).
  \item[Convex function] a function whose graph always lies below the chord joining two of its points; every local minimum of it is a global minimum (\emph{fungsi konveks}).
  \item[Coordinate vector] the list of coefficients that expresses a vector in a basis (\emph{koordinat}).
  \item[Correlation] a measure of the strength of the linear relation between two variables; it does not state cause and effect (\emph{korelasi}).
  \item[Countable] a set whose members can be numbered with nothing skipped; $\mathbb{Q}$ is countable, $\mathbb{R}$ is not (\emph{terhitung}).
  \item[Countable set] a set whose members can be listed as first, second, third, and so on with nothing left behind (\emph{himpunan terhitung}).
  \item[Counterexample] one example that breaks a statement of the form \"for all\" (\emph{lawan contoh}).
  \item[Covariance] a measure of two quantities moving together, computed from the product of their deviations (\emph{kovarians}).
  \item[Covariance] a measure of how far two variables move together; its unit is the product of their units (\emph{kovariansi}).
  \item[Cross-entropy] a measure of surprise $-\sum_i y_i\log \hat{p}_i$; the more confident the model is on a wrong answer, the larger the penalty (\emph{entropi silang}).
  \item[Cross-validation] assessing a model by dividing the data in turn into training and test parts (\emph{validasi silang}).
  \item[Cumulative distribution function] the function $F(x)=P(X\le x)$; a rising staircase whose height is the probability already accumulated up to $x$ (\emph{fungsi distribusi kumulatif}).
  \item[Data leakage] information from the test data that leaks into the training process without anyone intending it, so that test results look better than reality (\emph{kebocoran data}).
  \item[Data split] the division of data into training, validation, and test sets that do not leak into one another (\emph{pemisahan data}).
  \item[Defective] a matrix whose eigenvalues repeat but whose eigenspace is too narrow to be a basis (\emph{defective}).
  \item[Dense set] a set every point of which can be approached by members of that set other than the point itself (\emph{himpunan padat}).
  \item[Derivative] the limit of the difference quotient; the slope of the tangent line and at the same time the instantaneous rate of change (\emph{turunan}).
  \item[Determinant] one number that measures the factor by which volume changes and at the same time tests invertibility (\emph{determinan}).
  \item[Diagonalization] a change of basis until the matrix becomes diagonal, so that the work is reduced to multiplying numbers (\emph{diagonalisasi}).
  \item[Differential equation] a sentence that ties the rate of change of a quantity to the size of that quantity itself (\emph{persamaan diferensial}).
  \item[Differential privacy] the promise that one person's participation in the data changes the output distribution only slightly, measured by $\varepsilon$ (\emph{privasi diferensial}).
  \item[Digital signature] a value that only the holder of the private key can create but anyone can check with the public key (\emph{tanda tangan digital}).
  \item[Dimension] the number of vectors in a basis; the number belongs to the space, not to one basis that happens to be chosen (\emph{dimensi}).
  \item[Direct sum] the sum of two subspaces that meet only at the zero vector, so that every vector has one way of being written (\emph{jumlah langsung}).
  \item[Double integral] the sum of values over small cells of a region, then its limit as the cells shrink (\emph{integral lipat dua}).
  \item[Duality] formulating an optimization problem from the other side, whose optimal value gives a lower bound for the original problem (\emph{dualitas}).
  \item[Eigenspace] all the eigenvectors for one eigenvalue together with the zero vector; it is a subspace (\emph{ruang eigen}).
  \item[Eigenvalue] a scalar $\lambda$ such that $A\mathbf{v}=\lambda\mathbf{v}$ for some nonzero vector (\emph{nilai eigen}).
  \item[Elementary matrix] an identity matrix that has had one row operation applied to it, so that it becomes a record of that operation (\emph{matriks elementer}).
  \item[Elimination] replacing a system by a simpler system without changing its solution set (\emph{eliminasi}).
  \item[Epsilon-delta definition] the style of defining a limit by naming the output tolerance first, then the input bound that guarantees it (\emph{epsilon-delta}).
  \item[Equilibrium point] a state of a dynamical system that does not change when it is not disturbed (\emph{titik setimbang}).
  \item[Estimator] a rule that turns a sample into a guess at a parameter; the guess it produces is called the estimate (\emph{penaksir}).
  \item[Euler method] a numerical solution of a differential equation that steps along the tangent line; its error is first order in the step (\emph{metode Euler}).
  \item[Euler totient function] the function $\varphi(n)$: the count of numbers between $1$ and $n$ that are coprime with $n$ (\emph{fungsi phi}).
  \item[Expectation] the expected value $\mathbb{E}[X]$: the probability-weighted average, the balance point of a distribution (\emph{harapan}).
  \item[Expectation over transformation] an attack that maximizes the expected loss over the randomness of the defense, usually by averaging gradients over several random copies (\emph{EOT}).
  \item[Exponential distribution] the distribution of waiting time in a process without aging, with density $\lambda e^{-\lambda t}$ (\emph{distribusi eksponensial}).
  \item[Fast gradient sign method] a one-step attack (FGSM) that shifts every coordinate by the budget, in the direction of the sign of the loss derivative (\emph{metode tanda gradien}).
  \item[Field] a set with two operations that is closed, has identities and inverses, and satisfies the distributive law (\emph{lapangan}).
  \item[Floating-point system] binary scientific notation with a fixed number of digits; a number is stored as sign, exponent, and mantissa, so that the spacing between representable numbers is not uniform: dense near zero, sparse at the top end, and two different numbers can fall into the same place (\emph{sistem floating-point}).
  \item[Formal verification] a proof that a program or formula meets its specification, checked by a machine (\emph{verifikasi formal}).
  \item[Functional limit] $f(x)$ can be forced as close as we like to $L$ provided $x$ is close enough to $c$ but not equal to $c$ (\emph{limit fungsi}).
  \item[General solution] the family of all solutions of a differential equation, depending on arbitrary constants (\emph{solusi umum}).
  \item[Generalization] the ability to do problems never seen before, not problems already memorized (\emph{generalisasi}).
  \item[Global sensitivity] the largest shift in the answer to a query that one person can cause in the worst case (\emph{global sensitivity}).
  \item[Gradient] the vector of partial derivatives; the direction of steepest ascent of a function (\emph{gradien}).
  \item[Gradient descent] the parameter update rule that moves against the direction of the gradient (\emph{PGC, penurunan gradien}).
  \item[Harmonic series] the series $1+1/2+1/3+\cdots$; its terms go to zero but its sum diverges (\emph{deret harmonik}).
  \item[Heavy tail] a distribution whose probability of enormous values shrinks too slowly, so that one observation can outweigh all the rest (\emph{ekor berat}).
  \item[Hessian] the matrix of second partial derivatives; the concavity of a function is decided by its sign (\emph{Hessian}).
  \item[Hidden layer] a layer of units that touches neither the input nor the output of the network directly (\emph{lapisan tersembunyi}).
  \item[Hypothesis test] a procedure for deciding, with a stated risk of error, whether the data are enough to reject an initial claim (\emph{uji hipotesis}).
  \item[Implication] a statement of the form if $P$ then $Q$; it is false only when $P$ is true and $Q$ is false (\emph{implikasi}).
  \item[Independent] two events that give each other no news, $P(A\cap B)=P(A)P(B)$; knowing one does not change the probability of the other (\emph{independen}).
  \item[Independent and identically distributed] samples that are mutually independent and share one distribution; the most often used assumption, and the most often violated (\emph{i.i.d.}).
  \item[Induction] the domino pattern of proof: the base case holds, then each case knocks down the next (\emph{induksi}).
  \item[Inner product] a mapping of two vectors to one scalar that measures length and angle; in $\mathbb{R}^n$ its value is $\langle \mathbf{x},\mathbf{y}\rangle=\mathbf{x}^{\top}\mathbf{y}$ (\emph{hasilkali dalam}).
  \item[Integrality gap] the distance between the best value of the LP relaxation and the best value of the integer answer; it measures how far a continuous model misleads (\emph{celah integralitas}).
  \item[Integrating factor] a multiplier that turns the left side of a linear differential equation into a single derivative of a product (\emph{faktor integrasi}).
  \item[Interpolation] finding a function that passes through every data point; different from regression, which does not demand that (\emph{interpolasi}).
  \item[Isomorphism] a bijective mapping that preserves structure; two isomorphic spaces cannot be told apart algebraically (\emph{isomorfisme}).
  \item[Jacobian] the matrix of first partial derivatives of a vector mapping; its determinant measures the local volume scale factor (\emph{Jacobian}).
  \item[Jacobian determinant] the local volume scale factor when coordinates are changed; its absolute value multiplies area or volume in a multiple integral (\emph{determinan Jacobian}).
  \item[Kamus (glossary box)] a small box inside a chapter that explains one new term together with its analogy (\emph{kamus}).
  \item[Karush-Kuhn-Tucker conditions] necessary conditions for optimality in a problem with inequality constraints (\emph{kondisi KKT}).
  \item[Kernel, null space] the set of all vectors mapped to zero by a linear transformation (\emph{inti}).
  \item[Kernel, null space] the set of all inputs mapped to zero by a linear transformation; the same as the null space (\emph{kernel}).
  \item[Lagrange multiplier] an auxiliary variable that puts a price on a constraint; its value shows how much the objective could improve if the constraint were loosened (\emph{pengali Lagrange}).
  \item[Law of large numbers] the promise that the average of many observations sticks to the true mean, with the chance of missing it shrinking (\emph{hukum bilangan besar}).
  \item[Least squares] the choice of parameters that minimizes the sum of squared residuals when there is no exact solution (\emph{kuadrat terkecil}).
  \item[Linear map] a mapping that preserves addition and scalar multiplication (\emph{transformasi linear}).
  \item[Linear program] the problem of maximizing or minimizing a linear function with linear constraints (\emph{program linear}).
  \item[Lipschitz condition] the condition that the rate of change of a function is bounded by one constant; this is what guarantees a unique solution of a differential equation (\emph{Lipschitz}).
  \item[Logical connective] an operation that combines statements: and, or, not, if-then, if-and-only-if (\emph{penghubung logika}).
  \item[Logistic regression] a model for binary output that maps a linear combination to a probability through the logistic function (\emph{regresi logistik}).
  \item[Loss function] a number we want to make small while the model is trained; it measures how wrong a prediction is, not how right (\emph{fungsi kerugian}).
  \item[LP relaxation] replacing the integer condition by a real condition so that an integer problem can be solved as an LP (\emph{relaksasi LP}).
  \item[Lyapunov function] a positive-valued function that decreases along the trajectories of a system; its existence proves stability without solving the equations (\emph{fungsi Lyapunov}).
  \item[Machine epsilon] the distance between $1$ and the next machine number; the size of one step at scale one (\emph{epsilon mesin}).
  \item[Modular inverse] a number $b$ such that $ab\equiv 1 \pmod n$; the undo button of clock arithmetic (\emph{invers modular}).
  \item[Monotone sequence] a sequence that only rises or only falls, never turning back (\emph{barisan monoton}).
  \item[Multicollinearity] the situation in which explanatory variables are nearly linearly dependent on one another, so that regression coefficients become unstable (\emph{multikolinearitas}).
  \item[Multiplicity] the number of times an eigenvalue occurs as a root of the characteristic polynomial, which is not always equal to the dimension of its eigenspace (\emph{multiplisitas}).
  \item[Necessary condition] a condition that must hold for a statement to be true; without it the statement is impossible, but meeting it is not yet enough (\emph{syarat perlu}).
  \item[Nested cross-validation] two layers of folds: the outer layer assesses, the inner layer selects, so that model selection does not leak test data (\emph{validasi silang bersarang}).
  \item[Network intrusion detection system] a system that watches network traffic and decides whether a flow of packets is an attack (\emph{NIDS}).
  \item[Norm] a measure of the length of a vector; the $L_\infty$ norm takes the largest component (\emph{norma}).
  \item[Normal distribution] the symmetric bell distribution determined entirely by its mean and its variance (\emph{distribusi normal}).
  \item[Null hypothesis] the initial claim under test, usually \"there is no effect\"; it is rejected only when the data fit it very badly (\emph{hipotesis nol}).
  \item[Obfuscated gradient] a gradient handed to the attacker that does not point in the right direction, usually because a piece is not differentiable or is random (\emph{gradien dikaburkan}).
  \item[Odds ratio] the ratio of the probability of success to the probability of failure; the exponential of a logistic regression coefficient states how many times that ratio multiplies per unit (\emph{rasio peluang}).
  \item[Ordered field] a field whose members can be ordered and whose arithmetic rules respect that order (\emph{lapangan terurut}).
  \item[Ordinary differential equation] an equation that relates a function to its derivative with respect to one independent variable (\emph{persamaan diferensial biasa}).
  \item[Orientation] an agreement about which direction of rotation counts as positive; reversing it reverses the sign of an integral (\emph{orientasi}).
  \item[Orthogonal] mutually perpendicular in the sense that their inner product is zero (\emph{ortogonal}).
  \item[Orthogonal matrix] a square matrix whose columns are orthonormal; it preserves length and angle (\emph{matriks ortogonal}).
  \item[Orthonormal] an orthogonal set each of whose members has length one (\emph{ortonormal}).
  \item[Out-of-sample] data not used when the model was trained, so that they are honest for judging it (\emph{luar sampel}).
  \item[Overfitting] memorizing the training data including its noise, so that performance on new data falls (\emph{overfitting}).
  \item[Padding] extra content added to a packet before it is sent so that every packet has a fixed size, so that packet size no longer leaks its content (\emph{bantalan}).
  \item[Partition] a division of a set into pieces that do not overlap and that cover the whole set (\emph{partisi}).
  \item[Perceptron] the simplest linear classifier: the sign of the inner product of the weights and the input (\emph{perceptron}).
  \item[Poles and zeros] values of a complex variable at which a function blows up or vanishes; used in the analysis of dynamical systems and control (\emph{kutub dan nol}).
  \item[Polynomial time] the class of algorithms whose number of steps grows no faster than a polynomial in the size of the input (\emph{waktu polinomial}).
  \item[Power series] a series of the form $\sum a_n(x-c)^n$; it holds only inside its radius of convergence (\emph{deret pangkat}).
  \item[Power set] the set of all subsets, including the empty set and the set itself; its cardinality is $2^n$ (\emph{himpunan kuasa}).
  \item[Prediction interval] an interval for one observation that has not yet occurred; always wider than the interval for a mean (\emph{selang prediksi}).
  \item[Primitive root] a member that, raised repeatedly to powers, passes through every other member of the cyclic group before returning to one; in cryptography, a starting position that reaches every position of the key (\emph{akar primitif}).
  \item[Privacy budget] the amount $\varepsilon$ a system may spend; every query or training step draws down what is left (\emph{anggaran privasi}).
  \item[Probability] a weight on every possibility, nonnegative and summing to one (\emph{peluang}).
  \item[Probability amplitude] a complex number whose squared modulus gives a probability; different from a probability because two amplitudes can cancel each other (\emph{amplitudo}).
  \item[Problem space] the space of real objects an attacker may change, for example packets and timing; different from the feature space in which the model works (\emph{ruang masalah}).
  \item[Projection] the point in a subspace closest to a vector, that is, its perpendicular image (\emph{proyeksi}).
  \item[Proof term] an expression that is at once a proof and a value; checking a proof is the same as checking a type (\emph{term bukti}).
  \item[Proposition] a sentence that is either true or false, not a question or a command; in this book also the name of a theorem environment (\emph{proposisi}).
  \item[Pseudoinverse] a generalized inverse $A^{+}$ that still works when some directions have been flattened: what can be inverted is inverted, what is lost is left alone (\emph{pseudoinvers}).
  \item[P-value] the probability of obtaining a result at least as extreme as the one observed if the null hypothesis is true; not the probability that the null hypothesis is true (\emph{nilai-p}).
  \item[Quantifier] a marker for all ($\forall$) or there exists ($\exists$); the order in which they are used changes the meaning of a statement (\emph{kuantor}).
  \item[Quantization] replacing every value by the nearest multiple of a fixed step; the result is a staircase function whose derivative is zero almost everywhere (\emph{kuantisasi}).
  \item[Quantum computing era] the era when quantum computers are large enough to run Shor's algorithm on the keys used today (\emph{zaman kuantum}).
  \item[Radius of convergence] the greatest distance from the center at which a power series still converges (\emph{jari-jari konvergensi}).
  \item[Random variable] a function that maps the outcome of an experiment to a number (\emph{variabel acak}).
  \item[Rank] the number of independent directions that survive after mapping, that is, $\dim\operatorname{im}T$ (\emph{rank}).
  \item[Rank-nullity theorem] $\dim\ker T+\dim\operatorname{im}T=\dim V$ for a linear transformation on a finite-dimensional space (\emph{rank-nullity}).
  \item[Ratio test] a series converges absolutely if $\lim|a_{n+1}/a_n|<1$, and diverges if that limit is greater than one (\emph{uji rasio}).
  \item[Rearrangement] swapping the order of the terms of a series; in a series that converges absolutely the sum does not change, in one that does not, the sum can change (\emph{susunan ulang}).
  \item[Regression] drawing the function closest to all the points, not the function that passes through all of them (\emph{regresi}).
  \item[Regression line] the line (or plane) that minimizes the sum of squared residuals; equivalent to the orthogonal projection onto the column space (\emph{garis regresi}).
  \item[Regularization] a penalty on the size of the parameters so that the model is not too sensitive to the training data (\emph{regularisasi}).
  \item[Residual] the difference between an observed value and the value predicted by the model (\emph{residu}).
  \item[Residue class] the set of numbers that leave the same remainder when divided by $n$, like three in the morning and three in the afternoon (\emph{kelas sisa}).
  \item[Riemann integrability] the property of a function whose upper and lower sums can be squeezed until their difference vanishes (\emph{keterintegralan}).
  \item[Riemann integral] the limit of the sum of rectangle areas as the width of the partition shrinks (\emph{integral Riemann}).
  \item[ROC curve] a plot of the true rate against the false positive rate over all thresholds; the area under it is called AUROC (\emph{kurva ROC}).
  \item[Rounding] replacing a real number by the nearest machine number (\emph{pembulatan}).
  \item[Rounding error] the difference between a real number and its machine representative, which piles up in long sums (\emph{kesalahan pembulatan}).
  \item[Rule of three] if no event occurs in $n$ trials, the event rate may still be as high as about $3/n$ at 95 percent confidence (\emph{aturan tiga}).
  \item[Runge-Kutta] a family of numerical methods that evaluate the derivative several times per step for a smaller error (\emph{Runge-Kutta}).
  \item[Sample complexity] the number of examples needed for a model to generalize with a given confidence (\emph{kompleksitas sampel}).
  \item[Sensitivity] in differential privacy, the largest change in output that one person can cause (\emph{sensitivitas}).
  \item[Sequence] a row of numbers given a running order $a_1,a_2,\dots$; the order is part of the definition, not something added later (\emph{barisan}).
  \item[Series] the sum of the terms of a sequence without stopping, whose value is read as the limit of the partial sums (\emph{deret}).
  \item[Similar] two matrices that represent the same linear mapping in two different bases, $B=P^{-1}AP$ (\emph{serupa}).
  \item[Singular value] the square root of an eigenvalue of $A^{\top}A$; it measures the length of a principal axis of the mapping by the matrix $A$ (\emph{nilai singular}).
  \item[Singular value decomposition] writing a matrix as $U\Sigma V^{\top}$ with $U,V$ orthogonal (\emph{dekomposisi nilai singular, SVD}).
  \item[Slack variable] an extra variable that records the unused remainder, so that a no more than constraint becomes an equation (\emph{variabel kendur}).
  \item[Softmax] a mapping of a row of scores into probabilities that sum to one, $e^{z_i}/\sum_j e^{z_j}$ (\emph{softmax}).
  \item[Sorry] a Lean~4 marker that a proof has been postponed; the machine accepts it but records it as an axiom (\emph{sorry}).
  \item[Stability] the property of an equilibrium point or dynamical system that stays close to its state after a small disturbance (\emph{kestabilan}).
  \item[Standard error] the standard deviation of an estimator, not the standard deviation of the data; it measures the wobble of one guess (\emph{galat baku}).
  \item[Standard form] a uniform way of writing an optimization problem: one objective, constraints in the form no more than zero, and nonnegative variables (\emph{bentuk baku}).
  \item[Stochastic gradient descent] gradient descent that uses a small fraction of the data at every step; cheaper and noisier (\emph{penurunan gradien stokastik}).
  \item[Strong induction] induction that may use every earlier case, not only the nearest one (\emph{induksi kuat}).
  \item[Subsequence] a sequence obtained by discarding some terms without changing the order of the remaining terms (\emph{subbarisan}).
  \item[Subspace] a subset of a vector space that is closed under addition and scalar multiplication (\emph{subruang}).
  \item[Summary box] a box of bullet points at the end of a chapter that can be read on its own without the body of the chapter (\emph{ringkasan}).
  \item[Supremum, least upper bound] the least upper bound; the smallest number that is not smaller than every member of the set (\emph{supremum}).
  \item[Tactic] a step-by-step command in Lean~4 that changes the goal of a proof (\emph{taktik}).
  \item[Taylor series] writing a function as a sum of derivative terms at one point, together with a remainder that measures the truncation error (\emph{deret Taylor}).
  \item[Threat model] the complete list of an attacker's abilities: the size of the disturbance, what the attacker knows, and what the attacker wants (\emph{model ancaman}).
  \item[Transferability] the property of an adversarial example made for one model that also fools another model trained separately (\emph{transferabilitas}).
  \item[Truncation error] the difference that appears because an infinite process is replaced by a finite piece (\emph{galat pemotongan}).
  \item[Type class] a list of abilities a type must have in order to be usable in a given place (\emph{typeclass}).
  \item[Uniform continuity] continuity with one $\delta$ that works for every point of the set, not one $\delta$ per point (\emph{kekontinuan seragam}).
  \item[Uniform convergence] a sequence of functions that approaches its limit function at the same rate over the whole domain (\emph{kekonvergenan seragam}).
  \item[Uniform convergence] one rate for the whole domain, not one schedule per point (\emph{konvergensi seragam}).
  \item[Uniqueness] the guarantee that one initial state gives rise to only one solution trajectory (\emph{ketunggalan}).
  \item[Unit in the last place] the smallest step a machine can distinguish at the scale of one number (\emph{ULP}).
  \item[Unitary] a matrix that preserves the inner product, $U^{*}U=I$; the complex version of an orthogonal matrix (\emph{uniter}).
  \item[Upper limit] the largest value that a sequence still visits arbitrarily often near its end, written $\limsup_{n\to\infty}a_n$ (\emph{limit superior}).
  \item[Upper sum] a deliberately enlarged estimate of an integral, using the largest value on each piece (\emph{jumlah atas}).
  \item[Variance] a measure of spread around the mean; its square root is the standard deviation (\emph{variansi}).
  \item[Variance inflation factor] how much the variance of a coefficient is inflated by other explanatory variables carrying similar information; abbreviated VIF (\emph{varians inflasi}).
  \item[Vector space] a set with addition and scalar multiplication that is closed and satisfies eight axioms (\emph{ruang vektor}).
  \item[Vertex] a corner of the feasible region, that is, a point that is never the midpoint between two other feasible points (\emph{verteks}).
  \item[Wilson interval] a confidence interval for a proportion found by inverting a test, so that it stays sensible when the proportion approaches zero or one (\emph{interval Wilson}).
  \item[Wilson interval] a confidence interval for a proportion that stays sensible when the proportion approaches zero or one (\emph{Wilson}).
\end{description}

\chapter{Guidelines and Grading Criteria for Exercises}
\label{app:petunjuk}

The answer key is in Appendix~D, but its form is deliberately incomplete: final numbers
and sketches of the steps, not full proofs. The reason is consistency. A final number
without a trace of checking will teach the reader the very habit this book wants to
remove, and a proof served up whole removes the only useful exercise, namely finding for
yourself where an argument does not yet close. So use this key as a last comparison, not
as a place to start. What is provided here is a way to check your own work before you
open Appendix~D.

\section*{Three kinds of exercises and how to check them}

\begin{center}
\small
\begin{tabular}{@{}p{2.6cm}p{5.6cm}p{5.8cm}@{}}
\toprule
\textbf{Type} & \textbf{What is assessed} & \textbf{How to check} \\
\midrule
Proof & every step must start from a definition or a result already
  proved; no leaps are allowed & swap drafts with a classmate; ask her
  to point to the sentence she cannot derive herself \\
Computation & the final number and the intermediate steps & redo it with Python; compute the relative
  error against the exact result; check whether the answer makes sense in
  order of magnitude \\
Evaluating a claim & the quality of the argument, not a binary true-false & first write down what
  must be true for the claim to hold; look for one example that breaks it;
  only then decide ``true'', ``false'', or ``not enough evidence'' \\
\bottomrule
\end{tabular}
\end{center}

\section*{Verification tools per part}

\begin{center}
\small
\begin{tabular}{@{}llp{7.2cm}@{}}
\toprule
\textbf{Chapter} & \textbf{Tool} & \textbf{Typical check} \\
\midrule
1--5 & \texttt{python3} & table of values, factorization, partial sums; compare
  with an analytic estimate (e.g. $H_n\approx\ln n+\gamma$) \\
6--10 & \texttt{numpy}, \texttt{scipy.linalg} & rank, null space (\texttt{null\_space}),
  eigenvalues, SVD, least squares (\texttt{lstsq}); check
  $\|A\mathbf{v}-\lambda\mathbf{v}\|$ \\
11--16 & \texttt{scipy.integrate}, finite differences & quadrature vs antiderivative; rate of
  error at several $h$ or several $n$ \\
17--21 & \texttt{numpy.random}, \texttt{scipy.stats} & Monte Carlo simulation with a
  fixed \emph{seed}; compare the empirical probability with the exact probability and
  report its standard error \\
22 & \texttt{struct}, \texttt{float.hex} & bitfields, ULP, machine epsilon;
  compare the naive algorithm with Kahan \\
23--24 & \texttt{scipy.optimize} & \texttt{linprog}, \texttt{minimize},
  \texttt{milp}; check the KKT conditions and strong duality numerically \\
25--26 & \texttt{numpy} & gradient check by finite differences; training vs validation
  error curves at several regularization settings \\
27--30 & \texttt{numpy}, \texttt{scikit-learn} when available & confusion matrix from
  stored predictions; Wilson interval for proportions; accuracy under attack
  for several $\varepsilon$ and several \emph{seed} \\
\bottomrule
\end{tabular}
\end{center}

\section*{How to grade your own work}

\begin{enumerate}
  \item First write down what is given and what is to be proved.
        If neither is clear, the difficulty is not in the mathematics.
  \item For computation exercises, guess the answer before computing. A wrong guess
        is still useful: it shows where intuition misses.
  \item Compare the result by two different methods. One method can be right by
        accident; two methods that agree are rarely an accident.
  \item If the results differ greatly, look for the mistake in a simple case first
        (one variable, one step, $n=1$).
  \item Write down the limits of validity of the answer. An answer that is correct for a
        special case but presented as a general answer is the mistake that most often slips through.
\end{enumerate}

\chapter{Answer Keys}
\label{app:kunci}

Every number in this appendix was recomputed with Python before it was written down, and the
checking script ships with the manuscript sources, so a reader can run it independently. For
proof exercises only a sketch of two to five lines is given; that is enough to tell whether
the direction of your work is right, but not enough to be copied as a proof. Items whose
statement is itself flawed are marked plainly as ambiguous instead of being forced to have
an answer.

% Answer keys for Chapters 1--5. This file is \input directly inside the document.
% Every number comes from _wrap/kunci/cek.py.

\section*{Chapter 1: The Language of Mathematics and the Art of Proof}
\begin{enumerate}[label=\textbf{1.\arabic*}]
  \item $\exists \varepsilon>0\ \forall N\in\mathbb{N}\ \exists n\ge N:\ |a_{n}-L|\ge
        \varepsilon$. In words: there is one tolerance that is never attained,
        because however large $N$ is chosen, there is always a later term that
        deviates by at least that $\varepsilon$.

  \item The negation of ``every prime $>2$ is odd'' is ``some prime $>2$ is even''.
        That negation is \emph{false}, so the original statement is true and there is
        no counterexample. For the second claim, the negation is ``for every $n$,
        $n^{2}+n+41$ is prime''; the counterexample is $n=40$ with $1681=41^{2}$.

  \item $n^{2}+n=n(n+1)$. Cases: $n=2k$ gives $2k(2k+1)$, $n=2k+1$ gives
        $(2k+1)(2k+2)=2(2k+1)(k+1)$. Contraposition: if $n(n+1)$ is odd then $n$
        and $n+1$ are both odd, yet two consecutive numbers cannot both be odd.

  \item Check: $n=1$ gives $1$, $n=2$ gives $5$, $n=3$ gives $14$, matching
        $n(n+1)(2n+1)/6$. Induction step: add $(k+1)^{2}$ to both sides, then
        $\frac{k(k+1)(2k+1)}{6}+(k+1)^{2}=\frac{(k+1)(k+2)(2k+3)}{6}$.

  \item Strong induction on $n$. For $n=1$: $k=0$, $m=1$. For even $n$ write
        $n=2j$ with $j<n$; the induction hypothesis gives $j=2^{k}m$, so that
        $n=2^{k+1}m$ with $m$ still odd. The strong hypothesis is used exactly here:
        $j<n$, not $j=n-1$.

  \item (a) True, the ordinary induction principle. (b) False: take $P(n)$ to mean
        $n\le 10^{6}$; the counterexample is $n=10^{6}+1$. (c) False: $n=40$ gives
        $1681=41^{2}$. (d) False: $n=3$ is divisible by $3$ but not by $6$.

  \item The script finds $87$ primes out of $101$ values for $n=0$ through $100$. It
        fails at $40$, $41$, $44$, $49$, $56$, $65$, $76$, $81$, $82$, $84$, $87$, $89$,
        $91$, $96$. Only $5$ of the $14$ have the form $p-1$ or $p$: $40,41$ via
        $p=41$, $82$ via $83$, $89$, $96$ via $97$.

  \item $2^{1}=2>1$; $2^{2}=4=4$; $2^{3}=8<9$; $2^{4}=16=16$; $2^{5}=32>25$. Hence
        $n_{0}=5$. Induction: if $2^{k}>k^{2}$ and $k\ge 5$, then
        $2^{k+1}>2k^{2}\ge(k+1)^{2}$ because $k^{2}-2k-1>0$.

  \item 
The claim is false, because $n=40$ gives $40^{2}+40+41=1\,681=41^{2}$, which
        is not prime. Forty data points only support it, they do not prove it: one
        counterexample is enough to refute ``for every $n\ge0$''.

  \item False: $n=2$ gives $4\mid 4$ but $4\nmid 2$. The correct claim is:
        $4\mid n^{2}$ implies $2\mid n$. Proof by contraposition: if
        $n=2k+1$ then $n^{2}=4k^{2}+4k+1\equiv 1 \pmod 4$.
\end{enumerate}

\section*{Chapter 2: Sets, Relations, Functions, and Cardinality}
\begin{enumerate}[label=\textbf{2.\arabic*}]
  \item $x\in A\cap(B\cup C)\iff x\in A\wedge(x\in B\vee x\in C)\iff
        (x\in A\wedge x\in B)\vee(x\in A\wedge x\in C)$, using distributivity of
        $\wedge$ over $\vee$. Second: $x\in A\setminus(B\setminus C)\iff x\in
        A\wedge(\neg x\in B\vee x\in C)\iff x\in(A\setminus B)\cup(A\cap C)$.

  \item False. The smallest counterexample: $A=\{1\}$ with $R=\varnothing$; the empty
        relation is vacuously symmetric and transitive, but $(1,1)\notin R$. What is
        missing is reflexivity for elements that occur in no pair at all.

  \item Take $S_{1},S_{2}\subseteq A$ with $f(S_{1})=f(S_{2})$. For $s\in S_{1}$
        there is $t\in S_{2}$ with $f(s)=f(t)$, and injectivity gives $s=t\in S_{2}$.
        Without injectivity: $f$ constant on $\{1,2\}$ makes $f(\{1\})=f(\{2\})$.

  \item The direction $(\Leftarrow)$ is trivial. For $(\Rightarrow)$: if $a=b$, the
        equality $\{\{a\}\}=\{\{c\},\{c,d\}\}$ forces $\{c\}=\{c,d\}$, so $c=d=a$. If
        $a\ne b$, then $\{a\}$ must equal $\{c\}$, since $\{c,d\}$ has size
        two; so $a=c$ and then $b=d$.

  \item Define $g:f(A)\to A$ by choosing one preimage for each $y\in f(A)$;
        $g$ is injective, so $|f(A)|\le|A|$ by the previous lemma. If $f$ is not
        injective, two elements of $A$ give one value, so $|f(A)|<|A|$.

  \item $\binom{5}{2}=10$ partitions of type $2+3$ and $5$ of type $1+4$, altogether
        $15=2^{4}-1=S(5,2)$. The script counts partitions directly and also
        gives $15$; the total number of equivalence relations on $5$ elements is $52$.

  \item $\pi(7,5)=T_{12}+5=78+5=83$. For $1000$: $T_{44}=990\le 1000<T_{45}=1035$,
        so $s=44$, $n=10$, $m=34$; the check $\pi(34,10)=990+10=1000$.

  \item Cumulative through $k=9$ is $27$; level $k=10$ contains $1/9,\,3/7,\,7/3,\,
        9/1$ at positions $28$--$31$, so $7/3$ sits at position $30$. Because
        $\varphi(11)=10$, position $41$ is the last term of level $11$, namely
        $10/1$.

  \item Arrange $A_{n}=\{a_{n1},a_{n2},\dots\}$ and then run the Cantor diagonal over
        the pairs $(n,m)$: every member of $\bigcup A_{n}$ appears at least once,
        and discarding the repeated appearances does not increase the count.

  \item $g(x)=\tan\bigl(\pi(x-\tfrac12)\bigr)$ is a bijection from $(0,1)$ onto
        $\mathbb{R}$. Its domain must be restricted because $\tan$ is infinite
        at $\pm\tfrac12\pi$; on $[0,1]$ that formula is undefined at both ends.
\end{enumerate}

\section*{Chapter 3: Axioms of the Real Numbers and Supremum}
\begin{enumerate}[label=\textbf{3.\arabic*}]
  \item $a<b$ means $b-a>0$. Add $-a-b$ to both sides (closure under
        addition): $(-a)-(-b)=b-a>0$, that is $-b<-a$. The reverse direction is obtained
        by adding $a+b$. Hence multiplying by $-1$ reverses the direction.

  \item Since $a<b$, we have $2a<a+b$ and $a+b<2b$; multiply by $\tfrac12>0$ and
        we get $a<\frac{a+b}{2}<b$. Only the order axioms are used. Conversely
        $\sup\{q:q^{2}<2\}$ does not exist in $\mathbb{Q}$, so completeness is required.

  \item $\sup S=2$: $2$ is an upper bound, and $2-\varepsilon$ is not an upper bound because
        $2-1/n>2-\varepsilon$ for $n>1/\varepsilon$. $\inf S=1$ is attained at
        $n=1$, so a minimum exists; no maximum exists because $2\notin S$.

  \item True. The premise is the Archimedean property: for every $x>0$ there is $n$ with
        $1/n<x$. Consequently, if $m>0$ were the least, then $m/2>0$ is smaller
        than $m$, a contradiction.

  \item $q=1$: $h=1/4$, so $q+h=5/4$ with $(5/4)^{2}=1. 5625<2$. $q=7/5$:
        $h=1/120$, so $q+h=169/120\approx 1. 4083$ with square
        $1. 9834<2$. The second increment is $30$ times smaller: the closer $q$ is to
        $\sqrt2$, the shorter the correction.

  \item $|x^{3}-y^{3}|=|x-y|\,|x^{2}+xy+y^{2}|\le|x-y|\bigl(|x|^{2}+|x||y|+|y|^{2}\bigr)
        \le 12\,|x-y|$, because each term is at most $4$. The triangle inequality
        is used in the second step to split the absolute value of a sum.

  \item If $S$ is open and $x$ is a limit point of $S^{c}$ with $x\in S$, then there is
        $B(x,r)\subseteq S$, so that $B(x,r)\cap S^{c}=\varnothing$: a contradiction.
        Conversely, if $S$ is closed and $x\in S^{c}$, the absence of a ball inside
        $S^{c}$ makes $x$ a limit point of $S\subseteq S$, a contradiction.

  \item $x_{1}=7/4$; $x_{2}=97/56\approx 1. 732143$;
        $x_{3}=18817/10864\approx 1. 732050810$; $x_{4}$ and $x_{5}$ agree to
        $16$ digits, $1. 732050807568877$. Relative error after
        five steps: $1. 0363\times10^{-2}$, $5. 3145\times10^{-5}$,
        $1. 4121\times10^{-9}$, $9. 97\times10^{-19}$, $4. 97\times10^{-37}$.
        Ratios $e_{k+1}/e_{k}^{2}$: $0. 4330$, $0. 4949$, $0. 5000$, tending to $1/2$.

  \item $0$ lies in every $(-1/n,1/n)$. For $x\ne0$ choose $n>1/|x|$, then
        $x\notin(-1/n,1/n)$. A finite intersection of open sets stays open, whereas
        this infinite intersection equals $\{0\}$, which is not open.

  \item By density of $\mathbb{Q}$, in the interval $(a/\sqrt2,\,b/\sqrt2)$ there is
        $q\in\mathbb{Q}$ with $q\ne0$. Then $a<q\sqrt2<b$, and $q\sqrt2$ is irrational
        because if it were rational then $\sqrt2$ would be rational.
\end{enumerate}

\section*{Chapter 4: Sequences and Limits}
\begin{enumerate}[label=\textbf{4.\arabic*}]
  \item $|x_{n}-\tfrac23|=\Bigl|\frac{3(2n-1)-2(3n+4)}{3(3n+4)}\Bigr|
        =\frac{11}{9n+12}$. The condition is $n>\frac{11/\varepsilon-12}{9}$, so
        $N(\varepsilon)=\bigl\lfloor(11/\varepsilon-12)/9\bigr\rfloor$. For
        $\varepsilon=10^{-3}$: $N=1220$, while the smallest term that qualifies is
        $n=1221$ with error $9. 9991\times10^{-4}$.

  \item $\bigl||x_{n}|-|L|\bigr|\le|x_{n}-L|$ (the reverse triangle inequality),
        so the same $\varepsilon$ bound applies. The converse is false:
        $x_{n}=(-1)^{n}$ gives $|x_{n}|=1$ convergent, but $x_{n}$ diverges.

  \item For $x_{n}$: even terms $2+1/n\to2$, odd terms $1/n\to0$; hence
        $\limsup=2$, $\liminf=0$ (the maximum $2. 5$ at $n=2$ is not the limsup). For
        $y_{n}$: even $1+1/n\to1$, odd $-(1+1/n)\to-1$; hence $\limsup=1$,
        $\liminf=-1$. Numerically at $n=1000$: $1{,}001$ and $-1{,}001$.

  \item Partly right, wrong conclusion. It is true that $(1+1/n)^{n}$ increases and
        is bounded by $3$, so its limit exists and is at most $3$. But an upper bound
        is not the value of the limit: the limit is $e\approx 2. 7183$.

  \item $3=\bigl(3^{n}\bigr)^{1/n}\le(2^{n}+3^{n})^{1/n}\le 3\cdot 2^{1/n}$, and
        $2^{1/n}\to1$ because $(1+\delta)^{n}\ge 1+n\delta$ (Bernoulli) gives
        $\delta\le 2^{1/n}-1\le 1/n$. Both fences go to $3$.

  \item The error is $<10^{-6}$ first at $n=12$ ($2. 6146\times10^{-7}$) and
        $<10^{-12}$ at $n=21$. The ratio of successive errors settles at $0. 250000$, consistent with
        $|x_{n+1}-2|=|x_{n}-2|/(\sqrt{2+x_{n}}+2)\to|x_{n}-2|/4$.

  \item A Cauchy sequence is bounded: take $\varepsilon=1$ and get $N$; all terms after $N$
        lie in $(x_{N}-1,x_{N}+1)$, and the rest are finitely many. A bounded monotone
        sequence converges by the monotone convergence theorem, and a convergent sequence
        is Cauchy.

  \item $\varepsilon=10^{-3}$ gives the threshold $5/\varepsilon-2=4998$, so
        $N=4998$ is valid. Errors: $n=4998$ gives exactly $10^{-3}$ (not yet qualifying), $n=4999$
        gives $9. 9980\times10^{-4}$, $n=5000$ gives $9. 9960\times10^{-4}$.
        $N-1=4997$ is not valid because $n=4998>4997$ still has error exactly
        $\varepsilon$, not less.

  \item Suppose $x_{n}\not\to L$; then there is $\varepsilon>0$ and a subsequence
        $(x_{n_{k}})$ with $|x_{n_{k}}-L|\ge\varepsilon$. By
        Bolzano--Weierstrass that subsequence has a subsequence converging to $L'$ with
        $|L'-L|\ge\varepsilon$, contradicting the hypothesis.
\end{enumerate}

\section*{Chapter 5: Infinite Series}
\begin{enumerate}[label=\textbf{5.\arabic*}]
  \item $s_{20}=2. 4999451576$, whereas the sum is $5/2=2. 5$; the error is
        $5. 4842\times10^{-5}$. The bound $\frac{r^{21}}{1-r}$ gives exactly the same
        figure, $5. 4842\times10^{-5}$, because the tail is geometric with positive terms.

  \item $\frac{1}{n(n+2)}=\frac12\bigl(\frac1n-\frac1{n+2}\bigr)$, so that
        $s_{N}=\frac12\bigl(1+\frac12-\frac{1}{N+1}-\frac{1}{N+2}\bigr)$; its limit is
        $\frac12\cdot\frac32=\frac34$.

  \item Because $a_{n}\to0$, for $n\ge N$ we have $0\le a_{n}\le1$, so
        $a_{n}^{2}\le a_{n}$; the comparison test closes it. The converse is false:
        $a_{n}=1/n$ gives $\sum a_{n}^{2}$ convergent but $\sum a_{n}$ divergent.

  \item False. The argument breaks at ``smaller than a geometric term of ratio $\tfrac12$ for
        $n$ sufficiently large'': $1/n<(1/2)^{n}$ in fact fails, because $(1/2)^{n}$ shrinks
        far faster. The comparison test needs an inequality on the whole tail. The
        harmonic series diverges.

  \item (a) Convergent, ratio test: the ratio $\to\tfrac12<1$. (b) Divergent, ratio test:
        the ratio $2\bigl(\frac{n}{n+1}\bigr)^{2}\to2>1$. (c) Convergent, ratio test:
        $\bigl(\frac{n}{n+1}\bigr)^{n}\to\frac1e\approx0. 3679<1$. (d) Divergent,
        $p$-series with $p=\tfrac13\le1$.

  \item $s_{10^{6}}=1. 6449330668$ and $\pi^{2}/6=1. 6449340668$; the difference is
        $9. 9999950\times10^{-7}$, close to $\frac1N$. After the correction
        $\frac1N-\frac{1}{2N^{2}}$ the result is $1. 6449340668$ with error
        $2. 2\times10^{-16}$. Correct decimal digits: $5$ before, $15$ after.

  \item We need $\frac{1}{\sqrt{N+1}}<10^{-3}$, that is $N+1>10^{6}$, so $N=10^{6}$.
        At that $N$ the sum is $0. 6043986435$ while the limit is
        $0. 6048986434$; the actual error is $5. 0000\times10^{-4}$, exactly half the
        bound. The bound is loose by about a factor of two.

  \item $|b_{n}|\le M$ gives $|a_{n}b_{n}|\le M|a_{n}|$. Because $\sum M|a_{n}|$
        converges, the comparison test concludes that $\sum|a_{n}b_{n}|$ converges.

  \item After $K$ blocks, the positive terms used run up to denominator $6K-1$ and the
        negative terms up to $\frac{1}{2K}$, so
        $S_{K}=H_{6K}-\frac12H_{3K}-\frac12H_{K}\to\ln6-\frac12\ln3=\ln(2\sqrt3)
        \approx1. 2425$. Numerically $K=10^{6}$ gives $1. 2424531$. Its ratio to
        $\ln2$ is $1. 7925$, so it is not a rational multiple of $\ln2$.

  \item Write $p_{n}=\max(a_{n},0)$ and $q_{n}=\max(-a_{n},0)$; then
        $a_{n}=p_{n}-q_{n}$ and $|a_{n}|=p_{n}+q_{n}$. If both of those series
        converged, $\sum|a_{n}|$ would converge, contradicting conditional
        convergence. Hence both diverge, and those two unbounded piles are the fuel
        of the Riemann rearrangement theorem.
\end{enumerate}

% Answer keys for Chapters 6--15. This file is \input directly inside the document.
% Every number comes from the check scripts in the temporary working directory.

\section*{Chapter 6: Vector Spaces and Subspaces}
\begin{enumerate}[label=\textbf{6.\arabic*}]
  \item $W$ is the solution set of the homogeneous system $x-2y+z=0$, $x+y=0$, so it is
        closed under addition and scalar multiplication. From $x=-y$ and $z=3y$ we get
        $W=\operatorname{span}\{(-1,1,3)\}$: a one-vector basis, $\dim W=1$.

  \item The set $D=\{A:a_{11}=a_{22}=0\}$ is closed under addition and scalars,
        because both diagonal entries stay zero. Every one of its members is
        $\begin{pmatrix}0&b\\c&0\end{pmatrix}=bE_{12}+cE_{21}$, so
        $\{E_{12},E_{21}\}$ is a basis and $\dim D=2$.

  \item $(1+x)+(1+x^{2})=2+x+x^{2}$: the set is linearly dependent, hence
        $\dim\operatorname{span}S=2$. The vector $x$ is not in that span, because
        $a(1+x)+b(1+x^{2})$ demands $a=1$ and $a+b=0$ at the same time. So
        $\{1+x,\ 1+x^{2},\ x\}$ is a basis of $P_{2}$.

  \item Collect the coefficients: $(c_{1}+c_{3})\mathbf{v}_{1}+(c_{1}+c_{2})
        \mathbf{v}_{2}+(c_{2}+c_{3})\mathbf{v}_{3}=\mathbf{0}$. Independence of
        $\{\mathbf{v}_{1},\mathbf{v}_{2},\mathbf{v}_{3}\}$ gives three homogeneous
        equations with determinant $2\ne0$, so $c_{1}=c_{2}=c_{3}=0$.

  \item $U$ is the $xy$-plane, $W$ is the plane $x=y$. Their intersection is
        $U\cap W=\operatorname{span}\{(1,1,0)\}$, so $\dim(U\cap W)=1$ and
        $\dim(U+W)=2+2-1=3$. The formula fits: $U+W=\mathbb{R}^{3}$, and numpy
        gives the rank of the four spanning vectors together as $3$.

  \item Elimination puts pivots in columns $1,2,4$, so the rank is $3$; a basis of the
        column space is $(1,2,0)$, $(2,4,1)$, and $(4,9,1)$. Nullity $4-3=1$ with
        basis $\{(-1,-1,1,0)\}$; \texttt{matrix\_rank} and
        \texttt{null\_space} confirm both numbers.

  \item False. In $\mathbb{R}^{2}$ take $U_{1}=\operatorname{span}\{(1,0)\}$,
        $U_{2}=\operatorname{span}\{(0,1)\}$, $U_{3}=\operatorname{span}\{(1,1)\}$.
        All pairwise intersections are $\{\mathbf{0}\}$, yet
        $(1,1)=(1,0)+(0,1)$: the representation is not unique, and
        $\dim U_{1}+\dim U_{2}+\dim U_{3}=3>2=\dim\mathbb{R}^{2}$.

  \item Take a basis $\mathbf{u}_{1},\dots,\mathbf{u}_{n-1}$ of $U$. That set is
        linearly independent and has fewer than $n$ elements, so it does not yet span
        $V$; there is some $\mathbf{v}\in V$ outside its span. The extension theorem makes
        $\mathbf{u}_{1},\dots,\mathbf{u}_{n-1},\mathbf{v}$ a basis of $V$, so that
        $V=U\oplus\operatorname{span}\{\mathbf{v}\}$.

  \item A polynomial of degree at most $4$ that vanishes at $0$ and $1$ is divisible
        by $x(x-1)$, so $p=x(x-1)(a+bx+cx^{2})$: three free parameters,
        dimension $3$. Basis $\{x(x-1),\,x^{2}(x-1),\,x^{3}(x-1)\}$. The evaluation
        matrix $2\times5$ has rank $2$, so nullity $5-2=3$.

  \item The coordinates satisfy $c_{1}+c_{2}+c_{3}+c_{4}=4$, $c_{2}+c_{3}+c_{4}=2$,
        $c_{3}+c_{4}=-5$, $c_{4}=1$. From the bottom: $c_{4}=1$, $c_{3}=-6$,
        $c_{2}=7$, $c_{1}=2$. So $[q]_{B}=(2,7,-6,1)$, the same as the result of
        \texttt{numpy.linalg.solve}.
\end{enumerate}

\section*{Chapter 7: Basis, Dimension, and Linear Transformations}
\begin{enumerate}[label=\textbf{7.\arabic*}]
  \item Solve $c_{1}(1,1,0)+c_{2}(1,0,1)+c_{3}(0,1,1)=(2,-1,3)$, that is
        $c_{1}+c_{2}=2$, $c_{1}+c_{3}=-1$, $c_{2}+c_{3}=3$. The result is
        $c=(-1,3,0)$, so $[\mathbf{w}]_{B'}=(-1,3,0)$. Back:
        $-1(1,1,0)+3(1,0,1)+0=(2,-1,3)$.

  \item Let $\varphi$ be an isomorphism. Because $\varphi$ is bijective,
        $\mathbf{b}_{i}=\varphi^{-1}(\mathbf{e}_{i})$ is defined; because
        $\varphi^{-1}$ is linear and bijective, $\{\mathbf{b}_{i}\}$ is a basis and
        $\varphi=\varphi_{B}$. It is not unique because each basis $B$ gives a
        different isomorphism.

  \item $A$ has rank $2$ (pivots in columns $2,3$), nullity $1$. Kernel:
        $x_{2}=0$, $2x_{3}=0$, so $\ker A=\operatorname{span}\{(1,0,0)\}$.
        Range $\operatorname{span}\{(1,0,0),(0,2,0)\}$. This is the derivative matrix
        on $P_{2}$ in the basis $(1,x,x^{2})$: $D(1)=0$, $D(x)=1$,
        $D(x^{2})=2x$.

  \item $\operatorname{im}(S\circ T)\subseteq\operatorname{im}S$ gives
        rank$(S\circ T)\le$ rank $S$. Besides that,
        $\ker T\subseteq\ker(S\circ T)$, so nullity$(S\circ T)\ge$ nullity $T$;
        rank-nullity applied to both maps gives rank$(S\circ T)\le$ rank $T$.

  \item False. $\begin{pmatrix}1&0\\0&0\end{pmatrix}$ and
        $\begin{pmatrix}0&1\\0&0\end{pmatrix}$ both have rank $1$, but
        the second is nilpotent (trace $0$, square zero) while the first
        has trace $1$. Rank measures only the dimension of the range.

  \item $T$ sends $\mathbf{b}_{1},\mathbf{b}_{2},\mathbf{b}_{3}$ to $(1,0)$,
        $(3,1)$, $(3,2)$. Coordinates relative to $C$ are $(x,y)\mapsto(y,x-y)$,
        so the columns are $(0,1),(1,2),(2,1)$ and
        $[T]_{C\leftarrow B}=\begin{pmatrix}0&1&2\\1&2&1\end{pmatrix}$.
        For $\mathbf{v}=(1,2,3)$: $[\mathbf{v}]_{B}=(-1,-1,3)$ gives
        $(5,0)$, the same as $[T(\mathbf{v})]_{C}$.

  \item Take a basis $\mathbf{u}_{1},\dots,\mathbf{u}_{m}$ of $U$ with
        $m=\dim V$. It is linearly independent and has $\dim V$ elements, so it spans
        $V$; hence every $\mathbf{v}\in V$ lies in $U$, that is $U=V$.

  \item $B=P^{-1}AP$, so $B\mathbf{x}=\mathbf{0}\iff AP\mathbf{x}=\mathbf{0}
        \iff P\mathbf{x}\in\ker A$. The mapping $\mathbf{x}\mapsto P\mathbf{x}$ is
        an isomorphism from $\ker B$ onto $\ker A$, so the nullities agree. That
        equality is exactly the statement that the dimension of the kernel does not
        depend on the basis.

  \item The third row is the sum of the first two, so rank $2$ and nullity
        $4-2=2$. A nonzero kernel vector $(-2,-1,1,0)$: $F(-2,-1,1,0)=\mathbf{0}$.
        That direction shifts no feature at all, because
        $F(\mathbf{x}+t\mathbf{k})=F(\mathbf{x})$ for all $t$.
\end{enumerate}

\section*{Chapter 8: Linear Systems, Determinants, and Inverses}
\begin{enumerate}[label=\textbf{8.\arabic*}]
  \item The pivots that appear are $1$, $1$, $-1$; no row needs swapping, so
        $\det A=1\cdot1\cdot(-1)=-1$. Backward: $z=3$, $y=-1$, $x=2$. Cramer
        gives $\det A_{x}=-2$, $\det A_{y}=1$, $\det A_{z}=-3$; divided by $-1$
        the result is $(2,-1,3)$.

  \item Cofactor expansion gives $\det M=0$. The refuting vector is
        $\mathbf{v}=(1,-1,-1,1)$: every row gives $0$, for instance
        $0-1-2+3=0$. So $\det M=0$ without any expansion. numpy gives
        $5.3\times10^{-15}$, not zero, because elimination subtracts nearly equal
        numbers and the rounding error is counted too.

  \item Leibniz formula: $\det A=\sum_{\sigma}\operatorname{sgn}(\sigma)
        \prod_{i}a_{i\sigma(i)}$. If columns $j$ and $k$ agree, pair every
        $\sigma$ with $\sigma\circ(j\,k)$; the products of the terms are the same but
        the signs are opposite, so the sum of all terms is zero.

  \item Multiply $A\operatorname{adj}A=(\det A)I$ by the determinant:
        $\det A\cdot\det(\operatorname{adj}A)=(\det A)^{n}$, because
        $\det(cI)=c^{n}$. If $\det A\ne0$, divide by $\det A$. If
        $\det A=0$, both sides vanish for $n\ge2$, so the equality still holds.

  \item False as a general statement. Sensitivity is measured by
        $\kappa(A)=\lVert A\rVert\,\lVert A^{-1}\rVert$, not by $\det A$:
        $A=\varepsilon I$ has $\det A=\varepsilon^{n}$ very small yet
        $\kappa=1$. A small determinant is dangerous only when it comes with a
        large $\lVert A^{-1}\rVert$.

  \item $\lVert A\rVert_{1}=2+\varepsilon$ and
        $A^{-1}=\frac1\varepsilon\begin{pmatrix}1+\varepsilon&-1\\-1&1\end{pmatrix}$
        give $\lVert A^{-1}\rVert_{1}=(2+\varepsilon)/\varepsilon$. So
        $\kappa_{1}(A)=(2+\varepsilon)^{2}/\varepsilon$. The condition $\kappa>10^{6}$
        is $\varepsilon<4.000016\times10^{-6}$; at
        $\varepsilon=4\times10^{-6}$ numpy gives $1.000004\times10^{6}$.

  \item Adding a multiple of a row means multiplying on the left by
        $E=I+cE_{ij}$, and $\det E=1$, so $\det B=\det A$. Because $B=EA$,
        we have $B^{-1}=A^{-1}E^{-1}$ with $E^{-1}=I-cE_{ij}$.

  \item $A+AD=A(I+D)$, and $I+D$ is diagonal, so
        $\det(A+AD)=\det A\prod_{i}(1+d_{i})$. It is invertible exactly when
        $\det A\ne0$ and no $d_{i}=-1$. Example $\det A=5$ with
        $D=\operatorname{diag}(2,-\tfrac12,3)$: the left side is $30$, the right side
        $5\cdot3\cdot\tfrac12\cdot4=30$.

  \item With $A=\begin{pmatrix}1&2&3\\0&1&4\\5&6&0\end{pmatrix}$,
        $\det A=1$ and
        \[
          \operatorname{adj}A=A^{-1}=
          \begin{pmatrix}-24&18&5\\20&-15&-4\\-5&4&1\end{pmatrix},
        \]
        the same as \texttt{numpy.linalg.inv} up to machine error. For
        $\det A=7$:
        $\begin{pmatrix}2&1&0\\1&1&0\\0&0&7\end{pmatrix}^{-1}$ contains
        $1/7$ in the third entry, because division by $7$ cannot be avoided.
\end{enumerate}

\section*{Chapter 9: Eigenvalues and Diagonalization}
\begin{enumerate}[label=\textbf{9.\arabic*}]
  \item $\operatorname{tr}A_{1}=5$ and $\det A_{1}=4$, so $\lambda=4$ and
        $\lambda=1$. Eigenvectors: $(1,1)$ for $\lambda=4$, $(1,-2)$ for
        $\lambda=1$. Then $P=\begin{pmatrix}1&1\\1&-2\end{pmatrix}$ and
        $D=\operatorname{diag}(4,1)$; $P^{-1}A_{1}P=D$ is verified.

  \item $\lambda=5,2$ with eigenvectors $(2,1)$ and $(1,-1)$. For
        $k\ge1$,
        \[
          G^{k}=\frac13
          \begin{pmatrix}2\cdot5^{k}+2^{k}&2\cdot5^{k}-2\cdot2^{k}\\
          5^{k}-2^{k}&5^{k}+2\cdot2^{k}\end{pmatrix},
        \]
        so that
        $G^{10}=\begin{pmatrix}6510758&6509734\\3254867&3255891\end{pmatrix}$,
        the same as \texttt{matrix\_power}.

  \item $A\mathbf{v}=\lambda\mathbf{v}$ gives
        $A^{2}\mathbf{v}=\lambda^{2}\mathbf{v}$; induction gives
        $A^{k}\mathbf{v}=\lambda^{k}\mathbf{v}$. The converse fails:
        $A=\begin{pmatrix}0&-1\\1&0\end{pmatrix}$ has $A^{2}=-I$ with eigenvalue
        $-1$, whereas $A$ itself has no real eigenvalue.

  \item $0$ is an eigenvalue $\iff$ there is $\mathbf{v}\ne\mathbf{0}$ with
        $A\mathbf{v}=\mathbf{0}$ $\iff$ $\ker A\ne\{\mathbf{0}\}$ $\iff$ $A$ is not
        invertible. If $A$ is invertible and $A\mathbf{v}=\lambda\mathbf{v}$, then
        $\lambda\ne0$ and $A^{-1}\mathbf{v}=\lambda^{-1}\mathbf{v}$.

  \item False. $\begin{pmatrix}1&0\\0&1\end{pmatrix}$ and
        $\begin{pmatrix}1&1\\0&1\end{pmatrix}$ both have characteristic
        polynomial $(\lambda-1)^{2}$, but the second is defective:
        $\operatorname{rank}(J-I)=1$, so nullity $1<2$.

  \item $\lambda I-A$ and its transpose have the same determinant, because
        $\det X=\det X^{\top}$; hence $A$ and
        $A^{\top}$ have the same characteristic polynomial and the same eigenvalues.
        Their eigenvectors do not: $A$ and $A^{\top}$ have different eigenspaces as
        soon as $A$ is not symmetric.

  \item $K-2I=\begin{pmatrix}0&1&0\\0&0&0\\0&0&-1\end{pmatrix}$ has rank $2$,
        so nullity $=1$. The algebraic multiplicity of $\lambda=2$ is $2$,
        the geometric one is $1$; because $1<2$, $K$ cannot be diagonalized.

  \item $\lambda=1$ and $0.5$. The fixed point of $x=4y$ with $x+y=1$ is
        $(0.8;0.2)$. The error shrinks like $0.5^{k}$; because
        $0.5^{19}=1.91\times10^{-6}>\allowbreak10^{-6}$ and
        $0.5^{20}=9.54\times10^{-7}$, it takes $20$ steps.

  \item The printed matrix gives
        $\operatorname{tr}(B^{10})=(-19)+21+1024=1026$, and
        $1^{10}+1^{10}+2^{10}=1026$: they agree. For a diagonalizable $A$,
        $A^{k}=PD^{k}P^{-1}$ yields
        $\operatorname{tr}(A^{k})=\operatorname{tr}(D^{k})=\sum_{i}\lambda_{i}^{k}$.
        For defective $B$ that step is unavailable, because $B$ has no $P$ and $D$.
        Instead use the Schur triangular form $B=QTQ^{-1}$ with $T$ upper triangular
        and diagonal $\lambda_{i}=1,1,2$; then $B^{k}=QT^{k}Q^{-1}$, and $T^{k}$ is
        still upper triangular with diagonal $\lambda_{i}^{k}$, so
        $\operatorname{tr}(B^{k})=\operatorname{tr}(T^{k})
        =\sum_{i}\lambda_{i}^{k}=1026$. The trace is invariant under change of basis
        ($\operatorname{tr}XY=\operatorname{tr}YX$), which is why the result does not
        depend on $Q$. In short: diagonalization carries the proof only in the
        diagonalizable case; in the defective case the proof switches to the
        triangular form, the formula stays the same, and the factor $k$ that appears
        in the entries of $B^{k}$ never touches the diagonal.
\end{enumerate}

\section*{Chapter 10: Orthogonality, SVD, and Least Squares}
\begin{enumerate}[label=\textbf{10.\arabic*}]
  \item $(1,0,0)$ and $(0,1,1)$ are already orthogonal, so
        $\operatorname{proj}_{W}\mathbf{b}=1\cdot(1,0,0)+\frac{4}{2}(0,1,1)
        =(1,2,2)=\mathbf{b}$. Because $\mathbf{b}\in W$, the distance is $0$;
        \texttt{qr} gives $QQ^{\top}\mathbf{b}=(1,2,2)$ as well.

  \item $\mathbf{b}=\mathbf{p}+\mathbf{r}$ with $\mathbf{r}\perp W$, so
        $\mathbf{r}\perp\mathbf{p}$ and Pythagoras gives
        $\lVert\mathbf{b}\rVert^{2}=\lVert\mathbf{p}\rVert^{2}
        +\lVert\mathbf{r}\rVert^{2}$. If $\mathbf{p}$ is not an orthogonal
        projection, $\mathbf{r}$ is not perpendicular to $\mathbf{p}$; on an oblique
        basis in $\mathbb{R}^{2}$ a cross term $2\langle\mathbf{p},\mathbf{r}
        \rangle\ne0$ appears.

  \item Orthonormal columns: $(A^{\top}A)_{ij}=\mathbf{a}_{i}^{\top}\mathbf{a}_{j}
        =\delta_{ij}$, so $A^{\top}A=I$. The normal equations
        $A^{\top}A\hat{\mathbf{x}}=A^{\top}\mathbf{b}$ immediately give
        $\hat{\mathbf{x}}=A^{\top}\mathbf{b}$. All singular values of $A$ are
        $1$, so $\kappa(A)=1$.

  \item False. $A=\begin{pmatrix}1&0\\0&1\\1&1\end{pmatrix}$ gives
        $A^{\top}A=\begin{pmatrix}2&1\\1&2\end{pmatrix}$, which is invertible
        ($\det=3$), yet its rows are dependent: the third row is the sum of the first
        two. What invertibility of $A^{\top}A$ guarantees is only independence of the
        columns.

  \item Singular values $2$ and $0$, so rank $1$;
        $U=V=\frac{1}{\sqrt2}\begin{pmatrix}1&1\\1&-1\end{pmatrix}$ (up to
        sign). $A^{+}=\frac14\begin{pmatrix}1&1\\1&1\end{pmatrix}$, so
        $A^{+}\mathbf{b}=(1,1)$ and the residual is $(-1,1)$ with norm
        $\sqrt2=1.4142$; \texttt{pinv} agrees.

  \item 
Write $\mathbf{b}=A\hat{\mathbf{x}}+\mathbf{r}$ with
        $A\hat{\mathbf{x}}\in\operatorname{col}A$. For
        $\mathbf{h}\in\operatorname{col}A$ the function
        $t\mapsto\lVert\mathbf{b}-A(\hat{\mathbf{x}}+t\mathbf{h})\rVert^{2}$ is
        minimal at $t=0$, hence $\langle\mathbf{r},\mathbf{h}\rangle=0$; this is the
        ``necessary'' direction. Conversely, if
        $\langle\mathbf{r},\mathbf{h}\rangle=0$ for every
        $\mathbf{h}\in\operatorname{col}A$, then
        $\lVert\mathbf{b}-A\mathbf{x}\rVert^{2}=\lVert\mathbf{r}\rVert^{2}+\lVert
        A(\hat{\mathbf{x}}-\mathbf{x})\rVert^{2}\ge\lVert\mathbf{r}\rVert^{2}$, so
        $\hat{\mathbf{x}}$ minimises. Column independence is what makes
        $\hat{\mathbf{x}}$ unique.

  \item $\mathbf{u}_{1}=(1,1,1)/\sqrt3$; the remainder $(1,1,0)-\frac23(1,1,1)
        =\frac13(1,1,-2)$, so $\mathbf{u}_{2}=(1,1,-2)/\sqrt6$; then
        $\mathbf{u}_{3}=(1,-1,0)/\sqrt2$. $\lVert Q^{\top}Q-I\rVert$:
        $6.7\times10^{-16}$ (by hand), $3.1\times10^{-16}$ (\texttt{qr}).
        The signs may differ because each $\mathbf{u}_{i}$ is determined only up to
        $\pm$.

  \item $AA^{\top}\mathbf{x}=\mathbf{0}$ gives
        $\lVert A^{\top}\mathbf{x}\rVert^{2}=\mathbf{x}^{\top}AA^{\top}
        \mathbf{x}=0$, so $A^{\top}\mathbf{x}=\mathbf{0}$; independence of the rows
        makes $\mathbf{x}=\mathbf{0}$. Hence $AA^{\top}$ is square and invertible.
        For $\mathbf{x}=A^{\top}(AA^{\top})^{-1}\mathbf{b}$ we have
        $A\mathbf{x}=\mathbf{b}$ and
        $\mathbf{x}\in\operatorname{col}A^{\top}=(\ker A)^{\perp}$.

  \item With the last reading $314.3$: $\sum y=1540.5$,
        $\sum Ty=157161.7$, determinant $50$. Cramer:
        $a=(52030\cdot1540.5-510\cdot157161.7)/50=-5.04$ and
        $b=(5\cdot157161.7-510\cdot1540.5)/50=3.07$;
        \texttt{lstsq} agrees. The intercept shifts by $-10.10$, the slope by $+0.10$;
        the endpoint has leverage $0.6$, the largest of the five.

  \item $B$ has rank $1$, so $B^{+}$ gives
        $\mathbf{x}^{+}=\frac{1}{70}(1,2)\approx(0.01429;\,0.02857)$ and
        $\mathbf{k}=(-2,1)$. All solutions are $\mathbf{x}^{+}+t\mathbf{k}$;
        because $\mathbf{x}^{+}\perp\mathbf{k}$,
        $\lVert\mathbf{x}^{+}+t\mathbf{k}\rVert^{2}=\lVert\mathbf{x}^{+}
        \rVert^{2}+t^{2}\lVert\mathbf{k}\rVert^{2}$ is minimal at $t=0$.
\end{enumerate}

\section*{Chapter 11: Continuity and Uniform Continuity}
\begin{enumerate}[label=\textbf{11.\arabic*}]
  \item $\lvert(2x-1)-5\rvert=2\lvert x-3\rvert$, so $\delta=\varepsilon/2$ is
        valid and largest: at $x=3+\delta$ the value is exactly $\varepsilon$.
        The choice $\delta=\varepsilon/3$ is also valid, but it uses only a third of the
        room available.

  \item $\lvert x^{2}-c^{2}\rvert=\lvert x-c\rvert\,\lvert x+c\rvert$; with
        $\lvert x-c\rvert<1$ we get $\lvert x+c\rvert\le2\lvert c\rvert+1$,
        so $\delta=\min(1,\varepsilon/(2\lvert c\rvert+1))$. At $c=2$:
        $\delta=0.002$; at $c=100$: $\delta=4.975\times10^{-5}$.

  \item Partly true. $1/x$ is indeed continuous at every point of $(0,1]$. But the
        uniform conclusion is wrong: $(0,1]$ is not compact and $1/x$ is unbounded
        near $0$. The refuting pair $x_{n}=1/n$, $y_{n}=1/(n+1)$:
        $\lvert x_{n}-y_{n}\rvert\to0$ but the difference of the values is $1$.

  \item Given $\varepsilon$, take $\delta_{f}$ and $\delta_{g}$; with
        $\delta=\min(\delta_{f},\delta_{g})$,
        $\lvert(f+g)(x)-(f+g)(y)\rvert\le\varepsilon/2+\varepsilon/2$.
        For $fg$: $f=g=x$ on $\mathbb{R}$ with $x_{n}=n$,
        $y_{n}=n+1/n$ gives difference $2+1/n^{2}\ge2$.

  \item $\lvert\sqrt x-\sqrt y\rvert^{2}\le\lvert x-y\rvert$ (rationalize the
        denominator), so $\delta=\varepsilon^{2}$ works for all
        $x,y\ge0$. Only $x,y\ge0$ is used, so that the roots are defined;
        closedness is not needed because the proof never touches a limit point outside
        the domain.

  \item 
$(1+\delta)^{3}-1=\varepsilon$ gives
        $\delta=(1+\varepsilon)^{1/3}-1=0{,}0033223$. The choice
        $\delta=\varepsilon/3=0{,}0033333$ is actually too large:
        $(1+\varepsilon/3)^{3}-1-\varepsilon=\varepsilon^{2}/3+\varepsilon^{3}/27>0$
        for every $\varepsilon>0$, so it violates the bound and is not valid.

  \item With $\delta=\varepsilon/2$: the largest difference
        $2\delta-\delta^{2}=0.009975<\varepsilon$, so it is safe but not
        optimal. With $\delta=1-\sqrt{1-\varepsilon}=0.0050126$: the largest
        difference $1-(1-\delta)^{2}=\varepsilon$ exactly, attained at
        $x=1-\delta$.

  \item $\lvert f\rvert$ is continuous on the compact $[a,b]$, so it attains a minimum $m$;
        because $f\ne0$ everywhere, $m>0$ and $\lvert f\rvert\ge m$. For
        $f(x)=x$ on $(0,1)$ the conclusion fails: the infimum is $0$ and it is not
        attained.

  \item At $c\ne0$ take $q_{n}\to c$ rational and $r_{n}\to c$ irrational; the
        function values tend to $c$ and $-c$, which differ, so $f$ is not continuous. At $0$:
        $\lvert f(x)\rvert=\lvert x\rvert\to0=f(0)$. The function $\lvert f\rvert$
        is $x\mapsto\lvert x\rvert$, continuous everywhere; there is no
        contradiction because $\lvert f\rvert$ throws away the sign information.
\end{enumerate}

\section*{Chapter 12: Derivatives: One and Several Variables}
\begin{enumerate}[label=\textbf{12.\arabic*}]
  \item $\bigl((2+h)^{3}-8\bigr)/h=12+6h+h^{2}\to12$. For $\sqrt x$:
        $(\sqrt{2+h}-\sqrt2)/h=1/(\sqrt{2+h}+\sqrt2)\to1/(2\sqrt2)
        =0.3535534$. In both cases $h$ drops out of the denominator.

  \item $f'(c)>0$ makes $(f(c+h)-f(c))/h>f'(c)/2>0$ for $h$ small enough,
        so $f(c+h)>f(c)$. That is only a one-sided statement. The example
        $f(x)=x+2x^{2}\sin(1/x^{2})$ has $f'(0)=1>0$ but is not increasing on
        any neighbourhood of $0$: its derivative changes sign arbitrarily close
        to $0$.

  \item Differentiate $1/g$ first: $(1/g)'=-g'/g^{2}$; here the continuity of $g$
        is used, so that $g(x+h)\to g(x)\ne0$. Then
        $(f/g)'=f'\cdot(1/g)+f\cdot(1/g)'$, which simplifies to
        $(f'g-fg')/g^{2}$.

  \item False. $f(x,y)=xy/(x^{2}+y^{2})$ for $(x,y)\ne(0,0)$ and
        $f(0,0)=0$: both partial derivatives at the origin exist and equal $0$,
        but along $y=x$ the value is $\tfrac12$, so $f$ is not continuous
        there.

  \item $\nabla f=(3x^{2}y-y^{3},\,x^{3}-3xy^{2})$, at $(2,1)$ equal to
        $(11,2)$. The steepest ascent direction is $(11,2)/\sqrt{125}$, the maximum rate
        $\sqrt{125}=11.1803$. The central difference with $h=10^{-4}$ gives relative error
        $9.1\times10^{-10}$ and $1.0\times10^{-8}$, consistent with order $h^{2}$ in the
        second component.

  \item By the mean value theorem,
        $f(b)-f(a)=f'(\xi)(b-a)=0$ for every pair $a<b$ in the interval, so
        $f$ is constant. On a disconnected set that step fails because $\xi$
        need not lie in the set; for example $f=0$ on $(0,1)$ and $f=1$
        on $(2,3)$.

  \item At $(0,1)$: $J_{H}=\begin{pmatrix}1&1\\1&2\end{pmatrix}$; at
        $H(0,1)=(1,1)$: $J_{G}=\begin{pmatrix}2&1\\1&-1\end{pmatrix}$. The matrix
        product is $\begin{pmatrix}3&4\\0&-1\end{pmatrix}$, and the direct derivative
        gives the same matrix.

  \item $f_{x}=3x^{2}-3y$, $f_{y}=3y^{2}-3x$; both zero gives $y=x^{2}$ and
        $x=y^{2}$, so $x(x^{3}-1)=0$: the points $(0,0)$ and $(1,1)$. The Hessian
        $\begin{pmatrix}6x&-3\\-3&6y\end{pmatrix}$: at $(0,0)$ its determinant is
        $-9<0$ (saddle), at $(1,1)$ its determinant is $27>0$ with $f_{xx}=6>0$
        (local minimum, $f=-1$). If the determinant is zero, that test says nothing.

  \item $p_{2}(x)=1+\tfrac12h-\tfrac18h^{2}$ at $h=0.04$ gives $1.0198$,
        whereas $\sqrt{1.04}=1.0198039$; the error is
        $3.9\times10^{-6}$, the Lagrange remainder
        $\frac{3}{8}\xi^{-5/2}h^{3}/6\le4.0\times10^{-6}$. The bound is tight, not loose, because $\xi$ is close
        to $1$: the Lagrange remainder really does nearly equal the actual error.
\end{enumerate}

\section*{Chapter 13: Riemann Integrals and the Fundamental Theorem of Calculus}
\begin{enumerate}[label=\textbf{13.\arabic*}]
  \item $U_{n}=\frac1n\sum_{k=1}^{n}\frac kn=\frac{n+1}{2n}$ and
        $L_{n}=\frac{n-1}{2n}$, so $U_{n}-L_{n}=1/n\to0$: $f$ is integrable
        and $\int_{0}^{1}x\,dx=\tfrac12$.

  \item $\bigl\lvert\lvert f(x)\rvert-\lvert f(y)\rvert\bigr\rvert
        \le\lvert f(x)-f(y)\rvert$ (reverse triangle inequality), so
        integrability carries over. The bound comes from $-\lvert f\rvert\le f\le\lvert f\rvert$
        and monotonicity of the integral. The converse is false: $f=1$ on the rationals and
        $-1$ on the irrationals has $\lvert f\rvert=1$, which is integrable.

  \item False. The correct part: the upper and lower sums of a bounded function are always
        finite. What the proof does not support is the equality
        $\sup_{\mathcal{P}}L=\inf_{\mathcal{P}}U$. The Dirichlet function is bounded
        with $U-L=1$ for every partition.

  \item On a uniform partition each subinterval contributes oscillation
        $f(x_{k})-f(x_{k-1})$, so that
        $U-L\le\frac{b-a}{n}\sum_{k}\bigl(f(x_{k})-f(x_{k-1})\bigr)
        =\frac{(b-a)(f(b)-f(a))}{n}$. For a non-monotone function that telescoping fails
        because the signs of the oscillations are not uniform.

  \item $\int_{0}^{1}(3x^{2}-2x+1)\,dx=[x^{3}-x^{2}+x]_{0}^{1}=1$. Simpson
        with $n=100$ gives $1.0000000000$ and \texttt{quad} $1.0000000000$
        (error $1.1\times10^{-14}$); the three numbers differ by zero to machine
        precision.

  \item $\int_{1}^{10}x^{-2}dx=1-\tfrac1{10}=0.9$ and
        $\int_{1}^{\infty}x^{-2}dx=1$. Their difference is $0.1=\tfrac1{10}$ exactly,
        namely the tail $\int_{10}^{\infty}x^{-2}dx=1/10$; there is no need to recompute it,
        because the missing part is exactly that tail.

  \item If $f(x_{0})>0$, continuity gives $f>\tfrac12f(x_{0})$ on an interval of
        width $2\delta$ around $x_{0}$ (of width $\delta$ when $x_{0}$ is an
        endpoint), so the integral is at least
        $\tfrac12\delta f(x_{0})>0$: a contradiction. For $f$ that is merely integrable
        the statement fails: $f=1$ at a single point and $0$ elsewhere has integral
        $0$.

  \item $u=\sin x$: $\int_{0}^{1}u\,du=\tfrac12$. Through the identity:
        $\tfrac12\int_{0}^{\pi/2}\sin2x\,dx
        =\tfrac12\bigl[-\tfrac12\cos2x\bigr]_{0}^{\pi/2}=\tfrac12$. \texttt{quad}
        gives $0.5$ as well.

  \item On every subinterval there is a rational (value $1$) and an irrational (value
        $0$), so $U(\mathcal{P},D)=1$ and $L(\mathcal{P},D)=0$; the difference is $1$
        for all $\mathcal{P}$. The Riemann criterion demands that this difference can be
        made as small as we please, so $D$ is not integrable.

  \item The integral bounds give $\frac1{101}<\sum_{n=101}^{\infty}n^{-2}
        <\frac1{100}$, that is $9.9010\times10^{-3}$ up to
        $1.0000\times10^{-2}$. Their midpoint is $9.9505\times10^{-3}$,
        whereas the true value is $9.950167\times10^{-3}$; the relative error is
        $3.3\times10^{-5}$.
\end{enumerate}

\section*{Chapter 14: Series of Functions and the Exchange of Limits}
\begin{enumerate}[label=\textbf{14.\arabic*}]
  \item $\sup_{[0,1]}\lvert x^{n}\rvert=1$ for all $n$, attained at $x=1$.
        At $x=0.9$: $0.3487$ for $n=10$, $2.656\times10^{-5}$ for
        $n=100$, $1.748\times10^{-46}$ for $n=1000$. The first column is flat
        because the peak is always at the endpoint; the second column shrinks because $0.9<1$.

  \item For $n\ge N$, $\lVert f_{n}-f\rVert_{\infty}<\varepsilon$, so
        $\lVert f\rVert_{\infty}\le\lVert f_{N}\rVert_{\infty}+1$ and likewise
        for $\lVert f_{n}\rVert_{\infty}$. Pointwise convergence is not enough:
        $f_{n}=n\chi_{(0,1/n)}$ converges to $0$ at every point but
        $\lVert f_{n}\rVert_{\infty}=n$.

  \item The peak is at $x=1/\sqrt n$ with
        $\lVert f_{n}\rVert_{\infty}=1/(2\sqrt n)\to0$, and that bound holds
        for all $x$ at once, so the convergence is uniform. The condition
        $1/(2\sqrt n)<0.01$ requires $\sqrt n>50$, so $n=2501$.

  \item False. Pointwise reasoning gives one $N$ for each $x$, not
        one $N$ for all $x$. The counterexample is $x_{n}=2^{-1/n}$:
        $x_{n}^{n}=\tfrac12$ for every $n$, so
        $\sup_{[0,1)}\lvert x^{n}\rvert\ge\tfrac12$ and it does not tend to $0$.

  \item The error $\lvert e-S_{N}(1)\rvert$ for $N=5,10,15$:
        $1.615\times10^{-3}$, $2.731\times10^{-8}$,
        $5.077\times10^{-14}$; the bound $e/(N+1)!$ gives
        $3.775\times10^{-3}$, $6.810\times10^{-8}$,
        $1.299\times10^{-13}$, always larger. At $N=20$ the bound is already
        $5.32\times10^{-20}$, far below machine precision, so the measured error
        stops at $4.44\times10^{-16}$. The bound first drops below $10^{-16}$
        at $N=18$, and there the true error is unmeasurable as well.

  \item The sum is $(1-x)/(1-x)=1$ for $x\in[0,1)$, whereas at $x=1$ the
        sum is $0$. The limit is not continuous even though every partial sum is
        continuous; the uniform limit of continuous functions must be continuous, so
        the convergence is not uniform.

  \item Radii: $R=1$ for $\sum x^{n}/n^{2}$, $R=e$ for
        $\sum(n!/n^{n})x^{n}$, $R=\tfrac12$ for $\sum2^{n}x^{n}/n$. The ratio of the
        $1001$st and $1000$th terms: $0.998003$; $0.368063$ (so
        $R\approx2.7169$); $1.998002$. The slowest is $\sum x^{n}/n^{2}$,
        because its ratio tends to $1$ like $(1-1/n)^{2}$.

  \item $\lvert f_{n}g-fg\rvert=\lvert f_{n}-f\rvert\,\lvert g\rvert
        \le\lVert f_{n}-f\rVert_{\infty}\lVert g\rVert_{\infty}\to0$. Without
        boundedness of $g$: $f_{n}=x+1/n\to x$ uniformly on $(0,1)$, but
        $(x+1/n)/x=1+1/(nx)$ is unbounded, so $f_{n}g$ does not converge
        uniformly.

  \item The peak is at $x=1/\sqrt{2n+1}$. For $n=10$: height $1.3397$ at
        $x=0.2182$, $\int_{0}^{1}f_{10}=\tfrac{5}{11}=0.4545$. For
        $n=100$: height $4.2835$ at $x=0.07053$, the integral
        $\tfrac{50}{101}=0.4950$. The height follows $\sqrt{n/(2e)}$; the area
        tends to $\tfrac12$ because the peak gets narrower.

  \item Differentiating $\sum_{n\ge0}x^{n}=1/(1-x)$ gives
        $\sum_{n\ge1}nx^{n-1}=1/(1-x)^{2}$; at $x=\tfrac12$ we get
        $\sum_{n\ge1}n/2^{n}=2$. The tail remainder is $(N+2)/2^{N}$, so
        $N=39$ gives $7.46\times10^{-11}<10^{-10}$.
\end{enumerate}

\section*{Chapter 15: Multiple Integrals, Change of Variables, and Vector Field Theorems}
\begin{enumerate}[label=\textbf{15.\arabic*}]
  \item Cartesian: $\int_{-1}^{1}\int_{-\sqrt{1-x^{2}}}^{\sqrt{1-x^{2}}}
        (x^{2}+y^{2})\,dy\,dx=\pi/2$. Polar:
        $\int_{0}^{2\pi}\int_{0}^{1}r^{3}\,dr\,d\theta
        =2\pi\cdot\tfrac14=\pi/2=1.5708$. \texttt{dblquad} gives
        $1.5707963268$; in this test environment the Cartesian route requires $8379$
        evaluations, the polar route $441$.

  \item For $T(\mathbf{u})=A\mathbf{u}$ linear, $DT=A$ is constant, so
        its Jacobian determinant is $\det A$ and $\lvert\det A\rvert$ is the area
        scale factor. If $\det A\ne0$ the map is one-to-one on $D$, so
        $\iint_{T(D)}f=\lvert\det A\rvert\iint_{D}f\circ T$.

  \item Clockwise, each side is traversed in reverse, so the four terms of the example
        change sign: $-\tfrac13-2-\tfrac23+1=-2$, whereas the right-hand side stays
        $2$. The bottom and right sides that used to contribute $+\tfrac13$ and $+2$
        become negative; the left side that used to be $-1$ becomes $+1$.

  \item $x=au$, $y=bv$, $z=cw$ gives Jacobian $abc$ and the unit ball:
        $V=abc\cdot\frac{4\pi}{3}=\frac{4\pi abc}{3}$. At $a=b=c=1$
        the result is $4.1888$, the familiar formula for the volume of a ball.

  \item False. $\operatorname{curl}F=0$ does hold at every point where $F$
        is defined, but the domain has a hole so the curl test does not apply.
        The supporting number: $\oint_{\lvert\mathbf{r}\rvert=1}F\cdot d\mathbf{r}
        =2\pi\ne0$, so $F$ is not conservative.

  \item If the integral does not depend on the path, two paths $C_{1},C_{2}$ from
        $A$ to $B$ give the same integral, so the circulation around $C_{1}-C_{2}$ is zero.
        Conversely, if every closed curve has zero circulation, $C_{1}-C_{2}$ is
        closed and hence the two integrals agree.

  \item $\operatorname{curl}F=(0,0,2)$; the flux through $z=1-x^{2}-y^{2}$ with
        upward normal $=2\cdot\pi=2\pi$. The line integral along the boundary
        $x^{2}+y^{2}=1$, $z=0$:
        $\int_{0}^{2\pi}(-\sin t,\cos t,0)\cdot(-\sin t,\cos t,0)\,dt=2\pi$.
        Both are $6.2832$.

  \item $\operatorname{div}F=0$ away from the origin, so the multiple integral is $0$,
        whereas the flux through the unit sphere is $4\pi=12.5664$. Gauss's theorem
        requires a $C^{1}$ field on a region that contains the whole interior of the
        surface; that condition is violated at the enclosed origin.

  \item $\det J=r^{2}\sin\varphi$, zero at $\varphi=0$ and $\varphi=\pi$. At both
        poles the map stops being one-to-one: all $\theta$ give the same point.
        On $0<\varphi<\pi$ the map is one-to-one and $\det J>0$, so the change of
        variables formula remains valid even though the determinant is zero on the boundary.
\end{enumerate}

% Answer key for Chapters 16--23. This file is \input directly inside the document.
% Every number comes from the Python test scripts (see the report).

\section*{Chapter 16: Differential Equations and Dynamical Systems}
\begin{enumerate}[label=\textbf{16.\arabic*}]
  \item $y=Ce^{t^{2}}$ with $y(0)=1$ gives $y=e^{t^{2}}$, so $y(1)=e=2.718282$.
        Euler with $h=0.1$ uses $y_{n+1}=(1+0.02n)y_{n}$; the result is $2.334633$,
        error $3.8365\times10^{-1}$. That method lags behind because it uses
        the slope at the left end of each step.

  \item The difference $w=y_{1}-y_{2}$ satisfies $w'=-p(t)w$ and $w(t_{0})=0$. Multiply
        by $e^{\int p}$: $\bigl(we^{\int p}\bigr)'=0$, so $we^{\int p}$ is constant, and
        its value is zero at $t_{0}$. Hence $w\equiv0$. No solution formula is used,
        only the product rule and uniqueness of initial values.

  \item Not enough evidence. Linearization sees only the degree-one terms; the
        eigenvalues $\pm i$ say only that the linear part is a rotation. With $r'=r^{3}$
        the origin is a center for the linearization but an outward spiral, whereas
        $r'=-r^{3}$ gives it an inward spiral. So $\pm i$ does not determine stability.

  \item $E'=vv'+\omega_{0}^{2}\theta\theta'=v(-\omega_{0}^{2}\theta)+\omega_{0}^{2}\theta
        v=0$, so $E$ is constant. The equation $v^{2}/(2E)+\theta^{2}/(2E/\omega_{0}^{2})=1$
        is an ellipse, and $\theta=A\cos(\omega_{0}t+\varphi)$ closes one turn
        in $2\pi/\omega_{0}$.

  \item $\tau=RC=2.2\times10^{3}\cdot10^{-5}=0.022$ seconds. At $t=\tau$:
        $V=12(1-e^{-1})=7.585447$ volts. For $6$ volts, $t=-\tau\ln(1-\tfrac{6}{12})
        =0.022\cdot0.693147=15.249$ milliseconds.

  \item With $v=2\cos(\theta/2)$: $v^{2}/2=2\cos^{2}(\theta/2)$ and
        $1-\cos\theta=2\sin^{2}(\theta/2)$, and their sum is $E=2$. At $\theta=\pm\pi$
        the value is $v=0$, so the curve joins those two points; it separates
        oscillation from full rotation, and the time needed to travel it is infinite.

  \item $T_{0}=2\pi\sqrt{0.25/9.80665}=1.003205$ seconds.
        $T=T_{0}\cdot\tfrac{2}{\pi}K(\sin15^{\circ})=1.020669$ seconds; the difference is
        $1.74$ percent. An amplitude of $30^{\circ}$ lengthens the period, it does not
        shorten it.

  \item The factor per step is $1+h\lambda=1-100h$; the condition $|1-100h|<1$ gives
        $0<h<0.02$. Exactly at $h=0.02$ the factor is $-1$: the value changes sign
        with fixed magnitude $1$, neither decaying nor blowing up. That is the
        marginal stability boundary, not stability.

  \item Suppose $y$ is periodic and not constant. At its maximum point $y'=0$, so
        $f(y_{\max})=0$. The constant solution $y\equiv y_{\max}$ passes through that
        point, and uniqueness (Lipschitz) forces the solutions to be the same,
        contradicting non-constancy. So there is no non-constant periodic orbit.
\end{enumerate}

\section*{Chapter 17: Probability and Conditional Probability}
\begin{enumerate}[label=\textbf{17.\arabic*}]
  \item 
$P(A\cup B\cup C)=P(A\cup B)+P(C)-P\bigl((A\cup B)\cap C\bigr)$, then
        $P\bigl((A\cup B)\cap C\bigr)=P(A\cap C)+P(B\cap C)-P(A\cap B\cap C)$. The
        triple-overlap term is subtracted three times, once for each pair, and then
        added back once with a $+$ sign.

  \item False on general spaces. $P(A\mid B)=0$ only means $P(A\cap B)=0$. On a
        finite space with positive probabilities, that forces $A\cap B=\varnothing$.
        On a continuous space it does not: $A=\{1/2\}$ and $B=(0,1)$ uniform give
        $P(A\cap B)=0$ even though $1/2\in A\cap B$.

  \item From $10\,000$ people: $200$ are sick, $190$ correctly positive; with specificity
        $0.90$ there are $980$ false positives, so $190/1170=19/117\approx0.1624$.
        With $0.99$: only $98$ false positives remain, so $190/288=95/144
        \approx0.6597$: the precision rises fourfold.

  \item $P(A)=P(A\cap B)+P(A\cap B^{c})$, so
        $P(A\cap B^{c})=P(A)-P(A)P(B)=P(A)(1-P(B))=P(A)P(B^{c})$. For
        $A^{c},B^{c}$ use $P(A^{c}\cap B^{c})=1-P(A\cup B)$ and inclusion--exclusion,
        or apply the result above twice.

  \item $1-\prod_{k=0}^{29}(1-k/365)=0.706316$, whereas
        $1-e^{-435/365}=0.696320$; the difference is $0.009996$. The approximation uses
        $1-x\approx e^{-x}$ for $x$ up to $29/365$, and treats pair events
        as mutually independent when they are not.

  \item $A=(A\cap B)\cup(A\cap B^{c})$ and the two pieces are disjoint, because
        $B\cap B^{c}=\varnothing$. Additivity gives
        $P(A)=P(A\cap B)+P(A\cap B^{c})=P(A\mid B)P(B)+P(A\mid B^{c})P(B^{c})$,
        using the definition of conditional probability twice.

  \item The probability of winning by switching is $3/4=0.75$: the initial choice wins
        only when the car really is there ($1/4$). A simulation of $10^{4}$ games gives
        about $0.75$ with standard deviation $\sqrt{0.75\cdot0.25/10^{4}}\approx0.0043$.

  \item 
The correct figure is $0{,}9075$, not $0{,}9999$: with $0{,}1$ percent
        prevalence only $1$ person in $1000$ is ill, while the false positives among
        two positive tests are $999\cdot0{,}01^{2}\approx0{,}1$. Hence
        $0{,}99^{2}\cdot0{,}001/(0{,}99^{2}\cdot0{,}001+0{,}01^{2}\cdot0{,}999)=0{,}9075$.
        The assumptions: the two tests are conditionally independent and both
        sensitivity and specificity are exactly $0{,}99$.

  \item Not enough. Take $A$ the first coin shows heads, $B$ the second coin shows
        heads, $C$ both coins match. Every pair gives $1/4=P\cdot P$, but
        $P(A\cap B\cap C)=1/4$ whereas $P(A)P(B)P(C)=1/8$; the difference is $1/8$.
\end{enumerate}

\section*{Chapter 18: Random Variables and Distributions}
\begin{enumerate}[label=\textbf{18.\arabic*}]
  \item $c=1-(0.1+0.4)=0.5$; $\mathbb{E}[X]=-0.2+0+1.5=1.3$;
        $\mathbb{E}[X^{2}]=0.4+0+4.5=4.9$; $\operatorname{Var}(X)=4.9-1.69=3.21$.

  \item $(-\infty,b]=(-\infty,a]\cup(a,b]$ and the two pieces are disjoint, so
        $F(b)=F(a)+P(a<X\le b)$. For discrete variables $P(X=a)>0$ makes
        $P(a\le X\le b)$ larger; for continuous ones $P(X=a)=0$ so the two agree.

  \item It is true that $P(X=x)=0$ for every $x$ in a continuous distribution. The
        mistake is concluding ``impossible'': probability zero does not mean the empty
        set. Countable additivity gives probability zero to a countable union of points,
        but an interval contains uncountably many points and has probability one.

  \item Differentiate $\sum q^{k}=1/(1-q)$ twice: $\sum k(k-1)q^{k-2}=2/(1-q)^{3}$.
        Since $\sum k^{2}q^{k-1}=(1+q)/(1-q)^{3}$, with $q=1-p$ we get
        $\mathbb{E}[X^{2}]=p(1+q)/p^{3}=(1+q)/p^{2}=(2-p)/p^{2}$.
        Hence $\operatorname{Var}=(2-p)/p^{2}-1/p^{2}=(1-p)/p^{2}$; for $p=0.2$:
        $\mathbb{E}[X^{2}]=45$ and $\operatorname{Var}=20$.

  \item $\sigma=2$: $P(X>1)=1-\Phi(-1)=0.841345$;
        $P(1<X<5)=2\Phi(1)-1=0.682689$; the $0.2$ quantile is
        $3+2(-0.841621)=1.316758$. All three agree with
        \texttt{scipy.stats.norm} to six digits.

  \item $\sum_{k=0}^{40}e^{-2.5}2.5^{k}/k!=1-5.39\times10^{-35}$, that is $1$
        to every printed digit. $P(X\ge8)=0.004247$, whereas the binomial with
        $n=1000$, $p=0.0025$ gives $0.004191$; the binomial is smaller because
        its variance is $2.4938<2.5$ and it is bounded above by $n$.

  \item $P(X>s+t\mid X>s)=e^{-\lambda(s+t)}/e^{-\lambda s}=e^{-\lambda t}$.
        Geometric: $P(X>k)=q^{k}$, so the ratio is $q^{s+t}/q^{s}=q^{t}$. Uniform on
        the integers fails because the tail probability runs out, and the ratio depends
        on $s$.

  \item $f_{Y}(y)=\frac{1}{4\sqrt{y}}$ for $0<y<1$ (two preimages) and
        $\frac{1}{8\sqrt{y}}$ for $1\le y<9$ (one preimage). The integral is
        $\tfrac12+\tfrac12=1$. $\mathbb{E}[Y]=\frac14\int_{-1}^{3}x^{2}\,dx=7/3
        =2.3333$, the same as $\mathbb{E}[X^{2}]$, and via $f_{Y}$:
        $\int_{0}^{1}\frac{y}{4\sqrt{y}}\,dy+\int_{1}^{9}\frac{y}{8\sqrt{y}}\,dy
        =\tfrac16+\tfrac{13}6=\tfrac73$.

  \item Take $X$ uniform on $\{-1,0,1\}$ and $Y=|X|$. Then
        $\mathbb{E}[XY]=\mathbb{E}[X]=0$, so $\operatorname{Cov}=0$ and $\rho=0$.
        But $P(X=0,Y=0)=1/3$ whereas $P(X=0)P(Y=0)=1/9$: they are not independent.

  \item 
$\operatorname{se}(\bar X)=\sqrt{p(1-p)/(nR)}\le0{,}005$ gives
        $R\ge0{,}2275/(20\cdot0{,}000025)=455$. For the variance,
        $\operatorname{se}(\bar s^{2})=\sigma^{2}\sqrt{2/((n-1)R)}$ (normal
        approximation) with $\sigma^{2}=4{,}55$ and $n=20$, so
        $R\ge2\cdot4{,}55^{2}/(19\cdot0{,}02^{2})=5\,448{,}03$, that is $R=5\,449$.
        Using $2\cdot4{,}55^{2}/0{,}02^{2}$ alone drops the factor $n-1$ and gives a
        figure nineteen times too large.
  \end{enumerate}

\section*{Chapter 19: The Law of Large Numbers and the Central Limit Theorem}
\begin{enumerate}[label=\textbf{19.\arabic*}]
  \item $\{|(X_{n}+Y_{n})-(X+Y)|>\varepsilon\}$ is contained in
        $\{|X_{n}-X|>\varepsilon/2\}\cup\{|Y_{n}-Y|>\varepsilon/2\}$. The union bound
        gives a probability of at most the sum of two terms, each going to $0$.
        The $\varepsilon/2$ is needed so that two errors both below it add up to
        less than $\varepsilon$.

  \item $\mathbb{E}[X]=0$ and $\mathbb{E}[X^{2}]=2a^{2}\cdot\sigma^{2}/(2a^{2})=\sigma^{2}$,
        so $\operatorname{Var}(X)=\sigma^{2}$. Then
        $P(|X|\ge a)=\sigma^{2}/a^{2}$, exactly equal to the Chebyshev bound: the bound
        is attained, so it is sharp.

  \item False. The central limit theorem demands finite variance and iid observations; it
        also speaks only about the limit, not about $n=30$. In the lab section, the
        exponential with $n=30$ still gives $D=0.0219$ with $p=0.000$, and Cauchy
        gives a median $|\bar X_{n}|$ that settles near $1$.

  \item Chebyshev: $n\ge\sigma^{2}/(\varepsilon^{2}\delta)=0.25/(0.0025\cdot0.05)
        =2000$. The exact binomial computation gives $n=371$ (tail $0.0484$). The
        Chebyshev bound is five times looser.

  \item $\mathbb{E}[X_{n}]=n^{2}\cdot(1/n)=n\to\infty$, whereas for fixed
        $\varepsilon$ we have $P(|X_{n}|>\varepsilon)=1/n\to0$. So $X_{n}\to0$ in
        probability but its expectation explodes: no domination is available.

  \item Exact: $0.069545$. Without the correction, $z=(50-60)/\sqrt{42}=-1.5430$ gives
        $0.061411$, error $0.008134$. With the correction at $50.5$:
        $0.071340$, error $0.001795$. The continuity correction improves it fourfold.

  \item Median $|\bar X_{n}|$: $1.004$; $1.008$; $1.040$; $0.973$ for
        $n=10$ up to $10^{4}$. $P(|\bar X_{n}|>10)$: $0.0630$; $0.0680$;
        $0.0653$; $0.0580$, theory $0.0635$. The standard deviation $33.0$; $863.6$;
        $23.9$; $379.5$ does not shrink because the Cauchy variance is infinite.

  \item $\varphi_{\bar X_{n}}(t)=\bigl[\varphi(t/n)\bigr]^{n}
        =\bigl[e^{-|t|/n}\bigr]^{n}=e^{-|t|}$, so the average is again Cauchy for
        every $n$. Hence the se stays as large as the scale of the data itself, and
        ``mean $\pm$ two se'' does not narrow even as $n$ grows.

  \item $n=(1.96\cdot0.491971/0.0005)^{2}=3\,719\,214$; for half-width
        $0.001$ it is enough to take $929\,803$. The ratio is $4$, because $n\propto h^{-2}$:
        halve the width and the sample size rises fourfold.

  \item For $t<c$ and $\varepsilon=c-t$:
        $P(X_{n}\le t)\le P(X_{n}\le c-\varepsilon)+P(|X_{n}-c|>\varepsilon)\to0$;
        for $t>c$ use the complement. The constant condition cannot be dropped:
        take $X$ Rademacher and $X_{n}=(-1)^{n}X$. Every $X_{n}$ is still
        Rademacher (because $X$ is symmetric), so $X_{n}\rightsquigarrow X$; but for
        odd $n$ we have $|X_{n}-X|=2$ with probability $1$, so $X_{n}$ does not
        converge to $X$ in probability.
\end{enumerate}

\section*{Chapter 20: Estimation, Confidence Intervals, and Hypothesis Testing}
\begin{enumerate}[label=\textbf{20.\arabic*}]
  \item False. A p-value of $0.42$ means the data are not enough to contradict $p=0.5$;
        it does not prove $p=0.5$. The supplementary computations: a $95$ percent
        interval, for instance $(0.28;\,0.57)$, and the power of the test at the
        difference considered to matter.

  \item $\mathbb{E}[\hat p(1-\hat p)]=p-\bigl(p^{2}+p(1-p)/n\bigr)
        =p(1-p)(n-1)/n$. The unbiased estimator is
        $\frac{n}{n-1}\hat p(1-\hat p)$. For $p=0.3$, $n=20$: the expectation is
        $0.21\cdot19/20=0.1995$, and that estimator brings it back to $0.21$.

  \item Write $\hat\theta_{n}-\theta=\operatorname{bias}_{n}
        +(\hat\theta_{n}-\mathbb{E}\hat\theta_{n})$. The first term goes to $0$; the
        second goes to $0$ in probability by Chebyshev, because its variance goes to
        $0$. The MSE $=\operatorname{bias}^{2}+\operatorname{Var}$ going to $0$ is stronger,
        but only the two conditions above are used.

  \item 
$7/40$: Wilson $(0{,}0875;\,0{,}3195)$, width $0{,}2320$; Wald
        $(0{,}0572;\,0{,}2928)$, width $0{,}2355$. $13/20$: Wilson width $0{,}3860$,
        Wald $0.4181$. The $7/40$ interval is narrower, but that is not a sign of
        higher precision. Partly it is $n=40$; partly it is $\hat p=0.175$ being far
        from $0.5$, which shrinks $\hat p(1-\hat p)$: at $n=40$ with $\hat p=0.5$ the
        Wilson width is $0.2960$, still wider than $0.2320$. (At equal $n$, $7/20$ and
        $13/20$ both give $0.3860$, since $\hat p(1-\hat p)$ is symmetric about
        $1/2$.) What settles precision is not width but coverage: the Wald coverage
        at $p=0.175$, $n=40$ is only $0.921$, whereas Wilson gives $0.943$. A narrow
        interval near the edge is paid for with coverage below $95$ percent.

  \item $|k-15|>1.96\sqrt{7.5}=5.3676$, so reject when $k\le9$ or $k\ge21$.
        Its exact size is $0.042774$; its power at $p=0.7$ is only $0.588816$.
        With $n=100$ (reject when $k\le40$ or $k\ge60$) the power rises to $0.987502$.

  \item Wald: $|\hat p-p_{0}|\le z\sqrt{\hat p(1-\hat p)/n}$, which is equivalent to
        $p_{0}$ lying in $\hat p\pm z\sqrt{\hat p(1-\hat p)/n}$. A test that uses
        $\sqrt{p_{0}(1-p_{0})/n}$ uses a different denominator, so the set of $p_{0}$
        that are not rejected is not that Wald interval.

  \item 
Suppose the observed difference is $0.87$ hours with standard error $0.86$ hours.
``At the $5$ percent level the data do not reject equality of the two methods
        ($p=0.31$). The $95$ percent interval for the difference is
        $(-0.82;\,2.56)$ hours, so a $2$-hour difference is still a plausible value
        and even a $0.5$-hour difference cannot be ruled out. The power of this test
        against a difference of $0.5$ hours is only $0.09$: absence of evidence is
        not evidence of absence, and the honest conclusion is that these data do not
        decide the question.''

  \item Rule of three: $3/100=0.03$. The exact one-sided bound is
        $1-0.05^{1/100}=0.029513$; the Wilson upper bound is $0.036993$. The three
        differ because the rule of three is only an approximation and Wilson is two-sided.
        For the statement ``at most $5$ percent'' to be valid at the $95$ percent level:
        the exact bound is $0.0295$.

  \item $\bar x=19.125$, $s=1.082191$, $\operatorname{se}=0.312401$;
        $t_{11;0.95}=1.795885$ gives $(18.5640;\,19.6860)$ with width
        $1.1221$, whereas $95$ percent gives width $1.3752$. The width scales as
        $1/\sqrt{n}$, so halving it needs $n=48$.
\end{enumerate}

\section*{Chapter 21: Regression and Applied Statistics}
\begin{enumerate}[label=\textbf{21.\arabic*}]
  \item Not enough evidence. $R^{2}=0.99$ says only that the variance is explained;
        $s$ determines the error. On the cafe data $R^{2}=0.992821$ but
        $s=0.365148$ hours $=21.9$ minutes, and the half-width of the $95$ percent
        prediction interval at $x_{0}=5$ is $0.9101$ hours, that is $55$ minutes. What
        should be quoted: $s$ and a prediction interval.

  \item The normal equations are $X^{\top}e=0$. The first column of $X$ holds ones, so
        $\sum e_{i}=0$; the second column holds $x_{i}$, so $\sum x_{i}e_{i}=0$.
        Without the ones column (a line through the origin), $\sum e_{i}=0$ is no longer
        imposed and generally fails.

  \item $\hat b-\beta=\sum_{i}(x_{i}-\bar x)\varepsilon_{i}/S_{xx}$. The terms are
        independent and $\operatorname{Var}(\varepsilon_{i})=\sigma^{2}$, so
        $\operatorname{Var}(\hat b)=\sigma^{2}\sum(x_{i}-\bar x)^{2}/S_{xx}^{2}
        =\sigma^{2}/S_{xx}$. No matrix formula is used.

  \item $\hat b=0.7$, $\hat a=1.3$, $S_{xx}=10$, $s=0.516398$,
        $\operatorname{se}(\hat b)=0.163299$; with $t_{3}=3.182446$ the interval is
        $(0.1803;\,1.2197)$. It is wider than for the cafe data because $n=5$,
        $S_{xx}$ is small, and $s$ is large.

  \item 
In the cafe data $\bar x=5$, $n=9$, $S_{xx}=60$, so the leverage is
        $c(x_{0})=1/9+(x_{0}-5)^{2}/60$: $c(2)=0.261111$ and $c(9)=0.377778$.
        With $2.364624\sqrt{0.133333\,(1+c)}$ the half-widths are $0.9696$ for
        $x_{0}=2$ and $1.0135$ for $x_{0}=9$; they differ because
        $(x_{0}-\bar x)^{2}/S_{xx}$ is larger farther from the centre. That term
        shrinks like $1/n$, so the difference between the two intervals does vanish
        for large $n$: both tend to $1.96\sigma=0.716$ half-width. What never
        vanishes is the width itself, because the $+1$ term (the error of one new
        observation) remains: a prediction interval does not shrink to zero, it
        shrinks to the spread of the error.

  \item The column space of the big model contains the column space of the small model, so
        its projection is closest: the SSE cannot rise. $R^{2}$ may stay the same, but the
        divisor $n-p$ shrinks, so $R^{2}_{\text{adj}}$ falls: from $0.991795$ ($p=2$) to
        $0.990428$ ($p=3$).

  \item $\mathrm{VIF}=1/(1-0.81)=5.263158$, so
        $\operatorname{se}(\hat b_{1})$ grows by a factor $\sqrt{5.263158}$. With
        $x_{2}$: $\sqrt{2/(50\cdot0.19)}=0.458831$. Without $x_{2}$:
        $\sqrt{2/50}=0.2$. The ratio is $2.294157$.

  \item 
Group A: $(1;2)$, $(2;1{,}8)$, $(3;1{,}6)$; group B: $(5;4{,}0)$,
        $(6;3{,}8)$, $(7;3{,}6)$. Each group has slope $-0{,}2$, the pooled slope is
        $+0{,}4$. The valid decomposition adds contributions, not slopes:
        $b_{\text{gab}}=(S_{xy}^{\text{dalam}}+S_{xy}^{\text{antar}})/S_{xx}=-0{,}8/28+12/28=-0{,}028571+0{,}428571=0{,}4$,
        whereas the between-group slope on its own is $12/24=0{,}5$.

  \item ``Of the twelve models tried, the best gives $p=0.03$. After
        correcting for twelve comparisons (Bonferroni: $0.36$) the effect is no longer
        significant; therefore these data have not decided.''
\end{enumerate}

\section*{Chapter 22: Floating-Point Arithmetic and Error Analysis}
\begin{enumerate}[label=\textbf{22.\arabic*}]
  \item $\operatorname{ulp}(1)=2^{-52}=2.220\times10^{-16}$;
        $\operatorname{ulp}(1024)=2^{-42}=2.274\times10^{-13}$;
        $\operatorname{ulp}(0.1)=2^{-56}=1.388\times10^{-17}$;
        $\operatorname{ulp}(3\times10^{8})=2^{-24}=5.960\times10^{-8}$.
        $|x|u$ differs from the ulp by at most a factor of two, because the ulp is constant
        inside one binade. In $[1,2)$ there are $2^{52}$ numbers.

  \item $2^{-53}$ is exactly half an ulp of $1$, and rounding to even sends it back
        to $1$: $1+2^{-53}=1$. Since $2^{-52}$ is one full ulp,
        $1+2^{-52}>1$. In binary32, $u=2^{-24}$ and ulp$(1)=2^{-23}$, so the smallest
        positive number is $2^{-23}=1.192\times10^{-7}$.

  \item $\operatorname{fl}(x)$ is the nearest number, so
        $|\operatorname{fl}(x)-x|\le\tfrac12\operatorname{ulp}(x)$ by the definition of
        rounding. Equality cannot hold for all $x$: when $x$ is itself representable the
        error is $0$, and $\tfrac12$ ulp is attained only halfway between two numbers.

  \item False. Every operation is rounded correctly, but the errors accumulate along the
        algorithm. Example from this chapter: the sequential sum $1+10^{6}\cdot10^{-16}$
        stays $1$, whereas Kahan gives $1.0000000001$; the difference $10^{-10}$
        is far above $u$.

  \item With denominator $2^{54}$: $(0.1+0.2)+0.3=10808639105689192/2^{54}$ and
        $0.1+(0.2+0.3)=10808639105689190/2^{54}$. The difference is $2/2^{54}=2^{-53}$,
        exactly one ulp. \texttt{fsum} gives $0.6$; \texttt{float.hex} differs in the
        last digit.

  \item The naive formula gives $0.4996003611$; the double-angle
        identity gives $0.4999999999999996$ and $0.5-x^{2}/24+x^{4}/720$
        agrees. Its relative error is $-7.99\times10^{-4}$ and comes from
        cancellation in $1-\cos x$: $\cos x$ discards the digits the difference needs.

  \item Naive: the small root is $-1.0000000111\times10^{-4}$. The stable form
        $-2c/\bigl(b+\sqrt{b^{2}-4ac}\bigr)$ gives
        $-1.0000000100\times10^{-4}$, relative error $1.12\times10^{-9}$.
        The condition number of the subtraction is $\bigl(2\cdot10^{4}\bigr)/\bigl(2\cdot
        10^{-4}\bigr)=10^{8}$, so about eight digits are lost.

  \item At level $j$ there are $2^{k-j}$ additions, each with error at most
        $u$ times the sum of the absolute values of its inputs; summed over the whole tree,
        the error is $\le ku\sum|x_{i}|$ with $k=\log_{2}n$. That tree depth is what
        replaces $n$ in the bound for the sequential sum.

  \item Sequential: every $10^{-16}$ is below $\tfrac12\operatorname{ulp}(1)
        =1.11\times10^{-16}$, so the result stays $1$. Kahan:
        $1.0000000001$, the same as the exact value. Each Kahan step is also rounded,
        but the lost remainder is kept in $c$ and returned on the next step.

  \item $\kappa\approx1+\bar x^{2}/s^{2}$. For $1000$ values
        $999+\bigl((k\bmod7)-3\bigr)/4$ with $k=1,\dots,1000$:
        $\bar x=999.00075$, $s=0.4997$ (population standard deviation, matching
        the $1/n$ divisor in the problem),
        $\kappa\approx4.0\times10^{6}$. With $u_{32}=5.96\times10^{-8}$,
        $\kappa u\approx0.238$: almost all digits are lost. Measured: one pass
        errors by $25$ percent, two passes by $4.4\times10^{-8}$.
\end{enumerate}

\section*{Chapter 23: Convex Optimization and Duality}
\begin{enumerate}[label=\textbf{23.\arabic*}]
  \item $f$: Hessian $\begin{pmatrix}2&1\\1&2\end{pmatrix}$, eigenvalues $1$ and $3$,
        positive definite, strictly convex. $g$: $\begin{pmatrix}0&1\\1&0\end{pmatrix}$,
        eigenvalues $\pm1$, indefinite, neither of the two. $h$:
        $\begin{pmatrix}4&-4\\-4&6\end{pmatrix}$, eigenvalues $0.8769$ and $9.1231$,
        positive definite. None is singular semidefinite; for $x_{1}^{2}+2x_{1}x_{2}
        +x_{2}^{2}$ (eigenvalues $0$ and $4$) the minimum is not unique.

  \item For $x,y\in C_{1}\cap C_{2}$ and $t\in[0,1]$, the mixed point lies in
        $C_{1}$ and in $C_{2}$, hence in their intersection. Union: $[0,1]\cup[2,3]$ contains
        $1$ and $2$ but not $(1+2)/2=1.5$.

  \item $f_{i}(tx+(1-t)y)\le tf_{i}(x)+(1-t)f_{i}(y)\le tM(x)+(1-t)M(y)$ for each
        $i$, then take the maximum. For $\max\{x_{1}^{2},x_{2}\}$ at $(0,0)$ the gradient
        does not exist (a crease along $x_{2}=x_{1}^{2}$), but its subdifferential
        $\operatorname{conv}\{(0,0),(0,1)\}$ contains $\mathbf{0}$: a minimum. At $(1,1)$
        the subdifferential is $\operatorname{conv}\{(2,0),(0,1)\}$ with no zero, not a
        minimum.

  \item Without $x_{1}\ge1$: $x_{1}=0.8$, $x_{2}=1.6$, which violates it. So $x_{1}=1$,
        $x_{2}=1.5$, $p^{*}=3.25$. The multipliers are $\nu=-1.5$ for the equality and
        $\lambda=0.5\ge0$ for the inequality. \texttt{scipy.optimize.mini\-mize}
        gives the same value.

  \item $g(\lambda)=\min_{x}\bigl[x_{1}^{2}+x_{2}^{2}+\lambda(2-x_{1}-x_{2})\bigr]
        =2\lambda-\lambda^{2}/2$. Its derivative is $2-\lambda$, so the maximum is at
        $\lambda=2$ with $d^{*}=2=p^{*}$. The gap is zero because the constraint is affine
        and $x=(1.5;\,1.5)$ is strictly feasible: Slater holds.

  \item False. Strong duality only equates the values, it does not guarantee attainment.
        Example: minimize $e^{x}$ over $\mathbb{R}$. $p^{*}=0$ but it is not attained;
        the dual side is also $0$, so the gap is zero while the primal optimum does not
        exist.

  \item False. Complementary slackness states $\lambda_{i}g_{i}(\mathbf{x}^{*})=0$:
        an active constraint may have a zero multiplier, and an inactive constraint must
        have a zero multiplier. The implication runs one way, so ``active $\Rightarrow$
        positive'' is not valid.

  \item Without the bound, $(x_{1}-2)^{2}+(x_{2}-2)^{2}$ on $x_{1}+x_{2}=2$ is minimized
        at $x_{1}=x_{2}=1$. That satisfies $x_{1}\le c$ if and only if $c\ge1$, so
        at $c=1$ the multiplier is $\lambda=0$ and $p^{*}=2$; the constraint starts to
        bind only for $c<1$ (at $c=0.5$: $p^{*}=2.5$, $\lambda=2$).

  \item For $\lambda>0$ and $x_{k}\to-\infty$, the term $\lambda x_{k}^{3}$ dominates
        and is negative, so $g(\lambda)=-\infty$. Only $\lambda=0$ is finite:
        $g(0)=\min(x-2)^{2}=0$, hence $d^{*}=0$. Meanwhile $p^{*}=1$ at $x=1$ (the piece
        $[0,1]$), because at $x=-1$ the value is $9$. The gap is $1$, one hundred percent.

  \item The first-order inequality for a convex function gives
        $f(y)\ge f(x^{*})+\nabla f(x^{*})^{\top}(y-x^{*})=f(x^{*})$ for every $y$,
        without using strict convexity. That proof fails for $f(x)=x^{3}$ because
        $f$ is not convex; the condition violated is precisely convexity of $f$, at $x=0$
        where $\nabla f=0$ but it is not a minimum.
\end{enumerate}

% Answer keys for Chapters 24--32. This file is \input directly inside the document.
% Every number is tested by an assert script outside the repo (see the report).

\section*{Chapter 24: Linear and Integer Programming}
\begin{enumerate}[label=\textbf{24.\arabic*}]
  \item Maximize $30x_{1}+20x_{2}$ subject to $2x_{1}+x_{2}\le10$, $x_{1}+2x_{2}\le8$,
        $x\ge0$. The vertices $(0;0)$, $(5;0)$, $(0;4)$, $(4;2)$ give profit
        $0;\,150;\,80;\,160$. The continuous optimum is $(4;2)$ with profit $160$; integer
        enumeration also gives $160$, so here rounding down happens to be exact and
        the gap is $0$.

  \item For $x,y\in P$ and $\lambda\in[0,1]$:
        $A(\lambda x+(1-\lambda)y)=\lambda Ax+(1-\lambda)Ay\le\lambda b+(1-\lambda)b=b$,
        using $A$ linear and $\lambda\ge0$. On integer sets that argument collapses:
        a convex combination of two integer points need not be integer, for example $0$
        and $1$ are feasible while their midpoint is not a member of $\mathbb{Z}^{n}$.

  \item Profit $x_{1}+2x_{2}$ at the vertices: $0$; $4$; $6$; $6$. The optima $(0;3)$ and
        $(3;1{,}5)$ both equal $6$, and the whole segment $x_{1}+2x_{2}=6$ is optimal because
        the objective is constant there; the midpoint $(1{,}5;\,2{,}25)$ is not a vertex.
        The dual is $y=(0;1)$: the price of an hour is $1$, of plywood $0$.

  \item Initial table with basis $s_{1},s_{2},s_{3}=(4;3;3)$. First pivot: $x_{1}$ enters
        (reduced cost $3$), ratios $4$ and $3$, so $s_{2}$ leaves; value $9$.
        Second pivot: $x_{2}$ enters, $s_{1}$ leaves (ratio $1$), optimum $(3;1)$
        with $z=11$. The dual prices from the final table: $(2;1;0)$, and
        $4(2)+3(1)+3(0)=11$.

  \item The dual minimizes $10y_{1}+8y_{2}$ subject to $2y_{1}+y_{2}\ge30$,
        $y_{1}+2y_{2}\ge20$, $y\ge0$; the solution $y=(40/3;\,10/3)$ has value
        $160$, the same as the primal. Meaning: one extra kg of flour is worth
        $13{,}333$ and one extra oven hour $3{,}333$.

  \item False. Rounding down is not always feasible once there are equality
        constraints or constraints with sign $\ge$. Example: maximize $x_{1}+x_{2}$ with
        $2x_{1}+2x_{2}=5$; the continuous answer $(2{,}5;\,0)$ rounds to
        $(2;0)$, which gives $4\ne5$ and is therefore infeasible.

  \item True for the forward direction, and only for optimal pairs. Complementary
        slackness $y_{i}(b_{i}-\mathbf{a}_{i}^{\top}\mathbf{x})=0$ gives:
        $y_{i}>0$ forces that constraint to be active. The converse does not hold; at a
        degenerate point there is an active constraint with $y_{i}=0$, so a zero price
        does not mean the constraint is slack.

  \item Densities $2{,}0$; $1{,}75$; $1{,}667$; $1{,}5$. The first item fills up
        ($5$ kg, $10$), then $3/4$ of the second, total weight $8$: the LP bound
        is $15{,}25$. Enumeration gives $15$ (items $1$ and $3$). Gap $0{,}25$,
        that is $1{,}639$ percent of the bound.

  \item 
Capacity $7$: the first item fills the sack, then $1/2$ of the second, so
        the LP bound is $13{,}5$; the best integer answer is $13$ (items $1$ and
        $4$, weight $7$). The gap is $0{,}5$, i.e. $3{,}704$ percent. The gap
        changes because capacity $7$ rules out the pair $(1,3)$, whose weight
        $5+3=8$ used to fit, so the best integer answer moves to items $1$ and $4$;
        the fractional and integer optima do not move in step.

  \item $x_{v}=1/2$ for every $v$ is feasible ($x_{u}+x_{v}=1$) and gives $n/2$;
        every cover must contain $n-1$ vertices, since leaving out two vertices
        means the edge between them is uncovered. Ratio of LP to integer:
        $0{,}6667$; $0{,}625$; $0{,}5556$ for $n=4,5,10$. The ratio tends to
        $1/2$, so the relaxation bound can be off by a factor of two and rounding is not
        a general method.
\end{enumerate}

\section*{Chapter 25: Neural Networks and Backpropagation}
\begin{enumerate}[label=\textbf{25.\arabic*}]
  \item The bound is $\lfloor1/\gamma^{2}\rfloor$: for $\gamma=0{,}2$ we get $25$,
        for $\gamma=0{,}05$ we get $400$. That bound only counts mistakes,
        and each mistake consumes margin, which is why $n$ does not appear. For
        small $\gamma$ the bound blows up like $1/\gamma^{2}$ and is therefore useless.

  \item Write $\mathbf{W}=\mathbf{W}_{2}\mathbf{W}_{1}$ and
        $\mathbf{b}=\mathbf{W}_{2}\mathbf{b}_{1}+\mathbf{b}_{2}$; then the output is
        $\mathbf{W}\mathbf{x}+\mathbf{b}$, that is, one linear threshold. Adding the next
        identity layer only multiplies matrices, so the capacity does not
        increase at all.

  \item With $\mathbf{x}=(0{,}5;-1;\,0{,}25;2)$ and $y=0$: $\mathbf{z}_{1}=(2{,}9125;\,
        0{,}375)$, $a=0{,}8961350845$, $\bar{z}_{2}=0{,}8961350845$, hence
        $\bar{\mathbf{z}}_{1}=(0{,}7169081;\,-0{,}5376810)$ and
        \[
          \nabla_{\mathbf{W}_{1}}L=
          \begin{pmatrix}
            0{,}3584540 & -0{,}7169081 & 0{,}1792270 & 1{,}4338161\\
            -0{,}2688405 & 0{,}5376810 & -0{,}1344203 & -1{,}0753621
          \end{pmatrix},
        \]
        that is, the second row equals $-0{,}75$ times the first row. This item is
        ambiguous: both gates are open, so there is no row with zero gradient.

  \item Not enough evidence. Finite differences only check thirteen scalars at one
        point, and at $h=10^{-6}$ an error of $10^{-10}$ is exactly what is expected; a
        shared bug in both computations would not show up. What is needed: several points,
        both gate states, the cases $y=0$ and $y=1$, gradients with respect to the input, and
        an automatic differentiation comparison.

  \item False. The theorem is about density of continuous functions on a
        compact domain: there exist parameters that approximate, not that those parameters
        are reachable by optimization. Two things are not stated: how many units are
        needed, and the error on test data (generalization), not only on the target
        function.

  \item At $t=1$, $\hat m=g$ and $\hat v=g^{2}$, so the step is
        $\eta g/(|g|+\epsilon)=-0{,}01$ and $b_{2}$ becomes $0{,}06$. With a gradient
        $100$ times larger the step stays $-0{,}01$, because Adam only
        uses the ratio to the gradient scale; $b_{2}$ stays $0{,}06$.

  \item 
The output $w_{2}\max(\mathbf{W}_{1}\mathbf{x}+\mathbf{b}_{1},0)$ is a sum
        of two linear functions, each active on one half-plane, hence piecewise
        linear; two lines $z_{1}[i]=0$ cut the plane into at most four cells, each
        with its own linear formula. For $k$ units in general position, $k$ lines
        cut the plane into $1+k+k(k-1)/2$ regions, that is $2,4,7$ for $k=1,2,3$
        (not $8$: the third line adds only three regions, not four).

  \item $\nabla_{\mathbf{x}}L=(-0{,}0868370;\,0{,}1328096;\,-0{,}0970532;\,
        -0{,}1430257)$, so $\boldsymbol{\delta}=0{,}05\,\mathrm{sign}(\nabla)
        =(-0{,}05;0{,}05;-0{,}05;-0{,}05)$. The loss at $\mathbf{x}+\boldsymbol{\delta}$
        is $0{,}3186800$, whereas the linear prediction
        $0{,}2949123+0{,}0229863=0{,}3178985$; a difference of $7{,}81\times10^{-4}$ from
        curvature.

  \item 
Prediction: the second lab block uses ReLU, and ReLU saturates only on the
        negative side. Saturation table ($8$ units, random He $\mathbf{W}_{1}$,
        $\mathbf{b}_{2}=0$):
        \begin{center}
        \begin{tabular}{lrrrr}
        & range of $Z_{1}$ (initial) & dead units & final $L$ & accuracy \\ \hline
        $\mathbf{b}_{1}=0$   & $[-2.43;\,2.40]$   & $0/8$ & $0.0163$ & $1.0000$ \\
        $\mathbf{b}_{1}=+10$ & $[7.57;\,12.41]$   & $0/8$ & $0.3225$ & $0.8375$ \\
        $\mathbf{b}_{1}=-10$ & $[-12.43;\,-7.60]$ & $8/8$ & $0.6909$ & $0.5333$
        \end{tabular}
        \end{center}
        So $\mathbf{b}_{1}=+10$ indeed does not kill the hidden units ($Z_{1}$ far
        above zero passes through unchanged), but it is not free: $H\cdot\mathbf{w}_{2}$
        becomes very large, the output sigmoid sticks at $0$ or $1$, and $L$ starts
        at $4.04$; after $3000$ steps it stops at $L=0.3225$ with accuracy $0.8375$,
        worse than $\mathbf{b}_{1}=0$. For $\mathbf{b}_{1}=-10$ every unit is dead
        from the start ($Z_{1}$ below zero, the derivative exactly $0$, not merely
        small), $\partial L/\partial b_{1}=\mathbf{0}$ exactly, and training stops at
        $L=0.6909$ -- the label entropy -- with majority-class accuracy $0.5333$.
  \end{enumerate}

\section*{Chapter 26: Generalization, Regularization, and Training}
\begin{enumerate}[label=\textbf{26.\arabic*}]
  \item Not enough evidence, and what is wrong is equating training error with
        generalization error. Zero training error only means the training data has been
        memorized. The additional numbers needed: validation (or test) error on data not
        used for training, together with its interval width.

  \item With $\hat f=\hat f_{D}$ and $\bar f=\mathbb{E}_{D}\hat f$:
        $\mathbb{E}(y-\hat f)^{2}=\sigma^{2}+(\bar f-f)^{2}+\operatorname{Var}(\hat f)$.
        The cross term vanishes because $\mathbb{E}[\varepsilon]=0$ and $\varepsilon$
        is independent of $D$, so $\mathbb{E}[\varepsilon(\hat f-f)]=0$; independence
        is used precisely at this step, on $\varepsilon$, not on $\hat f$. The second
        cross term vanishes for a different reason: $\mathbb{E}[\hat f-\bar f]=0$ by
        the definition of $\bar f$, so $2(f-\bar f)\,\mathbb{E}[\hat f-\bar f]=0$
        without needing any independence.

  \item From $\frac{d}{dw}\bigl[\frac12\sum(y_{i}-wx_{i})^{2}\bigr]=0$ we get
        $w(\lambda)=\sum x_{i}y_{i}/(\sum x_{i}^{2}+\lambda)=64/(30+\lambda)$.
        Its derivative $-64/(30+\lambda)^{2}<0$ but never zero, so $w$ shrinks
        without touching zero. In L1, the subgradient at $w=0$ contains zero when
        $\lambda\ge|\sum x_{i}y_{i}|=64$; so $w=0$ exactly at $\lambda=64$.

  \item $\operatorname{MSE}(c)=(c-1)^{2}\mu^{2}+c^{2}\sigma^{2}/n$: for
        $c=0{,}95$ we get $0{,}01+0{,}0361=0{,}0461$; for $c=1{,}05$ we get
        $0{,}01+0{,}0441=0{,}0541$; for $c=1$ we get $0{,}04$. An equally large deviation
        is not an equally large loss, because the variance term grows like $c^{2}$, so
        exceeding the budget costs more than shrinking.

  \item Fix: that cross-validation number is an estimate of the error of the selection
        procedure, not of the selected model's error in the field, and it is
        optimistically biased because $\lambda$ was chosen on the same folds. Two
        things are needed: nested cross-validation (or a test set never seen during
        selection), and repetition with several partitions along with intervals.

  \item For $a<0$: with $u<0$ the derivative is $u-a-t$, zero at $u=a+t$, which is valid
        only when $a<-t$. For $u=0$ the subgradient is $[-a-t,\,-a+t]$, which contains
        zero when $|a|\le t$, so $u=0$ is optimal. Symmetric because the objective
        $\frac12(u-a)^{2}+t|u|$ is unchanged when $u$ and $a$ are both sign-flipped.

  \item The half-life is $\ln2/(\eta\lambda)=0{,}693147/0{,}002=346{,}57$ steps;
        a quarter of the initial value means two half-lives, $693{,}15$ steps, and
        $(1-\eta\lambda)^{693}=0{,}24973$. A numpy simulation with zero gradient
        gives the same numbers, because $w_{k}=w_{0}(1-\eta\lambda)^{k}$.

  \item The gap is $0{,}2406-0{,}1500=0{,}0906$; at $\lambda=0$ the gap is
        $0{,}3262-0=0{,}3262$, at $\lambda=10$ only $0{,}3395-0{,}3333=0{,}0062$.
        A small gap is not enough: at $\lambda=10$ the gap is smallest precisely because
        both are equally bad, close to the majority-class error $0{,}3419$.

  \item With $m$ hypotheses, $P(\exists i:|\hat R_{i}-R_{i}|>\epsilon)\le
        2m e^{-2n\epsilon^{2}}$, so $\epsilon=\sqrt{(\ln m+\ln(2/\delta))/(2n)}$.
        For $m=2^{10}$, $n=200$, $\delta=0{,}01$: $\epsilon=0{,}1749$. Halving it
        demands four times the $n$, namely $800$.
\end{enumerate}

\section*{Chapter 27: Honest Model Evaluation}
\begin{enumerate}[label=\textbf{27.\arabic*}]
  \item Not enough evidence. A difference of $0{,}4$ points on a single test set is not
        necessarily beyond noise; what is wrong is treating two paired proportions
        as independent. Two additional numbers: the size of the test set, and the count of
        paired disagreements (for McNemar's test) or the interval width.

  \item For hard labels $s,y\in\{0,1\}$ we have $(s-y)^{2}=1$ exactly at
        mistakes and $0$ at correct predictions, so the average is
        $(\mathrm{FP}+\mathrm{FN})/n$. For scores that are probabilities, $(s-y)^{2}$
        takes intermediate values, so the average is no longer a count of mistakes
        but a calibration penalty.

  \item $\mathbb{E}[(s-y)^{2}]=s^{2}-2s\pi+\pi$ with $\pi=\mathbb{E}[y]$.
        The derivative is $2s-2\pi$, zero at $s=\pi$; the minimum value is
        $\pi^{2}-2\pi^{2}+\pi=\pi(1-\pi)$, and the second derivative $2>0$ confirms
        that it is a minimum.

  \item Precision $180/225=0{,}8$; recall $180/200=0{,}9$; $F_{1}=0{,}847059$;
        accuracy $1935/2000=0{,}9675$; specificity $1755/1800=0{,}975$; balanced accuracy
        $(0{,}9+0{,}975)/2=0{,}9375$. \texttt{sklearn.metrics} gives the
        same numbers to six decimals.

  \item Wilson $95/100$: $[0{,}8882;\,0{,}9785]$, width $0{,}0902$; for
        $950/1000$: $[0{,}9347;\,0{,}9619]$, width $0{,}0272$, that is $3{,}318$
        times narrower even though $n$ is ten times larger. Width is not a measure of
        model quality: it depends on $n$ and on the data, and the two experiments are
        not paired.

  \item 
Out of $25$ pairs, $19$ are correct: the positive $0{,}95$ beats five
        negatives, $0{,}88$ also five (it sits above $0{,}75$), $0{,}70$ four,
        $0{,}60$ three, $0{,}40$ two. Hence $\mathrm{AUC}=19/25=0{,}76$;
        \texttt{roc\_auc\_score} reports $0{,}76$.

  \item Precision $=0{,}8\pi/(0{,}8\pi+0{,}01(1-\pi))$: at $\pi=0{,}002$ we get
        $0{,}138169$, at $\pi=0{,}2$ we get $0{,}952381$. The requirement precision
        $\ge0{,}5$ gives $0{,}4\pi\ge0{,}005-0{,}005\pi$, that is
        $\pi\ge0{,}005/0{,}405=0{,}012346$.

  \item Split by host, not by record: $60/20/20$ becomes $180$ training hosts,
        $60$ validation, $60$ test, each about $20$ records. Preprocessing is
        computed inside the training fold only. The choice among $40$ settings goes
        through nested cross-validation in the training part. The test set is touched
        once, at the end.

  \item False. AUC measures ranking, not usefulness at a particular threshold. With
        sensitivity $0{,}8$ and specificity $0{,}99$, at prevalence $0{,}002$ the
        precision is only $0{,}138169$: of ten alarms, $8{,}62$ are wrong, so about
        nine. The chosen
        threshold still has to be checked for precision and recall.
\end{enumerate}

\section*{Chapter 28: Formal Methods with Lean 4}
\begin{enumerate}[label=\textbf{28.\arabic*}]
  \item False. A file that is accepted means every theorem has a proof that is
        checked, but \texttt{sorry} passes with a warning, and
        \texttt{axiom} as well as \texttt{native\_decide} extend the trusted base.
        The checking command: \texttt{\#print axioms} on every theorem that is
        claimed.

  \item 
Helper lemma: for every $k$, $k(2k)+2(2k+1)=(k+1)\cdot2(k+1)$. That goal
        contains $k\cdot k$, so it is not a linear goal: \texttt{omega} will not
        close it on its own and a rewrite is needed first. Then \texttt{induction n
        with}: base case \texttt{rfl}; the \texttt{succ k ih} step uses \texttt{rw
        [oddSum, Nat.mul\_add, ih]}, which leaves exactly that helper lemma, then
        \texttt{exact bantu k}. The helper lemma itself is closed by
        \texttt{simp only [Nat.mul\_add, Nat.add\_mul, Nat.mul\_one, Nat.one\_mul]}
        followed by \texttt{omega}: \texttt{omega} cannot close
        $k\cdot(2k)+2(2k+1)=(k+1)\cdot2(k+1)$ on its own, because that goal contains
        $k\cdot k$ and is not linear (tried: \texttt{omega could not prove the goal}).

  \item Informal: for $n=0$ both sides are $0$; if $\sum_{i<n}2i=n(n-1)$, then
        adding $2n$ gives $n(n-1)+2n=(n+1)n$. The Lean version defines
        \texttt{evenSum : Nat -> Nat} recursively (\texttt{evenSum 0 = 0},
        \texttt{evenSum (n+1) = evenSum n + 2*n}) and states the form without
        truncated subtraction, $E(n)+n=n\cdot n$. The steps: \texttt{induction},
        base \texttt{rfl}, then \texttt{rw [evenSum]}. Note that
        \texttt{rw [evenSum, ih]} fails outright (\texttt{did not find instance of
        the pattern evenSum k + k}), because the goal afterwards reads
        $E(k)+2k+(k+1)=(k+1)(k+1)$. Insert the regrouping
        \texttt{have h : evenSum k + 2 * k + (k + 1) = (evenSum k + k) + (2 * k + 1) := by omega}
        first, continue with \texttt{rw [h, ih]}, and close with
        \texttt{simp only [Nat.mul\_add, Nat.add\_mul, Nat.mul\_one, Nat.one\_mul]}
        then \texttt{omega}.

  \item The sum of the first $37$ odd numbers $=37^{2}=1369$; the sum of the first $250$
        $=250^{2}=62500$. Script: \texttt{sum(2*i+1 for i in range(n))} gives
        both numbers, and both are perfect squares.

  \item        A direct check for all $n\le1000$ demands $1001$ cases, namely $n=0$
        through $n=1000$; the bound $n\ge1$ would give $1000$, but the exercise
        does not say that, and the \texttt{zero} base case above really does check
        $n=0$. An induction proof generates only two goals, \texttt{zero} and
        \texttt{succ}, then concludes for all $n$; so $1001$ against $2$.

  \item \texttt{Nat.add} is defined by recursion on the second argument, so
        $n+0$ reduces straight to $n$ and \texttt{rfl} closes it. On $0+n$
        the recursion gets stuck on the unknown $n$, so there is nothing
        to reduce and \texttt{induction n} is needed.

  \item 
Weak: for fixed $x$, $\exists\delta,\ |\delta|\le1\wedge f(x+\delta)=f(x)$;
        the witness $\delta=0$ makes both sides equal and \texttt{rfl} closes it.
        Strong: $\forall x\,\forall\delta,\ |\delta|\le1\to f(x+\delta)=f(x)$, and
        that is false for the identity $f$ at $\delta=1$.

  \item 
\texttt{decide}: no axioms, pure kernel computation. \texttt{omega}:
        \texttt{propext} and \texttt{Quot.sound}, because it proves goals by
        contradiction inside the integers; \texttt{Classical.choice} does not
        appear, and this chapter prints the same list. \texttt{sorry}:
        \texttt{sorryAx}, the axiom added specifically to close a skipped goal.

  \item First mistake: \texttt{delta} is already free in the context and then bound again
        by \texttt{exists}, so the reader of the statement determines its value,
        not the attacker; robustness is not implied. Second: there is no quantifier over $x$,
        so $x$ is free too. Correct version:
        $\forall x,\ \exists\delta,\ \lVert\delta\rVert\le1\wedge f(x+\delta)=f(x)$.

  \item \texttt{intro h} gives $h:P\wedge\neg P$. Its components \texttt{h.1} have type
        $P$ and \texttt{h.2} have type $\neg P$; then \texttt{exact h.2 h.1} closes
        the goal, because $\neg P$ is $P\to$ false.
\end{enumerate}

\section*{Chapter 29: Adversarial Machine Learning: Attacks and Defenses}
\begin{enumerate}[label=\textbf{29.\arabic*}]
  \item $\mathbf{s}_{A}=1$, $\mathbf{s}_{B}=3$, so the margin is already $-2$ (the example
        leans toward $B$); $\mathbf{d}=(1;-2;1)$ with $\lVert\mathbf{d}\rVert_{1}=4$
        and $\lVert\mathbf{d}\rVert_{2}=\sqrt6\approx2{,}449$. Hence
        $\varepsilon^{*}_{\infty}=2/4=0{,}5$ and
        $\varepsilon^{*}_{2}=2/2{,}449=0{,}8165$; the ratio is $\sqrt{8/3}=1{,}633$.

  \item The derivative of the objective with respect to $\boldsymbol{\delta}$ is $\mathbf{w}$, so
        the maximum is $\varepsilon\lVert\mathbf{w}\rVert_{2}$, attained uniquely at
        $\boldsymbol{\delta}=\varepsilon\mathbf{w}/\lVert\mathbf{w}\rVert_{2}$
        (Cauchy--Schwarz, with equality when parallel). In
        $\ell_\infty$ the maximum is $\varepsilon\lVert\mathbf{w}\rVert_{1}$ at
        $\varepsilon\,\mathrm{sign}(\mathbf{w})$, a vertex; the form differs
        because the dual norm differs.

  \item The error $\mathbf{d}^{\top}\boldsymbol{\delta}$ is linear on a compact convex
        set, so the maximum is at an extreme point. The extreme points of the
        $\ell_\infty$ ball are $\pm\varepsilon$ in each coordinate, and FGSM picks
        exactly that point; hence the PGD path, which always stays inside the ball,
        cannot exceed it.

  \item False. In this chapter a random signed perturbation at $\varepsilon=0{,}15$ only
        drops accuracy to $95{,}48$ percent, whereas PGD 40 steps
        brings it down to $10{,}72$ percent. Robustness to random noise
        says nothing about an attack that chooses the direction.

  \item 
Untargeted success (PGD with 40 steps, clean accuracy as the reference) for
        $A$: $14{,}02$, $53{,}16$, $89{,}06$ percent at $\varepsilon=
        0{,}05;0{,}10;0{,}15$. For $D$: $5{,}38$, $15{,}28$, $35{,}06$ percent. Over
        that range the gap actually widens ($8{,}64$, $37{,}88$, $54{,}00$
        percentage points), because $A$ is already close to the $100$ percent
        ceiling while $D$ is still far; the narrowing only happens outside that
        range, at $\varepsilon\ge0{,}20$, once both models run out of margin.

  \item One coordinate suffices: projection onto an interval is clipping, and
        $[x_{j}-\varepsilon,x_{j}+\varepsilon]\cap[0,1]$ equals
        $[\max(0,x_{j}-\varepsilon),\ \min(1,x_{j}+\varepsilon)]$. Clipping to the
        $\varepsilon$ box and then to $[0,1]$, or the other way round, yields the
        same interval because interval intersection is commutative; the Cartesian product
        then gives the same point.

  \item $D$ loses $97{,}99-96{,}48=1{,}51$ points of clean accuracy to
        gain $62{,}65-10{,}72=51{,}93$ points of robustness, so the price is
        $1{,}51/51{,}93=0{,}0291$ clean points per point of robustness. That ratio is not
        complete: it depends on one $\varepsilon$, on one attacker that is
        not an upper bound, and it ignores the clean accuracy of the comparison model.

  \item Protocol: attack with PGD with more steps and several random starts,
        with EOT over the defense's randomness, then repeat with attacks without
        gradients (imitation/transfer, finite-difference estimation). Report the curve of accuracy
        against $\varepsilon$, not a single point. Easiest symptom: success
        jumps as soon as the gradient is replaced by BPDA or EOT, while naive attacks
        report zero.

  \item White box $89{,}06$ percent, transfer $1-0{,}2379/0{,}9849=75{,}85$
        percent, a difference of $13{,}21$ points. That is not evidence that $B$ is safer:
        an attack built specifically for $B$ will be stronger, and in this chapter the
        white-box versus transfer difference does shrink when the attacker is adapted.
        This item is ambiguous: the lab output does not contain a PGD row for model $B$.

  \item Untargeted only needs a label change; targeted must reach a particular
        class, so the set of targets that must be met is far smaller. At
        $\varepsilon=0{,}15$ the target is reached $33{,}00$ percent of the time, whereas
        untargeted is $89{,}06$ percent; with $\varepsilon=0{,}30$ the target is reached
        $99{,}50$ percent. So the difficulty is a matter of budget, not of attack type.
\end{enumerate}

\section*{Chapter 30: The Limits of Attack: Obfuscated Gradients and NIDS Evasion}
\begin{enumerate}[label=\textbf{30.\arabic*}]
  \item Three checks: (a) repeat with adaptive attacks (BPDA, EOT), (b)
        report the curve of accuracy against $\varepsilon$, not a single point, (c) check
        gradient quality (cosine against finite differences or against the weight direction). If
        what is happening is gradient masking: success collapses only for attacks
        based on gradients, while random search and transfer still succeed, and
        the curve is flat.

  \item For $f(x)=w\cdot x+b$ the gradient is constant and the Hessian zero, so the
        Lipschitz constant is $L=0$ and the bound $L\delta$ gives zero error: BPDA is exact. With two
        ReLU layers, $f$ is piecewise linear so the gradient jumps at the
        kinks and the Lipschitz-gradient hypothesis fails; the bound is then
        meaningless.

  \item 
As the program shows, the label flips up to $q=0{,}6$ (score $-1{,}2$) and
        fails at $q\ge0{,}70$, when every coordinate rounds to zero so the score is
        $0$, a tie rather than a flip. Hence
        $q^{*}=2\max_{j}|x_{j}-\varepsilon|=2(\varepsilon+|\min_{j}x_{j}|)=0{,}70$.
        Analytic safety bound: the rounding error is at most $4q$, so the attack is
        only guaranteed to succeed when $q<(8\varepsilon-m)/4=0{,}345$; at $q=0{,}345$
        the label still flips, so the direction of the inequality matters.

  \item Probability of rejection $=P(\text{lebih dari }10\text{ dari }20\text{ hilang})
        =0{,}411901$, so what is answered is $588$ of $1000$ queries: $588/20=29$
        PGD steps with $m=20$. With $m=1$ the rejection probability is $0{,}5$, so $500$
        steps. What is paid: the gradient becomes noisier and its direction further from the true one.

  \item What can be predicted without running: clean accuracy is about the same,
        naive success stays near zero, and BPDA success is near the baseline for $q=0$
        and $q=0{,}05$ but falls for large $q$. The cosine no longer approaches one
        because the ReLU gate adds kinks on top of rounding, so the BPDA approximation
        only approaches the direction. The numbers for the third value of $q$ must be taken from execution,
        not invented.

  \item Binary flags and protocol numbers are assumed not to be changeable by the attacker:
        $7+3=10$ locked features, the remaining $40$ features can really be shifted (out of
        $50$). A $32\times32$ image with three channels has $3072$ values, so
        $3072/40=76{,}8$ times more.

  \item Daily retraining changes the parameters that a gradient-based attack uses,
        so old directions go stale quickly; a blind attack does not use those parameters, so
        it keeps working. According to the EOT lemma what must be checked is the expectation over
        transformations, and an adaptive attack must be defined against that defense
        itself, including its retraining schedule.

  \item Metric: success of the adaptive attack against $\varepsilon$, with
        at least $5$ seeds and an interval. The curve shape that accepts (a): the curve
        is shifted to the right or downward systematically at all $\varepsilon$,
        not only at one point. If after BPDA and EOT the success returns
        to the baseline, then only the measurement was broken, so (b).
\end{enumerate}

\section*{Chapter 31: Cryptography: Modular Integers to the Quantum Threat}
\begin{enumerate}[label=\textbf{31.\arabic*}]
  \item False. $2^{100}\equiv1\pmod{101}$ only says the order of $2$ divides $100$,
        that is, $101$ passes the Fermat test to base $2$; pseudoprimes (and
        Carmichael numbers) also pass. What that congruence proves is only $101\nmid2$
        and $2^{100}\equiv1$, not the absence of divisors.

  \item If $\gcd(a,n)=\gcd(b,n)=1$, Bezout's theorem gives $ax+ny=1$ and
        $bz+nw=1$; the product of the two gives $(ab)(xz)+n(\cdot)=1$, so $ab$ is a unit
        and $U(n)$ is closed. If $a$ is not a unit, $d=\gcd(a,n)>1$ and
        $b=n/d\ne0$ satisfies $ab=n(a/d)\equiv0$.

  \item $a^{2}-1=(a-1)(a+1)\equiv0\pmod p$; since $p$ is prime, Euclid's lemma gives
        $p\mid a-1$ or $p\mid a+1$, so $a\equiv\pm1$. For modulus $8$ that lemma
        fails because $8$ is not prime: the square roots of one are $1,3,5,7$, since
        $3^{2}=5^{2}=7^{2}=49\equiv1$.

  \item The classes $0\le a<pq$ that share a factor with $pq$ are $q$ multiples of $p$
        plus $p$ multiples of $q$, and $0$ is counted twice; so
        $\varphi(pq)=pq-p-q+1=(p-1)(q-1)$. For $p=1019$, $q=1031$:
        $\varphi=1018\cdot1030=1\,048\,540$.

  \item Extending Euclid backwards gives $17^{-1}\equiv2753\pmod{3120}$ (since
        $17\cdot2753=46\,801=15\cdot3120+1$), and
        $65\,537^{-1}\equiv60\,493\pmod{1\,048\,540}$. The division takes $4$ steps against $5$,
        so adding one digit to the modulus adds only a few
        steps.

  \item Fermat's little theorem: $a^{100}\equiv1\pmod{101}$ for $101\nmid a$. Hence
        $5^{100}\equiv1$, $3^{200}=(3^{100})^{2}\equiv1$, and
        $2^{1000}=(2^{100})^{10}\equiv1$. All three are checked by
        \texttt{pow(a,k,101)} and match.

  \item $A=5^{7}\equiv17$, $B=5^{15}\equiv19$, shared key $19^{7}\equiv15$
        (the same as $17^{15}$). If an eavesdropper replaces $A$ with $11$, Bob
        computes $11^{15}\equiv10$ while Alice still holds $15$; two different
        keys appear because the paired exponents are no longer a legitimate
        Diffie--Hellman pair.

  \item With $e=3$ and $m<n^{1/3}$ we have $m^{3}<n$, so $c=m^{3}$ with no
        modular reduction and the integer cube root returns $m$.
        For $n=1\,050\,589$ the bound is $n^{1/3}\approx101{,}66$, so the longest
        message is $m=101$. Random padding makes $m^{3}>n$ so that integer root
        does not exist.

  \item Two things must be true: (i) the parameters and implementation really
        meet the claimed security guarantee, that is, the reduction to a lattice
        problem that is still hard; (ii) the hardness assumption underneath still
        holds and the proof model covers the attacks that are actually used.
        Widespread adoption is not a condition for the claim to be true: a scheme
        used everywhere is still not proven secure if its assumption fails.
        The second sentence remains wrong: RSA is still embedded in legacy systems,
        and migration is done hybrid, so it takes years. (Post-quantum signatures
        already exist: ML-DSA, FIPS 204.)

  \item About $1{,}1774\sqrt{2^{48}}=1{,}1774\cdot2^{24}$, that is $19\,753\,662$
        samples for a one-half collision probability; for $32$ bits
        $1{,}1774\cdot2^{16}=77\,163$. The numbers $2^{24}$ and $2^{16}$ are rough
        square-root guesses, not the exact values.
\end{enumerate}

\section*{Chapter 32: Differential Privacy and the Mathematical Foundations of QML}
\begin{enumerate}[label=\textbf{32.\arabic*}]
  \item False as a sentence. $\varepsilon=0{,}5$ bounds the probability ratio per
        record, it does not remove leakage; it is also a worst-case guarantee. Two things
        that must be true: the mechanism really satisfies $(\varepsilon,\delta)$ for
        the adjacency relation and sensitivity used, and all outputs
        (queries, budget remainder, post-processing) are counted within that budget.

  \item The $d$-dimensional Laplace density is a product, so the ratio is
        $\prod_{j}\exp\bigl((|y_{j}-a'_{j}|-|y_{j}-a_{j}|)/b\bigr)$, and by the
        triangle inequality each factor is at most
        $\exp(|a_{j}-a'_{j}|/b)$. The product is
        $\exp(\sum_{j}|a_{j}-a'_{j}|/b)\le\exp(\Delta f/b)=e^{\varepsilon}$.
        The one-dimensional argument is replaced by the product of per-coordinate ratios, and
        $\ell_1$ takes the place of the single absolute value.

  \item Take $D\sim D'$ differing in one record. That record exists in only
        one of $D_{1},D_{2}$, so only one mechanism has its input changed;
        its ratio is bounded by $e^{\varepsilon_{i}}$, while the others have ratio $1$.
        Hence $\varepsilon=\max_{i}\varepsilon_{i}$ and
        $\delta=\max_{i}\delta_{i}$, not the sum.

  \item The probability of answering ``yes'' is $p+(1-p)/2$ if the truth is ``yes'' and $(1-p)/2$
        if the truth is ``no''; the ratio is $(1+p)/(1-p)$. Hence
        $p=(e^{\varepsilon}-1)/(e^{\varepsilon}+1)$, and for $\varepsilon=2$
        we get $p=0{,}761594$, that is $\tanh1$.

  \item Count query: $b=2$, so $P(|Y|>10)=e^{-5}=0{,}006738$. Mean query
        with $R=1$: $b=10^{-4}/0{,}5=2\times10^{-4}$; an error of $10$ people means
        $10^{-3}$ on the mean scale, which again gives $e^{-5}$, whereas an
        absolute error of $10$ on that scale gives $e^{-50000}\approx0$. This item is
        ambiguous about the units of the error.

  \item $H|0\rangle=(|0\rangle+|1\rangle)/\sqrt2$; $S$ gives
        $(|0\rangle+i|1\rangle)/\sqrt2$; then $H$ gives
        $\frac{1+i}{2}|0\rangle+\frac{1-i}{2}|1\rangle$. The probabilities are $1/2$ and $1/2$,
        the same as $H|0\rangle$, because measurement in the computational basis only sees the modulus
        of the amplitude. What distinguishes them is the relative phase, and the gate is $S^{\dagger}$
        (it multiplies $|1\rangle$ by $-i$). $Z$ after $H$ is no substitute: $HZH=X$,
        which swaps $|0\rangle$ with $|1\rangle$ instead of rotating the phase.

  \item $|\Phi^{+}\rangle=(|00\rangle+|11\rangle)/\sqrt2$. If the second qubit is measured
        first and the result is $0$, the first qubit collapses to $|0\rangle$, so the
        probability of the outcome $01$ is $0$; if the result is $1$, the first qubit collapses
        to $|1\rangle$, so the outcome $11$ is certain and the probability of $01$ is also $0$.
        Both show perfect correlation: the outcomes of the two qubits are always the same.

  \item $U^{\dagger}U=\operatorname{diag}(1,e^{-i\varphi})
        \operatorname{diag}(1,e^{i\varphi})=I$, so it is unitary for every
        $\varphi$. In the computational basis $|1\rangle$ only picks up a phase, its
        amplitude modulus stays the same, so probabilities do not change; but that relative phase
        changes the interference at the next $H$ gate, so the result can
        differ.

  \item 
$n=40$: $2^{40}=1{,}0995\times10^{12}$ amplitudes, times $16$ bytes, gives
        $16$ TiB (not GiB: $2^{40}$ bytes is already one TiB). $n=60$:
        $2^{60}=1{,}1529\times10^{18}$ amplitudes, that is $16$ EiB or
        $1{,}8447\times10^{19}$ bytes, far beyond any storage that exists.

  \item Three questions: whether all the amplitudes are really used, whether
        the output can be extracted with a number of measurements that is not exponential,
        and whether there is a classical comparison on the same task and data. What makes
        me accept: a proven separation or measurable scaling on that task,
        a polynomial number of measurements, and an advantage over a seriously tuned
        classical baseline.
\end{enumerate}



% ---------------------------- index ----------------------------
\cleardoublepage
\phantomsection
\addcontentsline{toc}{chapter}{Index}
\printindex

% ---------------------------- colophon and references ----------------------------
\chapter*{Colophon}
\addcontentsline{toc}{chapter}{Colophon}
\markboth{Colophon}{Colophon}

This manuscript was written in Indonesian, arranged into six parts and thirty-two
chapters, and built with \LaTeX{} (pdf\LaTeX{}, the \texttt{book} class). The entire
manuscript source in \LaTeX{} form ships with this book, so that every number and every
code fragment can be checked again by the reader; the thirteen-chapter version that
preceded it remains archived as a record of how the book was put together.

This book is an original piece of writing: its sentences, its order of presentation, its
worked examples, and its exercises were composed by the author. Its mathematical ideas
come from the books in the References, and the reader is urged to read them. No part of
the manuscript is copied verbatim from a source book; where a definition or a theorem
follows a source closely, that source is named in the chapter map.

Every number inside a computing Lab comes from a real execution on the author's machine
(Python 3, NumPy, SciPy) and can be repeated by the reader. The Lean~4 code fragments in
Chapter~28 were run with Lean 4.10.0 without Mathlib and their output is quoted as it
stood; library tactics such as \texttt{linarith} and \texttt{ring} were not run because
Mathlib is not installed, and that chapter states the boundary.

Cover and layout are deliberately plain, without photographs or raster images: every
figure inside the chapters is drawn as TikZ code, so that the reader can examine its
numbers and change them, and the list of figures is in the front matter. Seven study boxes
serve as the main visual markers --- \emph{For the young reader}, \emph{Note for the
instructor}, \emph{For the researcher}, \emph{Read aloud}, \emph{Work it out by hand},
\emph{Step recipe}, and \emph{Challenge} --- so that one and the same chapter can be
entered through different ways of learning.

\subsection*{License}
The text of this book is licensed under Creative Commons BY 4.0 (CC BY 4.0): it may be
shared, adapted, and used commercially as long as the author is credited and changes are
marked. The \texttt{LICENSE} and \texttt{CITATION.cff} files in the repository carry the
same terms. The Python and Lean~4 code fragments inside it are licensed under MIT, so that
they can be reused without friction.

\subsection*{How to cite}
Pangimpian, Afiq Andico. \emph{Why Machines Can Be Wrong: From Linear Algebra to
Machine-Learning Security}. First Edition. 2026.

\subsection*{Version and feedback}
Manuscript version 1.0, October 2026. Every error you find --- a wrong computation, a
reference that does not match, or a misleading sentence --- is more useful than praise;
send it to \texttt{afiqandico13@gmail.com} with the chapter and page number.

\chapter*{References}
\addcontentsline{toc}{chapter}{References}
\markboth{References}{References}

Every work below is used as a source of ideas, not as material to be copied. The chapter
numbers cited in the \emph{Map to the source books} refer to the chapter numbering of the
edition listed; a reader who wants to check a single claim need only open that edition.

\begin{sloppypar}
\begin{enumerate}[label={[\arabic*]},leftmargin=1.8em]
  \item W. Rudin, \emph{Principles of Mathematical Analysis}, 3rd ed.,
        McGraw-Hill, 1976.
  \item S. Abbott, \emph{Understanding Analysis}, 2nd ed., Springer, 2015.
  \item S. Axler, \emph{Linear Algebra Done Right}, 4th ed., Springer, 2024.
  \item H. Anton, \emph{Elementary Linear Algebra}, 11th ed., Wiley, 2013.
  \item K. Hoffman and R. Kunze, \emph{Linear Algebra}, 2nd ed.,
        Prentice-Hall, 1971.
  \item G. Strang, \emph{Introduction to Linear Algebra}, 6th ed.,
        Wellesley-Cambridge Press, 2023.
  \item J. Stewart, \emph{Calculus: Early Transcendentals}, 8th ed.,
        Cengage, 2016.
  \item J. Stewart, D. K. Clegg, S. Watson, \emph{Calculus: Early
        Transcendentals}, 9th ed., Cengage, 2020.
  \item M. Spivak, \emph{Calculus on Manifolds}, W. A. Benjamin, 1965.
  \item M. Boelkins, D. Austin, S. Schlicker, \emph{Active Calculus:
        Multivariable}, 2nd ed., open text, Grand Valley State University,
        \url{https://activecalculus.org/multi/}.
  \item J. K. Blitzstein and J. Hwang, \emph{Introduction to Probability},
        2nd ed., CRC Press, 2019.
  \item L. Wasserman, \emph{All of Statistics: A Concise Course in Statistical
        Inference}, Springer, 2004.
  \item D. Goldberg, \emph{What Every Computer Scientist Should Know About
        Floating-Point Arithmetic}, ACM Computing Surveys 23(1), 1991.
  \item S. Boyd and L. Vandenberghe, \emph{Convex Optimization}, Cambridge
        University Press, 2004.
  \item D. Bertsimas and J. N. Tsitsiklis, \emph{Introduction to Linear
        Optimization}, Athena Scientific, 1997.
  \item L. A. Wolsey, \emph{Integer Programming}, Wiley, 1998.
  \item I. Goodfellow, Y. Bengio, A. Courville, \emph{Deep Learning},
        MIT Press, 2016.
  \item M. A. Nielsen, \emph{Neural Networks and Deep Learning}, online book,
        Determination Press, 2015.
  \item J. Avigad and P. Massot, \emph{Mathematics in Lean}, version 4.19.0, June
        2025, \url{https://avigad.github.io/mathematics_in_lean/} (also available
        as an Indonesian-language edition).
  \item J. Avigad, L. de Moura, S. Kong, S. Ullrich, \emph{Theorem Proving in
        Lean 4}, Lean FRO, \url{https://lean-lang.org/theorem_proving_in_lean4/}.
  \item I. J. Goodfellow, J. Shlens, C. Szegedy, \emph{Explaining and Harnessing
        Adversarial Examples}, ICLR 2015 (arXiv:1412.6572).
  \item A. Madry, A. Makelov, L. Schmidt, D. Tsipras, A. Vladu,
        \emph{Towards Deep Learning Models Resistant to Adversarial Attacks},
        ICLR 2018 (arXiv:1706.06083).
  \item A. Athalye, N. Carlini, D. Wagner, \emph{Obfuscated Gradients Give a
        False Sense of Security: Circumventing Defenses to Adversarial
        Examples}, ICML 2018 (arXiv:1802.00420).
  \item C. Szegedy, W. Zaremba, I. Sutskever, J. Bruna, D. Erhan, I. Goodfellow,
        R. Fergus, \emph{Intriguing Properties of Neural Networks}, ICLR 2014
        (arXiv:1312.6199).
  \item N. Carlini and D. Wagner, \emph{Towards Evaluating the Robustness of
        Neural Networks}, IEEE Symposium on Security and Privacy 2017
        (arXiv:1608.04644).
  \item J. M. Cohen, E. Rosenfeld, J. Z. Kolter, \emph{Certified Adversarial
        Robustness via Randomized Smoothing}, ICML 2019 (arXiv:1902.02918).
  \item N. Moustafa and J. Slay, \emph{UNSW-NB15: A Comprehensive Data Set for
        Network Intrusion Detection Systems}, Military Communications and
        Information Systems Conference (Milcom), 2015, pp. 1--6.
  \item M. elShehaby and A. Matrawy, \emph{Evasion Adversarial Attacks Remain
        Impractical Against ML-based Network Intrusion Detection Systems,
        Especially Dynamic Ones}, arXiv:2306.05494, 2023.
  \item W. Diffie and M. E. Hellman, \emph{New Directions in Cryptography},
        IEEE Transactions on Information Theory 22(6), 1976, pp. 644--654.
  \item R. L. Rivest, A. Shamir, L. Adleman, \emph{A Method for Obtaining Digital
        Signatures and Public-Key Cryptosystems}, Communications of the ACM 21(2),
        1978, pp. 120--126.
  \item P. W. Shor, \emph{Polynomial-Time Algorithms for Prime Factorization and
        Discrete Logarithms on a Quantum Computer}, SIAM Journal on Computing
        26(5), 1997, pp. 1483--1509.
  \item C. Dwork, \emph{Differential Privacy}, International Colloquium on
        Automata, Languages and Programming (ICALP), LNCS 4052, Springer, 2006,
        pp. 1--12.
  \item M. Abadi, A. Chu, I. Goodfellow, H. B. McMahan, I. Mironov, K. Talwar,
        L. Zhang, \emph{Deep Learning with Differential Privacy}, ACM CCS 2016,
        pp. 308--318 (arXiv:1607.00133).
  \item NIST, \emph{FIPS 203: Module-Lattice-Based Key-Encapsulation Mechanism
        Standard}, National Institute of Standards and Technology, 13 August
        2024, \url{https://csrc.nist.gov/pubs/fips/203/final}.
\end{enumerate}
\end{sloppypar}


\end{document}
